% THIS IS NOT THE FILE YOU SHOULD PROCESS. IT IS THE "MAIN" FILE,
% BUT IT GETS INCLUDED BY ONE OF THE hott-xxx.tex FILES. THOSE ARE
% THE MAIN ONES.

% DOCUMENT CLASS
\documentclass[\OPTfontsize]{book}

% 中译本改动：正文已译成中文，加载 ctex 以支持 CJK 排版，用 Fandol 字体集
% （随 texlive 自带，不依赖系统安装的字体），跟 XeLaTeX/LuaLaTeX 都兼容；
% 必须紧跟在 \documentclass 之后、其余字体/编码相关宏包之前加载。
\usepackage[fontset=fandol]{ctex}

% 中译本改动：用 LuaLaTeX 编译时，xcolor（被 \pagecolor{white}/\nopagecolor
% 一类调用触发，见 front.tex/main.tex 的封面逻辑）内部直接用了 \pdfliteral/
% \pdfsave/\pdfrestore 这些 pdfTeX 原语，LuaTeX 并不原生提供，报
% "Undefined control sequence"（真实复现，编译到 setmath.tex 附近报了 9 次）。
% luatex85 包专门用于在 LuaTeX 下把这批 pdfTeX 原语重新定义回来，是这类
% "老宏包硬编码 pdfTeX 原语"问题的标准解法；只在 LuaTeX 下加载，其他引擎
% 下这个包本身会直接退出，不需要额外判断。
\usepackage{luatex85}

\PassOptionsToPackage{table}{xcolor}

\usepackage{etex} % We're running out of registers and dimensions, or some such

% PAGE GEOMETRY
\usepackage[papersize={\OPTpagesize},
  twoside,
  includehead,
  top=\OPTtopmargin,
  bottom=\OPTbottommargin,
  inner=\OPTinnermargin,
  outer=\OPToutermargin,
bindingoffset=\OPTbindingoffset]{geometry}

% HYPERLINKING AND PDF METADATA
\usepackage[backref=page,
  colorlinks,
  citecolor=linkcolor,
  linkcolor=linkcolor,
  urlcolor=linkcolor,
  unicode,
  pdfauthor={Univalent Foundations Program},
  pdftitle={Homotopy Type Theory: Univalent Foundations of Mathematics},
  pdfsubject={Mathematics},
pdfkeywords={type theory, homotopy theory, univalence axiom}]{hyperref}
\renewcommand{\backref}[1]{}
\renewcommand{\backrefalt}[4]{%
  \ifcase #1 %
  (No citations.)
  \or
  (Cited on page\ #2.)
  \else
  (Cited on pages\ #2.)
\fi}

% OTHER PACKAGES

% Use this package and stick \layout somewhere in the text to see
% page margins, text size and width etc. Useful for debugging page format.
\usepackage{layout}

%%% Because Germans have umlauts and Slavs have even stranger ways of mangling letters
\usepackage[utf8]{inputenc}

%%% For table {tab:theorems}
\usepackage{pifont}

%%% Multi-Columns for long lists of names
\usepackage{multicol}

% 中译本改动：amsmath/amsthm/stmaryrd/mathtools 都要放到 fontsetup（内部加载 unicode-math）之前
\usepackage{amsmath,amsthm,stmaryrd,mathtools}

%%% Set the fonts
% 改回 newcm：用 fontsetup 包（要求 Xe/LuaLaTeX，走 fontspec+unicode-math）
\usepackage[default]{fontsetup}

% 中译本改动：fontsetup[default] 正文本身默认已经是等高的现代体（lining）数字，
% 问题只出在 \textsc/\scshape（页眉里 \chaptermark/\sectionmark 就是拿 \textsc
% 包住章节号）：fspdefault.tex 给小型大写字母专门配了
% SmallCapsFeatures={Numbers=OldStyle}，数字进了 \textsc 就会变成会下沉出基线
% 的老式体。\addfontfeatures 只能改当前字形，改不动 sc 形状在 \setmainfont 时就
% 定死的 SmallCapsFeatures（实测验证过，追加特性无效）；必须重新整个
% \setmainfont 一次才会生效。下面这段就是 fspdefault.tex 里那条 \setmainfont
% 命令的原样复制，唯一改动是把 SmallCapsFeatures 里的 Numbers=OldStyle
% 换成 Numbers=Lining；SizeFeatures 分段保留，避免小字号（脚注等用到的
% NewCM08）跟着受影响。
\setmainfont[%
  SizeFeatures={%
{Size=-8, Font=NewCM08-Book.otf,
      ItalicFont=NewCM08-BookItalic.otf,%
      BoldFont=NewCM10-Bold.otf,%
      BoldItalicFont=NewCM10-BoldItalic.otf,%
      SlantedFont=NewCM08-Book.otf,%
      BoldSlantedFont=NewCM10-Bold.otf,%
      SmallCapsFeatures={Numbers=Lining}},
{Size=8, Font=NewCM08-Book.otf,
      ItalicFont=NewCM08-BookItalic.otf,%
      BoldFont=NewCM10-Bold.otf,%
      BoldItalicFont=NewCM10-BoldItalic.otf,%
      SlantedFont=NewCM08-Book.otf,%
      BoldSlantedFont=NewCM10-Bold.otf,%
      SmallCapsFeatures={Numbers=Lining}},
{Size= 9-, Font = NewCM10-Book.otf,
      ItalicFont=NewCM10-BookItalic.otf,%
      BoldFont=NewCM10-Bold.otf,%
      BoldItalicFont=NewCM10-BoldItalic.otf,%
      SlantedFont=NewCM10-Book.otf,%
      BoldSlantedFont=NewCM10-Bold.otf,%
      SmallCapsFeatures={Numbers=Lining}}%
  },%
  SmallCapsFeatures={Numbers=Lining},%
  BoldSlantedFont=NewCM10-Bold.otf,%
  SlantedFont=NewCM10-Book.otf,%
  BoldItalicFont=NewCM10-BoldItalic.otf,%
  BoldFont=NewCM10-Bold.otf,%
  ItalicFont=NewCM10-BookItalic.otf,%
  SlantedFeatures={FakeSlant=0.25},%
  BoldSlantedFeatures={FakeSlant=0.25},%
]{NewCM10-Book.otf}

% \Box（模态逻辑的"必然"算子）
\renewcommand{\Box}{□}

% \centerdot（macros.tex 里 \ct 道路拼接算子的底层符号，全书高频使用）原来
% 也是 amssymb 提供的，去掉后未定义，同样用字面 Unicode 字符代替。
\newcommand{\centerdot}{⋅}

% \ocircle（macros.tex 里 \modal 模态算子的底层符号）原来是 wasysym 提供
% 的，去掉后未定义，同样用字面 Unicode 字符（○ U+25CB）代替。
\newcommand{\ocircle}{○}

% \leadsto（"道路"记号 $a\leadsto b$，直接在正文里裸用，不经过 macros.tex）
% 原来是 latexsym/amsfonts 提供的，是 LaTeX 内核自带的"提示占位"桩函数
% （报错信息会提示"Load the latexsym or amsfonts package"），之前 amsmath
% 传递加载了 amsfonts 从而侧面提供了它；现在同样用字面 Unicode 字符
% （⇝ U+21DD）代替，跟 \Box 一样要用 \renewcommand（因为内核桩已经"定义"过
% 这个名字了，\newcommand 会报"已经 defined"）。
\renewcommand{\leadsto}{⇝}

% \square（macros.tex 里好几个 surreal-number "bar name" 宏，如 \tapbname/
% \tlbname/\tleb/\tltb，都用它拼出来）同样是 amssymb 提供的，去掉后未定义。
\newcommand{\square}{□}

\usepackage{ifpdf}

\usepackage{graphicx}
\DeclareGraphicsExtensions{.png}
\RequirePackage{bmpsize-base}
\RequirePackage{ifpdf}
\makeatletter
% 中译本改动：这段 DVI 模式下手工计算图片包围盒的 hack，是为真正的
% "dvips -> ps2pdf" 流程写的，假设的是 pdftex 的原语。ifpdf 包在 XeLaTeX/
% LuaLaTeX 下都可能判定成"不是 pdf 模式"从而误触发这段 hack，报
% "Undefined control sequence"（真实复现：\@bmpsize@read 报错，编译直接
% 中止；用 \ifdefined\XeTeXversion/\luatexversion 分别探测两种引擎又在
% LuaLaTeX 下报 "\ifdefined ... was incomplete"，怀疑是这两个宏在当前
% texlive 版本下的可用性问题，没有深究）。中译本只会用 XeLaTeX 或
% LuaLaTeX 编译，两者都能自己正确处理图片包围盒，不需要这段 hack，直接
% 无条件跳过最省事也最不容易出岔子。
\expandafter\@gobble{%
\AtBeginDocument{%
\let\Gin@ii@old=\Gin@ii
\def\Gin@ii[#1]#2{%
  \begingroup
    \let\@found\@empty
    \@for\@type:=\bmpsize@types\do{%
      \ifx\@found\@empty
        \@nameuse{bmpsize@read@\@type}{#2.\@type}%
        \ifbmpsize@ok
          \let\@found=\@type
        \fi
      \fi
      \ifx\@found\@empty
        \@nameuse{bmpsize@read@\@type}{#2}%
        \ifbmpsize@ok
          \let\@found=\@type
        \fi
      \fi
    }%
    \ifx\@found\@empty
      \Gin@ii@old[#1]{#2}%
    \else
      \Gin@ii@old[natwidth=\bmpsize@width bp,natheight=\bmpsize@height bp,#1]{#2}%
    \fi
  \endgroup
}%
}
}
\makeatother
 % for bounding boxes in dvi mode
\usepackage{comment}

\usepackage{fancyhdr} % To set headers and footers

\usepackage{nextpage} % So we can jump to odd-numbered pages

% amsmath/amsthm/stmaryrd/mathtools 已经在 fontsetup 之前加载过了，不再加载，避免 "Command 已经 defined" 冲突。
\usepackage{enumitem,xspace}
\usepackage{xstring} % For generating singluars and plurals in \backref

\usepackage{xcolor} % For colored cells in tables we need \cellcolor
\usepackage{wallpaper} % For the background image on the cover page

\usepackage{booktabs} % For nice tables
\usepackage{array} % For nice tables

\definecolor{linkcolor}{rgb}{\OPTlinkcolor}
\usepackage{aliascnt}
\usepackage[capitalize]{cleveref}
% 中译本改动：交换图箭头被画成实心黑块，根源定位到 xy-pic 的 pdf 画线驱动
% （xypdf.tex）里，线宽 \xP@lw 是直接拿 \fontdimen8\textfont3（数学字体族 3，
% 即扩展/定界符字体的``默认规则粗细''）当线宽用的。经典 TFM 数学字体（cmex10）
% 这个值是 0.4pt 左右；但 unicode-math（fontsetup 内部）接管数学字体族后，
% 从 OpenType MATH 表折算出的这个兼容量变成了 ~7.2pt，于是所有直线/斜线都
% 被画成十几倍粗的实心色块（复现：categories.tex 里的交换图）。这里手动把
% \fontdimen8\textfont3 改回经典值，只影响 xy 内部这一处线宽计算，不影响
% 正文数学排版本身用到的字体度量。必须放在 \setmathfont（fontsetup 内部）
% 之后、\usepackage{xy} 之前，且不依赖 pdf/cmtip 具体选哪个纸面驱动
% ——两者都读同一个 \fontdimen，同样受影响，同样被这一行修好。
\fontdimen8\textfont3=0.4pt\relax
\usepackage[all,2cell,cmtip]{xy}
\UseAllTwocells
\usepackage{braket} % used for \setof{ ... } macro

\usepackage{tikz}
\usetikzlibrary{decorations.pathmorphing,arrows}

\usepackage{etoolbox}           % hacking commands for TOC

\usepackage{mathpartir}         % for formal.tex appendix, section 3

\usepackage[numbered]{bookmark} % add chapter/section numbers to the toc in the pdf metadata

\input{macros}

%%%% Indexing
\usepackage{makeidx}
\makeindex

%%%% Header and footers
\pagestyle{fancyplain}
\setlength{\headheight}{15pt}
\renewcommand{\chaptermark}[1]{\markboth{\textsc{Chapter \thechapter. #1}}{}}
\renewcommand{\sectionmark}[1]{\markright{\textsc{\thesection\ #1}}}

\lhead[\fancyplain{}{{\thepage}}]%
{\fancyplain{}{\nouppercase{\rightmark}}}
\rhead[\fancyplain{}{\nouppercase{\leftmark}}]%
{\fancyplain{}{\thepage}}
\cfoot[]{}
\lfoot[]{}
\rfoot[]{}

%%%% Chapter & part style
\usepackage[raggedright]{titlesec}
\titleformat{\part}[display]{\fontsize{\OPTpartfont}{\OPTpartfont}\fontseries{m}\fontshape{sc}\selectfont}{\hfil\partname\ \Roman{part}}{\OPTpartskip}{\fontsize{\OPTparttitlefont}{\OPTparttitlefont}\fontseries{b}\fontshape{sc}\selectfont\hfil}
\titleformat{\chapter}[display]{\fontsize{\OPTchapterfont}{\OPTchapterfont}\fontseries{m}\fontshape{it}\selectfont}{\chaptertitlename\ \thechapter}{\OPTchapterskip}{\fontsize{\OPTchaptertitlefont}{\OPTchaptertitlefont}\fontseries{b}\fontshape{n}\selectfont}

% Patch bug in titlesec 2.10.1
\makeatletter
\@ifpackagelater{titlesec}{2016/03/21}{%
  % Package titlesec is on version 2.10.2 or higher, nothing to do
}{\@ifpackagelater{titlesec}{2016/03/15}{%
    % Package titlesec on version 2.10.1, patch accordingly
    \usepackage{etoolbox}%
    \patchcmd{\ttlh@hang}{\parindent\z@}{\parindent\z@\leavevmode}{}{}%
    \patchcmd{\ttlh@hang}{\noindent}{}{}{}%
  }{% Package titlesec is on version 2.10.0 or lower, nothing to do %
}}
\makeatother

% To avoid compiling stuff other than what you're working on right
% now, uncomment the following command and give your file as its
% argument.
%\includeonly{}

% For some reason \pagecolor overlays the cover image,
% unless we use it once before the document starts.
\definecolor{covercolor}{cmyk}{\OPTcovercolor}
\definecolor{covertext}{cmyk}{\OPTcovertextcolor}
\pagecolor{white}
\ifpdf
\nopagecolor
\fi

% Make a new page for each exercise, if the option is set
\ifdef{\OPTexerciseperpage}
{
  \BeforeBeginEnvironment{ex}{\newpage}
}{}

\begin{document}

% NB: This does not actually appear anywhere because we have
% a custom title page.
\title{Homotopy Type Theory: Univalent Foundations of Mathematics}
\author{The Univalent Foundations Program}

\frontmatter % Turn on arabic page numbers and unnumbered chapters

% Half-title, title, copyright page do not have displayed page numbers
\pagestyle{empty}

%%%%%%%%%%%%%%%%%%%% Chinese title, letter-spaced but not justified %%%%%%%%%%%%%%%%%%%%
\newcommand{\hottcjktitle}{同\hspace{0.4em}伦\hspace{0.4em}类\hspace{0.4em}型\hspace{0.4em}论}
\newcommand{\hottcjksubtitle}{数\hspace{0.3em}学\hspace{0.3em}的\hspace{0.3em}泛\hspace{0.3em}等\hspace{0.3em}基\hspace{0.3em}础}

%%%%%%%%%%%%%%%%%%%% Cover page %%%%%%%%%%%%%%%%%%%%

\newgeometry{noheadfoot,bindingoffset=-5pt,top=0pt,bottom=0pt,inner=0pt,outer=0pt}%
\ifOPTcover
\setcounter{page}{-1} % Otherwise we end up having two pages numbered 1
\newlength{\coverheight}
\setlength{\coverheight}{\OPTcoverheight}
\newlength{\coverwidth}
\setlength{\coverwidth}{\OPTcoverwidth}
\newcommand{\frontpage}{
\begin{minipage}[b][\coverheight][c]{\coverwidth}
\hbox{}\hfill
\begin{minipage}[b][\coverheight][t]{0.83\coverwidth}
\color{covertext}
\vspace{\OPTtopskip}
{\fontsize{\OPTcovertitlefont}{\OPTcovertitlefont}\fontseries{b}\selectfont%
  \hfill \hottcjktitle\par}\par
\vspace*{\OPTcovertitleskip}
{\fontsize{\OPTcoversubtitlefont}{\OPTcoversubtitlefont}\fontshape{it}\selectfont
\hfill \hottcjksubtitle}

\vfill

{\fontsize{\OPTcoverauthorfont}{\OPTcoverauthorfont}\fontseries{b}\fontshape{sc}\selectfont

\hfill The Univalent Foundations Program
\par
\vspace*{\OPTcoverauthorskip}
\hfill 普林斯顿高等研究院\par

\vspace*{\OPTbotskip}
}

\end{minipage}\hfill\hbox{}
\end{minipage}
}

\ThisLLCornerWallPaper{1.1}{\OPTfrontimage}
\pagecolor{covercolor}
\frontpage
\newpage
% Reset page counter, cover page does not count
\ifpdf
\nopagecolor
\else
\pagecolor{white}
\fi
\cleartooddpage
\else
\fi

%%%%%%%%%%%%%%%%%%%% Bastard page %%%%%%%%%%%%%%%%%%%%
\ifOPTbastard
\cleartooddpage
\hbox{}
\vspace{0.2\textwidth}
{\centering
\fontsize{\OPTbastardtitlefont}{\OPTbastardtitlefont}\fontshape{n}\selectfont%
\textbf{\hottcjktitle}\par
\vspace*{\OPTbastardtitleskip}
\fontsize{\OPTbastardsubtitlefont}{\OPTbastardsubtitlefont}\fontshape{n}\selectfont%
\textit{\hottcjksubtitle}\par
}
\else
\fi

%%%%%%%%%%%%%%%%%%%% Title page %%%%%%%%%%%%%%%%%%%%
\cleartooddpage
\hbox{}\vfill
{\centering
{\fontsize{\OPTtitletitlefont}{\OPTtitletitlefont}\fontseries{b}\selectfont%
\hottcjktitle}\par
\vspace*{\OPTtitletitleskip}
{\fontsize{\OPTtitlesubtitlefont}{\OPTtitlesubtitlefont}\fontshape{it}\selectfont%
\hottcjksubtitle}\par
\vspace*{\OPTtitleskip}
{\fontsize{\OPTtitleauthorfont}{\OPTtitleauthorfont}\fontshape{n}\selectfont%
The Univalent Foundations Program\par
\vspace*{\OPTtitleauthorskip}
普林斯顿高等研究院\par
}
\vspace*{\OPTtitleskip}
\vspace*{\OPTtitleskip}
\includegraphics[width=\OPTtitlewidth]{\OPThalftorus}\par
}

\vfill
\hbox{}

\clearpage
%%% Restore page style
\restoregeometry

%%%%%%%%%%%%%%%%%%%% Copyright page %%%%%%%%%%%%%%%%%%%%
\hbox{}
\vfill
\input{version.tex}
{\small
\noindent
\emph{``同伦类型论：数学的泛等基础''}\\
\copyright\ 2013 The Univalent Foundations Program

\medskip
\noindent
书籍版本：\texttt{\OPTversion}

\medskip
\noindent
MSC 2010 分类号：
\texttt{03-02},
\texttt{55-02},
\texttt{03B15}

\bigskip
\footnotesize

\noindent
本作品采用
\textbf{\emph{知识共享 署名-相同方式共享 3.0 未本地化版许可协议}}
授权。
%
如需查看该许可协议的副本，请访问
\url{http://creativecommons.org/licenses/by-sa/3.0/}。

\bigskip

\noindent
本书可在 \url{http://homotopytypetheory.org/book/} 免费获取。

\bigskip

\noindent
\emph{\textbf{\small 致谢}}

\medskip

\noindent
除了普林斯顿高等研究院的慷慨支持外，本书的部分贡献者还得到了以下机构和项目的部分或全额资助：
%
\begin{itemize}
\item 高等研究院成员协会：对高等研究院的一项资助 % Dan Grayson
% SLOVENIA
\item 斯洛文尼亚共和国研究活动局：  % Andrej's Slovenian agency
\href{http://www.sicris.si/search/prg.aspx?id=6120}{P1--0294}，
\href{http://www.sicris.si/search/prj.aspx?id=7109}{N1--0011}。

\item 美国空军科学研究办公室：
  FA9550-11-1-0143，以及 % Steve's ASFOR
  FA9550-12-1-0370。  % Bob's ASFOR
  {
    \setlength{\parskip}{0pt}
    \begin{quote}
      \noindent\scriptsize
      本材料部分基于 AFOSR 在上述资助项目下所支持的工作。
      本出版物中表达的任何观点、发现和结论或建议均为作者本人的观点，并不一定反映 AFOSR 的观点。
    \end{quote}
  }

\item 英国工程与物理科学研究理事会： % Thorsten and students
   \href{http://gow.epsrc.ac.uk/NGBOViewGrant.aspx?GrantRef=EP/G034109/1}{EP/G034109/1}， % Reusability and dependent types
   \href{http://gow.epsrc.ac.uk/NGBOViewGrant.aspx?GrantRef=EP/G03298X/1}{EP/G03298X/1}。 % Theory and Application of Induction Recursion

\item 欧盟第七框架计划，资助协议编号 243847（%
\href{http://wiki.portal.chalmers.se/cse/pmwiki.php/ForMath/ForMath/}{ForMath}）。 %% several Europeans, via Bas

\item 美国国家科学基金会：
  \href{http://www.nsf.gov/awardsearch/showAward.do?AwardNumber=1001191}{DMS-1001191}， %% Steve's NSF, including Chris and Kristina
  \href{http://www.nsf.gov/awardsearch/showAward.do?AwardNumber=1100938}{DMS-1100938}， %% Vladimir's NSF
  \href{http://www.nsf.gov/awardsearch/showAward.do?AwardNumber=1116703}{CCF-1116703}， %% Foundations and Applications of Higher-Dimensional Directed Type Theory
  以及
  \href{http://www.nsf.gov/awardsearch/showAward.do?AwardNumber=1128155}{DMS-1128155}。 %% IAS support for Mike Shulman, copied from our %% paper by Dan Licata
  {
    \setlength{\itemsep}{0pt}
    \begin{quote}
      \noindent\scriptsize
      本材料部分基于美国国家科学基金会在上述资助项目下所支持的工作。
      本材料中表达的任何观点、发现和结论或建议均为作者本人的观点，并不一定反映美国国家科学基金会的观点。
    \end{quote}
  }
\item 西蒙尼基金：对高等研究院的一项资助          %% Dan Grayson
  \end{itemize}


}
\cleartooddpage

%%% Local Variables:
%%% mode: latex
%%% TeX-master: "hott-online"
%%% End:
 %%% Title page and copyright

\cleartooddpage

% Add Preface to PDF Metadata but not printed TOC
\hypertarget{preface}{}
\bookmark[dest=preface]{前言}

% Preface and TOC have arabic numbers
\pagestyle{fancyplain}

\chapter*{前言}
\label{cha:preface}

% Uncomment this if you think Preface should appear in TOC
%\addcontentsline{toc}{chapter}{前言}
\subsection*{普林斯顿高等研究院 泛等基础特别学年}

2012--13 学年，在 Steve Awodey、Thierry Coquand 和 Vladimir Voevodsky 的组织下，普林斯顿高等研究院数学学院举办了一场关于数学的泛等基础的特别学年活动。以下是官方参与者名单。

\begin{multicols}{\OPTprefacecols}{
\begin{itemize}
\item[] Peter Aczel
\item[] Benedikt Ahrens
\item[] Thorsten Altenkirch
\item[] Steve Awodey
\item[] Bruno Barras
\item[] Andrej Bauer
\item[] Yves Bertot
\item[] Marc Bezem
\item[] Thierry Coquand
\item[] Eric Finster
\item[] Daniel Grayson
\item[] Hugo Herbelin
\item[] Andr\'e Joyal
\item[] Dan Licata
\item[] Peter Lumsdaine
\item[] Assia Mahboubi
\item[] Per Martin-L\"of
\item[] Sergey Melikhov
\item[] Alvaro Pelayo
\item[] Andrew Polonsky
\item[] Michael Shulman
\item[] Matthieu Sozeau
\item[] Bas Spitters
\item[] Benno van den Berg
\item[] Vladimir Voevodsky
\item[] Michael Warren
\item[] Noam Zeilberger
\end{itemize}
}
\end{multicols}

\noindent 此外，还有以下学生参与，他们的贡献同样宝贵。

\begin{multicols}{\OPTprefacecols}{
\begin{itemize}
\item[] Carlo Angiuli
\item[] Anthony Bordg
\item[] Guillaume Brunerie
\item[] Chris Kapulkin
\item[] Egbert Rijke
\item[] Kristina Sojakova
\end{itemize}
}
\end{multicols}

\noindent 另外，还有以下短期与长期访问学者（包括学生访客），他们对特别学年的贡献同样至关重要。

\begin{multicols}{\OPTprefacecols}{
\begin{itemize}
\item[] Jeremy Avigad
\item[] Cyril Cohen
\item[] Robert Constable
\item[] Pierre-Louis Curien
\item[] Peter Dybjer
\item[] Mart{\'\i}n Escard{\'o}
\item[] Kuen-Bang Hou
\item[] Nicola Gambino
\item[] Richard Garner
\item[] Georges Gonthier
\item[] Thomas Hales
\item[] Robert Harper
\item[] Martin Hofmann
\item[] Pieter Hofstra
\item[] Joachim Kock
\item[] Nicolai Kraus
\item[] Nuo Li
\item[] Zhaohui Luo
\item[] Michael Nahas
\item[] Erik Palmgren
\item[] Emily Riehl
\item[] Dana Scott
\item[] Philip Scott
\item[] Sergei Soloviev
\end{itemize}
}
\end{multicols}

\subsection*{关于本书}

我们起初并没有打算写一本书。
本书源于大家的一番共同努力：提出一种新的、人类可读可理解“非形式化类型论”写法，作为可由机器验证的形式证明的补充。
泛等基础与这样一种数学基础观紧密相连：数学应该基于一套能够在计算机证明助理中实现的基础。
虽然本书本身并不包含这样的形式化工作，但这里介绍的许多内容，实际上最初都是在证明助理中以完全形式化的方式完成的；
直到后来，我们才将其“去形式化”，整理成如今读者眼前的模样。这恰好将数学形式化中的通常次序颠倒了过来，实属一件奇事。

前述的每一位参与者，都以自己的方式为那一“特别学年”作出了贡献——也因此为本书作出了贡献；这些贡献或是思想，或是文字，或是实际的工作。
贯穿整整一年的那种合作精神，确实非同寻常。

\mentalpause

在此要特别感谢高等研究院；没有它，这本书显然根本不可能问世。
事实证明，这里正是孕育这个数学新分支的理想环境：思想活跃，气氛融洽，鼎力支持。
但愿这种独一无二的氛围还能有一丝留存于本书的字里行间，也留存在这一新兴领域的未来发展之中。

\bigskip

\begin{flushright}
The Univalent Foundations Program\\
普林斯顿高等研究院\\
普林斯顿，2013年4月
\end{flushright}

%%%%%%%%%%%%% end of scope of local macros
% Local Variables:
% TeX-master: "hott-online"
% End:

\cleartooddpage[\thispagestyle{empty}]

% Add TOC to PDF Metadata but not printed TOC
\hypertarget{toc}{}
\bookmark[dest=toc]{目录}

\setcounter{tocdepth}{1}        % chapters and sections for the toc
\tableofcontents
\setcounter{tocdepth}{2}        % chapters, sections, and subsections for the
% metadata of the pdf
\cleartooddpage[\thispagestyle{empty}]

\mainmatter % Turn on roman page numbers and numbered chapters

% Turn on headers and footers for mainmatter (must appear after \cleartooddpage)
\pagestyle{fancyplain}

\chapter*{引言}
\markboth{\textsc{引言}}{}
\addcontentsline{toc}{chapter}{引言}
\setcounter{page}{1}
\pagenumbering{arabic}

\emph{同伦类型论}是一个新兴的数学分支，它以一种出人意料的方式，将若干看似不同的领域汇聚到了一起。
这一学科基于近年来发现的\emph{同伦论}与\emph{类型论}之间的深刻联系。
同伦论脱胎于代数拓扑与同调代数，且与高阶范畴论有着千丝万缕的关联；而类型论则属于数理逻辑与理论计算机科学的范畴。
尽管二者之间的联系正处于密切研究之中，但人们日渐清楚地认识到，当下已有的发现不过是冰山一角；
要想真正理解这门学科，还需要投入更多的时间、付出更多艰辛的努力。
它所触及的话题跨度之大，看似天差地别：从球面的同伦群，到类型检查的算法，再到弱 $\infty$-群胚的定义，无不在其范围之内。

同伦类型论还为数学的根基问题带来了新的理念。
\index{foundations, univalent}%
一方面，我们有 Voevodsky （沃沃茨基）提出的精妙而优美的\emph{泛等公理（univalence axiom）}。
\index{univalence axiom}%
特别地，泛等公理蕴含了这样一条原则：同构的结构可以被等同。
这正是数学家们在日常工作中一直乐于采用的原则，尽管它与传统基础理论的“官方”教条并不相容。
另一方面，我们还有\emph{高阶归纳类型（higher inductive types）}，它们以直接的、逻辑的方式描述了同伦论中一些基本的空间与构造：球面、柱面、截断、局部化等。
二者在经典的集合论基础中都无法直接刻画，但在同伦类型论中结合起来后，它们使一种全新的“同伦类型逻辑”成为可能。
\index{foundations}%

由此浮现出一种新的数学基础的构想：它具有内在的同伦内容，以一种“在不变量意义下”理解数学对象的方式来组织数学；
与此同时，它还可以方便地在机器上实现，从而为实际从事数学研究的人提供切实的助力。
这就是\emph{泛等基础（Univalent Foundations）}计划。
本书旨在对泛等基础的基本内容作首次系统性阐述，并汇集这种新型推理方式的若干实例——但既不要求读者了解或学习任何形式逻辑，也无需使用任何计算机证明助理。

% This enlarges the page by one line in letter format. Use sparringly.
\OPTwidow

我们要强调，同伦类型论是一个年轻的领域，泛等基础更是尚在进行中的工作。
本书应当被视为该领域在写作之时一个局部的“快照”，而非对一座已落成大厦的精心阐释。
正如我们将在后文简略讨论的那样，同伦类型论的许多方面尚未被完全理解——有些甚至本书都未曾触及。
最终的理论几乎可以肯定不会与本书所述完全一样，但它\emph{至少}会同样强大、同样有力；因此我们相信，泛等基础终将成为集合论的一种可行替代方案，成为大多数数学家所从事的非形式化数学的“隐含基础”。

\subsection*{类型论}

类型论最初由 Bertrand Russell（罗素）\cite{Russell:1908}\index{Russell, Bertrand}发明，作为阻止当时正在研究的数学的逻辑基础中出现的悖论的一种手段。
在接下来的几十年里，许多人对其作了进一步发展，尤其是 Church（丘奇）~\cite{Church:1940tu,Church:1941tc}将其与他的 \textit{$\lambda$-演算}相结合。
尽管类型论通常不被视为经典数学的基础——集合论更为惯用——但它仍有大量应用，特别是在计算机科学与程序设计语言理论中~\cite{Pierce-TAPL}。
\index{programming}%
\index{type theory}%
\index{lambda-calculus@$\lambda$-calculus}%
Per Martin-L\"of（佩尔·马丁-洛夫）\cite{Martin-Lof-1972,Martin-Lof-1973,Martin-Lof-1979,martin-lof:bibliopolis}等人对 Church 的类型系统作了“直谓的”修改，这一系统现在通常被称为依值的、构造性的、直觉主义的类型论，或简称\emph{Martin-L\"of 类型论}。这正是我们在此所考虑的系统的基础；它最初是作为形式化构造性数学的严格框架而设计的。在下文中，我们经常用“类型论”来特指这个系统及其类似的系统，尽管类型论作为一门学科要宽广得多（参见 \cite{somma,kamar} 了解类型论的历史）。

在类型论中，与集合论不同，对象是依据一个原始的\emph{类型（type）}概念来分类的，这类似于编程语言中使用的数据类型。
这些结构精细的类型可以用来表达被分类对象的详细规格，从而产生关于这些对象的推理原则。
举一个非常简单的例子，积类型 $A\times B$ 的对象已知具有 $\pairr{a,b}$ 的形式，因此人们自然知道如何构造它们、如何分解它们。
类似地，函数类型 $A\to B$ 的对象可以由一个以类型 $A$ 的对象为参数的类型 $B$ 的对象给出，并且可以在类型 $A$ 的参数上求值。
所有对象的这种严格可预测的行为（与集合论更为自由的构成原则——允许异质集合——形成对比），正是类型论得以广泛用于验证计算机程序正确性的一个方面。
与类型构造相关的清晰推理原则，也构成了现代\emph{计算机证明助理}的基础，%
\index{proof!assistant}%
\indexsee{computer proof assistant}{proof assistant}
\index{mathematics!formalized}%
这类工具用于形式化数学并验证形式化证明的正确性。
下文还将回到类型论的这一方面。

然而，从数学角度理解类型论，一直以来都有一个困难：\emph{类型}这一基本概念与\emph{集合}的概念有所不同，而差异所在之处一直难以精确刻画。
我们相信，以下新想法是重大的一步：不把类型视为奇怪的集合（也许是不使用经典逻辑构造出的集合），而是从同伦论的视角将类型视为空间。
特别地，它解决了如下问题：类型元素的相等概念与集合元素的相等概念有何不同。

在同伦论中，我们关注空间
\index{topological!space}%
及其间的连续映射，
\index{function!continuous!in classical homotopy theory}%
并在同伦的意义下考察它们。
一对连续映射 $f : X \to Y$ 与 $g : X\to Y$ 之间的一个\emph{同伦（homotopy）}
\index{homotopy!topological}%
是一个满足 $H(x, 0) = f (x)$ 和 $H(x, 1) = g(x)$ 的连续映射 $H : X \times [0, 1] \to Y$。
同伦 $H$ 可以看作将 $f$ “连续形变”为 $g$ 的过程。
如果存在来回的连续映射，其复合分别同伦于相应的恒等映射，换言之，如果它们“在相差同伦的意义下”同构，则称空间 $X$ 和 $Y$ 是\emph{同伦等价（homotopy equivalent）}的，
\index{homotopy!equivalence!topological}%
记作 $\eqv X Y$。
同伦等价的空间具有相同的代数不变量（例如同调或基本群），并被称为具有相同的\emph{同伦类型（homotopy type）}。

\subsection*{同伦类型论}

同伦类型论（HoTT）从同伦的视角解释类型论。
在同伦类型论中，我们将类型视为“空间”（如同在同伦论中所研究的）或高阶群胚，而将逻辑构造（例如积 $A\times B$）视为在这些空间上的同伦不变构造。
如此一来，我们便能直接处理空间，无需先发展点集拓扑学（或任何组合替代物，如单纯集合论）。
为简要解释这一视角，先考虑类型论的基本概念，即\emph{项（term）} $a$ 属于\emph{类型} $A$，写作：
\[ a:A. \]
传统上，这个表达式被认为类似于：

\begin{center}
``$a$ 是集合 $A$ 的一个元素''。
\end{center}

然而，在同伦类型论中，我们将其理解为：

\begin{center}
``$a$ 是空间 $A$ 中的一个点''。
\end{center}

\index{continuity of functions in type theory@``continuity'' of functions in type theory}%
类似地，类型论中的每个函数 $f : A\to B$ 都被视为从空间 $A$ 到空间 $B$ 的一个连续映射。

我们必须强调，这里所说的``空间''是在纯同伦意义下（而非拓扑意义下）处理的。
例如，类型既没有``开子集''的概念，类型中元素的序列也没有``收敛''的概念。
我们只有``同伦意义下的''概念，比如点之间的道路和道路之间的同伦，这些概念在同伦论的其他模型（例如单纯集）中同样有意义。
因此，更准确的说法是：我们将类型视为\emph{$\infty$-群胚（$\infty$-groupoids）}\index{.infinity-groupoid@$\infty$-groupoid}；这一名称指的是同伦论中的``不变对象''——它们可以由拓扑空间、
\index{topological!space}%
单纯集或同伦论的任何其他模型来呈现。
不过，有时使用``空间''、``道路''这类拓扑词汇是方便的，只要我们记住其他拓扑概念并不适用即可。

（也难免让人想用\emph{同伦类型}
\index{homotopy!type}%
这个短语来指称这些对象，因为它暗示了双重解释：``（类型论中的）一个类型，从同伦角度看待''与``从同伦论角度考虑的一个空间''。
后一种解释与``同伦类型''的经典含义——即空间在模同伦等价意义下的\emph{等价类}——略有不同，尽管它确实保留了诸如``这两个空间具有相同的同伦类型''这类短语的意义。）

将类型解释为结构化对象（而非集合）的想法由来已久，且已被证明能够澄清类型论中诸多令人困惑的方面。
例如，将类型解释为层有助于解释类型论逻辑的直觉主义本质，而将其解释为偏等价关系或``域''则有助于解释其计算方面。
这也意味着，我们可以运用类型论推理来研究这些结构化对象，从而催生出范畴逻辑这一丰富的领域。
同伦解释符合同一模式：它澄清了类型论中\emph{相等}（或等同）的本质，并允许我们在同伦论研究中运用类型论推理。

同伦解释的关键新思想在于：两个对象 $a, b: A$（属于同一类型 $A$）在逻辑上的相等 $a = b$，可以理解为空间 $A$ 中从点 $a$ 到点 $b$ 存在一条道路 $p : a \leadsto b$。
这也意味着，如果两个函数 $f, g: A\to B$ 是同伦的，它们就可以被视为等同，因为同伦不过是 $B$ 中的一族（连续）道路 $p_x: f(x) \leadsto g(x)$——每个 $x:A$ 对应一条。
在类型论中，对于每个类型 $A$，都存在一个（曾经有些神秘的）类型 $\idtypevar{A}$，表示 $A$ 中两个对象的相等证明；在同伦类型论中，这恰恰就是\emph{道路空间} $A^I$，即所有从单位区间到 $A$ 的连续映射 $I\to A$。
\index{unit!interval}%
\index{interval!topological unit}%
\index{path!topological}%
\index{topological!path}%
这样，一个项 $p : \idtype[A]{a}{b}$ 就代表了 $A$ 中的一条道路 $p : a \leadsto b$。

同伦类型论的思想大约在 2006 年产生，源于 Awodey 和 Warren~\cite{AW} 与 Voevodsky~\cite{VV} 的独立工作，但其灵感来自 Hofmann 和 Streicher 更早的群胚解释~\cite{hs:gpd-typethy}。
事实上，正如 Grothendieck（格罗滕迪克）所提出、并正由两个领域的数学家深入开展的那样，高阶范畴论（特别是弱 $\infty$-群胚理论）如今已知与同伦论有着深刻的联系。
Awodey–Warren 和 Voevodsky 最初的语义模型使用了同伦论中若干广为人知的概念与技术，如 Quillen 模型范畴和 Kan\index{Kan complex} 单纯集\index{simplicial!sets}，这些概念如今也应用于高阶范畴论。
\index{Quillen model category}%
\index{model category}%

特别地，Voevodsky 在 Kan 单纯集中构建了一种类型论解释，并认识到该解释满足一个他称之为\emph{泛等性}（univalence）的额外关键性质。
这一性质此前在类型论中并未被考虑（尽管 Church 的命题外延性原理被证明是它的一个非常特殊的情形，Hofmann 和 Streicher 也曾在``宇宙外延性''的名下考虑过另一个特例）。
以新公理的形式将泛等性引入类型论，会产生深远的影响，其中许多是自然的、简化的且令人信服的。
泛等公理也进一步强化了类型论的同伦观点，因为它在单纯集模型及其他相关模型中成立，而在将类型视为集合的观点下则不成立。

\subsection*{泛等基础}

泛等公理的基本思想可以非常简要地解释如下。
在类型论中，可以存在一种类型，其元素本身也是类型；这样的类型被称为\emph{宇宙（universe）}，通常记作 $\UU$。
那些作为 $\UU$ 的项的类型，通常称为\emph{小类型}。\index{type!small}%
\index{small!type}%
与任何类型一样，$\UU$ 也有自己的相等类型 $\idtypevar{\UU}$，它表达了小类型之间的相等关系 $A = B$。
把类型设想为空间，则 $\UU$ 也是一个空间，其中的点是一个个空间；要理解它的相等类型，我们就要问：$\UU$ 中两个空间 $A$ 与 $B$ 之间的一条道路 $p : A \leadsto B$ 究竟是什么？
泛等公理断言，这样的道路对应于 $A$ 与 $B$ 之间的同伦等价 $\eqv A B$，大致如上文所述。
更精确地说：给定任意（小）类型 $A$ 和 $B$，除了原始的、表示 $A$ 与 $B$ 的相等证明的类型 $\idtype[\UU]AB$ 之外，还有一个由定义给出的、从 $A$ 到 $B$ 的等价的类型 $\texteqv AB$。
由于任何对象上的恒等映射都是一个等价，因此存在一个规范的映射，
\[\idtype[\UU]AB\to\texteqv AB.\]
泛等公理断言，这个映射本身就是一个等价。
冒着过度简化的风险，我们可以把它简明地表述为：

\begin{description}\index{univalence axiom}%
\item[泛等公理：] $\eqvspaced{(A = B)}{(\eqv A B)}$。
\end{description}

%
换言之，相等即等价。\index{identity}% 
特别地，人们可以说``等价的类型是等同的''。
然而，这种说法有一定的误导性：它听上去像是某种``骨架性''条件，将等价概念\emph{坍缩}为相等，而事实上泛等公理所做的是\emph{拓展}相等概念，使之与（未改变的）等价概念相吻合。

从同伦的观点看，泛等公理意味着具有相同同伦类型的空间在宇宙 $\UU$ 中由一条道路相连，这契合了将 $\UU$ 视为（小）空间的分类空间的直观。
但从逻辑的观点看，这是一个革命性的新思想：它表明同构的事物可以被等同！
数学家们当然习惯于在实践中将同构的结构视为等同，但他们通常是通过``记号滥用''\index{abuse!of notation}或其他非正式的手段来实现的，心里清楚所涉及的对象并非``真的''相同。
但在这一新的基础方案中，这样的结构可以被形式化地等同，其逻辑意义在于：涉及其中一个的每一个性质或构造，也适用于另一个。
事实上，这种等同现在是显式的，并且性质与构造可以沿着它系统地迁移。
此外，做出这种等同的\emph{不同方式}本身也构成一种结构，人们可以（而且应该！）将其纳入考量。

总而言之，对于宇宙 $\UU$ 中的点 $A$ 和 $B$（即小类型），泛等公理将以下三个概念等同起来：

\begin{itemize}
\item（逻辑上）$A$ 与 $B$ 的一个相等证明 $p:A=B$
\item（拓扑上）$\UU$ 中从 $A$ 到 $B$ 的一条道路 $p:A \leadsto B$
\item（同伦上）$A$ 与 $B$ 之间的一个等价 $p:\eqv A B$。
\end{itemize}

\subsection*{高阶归纳类型}\index{type!higher inductive}%

类型论的经典优势之一，在于它有一套简单而有效的方法来处理归纳定义的结构。
最简单且非平凡的归纳定义结构是自然数，它由零和后继函数归纳生成。
从这一表述出发，我们可以算法式地提取出刻画自然数的数学归纳原理。
更一般的归纳定义涵盖了各式各样的列表与良基树，每一种都由相应的``归纳原理''所刻画。
这包括了某些编程语言中使用的大多数数据结构；因此，类型论在对这些编程语言进行形式推理时大有用处。
如果从非常广义的角度来理解，归纳定义还包括像不相交并 $A+B$ 这样的例子，它可以被视为由两个单射 $A\to A+B$ 和 $B\to A+B$ ``归纳''生成的。此时的``归纳原理''就是``分情形讨论证明''，它刻画了不相交并。

在同伦论中，我们很自然地也会考虑那些``归纳定义空间''，它们不仅由一组\emph{点}生成，而且还由一组\emph{道路}和更高阶的道路生成。
经典上，这类空间被称为\emph{CW 复形（CW complexes）}。
\index{CW complex}%
例如，圆 $S^1$ 由一个点以及一条从该点出发回到自身的道路生成。
类似地，2-球面 $S^2$ 由一个点 $b$ 和一条从 $b$ 处的常值道路到其自身的二维道路生成；而环面 $T^2$ 则由一个点、两条从该点到自身的道路 $p$ 和 $q$，以及一条从 $p\ct q$ 到 $q\ct p$ 的二维道路生成。

通过将道路与同伦类型论中的相等等同起来，这类``归纳定义空间''就可以在类型论中用``归纳原理''来刻画，其方式完全类似于自然数和不交并这样的经典例子。
由此产生的 \emph{高阶归纳类型（higher inductive types）}
\index{type!higher inductive}%
为我们提供了一种直接的``逻辑''方式来对球面等熟悉空间进行推理，这种方式（与泛等性结合后）可以用于以纯粹形式化的方式，进行同伦论中那些熟悉的论证，例如计算球面的同伦群。
这样得到的证明，是经典同伦论思想与经典类型论思想的一次融合，为这两个学科都带来了新的见解。

而且，这还只是冰山一角：同伦论中的许多抽象构造，例如同伦余极限、纬悬、Postnikov（波斯特尼科夫）塔、局部化、完备化与谱化，也都可以表示为高阶归纳类型。
其中有许多在经典做法中是利用 Quillen（奎伦）的``小对象论证来构造的；这种论证可以看作是以有限的方式，通过算法来描述一个空间的无穷 CW 复形表示\index{presentation!of a space as a CW complex}，正如``零和后继''是对无穷自然数集合的一种有限算法\index{algorithm}描述一样。
通过小对象论证得到的空间是出了名的复杂和难以理解；而类型论的方法则可能简单得多——它无需任何显式构造，而是让人可以直接使用相应的``归纳原理''。
因此，泛等性与高阶归纳类型的结合，暗示着同伦论的实践在某种意义上可能迎来一场革命。

\subsection*{泛等基础中的集合}

\index{set|(}%

我们曾声称，泛等基础最终能够为``所有''数学提供一个基础；但到目前为止，我们讨论的还仅仅是同伦论这一分支。
当然，即使不借助同伦类型论的新特征，也已经有很多利用类型论来形式化数学的具体实例，
\index{mathematics!formalized}%
\index{theorem!Feit--Thompson}%
\index{theorem!odd-order}%
\index{Feit--Thompson theorem}%
\index{odd-order theorem}%
例如最近在 \Coq~\cite{gonthier} 中对 Feit–Thompson 奇阶定理的形式化工作。

但传统观点认为，数学建立在集合论之上，其含义是：所有数学对象与构造都可以编码进诸如 Zermelo–Fraenkel 集合论（ZF）这样的理论中。
\index{set theory!Zermelo--Fraenkel}%
\indexsee{Zermelo-Fraenkel set theory}{set theory}%
\indexsee{ZF}{set theory}%
\indexsee{ZFC}{set theory}%
然而，如今已经公认：对于集合论自身领域之外的大部分数学而言，ZF 集合中那种错综复杂的层次成员结构其实并无必要；一个更``结构性的''理论，比如 Lawvere\index{Lawvere}（劳维尔）的《集合范畴初等理论》~\cite{lawvere:etcs-long}，就已经足够。
\index{Elementary Theory of the Category of Sets}%

在泛等基础中，基本的对象是``同伦类型''而非集合，但我们可以\emph{定义}一类行为像集合的类型。
从同伦的角度看，它们可以看作是每个连通分支都可缩的空间，也就是说，是同伦等价于离散空间的那些空间。
\index{discrete!space}%
定理表明，由这样的``集合''所构成的范畴满足 Lawvere\index{Lawvere} 的公理（或与之相关的公理，具体取决于理论的细节）。
因此，任何能够用一种类似 ETCS 的理论来表示的数学（经验表明，这基本上就是所有数学），同样可以在泛等基础中得到表示。

这支持了如下主张：泛等基础至少不逊于现有的数学基础。
在泛等基础中工作的数学家，可以用一种熟悉的方式从集合出发来构造结构；而更一般的同伦类型则留在基础背景中待命，待到需要时便可取用。
正因如此，本书所选的大部分应用领域，都是泛等基础能做出一些\emph{崭新}贡献、从而区别于现有基础体系的领域。

不出所料，同伦论与范畴论便是其中的两个；但或许不那么显而易见的是，即使在集合论与实分析这样的学科中，泛等基础也能提供新颖有趣的见解。
例如，泛等公理允许我们将同构的结构视为等同，而高阶归纳类型则允许我们直接通过对象的泛性质来描述它们。
这样一来，我们通常就可以避免诉诸随意选出的代表元或超限的迭代构造。
事实上，即使是在泛等基础的集合内部，我们也可以用这样一种归纳的泛性质来刻画 ZF 集合论的研究对象。

\index{set|)}%

\subsection*{非形式化类型论}

\index{mathematics!formalized|(defstyle}%
\index{informal type theory|(defstyle}%
\index{type theory!informal|(defstyle}%
\index{type theory!formal|(}%
习惯于经典数学的学者，在学习类型论时常常会遇到一个困难：类型论通常以完全或部分形式化的演绎系统的形式呈现。
这种风格对于证明论研究非常有用，但在应用性的、非形式化的推理中却不太方便。
甚至对于大多数从事研究的数学家来说，这种风格也相当陌生——即使他们可能对数学基础感兴趣。
本书的目标之一，就是要发展一种在泛等基础上进行数学工作的非形式化风格，它既要保持严格与精确，又要更贴近日常数学的语言与表述风格。

在当今的数学中，人们通常以一种原则上可以形式化的方式来构造数学对象并进行推理——人们相信，这可以基于一种初等集合论系统（比如 ZFC）来完成，至少在具备足够的巧思与耐心的前提下。
在大多数情况下，人们甚至无需意识到这种可能性，因为它在很大程度上与``证明完全严谨''的要求相一致（即所有数学家通过教育和经验而直观理解的那种严谨）。
但人们确实需要学会对``非形式化的集合论''中的几个方面保持谨慎：使用过大或尚未成形而无法构成集合的搜集；选择公理及其等价形式；甚至（对本科生而言）反证法；等等。
采用像同伦类型论这样的新基础系统，作为非形式化推理的\emph{隐含形式基础}，将需要我们调整自己的一些本能与习惯。
本书旨在成为这种``新型数学''的一个例子：它仍然是非形式化的，但在原则上如今可以基于同伦类型论而非 ZFC 来形式化——当然，同样需要足够的巧思与耐心。

值得强调的是，在这个新系统中，这样的形式化能够带来真正的实用价值。
类型论的形式系统非常适合计算机系统，并且已经在现有的证明助理中得到了实现。
\index{proof!assistant}%
证明助理是一种计算机程序，它引导用户构建完全形式化的证明，只允许有效的推理步骤。
它还能提供一定程度的自动化，可以搜索定理库中的现有结论，甚至能从最终的（构造性）证明中提取出数值算法\index{algorithm} \index{extraction of algorithms}。

我们相信，泛等基础计划的这一特点将其与其他基础方案区分开来，并可能为一线数学家带来新的实用价值。
事实上，基于较早期类型论的证明助理，已经被用来形式化一些重大的数学证明，例如四色定理\index{theorem!four-color} \index{four-color theorem}以及 Feit--Thompson 定理。
泛等基础的计算机实现，当前同理论本身一样，尚在发展之中。
\index{proof!assistant}%
然而，即便是目前可用的实现（它们大多是对现有证明助理，如 \Coq 和 \Agda，进行小幅修改后的版本）也已经展现了自身的价值，不仅体现在对已知证明的形式化上，还体现在新证明的发现上。
事实上，本书中描述的许多证明，其\emph{首次}完成实际上是在证明助理中以完全形式化的形式进行的，直到现在才第一次被``去形式化''地呈现出来——这逆转了形式数学与非形式数学之间通常的关系。

我们可以想象，在不久的将来，数学家将有可能在由证明助理形式化的泛等基础系统内部工作，从而验证自己论文的正确性；并且，这样做会变得如同用 \TeX 为自己的论文排版一样自然。
%(Whether this proves to be the publishers' dream or their nightmare remains to be seen.) 
原则上，这对任何其他基础系统也同样成立，但我们相信，使用泛等基础能更切合实际地实现这一点；眼前的这部著作及其形式化的对应版本便是见证。

\index{type theory!formal|)}%
\index{informal type theory|)}%
\index{type theory!informal|)}%
\index{mathematics!formalized|)}%

\subsection*{构造性}

\index{mathematics!constructive|(}%

经典\index{mathematics!classical}基础与类型论之间最显著的区别之一，是\emph{证明相关性（proof relevance）}这一观念。根据这一观念，数学陈述、甚至其证明本身，都成为了第一类数学对象。
在类型论中，我们用类型来表示数学陈述，这些类型可以同时被视作数学构造和数学断言，这一概念也被称为\emph{命题即类型（propositions as types）}。
\index{proposition!as types}%
相应地，我们可以把一个项 $a : A$ 同时看作类型 $A$ 的一个元素（在同伦类型论中即空间 $A$ 的一个点），同时也是命题 $A$ 的一个证明。
举个例子：假设有集合 $A$ 和 $B$（离散空间），
\index{discrete!space}%
我们来考虑陈述``$A$ 同构于 $B$''。
在类型论中，这可以表达为：
\begin{narrowmultline*}
  \mathsf{Iso}(A,B) \defeq \narrowbreak
  \sm{f : A\to B}{g : B\to A}\Big(\big(\tprd{x:A} g(f(x)) = x\big) \times \big(\tprd{y:B}\, f(g(y)) = y\big)\Big).
\end{narrowmultline*}
%
将这里的类型构造子 $\Sigma, \Pi, \times$ 分别读作``存在''、``对所有''和``且''，就得到了``$A$ 与 $B$ 同构''这一命题的通常表述；另一方面，将它们读作和与积，得到的则是 $A$ 与 $B$ 之间\emph{所有同构的类型}！要证明 $A$ 与 $B$ 同构，就要构造一个证明 $p : \mathsf{Iso}(A,B)$，这等同于在 $A$ 与 $B$ 之间构造一个同构，亦即给出一对函数 $f, g$，并附上\emph{证明}表明它们的复合是相应的恒等映射。而后面的这些证明，反过来也无非是相应种类的同伦。这样一来，\emph{证明一个命题就等同于构造某个特定类型的元素}。
特别地，要证明一个形如``$A$ 且 $B$''的陈述，只需分别证明 $A$ 和 $B$，即给出类型 $A\times B$ 的一个元素。
而要证明 $A$ 蕴涵 $B$，只需找到 $A\to B$ 的一个元素，即从 $A$ 到 $B$ 的一个函数（它决定了从 $A$ 的证明到 $B$ 的证明的映射）。

“命题即类型”的逻辑是灵活的，它支持多种变体，例如只使用类型的某个子类来表示命题。
在同伦类型论中，存在一些自然的此类子类，其根源在于：所有类型构成的系统，正如经典同伦论中的空间一样，是按照其高阶同伦结构存在或退化的维度来“分层”的。
具体而言，Voevodsky 找到了\emph{同伦 $n$-类型（homotopy $n$-types）}的纯类型论定义，它对应于在维度 $n$ 以上没有非平凡同伦信息的空间。
（$0$-类型就是前面提到的、满足 Lawvere 公理\index{Lawvere}的“集合”。）
此外，利用高阶归纳类型，我们可以一般地将一个类型“截断”为 $n$-类型；在经典同伦论中，这对应于它的第 $n$ 个 Postnikov\index{Postnikov tower} 截面。\index{n-type@$n$-type}
对逻辑而言，特别重要的是同伦 $(-1)$-类型的情形，我们称之为\emph{命题}。
在经典视角下，每个 $(-1)$-类型要么是空的，要么是可缩的；我们将这两种可能性分别解释为真值“假”和“真”。

将所有类型作为命题，会得到一种非常“构造性”的逻辑观；更多相关内容参见~\cite{kolmogorov,TroelstraI,TroelstraII}。
例如，每一个关于某物存在的证明都自带足够的信息，足以实际找到所存在的对象；并且，每一个关于“$A$ 或 $B$”成立的证明，要么是 $A$ 成立的证明，要么是 $B$ 成立的证明。
因此，从每个证明中我们都能自动提取出一个算法；\index{algorithm} \index{extraction of algorithms} 这在计算机编程的应用中非常有用。

然而另一方面，这种逻辑确实偏离了传统数学对存在性证明的理解。
特别地，它并不能忠实表达某些重要的经典推理原则，例如选择公理和排中律。
为此，我们需要使用“$(-1)$-截断”逻辑，其中只有同伦 $(-1)$-类型才代表命题。

\index{axiom!of choice}%
更具体地说，一方面考虑\emph{选择公理}：“若对每个 $x: A$，存在一个 $y:B$ 使得 $R(x,y)$，则存在一个函数 $f : A\to B$，使得对所有 $x:A$ 都有 $R(x, f(x))$。”
纯“命题即类型”意义下的“存在”概念足够强，足以使该陈述直接得证——但它并不具备通常选择公理的所有推论。
然而，在 $(-1)$-截断逻辑中，这个陈述并非自动成立，而是一个强假设，具有与其在经典\index{mathematics!classical}集合论中的对应物相同的那些推论。

\index{excluded middle}%
\index{univalence axiom}%
另一方面，考虑\emph{排中律}：“对所有 $A$，要么 $A$ 成立，要么非 $A$ 成立。”
在纯“命题即类型”逻辑下解释该陈述，会得到一个与泛等公理不一致的命题。
因为证明“$A$”意味着展示它的一个元素，所以该假设将给出一种从每个非空类型中选取元素的统一方法——某种Hilbert 式的选择算子。
泛等性意味着，这样的选择算子所选出的 $A$ 的元素，必须在 $A$ 的所有自等价下保持不变，因为这些自等价被等同于自相等，而每个操作都必须尊重相等；但显然，有些类型存在无不动点的自同构，例如我们可以交换一个二元类型中的两个元素。
\index{automorphism!fixed-point-free}%
然而，“$(-1)$-截断排中律”虽非自动成立，却可以被一致地假定，且具有与经典数学中大体相同的推论。

换句话说，纯粹的``命题即类型''逻辑的``构造性''是上文提到的强算法意义上的，而默认的 $(-1)$-截断逻辑的``构造性''则是另一种意义上的（即 Heyting 以``直觉主义''为名形式化的那种逻辑）；
对于后者，我们可以自由地添加选择公理和排中律，从而得到一种可以称为``经典''的逻辑。
因此，同伦类型论同时兼容构造性与经典这两种逻辑观念，甚至远不止于此。
\index{logic!constructive vs classical}%
确实，同伦视角揭示出，经典逻辑与构造逻辑可以共存，分别作为一系列不同系统的两个端点，其间存在着无限多种可能性（即满足 $-1 < n < \infty$ 的同伦 $n$-类型）。
我们可以谈论``\LEM{n}''和``\choice{n}''，其中 $\choice{\infty}$ 是可证的，而 \LEM{\infty} 与泛等性不相容；
而 $\choice{-1}$ 和 $\LEM{-1}$ 则是经典数学家们所熟悉的版本（因此，在大多数情况下，当下标未给出时，假定其为 $(-1)$ 是恰当的）。
事实上，甚至可以有这样的有用系统：其中只有\emph{某些}类型满足这些进一步的``经典''原则，而一般类型仍保持``构造性''。\index{excluded middle}\index{axiom!of choice}%%

值得强调的是，泛等基础并不\emph{要求}使用构造性或直觉主义逻辑。\index{logic!intuitionistic}\index{logic!constructive} %
只需假定这两个原则成立（以其适当的、$(-1)$-截断的形式），大部分依赖排中律和选择公理的经典数学都可以在泛等基础中开展。
然而，类型论确实鼓励我们在不需要时避免使用这些原则，理由如下。

首先，每位数学家都知道，证明定理时使用的假设越少，定理就越强，因为它适用于更多的例子。
\choice{} 和 \LEM{} 的情况并无不同：类型论容许许多有趣的``非标准''模型，例如在层意象中，\choice{} 和 \LEM{} 这类经典性原则往往不成立。
\index{topos}%
同伦类型论在高阶意象中也容许类似的模型，如~\cite{ToenVezzosi02,Rezk05,lurie:higher-topoi} 中所研究的那些。
因此，如果我们避免使用这些原则，我们所证明的定理将在所有这些模型的内部均有效。

其次，类型论的另一个优点是它的可计算特性。
除了作为数学基础，类型论也是一种计算的形式理论，也可以视为一种强大的编程语言。
\index{programming}%
从这个角度看，系统的规则不能像集合论公理那样随意选取：它们之间必须有一种和谐，使得所有证明都能作为程序``执行''。
对于同伦类型论引入的新原理，例如泛等性和高阶归纳类型，从这个角度看，我们尚未完全理解，但其基本轮廓已逐渐显现；例如，可参见~\cite{lh:canonicity}。
然而，人们早就知道，像 \choice{} 和 \LEM{} 这样的原则从根本上就与可计算性相抵触，因为它们径直断言某些对象存在，却不给出任何计算它们的方法。
因此，要维持类型论作为计算理论的特性，避免使用这些原则是必要的。

幸运的是，构造性推理并不像看起来那么困难。
在某些情况下，只需重新表述一些定义，就可以使一个定理成为构造性的，同时其证明也会更加优雅。
而且，在泛等基础中，这种情况似乎更常发生。
例如：

\begin{enumerate}
\item 在集合论基础中，同伦论和范畴论中在多处需要借助选择公理来完成超限构造。
  但有了高阶归纳类型，我们就能直接且构造性地将这些构造编码。
  特别地，\cref{cha:homotopy} 中所有的``综合''同伦论都不需要 \LEM{} 或 \choice{}。
\item 在集合论基础中，命题``每个全忠实且本质满的函子都是范畴的等价''等价于选择公理。
  但在泛等公理下，它直接就是\emph{真的}；参见 \cref{cha:category-theory}。
\item 在集合论中，若要分别定义出能够典范地代表集合同构类与良序集同构类的``基数''与``序数''概念，需要经过各种迂回的处理——可能还要用到选择公理或良基公理。
  但有了泛等性与高阶归纳类型，我们可以直接对宇宙进行截断来获得这样的代表元；参见 \cref{cha:set-math}。
\item 在集合论基础中，将实数定义为柯西序列的等价类，要么需要排中律，要么需要（可数）选择公理，才能保证其行为良好。
  但有了高阶归纳类型，我们可以给出这个定义的一个版本，它行为良好且避免了任何选择原则；参见 \cref{cha:real-numbers}。
\end{enumerate}

当然，这些简化也完全可以视为一种证据，表明新方法最终未必是真正构造性的。
然而，我们要再次强调，阅读本书的读者完全不必在意或担心构造性问题。
关键在于，以上所有例子中，我们所给出的理论版本都具有独立的优势，无论是否假定 \LEM{} 与 \choice{} 可用。
倘若能够达到构造性，那将是一份额外的红利。\index{constructivity}%

既然讨论了添加泛等性、高阶归纳类型、\choice{} 和 \LEM{} 等新原则，人们可能会问，由此产生的系统是否仍然是一致的。
（类型论相对于集合论的一个原有优点是，可以通过证明论手段验证其一致性。）
与任何基础系统一样，一致性\index{consistency}是一个相对的问题：``相对于什么而言是一致的？''
简短的答案是，由于 Voevodsky 的工作~\cite{klv:ssetmodel}（关于高阶归纳类型，参见~\cite{ls:hits}），本书中考虑的所有构造与公理在 Kan\index{Kan complex} 复形范畴中都有一个模型。
因此，已知它们相对于 ZFC 是一致的（只需加入足够多的不可达基数
\index{inaccessible cardinal}\index{consistency}%
以提供我们所需的嵌套泛等宇宙）。
以更传统的类型论方式来解释这一一致性，是一项正在进行的工作（例如，参见~\cite{lh:canonicity,coquand2012constructive}）。

我们在 \cref{tab:pov} 中总结了关于类型论操作的不同观点。

\begin{table}[htb]
  \centering
  \OPTsmalltable
 \begin{tabular}{lllll}
    \toprule
       类型 && 逻辑 & 集合 & 同伦\\ \addlinespace[2pt]
    \midrule
       $A$ && 命题 & 集合 & 空间\\ \addlinespace[2pt]
       $a:A$ && 证明 & 元素 & 点 \\ \addlinespace[2pt]
       $B(x)$ && 谓词 & 集合族 & 纤维化 \\ \addlinespace[2pt]
       $b(x) : B(x)$ && 条件证明 & 元素族 & 截面\\ \addlinespace[2pt]
       $\emptyt, \unit$ && $\bot, \top$ & $\emptyset, \{ \emptyset \}$ & $\emptyset, *$\\ \addlinespace[2pt]
       $A + B$ && $A\vee B$ & 不交并 & 余积\\ \addlinespace[2pt]
       $A\times B$ && $A\wedge B$ & 序对的集合 & 积空间\\ \addlinespace[2pt]
       $A\to B$ && $A\Rightarrow B$ & 函数集合 & 函数空间\\ \addlinespace[2pt]
       $\sm{x:A}B(x)$ &&  $\exists_{x:A}B(x)$ & 不交和 & 全空间\\ \addlinespace[2pt]
       $\prd{x:A}B(x)$ &&  $\forall_{x:A}B(x)$ & 积 & 截面的空间\\ \addlinespace[2pt]
       $\mathsf{Id}_{A}$ && 相等关系 $=$ & $\setof{\pairr{x,x} | x\in A}$ & 道路空间 $A^I$ \\ \addlinespace[2pt]
    \bottomrule
  \end{tabular}
  \caption{类型论操作的不同观点比较}\label{tab:pov}
\end{table}

\index{mathematics!constructive|)}%

\subsection*{待解问题} 

\index{open!problem|(}%

对于那些有志于为这一数学新分支贡献力量的读者，知道这一领域存在许多有趣的待解问题，或许会令人鼓舞。

\index{univalence axiom!constructivity of}%
其中最紧迫的问题，或许就是 Voevodsky 在 \cite{Universe-poly} 中提出的关于泛等公理的``构造性''问题。
类型论的基本体系遵循 Gentzen 的自然演绎结构。
逻辑连接词由其引入规则定义，其消去规则则由计算规则来保证。
沿用这一模式，并运用 Tait 的可计算性方法（最初为分析 Gödel 的 Dialectica 解释而设计），可以证明类型论的\emph{典范化（normalization）}性质。
这一性质进而保证了若干重要性质，例如类型检查的可判定性（这是一条关键性质，因为类型检查对应于证明检查，而我们理应做到``看到证明时就能辨认出它''），以及所谓的``典范性\index{canonicity}性质''——即自然数类型的任意封闭项都可归约为一个数字。
当我们用一条公理（例如函数外延性公理或泛等公理）来扩展类型论时，上述最后一条性质以及引入/消去规则的统一结构便会丧失。
Voevodsky 针对加入泛等公理的类型论的典范性问题，提出了一个精确的数学猜想：给定一个自然数类型的封闭项，是否总能找到一个数字及一个证明，表明该项与该数字相等——其中这个相等证明本身可以使用泛等公理？
更一般地，一个重要问题是：能否为泛等公理提供构造性的证成？
如果再加入其他由同伦动机驱动的构造，例如高阶归纳类型，又会如何？
这些问题至今仍是悬而未决的，尽管目前人们正在发展各种方法以尝试寻求答案。

另一个基本问题是，在处理诸如自然数这样本质上是集合（即离散空间）的类型时所遇到的困难，
\index{discrete!space}%
它们只包含平凡的道路。
目前，同伦类型论实际上只能将空间刻画到同伦等价的程度，这意味着上述``离散空间''可能只是\emph{同伦等价}于离散空间。
从类型论的角度看，这意味着存在许多``相等''于自反性的道路，但并非在\emph{判断相等}的意义上与自反性相等（``判断相等''的含义见 \cref{sec:types-vs-sets}）。
虽然这种同伦不变性有其优势，但这些``无意义''的相等项确实给论证和构造引入了不必要的复杂性，因此，若能有一套系统的方法来消除或坍缩它们，将会十分便利。
% In some cases, the proliferation of such superfluous identity terms makes it very difficult or impossible to formulate what should be a straightforward concept, such as the definition of a (semi-)simplicial type.

一个更为专门但同样重要的问题，是同伦类型论与\emph{高阶意象（higher toposes）}%
\index{.infinity1-topos@$(\infty,1)$-topos}
研究之间的关系；后者目前正活跃于高阶范畴论与同伦论的交叉领域。
熟悉这两个领域的研究者越来越相信，二者之间存在着深刻的联系。
例如，泛等宇宙的概念应与\emph{对象分类子（object classifier）}的概念相一致，而高阶归纳类型则应是\emph{局部可表现性（local presentability）}的一种``初等''反映。
更一般地说，同伦类型论应该是 $(\infty,1)$-意象的``内部语言''，正如直觉主义高阶逻辑是普通 1-意象的内部语言一样。
然而，尽管存在这种普遍共识，细节仍有待厘清——尤其是一致性与严格性的问题尚需处理——而完成这项工作无疑将使我们对这两个概念都获得更深的洞见。

\index{mathematics!formalized}%
但迄今为止最广阔的待办工作，是在这一新系统中对日常数学进行持续的形式化。
近期在形式化同伦论与范畴论的一些基本事实方面取得的成效令人鼓舞；其中部分成果在 \cref{cha:homotopy,cha:category-theory} 中有所介绍。
不过显而易见，仍有大量工作有待完成。

\index{open!problem|)}%

同伦类型论社区维护着一个网站与集体博客 \url{http://homotopytypetheory.org}，并设有讨论邮件列表。
随时欢迎新来者！

\subsection*{如何阅读本书}

本书分为两个部分。
\cref{part:foundations}，``基础篇''，发展同伦类型论的基本概念。
这是构建具体学科发展的数学基础，也是理解泛等基础方法所必需的。对程序员而言，这就是``库代码''。
由于泛等基础是一种全新且不同的数学，其基本概念需要一些时间来适应；因此 \cref{part:foundations} 的篇幅相当可观。

\cref{part:mathematics}，``数学篇''，由四章组成，它们在 \cref{part:foundations} 基本概念的基础之上，展示我们在四个不同的数学领域中能够用泛等基础做的新事情：同伦论（\cref{cha:homotopy}）、范畴论（\cref{cha:category-theory}）、集合论（\cref{cha:set-math}）和实分析（\cref{cha:real-numbers}）。
\cref{part:mathematics} 的各章之间大体相互独立，尽管偶尔会用到在别章中证明的引理。

想要认真理解泛等基础并能在其中开展工作的读者，最终需要阅读并理解 \cref{part:foundations} 的大部分内容。
然而，对于只想对泛等基础及其能力略作了解的读者来说，在到达 \cref{part:mathematics} 的``实质内容''之前，必须先通读超过两百页的基础部分，这难免令人望而却步。
幸运的是，要阅读 \cref{part:mathematics} 的各章，并不需要 \cref{part:foundations} 的全部内容。
\cref{part:mathematics} 的每一章都以简短的概述开头，介绍其学科主题、泛等基础能为之贡献什么，以及来自 \cref{part:foundations} 的必要背景知识，因此有勇气的读者可以直接翻到他们最感兴趣的章节。
对于那些想比这更深入地理解 \cref{part:mathematics} 中一个或多个章节、但尚未准备好通读 \cref{part:foundations} 的读者，我们在此提供 \cref{part:foundations} 的简要小结，并注明其中哪些部分对 \cref{part:mathematics} 的哪些章节是必要的。

\cref{cha:typetheory} 讲述类型论的基本概念，不涉及任何同伦解释。
熟悉 Martin-L\"of 类型论的读者可以快速浏览，了解我们所用理论的具体细节。
然而，没有类型论经验的读者需要阅读 \cref{cha:typetheory}，因为类型论与集合论等其他基础之间存在着许多微妙的差异。

\cref{cha:basics} 引入了对类型论的同伦视角及支撑这一视角的基本概念，并描述了 \cref{cha:typetheory} 中类型论各组成部分的同伦行为。
它还引入了\emph{泛等公理（univalence axiom）}（\cref{sec:compute-universe}）——同伦类型论两项基本创新中的第一项。
因此，这部分内容非常基础，我们鼓励每位读者阅读，特别是 \crefrange{sec:equality}{sec:basics-equivalences}。

\cref{cha:logic} 描述我们如何在同伦类型论中表达逻辑，以及它与经典逻辑、构造逻辑和直觉主义逻辑之间的联系。
在此，我们定义了排中律、选择公理和命题大小重整公理（尽管在全书其余部分，我们在大多数情况下不需要假设它们成立），以及对于表达传统逻辑至关重要的\emph{命题截断}。
这一章是 \cref{cha:set-math,cha:real-numbers} 的必要背景，对 \cref{cha:category-theory} 较不重要，对 \cref{cha:homotopy} 则不太必要。

\cref{cha:equivalences,cha:induction} 详细研究了两个专题：等价（及相关概念）与广义归纳定义。
虽然这些主题本身很重要，并且能加深对同伦类型论的理解，但对 \cref{part:mathematics} 而言，它们在大多数情况下并非必需。
\cref{cha:equivalences} 中只有零星引理在个别地方被用到，而 \cref{sec:bool-nat,sec:strictly-positive,sec:generalizations} 中的一般性讨论则有助于为 \cref{cha:hits} 提供所需的直觉。
\cref{sec:generalizations} 中讨论的广义归纳定义也在 \cref{cha:set-math,cha:real-numbers} 的若干处被使用。

\cref{cha:hits} 引入同伦类型论的第二项基本创新——\emph{高阶归纳类型}——并给出许多例子。
高阶归纳类型是 \cref{cha:homotopy} 的主要研究对象，某些特定的高阶归纳类型在 \cref{cha:set-math,cha:real-numbers} 中也扮演重要角色。
它们对 \cref{cha:category-theory} 来说不太必要，尽管 \cref{sec:rezk} 中用到了其中一个例子。

最后，\cref{cha:hlevels} 讨论同伦 $n$-类型及相关概念，例如 $n$-连通类型。
这些概念对 \cref{cha:homotopy} 很重要，但在 \cref{part:mathematics} 的其余部分则不太重要，尽管其中某些引理在 $n=-1$ 的情形下被用于 \cref{sec:piw-pretopos}。

以上便是 \cref{part:foundations} 的全部。
如前所述，\cref{part:mathematics} 由四个在很大程度上互不相关的章节组成，每章描述泛等基础能为某个特定学科提供什么。

在 \cref{part:mathematics} 的各章中，\cref{cha:homotopy}（同伦论）也许是最激进的。
泛等基础对同伦论采取了一种截然不同的``综合''方法，其中同伦类型本身就是基本对象（即类型），而不是用拓扑空间或其他集合论模型构造出来的。
这使我们能够用新的证明风格来处理代数拓扑中的经典定理；我们在此给出了一些范例，从 $\pi_1(\Sn^1)=\Z$ 到 Freudenthal 纬悬定理。

在 \cref{cha:category-theory} (范畴论) 中，我们发展了基本的 (1-) 范畴论，并遵循泛等公理的原则，即\emph{相等即同构}。
这样做有一个令人愉快的效果：所有定义与构造均在范畴等价下自动保持不变——事实上，等价的范畴就是相等的，正如等价的类型就是相等的一样。
（这与高阶范畴论和高阶意象理论也有所联系。）

\cref{cha:set-math} (集合论) 研究的是泛等基础下的集合。
集合范畴具备其通常的性质，从而为任何不需要同伦或高阶范畴结构的数学提供了基础。
我们还会看到，泛等性使基数与序数变得更为自然，而高阶归纳类型则可产生一个满足 Zermelo--Fraenkel 集合论通常公理的累积层次。

在 \cref{cha:real-numbers} (实数) 中，我们总结了 Dedekind实数的构造，并注意到高阶归纳类型允许我们定义 Cauchy（柯西）实数，从而避开构造性数学中与之相关的一些问题。
接着，我们概述了处理 Conway（康威）超现实数的类似方法。

本书每一章都以一个“附注”小节作结，用以汇集历史评注、文献引用，并尽可能注明结果的归属。
我们还在每章末尾附有习题，以帮助读者熟悉如何在泛等基础下做数学。

最后，还需指出本书是由众多参与者大规模协作完成的。
我们已尽力确保术语与记号的统一，并将数学内容组织为逻辑上连贯的线性序列，但很可能仍存在一些瑕疵。
对于任何此类瑕疵，我们恳请读者谅解，并欢迎为下一版的改进提出建议。

% Local Variables:
% TeX-master: "hott-online"
% End:

\part{基础篇}
\label{part:foundations}

\chapter{类型论}
\label{cha:typetheory}

\section{类型论 vs. 集合论}
\label{sec:types-vs-sets}
\label{sec:axioms}

\index{type theory}
同伦类型论除了其他用途外，还是数学的一种基础语言，换言之，是 Zermelo--Fraenkel\index{set theory!Zermelo--Fraenkel} 集合论的一种替代方案。
然而，它在几个重要方面的运作方式与集合论不同，可能需要一些时间来适应。
若要仔细解释这些差异，就需要我们在这里比本书其余部分更为形式化。
正如引言所述，我们的目标是以\emph{非形式化}的方式来写类型论；但对惯用集合论的数学家而言，在开头给出更精确的说明，有助于避免一些常见的误解和错误。

我们注意到，集合论的基础有两个“层次”：一阶逻辑的演绎系统\index{first-order!logic}，以及在这个系统内部表述的特定理论（例如 ZFC）的公理。
因此，集合论不仅关乎集合，更关乎集合（第二层的对象）与命题（第一层的对象）之间的相互作用。

相比之下，类型论本身就是一个演绎系统：它不必在任何上层结构（例如一阶逻辑）内部来表述。
类型论的基本概念不是集合论中的那两个（集合与命题），而是只有一个：\emph{类型}。
命题（那些我们可以证明、反驳、假设、否定等等\footnote{容易混淆的是，将“命题（proposition）”一词与“定理（theorem）”作为同义词使用也是一种常见做法（可追溯至Euclid）。
  我们将遵循逻辑学家的用法：一个\emph{命题}是一个\emph{可以}被证明的陈述，而一个\emph{定理}\indexfoot{theorem}（或“引理”\indexfoot{lemma}、或“推论”\indexfoot{corollary}）则是这样一个\emph{已被}证明的陈述。
因此，“$0=1$”和它的否定“$\neg (0=1)$”都是命题，但只有后者是定理。}的陈述）被等同于特定类型，这种对应关系见第~\pageref{tab:pov} 页的 \cref{tab:pov}。
于是，\emph{证明定理}这一数学活动，就被等同于\emph{构造对象}这一数学活动的一个特例——在这里，构造的就是一个代表命题的类型的实例。

\index{deductive system}%
这引出了类型论与集合论之间的另一个差异，但为了解释它，我们必须先一般性地谈谈演绎系统。
非正式地说，一个演绎系统就是一套\define{规则}，
\indexdef{rule}%
用于推导所谓\define{判断}。
\indexdef{judgment}%
如果我们将演绎系统视为一个形式化的游戏\index{game!deductive system as}%
那么判断就是遵循游戏规则后所到达的“位置”。
我们也可以将演绎系统视作某种代数理论，在这种情况下，判断就是其中的元素（好比群中的元素），而演绎规则就是其中的运算（好比群的乘法）。
从逻辑的观点来看，判断可以被认为是“外部”的陈述，居于元理论之中，与理论自身的“内部”陈述相对。

在（作为集合论基础的）一阶逻辑的演绎系统中，只有一种判断：某个给定的命题有一个证明。
也就是说，每个命题 $A$ 都对应着一个判断“$A$ 有一个证明”，而所有判断都具有这种形式。
一条一阶逻辑的规则，例如“由 $A$ 和 $B$ 可推出 $A\wedge B$”，实际上是一条“证明构造”规则：它说的是，给定判断“$A$ 有一个证明”和“$B$ 有一个证明”，我们可以推导出“$A\wedge B$ 有一个证明”。
请注意，判断“$A$ 有一个证明”所处的层次，与命题 $A$ 本身——作为理论的一个内部陈述——是不同的。
% In particular, we cannot manipulate it to construct propositions such as ``if $A$ has a proof, then $B$ does not have a proof'' --- unless we are using our set-theoretic foundation as a meta-theory with which to talk about some other axiomatic system.

类型论的基本判断，类似于“$A$ 有一个证明”，写作“$a:A$”，读作“项 $a$ 具有类型 $A$”（the term $a$ has type $A$），或者更宽松地说“$a$ 是 $A$ 的一个元素”——在同伦类型论中，也常说“$a$ 是 $A$ 的一个点”。
\indexdef{term}%
\indexdef{element}%
\indexdef{point!of a type}%
当 $A$ 是一个代表命题的类型时，$a$ 可以称为 $A$ 之可证性的一个\emph{见证}\index{witness!to the truth of a proposition}，或其真理性的\emph{证据}\index{evidence, of the truth of a proposition}（有时甚至称为 $A$ 的一个\emph{证明}\index{proof}，但我们会尽量避免使用这种易混淆的术语）。
此时，在类型论中（对某个 $a$）可以推导出判断 $a:A$，恰对应于在一阶逻辑中可以推导出类似的判断“$A$ 有一个证明”（除去所设公理及数学编码方式的差异，这些差异将贯穿全书加以讨论）。

另一方面，如果类型 $A$ 更多地被当作一个集合而非一个命题来对待（尽管我们将会看到，这种区别有时是模糊的），那么“$a:A$”就可以看作类似于集合论中的陈述“$a\in A$”。
然而，二者有一个本质区别：“$a:A$”是一个\emph{判断}，而“$a\in A$”是一个\emph{命题}。
具体来说，在类型论内部工作时，我们不能做出像“如果 $a:A$，那么并非 $b:B$”这样的陈述，也不能“反驳”判断“$a:A$”。

理解这一点的一个好办法是：在集合论中，“属于”是两个预先存在的对象“$a$”和“$A$”之间可能成立也可能不成立的一种\emph{关系}；而在类型论中，我们不能孤立地谈论一个“元素” $a$：每个元素\emph{依其本质}就是某个类型的元素，并且该类型（一般而言）是唯一确定的。
因此，当我们非正式地说“令 $x$ 是一个自然数”时，在集合论中这是“令 $x$ 是一个东西，并假设 $x\in\nat$”的简写；而在类型论中，“令 $x:\nat$”则是一个原子语句：我们不能在不指明其类型的情况下引入一个变量。\index{membership}

初看之下，这似乎是一种不便的限制，但它可以说更接近于“令 $x$ 是一个自然数”的直观数学含义。
在实践中，似乎每当我们真正\emph{需要}“$a\in A$”作为一个命题而非判断时，总存在一个背景集合 $B$，使得已知 $a$ 是 $B$ 的元素，且 $A$ 是 $B$ 的子集。
这种情况在类型论中也容易表示，只需取 $a$ 为类型 $B$ 的一个元素，而 $A$ 为 $B$ 上的一个谓词；参见\cref{subsec:prop-subsets}。

类型论与集合论之间的最后一个区别在于对相等性的处理。
数学中熟悉的相等性概念是一个命题：例如，我们可以反驳一个等式，或者将一个等式作为假设。
由于在类型论中，命题就是类型，这意味着相等性也是一个类型：对于元素 $a,b:A$（即同时有 $a:A$ 和 $b:A$），我们有一个类型“$\id[A]ab$”。
（当然，在\emph{同伦}类型论中，这个相等命题可能表现出不熟悉的行为：参见\cref{sec:identity-types,cha:basics}及本书其余部分。）
当 $\id[A]ab$ 有实例时，我们说 $a$ 和 $b$ 是\define{（命题）相等的}。
\index{propositional!equality}%
\index{equality!propositional}%

然而，在类型论中还需要一种相等性\emph{判断}，它存在于与判断“$x:A$”相同的层面。\index{judgment}
\symlabel{defn:judgmental-equality}%
这被称为\define{判断相等性}
\indexdef{equality!judgmental}%
\indexdef{judgmental equality}%
或\define{定义相等性}，
\indexdef{equality!definitional}%
\indexsee{definitional equality}{equality, definitional}%
我们将其写作 $a\jdeq b : A$ 或简写为 $a \jdeq b$。
可以这样理解：它意味着“依定义相等”。
例如，如果我们用方程 $f(x)=x^2$ 定义一个函数 $f:\nat\to\nat$，那么表达式 $f(3)$ 就\emph{依定义}等于 $3^2$。
在理论内部，对“依定义相等”进行否定或假设是没有意义的；我们不能说“如果 $x$ 依定义等于 $y$，那么 $z$ 就不依定义等于 $w$”。
两个表达式是否依定义相等，只需展开定义便可判定；特别地，这是\emph{算法上可判定的}\index{algorithm}（尽管该算法必然是\emph{元理论的}，而非理论内部的）。\index{decidable!definitional equality}

随着类型论变得日益复杂，判断相等性也会比这里描述的更加微妙，但以此为起点已是一种很好的直觉。
换一个角度看，若我们将演绎系统视作一个代数理论，那么判断相等性无非就是该理论内部的相等性，类似于群元素之间的相等——唯一可能造成混淆之处在于，在类型论的演绎系统\emph{内部}，\emph{也}存在着一个对象（即类型``$a=b$''），它在系统内部扮演着``相等性''的角色。

我们之所以\emph{需要}判断相等性这个概念，是为了用它来控制另一种判断形式``$a:A$''。
举例来说，假设我们已经证明了 $3^2=9$，即对某个 $p$，我们已经推出了判断 $p:(3^2=9)$。
那么，同一个见证元 $p$ 理应同样算作 $f(3)=9$ 的证明，因为按照定义，$f(3)$ 就是 $3^2$。
表达这一点的最佳方式是一条规则：给定判断 $a:A$ 和 $A\jdeq B$，可推出判断 $a:B$。

因此，对我们而言，类型论将是一个基于以下两种判断形式的演绎系统：


\begin{center}
\medskip
\begin{tabular}{cl}
  \toprule
  判断 & 含义\\
  \midrule
  $a : A$       & ``$a$ 是类型 $A$ 的对象''\\
  $a \jdeq b : A$ & ``$a$ 和 $b$ 是类型 $A$ 的定义相等的对象''\\
  \bottomrule
\end{tabular}
\medskip
\end{center}


%
\symlabel{defn:defeq}%
当引入定义相等，亦即将一个对象定义为与另一个对象相等时，我们将使用符号``$\defeq$''。
因此，上面提到的函数 $f$ 的定义将写作 $f(x)\defeq x^2$。

由于判断无法组合成更复杂的陈述，符号``$:$''和``$\jdeq$''的结合力比任何其他符号都弱。%
\footnote{在形式化的\indexfoot{mathematics!formalized}类型论中，逗号和推断符（$\vdash$）的结合力可能更弱。
  例如，$x:A,y:B\vdash c:C$ 会被解析为 $((x:A),(y:B))\vdash (c:C)$。
  但在本书中，在\cref{cha:rules}之前，我们暂不使用这种记法。}
因此，举例来说，``$p:\id{x}{y}$''应解析为``$p:(\id{x}{y})$''，这是有意义的，因为``$\id{x}{y}$''是一个类型；而不应解析为``$\id{(p:x)}{y}$''，这是无意义的，因为``$p:x$''是一个判断，不能与任何东西相等。
类似地，``$A\jdeq \id{x}{y}$''只能被解析为``$A\jdeq(\id{x}{y})$''，尽管在这种极端情况下，为了辅助阅读理解，最好还是加上括号。
此外，后续我们还将采用一种常见的等式链记号——例如写 $a=b=c=d$ 来表示``$a=b$ 且 $b=c$ 且 $c=d$，因此 $a=d$''——并且我们也会在链中包含判断相等性。语境通常足以说明意图。

或许在这里也可以顺便指出，常见的数学记号``$f:A\to B$''（表示 $f$ 是从 $A$ 到 $B$ 的函数）可以看作一个类型判断，因为我们用``$A\to B$''作为从 $A$ 到 $B$ 的函数类型的记号（这是类型论的标准做法；参见 \cref{sec:pi-types}）。

\index{assumption|(defstyle}%
判断可以依赖于形如 $x:A$ 的\emph{假设}，其中 $x$ 是一个变量\indexdef{variable}，%
$A$ 是一个类型。
例如，在假设 $m,n : \nat$ 下，我们可以构造一个对象 $m + n : \nat$。
另一个例子是，假设 $A$ 是一个类型，$x,y : A$，并且 $p : \id[A]{x}{y}$，我们可以构造一个元素 $p^{-1} : \id[A]{y}{x}$。
所有这些假设的集合称为\define{上下文（context）}；%
\index{context}
从拓扑的观点来看，它可以被视为一个``参数\index{parameter!space}空间''。
事实上，严格来说，上下文必须是一个有序的假设列表，因为后面的假设可能依赖于前面的假设：只有在类型 $A$ 中出现的任何变量的假设都已给出\emph{之后}，才能做出假设 $x:A$。

如果假设 $x:A$ 中的类型 $A$ 代表一个命题，那么该假设就是\emph{假说}\indexdef{hypothesis}在类型论中的对应物：我们假设命题 $A$ 成立。
当类型被视为命题时，我们可以省略其证明的名称。
因此，在上面的第二个例子中，我们可以改为说：假设 $\id[A]{x}{y}$，可以证明 $\id[A]{y}{x}$。
然而，由于我们从事的是``证明相关''\index{mathematics!proof-relevant}的数学，我们会频繁地将证明本身作为对象来提及。
例如在上面的例子中，我们可能想要证明 $p^{-1}$ 连同传递性和自反性的证明一起表现得像一个群胚；参见 \cref{cha:basics}。

注意，在``假设''一词的上述含义下，我们可以假设一个命题相等（通过假设一个变量 $p:x=y$），但不能假设一个判断相等性 $x\jdeq y$，因为它不是一个可以拥有元素的类型。
不过，我们可以做另一件看起来有点像假设判断相等性的事情：如果我们有一个涉及变量 $x:A$ 的类型或元素，那么可以用某个特定的元素 $a:A$ 来\emph{替换} $x$，从而得到一个更具体的类型或元素。
我们有时会使用诸如``现在假设 $x\jdeq a$''这样的语言来指代这一替换过程，即使它严格来说并非上文所引入的\emph{假设}。
\index{assumption|)}%

同理，我们也不能\emph{证明}一个判断相等性，因为它不是一个可以展示见证元的类型。
尽管如此，我们有时会将判断相等性作为定理的一部分来陈述，例如``存在 $f:A\to B$ 使得 $f(x)\jdeq y$''。
这应被视为作出了两个独立的判断：首先，我们对某个元素 $f$ 作出判断 $f:A\to B$；然后，我们作出附加的判断 $f(x)\jdeq y$。

在本章剩余部分，我们将尝试对类型论做一个非正式的介绍，这足以满足本书的需要；我们会在\cref{cha:rules}中给出更形式化的描述。
除了一些相当明显的规则（例如判断相等的对象总是可以相互替换\index{substitution}），类型论的规则可以按\emph{类型形成子}进行分组。
每个类型形成子由一种构造类型的方法（可能利用先前构造的类型），连同该类型元素的构造与行为规则组成。
在大多数情况下，这些规则遵循相当可预测的模式，但我们不会在此尝试精确描述这一点；不过可以参见\cref{sec:finite-product-types}的开头以及\cref{cha:induction}。\index{type theory!informal}

\index{axiom!versus rules}%
\index{rule!versus axioms}%
本章介绍的类型论的一个重要方面是，它完全由\emph{规则}构成，没有任何\emph{公理}。
在用判断来描述演绎系统时，\emph{规则}允许我们从一组判断推导出另一个判断，而\emph{公理}则是初始给定的判断。
如果我们把演绎系统看作一个形式游戏，那么规则就是游戏的规则，而公理则是起始位置。
如果我们把演绎系统看作一个代数理论，那么规则就是理论中的运算，而公理则是该理论某个特定自由模型的\emph{生成元}。

在集合论中，唯一的规则是一阶逻辑的规则（例如允许我们从``$A$ 有一个证明''和``$B$ 有一个证明''推导出``$A\wedge B$ 有一个证明''的规则）：关于集合行为的所有信息都包含在公理中。
相比之下，在类型论中，通常\emph{规则}包含了所有的信息，而不需要公理。
例如，在\cref{sec:finite-product-types}中我们将看到，存在一条规则允许我们从``$a:A$''和``$b:B$''推导出判断``$(a,b):A\times B$''，而在集合论中，类似的陈述将是配对公理的（推论）。

仅以规则来表述类型论的一个优点是，规则是``过程性的''。
特别是，正是这一特性使类型论可能具有良好的计算性质（尽管并不能自动保证），例如``典范性''。\index{canonicity}
然而，尽管这种风格适用于传统的类型论，我们尚未完全理解如何以这种方式来表述\emph{同伦}类型论所需的一切。
特别地，在\cref{sec:compute-pi,sec:compute-universe,cha:hits}中，我们将不得不引入额外的公理——尤其是\emph{泛等公理}——来扩充本章所介绍的类型论规则。
但在本章中，我们仅限于讨论传统的、基于规则的类型论。

\section{函数类型}
\label{sec:function-types}

\index{type!function|(defstyle}%
\indexsee{function type}{type, function}%
给定类型 $A$ 和 $B$，我们可以构造类型 $A \to B$，其中的\define{函数}
\index{function|(defstyle}%
\indexsee{map}{function}%
\indexsee{mapping}{function}%
以 $A$ 为定义域、$B$ 为陪域。
我们有时也称函数为\define{映射}。
\index{domain!of a function}%
\index{codomain, of a function}%
\index{function!domain of}%
\index{function!codomain of}%
\index{functional relation}%
与集合论不同，在类型论中，我们不把函数定义为函数关系；相反，函数是一个原始概念。
我们通过规定以下内容来解释函数类型：可以对函数做什么操作，如何构造函数，以及函数会诱导出哪些相等关系。

给定一个函数 $f : A \to B$ 和定义域中的一个元素 $a : A$，我们可以\define{应用}
\indexdef{application!of function}%
\indexdef{function!application}%
\indexsee{evaluation}{application, of a function}
这个函数，从而得到陪域 $B$ 中的一个元素，记作 $f(a)$，并称之为 $f$ 在 $a$ 处的\define{值}。
\indexdef{value!of a function}%
在类型论中，省略括号\index{parentheses}而将 $f(a)$ 简记为 $f\,a$ 是很常见的做法，我们有时也会这样写。

但是，我们该如何构造 $A \to B$ 中的元素呢？有两种等价的方法：要么通过直接定义，要么使用 $\lambda$-抽象。
通过定义来引入函数
\indexdef{definition!of function, direct}%
指的是，我们为它取一个名字——比方说 $f$——然后通过给出一个方程
\begin{equation}
  \label{eq:expldef}
  f(x) \defeq \Phi
\end{equation}
来定义 $f : A \to B$，其中 $x$ 是一个
\index{variable}%
变量，而 $\Phi$ 是一个可能用到 $x$ 的表达式。
为了使这个定义有效，我们必须验证在假设 $x:A$ 时有 $\Phi : B$。

这样，我们就可以通过把表达式 $\Phi$ 中的变量 $x$ 替换成 $a$ 来计算 $f(a)$。
举一个例子，考虑函数 $f : \nat \to \nat$，它由 $f(x) \defeq x+x$ 定义。（我们将在 \cref{sec:inductive-types} 中定义 $\nat$ 和 $+$。）
那么 $f(2)$ 与 $2+2$ 判断相等。

如果我们不想为函数引入名字，可以使用
\define{$\lambda$-抽象}。
\index{lambda abstraction@$\lambda$-abstraction|defstyle}%
\indexsee{function!lambda abstraction@$\lambda$-abstraction}{$\lambda$-abstraction}%
\indexsee{abstraction!lambda-@$\lambda$-}{$\lambda$-abstraction}%
给定一个如上所述的类型为 $B$、可能用到 $x:A$ 的表达式 $\Phi$，我们写作 $\lam{x:A} \Phi$ 以表示由~\eqref{eq:expldef} 定义的同一个函数。
因此，我们有
\[ (\lamt{x:A}\Phi) : A \to B. \]
对于上一段的例子，我们有类型判断
\[ (\lam{x:\nat}x+x) : \nat \to \nat. \]
作为另一个例子，对于任意类型 $A$ 和 $B$ 以及任意元素 $y:B$，我们有一个\define{常值函数}
\indexdef{constant!function}%
\indexdef{function!constant}%
$(\lam{x:A} y): A\to B$。

我们通常会省略 $\lambda$-抽象中变量 $x$ 的类型标注，直接写作 $\lam{x}\Phi$，这是因为从函数 $\lam x \Phi$ 具有类型 $A\to B$ 这一判断，就可以推断出 $x:A$。
按约定，变量绑定符“$\lam{x}$”的\define{作用域}
\indexdef{variable!scope of}%
\indexdef{scope}%
为表达式中紧随其后的整个剩余部分，除非被括号\index{parentheses}隔开。
因此，举例来说，$\lam{x} x+x$ 应被解析为 $\lam{x} (x+x)$，而不是 $(\lam{x}x)+x$（而且，在此例中后者无论如何都会是类型错误的）。

另一种等价的记法是
\symlabel{mapsto}%
\[ (x \mapsto \Phi) : A \to B. \]
\symlabel{blank}%
我们有时也会在表达式 $\Phi$ 中使用空白记号“$\blank$”来代替变量，以表示一个隐式的 $\lambda$-抽象。
例如，$g(x,\blank)$ 就是 $\lam{y} g(x,y)$ 的另一种写法。

$\lambda$-抽象是一个函数，因此我们可以把它应用于参数 $a:A$。
于是我们得到下面的 \define{计算规则}\indexdef{computation rule!for function types}\footnote{这个等式的使用常被称为 \define{$\beta$-变换}
\indexsee{beta-conversion@$\beta $-conversion}{$\beta$-reduction}%
\indexsee{conversion!beta@$\beta $-}{$\beta$-reduction}%
或 \define{$\beta$-归约}。%
\index{beta-reduction@$\beta $-reduction|footstyle}%
\indexsee{reduction!beta@$\beta $-}{$\beta$-reduction}%
}，它是一个定义相等：
\[(\lamu{x:A}\Phi)(a) \jdeq \Phi'\]
其中 $\Phi'$ 是将表达式 $\Phi$ 中所有 $x$ 的出现都替换为 $a$ 后得到的结果。

沿用前面的例子，我们有
%
\[ (\lamu{x:\nat}x+x)(2) \jdeq 2+2. \]
%

注意，从任何函数 $f:A\to B$，我们都能构造一个 $\lambda$-抽象函数 $\lam{x} f(x)$。
因为按照定义它就是“将 $f$ 应用于其参数的函数”，所以我们认为它与 $f$ 定义相等：\footnote{这个等式的使用常被称为 \define{$\eta$-变换}
\indexsee{eta-conversion@$\eta $-conversion}{$\eta$-expansion}%
\indexsee{conversion!eta@$\eta $-}{$\eta$-expansion}%
或 \define{$\eta$-展开}。
\index{eta-expansion@$\eta $-expansion|footstyle}%
\indexsee{expansion, eta-@expansion, $\eta $-}{$\eta$-expansion}%
}
\[ f \jdeq (\lam{x} f(x)). \]
这个等式就是\define{函数类型的唯一性原理}\indexdef{uniqueness!principle!for function types}，因为它表明 $f$ 由其值唯一决定。

通过带显式参数的定义来引入函数，可以借助 $\lambda$-抽象归约为简单定义：也就是说，我们可以将形如
\[ f(x) \defeq \Phi \]
的 $f: A\to B$ 的定义，读作
\[ f \defeq \lamu{x:A}\Phi.\]

当我们进行涉及变量的计算时，若要将一个变量替换为也含有变量的表达式，就必须格外小心，因为我们希望保持表达式的绑定结构。这里所说的\emph{绑定结构}\indexdef{binding structure}，指的是由绑定符（比如 $\lambda$、$\Pi$ 和 $\Sigma$——后两个我们很快就会遇到）在变量被引入的位置与它被使用的位置之间生成的不可见链接。

举个例子，考虑定义为
\[ f(x) \defeq \lamu{y:\nat} x + y. \]
的 $f : \nat \to (\nat \to \nat)$。
现在，如果我们事先假设了 $y : \nat$，那么 $f(y)$ 是什么呢？如果只是简单地在定义 $f(x)$ 的表达式 ``$\lam{y}x+y$'' 中把所有 $x$ 都替换成 $y$，得到 $\lamu{y:\nat} y + y$，那就错了，因为这意味着 $y$ 被\define{捕获}了。
\indexdef{capture, of a variable}%
\indexdef{variable!captured}%
此前，被代入的\index{substitution} $y$ 指的是我们的假设，可现在它却指代了 $\lambda$-抽象的参数。于是，这种简单替换会破坏绑定结构，使我们得以进行语义上不正确的计算。

但是，在这个例子中，$f(y)$ 究竟\emph{是}什么呢？注意，约束变量（或称“哑变量”）
\indexdef{variable!bound}%
\indexdef{variable!dummy}%
\indexsee{bound variable}{variable, bound}%
\indexsee{dummy variable}{variable, bound}%
——比如表达式 $\lamu{y:\nat} x + y$ 中的 $y$——只在局部有意义，并且可以一致地替换为任何其他变量，同时保持绑定结构不变。事实上，$\lamu{y:\nat} x + y$ 被规定为与 $\lamu{z:\nat} x + z$ 判断相等\footnote{这个等式的使用常被称为 \define{$\alpha$-变换}。
\indexfoot{alpha-conversion@$\alpha $-conversion}
\indexsee{conversion!alpha@$\alpha$-}{$\alpha$-conversion}
}。由此可见，
$f(y)$ 与 $\lamu{z:\nat} y + z$ 判断相等，而这正是我们的答案。（除了 $z$，任何与 $y$ 不同的变量都可以使用，得到的结果都是相等的。）

当然，这对任何数学家来说都应该很熟悉：它与积分中的现象是同一回事——如果 $f(x) \defeq \int_1^2 \frac{dt}{x-t}$，那么 $f(t)$ 不是 $\int_1^2 \frac{dt}{t-t}$ 而是 $\int_1^2 \frac{ds}{t-s}$。
$\lambda$-抽象绑定一个哑变量的方式，与积分绑定积分变量的方式完全一样。

我们已经看到了如何定义单变量函数。定义多变量函数的一种方法是使用 Cartesian 积（这将在后面介绍）；一个参数为 $A$ 与 $B$、结果为 $C$ 的函数，其类型将是 $f : A \times B \to C$。
然而，还有另一种选择可以避免使用积类型，它被称为 \define{柯里化}
\indexdef{currying}%
\indexdef{function!currying of}%
（得名于数学家 Haskell Curry）。
\index{programming}%

柯里化的核心思想是：将接受两个输入 $a:A$ 和 $b:B$ 的函数，表示为这样一个函数——它只接受\emph{一个}输入 $a:A$，然后返回\emph{另一个函数}，而后者再接受第二个输入 $b:B$ 并给出最终结果。
也就是说，我们将双变量函数视为属于迭代函数类型 $f : A \to (B \to C)$。
我们也可以省略括号\index{parentheses}，将其写作 $f : A \to B \to C$，并以右结合\index{associativity!of function types}为默认约定。
这样，给定 $a : A$ 和 $b : B$，我们可以先将 $f$ 应用于 $a$，再将结果应用于 $b$，从而得到 $f(a)(b) : C$。
为了避免括号过多，我们约定可将 $f(a)(b)$ 写作 $f(a,b)$，尽管这里并没有涉及任何积类型。
当我们完全省略函数参数周围的括号时，我们用 $f\,a\,b$ 表示 $(f\,a)\,b$，此时的默认结合性为左结合，以保证 $f$ 能按正确的顺序应用于其参数。

我们带有显式参数的定义记号也可以推广到这种情况：我们可以通过给出一个等式
\[ f(x,y) \defeq \Phi\]
来定义一个命名函数 $f : A \to B \to C$，其中在给定 $x:A$ 和 $y:B$ 时有 $\Phi:C$。
利用 $\lambda$-抽象\index{lambda abstraction@$\lambda$-abstraction}，这对应着
\[ f \defeq \lamu{x:A}{y:B} \Phi, \]
也可以写作
\[ f \defeq x \mapsto y \mapsto \Phi. \]
我们还可以通过写出多个空位来隐式地对多个变量进行抽象，例如 $g(\blank,\blank)$ 表示 $\lam{x}{y} g(x,y)$。
对三个或更多参数的函数进行柯里化，是刚才所述内容的直接推广。
 
\index{type!function|)}%
\index{function|)}%

\section{宇宙与类型族}
\label{sec:universes}

到目前为止，我们一直非正式地使用“$A$ 是一个类型”这种说法。
为了使其更加精确，我们将引入\define{宇宙（universes）}。
\index{type!universe|(defstyle}%
\indexsee{universe}{type, universe}%
宇宙是其元素为类型的类型。
就像在朴素集合论中那样，我们或许希望有一个包含所有类型的宇宙 $\UU_\infty$，并且它包含自身（即 $\UU_\infty : \UU_\infty$）。
然而，与集合论中的情况一样，这是不可靠的；也就是说，我们可以由此推导出每一个类型（包括代表命题“假”的空类型，见 \cref{sec:coproduct-types}）都是\emph{有实例的}。
例如，利用集合的树形表示，我们可以直接编码 Russell（罗素）的悖论\index{paradox} \cite{coquand:paradox}。
%  or alternatively, in order to avoid the use of
% inductive types to define trees, we can follow Girard \cite{girard:paradox} and encode the Burali-Forti paradox,
% which shows that the collection of all ordinals cannot be an ordinal.

为了避免这个悖论，我们引入一个宇宙的层级结构
\indexsee{hierarchy!of universes}{type, universe}%
\[ \UU_0 : \UU_1 : \UU_2 : \cdots \]
其中每一个宇宙 $\UU_i$ 都是下一个宇宙 $\UU_{i+1}$ 的元素。
此外，我们假设我们的宇宙是\define{累积的（cumulative）}，
\indexdef{type!universe!cumulative}%
\indexdef{cumulative!universes}%
也就是说，第 $i$ 个宇宙的所有元素也都是第 $(i+1)^{\mathrm{st}}$ 个宇宙的元素，即若 $A:\UU_i$，则也有 $A:\UU_{i+1}$。
这样做很方便，但有一个略有不便的后果：元素不再具有唯一的类型；而且在其他一些此处无需深究的方面也略有些棘手；详见注释。

当我们说 $A$ 是一个类型时，我们的意思是它属于某个宇宙 $\UU_i$。
我们通常希望避免显式地提及层级
\indexdef{universe level}%
\indexsee{level}{universe level or $n$-type}%
\indexsee{type!universe!level}{universe level}%
$i$，而只是假设可以一致地指定层级；
因此我们可以省略层级，直接写作 $A:\UU$。
这样一来，我们甚至可以写出 $\UU:\UU$，它可以读作 $\UU_i:\UU_{i+1}$，因为我们隐去了指标。
这种书写宇宙的方式被称为\define{典型歧义（typical ambiguity）}。
\indexdef{typical ambiguity}%
它很方便，但也有一点危险，因为它允许我们写出看似有效、实则重现了自指悖论的证明。
如果对一个论证是否正确存在任何疑问，检查的方法就是尝试为其中出现的所有宇宙一致地分配层级。
当我们假设某个宇宙 \UU 时，我们将属于 \UU 的类型称为\define{小类型（small types）}。
\indexdef{small!type}%
\indexdef{type!small}%

为了对随给定类型 $A$ 变化的一组类型进行建模，我们使用陪域为宇宙的函数 $B : A \to \UU$。
这样的函数称为\define{类型族（families of types）}（有时也称为\emph{依值类型（dependent types）}）；
\indexsee{family!of types}{type, family of}%
\indexdef{type!family of}%
\indexsee{type!dependent}{type, family of}%
\indexsee{dependent!type}{type, family of}%
它们对应于集合论中使用的集合族。

\symlabel{fin}%
类型族的一个例子是有限集族 $\Fin : \nat \to \UU$，其中 $\Fin(n)$ 是一个恰好有 $n$ 个元素的类型。
（我们目前还不能\emph{定义}这个族 $\Fin$——事实上，我们甚至尚未引入其定义域 $\nat$——但我们很快就能做到；参见\cref{ex:fin}。）
我们可以用 $0_n,1_n,\dots,(n-1)_n$ 来记 $\Fin(n)$ 中的元素，并给这些元素加上下标，以强调当 $n$ 与 $m$ 不同时，$\Fin(n)$ 的元素与 $\Fin(m)$ 的元素是不同的，并且它们全都不同于普通的自然数（后者我们将在\cref{sec:inductive-types}中引入）。
\index{finite!sets, family of}%

一个更平凡（却非常重要）的类型族例子是\define{常值类型族（constant type family）}
\indexdef{constant!type family}%
\indexdef{type!family of!constant}%
，即取值恒为某个类型 $B:\UU$，自然也就是常值函数 $(\lam{x:A} B):A\to\UU$。

作为一个\emph{反}例，在我们这个版本的类型论中，不存在“$\lam{i:\nat} \UU_i$”这样的类型族。
的确，没有哪个宇宙足够大，能够充当其陪域。
此外，我们甚至不把诸宇宙 $\UU_i$ 的指标 $i$ 与类型论中的自然数 $\nat$ 视为等同（后者将在\cref{sec:inductive-types}中引入）。

\index{type!universe|)}%

\section{依值函数类型（\texorpdfstring{$\Pi$}{Π} 类型）}
\label{sec:pi-types}

\index{type!dependent function|(defstyle}%
\index{function!dependent|(defstyle}%
\indexsee{dependent!function}{function, dependent}%
\indexsee{type!Pi-@$\Pi$-}{type, dependent function}%
\indexsee{Pi-type@$\Pi$-type}{type, dependent function}%
在类型论中，我们经常使用函数类型的一个更一般的版本，称为\define{$\Pi$ 类型（$\Pi$-type）}或\define{依值函数类型（dependent function type）}。
$\Pi$ 类型的元素是一些函数，其陪域类型可随函数所应用的定义域元素而变化，这样的函数称为\define{依值函数（dependent functions）}。
之所以使用“$\Pi$ 类型”这个名称，是因为这种类型也可以被视为给定类型上的 Cartesian 积。

给定一个类型 $A:\UU$ 和一个类型族 $B:A \to \UU$，我们可以构造依值函数的类型 $\prd{x:A}B(x) : \UU$。
这个类型有许多替代记法，例如 \[ \tprd{x:A} B(x) \qquad \dprd{x:A}B(x) \qquad \lprd{x:A} B(x). \]
如果 $B$ 是一个常值族，那么这个依值乘积类型就是通常的函数类型：
\[\tprd{x:A} B \jdeq (A \to B).\]
事实上，所有关于 $\Pi$ 类型的构造都是普通函数类型上相应构造的推广。

\indexdef{definition!of function, direct}%
我们可以通过显式定义来引入依值函数：要定义 $f : \prd{x:A}B(x)$，其中 $f$ 是待定义的依值函数的名字，我们需要一个可能涉及变量 $x:A$ 的表达式 $\Phi : B(x)$，
\index{variable}%
并写作
\[ f(x) \defeq \Phi \qquad \mbox{for $x:A$}.\]
另一种方式是使用\define{$\lambda$-抽象（$\lambda$-abstraction）}。%
\index{lambda abstraction@$\lambda$-abstraction|defstyle}%
\begin{equation}
  \label{eq:lambda-abstraction}
  \lamu{x:A} \Phi \ :\ \prd{x:A} B(x).
\end{equation}
\indexdef{application!of dependent function}%
\indexdef{function!dependent!application}%
与非依值函数一样，我们可以将一个依值函数 $f : \prd{x:A}B(x)$ \define{应用（apply）}到一个参数 $a:A$ 上，从而得到一个元素 $f(a):B(a)$。
其等式与普通函数类型相同，即我们有以下计算规则：
\index{computation rule!for dependent function types}%
给定 $a:A$，我们有 $f(a) \jdeq \Phi'$ 且 $(\lamu{x:A} \Phi)(a) \jdeq \Phi'$，其中 $\Phi'$ 是通过将 $\Phi$ 中所有 $x$ 的出现替换为 $a$ 而得到的（并像往常一样避免变量捕获）。
类似地，对于任何 $f:\prd{x:A} B(x)$，我们有唯一性原理 $f\jdeq (\lam{x} f(x))$。
\index{uniqueness!principle!for dependent function types}%

作为一个例子，回忆 \cref{sec:universes}，那里有一个类型族 $\Fin:\nat\to\UU$，其值是标准有限集，元素为 $0_n,1_n,\dots,(n-1)_n : \Fin(n)$。
那么存在一个依值函数 $\fmax : \prd{n:\nat} \Fin(n+1)$，它返回每个非空有限类型的“最大”元素，即 $\fmax(n) \defeq n_{n+1}$。
\index{finite!sets, family of}%
正如 $\Fin$ 本身的情况，我们现在还不能定义 $\fmax$，但很快就可以；参见 \cref{ex:fin}。

另一类我们现在就能定义的重要依值函数类型，是那些在给定宇宙上\define{多态的（polymorphic）}函数。
\indexdef{function!polymorphic}%
\indexdef{polymorphic function}%
一个多态函数接受一个类型作为其参数之一，然后对该类型（或由它构造的其他类型）的元素进行操作。
\symlabel{idfunc}%
\indexdef{function!identity}%
\indexdef{identity!function}%
一个例子是多态恒等函数 $\idfunc : \prd{A:\UU} A \to A$，我们通过 $\idfunc{} \defeq \lam{A:\type}{x:A} x$ 来定义它。
（与 $\lambda$-抽象类似，除非被限定，否则 $\Pi$ 会自动将其作用域\index{scope}扩展到表达式的其余部分；因此 $\idfunc : \prd{A:\UU} A \to A$ 表示 $\idfunc : \prd{A:\UU} (A \to A)$。
这个约定在数学中虽不常见，但在类型论中很普遍。）

我们有时把依值函数的某些参数写作下标。
例如，我们可以用 $\idfunc[A](x) \defeq x$ 等价地定义多态恒等函数。
此外，如果某个参数能从上下文中推断出来，我们也可以将其完全省略。
例如，如果 $a:A$，那么写 $\idfunc(a)$ 是没有歧义的，因为要使 $\idfunc$ 能应用于 $a$，它就必须指 $\idfunc[A]$。

另一个不那么平凡的多态函数例子是``$\mathsf{swap}$''操作，它交换一个（柯里化的）两参数函数的参数顺序：
\[ \mathsf{swap} : \prd{A:\UU}{B:\UU}{C:\UU} (A\to B\to C) \to (B\to A \to C). \]
我们可以通过下式定义它：
\[ \mathsf{swap}(A,B,C,g) \defeq \lam{b}{a} g(a)(b). \]
我们也可以将类型参数等价地写作下标：
\[ \mathsf{swap}_{A,B,C}(g)(b,a) \defeq g(a,b). \]

注意，正如我们对普通函数所做的那样，我们使用柯里化来定义具有多个参数（例如 $\mathsf{swap}$）的依值函数。
然而，在依值的情形下，第二个定义域可能依赖于第一个定义域，并且陪域可能依赖于两者。
也就是说，给定 $A:\UU$ 以及类型族 $B : A \to \UU$ 和 $C : \prd{x:A}B(x) \to \UU$，我们可以构造具有两个参数的函数的类型 $\prd{x:A}{y : B(x)} C(x,y)$。
当 $B$ 为常值且等于 $A$ 时，我们可以简化写法，记为 $\prd{x,y:A}$；例如，$\mathsf{swap}$ 的类型也可以写成
\[ \mathsf{swap} : \prd{A,B,C:\UU} (A\to B\to C) \to (B\to A \to C). \]
最后，给定 $f:\prd{x:A}{y : B(x)} C(x,y)$ 以及参数 $a:A$ 和 $b:B(a)$，我们有 $f(a)(b) : C(a,b)$，如前所述，我们将其写作 $f(a,b) : C(a,b)$。

\index{type!dependent function|)}%
\index{function!dependent|)}%

\section{积类型}
\label{sec:finite-product-types}

给定类型 $A,B:\UU$，我们引入类型 $A\times B:\UU$，称之为它们的\define{Cartesian 积（cartesian product）}。
\indexsee{cartesian product}{type, product}%
\indexsee{type!cartesian product}{type, product}%
\index{type!product|(defstyle}%
\indexsee{product!of types}{type, product}%
我们还引入一种零元积类型，称为\define{单位类型（unit type）} $\unit : \UU$。
\indexsee{nullary!product}{type, unit}%
\indexsee{unit!type}{type, unit}%
\index{type!unit|(defstyle}%
我们希望 $A\times B$ 的元素是由 $a:A$ 和 $b:B$ 构成的对 $\tup{a}{b} : A \times B$，而 $\unit$ 的唯一元素是某个特定的对象 $\ttt : \unit$。
\indexdef{pair!ordered}%
然而，与集合论的做法不同——在集合论中，有序对先被定义为特定的集合，再汇总成 Cartesian 积——在类型论中，有序对与函数一样，都是原始概念。

\begin{rmk}\label{rmk:introducing-new-concepts}
  在类型论中，引入一种新类型有一个通用的模式。
  我们已经在前文见过这个模式\cref{sec:function-types,sec:pi-types}\footnote{上文对宇宙的描述是一个例外。}，因此值得强调一下它的通用形式。
  要规定一个类型，我们需要规定：
  \begin{enumerate}
  \item 如何通过\define{形成规则（formation rules）}来形成此类新类型。
    \indexdef{formation rule}%
    \index{rule!formation}%  
（例如，当 $A$ 是一个类型且 $B$ 是一个类型时，我们可以构造函数类型 $A \to B$。当 $A$ 是一个类型且对 $x:A$ 有 $B(x)$ 是一个类型时，我们可以构造依值函数类型 $\prd{x:A} B(x)$。）

  \item 如何构造该类型的元素。
    这些称为该类型的\define{构造子（constructors）}或\define{引入规则（introduction rules）}。
    \indexdef{constructor}%
    \indexdef{rule!introduction}%
    \indexdef{introduction rule}%
    （例如，函数类型有一个构造子：$\lambda$-抽象。回想一下，像 $f(x)\defeq 2x$ 这样的直接定义，等价于用 $\lambda$-抽象写成 $f\defeq \lam{x} 2x$。）

  \item 如何使用该类型的元素。
    这些称为该类型的\define{消去子（eliminators）}或\define{消去规则（elimination rules）}。
    \indexsee{rule!elimination}{eliminator}%
    \indexsee{elimination rule}{eliminator}%
    \indexdef{eliminator}%
    （例如，函数类型有一个消去子：函数应用。）

  \item
    一条\define{计算规则（computation rule）}\indexdef{computation rule}\footnote{也被称为\define{$\beta$-归约（$\beta $-reduction）}\index{beta-reduction@$\beta $-reduction|footstyle}}，它表达了消去子如何作用于构造子。
（例如，对于函数，计算规则断言 $(\lamu{x:A}\Phi)(a)$ 在判断相等性上等于将 $\Phi$ 中的 $x$ 替换为 $a$ 的结果。）

  \item
    一条可选的\define{唯一性原理（uniqueness principle）}\indexdef{uniqueness!principle}\footnote{也被称为\define{$\eta$-展开（$\eta $-expansion）}\index{eta-expansion@$\eta $-expansion|footstyle}}，它表达了进入或离开该类型的映射的唯一性。
    对于某些类型，唯一性原理刻画的是进入该类型的映射：它断言该类型的每个元素都由对其应用消去子所得的结果唯一确定，并且可以对这些结果应用构造子来重建该元素——从而表达了构造子如何作用于消去子，与计算规则对偶。
    （例如，对于函数，唯一性原理说任何函数 $f$ 在判断相等性上都等于``展开后的''函数 $\lamu{x} f(x)$，因此由其取值唯一确定。）
    对于另一些类型，唯一性原理则说，每一个\emph{从}该类型出发的映射（函数）都由某些数据唯一确定。（一个例子是 \cref{sec:coproduct-types} 中引入的余积类型，其唯一性原理在 \cref{sec:universal-properties} 中提到。）
    
    当唯一性原理不作为一条判断相等性规则时，它通常仍可从该类型的其他规则出发，作为一条\emph{命题}相等性得到证明。
    在这种情况下，我们称之为\define{命题唯一性原理（propositional uniqueness principle）}。
    \indexdef{uniqueness!principle, propositional}%
    \indexsee{propositional!uniqueness principle}{uniqueness principle, propositional}%
    （在后面的章节中，我们偶尔也会遇到\emph{命题计算规则（propositional computation rules）}。）
    \indexdef{computation rule!propositional}%
  \end{enumerate}
\cref{sec:syntax-more-formally} 中的推导规则正是据此组织和命名的；例如，可参见 \cref{sec:more-formal-pi}，其中每种可能性都得到了体现。
\end{rmk}

构造有序对的方式是显而易见的：给定 $a:A$ 和 $b:B$，我们可以构造 $(a,b):A\times B$。
类似地，构造 $\unit$ 的元素也有唯一的方式，即我们有 $\ttt:\unit$。
我们期望``$A\times B$ 的每个元素都是一个有序对''，这正是积类型的唯一性原理；我们并不将其作为类型论的一条规则来断言，但随后会将其作为一条命题相等性加以证明。

现在，我们该如何\emph{使用}这些对，也就是如何定义从积类型出发的函数呢？
我们先来考虑非依值函数 $f : A\times B \to C$ 的定义。
由于我们意图让 $A\times B$ 的元素仅为对，我们期望能够通过规定 $f$ 应用于对 $\tup{a}{b}$ 时的结果来定义这样的函数。
我们可以通过给出一个函数 $g : A \to B \to C$ 来规定这些结果。
于是，我们引入一条新的规则（积的消去规则），该规则说：对于任意这样的 $g$，我们可以定义一个函数 $f : A\times B \to C$，满足：
\[ f(\tup{a}{b}) \defeq g(a)(b). \]
我们在这里避免写作 $g(a,b)$，是为了强调 $g$ 并不是一个定义在积上的函数。
（不过，在本书后续部分，我们将会经常用 $g(a,b)$ 同时表示积上的函数和两个变量的柯里化函数。）
这个定义方程就是积类型的计算规则\index{computation rule!for product types}。

请注意，在集合论中，我们这样定义 $f$ 的依据是这样一个事实：$A\times B$ 中的每个元素都是有序对，因此只需在这些对上定义 $f$ 即可。
相比之下，类型论反转了这一情况：我们假定，只要规定了函数在对上的取值，该函数在 $A\times B$ 上就是良定义的；并且从这一点出发（更准确地说，从下面将要介绍的、针对依值函数的更一般版本出发），我们将能够\emph{证明} $A\times B$ 中的每个元素都是一个对。
从范畴论的角度来看，我们可以说，我们把积 $A\times B$ 定义为我们已经引入的``指数'' $B\to C$ 的左伴随。

举个例子，我们可以推导出\define{投影（projection）}
\indexsee{function!projection}{projection}%
\indexsee{component, of a pair}{projection}%
\indexdef{projection!from cartesian product type}%
函数
\symlabel{defn:proj}%
\begin{align*}
  \fst & :  A \times B \to A \\
  \snd & :  A \times B \to B
\end{align*}
其定义方程为：
\begin{align*}
  \fst(\tup{a}{b}) & \defeq  a \\
  \snd(\tup{a}{b}) & \defeq  b.
\end{align*}
%
\symlabel{defn:recursor-times}%
与其每次定义函数时都援引这一函数定义原则，另一种方法是只在一个最一般的情形下援引它一次，然后在所有其他情形中直接应用所得到的函数。
也就是说，我们可以定义类型为
\begin{equation}
  \rec{A\times B} : \prd{C:\UU}(A \to B \to C) \to A \times B \to C
\end{equation}
的函数，其定义方程为：
\[\rec{A\times B}(C,g,\tup{a}{b}) \defeq g(a)(b). \]
这样一来，我们就不必直接用定义方程来定义像 $\fst$ 和 $\snd$ 这样的函数，而是可以定义：
\begin{align*}
  \fst &\defeq \rec{A\times B}(A, \lam{a}{b} a)\\
  \snd &\defeq \rec{A\times B}(B, \lam{a}{b} b).
\end{align*}
我们称函数 $\rec{A\times B}$ 为积类型的 \define{递归子（recursor）}
\indexsee{recursor}{recursion principle}%
。``递归子''这个名字在这里有点不妥，因为并没有发生任何递归。这个名字来源于积类型是归纳类型一般框架下的一个退化例子；对于像自然数这样的类型，递归子才是真正递归的。我们也可以谈及 Cartesian 积的\define{递归原理（recursion principle）}，意思就是我们可以像上面那样，通过给出它在对上的值来定义一个函数 $f:A\times B\to C$。
\index{recursion principle!for cartesian product}%

递归子可以从投影函数推导出来，反之亦然，我们将此留作一个简单的练习。
% Ex: Derive from projections

\symlabel{defn:recursor-unit}%
单位类型也有一个递归子：
\[\rec{\unit} : \prd{C:\UU}C \to \unit \to C\]
其定义方程为：
\[ \rec{\unit}(C,c,\ttt) \defeq c. \]
虽然我们为了保持类型定义的模式一致性而包含了它，但 $\unit$ 的递归子完全没有用，因为我们本可以直接定义一个这样的函数，只需简单地忽略类型为 $\unit$ 的那个参数即可。

为了能在积类型上定义\emph{依值}函数，我们必须推广递归子。给定 $C: A \times B \to \UU$，我们可以通过提供一个函数
\narrowequation{
 g : \prd{x:A}\prd{y:B} C(\tup{x}{y})
}
来定义一个函数 $f : \prd{x : A \times B} C(x)$，其定义方程为：
\[ f(\tup x y) \defeq g(x)(y). \]
例如，用这种方法我们可以证明命题唯一性原理，即 $A\times B$ 中的每个元素都等于一个对。
\index{uniqueness!principle, propositional!for product types}%
具体来说，我们可以构造一个函数：
\[ \uniq{A\times B} : \prd{x:A \times B} (\id[A\times B]{\tup{\fst {(x)}}{\snd {(x)}}}{x}). \]
这里我们用到了相等类型，我们将在下文 \cref{sec:identity-types} 中介绍它。
不过，我们现在只需要知道，对于任何 $x:A$，存在一个自反性元素 $\refl{x} : \id[A]{x}{x}$。
据此，我们可以定义：
\label{uniquenessproduct}
\[ \uniq{A\times B}(\tup{a}{b}) \defeq \refl{\tup{a}{b}}. \]
这个构造之所以有效，是因为在 $x \defeq \tup{a}{b}$ 的情形下，我们可以利用投影的定义方程来计算：
\[ \tup{\fst(\tup{a}{b})}{\snd{(\tup{a}{b})}} \jdeq \tup{a}{b} \]
因此，
\[ \refl{\tup{a}{b}} : \id{\tup{\fst(\tup{a}{b})}{\snd{(\tup{a}{b})}}}{\tup{a}{b}} \]
是良类型的，因为等号两边是判断相等的。

更一般地说，以这种方式定义依值函数的能力意味着：要证明一个性质对于积类型的所有元素都成立，只需对它的规范元素——即有序对——证明该性质就够了。
当我们后面遇到像自然数这样的归纳类型时，与之类似的性质就是进行归纳证明的能力。
因此，如果我们仿照上面的做法，将这一原则在最一般的情形下应用一次，我们称所得到的函数为积类型的\define{归纳（induction）}：给定 $A,B : \UU$，我们有
\symlabel{defn:induction-times}%
\[ \ind{A\times B} : \prd{C:A \times B \to \UU}
\Parens{\prd{x:A}{y:B} C(\tup{x}{y})} \to \prd{x:A \times B} C(x) \]
其定义方程为
\[ \ind{A\times B}(C,g,\tup{a}{b}) \defeq g(a)(b). \]
类似地，我们也可以说，定义在有序对上的一个依值函数是从 Cartesian 积的\define{归纳原理（induction principle）}
\index{induction principle}%
\index{induction principle!for product}%
得到的。
容易看出，递归子只是归纳在类型族 $C$ 为常值时的特例。
因为归纳描述了如何使用一个积类型的元素，所以归纳也被称为\define{（依值）消去子（(dependent) eliminator）}，
\indexsee{eliminator!of inductive type!dependent}{induction principle}%
而递归则被称为\define{非依值消去子（non-dependent eliminator）}。
\indexsee{eliminator!of inductive type!non-dependent}{recursion principle}%
\indexsee{non-dependent eliminator}{recursion principle}%
\indexsee{dependent eliminator}{induction principle}%

% We can read induction propositionally as saying that a property which
% is true for all pairs holds for all elements of the product type.

单位类型的归纳实际上比递归子更有用：
\symlabel{defn:induction-unit}%
\[ \ind{\unit} : \prd{C:\unit \to \UU} C(\ttt) \to \prd{x:\unit}C(x)\]
其定义方程为
\[ \ind{\unit}(C,c,\ttt) \defeq c. \]
借助归纳，我们可以证明 $\unit$ 的命题唯一性原理，该原理断言 $\unit$ 中唯一的实例就是 $\ttt$。
也就是说，我们可以构造
\label{uniquenessunit}
\[\uniq{\unit} : \prd{x:\unit} \id{x}{\ttt} \]
通过定义方程：
\[\uniq{\unit}(\ttt) \defeq \refl{\ttt} \]
或等价地通过归纳：
\[\uniq{\unit} \defeq \ind{\unit}(\lamu{x:\unit} \id{x}{\ttt},\refl{\ttt}). \]

\index{type!product|)}%
\index{type!unit|)}%

\section{依值对类型 \texorpdfstring{($\Sigma$ 类型)}{(Σ 类型)}}
\label{sec:sigma-types}

\index{type!dependent pair|(defstyle}%
\indexsee{type!dependent sum}{type, dependent pair}%
\indexsee{type!Sigma-@$\Sigma$-}{type, dependent pair}%
\indexsee{Sigma-type@$\Sigma$-type}{type, dependent pair}%
\indexsee{sum!dependent}{type, dependent pair}%

正如我们将函数类型（\cref{sec:function-types}）推广为依值函数类型（\cref{sec:pi-types}）一样，将 \cref{sec:finite-product-types} 中引入的积类型加以推广，使得一个``对''中第二个分量的类型可以随第一个分量的选择不同而变化，这也常常很有用。
这被称为\define{依值对类型（dependent pair type）}，或\define{$\Sigma$ 类型（$\Sigma$-type）}，因为在集合论中，它对应于给定类型上的索引和（在余积或不相交并的意义上）。

给定一个类型 $A:\UU$ 以及一个类型族 $B : A \to \UU$，依值对类型写作 $\sm{x:A} B(x) : \UU$。
其他的记法还有
\[ \tsm{x:A} B(x) \hspace{2cm} \dsm{x:A}B(x) \hspace{2cm} \lsm{x:A} B(x). \]
与 $\lambda$-抽象和 $\Pi$ 类型这样的绑定构造类似，$\Sigma$ 也会自动作用于表达式的其余部分，除非被明确界定；例如，$\sm{x:A} B(x) \to C$ 表示 $\sm{x:A} (B(x) \to C)$。

\symlabel{defn:dependent-pair}%
\indexdef{pair!dependent}%
构造依值对类型元素的方式是通过配对：给定 $a:A$ 和 $b:B(a)$，我们有 $\tup{a}{b} : \sm{x:A} B(x)$。
如果 $B$ 是常值的，那么依值对类型就是通常的 Cartesian 积类型：
\[ \Parens{\sm{x:A} B} \jdeq (A \times B).\]
$\Sigma$ 类型上的所有构造都是积类型相应构造的直接推广，其中非依值的函数常常被依值函数所取代。

例如，$\Sigma$ 类型的递归原理%
\index{recursion principle!for dependent pair type}
说的是，要定义一个从 $\Sigma$ 类型$f : (\sm{x:A} B(x)) \to C$出发的非依值函数，我们给出一个函数
$g : \prd{x:A} B(x) \to C$，然后通过定义方程
\[ f(\tup{a}{b}) \defeq g(a)(b). \]
来定义 $f$。
\indexdef{projection!from dependent pair type}%
例如，我们可以推导出 $\Sigma$ 类型的第一投影：
\symlabel{defn:dependent-proj1}%
\begin{equation*}
  \fst : \Parens{\sm{x : A}B(x)} \to A
\end{equation*}
其定义方程为
\begin{equation*}
  \fst(\tup{a}{b}) \defeq a.
\end{equation*}
然而，由于一个配对
\narrowequation{
  (a,b):\sm{x:A} B(x)
}
的第二分量的类型是 $B(a)$，第二投影必须是一个\emph{依值}函数，其类型涉及第一投影函数：
\symlabel{defn:dependent-proj2}%
\[ \snd : \prd{p:\sm{x : A}B(x)}B(\fst(p)). \]
为此，我们需要 $\Sigma$ 类型的\emph{归纳}原理%
\index{induction principle!for dependent pair type}
（即``依值消去子''）。
该原理说，要构造一个从 $\Sigma$ 类型出发、取值于某个类型族$C : (\sm{x:A} B(x)) \to \UU$的依值函数，我们需要一个函数
\[ g : \prd{a:A}{b:B(a)} C(\tup{a}{b}). \]
然后我们可以导出一个函数
\[ f : \prd{p : \sm{x:A}B(x)} C(p) \]
其定义方程为\index{computation rule!for dependent pair type}
\[ f(\tup{a}{b}) \defeq g(a)(b).\]
取$C(p)\defeq B(\fst(p))$，我们可以定义
\narrowequation{
\snd : \prd{p:\sm{x : A}B(x)}B(\fst(p))
}
它满足如下自然的方程：
\[ \snd(\tup{a}{b})  \defeq  b. \]
为了确认这是正确的，我们注意到，利用 $\fst$ 的定义方程可得$B (\fst(\tup{a}{b})) \jdeq B(a)$，而确实有 $b : B(a)$。

我们可以将递归原理和归纳原理整合为 $\Sigma$ 的递归子：
\symlabel{defn:recursor-sm}%
\[ \rec{\sm{x:A}B(x)} : \dprd{C:\UU}\Parens{\tprd{x:A} B(x) \to C} \to
\Parens{\tsm{x:A}B(x)} \to C \]
其定义方程为
\[ \rec{\sm{x:A}B(x)}(C,g,\tup{a}{b}) \defeq g(a)(b) \]
以及相应的归纳算子：
\symlabel{defn:induction-sm}%
\begin{narrowmultline*}
  \ind{\sm{x:A}B(x)} : \narrowbreak
    \dprd{C:(\sm{x:A} B(x)) \to \UU}
    \Parens{\tprd{a:A}{b:B(a)} C(\tup{a}{b})}
    \to \dprd{p : \sm{x:A}B(x)} C(p)
\end{narrowmultline*}
其定义方程为
\[ \ind{\sm{x:A}B(x)}(C,g,\tup{a}{b}) \defeq g(a)(b). \]
和之前一样，当族 $C$ 为常量时，递归子是归纳算子的特例。

再举一个例子，考虑以下原理，其中 $A$ 和 $B$ 是类型，且 $R:A\to B\to \UU$：
\[ \ac : \Parens{\tprd{x:A} \tsm{y :B} R(x,y)} \to
\Parens{\tsm{f:A\to B} \tprd{x:A} R(x,f(x))}.
\]
我们可以将 $R$ 视为 $A$ 和 $B$ 之间的一个``证明相关关系''
\index{mathematics!proof-relevant}%
，其中 $R(a,b)$ 是 $a:A$ 与 $b:B$ 相关的见证的类型。
那么，$\ac$ 直观上说的是：如果我们有一个依值函数 $g$，它为每个 $a:A$ 指派一个依值对 $(b,r)$，其中 $b:B$ 且 $r:R(a,b)$，那么我们就能得到一个函数 $f:A\to B$ 和一个依值函数，后者为每个 $a:A$ 指派一个 $R(a,f(a))$ 的见证。
直觉告诉我们，只需将 $g$ 的值拆分为各自的分量即可。
事实上，利用我们刚刚定义的投影，我们可以如下定义：
\[ \ac(g) \defeq \Parens{\lamu{x:A} \fst(g(x)),\, \lamu{x:A} \snd(g(x))}. \]
为了验证这是良类型的，注意如果$g:\prd{x:A} \sm{y :B} R(x,y)$，我们有
\begin{align*}
\lamu{x:A} \fst(g(x)) &: A \to  B, \\
\lamu{x:A} \snd(g(x)) &: \tprd{x:A} R(x,\fst(g(x))).
\end{align*}
此外，类型$\prd{x:A} R(x,\fst(g(x)))$是将 $\ac$ 陪域中参与求和的类型族$\lamu{f:A\to B} \tprd{x:A} R(x,f(x))$应用于函数$\lamu{x:A} \fst(g(x))$所得的结果：
\[ \tprd{x:A} R(x,\fst(g(x))) \jdeq
\Parens{\lamu{f:A\to B} \tprd{x:A} R(x,f(x))}\big(\lamu{x:A} \fst(g(x))\big). \]
因此，我们得到
\[ \Parens{\lamu{x:A} \fst(g(x)),\, \lamu{x:A} \snd(g(x))} : \tsm{f:A\to B} \tprd{x:A} R(x,f(x))\]
这正是所要求的。

如果我们把 $\Pi$ 读作``对所有''，把 $\Sigma$ 读作``存在''，那么函数 $\ac$ 的类型表达的是：
\emph{如果对所有 $x:A$，都存在 $y:B$ 使得 $R(x,y)$ 成立，那么存在一个函数 $f : A \to B$，使得对所有 $x:A$ 我们都有 $R(x,f(x))$}。
由于这听起来像是选择公理的一个版本，函数 \ac 传统上被称为\define{类型论选择公理}，而且正如我们刚才所展示的，它可以直接从类型论的规则中证明，而无须作为公理引入。
\index{axiom!of choice!type-theoretic}%
然而，请注意这里实际上并不涉及任何``选择''，因为选择已经在前提中给出了：我们要做的仅仅是将其分解为两个函数——一个代表选择，另一个代表其正确性。
在\cref{sec:axiom-choice}中，我们将给出``选择公理''的另一种表述，它更接近通常的形式。

依值对类型常被用来定义各种数学结构的类型，这些结构通常由若干相互依值的数据构成。
举个简单的例子，假设我们想定义一个\define{原群}\indexdef{magma}为一个类型 $A$ 附带一个二元运算 $m:A\to A\to A$。
短语``附带（together with）''\index{together with}（及其同义词``配备（equipped with）''\index{equipped with}）的精确含义是：``一个原群''就是一个由类型 $A:\UU$ 和运算 $m:A\to A\to A$ 组成的\emph{对} $(A,m)$。
由于这个对的第二个分量 $m$ 的类型 $A\to A\to A$ 依赖于其第一个分量 $A$，这样的对属于一个依值对类型。
因此，定义``一个原群是一个类型 $A$ 附带一个二元运算 $m:A\to A\to A$''应理解为将\emph{原群的类型}定义为
\[ \mathsf{Magma} \defeq \sm{A:\UU} (A\to A\to A). \]
给定一个原群，我们用第一投影 $\proj1$ 提取它的底层类型（即``载体（carrier）''\index{carrier}），用第二投影 $\proj2$ 提取它的运算。
当然，由多于两个数据构成的结构需要用到迭代的对类型，其中的依赖关系可能只是部分的；例如，带点原群（即附带基点 $e:A$ 的原群 $(A,m)$）的类型是
\[ \mathsf{PointedMagma} \defeq \sm{A:\UU} (A\to A\to A) \times A. \]
我们通常还想为这样的结构施加公理，例如将一个带点原群变成一个幺半群或一个群。
这同样可以利用 $\Sigma$ 类型来实现；参见\cref{sec:pat}。

在本书后续部分，我们有时会明确写出这类定义，但最终我们信任读者能够自行将它们从英文描述翻译成 $\Sigma$ 类型。
我们通常也遵循常见的数学惯例，用同一个字母来表示这类结构及其载体（这相当于在记法中省略相应的投影函数）：也就是说，我们会谈论一个原群 $A$ 及其运算 $m:A\to A\to A$。

注意，$\mathsf{PointedMagma}$ 的典范元素具有形式 $(A,(m,e))$，其中 $A:\UU$，$m:A\to A\to A$，且 $e:A$。
由于这类迭代的 $\Sigma$ 类型频繁出现，我们采用通常的有序三元组、四元组等记法来表示（可能依值的）嵌套对，且约定为右结合。
也就是说，我们有 $(x,y,z) \defeq (x,(y,z))$ 与 $(x,y,z,w)\defeq (x,(y,(z,w)))$，等等。

\index{type!dependent pair|)}%

\section{余积类型}
\label{sec:coproduct-types}

给定 $A,B:\UU$，我们引入它们的\define{余积}\indexdef{coproduct}类型 $A+B:\UU$。
\indexsee{coproduct}{type, coproduct}%
\index{type!coproduct|(defstyle}%
\indexsee{disjoint!sum}{type, coproduct}%
\indexsee{disjoint!union}{type, coproduct}%
\indexsee{sum!disjoint}{type, coproduct}%
\indexsee{union!disjoint}{type, coproduct}%
这对应于集合论中的\emph{不交并}，我们有时也用这个名称来称呼它。
在类型论中，和函数类型、积类型的情况一样，余积必须作为一个基本构造来引入，因为此前并没有预先给定的``类型的并''的概念。
我们同时还引入一个零元的版本：\define{空类型 $\emptyt:\UU$}。
\indexsee{nullary!coproduct}{type, empty}%
\indexsee{empty type}{type, empty}%
\index{type!empty|(defstyle}%

构造 $A+B$ 的元素有两种方式：对于 $a:A$，可以构造 $\inl(a) : A+B$；对于 $b:B$，可以构造 $\inr(b):A+B$。
（名称 $\inl$ 和 $\inr$ 分别是``左嵌入（left injection）''和``右嵌入（right injection）''的缩写。）
空类型则没有任何构造其元素的方式。

\index{recursion principle!for coproduct}
要构造一个非依值函数 $f : A+B \to C$，我们需要两个函数 $g_0 : A \to C$ 和 $g_1 : B \to C$。
然后 $f$ 由以下定义方程定义：
\begin{align*}
  f(\inl(a)) &\defeq g_0(a), \\
  f(\inr(b)) &\defeq g_1(b).
\end{align*}
也就是说，函数 $f$ 是通过\define{分情形讨论}\indexdef{case analysis}来定义的。
%
和之前一样，我们可以推导出递归子：
\symlabel{defn:recursor-plus}%
\[ \rec{A+B} : \dprd{C:\UU}(A \to C) \to (B\to C) \to A+B \to C\]
其定义方程为
\begin{align*}
\rec{A+B}(C,g_0,g_1,\inl(a)) &\defeq g_0(a), \\
\rec{A+B}(C,g_0,g_1,\inr(b)) &\defeq g_1(b).
\end{align*}

\index{recursion principle!for empty type}
对于空类型，我们总是可以构造一个函数 $f : \emptyt \to C$，而无需给出任何定义方程，因为 $\emptyt$ 中根本没有元素可供定义 $f$。
因此，$\emptyt$ 的递归子是
\symlabel{defn:recursor-emptyt}%
\[\rec{\emptyt} : \tprd{C:\UU} \emptyt \to C,\]
它构造了从空类型到任何其他类型的典范函数。
在逻辑上，它对应于原则 \textit{ex falso quodlibet}。
\index{ex falso quodlibet@\textit{ex falso quodlibet}}

\index{induction principle!for coproduct}
要构造一个从余积出发的依值函数 $f:\prd{x:A+B}C(x)$，我们假设给定了类型族
$C: (A + B) \to \UU$，
并要求
\begin{align*}
  g_0 &: \prd{a:A} C(\inl(a)), \\
  g_1 &: \prd{b:B} C(\inr(b)).
\end{align*}
这样我们就得到了 $f$，其定义方程是：\index{computation rule!for coproduct type}
\begin{align*}
  f(\inl(a)) &\defeq g_0(a), \\
  f(\inr(b)) &\defeq g_1(b).
\end{align*}
我们将这一模式封装为余积的归纳原理：
\symlabel{defn:induction-plus}%
\begin{narrowmultline*}
  \ind{A+B} :
  \dprd{C: (A + B) \to \UU}
  \Parens{\tprd{a:A} C(\inl(a))} \to \narrowbreak
  \Parens{\tprd{b:B} C(\inr(b))} \to \tprd{x:A+B}C(x). 
\end{narrowmultline*}
和之前一样，当族 $C$ 是常值的时候，我们就得到了递归子。

\index{induction principle!for empty type}
空类型的归纳原理
\symlabel{defn:induction-emptyt}%
\[ \ind{\emptyt} : \prd{C:\emptyt \to \UU}{z:\emptyt} C(z) \]
为我们提供了一种从空类型出发定义平凡依值函数的方法。% In the presence of $\eta$-equality it is derivable
% from the recursor.
% ex

\index{type!coproduct|)}%
\index{type!empty|)}%

\section{布尔类型}
\label{sec:type-booleans}

\indexsee{boolean!type of}{type of booleans}%
\index{type!of booleans|(defstyle}%
布尔类型 $\bool:\UU$ 旨在恰好有两个元素
$\bfalse,\btrue : \bool$。显然，我们可以用余积
% this one results in a warning message just because it's on the same page as the previous entry 
% for {type!coproduct}, so it's not our fault
\index{type!coproduct}%
与单位\index{type!unit}类型把它构造出来，即 $\unit + \unit$。不过，
由于它很常用，我们在此给出显式规则。
实际上，我们将看到反过来也是可以的：我们可以从 $\Sigma$-类型和 $\bool$ 出发来推导出二元余积。

\index{recursion principle!for type of booleans}
要推导出一个函数 $f : \bool \to C$，我们需要 $c_0,c_1 : C$，并附加定义方程
\begin{align*}
  f(\bfalse) &\defeq c_0, \\
  f(\btrue)  &\defeq c_1.
\end{align*}
这个递归子对应着函数式编程中的 if-then-else 结构：
\symlabel{defn:recursor-bool}%
\[ \rec{\bool} : \prd{C:\UU}  C \to C \to \bool \to C \]
其定义方程是
\begin{align*}
  \rec{\bool}(C,c_0,c_1,\bfalse) &\defeq c_0, \\
  \rec{\bool}(C,c_0,c_1,\btrue)  &\defeq c_1.
\end{align*}

\index{induction principle!for type of booleans}
给定 $C : \bool \to \UU$，要推导一个依值函数
$f : \prd{x:\bool}C(x)$，我们需要 $c_0:C(\bfalse)$ 和 $c_1 : C(\btrue)$，在此情形下即可给出定义方程
\begin{align*}
  f(\bfalse) &\defeq c_0, \\
  f(\btrue)  &\defeq c_1.
\end{align*}
我们将此封装为归纳原理
\symlabel{defn:induction-bool}%
\[ \ind{\bool} : \dprd{C:\bool \to \UU}  C(\bfalse) \to C(\btrue)
\to \tprd{x:\bool} C(x) \]
其定义方程为
\begin{align*}
  \ind{\bool}(C,c_0,c_1,\bfalse) &\defeq c_0, \\
  \ind{\bool}(C,c_0,c_1,\btrue)  &\defeq c_1.
\end{align*}

举个例子，使用归纳原理我们可以推导出——正如我们所料—— $\bool$ 的每个元素要么是 $\btrue$，要么是 $\bfalse$。
和之前一样，为了表述这一点，我们需要用到尚未介绍的相等类型；但我们只需用到``每个对象都通过 $\refl{x}:x=x$ 等于自身''这一事实。
于是，我们构造一个 \begin{equation}\label{thm:allbool-trueorfalse}
  \prd{x:\bool}(x=\bfalse)+(x=\btrue),
\end{equation} 的元素，即一个函数，它能给每个 $x:\bool$ 指派一个相等 $x=\bfalse$ 或 $x=\btrue$。
我们使用 \bool 的归纳原理来定义这个元素，给定 $C(x) \defeq (x=\bfalse)+(x=\btrue)$；
两个输入分别是 $\inl(\refl{\bfalse}) : C(\bfalse)$ 和 $\inr(\refl{\btrue}):C(\btrue)$。
换句话说，我们所构造的 \eqref{thm:allbool-trueorfalse} 的元素是
\[ \ind{\bool}\big(\lam{x}(x=\bfalse)+(x=\btrue),\, \inl(\refl{\bfalse}),\, \inr(\refl{\btrue})\big). \]

我们曾提到，$\Sigma$ 类型可以类比为带索引的无交并，而余积则是二元无交并。
自然地，我们会期望，一个二元无交并 $A+B$ 可以构造为在两元素类型 \bool 上的带索引的无交并。
为此，我们需要一个类型族 $P:\bool\to\type$，使得 $P(\bfalse)\jdeq A$ 且 $P(\btrue)\jdeq B$。
事实上，我们能恰好通过 $\bool$ 的递归原理得到这样一个族。
\index{type!family of}%
（利用宇宙 $\UU$ 本身也是一个类型这一事实，通过归纳和递归来定义\emph{类型族（type families）}，是类型论中一个微妙且重要的方面。）
因此，我们本可以定义
\index{type!coproduct}%
\[ A + B \defeq \sm{x:\bool} \rec{\bool}(\UU,A,B,x) \]
其中
\begin{align*}
  \inl(a) &\defeq \tup{\bfalse}{a}, \\
  \inr(b) &\defeq \tup{\btrue}{b}.
\end{align*}
我们将其留作练习：从这一定义出发推导余积类型的归纳原理。
（另见 \cref{ex:sum-via-bool,sec:appetizer-univalence}。）

我们可以将同样的想法应用到积类型和 $\Pi$ 类型上：我们本可以定义
\[ A \times B \defeq \prd{x:\bool}\rec{\bool}(\UU,A,B,x). \]
此时，可以借助 \bool 的归纳原理来构造配对：
\[ \tup{a}{b} \defeq \ind{\bool}(\rec{\bool}(\UU,A,B),a,b) \]
而投影则是直接的应用：
\begin{align*}
  \fst(p) &\defeq p(\bfalse), \\
  \snd(p) &\defeq p(\btrue).
\end{align*}
为以这种方式定义的二元积推导归纳原理要稍微复杂一些，并且需要用到我们将在 \cref{sec:compute-pi} 中引入的函数外延性。
此外，我们也得不到相同的判断相等性；详见 \cref{ex:prod-via-bool}。
这是在用一种类型编码另一种类型时会反复出现的问题；我们将在 \cref{sec:htpy-inductive} 中回头讨论它。

我们偶尔会将 $\bool$ 的元素 $\bfalse$ 和 $\btrue$ 分别称为``假''和``真''。
然而，需要注意的是，与经典\index{mathematics!classical}数学不同，我们并不把 $\bool$ 的元素当作真值
\index{value!truth}%
也不将其当作命题。
（相反，我们将命题与类型等同起来；见 \cref{sec:pat}。）
特别地，类型 $A \to \bool$ 通常并非 $A$ 的幂集\index{power set}；它仅表示 $A$ 的那些``可判定''子集（见 \cref{cha:logic}）。
\index{decidable!subset}%

\index{type!of booleans|)}%

\section{自然数}
\label{sec:inductive-types}

\indexsee{type!of natural numbers}{natural numbers}%
\index{natural numbers|(defstyle}%
\indexsee{number!natural}{natural numbers}%
至此，我们已经有了通过抽象运算来构造新类型的规则，但要进行具体的数学工作，我们还需要一些具体的类型，例如数的类型。
其中最基本的便是自然数类型 $\nat : \UU$；一旦有了它，我们就能构造整数、有理数、实数等等（参见\cref{cha:real-numbers}）。

$\nat$ 的元素由 $0 : \nat$\indexdef{zero} 与后继（successor）\indexdef{successor}运算 $\suc : \nat \to \nat$ 构造而成。
在表示自然数时，我们采用常见的十进制记法：$1 \defeq \suc(0)$, $2 \defeq \suc(1)$, $3 \defeq \suc(2)$, \dots。

自然数的本质属性在于，我们可以通过递归来定义函数，并通过归纳来进行证明——此处``递归''与``归纳''这两个词具有更为熟悉的含义。
\index{recursion principle!for natural numbers}%
要通过递归构造一个从自然数出发的非依值函数 $f : \nat \to C$，只需提供一个起始点 $c_0 : C$ 以及一个``下一步''函数 $c_s : \nat \to C \to C$。
这样便得到了 $f$，其定义方程为\index{computation rule!for natural numbers}
\begin{align*}
  f(0) &\defeq c_0, \\
  f(\suc(n)) &\defeq c_s(n,f(n)).
\end{align*}
我们称 $f$ 是由\define{原始递归（primitive recursion）}定义的。
\indexdef{primitive!recursion}%
\indexdef{recursion!primitive}%

我们来看一个例子：如何定义一个将其参数加倍的函数。
此时取 $C\defeq \nat$。
我们首先需要给出 $\dbl(0)$ 的值，这很简单：令 $c_0 \defeq 0$。
接着，要计算自然数 $n$ 对应的 $\dbl(\suc(n))$，我们先计算 $\dbl(n)$ 的值，然后执行两次后继运算。
这由递推关系\index{recurrence} $c_s(n,y) \defeq \suc(\suc(y))$ 所刻画。
注意，$c_s$ 的第二个参数 $y$ 代表了\emph{递归调用（recursive call）}\index{recursive call} $\dbl(n)$ 的结果。

因此，以这种方式通过原始递归定义 $\dbl:\nat\to\nat$，我们得到如下定义方程：
\begin{align*}
  \dbl(0) &\defeq 0\\
  \dbl(\suc(n)) &\defeq \suc(\suc(\dbl(n))).
\end{align*}
这确实具有正确的计算行为：例如，我们有
\begin{align*}
  \dbl(2) &\jdeq \dbl(\suc(\suc(0)))\\
  & \jdeq c_s(\suc(0), \dbl(\suc(0))) \\
                 & \jdeq \suc(\suc(\dbl(\suc(0)))) \\
                 & \jdeq \suc(\suc(c_s(0,\dbl(0)))) \\
                 & \jdeq \suc(\suc(\suc(\suc(\dbl(0))))) \\
                 & \jdeq \suc(\suc(\suc(\suc(c_0)))) \\
                 & \jdeq \suc(\suc(\suc(\suc(0))))\\
                 &\jdeq 4.
\end{align*}
我们也可以通过原始递归来定义多变量函数，只需对函数进行柯里化，并允许 $C$ 自身是一个函数类型。
\indexdef{addition!of natural numbers}
例如，我们定义加法 $\add : \nat \to \nat \to \nat$，取 $C \defeq \nat \to \nat$ 及如下的``起始点''与``下一步''数据：
\begin{align*}
  c_0 & : \nat \to \nat \\
  c_0 (n) & \defeq n \\
  c_s & : \nat \to (\nat \to \nat) \to (\nat \to \nat) \\
  c_s(m,g)(n) & \defeq \suc(g(n)).
\end{align*}
由此我们得到了满足如下定义相等性的 $\add : \nat \to \nat \to \nat$：
\begin{align*}
  \add(0,n) &\jdeq n \\
  \add(\suc(m),n) &\jdeq \suc(\add(m,n)). 
\end{align*}
按惯例，我们将 $\add(m,n)$ 写作 $m+n$。
请读者自行验证 $2+2\jdeq 4$。
% ex: define multiplication and exponentiation.

与之前的例子一样，我们可以将原始递归原理打包为一个递归子（recursor）：
\[\rec{\nat}  : \dprd{C:\UU} C \to (\nat \to C \to C) \to \nat \to C \]
其定义方程为
\symlabel{defn:recursor-nat}%
\begin{align*}
\rec{\nat}(C,c_0,c_s,0)  &\defeq c_0, \\
\rec{\nat}(C,c_0,c_s,\suc(n)) &\defeq c_s(n,\rec{\nat}(C,c_0,c_s,n)).
\end{align*}
%ex derive rec from it
利用 $\rec{\nat}$，我们可以将 $\dbl$ 和 $\add$ 表示如下：
\begin{align}
\dbl &\defeq \rec\nat\big(\nat,\, 0,\, \lamu{n:\nat}{y:\nat} \suc(\suc(y))\big) \label{eq:dbl-as-rec}\\
\add &\defeq \rec{\nat}\big(\nat \to \nat,\, \lamu{n:\nat} n,\, \lamu{m:\nat}{g:\nat \to \nat}{n :\nat} \suc(g(n))\big).
\end{align}
当然，所有仅使用原始递归原理可定义的函数都是\emph{可计算的}。
（然而，高阶函数类型——即以其他函数为参数的函数——的存在，确实意味着我们可以定义比通常的原始递归函数更多的函数；例如参见\cref{ex:ackermann}。）
这在构造性数学中是合适的；
\index{mathematics!constructive}%
在\cref{sec:intuitionism,sec:axiom-choice}中，我们将看到如何扩充类型论，以便定义更一般的数学函数。

\index{induction principle!for natural numbers}
我们现在沿用与其他类型相同的方法，将原始递归推广至依值函数，从而得到一个\emph{归纳原理}。
因此，假设给定一个类型族 $C : \nat \to \UU$、一个元素 $c_0 : C(0)$ 以及一个函数 $c_s : \prd{n:\nat} C(n) \to C(\suc(n))$；那么我们可以构造 $f : \prd{n:\nat} C(n)$，其定义方程为：\index{computation rule!for natural numbers}
\begin{align*}
  f(0) &\defeq c_0, \\
  f(\suc(n)) &\defeq c_s(n,f(n)).
\end{align*}
我们也可以将其打包为单个函数
\symlabel{defn:induction-nat}%
\[\ind{\nat}  : \dprd{C:\nat\to \UU} C(0) \to \Parens{\tprd{n : \nat} C(n) \to C(\suc(n))} \to \tprd{n : \nat} C(n) \]
其定义方程为
\begin{align*}
\ind{\nat}(C,c_0,c_s,0)  &\defeq c_0, \\
\ind{\nat}(C,c_0,c_s,\suc(n)) &\defeq c_s(n,\ind{\nat}(C,c_0,c_s,n)).
\end{align*}
此处我们终于看到了与经典归纳证明概念的联系。
回顾一下，在类型论中，我们用类型来表示命题，并通过给出对应类型的实例来证明该命题。
特别地，自然数的一个\emph{性质}由一个类型族 $P:\nat\to\type$ 来表示。
从这个角度来看，上述归纳原理是说：如果我们能证明 $P(0)$，并且对于任何 $n$，我们都能在假设 $P(n)$ 成立的前提下证明 $P(\suc(n))$，那么对任意 $n$ 都有 $P(n)$。
这当然正是通常的自然数归纳证明原理。

\index{associativity!of addition!of natural numbers}
作为例子，考虑如何表示 $+$ 满足结合律的一个显式证明。
（我们实际上不会以这种风格书写证明，但它作为例子有助于理解归纳在类型论中是如何被形式表示的。）
为推导
\[\assoc : \prd{i,j,k:\nat} \id{i + (j + k)}{(i + j) + k}, \]
只需提供
\[ \assoc_0 :  \prd{j,k:\nat} \id{0 + (j + k)}{(0+ j) + k} \]
与
\begin{narrowmultline*}
  \assoc_s  : \prd{i:\nat} \left(\prd{j,k:\nat} \id{i + (j + k)}{(i + j) + k}\right)
   \narrowbreak
   \to \prd{j,k:\nat} \id{\suc(i) + (j + k)}{(\suc(i) + j) + k}.
\end{narrowmultline*}。
为推导 $\assoc_0$，回顾 $0+n \jdeq n$，因此 $0 + (j + k) \jdeq j+k \jdeq (0+ j) + k$。
于是我们可以直接令
\[ \assoc_0(j,k) \defeq \refl{j+k}. \]
对于 $\assoc_s$，回顾加法的定义给出 $\suc(m)+n \jdeq \suc(m+n)$，因此
\begin{align*}
   \suc(i) + (j + k)  &\jdeq \suc(i+(j+k)) \qquad\text{and}\\
   (\suc(i)+j)+k &\jdeq \suc((i+j)+k).
\end{align*}。
这样一来，$\assoc_s$ 的输出类型便等价地变为 $\id{\suc(i+(j+k))}{\suc((i+j)+k)}$。
但其输入（即``归纳假设''）
\index{hypothesis!inductive}%
\index{inductive!hypothesis}%
能给出 $\id{i+(j+k)}{(i+j)+k}$，因此只需援引以下事实：如果两个自然数相等，那么它们的后继也相等。
（我们将在 \cref{lem:map} 中利用相等类型的归纳原理证明这个显然的事实。）
我们将后一事实记作
$\apfunc{\suc} : %\prd{m,n:\nat}
(\id[\nat]{m}{n}) \to (\id[\nat]{\suc(m)}{\suc(n)})$，于是可以定义
\[\assoc_s(i,h,j,k) \defeq \apfunc{\suc}( %n+(j+k),(n+j)+k,
h(j,k)). \]
将这些与 $\ind{\nat}$ 组合起来，我们就得到了结合律的一个证明。

\index{natural numbers|)}%

\section{模式匹配与递归}
\label{sec:pattern-matching}

\index{pattern matching|(defstyle}%
\indexsee{matching}{pattern matching}%
\index{definition!by pattern matching|(}%
与迄今为止考虑过的类型相比，自然数引入了一层额外的微妙之处。
以余积为例，我们可以用递归子定义函数 $f:A+B\to C$：
\[ f \defeq \rec{A+B}(C, g_0, g_1) \]
也可以通过给出定义方程来定义：
\begin{align*}
  f(\inl(a)) &\defeq g_0(a)\\
  f(\inr(b)) &\defeq g_1(b).
\end{align*}
要从 $f$ 的前一种表达式得到后一种，只需使用递归子的计算规则。
反之，给定任意定义方程
\begin{align*}
  f(\inl(a)) &\defeq \Phi_0\\
  f(\inr(b)) &\defeq \Phi_1
\end{align*}
其中 $\Phi_0$ 和 $\Phi_1$ 是可能分别涉及变量
\index{variable}%
$a$ 与 $b$ 的表达式，我们也可以借助 $\lambda$-抽象\index{lambda abstraction@$\lambda$-abstraction}，用递归子等价地表达这些方程：
\[ f\defeq \rec{A+B}(C, \lam{a} \Phi_0, \lam{b} \Phi_1).\]
然而，对于自然数，像 $\dbl$ 这样的函数的``定义方程''：
\begin{align}
  \dbl(0) &\defeq 0 \label{eq:dbl0}\\
  \dbl(\suc(n)) &\defeq \suc(\suc(\dbl(n)))\label{eq:dblsuc}
\end{align}
其右边涉及了\emph{函数 $\dbl$ 自身}。
尽管如此，我们仍然希望以这些方程、而非 \eqref{eq:dbl-as-rec} 来作为 \dbl 的定义，因为它们要方便和易读得多。
解决方案是：将 \eqref{eq:dblsuc} 右边出现的表达式``$\dbl(n)$''理解为递归调用的结果，而在 $\dbl\defeq \rec{\nat}(\nat,c_0,c_s)$ 这种形式的定义中，递归调用的结果正是 $c_s$ 的第二个参数。

更一般地，如果我们有一个函数 $f:\nat\to C$ 的如下``定义''：
\begin{align*}
  f(0) &\defeq \Phi_0\\
  f(\suc(n)) &\defeq \Phi_s
\end{align*}
其中 $\Phi_0$ 是类型为 $C$ 的表达式，而 $\Phi_s$ 也是类型为 $C$ 的表达式，且可能涉及变量 $n$ 以及符号``$f(n)$''；那么我们可以将其转化为如下定义：
\[ f \defeq \rec{\nat}(C,\,\Phi_0,\,\lam{n}{r} \Phi_s') \]
其中 $\Phi_s'$ 是将 $\Phi_s$ 中所有出现的``$f(n)$''替换为新变量 $r$ 后所得的结果。

这种通过递归（或更一般地，通过归纳定义依值函数）来定义函数的方式非常方便，因此我们经常采用。
这被称为通过\define{模式匹配（pattern matching）}来定义。
当然，这与程序员定义递归函数的方式非常相似——其函数体中直接包含对自身的递归调用。
然而，与程序员不同，我们在能够进行何种递归调用方面受到了限制：为了使这样的定义能够用递归原理重新表达，被定义的函数 $f$ 在 $f(\suc(n))$ 的函数体中只能作为复合符号``$f(n)$''的一部分出现。
否则，我们可能会写出无意义的函数，例如
\begin{align*}
  f(0)&\defeq 0\\
  f(\suc(n)) &\defeq f(\suc(\suc(n))).
\end{align*}。
如果程序员写出这样的函数，对于任何正数输入，它都会永远调用自身，陷入无限循环而永不返回值。
然而在数学中，一个\emph{函数}要配得上这个名称，就必须为每个输入值关联唯一的输出值，因此上述做法是不可接受的。

当我们今后在 \cref{cha:induction,cha:hits,cha:real-numbers} 引入更复杂的归纳类型时，这一点将更加重要。
每当引入一种新的归纳定义时，我们总是首先推导其归纳原理。
然后，我们才引入一种恰当的``模式匹配''，并说明它可以作为归纳原理的简写形式。

\index{pattern matching|)}%
\index{definition!by pattern matching|)}%

\section{命题即类型}
\label{sec:pat}

\index{proposition!as types|(defstyle}%
\index{logic!propositions as types|(}%
正如引言中所说，在类型论中，要证明一个命题为真，就对应于展示与该命题对应的类型中的一个元素。
\index{evidence, of the truth of a proposition}%
\index{witness!to the truth of a proposition}%
\index{proof|(}
我们将该类型的元素视为该命题为真的\emph{证据}或\emph{见证}。
（有时它们甚至被称为\emph{证明}，但这种叫法容易引起误解，因此我们通常避免使用它。）

然而，我们通常不会显式地构造见证；
相反，我们以常规的数学叙述方式来呈现证明，使得它们可以被转化为类型的元素。
这与经典集合论中的推理并无不同：在那里，我们也不会期望看到使用谓词逻辑规则和集合论公理进行显式推导的过程。

不过，类型论关于证明的视角确实在一些重要方面有所不同。
类型论逻辑的基本原理是：命题不仅仅为真或假，而应被视为其所有可能见证的汇集。
依照这一观念，证明就不只是用来交流数学的工具，其本身也是数学对象，与数、映射、群等更为熟悉的对象具有同等地位。
因此，既然类型对可用的数学对象进行分类并规范它们之间的交互方式，那么命题无非就是特殊的类型——即其元素为证明的那些类型。

\index{propositional!logic}%
\index{logic!propositional}%
使这一对应得以成立的基本观察是：在命题上的\emph{逻辑}运算（用英语表述）与其对应见证类型上的\emph{类型论}运算之间，存在如下自然对应。
\index{false}%
\index{true}%
\index{conjunction}%
\index{disjunction}%
\index{implication}%

\begin{center}
\medskip
\begin{tabular}{ll}
  \toprule
  英语 & 类型论\\
  \midrule
  真 & $\unit$ \\
  假 & $\emptyt$ \\
  $A$ 且 $B$ & $A \times B$ \\
  $A$ 或 $B$ & $A + B$ \\
  若 $A$ 则 $B$ & $A \to B$ \\
  $A$ 当且仅当 $B$ & $(A \to B) \times (B \to A)$ \\
  非 $A$ &  $A \to \emptyt$ \\
  \bottomrule
\end{tabular}
\medskip
\end{center}

这一对应的要点在于：在每一种情形下，右侧类型的元素构造与使用规则，都对应着左侧命题的推理规则。
例如，证明形如``$A$ 且 $B$''的命题的基本方式是分别证明 $A$ 和 $B$；
而构造 $A\times B$ 的一个元素的基本方式则是给出一个对子 $(a,b)$，其中 $a$ 是 $A$ 的一个元素（或见证），$b$ 是 $B$ 的一个元素（或见证）。
并且，如果我们想利用``$A$ 且 $B$''来证明别的结论，我们可以自由地同时使用 $A$ 和 $B$；
这类似于 $A\times B$ 的归纳原理允许我们利用 $A$ 和 $B$ 的元素来构造从该类型出发的函数。

类似地，证明蕴含\index{implication}``若 $A$ 则 $B$''的基本方式是假设 $A$ 并证明 $B$；
而构造 $A\to B$ 的一个元素的基本方式，则是给出一个表示 $B$ 的元素（见证）的表达式，该表达式可以涉及类型 $A$ 的一个未指定的变量元素（见证）。
使用蕴含``若 $A$ 则 $B$''的基本方式则是：若已知 $A$，便可推出 $B$；
这类似于我们可以将函数 $f:A\to B$ 应用于 $A$ 的一个元素以产生 $B$ 的一个元素。

我们强烈建议读者自行练习，验证其他类型形成子的规则也能合理地翻译为逻辑规则。

特别值得注意的是，空类型 $\emptyt$ 对应着假\index{false}。
从逻辑上讲，我们将 $\emptyt$ 的一个实例称为\define{矛盾（contradiction）}：
\indexdef{contradiction}%
因此没有办法证明一个矛盾，%
\footnote{更准确地说，没有\emph{基本的}方法来证明一个矛盾，即 \emptyt 没有构造子。
如果我们的类型论不一致，那么将存在某种更复杂的方法来构造 \emptyt 的一个元素。}
而从矛盾出发可以推出任何东西。
我们还将类型 $A$ 的\define{否定（negation）}
\indexdef{negation}%
定义为
%
\begin{equation*}
  \neg A \ \defeq\ A \to \emptyt.
\end{equation*}
%
这样一来，$\neg A$ 的一个见证就是一个函数 $A \to \emptyt$；我们可以通过假设 $x : A$ 并推导出~$\emptyt$ 的一个元素来构造它。
\index{proof!by contradiction}%
\index{logic!constructive vs classical}
注意，尽管我们得到的是``构造性''逻辑（如引言中所讨论的），但这里所说的这种``通过反证法证明 $\neg A$''（即假设 $A$ 并推出矛盾，从而得出结论 $\neg A$）在构造性意义下是完全有效的：它仅仅是援引了``否定''的\emph{含义}。
而不被允许的那种``通过反证法证明 $A$''，是指假设 $\neg A$ 并推出矛盾，以此来证明 $A$。
构造性地，这样的论证只能让我们得出 $\neg\neg A$。读者可以验证，从 $\neg\neg A$（也就是从 $(A\to \emptyt)\to\emptyt$）并没有显而易见的方法能得到 $A$。

\mentalpause

上述将逻辑连接词翻译为类型形成操作的方式，被称为\define{命题即类型（propositions as types）}：它为我们提供了一种方法，将用语言陈述的命题及其证明，翻译为类型及其元素。
例如，假设我们想证明下面的重言式（``德摩根律''之一）：
\index{law!de Morgan's|(}%
\index{de Morgan's laws|(}%
\begin{equation}\label{eq:tautology1}
  \text{\emph{``If not $A$ and not $B$, then not ($A$ or $B$)''}.}
\end{equation}
对此，一段通常的文字证明可能如下所述。

\begin{quote}
  假设非 $A$ 且非 $B$，同时还假设 $A$ 或 $B$；我们将导出一个矛盾。
  有两种情形。
  若 $A$ 成立，则由非 $A$，我们得到矛盾。
  类似地，若 $B$ 成立，则由非 $B$，我们也得到矛盾。
  因此，在任何一种情形下我们都有矛盾，所以非（$A$ 或 $B$）。
\end{quote}

现在，根据上面给出的规则，对应于我们的重言式~\eqref{eq:tautology1} 的类型是
\begin{equation}\label{eq:tautology2}
  (A\to \emptyt) \times (B\to\emptyt) \to (A+B\to\emptyt)
\end{equation}
因此，我们应该能够将上述证明翻译为该类型的一个元素。

作为这种翻译如何运作的例子，让我们描述一下，一位阅读上述文字证明的数学家是如何在头脑中同时构造出~\eqref{eq:tautology2} 的一个元素的。
起始句``假设非 $A$ 且非 $B$''翻译为定义一个函数，其中隐式应用了其定义域中 Cartesian 积 $(A\to\emptyt)\times (B\to\emptyt)$ 的递归原理。
这引入了类型分别为 $A\to\emptyt$ 和 $B\to\emptyt$ 的未命名变量
\index{variable}%
（假设）
\index{hypothesis}。%
翻译到类型论时，我们必须给这些变量起名字；不妨称它们为 $x$ 和 $y$。
此时，\eqref{eq:tautology2} 的一个元素的部分定义可以写为
\[ f((x,y)) \defeq\; \Box\;:A+B\to\emptyt \]
其中类型为 $A+B\to\emptyt$ 的``洞'' $\Box$ 表示还有待完成的部分。
（我们也可以等价地写成 $f \defeq \rec{(A\to\emptyt)\times (B\to\emptyt)}(A+B\to\emptyt,\lam{x}{y} \Box)$，即使用递归子而非模式匹配。）
下一句``同时还假设 $A$ 或 $B$；我们将导出一个矛盾''表明要用一个函数定义来填充这个洞，并引入另一个未命名的假设 $z:A+B$，从而得到证明状态：
\[ f((x,y))(z) \defeq \;\Box\; :\emptyt. \]
现在，``有两种情形''表明要进行分情形讨论，即应用余积 $A+B$ 的递归原理。
如果用递归子来写，就是
\[ f((x,y))(z) \defeq \rec{A+B}(\emptyt,\lam{a} \Box,\lam{b}\Box,z) \]
而如果用模式匹配来写，则是
\begin{align*}
  f((x,y))(\inl(a)) &\defeq \;\Box\;:\emptyt\\
  f((x,y))(\inr(b)) &\defeq \;\Box\;:\emptyt.
\end{align*}
注意，在这两种写法中，我们现在都有两个类型为 $\emptyt$ 的``洞''需要填充，这对应于我们必须导出矛盾的两种情形。
最后，由 $a:A$ 和 $x:A\to\emptyt$ 得出矛盾，这仅仅是将函数 $x$ 应用于 $a$，另一种情形也类似。
\index{application!of hypothesis or theorem}%
（注意``应用函数''与``应用假设或定理''在措辞上的巧妙重合。）
因此，我们最终的定义是
\begin{align*}
  f((x,y))(\inl(a)) &\defeq x(a)\\
  f((x,y))(\inr(b)) &\defeq y(b).
\end{align*}

作为练习，请运用刚刚介绍的规则，对任意类型 $A$ 和 $B$ 给出 \[ ((A + B) \to \emptyt) \to (A \to \emptyt) \times (B \to \emptyt), \] 的一个元素，以此验证逆重言式\emph{``若非（$A$ 或 $B$），则（非 $A$）且（非 $B$）''}。

\index{logic!classical vs constructive|(}
然而，并非所有经典\index{mathematics!classical}逻辑的重言式都能在这种解释下成立。
例如，规则\emph{``若非（$A$ 且 $B$），则（非 $A$）或（非 $B$）''}并不成立：我们通常无法构造出对应类型
\[ ((A \times B) \to \emptyt) \to (A \to \emptyt) + (B \to \emptyt).\]
中的一个实例。
这反映出类型论中``命题即类型''这套``自然的''逻辑是\emph{构造性的}。
这意味着它不包含某些经典原则，例如排中律 (\LEM{})\index{excluded middle}或反证法，\index{proof!by contradiction}以及其他依赖于这些原则的定律，比如上面这个德摩根定律的实例。
\index{law!de Morgan's|)}%
\index{de Morgan's laws|)}%

从哲学上讲，构造逻辑之所以得名，是因为它将自身限制于那些可以\emph{能行地}完成的构造，也就是说，这些构造具有计算意义。
在不求过分精确的前提下，这意味着存在某种算法\index{algorithm}，它能一步步指明如何构造一个对象（作为一种特例，如何看出一个定理为真）。
这就必须把排中律排除在外，因为并不存在一个\emph{能行}的\index{effective!procedure}程序来判定一个命题是真还是假。

类型论逻辑的构造性意味着它具有内在的计算意义，这对计算机科学家也很有意义。
同时，这也意味着类型论提供了\emph{公理选择的自由（axiomatic freedom）}。\index{axiomatic freedom}
例如，尽管默认情况下不存在见证排中律的构造，但逻辑本身与排中律的存在是兼容的（参见\cref{sec:intuitionism}）。
因此，由于类型论并不\emph{否认}排中律，我们可以一致地将其作为假设加入，从而如通常那样不受限制地做数学。
从这个意义上说，类型论丰富而非限制了传统数学实践。

我们建议不熟悉构造逻辑的读者多做一些练习来熟悉它。相关建议参见\cref{ex:tautologies,ex:not-not-lem}。
\index{logic!classical vs constructive|)}

\mentalpause

到此为止，我们讨论的还只是命题逻辑。
\index{quantifier}%
\index{quantifier!existential}%
\index{quantifier!universal}%
\index{predicate!logic}%
\index{logic!predicate}%
现在我们来考虑\emph{谓词（predicate）}逻辑，其中除了像``且''和``或''这样的逻辑连接词，我们还有量词``存在''和``对所有''。
在这种情况下，类型扮演着双重角色：它们既作为命题，也作为传统意义上的类型，即我们进行量化所基于的论域。
一个关于类型 $A$ 的谓词表示为一个类型族 $P : A \to \UU$，它给每个元素 $a : A$ 指派一个类型 $P(a)$，这个类型对应于``$P$ 对 $a$ 成立''这个命题。现在我们通过解释量词来扩展上述翻译：

\begin{center}
  \medskip
  \begin{tabular}{ll}
    \toprule
    日常语言 & 类型论 \\
    \midrule
    对所有 $x:A$，$P(x)$ 成立 & $\prd{x:A} P(x)$ \\
    存在 $x:A$ 使得 $P(x)$ & $\sm{x:A}$ $P(x)$ \\
    \bottomrule
  \end{tabular}
  \medskip
\end{center}

和前面一样，我们可以证明（构造性）谓词逻辑的重言式会翻译成有实例的类型。
例如，\emph{若对所有 $x:A$，有 $P(x)$ 且 $Q(x)$，则（对所有 $x:A$，有 $P(x)$）且（对所有 $x:A$，有 $Q(x)$）} 翻译为
\[ (\tprd{x:A} P(x) \times Q(x)) \to (\tprd{x:A} P(x)) \times (\tprd{x:A} Q(x)). \]
这个重言式的一个非形式证明可以这样进行：

\begin{quote}
  假设对所有 $x$，有 $P(x)$ 且 $Q(x)$。
  首先，我们假设给定 $x$ 并证明 $P(x)$。
  根据假设，我们有 $P(x)$ 且 $Q(x)$，因此我们有 $P(x)$。
  其次，我们假设给定 $x$ 并证明 $Q(x)$。
  再次根据假设，我们有 $P(x)$ 且 $Q(x)$，因此我们有 $Q(x)$。
\end{quote}

第一句话通过为假设引入一个见证，开始将蕴含式定义为函数：\index{hypothesis}
\[ f(p) \defeq \;\Box\; : (\tprd{x:A} P(x)) \times (\tprd{x:A} Q(x)). \]
此时隐式地使用了配对构造子来生成积类型的元素，在这个例子中，``第一''和``第二''两个词对此略有提示：
\[ f(p) \defeq \Big( \;\Box\; : \tprd{x:A} P(x) \;,\; \Box\; : \tprd{x:A}Q(x) \;\Big). \]
``我们假设给定 $x$ 并证明 $P(x)$''这句话，现在表示以通常的方式定义一个\emph{依值}函数，为其输入引入一个变量 \index{variable}。%
由于这是在配对构造子内部，将其写成一个 $\lambda$-抽象就很自然了：\index{lambda abstraction@$\lambda$-abstraction}
\[ f(p) \defeq \Big( \; \lam{x} \;\big(\Box\; : P(x)\big) \;,\; \Box\; : \tprd{x:A}Q(x) \;\Big). \]
接着，``我们有 $P(x)$ 和 $Q(x)$''用到了假设，从而得到 $p(x) : P(x)\times Q(x)$；而``因此我们有 $P(x)$''则隐式地应用了恰当的投影：
\[ f(p) \defeq \Big( \; \lam{x} \proj1(p(x))  \;,\; \Box\; : \tprd{x:A}Q(x) \;\Big). \]
后面两句话以显而易见的方式填上了另一个空位：
\[ f(p) \defeq \Big( \; \lam{x} \proj1(p(x))  \;,\; \lam{x} \proj2(p(x)) \; \Big). \]
当然，我们用作例子的这些英文证明，远比数学家在实践中通常使用的证明冗长；它们更像是``证明导论''课上会用到的那种语言。
实践中的数学家已经学会了填补这些空白，因此在实践中我们可以省略大量细节，而且我们通常也会这样做。
然而，证明有效性的判据始终是：它们能够还原为对相应类型中某个元素的构造。

\symlabel{leq-nat}%
作为一个更具体的例子，考虑如何定义自然数的不等关系。
一个自然的定义是：$n\le m$ 当且仅当存在一个 $k:\nat$ 使得 $n+k=m$。
（这里再次用到了我们将在下一节介绍的相等类型，不过我们暂时不需要了解太多。）
根据命题即类型的对应关系，这会得到：
\[ (n\le m) \defeq \sm{k:\nat} (\id{n+k}{m}). \]
我们邀请读者根据这个定义来证明 $\le$ 的常见性质。
对于严格不等式，有几个自然的选项，例如
\[ (n<m) \defeq \sm{k:\nat} (\id{n+\suc(k)}{m}) \]
或者
\[ (n<m) \defeq (n\le m) \times \neg(\id{n}{m}). \]
前者在构造性数学中更为自然，但在本例中，由于 $\nat$ 具有``可判定相等性''（参见 \cref{sec:intuitionism,prop:nat-is-set}），它实际上等价于后者。
\index{decidable!equality}%

对于类型 $\sm{x:A} P(x)$，还有另一种解读方式。
既然它的一个实例是一个元素 $x:A$ 连同 $P(x)$ 成立的一个见证，我们就可以不把 $\sm{x:A} P(x)$ 看作命题``存在 $x:A$ 使得 $P(x)$''，而是将其看作``所有满足 $P(x)$ 的元素 $x:A$ 构成的类型''，即 $A$ 的一个``子类型''。
\index{subtype}%

我们将在 \cref{subsec:prop-subsets} 回到这种解读。
现在，我们注意到它允许我们在将类型定义为数学结构时纳入公理——这一点我们曾在 \cref{sec:sigma-types} 讨论过。
例如，假设我们想定义一个\define{半群}\index{semigroup}为一个配备了一个二元运算 $m:A\to A\to A$（即一个原群\index{magma}）的类型 $A$，并且满足对于所有 $x,y,z:A$ 都有 $m(x,m(y,z)) = m(m(x,y),z)$。
这后一个命题由类型
\[\prd{x,y,z:A} m(x,m(y,z)) = m(m(x,y),z),\]
来表示，所以半群的类型就是
\[ \semigroup \defeq \sm{A:\UU}{m:A\to A\to A} \prd{x,y,z:A} m(x,m(y,z)) = m(m(x,y),z), \]
即由所有半群组成的 $\mathsf{Magma}$ 的子类型。
从一个 $\semigroup$ 的实例中，我们可以通过应用恰当的投影，提取出载体 $A$、运算 $m$ 以及该公理的一个见证。
我们将在 \cref{sec:equality-of-structures} 回到这个例子。

还需注意的是，我们可以利用类型论中的宇宙来表示``高阶逻辑''——也就是说，我们可以对所有的命题或所有的谓词进行量化。
例如，给定 $A : \UU$ 和 $a,b : A$，我们可以将命题\emph{对于所有性质 $P : A \to \UU$，如果 $P(a)$ 那么 $P(b)$} 表示为
\[ \prd{P : A \to \UU} P(a) \to P(b) \]
然而，\emph{先验地}，这个命题位于一个比我们所量化的命题更高、也不同的宇宙中；即
\[ \Parens{\prd{P : A \to \UU_i} P(a) \to P(b)} : \UU_{i+1}. \]
我们将在 \cref{subsec:prop-subsets} 回到这个问题。

\mentalpause

我们在此描述的是一种``证明相关''的
\index{mathematics!proof-relevant}%
命题翻译方式，其中析取与存在性陈述的证明会携带一些信息。
例如，如果我们有 $A+B$ 的一个实例，并把它看作``$A$ 或 $B$''的一个见证，那么我们就能知道它究竟来自 $A$ 还是来自 $B$。
类似地，如果我们有 $\sm{x:A} P(x)$ 的一个实例，并把它看作``存在 $x:A$ 使得 $P(x)$''的一个见证，那么我们就能知道这个元素 $x$ 具体是什么（它就是给定实例的第一投影）。

由于这种逻辑具有证明相关的本性，我们可能会遇到``$A$ 当且仅当 $B$''的情形（回想一下，这意指 $(A\to B)\times (B\to A)$），但类型 $A$ 和 $B$ 却展现出不同的行为。
例如，很容易验证``$\mathbb{N}$ 当且仅当 $\unit$''，然而显然 $\mathbb{N}$ 和 $\unit$ 在重要方面有所不同。
``$\mathbb{N}$ 当且仅当 $\unit$''这个陈述仅仅告诉我们：\emph{当被视为一个命题时}，类型 $\mathbb{N}$ 所代表的命题与 $\unit$ 相同（此时，这个命题是真命题）。
我们有时会用 $A$ 与 $B$ 是\define{逻辑等价}这一说法来表达``$A$ 当且仅当 $B$''。
\indexdef{logical equivalence}%
\indexdef{equivalence!logical}%
这需要与 \cref{sec:basics-equivalences,cha:equivalences} 中将要引入的更强的\emph{类型等价}概念区分开：尽管 $\mathbb{N}$ 和 $\unit$ 是逻辑等价的，它们却不是等价的类型。

在 \cref{cha:logic} 中，我们将引入一类称为``命题''的类型，对于它们而言，等价与逻辑等价两个概念将重合。
利用这些类型，我们将对上述逻辑做出一处修改，这种修改在某些情况下是恰当的：它将抛弃析取与存在性陈述中所包含的额外信息。

最后我们注意到，``命题即类型''的对应关系反过来也成立，我们可以将任意类型 $A$ 本身视为一个命题，而证明它的方式就是展示 $A$ 的一个元素。
有时我们将这个命题表述为``$A$ 是\define{有实例的}''。
\indexdef{inhabited type}%
\indexsee{type!inhabited}{inhabited type}%
也就是说，当我们说 $A$ 是有实例的时，我们的意思是我们已经给出了 $A$ 的一个（具体）元素，只是我们选择不去为这个元素命名。
类似地，说 $A$ \emph{不是有实例的}，等同于给出 $\neg A$ 的一个元素。
特别地，空类型 $\emptyt$ 显然不是有实例的，因为 $\neg \emptyt \jdeq (\emptyt \to \emptyt)$ 有 $\idfunc[\emptyt]$ 作为实例。\footnote{这不应与类型论的一致性混淆，后者是一个\emph{元理论}主张，即在遵循类型论规则的前提下，不可能构造出 $\emptyt$ 的一个元素。\indexfoot{consistency}}

\index{proof|)}%
\index{proposition!as types|)}%
\index{logic!propositions as types|)}%

\section{相等类型}
\label{sec:identity-types}

\index{type!identity|(defstyle}%
\indexsee{identity!type}{type, identity}%
\indexsee{type!equality}{type, identity}%
\indexsee{equality!type}{type, identity}%
虽然之前介绍的那些构造可以看作是对标准集合论构造的推广，但我们处理相等性的方式，似乎为类型论所特有。
根据命题即类型的理念，两个相同类型的元素 $a,b:A$ 相等这一\emph{命题}，必须对应某个\emph{类型}。
由于这个命题依赖于 $a$ 和 $b$ 具体是什么，这些\define{相等类型}或\define{恒等类型}必须是依赖于 $A$ 的两个副本的类型族。

我们可以将这个族写作 $\idtypevar{A}:A\to A\to\type$（注意不要与恒等映射 $\idfunc[A]$ 混淆），于是 $\idtype[A]ab$ 就是代表 $a$ 与 $b$ 相等这一命题的类型。
不过，一旦我们熟悉了命题即类型的观点，使用标准的等号来表示这个类型也很方便；因此，``$\id{a}{b}$''也将是代表 $a$ 等于 $b$ 这一命题的\emph{类型} $\idtype[A]ab$ 的一种记法。
为了清晰起见，我们也可以写作``$\id[A]{a}{b}$''来指明类型 $A$。
如果我们有 $\id[A]{a}{b}$ 的一个元素，我们就可以说 $a$ 和 $b$ 是\define{相等}的，有时为了强调这不同于 \cref{sec:types-vs-sets} 中讨论的判断相等性 $a\jdeq b$，我们也说它们是\define{命题相等}的。
\indexdef{equality!propositional}%
\indexdef{propositional!equality}%

正如我们在 \cref{sec:pat} 中所提到的，``或''和``存在''的``命题即类型''版本可以包含比命题为真这一事实更多的信息；同样地，类型 $\id{a}{b}$ 也完全可以包含更多信息。
事实上，这正是同伦解释的基石：我们将 $\id{a}{b}$ 的见证视为空间 $A$ 中 $a$ 与 $b$ 之间的\emph{道路}\indexdef{path}或\emph{等价}。
正如空间中两点之间可以有多条道路，两个对象相等的见证也可以不止一个。
换一种说法，我们可以把 $\id{a}{b}$ 看作 $a$ 与 $b$ 的\emph{相等证明}\indexdef{identification}的类型，而将 $a$ 和 $b$ 确认为相等的方式也可以有多种。
我们将在 \cref{cha:basics} 回到这一解释；现在我们先聚焦于相等类型的基本规则。
和本章讨论的所有其他类型一样，它将有形成规则、引入规则、消去规则和计算规则，这些规则在形式上的行为完全相同。

形成规则是说：给定一个类型 $A:\UU$ 和两个元素 $a,b:A$，我们可以在同一个宇宙中形成类型 $(\id[A]{a}{b}):\UU$。
构造 $\id{a}{b}$ 的元素的基本方法是知道 $a$ 和 $b$ 相同。
因此，引入规则是一个依值函数
\[\refl{} : \prd{a:A} (\id[A]{a}{a})\]
称为\define{自反性}（reflexivity），
\indexdef{reflexivity!of equality}%
它说 $A$ 的每个元素都以指定的方式等于自身。
我们把 $\refl{a}$ 看作点 $a$ 处的
常值道路\indexdef{path!constant}\indexsee{loop!constant}{path, constant}。

特别地，这意味着如果 $a$ 和 $b$ 是\emph{判断相等的}，即 $a\jdeq b$，那么我们也得到一个元素 $\refl{a} : \id[A]{a}{b}$。
这是良类型的，因为 $a\jdeq b$ 意味着类型 $\id[A]{a}{b}$ 与 $\id[A]{a}{a}$ 也是判断相等的，而后者正是 $\refl{a}$ 的类型。

相等类型的归纳原理（即消去规则）是类型论中最微妙的部分之一，对同伦解释至关重要。
我们先考虑它的一个重要推论，即``相等者可以替换相等者''的原则，如下所示：
\index{indiscernibility of identicals}%
\index{equals may be substituted for equals}%

\begin{description}
\item[同一者的不可分辨性（Indiscernibility of identicals）：]
对于每个类型族
\[
C : A \to \UU
\]
存在一个函数
\[
f : \prd{x,y:A}{p:\id[A] x y} C(x) \to C(y)
\]
使得
\[
f(x,x,\refl{x}) \defeq \idfunc[C(x)].
\]
\end{description}

这个原理说的是，每个类型族 $C$ 都尊重相等：将 $C$ 应用到 $A$ 中\emph{相等}的元素上，也会在所得类型之间给出一个函数。所列等式表明，与自反性相关联的函数是恒等函数（我们之后将会看到，一般地，函数 $f(x,y,p): C(x) \to C(y)$ 总是类型之间的等价）。

同一者的不可分辨性可以被视为相等类型的一种递归原理，类似于前面为布尔值和自然数给出的递归原理。
正如 $\rec{\nat}$ 为任何其他某种类型的 $C$ 给出一个指定的映射 $\nat\to C$ 那样，同一者的不可分辨性也给出一个从 $\id[A] x y$ 到 $A$ 上某些自反二元关系的指定映射，即由某个一元谓词 $C(x)$ 给出的形如 $C(x) \to C(y)$ 的关系。
我们也可以针对更一般形式 $C(x, y)$ 的自反关系，表述一个更一般的递归原理。
然而，为了完全刻画相等类型，我们必须将这个递归原理推广为一个归纳原理；它不仅考虑从 $\id[A] x y$ 出发的映射，还要考虑其上的类型族。
换句话说，我们不仅考虑允许用相等者替换相等者，还要考虑等式的证据 $p$ 本身。
    
\subsection{道路归纳}

\index{generation!of a type, inductive|(}
相等类型的归纳原理被称为\define{道路归纳}（path induction），
\index{path!induction|(}%
\index{induction principle!for identity type|(}%
这是考虑到将在 \cref{cha:basics} 的引言中解释的同伦解释。它可以理解为：相等类型族是由形如 $\refl{x}: \id{x}{x}$ 的元素自由生成的。

\begin{description}
\item[道路归纳（Path induction）：] 
  给定一个类型族
  \[ C : \prd{x,y:A} (\id[A]{x}{y}) \to \UU \]
  以及一个函数
  \[ c :  \prd{x:A} C(x,x,\refl{x}),\]
  则存在一个函数
  \[ f : \prd{x,y:A}{p:\id[A]{x}{y}} C(x,y,p) \]
  使得
  \[ f(x,x,\refl{x}) \defeq c(x). \]
\end{description}

需要注意的是，就像积、余积、自然数等的归纳原理一样，道路归纳允许我们定义\emph{特定的}函数，这些函数展现出适当的计算行为。
正如我们有\emph{那个}由 $c_0:C$ 和 $c_s:\nat \to C \to C$ 经递归定义的函数 $f:\nat\to C$，而且它满足 $f(0)\jdeq c_0$ 和 $f(\suc(n))\jdeq c_s(n,f(n))$，我们也有\emph{那个}由 $c :  \prd{x:A} C(x,x,\refl{x})$ 经道路归纳定义的函数 $f : \dprd{x,y:A}{p:\id[A]{x}{y}} C(x,y,p)$，而且它满足 $f(x,x,\refl{x}) \jdeq c(x)$。

为了理解这个原理的含义，我们先考虑 $C$ 不依赖于 $p$ 的更简单情形。此时我们有 $C:A\to A\to \UU$，可以将其视为一个依赖于 $A$ 中两个元素的谓词。我们关心的是命题 $C(x,y)$ 对某个元素对 $x,y:A$ 何时成立。在这种情况下，道路归纳的假设是说，我们知道对所有 $x:A$ 都有 $C(x,x)$ 成立，也就是说，如果我们在 $(x, x)$ 这对元素处计算 $C$，会得到一个真命题——所以 $C$ 是一个\emph{自反关系}（reflexive relation）。而结论则告诉我们，只要 $\id{x}{y}$ 成立，$C(x,y)$ 就成立。这正是前面提到的、适用于自反关系的更一般的递归原理。

该规则的一般归纳形式则允许 $C$ 还依赖于见证 $x$ 与 $y$ 相等的 $p:\id{x}{y}$。在前提中，我们不仅将 $x, y$ 替换为 $x, x$，同时也将 $p$ 替换为自反性：要证明关于所有元素 $x, y$ 以及它们之间道路 $p : \id{x}{y}$ 的性质，只需考虑元素为 $x,x$ 且道路为 $\refl{x}: \id{x}{x}$ 的情形即可。如果我们仅仅把类型看作集合，这样做带来的好处不太明显；但由于在 $x$ 和 $y$ 之间可能存在许多不同的相等证明 $p : \id{x}{y}$，因此在考虑类型 $\id[A]{x}{y}$ 上的族时，对它们加以记录是有意义的。在 \cref{cha:basics} 中我们将会看到，这对于同伦解释至关重要。

若将道路归纳打包为一个函数，其形式如下：
\symlabel{defn:induction-ML-id}%
\begin{narrowmultline*}
  \indid{A} :  \dprd{C : \prd{x,y:A} (\id[A]{x}{y}) \to \UU}
  \Parens{\tprd{x:A} C(x,x,\refl{x})} \to 
  \narrowbreak
  \dprd{x,y:A}{p:\id[A]{x}{y}}   C(x,y,p)
\end{narrowmultline*}
并满足等式\index{computation rule!for identity types}
\[ \indid{A}(C,c,x,x,\refl{x}) \defeq c(x). \]
函数 $ \indid{A}$ 传统上称为 $J$。
\indexsee{J@$J$}{induction principle for identity type}我们将在 \cref{lem:transport} 中说明，同一者的不可分辨性是道路归纳的一个实例，并赋予它一个新的名称和记号。

\mentalpause

给定一个证明 $p : \id{a}{b}$，道路归纳要求我们将 $a$ 和 $b$ \emph{两者}都替换为同一个未知元素 $x$；因此，要为 $A$ 中所有相等的元素对定义族 $C$ 的元素，只需在``对角线''上定义它就足够了。然而，在一些证明中，利用等式 $p : \id{a}{b}$ 将所有出现的 $b$ 替换为 $a$（或者反之）会更简单，因为有时针对等式中提到的特定元素 $a$ 继续完成证明，要比为一个一般的未知元素 $x$ 完成证明更容易。这便引出了相等类型的第二种归纳原理，它指出类型族 $\id[A]{a}{x}$ 是由元素 $\refl{a} : \id{a}{a}$ 生成的。如下所示，这第二种原理与第一种是等价的；它只是有时提供了一个更方便的表述。

\index{path!induction based}%
\index{induction principle!for identity type!based}%

\begin{description}
\item[基于基点的道路归纳（Based path induction）：] 
  固定一个元素 $a:A$，并假设给定一个族
  \[ C : \prd{x:A} (\id[A]{a}{x}) \to \UU \]
  以及一个元素
  \[ c : C(a,\refl{a}). \]
  那么我们可以得到一个函数
  \[ f : \prd{x:A}{p:\id{a}{x}} C(x,p) \]
  使得
  \[ f(a,\refl{a}) \defeq c.\]
\end{description}

这里 $C(x,p)$ 是一个类型族，其中 $x$ 是 $A$ 的元素，而对于 $A$ 中固定的 $a$，$p$ 是相等类型 $\id[A]{a}{x}$ 的一个元素。基于基点的道路归纳原理是说，要为所有 $x$ 和 $p$ 定义这个族的一个元素，只需考虑 $x$ 是 $a$ 且 $p$ 是 $\refl{a} : \id{a}{a}$ 的情形。

打包为函数后，基于基点的道路归纳变为：
\symlabel{defn:induction-PM-id}%
\begin{align*}
  \indidb{A} :  \dprd{a:A}{C : \prd{x:A} (\id[A]{a}{x}) \to \UU}
  C(a,\refl{a}) \to \dprd{x:A}{p : \id[A]{a}{x}} C(x,p) 
\end{align*}
并满足等式
\[ \indidb{A}(a,C,c,a,\refl{a}) \defeq c. \]
%\[ g(x)(x,\refl{x}) \defeq d(x) \]

下面我们将证明，道路归纳与基于基点的道路归纳是等价的。正因如此，我们有时会不严格地将基于基点的道路归纳也简称为``道路归纳''，而依靠读者从证明的形式来推断所指的原理。

\begin{rmk}\label{rmk:the-only-path-is-refl}
  直观上，自然数的归纳原理表达的是这样一个事实：每个自然数要么是 $0$，要么具有 $\suc(n)$ 的形式（其中 $n$ 是某个自然数）。
  因此，如果我们能为这些情形（第二种情形带有归纳假设）证明某个性质，那么我们就为所有自然数证明了该性质。
  类似地，$A+B$ 的归纳原理表达的是：$A+B$ 中的每个元素要么形如 $\inl(a)$，要么形如 $\inr(b)$，依此类推。
  用同样的解读来理解道路归纳，我们可能会说，道路归纳表达的是这样一个事实：每条道路都具有 $\refl{a}$ 的形式，因此如果我们能为自反性道路证明某个性质，那么我们就为所有道路都证明了该性质。

  然而，这种理解放在道路的同伦解释语境下会相当令人困惑，因为两个元素 $a$ 和 $b$ 之间可能有许多种不同的识别方式，从而相等类型中可能有许多不同的元素！
  既然可能存在许多不同的道路，我们又怎么可能拥有一个断言``只有自反性才是道路''的归纳原理呢？

关键之处在于，归纳定义的不是相等\emph{类型}，而是相等\emph{族}。
具体来说，道路归纳说的是：当 $x, y$ 遍历 $A$ 的所有元素时，类型族 $(\id[A]{x}{y})$ 是由形如 $\refl{x}$ 的元素归纳定义的。
这意味着，对于任何依赖于相等族的一个\emph{一般}元素 $(x, y, p)$ 的族 $C(x, y, p)$，要给出它的一个元素，只需考虑形如 $(x, x, \refl{x})$ 的情形。
在同伦解释下，这表示由所有三元组 $(x, y, p)$ 构成的空间（其中 $x$ 和 $y$ 是道路 $p$ 的端点，即 $\Sigma$-类型 $\sm{x,y:A}(\id{x}{y})$）是由每个点 $x$ 处的常值环路\index{path!constant}归纳生成的。
我们将在 \cref{cha:basics} 中看到，在同伦论里，对应于 $\sm{x,y:A}(\id{x}{y})$ 的空间就是\emph{自由道路空间（free path space）}——即 $A$ 中端点可以自由变动的所有道路构成的空间——并且事实上，该空间中的任意一点都同伦于某个点处的常值环路，因为我们可以直接沿给定道路将其一个端点收回。
类似地，相应的事实也在类型论中成立：我们可以通过对 $p:x=y$ 作道路归纳来证明 $\id[\sm{x,y:A}(\id{x}{y})]{(x,y,p)}{(x,x,\refl{x})}$。

类似地，定基点道路归纳说的是，对于一个固定的 $a:A$，当 $y$ 遍历 $A$ 的所有元素时，类型族 $(\id[A]{a}{y})$ 是由元素 $\refl{a}$ 归纳定义的。
因此，对于任何依赖于该族的一个一般元素 $(y, p)$ 的族 $C(y, p)$，要给出它的一个元素，只需考虑情形 $(a, \refl{a})$。
从同伦角度看，这表达了以某个选定点为起点的所有道路构成的空间（即该点处的\emph{定基点道路空间（based path space）}，在类型论中即 $\sm{y:A} (\id{a}{y})$）可收缩为该选定点上的常值环路。
同样地，对应的事实也在类型论中成立：我们可以通过对 $p:a=y$ 作定基点道路归纳来证明 $\id[\sm{y:A}(\id{a}{y})]{(y,p)}{(a,\refl{a})}$。
另请注意，根据 \cref{sec:pat} 中所述将 $\Sigma$-类型视为子类型的解释，类型 $\sm{y:A}(\id{a}{y})$ 可以被看作``所有与 $a$ 相等的 $A$ 中元素构成的类型''，这是``单点\index{type!singleton}子集'' $\{a\}$ 的类型论版本。

无论是道路归纳还是定基点道路归纳，都没有提供一种方法来给出依赖于一条具有\emph{两个固定端点} $a$ 和 $b$ 的道路 $p$ 的族 $C(p)$ 的元素。
特别地，对于依赖于一条环路的族 $C: (\id[A]{a}{a}) \to \UU$，我们\emph{不能}应用道路归纳而只考虑 $C(\refl{a})$ 的情形，因此，我们不能证明所有环路都是自反性。
所以，归纳定义相等族，并不禁止在相等类型的特定实例中出现非自反性的道路。
换言之，一条道路 $p:\id{x}{x}$ 作为 $(\id{x}{x})$ 中的一个元素，可能不等于自反性，但序对 $(x,p)$ 作为 $\sm{y:A}(\id{x}{y})$ 中的元素，却仍然等于序对 $(x,\refl{x})$。

举一个拓扑学的例子：考虑有孔圆盘 \narrowequation{\setof{ (x,y) | 0 < x^2+y^2 < 2 }} 中的一条环路，它从 $(1,0)$ 出发，绕过位于 $(0,0)$ 的孔洞一圈后，回到 $(1,0)$。
如果我们把两个端点都固定在 $(1,0)$，那么这条环路无法在有孔圆盘内形变为常值道路，正如一根绕在杆子上的绳子，如果我们紧握两端，就无法把它拉回来一样。
然而，如果我们允许其中一个端点变动，这条环路就可以收缩为常值道路，正如我们如果只握住绳子的一端，就总能把它收回来一样。
\end{rmk}

\index{path!induction|)}%
\index{induction principle!for identity type|)}%
\index{generation!of a type, inductive|)}

\subsection{道路归纳与定基点道路归纳的等价性}

上面介绍的关于相等类型的两种归纳原理是等价的。
我们很容易看到，道路归纳可以从定基点道路归纳原理推导出来。
具体地，假设我们有道路归纳的前提：
\begin{align*}
C &: \prd{x,y:A}(\id[A]{x}{y}) \to \UU,\\
c &: \prd{x:A} C(x,x,\refl{x}).
\end{align*}
现在，给定一个元素 $x:A$，我们可以对以上两者进行实例化，得到
\begin{align*}
C' &: \prd{y:A} (\id[A]{x}{y}) \to \UU,  \\
C' &\defeq C(x), \\
c' &: C'(x,\refl{x}), \\
c' &\defeq c(x).
\end{align*}
显然，$C'$ 和 $c'$ 符合定基点道路归纳的前提，因此我们可以构造
\begin{equation*}
  g : \prd{y:A}{p : \id{x}{y}} C'(y,p)
\end{equation*}
其定义等式为
\[ g(x,\refl{x}) \defeq c'.\]
现在，我们注意到 $g$ 的余域就是 $C(x,y,p)$。
于是，释放我们对 $x:A$ 的假设，便可以推导出一个函数
\[ f : \prd{x,y:A}{p : \id[A]{x}{y}} C(x,y,p) \]
并满足所需的判断相等性 $f(x,x,\refl{x}) \judgeq g(x,\refl{x}) \defeq c' \defeq c(x)$。

这一事实的另一种证法是观察到：任何这样的 $f$ 都可以作为 $\indid{A}$ 的一个实例得到，所以只需用 $\indidb{A}$ 来定义 $\indid{A}$ 即可：
\[ \indid{A}(C,c,x,y,p) \defeq \indidb{A}(x,C(x),c(x),y,p). \]

另一个方向的证明则要棘手一些；我们并不清楚如何用道路归纳的某个特定实例来推导出定基点道路归纳的某个特定实例。
我们能做的是转而构造道路归纳的一个实例，让它一次性涵盖定基点道路归纳所有可能的实例。
定义
\begin{align*}
D &: \prd{x,y:A} (\id[A]{x}{y}) \to \UU, \\
D(x,y,p) &\defeq \prd{C : \prd{z:A} (\id[A]{x}{z}) \to \UU} C(x,\refl{x}) \to C(y,p).
\end{align*}
那么我们可以构造函数
\begin{align*}
d &: \prd{x : A} D(x,x,\refl{x}), \\
d &\defeq \lamu{x:A}\lamu{C:\prd{z:A}{p : \id[A]{x}{z}} \UU}\lam{c:C(x,\refl{x})} c
\end{align*}
从而利用道路归纳得到
\[ f : \prd{x,y:A}{p:\id[A]{x}{y}} D(x,y,p) \]
并满足 $f(x,x,\refl{x}) \defeq d(x)$。展开 $D$ 的定义，我们可以展开 $f$ 的类型：
\[ f : \prd{x,y:A}{p:\id[A]{x}{y}}{C : \prd{z:A} (\id[A]{x}{z}) \to \UU} C(x,\refl{x}) \to C(y,p). \]
现在，给定 $a:A$ 以及 $x:A$ 和 $p:\id[A]{a}{x}$，我们就能推导出定基点道路归纳的结论：
\[ f(a,x,p,C,c) : C(x,p). \]
注意，我们同时也得到了正确的定义相等。

另一种证法是观察到：任何定基点道路归纳的使用都是 $\indidb{A}$ 的一个实例，并定义
\begin{narrowmultline*}
\indidb{A}(a,C,c,x,p) \defeq \narrowbreak
\indid{A}
  \begin{aligned}[t]
    \big(
    &\big(\lamu{x,y:A}{p:\id[A]{x}{y}} \tprd{C : \prd{z:A} (\id[A]{x}{z}) \to \UU} C(x,\refl{x}) \to C(y,p) \big),\\
    &(\lamu{x:A}{C:\prd{z:A} (\id[A]{x}{z}) \to \UU}{d:C(x,\refl{x})} d),
     a, x, p\big)(C, c).
   \end{aligned}
\end{narrowmultline*}

需要注意的是，上面给出的构造用到了宇宙。
也就是说，如果想用 $C : \prd{x:A} (\id[A]{a}{x}) \to \UU_i$ 来建模 $\indidb{A}$，我们就需要使用带有
%
\[ D:\prd{x,y:A} (\id[A]{x}{y}) \to \UU_{i+1} \]
%
的 $\indid{A}$，因为 $D$ 对给定类型的所有 $C$ 进行了量化。
虽然这与我们对宇宙的定义是兼容的，但也有可能在不使用宇宙的情况下推导出 $\indidb{A}$：
我们可以证明 $\indid{A}$ 蕴涵 \cref{lem:transport,thm:contr-paths}，并且这两个原理联合起来可直接推出 $\indidb{A}$。
我们将细节留给读者，作为 \cref{ex:pm-to-ml}。

我们可以利用上述任一关于相等类型的表述来建立以下事实：相等是一种等价关系，每个函数都保持相等，并且每个类型族都尊重相等。
我们将这些细节留到下一章，在那里它们会在同伦类型论的语境下得到推导和解释。

\begin{rmk}\label{rmk:propeq-vs-jdeq}
  我们要强调，尽管命题相等有一些不熟悉的特征，但它才是同伦类型论中数学上真正意义上的``相等''。
  这一区分不属于判断相等性，后者更确切地说是类型论规则的一种元理论特征。
  例如，在\cref{sec:inductive-types}中证明的自然数加法结合律，就是一种\emph{命题}相等，而非判断相等性。
  交换律（\cref{ex:add-nat-commutative}）也是如此。
  即便像 $n+1=1+n$ 这样简单的交换性，对于一般的 $n$ 来说也不是判断相等性（不过，对于任何具体的 $n$，例如 $3+1\jdeq 1+3$，它是判断相等的，因为根据定义 $+$ 的计算规则，两者都与 $4$ 判断相等）。
  我们只能借助相等类型来证明这类事实，因为对 \nat 应用归纳原理时只能以类型（而非判断）作为输出。
\end{rmk}

\subsection{不等}
\label{sec:disequality}

最后，让我们也谈一谈\define{不等}，
\indexdef{disequality}%
它定义为相等的否定：%
\footnote{我们用 ``inequality'' 一词指代 $<$ 和 $\leq$。另外请注意，这里否定的是\emph{命题}相等类型。
当然，对判断相等性 $\jdeq$ 取否定是毫无意义的，因为判断本身不参与逻辑运算。}
%
\begin{equation*}
  (x \neq_A y) \ \defeq\ \lnot (\id[A]{x}{y}).
\end{equation*}
若 $x\neq y$，则称 $x$ 和 $y$ 是\define{不相等的}，
\indexdef{unequal}%
或\define{不相等的}。
%
就像否定一样，不等在此处的作用远不如它在经典\index{mathematics!classical}数学里那么重要。
例如，我们无法通过证明两个对象并非不相等来证明它们相等：那样会用到经典的双重否定律，参见\cref{sec:intuitionism}。

有时，用正面的方式来陈述不等是有用的。
例如，在~\cref{RD-inverse-apart-0}中我们将证明，一个实数 $x$ 有逆元当且仅当其与~$0$ 的距离是正的，这是一个比 $x \neq 0$ 更强的要求。

\index{type!identity|)}%

\sectionNotes

这里介绍的类型论是 Martin-L\"{o}f 直觉主义类型论~\cite{Martin-Lof-1972,Martin-Lof-1973,Martin-Lof-1979,martin-lof:bibliopolis} 的一个版本，后者本身又建立于 Brouwer \cite{beeson}、Heyting~\cite{heyting1966intuitionism}、Scott~\cite{scott70}、de Bruijn~\cite{deBruijn-1973}、Howard~\cite{howard:pat}、Tait~\cite{Tait-1966,Tait-1968} 以及 Lawvere~\cite{lawvere:adjinfound}\index{Lawvere} 的奠基性工作之上，并受其影响。
\index{proof!assistant}%
Martin-L\"{o}f 类型论的三种主要变体分别构成了 \NuPRL \cite{constable+86nuprl-book}、\Coq~\cite{Coq} 和 \Agda \cite{norell2007towards} 等类型论计算机实现的基础。本书给出的理论与这些表述在多个方面有所不同，其中有些差异对于同伦解释至关重要，另一些则是技术上的便利，或涉及尚未在同伦背景下研究的概念。

\index{type theory!intensional}%
\index{type theory!extensional}%
\index{intensional type theory}%
\index{extensional!type theory}%
最为显著的是，这里描述的类型论源自 Martin-L\"{o}f 类型论的\emph{内涵（intensional）}版本~\cite{Martin-Lof-1973}，而非其\emph{外延（extensional）}版本~\cite{Martin-Lof-1979}。外延理论不区分判断相等性和命题相等性，而内涵理论则将判断相等性视为纯粹定义性的，并允许对相等类型的一种更为宽泛的、\emph{证明相关的（proof-relevant）}解释——这对同伦解释至关重要。从同伦视角来看，外延类型论局限于同伦离散的集合（参见 \cref{sec:basics-sets}），而内涵理论则允许具有更高维结构的类型。\NuPRL 系统~\cite{constable+86nuprl-book} 是外延的，而 \Coq~\cite{Coq} 和 \Agda~\cite{norell2007towards} 都是内涵的。在内涵类型论中，还存在若干在相等证明的结构上有所不同的变体。我们在此依赖的是最宽松的解释，它承认对相等的\emph{证明相关}解释；而更受限的变体则会施加诸如\emph{相等证明唯一性（UIP）}~\cite{Streicher93}
\indexsee{UIP}{uniqueness of identity proofs}%
\index{uniqueness!of identity proofs}%（即断言任何两个相等证明都是判断相等的），以及\emph{K 公理}~\cite{Streicher93}
\index{axiom!Streicher's Axiom K}
（即断言相等的唯一证明是自反性，在判断相等的意义下）之类的限制。这些额外要求可以在 \Coq 和 \Agda 系统中选择性地施加。

%(In the terminology of \cref{cha:hlevels} such a type theory is about $0$-truncated types.)

内涵理论之间的另一个差异在于判断相等性的强度，尤其是关于函数类型的对象。在此，我们将唯一性原理\index{uniqueness!principle}（$\eta$-变换）$f \jdeq \lam{x} f(x)$ 作为一条判断相等原理纳入。例如，在 \cref{sec:univalence-implies-funext} 中用到这一原理来证明泛等性蕴含命题性的函数外延性。其他类型有时也会考虑唯一性原理。
例如，Cartesian 积 $A\times B$ 的唯一性原理\index{uniqueness!principle!for product types}就是我们在 \cref{sec:finite-product-types} 中构造的命题相等 $\uniq{A\times B}$ 的判断版本，它表述为 $u \jdeq (\proj1(u),\proj2(u))$。
这一原理以及依值对类型的相应版本都是合理的选择（尽管我们并未采纳），但我们不能将所有此类规则悉数纳入，因为相等类型所对应的唯一性原理会将所有高阶同伦结构抹平。
因此，我们\emph{不得不}将其排除，于是问题就变成在哪里划定界限。关于归纳类型，我们将在~\cref{sec:htpy-inductive} 进一步讨论这些问题。

对我们的目的而言，至关重要的一点是：函数的（命题）相等被取为\emph{外延的}（此处的``外延的''含义与上文不同！）。
这不是本章所述规则的推论；它将由 \cref{axiom:funext} 表达。
\index{function extensionality}%
这一决定对我们的目的意义重大，因为它明确了函数的相等性符合数学中的惯常期望。尽管我们将 \cref{axiom:funext} 作为公理纳入，但它可以从泛等公理与函数类型的唯一性原理\index{uniqueness!principle!for function types}（参见 \cref{sec:univalence-implies-funext}）推导出来，也可以从区间类型的存在性（参见 \cref{thm:interval-funext}）推导出来。

对于积、$\Sigma$ 类型、余积、自然数等归纳类型（参见 \cref{cha:induction}），归纳与递归的表述方式还有额外的选择。
\index{pattern matching}%
我们以\emph{归纳原理}为基础，将\emph{模式匹配（pattern matching）}视为由其派生；反过来，以模式匹配为基础、归纳原理为派生也是可以的，参见 \cref{cha:rules}。
后一种做法通常还会允许\emph{深层（deep）}模式匹配；参见~\cite{Coquand92Pattern}。
我们这样选择有几方面原因。
其一，在范畴语义中自然获得的就是归纳原理。
其二，界定允许哪些形式的（深层）模式匹配是相当棘手的；
例如，\Agda 的
\index{proof!assistant!Agda@\textsc{Agda}}%
模式匹配默认就能证明 K 公理，
\index{axiom!Streicher's Axiom K}%
尽管标志 \texttt{--without-K} 可以阻止这一点~\cite{CDP14}。
最后，对于\emph{高阶}归纳类型（参见 \cref{cha:hits}），深层模式匹配尚未得到充分理解。
因此，我们将只使用 \cref{sec:pattern-matching} 中所描述的那类模式匹配，它们直接等价于归纳原理的应用。

\index{proof!assistant!Coq@\textsc{Coq}}%
与 \Coq 的类型论不同，我们不包含原始的命题类型，而是如 \cref{sec:pat} 所讨论的，采纳\emph{命题即类型（propositions-as-types, PAT）}原则，将命题与类型等同。
这一原则最初由 de Bruijn~\cite{deBruijn-1973}、Howard~\cite{howard:pat}、Tait~\cite{Tait-1968} 和 Martin-L\"{o}f~\cite{Martin-Lof-1972} 提出。
（关于我们这一决定的更详细解释，见 \cref{subsec:pat?,subsec:hprops}。）

然而，我们的确包含了一个完整的累积宇宙层级，使得类型形成判断与相等判断成为宇宙的归属判断与相等判断的实例。
为了方便，我们将宇宙的对象直接视为类型，而非类型的代码；用 \cite{martin-lof:bibliopolis} 的术语来说，这意味着我们使用的是``Russell 风格宇宙''，而非``Tarski 风格宇宙''。
\index{type!universe!Tarski-style}%
\index{type!universe!Russell-style}%
另一种方案是使用 Tarski 风格宇宙，此时需要一个显式的强制转换函数\index{coercion, universe-raising}将宇宙的元素 $A:\UU$ 转化为类型 $\mathsf{El}(A)$，并约定在非形式化推理时省略此强制转换。

我们还将宇宙层级视为累积的：对于每个 $j\geq i$，$\UU_i$ 中的每个类型也属于 $\UU_j$。
形式上实现累积性有几种不同方法。最简单的是直接加入一条规则：若 $A:\UU_i$，则 $A:\UU_j$。
然而，这会带来一个恼人的后果：对于类型族 $B:A\to \UU_i$，我们无法推出 $B:A\to\UU_j$，尽管可以推出 $\lam{a} B(a) : A\to\UU_j$。
一种更精细的方案是引入由 $\UU_i<:\UU_j$ 生成的判断子类型关系 $<:$ 来解决此问题，但这会使类型论的研究变得更为复杂。
另一种方案是引入显式的强制转换函数 $\uparrow : \UU_i \to \UU_j$，同样可在非形式化推理时省略。

宇宙也不必以自然数为指标并按线性序排列。
在某些情况下，更合适的做法是仅假定每个宇宙都是某个更大宇宙的元素，同时满足一种``有向性''性质——即任意两个宇宙都共同被包含于某个更大的宇宙之中。
此外还有许多其他可能的变体，例如引入一个包含所有 $\UU_i$ 的宇宙 $\UU_{\omega}$（甚或是更高阶的``大基数''类型宇宙），或将整个层级结构内化为一个类型族 $\lam{i} \UU_i$。
后一种做法事实上已在 \Agda 中实现。

相等类型的道路归纳原理由 Martin-L\"{o}f~\cite{Martin-Lof-1973} 提出。
Martin-L\"of 类型论设定下的定基点道路归纳（based path induction）规则归功于 Paulin-Mohring~\cite{Moh93}；它可以看作是 \NuPRL~\cite[Section~8.1]{constable+86nuprl-book} 系统中关于假设判断的``逐点函数性（pointwise functionality）''
\index{pointwise!functionality}%
概念在内涵类型论中的推广。
Martin-L\"of 的规则蕴含 Paulin-Mohring 的规则，这一事实由 Streicher（使用 K 公理，见~\cref{sec:hedberg}）、Altenkirch 与 Goguen（如~\cref{sec:identity-types}），以及最终由 Hofmann（在不使用宇宙的情况下，如~\cref{ex:pm-to-ml}）所证明；详情参见~\cite[\S1.3 and Addendum]{Streicher93}。

\sectionExercises

\begin{ex}\label{ex:composition}
  给定函数 $f:A\to B$ 与 $g:B\to C$，定义它们的\define{复合（composite）}
  \indexdef{composition!of functions}%
  \indexdef{function!composition}%
  $g\circ f:A\to C$。
  \index{associativity!of function composition}%
  证明我们有 $h \circ (g\circ f) \jdeq (h\circ g)\circ f$。
\end{ex}

\begin{ex}\label{ex:pr-to-rec}
  仅使用投影映射，推导积类型 $A \times B$ 的递归原理 $\rec{A\times B} $，并验证定义相等性成立。
  对 $\Sigma$ 类型做同样的推导。
\end{ex}

\begin{ex}\label{ex:pr-to-ind}
  仅使用投影映射和命题唯一性原理 $\uniq{A\times B}$，推导积类型 $A \times B$ 的归纳原理 $\ind{A\times B}$。
  验证定义相等性成立。
  将 $\uniq{A\times B}$ 推广到 $\Sigma$ 类型，并对 $\Sigma$ 类型做同样的推导。
  \emph{(这需要 \cref{cha:basics} 中的概念。)}
\end{ex}

\begin{ex}\label{ex:iterator}
\index{iterator!for natural numbers}
假设仅给定自然数的\emph{迭代器（iterator）}
\[\ite : \prd{C:\UU} C \to (C \to C) \to \nat \to C \]
及其定义方程
\begin{align*}
\ite(C,c_0,c_s,0)  &\defeq c_0, \\
\ite(C,c_0,c_s,\suc(n)) &\defeq c_s(\ite(C,c_0,c_s,n)),
\end{align*}，
推导出一个类型与递归子 $\rec{\nat}$ 相同的函数。
利用 $\nat$ 的归纳原理证明，该函数在命题层面满足递归子的定义方程。
\end{ex}

\begin{ex}\label{ex:sum-via-bool}
\index{type!coproduct}%
证明：若定义 $A + B \defeq \sm{x:\bool} \rec{\bool}(\UU,A,B,x)$，则可给出 $\ind{A+B}$ 的一个定义，使得 \cref{sec:coproduct-types} 中所述的定义相等性成立。
\end{ex}

\begin{ex}\label{ex:prod-via-bool}
\index{type!product}%
证明：若定义 $A \times B \defeq \prd{x:\bool}\rec{\bool}(\UU,A,B,x)$，则可给出 $\ind{A\times B}$ 的一个定义，使得 \cref{sec:finite-product-types} 中所述的定义相等性在命题层面（即利用相等类型）成立。
\emph{(这需要用到 \cref{sec:compute-pi} 中引入的函数外延性公理。)}
\end{ex}

\begin{ex}\label{ex:pm-to-ml}
给出一个不使用宇宙的替代推导，从 $\indid{A}$ 导出 $\indidb{A}$。
  \emph{(使用后续章节的概念最为容易。)}
\end{ex}

\begin{ex}\label{ex:nat-semiring}
  \index{multiplication!of natural numbers}%
  使用 $\rec{\nat}$ 定义乘法和指数运算。
  仅使用 $\ind{\nat}$ 验证 $(\nat,+,0,\times,1)$ 构成一个半环\index{semiring}。
  你可能还需要用到相等的对称性与传递性，\cref{lem:opp,lem:concat}。
\end{ex}

\begin{ex}\label{ex:fin}
  \index{finite!sets, family of}%
  定义 \cref{sec:universes} 末尾提到的类型族 $\Fin : \nat \to \UU$，以及 \cref{sec:pi-types} 中提到的依值函数 $\fmax : \prd{n:\nat} \Fin(n+1)$。
\end{ex}

\begin{ex}\label{ex:ackermann}
  \indexdef{function!Ackermann}%
  \indexdef{Ackermann function}%
  证明 Ackermann 函数 $\ack : \nat \to \nat \to \nat$ 可以仅用 $\rec{\nat}$ 定义，且满足以下方程：
  \begin{align*}
    \ack(0,n) &\jdeq \suc(n), \\
    \ack(\suc(m),0) &\jdeq \ack(m,1), \\
    \ack(\suc(m),\suc(n)) &\jdeq \ack(m,\ack(\suc(m),n)).
  \end{align*}
\end{ex}

\begin{ex}\label{ex:neg-ldn}
  证明对任意类型 $A$，有 $\neg\neg\neg A \to \neg A$。
\end{ex}

\begin{ex}\label{ex:tautologies}
  使用命题即类型解释，推导以下重言式。
  \begin{enumerate}
  \item 若 $A$ 成立，则（若 $B$ 成立，则 $A$ 成立）。
  \item 若 $A$ 成立，则非（非 $A$）。
  \item 若（非 $A$ 或非 $B$）成立，则非（$A$ 且 $B$）。
  \end{enumerate}
\end{ex}

\begin{ex}\label{ex:not-not-lem}
  使用命题即类型，推导排中律的双重否定，即证明\emph{非（非（$P$ 或非 $P$））}。
\end{ex}

\begin{ex}\label{ex:without-K}
  为什么相等类型的归纳原理不允许我们构造一个函数 $f: \prd{x:A}{p:\id{x}{x}} (\id{p}{\refl{x}})$，使其满足定义方程
  \[ f(x,\refl{x}) \defeq \refl{\refl{x}} \quad ?\]？
\end{ex}

\begin{ex}\label{ex:subtFromPathInd}
  证明同一者的不可区分性可由道路归纳导出。  
\end{ex}

\begin{ex}\label{ex:add-nat-commutative}
  证明自然数的加法满足交换律：$\prd{i,j:\nat} (i+j=j+i)$。
\end{ex}

% Local Variables:
% TeX-master: "hott-online"
% End:

\chapter{同伦类型论}
\label{cha:basics}

同伦类型论的核心新思想是：类型可以被视为同伦论中的空间，或范畴论中的高维群胚。

\index{classical!homotopy theory|(}
\index{higher category theory|(}
我们先简要回顾同伦论与高维范畴论之间的联系。
在经典同伦论中，一个空间 $X$ 是一个配备了拓扑的点集，
\indexsee{space!topological}{topological space}
\index{topological!space}
而点 $x$ 与 $y$ 之间的一条道路则用一个连续映射 $p : [0,1] \to X$ 表示，其中 $p(0) = x$ 且 $p(1) = y$。
\index{path!topological}
\index{topological!path}
可以把这个函数理解为在每个"时刻"给出 $X$ 中的一个点。
对于许多目的而言，道路的严格相等（即逐点相等的函数）是一个过分精细的概念。
例如，我们可以定义道路的拼接运算（若 $p$ 是从 $x$ 到 $y$ 的道路，$q$ 是从 $y$ 到 $z$ 的道路，则拼接 $p \ct q$ 是从 $x$ 到 $z$ 的道路）和取逆运算（$\opp p$ 是从 $y$ 到 $x$ 的道路）。
然而，这些运算之间有一些自然的等式，对于严格相等而言并不成立：
例如，道路 $p \ct \opp p$（它随着时间从 $0$ 到 $1$，从 $x$ 走到 $y$，再沿原路返回）并不严格等于恒等道路（它始终静止在 $x$ 点）。

解决之道在于引入一种更粗略的道路相等概念，称为\emph{同伦}。
\index{homotopy!topological}
一对连续映射 $f : X_1 \to X_2$ 与 $g : X_1\to X_2$ 之间的一个同伦，是一个连续映射 $H : X_1
\times [0, 1] \to X_2$，满足 $H(x, 0) = f (x)$ 与 $H(x, 1) =
g(x)$。
特别地，对于从 $x$ 到 $y$ 的两条道路 $p$ 和 $q$，同伦是一个连续映射 $H : [0,1] \times [0,1] \rightarrow X$，
使得对所有 $s\in [0,1]$，有 $H(s,0) = p(s)$ 和 $H(s,1) = q(s)$。
在这种情况下，我们还要求对所有 $t\in [0,1]$，有 $H(0,t) = x$ 且 $H(1,t)=y$，
从而对每个 $t$，函数 $H(\blank,t)$ 仍然是一条从 $x$ 到 $y$ 的道路；
这类同伦被称为\emph{保端点的}或\emph{相对于端点的}。
在简单的情形中，我们可以把正方形 $[0,1]\times [0,1]$ 在 $H$ 下的像看作"填满"了 $p$ 与 $q$ 之间的"空间"，
尽管对于一般的 $X$ 来说这并不真正有意义；
更好的理解方式是把 $H$ 看作将 $p$ 连续形变为 $q$ 且不移动端点的过程。
由于 $[0,1]\times [0,1]$ 是 2 维的，我们也将 $H$ 称为\emph{道路之间的 2 维道路}。\index{path!2-}

例如，因为 $p \ct \opp p$ 沿同一条路线走出去又走回来，所以你可以将 $p \ct \opp p$ 连续地缩为恒等道路——例如，它不会被空间中的空洞"钩住"。
同伦是一种等价关系，且拼接、取逆等运算都与它相容。
此外，在某点 $x_0$ 处环路\index{loop}的同伦等价类（两条环路 $p$ 和 $q$ 在存在它们之间的\emph{保基点的}同伦时被视为相等，即如上所述的同伦 $H$ 并额外满足对所有 $t$ 有 $H(0,t) =
H(1,t) = x_0$）构成一个群，称为\emph{基本群}。\index{fundamental!group}
这个群是空间的一个\emph{代数不变量}，可用于探究两个空间是否\emph{同伦等价}
（即存在来回的连续映射，其复合均同伦于恒等映射），因为同伦等价的空间具有同构的基本群。

由于同伦本身是一种 2 维道路，自然就有 3 维的\emph{同伦之间的同伦}\index{path!3-}的概念，
进而还有\emph{同伦之间的同伦之间的同伦}，如此等等。
这个由点、道路、同伦、同伦之间的同伦……构成的无限塔，连同基本群等代数运算，是被称为（弱）\emph{$\infty$-群胚}的代数结构的一个实例。
一个 $\infty$-\emph{群胚}\index{.infinity-groupoid@$\infty$-groupoid}包含一组对象，
然后是对象之间的一组\emph{态射}\indexdef{morphism!in an .infinity-groupoid@in an $\infty$-groupoid}，
接着是\emph{态射之间的态射}，依此类推，并配有某种复杂的代数结构；
第 $k$ 层的态射称为\define{$k$-态射 ($k$-morphism)}\indexdef{k-morphism@$k$-morphism}。
每一层的态射都有单位、复合和逆运算，但这些运算具有所谓的"弱性"：
具体来说，它们满足群胚定律（复合的结合律、单位元是复合的单位、逆元可消去）仅在相差一个下一层态射的意义下成立，而这种弱性会引发进一步的结构。
例如，由于态射复合 $p \ct (q \ct r) = (p \ct q) \ct r$ 的结合律本身就是一个更高维的态射，我们需要额外的运算来关联各种结合性证明：
将 $p \ct
(q \ct (r \ct s))$ 用不同方式重新结合为 $((p \ct q) \ct r) \ct s$，就引出了 Mac Lane（麦克兰恩）的五边形\index{pentagon, Mac Lane}。
弱性还会造成不同层次之间非平凡的相互作用。

每个拓扑空间 $X$ 都有一个\emph{基本 $\infty$-群胚}
\index{.infinity-groupoid@$\infty$-groupoid!fundamental}
\index{fundamental!.infinity-groupoid@$\infty$-groupoid}，
其 $k$-态射就是 $X$ 中的 $k$ 维道路。
$\infty$-群胚的弱性直接对应于这样一个事实：道路仅在同伦意义下构成群，其中 $(k+1)$-道路充当 $k$-道路之间的同伦。
此外，将空间视为 $\infty$-群胚的观点保留了足够多的空间信息以开展同伦论：
基本 $\infty$-群胚的构造与一个 $\infty$-群胚作为空间的几何\index{geometric realization}实现构成伴随\index{adjoint!functor}关系，
且该伴随保持同伦论（这称为\emph{同伦假说/定理}，%
\index{hypothesis!homotopy}%
\index{homotopy!hypothesis}%
因为它是假说还是定理取决于如何定义 $\infty$-群胚）。
例如，你可以容易地定义一个 $\infty$-群胚的基本群；而如果你计算一个空间的基本 $\infty$-群胚的基本群，结果将与该空间基本群的经典定义相符。
正是由于这种对应关系，同伦论与高维范畴论紧密相连。

\index{classical!homotopy theory|)}%
\index{higher category theory|)}%

\mentalpause

现在，在同伦类型论中，每个类型都可以被看作具有 $\infty$-群胚的结构。
回想一下，对于任何类型 $A$ 和任意 $x, y : A$，我们都有一个相等类型 $\id[A]{x}{y}$，也写作 $\idtype[A]{x}{y}$ 或简写为 $x = y$。
从逻辑上讲，我们可以将 $x = y$ 的元素看作 $x$ 与 $y$ 相等的证据，或 $x$ 与 $y$ 的相等证明。
此外，类型论（不同于一阶逻辑等系统）允许我们将 $\id[A]{x}{y}$ 的这类元素本身也视为个体，从而成为进一步命题的主语。
因此，我们可以\emph{迭代}相等类型：可以构造相等证明 $p, q$ 之间的相等证明类型 $\id[{(\id[A]{x}{y})}]{p}{q}$，
以及类型 $\id[{(\id[{(\id[A]{x}{y})}]{p}{q})}]{r}{s}$，如此等等。
这个相等类型塔的结构，精确地对应于空间中的连续道路及其间的（高阶）同伦，或一个 $\infty$-群胚\index{.infinity-groupoid@$\infty$-groupoid}。

因此，我们通常将元素$p : \id[A]{x}{y}$称为从 $x$ 到 $y$ 的一条\define{道路（path）}；
\index{path}
我们称 $x$ 为其\define{起点（start point）}
\indexdef{start point of a path}
\indexdef{path!start point of}，
称 $y$ 为其\define{终点（end point）}。
\indexdef{end point of a path}
\indexdef{path!end point of}
两条起点和终点相同的道路$p,q : \id[A]{x}{y}$被称为\define{平行的（parallel）}，
\indexdef{parallel paths}
\indexdef{path!parallel}此时，元素$r : \id[{(\id[A]{x}{y})}]{p}{q}$可以看作一个同伦，或态射之间的态射；
我们通常将其称为\define{2-道路（2-path）}
\indexdef{path!2-}\indexsee{2-path}{path, 2-}%
或\define{2-维道路（2-dimensional path）}。
\index{dimension!of paths}%
\indexsee{2-dimensional path}{path, 2-}\indexsee{path!2-dimensional}{path, 2-}%
类似地，$\id[{(\id[{(\id[A]{x}{y})}]{p}{q})}]{r}{s}$是两个平行2-维道路之间的\define{3-维道路（3-dimensional path）}
\indexdef{path!3-}\indexsee{3-path}{path, 3-}\indexsee{3-dimensional path}{path, 3-}\indexsee{path!3-dimensional}{path, 3-}%
的类型，依此类推。
如果类型 $A$ 是"集合式的"，例如 \nat，那么这些迭代的相等类型将平淡无奇（参见\cref{sec:basics-sets}）；但在一般情况下，它们可以模拟非平凡的同伦类型。

%% (Obviously, the
%% notation ``$\id[A]{x}{y}$'' has its limitations here.  The style
%% $\idtype[A]{x}{y}$ is only slightly better in iterations:
%% $\idtype[{\idtype[{\idtype[A]{x}{y}}]{p}{q}}]{r}{s}$.)

同伦类型论与经典同伦论的一个重要区别在于，同伦类型论以如下意义提供了一种对空间的\emph{综合（synthetic）}描述。
\index{synthetic mathematics}%
\index{geometry, synthetic}%
\index{Euclid of Alexandria}%
综合几何是Euclid ~\cite{Euclid} 风格的几何学：从一些基本概念（点与直线）、构造（连接任意两点的直线）和公理（所有直角都相等）出发，通过逻辑推导出结论。
这与\emph{解析}几何形成对比，
\index{analytic mathematics}%
在解析几何中，点和直线等概念通过 $\R^n$ 中的笛卡尔坐标具体表示——直线是点的集合——而基本构造和公理由此表示导出。
经典同伦论是解析的（空间和道路由点构成），而同伦类型论是综合的：点、道路以及道路之间的道路是基本的、不可分割的、原始的概念。

此外，同伦类型论的一个令人惊叹之处是，所有基本构造和公理——所有高阶群胚结构——都自动从相等类型的归纳原理中产生。
回顾\cref{sec:identity-types}，该原理是说，如果

\begin{itemize}
\item 对于每个 $x,y:A$ 和每个 $p:\id[A]xy$，我们有一个类型 $D(x,y,p)$，且
\item 对于每个 $a:A$，我们有一个元素$d(a):D(a,a,\refl a)$，
\end{itemize}

那么

\begin{itemize}
\item 对于\emph{任意}两个元素 $x,y:A$ 和 $p:\id[A]xy$，都存在一个元素$\indid{A}(D,d,x,y,p):D(x,y,p)$，使得$\indid{A}(D,d,a,a,\refl a) \jdeq d(a)$。
\end{itemize}

换言之，给定依值函数
\begin{align*}
D & :\prd{x,y:A} (\id{x}{y}) \to \type\\
d & :\prd{a:A} D(a,a,\refl{a})
\end{align*}
存在一个依值函数
\[\indid{A}(D,d):\prd{x,y:A}{p:\id{x}{y}} D(x,y,p)\]
使得对每个 $a:A$ 有
\begin{equation}\label{eq:Jconv}
\indid{A}(D,d,a,a,\refl{a})\jdeq d(a)
\end{equation}。
通常，每次应用这个归纳规则时，我们要么不关心所定义的具体函数，要么立即给它另一个名字。

非正式地说，相等类型的归纳原理是说：如果我们想构造一个依赖于相等类型元素 $p:\id[A]xy$ 的对象（或证明一个陈述），那么只需在 $x$ 和 $y$ 判断相等且 $p$ 是自反元素 $\refl{x}:x=x$（判断相等）的特殊情况下进行构造（或证明）即可。
在非正式表述中，我们可以用"由归纳，只需假设……"这样的短语来表达。
这种"归约到自反性情形"的做法，类似于在自然数上作普通归纳证明时归约到"基础情况"和"归纳步骤"，也类似于在对不相交并或析取作分情形讨论时归约到"左情况"和"右情况"。\index{induction principle!for identity type}%

在关于自然数的归纳证明的语境中，"转换规则"~\eqref{eq:Jconv}可能不那么熟悉，但在与之相关的递归定义概念中，却有一个类似的想法。
如果序列\index{sequence} $(a_n)_{n\in \mathbb{N}}$是通过给定 $a_0$ 并用 $a_n$ 表示 $a_{n+1}$ 来定义的，那么结果序列的第 $0$ 项\emph{确实}就是给定的 $a_0$，并且给定的联系 $a_{n+1}$ 与 $a_n$ 的递推关系对结果序列成立。
这一点可能看起来过于显然，不值一提；但倘若我们将递归定义视为计算序列值的算法\index{algorithm}，那么它恰恰就是执行该算法的过程。
规则~\eqref{eq:Jconv} 与此类似：如果我们通过指定 $p$ 为 $\refl{x}:x=x$ 时的值来为所有 $p:x=y$ 定义对象 $f(p)$，那么我们指定的值事实上就是 $f(\refl{x})$ 的值。

这个归纳原理赋予每个类型以 $\infty$-群胚\index{.infinity-groupoid@$\infty$-groupoid}的结构，并使两个类型之间的每个函数具有相应群胚之间的 $\infty$-函子\index{.infinity-functor@$\infty$-functor}结构。
从数学的观点来看，这很有趣，因为它提供了一种处理 $\infty$-群胚的新方法。
从类型论的观点来看，这同样有趣，因为它揭示了与每个类型和函数相关联的新运算。
在本章剩余的部分里，我们将开始探索这种结构。

\section{类型是高阶群胚}
\label{sec:equality}

\index{type!identity|(}%
\index{path|(}%
\index{.infinity-groupoid@$\infty$-groupoid!structure of a type|(}%
我们现在从归纳原理出发，推导高阶群胚结构的初步内容。
我们从相等性的对称性开始，用拓扑学的语言来说，这意味着"道路可以反转"。

\begin{lem}\label{lem:opp}
  对每个类型 $A$ 以及任意 $x,y:A$，存在一个函数
  \begin{equation*}
    (x= y)\to(y= x)
  \end{equation*}
  记作 $p\mapsto \opp{p}$，使得对每个 $x:A$ 有 $\opp{\refl{x}}\jdeq\refl{x}$。
  我们称 $\opp{p}$ 为 $p$ 的\define{逆（inverse）}。
  \indexdef{path!inverse}%
  \indexdef{inverse!of path}%
  \index{equality!symmetry of}%a
  \index{symmetry!of equality}%
\end{lem}

鉴于这是我们第一次将某个陈述表述为"引理"或"定理"，不妨稍作停顿，思考一下这意味着什么。
回想一下，命题（可被证明的陈述）等同于类型，而引理和定理（已被证明的陈述）则等同于\emph{有实例的}类型。
因此，引理或定理的陈述应被转译为一个类型，如 \cref{sec:pat} 所示，而其证明则被转译为该类型的一个实例。
根据全称量词"对每个"的解释，与 \cref{lem:opp} 对应的类型是
\[ \prd{A:\UU}{x,y:A} (x= y)\to(y= x). \]
\cref{lem:opp} 的证明将在于构造该类型的一个元素，即对某个 $f$ 推导出判断 $f:\prd{A:\UU}{x,y:A} (x= y)\to(y= x)$。
然后，我们引入符号 $\opp{(\blank)}$ 来表示这个元素 $f$，其中参数 $A$、$x$ 和 $y$ 被省略并由上下文推断。
（如 \cref{sec:types-vs-sets} 所述，次要陈述"对每个 $x:A$，有 $\opp{\refl{x}}\jdeq\refl{x}$"应被视为一个独立的判断。）

\begin{proof}[第一次证明]
  假设给定 $A:\UU$，
  并令 $D:\prd{x,y:A}(x= y) \to \type$ 为由 $D(x,y,p)\defeq (y= x)$ 定义的类型族。
  换句话说，$D$ 是一个函数，它为任意 $x,y:A$ 和 $p:x=y$ 指派一个类型，即类型 $y=x$。
  于是我们有一个元素
  \begin{equation*}
    d\defeq \lam{x} \refl{x}:\prd{x:A} D(x,x,\refl{x}).
  \end{equation*}
  因此，相等类型的归纳原理给了我们一个元素
  \narrowequation{ \indid{A}(D,d,x,y,p): (y= x)}
  对每个 $p:(x= y)$。
  现在我们可以定义所需的函数 $\opp{(\blank)}$ 为 $\lam{p} \indid{A}(D,d,x,y,p)$，即令 $\opp{p} \defeq \indid{A}(D,d,x,y,p)$。
  转换规则~\eqref{eq:Jconv} 给出了 $\opp{\refl{x}}\jdeq \refl{x}$，如所需。
\end{proof}

我们以非常形式化的风格写出了这个证明，这在大家对相等类型的归纳规则还不熟悉时或许有所帮助。
若追求更形式化，我们可以说 \cref{lem:opp} 及其证明共同构成了判断
\begin{narrowmultline*}
  \lam{A}{x}{y}{p} \indid{A}((\lam{x}{y}{p} (y=x)), (\lam{x} \refl{x}), x, y, p)
  \narrowbreak : \prd{A:\UU}{x,y:A} (x= y)\to(y= x)
\end{narrowmultline*}
（以及一个额外的相等判断）。
然而，最终我们更倾向于使用更自然的语言，如下面这个等价的证明所示。

\begin{proof}[第二次证明]
  我们想为每对 $x,y:A$ 和每个 $p:x=y$ 构造一个元素 $\opp{p}:y=x$。
  由归纳，只需处理 $y$ 为 $x$ 且 $p$ 为 $\refl{x}$ 的情形。
  但在这种情形下，$p$ 的类型 $x=y$ 和我们要在其中构造 $\opp{p}$ 的类型 $y=x$ 都仅仅是 $x=x$。
  因此，在"自反性情形"下，我们可以直接定义 $\opp{\refl{x}}$ 为 $\refl{x}$。
  一般情形则由归纳原理得出，而转换规则 $\opp{\refl{x}}\jdeq\refl{x}$ 恰好就是我们在自反性情形中给出的证明。
\end{proof}

在接下来的几个证明中，我们将同时采用两种风格书写，以帮助读者适应后一种。
接下来，我们证明相等的传递性，换言之，也就是道路的"拼接"。

\begin{lem}\label{lem:concat}
  对每个类型 $A$ 以及每个 $x,y,z:A$，存在一个函数
  \begin{equation*}
  (x= y) \to   (y= z)\to (x=  z),
  \end{equation*}
  记作 $p \mapsto q \mapsto p\ct q$，使得对任意 $x:A$，有 $\refl{x}\ct \refl{x}\jdeq \refl{x}$。
  我们将 $p\ct q$ 称为 $p$ 与 $q$ 的\define{拼接（concatenation）}或\define{复合（composite）}。
  \indexdef{path!concatenation}%
  \indexdef{path!composite}%
  \indexdef{concatenation of paths}%
  \indexdef{composition!of paths}%
  \index{equality!transitivity of}%
  \index{transitivity!of equality}%
\end{lem}

请注意，我们选择的道路拼接记法与函数复合的顺序相反：由 $p:x=y$ 和 $q:y=z$ 得到 $p\ct q : x=z$，而由 $f:A\to B$ 和 $g:B\to C$ 得到 $g\circ f : A\to C$（参见 \cref{ex:composition}）。

\begin{proof}[第一种证法]
  所需函数的类型为 $\prd{x,y,z:A} (x= y) \to   (y= z)\to (x=  z)$。
  不过，我们将改为定义一个具有等价类型 $\prd{x,y:A} (x= y) \to \prd{z:A} (y= z)\to (x=  z)$ 的函数，这使得我们可以应用两次道路归纳。
  令 $D:\prd{x,y:A} (x=y) \to \type$ 为类型族
  \begin{equation*}
    D(x,y,p)\defeq \prd{z:A}{q:y=z} (x=z).
  \end{equation*}
  注意 $D(x,x,\refl x) \jdeq \prd{z:A}{q:x=z} (x=z)$。
  因此，为了对这个 $D$ 应用相等类型的归纳原理，我们需要一个类型为
  \begin{equation}\label{eq:concatD}
    \prd{x:A} D(x,x,\refl{x})
  \end{equation}
  的函数，也就是说，类型为
  \[ \prd{x,z:A}{q:x=z} (x=z). \]
  的函数。
  现在令 $E:\prd{x,z:A}{q:x=z}\type$ 为类型族 $E(x,z,q)\defeq (x=z)$。
  注意 $E(x,x,\refl x) \jdeq (x=x)$。
  于是，我们有如下函数
  \begin{equation*}
    e(x) \defeq \refl{x} : E(x,x,\refl{x}).
  \end{equation*}
  对 $E$ 应用相等类型的归纳原理，我们得到一个函数
  \begin{equation*}
    d : \prd{x,z:A}{q:x=z} E(x,z,q).
  \end{equation*}
  但 $E(x,z,q)\jdeq (x=z)$，故 $d$ 的类型为~\eqref{eq:concatD}。
  于是，我们可以利用这个函数 $d$，对 $D$ 应用相等类型的归纳原理，从而得到所需的类型为
  \begin{equation*}
    \prd{x,y:A} (x= y) \to \prd{z:A} (y= z)\to (x=  z)
  \end{equation*}
  的函数，因此也就得到了 $\prd{x,y,z:A} (y=z) \to (x=y) \to (x=z)$。
  两次归纳原理的计算规则给出：对任意 $x:A$，有 $\refl{x}\ct \refl{x}\jdeq \refl{x}$。
\end{proof}

\begin{proof}[第二种证法]
  我们要对每个 $x,y,z:A$ 以及每个 $p:x=y$ 和 $q:y=z$，构造 $x=z$ 的一个元素。
  对 $p$ 作归纳，只需假设 $y$ 是 $x$ 且 $p$ 是 $\refl{x}$。
  此时，$q$ 的类型 $y=z$ 即为 $x=z$。
  再对 $q$ 作归纳，只需假设 $z$ 是 $x$ 且 $q$ 是 $\refl{x}$。
  但此时，$x=z$ 即为 $x=x$，而我们已有 $\refl{x}:(x=x)$。
\end{proof}

读者很可能会觉得这个引理的证明过于曲折。
事实上，我们本可以在对 $p$ 归纳后就停下来，因为此时我们要构造的是 $x=z$ 的相等关系，而我们已经有了一个，即 $q$。
为何还要继续对 $q$ 做一次归纳？

答案在于，正如引言中所述，我们做的是\emph{证明相关的}数学。
\index{mathematics!proof-relevant}%
当我们证明一个引理时，我们是在定义某个类型的一个实例；在证明过程中究竟定义了哪个\emph{具体的}元素，这可能很重要，而不仅仅是该实例所属的类型（即引理的\emph{陈述}）。
\cref{lem:concat} 有三种显而易见的证明方式：可以对 $p$ 归纳，可以对 $q$ 归纳，也可以对两者都归纳。
如果用三种不同方式证明，我们将得到同一类型的三个不同元素。
不难证明这三个元素是相等的（见 \cref{ex:basics:concat}），但由于它们并非\emph{定义性}相等，仍可能有理由偏爱其中之一。

对于 \cref{lem:concat}，差异体现在计算规则上。
如果通过对 $p$ 作单次归纳来证明引理，则最终得到的计算规则形如 $\refl{y} \ct q \jdeq q$。
如果通过对 $q$ 作单次归纳来证明，则得到 $p\ct\refl{y}\jdeq p$，而像我们这样通过双重归纳来证明，则只能得到 $\refl{x}\ct\refl{x} \jdeq \refl{x}$。

\index{mathematics!formalized}%
在进行形式化数学时，不对称的计算规则有时会带来便利，因为它们能让计算机自动简化更多内容。
然而，在非形式化的数学中，甚至可以说在形式化的情形中，拥有一个行为不对称的拼接操作、并且需要记住哪一边是"特殊"的，可能会令人困惑。
对称地处理两边能使证明更加稳健；这就是我们如此证明的原因。
（诚然，这是一种风格上的选择。）

下表从"相等"、"同伦"和"高阶群胚"三个视角总结了迄今为止的工作。


\begin{center}
  \medskip
  \begin{tabular}{ccc}
    \toprule
    相等 & 同伦 & $\infty$-群胚\\
    \midrule
    自反性\index{equality!reflexivity of} & 常值道路 & 恒同态射\\
    对称性\index{equality!symmetry of} & 道路的逆 & 逆态射\\
    传递性\index{equality!transitivity of} & 道路的拼接 & 态射的复合\\
    \bottomrule
  \end{tabular}
  \medskip
\end{center}

在实践中，传递性常用于通过一串中间步骤来证明一个等式。
我们将使用诸如 $a=b=c=d$ 这样的常见记法。
如果中间表达式较长，或者我们希望指明每步等式的见证，则可以写作
\begin{align*}
  a &= b & \text{(by $p$)}\\ &= c &\text{(by $q$)} \\ &= d &\text{(by $r$)}.
\end{align*}
无论哪种情况，这种记法都表示 $(p\ct q)\ct r: (a=d)$ 的构造。
（为明确起见，我们选择左结合；不过鉴于下文 \cref{thm:omg}\ref{item:omg4}，选择哪种结合方式影响不大。）
假如恰好 $b$ 和 $c$ 是判断相等的，那么可以写作
\begin{align*}
  a &= b & \text{(by $p$)}\\ &\jdeq c \\ &= d &\text{(by $r$)}
\end{align*}
以表示 $p\ct r : (a=d)$ 的构造。
我们也遵循数学写作中的常见做法：不要求这种记法中的依据（"由 $p$"和"由 $r$"）给出所需的精确见证；
而是允许它们仅提及构造该见证时最关键（或最不显然）的要素。
例如，若"引理 A"陈述了对所有 $x$ 和 $y$ 有 $f(x)=g(y)$，则我们可以用"由引理 A"作为步骤 $f(a)=g(b)$ 的依据，相信读者能推断出我们是在将引理 A 应用于 $x\defeq a$ 和 $y\defeq b$。
若相信读者能够自行补出依据，也可将其完全省略。

现在，由于证明相关性，我们不能在证明了相等的"对称性"和"传递性"后就止步：
我们还需要知道这些作用于相等的\emph{操作}是行为良好的。
（这一问题在集合论中是不可见的，因为在那里对称性和传递性仅仅是相等的\emph{性质}，而非道路上的结构。）
从同伦论的观点来看，拼接和求逆只是高阶群胚结构的"第一层"——我们还需要关于这些操作的\index{coherence}相干律，以及高维的类似操作。
例如，我们需要知道拼接是\emph{结合的}，且求逆关于拼接给出\emph{逆}。

\begin{lem}\label{thm:omg}%[The $\omega$-groupoid structure of types]
  \index{associativity!of path concatenation}%
  \index{unit!law for path concatenation}%
  假设 $A:\type$，$x,y,z,w:A$，且 $p:x= y$，$q:y = z$，$r:z=w$。
  我们有以下性质：
  \begin{enumerate}
  \item $p= p\ct \refl{y}$ 且 $p = \refl{x} \ct p$。\label{item:omg1}
  \item $\opp{p}\ct p=  \refl{y}$ 与 $p\ct \opp{p}= \refl{x}$。\label{item:omg2}
  \item $\opp{(\opp{p})}= p$。\label{item:omg3}
  \item $p\ct (q\ct r)=  (p\ct q)\ct r$。\label{item:omg4}
  \end{enumerate}
\end{lem}

特别要注意，\ref{item:omg1}--\ref{item:omg4} 本身也是命题相等性，属于相等类型\emph{的}相等类型——例如对 $p,q:x=y$，有类型 $p=_{x=y}q$。
从拓扑上看，它们是\emph{道路的道路}，即同伦。
拓扑学中有一个熟知的事实：当我们将道路 $p$ 与逆道路 $\opp p$ 拼接时，我们并不会恰好得到一条常值道路（在类型论中对应于 $\refl{}$）——相反，我们得到的是从 $p\ct\opp p$ 到常值道路的一个同伦，或者说高阶道路。

\begin{proof}[\cref{thm:omg} 的证明]
  所有证明都使用相等类型的归纳原理。
  \begin{enumerate}
  \item \emph{第一个证明：}令 $D:\prd{x,y:A} (x=y) \to \type$ 为由
    \begin{equation*}
      D(x,y,p)\defeq (p= p\ct \refl{y}).
    \end{equation*}
    给出的类型族。
    那么 $D(x,x,\refl{x})$ 就是 $\refl{x}=\refl{x}\ct\refl{x}$。
    由于 $\refl{x}\ct\refl{x}\jdeq\refl{x}$，可得 $D(x,x,\refl{x})\jdeq (\refl{x}=\refl{x})$。
    于是我们得到一个函数
    \begin{equation*}
      d\defeq\lam{x} \refl{\refl{x}}:\prd{x:A} D(x,x,\refl{x}).
    \end{equation*}
    现在，相等类型的归纳原理对每个 $p:x= y$ 给出一个元素 $\indid{A}(D,d,x,y,p):(p= p\ct\refl{y})$。
    另一个等式同理可证。

    \mentalpause

    \noindent
    \emph{第二个证明：}对 $p$ 作归纳，只需考虑 $y$ 为 $x$ 且 $p$ 为 $\refl x$ 的情形。
    但此时，我们有 $\refl{x}\ct\refl{x}\jdeq\refl{x}$。
  \item \emph{第一个证明：}令 $D:\prd{x,y:A} (x=y) \to \type$ 为由
    \begin{equation*}
      D(x,y,p)\defeq (\opp{p}\ct p=  \refl{y}).
    \end{equation*}
    给出的类型族。
    那么 $D(x,x,\refl{x})$ 就是 $\opp{\refl{x}}\ct\refl{x}=\refl{x}$。
    由于 $\opp{\refl{x}}\jdeq\refl{x}$ 且 $\refl{x}\ct\refl{x}\jdeq\refl{x}$，可得 $D(x,x,\refl{x})\jdeq (\refl{x}=\refl{x})$。
    因此我们得到函数
    \begin{equation*}
      d\defeq\lam{x} \refl{\refl{x}}:\prd{x:A} D(x,x,\refl{x}).
    \end{equation*}
    现在，道路归纳对 $A$ 中的每个 $p:x= y$ 给出一个元素 $\indid{A}(D,d,x,y,p):\opp{p}\ct p=\refl{y}$。
    另一个等式同理可证。

    \mentalpause

    \noindent \emph{第二个证明：}通过归纳，只需考虑 $p$ 为 $\refl x$ 的情形。
    但此时，我们有 $\opp{p} \ct p \jdeq \opp{\refl x} \ct \refl x \jdeq \refl x$。

  \item \emph{第一个证明：}令 $D:\prd{x,y:A} (x=y) \to \type$ 为由
    \begin{equation*}
      D(x,y,p)\defeq (\opp{(\opp{p})}= p).
    \end{equation*}
    给出的类型族。
    那么 $D(x,x,\refl{x})$ 就是类型 $(\opp{(\opp{\refl x})}=\refl{x})$。
    但由于对每个 $x:A$ 有 $\opp{\refl{x}}\jdeq \refl{x}$，我们得到 $\opp{(\opp{\refl{x}})}\jdeq \opp{\refl{x}} \jdeq\refl{x}$，从而 $D(x,x,\refl{x})\jdeq(\refl{x}=\refl{x})$。
    因此我们得到函数
    \begin{equation*}
      d\defeq\lam{x} \refl{\refl{x}}:\prd{x:A} D(x,x,\refl{x}).
    \end{equation*}
    现在，道路归纳对每个 $p:x= y$ 给出一个元素 $\indid{A}(D,d,x,y,p):\opp{(\opp{p})}= p$。

    \mentalpause

    \noindent \emph{第二个证明：}通过归纳，只需考虑 $p$ 为 $\refl x$ 的情形。
    但此时，我们有 $\opp{(\opp{p})}\jdeq \opp{(\opp{\refl x})} \jdeq \refl x$。

  \item \emph{第一个证明：}令 $D_1:\prd{x,y:A} (x=y) \to \type$ 为由
    \begin{equation*}
      D_1(x,y,p)\defeq\prd{z,w:A}{q:y= z}{r:z= w} \big(p\ct (q\ct r)=  (p\ct q)\ct r\big).
    \end{equation*}
    给出的类型族。
    那么 $D_1(x,x,\refl{x})$ 是
    \begin{equation*}
      \prd{z,w:A}{q:x= z}{r:z= w} \big(\refl{x}\ct(q\ct r)= (\refl{x}\ct q)\ct r\big).
    \end{equation*}
    为构造此类型的一个元素，令 $D_2:\prd{x,z:A} (x=z) \to \type$ 为类型族
    \begin{equation*}
      D_2 (x,z,q) \defeq \prd{w:A}{r:z=w} \big(\refl{x}\ct(q\ct r)= (\refl{x}\ct q)\ct r\big).
    \end{equation*}
    那么 $D_2(x,x,\refl{x})$ 是
    \begin{equation*}
      \prd{w:A}{r:x=w} \big(\refl{x}\ct(\refl{x}\ct r)= (\refl{x}\ct \refl{x})\ct r\big).
    \end{equation*}
    为构造\emph{这个}类型的一个元素，令 $D_3:\prd{x,w:A} (x=w) \to \type$ 为类型族
    \begin{equation*}
      D_3(x,w,r) \defeq \big(\refl{x}\ct(\refl{x}\ct r)= (\refl{x}\ct \refl{x})\ct r\big).
    \end{equation*}
    那么 $D_3(x,x,\refl{x})$ 是
    \begin{equation*}
      \big(\refl{x}\ct(\refl{x}\ct \refl{x})= (\refl{x}\ct \refl{x})\ct \refl{x}\big)
    \end{equation*}
    它与类型 $(\refl{x} = \refl{x})$ 定义上相等，因此有实例 $\refl{\refl{x}}$。
    于是，应用三次道路归纳规则，我们就得到了整体所需类型的元素。

    \mentalpause

    \noindent \emph{第二个证明：}通过归纳，只需考虑 $p$、$q$ 和 $r$ 均为 $\refl x$ 的情形。
    但此时，我们有
    \begin{align*}
      p\ct (q\ct r)
      &\jdeq \refl{x}\ct(\refl{x}\ct \refl{x})\\
      &\jdeq \refl{x}\\
      &\jdeq (\refl{x}\ct \refl x)\ct \refl x\\
      &\jdeq (p\ct q)\ct r.
    \end{align*}
    因此，$\refl{\refl{x}}$ 是这个类型的实例。 \qedhere
  \end{enumerate}
\end{proof}

\begin{rmk}
  定义这些高阶道路还有其他方式。
  例如，在 \cref{thm:omg}\ref{item:omg4} 中，我们也可以只对一条或两条道路作归纳，而不必对全部三条都作归纳。
  每种选择都会产生一个\emph{定义上}不同的证明，但它们彼此之间都是相等的。
  任意两个具体证明之间的这种相等关系，同样可以用归纳来证明：把所涉及的所有道路都归约为自反性，然后观察到两个证明自身也都归约为自反性即可。
\end{rmk}

鉴于 \cref{thm:omg}\ref{item:omg4}，我们以后通常将 $(p\ct q)\ct r$ 写作 $p\ct q\ct r$，
类似地，用 $p\ct q\ct r\ct s$ 表示 $((p\ct q)\ct r)\ct s$，
依此类推。
为确定起见，我们选择左结合；但从实质上说，选哪一种结合方式并没有区别。
我们一般相信读者能在必要时自行补上 \cref{thm:omg}\ref{item:omg4} 的相应实例来重新加括号。

关于高阶群胚结构，我们其实还没有完全讲完：道路~\ref{item:omg1}--\ref{item:omg4} 还必须满足它们自身的高阶相干律 \index{coherence}，而这些相干律本身也是高阶道路，
\index{associativity!of path concatenation!coherence of}%
\index{globular operad}%
\index{operad}%
\index{groupoid!higher}%
如此"一直推到无穷"（这一点可以通过例如球状算畴（globular operad）的概念来严格表述）。
不过，对于大多数目的而言，并不需要将整个无穷维结构显式地写出。
同伦类型论的一个美妙之处在于，所有这些结构都可以仅仅从相等类型的归纳性质出发予以\emph{证明}，因此我们可以根据需要，将其中的任意多或任意少的部分显式写出。

具体来说，在本书中，我们并不需要那些用于精确描述诸如"所有高阶层面的相干结构"等概念时所涉及的复杂组合学。
除了普通的道路，我们还会用到道路的道路（即类型 $p =_{x=_A y} q$ 中的元素），如前所述，我们称之为\emph{2-道路}\index{path!2-}或\emph{2-维道路}；偶尔可能还会用到道路的道路的道路（即类型 $r = _{p =_{x=_A y} q} s$ 中的元素），我们称之为\emph{3-道路}\index{path!3-}或\emph{3-维道路}。
虽然可以定义一般的\emph{$n$-维道路}
\indexdef{path!n-@$n$-}%
\indexsee{n-path@$n$-path}{path, $n$-}%
\indexsee{n-dimensional path@$n$-dimensional path}{path, $n$-}%
\indexsee{path!n-dimensional@$n$-dimensional}{path, $n$-}%
（参见 \cref{ex:npaths}），但我们并不需要它。

然而，我们确实会用到高维道路的一种尤为重要且简单的情形，即起点与终点相同的情形。
在集合论中，命题 $a=a$ 是完全无趣的，但在同伦论中，从一个点到其自身的道路被称为\emph{环路}\index{loop}，并且承载着许多有趣的高阶结构。
因此，给定一个类型 $A$ 及其上的一个点 $a:A$，我们定义它的\define{环路空间（loop space）}
\index{loop space}%
$\Omega(A,a)$ 为类型 $\id[A]{a}{a}$。
如果点 $a$ 可以从上下文看出，我们有时也简写为 $\Omega A$。

由于 $\Omega A$ 中的任意两个元素都是具有相同起点和终点的道路，它们可以相互拼接；
因此我们有一个运算 $\Omega A\times \Omega A\to \Omega A$。
更一般地，$A$ 的高阶群胚结构赋予了 $\Omega A$ 类似的"高阶群"结构。

考虑 $A$ 的环路空间\emph{的}环路空间\index{loop space!iterated}\index{iterated loop space}也是有用的，它是在 $a$ 处的恒等环路上的2维环路空间。
这记作 $\Omega^2(A,a)$，在类型论中用类型 $\id[({\id[A]{a}{a}})]{\refl{a}}{\refl{a}}$ 表示。
虽然 $\Omega^2(A,a)$ 作为一个环路空间，又是一个"高阶群"，但由于它的元素是1维环路之间的2维环路，它现在还具有一些额外的结构。

\begin{thm}[Eckmann--Hilton]\label{thm:EckmannHilton}
  第二环路空间上的复合运算
  %
  \begin{equation*}
    \Omega^2(A)\times \Omega^2(A)\to \Omega^2(A)
  \end{equation*}
  是交换的：对任意 $\alpha, \beta:\Omega^2(A)$，都有 $\alpha\ct\beta = \beta\ct\alpha$。
  \index{Eckmann--Hilton argument}%
\end{thm}

\begin{proof}
首先观察到，1-环路的复合 $\Omega A\times \Omega A\to \Omega A$ 诱导出一个运算
\[
\star : \Omega^2(A)\times \Omega^2(A)\to \Omega^2(A)
\]
具体如下：考虑元素 $a, b, c : A$ 以及1-道路和2-道路，
%
\begin{align*}
  p &: a = b,       &       r &: b = c \\
  q &: a = b,       &       s &: b = c \\
  \alpha &: p = q,  &   \beta &: r = s
\end{align*}
%
如下图所示（道路用箭头表示）。
% Changed this to xymatrix in the name of having uniform source code,
% maybe the original using xy looked better (I think it was too big).
% It is commented out below in case you want to reinstate it.
\[
 \xymatrix@+5em{
   {a} \rtwocell<10>^p_q{\alpha}
   &
   {b} \rtwocell<10>^r_s{\beta}
   &
   {c}
 }
\]
分别复合上方和下方的1-道路，我们得到两条道路 $p\ct r,\ q\ct s : a = c$，然后它们之间存在一个"水平复合"
%
\begin{equation*}
  \alpha\hct\beta : p\ct r = q\ct s
\end{equation*}
%
其定义如下。
首先，通过对 $r$ 作道路归纳来定义 $\alpha \rightwhisker r : p\ct r = q\ct r$，使得
\[ \alpha \rightwhisker \refl{b} \jdeq \opp{\mathsf{ru}_p} \ct \alpha \ct \mathsf{ru}_q \]
其中 $\mathsf{ru}_p : p = p \ct \refl{b}$ 是来自 \cref{thm:omg}\ref{item:omg1} 的右单位律。
我们也可以通过对 $\alpha$ 或所有相关道路作归纳来定义 $\rightwhisker$，这将导致不同的判断相等，但就当前目的而言，通过对 $r$ 归纳来定义会使事情更简单。
类似地，通过对 $q$ 作道路归纳来定义 $q\leftwhisker \beta : q\ct r = q\ct s$，使得
\[ \refl{b} \leftwhisker \beta \jdeq \opp{\mathsf{lu}_r} \ct \beta \ct \mathsf{lu}_s \]
其中 $\mathsf{lu}_r$ 表示左单位律。
运算 $\leftwhisker$ 和 $\rightwhisker$ 称为\define{须接（whiskering）}\indexdef{whiskering}。
接下来，由于 $\alpha \rightwhisker r$ 和 $q\leftwhisker \beta$ 是可复合的2-道路，我们可以定义\define{水平复合（horizontal composition）}
\indexdef{horizontal composition!of paths}%
\indexdef{composition!of paths!horizontal}%
为：
\[
\alpha\hct\beta\ \defeq\ (\alpha\rightwhisker r) \ct (q\leftwhisker \beta).
\]
现在假设 $a \jdeq  b \jdeq  c$，从而所有1-道路 $p$, $q$, $r$, $s$ 都是 $\Omega(A,a)$ 的元素，并进一步假定 $p\jdeq q \jdeq r \jdeq s\jdeq \refl{a}$，因此 $\alpha:\refl{a} = \refl{a}$ 和 $\beta:\refl{a} = \refl{a}$ 可以按两种顺序复合。
在这种情况下，我们有
\begin{align*}
  \alpha\hct\beta
  &\jdeq (\alpha\rightwhisker\refl{a}) \ct (\refl{a}\leftwhisker \beta)\\
  &= \opp{\mathsf{ru}_{\refl{a}}} \ct \alpha \ct \mathsf{ru}_{\refl{a}} \ct \opp{\mathsf{lu}_{\refl a}} \ct \beta \ct \mathsf{lu}_{\refl{a}}\\
  &\jdeq \opp{\refl{\refl{a}}} \ct \alpha \ct \refl{\refl{a}} \ct \opp{\refl{\refl a}} \ct \beta \ct \refl{\refl{a}}\\
  &= \alpha \ct \beta.
\end{align*}
（回忆一下，根据道路归纳的计算规则，有 $\mathsf{ru}_{\refl{a}} \jdeq \mathsf{lu}_{\refl{a}} \jdeq \refl{\refl{a}}$。）
另一方面，我们可以类似地通过
\[
\alpha\hct'\beta\ \defeq\ (p\leftwhisker \beta)\ct (\alpha\rightwhisker s)
\]
定义另一个水平复合，并且类似地得到
\[
\alpha\hct'\beta = \beta\ct\alpha.
\]
\index{interchange law}%
但是，一般来说，两种定义水平复合的方式是一致的，即 $\alpha\hct\beta = \alpha\hct'\beta$，这可以通过对 $\alpha$ 和 $\beta$ 进行归纳，然后再对剩下的两条1-道路进行归纳来证明，从而将一切归约为自反性。
因此我们有
\[\alpha \ct \beta = \alpha\hct\beta = \alpha\hct'\beta = \beta\ct\alpha.
\qedhere
\]
\end{proof}

上述事实被称为 \emph{Eckmann--Hilton 论证}，源自经典同伦论。
在下文的 \cref{cha:homotopy} 中，正是借助它来证明一个类型的高阶同伦群总是交换\index{group!abelian}群。
证明中定义的须接和水平复合操作也是类型的 $\infty$-群胚结构的基本组成部分。
这些操作各自满足相应的定律（在相差高阶同伦的意义下），例如
\[ \alpha \rightwhisker (p\ct q) = (\alpha \rightwhisker p) \rightwhisker q \]
等等。
从现在起，我们相信读者能够在需要时运用道路归纳来定义更多此类操作并验证其性质。

正如这个例子所揭示的，高阶道路类型的代数结构远比每一层上的群胚结构复杂得多：各层次之间相互影响，产生更多运算和定律，如同同伦论中迭代环路空间的研究那样。
事实上，与经典同伦论中的做法一样，我们可以给出下面的一般性定义：

\begin{defn} \label{def:pointedtype}
  一个\define{带点类型（pointed type）}
  \indexsee{pointed!type}{type, pointed}%
  \indexdef{type!pointed}%
  $(A,a)$ 由一个类型 $A:\type$ 连同其中的一个点 $a:A$ 所组成；这个点称为其\define{基点（basepoint）}。
  \indexdef{basepoint}%
  我们记 $\pointed{\type} \defeq \sm{A:\type} A$ 为宇宙 $\type$ 中所有带点类型构成的类型。
\end{defn}

\begin{defn} \label{def:loopspace}
  给定一个带点类型 $(A,a)$，定义其\define{环路空间（loop space）}
  \indexdef{loop space}%
  为如下带点类型：
  \[\Omega(A,a)\defeq ((\id[A]aa),\refl a).\]
  其中的元素称为 $a$ 处的一个\define{环路（loop）}\indexdef{loop}。
  对于 $n:\N$，带点类型 $(A,a)$ 的\define{$n$ 重迭代环路空间（$n$-fold iterated loop space）}
  $\Omega^{n}(A,a)$
  \indexdef{loop space!iterated}%
  \indexsee{loop space!n-fold@$n$-fold}{loop space, iterated}%
  递归地定义为：
  \begin{align*}
    \Omega^0(A,a)&\defeq(A,a)\\
    \Omega^{n+1}(A,a)&\defeq\Omega^n(\Omega(A,a)).
  \end{align*}
  其中的元素称为 $a$ 处的一个\define{$n$-环路（$n$-loop）}
  \indexdef{loop!n-@$n$-}%
  \indexsee{n-loop@$n$-loop}{loop, $n$-}%
  或一个\define{$n$ 维环路（$n$-dimensional loop）}
  \indexsee{loop!n-dimensional@$n$-dimensional}{loop, $n$-}%
  \indexsee{n-dimensional loop@$n$-dimensional loop}{loop, $n$-}%
  。
\end{defn}

关于迭代环路空间，我们将在 \cref{cha:hlevels,cha:hits,cha:homotopy} 中进一步讨论。
\index{.infinity-groupoid@$\infty$-groupoid!structure of a type|)}%
\index{type!identity|)}
\index{path|)}%

\section{函数是函子}
\label{sec:functors}

\index{function|(}%
\index{functoriality of functions in type theory@``functoriality'' of functions in type theory}%
现在，我们希望论证函数 $f:A\to B$ 在道路上具有函子性。
在传统的类型论中，这等价于说函数保持相等。
\index{continuity of functions in type theory@``continuity'' of functions in type theory}%
从拓扑的角度看，这对应于说每个函数都是"连续"的，即保持道路。

\begin{lem}\label{lem:map}
  设 $f:A\to B$ 是一个函数。
  那么对于任意 $x,y:A$，存在一个操作
  \begin{equation*}
    \apfunc f : (\id[A] x y) \to (\id[B] {f(x)} {f(y)}).
  \end{equation*}
  此外，对于每个 $x:A$，有 $\apfunc{f}(\refl{x})\jdeq \refl{f(x)}$。
  \indexdef{application!of function to a path}%
  \indexdef{path!application of a function to}%
  \indexdef{function!application to a path of}%
  \indexdef{action!of a function on a path}%
\end{lem}

记号 $\apfunc f$ 既可以理解为 $f$ 对一条道路的\underline{应}用（\underline{ap}plication），也可以理解为 $f$ 在道路上的\underline{作}用（\underline{a}ction on \underline{p}aths）。

\begin{proof}[第一种证法]
  令 $D:\prd{x,y:A} (x=y) \to \type$ 为由下式定义的类型族：
  \[D(x,y,p)\defeq (f(x)= f(y)).\]
  于是我们有
  \begin{equation*}
    d\defeq\lam{x} \refl{f(x)}:\prd{x:A} D(x,x,\refl{x}).
  \end{equation*}
  由道路归纳，得到 $\apfunc f : \prd{x,y:A} (x=y) \to (f(x)=f(y))$。
  计算规则蕴含对每个 $x:A$ 有 $\apfunc f({\refl{x}})\jdeq\refl{f(x)}$。
\end{proof}

\begin{proof}[第二种证法]
  为对所有 $p:x=y$ 定义 $\apfunc{f}(p)$，由归纳法，只需假设 $p$ 是 $\refl{x}$。
  在此情形下，可以定义 $\apfunc f(p) \defeq \refl{f(x)}:f(x)= f(x)$。
\end{proof}

我们通常将 $\apfunc f (p)$ 简写为 $\ap f p$。
严格说来，这种写法有歧义，但通常不会引起混淆。
这与范畴论中的常见约定一致：用同一个符号表示函子作用于对象和作用于态射。

我们注意到 $\apfunc{}$ 以人们可能期望的各种方式呈现函子行为。

\begin{lem}\label{lem:ap-functor}
  对于函数 $f:A\to B$ 和 $g:B\to C$，以及道路 $p:\id[A]xy$ 和 $q:\id[A]yz$，我们有：
  \begin{enumerate}
  \item $\apfunc f(p\ct q) = \apfunc f(p) \ct \apfunc f(q)$。\label{item:apfunctor-ct}
  \item $\apfunc f(\opp p) = \opp{\apfunc f (p)}$。\label{item:apfunctor-opp}
  \item $\apfunc g (\apfunc f(p)) = \apfunc{g\circ f} (p)$。\label{item:apfunctor-compose}
  \item $\apfunc {\idfunc[A]} (p) = p$。
  \end{enumerate}
\end{lem}

\begin{proof}
  留给读者。
\end{proof}


\index{function|)}%

正如 \cref{thm:omg} 中的那些相等关系一样，\cref{lem:ap-functor} 中的这些相等关系本身也是道路，它们满足各自的相干律（可以用同样的方式证明），如此等等。

\section{依值类型族是纤维化}
\label{sec:fibrations}

\index{type!family of|(}%
\index{transport|(defstyle}%
由于\emph{依值类型的（dependently typed）}函数在类型论中至关重要，我们也将需要 \cref{lem:map} 对这种函数的一个版本。
然而，这一点陈述起来并不那么简单，因为如果 $f:\prd{x:A} B(x)$ 且 $p:x=y$，那么 $f(x):B(x)$ 和 $f(y):B(y)$ 是属于不同类型的元素，因此，先验地说，我们甚至无法问它们是否相等。
缺少的关键是，$p$ 本身为我们提供了一种联系类型 $B(x)$ 和 $B(y)$ 的方式。

我们已经在 \autoref{sec:identity-types} 中见过这一点，当时我们称之为"等同物的不可分辨性（indiscernibility of identicals）"。
\index{indiscernibility of identicals}%
现在，我们为它引入一个新的名称和记号，此后将一直使用。

\begin{lem}[Transport]\label{lem:transport}
  假设 $P$ 是 $A$ 上的一个类型族，且 $p:\id[A]xy$。
  那么存在一个函数 $\transf{p}:P(x)\to P(y)$。
\end{lem}

\begin{proof}[第一种证法]
  令 $D:\prd{x,y:A} (\id{x}{y}) \to \type$ 为由下式定义的类型族：
  \[D(x,y,p)\defeq P(x)\to P(y).\]
  那么我们有函数
  \begin{equation*}
    d\defeq\lam{x} \idfunc[P(x)]:\prd{x:A} D(x,x,\refl{x}),
  \end{equation*}
  因此，归纳原理给出 $\indid{A}(D,d,x,y,p):P(x)\to P(y)$，我们将其定义为 $\transf p$。
\end{proof}

\begin{proof}[第二种证法]
  由归纳法，只需假设 $p$ 是 $\refl x$。
  但在此情形下，我们可以取 $\transf{(\refl x)}:P(x)\to P(x)$ 为恒等函数。
\end{proof}

有时，需要在记号中标明传输操作发生在哪个类型族 $P$ 中。
此时，我们可以写作
\[\transfib P p \blank : P(x) \to P(y).\]

回想一下，$A$ 上的类型族 $P$ 可以看作 $A$ 中元素的一种性质：如果 $P(x)$ 有实例，我们就说该性质在 $A$ 中的 $x$ 处成立。
那么传输引理就表明 $P$ 保持相等，其含义是：如果 $x$ 等于 $y$，那么 $P(x)$ 成立当且仅当 $P(y)$ 成立。
事实上，我们后面会看到，如果 $x=y$，那么 $P(x)$ 和 $P(y)$ 实际上是\emph{等价}的。

从拓扑的角度看，传输引理可以视为纤维化中的一种"道路提升（path lifting）"操作。
\index{fibration}%
\indexdef{total!space}%
我们可以把一个类型族 $P:A\to \type$ 看作一个以 $A$ 为\emph{底空间}的\emph{纤维化（fibration）}，其中 $P(x)$ 是 $x$ 上的纤维，而 $\sm{x:A}P(x)$ 是这个纤维化的\define{全空间（total space）}，其第一投影为 $\sm{x:A}(P(x))\to A$。
纤维化的定义性质是：给定底空间 $A$ 中的一条道路 $p:x=y$ 以及 $x$ 上纤维中的一点 $u:P(x)$，我们可以将道路 $p$ 提升为全空间中的一条道路，使其始于 $u$（并且这个提升可以连续地进行）。
点 $\trans p u$ 可以看作这条提升后的道路的另一端点。
我们也可以在类型论中定义这条道路本身：

\begin{lem}[道路提升性质]\label{thm:path-lifting}
  \indexdef{path!lifting}%
  \indexdef{lifting!path}%
  设 $P:A\to\type$ 是 $A$ 上的一个类型族，并假设对某个 $x:A$ 有 $u:P(x)$。
  那么对任意 $p:x=y$，在 $\sm{x:A}P(x)$ 中存在
  \begin{equation*}
    \mathsf{lift}(u,p):(x,u)=(y,\trans{p}{u})
  \end{equation*}
  使得 $\ap{\proj1}{\mathsf{lift}(u,p)} = p$。
\end{lem}

\begin{proof}
  留给读者。
  我们将在 \cref{sec:compute-sigma} 中证明一个更一般的定理。
\end{proof}

在经典同伦论中，纤维化被定义为其道路的提升\emph{存在}的映射；而与之相对，我们刚刚说明了，在类型论中，每一个类型族都自带一个\emph{指定的}"道路提升函数"。
这符合构造性数学的哲学：除非将某物展示出来，否则不能证明它存在。
\index{continuity of functions in type theory@``continuity'' of functions in type theory}%
这也自动保证了道路提升是被"连续地"选择的，因为正如我们所见，类型论中的所有函数都是"连续"的。

\begin{rmk}
  尽管我们可以将类型族 $P:A\to \type$ 看作纤维化，但一般而言最好不要说"纤维化 $P:A\to\type$"这样的话，因为这种说法听起来像是我们在讨论一个以 $\type$ 为底、以 $A$ 为全空间的纤维化。
  再次强调，当将类型族 $P:A\to \type$ 视为纤维化时，它的底是 $A$，全空间是 $\sm{x:A} P(x)$。

  在谈论类型族时，我们偶尔也会用到其他拓扑学术语。
  例如，我们可能将依值函数 $f:\prd{x:A} P(x)$ 称为纤维化 $P$ 的一个\define{截面（section）}
  \indexdef{section!of a type family}；%
  我们也可能说某事\define{逐纤维地}
  \indexdef{fiberwise}%
  成立，意思是它对每个 $P(x)$ 都成立。
  例如，一个截面 $f:\prd{x:A} P(x)$ 就表明 $P$ 是"逐纤维地有实例的"。
\end{rmk}

\index{function!dependent|(}
现在我们可以证明 \cref{lem:map} 的依值版本了。
拓扑直觉是：给定 $f:\prd{x:A} P(x)$ 和一条道路 $p:\id[A]xy$，我们应该能够将 $f$ 应用于 $p$，从而在 $P$ 的全空间中得到一条"位于 $p$ 上方"的道路，如下图所示。

\begin{center}
  \begin{tikzpicture}[yscale=.5,xscale=2]
    \draw (0,0) arc (-90:170:8ex) node[anchor=south east] {$A$} arc (170:270:8ex);
    \draw (0,6) arc (-90:170:8ex) node[anchor=south east] {$\sm{x:A} P(x)$} arc (170:270:8ex);
    \draw[->] (0,5.8) -- node[auto] {$\proj1$} (0,3.2);
    \node[circle,fill,inner sep=1pt,label=left:{$x$}] (b1) at (-.5,1.4) {};
    \node[circle,fill,inner sep=1pt,label=right:{$y$}] (b2) at (.5,1.4) {};
    \draw[decorate,decoration={snake,amplitude=1}] (b1) -- node[auto,swap] {$p$} (b2);
    \node[circle,fill,inner sep=1pt,label=left:{$f(x)$}] (b1) at (-.5,7.2) {};
    \node[circle,fill,inner sep=1pt,label=right:{$f(y)$}] (b2) at (.5,7.2) {};
    \draw[decorate,decoration={snake,amplitude=1}] (b1) -- node[auto] {$f(p)$} (b2);
  \end{tikzpicture}
\end{center}

我们\emph{确实可以}从\cref{lem:map}得到这样的对象。
给定$f:\prd{x:A} P(x)$，我们可以令$f'(x)\defeq (x,f(x))$来定义一个非依值函数$f':A\to \sm{x:A} P(x)$，进而考虑$\ap{f'}{p} : f'(x) = f'(y)$。
由于 $\proj1 \circ f' \jdeq \idfunc[A]$，由\cref{lem:ap-functor}可得$\ap{\proj1}{\ap{f'}{p}} = p$；因此 $\ap{f'}{p}$ 确实在此意义上"位于" $p$ 的上方。
然而，仅凭 $\ap{f'}{p}$ 的\emph{类型}，并不容易看出它究竟位于 $A$ 中哪条具体道路（本例中即 $p$）的上方，而这一点有时很重要。

解决方法是使用迁移引理。
根据\cref{thm:path-lifting}，我们有一条从 $(x,u)$ 到 $(y,\trans p u)$ 的典范道路 $\mathsf{lift}(u,p)$，它位于 $p$ 的上方。
因此，任何从 $u:P(x)$ 到 $v:P(y)$ 且位于 $p$ 上方的道路，都应本质唯一地通过 $\mathsf{lift}(u,p)$ 进行分解，其分解因子是一条完全位于纤维 $P(y)$ 中、从 $\trans p u$ 到 $v$ 的道路。
因此，在相差等价的意义下，将"从 $u$ 到 $v$ 且位于 $p:x=y$ 上方的道路"定义为 $P(y)$ 中的道路 $\trans p u = v$ 是合理的。
而且，我们确实可以证明依值函数会产生这样的道路。

\begin{lem}[依值映射]\label{lem:mapdep}
  \indexdef{application!of dependent function to a path}%
  \indexdef{path!application of a dependent function to}%
  \indexdef{function!dependent!application to a path of}%
  \indexdef{action!of a dependent function on a path}%
  假设$f:\prd{x: A} P(x)$；那么我们有一个映射
  \[\apdfunc f : \prd{p:x=y}\big(\id[P(y)]{\trans p{f(x)}}{f(y)}\big).\]
\end{lem}

\begin{proof}[第一种证明]
  令$D:\prd{x,y:A} (\id{x}{y}) \to \type$为由下式定义的类型族：
  \begin{equation*}
    D(x,y,p)\defeq \trans p {f(x)}= f(y).
  \end{equation*}
  那么$D(x,x,\refl{x})$就是$\trans{(\refl{x})}{f(x)}= f(x)$。
  但由于$\trans{(\refl{x})}{f(x)}\jdeq f(x)$，我们得到$D(x,x,\refl{x})\jdeq (f(x)= f(x))$。
  因此，我们得到函数
  \begin{equation*}
    d\defeq\lam{x} \refl{f(x)}:\prd{x:A} D(x,x,\refl{x})
  \end{equation*}
  现在，道路归纳便对每个 $p:x=y$ 给出$\apdfunc f(p):\trans p{f(x)}= f(y)$。
\end{proof}

\begin{proof}[第二种证明]
  通过归纳，只需假设 $p$ 是 $\refl x$。
  但此时所求等式为$\trans{(\refl{x})}{f(x)}= f(x)$，判断成立。
\end{proof}

我们将把这种意义上"位于其他道路上方的道路"统称为\emph{依值道路（dependent paths）}。
\indexsee{dependent!path}{path, dependent}%
\index{path!dependent}%
从\cref{cha:hits}开始，它们将扮演越来越重要的角色。
在\cref{sec:computational}中我们会看到，对于少数几种特定的类型族，存在一些等价的、有时更方便的表示依值道路概念的方式。

现在回忆\cref{sec:pi-types}，非依值类型的函数 $f:A\to B$ 只是依值类型函数$f:\prd{x:A} P(x)$在 $P$ 为常值类型族（即 $P(x) \defeq B$）时的特例。
此时，$\apdfunc{f}$ 与 $\apfunc{f}$ 密切相关，这由以下引理保证：

\begin{lem}\label{thm:trans-trivial}
  如果 $P:A\to\type$ 对某个固定的 $B:\type$ 定义为 $P(x) \defeq B$，那么对于任意 $x,y:A$ 和 $p:x=y$ 以及 $b:B$，我们有一条道路
  \[ \transconst Bpb : \transfib P p b = b. \]
\end{lem}

\begin{proof}[第一种证明]
  固定一个 $b:B$，并令$D:\prd{x,y:A} (\id{x}{y}) \to \type$为由下式定义的类型族：
  \[ D(x,y,p) \defeq (\transfib P p b = b). \]
  则 $D(x,x,\refl x)$ 即$(\transfib P{\refl{x}}{b} = b)$，根据迁移的计算规则，它与 $(b=b)$ 判断相等。
  因此，我们有函数
  \[ d \defeq \lam{x} \refl{b} : \prd{x:A} D(x,x,\refl x). \]
  现在，道路归纳便给出
  \narrowequation{
    \prd{x,y:A}{p:x=y}(\transfib P p b = b)}
  类型的一个元素，如我们所愿。
\end{proof}

\begin{proof}[第二种证明]
  通过归纳，只需假设 $y$ 是 $x$ 且 $p$ 是 $\refl x$。
  但 $\transfib P {\refl x} b \jdeq b$，故此时所需证明的不过是 $b=b$，而 $\refl{b}$ 即可。
\end{proof}

因此，对于任意 $x,y:A$、$p:x=y$ 和 $f:A\to B$，通过分别与 $\transconst Bp{f(x)}$ 及其逆进行拼接，我们得到函数
\begin{align}
  \big(f(x) = f(y)\big) &\to \big(\trans{p}{f(x)} = f(y)\big)\label{eq:ap-to-apd}
  \qquad\text{and} \\
  \big(\trans{p}{f(x)} = f(y)\big) &\to \big(f(x) = f(y)\big).\label{eq:apd-to-ap}
\end{align}
实际上，这些函数是互逆的等价（在\cref{sec:basics-equivalences}将要引入的意义上），并且它们将 $\apfunc f (p)$ 与 $\apdfunc f (p)$ 联系起来。

\begin{lem}\label{thm:apd-const}
  对于 $f:A\to B$ 和 $p:\id[A]xy$，我们有
  \[ \apdfunc f(p) = \transconst B p{f(x)} \ct \apfunc f (p). \]
\end{lem}

\begin{proof}[第一种证明]
  令$D:\prd{x,y:A} (\id xy) \to \type$为由下式定义的类型族：
  \[ D(x,y,p) \defeq \big(\apdfunc f (p) = \transconst Bp{f(x)} \ct \apfunc f (p)\big). \]
  因此，我们有
  \[D(x,x,\refl x) \jdeq \big(\apdfunc f (\refl x) = \transconst B{\refl x}{f(x)} \ct \apfunc f ({\refl x})\big).\]
  但根据定义，出现在这个类型中的所有三条道路都是 $\refl{f(x)}$，所以我们有
  \[ \refl{\refl{f(x)}} : D(x,x,\refl x). \]
  因此，道路归纳便给出$\prd{x,y:A}{p:x=y} D(x,y,p)$类型的一个元素，即为所求。
\end{proof}

\begin{proof}[第二种证明]
  通过归纳，只需假设 $y$ 是 $x$ 且 $p$ 是 $\refl x$。
  此时需要证明的是$\refl{f(x)} = \refl{f(x)} \ct \refl{f(x)}$，判断成立。
\end{proof}

由于 $\apdfunc{f}$ 与 $\apfunc{f}$ 的类型不同，使用不同的记号通常会更清晰。
% We may sometimes use a notation $\apd f p$ for $\apdfunc{f}(p)$, which is similar to the notation $\ap f p$ for $\apfunc{f}(p)$.

\index{function!dependent|)}%

至此，我们希望读者已开始对基于相等类型的归纳证明有所体会。
从今往后，我们将不再同时给出两种风格的证明，而是选用最清晰、最便捷的（通常是第二种更简洁的）证明方式。
以下是关于迁移的其他几个有用引理；我们将其证明留给读者（选用任一风格均可）。

\begin{lem}\label{thm:transport-concat}
  给定 $P:A\to\type$，其中 $p:\id[A]xy$ 且 $q:\id[A]yz$，以及 $u:P(x)$，我们有
  \[ \trans{q}{\trans{p}{u}} = \trans{(p\ct q)}{u}. \]
\end{lem}

\begin{lem}\label{thm:transport-compose}
  对于函数 $f:A\to B$ 和类型族 $P:B\to\type$，以及任意的 $p:\id[A]xy$ 和 $u:P(f(x))$，我们有
  \[ \transfib{P\circ f}{p}{u} = \transfib{P}{\apfunc f(p)}{u}. \]
\end{lem}

\begin{lem}\label{thm:ap-transport}
  对于 $P,Q:A\to \type$ 和一族函数 $f:\prd{x:A} P(x)\to Q(x)$，以及任意的 $p:\id[A]xy$ 和 $u:P(x)$，我们有
  \[ \transfib{Q}{p}{f_x(u)} = f_y(\transfib{P}{p}{u}). \]
\end{lem}

\index{type!family of|)}%
\index{transport|)}

\section{同伦与等价}
\label{sec:basics-equivalences}

\index{homotopy|(defstyle}%

至此，我们已经了解了如何将相等类型 $\id[A]xy$ 视为类型 $A$ 中两个元素 $x$ 与 $y$ 之间的\emph{相等证明}、\emph{道路}或\emph{等价}。
现在，我们来探究\emph{函数}之间以及\emph{类型}之间合适的"等同"或"相同"概念。
在\cref{sec:compute-pi,sec:compute-universe}中，我们将看到同伦类型论允许我们将这些概念与相等类型的实例等同起来，但在此之前，我们需要先理解它们本身。

传统上，我们认为两个函数相同，如果它们在所有输入上取相等的值。
在"命题即类型"的解释下，这提示我们，两个（可能是依值的）函数 $f$ 和 $g$ 相同，当类型 $\prd{x:A} (f(x)=g(x))$ 是有实例的。
在同伦解释下，这个依值函数类型由\emph{连续}的道路或\emph{函子性}的等价构成，因此可视为\emph{同伦}或\emph{自然同构}的类型。\index{isomorphism!natural}
我们将采用拓扑学的术语。

\begin{defn} \label{defn:homotopy}
  设 $f,g:\prd{x:A} P(x)$ 是类型族 $P:A\to\type$ 的两个截面。
  从 $f$ 到 $g$ 的一个\define{同伦}（homotopy）
  是一个类型为
  \begin{equation*}
    (f\htpy g) \defeq \prd{x:A} (f(x)=g(x)).
  \end{equation*}
  的依值函数。
\end{defn}

注意，一个同伦并不等同于 $(f=g)$ 的一个相等证明。
不过，在\cref{sec:compute-pi}中，我们将引入一条公理，使得同伦与相等证明变得"等价"。

以下证明留给读者。

\begin{lem}\label{lem:homotopy-props}
  同伦关系在每个依值函数类型 $\prd{x:A} P(x)$ 上都是一个等价关系。
  即我们有以下类型的元素：
  \begin{gather*}
    \prd{f:\prd{x:A} P(x)} (f\htpy f)\\
    \prd{f,g:\prd{x:A} P(x)} (f\htpy g) \to (g\htpy f)\\
    \prd{f,g,h:\prd{x:A} P(x)} (f\htpy g) \to (g\htpy h) \to (f\htpy h).
  \end{gather*}
\end{lem}

% This is judgmental and is \cref{ex:composition}.
% \begin{lem}
%   Composition is associative and unital up to homotopy.
%   That is:
%   \begin{enumerate}
%   \item If $f:A\to B$ then $f\circ \idfunc[A]\htpy f\htpy \idfunc[B]\circ f$.
%   \item If $f:A\to B, g:B\to C$ and $h:C\to D$ then $h\circ (g\circ f) \htpy (h\circ g)\circ f$.
%   \end{enumerate}
% \end{lem}

\index{functoriality of functions in type theory@``functoriality'' of functions in type theory}%
\index{continuity of functions in type theory@``continuity'' of functions in type theory}%
正如类型论中的函数自动是"函子"一样，同伦也自动是\index{naturality of homotopies@``naturality'' of homotopies}%"自然变换"。
我们仅对非依值函数 $f,g:A\to B$ 陈述并证明这一点；在\cref{ex:dep-htpy-natural}中，请读者将其推广到依值函数的情形。

回忆一下，对于 $f:A\to B$ 和 $p:\id[A]xy$，我们可以记 $\ap f p$ 来表示 $\apfunc{f} (p)$。

\begin{lem}\label{lem:htpy-natural}
  假设 $H:f\htpy g$ 是函数 $f,g:A\to B$ 之间的一个同伦，并令 $p:\id[A]xy$。那么我们有
  \begin{equation*}
    H(x)\ct\ap{g}{p}=\ap{f}{p}\ct H(y).
  \end{equation*}
  我们也可以将其画成一个交换图：\index{diagram}
  \begin{align*}
    \xymatrix{
      f(x) \ar@{=}[r]^{\ap fp} \ar@{=}[d]_{H(x)} & f(y) \ar@{=}[d]^{H(y)} \\
      g(x) \ar@{=}[r]_{\ap gp} & g(y)
    }
  \end{align*}
\end{lem}

\begin{proof}
  通过归纳，只需假设 $p$ 是 $\refl x$。
  由于 $\apfunc{f}$ 和 $\apfunc g$ 在自反性上的计算，此时所需证明的是
  \[ H(x) \ct \refl{g(x)} = \refl{f(x)} \ct H(x). \]
  但由于等式两边都等于 $H(x)$，结论成立。
\end{proof}

\begin{cor}\label{cor:hom-fg}
  设 $H : f \htpy \idfunc[A]$ 是一个同伦，其中 $f : A \to A$。那么对于任意 $x : A$，我们有 \[ H(f(x)) = \ap f{H(x)}. \]
  % The above path will be denoted by $\com{H}{f}{x}$.
\end{cor}


\noindent
此处 $f(x)$ 表示 $f$ 对 $x$ 的普通函数应用，而 $\ap f{H(x)}$ 表示 $\apfunc{f}(H(x))$。


\begin{proof}
由 $H$ 的自然性，下面的道路图交换：
\begin{align*}
\xymatrix@C=3pc{
ffx \ar@{=}[r]^-{\ap f{Hx}} \ar@{=}[d]_{H(fx)} & fx \ar@{=}[d]^{Hx} \\
fx \ar@{=}[r]_-{Hx} & x
}
\end{align*}
也就是说，$\ap f{H x} \ct H x = H(f x) \ct H x$。
现在我们可以通过须接 $\opp{(H x)}$ 来消去 $H x$，从而得到
\[ \ap f{H x}
= \ap f{H x} \ct H x \ct \opp{(H x)}
= H(f x) \ct H x \ct \opp{(H x)}
= H(f x)
\]
如我们所愿（此处省略了一些结合律道路）。
\end{proof}

当然，如同函数的函子性（\cref{lem:ap-functor}），\cref{lem:htpy-natural}中的等式本身也是一条满足其自身相干律的道路，依此类推。

\index{homotopy|)}%

\index{equivalence|(}%
现在转向类型。从传统观点来看，我们可以说一个函数 $f:A\to B$ 是一个\emph{同构（isomorphism）}，如果存在一个函数 $g:B\to A$，使得两个复合 $f\circ g$ 和 $g\circ f$ 都点态相等于恒等函数；也就是说，满足 $f \circ g \htpy \idfunc[B]$ 和 $g\circ f \htpy \idfunc[A]$。
\indexsee{homotopy!equivalence}{equivalence}%
从同伦论的角度看，这应被称为\emph{同伦等价（homotopy equivalence）}；而从范畴论的角度看，则应称为\emph{（高阶）群胚的等价（equivalence of (higher) groupoids）}。
然而，在进行证明相关的数学时，
\index{mathematics!proof-relevant}%
其对应的类型
\begin{equation}
  \sm{g:B\to A} \big((f \circ g \htpy \idfunc[B]) \times (g\circ f \htpy \idfunc[A])\big)\label{eq:qinvtype}
\end{equation}
表现不佳。
例如，对于单个函数 $f:A\to B$，\eqref{eq:qinvtype} 中可能存在多个彼此不相等的居民。
（这与高阶范畴论中的一个观察密切相关：我们常常需要考虑\emph{伴随等价（adjoint equivalence）}\index{adjoint!equivalence}，而不仅仅是普通的等价。）
为此，我们转而赋予~\eqref{eq:qinvtype} 下面这个历史上准确但听起来略带贬义色彩的名字。

\begin{defn}\label{defn:quasi-inverse}
  对于函数 $f:A\to B$，$f$ 的一个\define{拟逆（quasi-inverse）}
  \indexdef{quasi-inverse}%
  \indexsee{function!quasi-inverse of}{quasi-inverse}%
  是一个三元组 $(g,\alpha,\beta)$，其中包含一个函数 $g:B\to A$ 以及两个同伦
$\alpha:f\circ g\htpy \idfunc[B]$ 和 $\beta:g\circ f\htpy \idfunc[A]$。
\end{defn}

\symlabel{qinv}
因此，~\eqref{eq:qinvtype} 就是\emph{$f$ 的拟逆构成的类型}；我们可以将其记作 $\qinv(f)$。

\begin{eg}\label{eg:idequiv}
  \index{identity!function}%
  \index{function!identity}%
  恒等函数 $\idfunc[A]:A\to A$ 有一个拟逆，就是 $\idfunc[A]$ 本身，连同由 $\alpha(y) \defeq \refl{y}$ 和 $\beta(x) \defeq \refl{x}$ 定义的同伦。
\end{eg}

\begin{eg}\label{eg:concatequiv}
  对于任意 $p:\id[A]xy$ 和 $z:A$，函数
  \begin{align*}
    (p\ct \blank)&:(\id[A]yz) \to (\id[A]xz) \qquad\text{and}\\
    (\blank \ct p)&:(\id[A]zx) \to (\id[A]zy)
  \end{align*}
  分别以 $(\opp p \ct \blank)$ 和 $(\blank \ct \opp p)$ 为拟逆；参见\cref{ex:equiv-concat}。
\end{eg}

\begin{eg}\label{thm:transportequiv}
  对于任意 $p:\id[A]xy$ 和 $P:A\to\type$，函数
  \[\transfib{P}{p}{\blank}:P(x) \to P(y)\]
  有一个拟逆 $\transfib{P}{\opp p}{\blank}$；这可由 \narrowbreak \cref{thm:transport-concat} 得出。
\end{eg}

\symlabel{basics-isequiv}\symlabel{basics:iso}
一般来说，我们仅在类型 $A$ 和 $B$ "表现得像集合"时（参见\cref{sec:basics-sets}），才使用\emph{同构（isomorphism）}一词
\index{isomorphism!of sets}
（以及类似词语，如\emph{双射（bijection）}，及相关记号 $A\cong B$）
\index{bijection}
。在这种情况下，类型~\eqref{eq:qinvtype} 不存在问题。
我们将\emph{等价（equivalence）}一词留给一个改进的概念 $\isequiv (f)$，它需满足以下性质：%


\begin{enumerate}
\item 对每个 $f:A\to B$，存在一个函数 $\qinv(f) \to \isequiv (f)$。\label{item:be1}
\item 类似地，对每个 $f$，我们有 $\isequiv (f) \to \qinv(f)$；因此二者是逻辑等价的（参见\cref{sec:pat}）。\label{item:be2}
\item 对于任意两个居民 $e_1,e_2:\isequiv(f)$，有 $e_1=e_2$。\label{item:be3}
\end{enumerate}


在\cref{cha:equivalences}中我们将会看到，有许多不同的 $\isequiv(f)$ 定义都满足这三条性质，但它们彼此都是等价的。
现在，为了让读者相信这样的定义是存在的，我们仅提及其中最简单的一个：
\begin{equation}\label{eq:isequiv-invertible}
  \isequiv(f) \;\defeq\;
  \Parens{\sm{g:B\to A} (f\circ g \htpy \idfunc[B])}
  \times
  \Parens{\sm{h:B\to A} (h\circ f \htpy \idfunc[A])}.
\end{equation}
我们现在就可以证明这个定义满足性质~\ref{item:be1}和~\ref{item:be2}。
要定义一个函数 $\qinv(f) \to \isequiv (f)$ 很容易：将 $(g,\alpha,\beta)$ 映到 $(g,\alpha,g,\beta)$ 即可。
反之，给定 $(g,\alpha,h,\beta)$，令 $\gamma$ 为如下的复合同伦
\[ g \overset{\beta}{\htpy} h\circ f\circ g \overset{\alpha}{\htpy} h, \]
即 $\gamma(x) \defeq \opp{\beta(g(x))} \ct \ap{h}{\alpha(x)}$。
现在，由 $\beta'(x) \defeq \gamma(f(x)) \ct \beta(x)$ 定义 $\beta':g\circ f\htpy \idfunc[A]$。
则有 $(g,\alpha,\beta'):\qinv(f)$。

对此定义证明性质~\ref{item:be3} 也不难，但需要先确定 Cartesian 积和依值对类型的相等类型，我们将在\cref{sec:compute-cartprod,sec:compute-sigma}中讨论这一点。
因此，我们也将其推迟；详见\cref{sec:biinv}。
此刻，主要应理解的是：存在一个表现良好的类型，我们可以将其读作"$f$ 是一个等价"；并且，我们可以通过给出 $f$ 的一个拟逆来证明 $f$ 是一个等价。
在实践中，这是证明一个函数是等价的最常用方法。

依照证明相关的理念，
\index{mathematics!proof-relevant}%
从 $A$ 到 $B$ 的一个\emph{等价（equivalence）}被定义为一个函数 $f:A\to B$ 连同 $\isequiv (f)$ 的一个居民，即它为等价的证明。
我们记 $(\eqv A B)$ 为从 $A$ 到 $B$ 的等价构成的类型，即类型
\begin{equation}\label{eq:eqv}
  (\eqv A B) \defeq \sm{f:A\to B} \isequiv(f).
\end{equation}
上述性质~\ref{item:be3} 保证了，如果两个等价作为函数是相等的（即它们作为 $A\to B$ 中的底层元素是相等的），那么它们作为等价也是相等的（参见 \cref{sec:compute-sigma}）。
因此，我们常常滥用记号，模糊等价与其底层函数之间的区别。
例如，如果我们有一个函数 $f:A\to B$，并且知道 $e:\isequiv(f)$，我们可以写成 $f:\eqv A B$，而不是 $\tup{f}{e}$。
反之，如果我们有一个等价 $g:\eqv A B$，那么给定 $a:A$，我们可以直接写 $g(a)$，而不是 $(\proj1 g)(a)$。

最后，我们指出如下观察：

\begin{lem}\label{thm:equiv-eqrel}
  类型等价是 \type 上的一个等价关系。
  具体来说：
  \begin{enumerate}
  \item 对任意 $A$，恒等函数 $\idfunc[A]$ 是一个等价；因此 $\eqv A A$。
  \item 对任意 $f:\eqv A B$，我们有一个等价 $f^{-1} : \eqv B A$。
  \item 对任意 $f:\eqv A B$ 和 $g:\eqv B C$，我们有 $g\circ f : \eqv A C$。
  \end{enumerate}
\end{lem}

\begin{proof}
  恒等函数显然是自身的拟逆；因此它是一个等价。

  如果 $f:A\to B$ 是一个等价，那么它有一个拟逆，设其为 $f^{-1}:B\to A$。
  则 $f$ 也是 $f^{-1}$ 的一个拟逆，所以 $f^{-1}$ 是一个 $B\to A$ 的等价。

  最后，给定 $f:\eqv A B$ 和 $g:\eqv B C$，设其拟逆分别为 $f^{-1}$ 和 $g^{-1}$，则对于任意 $a:A$，我们有 $f^{-1} g^{-1} g f a = f^{-1} f a = a$，而对于任意 $c:C$，我们有 $g f f^{-1} g^{-1} c = g g^{-1} c = c$。
  因此 $f^{-1} \circ g^{-1}$ 是 $g\circ f$ 的一个拟逆，从而 $g\circ f$ 是一个等价。
\end{proof}

\index{equivalence|)}%

\section{类型形成子的高阶群胚结构}
\label{sec:computational}

在 \cref{cha:typetheory} 中，我们介绍了多种形成新类型的方法：Cartesian 积、不交并、依值积、依值和等等。
在 \cref{sec:equality,sec:functors,sec:fibrations} 中，我们看到同伦类型论中的\emph{所有}类型都表现得像空间或高阶群胚。
本章剩余部分的目标是，在 \cref{cha:typetheory} 所定义的特定类型的案例中，明确这种高阶结构是如何表现的。

事实证明，对于许多类型 $A$，其相等类型 $\id[A]xy$ 可以用构造 $A$ 时所用的数据，在等价的意义下刻画出来。
例如，如果 $A$ 是一个 Cartesian 积 $B\times C$，并且 $x\jdeq (b,c)$，$y\jdeq(b',c')$，那么我们有一个等价
\begin{equation}\label{eq:prodeqv}
  \eqv{\big((b,c)=(b',c')\big)}{\big((b=b')\times (c=c')\big)}.
\end{equation}
用更传统的语言来说，两个有序对相等当且仅当它们的各分量分别相等（但等价~\eqref{eq:prodeqv} 所说的比这更多）。
相等类型的高阶结构也可以用这些等价来表达；例如，拼接两个关于有序对的相等性，对应于各分量上的逐对拼接。

类似地，当一个类型族 $P:A\to\type$ 是使用 \cref{cha:typetheory} 中的类型形成规则逐纤维地构建起来时，操作 $\transfib{P}{p}{\blank}$ 可以在同伦的意义下，根据构建 $P$ 所用数据的相应操作来刻画。
例如，如果 $P(x) \jdeq B(x)\times C(x)$，那么我们有
\[\transfib{P}{p}{(b,c)} = \big(\transfib{B}{p}{b},\transfib{C}{p}{c}\big).\]

最后，类型形成规则也具有函子性，并且如果一个函数 $f$ 是基于这种函子性构建的，那么 $\apfunc f$ 和 $\apdfunc f$ 这两个操作，可以根据构造 $f$ 所用数据的相应操作来计算。
例如，如果 $g:B\to B'$ 且 $h:C\to C'$，并且我们通过 $f(b,c)\defeq (g(b),h(c))$ 来定义 $f:B\times C \to B'\times C'$，那么模掉等价~\eqref{eq:prodeqv}，我们可以将 $\apfunc f$ 与 "$(\apfunc g,\apfunc h)$" 等同。

接下来的几节（\crefrange{sec:compute-cartprod}{sec:compute-nat}）将致力于为所有基本的类型形成规则陈述并证明这类定理，每个基本类型形成子对应一节。
在这里，我们遇到了当前可用的类型理论中一个表面上的不足之处；
正如在后续章节中将会明确的那样，如果这些关于相等类型、传输等等的刻画是\emph{判断相等性}\index{judgmental equality}，似乎会更方便且更直观。
然而，在 \cref{cha:typetheory} 所介绍的理论中，相等类型由其归纳原理为所有类型统一地定义，因此我们不能对不同的类型将它们"重新定义"为不同的东西。
因此，本章将要讨论的针对特定类型的刻画，大多是\emph{定理}，需要我们去发现和证明——如果可能的话。

实际上，\cref{cha:typetheory} 所描述的类型论不足以证明两类类型形成子——$\Pi$-类型与宇宙——所期望的那些定理。
正因如此，我们不得不向类型论中引入公理，以使这些"定理"成立。
从类型论的角度看，\emph{公理}（axiom，参见~\cref{sec:axioms}）是一个被声明居住于某个特定类型的"原子"元素；除了其所属类型的规则之外，没有任何其他规则支配其行为。
\index{axiom!versus rules}%

\index{function extensionality}%
\indexsee{extensionality, of functions}{function extensionality}
\index{univalence axiom}%
关于 $\Pi$-类型（\cref{sec:compute-pi}）的公理对于类型论研究者来说并不陌生：它被称为\emph{函数外延性}，其内容（粗略地说）是：如果两个函数在\cref{sec:basics-equivalences}的意义下是同伦的，那么它们相等。
然而，关于宇宙（\cref{sec:compute-universe}）的公理，则是同伦类型论中由 Voevodsky 提出的新贡献：它被称为\emph{泛等公理}，其内容（粗略地说）是：如果两个类型在\cref{sec:basics-equivalences}的意义下是等价的，那么它们相等。
我们已在导论中提到过这条公理；它将在本书中发挥重要作用。%
\footnote{我们选择将这些原理作为公理引入，但也存在其他可能的方式来构造一套使之成立的类型论。
  详见本章注记。}

需要特别注意的是，并非\emph{所有}的相等类型都能通过对类型构造进行归纳来"确定"。
反例包括大多数非平凡的高阶归纳类型（参见\cref{cha:hits,cha:homotopy}）。
例如，计算类型 $\Sn^n$（参见\cref{sec:circle}）的相等类型，就相当于计算球面的高阶同伦群——这是代数拓扑学中一个深刻而重要的研究领域。

\section{Cartesian 积类型}
\label{sec:compute-cartprod}

\index{type!product|(}%
给定类型 $A$ 和 $B$，考虑它们的 Cartesian 积类型 $A \times B$。
对于任意元素 $x,y:A\times B$ 以及一条道路$p:\id[A\times B]{x}{y}$，由函子性我们可以提取出道路$\ap{\proj1}p:\id[A]{\proj1(x)}{\proj1(y)}$与$\ap{\proj2}p:\id[B]{\proj2(x)}{\proj2(y)}$。
于是，我们得到一个函数
\begin{equation}\label{eq:path-prod}
  (\id[A\times B]{x}{y}) \to (\id[A]{\proj1(x)}{\proj1(y)}) \times (\id[B]{\proj2(x)}{\proj2(y)}).
\end{equation}

\begin{thm}\label{thm:path-prod}
  对于任意的 $x$ 和 $y$，函数~\eqref{eq:path-prod}是一个等价。
\end{thm}

从逻辑的角度看，这意味着两个序对相等，当且仅当它们逐个分量相等。
从范畴论的角度看，这意味着积群胚中的态射就是态射对。
从同伦论的角度看，这意味着积空间中的道路就是道路对。

\begin{proof}
  我们需要构造一个方向相反的函数：
  \begin{equation}
    (\id[A]{\proj1(x)}{\proj1(y)}) \times (\id[B]{\proj2(x)}{\proj2(y)}) \to (\id[A\times B]{x}{y}). \label{eq:path-prod-inverse}
  \end{equation}
  根据 Cartesian 积的归纳规则，我们可以假设 $x$ 和 $y$ 都是序对，即对某个 $a,a':A$ 和 $b,b':B$，有 $x\jdeq (a,b)$ 且 $y\jdeq (a',b')$。
  此时，我们需要的是一个形如
  \begin{equation*}
    (\id[A]{a}{a'}) \times (\id[B]{b}{b'}) \to \big(\id[A\times B]{(a,b)}{(a',b')}\big).
  \end{equation*}
  的函数。
  现在，通过对其定义域中的 Cartesian 积进行归纳，我们可以假设给定了 $p:a=a'$ 和 $q:b=b'$。
  接着，通过两次道路归纳，我们可以假设 $a\jdeq a'$ 且 $b\jdeq b'$，并且 $p$ 和 $q$ 都是自反性。
  但在这种情况下，我们有$(a,b)\jdeq(a',b')$，因此我们可以将输出也取为自反性。

  接下来需要证明~\eqref{eq:path-prod-inverse}是~\eqref{eq:path-prod}的拟逆。
  这是一系列简单的归纳，但必须按正确的顺序进行。

  在一个方向上，我们从$r:\id[A\times B]{x}{y}$出发。
  我们首先对 $r$ 作道路归纳，以便假设 $x\jdeq y$ 且 $r$ 是自反性。
  此时，由于 $\apfunc{\proj1}$ 和 $\apfunc{\proj2}$ 是通过道路归纳定义的，~\eqref{eq:path-prod}将 $r\jdeq \refl{x}$ 映射到序对$(\refl{\proj1x},\refl{\proj2x})$。
  现在，通过对 $x$ 进行归纳，可假设 $x\jdeq (a,b)$，于是这个序对就是 $(\refl a, \refl b)$。
  因此，根据定义，~\eqref{eq:path-prod-inverse}将其映射到 $\refl{(a,b)}$，这（在我们当前的假设下）就是 $r$。

  在另一个方向上，如果我们从$s:(\id[A]{\proj1(x)}{\proj1(y)}) \times (\id[B]{\proj2(x)}{\proj2(y)})$出发，那么首先对 $x$ 和 $y$ 进行归纳，假设它们是序对 $(a,b)$ 和 $(a',b')$；然后对$s:(\id[A]{a}{a'}) \times (\id[B]{b}{b'})$进行归纳，将其归约为序对 $(p,q)$，其中 $p:a=a'$ 且 $q:b=b'$。
  接着，对 $p$ 和 $q$ 进行归纳，可以假设它们分别是自反性 $\refl a$ 和 $\refl b$；在这种情况下，~\eqref{eq:path-prod-inverse}给出 $\refl{(a,b)}$，然后~\eqref{eq:path-prod}再给出$(\refl a,\refl b)\jdeq (p,q)\jdeq s$。
\end{proof}

特别地，我们证明了~\eqref{eq:path-prod}有一个逆~\eqref{eq:path-prod-inverse}，我们可以将其记作
\symlabel{defn:pairpath}
\[
\pairpath : (\id{\proj{1}(x)}{\proj{1}(y)}) \times (\id{\proj{2}(x)}{\proj{2}(y)}) \to (\id x y).
\]
注意，它的一个特例给出了积类型的命题唯一性原理\index{uniqueness!principle, propositional!for product types}：$z = (\proj1(z),\proj2(z))$。

不妨把 \pairpath 看作 $\id x y$ 的一个\emph{构造子}或\emph{引入规则}，这就好比 $A\times B$ 类型本身的"配对"构造子，它在给定 $a:A$ 和 $b:B$ 时构造出序对 $(a,b)$。
从这个视角看，~\eqref{eq:path-prod}的两个分量：
\begin{align*}
  \projpath{1} &: (\id{x}{y}) \to (\id{\proj{1}(x)}{\proj{1} (y)})\\
  \projpath{2} &: (\id{x}{y}) \to (\id{\proj{2}(x)}{\proj{2} (y)})
\end{align*}
就是\emph{消去}规则。
类似地，见证~\eqref{eq:path-prod-inverse}为~\eqref{eq:path-prod}之拟逆的两条同伦，分别由\emph{命题计算规则}：
\index{computation rule!propositional!for identities between pairs}%
\begin{align*}
  {\projpath{1}{(\pairpath(p, q)})}
  &= %_{(\id{\proj{1} x}{\proj{1} y})}
  {p} \\
  {\projpath{2}{(\pairpath(p,q)})}
  &= %_{(\id{\proj{2} x}{\proj{2} y})}
  {q}
\end{align*}
（针对 $p:\id{\proj{1} x}{\proj{1} y}$ 和 $q:\id{\proj{2} x}{\proj{2} y}$）
以及一条\emph{命题唯一性原理}：
\index{uniqueness!principle, propositional!for identities between pairs}%
\[
\id{r}{\pairpath(\projpath{1} (r), \projpath{2} (r)) }
\qquad\text{for } r : \id[A \times B] x y.
\]
构成。

我们也可以按分量刻画 $A\times B$ 中道路的自反性、逆和复合：
\begin{align*}
  {\refl{(z : A \times B)}}
  &= {\pairpath (\refl{\proj{1} z},\refl{\proj{2} z})} \\
  {\opp{p}}
  &= {\pairpath \big(\opp{\projpath{1} (p)},\, \opp{\projpath{2} (p)}\big)} \\
  {{p \ct q}}
  &= {\pairpath \big({\projpath{1} (p)} \ct {\projpath{1} (q)},\,{\projpath{2} (p)} \ct {\projpath{2} (q)}\big)}.
\end{align*}
或者，换种写法：
\begin{alignat*}{2}
  \projpath{i}(\refl{(z : A \times B)}) &= \refl{\proj{i} z} &\qquad (i=1,2)\\
  \pairpath(\opp p, \opp q) &= \opp{\pairpath(p,q)}\\
  \pairpath(p\ct q, p'\ct q') &= \pairpath(p,p') \ct \pairpath(q,q').
\end{alignat*}
所有这些等式都可以这样得到：先对给定的道路作道路归纳，然后归结为自反性。
对于 \cref{sec:equality} 中讨论的其余高阶群胚结构，这一点同样成立，尽管这会逐渐变得繁琐——因为需要插入足够多其他的一致性道路，才能写出一个能通过类型检查的等式。
例如，如果我们用 $\ctassoc(p,q,r)$ 表示 \cref{thm:omg}\ref{item:omg4} 中道路的逆，而用 $\pairct(p,q,p',q')$ 表示上面最后显示的那条道路，那么对于任何具有适当类型的 $u,v,z,w:A\times B$ 和 $p,q,r,p',q',r'$，我们有
\begin{equation*}
  \begin{array}{l}
    \pairct(p\ct q, r, p'\ct q', r') \\
    \ct\;  (\pairct(p,q,p',q') \rightwhisker \pairpath(r,r')) \\
    \ct\;  \ctassoc(\pairpath(p,p'),\pairpath(q,q'),\pairpath(r,r'))\\
    =
    \begin{array}[t]{l}
      \apfunc{\pairpath}({\pairpath(\ctassoc(p,q,r),\ctassoc(p',q',r'))})\\
      \ct\; \pairct(p, q\ct r, p', q'\ct r')\\
      \ct\; (\pairpath(p,p') \leftwhisker \pairct(q,r,q',r')).
    \end{array}
  \end{array}
\end{equation*}
幸运的是，我们永远不需要使用任何这类高维一致性。

\index{transport!in product types}%
我们现在考虑在类型族的逐点积上的传输。
给定类型族 $ A, B : Z \to \type$，我们略滥用记号，用 $A\times B:Z\to\type$ 表示由 $(A\times B)(z) \defeq A(z) \times B(z)$ 定义的类型族。
现在给定 $p : \id[Z]{z}{w}$ 和 $x : A(z) \times B(z)$，我们可以把 $x$ 沿着 $p$ 传输，从而得到 $A(w)\times B(w)$ 中的一个元素。

\begin{thm}\label{thm:trans-prod}
  在上述情形下，我们有
  \[
  \id[A(w) \times B(w)]
  {\transfib{A\times B}px}
  {(\transfib{A}{p}{\proj{1}x}, \transfib{B}{p}{\proj{2}x})}.
  \]
\end{thm}

\begin{proof}
  通过对 $p$ 作道路归纳，我们可以假设 $p$ 是自反性，此时我们有
  \begin{align*}
    \transfib{A\times B}px&\jdeq x\\
    \transfib{A}{p}{\proj{1}x}&\jdeq \proj1x\\
    \transfib{B}{p}{\proj{2}x}&\jdeq \proj2x.
  \end{align*}
  因此，只需证明 $x = (\proj1 x, \proj2x)$。
  但这就是积类型的命题唯一性原理，而如前文所述，它可由 \cref{thm:path-prod} 推出。
\end{proof}

最后，我们考虑 $\apfunc{}$ 在 Cartesian 积下的函子性。
假设给定类型 $A,B,A',B'$ 以及函数 $g:A\to A'$ 和 $h:B\to B'$；那么我们可以通过 $f(x) \defeq (g(\proj1x),h(\proj2x))$ 定义一个函数 $f:A\times B\to A'\times B'$。

\begin{thm}\label{thm:ap-prod}
  在上述情形下，给定 $x,y:A\times B$ 以及 $p:\proj1x=\proj1y$ 和 $q:\proj2x=\proj2y$，我们有
  \[ \id[(f(x)=f(y))]{\ap{f}{\pairpath(p,q)}} {\pairpath(\ap{g}{p},\ap{h}{q})}. \]
\end{thm}

\begin{proof}
  首先注意，上面的方程是良类型的。
  一方面，由于 $\pairpath(p,q):x=y$，我们有 $\ap{f}{\pairpath(p,q)}:f(x)=f(y)$。
  另一方面，由于 $\proj1(f(x))\jdeq g(\proj1x)$ 与 $\proj2(f(x))\jdeq h(\proj2x)$，我们也有 $\pairpath(\ap{g}{p},\ap{h}{q}):f(x)=f(y)$。

  现在，通过归纳，我们可以假设 $x\jdeq(a,b)$ 且 $y\jdeq(a',b')$，此时有 $p:a=a'$ 和 $q:b=b'$。
  于是，再通过道路归纳，我们可以假设 $p$ 和 $q$ 都是自反性，此时所求等式判断性地成立。
\end{proof}

\index{type!product|)}%

\section{\texorpdfstring{$\Sigma$}{Σ}-类型}
\label{sec:compute-sigma}

\index{type!dependent pair|(}%
设 $A$ 为一个类型，$P:A\to\type$ 为一个类型族。
回忆一下，$\Sigma$-类型，或称依值对类型，$\sm{x:A} P(x)$ 是 Cartesian 积类型的一种推广。
因此，我们预期它的高阶群胚结构也是对前一节内容的推广。
特别地，它的道路应当由一对道路组成，但要给出这些道路的正确类型，还需稍加思考。

假设我们在 $\sm{x:A}P(x)$ 中有一条道路 $p:w=w'$。
那么我们会得到 $\ap{\proj{1}}{p}:\proj{1}(w)=\proj{1}(w')$。
然而，我们不能直接问 $\proj{2}(w)$ 是否等于 $\proj{2}(w')$，因为它们不必处于同一个类型中。
但是，我们可以沿着道路 $\ap{\proj{1}}{p}$ 对 $\proj{2}(w)$ 进行传输\index{transport}，这确实能给出一个与 $\proj{2}(w')$ 类型相同的元素。
通过道路归纳，我们确实可以得到一条道路 $\trans{\ap{\proj{1}}{p}}{\proj{2}(w)}=\proj{2}(w')$。

回忆在 \cref{lem:mapdep} 之前的讨论，等式
\narrowequation{
  \trans{\ap{\proj{1}}{p}}{\proj{2}(w)}=\proj{2}(w')
}
可以看作是从 $\proj2(w)$ 到 $\proj2(w')$ 且居于 $A$ 中道路 $\ap{\proj1}{p}$ 之上的道路的类型。
\index{fibration}%
\index{total!space}%
因此，全空间中的一条道路 $w=w'$ 决定——且被决定于——$A$ 中的一条道路 $p:\proj1(w)=\proj1(w')$ 连同一条从 $\proj2(w)$ 到 $\proj2(w')$ 且居于 $p$ 之上的道路。这看来是合理的。

\begin{rmk}\label{rmk:sigma-equality-extraction}
  注意，如果我们有 $x:A$ 和 $u,v:P(x)$ 使得 $(x,u)=(x,v)$，这并不能推出 $u=v$——反例可见 \cref{ex:sigma-eq-components-neq}。
  我们所能得出的结论仅仅是：存在 $p:x=x$ 使得 $\trans p u = v$。
  这一点常常让刚接触类型论的人感到困惑，但从拓扑的观点来看却很自然：在纤维化的全空间中，两个恰好位于同一纤维内的点之间存在道路 $(x,u)=(x,v)$，这一事实并不意味着存在一条完全\emph{位于}该纤维内部的道路 $u=v$。
\end{rmk}

下面的定理说明这个过程也可以反过来。
由于它是 \cref{thm:path-prod} 的直接推广，我们在此就写得简略一些。

\begin{thm}\label{thm:path-sigma}
假设 $P:A\to\type$ 是类型 $A$ 上的类型族，且令 $w,w':\sm{x:A}P(x)$。那么存在等价
\begin{equation*}
\eqvspaced{(w=w')}{\dsm{p:\proj{1}(w)=\proj{1}(w')} \trans{p}{\proj{2}(w)}=\proj{2}(w')}.
\end{equation*}
\end{thm}

\begin{proof}
我们定义函数
\begin{equation*}
f : \prd{w,w':\sm{x:A}P(x)} (w=w') \to \dsm{p:\proj{1}(w)=\proj{1}(w')} \trans{p}{\proj{2}(w)}=\proj{2}(w')
\end{equation*}
由道路归纳给出，其中
\begin{equation*}
f(w,w,\refl{w})\defeq(\refl{\proj{1}(w)},\refl{\proj{2}(w)}).
\end{equation*}
我们要证明 $f$ 是一个等价。

反过来，我们定义
\begin{narrowmultline*}
  g : \prd{w,w':\sm{x:A}P(x)}
      \Parens{\sm{p:\proj{1}(w)=\proj{1}(w')}\trans{p}{\proj{2}(w)}=\proj{2}(w')}
      \to
      \narrowbreak
      (w=w')
\end{narrowmultline*}
通过首先对 $w$ 和 $w'$ 进行归纳，将其分别拆分为 $(w_1,w_2)$ 和
$(w_1',w_2')$，因此只需证明
\begin{equation*}
\Parens{\sm{p:w_1 = w_1'}\trans{p}{w_2}=w_2'} \to ((w_1,w_2)=(w_1',w_2')).
\end{equation*}
接下来，给定一对 $\sm{p:w_1 = w_1'}\trans{p}{w_2}=w_2'$，我们可以使用 $\Sigma$-归纳得到 $p : w_1 = w_1'$ 和 $q :
\trans{p}{w_2}=w_2'$。对 $p$ 进行归纳，我们有 $q :
\trans{(\refl{w_1})}{w_2}=w_2'$，此时只需证明
$(w_1,w_2)=(w_1,w_2')$。但 $\trans{(\refl{w_1})}{w_2} \jdeq w_2$，所以对 $q$ 进行归纳就将目标化为
$(w_1,w_2)=(w_1,w_2)$，而这可以通过 $\refl{(w_1,w_2)}$ 来证明。

接着我们证明对于所有 $w$、$w'$ 和 $r$ 都有 $f(g(r))=r$，其中 $r$ 的类型为
\[\dsm{p:\proj{1}(w)=\proj{1}(w')} (\trans{p}{\proj{2}(w)}=\proj{2}(w')).\]
首先，如 $g$ 的定义中那样，通过配对归纳拆开 $w$、$w'$ 和 $r$，然后使用两次道路归纳将 $r$ 的两个分量都归约为 \refl{}。于是只需证明
$f (g(\refl{w_1},\refl{w_2})) = (\refl{w_1},\refl{w_2})$，而由定义可知这成立。类似地，为证明对所有 $w$、$w'$ 和 $p : w = w'$ 都有 $g(f(p))=p$，我们可以对 $p$ 做道路归纳，然后做配对归纳以拆开 $w$，此时只需证明
$g(f (\refl{(w_1,w_2)})) = \refl{(w_1,w_2)}$，而这由定义即得。

因此，$f$ 有一个拟逆，从而是一个等价。
\end{proof}

正如在 Cartesian 积的情形中所做的那样，我们可以推出一个命题唯一性原理作为特例。

\begin{cor}\label{thm:eta-sigma}
  \index{uniqueness!principle, propositional!for dependent pair types}%
  对于 $z:\sm{x:A} P(x)$，我们有 $z = (\proj1(z),\proj2(z))$。
\end{cor}

\begin{proof}
  我们有 $\refl{\proj1(z)} : \proj1(z) = \proj1(\proj1(z),\proj2(z))$，故由 \cref{thm:path-sigma}，只需展示一条道路 $\trans{(\refl{\proj1(z)})}{\proj2(z)} = \proj2(\proj1(z),\proj2(z))$。
  但等式两边都判断相等于 $\proj2(z)$。
\end{proof}

类似于二元 Cartesian 积的情形，我们可以将
\cref{thm:path-sigma} 的逆向
视为引入形式（\pairpath{}{}），正向视为
消去形式（\projpath{1} 和 \projpath{2}），而该等价则给出了这些形式的命题计算规则和唯一性原理。

注意，\cref{thm:path-lifting} 中定义的提升道路 $\mathsf{lift}(u,p)$（即 $p:x=y$ 在 $u:P(x)$ 处的提升）可以等同于引入形式的特例
\[\pairpath(p,\refl{\trans p u}):(x,u) = (y,\trans p u).\]
\index{transport!in dependent pair types}%
它出现在关于 transport 在 $\Sigma$-类型上作用的定理陈述中，这也是二元 Cartesian 积情形的推广：

\begin{thm}\label{transport-Sigma}
  假设我们有类型族
  %
  \begin{equation*}
    P:A\to\type
    \qquad\text{and}\qquad
    Q:\Parens{\sm{x:A} P(x)}\to\type.
  \end{equation*}
  %
  那么我们可以构造 $A$ 上由下式定义的类型族
  \begin{equation*}
    x \mapsto \sm{u:P(x)} Q(x,u).
  \end{equation*}
  对任意道路 $p:x=y$ 和任意 $(u,z):\sm{u:P(x)} Q(x,u)$，我们有
  \begin{equation*}
    \trans{p}{u,z}=\big(\trans{p}{u},\,\trans{\pairpath(p,\refl{\trans pu})}{z}\big).
  \end{equation*}
\end{thm}

\begin{proof}
直接由道路归纳即得。
\end{proof}

我们把 \cref{thm:ap-prod} 的推广（参见 \cref{ex:ap-sigma}）的陈述和证明留给读者，并请读者逐分量地刻画 $\Sigma$-类型的自反性、逆和复合。

\index{type!dependent pair|)}%

\section{单位类型}
\label{sec:compute-unit}

\index{type!unit|(}%
平凡的情形有时也很重要，所以我们简要提一下单位类型~\unit。

\begin{thm}\label{thm:path-unit}
  对任意 $x,y:\unit$，我们有 $\eqv{(x=y)}{\unit}$。
\end{thm}

这个证明一开始可能会让人想先对 $x$ 和 $y$ 进行 $\unit$-归纳，从而将问题化为 $\eqv{(\ttt=\ttt)}{\unit}$。
但这样一来我们就会卡住，因为我们无法对 $p:\ttt=\ttt$ 进行道路归纳。
因此，我们转而尽可能一般地处理 $x$ 和 $y$，只在最后才通过归纳将它们归约为 $\ttt$。

\begin{proof}
  定义一个从 $(x=y)$ 到 $\unit$ 的函数是容易的：只需将所有元素都映到 $\ttt$。
  反过来，对于任意 $x,y:\unit$，由归纳法可以假设 $x\jdeq \ttt\jdeq y$。
  此时我们有 $\refl{\ttt}:x=y$，这就给出了一个常值函数 $\unit\to(x=y)$。

  为了证明这两个函数互逆，首先考虑一个元素 $u:\unit$。
  我们可以假设 $u\jdeq\ttt$，而这正是复合函数 $\unit \to (x=y)\to\unit$ 作用在 $u$ 上的结果。

  另一方面，给定 $p:x=y$。
  由道路归纳，我们可以假设 $x\jdeq y$ 且 $p$ 就是 $\refl x$。
  接着我们可以假设 $x$ 就是 $\ttt$，此时复合函数 $(x=y) \to \unit\to(x=y)$ 将 $p$ 映到 $\refl x$，也就是映回 $p$ 本身。
\end{proof}

因此，$\unit$ 的任意两个元素都是相等的。
我们将其留给读者，以引入规则、消去规则、计算规则和唯一性规则的形式表述此等价。
\index{transport!in unit type}%
$\unit$ 的 transport 引理就是常值类型族（\cref{thm:trans-trivial}）的 transport 引理。

\index{type!unit|)}%

\section{\texorpdfstring{$\Pi$}{Π}-类型与函数外延性公理}
\label{sec:compute-pi}

\index{type!dependent function|(}%
\index{type!function|(}%
\index{homotopy|(}%
给定一个类型 $A$ 和一个类型族 $B : A \to \type$，考虑依值函数类型 $\prd{x:A}B(x)$。
我们期望，在 $\prd{x:A} B(x)$ 中从 $f$ 到 $g$ 的道路类型 $f=g$，等价于逐点道路的类型：\index{pointwise!equality of functions}
\begin{equation}
  \eqvspaced{(\id{f}{g})}{\Parens{\prd{x:A} (\id[B(x)]{f(x)}{g(x)})}}.\label{eq:path-forall}
\end{equation}
从传统视角看，这意味着：两个在每个点上都相等的函数，作为函数也是相等的。
\index{continuity of functions in type theory@``continuity'' of functions in type theory}%
从拓扑视角看，这意味着：函数空间中的一条道路，等同于一个连续同伦。
\index{functoriality of functions in type theory@``functoriality'' of functions in type theory}%
而从范畴视角看，这意味着：函子范畴中的一个同构，就是一个自然的同构族。

然而，与前几节的情况不同，\cref{cha:typetheory} 中给出的基础类型论并不足以证明~\eqref{eq:path-forall}。
我们所能说的只是：存在某个函数
\begin{equation}\label{eq:happly}
  \happly : (\id{f}{g}) \to \prd{x:A} (\id[B(x)]{f(x)}{g(x)})
\end{equation}
它可以由道路归纳容易地定义。
因此，暂时假定：

\begin{axiom}[函数外延性]\label{axiom:funext}
  \indexsee{axiom!function extensionality}{function extensionality}%
  \indexdef{function extensionality}%
  对于任何 $A$、$B$、$f$ 和 $g$，函数~\eqref{eq:happly} 是一个等价。
\end{axiom}

我们将在后续章节中看到，这个公理既能从泛等性（参见 \cref{sec:compute-universe,sec:univalence-implies-funext}）推出，也能从区间类型（参见 \cref{sec:interval} 与 \cref{ex:funext-from-interval}）推出。

特别地，\cref{axiom:funext} 蕴含了~\eqref{eq:happly} 有一个拟逆
\[
\funext : \Parens{\prd{x:A} (\id{f(x)}{g(x)})} \to {(\id{f}{g})}.
\]
这个函数也常被称为"函数外延性"。
正如我们在 \cref{sec:compute-cartprod} 中对 $\pairpath$ 所做的那样，我们可以将 $\funext$ 视为类型 $\id f g$ 的一个\emph{引入规则}。
从这个视角看，$\happly$ 就是\emph{消去规则}，而那些见证 $\funext$ 作为 $\happly$ 之拟逆的同伦，则成为一个命题计算规则\index{computation rule!propositional!for identities between functions}
\[
\id{\happly({\funext{(h)}},x)}{h(x)} \qquad\text{for }h:\prd{x:A} (\id{f(x)}{g(x)})
\]
和一个命题唯一性原理\index{uniqueness!principle!for identities between functions}：
\[
\id{p}{\funext (x \mapsto \happly(p,{x}))} \qquad\text{for } p: \id f g.
\]

我们也可以计算 $\Pi$-类型中的恒等、逆与复合；它们简单地由逐点运算给出：\index{pointwise!operations on functions}
\begin{align*}
\refl{f} &= \funext(x \mapsto \refl{f(x)}) \\
\opp{\alpha} &= \funext (x \mapsto \opp{\happly (\alpha,x)})  \\
{\alpha} \ct \beta &= \funext (x \mapsto {\happly({\alpha},x) \ct \happly({\beta},x)}).
\end{align*}
第一个等式由 $\happly$ 的定义可得，第二、第三个则是简单的道路归纳。

由于当 $B$ 不依赖于 $x$ 时，非依值函数类型 $A\to B$ 是依值函数类型 $\prd{x:A} B(x)$ 的特例，上述所有讨论同样适用于非依值情形。
\index{transport!in function types}%
不过，在非依值情形下，transport 的规则要稍微简单一些。
给定一个类型 $X$、一条道路 $p:\id[X]{x_1}{x_2}$、两个类型族 $A,B:X\to \type$，以及一个函数 $f : A(x_1) \to B(x_1)$，我们有
\begin{align}\label{eq:transport-arrow}
  \transfib{A\to B}{p}{f} &=
  \Big(x \mapsto \transfib{B}{p}{f(\transfib{A}{\opp p}{x})}\Big)
\end{align}
其中 $A\to B$（滥用记号）表示由下式定义的类型族 $X\to \type$：
\[(A\to B)(x) \defeq (A(x)\to B(x)).\]
也就是说，当我们沿着道路 $p:x_1=x_2$ transport 一个函数 $f:A(x_1)\to B(x_1)$ 时，我们得到的函数是 $A(x_2)\to B(x_2)$：它将其参数沿着 $p$ 反向 transport（在类型族 $A$ 中），然后应用 $f$，再将结果沿着 $p$ 正向 transport（在类型族 $B$ 中）。
这容易由道路归纳证明。

\index{transport!in dependent function types}%
传输依值函数的情形与此类似，但要更复杂一些。
假设给定如前所述的 $X$ 与 $p$，类型族 $A:X\to \type$ 和 $B:\prd{x:X} (A(x)\to\type)$，以及一个依值函数 $f : \prd{a:A(x_1)} B(x_1,a)$。
那么对于 $a:A(x_2)$，我们有
\begin{narrowmultline*}
  \transfib{\Pi_A(B)}{p}{f}(a) = \narrowbreak
  \Transfib{\widehat{B}}{\opp{(\pairpath(\opp{p},\refl{ \trans{\opp p}{a} }))}}{f(\transfib{A}{\opp p}{a})}
\end{narrowmultline*}
其中 $\Pi_A(B)$ 与 $\widehat{B}$ 分别表示类型族：
\begin{equation}\label{eq:transport-arrow-families}
\begin{array}{rclcl}
\Pi_A(B) &\defeq& \big(x\mapsto \prd{a:A(x)} B(x,a) \big) &:& X\to \type\\
\widehat{B} &\defeq& \big(w \mapsto B(\proj1w,\proj2w) \big) &:& \big(\sm{x:X} A(x)\big) \to \type.
\end{array}
\end{equation}
这些公式看起来可能令人生畏，不必在意细节。基本思想与非依值函数类型的情形完全一样：将参数反向传输，应用函数，然后将结果再次正向传输。

现在回顾一下，对于一个一般的类型族 $P:X\to\type$，在 \cref{sec:functors} 中我们将从 $u:P(x)$ 到 $v:P(y)$、位于 $p:\id[X]xy$ 之上的\emph{依值道路（dependent paths）}类型定义为 $\id[P(y)]{\trans{p}{u}}{v}$。
当 $P$ 是一个函数类型族时，有一种等价的表示方式往往更为方便。
\index{path!dependent!in function types}

\begin{lem}\label{thm:dpath-arrow}
  给定类型族 $A,B:X\to\type$ 与 $p:\id[X]xy$，以及 $f:A(x)\to B(x)$ 和 $g:A(y)\to B(y)$，我们有一个等价
  \[ \eqvspaced{ \big(\trans{p}{f} = {g}\big) } { \prd{a:A(x)}  (\trans{p}{f(a)} = g(\trans{p}{a})) }. \]
  此外，若 $q:\trans{p}{f} = {g}$ 在此等价下对应于 $\widehat q$，则对于 $a:A(x)$，道路
  \[ \happly(q,\trans p a) : (\trans p f)(\trans p a) = g(\trans p a)\]
  等于拼接道路 $i\ct j\ct k$，其中
  \begin{itemize}
  \item $i:(\trans p f)(\trans p a) = \trans p {f (\trans {\opp p}{\trans p a})}$ 来自~\eqref{eq:transport-arrow}，
  \item $j:\trans p {f (\trans {\opp p}{\trans p a})} = \trans p {f(a)}$ 来自 \cref{thm:transport-concat,thm:omg}，且
  \item $k:\trans p {f(a)}= g(\trans p a)$ 就是 $\widehat{q}(a)$。
  \end{itemize}
\end{lem}

\begin{proof}
  通过道路归纳，可设 $p$ 为自反性，此时所需的等价便归结为函数外延性。
  第二条论断则由函数外延性的计算规则可得。
\end{proof}

一般说来，我们常常需要考虑若干条道路的拼接，而每条道路又都来自某个先前证明过的引理或假设的对象。如果像上面第二条论断那样，给拼接中的每条道路都单独起一个名字，描述起来会相当繁琐。
因此，我们采用一种惯例，按照熟悉的数学风格，以``带理由的等式链''来书写此类拼接，并允许省略那些读者容易自行补全的理由。
例如，从 \cref{thm:dpath-arrow} 出发的道路 $i\ct j\ct k$ 就可以写成这样：
  \begin{align*}
    (\trans p f)(\trans p a)
    &= \trans p {f (\trans {\opp p}{\trans p a})}
    \tag{by~\eqref{eq:transport-arrow}}\\
    &= \trans p {f(a)}\\
    &= g(\trans p a).
    \tag{by $\widehat{q}$}
  \end{align*}
在普通数学中，这样的等式链仅用于证明两个对象相等。
而我们在此对其加以强化：用它来描述连接二者的一个\emph{特定}道路。

照例，\cref{thm:dpath-arrow} 有一个针对依值函数的版本，形式相似但更为复杂。
\index{path!dependent!in dependent function types}

\begin{lem}\label{thm:dpath-forall}
  给定类型族 $A:X\to\type$ 和 $B:\prd{x:X} A(x)\to\type$，以及 $p:\id[X]xy$，还有 $f:\prd{a:A(x)} B(x,a)$ 和 $g:\prd{a:A(y)} B(y,a)$，我们有一个等价
  \[ \eqvspaced{ \big(\trans{p}{f} = {g}\big) } { \Parens{\prd{a:A(x)}  \transfib{\widehat{B}}{\pairpath(p,\refl{\trans pa})}{f(a)} = g(\trans{p}{a}) } } \]
  其中 $\widehat{B}$ 如~\eqref{eq:transport-arrow-families} 所述。
\end{lem}

我们将此引理的证明及适当计算规则的陈述留给读者。

\index{homotopy|)}%
\index{type!dependent function|)}%
\index{type!function|)}%

\section{宇宙与泛等公理}
\label{sec:compute-universe}

\index{type!universe|(}%
\index{equivalence|(}%
给定两个类型 $A$ 和 $B$，我们可以将它们视为某个宇宙类型 \type 中的元素，从而构成相等类型 $\id[\type]AB$。
如引言所述，\emph{泛等性（univalence）}就是将 $\id[\type]AB$ 与从 $A$ 到 $B$ 的等价类型 $(\eqv AB)$（我们在 \cref{sec:basics-equivalences} 中描述过）视为等同。
我们通过下述规范函数来实现这一等同。

\begin{lem}\label{thm:idtoeqv}
  对于类型 $A,B:\type$，存在一个特定的函数，
  \begin{equation}\label{eq:uidtoeqv}
    \idtoeqv : (\id[\type]AB) \to (\eqv A B),
  \end{equation}
  其定义见证明部分。
\end{lem}

\begin{proof}
  本可以直接通过对相等性做归纳来构造它，但下面的描述更为方便。
  \index{identity!function}%
  \index{function!identity}%
  注意到恒等函数 $\idfunc[\type]:\type\to\type$ 可以被看作一个以宇宙 \type 为指标的类型族；它给每个类型 $X:\type$ 指派类型 $X$ 自身。
  （若视其为纤维化，其全空间就是``带点类型''所构成的类型 $\sm{A:\type}A$；另见 \cref{sec:object-classification}。）
  于是，给定一条道路 $p:A =_\type B$，我们便有一个传输\index{transport}函数 $\transf{p}:A \to B$。
  我们断言 $\transf{p}$ 是一个等价。
  但通过归纳，只需假设 $p$ 是 $\refl A$，此时 $\transf{p} \jdeq \idfunc[A]$，而根据 \cref{eg:idequiv}，它是一个等价。
  因此，我们可以将 $\idtoeqv(p)$ 定义为 $\transf{p}$（连同上述证明它确为等价的证据）。
\end{proof}

我们希望能说 \idtoeqv 是一个等价。然而，正如函数类型中 $\happly$ 的情况那样，\cref{cha:typetheory} 中所描述的类型论并不足以保证这一点。于是，与我们处理函数外延性的方式相同，我们将这一性质表述为一个公理：Voevodsky 的\emph{泛等公理（univalence axiom）}。

\begin{axiom}[泛等]\label{axiom:univalence}
  \indexdef{univalence axiom}%
  \indexsee{axiom!univalence}{univalence axiom}%
  对于任意 $A,B:\type$，函数~\eqref{eq:uidtoeqv} 是一个等价。
\end{axiom}

因此，特别地，我们有
  \[
\eqv{(\id[\type]{A}{B})}{(\eqv A B)}.
\]

严格来说，泛等公理是关于一个特定宇宙类型 $\UU$ 的命题。
如果一个宇宙 $\UU$ 满足这条公理，我们就称它是\define{泛等的（univalent）}。
\indexdef{type!universe!univalent}%
\indexdef{univalent universe}%
除非另有说明（例如在 \cref{sec:univalence-implies-funext} 中），我们将假设\emph{所有}宇宙都是泛等的。

\begin{rmk}
  对泛等公理而言，至关重要的是我们使用了 \cref{sec:basics-equivalences} 中所述的``好''版本的 $\isequiv$ 来定义 $\eqv AB$，而不是（比如说）$\sm{f:A\to B} \qinv(f)$。详见 \cref{ex:qinv-univalence}。
\end{rmk}

特别地，泛等性意味着\emph{等价的类型可以被等同}。
如同前面几节所做的，将这个等价分解为以下几部分会很有用：
%
\symlabel{ua}


\begin{itemize}
\item {(\id[\type]{A}{B})} 的一个引入规则，记作 $\ua$（代表 ``univalence axiom''）：
  \[
  \ua : ({\eqv A B}) \to (\id[\type]{A}{B}).
  \]
\item 消去规则，即 $\idtoeqv$：
  \[
  \idtoeqv \jdeq \transfibf{X \mapsto X} : (\id[\type]{A}{B}) \to (\eqv A B).
  \]
\item 命题计算规则\index{computation rule!propositional!for univalence}：
  \[
  \transfib{X \mapsto X}{\ua(f)}{x} = f(x).
  \]
\item 命题唯一性原理\index{uniqueness!principle, propositional!for univalence}：
  对于任意 $p : \id A B$，
  \[
  \id{p}{\ua(\transfibf{X \mapsto X}(p))}.
  \]
\end{itemize}


%
我们还可以将宇宙中相等的自反性、拼接和逆，与等价上相应的运算等同起来：
\begin{align*}
  \refl{A} &= \ua(\idfunc[A]) \\
  \ua(f) \ct \ua(g) &= \ua(g\circ f) \\
  \opp{\ua(f)} &= \ua(f^{-1}).
\end{align*}
其中第一个等式之所以成立，是因为根据 $\idtoeqv$ 的定义有 $\idfunc[A] = \idtoeqv(\refl{A})$，而 $\ua$ 是 $\idtoeqv$ 的逆。
对于第二个，如果我们令 $p \defeq \ua(f)$ 且 $q\defeq \ua(g)$，那么利用 \cref{thm:transport-concat} 和 $\idtoeqv$ 的定义，我们有
\[ \ua(g\circ f) = \ua(\idtoeqv(q) \circ \idtoeqv(p)) = \ua(\idtoeqv(p\ct q)) = p\ct q\]
第三个的证明是类似的。

下面这个观察是 \cref{thm:transport-compose} 的一个特例，它在应用泛等公理时常常有用。

\begin{lem}\label{thm:transport-is-ap}
  对于任意类型族 $B:A\to\type$ 以及 $x,y:A$ 和一条道路 $p:x=y$ 和一个元素 $u:B(x)$，我们有
  \begin{align*}
    \transfib{B}{p}{u} &= \transfib{X\mapsto X}{\apfunc{B}(p)}{u}\\
    &= \idtoeqv(\apfunc{B}(p))(u).
  \end{align*}
\end{lem}

\index{equivalence|)}%
\index{type!universe|)}%

\section{相等类型}
\label{sec:compute-paths}

\index{type!identity|(}%
正如类型 $\id[A]{a}{a'}$ 在同构意义下被刻画、对每个 $A$ 都有一个单独的``定义''一样，对于道路之间的道路类型 $\id[{\id[A]{a}{a'}}]{p}{q}$ $p,q : \id[A]{a}{a'}$ 也没有简单的刻画方式。
然而，其他一般性的定理类确实可以推广到相等类型上，例如它们尊重等价这一事实。

\begin{thm}\label{thm:paths-respects-equiv}
  若 $f : A \to B$ 是一个等价，则对于所有 $a,a':A$，如下映射亦然：
  \[\apfunc{f} : (\id[A]{a}{a'}) \to (\id[B]{f(a)}{f(a')}).\]
\end{thm}

\begin{proof}
  令 $\opp f$ 为 $f$ 的一个拟逆，并配有同伦
  %
  \begin{equation*}
    \alpha:\prd{b:B} (f(\opp f(b))=b)
    \qquad\text{and}\qquad
    \beta:\prd{a:A} (\opp f(f(a)) = a).
  \end{equation*}
  %
  $\apfunc{f}$ 的拟逆本质上是
  \[\apfunc{\opp f} : (\id{f(a)}{f(a')}) \to (\id{\opp f(f(a))}{\opp f(f(a'))}).\]
  然而，为了从 $\apfunc{\opp f}(q)$ 得到 $\id[A]{a}{a'}$ 中的一个元素，我们必须在两侧拼接上道路 $\opp{\beta_a}$ 和 $\beta_{a'}$。
  为了证明这给出了 $\apfunc{f}$ 的一个拟逆，一方面我们必须证明：对于任意 $p:\id[A]{a}{a'}$，有
  \[ \opp{\beta_a} \ct \apfunc{\opp f}(\apfunc{f}(p)) \ct \beta_{a'} = p. \]
  这可由 $\apfunc{}$ 的函子性以及同伦的自然性得到，参见 \cref{lem:ap-functor,lem:htpy-natural}。
  另一方面，我们必须证明：对于任意 $q:\id[B]{f(a)}{f(a')}$，有
  \[ \apfunc{f}\big( \opp{\beta_a} \ct \apfunc{\opp f}(q) \ct \beta_{a'} \big) = q. \]
  这个证明略复杂一些，但每一步也同样是应用 \cref{lem:ap-functor,lem:htpy-natural}（或直接抵消逆道路）：
  \begin{align*}
    \apfunc{f}\big( \narrowamp\opp{\beta_a} \ct \apfunc{\opp f}(q) \ct \beta_{a'} \big) \narrowbreak
    &= \opp{\alpha_{f(a)}} \ct {\alpha_{f(a)}} \ct
    \apfunc{f}\big( \opp{\beta_a} \ct \apfunc{\opp f}(q) \ct \beta_{a'} \big)
    \ct \opp{\alpha_{f(a')}} \ct {\alpha_{f(a')}}\\
    &= \opp{\alpha_{f(a)}} \ct
    \apfunc f \big(\apfunc{\opp f}\big(\apfunc{f}\big( \opp{\beta_a} \ct \apfunc{\opp f}(q) \ct \beta_{a'} \big)\big)\big)
    \ct {\alpha_{f(a')}}\\
    &= \opp{\alpha_{f(a)}} \ct
    \apfunc f \big(\beta_a \ct \opp{\beta_a} \ct \apfunc{\opp f}(q) \ct \beta_{a'} \ct \opp{\beta_{a'}} \big)
    \ct {\alpha_{f(a')}}\\
    &= \opp{\alpha_{f(a)}} \ct
    \apfunc f (\apfunc{\opp f}(q))
    \ct {\alpha_{f(a')}}\\
    &= q.\qedhere
  \end{align*}
\end{proof}

因此，如果对于某个类型 $A$，我们能完全刻画 $\id[A]{a}{a'}$，那么类型 $\id[{\id[A]{a}{a'}}]{p}{q}$ 也就随之确定了。例如：

\begin{itemize}
\item 当 $p,q : \id[A \times B]{w}{w'}$ 时，道路 $p = q$ 等价于道路对
  \[\id[{\id[A]{\proj{1} w}{\proj{1} w'}}]{\projpath{1}{p}}{\projpath{1}{q}}
  \quad\text{and}\quad
  \id[{\id[B]{\proj{2} w}{\proj{2} w'}}]{\projpath{2}{p}}{\projpath{2}{q}}.
  \]
\item 当 $p,q : \id[\prd{x:A} B(x)]{f}{g}$ 时，道路 $p = q$ 等价于同伦
  \[\prd{x:A} (\id[f(x)=g(x)] {\happly(p)(x)}{\happly(q)(x)}).\]
\end{itemize}

\index{transport!in identity types}%
接下来我们考虑在以道路类型为纤维的类型族中的传输，即考虑 $C:A\to\type$，其中每个 $C(x)$ 都是一个相等类型。
最简单的情形是 $C(x)$ 本身就是 $A$ 中（可能固定了一个端点的）道路类型。

\begin{lem}\label{cor:transport-path-prepost}
  对于任意 $A$ 和 $a:A$，若 $p:x_1=x_2$，则有
  %
  \begin{align*}
    \transfib{x \mapsto (\id{a}{x})} {p} {q} &= q \ct p
    & &\text{for $q:a=x_1$,}\\
    \transfib{x \mapsto (\id{x}{a})} {p} {q} &= \opp {p} \ct q
    & &\text{for $q:x_1=a$,}\\
    \transfib{x \mapsto (\id{x}{x})} {p} {q} &= \opp{p} \ct q \ct p
    & &\text{for $q:x_1=x_1$.}
  \end{align*}
\end{lem}

\begin{proof}
  对 $p$ 作道路归纳，再利用复合的单位律。
\end{proof}

换言之，用 ${x \mapsto \id{c}{x}}$ 传输就是后复合，而用 ${x \mapsto \id{x}{c}}$ 传输则是逆变的前复合。
在范畴论中，这些可能令人熟悉，它们正是协变和逆变 Hom 函子 $\hom(c, {\blank})$ 与 $\hom({\blank},c)$ 的函子性作用。

类似地，我们可以证明下述 \cref{cor:transport-path-prepost} 的更一般形式，它与 \cref{thm:transport-compose} 相关。

\begin{thm}\label{thm:transport-path}
  对于 $f,g:A\to B$，满足条件 $p : \id[A]{a}{a'}$ 与 $q : \id[B]{f(a)}{g(a)}$，我们有
  \begin{equation*}
    \id[f(a') = g(a')]{\transfib{x \mapsto \id[B]{f(x)}{g(x)}}{p}{q}}
    {\opp{(\apfunc{f}{p})} \ct q \ct \apfunc{g}{p}}.
  \end{equation*}
\end{thm}

因为 $\apfunc{(x \mapsto x)}$ 是恒等函数，而 $\apfunc{(x \mapsto c)}$（此处 $c$ 是常量）将道路 $p$ 映为 $\refl{c}$，所以 \cref{cor:transport-path-prepost} 是一个特例。
更一般的版本允许 $B$ 是以 $A$ 为指标的类型族：

\begin{thm}\label{thm:transport-path2}
  设 $B : A \to \type$ 且 $f,g : \prd{x:A} B(x)$，并满足 $p : \id[A]{a}{a'}$ 与 $q : \id[B(a)]{f(a)}{g(a)}$，
  则我们有
  \begin{equation*}
    \transfib{x \mapsto \id[B(x)]{f(x)}{g(x)}}{p}{q} =
    \opp{(\apdfunc{f}(p))} \ct \apfunc{(\transfibf{B}{p})}(q) \ct \apdfunc{g}(p).
  \end{equation*}
\end{thm}

最后，如同在 \cref{sec:compute-pi} 中那样，对于相等类型族，依值道路还有另一种等价的刻画。
\index{path!dependent!in identity types}

\begin{thm}\label{thm:dpath-path}
  对于满足 $p:\id[A]a{a'}$ 的情形，其中 $q:a=a$ 且 $r:a'=a'$，我们有
  \[ \eqvspaced{ \big(\transfib{x\mapsto (x=x)}{p}{q} = r \big) }{ \big( q \ct p = p \ct r \big). } \]
\end{thm}

\begin{proof}
  对 $p$ 作道路归纳，然后利用以下事实：与单位等式 $q\ct 1 = q$ 和 $r = 1\ct r$ 的复合是一个等价。
\end{proof}

还存在一些涉及函数应用的更一般的等价，类似于 \cref{thm:transport-path,thm:transport-path2}。

\index{type!identity|)}%

\section{余积}
\label{sec:compute-coprod}

\index{type!coproduct|(}%
\index{encode-decode method|(}%
到目前为止，我们讨论过的类型形成子大多都属于所谓的\emph{负类型}。
\index{type!negative}\index{negative!type}%
\index{polarity}%
直观上说，这意味着它们的元素是由其在消去规则下的行为所决定的：一个（依值）对由其投影决定，而一个（依值）函数则由其取值决定。
负类型的相等类型几乎总能被直接刻画——连同其所有高阶结构——正如我们在 \crefrange{sec:compute-cartprod}{sec:compute-pi} 中所做的那样。
宇宙本身并不完全是负类型，但其相等类型的行为也类似：我们有一个直接的刻画（泛等性），并且也能描述其高阶结构。
当然，相等类型本身则是一个特例。

现在我们来考察第一个\emph{正类型}形成子的例子。
\index{type!positive}\index{positive!type}%
同样非正式地说，正类型是指由某些构造子所"呈现"的类型，其呈现的泛性质由消去规则来表达。
（从范畴论的角度看，正类型具有"映射出"的泛性质，而负类型则具有"映射入"的泛性质。）
由于一般来说，基于呈现的计算是不可计算问题，因此对于正类型，我们不能总是期望对其相等类型有一个直接的刻画。
不过，在许多具体情形中，刻画或部分刻画确实存在，并且可以通过我们借这个例子引入的一般方法来获得。

（技术上来说，我们选择的 Cartesian 积与 $\Sigma$ 类型的呈现也是正类型的呈现。
但由于这些类型同时允许一种与之相差无几的负呈现，它们的相等类型便已有直接的刻画，无需使用这里要介绍的方法。）

考虑余积类型 $A+B$，它由单射 $\inl:A\to A+B$ 与 $\inr:B\to A+B$ 所"呈现"。
直观上，我们期望 $A+B$ 中包含 $A$ 与 $B$ 的精确副本，且它们互不相交，因此应有
\begin{align}
  {(\inl(a_1)=\inl(a_2))}&\eqvsym {(a_1=a_2)} \label{eq:inlinj}\\
  {(\inr(b_1)=\inr(b_2))}&\eqvsym {(b_1=b_2)}\\
  {(\inl(a)= \inr(b))} &\eqvsym {\emptyt}. \label{eq:inlrdj}
\end{align}
我们如下证明这一点。
固定一个元素 $a_0:A$；我们将刻画类型族
\begin{equation}
  (x\mapsto (\inl(a_0)=x)) : A+B \to \type.\label{eq:sumcodefam}
\end{equation}
类似的论证可以对任意 $b_0:B$ 刻画相应的族 $x\mapsto (x = \inr(b_0))$。
合起来，这些刻画便蕴含~\eqref{eq:inlinj}--\eqref{eq:inlrdj}。

为了刻画~\eqref{eq:sumcodefam}，我们将定义一个类型族 $\code:A+B\to\type$，并证明 $\prd{x:A+B} (\eqv{(\inl(a_0)=x)}{\code(x)})$。
由于我们想由此得到结论~\eqref{eq:inlinj}，就应有 $\code(\inl(a)) = (a_0=a)$；
而由于我们也想得到结论~\eqref{eq:inlrdj}，就应有 $\code (\inr(b)) = \emptyt$。
关键的想法在于，我们可以利用 $A+B$ 的递归原理，通过下面两个方程来\emph{定义} $\code:A+B\to\type$：
\begin{align*}
  \code(\inl(a)) &\defeq (a_0=a),\\
  \code(\inr(b)) &\defeq \emptyt.
\end{align*}
这是证明技巧的一个非常简单的实例——在同伦类型论中进行同伦理论研究时，这一技巧颇为常用；
例如可参见\ \cref{sec:pi1-s1-intro,sec:general-encode-decode}。
%
现在我们可以证明：

\begin{thm}\label{thm:path-coprod}
  对所有 $x:A+B$，有 $\eqv{(\inl(a_0)=x)}{\code(x)}$。
\end{thm}

\begin{proof}
  以下证明的关键在于，我们对所有点 $x$ 同时进行论证，从而能够使用余积的消去原理。
  我们首先通过将自反性沿着 $p$ 输运来定义一个函数
  \[ \encode : \prd{x:A+B}{p:\inl(a_0)=x} \code(x) \]
  ：
  \[ \encode(x,p) \defeq \transfib{\code}{p}{\refl{a_0}}. \]
  注意 $\refl{a_0} : \code(\inl(a_0))$，因为根据 \code 的定义有 $\code(\inl(a_0))\jdeq (a_0=a_0)$。
  接下来，我们定义一个函数
  \[ \decode : \prd{x:A+B}{c:\code(x)} (\inl(a_0)=x). \]
  为了定义 $\decode(x,c)$，我们可以先利用 $A+B$ 的消去原理，根据 $x$ 是 $\inl(a)$ 形式还是 $\inr(b)$ 形式来分情况讨论。

  在第一种情况，即 $x\jdeq \inl(a)$ 时，有 $\code(x)\jdeq (a_0=a)$，于是 $c$ 是 $a_0$ 与 $a$ 之间的一个相等证明。
  因此 $\apfunc{\inl}(c):(\inl(a_0)=\inl(a))$，我们可以将 $\decode(\inl(a),c)$ 定义为此值。

  在第二种情况，即 $x\jdeq \inr(b)$ 时，有 $\code(x)\jdeq \emptyt$，于是 $c$ 属于空类型。
  因此，利用 $\emptyt$ 的消去规则，就能给出 $\decode(\inr(b),c)$ 的一个值。

  这样便完成了 \decode 的定义；我们现在证明，对所有 $x$，$\encode(x,{\blank})$ 和 $\decode(x,{\blank})$ 互为拟逆（quasi-inverses）。
  一方面，假设给定 $x:A+B$ 和 $p:\inl(a_0)=x$；我们想要证明
  \narrowequation{
    \decode(x,\encode(x,p)) = p.
  }
  但根据（基于起点的）道路归纳，只需考虑 $x\jdeq\inl(a_0)$ 且 $p\jdeq \refl{\inl(a_0)}$ 的情形：
  \begin{align*}
    \decode(x,\encode(x,p))
    &\jdeq \decode(\inl(a_0),\encode(\inl(a_0),\refl{\inl(a_0)}))\\
    &\jdeq \decode(\inl(a_0),\transfib{\code}{\refl{\inl(a_0)}}{\refl{a_0}})\\
    &\jdeq \decode(\inl(a_0),\refl{a_0})\\
    &\jdeq \apfunc{\inl}(\refl{a_0})\\
    &\jdeq \refl{\inl(a_0)}\\
    &\jdeq p.
  \end{align*}
  另一方面，设 $x:A+B$ 和 $c:\code(x)$；我们想要证明 $\encode(x,\decode(x,c))=c$。
  我们可以再次根据 $x$ 分情况讨论。
  如果 $x\jdeq\inl(a)$，那么 $c:a_0=a$ 且 $\decode(x,c)\jdeq \apfunc{\inl}(c)$，于是有
  \begin{align}
    \encode(x,\decode(x,c))
    &\jdeq \transfib{\code}{\apfunc{\inl}(c)}{\refl{a_0}}
    \notag\\
    &= \transfib{a\mapsto (a_0=a)}{c}{\refl{a_0}}
    \tag{by \cref{thm:transport-compose}}\\
    &= \refl{a_0} \ct c
    \tag{by \cref{cor:transport-path-prepost}}\\
    &= c. \notag
  \end{align}
  最后，如果 $x\jdeq \inr(b)$，那么 $c:\emptyt$，因此我们可以得出任意结论。
\end{proof}

\noindent
当然，如果我们固定的是 $b_0:B$ 而不是 $a_0:A$，也有对应的定理。

特别地，\cref{thm:path-coprod} 蕴涵着，对于任意 $a : A$ 和 $b : B$，存在函数
%
\[ \encode(\inl(a), {\blank}) : (\inl(a_0)=\inl(a)) \to (a_0=a)\]
%
与
%
\[ \encode(\inr(b), {\blank}) : (\inl(a_0)=\inr(b)) \to \emptyt. \]
%
第二个式子表明 "$\inl(a_0)$ 与 $\inr(b)$ 不相等"，即 $\inl$ 和 $\inr$ 的像是不交的。
在将相等类型视为命题的传统解读下，第一个式子只是 $\inl$ 的单射性。
而 \cref{thm:path-coprod} 的完整同伦表述则给出了更多信息：类型 $\inl(a_0)=\inl(a)$ 和 $a_0=a$ 实际上是等价的，类型 $\inr(b_0)=\inr(b)$ 和 $b_0=b$ 也是等价的。

\begin{rmk}\label{rmk:true-neq-false}
  特别地，由于二元类型 $\bool$ 等价于 $\unit+\unit$，我们有 $\bfalse\neq\btrue$。
\end{rmk}

这个证明展示了一种描述道路空间的通用方法，我们将会经常使用。
为了刻画一个道路空间，第一步是定义一个比较纤维化 "$\code$"，它能对道路给出更明确的描述。
有几种不同的方法可以证明这样的比较纤维化等价于道路（在 \cref{sec:pi1-s1-intro} 中我们对同一个结果展示了若干不同的证明）。
我们这里所使用的称为\define{编码-解码方法（encode-decode method）}：
\indexdef{encode-decode method}
其核心思想是，为纤维化的所有实例（即作为一个函数 $\prd{x:A+B} \code(x) \to (\inl(a_0)=x)$）普遍地定义 $\decode$，这样我们就可以用道路归纳来分析 $\decode(x,\encode(x,p))$。

\index{transport!in coproduct types}%
像往常一样，我们也可以刻画输运（transport）在余积类型中的作用。
给定一个类型 $X$、一条道路 $p:\id[X]{x_1}{x_2}$，以及两个类型族 $A,B:X\to\type$，我们有
\begin{align*}
  \transfib{A+B}{p}{\inl(a)} &= \inl (\transfib{A}{p}{a}),\\
  \transfib{A+B}{p}{\inr(b)} &= \inr (\transfib{B}{p}{b}),
\end{align*}
其中上标中的 $A+B$ 如通常一样系滥用记号，表示类型族 $x\mapsto A(x)+B(x)$。
证明是简单的道路归纳。

\index{encode-decode method|)}%
\index{type!coproduct|)}%

\section{自然数}
\label{sec:compute-nat}

\index{natural numbers|(}%
\index{encode-decode method|(}%
我们使用编码-解码方法来刻画自然数的道路空间，自然数也是一个正类型。
在这种情况下，我们不固定一个端点，而是直接刻画双侧的道路空间。
因此，相等证明的编码是一个类型族
\[\code:\N\to\N\to\type,\]
通过对 \N 作双重递归定义如下：
\begin{align*}
  \code(0,0) &\defeq \unit\\
  \code(\suc(m),0) &\defeq \emptyt\\
  \code(0,\suc(n)) &\defeq \emptyt\\
  \code(\suc(m),\suc(n)) &\defeq \code(m,n).
\end{align*}
我们还通过递归定义一个依值函数 $r:\prd{n:\N} \code(n,n)$，其中
\begin{align*}
  r(0) &\defeq \ttt\\
  r(\suc(n)) &\defeq r(n).
\end{align*}

\begin{thm}\label{thm:path-nat}
  对所有的 $m,n:\N$，有 $\eqv{(m=n)}{\code(m,n)}$。
\end{thm}

\begin{proof}
  我们通过传输定义 \[ \encode : \prd{m,n:\N} (m=n) \to \code(m,n) \] 为 $\encode(m,n,p) \defeq \transfib{\code(m,{\blank})}{p}{r(m)}$。
  （也可以通过道路归纳直接定义 $\encode$，但基于传输的定义往往使后续计算更为简便。）
  接下来，通过对 $m$, $n$ 的双重归纳定义 \[ \decode : \prd{m,n:\N} \code(m,n) \to (m=n) \]。
  当 $m$ 和 $n$ 均为 $0$ 时，我们需要一个函数 $\unit \to (0=0)$，将其定义为将所有输入映到 $\refl{0}$。
  当 $m$ 是后继而 $n$ 是 $0$，或者反过来时，定义域 $\code(m,n)$ 是 \emptyt，此时 \emptyt 的消去子即够用。
  当二者均为后继时，可将 $\decode(\suc(m),\suc(n))$ 定义为复合
  %
  \begin{narrowmultline*}
    \code(\suc(m),\suc(n))\jdeq\code(m,n)
    \xrightarrow{\decode(m,n)} \narrowbreak
    (m=n)
    \xrightarrow{\apfunc{\suc}}
    (\suc(m)=\suc(n)).
  \end{narrowmultline*}
  %
  接下来我们证明，对所有 $m,n$，$\encode(m,n)$ 与 $\decode(m,n)$ 互为拟逆。

  一方面，若从 $p:m=n$ 出发，则对 $p$ 作归纳后只需证明
  \[\decode(n,n,\encode(n,n,\refl{n}))=\refl{n}.\]
  但由于 $\encode(n,n,\refl{n}) \jdeq r(n)$，这归结为证明 $\decode(n,n,r(n)) =\refl{n}$。
  我们可以对 $n$ 作归纳来证明这一点。
  若 $n$ 判断等于 $0$，则根据 \decode 的定义有 $\decode(0,0,r(0)) =\refl{0}$。
  在后继的情形下，由归纳假设有 $\decode(n,n,r(n)) = \refl{n}$，故只需注意到 $\apfunc{\suc}(\refl{n}) \jdeq \refl{\suc(n)}$。

  另一方面，若从 $c:\code(m,n)$ 出发，则对 $m$ 与 $n$ 作双重归纳来推进。
  若二者均为 $0$，则 $\decode(0,0,c) \jdeq \refl{0}$，而 $\encode(0,0,\refl{0})\jdeq r(0) \jdeq \ttt$。
  因此，只需回顾 \cref{sec:compute-unit} 的结论：$\unit$ 的每个居民都等于 \ttt。
  若 $m$ 为 $0$ 而 $n$ 为后继，或者反过来，则 $c:\emptyt$，故结论自动成立。
  在二者均为后继的情形下，利用归纳假设，有
  \begin{multline*}
    \encode(\suc(m),\suc(n),\decode(\suc(m),\suc(n),c))\\
    \begin{aligned}
    &= \encode(\suc(m),\suc(n),\apfunc{\suc}(\decode(m,n,c)))\\
    &= \transfib{\code(\suc(m),{\blank})}{\apfunc{\suc}(\decode(m,n,c))}{r(\suc(m))}\\
    &= \transfib{\code(\suc(m),\suc({\blank}))}{\decode(m,n,c)}{r(\suc(m))}\\
    &= \transfib{\code(m,{\blank})}{\decode(m,n,c)}{r(m)}\\
    &= \encode(m,n,\decode(m,n,c))\\
    &= c
  \end{aligned}
  \end{multline*}
  （事实上这个证明比所需的更长；参见 \cref{ex:n-set}。）
\end{proof}

特别地，我们有
\begin{equation}\label{eq:zero-not-succ}
  \encode(\suc(m),0) : (\suc(m)=0) \to \emptyt
\end{equation}
它表明"$0$ 不是任何自然数的后继"。
我们还有复合
\begin{narrowmultline}\label{eq:suc-injective}
  (\suc(m)=\suc(n))
  \xrightarrow{\encode} \narrowbreak
  \code(\suc(m),\suc(n))
  \jdeq \code(m,n) \xrightarrow{\decode} (m=n)
\end{narrowmultline}
它表明函数 $\suc$ 是单射。
\index{successor}%

我们将在 \cref{cha:induction,cha:hits} 中研究更一般的正类型。
在 \cref{cha:homotopy} 中，我们将看到，此处用于刻画余积与 \nat 的相等类型的同一技术，也可用于计算球面的同伦群。

\index{encode-decode method|)}%
\index{natural numbers|)}%

\section{示例：结构的相等性}
\label{sec:equality-of-structures}

现在我们考虑一个例子，用以说明类型的群胚结构与类型构造子之间的相互作用。在引言中我们曾指出，泛等性的一个优点在于：两个同构的东西是可互换的，意即凡涉及其中一者的性质或构造，也适用于另一者。常见的"记号滥用"\index{abuse!of notation}在形式上便成为真命题。泛等性本身断言等价的类型是相等的，因而是可互换的——这涵盖了例如将同构的集合视为等同这一常见做法。更进一步，当我们将其他数学对象定义为配备了一定结构或性质的集合乃至一般类型时，便能借助泛等性为它们推演出正确的相等概念。我们将在 \cref{cha:category-theory} 中以一个重要的例子阐明此点：在那里，我们将以适当方式定义范畴论的基本概念，使得范畴的相等即为等价，函子的相等即为自然同构，等等。特别参见 \cref{sec:sip}。在本节中，我们则描述一个来自代数的非常简单的例子。

为简便起见，我们以\emph{半群}为例，其中半群是指配备了一个满足结合律的"乘法"运算的类型。同样的想法也适用于其他代数结构，如幺半群、群和环。
回顾 \cref{sec:sigma-types,sec:pat}，一种数学结构的定义应被理解为将该类结构的类型定义为某个特定的迭代 $\Sigma$ 类型。
对于半群的情形，就得到如下定义。

\begin{defn}
给定一个类型 $A$，以 $A$ 为\define{载体}\index{carrier}的\define{半群结构}
\indexdef{semigroup!structure}%
\index{structure!semigroup}%
\index{associativity!of semigroup operation}%
的类型 \semigroupstr{A} 定义为
\[
\semigroupstr{A} \defeq \sm{m:A \to A \to A} \prd{x,y,z:A} m(x,m(y,z)) = m(m(x,y),z).
\]
%
一个\define{半群}
\indexdef{semigroup}%
即为一个类型连同这样一个结构：
%
\[
\semigroup \defeq \sm{A:\type} \semigroupstr A
\]
\end{defn}

\noindent
在接下来的两个小节中，我们将阐述泛等性使得处理此类半群更为方便的两种方式。

\subsection{提升等价}

\index{lifting!equivalences}%
在非正式的讨论中，我们或许会说，集合 $A$ 与 $B$ 之间的一个双射"显然"能诱导出 $A$ 上的半群结构与 $B$ 上的半群结构之间的一个同构。在泛等性下，这确实是显然的，因为给定类型 $A$ 与 $B$ 间的一个等价，我们能从 $A$ 上的一个半群结构自动导出 $B$ 上的一个半群结构，并且可以证明此导出本身是半群结构间的一个等价。原因在于，\semigroupstrsym 是一个类型族，因此它在类型间的道路上有一个由 $\mathsf{transport}$ 给出的作用：
\[
\transfibf{\semigroupstrsym}{(\ua(e))} : \semigroupstr{A} \to \semigroupstr{B}.
\]
此外，这个映射还是一个等价，因为 $\transfibf{C}(\alpha)$ 总是一个等价，其逆为 $\transfibf{C}{(\opp \alpha)}$，参见 \cref{thm:transport-concat,thm:omg}。

\index{univalence axiom}%
虽然泛等公理确保了该映射的存在，但要计算它具体做了什么，我们需要用到前面章节中证明的关于 $\mathsf{transport}$ 的性质。设 $(m,a)$ 是 $A$ 上的一个半群结构，我们来考察由
\[
\transfib{\semigroupstrsym}{\ua(e)}{(m,a)}.
\]
给出的 $B$ 上的诱导半群结构。首先，因为 \semigroupstr{X} 被定义为一个 $\Sigma$ 类型，根据 \cref{transport-Sigma}，有
\[
\transfib{\semigroupstrsym}{\ua(e)}{(m,a)} = (m',a')
\]
其中 $m'$ 是 $B$ 上诱导出的乘法运算
\begin{flalign*}
& m' : B \to B \to B \\
& m'(b_1,b_2) \defeq \transfib{X \mapsto (X \to X \to X)}{\ua(e)}{m}(b_1,b_2)
\end{flalign*}
而 $a'$ 是 $m'$ 满足结合律的一个诱导证明。
再次根据 \cref{transport-Sigma}，有
\begin{equation}\label{eq:transport-semigroup-step1}
  \begin{aligned}
{}&a' :     \mathsf{Assoc}(B,m')\\
  &a' \defeq \transfib{(X,m) \mapsto \mathsf{Assoc}(X,m)}{(\pairpath(\ua(e),\refl{m'}))}{a},
  \end{aligned}
\end{equation}
其中 $\mathsf{Assoc}(X,m)$ 是类型 $\prd{x,y,z:X} m(x,m(y,z)) = m(m(x,y),z)$。
根据函数外延性，我们只需考察 $m'$ 作用于参数 $b_1,b_2 : B$ 时的行为。连续应用 \eqref{eq:transport-arrow} 两次，可知 $m'(b_1,b_2)$ 等于
%
\begin{narrowmultline*}
  \transfibf{X \mapsto X}\big(
      \ua(e), \narrowbreak
      m(\transfib{X \mapsto X}{\opp{\ua(e)}}{b_1},
        \transfib{X \mapsto X}{\opp{\ua(e)}}{b_2}
       )
   \big).
\end{narrowmultline*}
%
于是，因为 $\ua$ 是 $\transfibf{X\mapsto X}$ 的拟逆，上式等于
\[
e(m(\opp{e}(b_1), \opp{e}(b_2))).
\]
因此，给定 $B$ 中的两个元素，诱导乘法 $m'$ 先用等价 $e$ 将它们送到 $A$，在 $A$ 中将它们相乘，再由 $e$ 将结果带回 $B$，正如所预期的那样。

此外，我们虽不展示证明，但可以计算出 $m'$ 满足结合律的诱导证明（参见 \eqref{eq:transport-semigroup-step1}）等于这样一个函数：它将
$b_1,b_2,b_3 : B$ 映为由以下步骤给出的一条道路：
\begin{equation}
  \label{eq:transport-semigroup-assoc}
  \begin{aligned}
    m'(m'(b_1,b_2),b_3)
    &= e(m(\opp{e}(m'(b_1,b_2)),\opp{e}(b_3))) \\
    &= e(m(\opp{e}(e(m(\opp{e}(b_1),\opp{e}(b_2)))),\opp{e}(b_3))) \\
    &= e(m(m(\opp{e}(b_1),\opp{e}(b_2)),\opp{e}(b_3))) \\
    &= e(m(\opp{e}(b_1),m(\opp{e}(b_2),\opp{e}(b_3)))) \\
    &= e(m(\opp{e}(b_1),\opp{e}(e(m(\opp{e}(b_2),\opp{e}(b_3)))))) \\
    &= e(m(\opp{e}(b_1),\opp{e}(m'(b_2,b_3)))) \\
    &= m'(b_1,m'(b_2,b_3)).
\end{aligned}
\end{equation}
这些步骤用到了 $m$ 满足结合律的证明 $a$ 以及 $e$ 的逆律。从代数的视角看，探究一个运算满足结合律的证明的相等性可能显得奇怪，但如果我们把 $A$ 和 $B$ 视为具有非平凡同伦的一般空间，这就讲得通了。在 \cref{cha:logic} 中，我们将引入\emph{集合}的概念，即只有平凡同伦的类型；如果我们考虑集合上的半群结构，那么任何两个这样的结合律证明都将自动相等。

\subsection{半群的相等性}
\label{sec:equality-semigroups}

利用前面几节讨论的道路空间等式，我们可以探究两个半群何时相等。给定半群 $(A,m,a)$ 和 $(B,m',a')$，根据 \cref{thm:path-sigma}，道路类型
\narrowequation{
  (A,m,a) =_\semigroup (B,m',a')
}
等于由以下序对组成的类型
\begin{align*}
p_1 &: A =_{\type} B \qquad\text{and}\\
p_2 &: \transfib{\semigroupstrsym}{p_1}{(m,a)} = {(m',a')}.
\end{align*}
由泛等性，$p_1$ 是某个等价 $e$ 对应的 $\ua(e)$。由 \cref{thm:path-sigma}、函数外延性以及前文对类型族 $\semigroupstrsym$ 中传输的分析，$p_2$ 等价于一对证明，其中第一个证明表明
\begin{equation} \label{eq:equality-semigroup-mult}
\prd{y_1,y_2:B} e(m(\opp{e}(y_1), \opp{e}(y_2))) = m'(y_1,y_2)
\end{equation}
而第二个证明表明 $a'$ 等于在 \eqref{eq:transport-semigroup-assoc} 中由 $a$ 构造出的诱导结合律证明。但利用逆元的相消，\eqref{eq:equality-semigroup-mult} 等价于
\[
\prd{x_1,x_2:A} e(m(x_1, x_2)) = m'(e(x_1),e(x_2)).
\]
这表明 $e$ 与二元运算交换，意即将 $A$ 中的乘法（即 $m$）映为 $B$ 中的乘法（即 $m'$）。与 $a$ 和 $a'$ 相关的等式也可作类似的重排。因此，半群间的相等恰好由载体类型上的一个等价构成，且该等价与半群结构交换。

对于一般类型，结合律的证明被视为半群结构的一部分。然而，如果我们限制到类集合的类型（再次参见 \cref{cha:logic}），联系 $a$ 与 $a'$ 的等式平凡地成立。此外，在此情形下，集合间的等价恰好就是一个双射。于是，我们便得到了\emph{半群同构}\index{isomorphism!semigroup}的标准定义：载体集合上的一个保持乘法运算的双射。也可以使用范畴论的同构定义，即将\emph{半群同态}\index{homomorphism!semigroup}定义为一个保持乘法的映射，从而得出结论：半群的相等性等同于两个互逆的同态；但我们不在此展示细节，参见 \cref{sec:sip}。

结论是：借助泛等性，半群相等当且仅当它们作为代数结构同构。正如我们将在\cref{sec:sip}中看到的，这一结论可以推广到更一般的情形：在同伦类型论中，数学结构的所有构造都自动保持同构，无需任何繁琐的证明或符号滥用。

\section{泛性质}
\label{sec:universal-properties}

\index{universal!property|(}%
通过组合前几节描述的道路计算规则，我们可以证明，各种类型构造操作都满足预期的泛性质——在同伦视角下，这些性质体现为等价。
例如，给定类型 $X,A,B$，我们有一个函数
\index{type!product}%
\begin{equation}\label{eq:prod-ump-map}
  (X\to A\times B) \to (X\to A)\times (X\to B)
\end{equation}
其定义为 $f \mapsto (\proj1 \circ f, \proj2\circ f)$。

\begin{thm}\label{thm:prod-ump}
  \index{universal!property!of cartesian product}%
  \eqref{eq:prod-ump-map} 是一个等价。
\end{thm}

\begin{proof}
  我们如下定义其拟逆：将 $(g,h)$ 映到 $\lam{x}(g(x),h(x))$。
  （严格来说，这里我们使用了 Cartesian 积 $(X\to A)\times (X\to B)$ 的归纳原理，将其归约为序对的情形。
  今后我们将经常运用这一原理，不再一一说明。）

  现在给定 $f:X\to A\times B$，往返复合得到的函数为
  \begin{equation}
    \lam{x} (\proj1(f(x)),\proj2(f(x))).\label{eq:prod-ump-rt1}
  \end{equation}
  根据 \cref{thm:path-prod}，对任意 $x:X$ 有 $(\proj1(f(x)),\proj2(f(x))) = f(x)$。
  因此，由函数外延性，函数~\eqref{eq:prod-ump-rt1} 与 $f$ 相等。

  另一方面，给定 $(g,h)$，往返复合得到序对 $(\lam{x} g(x),\lam{x} h(x))$。
  根据函数的唯一性原理，它与 $(g,h)$（判断）相等。
\end{proof}

事实上，这一泛性质也有其依值类型版本。
假设给定类型 $X$ 和类型族 $A,B:X\to \type$，
那么我们有函数
\begin{equation}\label{eq:prod-umpd-map}
  \Parens{\prd{x:X} (A(x)\times B(x))} \to \Parens{\prd{x:X} A(x)} \times \Parens{\prd{x:X} B(x)}
\end{equation}
其定义如前，为 $f \mapsto (\proj1 \circ f, \proj2\circ f)$。

\begin{thm}\label{thm:prod-umpd}
  \eqref{eq:prod-umpd-map} 是一个等价。
\end{thm}

\begin{proof}
  留给读者。
\end{proof}

正如 $\Sigma$ 类型是 Cartesian 积的推广，它们也满足此泛性质的一个推广版本。
我们直接给出其依值类型版本：假设有类型 $X$ 和类型族 $A:X\to \type$ 以及 $P:\prd{x:X} A(x)\to\type$。
那么我们有函数
\index{type!dependent pair}%
\begin{equation}
  \label{eq:sigma-ump-map}
  \Parens{\prd{x:X}\dsm{a:A(x)} P(x,a)} \to
  \Parens{\sm{g:\prd{x:X} A(x)} \prd{x:X} P(x,g(x))}.
\end{equation}
注意，如果对某个 $B:X\to\type$ 有 $P(x,a) \defeq B(x)$，则~\eqref{eq:sigma-ump-map} 归约为~\eqref{eq:prod-umpd-map}。

\begin{thm}\label{thm:ttac}
  \index{universal!property!of dependent pair type}%
  \eqref{eq:sigma-ump-map} 是一个等价。
\end{thm}

\begin{proof}
  如前，我们定义拟逆，将 $(g,h)$ 映到函数 $\lam{x} (g(x),h(x))$。
  现在给定 $f:\prd{x:X} \sm{a:A(x)} P(x,a)$，往返复合得到的函数为
  \begin{equation}
    \lam{x} (\proj1(f(x)),\proj2(f(x))).\label{eq:prod-ump-rt2}
  \end{equation}
  对任意 $x:X$，根据 \cref{thm:eta-sigma}（$\Sigma$ 类型的唯一性原理），有
  %
  \begin{equation*}
    (\proj1(f(x)),\proj2(f(x))) = f(x).
  \end{equation*}
  %
  因此，由函数外延性，~\eqref{eq:prod-ump-rt2} 与 $f$ 相等。
  另一方面，给定 $(g,h)$，往返复合得到 $(\lam {x} g(x),\lam{x} h(x))$，如前所述，它与 $(g,h)$ 判断相等。
\end{proof}

\index{axiom!of choice!type-theoretic}
这一点值得注意，因为按命题即类型的解释，~\eqref{eq:sigma-ump-map} 就是"选择公理"。
如果将 $\Sigma$ 读作"存在"，将 $\Pi$（有时）读作"对所有"，则可以如下表述：

\begin{itemize}
\item $\prd{x:X} \sm{a:A(x)} P(x,a)$ 意为"对所有 $x:X$，存在一个 $a:A(x)$ 使得 $P(x,a)$"；
\item $\sm{g:\prd{x:X} A(x)} \prd{x:X} P(x,g(x))$ 意为"存在一个选择函数 $g:\prd{x:X} A(x)$，使得对所有 $x:X$ 有 $P(x,g(x))$"。
\end{itemize}

因此，\cref{thm:ttac} 表明：选择公理不仅"为真"，其前提实际上还等价于其结论。
（另一方面，经典\index{mathematics!classical}数学家可能会觉得~\eqref{eq:sigma-ump-map} 并不符合选择公理的通常含义，因为 $g$ 的值已经给定，不再有任何选择余地。
我们将在\cref{sec:axiom-choice}回到这一点。）

上述关于序对类型的泛性质是针对"映射入"的，这与范畴论中积的概念一致，读者对此应不陌生。
然而，序对类型也有一个针对"映射出"的泛性质，可能看起来不那么熟悉。
在 Cartesian 积的情形，非依值版本直接表达了 Cartesian 闭伴随\index{adjoint!functor}：
\[ \eqvspaced{\big((A\times B) \to C\big)}{\big(A\to (B\to C)\big)}.\]
其依值版本针对类型族 $C:A\times B\to \type$ 表述如下：
\[ \eqvspaced{\Parens{\prd{w:A\times B} C(w)}}{\Parens{\prd{x:A}{y:B} C(x,y)}}. \]
这里从右到左的函数就是 $A\times B$ 的归纳原理，而从左到右的函数则是在序对处求值。
我们留给读者证明它们互为拟逆。
对于 $\Sigma$ 类型也有相应的版本：
\begin{equation}
  \eqvspaced{\Parens{\prd{w:\sm{x:A} B(x)} C(w)}}{\Parens{\prd{x:A}{y:B(x)} C(x,y)}}.\label{eq:sigma-lump}
\end{equation}
同样，从右到左的函数是归纳原理。

另一些归纳原理也属于此类泛性质。
例如，道路归纳是下述等价的自右向左方向：
\index{type!identity}%
\index{universal!property!of identity type}%
\begin{equation}
  \label{eq:path-lump}
  \eqvspaced{\Parens{\prd{x:A}{p:a=x} B(x,p)}}{B(a,\refl a)}
\end{equation}
此式对任意 $a:A$ 和类型族 $B:\prd{x:A} (a=x) \to\type$ 成立。
然而，带有递归的归纳类型，例如自然数，其泛性质则更为复杂；参见 \cref{cha:induction}。

\index{type!limit}%
\index{type!colimit}%
\index{limit!of types}%
\index{colimit!of types}%
既然 \cref{thm:prod-ump} 表达了 Cartesian 积通常的泛性质（在同伦论的意义下），熟悉范畴论的读者自然会想了解类型的其他极限与余极限。
在 \cref{ex:coprod-ump} 中，我们请读者证明余积类型 $A+B$ 也具备预期的泛性质，而 $\unit$（终对象）和 $\emptyt$（始对象）的零元情形则较为容易。
\index{type!empty}%
\index{type!unit}%
\indexsee{initial!type}{type, empty}%
\indexsee{terminal!type}{type, unit}%

\indexdef{pullback}%
对于拉回，预期的显式构造依然可行：给定 $f:A\to C$ 和 $g:B\to C$，定义
\begin{equation}
  A\times_C B \defeq \sm{a:A}{b:B} (f(a)=g(b)).\label{eq:defn-pullback}
\end{equation}
在 \cref{ex:pullback} 中我们请读者验证这一点。
一些更一般的同伦极限可用类似方法构造，但对于余极限，我们将需要新的工具；见 \cref{cha:hits}。

\index{universal!property|)}%

%%%%%%%%%%%%%%%%%%%%%%%%%%%
\sectionNotes

相等类型及其归纳原理的定义源自 Martin-L\"of \cite{Martin-Lof-1973}。
\index{intensional type theory}%
\index{extensional!type theory}%
\index{type theory!intensional}%
\index{type theory!extensional}%
\index{reflection rule}%
正如在 \cref{cha:typetheory} 的注记中提到的，我们这里使用的相等类型属于\emph{内涵}类型论（intensional type theory），而非\emph{外延}类型论（extensional type theory）。
一般来说，若一种相等性概念会根据对象的具体定义方式来区分对象，则称其为"内涵的（intensional）"；若它对具有相同"外延"或"可观察行为"的对象不加区分，则称其为"外延的（extensional）"。
用 Frege（弗雷格）的术语来说，内涵相等性比较的是\emph{涵义（sense）}，而外延相等性仅比较\emph{指称（reference）}。
我们也可以说一种相等性比另一种"更外延"或"更不外延"，分别意味着它考虑的对象内涵方面更少或更多。

\emph{内涵}类型论之所以得名，是因为其\emph{判断相等性（judgmental equality）} $x\jdeq y$ 是一种内涵性很强的相等性：它本质上是在说，在展开函数的定义方程之后，$x$ 和 $y$"具有相同的定义"。
相比之下，命题相等类型 $\id xy$ 则更具外延性，即便是在 \cref{cha:typetheory} 中无公理的内涵类型论中也是如此：例如，我们可以通过归纳证明对所有 $m,n:\N$ 都有 $n+m=m+n$，但不能说对所有 $m,n:\N$ 都有 $n+m\jdeq m+n$，因为加法的\emph{定义}对其参数的处理是不对称的。

我们可以通过添加公理使内涵类型论中的相等类型变得更加外延，例如函数外延性（若两个函数在所有输入上行为相同，则无论其定义方式如何，这两个函数都相等）和泛等性（这可以视为宇宙的一种外延性：若两个类型在所有语境中行为相同，则它们相等）。
函数外延性公理，以及泛等性在命题特殊情形下的形式（"命题外延性"），早已出现在 Russell（罗素）与 Church（丘奇）最早的类型论中。

如前所述，\emph{外延}类型论还包含一条"反射规则（reflection rule）"，说的是若 $p:x=y$，则实际上 $x\jdeq y$。
因此，外延类型论之所以如此命名，正是因为它\emph{不容许}任何纯粹的\emph{内涵}相等性：反射规则迫使判断相等性与更具外延性的相等类型重合。
此外，从反射规则可以推出函数外延性（至少在存在关于函数的判断唯一性原理时）。
然而，反射规则同时也意味着所有高阶群胚结构都会坍缩（参见 \cref{ex:equality-reflection}），因此与泛等公理不相容（参见 \cref{thm:type-is-not-a-set}）。
因此，若将泛等性视为一种外延性质，则可以说内涵类型论所允许的相等类型比外延类型论所允许的更加外延。

\sectionNotes

相等性的对称性（求逆）与传递性（拼接）在类型论中的证明已广为人知。
这些性质使得每个类型（在同伦意义下）成为一个 1-群胚，\cite{hs:gpd-typethy} 利用这一事实给出了类型论的首个"同伦式"语义。

真正的同伦解释——即将相等类型视作道路空间，将类型族视作纤维化——归功于 \cite{AW}，他使用了 Quillen 模型范畴的形式体系。
在（严格）$\infty$-群胚\index{.infinity-groupoid@$\infty$-groupoid}中的一种解释也见于学位论文 \cite{mw:thesis}。
关于在类型论中构造 $\infty$-群胚的\emph{所有}高阶运算及连贯性，可参见~\cite{pll:wkom-type} 和~\cite{bg:type-wkom}。

\index{proof!assistant!Coq@\textsc{Coq}}%
诸如 $\transfib{P}{p}{\blank}$ 与 $\apfunc{f}$ 等操作，以及一种恰当的等价概念，最初由 Voevodsky 使用证明辅助工具 \Coq 在类型论中加以深入研究。
此后，人们又发现了"等价"的其他多种等价定义，在 \cref{cha:equivalences} 中作了比较。

\cref{sec:computational} 中描述的关于相等类型、迁移（transport）等的"计算性"解释，受到了~\cite{lh:canonicity} 的强调。
他们还描述了一种"1-截断的"类型论（见 \cref{cha:hlevels}），其中这些规则均为判断相等。
将此推广到完整的非截断理论的可能性，是当前的研究课题之一。

\index{function extensionality}%
朴素形式的函数外延性——即"如果两个函数逐点相等，那么它们就相等"——是类型论中一条常见的公理，最早可追溯到 \cite{PM2}。
一些更强的函数外延性形式曾在~\cite{garner:depprod} 中被讨论过。
我们所采用的版本——即在等价的意义下将函数类型的相等类型等同起来——最初由 Voevodsky 研究；他还证明了该版本可由朴素形式（以及泛等性；参见 \cref{sec:univalence-implies-funext}）推出。

\index{univalence axiom}%
泛等公理同样归功于 Voevodsky。
它最初的动机源于单纯集合模型中的语义考量；详见~\cite{klv:ssetmodel}。
受群胚模型启发，Hofmann 和 Streicher~\cite{hs:gpd-typethy} 提出了一个类似的公理，称之为"宇宙外延性"。
该公理使用了拟逆~\eqref{eq:qinvtype}，而非一个恰当的"等价"概念；因此，它仅对 1-类型的宇宙才是正确的（且仅在这种情况下才与泛等公理等价）（参见 \cref{defn:1type}）。

在本书所使用的类型论中，函数外延性与泛等性必须作为公理引入，即被断言为属于某个类型的元素，但并非按照该类型的规则构造而成。
这种做法虽尚可行，但也有若干不足。
例如，若能将类型论完全建立在规则之上而非断言公理，其在形式上的性质将更为良好。
另一个不便之处在于，\crefrange{sec:compute-cartprod}{sec:compute-nat} 中的定理仅仅是命题相等性（道路）或等价，因此每当需要经由它们来回转换时，都必须显式提及。
同伦类型论当前的一个研究方向是描述这样一个类型系统：其中的这些规则是\emph{判断}相等，从而一举解决上述两个问题。
迄今为止，这仅在若干简单情形下得以实现，不过像~\cite{lh:canonicity} 这样的初步结果也前景看好。
此外，还有其他潜在的途径可以将泛等性与函数外延性引入类型论，例如引入足够强大的"高阶商"或"高阶归纳-递归类型"概念。

\crefrange{sec:compute-coprod}{sec:compute-nat} 中诸如"$\inl$ 与 $\inr$ 是单的且不相交"这类简单结论，在类型论中是众所周知的；而构造 \encode 函数则是证明它们的常用方法。
我们所描述的、更精炼的方法——即在等价的意义下刻画正类型的整个相等类型——则属于较新的发展；例如可参见~\cite{ls:pi1s1}。

\index{axiom!of choice!type-theoretic}%
类型论选择公理~\eqref{eq:sigma-ump-map} 由 William Howard 在其关于"命题即类型"对应的原始论文~\cite{howard:pat} 中所注意到，其后 Martin-L\"of 在引入其依值类型论时对此作了进一步研究。在 Bourbaki 的集合论中，它被提及为一条"分配律"~\cite{Bourbaki}。
\index{Bourbaki}%

关于拉回以及更一般的同伦极限在同伦类型论中更全面（且形式化）的讨论，请参阅~\cite{AKL13}。
\index{limit!of types}有向图\index{graph}上的图表的极限是最容易形式化的一类一般极限；而范畴（或更一般地，$(\infty,1)$-范畴）上的图表则存在一个问题：
\index{.infinity1-category@$(\infty,1)$-category}%
\indexsee{category!.infinity1-@$(\infty,1)$-}{$(\infty,1)$-category}%
一般而言，（同伦连贯的）图表\index{diagram}这一概念会涉及无穷多个连贯条件。
解决这个问题是同伦类型论中一个重要的开放问题\index{open!problem}。

\sectionExercises

\begin{ex}\label{ex:basics:concat}
  证明 \cref{lem:concat} 的三种显然证法两两相等。
\end{ex}

\begin{ex}\label{ex:eq-proofs-commute}
  证明前一练习中所构造的三个证明之间的相等性构成一个交换三角形。
  换言之，若将拼接的三种定义分别记作 $(p \mathbin{\ct_1} q)$、$(p\mathbin{\ct_2} q)$ 和 $(p\mathbin{\ct_3} q)$，则由拼接得到的相等性
  \[(p\mathbin{\ct_1} q) = (p\mathbin{\ct_2} q) = (p\mathbin{\ct_3} q)\]
  等于相等性 $(p\mathbin{\ct_1} q) = (p\mathbin{\ct_3} q)$。
\end{ex}

\begin{ex}\label{ex:fourth-concat}
  给出 \cref{lem:concat} 的第四种不同的证明，并证明它与其他证明相等。
\end{ex}

\begin{ex}\label{ex:npaths}
  通过对 $n$ 归纳，同时定义类型 $A$ 中 \define{$n$-维道路（$n$-dimensional path）}\index{path!n-@$n$-} 的一般概念，以及此类道路的边界类型。
\end{ex}

\begin{ex}\label{ex:ap-to-apd-equiv-apd-to-ap}
  证明函数~\eqref{eq:ap-to-apd} 与~\eqref{eq:apd-to-ap} 互为逆等价。
  % and that they take $\apfunc f(p)$ to $\apdfunc f (p)$ and vice versa. (that was \cref{thm:apd-const})
\end{ex}

\begin{ex}\label{ex:equiv-concat}
  证明若 $p:x=y$，则函数 $(p\ct \blank):(y=z) \to (x=z)$ 是一个等价。
\end{ex}

\begin{ex}\label{ex:ap-sigma}
  陈述并证明 \cref{thm:ap-prod} 从 Cartesian 积到 $\Sigma$-类型的一个推广形式。
\end{ex}

\begin{ex}\label{ex:ap-coprod}
  陈述并证明 \cref{thm:ap-prod} 关于余积的一个类似结论。
\end{ex}

\begin{ex}\label{ex:coprod-ump}
  \index{universal!property!of coproduct}%
  证明余积具有所期望的泛性质，
  \[ \eqv{(A+B \to X)}{(A\to X)\times (B\to X)}. \]
  你能将其推广为涉及依值函数的等价吗？
\end{ex}

\begin{ex}\label{ex:sigma-assoc}
  证明 $\Sigma$-类型是"结合"的，
  \index{associativity!of Sigma-types@of $\Sigma$-types}%
  即对任意 $A:\UU$ 以及类型族 $B:A\to\UU$ 和 $C:(\sm{x:A} B(x))\to\UU$，有
  \[\eqvspaced{\Parens{\sm{x:A}{y:B(x)} C(\pairr{x,y})}}{\Parens{\sm{p:\sm{x:A}B(x)} C(p)}}. \]
\end{ex}

\begin{ex}\label{ex:pullback}
  一个（同伦）\define{交换方块（commutative square）}
  \indexdef{commutative!square}%
  \begin{equation*}
  \vcenter{\xymatrix{
      P\ar[r]^h\ar[d]_k &
      A\ar[d]^f\\
      B\ar[r]_g &
      C
      }}
  \end{equation*}
  由如图所示的函数 $f$、$g$、$h$、$k$ 以及一条道路 $f \circ h= g \circ k$ 组成。
  注意，这恰好是~\eqref{eq:defn-pullback} 中所定义的拉回 $(P\to A) \times_{P\to C} (P\to B)$ 中的一个元素。
  一个交换方块称为（同伦）\define{拉回方块}
  \indexdef{pullback}%
  如果对任意 $X$，诱导映射
  \[ (X\to P) \to (X\to A) \times_{(X\to C)} (X\to B) \]
  是一个等价。
  证明在~\eqref{eq:defn-pullback} 中定义的拉回 $P \defeq A\times_C B$ 构成一个拉回方块的角。
\end{ex}

\begin{ex}\label{ex:pullback-pasting}
  假设给定两个交换方块
  \begin{equation*}
    \vcenter{\xymatrix{
        A\ar[r]\ar[d] &
        C\ar[r]\ar[d] &
        E\ar[d]\\
        B\ar[r] &
        D\ar[r] &
        F
      }}
  \end{equation*}
  且右边的方块是拉回方块。证明：左边的方块是拉回方块当且仅当外层的矩形是拉回方块。
\end{ex}

\begin{ex}\label{ex:eqvboolbool}
  证明 $\eqv{(\eqv\bool\bool)}{\bool}$。
\end{ex}

\begin{ex}\label{ex:equality-reflection}
  假设我们在类型论中加入\emph{相等反射规则（equality reflection rule）}，其规定为：若存在元素 $p:x=y$，则实际上 $x\jdeq y$。
  证明对任意 $p:x=x$ 有 $p\jdeq \refl{x}$。
  （这意味着每个类型在 \cref{sec:basics-sets} 将要引入的意义下都是一个\emph{集合}；参见 \cref{sec:hedberg}。）
\end{ex}

\begin{ex}\label{ex:strengthen-transport-is-ap}
  证明在不使用函数外延性的前提下，可将 \cref{thm:transport-is-ap} 加强为
  \[\transfib{B}{p}{{\blank}} =_{B(x)\to B(y)} \idtoeqv(\apfunc{B}(p))\]
  （在此类及其他类似情形中，之所以选择表面上较弱的表述，是出于可读性和一致性的考虑。）
\end{ex}

\begin{ex}\label{ex:strong-from-weak-funext}
  假设我们拥有的不是函数外延性（\cref{axiom:funext}），而仅仅是存在一个元素
  \[ \funext : \prd{A:\UU}{B:A\to \UU}{f,g:\prd{x:A} B(x)} (f\htpy g) \to (f=g) \]
  （不假定它与 $\happly$ 有何关联）。
  证明实际上仅此一项就足以推出完整的函数外延性公理（即 $\happly$ 是一个等价）。
  此结论归功于 Voevodsky；其证明较为棘手，可能需要用到后面章节的概念。
\end{ex}

\begin{ex}\label{ex:equiv-functor-types}\
  \begin{enumerate}
  \item 证明如果 $\eqv{A}{A'}$ 且 $\eqv{B}{B'}$，则 $\eqv{(A\times B)}{(A'\times B')}$。
  \item 给出该结论的两种证明，一种使用泛等性，另一种不使用，并证明这两个证明相等。
  \item 对其他类型形成子：$\Sigma$、$\to$、$\Pi$ 和 $+$，陈述并证明类似的结果。
\end{enumerate}
\end{ex}

\begin{ex}\label{ex:dep-htpy-natural}
  陈述并证明 \cref{lem:htpy-natural} 关于依值函数的一个版本。
\end{ex}

\begin{ex}\label{ex:sigma-eq-components-neq}
  我们已经看到 $\bfalse \neq \btrue$（\cref{rmk:true-neq-false}）。证明：尽管如此，$(\bool, \bfalse) =_{\sm{A:\type}A} (\bool, \btrue)$（参见 \cref{rmk:sigma-equality-extraction}）。
\end{ex}

% Local Variables:
% TeX-master: "hott-online"
% End:

\chapter{集合与逻辑}
\label{cha:logic}

类型论，无论形式的还是非形式的，都是一套处理类型及其元素的规则。
然而，当用自然语言进行非形式的数学写作时，我们通常会使用一些熟悉的词汇，特别是``与''``或''等逻辑连接词，以及``对所有''``存在''等逻辑量词。
与集合论不同，类型论提供了不止一种方式来将这些英语短语理解为对类型的操作。
这种潜在的歧义需要加以消除，途径是制定局部或全局的约定、为非形式数学引入新的标注，或两者兼用。
这需要一定的适应过程，但好处在于：由于类型论允许对逻辑进行更精细的分析，我们能够比在集合论基础下更忠实地表达数学，减少``滥用语言''的情况。
本章将解释其中涉及的问题，并论证我们所做选择的合理性。

\section{集合与 \texorpdfstring{$n$}{n}-类型}
\label{sec:basics-sets}

\index{set|(defstyle}%

为了解释类型论的逻辑与集合论的逻辑之间的联系，在类型论中引入``集合''的概念是有益的。
一般而言，类型的行为类似于空间或高阶群胚，但其中有一个子类，其行为更接近于传统集合论系统中的集合。
从范畴论的角度看，我们可以考虑\emph{离散}群胚，它们由一组对象确定，且仅有恒同态射作为高阶态射；从拓扑学的角度看，我们可以考虑具有离散拓扑的空间。\index{discrete!space}
更一般地，我们可以考虑与这类结构\emph{等价}的群胚或空间；由于类型论中的一切都是在同伦意义下进行的，我们无法分辨它们的差异。

直觉上，如果一个类型不包含高阶同伦信息，我们就会预期它在这种意义下``是一个集合''：任何两条平行道路都是（在同伦意义下）相等的，所有更高维度的平行道路亦然。
幸运的是，由于同伦类型论中的一切构造都自动具有``函子性''/``连续性''，
\index{continuity of functions in type theory@``continuity'' of functions in type theory}%
\index{functoriality of functions in type theory@``functoriality'' of functions in type theory}%
实际上只需在最底层提出这一要求就足够了。

\begin{defn}\label{defn:set}
  一个类型 $A$ 是一个\define{集合（set）}，如果对于所有 $x,y:A$ 和所有 $p,q:x=y$，都有 $p=q$。
\end{defn}

更精确地说，命题 $\isset(A)$ 被定义为如下类型：
\[ \isset(A) \defeq \prd{x,y:A}{p,q:x=y} (p=q). \]

如 \cref{sec:types-vs-sets} 所述，同伦类型论中的集合不同于 ZF 集合论中的集合，因为不存在全局的``属于''谓词 $\in$。
它们更接近于结构数学和范畴论中使用的集合，其元素是``抽象的点''，我们通过函数和关系来赋予其结构。
这足以使它们成为大多数基于集合的数学的基础系统；我们将在 \cref{cha:set-math} 中看到一些例子。

哪些类型是集合？
在 \cref{cha:hlevels} 中我们将深入研究这个问题更一般的形式，但目前我们可以先看一些简单的例子。

\begin{eg}\label{eg:isset-unit}
  类型 \unit 是一个集合。
  因为根据 \cref{thm:path-unit}，对于任意 $x,y:\unit$，类型 $(x=y)$ 等价于 \unit。
  由于 \unit 的任意两个元素都相等，这意味着 $x=y$ 的任意两个元素也相等。
\end{eg}

\begin{eg}\label{eg:isset-empty}
  类型 $\emptyt$ 是一个集合，因为给定任意 $x,y:\emptyt$，根据 $\emptyt$ 的归纳原理，我们可以推导出任何结论。
\end{eg}

\begin{eg}\label{thm:nat-set}
  自然数类型 \nat 也是一个集合。
  这由 \cref{thm:path-nat} 可得，因为所有相等类型 $\id[\nat]xy$ 都等价于 \unit 或 $\emptyt$，而 \unit 或 $\emptyt$ 的任意两个实例都是相等的。
  我们将在 \cref{cha:hlevels} 中看到该事实的另一个证明。
\end{eg}

我们迄今考虑的大多数类型形成操作也都保持集合性质。

\begin{eg}\label{thm:isset-prod}
  如果 $A$ 和 $B$ 是集合，那么 $A\times B$ 也是集合。
  给定 $x,y:A\times B$ 和 $p,q:x=y$，根据 \cref{thm:path-prod} 有 $p= \pairpath(\projpath1(p),\projpath2(p))$ 和 $q= \pairpath(\projpath1(q),\projpath2(q))$。
  但由于 $A$ 是集合，可知 $\projpath1(p)=\projpath1(q)$；同理，因为 $B$ 是集合，可知 $\projpath2(p)=\projpath2(q)$。因此 $p=q$。

  类似地，如果 $A$ 是集合且 $B:A\to\type$ 满足每个 $B(x)$ 都是集合，那么 $\sm{x:A} B(x)$ 也是集合。
\end{eg}

\begin{eg}\label{thm:isset-forall}
  如果 $A$ 是\emph{任意}类型，且 $B:A\to \type$ 满足每个 $B(x)$ 都是集合，那么类型 $\prd{x:A} B(x)$ 也是一个集合。
  为此假设 $f,g:\prd{x:A} B(x)$ 且 $p,q:f=g$。
  根据函数外延性，我们有
  %
  \begin{equation*}
    p = {\funext (x \mapsto \happly(p,x))}
    \quad\text{and}\quad
    q = {\funext (x \mapsto \happly(q,x))}.
  \end{equation*}
  %
  但对任意 $x:A$，我们有
  %
  \begin{equation*}
   \happly(p,x):f(x)=g(x)
   \qquad\text{and}\qquad
   \happly(q,x):f(x)=g(x),
  \end{equation*}
  %
  由于 $B(x)$ 是集合，我们得到 $\happly(p,x) = \happly(q,x)$。
  现在再次使用函数外延性，依值函数 $(x \mapsto \happly(p,x))$ 与 $(x \mapsto \happly(q,x))$ 相等，因而（应用 $\apfunc{\funext}$）$p$ 与~$q$ 也相等。
\end{eg}

更多例子，参见 \cref{ex:isset-coprod,ex:isset-sigma}。关于同伦类型论中所有集合构成的子系统（范畴）的更系统研究，请参见~\cref{cha:set-math}。

\index{n-type@$n$-type|(}%

集合仅仅是所谓\emph{同伦 $n$-类型}阶梯上的第一级。
下一级由\emph{$1$-类型}构成，它们在范畴论中类似于 $1$-群胚。
集合（我们也可称之为\emph{$0$-类型}）的定义性质是其中没有非平凡的道路。
类似地，$1$-类型的定义性质是在道路之间没有非平凡的道路：

\begin{defn}\label{defn:1type}
  一个类型 $A$ 是\define{1-类型（1-type）}
  \indexdef{1-type}%
  当且仅当对所有 $x,y:A$ 以及 $p,q:x=y$ 和 $r,s:p=q$，都有 $r=s$。
\end{defn}

类似地，我们可以定义 $2$-类型、$3$-类型等等。
我们将在 \cref{cha:hlevels} 中归纳地定义 $n$-类型的一般概念，并研究不同 $n$ 值对应的 $n$-类型之间的关系。

然而，目前先了解两个事实是有益的。
第一，这些层级是向上封闭的：如果 $A$ 是 $n$-类型，那么 $A$ 也是 $(n+1)$-类型。
例如：
%%  will make precise the sense in which this
%% ``suffices for all higher levels'', but as an example, we observe that
%% it suffices for the next level up.

\begin{lem}\label{thm:isset-is1type}
  如果 $A$ 是集合（即 $\isset(A)$ 有实例），那么 $A$ 是 $1$-类型。
\end{lem}

\begin{proof}
  假设 $f:\isset(A)$；则对任意 $x,y:A$ 和 $p,q:x=y$，我们有 $f(x,y,p,q):p=q$。
  固定 $x$，$y$ 和 $p$，并定义 $g: \prd{q:x=y} (p=q)$ 为 $g(q) \defeq f (x,y,p,q)$。
  那么对任意 $r:q=q'$，我们有 $\apdfunc{g}(r) : \trans{r}{g(q)} = g(q')$。
  因此，由 \cref{cor:transport-path-prepost}，我们得到 $g(q) \ct r = g(q')$。

  特别地，假设给定如 \cref{defn:1type} 所示的 $x,y,p,q$ 和 $r,s:p=q$，并按上述定义 $g$。
  则 $g(p) \ct r = g(q)$ 且 $g(p) \ct s = g(q)$，故由消去律可得 $r=s$。
\end{proof}

第二，这种按层级对类型的划分并非退化，也就是说并非所有类型都是集合：

\begin{eg}\label{thm:type-is-not-a-set}
  \index{type!universe}%
  宇宙 \type 不是集合。
  为证明这一点，只需举出一个类型 $A$ 和一条道路 $p:A=A$，使得 $p$ 不等于 $\refl A$。
  取 $A=\bool$，并令 $f:A\to A$ 定义为 $f(\bfalse)\defeq \btrue$ 和 $f(\btrue)\defeq \bfalse$。
  那么对所有 $x$（通过简单的分情形讨论）都有 $f(f(x))=x$，因此 $f$ 是一个等价。
  于是，由泛等性，$f$ 给出了一条道路 $p:A=A$。

  如果 $p$ 等于 $\refl A$，那么（再次由泛等性）$f$ 将等于 $A$ 上的恒等函数。
  但这将意味着 $\bfalse=\btrue$，与 \cref{rmk:true-neq-false} 矛盾。
\end{eg}

在 \cref{cha:hits,cha:homotopy} 中我们将证明，对任意 $n$，都存在不是 $n$-类型的类型。

注意，$A$ 是 $1$-类型恰好意味着对任意 $x,y:A$，相等类型 $\id[A]xy$ 是集合。
（因此，\cref{thm:isset-is1type} 可以等价地理解为：集合的相等类型也是集合。）
这将是我们将在 \cref{cha:hlevels} 中给出的 $n$-类型递归定义的基础。

我们也可以将这一刻画从集合``向下''延伸。
也就是说，一个类型 $A$ 是集合当且仅当对任意 $x,y:A$，$\id[A]xy$ 中的任意两个元素都相等。
由于集合也就是 $0$-类型，因此自然称一个具有后一种性质（其任意两个元素相等）的类型为\emph{$(-1)$-类型}。
这类类型可视为\emph{狭义下的命题}，对它们的研究正是通常所谓的``逻辑''；这将占据本章剩余部分的内容。

\index{n-type@$n$-type|)}%
\index{set|)}%

\section{命题即类型？}
\label{subsec:pat?}

\index{proposition!as types|(}%
\index{logic!constructive vs classical|(}%
\index{anger|(}%
到目前为止，我们一直遵循着\cref{sec:pat}中描述的直接的``命题即类型''哲学。
按照这种哲学，像``存在 $x:A$ 使得 $P(x)$''这样的自然语言短语，被解释为对应的类型，例如 $\sm{x:A} P(x)$；
而一个命题的证明，则被视为判定某个特定元素居留于该类型。
然而，我们也已经看到，由这种解读所产生的``逻辑''，在某些方面会让熟悉经典数学的数学家感到陌生。
例如，在\cref{thm:ttac}中我们看到，形如
\index{axiom!of choice!type-theoretic}%
\begin{equation}\label{eq:english-ac}
  \parbox{\textwidth-2cm}{``If for all $x:X$ there exists an $a:A(x)$ such that $P(x,a)$, then there exists a function $g:\prd{x:X} A(x)$ such that for all $x:X$ we have $P(x,g(x))$'',}
\end{equation}
的陈述——它看上去很像经典的\index{mathematics!classical}\emph{选择公理（axiom of choice）}——在这种解读下总是真的。
这是命题即类型逻辑一个值得注意且通常很有用的特点，但它同时也表明了这种逻辑与逻辑的经典解释之间存在显著差异：
在经典解释下，选择公理并非逻辑真理，而是一个额外的``公理''。

另一方面，我们现在也可以证明，那些形如经典\emph{双重否定律（law of double negation）}和\emph{排中律（law of excluded middle）}的对应陈述，与泛等公理是不相容的。
\index{univalence axiom}%

\begin{thm}\label{thm:not-dneg}
  \index{double negation, law of}%
  并非对于所有 $A:\UU$ 都有 $\neg(\neg A) \to A$。
\end{thm}

\begin{proof}
  回顾 $\neg A \jdeq (A\to\emptyt)$。
  我们也将``并非……''读作算子 $\neg$。
  因此，为了证明这个陈述，只需假设给定某个$f:\prd{A:\UU} (\neg\neg A \to A)$，然后构造 \emptyt 的一个元素即可。

  接下来的证明思路是注意到，$f$ 和类型论中的任何函数一样，是``连续的''。
  \index{continuity of functions in type theory@``continuity'' of functions in type theory}%
  \index{functoriality of functions in type theory@``functoriality'' of functions in type theory}%
  由泛等公理可知，这意味着 $f$ 对于类型的等价是\emph{自然的}。
  由此，再加上一个无不动点的自等价\index{automorphism!fixed-point-free}，我们就能导出矛盾。

  令 $e:\eqv\bool\bool$ 为由 $e(\btrue)\defeq\bfalse$ 和 $e(\bfalse)\defeq\btrue$ 定义的等价，如\cref{thm:type-is-not-a-set}所示。
  令 $p:\bool=\bool$ 为由 $e$ 经泛等公理给出的对应道路，即 $p\defeq \ua(e)$。
  那么我们有 $f(\bool) : \neg\neg\bool \to\bool$，并且
  \[\apd f p : \transfib{A\mapsto (\neg\neg A \to A)}{p}{f(\bool)} = f(\bool).\]
  因此，对于任意 $u:\neg\neg\bool$，我们有
  \[\happly(\apd f p,u) : \transfib{A\mapsto (\neg\neg A \to A)}{p}{f(\bool)}(u) = f(\bool)(u).\]

  现在根据\eqref{eq:transport-arrow}，将 $f(\bool):\neg\neg\bool\to\bool$ 沿着 $p$ 在类型族${A\mapsto (\neg\neg A \to A)}$中传输，等于如下函数：
  它先将参数沿着 $\opp p$ 在类型族 $A\mapsto \neg\neg A$ 中传输，然后应用 $f(\bool)$，最后再将结果沿着 $p$ 在类型族 $A\mapsto A$ 中传输：
  \begin{narrowmultline*}
    \transfib{A\mapsto (\neg\neg A \to A)}{p}{f(\bool)}(u) =
    \narrowbreak
    \transfib{A\mapsto A}{p}{f(\bool) (\transfib{A\mapsto \neg\neg
        A}{\opp{p}}{u})}.
  \end{narrowmultline*}
  %
  然而，由函数外延性可知，任意两点 $u,v:\neg\neg\bool$ 都是相等的，因为对于任意 $x:\neg\bool$，我们有 $u(x):\emptyt$，从而可以推导出任何结论，特别是 $u(x)=v(x)$。
  因此，我们有$\transfib{A\mapsto \neg\neg A}{\opp{p}}{u} = u$，于是从 $\happly(\apd f p,u)$ 我们得到一个等式
  \[ \transfib{A\mapsto A}{p}{f(\bool)(u)} = f(\bool)(u).\]
  最后，如\cref{sec:compute-universe}所讨论的，沿着道路 $p\jdeq \ua(e)$ 在类型族 $A\mapsto A$ 中传输，等价于应用等价 $e$；因此我们有
  \begin{equation}
    e(f(\bool)(u)) = f(\bool)(u).\label{eq:fpaut}
  \end{equation}
  %
  然而，我们也可以证明
  \begin{equation}
    \prd{x:\bool} \neg(e(x)=x).\label{eq:fpfaut}
  \end{equation}
  这可以通过对 $x$ 进行分情形讨论得到：两种情形都能直接从 $e$ 的定义以及 $\bfalse\neq\btrue$（\cref{rmk:true-neq-false}）推出。
  因此，将\eqref{eq:fpfaut}应用于 $f(\bool)(u)$ 和\eqref{eq:fpaut}，我们就得到了 $\emptyt$ 的一个元素。
\end{proof}

\begin{rmk}
  \index{choice operator}%
  \indexsee{operator!choice}{choice operator}%
  特别地，这意味着不可能存在一个Hilbert 风格的``选择算子''，能够为每个非空类型选出一个元素。
  关键在于，任何这样的算子都不可能是\emph{自然的}，而在泛等公理下，所有作用于类型的函数对于等价都必须是自然的。
\end{rmk}

\begin{rmk}
  然而，对于任何 $A$，$\neg\neg\neg A \to \neg A$ 仍然成立；参见\cref{ex:neg-ldn}。
\end{rmk}

\begin{cor}\label{thm:not-lem}
  \index{excluded middle}%
  并非对于所有 $A:\UU$ 都有 $A+(\neg A)$。
\end{cor}

\begin{proof}
  假设我们有 $g:\prd{A:\UU} (A+(\neg A))$。
  我们将证明此时有 $\prd{A:\UU} (\neg\neg A \to A)$，从而可以应用 \cref{thm:not-dneg}。
  于是假设 $A:\UU$ 且 $u:\neg\neg A$；我们想要构造 $A$ 的一个元素。

  根据 $g(A):A+(\neg A)$，通过分情形讨论，我们可以假设要么 $g(A)\jdeq \inl(a)$ 对某个 $a:A$ 成立，要么 $g(A)\jdeq \inr(w)$ 对某个 $w:\neg A$ 成立。
  在第一种情形中，我们有 $a:A$；在第二种情形中，我们有 $u(w):\emptyt$，从而可以得出我们想要的任何结果（比如 $A$）。
  因此，在两种情形下我们都得到了 $A$ 的一个元素，如所需。
\end{proof}

因此，如果我们想假定泛等公理（我们当然想这样做），同时仍保留经典\index{mathematics!classical}推理的选项（这也是可取的），我们就不能使用未经修改的命题即类型原则来将\emph{所有}非形式的数学陈述解释到类型论中，因为那样排中律将为假。
然而，我们也不想完全抛弃命题即类型，因为它有许多优良性质（例如简单性、构造性以及可计算性）。
我们现在讨论一种对命题即类型的修改方案，以解决这些问题；在 \cref{subsec:when-trunc} 中，我们将回到何时使用哪种逻辑的问题。
\index{anger|)}%
\index{proposition!as types|)}%
\index{logic!constructive vs classical|)}%

\section{命题}
\label{subsec:hprops}

\index{logic!of mere propositions|(}%
\index{mere proposition|(defstyle}%
\indexsee{proposition!mere}{mere proposition}%
我们已经看到，命题即类型逻辑既有优点也有缺点。
两者有一个共同的原因：当类型被视作命题时，它们可以包含比单纯的真或假更多的信息，并且所有对它们的``逻辑''构造都必须尊重这些额外信息。
这表明，我们可以通过将注意力限制在那些\emph{不}包含任何超出真值的信息的类型上，并仅将这些类型视为逻辑命题，从而获得一种更传统的逻辑。

这样的类型 $A$ 如果有实例则视为``真''，如果从其实例可导出矛盾（即如果 $\neg A \jdeq (A\to\emptyt)$ 是有实例的）则视为``假''。
\index{inhabited type}%
为了获得更传统的逻辑，我们希望避免将这样的类型当作逻辑命题：给出其一个元素会提供比仅知道该类型有实例更多的信息。
例如，如果给定 \bool 的一个元素，那么我们得到的信息就比仅仅知道 \bool 包含某个元素要多。
事实上，我们正好多得到了\emph{一比特}\index{bit}的信息：我们知道所给的是 \bool 的\emph{哪一个}元素。
相比之下，如果给定 \unit 的一个元素，那么我们得到的信息不比仅知道 \unit 包含一个元素更多，因为 \unit 的任意两个元素都彼此相等。
这启发我们做出如下定义。

\begin{defn}\label{defn:isprop}
  一个类型 $P$ 称为\define{命题（mere proposition）}，如果对于所有 $x,y:P$ 都有 $x=y$。
\end{defn}

注意，由于我们仍然\emph{在}类型论中做数学，所以这也是一个\emph{在}类型论中的定义，这意味着它本身是一个类型——或者更准确地说，是一个类型族。
具体来说，对于任意 $P:\type$，类型 $\isprop(P)$ 定义为
\[ \isprop(P) \defeq \prd{x,y:P} (x=y). \]
因此，断言``$P$ 是一个命题''就是要给出 $\isprop(P)$ 的一个实例，即一个依值函数，它通过道路将 $P$ 的任意两个元素连接起来。
这个函数的连续性/自然性意味着，不仅 $P$ 的任意两个元素相等，而且 $P$ 也不包含任何高阶同伦。

\begin{lem}\label{thm:inhabprop-eqvunit}
  如果 $P$ 是一个命题且 $x_0:P$，则有 $\eqv P \unit$。
\end{lem}

\begin{proof}
  定义 $f:P\to\unit$ 为 $f(x)\defeq \ttt$，定义 $g:\unit\to P$ 为 $g(u)\defeq x_0$。
  结论由下一个引理以及 \unit 是一个命题（根据 \cref{thm:path-unit}）这一观察得出。
\end{proof}

\begin{lem}\label{lem:equiv-iff-hprop}
  如果 $P$ 和 $Q$ 都是命题，且 $P\to Q$ 与 $Q\to P$ 成立，那么 $\eqv P Q$。
\end{lem}

\begin{proof}
  假设给定 $f:P\to Q$ 和 $g:Q\to P$。
  那么对于任意 $x:P$，由于 $P$ 是命题，我们有 $g(f(x))=x$。
  类似地，对于任意 $y:Q$，由于 $Q$ 是命题，我们有 $f(g(y))=y$；因此 $f$ 和 $g$ 互为拟逆。
\end{proof}

也就是说，正如在 \cref{sec:pat} 中所承诺的，如果两个命题是逻辑等价的，那么它们就是等价的。

在同伦论中，一个与 \unit 同伦等价的空间被称为\emph{可缩的}。
因此，任何有实例的命题都是可缩的（另见\cref{sec:contractibility}）。
另一方面，空类型 \emptyt 也是一个（平凡地）命题。
至少在经典\index{mathematics!classical}数学中，只有这两种可能性。

命题也称为\emph{子终对象}（从范畴论角度）、\emph{子单点集}（从集合论角度），或\emph{h-命题}。
\indexsee{object!subterminal}{mere proposition}%
\indexsee{subterminal object}{mere proposition}%
\indexsee{subsingleton}{mere proposition}%
\indexsee{h-proposition}{mere proposition}
\cref{sec:basics-sets}中的讨论提示我们也可以称其为\emph{$(-1)$-类型}；我们将在\cref{cha:hlevels}回到这一点。
形容词``mere''强调的是：尽管任何类型都可以被视作一个命题（我们通过给出它的一个实例来做到这一点），但一个仅仅是命题的类型，不能被有用地视作比命题\emph{更多}的东西：在其为真的见证中不包含任何额外信息。

注意，一个类型 $A$ 是集合当且仅当对任意 $x,y:A$，相等类型 $\id[A]xy$ 都是一个命题。
另一方面，通过复制并简化\cref{thm:isset-is1type}的证明，我们得到：

\begin{lem}\label{thm:prop-set}
  每个命题都是集合。
\end{lem}

\begin{proof}
  假设 $f:\isprop(A)$；于是对任意 $x,y:A$ 有 $f(x,y):x=y$。固定 $x:A$ 并定义$g(y)\defeq f(x,y)$。那么对任意 $y,z:A$ 和 $p:y=z$，我们有$\apd
  g p : \trans{p}{g(y)}={g(z)}$。因此由\cref{cor:transport-path-prepost}，我们得到$g(y)\ct p = g(z)$，这就是说$p=\opp{g(y)}\ct g(z)$。于是，对任意 $p,q:x=y$，我们有$p = \opp{g(x)}\ct g(y) = q$。
\end{proof}

特别地，这意味着：

\begin{lem}\label{thm:isprop-isprop}\label{thm:isprop-isset}
  对任意类型 $A$，类型 $\isprop(A)$ 和 $\isset(A)$ 都是命题。
\end{lem}

\begin{proof}
  假设 $f,g:\isprop(A)$。由函数外延性，要证明 $f=g$，只需证明对任意 $x,y:A$ 有$f(x,y)=g(x,y)$。而 $f(x,y)$ 和 $g(x,y)$ 都是 $A$ 中的道路，它们之所以相等，是因为由 $f$ 或 $g$ 可知 $A$ 是一个命题，从而由\cref{thm:prop-set}可知它是一个集合。
  类似地，假设 $f,g:\isset(A)$，即对任意 $a,b:A$ 和 $p,q:a=b$，我们有$f(a,b,p,q):p=q$和$g(a,b,p,q):p=q$。但由于 $A$ 是一个集合（由 $f$ 或 $g$ 可知），故而是一个 $1$-类型，由此可得$f(a,b,p,q)=g(a,b,p,q)$；于是由函数外延性有 $f=g$。
\end{proof}

到目前为止我们还见过另一个例子：\cref{sec:basics-equivalences}中的条件~\ref{item:be3}断言，对任意函数 $f$，类型 $\isequiv (f)$ 应该是一个命题。

\index{logic!of mere propositions|)}%
\index{mere proposition|)}%

\section{经典逻辑与直觉主义逻辑}
\label{sec:intuitionism}

\index{logic!constructive vs classical|(}%
\index{denial|(}%
有了命题的概念，我们现在可以给出同伦类型论中\define{排中律（law of excluded middle）}
\indexdef{excluded middle}%
\indexsee{axiom!excluded middle}{excluded middle}%
\indexsee{law!of excluded middle}{excluded middle}%
的正确定义：
\begin{equation}
  \label{eq:lem}
  \LEM{}\;\defeq\;
  \prd{A:\UU} \Big(\isprop(A) \to (A + \neg A)\Big).
\end{equation}
类似地，\define{双重否定律（law of double negation）}
\indexdef{double negation, law of}%
\indexdef{axiom!double negation}%
\indexdef{law!of double negation}%
是
\begin{equation}
  \label{eq:ldn}
  % \mathsf{DN}\;\defeq\;
  \prd{A:\UU} \Big(\isprop(A) \to (\neg\neg A \to A)\Big).
\end{equation}
二者也很容易看出是等价的——参见\cref{ex:lem-ldn}——因此从现在起我们通常只提 \LEM{}。

\LEM{} 的这一表述避免了\cref{thm:not-dneg,thm:not-lem}中的``悖论''，因为 \bool 不是一个命题。
为了将其与更一般的``命题即类型''表述区分开来，我们将后者重命名为：
\symlabel{lem-infty}
\begin{equation*}
  \LEM\infty \defeq \prd{A:\UU} (A + \neg A).
\end{equation*}
为强调起见，正确的版本~\eqref{eq:lem}
可以记作 $\LEM{-1}$；
另见\cref{ex:lemnm}。
虽然 $\LEM{}$
不是\cref{cha:typetheory}中所描述的基础类型论的推论，但可以一致地将其假定为一公理（不像它的 $\infty$ 版本）。
例如，我们将在\cref{sec:wellorderings}中假设它。

然而，不使用 \LEM{} 也能走很远，这一点可能会令人惊讶。
通常，只需对定义或定理稍加改写，就能避免诉诸排中律。
虽然这有时需要一点时间适应，但这番周折往往是值得的，能带来更优雅、更一般的证明。
我们在引论中已经讨论了这样做的一些益处。

例如，在经典\index{mathematics!classical}数学中，双重否定常被不必要地使用。
一个非常简单的例子是通常的假设：集合 $A$ 是“非空的”，其字面意思是“$A$ \emph{不}包含\emph{任何}元素”这一情形不成立。
而几乎总是，其真正含义是 $A$ \emph{确实}包含至少一个元素这一正面断言；去掉这层双重否定，我们就能让陈述减少对 \LEM{} 的依赖。
回想一下，当我们把 $A$ 本身作为命题来断言时（即构造出 $A$ 的一个元素，通常不命名），我们就说类型 $A$ 是\emph{有实例的（inhabited）}
\index{inhabited type}%
。
因此，在将经典证明转译为构造逻辑时，我们常常将“非空的”替换为“有实例的”（尽管有时必须替换为“仅知有实例的”；参见\cref{subsec:prop-trunc}）。

类似地，在经典数学中，不必要的反证法也并不罕见。
\index{proof!by contradiction}%
当然，经典反证法的形式依赖双重否定律：我们先假设 $\neg A$ 并推导出矛盾，从而得到 $\neg \neg A$，再由双重否定律得到 $A$。
然而，从 $\neg A$ 推导出矛盾的过程，常常稍作改写即可给出 $A$ 的一个直接证明，从而无需使用 \LEM{}。

同样需要注意的是，如果目标是证明一个\emph{否定}\index{negation}，那么“反证法”并不涉及 \LEM{}。
事实上，$\neg A$ 根据定义就是类型 $A\to\emptyt$，所以证明 $\neg A$ 无非是在假设 $A$ 的前提下证明矛盾（$\emptyt$）。
类似地，双重否定律对否定命题确实成立：$\neg\neg\neg A \to \neg A$。
通过练习，我们会逐渐学会更仔细地区分否定命题与非否定命题，并注意到何时用到了 \LEM{}、何时没有。

因此，与表面看来相反，“构造性地”做数学通常并不意味着放弃重要定理，而是要找到表述定义的最佳方式，使重要定理能够构造性地证明。
也就是说，我们在初步研究一个主题时可以自由使用 \LEM{}，但一旦对该主题有了更深入的理解，就有望改进其定义和证明，从而不再依赖那条公理。
% For instance, the theory of ordinal numbers, which classically makes heavy use of \LEM{}, works quite well constructively once we choose the correct definition of ``ordinal''; see \cref{sec:ordinals}.
这一观察在同伦类型论中尤为显著：泛等性和高阶归纳类型这些强大工具使我们能够构造性地应对许多传统上需要经典\index{mathematics!classical}推理的问题。
我们将在\cref{part:mathematics}中看到若干这样的例子。
% For instance, none of the ``synthetic'' homotopy theory we will develop in \cref{cha:homotopy} requires \LEM{} --- despite the fact that classical homotopy theory (formulated using topological spaces or simplicial sets) makes heavy use of them (as well as the axiom of choice).

还值得提及的是，即便在构造性数学中，排中律对\emph{某些}命题也可以成立。
这类命题传统上称为\emph{可判定的（decidable）}。

\begin{defn}\label{defn:decidable-equality}
  \mbox{}
  \begin{enumerate}
  \item 若 $A+\neg A$，则称类型 $A$ 为\define{可判定的（decidable）}
    \indexdef{decidable!type}%
    \indexdef{type!decidable}%
    。
  \item 类似地，若 \narrowequation{\prd{a:A} (B(a)+\neg B(a))}， 则称类型族 $B:A\to \type$ 是\define{可判定的（decidable）}
    \indexdef{decidable!type family}%
    \indexdef{type!family of!decidable}%
    。 \label{item:decidable-equality2}
  \item 特别地，若 \narrowequation{\prd{a,b:A} ((a=b) + \neg(a=b))}， 则称 $A$ 具有\define{可判定相等性（decidable equality）}
    \indexdef{decidable!equality}%
    \indexsee{equality!decidable}{decidable equality}%
    。
  \end{enumerate}
\end{defn}

因此，$\LEM{}$ 恰好断言了所有命题都是可判定的，从而所有命题族也是可判定的。
特别地，$\LEM{}$ 蕴含所有集合（在\cref{sec:basics-sets}的意义下）都具有可判定相等性。
这种意义上的可判定相等性是非常强的性质；参见\cref{thm:hedberg}。

\index{denial|)}%
\index{logic!constructive vs classical|)}%

\section{子集与命题大小重整}
\label{subsec:prop-subsets}

\index{mere proposition|(}%

作为命题有用性的另一个例子，我们来讨论子集（以及更一般的子类型）。
假设 $P:A\to\type$ 是一个类型族，且每个类型 $P(x)$ 都被视为一个命题。
那么 $P$ 本身就是 $A$ 上的一个\emph{谓词（predicate）}，或者说 $A$ 的元素的一个\emph{性质（property）}。

在集合论中，给定集合$A$上的一个谓词$P$，我们便可以构造子集$\setof{x\in A | P(x)}$。
如\cref{sec:pat}中所简要提及的，在类型论中与之最直接对应的类比便是$\Sigma$-类型 $\sm{x:A} P(x)$。
当然，$\sm{x:A} P(x)$的一个实例是一个对子$(x,p)$，其中$x:A$且$p$是$P(x)$的一个证明。
然而，对于一般的$P$，一个元素$a:A$可能会产生$\sm{x:A} P(x)$中多个不同的元素——如果命题$P(a)$有多个不同的证明。
这就有悖于我们通常对子集的直觉。
但倘若$P$是一个\emph{纯}命题（mere proposition），那么这种情况便不会发生。

\begin{lem}\label{thm:path-subset}
  设$P:A\to\type$是一个类型族，且对任意$x:A$，$P(x)$均为一个命题。
  若$u,v:\sm{x:A} P(x)$满足$\proj1(u) = \proj1(v)$，则有$u=v$。
\end{lem}

\begin{proof}
  假设$p:\proj1(u) = \proj1(v)$。
  根据\cref{thm:path-sigma}，为证明$u=v$，只需证$\trans{p}{\proj2(u)} = \proj2(v)$。
  然而，$\trans{p}{\proj2(u)}$与$\proj2(v)$都是$P(\proj1(v))$的元素，而$P(\proj1(v))$是命题；因此它们是相等的。
\end{proof}

举例来说，回忆我们在\cref{sec:basics-equivalences}中定义了
\[(\eqv A B) \;\defeq\; \sm{f:A\to B} \isequiv (f),\]
其中每个类型$\isequiv (f)$都被假定为一个命题。
由此可知，若两个等价的底层函数相等，则它们作为等价而言也是相等的。

\label{defn:setof}%
此后，若$P:A\to \type$是一个由命题组成的族（即每个$P(x)$都是命题），我们可以将
%
\begin{equation}
  \label{eq:subset}
  \setof{x:A | P(x)}
\end{equation}
%
作为$\sm{x:A} P(x)$的另一种记法。
（在技术上，没有理由不能将此记法也用于一般的$P$，但这种用法可能会因带来意外的隐含联想而引起混淆。）
若$A$是一个集合，我们称\eqref{eq:subset}为$A$的一个\define{子集}
\indexdef{subset}%
\indexdef{type!subset}%
；对于一般的$A$，我们或可称其为\define{子类型}。
\indexdef{subtype}%
我们也可以直接将$P$本身称作$A$的一个\emph{子集}或\emph{子类型}；这实际上更为正确，因为孤立地看类型~\eqref{eq:subset}并不记得它与$A$的关系。

\symlabel{membership}
对于这样的$P$和$a:A$，我们可以记作$a\in P$或$a\in \setof{x:A | P(x)}$以指代命题$P(a)$。
若它成立，则我们称$a$是$P$的一个\define{元素}。
\symlabel{subset}
类似地，若$\setof{x:A | Q(x)}$是$A$的另一个子集，则说$P$\define{被包含}
\indexdef{containment!of subsets}%
\indexdef{inclusion!of subsets}%
于$Q$中，并记作$P\subseteq Q$，若有$\prd{x:A}(P(x)\rightarrow Q(x))$。

作为子类型的更多例子，我们可以定义某个宇宙\UU{}中所有集合与所有命题构成的“子宇宙”：
\symlabel{setU}\symlabel{propU}
\begin{align*}
  \setU &\defeq \setof{A:\UU | \isset(A) },\\
  \propU &\defeq \setof{A:\UU | \isprop(A) }.
\end{align*}
$\setU$的一个元素是一个类型$A:\UU$连同证据$s:\isset(A)$，$\propU$亦同理。
\cref{thm:path-subset}蕴含着$\id[\setU]{(A,s)}{(B,t)}$等价于$\id[\UU]AB$（因而也等价于$\eqv AB$）。
因此，我们经常会滥用记号，简写$A:\setU$而非$(A,s):\setU$。
若无需指定所讨论的具体宇宙，我们也可以省略下标\UU。

回忆一下，对于任意两个宇宙$\UU_i$和$\UU_{i+1}$，若$A:\UU_i$则亦有$A:\UU_{i+1}$。
因此，对于任意$(A,s):\set_{\UU_i}$，我们也有$(A,s):\set_{\UU_{i+1}}$，对$\prop_{\UU_i}$亦类似，从而得到自然的映射
\begin{align}
  \set_{\UU_i} &\to \set_{\UU_{i+1}}\label{eq:set-up},\\
  \prop_{\UU_i} &\to \prop_{\UU_{i+1}}.\label{eq:prop-up}
\end{align}
映射~\eqref{eq:set-up}不可能是一个等价，否则我们将重现康托尔集合论中为人熟知的自指悖论\index{paradox}。
然而，尽管在目前所给出的类型论中，~\eqref{eq:prop-up}并非自动成为一个等价，但假定其为一个等价是一致的。
也就是说，我们可以考虑向类型论中添加如下公理。

\begin{axiom}[命题大小重整]
  \indexsee{axiom!propositional resizing}{propositional resizing}%
  \indexdef{propositional!resizing}%
  \indexsee{resizing!propositional}{propositional resizing}%
  映射$\prop_{\UU_i} \to \prop_{\UU_{i+1}}$是一个等价。
\end{axiom}

我们将此公理称为\define{命题大小重整}，
因为它意味着宇宙$\UU_{i+1}$中的任何命题均可被“重整大小”为较小宇宙$\UU_i$中与之等价的一个命题。
若$\UU_{i+1}$满足\LEM{}，则此公理自动成立（参见\cref{ex:lem-impred}）。
我们一般不假定此公理成立，不过在有些地方会将其用作明确的假设。
它是针对命题的一种\emph{非直谓性}形式，而避免使用它，则称类型论保持\emph{直谓的}。
\indexsee{impredicativity!for mere propositions}{propositional resizing}%
\indexdef{mathematics!predicative}%

在实践中，我们最常需要的是一个略有不同的陈述：所考虑的宇宙 \UU 包含一个能``囊括所有命题''的类型。
换言之，我们需要一个类型 $\Omega:\UU$ 以及一族以 $\Omega$ 为指标的命题，使得该族在等价的意义下包含所有命题。
如果 \UU 不是最小的宇宙 $\UU_0$，那么上述命题大小重整即可推出此陈述，因为此时我们可以定义 $\Omega\defeq \prop_{\UU_0}$。

非直谓性的一个用途是定义幂集。
将集合 $A$ 的\define{幂集（power set）}\indexdef{power set}定义为 $A\to\propU$ 是很自然的；但在缺乏非直谓性的情况下，这个定义（哪怕仅就等价而言）也依赖于宇宙 \UU 的选取。
然而，有了命题大小重整，我们就可以将幂集定义为
\symlabel{powerset}%
\[ \power A \defeq (A\to\Omega),\]
这样它就不依赖于 $\UU$ 了。
另见 \cref{subsec:piw}。

\section{命题的逻辑}
\label{subsec:logic-hprop}

\index{logic!of mere propositions|(}%
我们曾在 \cref{sec:types-vs-sets} 提到，与仅有一个基本概念（类型）的类型论不同，集合论基础有两个基本概念：集合与命题。
因此，一位经典\index{mathematics!classical}数学家习惯于分开处理这两种对象。

在类型论中也可以恢复类似的二分法：由那些是\emph{命题}的类型（及类型族）来扮演集合论中命题的角色。
在许多情况下，只需将相应的类型形成子限制到命题上，就可以在这种逻辑中表示出逻辑连接词和量词。
当然，这要求我们知道所讨论的类型形成子能保持命题性。

\begin{eg}
  如果 $P$ 和 $Q$ 是命题，那么 $P\times Q$ 也是命题。
  这很容易证明，只需利用积类型中道路的刻画，类似于 \cref{thm:isset-prod} 但更简单。
  因此，连接词``且''保持命题性。
\end{eg}

\begin{eg}\label{thm:isprop-forall}
  如果 $A$ 是任意类型，且 $P:A\to \type$ 满足对一切 $x:A$，类型 $P(x)$ 都是命题，那么 $\prd{x:A} P(x)$ 也是命题。
  证明思路类似于 \cref{thm:isset-forall} 但更简单：给定 $f,g:\prd{x:A} P(x)$，对任意 $x:A$，由于 $P(x)$ 是命题，我们有 $f(x)=g(x)$。
  然后根据函数外延性，即得 $f=g$。

  特别地，如果 $P$ 是命题，那么无论 $A$ 是什么，$A\to P$ 也是命题。
  更进一步，由于 \emptyt 是命题，所以 $\neg A \jdeq (A\to\emptyt)$ 也是命题。
  \index{quantifier!universal}%
  因此，连接词``蕴含''和``非''保持命题性，量词``对所有''也是如此。
\end{eg}

另一方面，并非所有类型形成子都能保持命题。
即使 $P$ 和 $Q$ 是命题， $P+Q$ 在一般情况下也未必是。
例如， \unit 是一个命题，但 $\bool=\unit+\unit$ 不是。
从逻辑上讲， $P+Q$ 是一种``纯粹构造性''的``或''：它的一个见证会额外包含\emph{哪个}析取支为真的信息。
有时这非常有用，但如果我们想要一种更经典的、能保持命题的``或''，就需要一种方法来``截断''这个类型，通过遗忘这些额外信息将其变为一个命题。

\index{quantifier!existential}%
$\Sigma$-类型 $\sm{x:A} P(x)$（其中 $A$ 是任意类型）也存在同样的问题。
这是``存在一个 $x:A$ 使得 $P(x)$''的纯粹构造性解释，它记住了见证 $x$，因此即使每个类型 $P(x)$ 都是命题，该类型本身通常也不是。
（回想一下，我们在 \cref{subsec:prop-subsets} 中曾指出， $\sm{x:A} P(x)$ 也可以看作``满足 $P(x)$ 的那些 $x:A$ 构成的子集''。）

\section{命题截断}
\label{subsec:prop-trunc}

\index{truncation!propositional|(defstyle}%
\indexsee{type!squash}{truncation, propositional}%
\indexsee{squash type}{truncation, propositional}%
\indexsee{bracket type}{truncation, propositional}%
\indexsee{type!bracket}{truncation, propositional}%
\emph{命题截断（propositional truncation）}，也称为\emph{$(-1)$-截断}、\emph{括号类型（bracket type）}或\emph{挤压类型（squash type）}，是一个额外的类型形成子，它能将一个类型``挤压''或``截断''为一个命题，从而忘掉该类型中居民所包含的、除其存在以外的所有信息。

更准确地说，对于任意类型 $A$，都存在一个类型 $\brck{A}$。
它有两个构造子：

\begin{itemize}
\item 对于任意 $a:A$，我们有 $\bproj a : \brck A$。
\item 对于任意 $x,y:\brck A$，我们有 $x=y$。
\end{itemize}

第一个构造子意味着，如果 $A$ 是有实例的，那么 $\brck A$ 也是有实例的。
第二个构造子确保 $\brck A$ 是一个命题；通常我们不给这个事实的见证命名。

\index{recursion principle!for truncation}%
$\brck A$ 的递归原理如下：

\begin{itemize}
\item 如果 $B$ 是一个命题，并且我们有 $f:A\to B$，那么就存在一个诱导出的 $g:\brck A \to B$，使得对所有 $a:A$，都有 $g(\bproj a) \jdeq f(a)$。
\end{itemize}

换句话说，任何（从 $A$ 的有实例性）能推出的命题，也都能从 $\brck A$ 推出。
因此，作为命题的 $\brck A$，所包含的信息并不比 $A$ 的有实例性更多。
（$\brck A$ 也有一个归纳原理，但它不是特别有用；参见 \cref{ex:prop-trunc-ind}。）

在 \cref{ex:lem-brck,ex:impred-brck,sec:hittruncations} 中，我们将描述如何用更一般的概念来构造 $\brck{A}$。
目前，我们只是将其作为 \cref{cha:typetheory} 中那些规则之外的一条附加规则来使用。

有了命题截断，我们就可以将``命题的逻辑''扩展到涵盖析取与存在量词。
具体来说，$\brck{A+B}$ 是``$A$ 或 $B$''的命题版本，它并不``记住''哪个析取支为真的信息。

截断的递归原理意味着，\emph{当我们试图证明一个命题时}，仍然可以在 $\brck{A+B}$ 上进行分情形讨论。
也就是说，假设我们有 $u:\brck{A+B}$，并且我们试图证明一个命题 $Q$。
换句话说，我们试图定义 $\brck{A+B} \to Q$ 的一个元素。
由于 $Q$ 是一个命题，根据命题截断的递归原理，我们只需构造一个函数 $A+B\to Q$ 即可。
于是现在我们可以对 $A+B$ 进行分情形讨论了。

类似地，对于类型族 $P:A\to\type$，我们可以考虑 $\brck{\sm{x:A} P(x)}$，它是``存在 $x:A$ 使得 $P(x)$''的命题版本。
与析取的情形一样，通过结合截断与 $\Sigma$ 类型的归纳原理，如果我们有一个类型为 $\brck{\sm{x:A} P(x)}$ 的假设，那么\emph{当试图证明一个命题时}，我们可以引入新的假设 $x:A$ 和 $y:P(x)$。
换句话说，如果我们知道存在某个 $x:A$ 使得 $P(x)$，但我们手头并没有一个具体的这样的 $x$，那么我们仍然可以自由地使用这样的 $x$，只要我们不是在试图构造任何可能依赖于 $x$ 的具体取值的对象。
要求陪域是一个命题就表达了结果对见证的独立性，因为这样一个类型的所有可能的实例都必须相等。

在\cref{cha:real-numbers,cha:set-math}所讨论的集合层面的数学中，我们主要处理集合与命题，此时使用传统的逻辑记号来专门指称``命题截断逻辑''会很方便。

\begin{defn} \label{defn:logical-notation}
  我们通过截断来定义\define{传统逻辑记号}
  \indexdef{implication}%
  \indexdef{traditional logical notation}%
  \indexdef{logical notation, traditional}%
  \index{quantifier}%
  \indexsee{existential quantifier}{quantifier, existential}%
  \index{quantifier!existential}%
  \indexsee{universal!quantifier}{quantifier, universal}%
  \index{quantifier!universal}%
  \indexdef{conjunction}%
  \indexdef{disjunction}%
  \indexdef{true}%
  \indexdef{false}%
  ，其中 $P$ 和 $Q$ 表示命题（或命题族）：
  {\allowdisplaybreaks
  \begin{align*}
    \top            &\ \defeq \ \unit \\
    \bot            &\ \defeq \ \emptyt \\
    P \land Q       &\ \defeq \ P \times Q \\
    P \Rightarrow Q &\ \defeq \ P \to Q \\
    P \Leftrightarrow Q &\ \defeq \ P = Q \\
    \neg P          &\ \defeq \ P \to \emptyt \\
    P \lor Q        &\ \defeq \ \brck{P + Q} \\
    \fall{x : A} P(x) &\ \defeq \ \prd{x : A} P(x) \\
    \exis{x : A} P(x) &\ \defeq \ \Brck{\sm{x : A} P(x)}
  \end{align*}}
\end{defn}

符号 $\land$ 和 $\lor$ 在同伦论中也用于表示带点空间的缩积与楔，我们将在\cref{cha:hits}引入它们。严格来说这造成了潜在的冲突，但通常不会引起混淆。

类似地，当如\cref{subsec:prop-subsets}那样讨论子集时，我们也会使用传统的记号来表示交、并、补：
\indexdef{intersection!of subsets}%
\symlabel{intersection}%
\indexdef{union!of subsets}%
\symlabel{union}%
\indexdef{complement, of a subset}%
\symlabel{complement}%
\begin{align*}
  \setof{x:A | P(x)} \cap \setof{x:A | Q(x)}
  &\defeq \setof{x:A | P(x) \land Q(x)},\\
  \setof{x:A | P(x)} \cup \setof{x:A | Q(x)}
  &\defeq \setof{x:A | P(x) \lor Q(x)},\\
  A \setminus \setof{x:A | P(x)}
  &\defeq \setof{x:A | \neg P(x)}.
\end{align*}
当然，在没有排中律的情况下，后者并非通常意义上的``补集''：对于 $A$ 的任意子集 $B$，我们未必有 $B \cup (A\setminus B) = A$。

\index{truncation!propositional|)}%
\index{mere proposition|)}%
\index{logic!of mere propositions|)}%

\section{选择公理}
\label{sec:axiom-choice}

\index{axiom!of choice|(defstyle}%
\index{denial|(}%
现在我们可以恰当表述同伦类型论中的选择公理了。
假设给定一个类型 $X$ 以及类型族
%
\begin{equation*}
  A:X\to\type
  \qquad\text{and}\qquad
  P:\prd{x:X} A(x)\to\type,
\end{equation*}
%
并且满足以下条件：

\begin{itemize}
\item $X$ 是一个集合；
\item 对每个 $x:X$，$A(x)$ 是一个集合；
\item 对每个 $x:X$ 和 $a:A(x)$，$P(x,a)$ 是一个命题。
\end{itemize}

\define{选择公理（axiom of choice）}
$\choice{}$ 断言，在如下假设之下，
\begin{equation}\label{eq:ac}
  \Parens{\prd{x:X} \Brck{\sm{a:A(x)} P(x,a)}}
  \to
  \Brck{\sm{g:\prd{x:X} A(x)} \prd{x:X} P(x,g(x))}.
\end{equation}
当然，这是~\eqref{eq:english-ac} 的直接翻译，其中我们把``存在 $x:A$ 使得 $B(x)$''读作 $\brck{\sm{x:A}B(x)}$，所以也可以用熟悉的逻辑记号将这一陈述写作
\begin{narrowmultline*}
  \textstyle
  \Big(\fall{x:X}\exis{a:A(x)} P(x,a)\Big)
  \Rightarrow \narrowbreak
  \Big(\exis{g : \prd{x:X} A(x)} \fall{x : X} P(x,g(x))\Big).
\end{narrowmultline*}
%
特别要注意，命题截断在这里出现了两次。
定义域中的截断意味着我们假设：对于每个 $x$，都存在某个 $a:A(x)$ 使得 $P(x,a)$ 成立，但这些值并未以任何已知方式被选定或指定。
值域中的截断则意味着我们得出结论：存在某个函数 $g$，但这个函数并未以任何已知方式被确定或指定。

事实上，根据 \cref{thm:ttac}，这个公理也可以用一个更简单的形式来表达。

\begin{lem}\label{thm:ac-epis-split}
  选择公理~\eqref{eq:ac} 等价于下述陈述：对于任意集合 $X$ 以及任意的 $Y:X\to\type$，如果每个 $Y(x)$ 都是集合，则有
  \begin{equation}
    \Parens{\prd{x:X} \Brck{Y(x)}}
    \to
    \Brck{\prd{x:X} Y(x)}.\label{eq:epis-split}
  \end{equation}
\end{lem}

这对应于经典\index{mathematics!classical}选择公理的一个众所周知的等价形式，即``一族非空集合的 Cartesian 积是非空的''。

\begin{proof}
  由 \cref{thm:ttac}，~\eqref{eq:ac} 的值域等价于
  \[\Brck{\prd{x:X} \sm{a:A(x)} P(x,a)}.\]
  因此，~\eqref{eq:ac} 等价于将~\eqref{eq:epis-split} 应用于以下情形：\narrowequation{Y(x) \defeq \sm{a:A(x)} P(x,a).}
  （根据 \cref{thm:isset-prod,thm:prop-set}，这是一个集合。）
  反过来，~\eqref{eq:epis-split} 等价于将~\eqref{eq:ac} 应用于 $A(x)\defeq Y(x)$ 和 $P(x,a)\defeq\unit$ 的情形。
  这样，二者逻辑等价。
  由于两者都是命题，根据 \cref{lem:equiv-iff-hprop}，它们作为类型也是等价的。
\end{proof}

与 \LEM{} 类似，等价形式~\eqref{eq:ac} 和~\eqref{eq:epis-split} 并非我们基本类型论的推论，但可以相容地将其作为公理来假设。

\begin{rmk}
  不难证明，~\eqref{eq:epis-split} 的右边总是蕴含左边。
  由于两边都是命题，根据 \cref{lem:equiv-iff-hprop}，选择公理也就等价于要求一个等价
  \[ \eqv{\Parens{\prd{x:X} \Brck{Y(x)}}}{\Brck{\prd{x:X} Y(x)}} \]
  这揭示了一个常见的误区：虽然依值函数类型保持命题性（\cref{thm:isprop-forall}），但它们并不与截断交换：$\brck{\prd{x:A} P(x)}$ 通常并不等价于 $\prd{x:A} \brck{P(x)}$。
  如果我们假设选择公理，它告诉我们这对\emph{集合}是成立的；而正如我们下面将会看到的，在一般情况下它并不成立。
\end{rmk}

选择公理中对类型必须是集合的限制，可以在一定程度上放宽。
例如，我们可以允许~\eqref{eq:ac} 中的 $A$ 和 $P$，或者~\eqref{eq:epis-split} 中的 $Y$，是任意的类型族；这样会得到一个看似更强但仍然相容的陈述。
我们也可以用将在 \cref{cha:hlevels} 中考虑的更一般的 $n$-截断来替代命题截断，从而得到一族公理 $\choice n$，它们在~\eqref{eq:ac} 与~\cref{thm:ttac} 之间插值；前者简称为 \choice{}（为强调起见也可记为 $\choice{-1}$），后者我们将称为 $\choice\infty$。
亦可参见 \cref{ex:acnm,ex:acconn}。
但要注意，我们不能放宽 $X$ 必须是集合这一要求。

\begin{lem}\label{thm:no-higher-ac}
  存在一个类型 $X$ 和一个类型族 $Y:X\to \type$，使得每个 $Y(x)$ 都是集合，但条件~\eqref{eq:epis-split} 为假。
\end{lem}

\begin{proof}
  定义 $X\defeq \sm{A:\type} \brck{\bool = A}$，并设 $x_0 \defeq (\bool, \bproj{\refl{\bool}}) : X$。
  根据 $\Sigma$ 类型中道路的刻画、$\brck{A=\bool}$ 是命题这一事实以及泛等性，对任意 $(A,p),(B,q):X$，我们有 $\eqv{(\id[X]{(A,p)}{(B,q)})}{(\eqv AB)}$。
  特别地，$\eqv{(\id[X]{x_0}{x_0})}{(\eqv \bool\bool)}$，因此正如 \cref{thm:type-is-not-a-set} 中所述，$X$ 不是一个集合。

  另一方面，如果 $(A,p):X$，那么 $A$ 是一个集合。这一结论可以通过对 $p:\brck{\bool=A}$ 施行截断归纳，并利用 $\bool$ 是集合的事实得出。
  由于只要 $A$ 和 $B$ 是集合，$\eqv A B$ 就也是集合，因此对任意 $x_1,x_2:X$，$\id[X]{x_1}{x_2}$ 都是集合，也就是说 $X$ 是一个 1-类型。
  特别地，如果我们通过 $Y(x) \defeq (x_0=x)$ 来定义 $Y:X\to\UU$，那么每个 $Y(x)$ 都是集合。

  现在根据定义，对于任意 $(A,p):X$，我们有 $\brck{\bool=A}$，从而有 $\brck{x_0 = (A,p)}$。
  因此，我们得到 $\prd{x:X} \brck{Y(x)}$。
  如果条件~\eqref{eq:epis-split} 对这个 $X$ 和 $Y$ 成立，那么我们将有 $\brck{\prd{x:X} Y(x)}$。
  由于我们正试图推导出一个矛盾（$\emptyt$），而矛盾本身是一个命题，因此我们可以假设 $\prd{x:X} Y(x)$，即 $\prd{x:X} (x_0=x)$。
  但这意味着 $X$ 是一个命题，因而也是一个集合，从而产生了矛盾。
\end{proof}

\index{denial|)}%
\index{axiom!of choice|)}%

\section{唯一选择原理}
\label{sec:unique-choice}

\index{unique!choice|(defstyle}%
\indexsee{axiom!of choice!unique}{unique choice}%

以下观察虽然平凡，但非常有用。

\begin{lem}\label{thm:prop-equiv-trunc}
  如果 $P$ 是一个命题，那么 $\eqv P {\brck P}$。
\end{lem}

\begin{proof}
  当然，由定义我们有 $P\to \brck{P}$。
  并且由于 $P$ 是一个命题，将 $\brck P$ 的泛性质应用于 $\idfunc[P] :P\to P$，便得到 $\brck P \to P$。
  根据 \cref{lem:equiv-iff-hprop}，这两个函数互为拟逆。
\end{proof}

其重要推论之一如下。

\begin{cor}[唯一选择原理]\label{cor:UC}
  假设给定一个类型族 $P:A\to \type$，满足
  \begin{enumerate}
  \item 对每个 $x$，类型 $P(x)$ 都是一个命题，并且
  \item 对每个 $x$，我们有 $\brck {P(x)}$。
  \end{enumerate}
  那么我们可以得到 $\prd{x:A} P(x)$。
\end{cor}

\begin{proof}
  由上述两条假设及前一引理立即可得。
\end{proof}

这条推论还体现了一种非常有用的推理技巧。
具体来说，假设我们知道 $\brck A$，并且希望利用这一点来构造另一个类型 $B$ 的元素。
我们自然希望在构造 $B$ 的元素时能使用 $A$ 的元素，但这仅在 $B$ 是命题时才被允许，这样我们才能应用命题截断 $\brck A$ 的归纳原理；在一般情形下，我们最多只能期望证明 $\brck B$。
%
另一种做法是，我们可以为 $B$ 补充额外的数据，来\emph{唯一地}刻画我们希望构造的对象。
确切地说，我们定义一个谓词 $Q:B\to\type$，使得 $\sm{x:B} Q(x)$ 是一个命题。
这样，从 $A$ 的一个元素可以构造出满足 $Q(b)$ 的元素 $b:B$，因此从 $\brck A$ 出发可以构造 $\brck{\sm{x:B} Q(x)}$；又因为 $\brck{\sm{x:B} Q(x)}$ 等价于 $\sm{x:B} Q(x)$，即可从中投影出 $B$ 的一个元素。
这方面的例子可参见 \cref{ex:decidable-choice}。

集合论数学中也会出现类似的问题，尽管其表现方式略有不同。
如果我们试图定义一个函数 $f: A \to B$，并且对于给定的 $a : A$，我们只能证明存在某个 $b : B$，那么我们的工作尚未完成，因为我们需要具体指定一个 $B$ 的元素，而不仅仅是证明其存在性。
当然，一种方案是像我们在类型论中所做的那样，细化论证以证明 $b:B$ 的唯一存在性。
但在集合论中，通过应用选择公理——由它替我们选取所需的元素——往往可以更简单地回避这一问题。
然而在同伦类型论中，且不论是否有意避免选择公理，可用的选择公理形式本身适用性就更低，因为它们要求选择的作用域是一个\emph{集合}。
因此，如果 $A$ 不是一个集合（例如可能是宇宙 $\UU$），就不存在一致的选择公理形式，能让我们在定义 $f(a)$ 时，为每个 $a : A$ 简单地选取一个 $B$ 的元素来使用。

\index{unique!choice|)}%

\section{命题何时被截断？}
\label{subsec:when-trunc}

\index{logic!of mere propositions|(}%
\index{mere proposition|(}%
\index{logic!truncated}%

乍看之下，截断版本的 $+$ 和 $\Sigma$ 似乎比未截断的版本更接近``或''与``存在''在非形式数学中的含义。
当然，它们更接近作为形式集合论基础的一阶逻辑中\index{first-order!logic}``或''与``存在''的\emph{精确}含义，因为后者并不试图保留命题为真的任何见证。
然而，一个可能令人惊讶的事实是：非形式数学的实践往往更准确地由未截断的形式来描述。

\index{prime number}%
例如，考虑这样一个陈述：``每个素数要么是 $2$，要么是奇数''。
做研究的数学家在使用这一事实时不会感到任何不妥：他们不仅用它来证明关于素数的\emph{定理}，还用它来对素数进行\emph{构造}——在 $2$ 的情形做一件事，在奇素数的情形做另一件事。
构造的最终结果不仅仅是某个陈述的真值，而是一个可能依赖于该素数奇偶性的数据。
因此，从类型论的视角来看，这类构造很自然地要用余积类型``$(p=2)+(p\text{ is odd})$''的归纳原理来表述，而不是它的命题截断。

诚然，这不是一个最理想的例子，因为``$p=2$''与``$p$ 是奇数''是互斥的，因此 $(p=2)+(p\text{ is odd})$ 实际上已经是一个命题，从而等价于它的截断（参见 \cref{ex:disjoint-or}）。
更有说服力的例子来自存在量词。
常见的做法是：先证明一个形如``存在 $x$ 使得……''的定理，而后在后续讨论中提及``定理 Y 中构造的那个 $x$''（注意定冠词的使用）。
此外，在推导这个 $x$ 的更多性质时，可能会用到诸如``根据定理 Y 证明中对 $x$ 的构造''这样的短语。

一个非常常见的例子是``$A$ 同构于 $B$''。严格来说，这只是说 $A$ 与 $B$ 之间存在\emph{某个}同构。但在证明此类命题时，几乎总是如此：要么给出一个具体的同构，要么证明某个已知的映射是同构，而后续推理往往取决于具体给出的是哪个同构。

受过集合论训练的数学家常常会为这种``滥用语言''而感到一丝内疚。\index{abuse!of language}
我们可能会试图为此道歉，将它们从最终稿中剔除，或者用``典范的''这类模糊词语来含糊其辞。
问题还因以下事实而加剧：在形式化的集合论中，技术上根本没有办法``构造''对象——我们只能证明具有某些性质的对象存在。
因此，类型论中的\emph{非截断逻辑}捕捉到了非形式数学中一些被集合论重构所掩盖的常见做法。
（这与泛等公理佐证``将同构对象视为相同''这一常见做法——尽管它在形式上缺乏依据——是类似的。）

另一方面，\emph{截断逻辑}有时又是不可或缺的。
我们在 \LEM{} 和 \choice{} 的陈述中已经看到了这一点；本书后续还会出现其他例子。
于是，我们面临这样一个问题：在书写非形式化的类型论时，当我们使用``或''和``存在''（以及``存在一个''``我们有''这类常见同义词）这些词时，究竟应当意指什么？

达成普遍共识或许不太可能。
也许根据所从事的数学门类不同，某一种约定可能更为适用——又或许，约定如何选择根本无关紧要。
在这种情况下，在数学论文开头加一段说明，告知读者所采用的语言约定，可能就足够了。
然而，即便选定了一个总体约定，另一种逻辑通常仍会至少偶尔出现，因此我们需要一种方式来指称它。
更一般地，可以考虑用其他行为类似的类型运算来替换命题截断，例如双重否定运算 $A\mapsto \neg\neg A$，或将在 \cref{cha:hlevels} 中讨论的 $n$-截断。
作为表述方式的一种尝试，在后续内容中，我们将偶尔使用\emph{副词}\index{adverb}来表示命题截断这类``模态''的应用。

例如，如果非截断逻辑是默认约定，我们可以用副词 \define{仅仅}
\indexdef{merely}%
来表示命题截断。
于是，短语

\begin{center}
  ``仅仅存在一个 $x:A$ 使得 $P(x)$''
\end{center}

表示类型 $\brck{\sm{x:A} P(x)}$。
类似地，我们将称类型 $A$ 是\define{仅知有实例}的，
\indexdef{merely!inhabited}%
\indexdef{inhabited type!merely}%
意思是它的命题截断 $\brck A$ 是有实例的（即我们有它的一个未命名元素）。
请注意，这是对副词``merely''在非形式数学英语中用法的\emph{定义}\index{definition!of adverbs}，
正如我们为``群''和``环''等名词\index{noun}，以及``正则的''和``正规的''等形容词\index{adjective}赋予精确数学含义一样。
我们并非声称``merely''在词典中的定义就是指命题截断；选择这个词只是意在提醒数学家读者：一个命题所包含的信息``仅仅''是一个真值，再无其他。

另一方面，如果截断逻辑是当前的默认约定，我们可以用\define{纯粹地}
\indexdef{purely}%
或\define{构造性地}这样的副词来标示截断的缺失，从而

\begin{center}
``纯粹地存在一个 $x:A$ 使得 $P(x)$''
\end{center}

将表示类型 $\sm{x:A} P(x)$。
即使无截断已是默认约定，我们也可以用``purely''或``actually''来强调这一点。

在本书中，出于若干原因，我们将继续以非截断逻辑作为默认约定。

\begin{enumerate}[label=(\arabic*)]
\item 我们想鼓励初学者去尝试使用它，而不是因为更熟悉就固守截断逻辑。
\item 在类型论中将截断逻辑作为默认，会面临与集合论基础同类的``语言滥用''\index{abuse!of language}问题，而未截断逻辑则能避免这些问题。
  举例来说，我们将 ``$\eqv A B$'' 定义为 $A$ 与 $B$ 之间等价的类型，而非其命题截断，这意味着证明形如 ``$\eqv A B$'' 的定理，字面上就是构造一个具体的此类等价。
  这个具体的等价之后可以引用。
\item 我们想强调，``命题''这一概念并非类型论的基本组成部分。
  正如我们将在 \cref{cha:hlevels} 中看到的，命题仅仅是无限阶梯中的第二级，此外还存在许多完全不在这个阶梯上的其他模态。
\item 许多在经典意义下是命题的陈述，在同伦类型论中不再如此。
  当然，其中最重要的便是相等。
\item 另一方面，同伦类型论最有趣的观察之一是，有数量惊人的类型\emph{自动}成为命题，或者稍加修改即可如此，而无需任何截断。
  （参见 \cref{thm:isprop-isprop,cha:equivalences,cha:hlevels,cha:category-theory,cha:set-math}。）
  因此，尽管这些类型除真值外不含任何数据，我们仍然可以用它们来构造未截断的对象，因为无需使用命题截断的归纳原理。
  如果默认对所有陈述都应用命题截断，这一有用事实的表达会变得更为笨拙。
\item 最后，截断对于本书中的大部分数学来说用处不大，因此当它们出现时明确标注出来更为简便。
\end{enumerate}

\index{mere proposition|)}%
\index{logic!of mere propositions|)}%

\section{可缩性}
\label{sec:contractibility}

\index{type!contractible|(defstyle}%
\index{contractible!type|(defstyle}%

在 \cref{thm:inhabprop-eqvunit} 中我们观察到，一个有实例的命题必然等价于 $\unit$，
\index{type!unit}%
反过来也成立，这不难看出。
具有这一性质的类型称为\emph{可缩的}。
可缩性的另一个等价定义，有时也较为方便，如下所述。

\begin{defn}\label{defn:contractible}
  一个类型 $A$ 是\define{可缩的（contractible）}，
  或称\define{单点集}，
  \indexdef{type!singleton}%
  \indexsee{singleton type}{type, singleton}%
  如果存在 $a:A$，称为\define{收缩中心}，
  \indexdef{center!of contraction}%
  使得对所有 $x:A$ 都有 $a=x$。
  我们将指定的道路 $a=x$ 记作 $\contr_x$。
\end{defn}

换言之，类型 $\iscontr(A)$ 定义为
\[ \iscontr(A) \defeq \sm{a:A} \prd{x:A}(a=x). \]
注意，在通常的命题即类型解读下，我们可以将 $\iscontr(A)$ 读作``$A$ 恰好包含一个元素''，或更精确地说，``$A$ 包含一个元素，且 $A$ 的每个元素都等于该元素''。

\begin{rmk}
  我们还可以从拓扑角度将 $\iscontr(A)$ 读作``存在点 $a:A$，使得对所有 $x:A$，存在一条从 $a$ 到 $x$ 的道路''。
  注意，对熟悉经典数学的人而言，这听起来更像是\emph{连通性}而非可缩性的定义。
  \index{continuity of functions in type theory@``continuity'' of functions in type theory}%
  \index{functoriality of functions in type theory@``functoriality'' of functions in type theory}%
  关键在于，此句中``存在''的含义是连续的/自然的。

  表达连通性的更好方式是 $\sm{a:A}\prd{x:A} \brck{a=x}$。
  如果假定 $A$ 是带点的，这确实是正确的——参见 \cref{thm:connected-pointed} 后的备注——但一般而言，一个类型可以不带基点而连通。
  在 \cref{sec:connectivity} 中，我们将把连通性定义为一般 $n$-连通性概念在 $n=0$ 时的特例，并且在 \cref{ex:connectivity-inductively} 中，请读者证明该定义等价于同时拥有 $\brck{A}$ 与 $\prd{x,y:A} \brck{x=y}$。
\end{rmk}

\begin{lem}\label{thm:contr-unit}
  对于类型 $A$，以下逻辑等价。
  \begin{enumerate}
  \item $A$ 是按 \cref{defn:contractible} 定义的可缩的。\label{item:contr}
  \item $A$ 是一个命题，并且存在点 $a:A$。\label{item:contr-inhabited-prop}
  \item $A$ 等价于 \unit。\label{item:contr-eqv-unit}
  \end{enumerate}
\end{lem}

\begin{proof}
  如果 $A$ 可缩，那么它显然有一个点 $a:A$（即收缩中心），并且对于任意 $x,y:A$ 都有 $x=a=y$；因此 $A$ 是一个命题。
  反过来，如果我们有一个点 $a:A$ 并且 $A$ 是一个命题，那么对于任意 $x:A$ 都有 $x=a$；因此 $A$ 可缩。
  而我们已在 \cref{thm:inhabprop-eqvunit} 中证明了~\ref{item:contr-inhabited-prop}$\Rightarrow$\ref{item:contr-eqv-unit}，其逆命题则可由单位类型~\unit 显然满足性质~\ref{item:contr-inhabited-prop} 得出。
\end{proof}

\begin{lem}\label{thm:isprop-iscontr}
  对于任何类型 $A$，类型 $\iscontr(A)$ 是一个命题。
\end{lem}

\begin{proof}
  假设给定 $c,c':\iscontr(A)$。
  我们可设 $c\jdeq(a,p)$ 与 $c'\jdeq(a',p')$，其中 $a,a':A$，且有 $p:\prd{x:A} (a=x)$ 与 $p':\prd{x:A} (a'=x)$。
  根据 $\Sigma$-类型中道路的特征，要证明 $c=c'$，只需给出 $q:a=a'$ 使得 $\trans{q}{p}=p'$。
  %
  我们取 $q\defeq p(a')$。
  现在，由于 $A$ 可缩（根据 $c$ 或 $c'$），由 \cref{thm:contr-unit} 可知它是一个命题。
  因此由 \cref{thm:prop-set,thm:isprop-forall}，$\prd{x:A}(a'=x)$ 也是命题；于是 $\trans{q}{p}=p'$ 自动成立。
\end{proof}

\begin{cor}\label{thm:contr-contr}
  如果 $A$ 可缩，那么 $\iscontr(A)$ 也可缩。
\end{cor}

\begin{proof}
  由 \cref{thm:isprop-iscontr} 与 \cref{thm:contr-unit}\ref{item:contr-inhabited-prop} 即得。
\end{proof}

如同命题一样，可缩类型也被许多类型构造子所保持。
例如，我们有：

\begin{lem}\label{thm:contr-forall}
  如果 $P:A\to\type$ 是一个类型族，且每个 $P(a)$ 都可缩，那么 $\prd{x:A} P(x)$ 也可缩。
\end{lem}

\begin{proof}
  根据 \cref{thm:isprop-forall}，$\prd{x:A} P(x)$ 是一个命题，因为每个 $P(x)$ 都是命题。
  但它也有一个元素，即把每个 $x:A$ 映到 $P(x)$ 的收缩中心的函数。
  于是由 \cref{thm:contr-unit}\ref{item:contr-inhabited-prop} 可知，$\prd{x:A} P(x)$ 可缩。
\end{proof}

\index{function extensionality}%
（事实上，\cref{thm:contr-forall} 的陈述等价于函数外延性公理。
参见~\cref{sec:univalence-implies-funext}。）

显然，如果 $A$ 等价于 $B$ 且 $A$ 可缩，那么 $B$ 也可缩。
更一般地，只需 $B$ 是 $A$ 的一个\emph{收缩核（retract）}即可。
按定义，一个\define{收缩}
\indexdef{retraction}%
\indexdef{function!retraction}%
是一个函数 $r : A \to B$，使得存在一个函数 $s : B \to A$（称为它的\define{截面}，
\indexdef{section}%
\indexdef{function!section}%）
以及一个同伦 $\epsilon:\prd{y:B} (r(s(y))=y)$；此时我们称 $B$ 是 $A$ 的一个\define{收缩核}%
\indexdef{retract!of a type}。

\begin{lem}\label{thm:retract-contr}
  如果 $B$ 是 $A$ 的一个收缩核，并且 $A$ 是可缩的，那么 $B$ 也是可缩的。
\end{lem}

\begin{proof}
  设 $a_0 : A$ 为收缩中心。
  我们断言 $b_0 \defeq r(a_0) : B$ 是 $B$ 的一个收缩中心。
  令 $b : B$；我们需要一条道路 $b = b_0$。
  但我们有 $\epsilon_b : r(s(b)) = b$ 与 $\contr_{s(b)} : s(b) = a_0$，因此通过复合可得
  \[ \opp{\epsilon_b} \ct \ap{r}{\contr_{s(b)}} : b = r(a_0) \jdeq b_0. \qedhere\]
\end{proof}

可缩类型乍看可能并不十分有趣，因为它们都等价于 \unit。
这一概念之所以有用，原因之一在于，有时一组各自非平凡的数据，整体上却会形成一个可缩的类型。
一个重要的例子就是带一个自由端点的道路空间。
正如我们将在 \cref{sec:identity-systems} 中看到的，这个事实本质上封装了相等类型中基于一点的道路归纳原理。

\begin{lem}\label{thm:contr-paths}
  对任意 $A$ 与任意 $a:A$，类型 $\sm{x:A} (a=x)$ 是可缩的。
\end{lem}

\begin{proof}
  我们选择点 $(a,\refl a)$ 作为中心。
  现在假设 $(x,p):\sm{x:A}(a=x)$；我们必须证明 $(a,\refl a) = (x,p)$。
  根据 $\Sigma$-类型中道路的特征，只需给出 $q:a=x$，使得 $\trans{q}{\refl a} = p$。
  但我们可取 $q\defeq p$，此时 $\trans{q}{\refl a} = p$ 由道路类型中输运的特征可得。
\end{proof}

出现这种情况时，我们便可在等价意义下简化复杂的构造，所依据的非正式原则是：可缩的数据可以自由忽略。
这一原则体现为许多引理，其中大部分留给读者；下面是一例。

\begin{lem}\label{thm:omit-contr}
  设 $P:A\to\type$ 是一个类型族。
  \begin{enumerate}
  \item 如果每个 $P(x)$ 都是可缩的，那么 $\sm{x:A} P(x)$ 等价于 $A$。\label{item:omitcontr1}
  \item 如果 $A$ 是可缩的且收缩中心为 $a$，那么 $\sm{x:A} P(x)$ 等价于 $P(a)$。\label{item:omitcontr2}
  \end{enumerate}
\end{lem}

\begin{proof}
  对于情况~\ref{item:omitcontr1}，我们证明 $\proj1:\sm{x:A} P(x) \to A$ 是一个等价。
  取其拟逆为 $g(x)\defeq (x,c_x)$，其中 $c_x$ 是 $P(x)$ 的收缩中心。
  复合 $\proj1 \circ g$ 显然是 $\idfunc[A]$，而反向复合利用每个 $P(x)$ 的收缩性，同伦于恒等映射。

  我们将情况~\ref{item:omitcontr2} 的证明留给读者（参见 \cref{ex:omit-contr2}）。
\end{proof}

可缩类型之所以有趣的另一个原因在于，它们将 \cref{sec:basics-sets} 中提到的 $n$-类型的阶梯向下又延伸了一级。

\begin{lem}\label{thm:prop-minusonetype}
  一个类型 $A$ 是命题，当且仅当对所有 $x,y:A$，类型 $\id[A]xy$ 都是可缩的。
\end{lem}

\begin{proof}
  对于“如果”方向，我们只需注意到任何可缩类型都是有实例的。
  对于“仅当”方向，我们在 \cref{subsec:hprops} 中已观察到每个命题都是集合，因而每个类型 $\id[A]xy$ 都是命题。
  但它也是有实例的（因为 $A$ 是命题），于是由 \cref{thm:contr-unit}\ref{item:contr-inhabited-prop} 可知它是可缩的。
\end{proof}

因此，可缩类型也可以称为 \define{$(-2)$-类型（$(-2)$-types）}。
它们是 $n$-类型阶梯的最底层，也将是 \cref{cha:hlevels} 中递归定义 $n$-类型时的基础情形。

\index{type!contractible|)}%
\index{contractible!type|)}%

\sectionNotes

在类型论中可以定义集合、命题和可缩类型，且所有高阶同伦都如 \cref{sec:basics-sets,subsec:hprops,sec:contractibility} 中那样自动得到处理——这一事实最早由 Voevodsky 注意到。
事实上，他通过归纳定义了整个 $n$-类型的层级，正如我们将在 \cref{cha:hlevels} 中所做的那样。\index{n-type@$n$-type!definable in type theory}%

\cref{thm:not-dneg,thm:not-lem} 本质上依赖于 Hedberg 的一个经典定理，
\index{Hedberg's theorem}%
\index{theorem!Hedberg's}%
我们将在 \cref{sec:hedberg} 中加以证明。
“命题即类型”形式的 \LEM{} 与泛等性相矛盾，这一结论是由 Martín Escardó（马丁·埃斯卡多）在 \Agda 邮件列表上指出的。
而我们给出的 \cref{thm:not-dneg} 的证明则归功于 Thierry Coquand。

命题截断最初由 Constable~\cite{Con85} 于 1983 年作为“子集”类型和“商”类型的一种应用，引入到 \NuPRL 的外延类型论
\index{proof!assistant!\NuPRL}%
之中。在 \NuPRL 类型论~\cite{constable+86nuprl-book} 中，本书所称的“命题截断”当时被称为“压平”。
此后，~\cite{ab:bracket-types} 仍在外延类型论的框架下，给出了直接刻画命题截断的规则。而在同伦类型论中的内涵版本，则由 Voevodsky 利用非直谓\index{impredicative!truncation}量化构造，随后 Lumsdaine 又用高阶归纳类型加以构造（参见 \cref{sec:hittruncations}）。

\index{propositional!resizing}%
Voevodsky~\cite{Universe-poly} 提出了 \cref{subsec:prop-subsets} 中所讨论的那一类大小重整规则。
\index{axiom!of reducibility}\index{resizing}%
这些规则显然与 Russell 在他与 Whitehead 合著的《数学原理》~\cite{PM2} 中提出的那个备受争议的\emph{可化归公理（axiom of reducibility）}密切相关。\index{Russell, Bertrand}

副词“纯粹地（purely）”用于指代未截断的逻辑，这一用法是在援引编程语言中以单子模态来建模副作用（effects）的做法；参见 \cref{sec:modalities} 及 \cref{cha:hlevels} 的附注。

相对于类型论，逻辑可以有多种不同的处理方式。
例如，除了 \cref{sec:pat} 中描述的单纯的“命题即类型”逻辑，以及 \cref{subsec:logic-hprop} 描述的仅使用命题的另一种方式外，人们还可以引入一种独立的命题 ``sort''，其行为在某些方面类似于类型，但不与类型等同。
此方法被逻辑充实类型论~\cite{aczel2002collection}、对意象（topos）及相关范畴内部语言的某些表述（例如~\cite{jacobs1999categorical,elephant}）以及证明助理 \Coq\index{proof!assistant!Coq@\textsc{Coq}} 所采用。
这种处理方法更一般，但功能也较弱。
举例来说，唯一选择原理（\cref{sec:unique-choice}）在 \Coq~\cite{Spiwack} 中所谓的 setoid 构成的范畴里、在逻辑充实类型论~\cite{aczel2002collection} 中、以及在最小类型论~\cite{maietti2005toward} 中都是不成立的。\index{setoid}
因此，泛等公理使我们的类型论更接近意象的内部逻辑；另见 \cref{cha:set-math}。\index{topos}

Martin-L\"of~\cite{martin2006100} 对选择公理的历史进行了讨论。当然，构造性数学和直觉主义数学有着悠久而复杂的历史，我们在此不再深究；例如可参阅~\cite{TroelstraI,TroelstraII}。

\sectionExercises

\begin{ex}\label{ex:equiv-functor-set}
  证明若 $\eqv A B$ 且 $A$ 是一个集合，则 $B$ 亦然。
\end{ex}

\begin{ex}\label{ex:isset-coprod}
  证明若 $A$ 和 $B$ 都是集合，则 $A+B$ 也是集合。
\end{ex}

\begin{ex}\label{ex:isset-sigma}
  证明若 $A$ 是一个集合，且 $B:A\to \type$ 是一个类型族，使得对所有 $x:A$，$B(x)$ 都是集合，则 $\sm{x:A} B(x)$ 也是集合。
\end{ex}

\begin{ex}\label{ex:prop-endocontr}
  证明 $A$ 是一个命题当且仅当 $A\to A$ 是可缩的。
\end{ex}

\begin{ex}\label{ex:prop-inhabcontr}
  证明 $\eqv{\isprop(A)}{(A\to\iscontr(A))}$。
\end{ex}

\begin{ex}\label{ex:lem-mereprop}
  证明若 $A$ 是一个命题，则 $A+(\neg A)$ 也是命题。
  因此，在~\eqref{eq:lem} 处无需插入命题截断。
\end{ex}

\begin{ex}\label{ex:disjoint-or}
  更一般地，证明若 $A$ 和 $B$ 都是命题且 $\neg(A\times B)$，则 $A+B$ 也是命题。
\end{ex}

% \begin{ex}\label{ex:hprop-iff-equiv}
%   Show that if $A$ and $B$ are mere propositions such that $A\to B$ and $B\to A$, then $\eqv A B$.
% \end{ex}

% \begin{ex}\label{ex:isprop-isprop}
%   Show that for any type $A$, the types $\isprop(A)$ and $\isset(A)$ are mere propositions.
% \end{ex}

% \begin{ex}\label{ex:prop-eqvtrunc}
%   Show that if $A$ is already a mere proposition, then $\eqv A{\brck{A}}$.
% \end{ex}

\begin{ex}\label{ex:brck-qinv}
  假设某个类型 $\isequiv(f)$ 满足\cref{sec:basics-equivalences}的条件~\ref{item:be1}--\ref{item:be3}，证明类型 $\brck{\qinv(f)}$ 满足同样的条件，并且等价于 $\isequiv(f)$。
\end{ex}

\begin{ex}\label{ex:lem-impl-prop-equiv-bool}
  证明：若排中律成立，则类型$\prop \defeq \sm{A:\type} \isprop(A)$等价于 \bool。
\end{ex}

\begin{ex}\label{ex:lem-impred}
  证明：若 $\UU_{i+1}$ 满足排中律，则典范包含$\prop_{\UU_i} \to \prop_{\UU_{i+1}}$是一个等价。
\end{ex}

\begin{ex}\label{ex:not-brck-A-impl-A}
  证明：并非对所有的 $A:\type$ 都有 $\brck{A} \to A$。
  （不过，对于某些特定的类型可以有 $\brck{A}\to A$。
  \cref{ex:brck-qinv}蕴含 $\qinv(f)$ 就是这样的类型。）
\end{ex}

\begin{ex}\label{ex:lem-impl-simple-ac}
  \index{axiom!of choice}%
  证明：若排中律成立，则对所有的 $A:\type$，有 $\bbrck{(\brck A \to A)}$。
  （这个性质是选择公理的一种极简形式，在没有排中律的情况下可能不成立；参见~\cite{krausgeneralizations}。）
\end{ex}

\begin{ex}\label{ex:naive-lem-impl-ac}
  我们在\cref{thm:not-lem}中证明了，下述朴素形式的排中律与泛等性不相容：
  \[ \prd{A:\type} (A+(\neg A)) \]
  在没有泛等性的情况下，这条公理是一致的。
  然而，证明它蕴含选择公理~\eqref{eq:ac}。
\end{ex}

\begin{ex}\label{ex:lem-brck}
  证明：假设排中律成立，双重否定 $\neg \neg A$ 与命题截断 $\brck A$ 具有相同的递归原理，但其计算规则是命题性的而非判断性的。
  换言之，证明：假设排中律成立，若 $B$ 是一个命题且有 $f:A\to B$，则存在导出的 $g:\neg\neg A \to B$ 使得对所有 $a:A$ 都有 $g(\bproj a) = f(a)$。
  由此推出（假设排中律成立）$\eqv{\neg\neg A}{\brck{A}}$。
  因此，在排中律下，命题截断可以被定义出来，而不必作为一个单独的类型构造子。
\end{ex}

\begin{ex}\label{ex:impred-brck}\ 
  \index{propositional!resizing}%
  \begin{enumerate}
  \item 证明：对任何 $A : \UU$，类型
    \[\prd{P:\prop_{\UU}} \Parens{(A\to P)\to P}\]
    与 $\brck A$ 具有相同的递归原理（至少相对于 $\UU$ 中的命题而言），\emph{并}具有相同的判断性计算规则。
  \item 上一部分考虑的类型并不位于同一个宇宙 $\UU$ 中，而且它的递归原理只适用于 $\UU$ 中的命题。
    证明：如果假设命题大小重整，那么可以定义一个确实位于同一宇宙 $\UU$ 中的类型，它满足与 $\brck A$ 相同的递归原理，尽管其计算规则只能是命题性的。
    因此，在这种情况下我们同样可以定义命题截断。
  \end{enumerate}
\end{ex}

\begin{ex}\label{ex:lem-impl-dn-commutes}
  假设排中律成立，证明双重否定与集合上命题的全称量化可交换。
  也就是说，证明：若 $X$ 是一个集合且每个 $Y(x)$ 都是命题，则排中律蕴含
  \begin{equation}
    \eqv{\Parens{\prd{x:X} \neg\neg Y(x)}}{\Parens{\neg\neg \prd{x:X} Y(x)}}.\label{eq:dnshift}
  \end{equation}
  注意，如果我们改为假设每个 $Y(x)$ 都是集合，那么~\eqref{eq:dnshift}将等价于选择公理~\eqref{eq:epis-split}。
\end{ex}

\begin{ex}\label{ex:prop-trunc-ind}
  \index{induction principle!for truncation}%
  证明\cref{subsec:prop-trunc}中给出的命题截断规则足以推出下述归纳原理：对于任意类型族 $B:\brck A \to \type$，若每个 $B(x)$ 都是命题，且对每个 $a:A$ 都有 $B(\bproj a)$，则对每个 $x:\brck A$ 都有 $B(x)$。
\end{ex}

\begin{ex}\label{ex:lem-ldn}
  证明排中律~\eqref{eq:lem}与双重否定律~\eqref{eq:ldn}是逻辑等价的。
\end{ex}

\begin{ex}\label{ex:decidable-choice}
  设 $P:\nat\to\type$ 是一个可判定的命题族（参见\cref{defn:decidable-equality}\ref{item:decidable-equality2}）。
  证明：
  \[ \Brck{\sm{n:\nat} P(n)} \;\to\; \sm{n:\nat}P(n).\]
\end{ex}

\begin{ex}\label{ex:omit-contr2}
  证明\cref{thm:omit-contr}\ref{item:omitcontr2}：如果 $A$ 是可缩的，且收缩中心为 $a$，那么 $\sm{x:A} P(x)$ 等价于 $P(a)$。
\end{ex}

\begin{ex}\label{ex:isprop-equiv-equiv-bracket}
  证明$\isprop(P) \simeq (P \simeq \brck P)$。
\end{ex}

\begin{ex}\label{ex:finite-choice}
  如同在经典集合论中一样，选择公理的有限版本是一个定理。证明当 $X$ 是有限类型 $\Fin(n)$（定义见\cref{ex:fin}）时，选择公理\eqref{eq:ac}成立。
\end{ex}

\begin{ex}\label{ex:decidable-choice-strong}
  证明：如果 $P:\nat\to\type$ 是任意可判定族，则\cref{ex:decidable-choice}的结论成立。
\end{ex}

\begin{ex}\label{ex:n-set}
  通过先证明对所有 $m,n$，$\code(m,n)$ 都是命题，来简化\cref{thm:path-nat}的证明。
\end{ex}

% Local Variables:
% TeX-master: "hott-online"
% End:

\chapter{等价}
\label{cha:equivalences}

现在我们要更详细地研究\emph{类型的等价（equivalence of types）}这一概念，该概念在 \cref{sec:basics-equivalences} 中已有简要介绍。
具体来说，我们将给出几种不同的方式来定义一个类型 $\isequiv(f)$，使其满足那里提到的性质。
回想一下，我们希望 $\isequiv(f)$ 具有以下性质，现重述如下：

\begin{enumerate}
\item $\qinv(f) \to \isequiv (f)$。\label{item:beb1}
\item $\isequiv (f) \to \qinv(f)$。\label{item:beb2}
\item $\isequiv(f)$ 是一个命题。\label{item:beb3}
\end{enumerate}

这里 $\qinv(f)$ 表示 $f$ 的拟逆所构成的类型：
\begin{equation*}
  \sm{g:B\to A} \big((f \circ g \htpy \idfunc[B]) \times (g\circ f \htpy \idfunc[A])\big).
\end{equation*}
利用函数外延性可知，$\qinv(f)$ 等价于类型
\begin{equation*}
  \sm{g:B\to A} \big((f \circ g = \idfunc[B]) \times (g\circ f = \idfunc[A])\big).
\end{equation*}

我们将定义三种不同的类型，它们都满足性质~\ref{item:beb1}--\ref{item:beb3}，分别称为：

\begin{itemize}
\item 半伴随等价，
\item 双可逆映射，
  \index{function!bi-invertible}
  以及
\item 可缩函数。
\end{itemize}

我们还将证明所有这些类型都是彼此等价的。
这些名称故意取得有些冗长，因为一旦我们知道它们彼此等价并且满足性质~\ref{item:beb1}--\ref{item:beb3}，
我们便会直接使用"等价"一词，而无需指明具体选择了哪一种定义。
不过，为了本章进行比较的需要，我们仍然需要为每种定义保留不同的名称。

然而，在具体考察不同的等价概念之前，我们先略微解释一下，为什么需要一个不同于拟逆性的概念。

\section{拟逆}
\label{sec:quasi-inverses}

\index{quasi-inverse|(}%
我们曾说过，$\qinv(f)$ 并不令人满意，因为它不是一个命题；而我们更希望一个给定的函数最多只能以一种方式``成为等价''。
然而，我们尚未给出任何证据表明 $\qinv(f)$ 不是一个命题。
在本节中，我们将展示一个具体的反例。

\begin{lem}\label{lem:qinv-autohtpy}
  如果 $f:A\to B$ 满足 $\qinv (f)$ 是有实例的，那么
  \[\eqv{\qinv(f)}{\Parens{\prd{x:A}(x=x)}}.\]
\end{lem}

\begin{proof}
  根据假设，$f$ 是一个等价；也就是说，我们有 $e:\isequiv(f)$，从而 $(f,e):\eqv A B$。
  由泛等性，$\idtoeqv:(A=B) \to (\eqv A B)$ 是一个等价，因此我们可以假设 $(f,e)$ 具有 $\idtoeqv(p)$ 的形式，其中 $p:A=B$。
  然后利用道路归纳，可以假设 $p$ 就是 $\refl{A}$，此时 $f$ 即为 $\idfunc[A]$。
  这样，我们只需证明 $\eqv{\qinv(\idfunc[A])}{(\prd{x:A}(x=x))}$。
  现在根据定义，我们有
  \[ \qinv(\idfunc[A]) \jdeq
  \sm{g:A\to A} \big((g \htpy \idfunc[A]) \times (g \htpy \idfunc[A])\big).
  \]
  由函数外延性，这等价于
  \[ \sm{g:A\to A} \big((g = \idfunc[A]) \times (g = \idfunc[A])\big).
  \]
  并且根据 \cref{ex:sigma-assoc}，上式又等价于
  \[ \sm{h:\sm{g:A\to A} (g = \idfunc[A])} (\proj1(h) = \idfunc[A])
  \]
  然而，根据 \cref{thm:contr-paths}，$\sm{g:A\to A} (g = \idfunc[A])$ 是以 $(\idfunc[A],\refl{\idfunc[A]})$ 为中心的可缩类型；因此由 \cref{thm:omit-contr} 可知，这个类型等价于 $\idfunc[A] = \idfunc[A]$。
  而再次利用函数外延性，$\idfunc[A] = \idfunc[A]$ 等价于 $\prd{x:A} x=x$。
\end{proof}

\noindent
我们指出，\cref{ex:qinv-autohtpy-no-univalence} 要求给出上述引理的一个避免使用泛等性的证明。

因此，我们需要的是一个类型 $A$，使得 $\prd{x:A}(x=x)$ 有非平凡的元素。
将 $A$ 视为一个高阶群胚，那么 $\prd{x:A}(x=x)$ 的一个实例就是从 $A$ 的恒等函子到其自身的一个自然变换\index{natural!transformation}。
这类变换被称为构成\define{范畴的中心（center of a category）}，
\index{center!of a category}%
\index{category!center of}%
因为自然性公理要求它们与所有态射交换。
在经典情形下，如果 $A$ 就是一个被视作单对象群胚的群，那么这给出的恰好就是通常群论意义上的中心。
这为下面的引理提供了一些动机。

\begin{lem}\label{lem:autohtpy}
  假设我们有一个类型 $A$，其中 $a:A$ 且 $q:a=a$，满足以下条件：
  \begin{enumerate}
  \item 类型 $a=a$ 是一个集合。\label{item:autohtpy1}
  \item 对所有 $x:A$，我们有 $\brck{a=x}$。\label{item:autohtpy2}
  \item 对所有 $p:a=a$，我们有 $p\ct q = q \ct p$。\label{item:autohtpy3}
  \end{enumerate}
  那么存在 $f:\prd{x:A} (x=x)$ 使得 $f(a)=q$。
\end{lem}

\begin{proof}
  设 $g:\prd{x:A} \brck{a=x}$ 如~\ref{item:autohtpy2} 所给出。
  首先我们观察到，每个类型 $\id[A]xy$ 都是一个集合。
  因为"是集合"是一个命题，所以我们可以应用命题截断的归纳原理，并假设 $g(x)=\bproj p$ 且 $g(y)=\bproj{p'}$，其中 $p:a=x$ 且 $p':a=y$。
  在这种情况下，与 $p$ 和 $\opp{p'}$ 复合，便得到一个等价 $\eqv{(x=y)}{(a=a)}$。
  而 $(a=a)$ 根据~\ref{item:autohtpy1} 是一个集合，所以 $(x=y)$ 也是一个集合。

  现在，我们本想通过给每个 $x$ 指派路径 $\opp{g(x)}
  \ct q \ct g(x)$ 来定义 $f$，但这行不通。
  因为 $g(x)$ 并不是 $a=x$ 的元素，而是 $\brck{a=x}$ 的元素；并且类型 $(x=x)$ 可能不是一个命题，所以我们不能对命题截断使用归纳原理。
  我们可以转而应用 \cref{sec:unique-choice} 中提到的技巧：通过唯一性来刻画我们希望构造的对象。
  对每个 $x:A$，定义类型
  \[ B(x) \defeq \sm{r:x=x} \prd{s:a=x} (r = \opp s \ct q\ct s).\]
  我们断言，对于每个 $x:A$，$B(x)$ 都是一个命题。
  由于这个断言本身也是一个命题，我们可以再次对截断应用归纳原理，并假设 $g(x) = \bproj p$，其中 $p:a=x$。
  现在假设给定 $B(x)$ 中的 $(r,h)$ 和 $(r',h')$；那么我们有
  \[ h(p) \ct \opp{h'(p)} : r = r'. \]
  剩下要证明的是，当沿着这个等式传输时，$h$ 与 $h'$ 是等同的。根据在相等类型和函数类型中的传输（\cref{sec:compute-paths,sec:compute-pi}），这归结为要证：对于任何 $s:a=x$，
  \[ h(s) = h(p) \ct \opp{h'(p)} \ct h'(s) \]
  成立。
  但这个等式的两边都是 $(x=x)$ 中元素之间的等式，因此由我们上面的观察（$(x=x)$ 是集合）可知其成立。

  这样，每个 $B(x)$ 都是一个命题；接下来我们断言 $\prd{x:A} B(x)$ 成立。
  给定 $x:A$，我们现在可以调用命题截断的归纳原理，假设 $g(x) = \bproj p$，其中 $p:a=x$。
  我们定义 $r \defeq \opp p \ct q \ct p$；为了构造 $B(x)$ 的实例，剩下要证明对于任何 $s:a=x$，都有 $r = \opp s \ct q \ct s$。
  对道路进行变形，这归结为要证明 $q\ct (p\ct \opp s) = (p\ct \opp s) \ct q$。
  而这正好是~\ref{item:autohtpy3} 的一个实例。
\end{proof}

\begin{thm}\label{thm:qinv-notprop}
  存在类型 $A$ 和 $B$ 以及函数 $f:A\to B$，使得 $\qinv(f)$ 不是一个命题。
\end{thm}

\begin{proof}
  我们只需展示一个类型 $X$，使得 $\prd{x:X} (x=x)$ 不是一个命题即可。
  如同在 \cref{thm:no-higher-ac} 的证明中那样定义 $X\defeq \sm{A:\type} \brck{\bool=A}$。
  此时，我们只需展示一个 $f:\prd{x:X} (x=x)$，使其不等于 $\lam{x} \refl{x}$。

  令 $a \defeq (\bool,\bproj{\refl{\bool}}) : X$，并取 $q:a=a$ 为对应于非恒等的等价 $e:\eqv\bool\bool$ 的道路；这个等价 $e$ 由 $e(\bfalse)\defeq\btrue$ 与 $e(\btrue)\defeq\bfalse$ 定义。
  我们想要应用 \cref{lem:autohtpy} 来构造 $f$。
  根据 $X$ 的定义、子集类型中的相等性（\cref{subsec:prop-subsets}）以及泛等性，我们得到 $\eqv{(a=a)}{(\eqv{\bool}{\bool})}$，这是一个集合，因此条件~\ref{item:autohtpy1} 成立。
  类似地，根据 $X$ 的定义以及子集类型中的相等性，我们得到条件~\ref{item:autohtpy2}。
  最后，根据 \cref{ex:eqvboolbool} 可知，每个等价 $\eqv\bool\bool$ 都等于 $\idfunc[\bool]$ 或 $e$，于是通过分四种情形讨论，我们便能证明~\ref{item:autohtpy3}。

  这样，我们就得到了 $f:\prd{x:X} (x=x)$，满足 $f(a) = q$。
  由于 $e$ 不等于 $\idfunc[\bool]$，因此 $q$ 不等于 $\refl{a}$，从而 $f$ 不等于 $\lam{x} \refl{x}$。
  所以，$\prd{x:X} (x=x)$ 不是一个命题。
\end{proof}

更一般地，\cref{lem:autohtpy} 蕴含了，任何艾伦伯格–麦克兰恩空间 $K(G,1)$（其中 $G$ 是一个非平凡交换\index{group!abelian}群）都将给出一个反例；见 \cref{cha:homotopy}。
我们之前用过的类型 $X$，实际上等价于 $K(\mathbb{Z}_2,1)$。
在 \cref{cha:hits} 我们将看到，圆 $\Sn^1 = K(\mathbb{Z},1)$ 是另一个容易描述的例子。

我们现在转向描述更好的等价概念。

\index{quasi-inverse|)}%

%%%%%%%%%%%%%%%%%%%%%%%%%%%%%%%%%%%%%%
\section{半伴随等价}
\label{sec:hae}
%%%%%%%%%%%%%%%%%%%%%%%%%%%%%%%%%%%%%%

\index{equivalence!half adjoint|(defstyle}%
\index{half adjoint equivalence|(defstyle}%
\index{adjoint!equivalence!of types, half|(defstyle}%

在 \cref{sec:quasi-inverses} 中，我们通过丢弃一个可缩类型的做法，得出了 $\qinv(f)$ 等价于 $\prd{x:A} (x=x)$ 的结论。
粗略来说，类型 $\qinv(f)$ 包含三个数据：$g$、$\eta$ 和 $\epsilon$。
当 $f$ 是一个等价时，其中两个数据（$g$ 和 $\eta$）可以被证明共同构成一个可缩类型。
但问题在于，移除这两个数据后，还会剩下一个数据（$\epsilon$）。
为了解决这个问题，我们的想法是添加一个\emph{额外的}数据，让它与 $\epsilon$ 一起形成一个可缩类型。

\begin{defn}\label{defn:ishae}
  一个函数 $f:A\to B$ 是一个\define{半伴随等价（half adjoint equivalence）}，
  如果存在 $g:B\to A$ 以及同伦 $\eta: g \circ f \htpy \idfunc[A]$ 和 $\epsilon:f \circ g \htpy \idfunc[B]$，
  使得存在一个同伦
  \[\tau : \prd{x:A} \map{f}{\eta x} = \epsilon(fx).\]
\end{defn}

由此我们得到一个类型 $\ishae(f)$，它定义为
\begin{equation*}
  \sm{g:B\to A}{\eta: g \circ f \htpy \idfunc[A]}{\epsilon:f \circ g \htpy \idfunc[B]} \prd{x:A} \map{f}{\eta x} = \epsilon(fx).
\end{equation*}
注意在上述定义中，连接 $\eta$ 和 $\epsilon$ 的相干性\index{coherence}条件只涉及 $f$。
我们也可以转而考虑一个涉及 $g$ 的类似相干性条件：
\[\upsilon : \prd{y:B} \map{g}{\epsilon y} = \eta(gy)\]
并由此得到一个类似的定义 $\ishae'(f)$。

幸运的是，这两个条件互相蕴含：

\begin{lem}\label{lem:coh-equiv}
对于函数 $f : A \to B$ 和 $g:B\to A$，以及同伦 $\eta: g \circ f \htpy \idfunc[A]$ 和 $\epsilon:f \circ g \htpy \idfunc[B]$，以下两个条件逻辑等价：
\begin{itemize}
\item $\prd{x:A} \map{f}{\eta x} = \epsilon(fx)$
\item $\prd{y:B} \map{g}{\epsilon y} = \eta(gy)$
\end{itemize}
\end{lem}

\begin{proof}
只需证明一个方向；另一个方向可以通过将 $A$、$f$ 和 $\eta$ 分别替换为 $B$、$g$ 和 $\epsilon$ 得到。
设 $\tau : \prd{x:A}\;\map{f}{\eta x} = \epsilon(fx)$。
固定 $y : B$。
利用 $\epsilon$ 的自然性并应用 $g$，我们得到如下道路的交换图：
\[\uppercurveobject{{ }}\lowercurveobject{{ }}\twocellhead{{ }}
  \xymatrix@C=3pc{gfgfgy \ar@{=}^-{gfg(\epsilon y)}[r] \ar@{=}_{g(\epsilon (fgy))}[d] & gfgy \ar@{=}^{g(\epsilon y)}[d] \\ gfgy \ar@{=}_{g(\epsilon y)}[r] & gy
  }\]
在图的左侧使用 $\tau(gy)$，我们得到
\[\uppercurveobject{{ }}\lowercurveobject{{ }}\twocellhead{{ }}
  \xymatrix@C=3pc{gfgfgy \ar@{=}^-{gfg(\epsilon y)}[r] \ar@{=}_{gf(\eta (gy))}[d] & gfgy \ar@{=}^{g(\epsilon y)}[d] \\ gfgy \ar@{=}_{g(\epsilon y)}[r] & gy
  }\]
利用 $\eta$ 与 $g \circ f$ 的交换性（\cref{cor:hom-fg}），我们有
\[\uppercurveobject{{ }}\lowercurveobject{{ }}\twocellhead{{ }}
  \xymatrix@C=3pc{gfgfgy \ar@{=}^-{gfg(\epsilon y)}[r] \ar@{=}_{\eta (gfgy)}[d] & gfgy \ar@{=}^{g(\epsilon y)}[d] \\ gfgy \ar@{=}_{g(\epsilon y)}[r] & gy
  }\]
然而，根据 $\eta$ 的自然性，我们还有
\[\uppercurveobject{{ }}\lowercurveobject{{ }}\twocellhead{{ }}
  \xymatrix@C=3pc{gfgfgy \ar@{=}^-{gfg(\epsilon y)}[r] \ar@{=}_{\eta (gfgy)}[d] & gfgy \ar@{=}^{\eta(gy)}[d] \\ gfgy \ar@{=}_{g(\epsilon y)}[r] & gy
  }\]
因此，除右侧同伦外全部消去，便得到 $g(\epsilon y) = \eta(g y)$，如所需。
\end{proof}

但重要的是，我们不在 $\ishae (f)$ 的定义中\emph{同时}包含 $\tau$ 和 $\upsilon$（因此得名"\emph{半}伴随等价"（half adjoint equivalence））。
如果同时包含，那么在消去可缩类型之后，我们仍然会剩下一项数据——除非我们再添加一个更高阶的相干条件。
一般来说，如果在奇数个相干条件之后截断，我们期望会得到一个表现良好的类型。

显然，我们有 $\ishae(f) \to\qinv(f)$：只需忘记相干数据。
反方向则是同伦论与范畴论中一个标准论证的变体。

\begin{thm}\label{thm:equiv-iso-adj}
对任意 $f:A\to B$，我们有 $\qinv(f)\to\ishae(f)$。
\end{thm}

\begin{proof}
假设 $(g,\eta,\epsilon)$ 是 $f$ 的一个拟逆。
我们需要提供一个四元组 $(g',\eta',\epsilon',\tau)$ 来见证 $f$ 是半伴随等价。
要定义 $g'$ 和 $\eta'$，我们可以做出显然的选择：令 $g' \defeq g$，$\eta'\defeq \eta$。
然而，在定义 $\epsilon'$ 时，我们需要考虑 $\tau$ 的构造，因此不能想当然地取 $\epsilon'$ 为 $\epsilon$。
取而代之，我们取
\begin{equation*}
\epsilon'(b) \defeq \opp{\epsilon(f(g(b)))}\ct (\ap{f}{\eta(g(b))}\ct \epsilon(b)).
\end{equation*}
现在我们需要找到
\begin{equation*}
\tau(a): \ap{f}{\eta(a)}=\opp{\epsilon(f(g(f(a))))}\ct (\ap{f}{\eta(g(f(a)))}\ct \epsilon(f(a))).
\end{equation*}
首先注意，由 \cref{cor:hom-fg} 我们有
%$\eta(g(f(a)))\ct\eta(a)=\ap{g}{\ap{f}{\eta(a)}}\ct\eta(a)$ and hence it follows that
$\eta(g(f(a)))=\ap{g}{\ap{f}{\eta(a)}}$。
因此，我们可以应用
\cref{lem:htpy-natural} 来计算
\begin{align*}
\ap{f}{\eta(g(f(a)))}\ct \epsilon(f(a))
& = \ap{f}{\ap{g}{\ap{f}{\eta(a)}}}\ct \epsilon(f(a))\\
& = \epsilon(f(g(f(a))))\ct \ap{f}{\eta(a)}
\end{align*}
由此我们便得到了所需的道路 $\tau(a)$。
\end{proof}

结合 \cref{lem:coh-equiv}（或对证明作对称化处理），我们也有 $\qinv(f)\to\ishae'(f)$。

尚需证明 $\ishae(f)$ 是一个命题。
为此，我们需要先知道等价映射的纤维是可缩的。

\begin{defn}\label{defn:homotopy-fiber}
  一个映射 $f:A\to B$ 在点 $y:B$ 处的 \define{纤维}
  \indexdef{fiber}%
  \indexsee{function!fiber of}{fiber}%
  定义为：
  \[ \hfib f y \defeq \sm{x:A} (f(x) = y).\]
\end{defn}

在同伦论中，这被称为 $f$ 的\emph{同伦纤维（homotopy fiber）}。
\cref{sec:computational} 中的道路引理导出了纤维中道路的如下刻画：

\begin{lem}\label{lem:hfib}
  对于任意 $f : A \to B$，$y : B$，以及 $(x,p),(x',p') : \hfib{f}{y}$，我们有
  \[ \big((x,p) = (x',p')\big) \eqvsym \Parens{\sm{\gamma : x = x'} f(\gamma) \ct p' = p} \qedhere\]
\end{lem}

\begin{thm}\label{thm:contr-hae}
  如果 $f:A\to B$ 是一个半伴随等价，那么对于任意 $y:B$，纤维 $\hfib f y$ 都是可缩的。
\end{thm}

\begin{proof}
  令 $(g,\eta,\epsilon,\tau) : \ishae(f)$，并固定 $y : B$。
  我们选取 $(gy, \epsilon y)$ 作为 $\hfib{f}{y}$ 的收缩中心。
  现在任取 $(x,p) : \hfib{f}{y}$；我们想要构造一条从 $(gy, \epsilon y)$ 到 $(x,p)$ 的道路。
  根据 \cref{lem:hfib}，只需给出一条道路 $\gamma : \id{gy}{x}$ 使得 $\ap f\gamma \ct p = \epsilon y$ 即可。
  我们令 $\gamma \defeq \opp{g(p)} \ct \eta x$。
  于是有
  \begin{align*}
    f(\gamma) \ct p & = \opp{fg(p)} \ct f (\eta x) \ct p \\
    & = \opp{fg(p)} \ct \epsilon(fx) \ct p \\
    & = \epsilon y
  \end{align*}
  其中第二个等式由 $\tau x$ 推出，第三个等式是 $\epsilon$ 的自然性。
\end{proof}

现在我们来定义封装可缩数据对的类型。
下面的类型将拟逆 $g$ 与其中一条同伦组合在一起。

\begin{defn}\label{defn:linv-rinv}
  给定一个函数 $f:A\to B$，我们定义类型
    \begin{align*}
      \linv(f) &\defeq \sm{g:B\to A} (g\circ f\htpy \idfunc[A])\\
      \rinv(f) &\defeq \sm{g:B\to A} (f\circ g\htpy \idfunc[B])
    \end{align*}
  分别为 $f$ 的 \define{左逆}
  \indexdef{left!inverse}%
  \indexdef{inverse!left}%
  与 \define{右逆}
  \indexdef{right!inverse}%
  \indexdef{inverse!right}%
  。
  如果 $\linv(f)$ 有实例，我们就称 $f$ 是 \define{左可逆的}
  \indexdef{function!left invertible}%
  \indexdef{function!right invertible}%
  ；类似地，如果 $\rinv(f)$ 有实例，则称 $f$ 是 \define{右可逆的}
  \indexdef{left!invertible function}%
  \indexdef{right!invertible function}%
  。
\end{defn}

\begin{lem}\label{thm:equiv-compose-equiv}
  如果 $f:A\to B$ 有一个拟逆，那么 \begin{align*}
    (f\circ \blank) &: (C\to A) \to (C\to B)\\
    (\blank\circ f) &: (B\to C) \to (A\to C).
  \end{align*} 亦然。
\end{lem}

\begin{proof}
  如果 $g$ 是 $f$ 的一个拟逆，那么 $(g\circ \blank)$ 和 $(\blank\circ g)$ 就分别是 $(f\circ \blank)$ 和 $(\blank\circ f)$ 的拟逆。
\end{proof}

\begin{lem}\label{lem:inv-hprop}
  如果 $f : A \to B$ 有一个拟逆，那么类型 $\rinv(f)$ 和 $\linv(f)$ 都是可缩的。
\end{lem}

\begin{proof}
  由函数外延性，我们有
  \[\eqv{\linv(f)}{\sm{g:B\to A} (g\circ f = \idfunc[A])}.\]
  但这就是 $(\blank\circ f)$ 在 $\idfunc[A]$ 上的纤维，因此由 \cref{thm:equiv-compose-equiv,thm:equiv-iso-adj,thm:contr-hae}，它是可缩的。
  类似地，$\rinv(f)$ 等价于 $(f\circ \blank)$ 在 $\idfunc[B]$ 上的纤维，因而也是可缩的。
\end{proof}

接下来我们定义那些将另一个同伦与额外的相干数据组合在一起的类型。\index{coherence}%

\begin{defn}\label{defn:lcoh-rcoh}
对于 $f : A \to B$，一个左逆 $(g,\eta) : \linv(f)$，以及一个右逆 $(g,\epsilon) : \rinv(f)$，我们记
\begin{align*}
\lcoh{f}{g}{\eta} & \defeq \sm{\epsilon : f\circ g \htpy \idfunc[B]} \prd{y:B} g(\epsilon y) = \eta (gy), \\
\rcoh{f}{g}{\epsilon} & \defeq \sm{\eta : g\circ f \htpy \idfunc[A]} \prd{x:A} f(\eta x) = \epsilon (fx).
\end{align*}
\end{defn}

\begin{lem}\label{lem:coh-hfib}
对于任何 $f,g,\epsilon,\eta$，我们有
\begin{align*}
\lcoh{f}{g}{\eta} & \eqvsym {\prd{y:B} \id[\hfib{g}{gy}]{(fgy,\eta(gy))}{(y,\refl{gy})}}, \\
\rcoh{f}{g}{\epsilon} & \eqvsym {\prd{x:A} \id[\hfib{f}{fx}]{(gfx,\epsilon(fx))}{(x,\refl{fx})}}.
\end{align*}
\end{lem}

\begin{proof}
利用 \cref{lem:hfib}。
\end{proof}

\begin{lem}\label{lem:coh-hprop}
  如果 $f$ 是一个半伴随等价，那么对于任何 $(g,\epsilon) : \rinv(f)$，类型 $\rcoh{f}{g}{\epsilon}$ 都是可缩的。
\end{lem}

\begin{proof}
  由 \cref{lem:coh-hfib} 以及依值函数类型保持可缩空间这一事实，只需证明对每个 $x:A$，类型 $\id[\hfib{f}{fx}]{(gfx,\epsilon(fx))}{(x,\refl{fx})}$ 是可缩的。
  但由 \cref{thm:contr-hae}，$\hfib{f}{fx}$ 是可缩的，而任何可缩空间的道路空间自身也是可缩的。
\end{proof}

\begin{thm}\label{thm:hae-hprop}
  对于任何 $f : A \to B$，类型 $\ishae(f)$ 是一个命题。
\end{thm}

\begin{proof}
  由 \cref{ex:prop-inhabcontr}，可以假设 $f$ 是一个半伴随等价，并证明 $\ishae(f)$ 是可缩的。
  现在根据 $\Sigma$ 类型的结合律（\cref{ex:sigma-assoc}），类型 $\ishae(f)$ 等价于
  \[\sm{u : \rinv(f)} \rcoh{f}{\proj{1}(u)}{\proj{2}(u)}.\]
  但由 \cref{lem:inv-hprop,lem:coh-hprop} 以及 $\Sigma$ 类型保持可缩性这一事实，后一类型也是可缩的。
\end{proof}

这样，我们就证明了 $\ishae(f)$ 具备了类型 $\isequiv(f)$ 所需的全部三个理想性质。
在接下来的两节中，我们将考虑另外几种可能性。

\index{equivalence!half adjoint|)}%
\index{half adjoint equivalence|)}%
\index{adjoint!equivalence!of types, half|)}%

\section{双可逆映射}
\label{sec:biinv}

\index{function!bi-invertible|(defstyle}%
\index{bi-invertible function|(defstyle}%
\index{equivalence!as bi-invertible function|(defstyle}%

使用 \cref{sec:hae} 中引入的语言，我们可以将 \cref{sec:basics-equivalences} 中提出的定义重述如下。

\begin{defn}\label{defn:biinv}
  我们称 $f:A\to B$ 是\define{双可逆的（bi-invertible）}，
  如果它同时有一个左逆和一个右逆：
  \[ \biinv (f) \defeq \linv(f) \times \rinv(f). \]
\end{defn}

在 \cref{sec:basics-equivalences} 中，我们证明了 $\qinv(f)\to\biinv(f)$ 以及 $\biinv(f)\to\qinv(f)$。
剩下需要证明的是以下定理。

\begin{thm}\label{thm:isprop-biinv}
  对于任何 $f:A\to B$，类型 $\biinv(f)$ 是一个命题。
\end{thm}

\begin{proof}
  可以假设 $f$ 是双可逆的，并证明 $\biinv(f)$ 是可缩的。
  但由于 $\biinv(f)\to\qinv(f)$，根据 \cref{lem:inv-hprop}，在这种情况下 $\linv(f)$ 和 $\rinv(f)$ 都是可缩的，而可缩类型的积是可缩的。
\end{proof}

注意，这也符合我们在 \cref{sec:hae} 开头提出的建议：我们将 $g$ 和 $\eta$ 结合成一个可缩的类型，并添加一项额外数据，该数据再与 $\epsilon$ 结合成另一个可缩的类型。
区别在于，我们不是添加一个\emph{高阶}数据（一条 2 维道路）来与 $\epsilon$ 结合，而是添加了一个\emph{低阶}数据（一个独立于左逆的右逆）。

\begin{cor}\label{thm:equiv-biinv-isequiv}
  对于任何 $f:A\to B$，我们有 $\eqv{\biinv(f)}{\ishae(f)}$。
\end{cor}

\begin{proof}
  我们有 $\biinv(f) \to \qinv(f) \to \ishae(f)$ 和 $\ishae(f) \to \qinv(f) \to \biinv(f)$。
  由于 $\ishae(f)$ 和 $\biinv(f)$ 都是命题，该等价由 \cref{lem:equiv-iff-hprop} 得出。
\end{proof}

\index{function!bi-invertible|)}%
\index{bi-invertible function|)}%
\index{equivalence!as bi-invertible function|)}%

\section{具有可缩纤维的映射}
\label{sec:contrf}

\index{function!contractible|(defstyle}%
\index{contractible!function|(defstyle}%
\index{equivalence!as contractible function|(defstyle}%

注意，我们关于 $\ishae(f)$ 和 $\biinv(f)$ 的证明，本质上使用了"等价映射的纤维是可缩的"这一事实。
事实上，这一性质本身就可以作为等价的充分定义。

\begin{defn}[可缩映射] \label{defn:equivalence}
  一个映射 $f:A\to B$ 被称为\define{可缩的（contractible）}，
  如果对所有 $y:B$，纤维 $\hfib f y$ 都是可缩的。
\end{defn}

因此，类型 $\iscontr(f)$ 定义为
\begin{align}
  \iscontr(f) &\defeq \prd{y:B} \iscontr(\hfib f y)\label{eq:iscontrf}
  % \\
  % &\defeq \prd{y:B} \iscontr (\setof{x:A | f(x) = y}).
\end{align}
注意在 \cref{sec:contractibility} 中我们定义了\emph{类型}可缩的含义。
这里我们定义的是\emph{映射}可缩的含义。
我们的术语遵循同伦论的一般惯例：如果一个映射的所有（同伦）纤维都具有某个性质，我们就说这个映射具有该性质。
因此，一个类型 $A$ 可缩当且仅当映射 $A\to\unit$ 可缩。
从 \cref{cha:hlevels} 开始，我们也将可缩的映射和类型称为\emph{$(-2)$-截断的}。

我们已经在 \cref{thm:contr-hae} 中证明了 $\ishae(f) \to \iscontr(f)$。
反过来：

\begin{thm}\label{thm:lequiv-contr-hae}
对任意 $f:A\to B$，有 ${\iscontr(f)} \to {\ishae(f)}$。
\end{thm}

\begin{proof}
令 $P : \iscontr(f)$。我们如下定义逆映射 $g : B \to A$：将每个 $y : B$ 送到 $y$ 处纤维的收缩中心：
\[ g(y) \defeq \proj{1}(\proj{1}(Py)). \]
于是可以如下定义同伦 $\epsilon$：将 $y$ 映为 $g(y)$ 确实属于 $y$ 处纤维的见证：
\[ \epsilon(y) \defeq \proj{2}(\proj{1}(P y)). \]
剩下还需要定义 $\eta$ 和 $\tau$。这当然相当于给出 $\rcoh{f}{g}{\epsilon}$ 的一个元素。由 \cref{lem:coh-hfib}，这等价于对每个 $x:A$，在 $f$ 在 $fx$ 处的纤维中给出一条从 $(gfx,\epsilon(fx))$ 到 $(x,\refl{fx})$ 的道路。但这一点很容易：对于任意 $x : A$，类型 $\hfib{f}{fx}$
由假设是可缩的，因此这样一条道路必然存在。我们可以显式地将其构造为
\[\opp{\big(\proj{2}(P(fx))(gfx,\epsilon(fx))\big)} \ct \big(\proj{2}(P(fx)) (x,\refl{fx})\big). \qedhere \]
\end{proof}

也很容易看出：

\begin{lem}\label{thm:contr-hprop}
  对任意 $f$，类型 $\iscontr(f)$ 是一个命题。
\end{lem}

\begin{proof}
  由 \cref{thm:isprop-iscontr}，每个类型 $\iscontr (\hfib f y)$ 都是命题。
  因此，由 \cref{thm:isprop-forall}，\eqref{eq:iscontrf} 也是命题。
\end{proof}

\begin{thm}\label{thm:equiv-contr-hae}
  对任意 $f:A\to B$，有 $\eqv{\iscontr(f)}{\ishae(f)}$。
\end{thm}

\begin{proof}
  我们已经建立了逻辑等价 ${\iscontr(f)} \Leftrightarrow {\ishae(f)}$，并且两者都是命题（\cref{thm:contr-hprop,thm:hae-hprop}）。
  因此 \cref{lem:equiv-iff-hprop} 适用。
\end{proof}

通常我们通过展示拟逆来证明一个函数是等价，但有时上述定义更为方便。
例如，它意味着在证明一个函数是等价时，我们可以自由地假设其陪域是有实例的。

\begin{cor}\label{thm:equiv-inhabcod}
  如果 $f:A\to B$ 满足 $B\to \isequiv(f)$，那么 $f$ 是一个等价。
\end{cor}

\begin{proof}
  要证明 $f$ 是等价，只需证明对任意 $y:B$，$\hfib f y$ 是可缩的。
  但如果有 $e:B\to \isequiv(f)$，那么给定任意这样的 $y$，我们有 $e(y):\isequiv(f)$，因此 $f$ 是等价，从而 $\hfib f y$ 可缩，如所欲证。
\end{proof}

\index{function!contractible|)}%
\index{contractible!function|)}%
\index{equivalence!as contractible function|)}%

\section{关于等价的定义}
\label{sec:concluding-remarks}

\indexdef{equivalence}
我们已经证明等价的三种定义都满足三条所期望的性质，并且彼此等价：
\[ \iscontr(f) \eqvsym \ishae(f) \eqvsym \biinv(f). \]
（等价还有更多可能的定义，但我们只考虑这三种。
更多定义可参见 \cref{ex:brck-qinv} 和本章的习题。）
因此，我们可以选择其中任意一种作为 $\isequiv (f)$ 的"那个"定义。
为明确起见，我们选择定义
\[ \isequiv(f) \defeq \ishae(f).\]
\index{mathematics!formalized}%
这一选择对形式化有利，因为 $\ishae(f)$ 包含了最直接有用的数据。
另一方面，对于其他目的，$\biinv(f)$ 往往更容易处理，因为它不包含 $2$-维道路，而且它的两个对称部分可以独立处理。
不过，就本书而言，具体选择哪一种差别不大。

在本章的剩余部分，我们将研究等价的其他一些性质与刻画。
\index{equivalence!properties of}%

\section{满射与嵌入}
\label{sec:mono-surj}

\index{set}
当 $A$ 和 $B$ 都是集合且 $f:A\to B$ 是一个等价时，我们也称它为\define{同构（isomorphism）}
\indexdef{isomorphism!of sets}%
或\define{双射（bijection）}。
\indexdef{bijection}%
\indexsee{function!bijective}{bijection}%
（对于不是集合的类型，我们避免使用这些词，因为在同伦论和高阶范畴论中，它们通常表示比同伦等价更严格的"相同"概念。）
在集合论中，一个函数是双射当且仅当它既是单射又是满射。
在类型论中同样如此，只要我们适当地表述这些条件。
为了清晰起见，在处理不是集合的类型时，我们将用\emph{嵌入（embedding）}一词代替单射。

\begin{defn}\label{defn:surj-emb}
  设 $f:A\to B$。
  \begin{enumerate}
  \item 若对每个 $b:B$ 都有 $\brck{\hfib f b}$，则称 $f$ 是\define{满的（surjective）}
    \indexsee{surjective!function}{function, surjective}%
    \indexdef{function!surjective}%
    （或称\define{满射（surjection）}）
    \indexsee{surjection}{function, surjective}%
    。
  \item 若对每对 $x,y:A$，函数 $\apfunc f : (\id[A]xy) \to (\id[B]{f(x)}{f(y)})$ 都是一个等价，则称 $f$ 是一个\define{嵌入（embedding）}
    \indexdef{function!embedding}%
    \indexsee{embedding}{function, embedding}%
    。
  \end{enumerate}
\end{defn}

换言之，$f$ 是满射当且仅当 $f$ 的每个纤维都仅知有实例，或者说，对于所有 $b:B$，都仅知存在一个 $a:A$ 使得 $f(a)=b$。
用传统逻辑记号来说，$f$ 是满射当 $\fall{b:B}\exis{a:A} (f(a)=b)$。
这必须与更强的断言 $\prd{b:B}\sm{a:A} (f(a)=b)$ 区分开来；若后者成立，则称 $f$ 是一个\define{可分裂满射（split surjection）}。
\indexsee{split!surjection}{function, split surjective}%
\indexsee{surjection!split}{function, split surjective}%
\indexsee{surjective!function!split}{function, split surjective}%
\indexdef{function!split surjective}%
（由于后一个类型等价于 $\sm{g:B\to A}\prd{b:B} (f(g(b))=b)$，是可分裂满射就等于是\emph{收缩}，如 \cref{sec:contractibility} 所定义。）
\index{retraction}%
\index{function!retraction}%

\cref{sec:axiom-choice} 中的选择公理确切地说就是：\emph{集合之间}的每个满射都是可分裂的。
然而，在泛等公理成立的前提下，说\emph{所有}满射都可分裂是根本不成立的。
在 \cref{thm:no-higher-ac} 中，我们构造了一个类型族 $Y:X\to \type$，使得 $\prd{x:X} \brck{Y(x)}$ 但 $\neg \prd{x:X} Y(x)$；
对于任何这样的族，其第一投影 $(\sm{x:X} Y(x)) \to X$ 就是一个不可分裂的满射。

如果 $A$ 和 $B$ 都是集合，那么由 \cref{lem:equiv-iff-hprop}，$f$ 是嵌入当且仅当
\begin{equation}
  \prd{x,y:A} (\id[B]{f(x)}{f(y)}) \to (\id[A]xy).\label{eq:injective}
\end{equation}
在这种情况下，我们说 $f$ 是\define{单的（injective）}，
\indexsee{injective function}{function, injective}%
\indexdef{function!injective}%
或称它是一个\define{单射（injection）}。
\indexsee{injection}{function, injective}%
对于不是集合的类型，我们避免使用这些词，因为它们可能被解读为 \eqref{eq:injective}，这对于非集合来说是一个性质不佳的概念。
同样成立的是，集合之间的函数是满射当且仅当它在某种合适意义下是\emph{满态射（epimorphism）}，但这对于更一般的类型也不成立，并且满射性通常是更重要的概念。

\begin{thm}\label{thm:mono-surj-equiv}
  一个函数 $f:A\to B$ 是等价，当且仅当它同时是满射和嵌入。
\end{thm}

\begin{proof}
  如果 $f$ 是一个等价，那么每个 $\hfib f b$ 都是可缩的，从而 $\brck{\hfib f b}$ 也是可缩的，所以 $f$ 是满射。
  并且我们在 \cref{thm:paths-respects-equiv} 中已经证明过，任何等价都是嵌入。

  反之，假设 $f$ 既是满射又是嵌入。
  令 $b:B$；我们来证明 $\sm{x:A}(f(x)=b)$ 是可缩的。
  因为 $f$ 是满射，仅知存在 $a:A$ 使得 $f(a)=b$。
  因此，$f$ 在 $b$ 处的纤维是有实例的；剩下只需证明它是一个命题。
  为此，假设给定 $x,y:A$ 以及 $p:f(x)=b$ 和 $q:f(y)=b$。
  那么由于 $\apfunc f$ 是一个等价，存在 $r:x=y$ 使得 $\apfunc f (r) = p \ct \opp q$。
  然而，利用 $\Sigma$-类型中道路的表征，后一个等式可重新整理为 $\trans{r}{p} = q$。
  于是，连同 $r$ 一起，它就在 $f$ 在 $b$ 处的纤维中给出了 $(x,p) = (y,q)$。
\end{proof}

\begin{cor}
  对于任意 $f:A\to B$，我们有
  \[ \isequiv(f) \eqvsym (\mathsf{isEmbedding}(f) \times \mathsf{isSurjective}(f)).\]
\end{cor}

\begin{proof}
  "是满射"和"是嵌入"都是命题；现在应用 \cref{lem:equiv-iff-hprop} 即可。
\end{proof}

当然，我们不能把这当作"等价"的定义，因为嵌入的定义本身就依赖于等价。
不过，这一刻画仍然有用；参见 \cref{sec:whitehead}。
我们将在 \cref{cha:hlevels} 中对其加以推广。

% \section{Fiberwise equivalences}
\section{等价的闭包性质}
\label{sec:equiv-closures}
\label{sec:fiberwise-equivalences}
\index{equivalence!properties of}%

% We end this chapter by observing some important closure properties of equivalences.
我们已经在 \cref{thm:equiv-eqrel} 中看到，等价在复合运算下是封闭的。
此外，还有：

\begin{thm}[2-out-of-3 性质]\label{thm:two-out-of-three}
  \index{2-out-of-3 property}%
  设 $f:A\to B$ 和 $g:B\to C$。
  如果 $f$、$g$ 和 $g\circ f$ 三者中有任意两个是等价，那么第三个也是等价。
\end{thm}

\begin{proof}
  若 $g\circ f$ 和 $g$ 是等价，则 $\opp{(g\circ f)} \circ g$ 是 $f$ 的一个拟逆。
  一方面，我们有 $\opp{(g\circ f)} \circ g \circ f \htpy \idfunc[A]$；另一方面，则有
  \begin{align*}
    f \circ \opp{(g\circ f)} \circ g
    &\htpy \opp g \circ g \circ f \circ \opp{(g\circ f)} \circ g\\
    &\htpy \opp g \circ g\\
    &\htpy \idfunc[B].
  \end{align*}
  类似地，若 $g\circ f$ 和 $f$ 是等价，则 $f\circ \opp{(g\circ f)}$ 是 $g$ 的一个拟逆。
\end{proof}

这是同伦论中关于等价的一个标准闭包条件。
同样广为人知的是，等价对收缩核也封闭，具体含义如下。

\index{retract!of a function|(defstyle}%

\begin{defn}\label{defn:retract}
  称一个函数 $g:A\to B$ 为函数 $f:X\to Y$ 的\define{收缩核}，
  如果存在一个图表
  \begin{equation*}
  \xymatrix{
    {A} \ar[r]^{s} \ar[d]_{g}
    &
    {X} \ar[r]^{r} \ar[d]_{f}
    &
    {A} \ar[d]^{g}
    \\
    {B} \ar[r]_{s'}
    &
    {Y} \ar[r]_{r'}
    &
    {B}
  }
\end{equation*}
  并满足以下条件：
  \begin{enumerate}
    \item 一个同伦 $R:r\circ s \htpy \idfunc[A]$。
    \item 一个同伦 $R':r'\circ s' \htpy\idfunc[B]$。
    \item 一个同伦 $L:f\circ s\htpy s'\circ g$。
    \item 一个同伦 $K:g\circ r\htpy r'\circ f$。
    \item 对每个 $a:A$，存在一条道路 $H(a)$，见证如下方块的交换性：
      \begin{equation*}
  \xymatrix@C=3pc{
    {g(r(s(a)))} \ar@{=}[r]^-{K(s(a))} \ar@{=}[d]_{\ap g{R(a)}}
    &
    {r'(f(s(a)))} \ar@{=}[d]^{\ap{r'}{L(a)}}
    \\
    {g(a)} \ar@{=}[r]_-{\opp{R'(g(a))}}
    &
    {r'(s'(g(a)))}
  }
\end{equation*}
  \end{enumerate}
\end{defn}

回顾 \cref{sec:contractibility}，其中我们定义了何为一个类型是另一个类型的收缩核。
那是上述定义在 $B$ 和 $Y$ 均为 $\unit$ 时的特例。
反过来，正如可缩性的情况一样，映射的收缩核会诱导其纤维的收缩核。

\begin{lem}\label{lem:func_retract_to_fiber_retract}
  如果函数 $g:A\to B$ 是函数 $f:X\to Y$ 的收缩核，那么对每个 $b:B$，
  $\hfib{g}b$ 是 $\hfib{f}{s'(b)}$ 的收缩核，
  其中 $s':B\to Y$ 如 \cref{defn:retract} 中所述。
\end{lem}

\begin{proof}
  假设 $g:A\to B$ 是 $f:X\to Y$ 的收缩核。那么对任意 $b:B$，我们有两个函数
  \begin{align*}
\varphi_b &:\hfiber{g}b\to\hfib{f}{s'(b)}, &
\varphi_b(a,p) & \defeq \pairr{s(a),L(a)\ct s'(p)},\\
\psi_b &:\hfib{f}{s'(b)}\to\hfib{g}b, &
\psi_b(x,q) &\defeq \pairr{r(x),K(x)\ct r'(q)\ct R'(b)}.
\end{align*}
  然后我们有 $\psi_b(\varphi_b({a,p}))\equiv\pairr{r(s(a)),K(s(a))\ct r'(L(a)\ct s'(p))\ct R'(b)}$。
  我们断言：对所有 $b:B$，$\psi_b$ 是一个以 $\varphi_b$ 为截面的收缩，
  也就是说，对所有 $(a,p):\hfib g b$，有 $\psi_b(\varphi_b({a,p}))= \pairr{a,p}$。
  换言之，我们希望证明
  \begin{equation*}
\prd{b:B}{a:A}{p:g(a)=b} \psi_b(\varphi_b({a,p}))= \pairr{a,p}.
\end{equation*}
  通过调换前两个 $\Pi$ 的位置并应用 \cref{thm:omit-contr} 的一种形式，上式等价于
  \begin{equation*}
\prd{a:A}\psi_{g(a)}(\varphi_{g(a)}({a,\refl{g(a)}}))=\pairr{a,\refl{g(a)}}.
\end{equation*}
  对任意 $a$，由 \cref{thm:path-sigma}，此配对的相等等价于一对相等。
  第一个分量由 $R(a):r(s(a))= a$ 可知相等，因此我们只需证明
  \begin{equation*}
\trans{R(a)}{K(s(a))\ct r'(L(a))\ct R'(g(a))} = \refl{g(a)}.
\end{equation*}
  但这个输运计算为 $\opp{g(R(a))}\ct K(s(a))\ct r'(L(a))\ct R'(g(a))$，所以所需的道路由 $H(a)$ 给出。
\end{proof}

\begin{thm}\label{thm:retract-equiv}
  如果 $g$ 是某个等价 $f$ 的收缩核，那么 $g$ 也是一个等价。
\end{thm}

\begin{proof}
  由 \cref{lem:func_retract_to_fiber_retract}，$g$ 的每个纤维都是 $f$ 的某个纤维的收缩核。
  因此，由 \cref{thm:retract-contr}，如果 $f$ 的所有纤维都是可缩的，那么 $g$ 的所有纤维也是可缩的。
\end{proof}

\index{retract!of a function|)}%

\index{fibration}%
\index{total!space}%
最后，我们说明逐纤维等价可以用全空间的等价来刻画。
为解释相关术语，回顾 \cref{sec:fibrations}：一个类型族 $P:A\to\type$ 可以看作以 $A$ 为底空间的纤维化，
其全空间为 $\sm{x:A} P(x)$，而该纤维化即投影映射 $\proj1:\sm{x:A} P(x) \to A$。
从这个角度看，给定两个类型族 $P,Q:A\to\type$，我们可以称一个函数 $f:\prd{x:A} (P(x)\to Q(x))$ 为\define{逐纤维映射}或\define{逐纤维变换}。
\indexsee{transformation!fiberwise}{fiberwise transformation}%
\indexsee{function!fiberwise}{fiberwise transformation}%
\index{fiberwise!transformation|(defstyle}%
\indexsee{fiberwise!map}{fiberwise transformation}%
\indexsee{map!fiberwise}{fiberwise transformation}
这样的映射会诱导一个全空间之间的函数：

\begin{defn}\label{defn:total-map}
  给定类型族 $P,Q:A\to\type$ 和一个映射 $f:\prd{x:A} P(x)\to Q(x)$，定义
  \begin{equation*}
    \total f  \defeq \lam{w}\pairr{\proj{1}w,f(\proj{1}w,\proj{2}w)} : \sm{x:A}P(x)\to\sm{x:A}Q(x).
  \end{equation*}
\end{defn}

\begin{thm}\label{fibwise-fiber-total-fiber-equiv}
  设 $f$ 是类型 $A$ 上两个族 $P$ 和 $Q$ 之间的一个逐纤维变换，并令 $x:A$ 和 $v:Q(x)$。
  那么我们有如下等价
  \begin{equation*}
\eqv{\hfib{\total{f}}{\pairr{x,v}}}{\hfib{f(x)}{v}}.
\end{equation*}
\end{thm}

\begin{proof}
  我们计算如下：
  \begin{align}
  \hfib{\total{f}}{\pairr{x,v}}
  & \jdeq \sm{w:\sm{x:A}P(x)}\pairr{\proj{1}w,f(\proj{1}w,\proj{2}w)}=\pairr{x,v}
  \notag \\
  & \eqv{}{} \sm{a:A}{u:P(a)}\pairr{a,f(a,u)}=\pairr{x,v}
  \tag{by~\cref{ex:sigma-assoc}} \\
  & \eqv{}{} \sm{a:A}{u:P(a)}{p:a=x}\trans{p}{f(a,u)}=v
  \tag{by \cref{thm:path-sigma}} \\
  & \eqv{}{} \sm{a:A}{p:a=x}{u:P(a)}\trans{p}{f(a,u)}=v
  \notag \\
  & \eqv{}{} \sm{u:P(x)}f(x,u)=v
  \tag{$*$}\label{eq:uses-sum-over-paths} \\
  & \jdeq \hfib{f(x)}{v}. \notag
\end{align}
  等价~\eqref{eq:uses-sum-over-paths} 由 \cref{thm:omit-contr,thm:contr-paths,ex:sigma-assoc} 得出。
\end{proof}

我们称一个逐纤维变换 $f:\prd{x:A} P(x)\to Q(x)$ 为\define{逐纤维等价}%
\indexdef{fiberwise!equivalence}%
\indexdef{equivalence!fiberwise}
如果每个 $f(x):P(x) \to Q(x)$ 都是一个等价。

\begin{thm}\label{thm:total-fiber-equiv}
假设 $f$ 是类型 $A$ 上两个族 $P$ 和 $Q$ 之间的一个逐纤维变换。
那么 $f$ 是逐纤维等价当且仅当 $\total{f}$ 是一个等价。
\end{thm}

\begin{proof}
令 $f$、$P$、$Q$ 和 $A$ 如定理所述。
由 \cref{fibwise-fiber-total-fiber-equiv} 可知，对所有的 $x:A$ 和 $v:Q(x)$，$\hfib{\total{f}}{\pairr{x,v}}$ 可缩当且仅当 $\hfib{f(x)}{v}$ 可缩。
因此，对所有 $w:\sm{x:A}Q(x)$ 而言 $\hfib{\total{f}}{w}$ 可缩，当且仅当对所有 $x:A$ 和 $v:Q(x)$ 而言 $\hfib{f(x)}{v}$ 可缩。
\end{proof}

\index{fiberwise!transformation|)}%

\section{对象分类子}
\label{sec:object-classification}

在类型论中，\emph{类型族}是一个基本概念，即一个函数 $B:A\to\type$。
我们已经看到，这样的族在某种程度上类似于同伦论中的\emph{纤维化}，其中的纤维化即为投影 $\proj1:\sm{a:A} B(a) \to A$。
同伦论中的一个基本事实是：每个映射都等价于一个纤维化。
有了泛等性，我们也可以在类型论中证明同样的结论。

\begin{lem}\label{thm:fiber-of-a-fibration}
  对于任何类型族 $B:A\to\type$，$\proj1:\sm{x:A} B(x) \to A$ 在 $a:A$ 上的纤维等价于 $B(a)$：
  \[ \eqv{\hfib{\proj1}{a}}{B(a)} \]
\end{lem}

\begin{proof}
  我们有
  \begin{align*}
    \hfib{\proj1}{a} &\defeq \sm{u:\sm{x:A} B(x)} \proj1(u)=a\\
    &\eqvsym \sm{x:A}{b:B(x)} (x=a)\\
    &\eqvsym \sm{x:A}{p:x=a} B(x)\\
    &\eqvsym B(a)
  \end{align*}
  这里利用了相等类型的左泛性质。
\end{proof}

\begin{lem}\label{thm:total-space-of-the-fibers}
  对于任意函数 $f:A\to B$，我们有 $\eqv{A}{\sm{b:B}\hfib{f}{b}}$。
\end{lem}

\begin{proof}
  我们有
  \begin{align*}
    \sm{b:B}\hfib{f}{b} &\defeq \sm{b:B}{a:A} (f(a)=b)\\
    &\eqvsym \sm{a:A}{b:B} (f(a)=b)\\
    &\eqvsym A
  \end{align*}
  这里利用了 $\sm{b:B} (f(a)=b)$ 可缩这一事实。
\end{proof}

\begin{thm}\label{thm:nobject-classifier-appetizer}
对任意类型 $B$，存在等价
\begin{equation*}
\chi:\Parens{\sm{A:\type} (A\to B)}\eqvsym (B\to\type).
\end{equation*}
\end{thm}

\begin{proof}
我们需要构造拟逆
\begin{align*}
\chi & : \Parens{\sm{A:\type} (A\to B)}\to B\to\type\\
\psi & : (B\to\type)\to\Parens{\sm{A:\type} (A\to B)}.
\end{align*}
定义 $\chi$ 为 $\chi((A,f),b)\defeq\hfiber{f}b$，定义 $\psi$ 为 $\psi(P)\defeq\Pairr{(\sm{b:B} P(b)),\proj1}$。
现在需要验证 $\chi\circ\psi\htpy\idfunc{}$ 以及 $\psi\circ\chi \htpy\idfunc{}$。
\begin{enumerate}
\item 设 $P:B\to\type$。
  由 \cref{thm:fiber-of-a-fibration}，
$\hfiber{\proj1}{b}\eqvsym P(b)$ 对任意 $b:B$ 成立，因此立即得到 $P\htpy\chi(\psi(P))$。
\item 设 $f:A\to B$ 为一个函数。我们需要找到一条道路
\begin{equation*}
\Pairr{\tsm{b:B} \hfiber{f}b,\,\proj1}=\pairr{A,f}.
\end{equation*}
首先注意，由 \cref{thm:total-space-of-the-fibers}，我们有
$e:\sm{b:B} \hfiber{f}b\eqvsym A$ 其中 $e(b,a,p)\defeq a$ 且 $e^{-1}(a)
\defeq(f(a),a,\refl{f(a)})$。
根据 \cref{thm:path-sigma}，只需证明 $\trans{(\ua(e))}{\proj1} = f$。
但由泛等公理的计算规则以及 \eqref{eq:transport-arrow}，我们有 $\trans{(\ua(e))}{\proj1} = \proj1\circ e^{-1}$，而 $e^{-1}$ 的定义立即给出 $\proj1 \circ e^{-1} \jdeq f$。\qedhere
\end{enumerate}
\end{proof}

\noindent
\indexdef{object!classifier}%
\indexdef{classifier!object}%
\index{.infinity1-topos@$(\infty,1)$-topos}%
特别地，这意味着在高阶意象理论的意义下我们有一个\emph{对象分类子}。
回顾 \cref{def:pointedtype}，$\pointed\type$ 表示带点类型 $\sm{A:\type} A$ 构成的类型。

\begin{thm}\label{thm:object-classifier}
设 $f:A\to B$ 为一个函数。则图表
\begin{equation*}
  \vcenter{\xymatrix{
      A\ar[r]^-{\vartheta_f} \ar[d]_{f} &
      \pointed{\type}\ar[d]^{\proj1}\\
      B\ar[r]_{\chi_f} &
      \type
      }}
\end{equation*}
是一个\emph{拉回}\index{pullback}方块（参见 \cref{ex:pullback}）。
其中函数 $\vartheta_f$ 定义为
\begin{equation*}
 \lam{a} \pairr{\hfiber{f}{f(a)},\pairr{a,\refl{f(a)}}}.
\end{equation*}
\end{thm}

\begin{proof}
注意到我们有等价
\begin{align*}
A & \eqvsym \sm{b:B} \hfiber{f}b\\
& \eqvsym \sm{b:B}{X:\type}{p:\hfiber{f}b= X} X\\
& \eqvsym \sm{b:B}{X:\type}{x:X} \hfiber{f}b= X\\
& \eqvsym \sm{b:B}{Y:\pointed{\type}} \hfiber{f}b = \proj1 Y\\
& \jdeq B\times_{\type}\pointed{\type}
\end{align*}
由此得到一个复合等价 $e:A\eqvsym B\times_\type\pointed{\type}$。
我们可以逐步展示这个复合等价的作用如下：
\begin{align*}
a & \mapsto \pairr{f(a),\; \pairr{a,\refl{f(a)}}}\\
& \mapsto \pairr{f(a), \; \hfiber{f}{f(a)}, \; \refl{\hfiber{f}{f(a)}}, \; \pairr{a,\refl{f(a)}}}\\
& \mapsto \pairr{f(a), \; \hfiber{f}{f(a)}, \; \pairr{a,\refl{f(a)}}, \; \refl{\hfiber{f}{f(a)}}}\\
& \mapsto \pairr{f(a), \; \pairr{\hfiber{f}{f(a)}, \; \pairr{a,\refl{f(a)}}}, \; \refl{\hfiber{f}{f(a)}}}.
\end{align*}
因此，我们得到同伦 $f\htpy\proj1\circ e$ 和 $\vartheta_f\htpy \proj2\circ e$。
\end{proof}

\section{泛等性蕴含函数外延性}
\label{sec:univalence-implies-funext}

\index{function extensionality!proof from univalence}%
在本章的最后一节，我们将证明泛等公理蕴含函数外延性。因此，在本节中我们\emph{不使用}函数外延性公理。
证明分为两步。首先，在 \cref{uatowfe} 中我们将证明泛等公理蕴含一种弱形式的函数外延性，即下面 \cref{weakfunext} 所定义的弱函数外延性原理。
而弱函数外延性原理又能推出通常的函数外延性，且这一步推导不需要泛等公理（\cref{wfetofe}）。

\index{univalence axiom}%
设 $\type$ 为一个宇宙；我们会明确指出在何处假设它是泛等的。

\begin{defn}\label{weakfunext}
\define{弱函数外延性原理}
\indexdef{function extensionality!weak}%
断言对于任意类型 $A$ 上的任意类型族 $P:A\to\type$，存在一个函数
\begin{equation*}
\Parens{\prd{x:A}\iscontr(P(x))} \to\iscontr\Parens{\prd{x:A}P(x)}
\end{equation*}
\end{defn}

下面的引理在已有函数外延性时容易证明；但这里的要点是，它也可以单独从泛等公理推出，而不需额外假设函数外延性。

\begin{lem} \label{UA-eqv-hom-eqv}
假设 $\type$ 是泛等的，则对任意 $A,B,X:\type$ 和任意 $e:\eqv{A}{B}$，存在一个等价
\begin{equation*}
\eqv{(X\to A)}{(X\to B)}
\end{equation*}
其底层映射就是对 $e$ 的底层函数做后复合。
\end{lem}

\begin{proof}
  % Immediate by induction on $\eqv{}{}$ (see \cref{thm:equiv-induction}).
  如同在 \cref{lem:qinv-autohtpy} 的证明中，我们可以假设 $e = \idtoeqv(p)$，其中 $p:A=B$。
  然后通过道路归纳，可设 $p$ 为 $\refl{A}$，于是 $e = \idfunc[A]$。
  但此时与 $e$ 的后复合就是恒等映射，因此是一个等价。
\end{proof}

\begin{cor}\label{contrfamtotalpostcompequiv}
设 $P:A\to\type$ 为一族可缩类型，即 \narrowequation{\prd{x:A}\iscontr(P(x)).}
则投影 $\proj{1}:(\sm{x:A}P(x))\to A$ 是一个等价。假设 $\type$ 是泛等的，与 $\proj{1}$ 的后复合立即给出一个等价
\begin{equation*}
\alpha : \eqv{\Parens{A\to\sm{x:A}P(x)}}{(A\to A)}.
\end{equation*}
\end{cor}

\begin{proof}
  根据 \cref{thm:fiber-of-a-fibration}，对 $\proj{1}:(\sm{x:A}P(x))\to A$ 和 $x:A$ 我们有等价
  \begin{equation*}
    \eqv{\hfiber{\proj{1}}{x}}{P(x)}.
  \end{equation*}
  因此当每个 $P(x)$ 可缩时，$\proj{1}$ 是一个等价。断言现在是 \cref{UA-eqv-hom-eqv} 的推论。
\end{proof}

特别地，上述等价在 $\idfunc[A]$ 处的同伦纤维是可缩的。因此，要证明泛等公理蕴含弱函数外延性，我们只需证明依值函数类型 $\prd{x:A}P(x)$ 是 $\hfiber{\alpha}{\idfunc[A]}$ 的收缩核。

\begin{thm}\label{uatowfe}
在泛等宇宙 $\type$ 中，设 $P:A\to\type$ 是一族可缩类型，并令 $\alpha$ 为 \cref{contrfamtotalpostcompequiv} 中的函数。
则 $\prd{x:A}P(x)$ 是 $\hfiber{\alpha}{\idfunc[A]}$ 的收缩核。从而 $\prd{x:A}P(x)$ 是可缩的。换言之，泛等公理蕴含弱函数外延性原理。
\end{thm}

\begin{proof}
定义函数
\begin{align*}
  \varphi &: (\tprd{x:A}P(x))\to\hfiber{\alpha}{\idfunc[A]},\\
  \varphi(f) &\defeq (\lam{x} (x,f(x)),\refl{\idfunc[A]}),
\intertext{and}
  \psi &: \hfiber{\alpha}{\idfunc[A]}\to \tprd{x:A}P(x), \\
  \psi(g,p) &\defeq \lam{x} \trans {\happly (p,x)}{\proj{2} (g(x))}.
\end{align*}
则由依值函数类型的唯一性原理，$\psi(\varphi(f))=\lam{x} f(x)$ 就是 $f$。
\end{proof}

我们现在来证明，弱函数外延性蕴含通常的函数外延性。
回想 \eqref{eq:happly} 中的函数 $\happly (f,g) : (f = g)\to(f\htpy g)$，它将函数的相等转换为同伦。
下面的证明不使用泛等公理。

\begin{thm}\label{wfetofe}
  \index{function extensionality}%
弱函数外延性蕴含函数外延性 \cref{axiom:funext}。
\end{thm}

\begin{proof}
我们想要证明
\begin{equation*}
\prd{A:\type}{P:A\to\type}{f,g:\prd{x:A}P(x)}\isequiv(\happly (f,g)).
\end{equation*}
根据 \cref{thm:total-fiber-equiv}，一个逐纤维映射诱导全空间上的等价，当且仅当它逐纤维为等价。因此，我们只需证明由 $\lam{g:\prd{x:A}P(x)} \happly (f,g)$ 诱导的、类型为
\begin{equation*}
\Parens{\sm{g:\prd{x:A}P(x)}(f= g)} \to \sm{g:\prd{x:A}P(x)}(f\htpy g)
\end{equation*}
的函数是等价。
由于 \cref{thm:contr-paths}，左边的类型是可缩的，故只需证明右边的类型：
\begin{equation}\label{eq:uatofesp}
\sm{g:\prd{x:A}P(x)}\prd{x:A}f(x)= g(x)
\end{equation}
是可缩的。
根据 \cref{thm:ttac}，该类型等价于
\begin{equation}\label{eq:uatofeps}
\prd{x:A}\sm{u:P(x)}f(x)= u.
\end{equation}
\cref{thm:ttac} 的证明虽然用到了函数外延性，但仅涉及其中的一个复合。
因此，在不假设函数外延性的前提下，我们仍可断言：~\eqref{eq:uatofesp} 是~\eqref{eq:uatofeps} 的一个收缩核\index{retract!of a type}。
而~\eqref{eq:uatofeps} 是可缩类型的积，根据弱函数外延性原理，它是可缩的；因此~\eqref{eq:uatofesp} 也是可缩的。
\end{proof}

\sectionNotes

在代数拓扑中，一个众所周知的事实是：配备了拟逆的连续映射空间，其同伦类型并非"同伦等价空间"应有的同伦类型。
在代数拓扑中，"同伦等价空间" $(\eqv ab)$ 通常被简单地定义为函数空间 $(A\to B)$ 的一个子空间，由那些作为同伦等价的函数组成。
在类型论中，与之最接近的是 $\sm{f:A\to B} \brck{\qinv(f)}$；参见 \cref{ex:brck-qinv}。

在同伦类型论中，等价的第一个定义是我们称为 $\iscontr(f)$ 的那个，它由 Voevodsky 提出。
其他定义随后被多人相继发现。
关于伴随等价\index{adjoint!equivalence}的基本定理，例如 \cref{lem:coh-equiv,thm:equiv-iso-adj}，是改编自高阶范畴论和同伦论中的标准事实。
将双可逆性用作等价的定义，是由 André Joyal 提出的。

\cref{sec:mono-surj,sec:equiv-closures} 中讨论的等价性质在同伦论中是众所周知的。
其中大多数在类型论中的首次证明由 Voevodsky 完成。

每个函数都等价于一个纤维化，这是同伦论中的一个标准事实。
$(\infty,1)$-范畴
\index{.infinity1-category@$(\infty,1)$-category}%
理论中的对象分类子
\index{object!classifier}%
\index{classifier!object}%
概念（即 \cref{thm:nobject-classifier-appetizer} 的范畴论类比）是由 Rezk（雷兹克）提出的（参见~\cite{Rezk05,lurie:higher-topoi}）。

最后，泛等公理蕴含函数外延性（\cref{sec:univalence-implies-funext}）这一事实归功于 Voevodsky。
我们给出的证明是他的证明的简化版本。
\cref{ex:funext-from-nondep} 同样归功于 Voevodsky。

\sectionExercises

\begin{ex}\label{ex:two-sided-adjoint-equivalences}
  考虑函数 $f:A\to B$ 的"双侧伴随等价\index{adjoint!equivalence}数据"所构成的类型，
  \begin{narrowmultline*}
    \sm{g:B\to A}{\eta: g \circ f \htpy \idfunc[A]}{\epsilon:f \circ g \htpy \idfunc[B]}
    \narrowbreak
    \Parens{\prd{x:A} \map{f}{\eta x} = \epsilon(fx)} \times
    \Parens{\prd{y:B} \map{g}{\epsilon y} = \eta(gy) }.
  \end{narrowmultline*}
  根据 \cref{lem:coh-equiv}，我们知道若 $f$ 是一个等价，则该类型是有实例的。
  试给出该类型的一个类似于 \cref{lem:qinv-autohtpy} 的刻画。

  你能否给出一个例子，说明该类型通常不是一个命题？
  （在 \cref{cha:hits} 之后这会更容易。）
\end{ex}

\begin{ex}\label{ex:symmetric-equiv}
  证明对于任意 $A,B:\UU$，下述类型等价于 $\eqv A B$。
  \begin{equation*}
    \sm{R:A\to B\to \type}
    \Parens{\prd{a:A} \iscontr\Parens{\sm{b:B} R(a,b)}} \times
    \Parens{\prd{b:B} \iscontr\Parens{\sm{a:A} R(a,b)}}.
  \end{equation*}
  你能否由此提取出一种满足 $\isequiv(f)$ 三项要求的类型定义？
\end{ex}

\begin{ex} \label{ex:qinv-autohtpy-no-univalence}
  试在不使用泛等公理的前提下，重新表述 \cref{lem:qinv-autohtpy} 的证明。
\end{ex}

\begin{ex}[不稳定的八面体公理]\label{ex:unstable-octahedron}
  \index{axiom!unstable octahedral}%
  \index{octahedral axiom, unstable}%
  假设 $f:A\to B$，$g:B\to C$，$b:B$。
  \begin{enumerate}
  \item 证明存在一个自然映射 $\hfib{g\circ f}{g(b)} \to \hfib{g}{g(b)}$，其在 $(b,\refl{g(b)})$ 上的纤维等价于 $\hfib f b$。
  \item 证明 $\eqv{\hfib{g\circ f}{c}}{\sm{w:\hfib{g}{c}} \hfib f {\proj1 w}}$。
  \end{enumerate}
\end{ex}

\begin{ex}\label{ex:2-out-of-6}
  \index{2-out-of-6 property}%
  证明等价满足\emph{2-out-of-6 性质（2-out-of-6 property）}：给定 $f:A\to B$，$g:B\to C$ 和 $h:C\to D$，如果 $g\circ f$ 和 $h\circ g$ 是等价，那么 $f$，$g$，$h$ 以及 $h\circ g\circ f$ 也都是等价。
  利用这一点给出 \cref{thm:paths-respects-equiv} 的一个更高层次的证明。
\end{ex}

\begin{ex}\label{ex:qinv-univalence}
  对于 $A,B:\UU$，以显然的方式通过道路归纳定义
  \[ \mathsf{idtoqinv}_{A,B} :(A=B) \to \sm{f:A\to B}\qinv(f) \]
  。
  令 \textbf{\textsf{qinv}-泛等性}表示泛等公理的一种修改形式，它断言对于所有 $A,B:\UU$，函数 $\mathsf{idtoqinv}_{A,B}$ 都有一个拟逆。
  \begin{enumerate}
  \item 证明在 \cref{sec:univalence-implies-funext} 中函数外延性的证明里，可以用 \qinv-泛等性代替泛等公理。
  \item 证明在 \cref{thm:qinv-notprop} 的证明里，可以用 \qinv-泛等性代替泛等公理。
  \item 证明 \qinv-泛等性是不一致的（即可以构造出 $\emptyt$ 的一个实例）。
    由此可见，在陈述泛等公理时，使用 $\isequiv$ 的"好"版本是至关重要的。
  \end{enumerate}
\end{ex}

\begin{ex}\label{ex:embedding-cancellable}
  证明一个函数 $f:A\to B$ 是嵌入，当且仅当以下两个条件成立：
  \begin{enumerate}
  \item $f$ 是\emph{左可消的}，即对任意 $x,y:A$，若 $f(x)=f(y)$ 则 $x=y$。\label{item:ex:ec1}
  \item 对任意 $x:A$，映射 $\apfunc f: \Omega(A,x) \to \Omega(B,f(x))$ 是一个等价。\label{item:ex:ec2}
  \end{enumerate}
  （特别地，若 $A$ 是一个集合，则 $f$ 是嵌入当且仅当它是左可消的，并且对所有的 $x:A$，$\Omega(B,f(x))$ 都是可缩的。）
  请举例说明条件~\ref{item:ex:ec1} 与条件~\ref{item:ex:ec2} 互不蕴含。
\end{ex}

\begin{ex}\label{ex:cancellable-from-bool}
  证明左可消函数 $\bool\to B$ 的类型（见 \cref{ex:embedding-cancellable}）等价于 $\sm{x,y:B}(x\neq y)$。
  类似地，请给出嵌入 $\bool\to B$ 的类型的一个显式刻画。
\end{ex}

\begin{ex}\label{ex:funext-from-nondep}
  \textbf{朴素非依值函数外延性公理}说的是：对于 $A,B:\type$ 和 $f,g:A\to B$，存在一个函数 $(\prd{x:A} f(x)=g(x)) \to (f=g)$。
  \indexdef{function extensionality!non-dependent}%
  修改 \cref{sec:univalence-implies-funext} 中的论证，以此证明该公理蕴含完整的函数外延性公理（\cref{axiom:funext}）。
\end{ex}

% Local Variables:
% TeX-master: "hott-online"
% End:

\chapter{归纳}
\label{cha:induction}

在 \cref{cha:typetheory} 中，我们介绍了许多从已有类型构造新类型的方法。
除了（依值）函数类型和宇宙之外，所有这些规则都是\emph{归纳定义（inductive definition）}\index{definition!inductive}\index{inductive!definition}这一普遍概念的特殊情形。
在本章中，我们将更一般地研究归纳定义。

\section{归纳类型简介}
\label{sec:bool-nat}

\index{type!inductive|(}%
\index{generation!of a type, inductive|(}
直观上，一个\emph{归纳类型（inductive type）} $X$ 可以理解为由某个有限的\emph{构造子（constructors）}集合“自由生成”的类型；其中每个构造子都是一个以 $X$ 为陪域的函数（可以带有若干个参数）。
这也包括零个参数的函数，它们就是 $X$ 的元素。

在描述一个具体的归纳类型时，我们逐条列出其构造子。
例如，\cref{sec:type-booleans} 中的类型 \bool 是由以下两个构造子归纳生成的：

\begin{itemize}
\item $\bfalse:\bool$
\item $\btrue:\bool$
\end{itemize}

类似地，$\unit$ 由以下一个构造子归纳生成：

\begin{itemize}
\item $\ttt:\unit$
\end{itemize}

而 $\emptyt$ 则根本没有构造子。
构造子带参数的一个例子是余积 $A+B$，它由两个构造子生成：

\begin{itemize}
\item $\inl:A\to A+B$
\item $\inr:B\to A+B$。
\end{itemize}

构造子带多个参数的一个例子是 Cartesian 积 $A\times B$，它由一个构造子生成：

\begin{itemize}
\item $\pairr{\blank, \blank} : A\to B \to A\times B$。
\end{itemize}

至关重要的是，我们还允许归纳类型的构造子接受来自正在定义的归纳类型本身的参数。
例如，自然数类型 $\nat$ 具有构造子：

\begin{itemize}
\item $0:\nat$
\item $\suc:\nat\to\nat$。
\end{itemize}

\symlabel{lst}
\index{type!of lists}%
\indexsee{list}{type of lists}%
另一个有用的例子是由某类型 $A$ 的元素组成的有限列表类型 $\lst A$，它具有构造子：

\begin{itemize}\label{list_constructors}
\item $\nil:\lst A$
\item $\cons:A\to \lst A \to \lst A$。
\end{itemize}

\index{free!generation of an inductive type}%
直观上，我们应该将一个归纳类型理解为其构造子\emph{自由生成的}类型。
也就是说，归纳类型的元素恰好就是从无开始、反复应用构造子所能得到的一切。
（我们将在 \cref{sec:identity-systems,cha:hits} 看到，对于更一般的归纳定义，这个概念需要稍作修改，但目前这样理解已经足够。）
例如，对于 \bool，我们应当认为其元素只有 $\bfalse$ 和 $\btrue$。
类似地，对于 \nat，我们应当认为每个元素要么是 $0$，要么是对某个“先前构造的”自然数应用 $\suc$ 得到的。

\index{induction principle}%
然而，我们并非直接断言此类性质，而是通过\emph{归纳原理（induction principle）}，也称为\emph{（依值）消去规则（(dependent) elimination rule）}来表达它们。
我们在 \cref{cha:typetheory} 已经见过这些原理。
例如，\bool 的归纳原理是：
%

\begin{itemize}
\item 要证明关于 \bool 的\emph{所有}元素的命题 $E : \bool \to \type$，只需对 $\bfalse$ 和 $\btrue$ 分别证明即可，亦即给出证明 $e_0 : E(\bfalse)$ 和 $e_1 : E(\btrue)$。
\end{itemize}

此外，将所得的证明 $\ind\bool(E,e_0,e_1): \prd{b : \bool}E(b)$ 作用于构造子 $\bfalse$ 和 $\btrue$ 时，其行为符合预期；这一原理由\emph{计算规则（computation rules）}\index{computation rule!for type of booleans}表达：

\begin{itemize}
\item 我们有 $\ind\bool(E,e_0,e_1,\bfalse) \jdeq e_0$。
\item 我们有 $\ind\bool(E,e_0,e_1,\btrue) \jdeq e_1$。
\end{itemize}

\index{case analysis}%
因此，布尔类型 $\bool$ 的归纳原理允许我们进行\emph{分情形讨论（case analysis）}。
由于两个构造子都不带任何参数，这对布尔类型来说已经足够了。

\index{natural numbers}%
然而，对于自然数，仅靠分情形讨论通常是不够的：在对应于归纳步骤 $\suc(n)$ 的情形中，我们还需要假设待证的命题对 $n$ 已经成立。
这就得到了如下的归纳原理：

\begin{itemize}
\item 要证明关于\emph{所有}自然数的命题 $E : \nat \to \type$，只需对 $0$ 和 $\suc(n)$ 分别证明即可，其中在证明 $\suc(n)$ 时可以假设命题对 $n$ 成立，亦即我们构造 $e_z : E(0)$ 和 $e_s : \prd{n : \nat} E(n) \to E(\suc(n))$。
\end{itemize}

与布尔类型的情况类似，对于函数 $\ind\nat(E,e_z,e_s) : \prd{x:\nat} E(x)$，我们也有相应的计算规则：
\index{computation rule!for natural numbers}%

\begin{itemize}
\item $\ind\nat(E,e_z,e_s,0) \jdeq e_z$。
\item 对任意 $n : \nat$，有 $\ind\nat(E,e_z,e_s,\suc(n)) \jdeq e_s(n,\ind\nat(E,e_z,e_s,n))$。
\end{itemize}

因此，依值函数 $\ind\nat(E,e_z,e_s)$ 可以理解为根据参数 $x : \nat$ 递归定义的；这一定义通过函数 $e_z$ 和 $e_s$ 实现，我们称它们为\define{递归式（recurrences）}\indexdef{recurrence}。
当 $x$ 为零\index{zero}时，该函数直接返回 $e_z$。
当 $x$ 是另一个自然数 $n$ 的后继\index{successor}时，结果是这样得到的：取递归式 $e_s$，并将具体的前驱\index{predecessor} $n$ 以及递归调用的值 $\ind\nat(E,e_z,e_s,n)$ 代入其中。

上述所有例子的归纳原理都具有这种家族相似性。
在 \cref{sec:strictly-positive} 中，我们将讨论“归纳定义”的一般概念，以及如何为其推导出合适的\emph{归纳原理}；但首先，我们来探究归纳定义之间的若干共通之处。

\index{recursion principle}%
例如，我们曾在 \cref{cha:typetheory} 的每一个例子里指出，从归纳原理可以推导出一个\emph{递归原理}，其中值域是简单类型（而非类型族）。
归纳原理和递归原理可能看起来有些奇怪，因为它们似乎只给出了函数的\emph{存在性}，而没有刻画出其唯一性。
然而事实上，归纳原理本身已经足够强，足以证明其自身的\emph{唯一性原理}\index{uniqueness!principle, propositional!for functions on N@for functions on $\nat$}，如下述定理所示。

\begin{thm}\label{thm:nat-uniq}
设 $f,g : \prd{x:\nat} E(x)$ 为两个函数，它们在命题相等的意义下满足递归式
%
\begin{equation*}
  e_z : E(0)
  \qquad\text{and}\qquad
  e_s : \prd{n : \nat} E(n) \to E(\suc(n))
\end{equation*}
%
也就是说，满足
\begin{equation*}
  \id{f(0)}{e_z}
  \qquad\text{and}\qquad
  \id{g(0)}{e_z}
\end{equation*}
以及
\begin{gather*}
  \prd{n : \nat} \id{f(\suc(n))}{e_s(n, f(n))},\\
  \prd{n : \nat} \id{g(\suc(n))}{e_s(n, g(n))}.
\end{gather*}
那么 $f$ 和 $g$ 相等。
\end{thm}

\begin{proof}
我们在类型族 $D(x) \defeq \id{f(x)}{g(x)}$ 上使用归纳法。对于基始情形，有
\[ f(0) = e_z = g(0). \]
对于归纳情形，假设 $n : \nat$ 且 $f(n) = g(n)$。那么
\[ f(\suc(n)) = e_s(n, f(n)) = e_s(n, g(n)) = g(\suc(n)). \]
第一个和最后一个等式由 $f$ 和 $g$ 所满足的条件得出。中间的等式则由归纳假设以及函数应用保持相等性这一事实得出。这样我们就得到了 $f$ 与 $g$ 的逐点相等；再由函数外延性即可完成证明。
\end{proof}

注意，唯一性原理甚至适用于那些仅\emph{在命题相等意义下}满足递归式的函数，亦即一条道路。
当然，由归纳原理直接得到的那个特定函数是在判断相等性\index{judgmental equality}的意义下满足这些递归式的；我们将在 \cref{sec:htpy-inductive} 回到这一点。
另一方面，定理本身只断言了函数之间的一个命题相等（亦见 \cref{ex:same-recurrence-not-defeq}）。
从同伦的观点来看，很自然地会问这条道路是否\emph{相容}，即等式 $f=g$ 在高阶道路意义下是否唯一；在 \cref{sec:initial-alg} 我们将会看到事实的确如此。

当然，类似的函数唯一性定理通常也能为其它归纳类型表述并证明。
在下一节中，我们将展示这一唯一性性质如何与泛等性结合，从而推出像自然数这样的归纳类型完全由其引入规则、消去规则和计算规则所刻画。

\index{type!inductive|)}%
\index{generation!of a type, inductive|)}%

%%%%%%%%%%%%%%%%%%%%%%%%%%%%%%%%%%%%%%%%%%%%%%%%%%%%%%%%%%

\section{归纳类型的唯一性}
\label{sec:appetizer-univalence}

\index{uniqueness!of identity types}%
我们已将“那个”自然数定义为一个特定的类型 \nat，并配有特定的归纳构造子 $0$ 和 $\suc$。
然而，根据前一节描述的类型论中归纳定义的一般原理，我们完全可以用相同的方式定义\emph{另一个}类型。
\index{natural numbers!isomorphic definition of}
也就是说，假设我们令 $\natp$ 是由以下构造子生成的归纳类型：

\begin{itemize}
\item $\zerop:\natp$
\item $\sucp:\natp\to\natp$。
\end{itemize}

那么 $\natp$ 将具有与 $\nat$ 形式上相同的归纳原理和递归原理。
当要为所有这些“新”自然数证明一个命题 $E : \natp \to \type$ 时，只需给出证明 $e_z : E(\zerop)$ 和 \narrowequation{e_s : \prd{n : \natp} E(n) \to E(\sucp(n))。}
而函数 $\rec\natp(E,e_z,e_s) : \prd{n:\natp} E(n)$ 则具有如下计算规则：

\begin{itemize}
\item $\rec\natp(E,e_z,e_s,\zerop) \jdeq e_z$，
\item $\rec\natp(E,e_z,e_s,\sucp(n)) \jdeq e_s(n,\rec\natp(E,e_z,e_s,n))$ 对任意 $n : \natp$ 成立。
\end{itemize}

但是，$\nat$ 和 $\natp$ 之间有什么关系呢？

这并非仅仅是一个学术问题，因为许多其他地方也能找到``像''自然数的结构。
\index{type!of lists}%
例如，我们可以将自然数与单元素类型上的列表视为等同（这可以说是最古老的呈现方式——见于洞穴墙壁上\index{caves, walls of}），与非负整数、有理数和实数的某些子集等等视为等同。
而从编程\index{programming}的角度看，自然数的这种``一元''表示效率非常低，因此我们有时可能更倾向于使用二进制表示。
我们希望能够在所有这些``自然数''版本之间相互等同，以便将构造和结果从一个版本转移到另一个版本。

当然，如果自然数的两个版本满足完全相同的归纳原理，那么它们导出的结构也完全相同。
譬如，回顾在\cref{sec:inductive-types}中定义的函数 $\dbl$ 的例子。为我们的新自然数定义一个类似的函数，只需依样复制并加上撇号即可：
\[ \dblp \defeq \rec\natp(\natp, \; \zerop, \;  \lamu{n:\natp}{m:\natp} \sucp(\sucp(m))). \]
这看起来简单，却有一个明显的缺点：会导致大量重复。
不仅函数需要复制，所有的引理及其证明也得复制。
例如，像$\prd{n : \nat} \dbl(\suc(n))=\suc(\suc(\dbl(n)))$这样一个简单的结果，连同其归纳证明，也必须被``加上撇号''。

在传统数学中，人们直接宣称 $\nat$ 和 $\natp$ 显然``相同''，并且可以在需要时互相替换。
这通常没什么问题，但它把相当多的东西扫到了地毯下面，拉大了非形式化数学与其精确描述之间的差距。
在同伦类型论中，我们可以做得更好。

首先注意到我们有如下可定义的映射：

\begin{itemize}
\item $f \defeq \rec\nat(\natp, \; \zerop, \;  \lamu{n:\nat} \sucp)
       : \nat \to\natp$,
\item $g \defeq \rec\natp(\nat, \; 0, \;  \lamu{n:\natp} \suc)
       : \natp \to\nat$。
\end{itemize}

由于 $g$ 和 $f$ 的复合满足与 $\nat$ 上恒等函数相同的递归式，由\cref{thm:nat-uniq}可得$\prd{n : \nat} \id{g(f(n))}{n}$，而同一定理的``加撇''版本则给出$\prd{n : \natp} \id{f(g(n))}{n}$。
因此，$f$ 和 $g$ 互为拟逆，故有 $\eqv{\nat}{\natp}$。
现在，我们可以沿着这个等价，将 $\nat$ 上的函数直接转移到 $\natp$ 上（反之亦然），例如：
\[ \dblp \defeq \lamu{n:\natp} f(\dbl(g(n))). \]
证明这个版本的 $\dblp$ 与早先定义的相等，是一个简单的练习。

当然，这没什么好惊讶的；数学家正是通过这样的同构来设想如何将 $\nat$ 与 $\natp$``等同''起来的。
\index{isomorphism!transfer across}%
然而，通过同构进行``转移''的机制，取决于被转移的对象；它并不总是像用 $f$ 和 $g$ 对单个函数做前复合与后复合那么简单。
考虑，例如，这样一个简单的引理：
\[\prd{n : \natp} \dblp(\sucp(n))=\sucp(\sucp(\dblp(n))).\]
插入正确的 $f$ 和 $g$，比起直接对 $n:\natp$ 重新进行归纳证明，也好不了太多。

\index{univalence axiom}%
正是在此处，泛等公理发挥了作用：由于 $\eqv{\nat}{\natp}$，我们还有 $\id[\type]{\nat}{\natp}$，即 $\nat$ 和 $\natp$ 作为类型是\emph{相等}的。
现在，关于相等类型的归纳原理保证了任何涉及 $\nat$ 的构造或证明，都能以同样的方式自动转移到 $\natp$ 上。
我们只需将该函数或定理的类型视为一个以类型为指标的类型族 $P:\type\to\type$，并将给定的对象看作 $P(\nat)$ 中的一个元素，然后沿着道路 $\id \nat\natp$ 进行输运即可。
这涉及的额外开销要小得多。

为简单起见，我们是在 \nat 和 \natp 两个类型的定义\emph{看起来相同}的情况下描述这一方法的。
然而，实践中更常见的情形是，两个定义并非字面上相同，但其中一个归纳原理蕴含另一个。
\index{type!unit}%
\index{natural numbers!encoded as list(unit)@encoded as $\lst\unit$}%
例如，考虑单元素类型上的列表类型 $\lst\unit$，它由以下两者生成：

\begin{itemize}
\item 一个元素 $\nil:\lst\unit$，以及
\item 一个函数 $\cons:\unit \times \lst\unit \to\lst\unit$。
\end{itemize}

这与 \nat 的定义并不相同，也不会产生相同的归纳原理。
$\lst\unit$ 的归纳原理是说，对于任意 $E:\lst\unit\to\type$ 以及递归式数据 $e_\nil:E(\nil)$ 和 $e_\cons : \prd{u:\unit}{\ell:\lst\unit} E(\ell) \to E(\cons(u,\ell))$，存在 $f:\prd{\ell:\lst\unit} E(\ell)$ 使得 $f(\nil)\jdeq e_\nil$ 且 $f(\cons(u,\ell))\jdeq e_\cons(u,\ell,f(\ell))$。
（我们将在\cref{sec:strictly-positive}看到如何推导一个归纳定义的归纳原理。）

现在假设我们定义 $0'' \defeq \nil: \lst\unit$，并用 $\suc''(\ell) \defeq \cons(\ttt,\ell)$ 定义 $\suc'':\lst\unit\to\lst\unit$。
那么，对于任意 $E:\lst\unit\to\type$ 以及 $e_0:E(0'')$ 和 $e_s:\prd{\ell:\lst\unit} E(\ell) \to E(\suc''(\ell))$，我们可以定义
\begin{align*}
  e_\nil &\defeq e_0\\
  e_\cons(\ttt,\ell,x) &\defeq e_s(\ell,x).
\end{align*}
（在定义 $e_\cons$ 时，我们利用了 \unit 的归纳原理来假设 $u$ 就是 $\ttt$。）
现在我们可以应用 $\lst\unit$ 的归纳原理，得到 $f:\prd{\ell:\lst\unit} E(\ell)$ 使得
\begin{gather*}
  f(0'') \jdeq f(\nil) \jdeq e_\nil \jdeq e_0\\
  f(\suc''(\ell)) \jdeq f(\cons(\ttt,\ell)) \jdeq e_\cons(\ttt,\ell,f(\ell)) \jdeq e_s(\ell,f(\ell)).
\end{gather*}
因此，$\lst\unit$ 满足与 $\nat$ 相同的归纳原理，从而（由与上述相同的论证）等于 $\nat$。

最后，这些结论并不局限于自然数：它们适用于任何归纳类型。
如果我们有一个归纳定义的类型 $W$，以及另一个类型 $W'$，满足与 $W$ 相同的归纳原理，那么可以推出 $\eqv{W}{W'}$，从而 $W=W'$。
我们分别利用 $W$ 和 $W'$ 导出的递归原理来构造映射 $W\to W'$ 和 $W'\to W$，然后利用各自的归纳原理证明两个复合都等于恒等映射。
例如，在\cref{cha:typetheory}中我们看到，余积 $A+B$ 也可以定义为$\sm{x:\bool} \rec{\bool}(\UU,A,B,x)$。
后者满足与前者相同的归纳原理；因此它们是典范等价的。

当然，这非常类似于范畴论中一个熟悉的事实：如果两个对象具有相同的\emph{泛性质}，那么它们等价。
在\cref{sec:initial-alg}中我们将看到，归纳类型实际上确实具有一个泛性质，因此这是那个一般原理的一个体现。
\index{universal property}

\section{\texorpdfstring{$\w$}{W}-类型}
\label{sec:w-types}

归纳类型的概念非常宽泛，这固然有利于它们的有用性和适用性，但也使得它们难以作为一个整体来研究。
幸运的是，它们都可以形式化地归约到几个特殊情形。
讨论这一归约过程超出了本书的范围——反正它对实践中使用类型论的数学家来说也无关紧要——但我们将花少许时间讨论一个我们尚未遇到的基本特殊情形。
这些就是 Martin-L{\"o}f 的\emph{$\w$-类型（$\w$-types）}，也称为\emph{良基树（well-founded trees）类型}。
\index{tree, well-founded}%
$\w$-类型是自然数、列表和二叉树等类型的推广，它足够宽泛，能涵盖\emph{任意}归纳类型的``递归''层面。

一个特定的 $\w$-类型通过给出两个参数 $A : \type$ 和 $B : A \to \type$ 来确定，在这种情况下，所得的 $\w$-类型记作 $\wtype{a:A} B(a)$。
类型 $A$ 表示 $\wtype{a :A} B(a)$ 的\emph{标签（labels）}类型，其作用类似于构造子（不过，我们将``构造子''一词保留给归纳定义中实际产生的函数）。
例如，当用 $\w$-类型定义自然数时，%
\index{natural numbers!encoded as a W-type@encoded as a $\w$-type}
类型 $A$ 将是包含两个元素 $\bfalse$ 和 $\btrue$ 的类型 $\bool$，因为恰好有两种方式得到一个自然数——它要么是零，要么是另一个自然数的后继。

类型族 $B : A \to \type$ 用于记录标签的\emph{元数（arity）}\index{arity}：一个标签 $a : A$ 将接受一族以 $B(a)$ 为指标的归纳参数。
因此，我们可以把 $B(a)$ 看作是标签 $a$ 的参数``个数''。
这些参数通过函数 $f : B(a) \to \wtype{a :A} B(a)$ 表示，其含义是，对于任何 $b : B(a)$，$f(b)$ 是该标签 $a$ 的``第 $b$ 个''参数。
因此，$\w$-类型 $\wtype{a :A} B(a)$ 可以被看作是良基树的类型，其中节点由 $A$ 中的元素标记，并且每个标记为 $a : A$ 的节点都有 $B(a)$ 条分支。

在自然数的例子中，标签 $\bfalse$ 的元数（arity）为 $0$，因为它构造的是常量零\index{zero}；标签 $\btrue$ 的元数为 $1$，因为它构造的是其参数的后继\index{successor}。
我们可以通过对 $\bool$ 做简单的消去来刻画这一点，定义一个到类型宇宙的函数 $\rec\bool(\bbU,\emptyt,\unit)$；对于 $\bfalse$，该函数返回空类型 $\emptyt$，对于 $\btrue$，则返回单位\index{type!unit}类型 $\unit$。
因此我们可以定义
\symlabel{natw}
\index{type!unit}%
\index{type!empty}%
\[ \natw \defeq \wtype{b:\bool} \rec\bool(\bbU,\emptyt,\unit,b) \]
其中上标 $\mathbf{w}$ 用于区分这个版本的自然数与之前使用的版本。
类似地，我们可以将 $A$ 上的列表\index{type!of lists}类型定义为一个具有 $\unit + A$ 个标签的 $\w$-类型：一个零元标签对应空列表，加上每个 $a : A$ 的一个一元标签，对应将 $a$ 附加到列表头部：
\[ \lst A \defeq \wtype{x: \unit + A} \rec{\unit + A}(\bbU, \; \emptyt, \; \lamu{a:A} \unit,\; x). \]
%
\index{W-type@$\w$-type}%
\indexsee{type!W-@$\w$-}{$\w$-type}%
一般而言，由 $A : \type$ 和 $B : A \to \type$ 所指定的 $\w$-类型 $\wtype{x:A} B(x)$，是由以下构造子生成的归纳类型：

\begin{itemize}
\item \label{defn:supp}
  $\supp : \prd{a:A} \Big(B(a) \to \wtype{x:A} B(x)\Big) \to \wtype{x:A} B(x)$.
\end{itemize}

%
构造子 $\supp$（上确界\index{supremum!constructor of a W-type@constructor of a $\w$-type}的缩写）接受一个标签 $a : A$ 以及一个函数 $f : B(a) \to \wtype{x:A} B(x)$（表示 $a$ 的参数），并构造出 $\wtype{x:A} B(x)$ 的一个新元素。
利用前面对自然数的 $\w$-类型编码，我们可以定义例如
\begin{equation*}
\zerow \defeq \supp(\bfalse, \; \lamu{x:\emptyt} \rec\emptyt(\natw,x)).
\end{equation*}
换句话说，我们使用标签 $\bfalse$ 来构造 $\zerow$。
而 $\rec\bool(\bbU,\emptyt,\unit, \bfalse)$ 会计算为 $\emptyt$，这符合预期，因为 $\bfalse$ 是一个零元标签。
因此，我们需要构造一个函数 $f : \emptyt \to \natw$，它表示提供给 $\bfalse$ 的（零个）参数。这当然是平凡的，如下所示，对 $\emptyt$ 做简单的消去即可。
类似地，我们可以定义 $1^{\mathbf{w}}$ 和一个后继函数 $\sucw$
\begin{align*}
1^{\mathbf{w}} &\defeq \supp(\btrue, \; \lamu{x:\unit} 0^{\mathbf{w}}) \\
\sucw&\defeq \lamu{n:\natw} \supp(\btrue, \; \lamu{x:\unit} n).
\end{align*}

\index{induction principle!for W-types@for $\w$-types}%
对于 $\w$-类型，我们有如下的归纳原理：

\begin{itemize}
\item 要证明关于 $\w$-类型 $\wtype{x:A} B(x)$ 中\emph{所有}元素的某个命题 $E : \big(\wtype{x:A} B(x)\big) \to \type$，只需对形如 $\supp(a,f)$ 的元素证明它，同时假设它对所有满足 $b : B(a)$ 的 $f(b)$ 成立。
换言之，只需给出一个证明
\begin{equation*}
e : \prd{a:A}{f : B(a) \to \wtype{x:A} B(x)}{g : \prd{b : B(a)} E(f(b))} E(\supp(a,f))
\end{equation*}
\end{itemize}

\index{variable}%
变量 $g$ 代表了我们的归纳假设，即 $a$ 的所有参数都满足 $E$。
为了表述这一点，我们对 $B(a)$ 类型的所有元素进行量化，因为每个 $b : B(a)$ 对应 $a$ 的一个参数 $f(b)$。

对于编码为 $\w$-类型的自然数，我们如何定义函数 $\dbl$？
我们希望使用 $\natw$ 的递归原理，其值域是 $\natw$ 本身。
因此需要构造一个合适的函数
\[e : \prd{a : \bool}{f : B(a) \to \natw}{g : B(a) \to \natw} \natw\]
它将代表 $\dbl$ 函数的递归式；为简便起见，我们将类型族 $\rec\bool(\bbU,\emptyt,\unit)$ 记作 $B$。

显然，$e$ 将是一个以 $a : \bool$ 作为其第一个参数的函数。
下一步是对 $a$ 进行分情形讨论，根据它是 $\bfalse$ 还是 $\btrue$ 来分别处理。
这提示了如下的形式
\[ e \defeq \lamu{a:\bool} \ind\bool(C,e_0,e_1,a) \]
其中
\[C \defeq \lamu{a:\bool} \prd{f : B(a) \to \natw}{g : B(a) \to \natw} \natw.\]
如果 $a$ 是 $\bfalse$，则类型 $B(a)$ 变为 $\emptyt$。
因此，给定 $f : \emptyt \to \natw$ 和 $g : \emptyt \to \natw$，我们希望构造一个 $\natw$ 的元素。
由于标签 $\bfalse$ 代表 $\emptyt$，它需要零个归纳参数，变量 $f$ 和 $g$ 无关紧要。
我们返回 $\zerow$\index{zero} 作为结果：
\[ e_0 \defeq \lamu{f:\emptyt \to \natw}{g:\emptyt \to \natw} \zerow. \]
类似地，如果 $a$ 是 $\btrue$，则类型 $B(a)$ 变为 $\unit$。
由于标签 $\btrue$ 代表后继\index{successor}运算，它需要一个归纳参数——即前驱\index{predecessor}——由变量 $f : \unit \to \natw$ 表示。
在前驱上的递归调用的值由变量 $g : \unit \to \natw$ 表示。
于是，取这个值（即 $g(\ttt)$）并应用后继函数两次，便得到想要的结果：
\begin{equation*}
e_1 \defeq \; \lamu{f:\unit \to \natw}{g:\unit \to \natw} \sucw(\sucw(g(\ttt))).
\end{equation*}
综上，我们得到
\[ \dbl \defeq \rec\natw(\natw, e) \]
其中 $e$ 如上所定义。

\symlabel{defn:recursor-wtype}
对于函数 $\rec{\wtype{x:A} B(x)}(E,e) : \prd{w : \wtype{x:A} B(x)} E(w)$，其对应的计算规则如下。
\index{computation rule!for W-types@for $\w$-types}%

\begin{itemize}
\item
  对于任意 $a : A$ 和 $f : B(a) \to \wtype{x:A} B(x)$，有
  \begin{equation*}
    \rec{\wtype{x:A} B(x)}(E,e,\supp(a,f)) \jdeq
    e(a,f,\big(\lamu{b:B(a)} \rec{\wtype{x:A} B(x)}(E,e,f(b))\big)).
  \end{equation*}
\end{itemize}

换言之，函数 $\rec{\wtype{x:A} B(x)}(E,e)$ 满足递归式 $e$。

根据上述计算规则，函数 $\dbl$ 的行为符合预期：
\begin{align*}
\dbl(\zerow) & \jdeq \rec\natw(\natw, e, \supp(\bfalse, \; \lamu{x:\emptyt} \rec\emptyt(\natw,x))) \\
& \jdeq e(\bfalse, \big(\lamu{x:\emptyt} \rec\emptyt(\natw,x)\big),
   \big(\lamu{x:\emptyt} \dbl(\rec\emptyt(\natw,x))\big)) \\
 & \jdeq e_0(\big(\lamu{x:\emptyt} \rec\emptyt(\natw,x)\big), \big(\lamu{x:\emptyt} \dbl(\rec\emptyt(\natw,x))\big))\\
 & \jdeq \zerow \\
 \intertext{and}
\dbl(1^{\mathbf{w}}) & \jdeq \rec\natw(\natw, e, \supp(\btrue, \; \lamu{x:\unit} \zerow)) \\
& \jdeq e(\btrue, \big(\lamu{x:\unit} \zerow\big), \big(\lamu{x:\unit} \dbl(\zerow)\big)) \\
 & \jdeq e_1(\big(\lamu{x:\unit} \zerow\big), \big(\lamu{x:\unit} \dbl(\zerow)\big)) \\
 & \jdeq \sucw(\sucw\big((\lamu{x:\unit} \dbl(\zerow))(\ttt)\big))\\
 & \jdeq \sucw(\sucw(\zerow))
\end{align*}
依此类推。

正如对于自然数一样，我们也能证明 $\w$-类型的唯一性定理：

\begin{thm}\label{thm:w-uniq}
  \index{uniqueness!principle, propositional!for functions on W-types@for functions on $\w$-types}%
设 $g,h : \prd{w:\wtype{x:A}B(x)} E(w)$ 为两个满足递归式
%
\begin{equation*}
  e : \prd{a,f} \Parens{\prd{b : B(a)} E(f(b))} \to  E(\supp(a,f)),
\end{equation*}
%
的函数，即在命题意义上满足该递归式，亦即
%
\begin{gather*}
 \prd{a,f} \id{g(\supp(a,f))} {e(a,f,\lamu{b:B(a)} g(f(b)))}, \\
 \prd{a,f} \id{h(\supp(a,f))}{e(a,f,\lamu{b:B(a)} h(f(b)))}.
\end{gather*}
那么 $g$ 与 $h$ 相等。
\end{thm}

\section{归纳类型是初始代数}
\label{sec:initial-alg}

\indexsee{initial!algebra characterization of inductive types}{homotopy-initial}%

如前所述，归纳类型也具有范畴论意义上的泛性质。
它们是\emph{同伦初始代数（homotopy-initial algebras）}：在由指定的构造子决定的``代数''范畴中，它们是在同伦相干的意义下的初始对象。
举一个简单的例子，考虑自然数。
这里合适的``代数''是指一个类型，它配备了与 $\nat$ 的构造子所赋予的相同结构。

\index{natural numbers!as homotopy-initial algebra}


\begin{defn}\label{defn:nalg}
  一个\define{$\nat$-代数（$\nat$-algebra）}
  \indexdef{N-algebra@$\nat$-algebra}%
  \indexdef{algebra!N-@$\nat$-}%
  是一个类型 $C$ 连同两个元素 $c_0 : C$ 与 $c_s : C \to C$。这类代数的类型为
\begin{equation*}
\nalg \defeq \sm {C : \type} C \times (C \to C).
\end{equation*}
\end{defn}

\begin{defn}\label{defn:nhom}
  在两个 $\nat$-代数 $(C,c_0,c_s)$ 与 $(D,d_0,d_s)$ 之间的一个\define{$\nat$-同态（$\nat$-homomorphism）}
  \indexdef{N-homomorphism@$\nat$-homomorphism}%
  \indexdef{homomorphism!N-@$\nat$-}%
  是一个函数 $h : C \to D$，满足 $h(c_0) = d_0$ 且对所有 $c : C$ 有 $h(c_s(c)) = d_s(h(c))$。这类同态的类型为
\begin{narrowmultline*}
\nhom((C,c_0,c_s),(D,d_0,d_s)) \defeq \narrowbreak
 \dsm {h : C \to D} (\id{h(c_0)}{d_0}) \times \tprd{c:C} (\id{h(c_s(c))}{d_s(h(c))}).
\end{narrowmultline*}
\end{defn}

这样我们就得到了一个由 $\nat$-代数与 $\nat$-同态构成的范畴，而我们的断言是：$\nat$ 是这个范畴的初始对象。
范畴论学者会立刻认出，这就是范畴中\emph{自然数对象（natural numbers object）}的定义。

当然，由于我们的类型表现得像 $\infty$-群胚\index{.infinity-groupoid@$\infty$-groupoid}，我们实际上有一个 $\nat$-代数的 $(\infty,1)$-范畴\index{.infinity1-category@$(\infty,1)$-category}，因此我们应该要求 $\nat$ 在适当的 $(\infty,1)$-范畴意义下是初始的。
幸运的是，我们无需定义 $(\infty,1)$-范畴也能表述这一点。

\begin{defn}
  \index{universal!property!of natural numbers}%
  一个 $\nat$-代数 $I$ 被称为\define{同伦初始的（homotopy-initial）}，
  \indexdef{homotopy-initial!N-algebra@$\nat$-algebra}%
  \indexdef{N-algebra@$\nat$-algebra!homotopy-initial (h-initial)}%
  或简称为\define{h-初始的（h-initial）}，
  \indexsee{h-initial}{homotopy-initial}%
  如果对于任何其他 $\nat$-代数 $C$，从 $I$ 到 $C$ 的 $\nat$-同态类型是可缩的。也就是说，
\begin{equation*}
\ishinitn(I) \defeq \prd{C : \nalg} \iscontr(\nhom(I,C)).
\end{equation*}
\end{defn}

当 h-初始代数存在时，它们是唯一的——不仅是范畴论中常见的``在同构意义下''唯一，更由泛等公理在``相等''的意义下唯一。

\begin{thm}
  任何两个 h-初始的 $\nat$-代数都是相等的。
  因此，h-初始 $\nat$-代数的类型是一个命题。
\end{thm}

\begin{proof}
  设 $I$ 和 $J$ 是 h-初始的 $\nat$-代数。
  那么 $\nhom(I,J)$ 是可缩的，因而其中存在某个 $\nat$-同态 $f:I\to J$；同理我们也有一个 $\nat$-同态 $g:J\to I$。
  现在复合 $g\circ f$ 是一个从 $I$ 到 $I$ 的 $\nat$-同态，$\idfunc[I]$ 也是；而 $\nhom(I,I)$ 是可缩的，所以 $g\circ f = \idfunc[I]$。
  类似地，$f\circ g = \idfunc[J]$。
  因此 $\eqv IJ$，从而 $I=J$。由于``可缩''是一个命题，并且依值函数类型保持命题，所以``是 h-初始的''本身也是一个命题。因此，关于 $I$（或 $J$）为 h-初始的任何两个证明都必然相等，这就完成了证明。
\end{proof}

我们现在有以下定理。

\begin{thm}\label{thm:nat-hinitial}
$\nat$-代数 $(\nat, \emptyt, \suc)$ 是同伦初始的。
\end{thm}

\begin{proof}[证明概要]
  固定任意一个 $\nat$-代数 $(C,c_0,c_s)$。
  $\nat$ 的递归原理给出一个函数 $f:\nat\to C$，定义为
  \begin{align*}
    f(0) &\defeq c_0\\
    f(\suc(n)) &\defeq c_s(f(n)).
  \end{align*}
  这两个等式使得 $f$ 成为一个 $\nat$-同态，我们可以将其取为 $\nhom(\nat,C)$ 的收缩中心。
  唯一性定理（\cref{thm:nat-uniq}）则意味着任何其他 $\nat$-同态都等于 $f$。
\end{proof}

为了将其置于更一般的上下文中，考虑\emph{自函子的代数（algebra for an endofunctor）}这一概念会有所帮助。\index{algebra!for an endofunctor}
注意，将一个类型 $C$ 变成 $\nat$-代数，等同于给出一个函数 $c:C+\unit\to C$；而一个函数 $f:C\to D$ 是一个 $\nat$-同态，当且仅当 $f \circ c \htpy d \circ (f+\unit)$。
用范畴论的语言来说，这意味着 $\nat$-代数就是类型范畴上的自函子 $F(X)\defeq X+1$ 的代数。

\indexsee{functor!polynomial}{endofunctor, polynomial}%
\indexsee{polynomial functor}{endofunctor, polynomial}%
\indexdef{endofunctor!polynomial}%
\index{W-type@$\w$-type!as homotopy-initial algebra}
考虑一个更一般的情形，对于给定的 $A : \type$ 和 $B : A \to \type$，考察与之关联的 $\w$ 类型。
在这种情况下，我们有一个关联的\define{多项式函子}：
\begin{equation}
\label{eq:polyfunc}
P(X) = \sm{x : A} (B(x) \rightarrow X).
\end{equation}
实际上，这个指派只是在同伦意义下具有函子性，但这不影响接下来的讨论。
根据定义，一个\define{$P$-代数（$P$-algebra）}
\indexdef{algebra!for a polynomial functor}%
\indexdef{algebra!W-@$\w$-}%
就是一个类型 $C$ 配上一个函数 $s_C :  PC \rightarrow C$。
由 $\Sigma$ 类型的泛性质，这等价于给出一个函数 $\prd{a:A} (B(a) \to C) \to C$。
我们也称这样的对象为关于 $A$ 和 $B$ 的\define{$\w$-代数（$\w$-algebra）}，
\indexdef{W-algebra@$\w$-algebra}%
并记
\symlabel{walg}
\begin{equation*}
\walg(A,B) \defeq \sm {C : \type} \prd{a:A} (B(a) \to C) \to C.
\end{equation*}

类似地，对于 $P$-代数 $(C,s_C)$ 和 $(D,s_D)$，它们之间的一个\define{同态}
\indexdef{homomorphism!of algebras for a functor}%
$(f, s_f) : (C, s_C) \rightarrow (D, s_D)$ 由一个函数 $f : C \rightarrow D$ 和映射 $PC \rightarrow D$ 之间的一个同伦构成：
\[
s_f :  f \circ s_C \, = s_{D} \circ Pf,
\]
其中 $Pf : PC\rightarrow PD$ 是 $P$ 对 $f: C \rightarrow D$ 施加作用的结果（该作用易于定义）。这样的代数同态可以示意性地表示为：
\[
\xymatrix{
 PC \ar[d]_{s_C} \ar[r]^{Pf}  \ar@{}[dr]|{s_f} &  PD \ar[d]^{s_D}\\
C \ar[r]_{f}   & D }
\]
就元素而言，$f$ 是一个 $P$-同态（或称为\define{$\w$-同态（$\w$-homomorphism）}）当且仅当
\indexdef{W-homomorphism@$\w$-homomorphism}%
\indexdef{homomorphism!W-@$\w$-}%
\[f(s_C(a,h)) = s_D(a,f \circ h).\]
我们有 $\w$-同态的类型：
\symlabel{whom}
\begin{equation*}
  \whom_{A,B}((C, s_C),(D,s_D)) \defeq \sm{f : C \to D} \prd{a:A}{h:B(a)\to C} \id{f(s_C(a,h))}{s_D(a, f\circ h)}
\end{equation*}

\index{universal!property!of $\w$-type}%
最后，一个 $P$-代数 $(C, s_C)$ 称为\define{同伦初始的}
\indexdef{homotopy-initial!algebra for a functor}%
\indexdef{homotopy-initial!W-algebra@$\w$-algebra}%
，如果对于每一个 $P$-代数 $(D, s_D)$，所有代数同态的类型 $(C, s_C) \rightarrow (D, s_D)$ 是可缩的。
也就是说，
\begin{equation*}
\ishinitw(A,B,I) \defeq \prd{C : \walg(A,B)} \iscontr(\whom_{A,B}(I,C)).
\end{equation*}
%
现在，与 \cref{thm:nat-hinitial} 类似的定理是：

\begin{thm}\label{thm:w-hinit}
对于任意类型 $A : \type$ 和类型族 $B : A \to \type$，$\w$-代数 $(\wtype{x:A}B(x), \supp)$ 是同伦初始的。
\end{thm}

\begin{proof}[证明概要]
假设我们有 $A : \type$ 和 $B : A \to \type$，
并考虑关联的多项式函子 $P(X)\defeq\sm{x:A}(B(x)\to X)$。
令 $W \defeq \wtype{x:A}B(x)$。那么利用
\cref{sec:w-types} 中的 $\w$ 引入规则，我们有一个结构映射 $s_W\defeq\supp: PW \rightarrow W$。
我们希望证明代数 $(W, s_W)$ 是
同伦初始的。因此，考虑另一个代数 $(C,s_C)$，并证明从 $(W, s_W)$ 到 $(C, s_C)$ 的 $\w$-同态的类型 $T\defeq \whom_{A,B}((W, s_W),(C,s_C)) $
是可缩的。为此，
观察到 $\w$-消去规则和 $\w$-计算
规则允许我们定义一个 $\w$-同态 $(f, s_f) : (W, s_W) \rightarrow (C, s_C)$，
从而证明 $T$ 是有实例的。此外还需要证明，对于每一个 $\w$-同态 $(g, s_g) : (W, s_W) \rightarrow (C, s_C)$，存在一个恒等证明
\begin{equation}
\label{equ:prequired}
p :  (f,s_f) = (g,s_g).
\end{equation}
这里用到如下一般性事实：形如 $(f,s_f) = (g,s_g) $ 的类型
等价于我们所谓的从 $f$ 到 $g$ 的 \define{代数 $2$-胞（algebra $2$-cells）}
\indexdef{algebra!2-cell}%
的类型，其典范
元素形如 $(e, s_e)$，其中 $e : f=g$，而 $s_e$ 是一个更高阶的恒等证明，连接如下粘合图表所表示的恒等证明：
\[
\xymatrix{
PW \ar@/^1pc/[r]^{Pg}   \ar[d]_{s_W}   \ar@/_1pc/[r]_{Pf} \ar@{}[r]|{Pe}
& PC \ar[d]^{s_C}  \\
W  \ar@/_1pc/[r]_f  \ar@{}[r]^{s_f} & C } \qquad
\xymatrix{
PW \ar@/^1pc/[r]^{Pg}   \ar[d]_{s_W} \ar@{}[r]_(.52){s_g}  & PC \ar[d]^{s_C}  \\
W \ar@/^1pc/[r]^g  \ar@/_1pc/[r]_f  \ar@{}[r]|{e} & C }
\]
鉴于这一事实，要证明存在如~\eqref{equ:prequired} 所示的元素，只需
证明存在一个从 $f$ 到 $g$ 的代数 2-胞 $(e,s_e)$ 即可。
恒等证明 $e : f=g$ 现在可以通过函数外延性和
$\w$-消去规则来构造，以保证所需的恒等
证明 $s_e$ 的存在性。
\end{proof}

%%%%%%%%%%%%%%%%%%%%%%%%%%%%%%%%%%%%%%%%%%%%%%%%%%%%%%%%%%

\section{同伦归纳类型}
\label{sec:htpy-inductive}

在 \cref{sec:w-types} 中，我们展示了如何将自然数编码为 $\w$ 类型，即
\begin{align*}
\natw & \defeq \wtype{b:\bool} \rec\bool(\bbU,\emptyt,\unit,b), \\
\zerow & \defeq \supp(\bfalse, (\lamu{x:\emptyt} \rec\emptyt(\natw,x))), \\
\sucw & \defeq \lamu{n:\natw} \supp(\btrue, (\lamu{x:\unit} n)).
\end{align*}
我们还展示了如何利用递归原理在 $\natw$ 上定义 $\dbl$ 函数。
然而，当涉及到归纳原理时，这种编码方式就不再令人满意了：给定 $E : \natw \to \type$ 以及递归式 $e_z : E(\zerow)$ 和 $e_s : \prd{n : \natw}  E(n) \to E(\sucw(n))$，我们只能构造一个依值函数 $r(E,e_z,e_s) : \prd{n : \natw} E(n)$，它\emph{仅在命题意义上}满足给定的递归式，即只相差一条道路。
这意味着，自然数的计算规则（它给出的是判断相等）无法从 $\w$ 类型的规则中以任何明显的方式推导出来。

\index{type!homotopy-inductive}%
\index{homotopy-inductive type}%
如果我们不考虑传统的归纳类型，而是考虑\emph{同伦归纳类型（homotopy-inductive types）}，那么这个问题就不复存在了。在同伦归纳类型中，所有计算规则都表述为道路层面的相等，也就是说，符号 $\jdeq$ 被替换为 $=$。例如，同伦版本 $\w$ 类型 $\mathsf{W^h}$ 的计算规则变为：
\index{computation rule!propositional}%


\begin{itemize}
\item 对于任意 $a : A$ 与 $f : B(a) \to \wtypeh{x:A} B(x)$，我们有
\begin{equation*}
  \rec{\wtypeh{x:A} B(x)}(E,e,\supp(a,f)) = e\Big(a,f,\big(\lamu{b:B(a)} \rec{\wtypeh{x:A} B(x)}(E,f(b))\big)\Big)
\end{equation*}
\end{itemize}

同伦归纳类型在计算性质上有一个明显的劣势——任何使用归纳原理构造的函数，其行为现在都只能命题地表征。
但仍有诸多其他考量促使我们去研究同伦归纳类型。
例如，尽管我们在 \cref{sec:initial-alg} 中说明了归纳类型是同伦初始代数，但并非每个同伦初始代数都是归纳类型（即满足相应的归纳原理）——而每个同伦初始代数\emph{都}是一个同伦归纳类型。
类似地，当涉及的类型中有一个（或两个）只是同伦归纳类型时，我们可能仍想应用来自 \cref{sec:appetizer-univalence} 的唯一性论证——比方说，用于证明用 $\w$-类型编码的 $\nat$ 等价于通常的 $\nat$。

此外，同伦归纳类型这一概念现在是类型论\emph{内部}的了。
例如，这意味着我们可以形成由所有自然数对象构成的类型，并对其作出断言。
对于 $\w$-类型的情形，我们可以将同伦 $\w$-类型 $\wtype{x:A} B(x)$ 刻画为：配备了上确界函数以及满足恰当（命题的）计算规则的归纳原理的任何类型：
\begin{multline*}
\w_d(A,B) \defeq \sm{W : \type}
                 \sm{\supp : \prd {a} (B(a) \to W) \to W}
                 \prd{E : W \to \type} \\
                 \prd{e : \prd{a,f} (\prd{b : B(a)} E(f(b))) \to E(\supp(a,f))}
                 \sm{\ind{} : \prd{w : W} E(w)}
                 \prd{a,f} \\
                 \ind{}(\supp(a,f)) = e(a,\lamu{b:B(a)} \ind{}(f(b))).
\end{multline*}
在 \cref{cha:hits} 中，我们将看到命题计算规则值得考虑的其他一些原因。

本节中，我们将陈述关于同伦归纳类型的一些基本事实。
我们省略大部分证明，因为它们技术性较强。

\begin{thm}
  对于任意 $A : \type$ 与 $B : A \to \type$，类型 $\w_d(A,B)$ 是一个命题。
\end{thm}

事实上，存在一种使用递归原理外加某些\emph{唯一性}与\emph{相干性}定律的等价刻画。首先给出递归原理：
%


\begin{itemize}
\item 当构造从 $\w$-类型 $\wtypeh{x:A} B(x)$ 到类型 $C$ 的函数时，只要对 $\supp(a,f)$ 给出其值就足够了——这里假设我们已经得到了所有 $b : B(a)$ 对应的 $f(b)$ 的值。
换言之，只需构造一个函数
\begin{equation*}
  c : \prd{a:A} (B(a) \to C) \to C.
\end{equation*}
\end{itemize}


\index{computation rule!propositional}%
对于 $\rec{\wtypeh{x:A} B(x)}(C,c) : (\wtype{x:A} B(x)) \to C$ 的相应计算规则如下：


\begin{itemize}
\item 对于任意 $a : A$ 与 $f : B(a) \to \wtypeh{x:A} B(x)$，我们有一个见证 $\beta(C,c,a,f)$，表明下述等式成立：
\begin{equation*}
  \rec{\wtypeh{x:A} B(x)}(C,c,\supp(a,f)) =
  c(a,\lamu{b:B(a)} \rec{\wtypeh{x:A} B(x)}(C,c,f(b))).
\end{equation*}
\end{itemize}

此外，我们断言下述唯一性原理，它表明任何两个由相同递归式定义的函数都是相等的：
\index{uniqueness!principle, propositional!for homotopy W-types@for homotopy $\w$-types}%


\begin{itemize}
\item 设给定 $C : \type$ 与 $c : \prd{a:A} (B(a) \to C) \to C$。令 $g,h : (\wtypeh{x:A} B(x)) \to C$ 为两个函数，它们满足命题相等意义下的递归式 $c$，即我们有
\begin{align*}
  \beta_g &: \prd{a,f} \id{g(\supp(a,f))}{c(a,\lamu{b: B(a)} g(f(b)))}, \\
  \beta_h &: \prd{a,f} \id{h(\supp(a,f))}{c(a,\lamu{b: B(a)} h(f(b)))}.
\end{align*}
那么 $g$ 与 $h$ 相等，即存在一个 $\alpha(C,c,f,g,\beta_g,\beta_h)$，其类型为 $g = h$。
\end{itemize}

\index{coherence}%
回想一下，当我们拥有的是归纳原理而不仅仅是递归原理时，这个命题唯一性原理（propositional uniqueness principle）是可以导出的（\cref{thm:w-uniq}）。
但如果只有递归原理，唯一性原理就不再可导出了——事实上，这条陈述甚至都不成立（练习）。因此，我们将其假定为一条公理。
我们还假定了以下相干性（coherence）\index{coherence}定律，它告诉我们唯一性的证明在典范元上的行为方式：


\begin{itemize}
\item
对于任意的 $a : A$ 和 $f : B(a) \to C$，下述图表在命题意义下交换：
\[\xymatrix{
  g(\supp(a,f)) \ar_{\alpha(\supp(a,f))}[d] \ar^-{\beta_g}[r] & c(a,\lamu{b:B(a)} g(f(b)))
  \ar^{c(a,\blank)(\funext (\lam{b} \alpha(f(b))))}[d] \\
  h(\supp(a,f)) \ar_-{\beta_h}[r] & c(a,\lamu{b: B(a)} h(f(b))) \\
}\]
其中 $\alpha$ 是道路 $\alpha(C,c,f,g,\beta_g,\beta_h) : g = h$ 的缩写。
\end{itemize}

将这些数据综合起来，便得到了 $\wtype{x:A} B(x)$ 的另一种刻画：它是一个具有上确界函数的类型，满足简单的消去、计算、唯一性与相干性规则：
\begin{multline*}
\w_s(A,B) \defeq \sm{W : \type}
                       \sm{\supp : \prd {a} (B(a) \to W) \to W}
                       \prd{C : \type}
                       \prd{c : \prd{a} (B(a) \to C) \to C}\\
                       \sm{\rec{} : W \to C}
                       \sm{\beta : \prd{a,f} \rec{}(\supp(a,f)) = c(a,\lamu{b: B(a)} \rec{}(f(b)))} \narrowbreak
                       \prd{g : W \to C}
                       \prd{h : W \to C}
                       \prd{\beta_g : \prd{a,f} g(\supp(a,f)) = c(a,\lamu{b: B(a)} g(f(b)))} \\
                       \prd{\beta_h : \prd{a,f} h(\supp(a,f)) = c(a,\lamu{b: B(a)} h(f(b)))}
                       \sm{\alpha : \prd {w : W} g(w) = h(w)}
                       \prd{a,f} \\
                       \alpha(\supp(a,f)) \ct \beta_h = \beta_g \ct c(a,-)(\funext \; \lam{b} \alpha(f(b)))
\end{multline*}

\begin{thm}
对于任意 $A : \type$ 和 $B : A \to \type$，类型 $\w_s (A,B)$ 是一个命题。
\end{thm}

最后，我们还有第三种非常简洁的刻画：将 $\wtype{x:A} B(x)$ 视为一个同伦初始的 $\w$-代数：
\begin{equation*}
\w_h(A,B) \defeq \sm{I : \walg(A,B)} \ishinitw(A,B,I).
\end{equation*}

\begin{thm}
对于任意 $A : \type$ 和 $B : A \to \type$，类型 $\w_h (A,B)$ 是一个命题。
\end{thm}

事实上，$\w$-类型的这三种刻画是等价的：


\begin{lem}\label{lem:homotopy-induction-times-3}
对于任意 $A : \type$ 和 $B : A \to \type$，我们有
\[ \w_d(A,B) \eqvsym \w_s(A,B) \eqvsym \w_h(A,B) \]
\end{lem}

更进一步，我们有如下定理，它改进了 \cref{thm:w-hinit}：

\begin{thm}
满足 $\w$-类型的形成、引入、消去与命题计算规则的类型，正是同伦初始的 $\w$-代数。
\end{thm}

%%%%%

\begin{proof}[证明概要]
%%%%%
回顾 \cref{thm:w-hinit} 的证明可知，要确立 $\wtype{x:A}B(x)$ 的同伦初始性（h-initial），仅仅用到其\emph{命题}计算规则就已足够。
为证其逆，假设与 $A : \type$ 和 $B : A \to \UU$ 关联的多项式函子（polynomial functor）有一个同伦初始代数 $(W,s_W)$；我们来证明 $W$ 满足 $\w$-类型的命题规则。
$\w$-引入规则很简单：对于 $a : A$ 和 $t : B(a) \rightarrow W$，我们定义 $\supp(a,t) : W$ 为将结构映射 $s_W : PW \rightarrow W$ 应用于 $(a,t) : PW$ 的结果。
对于 $\w$-消去规则，我们假设其前提成立，特别是有 $C' : W \to \type$。
利用其他前提，可以证明类型 $C \defeq \sm{ w : W} C'(w)$ 可以配备一个结构映射 $s_C : PC \rightarrow C$。根据 $W$ 的同伦初始性，我们得到一个代数同态 $(f, s_f) : (W, s_W) \rightarrow (C, s_C)$。此外，第一个投影 $\proj1 : C \rightarrow W$ 也能配备同态结构，因此我们得到如下形式的图表：
\[
\xymatrix{
PW \ar[r]^{Pf} \ar[d]_{s_W}  & PC \ar[d]^{s_C}  \ar[r]^{P \proj1}  & PW  \ar[d]^{s_W}  \\
W \ar[r]_f  & C \ar[r]_{\proj1}  & W.}
\]
但恒等函数 $1_W : W \rightarrow W$ 具有典范的代数同态结构，所以根据从 $(W,s_W)$ 到自身的同态类型的可缩性，必须存在一个恒等证明，表明 $(f,s_f)$ 与 $(\proj1, s_{\proj1})$ 的复合等于 $(1_W, s_{1_W})$。这特别意味着存在一个恒等证明 $p :  \proj1 \circ f = 1_W$。

由于 $(\proj2 \circ f) w : C( (\proj1 \circ f) w)$，我们可以定义
\[
\rec{}(w,c) \defeq
p_{\, * \,}( ( \proj2 \circ  f)   w )   : C(w)
\]
其中输运 $p_{\, * \,}$ 是关于族
\[
\lamu{u}C\circ u : (W\to W)\to W\to \UU.
\] 的。
命题 $\w$-计算规则的验证是一系列计算，涉及形如 $p_{\, * \,}$ 的操作的自然性性质。
\end{proof}


%%%%%

\index{natural numbers!encoded as a W-type@encoded as a $\w$-type}%
最后，如我们所愿，可以将同伦自然数编码为同伦 $\w$-类型：

\begin{thm}
带有命题计算规则的自然数规则，可以从带有命题计算规则的 $\w$-类型规则中推导出来。
\end{thm}

%%%%%%%%%%%%%%%%%%%%%%%%%%%%%%%%%%%%%%%%%%%%%%%%%%%%%%%%%%

\section{归纳定义的一般语法}
\label{sec:strictly-positive}

\index{type!inductive|(}%
\indexsee{inductive!type}{type, inductive}%

到目前为止，我们只讨论了特定的归纳类型：$\emptyt$、$\unit$、$\bool$、$\nat$、余积、积、$\Sigma$-类型、$\w$-类型等等。
然而，类型论的一个重要方面是能够定义\emph{新的}归纳类型，而不仅仅局限于某些特定的固定列表。
但要做到这一点，我们需要知道什么样的“归纳定义”是有效或合理的。

要知道并非所有“看起来像归纳定义”的东西都有意义，考虑下面这个类型 $C$ 的“构造子”：

\begin{itemize}
\item $g:(C\to \nat) \to C$.
\end{itemize}

此类类型 $C$ 的递归原理大概会说：给定一个类型 $P$，要构造一个函数 $f:C\to P$，只需考虑输入 $c:C$ 是形如 $g(\alpha)$ 的情况，其中 $\alpha:C\to\nat$。此外，我们期望能够以某种方式使用“递归数据”，即 $f$ 应用于 $\alpha$ 的结果。然而，“将 $f$ 应用于 $\alpha$”究竟是什么意思并不清楚，因为两者都是以 $C$ 为定义域的函数。

我们可以为 $C$ 写下一条“递归原理”，即（无根据地）假设存在某种方法能将 $f$ 应用于 $\alpha$ 并获得一个函数 $P\to\nat$。那么，递归规则的输入会要求一个类型 $P$ 连同函数
\begin{equation}
  h:(C\to\nat) \to (P\to\nat) \to P\label{eq:fake-recursor}
\end{equation}
其中 $h$ 的两个参数是 $\alpha$ 和“将 $f$ 应用于 $\alpha$ 的结果”。但是，所得的函数 $f:C\to P$ 的计算规则会是什么样呢？观察其他计算规则，我们会期望类似于“$f(g(\alpha)) \jdeq h(\alpha,f(\alpha))$”这样的东西，但如前所述，“$f(\alpha)$”没有意义。$C$ 的归纳原理就更成问题了；甚至连如何写下其前提都不清楚。

另一方面，我们可以为 $C$ 写下另一条不同的“递归原理”，即忽略 $\alpha$ 定义域中 $C$ 的“递归”存在，仅将其视为一簇自然数的指标类型。在这种情况下，输入会要求一个类型 $P$ 连同函数
\begin{equation*}
  h:(C\to \nat) \to P,
\end{equation*}
因此，递归原理的类型就是 $\rec{C}:\prd{P:\UU} ((C\to \nat) \to P) \to C\to P$，归纳原理也类似。现在可以写出一条计算规则，即 $\rec{C}(P,h,g(\alpha))\jdeq h(\alpha)$。然而，事实证明，存在一个具有此种递归子和计算规则的类型 $C$ 会导致矛盾。参见 \cref{ex:loop,ex:loop2,ex:inductive-lawvere,ex:ilunit} 了解其证明及其他变体。

上面的例子提示了归纳定义的一条限制：所有构造子的定义域都必须是所定义类型的\emph{协变函子}\index{functor!covariant}\index{covariant functor}，这样才能对它们“应用 $f$”，以得到“递归调用”的结果。
换句话说，如果把要定义的类型的所有出现都替换成一个变量
\index{variable!type}%
$X:\type$，那么每个构造子的定义域
\index{domain!of a constructor}%
必须是一个表达式，并且能够成为 $X$ 的一个协变函子。
我们目前考虑过的所有例子都满足这条限制。
例如，对于构造子 $\inl:A\to A+B$，对应的函子是取值于 $A$ 的常值函子（即 $X\mapsto A$）；
而对于构造子 $\suc:\nat\to\nat$，对应的函子是恒等函子（$X\mapsto X$）。

然而，这条必要条件并不是充分的。
协变性确实能防止归纳类型出现在单层函数类型的左侧，就像前面那个“构造子” $g$ 的参数 $C\to\nat$ 那样——因为那会产生一个逆变函子\index{functor!contravariant}\index{contravariant functor}，而不是协变函子。
但由于两个逆变函子的复合是协变的，因此像 $((X\to \nat)\to \nat)$ 这样的\emph{双重}函数类型又成为协变的。
这使我们能够重现 Cantor 风格的悖论\index{paradox}。

例如，考虑如下的“归纳类型” $D$，它具有如下构造子：

\begin{itemize}
\item $k:((D\to\prop)\to\prop)\to D$。
\end{itemize}

假设这样的类型存在，我们定义函数
\begin{align*}
  r&:D\to (D\to\prop)\to\prop,\\
  f&:(D\to\prop) \to D,\\
  p&:(D\to \prop) \to (D\to\prop)\to \prop,\\
  \intertext{by}
  r(k(\theta)) &\defeq \theta,\\
  f(\delta) &\defeq k(\lam{x} (x=\delta)),\\
  p(\delta) &\defeq \lam{x} \delta(f(x)).
\end{align*}
其中 $r$ 是利用 $D$ 的递归原理定义的，而 $f$ 和 $p$ 是显式定义的。
那么对任意 $\delta:D\to\prop$，有 $r(f(\delta)) = \lam{x}(x=\delta)$。

因此特别地，如果 $f(\delta)=f(\delta')$，那么我们有一条道路 $s:(\lam{x}(x=\delta)) = (\lam{x}(x=\delta'))$。
于是，$\happly(s,\delta) : (\delta=\delta) = (\delta=\delta')$，所以特别地有 $\delta=\delta'$ 成立。
从而，$f$ 是“单射”（尽管\emph{先验地} $D$ 可能不是一个集合）。
这已然颇为可疑——我们有了从 $D$ 的“幂集”\index{power set}到 $D$ 的“单射”——再稍加处理，便可将其化为矛盾。

假设给定 $\theta:(D\to\prop)\to\prop$，并定义 $\delta:D\to\prop$ 为
\begin{equation}
  \delta(d) \defeq \exis{\gamma:D\to\prop} (f(\gamma) = d) \times \theta(\gamma).\label{eq:Pinj}
\end{equation}
我们断言 $p(\delta)=\theta$。
根据函数外延性，只需证明对任意 $\gamma:D\to\prop$ 都有 $p(\delta)(\gamma) =_{\prop} \theta(\gamma)$。
而根据泛等性，为此只需证明两者互相蕴含即可。
现在由 $p$ 的定义，我们有
\begin{align*}
  p(\delta)(\gamma) &\jdeq \delta(f(\gamma))\\
  &\jdeq \exis{\gamma':D\to\prop} (f(\gamma') = f(\gamma)) \times \theta(\gamma').
\end{align*}
若 $\theta(\gamma)$ 成立，则此式显然成立，因为我们可以取 $\gamma'\defeq \gamma$。
另一方面，如果我们有 $\gamma'$ 满足 $f(\gamma') = f(\gamma)$ 且 $\theta(\gamma')$，那么由 $f$ 是单射可得 $\gamma'=\gamma$，因此也有 $\theta(\gamma)$。

这就完成了 $p(\delta)=\theta$ 的证明。
于是，每个元素 $\theta:(D\to\prop)\to\prop$ 都是某个元素 $\delta:D\to\prop$ 在 $p$ 下的像。
然而，如果我们用经典的对角线方法定义 $\theta$：
\[ \theta(\gamma) \defeq \neg p(\gamma)(\gamma) \quad\text{for all $\gamma:D\to\prop$} \]
那么由 $\theta = p(\delta)$ 可以推出 $p(\delta)(\delta) = \neg p(\delta)(\delta)$。
这是一个矛盾：没有一个命题可以等价于它的否定。
（假设 $P\Leftrightarrow \neg P$，如果 $P$，那么 $\neg P$，于是得到 $\emptyt$；因此 $\neg P$，但又推出 $P$，于是得到 $\emptyt$。）

\begin{rmk}
  这里需要处理宇宙大小的问题。
  一般来说，一个归纳类型必须存在于一个已经包含了其定义中所有类型的宇宙中。
  因此，如果在 $D$ 的定义中，模糊的记号 \prop 指的是 $\prop_{\UU}$，那么我们并没有 $D:\UU$，而只有对某个更大的满足 $\UU:\UU'$ 的宇宙 $\UU'$，有 $D:\UU'$。
  \index{mathematics!predicative}%
  \indexsee{impredicativity}{mathematics, predicative}%
  \indexsee{predicative mathematics}{mathematics, predicative}%
  因此，在直谓理论中，~\eqref{eq:Pinj} 的右侧属于 $\prop_{\UU'}$ 而非 $\prop_{\UU}$。
  所以这个矛盾的推导确实需要用到\cref{subsec:prop-subsets}中提到的命题大小重整公理。
  \index{propositional!resizing}%
\end{rmk}

\index{consistency}%
这个反例表明，我们应该禁止一个归纳类型出现在其构造子的定义域中的箭头左侧，即使这种出现嵌套在其他箭头中以至于最终是协变的。
（类似地，我们也禁止它出现在依值函数类型的定义域中。）
这种限制被称为\define{严格正性（strict positivity）}
\indexdef{strict!positivity}%
\indexsee{positivity, strict}{strict positivity}%
（普通的“正性”本质上是协变性），且事实证明它是充分的。

\index{constructor}%
总而言之，类型 $W$ 的一个有效归纳定义由一列\emph{构造子（constructor）}构成。
每个构造子具有一个函数类型：它取若干（可能为零）个输入（这些输入可能相互依赖），并返回 $W$ 的一个元素。
最后，我们允许 $W$ 自身出现在其构造子的输入类型中，但仅限于\emph{严格正的（strictly positive）}出现。
这实质上意味着，构造子的每个参数要么是一个不涉及 $W$ 的类型，要么是某种迭代的函数类型，且其最终的上域为 $W$。
例如，下面就是一个合法的构造子类型：
\begin{equation}
  c:(A\to W) \to (B\to C \to W) \to D \to W \to W.\label{eq:example-constructor}
\end{equation}
所有这些函数类型同样可以是依值函数（$\Pi$-类型）。%
\footnote{用 \cref{sec:initial-alg} 的语言来说，严格正性的条件确保了相关的自函子是多项式函子。\indexfoot{endofunctor!polynomial}\indexfoot{algebra!initial}\indexsee{algebra!initial}{homotopy-initial} 范畴论中众所周知，并非\emph{所有}自函子都具有初始代数；限制于多项式函子确保了相容性。 对这一条件可以有多种放宽方案，但在本书中，我们将限定于此处定义的严格正性。}

注意，我们要求归纳定义由\emph{有限}个构造子给出。
这仅仅是因为我们必须将其写在纸面上。
如果我们想要一个表现得仿佛拥有无限多个构造子的归纳类型，只需用某个无限类型来参数化一个构造子即可。
例如，像 $\nat \to W \to W$ 这样的构造子，可以视为等价于可数多个形如 $W\to W$ 的构造子。
（当然，这种无限性现在是类型理论\emph{内部}的，但对任何基础系统而言都理应如此。）
类似地，如果我们想要一个取“无限多个参数”的构造子，可以允许它取一个以某个无限类型参数化的参数族，例如 $(\nat\to W) \to W$，它取一个 $W$ 的元素的无限序列\index{sequence}。

\index{recursion principle!for an inductive type}%
那么，一旦有了这样一个归纳定义，可以用它来做什么呢？
首先，有一个\define{递归原理（recursion principle）}，它表明：为了定义一个函数 $f:W\to P$，只需考虑输入 $w:W$ 来自某个构造子的情形，并可在该构造子的输入上递归调用 $f$。
对于示例构造子~\eqref{eq:example-constructor}，我们要求 $P$ 配备一个类型为
\begin{narrowmultline}\label{eq:example-rechyp}
  d : (A\to W) \to (A\to P) \to (B\to C\to W) \to
  \narrowbreak
  (B\to C \to P) \to D \to W \to P \to P.
\end{narrowmultline}
的函数。
在这些假设下，递归原理给出 $f:W\to P$，并且它以显然的方式“保持构造子数据”——这就是计算规则，其中我们利用了输入的协变性\index{covariance}。
\index{computation rule!for inductive types}%
例如，在例子~\eqref{eq:example-constructor}中，计算规则说的是，对任意 $\alpha:A\to W$，$\beta:B\to C\to W$，$\delta:D$，以及 $\omega:W$，我们有
\begin{equation}
  f(c(\alpha,\beta,\delta,\omega)) \jdeq d(\alpha,f\circ \alpha,\beta, f\circ \beta, \delta, \omega,f(\omega)).\label{eq:example-comp}
\end{equation}
% As we have before in particular cases, when defining a particular function $f$, we may write these rules with $\defeq$ as a way of specifying the data $d$, and say that $f$ is defined by them.

\index{induction principle!for an inductive type}%
一般归纳类型 $W$ 的\define{归纳原理（induction principle）}只稍微复杂一点。
当然，我们从类型族 $P:W\to\type$ 出发，要求它配备“位于”$W$ 的构造子数据“之上”的构造子数据。
这意味着，上述诸如 $A\to P$ 这样的“递归调用”参数必须被替换为类型如 $\prd{a:A} P(\alpha(a))$ 的依值函数。
在~\eqref{eq:example-constructor} 的完整例子中，归纳原理对应的假设将要求
\begin{multline}\label{eq:example-indhyp}
d : \prd{\alpha:A\to W}\Parens{\prd{a:A} P(\alpha(a))} \to \narrowbreak
\prd{\beta:B\to C\to W} \Parens{\prd{b:B}{c:C} P(\beta(b,c))} \to\\
\prd{\delta:D}
\prd{\omega:W} P(\omega) \to
P(c(\alpha,\beta,\delta,\omega)).
\end{multline}
相应的计算规则形式上与~\eqref{eq:example-comp} 相同。
当然，递归原理是归纳原理当 $P$ 为常值族时的特例。
如前所述，归纳原理也称为\define{消去子（eliminator）}，而递归原理称为\define{非依值消去子（non-dependent eliminator）}。

如 \cref{sec:pattern-matching} 所述，我们也允许隐式地调用归纳与递归原理，为归纳原理的每个前提所对应的表达式写下带 $\defeq$ 的定义方程。
这称为通过（依值）\define{模式匹配（pattern matching）}来给出定义。
\index{pattern matching}%
\index{definition!by pattern matching}%
在我们这个贯穿始终的例子中，这意味着可以通过下式来定义 $f:\prd{w:W} P(w) $：
\[ f(c(\alpha,\beta,\delta,\omega)) \defeq \cdots \]
其中 $\alpha:A\to W$、$\beta:B\to C\to W$、$\delta:D$ 和 $\omega:W$ 是绑定在右侧表达式中的变量。
\index{variable}%
此外，右侧表达式可能涉及对 $f$ 的递归调用，形如 $f(\alpha(a))$、$f(\beta(b,c))$ 和 $f(\omega)$。
当用归纳原理重新表述这个定义时，我们将这类递归调用分别替换为新的变量 $\bar\alpha(a)$、$\bar\beta(b,c)$ 和 $\bar\omega$。
\begin{align*}
  \bar\alpha &: \prd{a:A} P(\alpha(a))\\
  \bar\beta &: \prd{b:B}{c:C} P(\beta(b,c))\\
  \bar\omega &: P(\omega).
\end{align*}
\symlabel{defn:induction-wtype}%
这样，我们可以写：
\[ f \defeq \ind{W}(P,\, \lam{\alpha}{\bar\alpha}{\beta}{\bar\beta}{\delta}{\omega}{\bar\omega} \cdots ) \]
其中 $\ind{W}$ 的第二个参数具有~\eqref{eq:example-indhyp} 的类型。

我们不准备尝试形式化地描述有效归纳定义的语法，及其所产生的归纳与递归原理和模式匹配规则。
这当然可以做到（事实上，如果要实现一个计算机证明助理，这是必须做的），但不能提供更多洞见。
经过一定练习后，便能自动为任何归纳定义推导出归纳与递归原理，乃至不假思索地运用它们。

\section{归纳类型的推广}
\label{sec:generalizations}

\index{type!inductive!generalizations}%
归纳类型的概念在类型论中已被研究多年，并有许多许多推广形式：归纳类型族、相互归纳类型、归纳-归纳（inductive-inductive）类型、归纳-递归（inductive-recursive）类型等等。
本节我们将概述其中的一部分，有几种将在本书后续内容中用到。
（在 \cref{cha:hits} 中，我们将更深入地研究一种非常不同的归纳类型推广，它是 \emph{同伦}类型论所特有的。）

这些推广大多涉及同时归纳定义多个类型。
一个我们已经见过的、非常简单的例子是余积 $A+B$。
如果每次想要考虑两个类型的余积时，都必须为 $\nat+\nat$、$\nat+\bool$、$\bool+\bool$ 等等分别写下独立的归纳定义，那将非常繁琐。
相反，我们做一个定义，其中 $A$ 和 $B$ 是代表类型的变量；
\index{variable!type}%
在类型论中，它们被称为\define{参数（parameters）}。%
\indexdef{parameter!of an inductive definition}
因此，严格地说，这样定义的结果不是一个单一类型，而是一个类型族 $+ : \type\to\type\to\type$，它取两个类型作为输入并产生它们的余积。
类似地，列表\index{type!of lists}类型 $\lst A$ 也是一个以类型 $A$ 为参数的类型族 $\lst{\blank}:\type\to\type$。

在数学中，这类事情显而易见，不值一提，但我们在这里提出来是为了与下一个例子作对比。
注意，每个类型 $A+B$ 都是\emph{各自独立地}归纳定义的，每个类型 $\lst A$ 也是如此。
\index{type!family of!inductive}%
\index{inductive!type family}%
相比之下，我们也可以考虑\emph{共同}归纳定义整个类型族 $B:A\to\type$。
区别在于，现在的构造子可能会改变指标 $a:A$，其结果是，我们不能说单个类型 $B(a)$ 是归纳定义的，只能说整个类型族是归纳定义的。

\index{type!of vectors}%
\index{vector}%
标准的例子是\emph{指定长度的列表}类型，传统上称为\define{向量（vectors）}。
我们固定一个参数类型 $A$，并对 $n:\nat$ 定义一个由以下构造子生成的类型族 $\vect n A$：

\begin{itemize}
\item 一个长度为零的向量 $\nil:\vect 0 A$，
\item 一个函数 $\cons:\prd{n:\nat} A\to \vect n A \to \vect{\suc (n)} A$。
\end{itemize}

与列表不同，向量（其元素取自固定类型 $A$）形成一个以长度为索引的类型族。
$A$ 是一个参数，而我们称 $n:\nat$ 是这个归纳族的\define{指标（index）}。
\indexdef{index of an inductive definition}%
单个类型如 $\vect3A$ 并不是归纳定义的：构造 $\vect3A$ 中元素的构造子从该族中的另一个类型取输入，例如 $\cons:A \to \vect2A \to \vect3A$。

\index{induction principle!for type of vectors}
\index{vector!induction principle for}
具体来说，归纳原理也必须涉及整个类型族；因此，前提与结论需对指标进行适当的量化。
在向量类型的情形下，其归纳原理是说，给定一个类型族 $C:\prd{n:\nat} \vect n A \to \type$，以及

\begin{itemize}
\item 一个元素 $c_\nil : C(0,\nil)$，和
\item 一个函数 \narrowequation{c_\cons : \prd{n:\nat}{a:A}{\ell:\vect n A} C(n,\ell) \to C(\suc(n),\cons(a,\ell))}，
\end{itemize}

则存在一个函数 $f:\prd{n:\nat}{\ell:\vect n A} C(n,\ell)$ 使得
\begin{align*}
  f(0,\nil) &\jdeq c_\nil\\
  f(\suc(n),\cons(a,\ell)) &\jdeq c_\cons(n,a,\ell,f(\ell)).
\end{align*}

\index{predicate!inductive}%
\index{inductive!predicate}%
归纳族的一个用途是归纳地定义\emph{谓词}。
例如，我们可以将谓词 $\mathsf{iseven}:\nat\to\type$ 定义为一个以 $\nat$ 为指标的归纳族，其构造子如下：

\begin{itemize}
\item 一个元素 $\mathsf{even}_0 : \mathsf{iseven}(0)$，
\item 一个函数 $\mathsf{even}_{ss} : \prd{n:\nat} \mathsf{iseven}(n) \to \mathsf{iseven}(\suc(\suc(n)))$。
\end{itemize}

换句话说，我们规定 $0$ 是偶数，并且如果 $n$ 是偶数，那么 $\suc(\suc(n))$ 也是偶数。
这些构造子``显然''没有提供任何办法来构造，比如说，$\mathsf{iseven}(1)$ 的一个元素；又因为 $\mathsf{iseven}$ 理应由这些构造子自由生成，所以这样的元素必然不存在。
（不过，严格证明 $\neg \mathsf{iseven}(1)$ 也并非完全平凡。）
$\mathsf{iseven}$ 的归纳原理是说，要证明关于所有偶数自然数的命题，只需对 $0$ 证明它，并验证在加二后仍然成立即可。

\index{mathematics!formalized}%
归纳定义的谓词在数学形式化与软件验证中被大量使用。
但本书中不会过多使用它们，仅在 \cref{sec:ordinals,sec:compactness-interval} 中有几处例外。

\index{type!mutual inductive}%
\index{mutual inductive type}%
另一个重要的特例是当归纳族的指标类型有限时。
此时，我们可以等价地将其表述为由\emph{相互归纳}定义的有限个类型。
例如，我们可以通过相互归纳来定义偶数自然数类型 $\mathsf{even}$ 和奇数自然数类型 $\mathsf{odd}$，其中 $\mathsf{even}$ 由构造子

\begin{itemize}
\item $0:\mathsf{even}$ 和
\item $\mathsf{esucc} : \mathsf{odd}\to\mathsf{even}$
\end{itemize}

生成，而 $\mathsf{odd}$ 则由一个构造子

\begin{itemize}
\item $\mathsf{osucc} : \mathsf{even}\to \mathsf{odd}$
\end{itemize}

生成。
注意 $\mathsf{even}$ 和 $\mathsf{odd}$ 是简单类型（而非类型族），但它们的构造子可以互相引用。
如果我们将这个定义表述为一个归纳类型族 $\mathsf{paritynat} : \bool \to \type$，其中 $\mathsf{paritynat}(\bfalse)$ 和 $\mathsf{paritynat}(\btrue)$ 分别代表 $\mathsf{even}$ 和 $\mathsf{odd}$，那么它的构造子会是：

\begin{itemize}
\item $0 : \mathsf{paritynat}(\bfalse)$，
\item $\mathsf{esucc} : \mathsf{paritynat}(\btrue) \to \mathsf{paritynat}(\bfalse)$，
\item $\mathsf{osucc} : \mathsf{paritynat}(\bfalse) \to \mathsf{paritynat}(\btrue)$。
\end{itemize}

当将其显式表述为相互归纳定义时，$\mathsf{even}$ 和 $\mathsf{odd}$ 的归纳原理是说，给定 $C:\mathsf{even}\to\type$ 和 $D:\mathsf{odd}\to\type$，以及

\begin{itemize}
\item $c_0 : C(0)$，
\item $c_s : \prd{n:\mathsf{odd}} D(n) \to C(\mathsf{esucc}(n))$，
\item $d_s : \prd{n:\mathsf{even}} C(n) \to D(\mathsf{osucc}(n))$，
\end{itemize}

则存在 $f:\prd{n:\mathsf{even}} C(n)$ 和 $g:\prd{n:\mathsf{odd}}D(n)$ 使得
\begin{align*}
  f(0) &\jdeq c_0\\
  f(\mathsf{esucc}(n)) &\jdeq c_s(g(n))\\
  g(\mathsf{osucc}(n)) &\jdeq d_s(f(n)).
\end{align*}
特别地，正如我们只能对归纳族``整体''进行归纳一样，我们必须同时对 $\mathsf{even}$ 和 $\mathsf{odd}$ 进行归纳。
本书同样不会过多使用相互归纳定义。

\index{type!inductive-inductive}%
\index{inductive-inductive type}%
一种更为激进的进一步推广是，允许定义一个类型族 $B:A\to \type$，其中不仅是类型 $B(a)$，连类型 $A$ 本身也是同一个大的归纳定义的一部分。
换言之，我们不仅像在归纳族中那样，为 $B(a)$ 指定可以接受其他 $B(a')$ 作为输入的构造子；同时也为 $A$ 本身指定构造子，而这些构造子可以接受 $B(a)$ 作为输入。
这可以看作是一种归纳族，其指标与被索引的类型是同时归纳定义的；也可以看作是一种相互归纳定义，其中一个类型可以依赖于另一个。
更复杂的依赖结构也是可能的。
一般而言，这些就称为\define{归纳-归纳定义}。
在本书的大部分内容中，我们不会使用它们，但其高阶变体（参见 \cref{cha:hits}）将在 \cref{cha:real-numbers} 的几个实验性例子中出现。

\index{type!inductive-recursive}%
\index{inductive-recursive type}%
我们要提及的最后一个推广是\define{归纳-递归定义}，在这种定义中，我们同时归纳地定义一个类型以及该类型上的一个\emph{递归}函数。
具体来说，我们先固定一个已知类型 $P$，然后为归纳类型 $A$ 给出构造子，同时利用这些构造子为 $A$ 带来的递归原理定义一个函数 $f:A\to P$——特别之处在于，$A$ 的构造子也被允许引用 $f$ 的值。
我们目前还不知道如何从同伦的视角证立这类定义，因此本书将不会使用它们。

\index{type!inductive|)}%

\section{相等类型与恒等系统}
\label{sec:identity-systems}

\index{type!identity!as inductive}%
现在我们要指出，在同伦类型论中处于核心地位的\emph{相等类型}，同样也可以被认为是归纳定义的。
具体来说，它们是一种带指标的``归纳族''，即 \cref{sec:generalizations} 意义下的归纳族。
事实上，将相等类型描述为归纳族有\emph{两种}方式，由此便产生了 \cref{cha:typetheory} 中所述的那两种归纳原理：道路归纳与基于点的道路归纳。

在两种定义中，类型 $A$ 都是一个参数。
对于第一种定义，我们归纳地定义一个族 $=_A : A\to A\to \type$，它有两个属于 $A$ 的指标，并由以下构造子给出：

\begin{itemize}
\item 对任意 $a:A$，有一个元素 $\refl a : a=_A a$。
\end{itemize}

通过与其它归纳族作类比，我们可以从这个定义中提取出归纳原理。
该原理陈述如下：给定任意 \narrowequation{C:\prd{a,b:A} (a=_A b) \to \type,} 以及 $d:\prd{a:A} C(a,a,\refl{a})$，则存在 \narrowequation{f:\prd{a,b:A}{p:a=_A b} C(a,b,p)} 使得 $f(a,a,\refl a)\jdeq d(a)$。
这正是相等类型的道路归纳原理。

对于第二种定义，我们将一个元素 $a_0:A$ 与 $A:\type$ 一并视为参数，然后归纳地定义一族 $(a_0 =_A \blank):A\to \type$，该族仅有\emph{一个}属于 $A$ 的指标，并由以下构造子给出：

\begin{itemize}
\item 一个元素 $\refl{a_0} : a_0 =_A a_0$。
\end{itemize}

请注意，由于 $a_0:A$ 被固定为参数，构造子 $\refl{a_0}$ 在归纳定义中并非以函数形式出现，而仅仅是一个元素。
该定义对应的归纳原理是说：给定 $C:\prd{b:A} (a_0 =_A b) \to \type$ 以及一个元素 $d:C(a_0,\refl{a_0})$，则存在 $f:\prd{b:A}{p:a_0 =_A b} C(b,p)$ 使得 $f(a_0,\refl{a_0})\jdeq d$。
这正是相等类型的基于点的道路归纳原理。

将相等类型视为归纳类型的观点，历史上曾引发一些困惑，原因在于 \cref{sec:bool-nat} 中提到的一种直觉：归纳类型的所有元素都应当通过反复应用其构造子来得到。
对于像 \bool 和 \nat 这类普通的归纳类型，情况确实如此：我们在 \cref{thm:allbool-trueorfalse} 中看到，\bool 的每个元素要么是 $\bfalse$，要么是 $\btrue$；类似地，我们也能证明 \nat 的每个元素要么是 $0$，要么是某个后继。

然而，对于相等类型，这一点却\emph{不}成立：它只有一个构造子 $\refl{}$，但并非每条道路都等于常值道路。
更准确地说，我们仅使用相等类型的归纳原理（无论哪一种）都无法证明：$a=_A a$ 的每个居民都等于 $\refl a$。
为了真正展示一个反例，我们需要某种额外的原则，例如泛等公理——回想一下，在 \cref{thm:type-is-not-a-set} 中，我们正是利用泛等公理展示了一条特殊的道路 $\bool=_{\type}\bool$，它不等于 $\refl{\bool}$。

\index{free!generation of an inductive type}%
\index{generation!of a type, inductive|(}%
关键在于，正如对同伦初始代数的研究所证实的那样，一个归纳定义应当被视为由其构造子\emph{自由生成}的。
当然，一个自由生成的结构所包含的元素可以多于其生成元：例如，由两个符号 $x$ 和 $y$ 生成的自由群中，不仅包含 $x$ 和 $y$，还包含诸如 $xy$、$yx^{-1}y$ 以及 $x^3y^2x^{-2}yx$ 这样的单词。
一般而言，自由结构的元素不仅通过应用生成元得到，也通过应用环境结构中的运算而得到——比如，如果我们讨论的是自由群，那就是群运算。

就归纳类型而言，我们所讨论的是自由生成的\emph{类型}——那么，一个类型的结构所具有的``运算''指的是什么呢？

如果像传统类型论那样，将类型视作\emph{集合}，那么类型上并没有这样的运算，因此我们自然期望归纳类型中除了其构造子所产生的元素之外，不含其他元素。

在同伦类型论中，我们将类型视作\emph{空间}或 $\infty$-群胚，%
\index{.infinity-groupoid@$\infty$-groupoid}
此时\emph{道路}上便有了许多运算（拼接、取逆等等）——这将在 \cref{cha:hits} 中变得重要——但\emph{对象}（元素）上仍然没有运算。
因此，对我们而言，诸如``$\bool$ 的每个元素要么是 $\bfalse$ 要么是 $\btrue$''、``$\nat$ 的每个元素要么是 $0$ 要么是某个后继''这样的陈述依然成立。

然而，正如我们在 \cref{cha:basics} 中看到的那样，将类型视为 $\infty$-群胚也意味着要将函数视作函子，这包括了类型族 $B:A\to\type$。
于是，被视为归纳类型族的相等类型 $(a_0 =_A \blank)$，实际上是一个\emph{自由生成的函子} $A\to\type$。
具体地说，它是由一个元素 $\refl{a_0}: F(a_0)$ 自由生成的函子 $F:A\to\type$。
而函子确实在对象上有运算，即 $A$ 的态射（道路）的作用。

范畴论中的\emph{米田引理}\index{Yoneda!lemma}告诉我们，对于任意范畴 $A$ 和对象 $a_0$，由 $F(a_0)$ 中的一个元素自由生成的函子，正是可表函子 $\hom_A(a_0,\blank)$。
因此，相等类型 $(a_0 =_A \blank)$ 应当就是这个可表函子，而我们确实也是这样看待它的：$(a_0 =_A b)$ 就是 $A$ 中从 $a_0$ 到 $b$ 的态射（道路）所构成的空间。

\index{generation!of a type, inductive|)}

\mentalpause

将相等类型视为归纳族的理由之一，是为了能够应用\cref{sec:appetizer-univalence,sec:htpy-inductive}中的唯一性原理。
具体地说，我们可以通过在 $A\times A$ 上给出另一个满足相同归纳原理的类型族，从而在等价的意义下刻画类型 $A$ 的相等类型族。
这就引出了下面的定义和定理。

\indexsee{system, identity}{identity system}%
\index{identity!system!at a point|(defstyle}%

\begin{defn}\label{defn:identity-systems}
  设 $A$ 是一个类型，$a_0:A$ 是其元素。
  \begin{itemize}
  \item $(A,a_0)$ 上的一个 \define{带点谓词（pointed predicate）}
    \indexdef{predicate!pointed}%
    \indexdef{pointed!predicate}%
    是一个配备了元素 $r_0:R(a_0)$ 的类型族 $R:A\to\type$。
  \item 对于带点谓词 $(R,r_0)$ 与 $(S,s_0)$，一个映射族 $g:\prd{b:A} R(b) \to S(b)$ 称为 \define{带点的}，如果 $g(a_0, r_0)=s_0$。
    我们有
    \[ \mathsf{ppmap}(R,S) \defeq \sm{g:\prd{b:A} R(b) \to S(b)} (g(a_0, r_0)=s_0).\]
  \item 一个 \define{在 $a_0$ 处的恒等系统（identity system at $a_0$）}
    是一个满足如下条件的带点谓词 $(R,r_0)$：对于任意的类型族 $D:\prd{b:A} R(b) \to \type$ 与 $d:D(a_0,r_0)$，都存在一个函数 $f:\prd{b:A}{r:R(b)} D(b,r)$ 使得 $f(a_0,r_0)=d$。
\end{itemize}
\end{defn}

\begin{thm}\label{thm:identity-systems}
  对于 $(A,a_0)$ 上的带点谓词 $(R,r_0)$，以下条件逻辑等价。
  \begin{enumerate}
  \item $(R,r_0)$ 是一个在 $a_0$ 处的恒等系统。\label{item:identity-systems1}
  \item 对于任意的带点谓词 $(S,s_0)$，类型 $\mathsf{ppmap}(R,S)$ 是可缩的。\label{item:identity-systems2}
  \item 对于任意的 $b:A$，函数 $\transfib{R}{\blank}{r_0} : (a_0 =_A b) \to R(b)$ 是一个等价。\label{item:identity-systems3}
  \item 类型 $\sm{b:A} R(b)$ 是可缩的。\label{item:identity-systems4}
  \end{enumerate}
\end{thm}

注意，\ref{item:identity-systems1}$\Leftrightarrow$\ref{item:identity-systems2}$\Leftrightarrow$\ref{item:identity-systems3} 之间的等价是 \cref{lem:homotopy-induction-times-3} 的一个版本，它针对的是被视为在 $A$ 的一个元素上变化的归纳族的相等类型 $a_0 =_A \blank$。
当然，条件\ref{item:identity-systems2}--\ref{item:identity-systems4}都是命题，因此逻辑等价蕴含真正的等价。
（条件\ref{item:identity-systems1}也是一个命题，但我们在此不证明这一点。）
另外需要注意，与条件\ref{item:identity-systems1}--\ref{item:identity-systems3}不同，陈述\ref{item:identity-systems4}并未提及 $a_0$ 或 $r_0$。

\begin{proof}
  首先，假设~\ref{item:identity-systems1} 并令 $(S,s_0)$ 为一个带点谓词。
  定义 $D(b,r) \defeq S(b)$ 和 $d\defeq s_0: S(a_0) \jdeq D(a_0,r_0)$。
  既然 $R$ 是一个恒等系统，我们有 $f:\prd{b:A} R(b) \to S(b)$ 满足 $f(a_0,r_0) = s_0$；因此 $\mathsf{ppmap}(R,S)$ 是有实例的。
  现在假设 $(f,f_r),(g,g_r) : \mathsf{ppmap}(R,S)$，定义 $D(b,r) \defeq (f(b,r) = g(b,r))$，并令 $d \defeq f_r \ct \opp{g_r} : f(a_0,r_0) = s_0 = g(a_0,r_0)$。
  那么同样由于 $R$ 是恒等系统，我们有 $h:\prd{b:A}{r:R(b)} D(b,r)$ 使得 $h(a_0,r_0) = f_r \ct \opp{g_r}$ 成立。
  根据函数外延性以及 $\Sigma$ 类型和道路类型中道路的刻画，这些数据给出一个等式 $(f,f_r) = (g,g_r)$。
  因此，$\mathsf{ppmap}(R,S)$ 是一个有实例的命题，从而可缩；故~\ref{item:identity-systems2} 成立。

  现在假设~\ref{item:identity-systems2}，并定义 $S(b) \defeq (a_0=b)$ 满足 $s_0 \defeq \refl{a_0}:S(a_0)$。
  那么 $(S,s_0)$ 是一个带点谓词，且 $\lamu{b:B}{p:a_0=b} \transfib{R}{p}{r} : \prd{b:A} S(b) \to R(b)$ 是从 $S$ 到 $R$ 的带点的映射族。
  由假设，$\mathsf{ppmap}(R,S)$ 是可缩的，因而是有实例的，故也存在一个从 $R$ 到 $S$ 的带点的映射族。
  两个方向的复合分别是从 $R$ 到 $R$ 以及从 $S$ 到 $S$ 的带点的映射族，由于 $\mathsf{ppmap}(R,R)$ 和 $\mathsf{ppmap}(S,S)$ 都是可缩的，故它们等于各自的恒等映射族。
  于是~\ref{item:identity-systems3} 成立。

  现在假设~\ref{item:identity-systems3}，条件~\ref{item:identity-systems4} 可由 \cref{thm:contr-paths} 得出，这里用到了 $\Sigma$ 类型保持等价这一事实（即 \cref{thm:total-fiber-equiv} 的``若''方向）。

  最后，假定~\ref{item:identity-systems4}，并令 $D:\prd{b:A} R(b)\to  \type$ 和 $d:D(a_0,r_0)$。
  我们可以将 $D$ 等价地表达为一个族 $D':(\sm{b:A} R(b)) \to \type$。
  现在，因为 $\sm{b:A} R(b)$ 是可缩的，我们有
  \[p:\prd{u:\sm{b:A} R(b)} (a_0,r_0) = u. \]
  此外，由于可缩类型的道路类型仍是可缩的，我们有 $p((a_0,r_0)) = \refl{(a_0,r_0)}$。
  定义 $f(u) \defeq \transfib{D'}{p(u)}{d}$，由此得到 $f:\prd{u:\sm{b:A} R(b)} D'(u)$，或等价地 $f:\prd{b:A}{r:R(b)} D(b,r)$。
  最后，我们有
  \[f(a_0,r_0) \jdeq \transfib{D'}{p((a_0,r_0))}{d} = \transfib{D'}{\refl{(a_0,r_0)}}{d} = d.\]
  因此，~\ref{item:identity-systems1} 成立。
\end{proof}

\index{identity!system!at a point|)}%

对于被视为在 $A$ 的两个元素上变化的族的相等类型 $=_A$，我们可以推导出类似的结果。

\index{identity!system|(defstyle}%

\begin{defn}
  一个在类型 $A$ 上的 \define{恒等系统（identity system）}
  是一个族 $R:A\to A\to \type$ 配备一个函数 $r_0:\prd{a:A} R(a,a)$，使得对于任何类型族 $D:\prd{a,b:A} R(a,b) \to \type$ 和 $d:\prd{a:A} D(a,a,r_0(a))$，都存在一个函数 $f:\prd{a,b:A}{r:R(a,b)} D(a,b,r)$ 满足对所有 $a:A$ 有 $f(a,a,r_0(a))=d(a)$。
\end{defn}

\begin{thm}\label{thm:ML-identity-systems}
  对于配备 $r_0:\prd{a:A} R(a,a)$ 的 $R:A\to A\to\type$，以下条件逻辑等价。
  \begin{enumerate}
  \item $(R,r_0)$ 是 $A$ 上的一个恒等系统。\label{item:MLis1}
  \item 对于所有 $a_0:A$，带点谓词 $(R(a_0),r_0(a_0))$ 是 $a_0$ 处的一个恒等系统。\label{item:MLis2}
  \item 对于任何 $S:A\to A\to\type$ 和 $s_0:\prd{a:A} S(a,a)$，类型
    \[ \sm{g:\prd{a,b:A} R(a,b) \to S(a,b)} \prd{a:A} g(a,a,r_0(a)) = s_0(a) \]
    是可缩的。\label{item:MLis3}
  \item 对于任何 $a,b:A$，映射 $\transfib{R(a)}{\blank}{r_0(a)} : (a =_A b) \to R(a,b)$ 是一个等价。\label{item:MLis4}
  \item 对于任何 $a:A$，类型 $\sm{b:A} R(a,b)$ 是可缩的。\label{item:MLis5}
  \end{enumerate}
\end{thm}

\begin{proof}
  \ref{item:MLis1}$\Leftrightarrow$\ref{item:MLis2} 的证明完全类似于相等类型的道路归纳原理与基于点的道路归纳原理之间的等价性证明；参见 \cref{sec:identity-types}。
  与\ref{item:MLis4}和\ref{item:MLis5}的等价则由 \cref{thm:identity-systems} 得出，而\ref{item:MLis3}是直接的。
\end{proof}

\index{identity!system|)}%

这个刻画之所以有趣，原因之一是它提供了陈述泛等性和函数外延性的另一种方式。
\index{univalence axiom}%
泛等公理针对宇宙 \UU 所说的正是：类型族
\[ (\eqv{\blank}{\blank}) : \UU\to\UU\to\UU \]
配上 $\idfunc : \prd{A:\UU} (\eqv AA)$ 后满足~\cref{thm:ML-identity-systems}\ref{item:MLis4}。
因此，它等价于~\ref{item:MLis1}的对应形式，我们可以将其陈述如下。

\begin{cor}[等价归纳]\label{thm:equiv-induction}
  \index{induction principle!for equivalences}%
  \index{equivalence!induction}%
  给定任何类型族 \narrowequation{D:\prd{A,B:\UU} (\eqv AB) \to \type} 和函数 $d:\prd{A:\UU} D(A,A,\idfunc[A])$，存在 \narrowequation{f : \prd{A,B:\UU}{e:\eqv AB} D(A,B,e)} 使得对所有 $A:\UU$ 有 $f(A,A,\idfunc[A]) = d(A)$。
\end{cor}

换言之，要证明关于所有等价的性质，只需对恒等映射证明它即可。
我们已经在 \cref{lem:qinv-autohtpy} 中使用过这个原理（但未作一般性陈述）。

类似地，函数外延性说的是：对于任何 $B:A\to\type$，类型族
\[ (\blank\htpy\blank) : \Parens{\prd{a:A} B(a)} \to \Parens{\prd{a:A} B(a)} \to \type
\]
配上 $\lamu{f:\prd{a:A} B(a)}{a:A} \refl{f(a)}$ 后满足~\cref{thm:ML-identity-systems}\ref{item:MLis4}。
因此，它也等价于~\ref{item:MLis1}的对应形式。

\begin{cor}[同伦归纳]\label{thm:htpy-induction}
  \index{induction principle!for homotopies}%
  \index{homotopy!induction}%
  给定任何 \narrowequation{D:\prd{f,g:\prd{a:A} B(a)} (f\htpy g) \to \type} 和 $d:\prd{f:\prd{a:A} B(a)} D(f,f,\lam{x}\refl{f(x)})$，存在
  %
  \begin{equation*}
    k:\prd{f,g:\prd{a:A} B(a)}{h:f\htpy g} D(f,g,h)
  \end{equation*}
  %
  使得对所有 $f$ 有 $k(f,f,\lam{x}\refl{f(x)}) = d(f)$。
\end{cor}

\sectionNotes

归纳定义在数学中历史悠久，可以说至少可追溯至 Frege（弗雷格）和 Peano（皮亚诺）的自然数公理。\index{Frege}\index{Peano} %
更为一般的“归纳谓词”也并不罕见，但在集合论基础中，它们通常被显式构造出来——或是作为某个适当的子集类的交，或是沿着序数进行超限迭代——而非被视为一个基本的概念。

在类型论中，归纳定义的特例可追溯至 Martin-L\"of（马丁-洛夫）的原始论文：\cite{martin-lof-hauptsatz} 提出了一种归纳定义的谓词与关系的一般性概念；归纳类型的概念在 Martin-L\"of 最早的类型论论文 \cite{Martin-Lof-1973} 中已经出现（但仅有具体实例，并未作为一般性概念提出）；
随后在 \cite{Martin-Lof-1979} 中，才作为与 $\w$-类型相关的一般性概念出现。\index{Martin-L\"of}%

归纳类型的一般概念由 Constable 与 Mendler 于 1985 年引入 \cite{DBLP:conf/lop/ConstableM85}。内涵类型论中归纳类型的一般图式则由
\cite{PfenningPaulinMohring, CoquandPaulin, Dybjer:1991} 提出。

归纳-递归定义的概念出现于 \cite{Dybjer:2000}。一个重要的类型论概念是树类型（这是 Post 系统概念在类型论中的一般表述），它出现于 \cite{PeterssonSynek}。

自然数作为 $\nat$-代数范畴中始对象的\emph{泛性质}，归功于 Lawvere（劳维尔）\cite{lawvere:adjinfound}\index{Lawvere}。
后来，\cite{mp:wftrees} 将此推广，把 $\w$-类型描述为多项式自函子的初始代数。\index{endofunctor!algebra for}
此类泛性质与相应归纳原理之间的相干同伦等价性（\cref{sec:initial-alg,sec:htpy-inductive}）归功于 \cite{ags:it-hott}。

关于类型论的同伦论语义中归纳类型的具体构造，可参见 \cite{klv:ssetmodel,mvdb:wtypes,ls:hits}。

\sectionExercises

\begin{ex}\label{ex:ind-lst}
  根据 $\lst{A}$ 在 \cref{sec:bool-nat} 中作为归纳类型的定义，推导其归纳原理。\index{type!of lists}
\end{ex}

\begin{ex}\label{ex:same-recurrence-not-defeq}
  构造自然数上的两个函数，使它们在判断相等意义下满足同一递归式\index{recurrence} $(e_z, e_s)$，但彼此并不判断相等。
\end{ex}

\begin{ex}\label{ex:one-function-two-recurrences}
  在同一类型 $E$ 上构造两个不同的递归式 $(e_z,e_s)$，使得同一函数 $f:\nat\to E$ 在判断相等意义下同时满足二者。
\end{ex}

\begin{ex}\label{ex:bool}
  证明：对任意类型族 $E : \bool \to \type$，归纳算子
  \[ \ind{\bool}(E) : \big(E(\bfalse) \times E(\btrue)\big) \to \prd{b : \bool} E(b) \]
  都是一个等价。
\end{ex}

\begin{ex}\label{ex:ind-nat-not-equiv}
  证明：对于 $\nat$ 而言，与 \cref{ex:bool} 类似的命题不成立。
\end{ex}

\begin{ex}\label{ex:no-dep-uniqueness-failure}
  证明：如果我们对 $\w$-类型假设的是简单消去（而非依值消去），则唯一性性质（\cref{thm:w-uniq} 的类比）不再成立。
  亦即，举出一个满足 $\w$-类型递归原理的类型，但其中的函数并不由其递归式\index{recurrence}唯一确定。
\end{ex}

\begin{ex}\label{ex:loop}
  假设在 \cref{sec:strictly-positive} 开头处类型 $C$ 的“归纳定义”中，我们将类型 \nat 替换为 \emptyt。
  类似于 \eqref{eq:fake-recursor}，我们可以为此类型考虑一个递归原理，其假设为
  \[ h:(C\to\emptyt) \to (P\to\emptyt) \to P. \]
  试证：即使没有计算规则，这个递归原理也是不一致的，即它允许我们构造出 \emptyt 中的一个元素。
\end{ex}

\begin{ex}\label{ex:loop2}
  现在考虑一个“归纳类型” $D$，它只有一个构造子 $\mathsf{scott}:(D\to D) \to D$。
  \index{Scott}%
  \cref{sec:strictly-positive} 中为 $C$ 提议的第二个递归子导出 $D$ 的如下递归子：
  \[ \rec{D} : \prd{P:\UU} ((D\to D) \to (D\to P)\to P) \to D \to P \]
  其计算规则为 $\rec{D}(P,h,\mathsf{scott}(\alpha)) \jdeq h(\alpha,(\lam{d} \rec{D}(P,h,\alpha(d))))$。
  试证这同样会导致矛盾。
\end{ex}

\begin{ex}\label{ex:inductive-lawvere}
  设 $A$ 为任意类型，并考虑类型 $L_A$ 的一个一般性“归纳定义”，其构造子为 $\mathsf{lawvere}:(L_A\to A) \to L_A$。
  \cref{sec:strictly-positive} 中为 $C$ 提议的第二个递归子导出 $L_A$ 的如下递归子：
  \[ \rec{L_A} : \prd{P:\UU} ((L_A\to A) \to P) \to L_A\to P \]
  其计算规则为 $\rec{L_A}(P,h,\mathsf{lawvere}(\alpha)) \jdeq h(\alpha)$。
  利用此递归子，证明 $A$ 具有\define{不动点性质（fixed-point property）}，即对任意函数 $f:A\to A$，均存在 $a:A$ 使得 $f(a)=a$。
  \index{Lawvere}%
  \index{fixed-point property}%
  特别地，若 $A$ 是一个不具有不动点性质的类型（例如 $\emptyt$、$\bool$ 或 $\nat$），则 $L_A$ 是不一致的。
\end{ex}

\begin{ex}\label{ex:ilunit}
  承接 \cref{ex:inductive-lawvere}，考虑 $L_{\unit}$；由于 $\unit$ 确实具有不动点性质，$L_{\unit}$ 并非显然不一致。
  试为 $L_{\unit}$ 表述一个归纳原理及其计算规则（类似于其递归子），并利用此归纳原理证明 $L_{\unit}$ 是可缩的。
\end{ex}

\begin{ex}\label{ex:empty-inductive-type}
在 \cref{sec:bool-nat} 中，我们定义了某个类型 $A$ 中元素的有限列表类型 $\lst A$。
现考虑类似的归纳定义，定义类型 $\lost A$，其唯一的构造子为
\[ \cons: A \to \lost A \to \lost A. \]
证明 $\lost A$ 等价于 $\emptyt$。
\end{ex}

\begin{ex}\label{ex:Wprop}
  设 $A$ 是一个命题，且 $B:A\to \UU$。
  \begin{enumerate}
  \item 证明 $\wtype{a:A} B(a)$ 是一个命题。
  \item 证明 $\wtype{a:A} B(a)$ 等价于 $\sm{a:A} \neg B(a)$。
  \item 不使用 $\wtype{a:A} B(a)$，证明 $\sm{a:A} \neg B(a)$ 是 \cref{sec:htpy-inductive} 意义下的同伦 $\w$-类型 $\wtypeh{a:A} B(a)$。
  \end{enumerate}
\end{ex}

\begin{ex}\label{ex:Wbounds}
  设 $A:\UU$ 且 $B:A\to \UU$。
  \begin{enumerate}
  \item 证明 $\Parens{\sm{a:A} \neg B(a)} \to \Parens{\wtype{a:A} B(a)}$。
  \item 证明 $\Parens{\wtype{a:A} B(a)} \to \Parens{\neg \prd{a:A} B(a)}$。
  \end{enumerate}
\end{ex}

\begin{ex}\label{ex:Wdec}
  设 $A:\UU$，并假定 $B:A\to \UU$ 是可判定的，即 $\prd{a:A} (B(a)+\neg B(a))$（见 \cref{defn:decidable-equality}）。
  证明 $\Parens{\wtype{a:A} B(a)} \to \Parens{\sm{a:A} \neg B(a)}$。
\end{ex}

\begin{ex}\label{ex:Wbounds-loose}
  证明下列命题是逻辑等价的。
  \begin{enumerate}
  \item 对任意 $A:\set$ 和 $B:A\to \prop$，有 $\Parens{\wtype{a:A} B(a)} \to \Brck{\sm{a:A} \neg B(a)}$。\label{item:Wbounds-loose-Sigma}
  \item 对任意 $A:\set$ 和 $B:A\to \prop$，有 $\Parens{\neg \prd{a:A} B(a)} \to \Brck{\wtype{a:A} B(a)}$。\label{item:Wbounds-loose-Pi}
  \item 排中律成立（如 \cref{sec:intuitionism} 中所述）。
  \end{enumerate}
  类似地，利用 \cref{thm:not-lem}，证明：若假设 \ref{item:Wbounds-loose-Sigma} 或 \ref{item:Wbounds-loose-Pi} 中的蕴含式对所有 $A:\UU$ 和 $B:A\to \UU$ 成立，将导致不一致。
\end{ex}

\begin{ex}\label{ex:Wimpred}
  对于 $A:\UU$ 和 $B:A\to \UU$，定义
  \[ W'_{A,B} \defeq \prd{R:\UU} \Parens{\prd{a:A} (B(a) \to R) \to R} \to R \]
  $W'_{A,B}$ 称为 $\wtype{a:A} B(a)$ 的\define{非直谓编码（impredicative encoding）}。
  \index{impredicative!encoding of a W-type@encoding of a $\w$-type}%
  \index{W-type@$\w$-type!impredicative encoding of}%
  注意，与 $\wtype{a:A} B(a)$ 不同，它位于比 $A$ 和 $B$ 更高的宇宙中。
  \begin{enumerate}
  \item 证明 $W'_{A,B}$ 逻辑等价于 $\wtype{a:A} B(a)$（逻辑等价的定义见 \cref{sec:pat}）。
  \item 证明 $W'_{A,B}$ 蕴含 $\neg\neg \sm{a:A} \neg B(a)$。
  \item 不使用 $\wtype{a:A} B(a)$，证明 $W'_{A,B}$ 满足与 $\wtype{a:A} B(a)$ 相同的\emph{递归}原理，可用于定义取值于宇宙 $\UU$（$W'_{A,B}$ 自身不属于该宇宙）中类型的函数。
  \item 利用\LEM，给出一个 $A:\UU$ 和 $B:A\to \UU$ 的例子，使得 $W'_{A,B}$ 不与 $\wtype{a:A} B(a)$ 等价。
  \end{enumerate}
\end{ex}

\begin{ex}\label{ex:no-nullary-constructor}
  证明对于任意 $A:\UU$ 和 $B:A\to\UU$，有
  \[ \eqv{\neg\Parens{\wtype{a:A} B(a)}}{\neg\Parens{\sm{a:A} \neg B(a)}}. \]
  换言之，$\wtype{a:A} B(a)$ 为空当且仅当它没有零元构造子。
  （与 \cref{ex:empty-inductive-type} 比较。）
\end{ex}

% Local Variables:
% TeX-master: "hott-online"
% End:

\chapter{高阶归纳类型}
\label{cha:hits}

\index{type!higher inductive|(}%
\indexsee{inductive!type!higher}{type, higher inductive}%
\indexsee{higher inductive type}{type, higher inductive}%

\section{引言}
\label{sec:intro-hits}

\index{generation!of a type, inductive|(}

正如我们在\cref{cha:induction}中讨论的一般归纳类型那样，\emph{高阶归纳类型（higher inductive types）}也是一个用于定义由某些构造子生成新类型的一般框架。
但与普通归纳类型不同，在定义高阶归纳类型时，我们可以有一些``构造子''，它们不仅能生成该类型的\emph{点}，还能生成该类型中的\emph{道路}乃至更高维的道路。
\index{type!circle}%
\indexsee{circle type}{type,circle}%
例如，我们可以考虑由以下构造子生成的高阶归纳类型 $\Sn^1$：

\begin{itemize}
\item 一个点 $\base:\Sn^1$，以及
\item 一条道路 $\lloop : {\id[\Sn^1]\base\base}$。
\end{itemize}

这与例如 $\bool$ 的定义完全类似——$\bool$ 由
\begin{itemize}
\item 一个点 $\bfalse:\bool$ 与
\item 一个点 $\btrue:\bool$
\end{itemize}
生成，或者 $\nat$ 由
\begin{itemize}
\item 一个点 $0:\nat$ 与
\item 一个函数 $\suc:\nat\to\nat$
\end{itemize}
生成。

当我们把类型看作高阶群胚时，这种更一般的``生成''概念便十分自然：
既然高阶群胚是一种兼有点、道路和高维道路的``多排序对象''，我们就应当允许在所有维数上都有``生成元''。

我们将那些普通的构造子（例如 $\base$）称为\define{点构造子（point constructors）}
\indexdef{constructor!point}%
\indexdef{point!constructor}%
或\emph{普通构造子}，而将其他构造子（例如 $\lloop$）称为\define{道路构造子（path constructors）}
\indexdef{constructor!path}%
\indexdef{path!constructor}%
或\emph{高阶构造子}。
每个道路构造子都必须指明该道路的起点和终点，我们分别称之为其\define{源（source）}
\indexdef{source!of a path constructor}%
与\define{靶（target）}；
\indexdef{target!of a path constructor}%
对于 $\lloop$ 而言，源和靶都是 $\base$。

注意，像 $\lloop$ 这样的道路构造子会生成恒等类型中的一个\emph{新}居民，它（至少在\emph{先验}意义上）不等于任何先前已存在的此类居民。
特别地，$\lloop$ 并非\emph{先验地}等于 $\refl{\base}$（不过要证明它们确实不相等需要稍加思考，参见\cref{thm:loop-nontrivial}）。
这也正是 $\Sn^1$ 与普通归纳类型 \unit 的区别所在。

关于这一推广，有几个要点需要说明。

\index{free!generation of an inductive type}%
首先，``生成''一词应当认真对待，就如同群可以由某个集合自由生成一样。
具体来说，由于高阶群胚带有道路和高维道路上的\emph{运算}，当这样一个对象由某些构造子``生成''时，这些运算会产生更多并非直接来自构造子本身的道路。
例如，在高阶归纳类型 $\Sn^1$ 中，构造子 $\lloop$ 并不是从 $\base$ 到 $\base$ 的唯一非平凡道路；
我们还有``$\lloop\ct\lloop$''、``$\lloop\ct\lloop\ct\lloop$''等等，以及 $\opp{\lloop}$ 等，它们全都各不相同。
这可能看起来显而易见，甚至不值一提，但它偏离了``普通''归纳类型的行为：
在普通归纳类型中，除了构造子直接``放入''的内容之外，不应期待看到其他东西。

其次，这种生成实际上是\emph{自由}生成：高阶归纳类型严格来说并不允许我们施加``公理''，例如强行要求``$\lloop\ct\lloop$''等于 $\refl{\base}$。
然而，在 $\infty$-群胚的世界里，%
\index{.infinity-groupoid@$\infty$-groupoid}
``自由生成''与``展示''之间几乎没有差别，
\index{presentation!of an infinity-groupoid@of an $\infty$-groupoid}%
\index{generation!of an infinity-groupoid@of an $\infty$-groupoid}%
因为我们可以通过添加一个新的二维生成元（例如在 $\base=\base$ 中添加一条道路 $\lloop\ct\lloop = \refl{\base}$）来使两条道路\emph{在同伦意义下}相等。
当然，我们随后还需要考虑这个新生成元是否应当满足它自己的``公理''，依此类推；但原则上，任何``展示''都可以通过将公理转化为构造子来变成``自由''生成。
正如我们将要看到的，通过添加``截断构造子''，我们同样可以用高阶归纳类型来表达诸如群展示之类的经典概念。

第三，尽管高阶归纳类型包含了能在该类型中生成\emph{道路}的“构造子”，它仍然是对\emph{单个}类型的归纳定义。
特别地，正如我们将要看到的，具有泛性质（通常由归纳原理表述）的是高阶归纳类型本身，而\emph{不是}它的相等类型。
高阶归纳类型的相等类型保留了所有相等类型通常的归纳原理（即道路归纳），并不会获得任何新的归纳原理。

因此，以具体方式确定高阶归纳类型的相等类型可能并非易事；这与\cref{cha:basics}中的情况不同，在那里我们能够对所有传统类型形成操作下相等类型的行为给出明确描述。
例如，在 $\Sn^1$ 中，是否存在从 $\base$ 到 $\base$ 的道路，它不仅仅是 $\lloop$ 及其逆的复合？
直觉上答案似乎是否定的（事实上确实如此），但证明这一点绝非易事。
实际上，这类问题迅速将我们引向诸如计算球面同伦群这样的难题——这是代数拓扑中长期悬而未决的问题，目前尚未知简单的公式。
同伦类型论为这类问题带来了强大的新视角，但也要求类型论本身变得与这些问题的答案一样复杂。

\index{dimension!of path constructors}%
第四，构造子的“维度”（即它们输出的是点、道路、道路之间的道路等等）与结果类型在哪些维度具有非平凡同伦没有直接联系。
举个简单的例子：如果归纳类型 $B$ 有一个类型为 $A\to B$ 的构造子，那么 $A$ 中的任何道路和高阶道路都会在 $B$ 中产生相应的道路和高阶道路，即使该构造子根本不是“高阶”构造子。
对于高阶构造子同样如此：具有一个类型为 $A\to (\id[B]xy)$ 的构造子，不仅意味着 $A$ 中的点会在 $B$ 中生成从 $x$ 到 $y$ 的道路，还意味着 $A$ 中的道路会生成这些道路之间的道路，依此类推。
我们将看到，这种可能性正是高阶归纳类型如此强大的主要原因。

另一方面，即便构造子的输入中\emph{没有}高阶类型，也有可能生成“意想不到的”高阶道路。
例如，在由以下构造生成的二维球面 $\Sn^2$ 中：
\symlabel{s2a}
\index{type!2-sphere}%

\begin{itemize}
\item 一个点 $\base:\Sn^2$，以及
\item 一条在 ${\base=\base}$ 中的 2-维道路$\surf:\refl{\base} = \refl{\base}$，
\end{itemize}

存在一条从 $\refl{\refl{\base}}$ 到自身的非平凡\emph{3-维道路}。
拓扑学家会认出这条道路是\emph{Hopf 纤维化}的一种体现。
从范畴论的观点来看，这与前文提到的 $\Sn^1$ 不仅包含 $\lloop$ 还包含 $\lloop\ct\lloop$ 等现象属于同类：只不过在\emph{高阶}群胚中，存在能够提升维度的\emph{运算}。
事实上，我们在\cref{sec:equality}中已经见过许多这类运算：结合律和单位律不仅仅是性质，还是运算——它们的输入是 1-道路，输出是 2-道路。

\index{generation!of a type, inductive|)}%

% In US Trade format it wants a page break here but then it stretches the above itemize,
% so we give it some stretchable space to use if it wants to.
\vspace*{0pt plus 20ex}

\section{归纳原理与依值道路}
\label{sec:dependent-paths}

当我们描述诸如圆周这样的高阶归纳类型是由某些构造子生成时，必须通过给出类似于\cref{cha:typetheory}中基本类型构造子的规则来解释这意味着什么。
构造子本身给出了\emph{引入（introduction）}规则，但要解释\emph{消去（elimination）}规则（即归纳原理和递归原理），则需要多费些思量。
在本书中，我们不打算给出一个一般性的形式化表述，来界定什么算作“高阶归纳定义”以及如何从这样的定义中提取消去规则——事实上，这是一个微妙的问题，也是当前研究的课题。
相反，我们将依赖一些一般性的非正式讨论和大量例子。

\index{type!circle}%
\index{recursion principle!for S1@for $\Sn^1$}%
递归原理通常比较容易描述：给定一个类型，只要它配备了与所讨论的高阶归纳类型的构造子所赋予的相同结构，就存在一个将构造子映射到该结构的函数。
例如，就 $\Sn^1$ 而言，其递归原理说的是：给定任意配备了基点 $b:B$ 和道路 $\ell:b=b$ 的类型 $B$，存在一个函数 $f:\Sn^1\to B$，使得 $f(\base)=b$ 且 $\apfunc f (\lloop) = \ell$。

\index{computation rule!for S1@for $\Sn^1$}%
\index{equality!definitional}%
后两个等式就是\emph{计算规则}。
\index{computation rule!for higher inductive types|(}%
\index{computation rule!propositional|(}%
然而，这里存在一个问题：这些计算规则应该是\index{judgmental equality}判断相等性还是命题相等性（道路）？
对于普通的归纳类型，我们曾毫不犹豫地将其设为判断相等性，尽管我们在\cref{cha:induction}中看到，即使将其设为命题相等性，在等价的意义下仍会得到相同的类型。
在普通情形下，可以论证计算规则在引言所描述的直观意义下确实是\emph{定义性}相等。

\index{equality!judgmental}%
但对于高阶归纳类型，这一点就不那么清晰了。%, and it is likewise less clear to what extent these equalities can be made judgmental in the known set-theoretic models.
此外，由于操作 $\apfunc f$ 并非类型论真正的基本组成部分，而是我们利用相等类型的归纳原理\emph{定义}出来的（而且我们本可选择其他等价的方式来定义），在一条\emph{判断}相等性中显式地引用它似乎并不合适。
判断相等性是演绎系统的一部分，而演绎系统不应该依赖于我们可能在该系统\emph{内部}做出的特定定义选择。
同时还需要考虑语义和实现方面的问题；详见注记。

不过，将高阶归纳类型的\emph{点}构造子的计算规则设为判断相等性，似乎是没有问题的。
在上面的例子中，这就意味着我们有 $f(\base)\jdeq b$。
这一选择有利于从计算的角度理解高阶归纳类型。
同时，它也极大地简化了我们的工作，因为否则第二个计算规则 $\apfunc f (\lloop) = \ell$ 甚至无法作为命题相等性而良类型；我们将不得不在等式的一边或另一边复合上 $f(\base)$ 与 $b$ 之间指定的那个相等证明。
（当然，当我们讨论更高维的道路时，这类问题最终确实会出现，但在此我们无需对此过于担忧。
亦可参见\cref{sec:hubs-spokes}。）
因此，我们将点构造子的计算规则取为判断相等性，而将道路以及更高维道路的计算规则取为命题相等性。%
\footnote{特别地，用\cref{sec:types-vs-sets}的语言来说，这意味着我们的高阶归纳类型是一种混合体：既有\emph{规则}（规定如何引入此类类型及其元素、其归纳原理、以及其点构造子的计算规则），也有\emph{公理}（道路构造子的计算规则，它们断言某些相等类型\index{inhabited}有实例，且给出实例的项未另行指定）。
我们可以期望，将来能够出现一种更好的类型理论，在其中，高阶归纳类型——就像泛等性一样——将仅通过规则而不借助公理来呈现。%
\indexfoot{axiom!versus rules}%
\indexfoot{rule!versus axioms}%
}

\begin{rmk}\label{rmk:defid}
回忆一下，对于普通的归纳类型，我们将递归定义的函数的计算规则不仅视为判断相等性，而且视作\emph{定义性}相等，因此可以使用符号 $\defeq$ 来表示它们。
例如，截断的前驱\index{predecessor!function, truncated}函数 $p:\nat\to\nat$ 就是由 $p(0)\defeq 0$ 和 $p(\suc(n))\defeq n$ 定义的。
对于高阶归纳类型，这类记号对点构造子是合理的（例如 $f(\base)\defeq b$），但对于道路构造子则可能产生误导，因为像 $\ap f \lloop = \ell$ 这样的相等性并非判断相等性。
因此，我们混合使用记法，将这种``通过定义的命题相等性''写作 $\ap f \lloop \defid \ell$。
\end{rmk}


\index{computation rule!for higher inductive types|)}%
\index{computation rule!propositional|)}%

\index{type!circle|(}%
\index{induction principle!for S1@for $\Sn^1$}%
那么，归纳原理（依值消去子）又是怎样的呢？
回想一下，对于一个普通的归纳类型 $W$，为了用归纳法证明 $\prd{x:W} P(x)$，我们必须为 $W$ 的每个构造子指定一个 $P$ 上的操作，该操作作用于 $W$ 中该构造子``上方''的``纤维''上。
例如，如果 $W$ 是自然数 $\nat$，那么要归纳地证明 $\prd{x:\nat} P(x)$，我们必须指定


\begin{itemize}
\item 一个位于构造子 $0:\nat$ 上方纤维中的元素 $b:P(0)$，以及
\item 对于每个 $n:\nat$，一个函数 $P(n) \to P(\suc(n))$。
\end{itemize}


第二个条件可以看作一个位于构造子 $\suc:\nat\to\nat$ \emph{上方}的``$P\to P$''函数，这推广了 $b:P(0)$ 位于构造子 $0:\nat$ 上方的方式。

因此，类比之下，为了证明 $\prd{x:\Sn^1} P(x)$，我们应当指定

\begin{itemize}
\item 在构造子 $\base:\Sn^1$ 上方纤维中的一个元素 $b:P(\base)$，以及
\item 一条从 $b$ 到 $b$、``落在构造子 $\lloop:\base=\base$ 上方''的道路。
\end{itemize}

注意，尽管 $\Sn^1$ 还包含 $\lloop$ 以外的其他道路（例如 $\refl{\base}$ 和 $\lloop\ct\lloop$），我们只需指定一条落在构造子\emph{本身}上方的道路。
这表达了 $\Sn^1$ 由其构造子``自由生成''的直观。

然而，问题在于一条道路``落在''另一条道路``上方''究竟是什么意思。
它当然\emph{不}意味着仅仅是一条 $b=b$ 的道路，因为那将是纤维 $P(\base)$ 中的道路（从拓扑上说，是一条落在 $\base$ 处\emph{常值}道路上方的道路）。
不过，实际上我们已经在\cref{cha:basics}中回答了这个问题：在\cref{lem:mapdep}之前的讨论中，我们得出结论，一条从 $u:P(x)$ 到 $v:P(y)$ 且落在 $p:x=y$ 上方的道路，可以表示为纤维 $P(y)$ 中的道路 $\trans p u = v$。
由于本章将大量使用这样的\define{依值道路}
\index{path!dependent}%
，我们为此引入一个专门的记号：
\begin{equation}
  (\dpath P p u v) \defeq (\transfib{P} p u = v).\label{eq:dpath}
\end{equation}

\begin{rmk}
定义依值道路还有其他可能的方式。
例如，除了 $\trans p u = v$，我们也可以考虑$u = \trans{(\opp p)}{v}$。
我们还可以将其视为更一般的``异质相等''的特例，
\index{heterogeneous equality}%
\index{equality!heterogeneous}%
或者直接定义为一个归纳类型族。
所有这些定义都给出等价的类型，因此在这个意义上选择哪一个并不太重要。
然而，选择 $\trans p u = v$ 作为定义，使我们最容易推导出关于依值道路的其他结论，例如 $\apdfunc{f}$ 会给出这样的道路，或者我们可以利用\cref{sec:computational}中的传输引理在特定的类型族中计算它们。
\end{rmk}

有了依值道路的概念，我们现在可以更精确地表述 $\Sn^1$ 的归纳原理：给定 $P:\Sn^1\to\type$ 以及
\begin{itemize}
\item 一个元素 $b:P(\base)$，和
\item 一条道路 $\ell : \dpath P \lloop b b$，
\end{itemize}
则存在一个函数$f:\prd{x:\Sn^1} P(x)$，满足 $f(\base)\jdeq b$ 且 $\apd f \lloop = \ell$。
与非依值的情形一样，我们说通过 $f(\base)\defeq b$ 与 $\apd f \lloop \defid \ell$ 来定义 $f$。

\begin{rmk}\label{rmk:varies-along}
在非正式地描述这一归纳原理的某个应用时，我们将其视为把目标``对所有 $x:\Sn^1$ 有 $P(x)$''拆分为两种情形，并有时分别以``当 $x$ 是 $\base$ 时''和``当 $x$ 沿 $\lloop$ 变化时''这样的说法来引入。
``沿一条道路变化''并没有特定的数学含义：它只是一种方便的说法，用以指示证明中相应部分的开始；参见\cref{thm:S1-autohtpy}中的例子。
\index{vary along a path constructor}%
\end{rmk}

从拓扑上说，$\Sn^1$ 的归纳原理可以如\cref{fig:topS1ind}所示来可视化。给定圆上的一个纤维化（图中是环面），定义该纤维化的一个截面，就等同于给出基点 $\base$ 上方纤维中的一个点 $b$，以及一条落在 $\lloop$ 上方、从 $b$ 到 $b$ 的道路。按照我们对依值道路的定义，这一点的类型论解释如\cref{fig:ttS1ind}所示：落在 $\lloop$ 上方、从 $b$ 到 $b$ 的道路，由基点上方纤维中从 $\trans \lloop b$ 到 $b$ 的道路来表示。

\begin{figure}
  \centering
  \begin{tikzpicture}
    \draw (0,0) ellipse (3 and .5);
    \draw (0,3) ellipse (3.5 and 1.5);
    \begin{scope}[yshift=4]
      \clip (-3,3) -- (-1.8,3) -- (-1.8,3.7) -- (1.8,3.7) -- (1.8,3) -- (3,3) -- (3,0) -- (-3,0) -- cycle;
      \draw[clip] (0,3.5) ellipse (2.25 and 1);
      \draw (0,2.5) ellipse (1.7 and .7);
    \end{scope}
    \node (P) at (4.5,3) {$P$};
    \node (S1) at (4.5,0) {$\Sn^1$};
    \draw[->>,thick] (P) -- (S1);
    \node[fill,circle,inner sep=1pt,label={below right:$\base$}] at (0,-.5) {};
    \node at (-2.6,.6) {$\lloop$};
    \node[fill,circle,\OPTblue,inner sep=1pt] (b) at (0,2.3) {};
    \node[\OPTblue] at (-.2,2.1) {$b$};
      \begin{scope}
        \draw[\OPTblue] (b) to[out=180,in=-150] (-2.7,3.5) to[out=30,in=180] (0,3.35);
        \draw[\OPTblue,dotted] (0,3.35) to[out=0,in=175] (1.4,4.35);
        \draw[\OPTblue] (1.4,4.35) to[out=-5,in=90] (2.5,3) to[out=-90,in=0,looseness=.8] (b);
      \end{scope}
      \node[\OPTblue] at (-2.2, 3.3) {$\ell$};
  \end{tikzpicture}
  \caption{$\Sn^1$ 的拓扑归纳原理}
  \label{fig:topS1ind}
\end{figure}

\begin{figure}
  \centering
  \begin{tikzpicture}
    \draw (0,0) ellipse (3 and .5);
    \draw (0,3) ellipse (3.5 and 1.5);
    \begin{scope}[yshift=4]
      \clip (-3,3) -- (-1.8,3) -- (-1.8,3.7) -- (1.8,3.7) -- (1.8,3) -- (3,3) -- (3,0) -- (-3,0) -- cycle;
      \draw[clip] (0,3.5) ellipse (2.25 and 1);
      \draw (0,2.5) ellipse (1.7 and .7);
    \end{scope}
    \node (P) at (4.5,3) {$P$};
    \node (S1) at (4.5,0) {$\Sn^1$};
    \draw[->>,thick] (P) -- (S1);
    \node[fill,circle,inner sep=1pt,label={below right:$\base$}] at (0,-.5) {};
    \node at (-2.6,.6) {$\lloop$};
    \node[fill,circle,\OPTblue,inner sep=1pt] (b) at (0,2.3) {};
      \node[\OPTblue] at (-.3,2.3) {$b$};
      \node[fill,circle,\OPTpurple,inner sep=1pt] (tb) at (0,1.8) {};
      % \draw[\OPTpurple,dashed] (b) to[out=0,in=0,looseness=5] (0,4) to[out=180,in=180] (tb);
      \draw[\OPTpurple,dashed] (b) arc (-90:90:2.9 and 0.85) arc (90:270:2.8 and 1.1);
      \begin{scope}
        \clip (b) -- ++(.1,0) -- (.1,1.8) -- ++(-.2,0) -- ++(0,-1) -- ++(3,2) -- ++(-3,0) -- (-.1,2.3) -- cycle;
        \draw[\OPTred,dotted,thick] (.2,2.07) ellipse (.2 and .57);
        \begin{scope}
          % \draw[clip] (b) -- ++(.1,0) |- (tb) -- ++(-.2,0) -- ++(0,-1) -| ++(3,3) -| (b);
          \clip (.2,0) rectangle (-2,3);
          \draw[\OPTred,thick] (.2,2.07) ellipse (.2 and .57);
        \end{scope}
      \end{scope}
      \node[\OPTred] at (1,1.2) {$\ell: \trans \lloop b=b$};
  \end{tikzpicture}
  \caption{$\Sn^1$ 的类型论归纳原理}
  \label{fig:ttS1ind}
\end{figure}

当然，我们期望能够从归纳原理推导出递归原理，只需取 $P$ 为一个常值类型族即可。
事实上确实如此，不过，从关于 $\lloop$ 的依值计算规则（涉及 $\apdfunc f$）推导出其非依值的版本（涉及 $\apfunc f$），这一步比预想的要棘手一些。

\begin{lem}\label{thm:S1rec}
  \index{recursion principle!for S1@for $\Sn^1$}%
  \index{computation rule!for S1@for $\Sn^1$}%
  若 $A$ 是一个类型，且给定 $a:A$ 与 $p:\id[A]aa$，则存在函数 $f:\Sn^1\to{}A$ 满足
  \begin{align*}
    f(\base)&\defeq a \\
    \apfunc f(\lloop)&\defid p.
  \end{align*}
\end{lem}

\begin{proof}
  我们希望将 $\Sn^1$ 的归纳原理应用于常值类型族，即 $(\lam{x} A): \Sn^1\to \UU$。
  为此所需的前提是：$(\lam{x} A)(\base) \jdeq A$ 中的一个点（即 $a:A$），以及一条 $\dpath {x \mapsto A}{\lloop} a a$ 中的依值道路（或等价地说，$\transfib{x \mapsto A}{\lloop} a = a$）。
  后一个类型并非 $p$ 所在的类型 $\id[A]aa$，但它与后者等价，因为根据 \cref{thm:trans-trivial}，我们有 $\transconst{A}{\lloop}{a} : \transfib{x \mapsto A}{\lloop} a= a$。
  因此，给定 $a:A$ 与 $p:a=a$，我们可以考虑复合道路
  \[\transconst{A}{\lloop}{a} \ct p:(\dpath {x \mapsto A}\lloop aa).\]
  应用归纳原理，我们得到 $f:\Sn^1\to A$，使得
  \begin{align}
    f(\base) &\jdeq a \qquad\text{and}\label{eq:S1recindbase}\\
    \apdfunc f(\lloop) &= \transconst{A}{\lloop}{a} \ct p.\label{eq:S1recindloop}
  \end{align}
  接下来需要推导等式 $\apfunc f(\lloop)=p$。
  然而，根据 \cref{thm:apd-const}，我们有
  \[\apdfunc f(\lloop) = \transconst{A}{\lloop}{f(\base)} \ct \apfunc f(\lloop).\]
  将此式与 \eqref{eq:S1recindloop} 结合，并消去其中出现的 $\transconstf$（由 \eqref{eq:S1recindbase}，二者相同），我们得到 $\apfunc f(\lloop)=p$。
\end{proof}

% Similarly, in this case we speak of defining $f$ by $f(\base)\defeq a$ and $\ap f \lloop \defid p$.
我们也有相应的唯一性原理。

\begin{lem}\label{thm:uniqueness-for-functions-on-S1}
  \index{uniqueness!principle, propositional!for functions on the circle}%
  若 $A$ 是一个类型，且 $f,g:\Sn^1\to{}A$ 是两个映射，并给定两个等式 $p,q$：
  \begin{align*}
    p:f(\base)&=_Ag(\base),\\
    q:\map{f}\lloop&=^{\lam{x} x=_Ax}_p\map{g}\lloop.
  \end{align*}
  那么对于所有 $x:\Sn^1$，有 $f(x)=g(x)$。
\end{lem}

\begin{proof}
  我们对类型族 $P(x)\defeq(f(x)=g(x))$ 应用 $\Sn^1$ 的归纳原理。
  当 $x$ 为 $\base$ 时，$p$ 恰好就是我们所需的。
  当 $x$ 沿 $\lloop$ 变化时，我们需要
  \(p=^{\lam{x} f(x)=g(x)}_{\lloop} p,\)
  根据 \cref{thm:transport-path,thm:dpath-path}，这可以约化为 $q$。
\end{proof}

\index{universal!property!of S1@of $\Sn^1$}%
这两个引理即可推出圆的预期泛性质：

\begin{lem}\label{thm:S1ump}
  对任何类型 $A$，都有一个自然的等价
  \[ (\Sn^1 \to A) \;\eqvsym\;
  \sm{x:A} (x=x).
  \]
\end{lem}

\begin{proof}
  一方面，我们有一个典范函数 $f:(\Sn^1 \to A) \to \sm{x:A} (x=x)$，由 $f(g) \defeq (g(\base),\ap g \lloop)$ 定义。
  另一方面，我们有一个函数 $g:\sm{x:A} (x=x) \to (\Sn^1 \to A)$，它将一对 $(b,\ell)$ 映为由圆的递归原理给出的函数 $\Sn^1 \to A$。

  现在，根据递归原理的计算规则，有 $f \circ g \htpy \idfunc$。
  至于 $g \circ f \htpy \idfunc$，则由唯一性原理得到，因为 \((g \circ f)(\lloop) =^{\lam{x} x=_Ax}_{\refl{\base}} \lloop\)，而这一等式同样是由圆的递归原理的计算规则得到的。
  因此，$f$ 有一个拟逆，从而是一个等价。
\end{proof}

\index{type!circle|)}%

如同在 \cref{sec:htpy-inductive} 中那样，我们可以证明 \cref{thm:S1ump} 的结论等价于拥有带命题计算规则的归纳原理。
其他高阶归纳类型也满足类似于 \cref{thm:S1rec,thm:S1ump} 的引理；我们通常将其证明留给读者。
现在，我们来看许多例子。

\section{区间}
\label{sec:interval}

\index{type!interval|(defstyle}%
\indexsee{interval!type}{type, interval}%
\define{区间（interval）}，记作 $\interval$，或许是比圆还要简单的一种高阶归纳类型。
它由下述构造子生成：

\begin{itemize}
\item 一个点 $\izero:\interval$，
\item 一个点 $\ione:\interval$，以及
\item 一条道路 $\seg : \id[\interval]\izero\ione$。
\end{itemize}


\index{recursion principle!for interval type}%
区间类型的递归原理断言：给定一个类型 $B$，连同


\begin{itemize}
\item 一个点 $b_0:B$，
\item 一个点 $b_1:B$，以及
\item 一条道路 $s:b_0=b_1$，
\end{itemize}


那么存在函数 $f:\interval\to B$，使得 $f(\izero)\jdeq b_0$，$f(\ione)\jdeq b_1$，并且 $\ap f \seg = s$。
\index{induction principle!for interval type}%
其归纳原理断言：给定 $P:\interval\to\type$，连同


\begin{itemize}
\item 一个点 $b_0:P(\izero)$，
\item 一个点 $b_1:P(\ione)$，以及
\item 一条道路 $s:\dpath{P}{\seg}{b_0}{b_1}$，
\end{itemize}


那么存在函数 $f:\prd{x:\interval} P(x)$，使得 $f(\izero)\jdeq b_0$，$f(\ione)\jdeq b_1$，并且 $\apd f \seg = s$。

如果纯粹在同伦意义下考虑，区间其实并不有趣：

\begin{lem}\label{thm:contr-interval}
  类型 $\interval$ 是可缩的。
\end{lem}

\begin{proof}
  我们证明对于所有 $x:\interval$，有 $x=_{\interval}\ione$。换言之，我们需要一个函数 $f$，其类型为 $\prd{x:\interval}(x=_{\interval}\ione)$。我们按如下方式开始定义 $f$：
  \begin{alignat*}{2}
    f(\izero)&\defeq \seg  &:\izero&=_{\interval}\ione,\\
    f(\ione)&\defeq \refl\ione &:\ione &=_{\interval}\ione.
  \end{alignat*}
  尚需定义 $\apd{f}\seg$，其类型必须为 $\seg =_{\seg}^{\lam{x} x=_{\interval}\ione}\refl \ione$。
  根据定义，该类型就是 $\trans\seg\seg=_{\ione=_{\interval}\ione}\refl\ione$，而它又等价于 $\rev\seg\ct\seg=\refl\ione$。
  但该类型中有一个典范元素，即道路之逆确为逆的证明。
\end{proof}

然而，从类型论的角度看，区间仍然具有一些有趣的特征，就像经典同伦论中的拓扑区间一样。
例如，它使我们能够给出函数外延性的一个简单证明。
（当然，如同在 \cref{sec:univalence-implies-funext} 中一样，在接下来的证明过程中，我们暂时搁置对函数外延性公理的总体假设。）

\begin{lem}\label{thm:interval-funext}
  \index{function extensionality!proof from interval type}%
  如果 $f,g:A\to{}B$ 是两个函数，且对于每个 $x:A$ 都有 $f(x)=g(x)$，那么在类型 $A\to{}B$ 中有 $f=g$。
\end{lem}

\begin{proof}
  记我们已有的证明为 $p:\prd{x:A}(f(x)=g(x))$。对于所有 $x:A$，我们如下定义一个函数
  $\widetilde{p}_x:\interval\to{}B$：
  \begin{align*}
    \widetilde{p}_x(\izero) &\defeq f(x), \\
    \widetilde{p}_x(\ione) &\defeq g(x), \\
    \map{\widetilde{p}_x}\seg &\defid p(x).
  \end{align*}
  现在我们定义 $q:\interval\to(A\to{}B)$ 如下：
  \[q(i)\defeq(\lam{x} \widetilde{p}_x(i))\]
  那么 $q(\izero)$ 是函数 $\lam{x} \widetilde{p}_x(\izero)$，它等于 $f$，因为 $\widetilde{p}_x(\izero)$ 定义为 $f(x)$。
  类似地，我们有 $q(\ione)=g$，从而
  \[\map{q}\seg:f=_{(A\to{}B)}g \qedhere\]
\end{proof}

在 \cref{ex:funext-from-interval} 中，我们请读者补全从 \cref{thm:interval-funext} 出发证明完整函数外延性公理的论证。

\index{type!interval|)}%

\section{圆与球面}
\label{sec:circle}

\index{type!circle|(}%
我们已经讨论过圆 $\Sn^1$，它是一个高阶归纳类型，由以下构造子生成：
\begin{itemize}
\item 一个点 $\base:\Sn^1$，以及
\item 一条道路 $\lloop : {\id[\Sn^1]\base\base}$。
\end{itemize}

\index{induction principle!for S1@for $\Sn^1$}%
其归纳原理断言：给定 $P:\Sn^1\to\type$，连同 $b:P(\base)$ 和 $\ell :\dpath P \lloop b b$，我们有 $f:\prd{x:\Sn^1} P(x)$ 满足 $f(\base)\jdeq b$ 且 $\apd f \lloop = \ell$。
其非依值的递归原理断言：给定 $B$ 及 $b:B$ 和 $\ell:b=b$，我们有 $f:\Sn^1\to B$ 满足 $f(\base)\jdeq b$ 且 $\ap f \lloop = \ell$。

我们注意到圆是非平凡的。

\begin{lem}\label{thm:loop-nontrivial}
  $\lloop\neq\refl{\base}$。
\end{lem}

\begin{proof}
  假设 $\lloop=\refl{\base}$。
  那么由于对于任意类型 $A$ 及其中的点 $x:A$ 和道路 $p:x=x$，都存在函数 $f:\Sn^1\to A$ 定义为 $f(\base)\defeq x$ 且 $\ap f \lloop \defid p$，我们有
  \[p = f(\lloop) = f(\refl{\base}) = \refl{x}.\]
  但这意味着每个类型都是集合，而我们已经知道并非如此（参见 \cref{thm:type-is-not-a-set}）。
\end{proof}

圆还具有下面这个有趣的性质，可用作反例的来源。

\begin{lem}\label{thm:S1-autohtpy}
  存在 $\prd{x:\Sn^1} (x=x)$ 中的一个元素，它不等于 $x\mapsto \refl{x}$。
\end{lem}

\begin{proof}
  我们通过 $\Sn^1$-归纳来定义 $f:\prd{x:\Sn^1} (x=x)$。
  当 $x$ 为 $\base$ 时，令 $f(\base)\defeq \lloop$。
  现在当 $x$ 沿着 $\lloop$ 变化时（参见 \cref{rmk:varies-along}），我们需要证明 $\transfib{x\mapsto x=x}{\lloop}{\lloop} = \lloop$。
  然而，在 \cref{sec:compute-paths} 中我们观察到 $\transfib{x\mapsto x=x}{p}{q} = \opp{p} \ct q \ct p$，所以我们必须证明 $\opp{\lloop} \ct \lloop \ct \lloop = \lloop$。
  但消去一个逆后，这是显然的。

  为证明 $f\neq (x\mapsto \refl{x})$，只需证明 $f(\base) \neq \refl{\base}$。
  但 $f(\base)=\lloop$，这正是上一个引理的结论。
\end{proof}

例如，这使我们能够推广 \cref{thm:type-is-not-a-set}，证明任何包含圆的宇宙都不可能是 1-类型。

\begin{cor}
  如果类型 $\Sn^1$ 属于某个宇宙 \type，那么 \type 不是 1-类型。
\end{cor}

\begin{proof}
  根据泛等性，\type 中的类型 $\Sn^1=\Sn^1$ 等价于 $\Sn^1$ 的自等价类型 $\eqv{\Sn^1}{\Sn^1}$，因此只需证明 $\eqv{\Sn^1}{\Sn^1}$ 不是集合即可。
  \index{automorphism!of S1@of $\Sn^1$}%
  为此，只需证明其相等类型 $\id[(\eqv{\Sn^1}{\Sn^1})]{\idfunc[\Sn^1]}{\idfunc[\Sn^1]}$ 不是一个命题。
  由于``是等价''这一性质本身是一个命题，该类型等价于 $\id[(\Sn^1\to\Sn^1)]{\idfunc[\Sn^1]}{\idfunc[\Sn^1]}$。
  而根据函数外延性，后者又等价于 $\prd{x:\Sn^1} (x=x)$；正如我们在 \cref{thm:S1-autohtpy} 中所见，它包含两个不相等的元素。
\end{proof}

\index{type!circle|)}%

\index{type!2-sphere|(}%
\indexsee{sphere type}{type, sphere}%
我们之前也提到过，2-球面 $\Sn^2$ 应当是由以下构造子生成的高阶归纳类型：
\symlabel{s2b}


\begin{itemize}
\item 一个点 $\base:\Sn^2$，以及
\item 一条在 ${\base=\base}$ 中的 2-维道路 $\surf:\refl{\base} = \refl{\base}$。
\end{itemize}


\index{recursion principle!for S2@for $\Sn^2$}%
$\Sn^2$ 的递归原理并不难陈述：给定类型 $B$ 及其中的点 $b:B$ 和一条 2-维道路 $s:\refl b = \refl b$，则存在函数 $f:\Sn^2\to B$，满足 $f(\base)\jdeq b$ 且 $\aptwo f \surf = s$。
此处，`$\aptwo f \surf$` 指的是将 $f$ 的函子作用推广到 2-维道路上的结果；其精确定义如下。

\begin{lem}\label{thm:ap2}
  给定 $f:A\to B$ 以及 $x,y:A$、$p,q:x=y$ 和 $r:p=q$，我们有一条道路 $\aptwo f r : \ap f p = \ap f q$。
\end{lem}

\begin{proof}
  通过道路归纳，可以假设 $p\jdeq q$ 且 $r$ 为自反性。此时可定义 $\aptwo f {\refl p} \defeq \refl{\ap f p}$。
\end{proof}

为了陈述一般的归纳原理，我们需要这个引理在依值函数情形下的版本，而这又要求定义依值 2-维道路的概念。
和之前一样，定义方式有很多种；其中一种是通过传输的 2-维版本来实现。

\begin{lem}\label{thm:transport2}
  给定 $P:A\to\type$ 以及 $x,y:A$、$p,q:x=y$ 和 $r:p=q$，对任意 $u:P(x)$，有 $\transtwo r u : \trans p u = \trans q u$。
\end{lem}

\begin{proof}
  使用道路归纳。
\end{proof}

现在，假设给定 $x,y:A$、$p,q:x=y$、$r:p=q$，以及点 $u:P(x)$、$v:P(y)$ 和依值道路 $h:\dpath P p u v$、$k:\dpath P q u v$。
根据依值道路的定义，这意味着 $h:\trans p u = v$ 且 $k:\trans q u = v$。
因此，将 $r$ 上方的依值 2-道路类型定义为如下形式是合理的：
\[ (\dpath P r h k )\defeq (h = \transtwo r u \ct k). \]
现在我们可以陈述 \cref{thm:ap2} 的依值版本了。

\begin{lem}\label{thm:apd2}
  给定 $P:A\to\type$、$x,y:A$、$p,q:x=y$、$r:p=q$ 以及函数 $f:\prd{x:A} P(x)$，有
  $\apdtwo f r : \dpath P r {\apd f p}{\apd f q}$。
\end{lem}

\begin{proof}
  道路归纳。
\end{proof}

\index{induction principle!for S2@for $\Sn^2$}%
现在可以陈述 $\Sn^2$ 的归纳原理：给定 $P:\Sn^2\to\type$、$b:P(\base)$ 以及 $s:\dpath Q \surf {\refl b}{\refl b}$，其中 $Q\defeq\lam{p} \dpath P p b b$，则存在函数 $f:\prd{x:\Sn^2} P(x)$，使得 $f(\base)\jdeq b$ 且 $\apdtwo f \surf = s$。

\index{type!2-sphere|)}%

当然，随着维数增加，这种显式方法会越来越复杂。
因此，若要定义对一切 $n$ 的 $n$-球面，就需要更系统化的思路。
一种方法是直接处理 $n$-维环路\index{loop!n-@$n$-}，而非一般的 $n$-维道路\index{path!n-@$n$-}。

\index{type!pointed}%
回忆 \cref{sec:equality} 中\emph{带点类型} $\type_*$ 以及 $n$-重环路空间\index{loop space!iterated} $\Omega^n : \type_* \to \type_*$ (\cref{def:pointedtype,def:loopspace}) 的定义。现在可以将 $n$-球面 $\Sn^n$ 定义为由以下生成元生成的高阶归纳类型：
\index{type!n-sphere@$n$-sphere}%


\begin{itemize}
\item 一个点 $\base:\Sn^n$，以及
\item 一条 $n$-环路 $\lloop_n : \Omega^n(\Sn^n,\base)$。
\end{itemize}


为了写出这种表示下的归纳原理，需要定义``依值 $n$-环路\indexdef{loop!dependent n-@dependent $n$-}''的概念，以及依值函数在 $n$-环路上的作用。
我们将此留给读者（参见 \cref{ex:nspheres}）；在下一节中，我们将讨论另一种有时更易处理的定义球面的方法。

\section{纬悬}
\label{sec:suspension}

\indexsee{type!suspension of}{suspension}%
\index{suspension|(defstyle}%
一个类型 $A$ 的\define{纬悬}是一种将 $A$ 的点变成道路（从而将 $A$ 中的道路变成 2-道路，依此类推）的泛用方法。
它是一个类型 $\susp A$，由以下生成元定义：\footnote{这里与依值对类型的记号有一个不幸的冲突，当然后者也用 $\Sigma$ 表示。
  不过，通常上下文可以区分。}


\begin{itemize}
\item 一个点 $\north:\susp A$，
\item 一个点 $\south:\susp A$，以及
\item 一个函数 $\merid:A \to (\id[\susp A]\north\south)$。
\end{itemize}

这些名称意在让人联想到某种``地球仪''，它有一个北极、一个南极，以及以类型 $A$ 为指标的、连接两极的经线。
\indexdef{pole}%
\indexdef{meridian}%
事实上，稍后我们会看到，如果 $A=\Sn^1$，那么它的纬悬就等价于通常球面 $\Sn^2$ 的表面。

\index{recursion principle!for suspension}%
$\susp A$ 的递归原理是说，给定一个类型 $B$ 以及以下数据：

\begin{itemize}
\item 点 $n,s:B$，以及
\item 函数 $m:A \to (n=s)$，
\end{itemize}

那么存在一个函数 $f:\susp A \to B$，使得 $f(\north)\jdeq n$ 且 $f(\south)\jdeq s$，并且对于所有 $a:A$，有 $\ap f {\merid(a)} = m(a)$。
\index{induction principle!for suspension}%
类似地，归纳原理是说，给定 $P:\susp A \to \type$ 以及以下数据：

\begin{itemize}
\item 点 $n:P(\north)$，
\item 点 $s:P(\south)$，以及
\item 对每个 $a:A$，一条道路 $m(a):\dpath P{\merid(a)}ns$，
\end{itemize}

那么存在一个函数 $f:\prd{x:\susp A} P(x)$，使得 $f(\north)\jdeq n$ 且 $f(\south)\jdeq s$，并且对每个 $a:A$ 有 $\apd f {\merid(a)} = m(a)$。

关于纬悬，我们的第一个观察是：它给出了定义圆的另一种方式。

\begin{lem}\label{thm:suspbool}
  \index{type!circle}%
  $\eqv{\susp\bool}{\Sn^1}$。
\end{lem}

\begin{proof}
  按递归原理定义 $f:\susp\bool\to\Sn^1$：令 $f(\north)\defeq \base$，$f(\south)\defeq\base$，而令 $\ap f{\merid(\bfalse)}\defid\lloop$，但 $\ap f{\merid(\btrue)} \defid \refl{\base}$。
  按递归原理定义 $g:\Sn^1\to\susp\bool$：令 $g(\base)\defeq \north$，且 $\ap g \lloop \defid \merid(\bfalse) \ct \opp{\merid(\btrue)}$。
  接下来我们证明 $f$ 和 $g$ 互为拟逆。

  首先，我们对所有 $x:\susp \bool$ 归纳证明 $g(f(x))=x$。
  若 $x\jdeq\north$，则 $g(f(\north)) \jdeq g(\base)\jdeq \north$，因此我们有 $\refl{\north} : g(f(\north))=\north$。
  若 $x\jdeq\south$，则 $g(f(\south)) \jdeq g(\base)\jdeq \north$，我们选择等式 $\merid(\btrue) : g(f(\south)) = \south$。
  还需证明：对于任意 $y:\bool$，当 $x$ 沿 $\merid(y)$ 变化时，这些等式保持不变；也就是说，将 $\refl{\north}$ 沿 $\merid(y)$ 输运后得到 $\merid(\btrue)$。
  根据道路空间和拉回纤维化中的输运，这意味着我们需要证明：
  \[ \opp{\ap g {\ap f {\merid(y)}}} \ct \refl{\north} \ct \merid(y) = \merid(\btrue). \]
  当然，我们可以约去 $\refl{\north}$。
  现在对 $\bool$ 做归纳，可以假设 $y\jdeq \bfalse$ 或 $y\jdeq \btrue$。
  若 $y\jdeq \bfalse$，则我们有
  \begin{align*}
    \opp{\ap g {\ap f {\merid(\bfalse)}}} \ct \merid(\bfalse)
    &= \opp{\ap g {\lloop}} \ct \merid(\bfalse)\\
    &= \opp{(\merid(\bfalse) \ct \opp{\merid(\btrue)})} \ct \merid(\bfalse)\\
    &= \merid(\btrue) \ct \opp{\merid(\bfalse)} \ct \merid(\bfalse)\\
    &= \merid(\btrue)
  \end{align*}
  而若 $y\jdeq \btrue$，则我们有
  \begin{align*}
    \opp{\ap g {\ap f {\merid(\btrue)}}} \ct \merid(\btrue)
    &= \opp{\ap g {\refl{\base}}} \ct \merid(\btrue)\\
    &= \opp{\refl{\north}} \ct \merid(\btrue)\\
    &= \merid(\btrue).
  \end{align*}
  于是，对所有 $x:\susp \bool$，我们有 $g(f(x))=x$。

  现在，我们对所有 $x:\Sn^1$ 归纳证明 $f(g(x))=x$。
  若 $x\jdeq \base$，则 $f(g(\base))\jdeq f(\north)\jdeq\base$，因此我们有 $\refl{\base} : f(g(\base))=\base$。
  还需证明：当 $x$ 沿 $\lloop$ 变化时，这个等式保持不变；也就是说，它沿 $\lloop$ 输运后变为自身。
  再次根据道路空间和拉回纤维化中的输运，这意味着要证明：
  \[ \opp{\ap f {\ap g {\lloop}}} \ct \refl{\base} \ct \lloop = \refl{\base}.\]
  然而，我们有
  \begin{align*}
    \ap f {\ap g {\lloop}} &= \ap f {\merid(\bfalse) \ct \opp{\merid(\btrue)}}\\
    &= \ap f {\merid(\bfalse)} \ct \opp{\ap f {\merid(\btrue)}}\\
    &= \lloop \ct \refl{\base}
  \end{align*}
  故容易得证。
\end{proof}

从拓扑上看，两点空间 \bool 也被称为\emph{0维球面} $\Sn^0$。
（例如，它是 $\mathbb{R}^1$ 中到原点距离为 $1$ 的点构成的空间，正如拓扑 $1$-球面是 $\mathbb{R}^2$ 中到原点距离为 $1$ 的点构成的空间。）
因此，\cref{thm:suspbool} 可以颇具暗示地表述为 $\eqv{\susp\Sn^0}{\Sn^1}$。
\index{type!n-sphere@$n$-sphere|defstyle}%
\indexsee{n-sphere@$n$-sphere}{type, $n$-sphere}%
事实上，这种模式可以继续下去：我们可以用归纳法定义所有球面：
\begin{equation}\label{eq:Snsusp}
  \Sn^0 \defeq \bool
  \qquad\text{and}\qquad
  \Sn^{n+1} \defeq \susp \Sn^n.
\end{equation}
我们甚至可以从更低一维开始，定义 $\Sn^{-1}\defeq \emptyt$，并观察到 $\eqv{\susp\emptyt}{\bool}$。

要严格证明这与上一节中 $\Sn^n$ 的定义相一致，需要将后一种定义表述得更明确。
不过，我们可以证明这个递归定义具有我们期望另一种定义也具有的泛性质。
若 $(A,a_0)$ 和 $(B,b_0)$ 是带点类型（基点常被隐去），令 $\Map_*(A,B)$ 表示带基点映射的类型：
\index{based map}
\symlabel{based-maps}
\[ \Map_*(A,B) \defeq \sm{f:A\to B} (f(a_0)=b_0). \]
注意，任何类型 $A$ 都对应一个带点类型 $A_+ \defeq A+\unit$，其基点为 $\inr(\ttt)$；这称为\emph{附加一个不交基点}。
\indexdef{basepoint!adjoining a disjoint}%
\index{disjoint!basepoint}%
\index{adjoining a disjoint basepoint}%

\begin{lem}
  对于一个类型 $A$ 和一个带点类型 $(B,b_0)$，我们有
  \[ \eqv{\Map_*(A_+,B)}{(A\to B)} \]
\end{lem}

注意，等式右边是从 $A$ 到 $B$ 的普通的\emph{无基点}函数的类型。

\begin{proof}
  从左到右：给定 $f:A_+ \to B$ 满足 $p:f(\inr(\ttt)) = b_0$，我们得到 $f\circ \inl : A \to B$。
  从右到左：给定 $g:A\to B$，我们定义 $g':A_+ \to B$，规定 $g'(\inl(a))\defeq g(a)$ 且 $g'(\inr(u)) \defeq b_0$。
  我们留给读者验证这两者是拟互逆的操作。
\end{proof}

特别地，注意 $\eqv{\bool}{\unit_+}$。
因此，对于任意带点类型 $B$，有
\[{\Map_*(\bool,B)} \eqvsym {(\unit \to B)}\eqvsym B.\]
%
现在回忆一下，环路空间\index{loop space}运算 $\Omega$ 作用于带点类型上，其定义为 $\Omega(A,a_0) \defeq (\id[A]{a_0}{a_0},\refl{a_0})$。
我们也可通过 $\susp(A,a_0)\defeq (\susp A,\north)$ 让纬悬 $\susp$ 作用于带点类型。

\begin{lem}\label{lem:susp-loop-adj}
  \index{universal!property!of suspension}%
  对于带点类型 $(A,a_0)$ 和 $(B,b_0)$，有
  \[ \eqv{\Map_*(\susp A, B)}{\Map_*(A,\Omega B)}.\]
\end{lem}


\addtocounter{thm}{1}           % Because we removed a numbered equation in commit 8f54d16


\begin{proof}
我们首先观察到以下等价链：
\begin{align*}
\Map_*(\susp A, B) & \defeq \sm{f:\susp A\to B} (f(\north)=b_0) \\
                   & \eqvsym \sm{f:\sm{b_n : B}{b_s : B} (A \to (b_n = b_s))} (\fst(f)=b_0) \\
                   & \eqvsym \sm{b_n : B}{b_s : B} \big(A \to (b_n = b_s)\big) \times (b_n=b_0) \\
                   & \eqvsym \sm{p : \sm{b_n : B} (b_n=b_0)}{b_s : B} (A \to (\fst(p) = b_s)) \\
                   & \eqvsym \sm{b_s : B} (A \to (b_0 = b_s))
\end{align*}
第一个等价依据的是纬悬的泛性质，即
\[ \Parens{\susp A \to B} \eqvsym \Parens{\sm{b_n : B} \sm{b_s : B} (A \to (b_n = b_s)) } \]
其中由右向左的函数由递归子给出（参见 \cref{ex:susp-lump}）。
第二个和第三个等价依据的是 \cref{ex:sigma-assoc}，外加各分量的重排。
最后，最后一个等价由 \cref{thm:omit-contr} 推出，因为根据 \cref{thm:contr-paths}，$\sm{b_n : B} (b_n=b_0)$ 是可缩的，以 $(b_0, \refl{b_0})$ 为中心。

现在通过以下等价链来完成证明：
\begin{align*}
  \sm{b_s : B} (A \to (b_0 = b_s))
  &\eqvsym \sm{b_s : B}{g:A \to (b_0 = b_s)}{q:b_0 = b_s} (g(a_0) = q)\\
  &\eqvsym \sm{r : \sm{b_s : B}(b_0 = b_s)}{g:A \to (b_0 = \proj1(r))} (g(a_0) = \proj2(r))\\
  &\eqvsym \sm{g:A \to (b_0 = b_0)} (g(a_0) = \refl{b_0})\\
  &\jdeq \Map_*(A,\Omega B).
\end{align*}
与前面类似，第一个和最后一个等价由 \cref{thm:omit-contr,thm:contr-paths} 给出，第二个由 \cref{ex:sigma-assoc} 及分量的重排给出。
\end{proof}

\index{type!n-sphere@$n$-sphere|defstyle}%
特别地，对于如 \eqref{eq:Snsusp} 中那样定义的球面，我们有
\index{universal!property!of Sn@of $\Sn^n$}%
\[ \Map_*(\Sn^n,B) \eqvsym \Map_*(\Sn^{n-1}, \Omega B) \eqvsym \cdots \eqvsym \Map_*(\bool,\Omega^n B) \eqvsym \Omega^n B. \]
因此，这些球面 $\Sn^n$ 具备我们期望它们应具备的泛性质，即对于如 \cref{sec:circle} 中直接用 $n$ 重环路空间\index{loop space!iterated} 定义的球面所期望的那样的泛性质。

\index{suspension|)}%

\section{胞腔复形}
\label{sec:cell-complexes}

\index{cell complex|(defstyle}%
\index{CW complex|(defstyle}%
在经典拓扑中，\emph{胞腔复形}是通过沿边界逐个粘贴圆盘得到的空间。
若 $n$ 维圆盘\index{disc}的边界被限制在维数严格小于 $n$ 的圆盘（即 $(n-1)$-骨架）内，则称之为\emph{CW 复形}。\index{skeleton!of a CW-complex}

任何有限 CW 复形都可以呈现为一个高阶归纳类型：将 $n$ 维圆盘转化为 $n$ 维道路，并将粘贴映射\index{attaching map}的像划分为道路构造子的源\index{source!of a path constructor}与靶\index{target!of a path constructor}（源和靶均写成低维道路的复合）。
我们在 \cref{sec:circle} 中对 $\Sn^1$ 和 $\Sn^2$ 的显式定义就具有这种形式。

\index{torus}%
另一个例子是环面 $T^2$，它由以下元素生成：


\begin{itemize}
\item 一点 $b:T^2$，
\item 一条道路 $p:b=b$，
\item 另一条道路 $q:b=b$，以及
\item 一条 2-道路 $t: p\ct q = q \ct p$。
\end{itemize}


要看出这是一个环面，最简单的方法或许是从一个矩形出发，它有四个角 $a,b,c,d$、四条边 $p,q,r,s$，其内部显然是一条从 $p\ct q$ 到 $r\ct s$ 的 2-道路 $t$：
\begin{equation*}
  \xymatrix{
      a\ar@{=}[r]^p\ar@{=}[d]_r \ar@{}[dr]|{\Downarrow t} &
      b\ar@{=}[d]^q\\
      c\ar@{=}[r]_s &
      d
      }
\end{equation*}
现在，将边 $r$ 与 $q$ 等同，边 $s$ 与 $p$ 等同，从而四个角也相互等同。
从拓扑上看，这一等同操作产生了环面。

\index{induction principle!for torus}%
\index{torus!induction principle for}%
环面的归纳原理是我们迄今为止写出的归纳原理中最棘手的一个。
给定 $P:T^2\to\type$，要给出一个截面 $\prd{x:T^2} P(x)$，我们需要：


\begin{itemize}
\item 一个点 $b':P(b)$，
\item 一条道路 $p' : \dpath P p {b'} {b'}$，
\item 一条道路 $q' : \dpath P q {b'} {b'}$，以及
\item 一条连接“复合” $p'\ct q'$ 与 $q'\ct p'$ 且位于 $t$ 之上的 2-道路 $t'$。
\end{itemize}


为了使最后这一项数据有意义，我们需要为依值道路定义一个复合运算，但这并不难定义。
然后，归纳原理给出一个函数 $f:\prd{x:T^2} P(x)$，使得 $f(b)\jdeq b'$，$\apd f {p} = p'$，$\apd f {q} = q'$，以及类似“$\apdtwo f t = t'$”的等式。
然而，照现在的写法它还不是良类型的，首先因为等式 $\apd f {p} = p'$ 和 $\apd f {q} = q'$ 并非判断相等，其次因为 $\apdfunc f$ 仅在同伦意义下保持道路拼接。
我们把细节留给读者（参见 \cref{ex:torus}）。

当然，环面的另一种定义是 $T^2 \defeq \Sn^1 \times \Sn^1$（在 \cref{ex:torus-s1-times-s1} 中我们请读者验证两者等价）。
\index{Klein bottle}%
\index{projective plane}%
然而，这种基于胞腔复形的定义很容易推广到其他没有此类描述的空间，例如Klein 瓶、射影平面等。
不过，写出归纳原理确实会变得越来越困难，这需要我们定义依值 $n$-道路的概念，以及 $\apdfunc{}$ 作用于 $n$-道路的方式。
幸运的是，一旦我们掌握了球面，就有办法绕过这个问题。

\section{Hubs and spokes}
\label{sec:hubs-spokes}

\indexsee{spoke}{hub and spoke}%
\index{hub and spoke|(defstyle}%

在拓扑学中，通常通过沿$(n-1)$维边界球面粘贴$n$维圆盘来构建CW复形。
\index{attaching map}%
然而，另一种表述方式是粘入$(n-1)$维球面的\emph{锥（cone）}\index{cone!of a sphere}。
也就是说，我们可以将圆盘\index{disc}看作由一个锥点（或称``轮毂''）和经线\index{meridian}%（或称``辐条''）构成，这些经线将该点连续地连接到边界上的每一点，如\cref{fig:hub-and-spokes}所示。

\begin{figure}
  \centering
  \begin{tikzpicture}
    \draw (0,0) circle (2cm);
    \foreach \x in {0,20,...,350}
      \draw[\OPTblue] (0,0) -- (\x:2cm);
    \node[\OPTblue,circle,fill,inner sep=2pt] (hub) at (0,0) {};
  \end{tikzpicture}
  \caption{一个由轮毂和辐条构成的2维圆盘}
  \label{fig:hub-and-spokes}
\end{figure}

我们可以利用这一想法，将包含$n>1$维道路构造子的高阶归纳类型，用仅包含1维道路构造子的高阶归纳类型来表达。
关键在于，$n$维道路可以表示为一族以某个$(n-1)$维对象为参数的1维道路的连续族。
用于此目的的最简单的$(n-1)$维对象是$(n-1)$维球面，尽管在某些情况下使用其他对象可能更合适。
（回想在\cref{sec:suspension}中，我们能够利用悬垂来归纳地定义球面，而悬垂只涉及1维道路构造子。
事实上，悬垂也可以看作这一想法的一个实例，因为它包含了一族以被悬垂的类型为参数的1维道路。）

\index{torus}
例如，上一节中的环面$T^2$也可以定义为由以下元素生成：


\begin{itemize}
\item 一个点 $b:T^2$，
\item 一条道路 $p:b=b$，
\item 另一条道路 $q:b=b$，
\item 一个点 $h:T^2$，以及
\item 对于每个 $x:\Sn^1$，一条道路 $s(x) : f(x)=h$，其中 $f:\Sn^1\to T^2$ 由 $f(\base)\defeq b$ 和 $\ap f \lloop \defid p \ct q \ct \opp p \ct \opp q$ 定义。
\end{itemize}


这一版本的环面的归纳原理表明，给定 $P:T^2\to\type$，要定义一个截面 $\prd{x:T^2} P(x)$，需要：


\begin{itemize}
\item 一个点 $b':P(b)$，
\item 一条道路 $p' : \dpath P p {b'} {b'}$，
\item 一条道路 $q' : \dpath P q {b'} {b'}$，
\item 一个点 $h':P(h)$，以及
\item 对于每个 $x:\Sn^1$，一条道路 $\dpath {P}{s(x)}{g(x)}{h'}$，其中 $g:\prd{x:\Sn^1} P(f(x))$ 由 $g(\base)\defeq b'$ 和 $\apd g \lloop \defid t(p' \ct q' \ct \opp{(p')} \ct \opp{(q')})$ 定义。
  在后一项中，$\ct$ 表示依值道路的拼接，而 $t:\eqv{(\dpath{P}{\ap f \lloop}{b'}{b'})}{(\dpath{P\circ f}{\lloop}{b'}{b'})}$ 的定义留给读者。
\end{itemize}


注意，这里不需要依值2-道路或$\apdtwofunc{}$。计算规则也留待读者写出。

\begin{rmk}\label{rmk:spokes-no-hub}
有人可能会质疑引入轮毂点$h$的必要性；为什么不能像\cref{fig:spokes-no-hub}所示的那样，简单地添加道路将圆盘的边界连续地连接到边界\emph{上}的某个点呢？
然而，若不作进一步修改，这是行不通的。
因为如果给定某个 $f:\Sn^1 \to X$，我们给出一个道路构造子将每个 $f(x)$ 连接到 $f(\base)$，那么最终得到的结构会更像\cref{fig:spokes-no-hub-ii}中的图形：一个锥，其顶点被扭转并粘合到了底面的某一点上。
问题在于，这样指定的从 $f(\base)$ 到自身的道路可能不是自反性。
我们可以添加一个2维道路构造子来确保这一点，从而补救该问题，但使用一个独立的轮毂点则无需任何高于~1维的道路构造子。
\end{rmk}

\begin{figure}
  \centering
  \begin{minipage}{2in}
    \begin{center}
      \begin{tikzpicture}
        \draw (0,0) circle (2cm);
        \clip (0,0) circle (2cm);
        \foreach \x in {0,15,...,165}
        \draw[\OPTblue] (0,-2cm) -- (\x:4cm);
      \end{tikzpicture}
    \end{center}
    \caption{无轮毂辐条}
    \label{fig:spokes-no-hub}
  \end{minipage}
  \qquad
  \begin{minipage}{2in}
    \begin{center}
      \begin{tikzpicture}[xscale=1.3]
        \draw (0,0) arc (-90:90:.7cm and 2cm) ;
        \draw[dashed] (0,4cm) arc (90:270:.7cm and 2cm) ;
        \draw[\OPTblue] (0,0) to[out=90,in=0] (-1,1) to[out=180,in=180] (0,0);
        \draw[\OPTblue] (0,4cm) to[out=180,in=180,looseness=2] (0,0);
        \path (0,0) arc (-90:-60:.7cm and 2cm) node (a) {};
        \draw[\OPTblue] (a.center) to[out=120,in=10] (-1.2,1.2) to[out=190,in=180] (0,0);
        \path (0,0) arc (-90:-30:.7cm and 2cm) node (b) {};
        \draw[\OPTblue] (b.center) to[out=150,in=20] (-1.4,1.4) to[out=200,in=180] (0,0);
        \path (0,0) arc (-90:0:.7cm and 2cm) node (c) {};
        \draw[\OPTblue] (c.center) to[out=180,in=30] (-1.5,1.5) to[out=210,in=180] (0,0);
        \path (0,0) arc (-90:30:.7cm and 2cm) node (d) {};
        \draw[\OPTblue] (d.center) to[out=190,in=50] (-1.7,1.7) to[out=230,in=180] (0,0);
        \path (0,0) arc (-90:60:.7cm and 2cm) node (e) {};
        \draw[\OPTblue] (e.center) to[out=200,in=70] (-2,2) to[out=250,in=180] (0,0);
        \clip (0,0) to[out=90,in=0] (-1,1) to[out=180,in=180] (0,0);
        \draw (0,4cm) arc (90:270:.7cm and 2cm) ;
      \end{tikzpicture}
    \end{center}
    \caption{无轮毂辐条 II}
    \label{fig:spokes-no-hub-ii}
  \end{minipage}
\end{figure}

\begin{rmk}
  \index{computation rule!propositional}%
  还需注意，这种将高阶道路“转化”为 1-道路的做法，并不能保持这些道路的\emph{判断式}计算规则，尽管它确实能保持\emph{命题式}计算规则。
\end{rmk}

\index{cell complex|)}%
\index{CW complex|)}%

\index{hub and spoke|)}%

\section{推出}
\label{sec:colimits}

\index{type!limit}%
\index{type!colimit}%
\index{limit!of types}%
\index{colimit!of types}%
从范畴论的观点来看，任何基础系统的一个重要方面都在于能够构造极限与余极限。
在集合论基础中，这些是集合的极限与余极限；而在我们的理论里，它们则是\emph{类型}的极限与余极限。
我们已经在 \cref{sec:universal-properties} 中看到，Cartesian 积类型具有作为类型范畴积的正确泛性质，并且在 \cref{ex:coprod-ump} 中看到，余积类型同样具有其预期的泛性质。

如 \cref{sec:universal-properties} 所述，更一般的极限可以利用相等类型与 $\Sigma$ 类型来构造，例如 $f:A\to C$ 与 $g:B\to C$ 的\index{pullback}拉回是 $\sm{a:A}{b:B} (f(a)=g(b))$（参见 \cref{ex:pullback}）。
然而，构造更一般的\emph{余极限}则需要将来自不同类型的元素等同起来，而高阶归纳类型恰好适用于此。
由于我们所有的构造都是同伦不变的，所有的余极限必然都是\emph{同伦余极限}；但为简洁起见，我们省略这一随处可见的修饰语。

本节我们讨论\emph{推出}，它大概是最简单的、也是最有用的余极限之一。
事实上，人们预期所有有限余极限（对于一个恰当的同伦意义上的“有限”定义而言）都可以由推出与有限余积构造出来。
使用高阶归纳类型直接构造更一般的余极限也是可能的，但这有些技术性，且并不完全令人满意，因为我们尚无一个良好的、完全一般的同伦相干图表概念。

\indexsee{type!pushout of}{pushout}%
\index{pushout|(defstyle}%
\index{span}%
给定类型与函数的一个伸展：
\[\Ddiag=\;\vcenter{\xymatrix{C \ar^g[r] \ar_f[d] & B \\ A & }}\]
该伸展的\define{推出}是以下数据所呈现的高阶归纳类型 $A\sqcup^CB$：
\begin{itemize}
\item 一个函数 $\inl:A\to A\sqcup^CB$，
\item 一个函数 $\inr:B \to A\sqcup^CB$，以及
\item 对每个 $c:C$，一条道路 $\glue(c):(\inl(f(c))=\inr(g(c)))$。
\end{itemize}

换句话说，$A\sqcup^CB$ 是 $A$ 与 $B$ 的不交并，再加上对每个 $c:C$ 的一个证据，证明 $f(c)$ 与 $g(c)$ 相等。
递归原理表明，若 $D$ 是另一个类型，我们可以通过定义以下数据来给出映射 $s:A\sqcup^CB\to{}D$：
\begin{itemize}
\item 对每个 $a:A$，$s(\inl(a)):D$ 的值，
\item 对每个 $b:B$，$s(\inr(b)):D$ 的值，以及
\item 对每个 $c:C$，$\mapfunc{s}(\glue(c)):s(\inl(f(c)))=s(\inr(g(c)))$ 的值。
\end{itemize}

我们将归纳原理的表述留给读者。
它还蕴含着下述唯一性原理：如果 $s,s':A\sqcup^CB\to{}D$ 是两个映射，使得
\index{uniqueness!principle, propositional!for functions on a pushout}%
\begin{align*}
  s(\inl(a))&=s'(\inl(a))\\
  s(\inr(b))&=s'(\inr(b))\\
  \mapfunc{s}(\glue(c))&=\mapfunc{s'}(\glue(c))
  \qquad\text{(modulo the previous two equalities)}
\end{align*}
对每个 $a,b,c$ 成立，那么 $s=s'$。

为了表述推出的泛性质，我们引入下述定义。

\begin{defn}\label{defn:cocone}
  给定一个伸展 $\Ddiag= (A \xleftarrow{f} C \xrightarrow{g} B)$ 和一个类型 $D$，一个\define{在伸展 $\Ddiag$ 下、以 $D$ 为顶点的上锥（cocone under $\Ddiag$ with vertex $D$）}
  \indexdef{cocone}%
  \index{vertex of a cocone}%
  由函数 $i:A\to{}D$ 和 $j:B\to{}D$ 以及一个同伦 $h : \prd{c:C} (i(f(c))=j(g(c)))$ 组成：
  \[\uppercurveobject{{ }}\lowercurveobject{{ }}\twocellhead{{ }}
  \xymatrix{C \ar^g[r] \ar_f[d] \drtwocell{^h} & B \ar^j[d] \\ A \ar_i[r] & D
  }\]
  我们用 $\cocone{\Ddiag}{D}$ 表示所有这类上锥的类型，即
  \[ \cocone{\Ddiag}{D} \defeq
  \sm{i:A\to D}{j:B\to D} \prd{c:C} (i(f(c))=j(g(c))).
  \]
\end{defn}

显然，存在一个在 $\Ddiag$ 下以 $A\sqcup^C B$ 为顶点的典范上锥，它由 $\inl$、$\inr$ 和 $\glue$ 组成。
\[\uppercurveobject{{ }}\lowercurveobject{{ }}\twocellhead{{ }}
\xymatrix{C \ar^g[r] \ar_f[d] \drtwocell{^\glue\ \ } & B \ar^{\inr}[d] \\
  A \ar_-\inl[r] & A\sqcup^CB }\]
下述引理说明，这个上锥是万有上锥。

\begin{lem}\label{thm:pushout-ump}
  \index{universal!property!of pushout}%
  对任意类型 $E$，存在一个等价
  \[ (A\sqcup^C B \to E) \;\eqvsym\; \cocone{\Ddiag}{E}. \]
\end{lem}

\begin{proof}
  考虑任意类型 $E:\type$。
  有一个典范函数 $c_\sqcup$，其定义为
  \[\function{(A\sqcup^CB\to{}E)}{\cocone{\Ddiag}{E}}
  {t}{(t\circ{}\inl,t\circ{}\inr,\mapfunc{t}\circ{}\glue)}\]
  我们非正式地将这个函数记作 $t\mapsto\composecocone{t}c_\sqcup$。
  我们将证明这是一个等价。

  首先，给定一个 $c=(i,j,h):\cocone{\mathscr{D}}{E}$，我们需要构造一个从 $A\sqcup^CB$ 到 $E$ 的映射 $\mathsf{s}(c)$。
  \[\uppercurveobject{{ }}\lowercurveobject{{ }}\twocellhead{{ }}
  \xymatrix{C \ar^g[r] \ar_f[d] \drtwocell{^h} & B \ar^{j}[d] \\
    A \ar_-{i}[r] & E }\]
  映射 $\mathsf{s}(c)$ 按以下方式定义
  \begin{align*}
    \mathsf{s}(c)(\inl(a))&\defeq i(a),\\
    \mathsf{s}(c)(\inr(b))&\defeq j(b),\\
    \mapfunc{\mathsf{s}(c)}(\glue(x))&\defid h(x).
  \end{align*}
我们定义了一个映射
\[\function{\cocone{\Ddiag}{E}}{(A\sqcup^CB\to{}E)}{c}{\mathsf{s}(c)}\]
现在需要证明这个映射是
$t\mapsto{}\composecocone{t}c_\sqcup$
的逆。
一方面，如果 $c=(i,j,h):\cocone{\Ddiag}{E}$，我们有
\begin{align*}
  \composecocone{\mathsf{s}(c)}c_\sqcup & =
  (\mathsf{s}(c)\circ\inl,\mathsf{s}(c)\circ\inr,
  \mapfunc{\mathsf{s}(c)}\circ\glue) \\
  & = (\lamu{a:A} \mathsf{s}(c)(\inl(a)),\;
  \lamu{b:B} \mathsf{s}(c)(\inr(b)),\;
  \lamu{x:C} \mapfunc{\mathsf{s}(c)}(\glue(x))) \\
  & = (\lamu{a:A} i(a),\;
  \lamu{b:B} j(b),\;
  \lamu{x:C} h(x)) \\
  & \jdeq (i, j, h) \\
  & = c.
\end{align*}
%
另一方面，如果 $t:A\sqcup^CB\to{}E$，我们希望证明
$\mathsf{s}(\composecocone{t}c_\sqcup)=t$。
对于 $a:A$，我们有
\[\mathsf{s}(\composecocone{t}c_\sqcup)(\inl(a))=t(\inl(a))\]
因为 $\composecocone{t}c_\sqcup$ 的第一个分量是 $t\circ\inl$。
同样地，对于 $b:B$ 我们有
\[\mathsf{s}(\composecocone{t}c_\sqcup)(\inr(b))=t(\inr(b))\]
而对于 $x:C$ 我们有
\[\mapfunc{\mathsf{s}(\composecocone{t}c_\sqcup)}(\glue(x))
=\mapfunc{t}(\glue(x))\]
因此 $\mathsf{s}(\composecocone{t}c_\sqcup)=t$。

这就证明了 $c\mapsto\mathsf{s}(c)$ 是 $t\mapsto{}\composecocone{t}c_\sqcup$ 的拟逆，符合要求。
\end{proof}

许多标准的同伦论构造都可以表示为（同伦）推出。


\begin{itemize}
\item 伸展 $\unit \leftarrow A \to \unit$ 的推出是\define{纬悬（suspension）} $\susp A$（参见 \cref{sec:suspension}）。%
  \index{suspension}
\symlabel{join}
\item $A \xleftarrow{\proj1} A\times B \xrightarrow{\proj2} B$ 的推出称为 $A$ 和 $B$ 的\define{联结（join）}，记作 $A*B$。%
  \indexdef{join!of types}
\item $\unit \leftarrow A \xrightarrow{f} B$ 的推出是 $f$ 的\define{锥（cone）}或\define{余纤维（cofiber）}。%
  \indexdef{cone!of a function}%
  \indexsee{mapping cone}{cone of a function}%
  \indexdef{cofiber of a function}%
\symlabel{wedge}
\item 如果 $A$ 和 $B$ 分别配有基点 $a_0:A$ 和 $b_0:B$，那么 $A \xleftarrow{a_0} \unit \xrightarrow{b_0} B$ 的推出是\define{楔（wedge）} $A\vee B$。%
  \indexdef{wedge}
\symlabel{smash}
\item 如果 $A$ 和 $B$ 如前所述带点，由 $f(\inl(a))\defeq (a,b_0)$ 和 $f(\inr(b))\defeq (a_0,b)$ 定义 $f:A\vee B \to A\times B$，并满足 $\ap f \glue \defid \refl{(a_0,b_0)}$。
  那么 $f$ 的锥称为\define{缩积（smash product）} $A\wedge B$。%
  \indexdef{smash product}
\end{itemize}


我们将在 \cref{cha:hlevels,cha:homotopy} 中进一步讨论推出。

\begin{rmk}
  如 \cref{subsec:prop-trunc} 所述，符号 $\wedge$ 和 $\vee$ 既用于表示带点空间的缩积和楔，也用于逻辑中的``与''和``或''。
  由于同伦类型论中的类型既可以表现得像空间，也可以表现得像命题，从技术上讲，存在潜在的冲突——但由于它们很少同时扮演两种角色，语境通常能消除歧义。
  此外，缩积和楔仅适用于\emph{带点}空间，而唯一带点的命题是 $\top\jdeq\unit$——并且无论 $\wedge$ 和 $\vee$ 取哪种含义，都有 $\unit\wedge \unit = \unit$ 和 $\unit\vee\unit=\unit$。
\end{rmk}

\index{pushout|)}%

\begin{rmk}
  需要注意的是，余极限一般不保持截断性质。
  例如，$\Sn^0$ 和 \unit 都是集合，但 $\unit \leftarrow \Sn^0 \to \unit$ 的推出是 $\Sn^1$，而 $\Sn^1$ 不是集合。
  因此，如果我们对 $n$-类型范畴（特别是集合范畴）中的余极限感兴趣，就需要以某种方式对余极限进行``截断''。
  我们将在 \cref{sec:hittruncations,cha:hlevels,cha:set-math} 中回到这一点。
\end{rmk}

\section{截断}
\label{sec:hittruncations}

\index{truncation!propositional|(}%
在 \cref{subsec:prop-trunc} 中，我们将命题截断作为一种新的类型形成操作引入；
现在我们可以观察到，它可以作为高阶归纳类型的一个特例而得到。
这将理解截断的问题化归为理解高阶归纳类型的问题，而后者至少是可以通过系统方法来处理的。
这一点也很有趣，因为它提供了我们的第一个真正\emph{递归}的高阶归纳类型的例子，因为其构造子的输入取自正被定义的类型本身（就像后继函数 $\suc:\nat\to\nat$ 那样）。

设 $A$ 是一个类型；我们定义其命题截断 $\brck A$ 为由以下生成的高阶归纳类型：

\begin{itemize}
\item 一个函数 $\bprojf : A \to \brck A$，以及
\item 对于每个 $x,y:\brck A$，一条道路 $x=y$。
\end{itemize}

注意，第二个构造子所定义的正是 $\brck A$ 为纯粹命题这一断言。
因此，$\brck A$ 的定义可以这样理解：$\brck A$ 由一个函数 $A\to\brck A$ 连同它是纯粹命题这一事实\emph{自由生成}。

这个高阶归纳定义的递归原理不难写出：给定任意类型 $B$ 及以下条件

\begin{itemize}
\item 一个函数 $g:A\to B$，以及
\item 对于任意 $x,y:B$，都有一条道路 $x=_B y$，
\end{itemize}

则存在函数 $f:\brck A \to B$，使得

\begin{itemize}
\item 对所有 $a:A$，有 $f(\bproj a) \jdeq g(a)$，且
\item 对于任意 $x,y:\brck A$，函数 $\apfunc f$ 将 $\brck A$ 中指定的道路 $x=y$ 映到 $B$ 中指定的道路 $f(x) = f(y)$（在命题层面）。
\end{itemize}

\index{recursion principle!for truncation}%
这些正是我们在 \cref{subsec:prop-trunc} 中为命题截断的递归原理所陈述的前提——即一个函数 $A\to B$，其中 $B$ 是纯粹命题——且结论的第一部分也与我们那里的陈述完全一致。
第二部分（$\apfunc f$ 的作用）之前并未提及，但在此情形下它实际上是空洞的，因为 $B$ 是纯粹命题，所以其中的\emph{任意}两条道路都自动相等。

\index{induction principle!for truncation}%
$\brck A$ 同样有归纳原理：给定任意 $B:\brck A \to \type$ 及以下条件

\begin{itemize}
\item 一个函数 $g:\prd{a:A} B(\bproj a)$，以及
\item 对于任意 $x,y:\brck A$，$u:B(x)$ 和 $v:B(y)$，都有一条依值道路 $q:\dpath{B}{p(x,y)}{u}{v}$，其中 $p(x,y)$ 是来自 $\brck A$ 第二个构造子的道路，
\end{itemize}

则存在 $f:\prd{x:\brck A} B(x)$，使得对 $a:A$ 有 $f(\bproj a)\jdeq g(a)$，并且还有另一条计算规则。
然而，由于任意两个纯粹命题之间（在同伦意义下）至多只有一个函数，这个归纳原理其实并不太有用（另见 \cref{ex:prop-trunc-ind}）。

\index{truncation!propositional|)}%
\index{truncation!set|(}%

\index{set|(}%
不过，我们可以将这一想法推广，对任意 $n$，构造出类似的落在 $n$-类型中的截断。
例如，我们可以将\emph{$0$-截断} $\trunc0A$ 定义为由以下内容生成：

\begin{itemize}
\item 一个函数 $\tprojf0 : A \to \trunc0 A$，以及
\item 对于每对 $x,y:\trunc0A$ 及每对道路 $p,q:x=y$，都有一条道路 $p=q$。
\end{itemize}

那么，$\trunc0A$ 就是由一个函数 $A\to \trunc0A$ 连同``$\trunc0A$ 是集合''这一断言自由生成的。
它的一个自然的归纳原理是说：给定 $B:\trunc0 A \to \type$ 及以下条件

\begin{itemize}
\item 一个函数 $g:\prd{a:A} B(\tproj0a)$，以及
\item 对于任意 $x,y:\trunc0A$，$z:B(x)$ 和 $w:B(y)$，以及每对道路 $p,q:x=y$ 满足 $r:\dpath{B}{p}{z}{w}$ 和 $s:\dpath{B}{q}{z}{w}$，都有一条 $2$-道路 $v:\dpath{\dpath{B}{-}{z}{w}}{u(x,y,p,q)}{r}{s}$，其中 $u(x,y,p,q):p=q$ 得自 $\trunc0A$ 的第二个构造子，
\end{itemize}

则存在 $f:\prd{x:\trunc0A} B(x)$，使得对所有 $a:A$ 有 $f(\tproj0a)\jdeq g(a)$，并且 $\apdtwo{f}{u(x,y,p,q)}$ 就是上面指定的那条 $2$-道路。
（与命题截断的情形一样，后一个条件实际上是无趣的。）
然而，由此出发，我们可以证明一个更有用的归纳原理。

\begin{lem}\label{thm:trunc0-ind}
  假设给定 $B:\trunc0 A \to \type$ 及 $g:\prd{a:A} B(\tproj0a)$，并且每个 $B(x)$ 都是集合。
  则存在 $f:\prd{x:\trunc0A} B(x)$，使得对所有 $a:A$ 都有 $f(\tproj0a)\jdeq g(a)$。
\end{lem}

\begin{proof}
  只需对如上所述的任意 $x,y,z,w,p,q,r,s$ 构造一条 $2$-道路 $v:\dpath{B}{u(x,y,p,q)}{r}{s}$ 即可。
  然而，根据依值 $2$-道路的定义，这是纤维 $B(y)$ 中的一条普通 $2$-道路。
  由于 $B(y)$ 是集合，任意两条平行道路之间都存在 $2$-道路。
\end{proof}

这蕴含了所期望的泛性质。

\begin{lem}\label{thm:trunc0-lump}
  \index{universal!property!of truncation}%
  对于任意集合 $B$ 和任意类型 $A$，与 $\tprojf0:A\to \trunc0A$ 复合确定了一个等价
  \[ \eqvspaced{(\trunc0A\to B)}{(A\to B)}. \]
\end{lem}

\begin{proof}
  将 \cref{thm:trunc0-ind} 应用于常值族这一特殊情形，得到一个从右到左的映射，它是从左到右``与 $\tprojf0$ 复合''映射的右逆。
  为证明它也是左逆，任取 $h:\trunc0A\to B$，并将 \cref{thm:trunc0-ind} 应用于复合 $a\mapsto h(\tproj0a)$ 来定义 $h':\trunc0A\to B$。
  于是有 $h'(\tproj0a)=h(\tproj0a)$。

  然而，由于 $B$ 是集合，对于任意 $x:\trunc0A$，类型 $h(x)=h'(x)$ 是纯粹命题，因而也是集合。
  于是，由 \cref{thm:trunc0-ind}，对任意 $a:A$ 有 $h'(\tproj0a)=h(\tproj0a)$ 这一观察蕴含了对任意 $x:\trunc0A$ 都有 $h(x)=h'(x)$，从而 $h=h'$。
\end{proof}

\index{limit!of sets}%
\index{colimit!of sets}%
例如，这使我们能够构造集合的余极限。
我们已经看到，如果 $A \xleftarrow{f} C \xrightarrow{g} B$ 是集合的一个伸展，那么推出 $A\sqcup^C B$ 可能不再是集合。
（例如，若 $A$ 和 $B$ 都是 \unit 而 $C$ 是 \bool，则推出是 $\Sn^1$。）
然而，我们可以通过截断，构造出一个作为集合的推出，并且它对其他集合具有预期的泛性质。

\begin{lem}\label{thm:set-pushout}
  \index{universal!property!of pushout}%
  设 $A \xleftarrow{f} C \xrightarrow{g} B$ 是集合的一个伸展\index{span}。
  那么对任意集合 $E$，存在一个典范等价
  \[ \Parens{\trunc0{A\sqcup^C B} \to E} \;\eqvsym\; \cocone{\Ddiag}{E}. \]
\end{lem}

\begin{proof}
  组合 \cref{thm:pushout-ump,thm:trunc0-lump} 中的等价即得。
\end{proof}

我们将 $\trunc0{A\sqcup^C B}$ 称为 $f$ 与 $g$ 的 \define{集合推出}
\indexdef{set-pushout}%
\index{pushout!of sets}，
以区别于（同伦）推出 $A\sqcup^C B$。
或者，我们也可以修改 \cref{sec:colimits} 中推出的定义，直接加入 $0$-截断构造子，从而避免事后再做截断。
类似的说明适用于任何一种集合余极限；我们将在 \cref{cha:set-math} 中进一步探讨这一点。

然而，尽管上述 $0$-截断的定义是可行的——它给出了我们期望的结果，并且是相容的——但它存在一些问题。
首先，它不能很好地融入高阶归纳类型的一般理论。
一般来说，直接处理像 $\trunc0A$ 的第二个构造子那样的构造子是很棘手的，因为它的\emph{输入}不仅涉及被定义类型的元素，还涉及其中的道路。

不过，这个问题可以比较容易地解决。
回想一下，在 \cref{sec:bool-nat} 中我们提到，可以通过让一个构造子接受一个类型为 $\nat\to W$ 的单个参数，从而允许归纳类型 $W$ 的构造子接受“无限多个”类型为 $W$ 的参数。
这背后有一个一般性原则：要为输入形式古怪的构造子建模，可以使用一个辅助的归纳类型（例如 \nat）来参数化它们，从而将输入化为一个以归纳类型为定义域的简单函数。

对于 $0$-截断，我们可以考虑辅助的\emph{高阶}归纳类型 $S$，它由两个点 $a,b:S$ 和两条道路 $p,q:a=b$ 生成。
那么，$\trunc 0A$ 那个略显可疑的构造子就可以替换为下面这个毫无问题的构造子：


\begin{itemize}
\item 对每个 $f:S\to \trunc 0A$，一条道路 $\apfunc{f}(p) = \apfunc{f}(q)$。
\end{itemize}


因为给出一个从 $S$ 出发的映射，相当于给出两个点以及它们之间的两条平行道路，所以这导出了相同的归纳原理。

\index{set|)}%

\index{truncation!set|)}%
\index{truncation!n-truncation@$n$-truncation}%
然而，我们当前 $0$-截断定义的一个更严重的问题是，它不太容易推广。
如果我们想为所有 $n:\nat$ 统一地描述一种将类型“$n$-截断”到 $n$-类型的定义方式，那么当前的方法是行不通的，因为第二个构造子所需的参数个数会随 $n$ 的增大而增加。
因此，在 \cref{sec:truncations} 中，我们将基于上面引入的类型 $S$ 等价于圆 $\Sn^1$ 这一观察，采用不同的想法来构造这些截断。
这一方法将 $0$-截断作为特例包含在内，并且满足 \cref{thm:trunc0-ind,thm:trunc0-lump} 的推广版本。

\section{商}
\label{sec:set-quotients}

集合的一种特别重要的余极限是按\emph{关系}取的\emph{商}。
也就是说，设 $A$ 是一个集合，且 $R:A\times A \to \prop$ 是一个命题族（一个\define{纯粹关系}）。
\indexdef{relation!mere}%
\indexdef{mere relation}%
它的商应该是两个投影
\[ \tsm{a,b:A} R(a,b) \rightrightarrows A. \]
的集合余等子。
我们也可以直接将其描述为高阶归纳类型 $A/R$，它由以下内容生成：
\index{set-quotient|(defstyle}%
\indexsee{quotient of sets}{set-quotient}%
\indexsee{type!quotient}{set-quotient}%


\begin{itemize}
\item 一个函数 $q:A\to A/R$；
\item 对每个满足 $R(a,b)$ 的 $a,b:A$，一个相等 $q(a)=q(b)$；以及
\item $0$-截断构造子：对所有 $x,y:A/R$ 和 $r,s:x=y$，有 $r=s$。
\end{itemize}


我们有时会将这个高阶归纳类型 $A/R$ 称为 $A$ 按 $R$ 的 \define{集合商}，以强调它按定义产生一个集合。
（在同伦论中存在更一般的“商”的概念，但它们大多超出了本书的范围。不过，在 \cref{sec:rezk} 中，我们将考虑类型按 $1$-群胚取“商”的情形，这是比集合商更高一层的构造。）

\begin{rmk}\label{rmk:quotient-of-non-set}
  对于集合商的定义及其大多数性质而言，$A$ 实际上并不一定需要是集合。
  不过，集合的情形通常最受关注。
\end{rmk}

\begin{lem}\label{thm:quotient-surjective}
  函数 $q:A\to A/R$ 是满的。
\end{lem}

\begin{proof}
  我们需要证明，对于任何 $x:A/R$，都仅仅存在一个 $a:A$ 使得 $q(a)=x$。
  我们使用 $A/R$ 的归纳原理。
  第一种情况是平凡的：如果 $x$ 是 $q(a)$，那么当然仅仅存在一个 $a$ 使得 $q(a)=q(a)$。
  由于目标是一个命题，它自动尊重所有道路构造子，因此我们就证完了。
\end{proof}

现在我们可以证明，集合商具有所期望的（集合）余等子的泛性质。

\begin{lem}\label{thm:quotient-ump}
  对于任何集合 $B$，与 $q$ 预复合给出一个等价：
  \[ \eqvspaced{(A/R \to B)}{\Parens{\sm{f:A\to B} \prd{a,b:A} R(a,b) \to (f(a)=f(b))}}.\]
\end{lem}

\begin{proof}
  $\blank\circ q$ 的拟逆（从右到左）就是 $A/R$ 的递归原理。
  也就是说，给定 $f:A\to B$ 使得
  \narrowequation{\prd{a,b:A} R(a,b) \to (f(a)=f(b)),} 我们通过 $\bar f(q(a))\defeq f(a)$ 来定义 $\bar f:A/R\to B$。
  这个定义方程说的正是 $(f\mapsto \bar f)$ 是 $(\blank\circ q)$ 的一个右逆。

  要证明它同时也是左逆，我们必须证明对于任意 $g:A/R\to B$ 和 $x:A/R$ 有 $g(x) = \overline{g\circ q}(x)$。
  然而，根据 \cref{thm:quotient-surjective}，仅仅存在 $a$ 使得 $q(a)=x$。
  由于我们欲证的相等是一个命题，我们可以假设确实存在这样的 $a$，此时便有 $g(x) = g(q(a)) = \overline{g\circ q}(q(a)) = \overline{g\circ q}(x)$。
\end{proof}

当然，经典数学中通常考虑的情况是 $R$ 为一个\define{等价关系（equivalence relation）}，即我们有
\indexdef{relation!equivalence}%
\indexsee{equivalence!relation}{relation, equivalence}%
%


\begin{itemize}
\item \define{自反性（reflexivity）}：$\prd{a:A} R(a,a)$，
  \indexdef{reflexivity!of a relation}%
  \indexdef{relation!reflexive}%
\item \define{对称性（symmetry）}：$\prd{a,b:A} R(a,b) \to R(b,a)$，以及
  \indexdef{symmetry!of a relation}%
  \indexdef{relation!symmetric}%
\item \define{传递性（transitivity）}：$\prd{a,b,c:C} R(a,b) \times R(b,c) \to R(a,c)$。
  \indexdef{transitivity!of a relation}%
  \indexdef{relation!transitive}%
\end{itemize}


%
此时，集合商 $A/R$ 将具备更多良好的性质，我们将在 \cref{sec:piw-pretopos} 中看到：例如，我们有 $R(a,b) \eqvsym (\id[A/R]{q(a)}{q(b)})$。
\symlabel{equivalencerelation}
我们通常将等价关系 $R(a,b)$ 写作中缀形式 $a\eqr b$。

等价关系的商也可以用其他方式构造。
集合论的方法是考虑等价类构成的集合，作为 $A$ 的幂集\index{power set}的一个子集。
我们也可以在类型论中模仿这种``非直谓''的构造。
\index{impredicative!quotient}

\begin{defn}
  一个谓词 $P:A\to\prop$ 是关系 $R : A \times A \to \prop$ 的一个\define{等价类（equivalence class）}
  \indexdef{equivalence!class}%
  如果仅仅存在一个 $a:A$，使得对所有的 $b:A$ 有 $\eqv{R(a,b)}{P(b)}$。
\end{defn}

由于 $R$ 和 $P$ 都是命题，等价 $\eqv{R(a,b)}{P(b)}$ 与蕴含 $R(a,b) \to P(b)$ 和 $P(b) \to R(a,b)$ 是一回事。
当然，对于任何 $a:A$，我们有典范的等价类 $P_a(b) \defeq R(a,b)$。

\begin{defn}\label{def:VVquotient}
  我们定义
  \begin{equation*}
    A\sslash R \defeq \setof{ P:A\to\prop | P \text{ is an equivalence class of } R}.
  \end{equation*}
  函数 $q':A\to A\sslash R$ 定义为 $q'(a) \defeq P_a$。
\end{defn}

\begin{thm}
  对于 $A$ 上的任意等价关系 $R$，类型 $A\sslash R$ 等价于集合商 $A/R$。
\end{thm}

\begin{proof}
  首先注意，如果 $R(a,b)$，那么由于 $R$ 是等价关系，对于任何 $c:A$ 我们有 $R(a,c) \Leftrightarrow R(b,c)$。
  因此，根据泛等性有 $R(a,c) = R(b,c)$，再由函数外延性即得 $P_a=P_b$，也就是说 $q'(a)=q'(b)$。
  于是，由 \cref{thm:quotient-ump} 我们得到一个诱导映射 $f:A/R \to A\sslash R$，满足 $f\circ q = q'$。

  我们证明 $f$ 既是单的又是满的，从而是一个等价。
  满射性直接由 $q'$ 是满的得出，而 $q'$ 的满射性本质上由 $A\sslash R$ 的定义保证。
  对于单射性，如果 $f(x)=f(y)$，为了证明命题 $x=y$，根据 $q$ 的满射性，我们可以假设存在 $a,b:A$ 使得 $x=q(a)$ 且 $y=q(b)$。
  那么对于任意 $c:A$ 有 $R(a,c) = f(q(a))(c) = f(q(b))(c) = R(b,c)$，特别地有 $R(a,b) = R(b,b)$。
  但由于 $R$ 是等价关系，$R(b,b)$ 是有实例的，故而 $R(a,b)$ 也有实例。
  因此 $q(a)=q(b)$，即 $x=y$。
\end{proof}

在 \cref{subsec:quotients} 中，我们将给出这个定理的另一个证明。
请注意，与 $A/R$ 不同，构造 $A\sslash R$ 会提升宇宙层级：如果 $A:\UU_i$ 且 $R:A\to A\to \prop_{\UU_i}$，那么在定义 $A\sslash R$ 时，我们也必须使用 $\prop_{\UU_i}$ 来包含所有等价类，因此 $A\sslash R : \UU_{i+1}$。
当然，如果我们假设 \cref{subsec:prop-subsets} 中的命题大小重整公理，就可以避免这个问题。

\begin{rmk}\label{defn-Z}
前面的两种构造提供了一般性的商构造，但在某些具体情形下，可能存在更简单的构造。
例如，我们可以将整数 \Z 定义为一个集合商
\indexdef{integers}%
\indexdef{number!integers}%
%
\[ \Z \defeq (\N \times \N)/{\eqr} \]
%
其中 $\eqr$ 是由下式定义的等价关系：
%
\[ (a,b) \eqr (c,d) \defeq (a + d = b + c). \]
%
换言之，一对数 $(a,b)$ 代表整数 $a - b$。
然而，在这种情况下，等价类存在\emph{典范代表元（canonical representatives）}：形如 $(n,0)$ 或 $(0,n)$ 的那些对。
\end{rmk}

以下引理表明，当出现这类情况时，我们就不再需要前面任何一种通用的商构造了。
（一个函数 $r:A\to A$ 如果满足 $r\circ r = r$，则称之为\define{幂等的（idempotent）}
\indexdef{function!idempotent}%
\indexdef{idempotent!function}%
。）

\begin{lem}\label{lem:quotient-when-canonical-representatives}
  设 $\eqr$ 是集合 $A$ 上的一个关系，并且存在一个幂等函数 $r: A \to A$，使得对一切 $x, y: A$ 有 $\eqv{(r(x) = r(y))}{(x \eqr y)}$。
  （这蕴含 $\eqr$ 是一个等价关系。）
  那么类型
  %
  \begin{equation*}
    (A/{\eqr}) \defeq \Parens{\sm{x : A} r(x) = x}
  \end{equation*}
  %
  满足 $A$ 关于 $\eqr$ 的集合商的泛性质，因此与之等价。
  换句话说，存在一个映射 $q : A \to (A/{\eqr})$，使得对于任意集合 $B$，与 $q$ 的预复合诱导出一个等价
  %
  \begin{equation}
    \label{eq:quotient-when-canonical}
    \Parens{(A/{\eqr}) \to B} \eqvsym \Parens{\sm{g : A \to B} \prd{x, y : A} (x \eqr y) \to (g(x) = g(y))}.
  \end{equation}
\end{lem}

\begin{proof}
  令 $i : \prd{x : A} r(r(x)) = r(x)$ 见证 $r$ 的幂等性。
  映射 $q : A \to (A/{\eqr})$ 由 $q(x) \defeq (r(x), i(x))$ 定义。
  注意，由于 $A$ 是集合，我们有 $q(x)=q(y)$ 当且仅当 $r(x)=r(y)$，从而（根据假设）当且仅当 $x \eqr y$。
  我们在\eqref{eq:quotient-when-canonical}中如下定义从左到右的映射 $e$：
  \[ e(f) \defeq (f \circ q, \nameless), \]
  %
  其中下划线 $\nameless$ 表示如下的证明：如果 $x, y : A$ 且 $x \eqr y$，那么如上所述 $q(x)=q(y)$，因此 $f(q(x)) = f(q(y))$。
  为证明 $e$ 是一个等价，考虑反方向的映射 $e'$，其定义为
  %
  \[ e'(g, s) (x,p) \defeq g(x). \]
  %
  给定任意 $f : (A/{\eqr}) \to B$，有
  %
  \[ e'(e(f))(x, p) \jdeq f(q(x)) \jdeq f(r(x), i(x)) = f(x, p) \]
  %
  最后一个等式成立是因为 $p : r(x) = x$，并且由于 $A$ 是集合，我们有 $(x,p) = (r(x), i(x))$。
  类似地，我们计算
  %
  \[ e(e'(g, s)) \jdeq e(g \circ \proj{1}) \jdeq (g \circ \proj{1} \circ q, {\nameless}). \]
  %
  由于 $B$ 是集合，我们不必考虑 $\nameless$ 的部分，而对于第一个分量，我们有
  %
  \[ g(\proj{1}(q(x))) \jdeq g(r(x)) = g(x), \]
  %
  其中最后一个等式成立是因为 $r(x) \eqr x$，而由假设 $s$，$g$ 尊重 $\eqr$。
\end{proof}

\begin{cor}\label{thm:retraction-quotient}
  设 $p:A\to B$ 是集合之间的一个收缩。
  那么 $B$ 是 $A$ 模去由下式定义的等价关系 $\eqr$ 的商：
  \[ (a_1 \eqr a_2) \defeq (p(a_1) = p(a_2)). \]
\end{cor}

\begin{proof}
  设 $s:B\to A$ 是 $p$ 的一个截面。
  那么 $s\circ p : A\to A$ 是一个幂等函数，它关于这个 $\eqr$ 满足 \cref{lem:quotient-when-canonical-representatives} 的条件，并且 $s$ 诱导了从 $B$ 到其不动点集的一个同构。
\end{proof}

\begin{rmk}\label{Z-quotient-by-canonical-representatives}
\cref{lem:quotient-when-canonical-representatives} 可应用于 $\Z$，其中幂等函数 $r : \N \times \N \to \N \times \N$ 定义为
%
\begin{equation*}
  r(a, b) =
  \begin{cases}
    (a - b, 0) & \text{if $a \geq b$,} \\
    (0, b - a) & \text{otherwise.}
  \end{cases}
\end{equation*}
%
（即使构造性地看，这也是一个合法的定义，因为 $\N$ 上的关系 $\geq$ 是可判定的。）
这样，一个非负整数典范地由 $(k, 0)$ 代表，一个非正整数由 $(0, m)$ 代表，其中 $k,m:\N$。
这种分情况处理的方式蕴含着下面这个关于整数的``归纳原理''，它在 \cref{cha:homotopy} 中会很有用。
\index{natural numbers}%
（照例，我们将自然数 $n$ 与相应的非负整数等同，即等同于 $(n,0):\N\times\N$ 在 $\Z$ 中的像。）
\end{rmk}

\begin{lem}\label{thm:sign-induction}
  \index{integers!induction principle for}%
  \index{induction principle!for integers}%
  设 $P:\Z\to\type$ 是一个类型族，并且我们已有
  \begin{itemize}
  \item $d_0: P(0)$，
  \item $d_+: \prd{n:\N} P(n) \to P(\suc(n))$，以及
  \item $d_- : \prd{n:\N} P(-n) \to P(-\suc(n))$。
  \end{itemize}
  那么存在 $f:\prd{z:\Z} P(z)$，使得
  \begin{itemize}
    \item $f(0) = d_0$，
    \item $f(\suc(n)) = d_+(n,f(n))$ 对所有 $n:\N$ 成立，且
    \item $f(-\suc(n)) = d_-(n,f(-n))$ 对所有 $n:\N$ 成立。
  \end{itemize}
\end{lem}

\begin{proof}
  在此证明中，我们令 $\Z$ 表示 $\sm{x:\N\times\N}(r(x)=x)$，其中 $r$ 是上面给出的幂等函数。
  （然后我们可以将结果传输到任何与 $\Z$ 等价的定义上。）
  令 $q:\N\times\N\to\Z$ 为商映射，如 \cref{lem:quotient-when-canonical-representatives} 中那样由 $q(x) = (r(x),i(x))$ 定义。
  现在定义 $Q\defeq P\circ q:\N\times \N \to \type$。
  通过将给定的数据沿适当的等式传输，我们得到
  \begin{align*}
    d'_0 &: Q(0,0)\\
    d'_+ &: \prd{n:\N} Q(n,0) \to Q(\suc(n),0)\\
    d'_- &: \prd{n:\N} Q(0,n) \to Q(0,\suc(n)).
  \end{align*}
  另请注意，由于 $q(n,m) = q(\suc(n),\suc(m))$，我们有一个诱导的等价
  \[e_{n,m}:\eqv{Q(n,m)}{Q(\suc(n),\suc(m))}.\]
  然后我们可以通过对 $x$ 进行双重归纳来构造 $g:\prd{x:\N\times \N} Q(x)$：
  \begin{align*}
    g(0,0) &\defeq d'_0,\\
    g(\suc(n),0) &\defeq d'_+(n,g(n,0)),\\
    g(0,\suc(m)) &\defeq d'_-(m,g(0,m)),\\
    g(\suc(n),\suc(m)) &\defeq e_{n,m}(g(n,m)).
  \end{align*}
  现在我们有 $\proj1 : \Z \to \N\times\N$，且满足 $q\circ \proj1 = \idfunc$。
  因此，特别地，我们有 $Q\circ \proj1 = P$，从而得到一族等价 $s:\prd{z:\Z} \eqv{Q(\proj1(z))}{P(z)}$。
  于是，我们可以定义 $f(z) = s(z,g(\proj1(z)))$ 来得到 $f:\prd{z:\Z} P(z)$，并验证所要求的等式。
\end{proof}

我们有时会用模式匹配语法来表示从 \cref{thm:sign-induction} 得到的函数 $f:\prd{z:\Z} P(z)$，它包含 $d_0$、$d_+$ 和 $d_-$ 这三种情形：
\begin{align*}
  f(0) &\defid d_0\\
  f(\suc(n)) &\defid d_+(n,f(n))\\
  f(-\suc(n)) &\defid d_-(n,f(-n))
\end{align*}
正如我们在处理高阶归纳类型的道路构造子时所做的那样，我们使用 $\defid$ 而非 $\defeq$，以表明由 \cref{thm:sign-induction} 蕴含的``计算''规则仅仅是命题相等性。
例如，通过这种方式，我们可以为任意整数 $n$ 定义一个环路的 $n$ 重拼接。

\begin{cor}\label{thm:looptothe}
  \indexdef{path!concatenation!n-fold@$n$-fold}%
  设 $A$ 是一个类型，$a:A$ 且 $p:a=a$。
  存在一个函数 $\prd{n:\Z} (a=a)$，记作 $n\mapsto p^n$，其定义如下：
  \begin{align*}
    p^0 &\defid \refl{a}\\
    p^{n+1} &\defid p^n \ct p
    & &\text{for $n\ge 0$}\\
    p^{n-1} &\defid p^n \ct \opp p
    & &\text{for $n\le 0$.}
  \end{align*}
\end{cor}

关于整数，我们将在 \cref{sec:free-algebras,sec:field-rati-numb} 作进一步讨论。

\index{set-quotient|)}%

\section{代数}
\label{sec:free-algebras}

高阶归纳类型不仅能用来构造球面、胞腔复形等高维对象，即便仅限于集合的层面，它们也非常有用。
在 \cref{thm:set-pushout} 中我们已经见过一个例子：它们使我们能够构造任意集合图表的余极限，而这在 \cref{cha:typetheory} 的基础类型论中是无法做到的。
当我们研究带有代数结构的集合时，高阶归纳类型同样极为有用。

作为本节贯穿的例子，我们考虑\emph{群（groups）}——它对大多数数学家来说都很熟悉，并能展现本质的现象（后续章节中也会用到群）。
不过，我们所说的大部分内容同样适用于任何类型的代数结构。

\index{monoid|(}%

\begin{defn}
  一个\define{幺半群（monoid）}
  \indexdef{monoid}%
  是一个集合 $G$，并配备以下结构：
  \begin{itemize}
  \item 一个\emph{乘法（multiplication）}
    \indexdef{multiplication!in a monoid}%
    \indexdef{multiplication!in a group}%
    函数 $G\times G\to G$，以中缀形式记作 $(x,y) \mapsto x\cdot y$；以及
  \item 一个\emph{单位元（unit）}
    \indexdef{unit!of a monoid}%
    \indexdef{unit!of a group}%
    $e:G$；使得
  \item 对任意 $x:G$，有 $x\cdot e = x$ 且 $e\cdot x = x$；并且
  \item 对任意 $x,y,z:G$，有 $x\cdot (y\cdot z) = (x\cdot y)\cdot z$。
    \index{associativity!in a monoid}%
    \index{associativity!in a group}%
  \end{itemize}
  一个\define{群（group）}
  \indexdef{group}%
  是一个幺半群 $G$，并额外配备：
  \begin{itemize}
  \item 一个\emph{取逆（inversion）}函数 $i:G\to G$，记作 $x\mapsto \opp x$；使得
    \index{inverse!in a group}%
  \item 对任意 $x:G$，有 $x\cdot \opp x = e$ 且 $\opp x \cdot x = e$。
  \end{itemize}
\end{defn}

\begin{rmk}\label{rmk:infty-group}
请注意，我们要求群必须是一个集合。
我们可以考虑一种更一般的``$\infty$-群''%
\index{.infinity-group@$\infty$-group}
的概念，它不必是集合，但这会使我们离题太远，目前并不合适。
根据我们目前的定义，可以预期由此产生的``群论''其行为将类似于集合论数学中的群论（需要注意的是，除非我们假设 \LEM{}，否则它将会是``构造性''的群论）。\index{mathematics!constructive}
\end{rmk}

\begin{eg}
  自然数集 \N 在加法下构成一个幺半群，单位元是 $0$；在乘法下也构成一个幺半群，单位元是 $1$。
  如果我们按通常的方式定义整数集 \Z 上的算术运算，那么如所预期的那样，它们在加法下构成一个群，在乘法下构成一个幺半群（当然，还是一个环）。
  例如，如果 $u, v \in \Z$ 分别由 $(a,b)$ 和 $(c,d)$ 表示，那么 $u + v$ 由 $(a + c, b + d)$ 表示，$-u$ 由 $(b, a)$ 表示，而 $u v$ 由 $(a c + b d, a d + b c)$ 表示。
\end{eg}

\begin{eg}\label{thm:homotopy-groups}
  在 \cref{sec:equality} 中我们基本上已经观察到：若 $(A,a)$ 是一个带点类型，则其环路空间\index{loop space} $\Omega(A,a)\defeq (\id[A]aa)$ 具有群的所有结构，只不过它一般不是集合。
  它应该是 \cref{rmk:infty-group} 中提到的那种意义上的“$\infty$-群”，但我们也可以通过截断把它变成一个群。
  具体来说，我们定义 $A$ 的以 $a:A$ 为基点的\define{基本群（fundamental group）}
  \indexsee{group!fundamental}{fundamental group}%
  \indexdef{fundamental!group}%
  为
  \[\pi_1(A,a)\defeq \trunc0{\Omega(A,a)}.\]
  它继承了群结构；例如，乘法 $\pi_1(A,a) \times \pi_1(A,a) \to \pi_1(A,a)$ 是由道路的拼接通过截断的双重归纳定义的。

  更一般地，$(A,a)$ 的\define{$n$ 阶同伦群（$n^{\mathrm{th}}$ homotopy group）}
  \index{homotopy!group}%
  \indexsee{group!homotopy}{homotopy group}%
  是 $\pi_n(A,a)\defeq \trunc0{\Omega^n(A,a)}$。
  \index{loop space!iterated}%
  于是，对于 $n\ge 1$ 有 $\pi_n(A,a) = \pi_1(\Omega^{n-1}(A,a))$，所以它也是一个群。
  （当 $n=0$ 时，$\pi_0(A) \jdeq \trunc0 A$，它不是一个群。）
  此外，Eckmann--Hilton 论证 \index{Eckmann--Hilton argument} （\cref{thm:EckmannHilton}）表明，若 $n\ge 2$，则 $\pi_n(A,a)$ 是一个\emph{交换}\index{group!abelian}群，即对所有 $x,y$ 有 $x\cdot y = y\cdot x$。
  \cref{cha:homotopy} 将主要研究这些群。
\end{eg}

\index{algebra!free}%
\index{free!algebraic structure}%
群论中的一个重要概念是由一个集合生成的\emph{自由群}，或更一般地，由一组生成元和关系\index{generator!of a group}表现的群。
类型论中一个熟知的事实是，\emph{某些}自由代数对象可以用\emph{普通}归纳类型来定义。
\symlabel{lst-freemonoid}%
\indexdef{type!of lists}%
\indexsee{list type}{type, of lists}%
\index{monoid!free|(}%
例如，集合 $A$ 上的自由幺半群可以等同于 $A$ 中元素的\emph{有限列表}\index{finite!lists, type of}的类型 $\lst A$，它由以下方式归纳生成：

\begin{itemize}
\item 一个构造子 $\nil:\lst A$，以及
\item 对每个 $\ell:\lst A$ 和 $a:A$，一个元素 $\cons(a,\ell):\lst A$。
\end{itemize}

我们有一个明显的包含映射 $\eta : A\to \lst A$，由 $a\mapsto \cons(a,\nil)$ 定义。
$\lst A$ 上的幺半群运算是列表的拼接，递归地定义为
\begin{align*}
  \nil \cdot \ell &\defeq \ell\\
  \cons (a,\ell_1) \cdot \ell_2 &\defeq \cons(a, \ell_1\cdot\ell_2).
\end{align*}
利用 $\lst A$ 的归纳原理不难证明，$\lst A$ 是一个集合，并且列表的拼接是结合的
\index{associativity!of list concatenation}%
且以 $\nil$ 为单位元。
因此，$\lst A$ 是一个幺半群。

\begin{lem}\label{thm:free-monoid}
  \indexsee{free!monoid}{monoid, free}%
  对任意集合 $A$，类型 $\lst A$ 是 $A$ 上的自由幺半群。
  换言之，对任意幺半群 $G$，与 $\eta$ 复合是一个等价
  \[ \eqv{\hom_{\mathrm{Monoid}}(\lst A,G)}{(A\to G)}, \]
  其中 $\hom_{\mathrm{Monoid}}(\blank,\blank)$ 表示幺半群同态（保持乘法和单位元的函数）的集合。
  \indexdef{homomorphism!monoid}%
  \indexdef{monoid!homomorphism}%
\end{lem}

\begin{proof}
  给定 $f:A\to G$，我们递归地定义 $\bar{f}:\lst A \to G$：
  \begin{align*}
    \bar{f}(\nil) &\defeq e\\
    \bar{f}(\cons(a,\ell)) &\defeq f(a) \cdot \bar{f}(\ell).
  \end{align*}
  由归纳法不难证明 $\bar{f}$ 是一个幺半群同态，并且映射 $f\mapsto \bar f$ 是 $(\blank\circ \eta)$ 的一个拟逆；详见 \cref{ex:free-monoid}。
\end{proof}

\index{monoid!free|)}%

自由幺半群的这种构造之所以可能，本质上是因为自由幺半群的元素具有可计算的典范形式（即有限列表）。
然而，其他自由（以及展示的）代数结构——例如群——的元素通常没有\emph{可计算的}典范形式。
举例来说，群展示中字的相等性在算法上\index{algorithm}是不可判定的。
尽管如此，我们仍然可以将自由代数对象描述为\emph{高阶}归纳类型，只需将所有公理化等式声明为道路构造子。

\indexsee{free!group}{group, free}%
\index{group!free|(}%
例如，令 $A$ 是一个集合，并定义高阶归纳类型 $\freegroup{A}$，它具有以下生成子。

\begin{itemize}
\item 一个函数 $\eta:A\to \freegroup{A}$。
\item 一个函数 $m: \freegroup{A} \times \freegroup{A} \to \freegroup{A}$。
\item 一个元素 $e:\freegroup{A}$。
\item 一个函数 $i:\freegroup{A} \to \freegroup{A}$。
\item 对每个 $x,y,z:\freegroup{A}$，一个等式 $m(x,m(y,z)) = m(m(x,y),z)$。
\item 对每个 $x:\freegroup{A}$，等式 $m(x,e) = x$ 和 $m(e,x) = x$。
\item 对每个 $x:\freegroup{A}$，等式 $m(x,i(x)) = e$ 和 $m(i(x),x) = e$。
\item $0$-截断构造子：对任意 $x,y:\freegroup{A}$ 和 $p,q:x=y$，有 $p=q$。
\end{itemize}

第一个构造子给出从 $A$ 到 $\freegroup{A}$ 的映射。
接下来的三个构造子为 $\freegroup{A}$ 赋予群的运算：乘法、单位元和取逆。
随后的三个构造子断言群的公理：结合律\index{associativity}、单位元律和逆元律。
最后，最后一个构造子断言 $\freegroup{A}$ 是一个集合。

因此，$\freegroup{A}$ 是一个群。
同样不难证明：

\begin{thm}
  \index{universal!property!of free group}%
  $\freegroup{A}$ 是集合 $A$ 上的自由群。
  换言之，对任意（集合意义上的）群 $G$，与 $\eta:A\to \freegroup{A}$ 复合确定了一个等价
  \[ \hom_{\mathrm{Group}}(\freegroup{A},G) \eqvsym (A\to G) \]
  其中 $\hom_{\mathrm{Group}}(\blank,\blank)$ 表示两个群之间所有群同态所成的集合。
  \indexdef{group!homomorphism}%
  \indexdef{homomorphism!group}%
\end{thm}

\begin{proof}
  高阶归纳类型 $\freegroup{A}$ 的递归原理\emph{恰恰}是说：若 $G$ 是一个群且给定 $f:A\to G$，则我们得到 $\bar{f}:\freegroup{A} \to G$。
  其计算规则给出 $\bar{f}\circ \eta \jdeq f$，并且 $\bar f$ 是群同态。
  因此，$(\blank\circ \eta) :  \hom_{\mathrm{Group}}(\freegroup{A},G) \to (A\to G)$ 有一个右逆。
  利用 $\freegroup{A}$ 的归纳原理不难证明它同时也是左逆。
\end{proof}

\index{acceptance}
我们不妨退后一步，回顾一下刚才做了什么。
我们证明了任意集合上的自由群存在，而\emph{无需}给出它的显式构造。
本质上，我们所要做的仅仅是将它应当满足的泛性质写下来。
在集合论中，要达到类似的结果，需要诉诸伴随函子定理\index{adjoint!functor theorem}之类的``黑箱''；而类型论则将此类构造内建于数学基础之中。

当然，有时自由代数结构的具体描述也是有用的。
对于自由群，我们可以利用商来给出这样的描述。
考虑 $\lst{A+A}$，其中在 $A+A$ 里我们把 $\inl(a)$ 写成 $a$，把 $\inr(a)$ 写成 $\hat{a}$（意在表示 $a$ 的形式逆）。
$\lst{A+A}$ 中的元素就是集合 $A$ 上自由群的\emph{字}。

\begin{thm}
  令 $A$ 为一个集合，并令 $\freegroupx{A}$ 是 $\lst{A+A}$ 模去下述关系所得的集合商。
  \begin{align*}
    (\dots,a_1,a_2,\widehat{a_2},a_3,\dots) &=
    (\dots,a_1,a_3,\dots)\\
    (\dots,a_1,\widehat{a_2},a_2,a_3,\dots) &=
    (\dots,a_1,a_3,\dots).
  \end{align*}
  那么 $\freegroupx{A}$ 也是集合 $A$ 上的自由群。
\end{thm}

\begin{proof}
  首先证明 $\freegroupx{A}$ 是一个群。
  我们已经知道 $\lst{A+A}$ 是一个幺半群；我们断言这个幺半群结构可以下降到商上。
  我们通过双商递归定义 $\freegroupx{A} \times \freegroupx{A} \to \freegroupx{A}$；只需验证给定的关系所生成的等价关系在列表的拼接操作下保持不变。
  类似地，我们通过商归纳证明结合律和单位元律。

  为了在 $\freegroupx{A}$ 中定义逆元，我们首先通过列表递归定义 $\mathsf{reverse}:\lst B\to\lst B$：
  \begin{align*}
    \mathsf{reverse}(\nil) &\defeq \nil,\\
    \mathsf{reverse}(\cons(b,\ell))&\defeq \mathsf{reverse}(\ell)\cdot \cons(b,\nil).
  \end{align*}
  现在，我们通过商递归定义 $i:\freegroupx{A}\to \freegroupx{A}$，它在列表 $\ell:\lst{A+A}$ 上的作用是交换 $A$ 的两个副本并将列表反转。
  这一操作保持关系不变，因此可以下降到商。
  并且，我们可以通过归纳证明对于 $x:\freegroupx{A}$ 有 $i(x) \cdot x = e$。
  首先，商归纳允许我们假设 $x$ 来自某个 $\ell:\lst{A+A}$，然后对列表进行归纳；如果将商映射记作 $q:\lst{A+A}\to \freegroupx{A}$，那么归纳情形为
  \begin{align*}
    i(q(\nil)) \ct q(\nil) &= q(\nil) \ct q(\nil)\\
    &= q(\nil)\\
    i(q(\cons(a,\ell))) \ct q(\cons(a,\ell)) &= i(q(\ell)) \ct q(\cons(\hat{a},\nil)) \ct q(\cons(a,\ell))\\
    &= i(q(\ell)) \ct q(\cons(\hat{a},\cons(a,\ell)))\\
    &= i(q(\ell)) \ct q(\ell)\\
    &= q(\nil). \tag{by the inductive hypothesis}
  \end{align*}
  （这里我们省略了关于列表拼接等行为的若干相当明显的引理。）

  这就完成了 $\freegroupx{A}$ 是一个群的证明。
  现在，若 $G$ 是任意一个群且给定函数 $f:A\to G$，我们可以定义 $A+A\to G$ 在 $A$ 的第一个副本上取 $f$，在第二个副本上取 $f$ 与 $G$ 的取逆映射的复合。
  由于 $G$ 是幺半群，我们得到一个幺半群同态 $\lst{A+A} \to G$。
  又因 $G$ 是一个群，这个映射保持上述关系，因此可以下降为映射 $\freegroupx{A}\to G$。
  不难证明这是一个群同态，并且是在 $A$ 上限制为 $f$ 的唯一同态。
\end{proof}

\index{monoid|)}%

如果 $A$ 具有可判定相等性\index{decidable!equality}（例如当我们假设排中律时），那么定义 $\freegroupx{A}$ 的商可以像在 \cref{lem:quotient-when-canonical-representatives} 中那样通过一个幂等元得到。
我们称一个字（回顾一下，它只是 $\lst{A+A}$ 的一个元素）为\define{既约的}
\indexdef{reduced word in a free group}
如果它不包含形如 $(a,\hat a)$ 或 $(\hat a,a)$ 的相邻对。
当 $A$ 具有可判定相等性时，不难定义字的\define{约化}
\index{reduction!of a word in a free group}%
它是一个幂等元，生成相应的商；我们将细节留给读者。

如果 $A\defeq \unit$（具有可判定相等性），则既约字要么全由 $\ttt$ 组成，要么全由 $\hat{\ttt}$ 组成。
因此，$\unit$ 上的自由群（free group）等价于整数 \Z：$0$ 对应 $\nil$，正整数 $n$ 对应由 $n$ 个 $\ttt$ 构成的既约字，负整数 $(-n)$ 对应由 $n$ 个 $\hat{\ttt}$ 构成的既约字。
当然，也可以直接证明 \Z 满足 $\freegroup{\unit}$ 的泛性质。

\begin{rmk}\label{thm:freegroup-nonset}
  在构造 $\freegroup{A}$ 与 $\freegroupx{A}$ 及证明其泛性质的全过程中，我们从未用到 $A$ 是集合这一假设。
  因此，我们实际上可以在任意类型上构造自由群。
  比较泛性质，即可得出 $\eqv{\freegroup{A}}{\freegroup{\trunc0A}}$。
\end{rmk}

\index{group!free|)}%

\index{algebra!colimits of}%
我们也可以用高阶归纳类型来构造代数对象的余极限。
例如，设 $f:G\to H$ 和 $g:G\to K$ 为群同态。
它们在群范畴中的推出称为\define{融合自由积}
\indexdef{amalgamated free product}%
\indexdef{free!product!amalgamated}%
$H *_G K$，它可以构造为由以下条目生成的高阶归纳类型：


\begin{itemize}
\item 函数 $h:H\to H *_G K$ 和 $k:K\to H *_G K$。
\item 群的运算与公理，与 $\freegroup{A}$ 的定义中相同。
\item 断言 $h$ 与 $k$ 为群同态的公理。
\item 对任意 $x:G$，有 $h(f(x)) = k(g(x))$。
\item $0$-截断构造子。
\end{itemize}


另一方面，它也可以显式地构造为 $\lst{H+K}$ 关于以下关系的集合商（set-quotient）：\begin{align*}
  (\dots, x_1, x_2, \dots) &= (\dots, x_1\cdot x_2, \dots)
  & &\text{for $x_1,x_2:H$}\\
  (\dots, y_1, y_2, \dots) &= (\dots, y_1\cdot y_2, \dots)
  & &\text{for $y_1,y_2:K$}\\
  (\dots, 1_G, \dots) &= (\dots, \dots) &&  \\
  (\dots, 1_H, \dots) &= (\dots, \dots) &&  \\
  (\dots, f(x), \dots) &= (\dots, g(x), \dots)
  & &\text{for $x:G$.}
\end{align*}
我们将证明留给读者。
在 $G$ 为平凡群的特殊情形下，最后一条关系不再需要，由此得到\define{自由积}
\indexdef{free!product}%
$H*K$，即群范畴中的余积。
（遗憾的是，此记法与类型的\emph{联结}记法冲突，如 \cref{sec:colimits} 所述，但通常可由上下文消除歧义。）

\index{presentation!of a group}%
注意，由\emph{展示}给出的群可视为余极限的特例。
设给定集合（或更一般地，类型）$A$ 及一对函数 $R\rightrightarrows \freegroup{A}$。
将 $R$ 视为``关系''的类型，其中这对函数为每个关系指定它所设定为相等的两个字。
例如，在展示 $\langle a \mid a^2 = e \rangle$ 中，可取 $A\defeq \unit$、$R\defeq \unit$，两个态射 $R\rightrightarrows \freegroup{A}$ 分别选取列表 $(a,a)$ 与空列表 $\nil$。
于是，由自由群的泛性质，可得一对群同态 $\freegroup{R} \rightrightarrows \freegroup{A}$。
它们在群范畴中的余等子（可仿照推出构造）即为此展示所\emph{展示}的群。

\mentalpause

注意，所有这类构造仅适用于\emph{代数理论}\index{theory!algebraic}，即公理为（全称量化的）方程的理论，方程中涉及的变量、常量与运算均来自给定的\emph{签名}\index{signature!of an algebraic theory}。
对其加以修改后，也可应用于所谓的\emph{本质代数理论}\index{theory!essentially algebraic}：这类理论中的运算为偏定义的，其定义域由先前运算之间的等式所确定。
例如，它们不适用于域的理论：域理论中的``求逆''运算是偏定义的，其定义域 $\setof{x | x \mathrel{\#} 0}$ 由先前运算之间的一个\emph{远离关系} $\#$ 所确定，参见 \cref{RD-inverse-apart-0}。
事实上，域范畴没有始对象，这是众所周知的。
\index{initial!field}%

另一方面，这些构造同样适用于\emph{无穷元}\index{infinitary!algebraic theory}代数理论，其``运算''可接受无穷多个输入。
在此情形下，除非假设选择公理，否则自由代数或代数余极限可能无法表示为简单的商。
这意味着高阶归纳类型显著增强了构造类型论（未必体现在证明论强度上，而是实用效力上），且在某些方面确实比（不带选择公理的）Zermelo--Fraenkel\index{set theory!Zermelo--Fraenkel} 集合论更强~\cite{blass:freealg}。
% We will see an example of this in \cref{sec:ordinals}.

\section{展平引理}
\label{sec:flattening}

正如我们将在\cref{cha:homotopy}中看到的，当高阶归纳类型与泛等性相结合时，会产生令人惊叹的结果。
其主要途径是：若 $W$ 是一个高阶归纳类型且 \UU 是一个类型宇宙，那么我们可以利用 $W$ 的递归原理来定义一个类型族 $P:W\to \UU$。
当涉及 $W$ 递归原理中处理道路构造子的条款时，就需要在 \UU 中提供道路，而这正是泛等性派上用场的地方。

例如，假设我们有一个类型 $X$ 和一个自等价 $e:\eqv X X$。
那么我们可以利用 $\Sn^1$ 递归来定义一个类型族 $P:\Sn^1 \to \UU$：
\begin{equation*}
  P(\base) \defeq X
  \qquad\text{and}\qquad
  \ap P\lloop \defid \ua(e).
\end{equation*}
这样一来，$X$ 便作为 $P$ 在基点 $\base$ 上的纤维 $P(\base)$ 出现。
自等价 $e$ 在 $P$ 中稍显隐蔽，但下面的引理指出，可以通过沿 \lloop 的传输将其提取出来。

\begin{lem}\label{thm:transport-is-given}
  给定 $B:A\to\type$ 和 $x,y:A$，以及一条道路 $p:x=y$ 和一个等价 $e:\eqv{B(x)}{B(y)}$ 使得 $\ap{B}p = \ua(e)$，则对任意 $u:B(x)$ 有
  \begin{align*}
    \transfib{B}{p}{u} &= e(u).
  \end{align*}
\end{lem}

\begin{proof}
  应用\cref{thm:transport-is-ap}，可得
  \begin{align*}
    \transfib{B}{p}{u} &= \idtoeqv(\ap{B}p)(u)\\
    &= \idtoeqv(\ua(e))(u)\\
    &= e(u).\qedhere
  \end{align*}
\end{proof}

我们此前已经见过通过递归定义类型族的例子：在\cref{sec:compute-coprod,sec:compute-nat}中，我们用它刻画了（普通）归纳类型的相等类型。
在\cref{cha:homotopy}中，我们将使用类似的思想来计算高阶归纳类型的同伦群。

在本节中，我们将描述一个关于这类类型族的一般性引理，该引理在后续内容中将很有用。
我们称它为\define{展平引理}：
\indexdef{flattening lemma}%
\indexdef{lemma!flattening}%
它断言，如果 $P:W\to\UU$ 是按上述方式递归定义的，那么它的全空间 $\sm{x:W} P(x)$ 等价于一个``展平了的''高阶归纳类型，后者的构造子可以从 $W$ 的构造子以及 $P$ 的定义推导出来。
（从范畴论的观点看，$\sm{x:W} P(x)$ 是 $P$ 的``Grothendieck（格罗滕迪克）构造\index{Grothendieck construction}''，而展平引理表述了它作为``松弛余极限\index{lax colimit}（lax colimit）''的泛性质。
不过，由于同伦类型论中的类型（比如 $W$）在范畴上对应于 $\infty$-群胚（因为所有道路都是可逆的），所以在这种情况下，松弛余极限与伪余极限是一回事。）

我们在这里证明展平引理的一个通用情况，它可以直接推出许多特例，并提示证明其他情况的方法。
假设有 $A,B:\type$ 和 $f,g:B\to{}A$，并且高阶归纳类型 $W$ 是由以下构造生成的：


\begin{itemize}
\item $\cc:A\to{}W$，以及
\item $\pp:\prd{b:B} (\cc(f(b))=_W\cc(g(b)))$。
\end{itemize}


这样，$W$ 就是 $f$ 和 $g$ 的\define{（同伦）余等子（(homotopy) coequalizer）}。
\indexdef{coequalizer}%
\indexdef{type!coequalizer}%
利用二元和（余积）与依值和（$\Sigma$ 类型），许多有趣的非递归高阶归纳类型都可以用这种形式表示。
所有的点构造子都必须打包进类型 $A$，而所有的道路构造子则打包进类型 $B$。
例如：

\begin{itemize}
\item 圆 $\Sn^1$ 可以通过取 $A\defeq \unit$ 和 $B\defeq \unit$ 来表示，并令 $f$ 和 $g$ 为恒等映射。
\item $j:X\to Y$ 与 $k:X\to Z$ 的推出可以通过取 $A\defeq Y+Z$ 和 $B\defeq X$ 来表示，并令 $f\defeq \inl \circ j$ 以及 $g\defeq \inr\circ k$。
\end{itemize}


现在进一步假设：

\begin{itemize}
\item $C:A\to\type$ 是 $A$ 上的一个类型族，且
\item $D:\prd{b:B}\eqv{C(f(b))}{C(g(b))}$ 是 $B$ 上的一个等价族。
\end{itemize}


递归地定义类型族 $P : W\to\type$ 如下：
\begin{align*}
  P(\cc(a)) &\defeq C(a)\\
  \map{P}{\pp(b)} &\defid \ua(D(b)).
\end{align*}
令 \Wtil 为由以下构造子生成的高阶归纳类型：

\begin{itemize}
\item $\cct:\prd{a:A} C(a) \to \Wtil$ 与
\item $\ppt:\prd{b:B}{y:C(f(b))} (\cct(f(b),y)=_{\Wtil}\cct(g(b),D(b)(y)))$。
\end{itemize}

展平引理如下：

\begin{lem}[展平引理]\label{thm:flattening}
  在上述设定下，我们有
  \[ \eqvspaced{\Parens{\sm{x:W} P(x)}}{\widetilde{W}}. \]
\end{lem}

\index{universal!property!of dependent pair type}%
如上所述，这一等价可视为表达了 $\sm{x:W} P(x)$ 作为 $P$ 在 $W$ 上的``松弛\index{lax colimit}余极限''所具有的泛性质。
同时，它也可视作余极限的\emph{稳定性与下降性}性质的一部分，该性质刻画了高阶意象。%
\index{.infinity1-topos@$(\infty,1)$-topos}%
\index{stability!and descent}%

\cref{thm:flattening} 的证明将占据本节剩余部分。
它略显技术性，初次阅读时可以选择跳过。
但它也是``证明相关数学''的一个很好的范例，
\index{mathematics!proof-relevant}%
因此建议在第二次阅读时仔细研读。

证明思路在于说明 $\sm{x:W} P(x)$ 具有与 \Wtil 相同的泛性质。
我们首先证明它配有与构造子 $\cct$ 和 $\ppt$ 相对应的类似物。

\begin{lem}\label{thm:flattening-cp}
  存在函数
  \begin{itemize}
  \item $\cct':\prd{a:A} C(a) \to \sm{x:W} P(x)$ 与
  \item $\ppt':\prd{b:B}{y:C(f(b))} \Big(\cct'(f(b),y)=_{\sm{w:W}P(w)}\cct'(g(b),D(b)(y))\Big)$。
  \end{itemize}
\end{lem}

\begin{proof}
  第一部分很容易：定义 $\cct'(a,x) \defeq (\cc(a),x)$，由定义即得 $P(\cc(a))\jdeq C(a)$。
  对于第二部分，假设给定 $b:B$ 和 $y:C(f(b))$；我们需要给出一个相等
  \[ (\cc(f(b)),y) = (\cc(g(b)),D(b)(y)). \]
  由于已有 $\pp(b):\cc(f(b))=\cc(g(b))$，由 $\Sigma$-类型中的相等可知，只需给出一个相等 $\trans{\pp(b)}{y} = D(b)(y)$。
  而这可由 \cref{thm:transport-is-given} 得出，其中用到了 $P$ 的定义。
\end{proof}

下面的引理说明，要在 $\sm{w:W} P(w)$ 上定义一个类型族的截面，只需给出与 \Wtil 情形中类似的数据。

\begin{lem}\label{thm:flattening-rect}
  假设 $Q:\big(\sm{x:W} P(x)\big) \to \type$ 是一个类型族，并且我们有
  \begin{itemize}
  \item $c : \prd{a:A}{x:C(a)} Q(\cct'(a,x))$ 以及
  \item $p : \prd{b:B}{y:C(f(b))} \Big(\trans{\ppt'(b,y)}{c(f(b),y)} = c(g(b),D(b)(y))\Big)$。%_{Q(\cct'(g(b),D(b)(y)))}
  \end{itemize}
  那么存在 $k:\prd{z:\sm{w:W} P(w)} Q(z)$ 使得 $k(\cct'(a,x)) \jdeq c(a,x)$ 成立。
\end{lem}

\begin{proof}
  假设给定 $w:W$ 和 $x:P(w)$；我们需要给出一个元素 $k(w,x):Q(w,x)$。
  对 $w$ 作归纳，只需考虑两种情形。
  当 $w\jdeq \cc(a)$ 时，有 $x:C(a)$，于是 $c(a,x):Q(\cc(a),x)$ 即为所求。
  （这部分定义也确保了所述计算规则成立。）

  现在需要证明，对任意 $b:B$，该定义沿 $\pp(b)$ 传输后保持不变。
  由于对所有 $w:W$ 所定义的是一个类型为 $\prd{x:P(w)} Q(w,x)$ 的函数，由 \cref{thm:dpath-forall} 可知，只需证明对任意 $y:C(f(b))$ 有
  \[ \transfib{Q}{\pairpath(\pp(b),\refl{\trans{\pp(b)}{y}})}{c(f(b),y)} = c(g(b),\trans{\pp(b)}{y}). \]
  令 $q:\trans{\pp(b)}{y} = D(b)(y)$ 为由 \cref{thm:transport-is-given} 得到的道路。
  那么有
  \begin{align}
    c(g(b),\trans{\pp(b)}{y})
    &= \transfib{x\mapsto Q(\cct'(g(b),x))}{\opp{q}}{c(g(b),D(b)(y))}
    \tag{by $\opp{\apdfunc{x\mapsto c(g(b),x)}(\opp q)}$} \\
    &= \transfib{Q}{\apfunc{x\mapsto \cct'(g(b),x)}(\opp q)}{c(g(b),D(b)(y))}
    \tag{by \cref{thm:transport-compose}}.
  \end{align}
  因此，只需证明
  \begin{multline*}
    \Transfib{Q}{\pairpath(\pp(b),\refl{\trans{\pp(b)}{y}})}{c(f(b),y)} = {}\\
    \Transfib{Q}{\apfunc{x\mapsto \cct'(g(b),x)}(\opp q)}{c(g(b),D(b)(y))}.
  \end{multline*}
  将右端的传输移到另一侧，并将两个传输合并，上式等价于
  %
  \begin{narrowmultline*}
    \Transfib{Q}{\pairpath(\pp(b),\refl{\trans{\pp(b)}{y}}) \ct
      \apfunc{x\mapsto \cct'(g(b),x)}(q)}{c(f(b),y)} =
    \narrowbreak
    c(g(b),D(b)(y)).
  \end{narrowmultline*}
  %
  然而，有
  \begin{multline*}
    \pairpath(\pp(b),\refl{\trans{\pp(b)}{y}}) \ct \apfunc{x\mapsto \cct'(g(b),x)}(q)
    = {} \\
    \pairpath(\pp(b),\refl{\trans{\pp(b)}{y}}) \ct \pairpath(\refl{\cc(g(b))},q)
    = \pairpath(\pp(b),q)
    = \ppt'(b,y)
  \end{multline*}
  于是构造由类型为 \[ \transfib{Q}{\ppt'(b,y)}{c(f(b),y)} = c(g(b),D(b)(y)). \qedhere \] 的假设 $p(b,y)$ 完成。
\end{proof}

\cref{thm:flattening-rect} \emph{几乎}赋予了 $\sm{w:W}P(w)$ 与 \Wtil 相同的归纳原理。
所缺少的是相等 $\apdfunc{k}(\ppt'(b,y)) = p(b,y)$。
为了证明这一点，需要分析 \cref{thm:flattening-rect} 的证明，而它当然就是 $k$ 的定义。

这应当是可行的，但实际上我们只需要非依值递归原理的计算规则。
因此，我们现在给出递归子的一种稍简化的直接构造，并证明其计算规则。

\begin{lem}\label{thm:flattening-rectnd}
  假设 $Q$ 是一个类型，并且我们有
  \begin{itemize}
  \item $c : \prd{a:A} C(a) \to Q$ 以及
  \item $p : \prd{b:B}{y:C(f(b))} \Big(c(f(b),y) =_Q c(g(b),D(b)(y))\Big)$。
  \end{itemize}
  那么存在 $k:\big(\sm{w:W} P(w)\big) \to Q$ 使得 $k(\cct'(a,x)) \jdeq c(a,x)$ 成立。
\end{lem}

\begin{proof}
  与 \cref{thm:flattening-rect} 中一样，我们对 $w:W$ 作归纳来定义 $k(w,x)$。
  当 $w\jdeq \cc(a)$ 时，定义 $k(\cc(a),x)\defeq c(a,x)$。
  现在由 \cref{thm:dpath-arrow} 可知，只需对 $b:B$ 和 $y:C(f(b))$ 考虑复合道路
  \begin{equation}\label{eq:flattening-rectnd}
    \transfib{x\mapsto Q}{\pp(b)}{c(f(b),y)}
    = c(g(b),\transfib{P}{\pp(b)}{y})
  \end{equation}
  %
  它定义为如下复合
  %
  \begin{align}
    \transfib{x\mapsto Q}{\pp(b)}{c(f(b),y)}
    &= c(f(b),y) \tag{by \cref{thm:trans-trivial}}\\
    &= c(g(b),D(b)(y)) \tag{by $p(b,y)$}\\
    &= c(g(b),\transfib{P}{\pp(b)}{y}). \tag{by \cref{thm:transport-is-given}}
  \end{align}
  如前所述，计算规则 $k(\cct'(a,x)) \jdeq c(a,x)$ 由定义即得。
\end{proof}

要得到第二个计算规则，我们需要下面的引理。

\begin{lem}\label{thm:ap-sigma-rect-path-pair}
  设 $Y:X\to\type$ 是一个类型族，并令 $k:(\sm{x:X}Y(x)) \to Z$ 由 $k(x,y) \defeq d(x)(y)$ 按分量定义，其中 $d:\prd{x:X} Y(x)\to Z$ 是一个柯里化函数。
  那么对任意 $s:\id[X]{x_1}{x_2}$ 以及任意 $y_1:Y(x_1)$ 和 $y_2:Y(x_2)$，若还给定道路 $r:\trans{s}{y_1}=y_2$，则道路
  \[\apfunc k (\pairpath(s,r)) :k(x_1,y_1) = k(x_2,y_2)\]
  等于复合
  \begin{align}
    k(x_1,y_1)
    &\jdeq d(x_1)(y_1) \notag\\
    &= \transfib{x\mapsto Z}{s}{d(x_1)(y_1)}
    \tag{by $\opp{\text{(\cref{thm:trans-trivial})}}$}\\
    &= \transfib{x\mapsto Z}{s}{d(x_1)(\trans{\opp s}{\trans{s}{y_1}})}
    \notag\\
    &= \big(\transfib{x\mapsto (Y(x)\to Z)}{s}{d(x_1)}\big)(\trans{s}{y_1})
    \tag{by~\eqref{eq:transport-arrow}}\\
    &= d(x_2)(\trans{s}{y_1})
    \tag{by $\happly(\apdfunc{d}(s))(\trans{s}{y_1}$}\\
    &= d(x_2)(y_2)
    \tag{by $\apfunc{d(x_2)}(r)$}\\
    &\jdeq k(x_2,y_2).
    \notag
  \end{align}
\end{lem}

\begin{proof}
  对 $s$ 和 $r$ 作道路归纳后，两边都归约为自反性。
\end{proof}

初看之下可能会觉得惊讶：\cref{thm:ap-sigma-rect-path-pair} 的陈述如此复杂，而证明却如此简单。
之所以如此复杂，是为了确保陈述是良类型的：$\apfunc k (\pairpath(s,r))$ 与声称与之相等的复合道路，必须具有相同的起点和终点。
一旦做到这一点，证明便可由道路归纳轻松完成。

\begin{lem}\label{thm:flattening-rectnd-beta-ppt}
  在 \cref{thm:flattening-rectnd} 的情形下，有 $\apfunc{k}(\ppt'(b,y)) = p(b,y)$。
\end{lem}

\begin{proof}
  回顾 $\ppt'(b,y) \defeq \pairpath(\pp(b),q)$，其中 $q:\trans{\pp(b)}{y} = D(b)(y)$ 来自 \cref{thm:transport-is-given}。
  因此，由于 $k$ 是按分量定义的，我们可以由 \cref{thm:ap-sigma-rect-path-pair} 计算 $\apfunc{k}(\ppt'(b,y))$，其中
  \begin{align*}
    x_1 &\defeq \cc(f(b)) & y_1 &\defeq y\\
    x_2 &\defeq \cc(g(b)) & y_2 &\defeq D(b)(y)\\
    s &\defeq \pp(b)      &   r &\defeq q.
  \end{align*}
  柯里化函数 $d:\prd{w:W} P(w) \to Q$ 是对 $w:W$ 作归纳定义的；
  要应用 \cref{thm:ap-sigma-rect-path-pair}，需要理解 $\apfunc{d(x_2)}(r)$ 和 $\happly(\apdfunc{d}(s),\trans s {y_1})$。

  对于第一个，由于 $d(\cc(a),x)\jdeq c(a,x)$，有
  \[ \apfunc{d(x_2)}(r) \jdeq \apfunc{c(g(b),-)}(q). \]
  对于第二个，$W$ 的归纳原理的计算规则告诉我们，$\apdfunc{d}(\pp(b))$ 等于复合道路~\eqref{eq:flattening-rectnd} 经 \cref{thm:dpath-arrow} 的等价转换后所得。
  因此，\cref{thm:dpath-arrow} 中给出的计算规则蕴含 $\happly(\apdfunc{d}(\pp(b)),\trans {\pp(b)}{y})$ 等于复合
  \begin{align}
    \big(\trans{\pp(b)}{c(f(b),-)}\big)(\trans {\pp(b)}{y})
    &= \trans{\pp(b)}{c(f(b),\trans{\opp {\pp(b)}}{\trans {\pp(b)}{y}})}
    \tag{by~\eqref{eq:transport-arrow}}\\
    &= \trans{\pp(b)}{c(f(b),y)}
    \notag \\
    &= c(f(b),y)
    \tag{by \cref{thm:trans-trivial}}\\
    &= c(g(b),D(b)(y))
   \tag{by $p(b,y)$}\\
    &= c(g(b),\trans{\pp(b)}{y}).
    \tag{by $\opp{\apfunc{c(g(b),-)}(q)}$}
  \end{align}
  最后，将 $\apfunc{d(x_2)}(r)$ 和 $\happly(\apdfunc{d}(s),\trans s {y_1})$ 的值代入 \cref{thm:ap-sigma-rect-path-pair}，可见所有道路成对消去，只剩 $p(b,y)$。
\end{proof}

现在我们终于可以证明展平引理了。

\begin{proof}[Proof of \cref{thm:flattening}]
  我们利用 \Wtil 的递归原理，以 $\cct'$ 和 $\ppt'$ 作为输入数据来定义 $h:\Wtil \to \sm{w:W}P(w)$。
  类似地，我们利用 \cref{thm:flattening-rectnd} 的递归原理，以 $\cct$ 和 $\ppt$ 作为输入数据来定义 $k:(\sm{w:W}P(w)) \to \Wtil$。

  一方面，我们需要证明对于任意 $z:\Wtil$，有 $k(h(z))=z$。
  对 $z$ 作归纳，只需考虑 \Wtil 的两个构造子。
  而根据定义，我们有
  \[k(h(\cct(a,x))) \jdeq k(\cct'(a,x)) \jdeq \cct(a,x)\]
  类似地，利用 $\Wtil$ 和 \cref{thm:flattening-rectnd-beta-ppt} 的命题计算规则，可得
  \[\ap k{\ap h{\ppt(b,y)}} = \ap k{\ppt'(b,y)} = \ppt(b,y) \]

  另一方面，我们需要证明对于任意 $z:\sm{w:W}P(w)$，有 $h(k(z))=z$。
  但这本质上是相同的，只需使用 \cref{thm:flattening-rect} 对 $\sm{w:W}P(w)$ 进行``归纳''，并应用相同的计算规则即可。
\end{proof}

\section{高阶归纳定义的一般语法}
\label{sec:naturality}

在 \cref{sec:strictly-positive} 中，我们讨论了使一个假定的``归纳定义''可接受的条件，即类型在其构造子中的所有归纳出现都必须是``严格正的''。\index{strict!positivity}
在本节中，我们讨论\emph{高阶}归纳定义所需的附加条件。
寻找合法高阶归纳定义的一般语法描述，是当前的一个研究课题，迄今为止提出的所有方案本质上都有一定的技术性；因此，我们只给出一般性的描述而非精确的定义。
所幸，在实践中这些极端情况似乎从未出现。

与普通的归纳定义一样，高阶归纳定义由一列\emph{构造子}指定，每个构造子都是一个（依值）函数。
为简化起见，我们可以要求每个构造子的输入满足与普通归纳类型构造子相同的条件。
具体来说，这些输入只能以严格正的方式包含正在被定义的类型。
注意，这就排除了诸如\cref{sec:hittruncations}中所给出的 $0$-截断那样的定义，因为其构造子的输入不仅包含正在定义的归纳类型，还包含了它的恒等类型。
或许可以扩展语法以允许这类定义；不过，在\cref{sec:truncations}中，我们将给出 $0$-截断的另一种构造，其构造子确实满足这一更具限制性的条件。

那么，普通归纳定义与高阶归纳定义之间的唯一区别在于，构造子的\emph{输出}类型可以不是正在定义的类型（比如 $W$），而是它的某个恒等类型，例如 $\id[W]uv$，或者更一般地，是一个迭代恒等类型，例如 $\id[({\id[W]uv})]pq$。
因此，在给出一个高阶归纳定义时，我们不仅需要指定每个构造子的输入，还需要指定那些决定所构造道路之\index{source!of a path constructor}源和\index{target!of a path constructor}靶的表达式 $u$ 和 $v$（或者 $u$, $v$, $p$, $q$ 等等）。

重要的是，这些表达式可以引用 $W$ 的\emph{其他}构造子。
例如，在 $\Sn^1$ 的定义中，构造子 $\lloop$ 的 $u$ 和 $v$ 都是前一个构造子 $\base$。
为了理解这一点，我们要求高阶归纳类型的构造子必须\emph{按顺序}指定，并且允许每个构造子的源和靶表达式 $u$ 和 $v$ 引用之前的构造子，但不能引用之后的构造子。
（当然，在实践中任何归纳定义的构造子都会以某种顺序写出，但对于普通归纳类型而言，该顺序无关紧要。）

注意，这个顺序并不一定是``维度''的顺序：原则上，一个 $1$ 维道路构造子可能引用一个 $2$ 维的构造子，因此需要排在它之后。
然而，我们并未赋予 $0$ 维构造子（点构造子）任何引用先前构造子的方式，所以它们大可以全都放在最前面。
而如果我们使用中心辐条构造（hub-and-spoke construction，见\cref{sec:hubs-spokes}）将所有构造子都归约为点和 $1$-道路，那么我们可以假定所有点构造子最先出现，然后是所有 $1$-道路构造子——但 $1$-道路构造子之间的顺序依然要紧。

剩下的问题是：$u$ 和 $v$ 可以是哪种表达式？
我们可能希望它们可以是涉及先前构造子的任意表达式。
然而，下面的例子表明，对这一想法的朴素处理并不奏效。

\begin{eg}\label{eg:unnatural-hit}
  考虑一个函数族 $f:\prd{X:\type} (X\to X)$。
  当然，$f_X$ 对所有的 $X$ 可能就是 $\idfunc[X]$，但其他的 $f$ 也可能存在。
  例如，$f_{\bool}:\bool\to\bool$ 完全可以是非恒等{\bf 自同构}\index{automorphism!of 2, nonidentity@of $\bool$, nonidentity}（参见 \cref{ex:unnatural-endomorphisms}）。

  现在假设我们尝试定义一个高阶归纳类型 $K$，它由以下生成元生成：
  \begin{itemize}
  \item 两个元素 $a,b:K$，以及
  \item 一条道路 $\sigma:f_K(a)=f_K(b)$。
  \end{itemize}
  $K$ 的归纳原理会是什么样子？
  我们需要假设一个类型族 $P:K\to\type$，当然还需要 $x:P(a)$ 和 $y:P(b)$。
  剩下的数据应当是 $P$ 中的一条依值道路，它位于 $\sigma$ 上方，因此必须将 $P(f_K(a))$ 中的某个元素连接到 $P(f_K(b))$ 中的某个元素。
  但这些元素可能是什么呢？
  我们知道 $P(a)$ 和 $P(b)$ 分别有实例 $x$ 和 $y$，但这丝毫不能告诉我们 $P(f_K(a))$ 和 $P(f_K(b))$ 的情况。
\end{eg}

显然，为了使定义合理，必须对 $u$ 和 $v$ 附加某些条件。
看起来，正如每个构造子的定义域（除其他要求外）必须是一个\emph{协变函子}一样，对表达式 $u$ 和 $v$ 施加的合适条件是它们应当定义\emph{自然变换}。
精确地表述这一要求已超出本书的范围；但非正式地说，它的意思是 $u$ 和 $v$ 只能涉及那些被所有类型间的函数所保持的操作。

例如，$u$ 和 $v$ 中允许涉及道路的拼接，就像 \cref{sec:cell-complexes} 中环面的最后一个构造子那样，因为类型论中的所有函数（在同伦意义下）都保持道路拼接。
但是，不允许它们涉及像 \cref{eg:unnatural-hit} 中函数 $f$ 那样的操作，因为 $f$ 不一定是自然的：可能存在某个函数 $g:X\to Y$，使得 $f_Y \circ g \neq g\circ f_X$。
（泛等性意味着 $f_X$ 必须对所有\emph{等价}保持自然性，但对那些不是等价的函数则不一定。）

自然性的直觉仅为判断一个高阶归纳定义在什么情况下可允许提供了粗略的指南。
即便能在本书中给出此类定义可允许形式的精确规范，这样的规范也很可能迅速过时，因为这一理论的新扩展正在不断被探索。
例如，\cref{sec:circle} 中提到的用``依值 $n$-环路\index{loop!dependent n-@dependent $n$-}''来表示 $n$-球面，以及 \cref{cha:real-numbers} 中使用的``高阶归纳-递归定义''，都是本书撰写期间引入的创新。
我们鼓励读者去实验——但务必谨慎。

\sectionNotes

高阶归纳类型的一般概念诞生于 2011 年 Oberwolfach 会议上 Andrej Bauer、Peter Lumsdaine、Mike Shulman 和 Michael Warren 之间的讨论，尽管早期工作中已有一些特殊情形的线索。
随后，Guillaume Brunerie 和 Dan Licata 对一般理论作出了重要贡献，尤其是找到了在计算机证明助手中表示它们并用其开展同伦论的便捷方法（参见 \cref{cha:homotopy}）。

关于高阶归纳类型的语法及其在高阶范畴模型中的语义的一般讨论，可见于 \cite{ls:hits}。
与普通的归纳类型类似，高阶归纳类型的模型可以通过超限迭代过程来构造；一个通俗的说法是：普通归纳类型描述的是\emph{自由}单子，而高阶归纳类型描述的是单子的\emph{展示}。\index{monad}
引入道路构造子也涉及模型范畴论中``右同伦''（用道路空间定义）与``左同伦''（用柱体定义）之间的等价——这一等价通常只在同伦意义下成立，这为更倾向于对道路构造子采用命题计算规则提供了语义上的理由。

另一个（暂时的）理由来自现有计算机实现的局限。
像 \Coq 和 \Agda 这样的证明助手\index{proof!assistant}内置了普通的归纳类型，但尚未内置高阶归纳类型。
我们当然可以通过假设大量公理来引入它们，但这只能得到命题层面的计算规则。
然而，Dan Licata 有一个技巧，可以利用私有数据类型来实现高阶归纳类型；这样可以获得点构造子的判断规则，但对道路构造子则不行。

\sectionNotes

在 \cref{sec:circle,sec:suspension} 中利用环路空间和纬悬对高维球面所作的类型论描述，主要归功于 Brunerie 和 Licata；Hou 则给出了使用 $n$ 维道路\index{path!n-@$n$-}进行替代描述的类型论版本。
将高阶道路化约为带枢纽和辐条的 1 维道路（\cref{sec:hubs-spokes}）的方法，归功于 Lumsdaine 和 Shulman。
将截断描述为高阶归纳类型的方法源于 Lumsdaine；$(-1)$-截断则与~\cite{ab:bracket-types} 的“方括号类型”密切相关。
展平引理首次以一般形式表述是由 Brunerie 完成的。

\index{set-quotient}
商类型在外延类型论中不成问题，例如在 \NuPRL~\cite{constable+86nuprl-book} 中。
人们通常通过转向扩展的 setoid 系统来引入它们。\index{setoid}
然而，商类型在内涵类型论（同伦类型论的起点）中是一个更棘手的问题，因为我们不能在不规定其行为方式的情况下，简单地添加新的命题相等性。
针对此问题的一些解决方案已被研究过~\cite{hofmann:thesis,Altenkirch1999,altenkirch+07ott}，并且人们已经考虑了数种不同的商类型概念。
使用高阶归纳类型来构造集合商，为我们所采用的特定方法（类似于先前考虑过的某些方法）提供了依据，因为它是一般机制的特例。
我们的构造尚未为所有与商相关的计算问题提供新的解决方案，因为一般而言我们仍然缺乏对高阶归纳类型的良好计算理解——但这确实意味着，当前关于高阶归纳类型计算解释的研究工作也同样适用于商。
当然，基于等价类构造商是标准的集合论思想，也是初等意象理论中熟知的内容；它在类型论中的使用（至少对于命题而言，这依赖于泛等公理）由 Voevodsky（沃沃茨基）提出。
内涵类型论中的商类型蕴含函数外延性这一事实由~\cite{hofmann:thesis} 证明，其灵感来源于~\cite{carboni} 关于正合完备化的工作；\cref{thm:interval-funext} 是对这类论证的改编。

\sectionExercises

\begin{ex}\label{ex:torus}
  定义依值道路的拼接，证明依值函数的应用保持道路拼接，并写出环面 $T^2$ 的精确归纳原理及其计算规则。\index{torus}
\end{ex}

\begin{ex}\label{ex:suspS1}
  利用 \cref{sec:circle} 中基于 $\base$ 和 $\surf$ 给出的 $\Sn^2$ 的显式定义，证明 $\eqv{\susp \Sn^1}{\Sn^2}$。
\end{ex}

\begin{ex}\label{ex:torus-s1-times-s1}
  证明 \cref{sec:cell-complexes} 中定义的环面 $T^2$ 等价于 $\Sn^1\times \Sn^1$。
  （警告：其道路代数相当困难。）
\end{ex}

\begin{ex}\label{ex:nspheres}
  定义依值 $n$-环路\index{loop!dependent n-@dependent $n$-}，以及依值函数在 $n$-环路上的作用，并根据 \cref{sec:circle} 末尾的定义，写出 $n$ 维球面的归纳原理。
\end{ex}

\begin{ex}\label{ex:susp-spheres-equiv}
  利用 \cref{sec:circle} 中使用 $\Omega^n$ 给出的 $\Sn^n$ 的定义，证明 $\eqv{\susp \Sn^n}{\Sn^{n+1}}$。
\end{ex}

\begin{ex}\label{ex:spheres-make-U-not-2-type}
  证明如果类型 $\Sn^2$ 属于某个宇宙 \type，则 \type 不是 2-类型。
\end{ex}

\begin{ex}\label{ex:monoid-eq-prop}
  证明如果 $G$ 是一个幺半群且 $x:G$，则 $\sm{y:G}((x\cdot y = e) \times (y\cdot x =e))$ 是一个命题。
  由此，利用唯一选择原理（\cref{cor:UC}），可以得出如下等价定义：群是满足对于每个 $x:G$，纯存在一个 $y:G$ 使得 $x\cdot y = e$ 且 $y\cdot x=e$ 的幺半群。
\end{ex}

\begin{ex}\label{ex:free-monoid}
  证明如果 $A$ 是一个集合，则 $\lst A$ 是一个幺半群。
  然后完成 \cref{thm:free-monoid} 的证明。\index{monoid!free}
\end{ex}

\begin{ex}\label{ex:unnatural-endomorphisms}
  假设 \LEM{}，构造一个函数族 $f:\prd{X:\type}(X\to X)$，使得 $f_{\bool}:\bool\to\bool$ 是非恒等自同构。\index{automorphism!of 2, nonidentity@of $\bool$, nonidentity}
\end{ex}

\begin{ex}\label{ex:funext-from-interval}
  证明在 \cref{thm:interval-funext} 中构造的映射实际上是 $\happly$ 的一个拟逆，从而表明区间类型蕴含完整的函数外延性公理。
  （你可能需要使用 \cref{ex:strong-from-weak-funext}。）
\end{ex}

\begin{ex}\label{ex:susp-lump}
  证明纬悬的泛性质：
  \[ \Parens{\susp A \to B} \eqvsym \Parens{\sm{b_n : B} \sm{b_s : B} (A \to (b_n = b_s)) } \]
\end{ex}

\begin{ex}\label{ex:alt-integers}
  证明 $\eqv{\Z}{\N+\unit+\N}$。
  证明如果我们把 $\Z$ 定义为 $\N+\unit+\N$，那么我们能够以判断式计算规则得到 \cref{thm:sign-induction}。
\end{ex}

\begin{ex}\label{ex:trunc-bool-interval}
  证明我们也可以通过使用 $\brck \bool$ 而非 $\interval$ 来证明 \cref{thm:interval-funext}。
\end{ex}

\index{type!higher inductive|)}%

% Local Variables:
% TeX-master: "hott-online"
% End:

\chapter{同伦 \texorpdfstring{$n$}{n}类型}
\label{cha:hlevels}

\index{n-type@$n$-type|(}%
\indexsee{h-level}{$n$-type}

同伦论的一个基本概念是\emph{同伦 $n$-类型（homotopy $n$-type）}：一个在维度 $n$ 以上不含任何有趣同伦的空间。
例如，同伦 $0$-类型本质上就是集合，不包含非平凡的道路；而同伦 $1$-类型可能包含非平凡的道路，但不包含道路之间的非平凡道路。
同伦 $n$-类型也称为\emph{$n$-截断空间（$n$-truncated spaces）}。
我们已经在 \cref{sec:basics-sets} 中提到过这个概念；本章的第一个目标就是在同伦类型论中给出它的精确定义。

与截断性对偶的概念是连通性：如果一个空间在维度 $n$ 及\emph{以下}不含任何有趣的同伦，则称其为\emph{$n$-连通的（$n$-connected）}。
例如，一个空间是 $0$-连通的（也简称为"连通的"），当且仅当它只有一个连通分支；是 $1$-连通的（也称为"单连通的"），当且仅当它也没有非平凡环路（尽管它可能有环路之间的非平凡高阶环路\index{loop!n-@$n$-}）。

将这两个概念都扩展到映射上，便能最清楚地看到截断性与连通性之间的对偶关系。
如果映射的所有纤维都是 $n$-截断的（或 $n$-连通的），我们就称该映射是\emph{$n$-截断的}或\emph{$n$-连通的}。
那么，$n$-连通映射和 $n$-截断映射就构成了一个\emph{正交分解系统}
\index{orthogonal factorization system}
\indexsee{factorization!system, orthogonal}{orthogonal factorization system}
中的两个映射类，也就是说，每个映射都可以唯一地分解为一个 $n$-连通映射后接一个 $n$-截断映射。

在 $n={-1}$ 的情形，$n$-截断映射就是嵌入，而 $n$-连通映射就是满射，正如在 \cref{sec:mono-surj} 中所定义的那样。
因此，$n$-连通分解系统是集合间函数标准像分解（即分解为满射后接单射）的大规模推广。
在本章末尾，我们会简要概述一个更为一般的理论：任何类型论的\emph{模态}都会引出一个类似的分解系统。

\section{\texorpdfstring{$n$}{n}类型的定义}
\label{sec:n-types}

正如在 \cref{sec:basics-sets,sec:contractibility} 中提到的，从零以下两个层级开始定义 $n$-类型是比较方便的：$(-1)$-类型就是命题，$(-2)$-类型则是可缩的类型。

\begin{defn}\label{def:hlevel}
  对于 $n \geq -2$，我们通过递归定义谓词 $\istype{n} : \type \to \type$ 如下：
  \[ \istype{n}(X) \defeq
  \begin{cases}
    \iscontr(X) & \text{ if } n = -2, \\
    \prd{x,y : X} \istype{n'}(\id[X]{x}{y}) & \text{ if } n = n'+1.
  \end{cases}
  \]
  如果 $\istype{n}(X)$ 有实例，我们就说 $X$ 是一个\define{$n$-类型（$n$-type）}，有时也说它是\emph{$n$-截断的}，
  \indexdef{n-type@$n$-type}%
  \indexsee{n-truncated@$n$-truncated!type}{$n$-type}%
  \indexsee{type!n-type@$n$-type}{$n$-type}%
  \indexsee{type!n-truncated@$n$-truncated}{$n$-type}。%
\end{defn}

\begin{rmk}
  \cref{def:hlevel} 中的数字 $n$ 取遍所有大于等于 $-2$ 的整数。
  我们可以形式化地理解这一点：可以定义一个这样的整数的类型 $\Z_{{\geq}-2}$（其归纳原理与 $\nat$ 相同），或者对每个 $k : \nat$ 定义一个谓词 $\istype{(k-2)}$。
  无论哪种方式，我们都可以通过对 $n$ 进行归纳来证明关于 $n$-类型的定理，其中 $n = -2$ 是基础情形。
\end{rmk}

\begin{eg}
  \index{set}
  我们在 \cref{thm:prop-minusonetype} 中看到，$X$ 是 $(-1)$-类型当且仅当它是一个命题。
  因此，$X$ 是 $0$-类型当且仅当它是一个集合。
\end{eg}

我们也已经看到存在不是集合的类型（\cref{thm:type-is-not-a-set}）。
然而，到目前为止，我们还没有对任何 $n>0$ 证明存在不是 $n$-类型的类型。
不过在 \cref{cha:homotopy} 中，我们将证明 $(n+1)$-球面 $\Sn^{n+1}$ 不是 $n$-类型。
（Kraus（克劳斯）也证明了第 $n$ 层嵌套的泛等宇宙也不是 $n$-类型，且未使用任何高阶归纳类型。）
此外，在 \cref{sec:whitehead} 中，我们将给出一个对\emph{任何}（有限）数字 $n$ 都不是 $n$-类型的例子。

我们先证明 $n$-类型在某些运算和构造子下封闭，以此开始 $n$-类型的一般理论。

\begin{thm}\label{thm:h-level-retracts}
  \index{retract!of a type}%
  \index{retraction}%
 设 $p : X \to Y$ 是一个收缩。对任意 $n\geq -2$，若 $X$ 是 $n$-类型，则 $Y$ 也是 $n$-类型。
\end{thm}

\begin{proof}
 我们对 $n$ 进行归纳。
 基础情形 $n=-2$ 由 \cref{thm:retract-contr} 处理。

 对于归纳步骤，假设任何 $n$-类型的收缩核仍是 $n$-类型，并且 $X$ 是一个 $\nplusone$-类型。
 任取 $y, y' : Y$；我们需要证明 $\id{y}{y'}$ 是一个 $n$-类型。
 令 $s$ 为 $p$ 的一个截面，并令 $\epsilon$ 为一个同伦 $\epsilon : p \circ s \htpy 1$。
 由于 $X$ 是一个 $\nplusone$-类型，$\id[X]{s(y)}{s(y')}$ 是一个 $n$-类型。
 我们断言 $\id{y}{y'}$ 是 $\id[X]{s(y)}{s(y')}$ 的一个收缩核。
 对于截面，我们取
 \[ \apfunc s : (y=y') \to (s(y)=s(y')). \]
 对于收缩，我们定义 $t:(s(y)=s(y'))\to(y=y')$ 为
 \[ t(q) \defeq  \opp{\epsilon_y} \ct \ap p q \ct \epsilon_{y'}.\]
 为了证明 $t$ 是 $\apfunc s$ 的一个收缩，我们需要证明对于任意 $r:y=y'$，有
 \[ \opp{\epsilon_y} \ct \ap p {\ap sr} \ct \epsilon_{y'} = r \]
 但这由 \cref{lem:htpy-natural} 可得。
\end{proof}

作为直接推论，我们得到 $n$-类型在等价下的稳定性（这从泛等性也能直接推出）：

\begin{cor}\label{cor:preservation-hlevels-weq}
 如果 $\eqv{X}{Y}$ 且 $X$ 是一个 $n$-类型，那么 $Y$ 也是。
\end{cor}

此外，回忆 \cref{sec:mono-surj} 中引入的嵌入概念。

\begin{thm}\label{thm:isntype-mono}
  \index{function!embedding}
  若 $f:X\to Y$ 是一个嵌入，且对于某个 $n\ge -1$，$Y$ 是一个 $n$-类型，那么 $X$ 也是。
\end{thm}

\begin{proof}
  任取 $x,x':X$；我们需要证明 $\id[X]{x}{x'}$ 是一个 $\nminusone$-类型。
  但由于 $f$ 是嵌入，我们有 $(\id[X]{x}{x'}) \eqvsym (\id[Y]{f(x)}{f(x')})$，而由假设后者是一个 $\nminusone$-类型。
\end{proof}

注意，这一定理在 $n=-2$ 时不成立：映射 $\emptyt \to \unit$ 是一个嵌入，但 $\unit$ 是一个 $(-2)$-类型而 $\emptyt$ 不是。

\begin{thm}\label{thm:hlevel-cumulative}
 $n$-类型的层级在如下意义下是累积的：
   给定一个数 $n \geq -2$，如果 $X$ 是一个 $n$-类型，那么它也是一个 $\nplusone$-类型。
\end{thm}

\begin{proof}
 我们对 $n$ 进行归纳。

 当 $n = -2$ 时，我们需要证明一个可缩类型，比如 $A$，具有可缩的道路空间。
       设 $a_0: A$ 是 $A$ 的收缩中心，并设 $x, y : A$。我们证明 $\id[A]{x}{y}$
       是可缩的。
       由 $A$ 的可缩性，我们有一条道路 $\contr_x \ct \opp{\contr_y} : x = y$，我们选择它作为
       $\id{x}{y}$ 的收缩中心。
       给定任意 $p : x = y$，我们需要证明 $p = \contr_x \ct \opp{\contr_y}$。
           根据道路归纳，只需证明
       $\refl{x} = \contr_x \ct \opp{\contr_x}$，而这是平凡的。

 对于归纳步骤，我们需要证明：若 $X$ 是 $\nplusone$-类型，则
           $\id[X]{x}{y}$ 是 $\nplusone$-类型。将归纳假设应用于 $\id[X]{x}{y}$
           即得所需结论。
\end{proof}

% \section{Preservation under constructors}
% \label{sec:ntype-pres}

我们现在证明，大多数类型形成运算都保持 $n$-类型。

\begin{thm}\label{thm:ntypes-sigma}
 设 $n \geq -2$，且令 $A : \type$ 与 $B : A \to \type$。
 如果 $A$ 是一个 $n$-类型，并且对于所有 $a : A$，$B(a)$ 都是 $n$-类型，那么
 $\sm{x : A} B(x)$ 也是 $n$-类型。
\end{thm}

\begin{proof}
 我们对 $n$ 进行归纳。

 当 $n = -2$ 时，我们选择 $\sm{x : A} B(x)$ 的收缩中心为偶对
       $(a_0, b_0)$，其中 $a_0 : A$ 是 $A$ 的收缩中心，而 $b_0 : B(a_0)$ 是 $B(a_0)$ 的收缩中心。
       给定 $\sm{x : A} B(x)$ 中任意其他元素 $(a,b)$，我们分别利用 $A$ 和 $B(a_0)$ 的可缩性
       提供一条道路 $\id{(a, b)}{(a_0,b_0)}$。

 对于归纳步骤，假设 $A$ 是 $\nplusone$-类型，并且
         对于任意 $a : A$，$B(a)$ 是 $\nplusone$-类型。我们证明 $\sm{x : A} B(x)$ 是 $\nplusone$-类型：
      固定 $\sm{x : A} B(x)$ 中的 $(a_1, b_1)$ 和 $(a_2,b_2)$，
     我们证明 $\id{(a_1, b_1)}{(a_2,b_2)}$ 是一个 $n$-类型。
      根据 \cref{thm:path-sigma}，我们有
      \[ \eqvspaced{(\id{(a_1, b_1)}{(a_2,b_2)})}{\sm{p : \id{a_1}{a_2}} (\id[B(a_2)]{\trans{p}{b_1}}{b_2})} \]
   并且由 $n$-类型在等价下的保持性（\cref{cor:preservation-hlevels-weq}），
   只需证明后者是 $n$-类型。这由归纳假设可得。
\end{proof}

作为一个特例，如果 $A$ 和 $B$ 是 $n$-类型，那么 $A\times B$ 也是。
另请注意，\cref{thm:hlevel-cumulative} 蕴含：如果 $A$ 是 $n$-类型，那么对于任意 $x,y:A$，$\id[A]xy$ 也是 $n$-类型。
将此与 \cref{thm:ntypes-sigma} 结合，我们看到对于任意 $n$-类型之间的函数 $f:A\to C$ 和 $g:B\to C$，它们的拉回\index{pullback}
\[ A\times_C B \defeq \sm{x:A}{y:B} (f(x)=g(y)) \]
（参见 \cref{ex:pullback}）也是 $n$-类型。
更一般地，$n$-类型在所有 \emph{极限} 下封闭。

\begin{thm}\label{thm:hlevel-prod}
 设 $n\geq -2$，且令 $A : \type$ 与 $B : A \to \type$。
 如果对于所有 $a : A$，$B(a)$ 是 $n$-类型，那么 $\prd{x : A} B(x)$ 也是 $n$-类型。
\end{thm}

\begin{proof}
  我们对 $n$ 进行归纳。
  当 $n = -2$ 时，结果直接就是 \cref{thm:contr-forall}。

  对于归纳步骤，假设结论对 $n$-类型成立，并且每个 $B(a)$ 是 $\nplusone$-类型。
  令 $f, g : \prd{a:A}B(a)$。
  我们需要证明 $\id{f}{g}$ 是 $n$-类型。
  根据函数外延性以及 $n$-类型在等价下的封闭性，只需证明 $\prd{a : A} (\id[B(a)]{f(a)}{g(a)})$ 是 $n$-类型。
  这由归纳假设可得。
\end{proof}

作为上述定理的一个特例，如果 $B$ 是一个 $n$-类型，那么函数空间 $A \to B$ 也是 $n$-类型。
现在我们可以推广 \cref{cha:basics} 中的观察，即 $\isset(A)$ 和 $\isprop(A)$ 是命题：

\begin{thm}\label{thm:isaprop-isofhlevel}
 对于任意 $n \geq -2$ 和任意类型 $X$，类型 $\istype{n}(X)$ 都是一个命题。
\end{thm}

\begin{proof}
  我们对 $n$ 进行归纳。

 对于基情形，需要证明对于任意 $X$，类型 $\iscontr(X)$ 是一个命题。
 这正是 \cref{thm:isprop-iscontr}。

对于归纳步骤，我们需要证明
\[\prd{X : \type} \isprop (\istype{n}(X)) \to \prd{X : \type} \isprop (\istype{\nplusone}(X)). \]
为了证明此蕴含关系的结论，我们需要证明：对于任意类型 $X$，类型
\[\prd{x, x' : X}\istype{n}(x = x')\]
是一个命题。根据 \cref{thm:isprop-forall} 或 \cref{thm:hlevel-prod}，只需证明对于任意 $x, x' : X$，类型 $\istype{n}(x =_X x')$ 是一个命题。
而这由对类型 $(x =_X x')$ 施加归纳假设即可得出。
\end{proof}

最后，我们证明 $n$-类型构成的类型本身是一个 $\nplusone$-类型。
我们将其定义为：
\symlabel{universe-of-ntypes}
\[\ntype{n} \defeq \sm{X : \type} \istype{n}(X). \]
必要时，我们可以写作 $\ntypeU{n}$ 以指定宇宙 $\UU$。
特别地，如 \cref{cha:basics} 中所定义，我们有 $\prop \defeq \ntype{(-1)}$ 与 $\set \defeq \ntype{0}$。
需要注意的是，正如对于 \prop 和 \set 一样，由于 $\istype{n}(X)$ 是一个命题，根据 \cref{thm:path-subset}，对于任何 $(X,p), (X',p'):\ntype{n}$ 我们有
\begin{align*}
  \Big(\id[\ntype{n}]{(X, p)}{(X', p')}\Big) &\eqvsym (\id[\type] X X')\\
  &\eqvsym (\eqv{X}{X'}).
\end{align*}

\begin{thm}\label{thm:hleveln-of-hlevelSn}
 对于任意 $n \geq -2$，类型 $\ntype{n}$ 是一个 $\nplusone$-类型。
\end{thm}

\begin{proof}%[Proof of \cref{thm:hleveln-of-hlevelSn}]
  令 $(X, p), (X', p') : \ntype{n}$；我们需要证明 $\id{(X, p)}{(X', p')}$ 是一个 $n$-类型。
  根据上述观察，这个类型等价于 $\eqv{X}{X'}$。
  接下来，我们注意到投影
  \[(\eqv{X}{X'}) \to (X \rightarrow X').\]
  是一个嵌入，因此如果 $n\geq -1$，那么由 \cref{thm:isntype-mono}，只需证明 $X \rightarrow X'$ 是一个 $n$-类型。
  但由于 $n$-类型在箭头类型下是封闭的，这就化归到 $X'$ 是 $n$-类型这一假设。

  在 $n=-2$ 的情形下，上述论证表明 $\eqv{X}{X'}$ 是一个 $(-1)$-类型——但它也是有实例的，因为任何两个可缩的类型都等价于 \unit，从而彼此等价。
  因此，$\eqv{X}{X'}$ 也是一个 $(-2)$-类型。
\end{proof}

\section{恒等证明唯一性与 Hedberg 定理}
\label{sec:hedberg}

\index{set|(}%

在 \cref{sec:basics-sets} 中，我们定义了什么是\emph{集合（set）}：如果对于任意 $x, y : X$ 以及 $p, q : x =_X y$ 都有 $p = q$，则称类型 $X$ 是一个集合。
在传统类型理论中，这一性质被称为\define{恒等证明唯一性（uniqueness of identity proofs, UIP）}。
\indexdef{uniqueness!of identity proofs}%
我们也已经看到，这等价于它是上一节意义下的 $0$-类型。
下面给出另一种等价的刻画，涉及 Streicher 的``K 公理'' \cite{Streicher93}：

\begin{thm}\label{thm:h-set-uip-K}
 一个类型 $X$ 是集合，当且仅当它满足\define{K 公理（Axiom K）}：
 \indexdef{axiom!Streicher's Axiom K}%
 即，对于所有 $x : X$ 和 $p : (x =_A x)$，都有 $p = \refl{x}$。
\end{thm}

\begin{proof}
 显然，K 公理是 UIP 的一个特例。
 反之，若 $X$ 满足 K 公理，任取 $x, y : X$ 和 $p, q : (\id{x}{y})$；我们要证明 $p=q$。
 但对 $q$ 进行归纳恰好将这一目标归约为 K 公理。
\end{proof}

我们要强调，\emph{我们}并没有将 UIP 或 K 原则作为公理来假设！
它们只是某些类型可能具有、也可能不具有的性质（即等价于该类型为集合）。
回顾 \cref{thm:type-is-not-a-set}，\emph{并非}所有类型都是集合。

下面的定理提供了另一种证明类型是集合的有用方法。

\begin{thm}\label{thm:h-set-refrel-in-paths-sets}
  \index{relation!reflexive}%
  设 $R$ 是类型 $X$ 上的一个自反的\index{reflexivity!of a relation}纯粹关系，且蕴含恒等。
  那么 $X$ 是一个集合，并且对于所有 $x,y:X$，$R(x,y)$ 都等价于 $\id[X]{x}{y}$。
\end{thm}

\begin{proof}
  令 $\rho : \prd{x:X} R(x,x)$ 见证 $R$ 的自反性，并令 \narrowequation{f : \prd{x,y:X} R(x,y) \to (\id[X]{x}{y})} 见证 $R$ 蕴含恒等。
  我们首先注意到，定理中的两个陈述是等价的。
  因为一方面，如果 $X$ 是一个集合，那么 $\id[X]xy$ 就是一个命题，而根据假设，它与命题 $R(x,y)$ 逻辑等价，因此也必然与 $R(x,y)$ 等价。
  另一方面，如果 $\id[X]xy$ 等价于 $R(x,y)$，那么与后者一样，它对所有 $x,y:X$ 都是命题，因此 $X$ 是一个集合。

  我们给出这个定理的两种证明。
  第一种直接证明 $X$ 是集合；第二种直接证明 $R(x,y)\eqvsym (x=y)$。

  \emph{第一种证明：}我们证明 $X$ 是集合。
  思路与 \cref{thm:prop-set} 相同：函数 $f$ 对其参数 $x$ 和 $y$ 必须是"连续的"。
  不过，由于我们需要处理类型为 $R(x,y)$ 的额外参数，符号上会稍微复杂一些。

  首先，对任意 $x:X$ 和 $p:\id[X]xx$，考虑 $\apdfunc{f(x)}(p)$。
  这是从 $f(x,x)$ 到自身的一条依赖道路。
  由于 $f(x,x)$ 仍然是一个函数 $R(x,x) \to (\id[X]xx)$，根据 \cref{thm:dpath-arrow}，对任意 $r:R(x,x)$ 我们得到一条道路
  \[\trans{p}{f(x,x,r)} = f(x,x,\trans{p}r).
  \]
  在等式的左边，是恒等类型中的传输，也就是拼接。
  而在右边，由于 $r$ 和 $\trans{p}r$ 都位于命题 $R(x,x)$ 中，我们有 $\trans{p}r = r$。
  因此，将 $r$ 替换为 $\rho(x)$，我们得到
  \[ f(x,x,\rho(x)) \ct p = f(x,x,\rho(x)). \]
  通过消去，得到 $p=\refl{x}$。
  所以 $X$ 满足 K 公理，从而是集合。

  \emph{第二种证明：}我们证明每个 $f(x,y) : R(x,y) \to \id[X]{x}{y}$ 都是一个等价。
  根据 \cref{thm:total-fiber-equiv}，只需证明 $f$ 诱导了一个全空间之间的等价：
  \begin{equation*}
    \eqv{\Parens{\sm{y:X}R(x,y)}}{\Parens{\sm{y:X}\id[X]{x}{y}}}.
  \end{equation*}
  由 \cref{thm:contr-paths}，右边的类型是可缩的，因此我们只需证明左边的类型也是可缩的。
  我们取偶对 $\pairr{x,\rho(x)}$ 作为收缩的中心。
  余下只需证明，对每个 ${y:X}$ 和每个 ${H:R(x,y)}$，都有
  \begin{equation*}
    \id{\pairr{x,\rho(x)}}{\pairr{y,H}}.
  \end{equation*}
  但由于 $R(x,y)$ 是一个命题，根据 \cref{thm:path-sigma}，我们只需证明 $\id[X]{x}{y}$，而这可以由 $f(H)$ 给出。
\end{proof}

\begin{cor}\label{notnotstable-equality-to-set}
  如果一个类型 $X$ 满足以下性质：对任意 $x,y:X$，$\neg\neg(x=y)\to(x=y)$，那么 $X$ 是一个集合。
\end{cor}

另一种证明类型是集合的便捷方法如下。
回顾 \cref{sec:intuitionism}，一个类型 $X$ 被称为具有\emph{可判定相等性（decidable equality）}
\index{decidable!equality|(}%
如果对所有 $x, y : X$，我们有
\[(x =_X y) + \neg (x =_X y).\]
\index{continuity of functions in type theory@``continuity'' of functions in type theory}%
\index{functoriality of functions in type theory@``functoriality'' of functions in type theory}%
这是一个非常强的条件：它意味着，当道路 $x=y$ 存在时，可以关于 $x$ 和 $y$ "连续地"（或"可计算地"或"函子性地"）选取该道路。
结合 \cref{thm:h-set-refrel-in-paths-sets} 与下面的引理，这蕴含 $X$ 是一个集合。

\begin{lem}\label{lem:hedberg-helper}
对于任意类型 $A$，我们有 $(A+\neg A)\to(\neg\neg A\to A)$。
\end{lem}

\begin{proof}
这在 \cref{thm:not-lem} 中已经基本证明过了，但我们再重复一下论证过程。
假设 $x:A+\neg A$。我们需要考虑两种情况。
如果 $x$ 是某个 $a:A$ 的左注入 $\inl(a)$，那么我们有一个常值函数 $\neg\neg A
\to A$，它将所有元素都映射到 $a$。如果 $x$ 是某个 $t:\neg A$ 的右注入 $\inr(t)$，
那么对于每个 $g:\neg\neg A$，都有 $g(t):\emptyt$。因此我们可以使用
\emph{ex falso quodlibet}，即 $\rec{\emptyt}$，来为任意 $g:\neg\neg A$ 获得一个 $A$ 中的元素。
\end{proof}

\index{anger}


\begin{thm}[Hedberg]\label{thm:hedberg}
  \index{Hedberg's theorem}%
  \index{theorem!Hedberg's}%
  如果 $X$ 有可判定相等性，那么 $X$ 是一个集合。
\end{thm}

\begin{proof}
如果 $X$ 有可判定相等性，那么对于任意 $x,y:X$，$\neg\neg(x=y)\to(x=y)$ 成立。因此，Hedberg 定理可由
\cref{notnotstable-equality-to-set} 推出。
\end{proof}

当然，这个定理与 \cref{thm:not-lem} 有密切的联系。
被 \cref{thm:not-lem} 所否定的命题 \LEM{\infty}，显然蕴含了每个类型都具有可判定相等性，从而都是集合——然而我们知道实际情况并非如此。
\index{excluded middle}%
请注意，从 \cref{sec:intuitionism} 引入的一致公理 \LEM{}，仅蕴含每个类型具有\emph{命题意义下的可判定相等性}，即对于任意 $A$，我们有
\indexdef{equality!merely decidable}%
\indexdef{merely!decidable equality}%
\[ \prd{a,b:A} (\brck{a=b} + \neg\brck{a=b}). \]

\index{decidable!equality|)}%

作为 \cref{thm:hedberg} 的一个应用示例，回顾 \cref{thm:nat-set}，我们曾利用 \cref{sec:compute-nat} 中对自然数相等类型的刻画，观察到 $\nat$ 是一个集合。
一个更传统的证明该定理的方法，仅使用 \eqref{eq:zero-not-succ} 和 \eqref{eq:suc-injective}，而不是 \cref{thm:path-nat} 的完整刻画，并用 \cref{thm:hedberg} 来填补空白。

\begin{thm}\label{prop:nat-is-set}
 自然数类型 $\nat$ 具有可判定相等性，因此是一个集合。
\end{thm}

\begin{proof}
  给定 $x, y : \nat$，我们通过对 $x$ 作归纳、对 $y$ 作分情形讨论来证明 $(x=y)+\neg(x=y)$。
  若 $x \jdeq 0$ 且 $y \jdeq 0$，我们取 $\inl(\refl0)$。
  若 $x \jdeq 0$ 且 $y \jdeq \suc(n)$，则由 \eqref{eq:zero-not-succ} 可得 $\neg (0 = \suc (n))$。

  对于归纳步骤，设 $x \jdeq \suc (n)$。
  若 $y \jdeq 0$，我们再次使用 \eqref{eq:zero-not-succ}。
  最后，若 $y \jdeq \suc (m)$，归纳假设给出 $(m = n)+\neg(m = n)$。
  在第一种情形中，若 $p:m=n$，则 $\ap \suc p:\suc(m)=\suc(n)$。
  而在第二种情形中，\eqref{eq:suc-injective} 推出 $\neg(\suc(m)=\suc(n))$。
\end{proof}

\index{set|)}%

\index{axiom!Streicher's Axiom K!generalization to n-types@generalization to $n$-types}%
尽管 Hedberg 定理似乎专属于集合（$0$-类型），但 K 公理可以自然地推广到 $n$-类型。
注意，普通的 K 公理（作为类型 $X$ 的一个性质）断言：对所有 $x:X$，环路空间 $\Omega(X,x)$（参见 \cref{def:loopspace}）是可缩的。
由于 $\Omega(X,x)$ 总是有实例的（由 $\refl{x}$ 给出），这等价于它是一个命题（一个 $(-1)$-类型）。
因为 $0 = (-1)+1$，这启发我们如下推广。

\begin{thm}\label{thm:hlevel-loops}
  对于任意 $n\geq -1$，类型 $X$ 是一个 $\nplusone$-类型，当且仅当对所有 $x : X$，类型 $\Omega(X, x)$ 是一个 $n$-类型。
\end{thm}

在证明此定理前，我们先证明一个辅助引理：

\begin{lem}\label{lem:hlevel-if-inhab-hlevel}
  给定 $n \geq -1$ 和 $X : \type$。
  如果已知 $X$ 有实例便能推出 $X$ 是 $n$-类型，那么 $X$ 本身就是 $n$-类型。
\end{lem}

\begin{proof}
  令 $f : X \to \istype{n}(X)$ 为给定的映射。
  我们需要证明，对任意 $x, x' : X$，类型 $\id{x}{x'}$ 是一个 $\nminusone$-类型。
  而此时 $f(x)$ 表明 $X$ 是一个 $n$-类型，因而其所有道路空间都是 $\nminusone$-类型。
\end{proof}

\begin{proof}[Proof of \cref{thm:hlevel-loops}]
  “仅当”方向是显然的，因为 $\Omega(X,x)\defeq (\id[X]xx)$。
  反过来，为证 $X$ 是一个 $\nplusone$-类型，我们需要证明：对任意 $x, x' : X$，类型 $\id{x}{x'}$ 是一个 $n$-类型。
  根据 \cref{lem:hlevel-if-inhab-hlevel}，只需构造一个映射
  \[ (\id{x}{x'}) \to \istype{n}(\id{x}{x'}). \]
  由道路归纳，只需考虑 $x\jdeq x'$ 的情形，此时由 $\Omega(X, x)$ 是 $n$-类型的假设即得证。
\end{proof}

\index{whiskering}
通过归纳和一些略显巧妙的须接（whiskering）操作，我们可以将 K 性质推广到 $n>0$ 的情形。

\begin{thm}\label{thm:ntype-nloop}
  \index{loop space!iterated}%
  对于每个 $n\ge -1$，类型 $A$ 是 $n$-类型当且仅当对所有的 $a:A$，$\Omega^{n+1}(A,a)$ 都是可缩的。
\end{thm}

\begin{proof}
  回顾 $\Omega^0(A,a) = (A,a)$，$n=-1$ 的情形即为 \cref{ex:prop-inhabcontr}。
  $n=0$ 的情形是 \cref{thm:h-set-uip-K}。
  现在我们使用归纳法；假设命题对 $n:\N$ 成立。
  根据 \cref{thm:hlevel-loops}，$A$ 是 $(n+1)$-类型当且仅当对所有 $a:A$，$\Omega(A,a)$ 都是 $n$-类型。
  由归纳假设，后者等价于：对所有 $p:\Omega(A,a)$，$\Omega^{n+1}(\Omega(A,a),p)$ 都是可缩的。

  由于 $\Omega^{n+2}(A,a) \defeq \Omega^{n+1}(\Omega(A,a),\refl{a})$ 且 $\Omega^{n+1} = \Omega^n \circ \Omega$，只需证明在带点类型 $\pointed\type$ 中 $\Omega(\Omega(A,a),p)$ 与 $\Omega(\Omega(A,a),\refl{a})$ 相等。
  为此，只需给出一个等价
  \[ g : \Omega(\Omega(A,a),p) \eqvsym \Omega(\Omega(A,a),\refl{a}) \]
  使其将基点 $\refl{p}$ 映射到基点 $\refl{\refl{a}}$。
  对于 $q:p=p$，定义 $g(q):\refl{a} = \refl{a}$ 为如下复合：
  \[ \refl{a} = p\ct \opp p \overset{q}{=} p\ct\opp p = \refl{a}, \]
  其中标记为“$q$”的道路实际上是 $\apfunc{\lam{r} r\ct\opp p} (q)$。
  那么 $g$ 是一个等价，因为它是以下一系列等价的复合：
  \[ (p=p) \xrightarrow{\apfunc{\lam{r} r\ct\opp p}} (p\ct \opp p = p\ct \opp p) \xrightarrow{i\ct - \ct \opp i} (\refl{a} = \refl{a}). \]
  这里用到了 \cref{eg:concatequiv,thm:paths-respects-equiv}，其中 $i:\refl{a} = p\ct \opp p$ 是典范的相等。
  显然有 $g(\refl{p}) = \refl{\refl{a}}$。
\end{proof}

\section{截断}
\label{sec:truncations}

\indexsee{n-truncation@$n$-truncation}{truncation}%
\index{truncation!n-truncation@$n$-truncation|(defstyle}%

在 \cref{subsec:prop-trunc} 中，我们引入了命题截断，它给出了一个类型到“仅仅是一个命题”（即 $(-1)$-类型）的“最佳逼近”。
在 \cref{sec:hittruncations} 中，我们将此截断构造为一个高阶归纳类型，并给出了一种将其推广到 0-截断的方法。
现在我们来解释一种更好的推广，它可以将任何类型截断为任意 $n\geq -2$ 的 $n$-类型；在经典同伦论中，这被称为它的 \define{$n$ 阶 Postnikov 截面（$n^{\mathrm{th}}$ Postnikov section）}。\index{Postnikov tower}

其思路是利用 \cref{thm:ntype-nloop} 与 \cref{lem:susp-loop-adj}。
前者指出，$A$ 是 $n$-类型当且仅当对所有 $a:A$，$\Omega^{n+1}(A,a)$
\index{loop space!iterated}%
是可缩的；后者则意味着
\narrowequation{\Omega^{n+1}(A,a) \eqvsym \Map_{*}(\Sn^{n+1},(A,a)),}
其中 $\Sn^{n+1}$ 配备了一个基点，我们不妨称之为 \base。
然而，$\Map_*(\Sn^{n+1},(A,a))$ 的可缩性是我们可以通过直接给出道路构造子来保证的。

\index{hub and spoke}%
我们将使用类似于 \cref{sec:hubs-spokes} 中的“毂辐构造（hub and spoke）”。
因此，对 $n\ge -1$，我们取 $\trunc nA$ 为由以下生成子生成的高阶归纳类型：

\begin{itemize}
\item 一个函数 $\tprojf n : A \to \trunc n A$，
\item 对于每个 $r:\Sn^{n+1} \to \trunc n A$，一个\emph{轮毂}点 $h(r):\trunc n A$，以及
\item 对于每个 $r:\Sn^{n+1} \to \trunc n A$ 和每个 $x:\Sn^{n+1}$，一条\emph{辐条}道路 $s_r(x):r(x) = h(r)$。
\end{itemize}

\noindent
有了这些构造子，就足以证明：

\begin{lem}
  $\trunc n A$ 是一个 $n$-类型。
\end{lem}

\begin{proof}
  根据 \cref{thm:ntype-nloop}，只需证明对所有 $b:\trunc nA$，$\Omega ^{n+1}(\trunc nA,b)$ 是可缩的。
  而由 \cref{lem:susp-loop-adj}，这等价于 \narrowequation{\Map_*(\Sn^{n+1},(\trunc nA,b)).}
  我们以后者为收缩中心：取值为 $b$ 的常值函数 $c_b:\Sn^{n+1} \to \trunc nA$，连同道路 $\refl b : c_b(\base) = b$。

  现在，$\Map_*(\Sn^{n+1},(\trunc nA,b))$ 中的任意元素由一个映射 $r:\Sn^{n+1} \to \trunc n A$ 和一条道路 $p:r(\base)=b$ 组成。
  根据函数外延性，要证明 $r = c_b$，只需对每个 $x:\Sn^{n+1}$ 给出一条道路 $r(x)=c_b(x) \jdeq b$。
  我们将其取为复合 $s_r(x) \ct \opp{s_r(\base)} \ct p$，其中 $s_r(x)$ 是在 $x$ 处的辐条。


  最后，我们必须证明，当沿着这个相等 $r=c_b$ 输运时，道路 $p$ 会变成 $\refl b$。
  根据道路类型中的输运，这意味着我们需要
  \[\opp{(s_r(\base) \ct \opp{s_r(\base)} \ct p)} \ct p = \refl b.\]
  而这直接由道路运算可得。
\end{proof}

（此构造在 $n=-2$ 时失败，但在那种情况下，我们可以对所有 $A$ 简单地定义 $\trunc{-2}{A}\defeq \unit$。
从现在起，我们假设 $n\ge -1$。）

\index{induction principle!for truncation}%
为了证明 $n$-截断所期望的泛性质，我们需要归纳原理。
我们像往常一样从构造子中提取它；该原理是说，给定 $P:\trunc nA\to\type$ 以及


\begin{itemize}
\item 对于每个 $a:A$，一个元素 $g(a) : P(\tproj na)$，
\item 对于每个 $r:\Sn^{n+1} \to \trunc n A$ 和 $r':\prd{x:\Sn^{n+1}} P(r(x))$，一个元素 $h'(r,r'):P(h(r))$，
\item 对于每个 $r:\Sn^{n+1} \to \trunc n A$ 和 $r':\prd{x:\Sn^{n+1}} P(r(x))$，以及每个 $x:\Sn^{n+1}$，一条依值道路
$\dpath{P}{s_r(x)}{r'(x)}{h'(r,r')}$，
\end{itemize}


则存在一个截面 $f:\prd{x:\trunc n A} P(x)$ 使得对所有 $a:A$ 有 $f(\tproj n a) \jdeq g(a)$。
为使其更为有用，我们将其重新表述如下。

\begin{thm}\label{thm:truncn-ind}
  对于任意类型族 $P:\trunc n A \to \type$，若每个 $P(x)$ 都是 $n$-类型，则对任意函数 $g : \prd{a:A} P(\tproj n a)$，存在一个截面 $f:\prd{x:\trunc n A} P(x)$ 使得对所有 $a:A$ 有 $f(\tproj n a)\defeq g(a)$。
\end{thm}

\begin{proof}
  只需构造上面列出的第二和第三项数据即可，因为 $g$ 的类型恰好就是第一项数据。
  给定 $r:\Sn^{n+1} \to \trunc n A$ 和 $r':\prd{x:\Sn^{n+1}} P(r(x))$，我们有 $h(r):\trunc n A$ 和 $s_r :\prd{x:\Sn^{n+1}} (r(x) =
h(r))$。
  定义 $t:\Sn^{n+1} \to P(h(r))$ 为 $t(x) \defeq \trans{s_r(x)}{r'(x)}$。
  由于 $P(h(r))$ 是 $n$-截断的，故存在一个点 $u:P(h(r))$ 和一个收缩 $v:\prd{x:\Sn^{n+1}} (t(x) = u)$。
  定义 $h'(r,r') \defeq u$，这给出了第二项数据。
  那么（回忆依值道路的定义），$v$ 的类型恰好就是第三项数据所要求的类型。
\end{proof}

特别地，如果 $E$ 是某个 $n$-类型，我们可以考虑在 $A$ 的每一点处取值恒为 $E$ 的常值类型族。
\symlabel{extend}
\index{recursion principle!for truncation}%
于是，每个映射 $f:A\to{}E$ 都可以延拓为一个映射 $\extend{f}:\trunc nA\to{}E$，其定义为 $\extend{f}(\tproj na)\defeq f(a)$；这就是 $\trunc n A$ 的\emph{递归原理}。

归纳原理还蕴含了这种形式的函数的唯一性原理。
\index{uniqueness!principle, propositional!for functions on a truncation}%
即，如果 $E$ 是 $n$-类型，且 $g,g':\trunc nA\to{}E$ 满足对每个 $a:A$ 有 $g(\tproj na)=g'(\tproj na)$，那么对所有 $x:\trunc nA$ 有 $g(x)=g'(x)$，因为类型 $g(x)=g'(x)$ 是 $n$-类型。
因此，$g=g'$。
（事实上，当 $E$ 是 $\nplusone$-类型时，此唯一性原理更一般地也成立。）
由此我们得到以下泛性质。

\begin{lem}[截断的泛性质]\label{thm:trunc-reflective}
  \index{universal!property!of truncation}%
  设 $n\ge-2$，$A:\type$ 且 $B:\typele{n}$。则以下映射是一个等价：
  \[\function{(\trunc nA\to{}B)}{(A\to{}B)}{g}{g\circ\tprojf n}\]
\end{lem}

\begin{proof}
  给定 $B$ 是 $n$-截断的，任何 $f:A\to{}B$ 都可以延拓为一个映射 $\extend{f}:\trunc nA\to{}B$。
  映射 $\extend{f}\circ\tprojf n$ 等于 $f$，因为根据定义，对于每个 $a:A$ 我们有 $\extend{f}(\tproj na)=f(a)$。
  而映射 $\extend{g\circ\tprojf n}$ 等于 $g$，因为它们都将 $\tproj na$ 映到 $g(\tproj na)$。
\end{proof}

用范畴论的语言来说，这表明 $n$-类型构成了类型范畴的一个\emph{反射子范畴}。
\index{reflective!subcategory}%
（要完全精确地陈述这一点，需要使用 $(\infty,1)$-范畴的语言。）
\index{.infinity1-category@$(\infty,1)$-category}%
特别地，这意味着 $n$-截断具有函子性：给定 $f:A\to B$，对复合 $A\xrightarrow{f} B \to \trunc n B$ 应用递归原理，就得到一个映射 $\trunc n f: \trunc n A \to \trunc n B$。
根据定义，我们有一个同伦
\begin{equation}
  \mathsf{nat}^f_n : \prd{a:A} \trunc n f(\tproj n a) = \tproj n {f(a)},\label{eq:trunc-nat}
\end{equation}
表达了映射 $\tprojf n$ 的\emph{自然性}。

唯一性蕴含了函子律，例如 $\trunc n {g\circ f} = \trunc n g \circ \trunc n f$ 和 $\trunc n{\idfunc[A]} = \idfunc[\trunc n A]$，以及相应的相干律。
我们还有高阶函子性，例如：

\begin{lem}\label{thm:trunc-htpy}
  给定 $f,g:A\to B$ 以及同伦 $h:f\htpy g$，存在一个诱导的同伦 $\trunc n h : \trunc n f \htpy \trunc n g$，使得复合
  \begin{equation}
    \xymatrix@C=3.6pc{\tproj n{f(a)} \ar@{=}[r]^-{\opp{\mathsf{nat}^f_n(a)}} &
      \trunc n f(\tproj n a) \ar@{=}[r]^-{\trunc n h(\tproj na)} &
      \trunc n g(\tproj n a) \ar@{=}[r]^-{\mathsf{nat}^g_n(a)} &
      \tproj n{g(a)}}\label{eq:trunc-htpy}
  \end{equation}
  等于 $\apfunc{\tprojf n}(h(a))$。
\end{lem}

\begin{proof}
  首先，我们确实有一个分量为 $\apfunc{\tprojf n}(h(a)) : \tproj n{f(a)} = \tproj n{g(a)}$ 的同伦。
  将其与路径 $\tproj n{f(a)} = \trunc n f(\tproj n a)$ 和 $\tproj n{g(a)} = \trunc n g(\tproj n a)$（分别源自 $\trunc n f$ 和 $\trunc ng$ 的定义）在两侧复合，便得到同伦 $(\trunc n f \circ \tprojf n) \htpy (\trunc n g \circ \tprojf n)$，进而由函数外延性得到一个相等。
  但由于 $(\blank\circ \tprojf n)$ 是一个等价，必然存在一条道路 $\trunc nf = \trunc ng$ 诱导出该同伦，而函数外延性的相干律蕴含了~\eqref{eq:trunc-htpy}。
\end{proof}

关于反射子范畴，以下观察也是标准的。

\begin{cor}
  一个类型 $A$ 是 $n$-类型，当且仅当 $\tprojf n : A \to \trunc n A$ 是一个等价。
\end{cor}

\begin{proof}
  “当”的部分由 $n$-类型在等价下封闭可得。
  反之，若 $A$ 是 $n$-类型，我们可以定义 $\ext(\idfunc[A]):\trunc n A\to{}A$。
  那么根据定义我们有 $\ext(\idfunc[A])\circ\tprojf n=\idfunc[A]:A\to{}A$。
  为了证明
  $\tprojf n\circ\ext(\idfunc[A])=\idfunc[\trunc nA]$，
  我们只需证明
  $\tprojf n\circ\ext(\idfunc[A])\circ\tprojf n=
  \idfunc[\trunc nA]\circ\tprojf n$。
  这一点同样成立：
  \[\raisebox{\depth-\height+1em}{\xymatrix{
    A \ar^-{\tprojf n}[r] \ar_{\idfunc[A]}[rd] &
    \trunc nA \ar^>>>{\ext(\idfunc[A])}[d] \ar@/^40pt/^{\idfunc[\trunc nA]}[dd] \\
    & A \ar_{\tprojf n}[d] \\
    & \trunc nA}}
  \qedhere\]
\end{proof}

$n$-类型的范畴还有一些并非所有反射子范畴都具有的特殊性质。
例如，反射子 $\trunc n-$ 保持有限积。

\begin{thm}\label{cor:trunc-prod}
  对于任意类型 $A$ 和 $B$，诱导的映射 $\trunc n{A\times B} \to \trunc nA \times \trunc nB$ 是一个等价。
\end{thm}

\begin{proof}
  只需证明 $\trunc nA \times \trunc nB$ 具有与 $\trunc n{A\times B}$ 相同的泛性质。
  于是，设 $C$ 是一个 $n$-类型；我们有如下等价：
  \begin{align*}
    (\trunc nA \times \trunc nB \to C)
    &= (\trunc nA \to (\trunc nB \to C))\\
    &= (\trunc nA \to (B \to C))\\
    &= (A \to (B \to C))\\
    &= (A \times B \to C)
  \end{align*}
  这里用到了 $\trunc nB$ 和 $\trunc nA$ 的泛性质，以及 $C$ 是 $n$-类型时 $B\to C$ 也是 $n$-类型这一事实。
  不难验证，此等价由与 $\tprojf n \times \tprojf n$ 复合给出，恰如所需。
\end{proof}

下面这个关于依赖和的事实也常常很有用。

\begin{thm}\label{thm:trunc-in-truncated-sigma}
设 $P:A\to\type$ 为一个类型族。则存在等价
\begin{equation*}
\eqv{\Trunc n{\sm{x:A}\trunc n{P(x)}}}{\Trunc n{\sm{x:A}P(x)}}.
\end{equation*}
\end{thm}

\begin{proof}
我们多次利用 $n$-截断的归纳原理，来构造函数
\begin{align*}
\varphi & : \Trunc n{\sm{x:A}\trunc n{P(x)}}\to\Trunc n{\sm{x:A}P(x)}\\
\psi & : \Trunc n{\sm{x:A}P(x)}\to \Trunc n{\sm{x:A} \trunc n{P(x)}}
\end{align*}
以及同伦 $H:\varphi\circ\psi\htpy \idfunc$ 与 $K:\psi\circ\varphi\htpy
\idfunc$，以表明它们互为拟逆。
定义 $\varphi$ 为 $\varphi(\tproj n{\pairr{x,\tproj nu}})\defeq\tproj n{\pairr{x,u}}$。
定义 $\psi$ 为 $\psi(\tproj n{\pairr{x,u}})\defeq\tproj n{\pairr{x,\tproj nu}}$。
然后定义 $H(\tproj n{\pairr{x,u}})\defeq \refl{\tproj n{\pairr{x,u}}}$ 与
$K(\tproj n{\pairr{x,\tproj nu}})\defeq \refl{\tproj n{\pairr{x,\tproj nu}}}$。
\end{proof}

\begin{cor}\label{thm:refl-over-ntype-base}
  若 $A$ 是一个 $n$-类型，且 $P:A\to\type$ 是任意类型族，则
  \[ \eqv{\sm{a:A} \trunc n{P(a)}}{\Trunc n{\sm{a:A}P(a)}} \]
\end{cor}

\begin{proof}
  如果 $A$ 是 $n$-类型，那么上式左边的类型本身就已经是 $n$-类型，因此等价于它的 $n$-截断；于是本推论由 \cref{thm:trunc-in-truncated-sigma} 可得。
\end{proof}

我们可以用与 \cref{sec:compute-coprod,sec:compute-nat} 中处理余积和自然数时相同的方法，来刻画截断的路径空间（该方法之后在 \cref{cha:homotopy} 计算同伦群时也会用到）。
不出所料，$A$ 的 $(n+1)$-截断中的路径空间，就是 $A$ 的路径空间的 $n$-截断。
确切地说，对于任意 $x,y:A$，存在一个典范映射
\begin{equation}
  f:\ttrunc n{x=_Ay}\to \Big(\tproj {n+1}x=_{\trunc{n+1}A}\tproj {n+1}y\Big)\label{eq:path-trunc-map}
\end{equation}
其定义为
\[f(\tproj n{p})\defeq \apfunc{\tprojf {n+1}}(p). \]
这个定义使用了 $\truncf n$ 的递归原理，这是可行的，因为 $\trunc {n+1}A$ 是 $(n+1)$-截断的，因此 $f$ 的余域是 $n$-截断的。

\begin{thm} \label{thm:path-truncation}
  对于任意 $A$、$x,y:A$ 以及 $n\ge -2$，映射~\eqref{eq:path-trunc-map} 是一个等价；因此我们有
  \[ \eqv{\ttrunc n{x=_Ay}}{\Big(\tproj {n+1}x=_{\trunc{n+1}A}\tproj {n+1}y\Big)}. \]
\end{thm}

\begin{proof}
  证明是编码-解码方法的一个简单应用：
  如同以往的情形一样，我们无法直接定义映射~\eqref{eq:path-trunc-map} 的一个拟逆，因为我们无法对 $\tproj {n+1}x$ 和 $\tproj {n+1}y$ 之间的相等进行归纳。
  因此，我们转而将其类型推广，以便处理类型 $\trunc{n+1}A$ 中的一般元素，而不仅仅是 $\tproj {n+1}x$ 和 $\tproj {n+1}y$。
  如下定义 $P:\trunc {n+1}A\to\trunc {n+1}A\to\typele{n}$：
  \[P(\tproj {n+1}x,\tproj {n+1}y)\defeq \trunc n{x=_Ay}\]
  这个定义是良定义的，因为 $\trunc n{x=_Ay}$ 是 $n$-截断的，而根据 \cref{thm:hleveln-of-hlevelSn}，$\typele{n}$ 是 $(n+1)$-截断的。
  现在，对于每一对 $u,v:\trunc{n+1}A$，都有一个映射
  \[\decode:P(u,v) \to \big(u=_{\trunc{n+1}A}v\big)\]
  它对 $u=\tproj {n+1}x$、$v=\tproj {n+1}y$ 以及 $p:x=y$ 定义为
  \[\decode(\tproj n{p})\defeq \apfunc{\tprojf{n+1}} (p).\]
  由于 $\decode$ 的余域是 $n$-截断的，我们只需对 $u$ 和 $v$ 是这种形式的情形定义它即可，而这与之前的定义相同。
  我们还通过对 $u$ 作归纳来定义一个函数
  \[ r : \prd{u:\trunc{n+1} A} P(u,u) \]
  其中 $r(\tproj{n+1} x) \defeq \tproj n {\refl x}$。

  现在我们可以定义一个逆映射
  \[\encode: (u=_{\trunc{n+1}A}v) \to P(u,v)\]
  定义为
  \[\encode(p) \defeq \transfib{v\mapsto P(u,v)}{p}{r(u)}. \]
  为了证明复合映射
  \[ (u=_{\trunc{n+1}A}v) \xrightarrow{\encode} P(u,v) \xrightarrow{\decode} (u=_{\trunc{n+1}A}v) \]
  是恒等函数，根据道路归纳，只需对 $\refl u : u=u$ 的情形进行验证；此时我们需要知道的是 $\decode(r(u)) = \refl{u}$。
  但由于这是一个 $(n-1)$-类型，因而也是一个 $(n+1)$-类型，我们可以假定 $u\jdeq \tproj {n+1} x$，而在此假定下，结论由 $r$ 和 $\decode$ 的定义即得。
  最后，为了证明
  \[ P(u,v) \xrightarrow{\decode} (u=_{\trunc{n+1}A}v) \xrightarrow{\encode} P(u,v) \]
  是恒等函数，由于该目标同样是一个 $(n-1)$-类型，我们可以假定 $u=\tproj {n+1}x$ 且 $v=\tproj {n+1}y$，且我们考虑的是对应于某个 $p:x=y$ 的 $\tproj n p:P(\tproj{n+1}x,\tproj{n+1}y)$。
  那么我们有
  \begin{align*}
    \encode(\decode(\tproj n p)) &= \encode(\apfunc{\tprojf{n+1}}(p))\\
    &= \transfib{v\mapsto P(\tproj{n+1}x,v)}{\apfunc{\tprojf{n+1}}(p)}{\tproj n {\refl x}}\\
    &= \transfib{y\mapsto \trunc n{x=y}}{p}{\tproj n {\refl x}}\\
    &= \tproj n {\transfib{y \mapsto (x=y)}{p}{\refl x}}\\
    &= \tproj n p,
  \end{align*}
  这里用到了 \cref{thm:transport-compose,thm:ap-transport}。
  （或者，我们也可以对 $p$ 进行道路归纳；则所需的相等在判断意义下成立。）
  这就完成了 \decode 和 \encode 是拟逆的证明。
  所陈述的结果便是 $u=\tproj {n+1}x$ 且 $v=\tproj {n+1}y$ 时的特例。
\end{proof}

\begin{cor}
  设 $n\ge-2$，且 $(A,a)$ 是一个带点类型。那么
  \[\ttrunc n{\Omega(A,a)}=\Omega\mathopen{}\left(\trunc{n+1}{(A,a)}\right)\]
\end{cor}

\begin{proof}
  这是前一个引理在 $x=y=a$ 时的特例。
\end{proof}

\begin{cor}
  设 $n\ge -2$，$k\ge 0$，且 $(A,a)$ 是一个带点类型。那么
  \[\ttrunc n{\Omega^k(A,a)} = \Omega^k\mathopen{}\left(\trunc{n+k}{(A,a)}\right). \]
\end{cor}

\begin{proof}
  对 $k$ 进行归纳，利用 $\Omega^k$ 的递归定义即可。
\end{proof}

我们还注意到“截断是累积的”：如果我们将一个类型先截断为一个 $n$-类型，再截断为一个 $k$-类型（其中 $k\le n$），那么我们不妨一开始就直接将其截断为一个 $k$-类型。

\begin{lem} \label{lem:truncation-le}
  设 $k,n\ge-2$ 且 $k\le{}n$，$A:\type$。那么
  $\trunc k{\trunc nA}=\trunc kA$。
\end{lem}

\begin{proof}
  我们定义两个映射 $f:\trunc k{\trunc nA}\to\trunc kA$ 和
  $g:\trunc kA\to\trunc k{\trunc nA}$ 如下：
  %
  \[
   f(\tproj k{\tproj na}) \defeq \tproj ka
   \qquad\text{and}\qquad
   g(\tproj ka) \defeq \tproj k{\tproj na}.
  \]
  %
  映射 $f$ 是良定义的，因为 $\trunc kA$ 是 $k$-截断的，并且也是 $n$-截断的（因为 $k\le{}n$）；映射 $g$ 是良定义的，因为 $\trunc k{\trunc nA}$ 是 $k$-截断的。

  复合映射 $f\circ{}g:\trunc kA\to\trunc kA$ 满足
  $(f\circ{}g)(\tproj ka)=\tproj ka$，因此 $f\circ{}g=\idfunc[\trunc kA]$。
  类似地，我们有 $(g\circ{}f)(\tproj k{\tproj na})=\tproj k{\tproj na}$，因此 $g\circ{}f=\idfunc[\trunc k{\trunc nA}]$。
\end{proof}

% \begin{lem}
%   We have $\trunc n{\unit}=\unit$.
% \end{lem}
% \begin{proof}
%   Indeed, $\unit$ is $n$-truncated for every $n$ hence $\trunc n{\unit}=\unit$ by
%   \cref{reflectPequiv}.
% \end{proof}

\index{truncation!n-truncation@$n$-truncation|)}%

\section{$n$-类型的余极限}
\label{sec:pushouts}

回想一下，在 \cref{sec:colimits} 中，我们曾用高阶归纳类型定义了类型的推出，并证明了其泛性质。
一般来说，$n$-\type{} 的（同伦）余极限可能不再是 $n$-\type{}（关于一个极端的反例，可见 \cref{ex:s2-colim-unit}）。
然而，如果对其进行 $n$-截断，我们便能得到一个 $n$-\type{}，并且它对于其他 $n$-\type{} 而言满足正确的泛性质。

本节我们将对推出这一情形证明上述结论，因为推出是余极限中最重要且非平凡的情形。
我们先回顾 \cref{sec:colimits} 中给出的如下定义。

\begin{defn}
  一个 \define{伸展（span）}% in $\P$
  \indexdef{span}%
  由一个五元组 $\Ddiag=(A,B,C,f,g)$ 构成，其中 % $A,B,C:\P$ and，并有映射 $f:C\to{}A$ 与 $g:C\to{}B$。
  \[\Ddiag=\quad\vcenter{\xymatrix{C \ar^g[r] \ar_f[d] & B \\ A & }}\]
\end{defn}

\begin{defn}
  给定一个伸展 $\Ddiag=(A,B,C,f,g)$ 和一个类型 $D$，%$D:\P$, a
  \define{上锥（cocone under $\Ddiag$ with base $D$）} 是一个三元组 $(i, j, h)$，
  \index{cocone}%
  其中 $i:A\to{}D$， $j:B\to{}D$，并且 $h : \prd{c:C}i(f(c))=j(g(c))$：
  \[\uppercurveobject{{ }}\lowercurveobject{{ }}\twocellhead{{ }}
  \xymatrix{C \ar^g[r] \ar_f[d] \drtwocell{^h} & B \ar^j[d] \\ A \ar_i[r] & D
  }\]
  我们用 $\cocone{\Ddiag}{D}$ 表示所有这类上锥构成的类型。
\end{defn}

上锥的类型具有（协变）函子性。
例如，给定 $D$， $E$% $D,E:\P$
以及一个映射 $t:D\to{}E$，则存在一个映射
  \[\function{\cocone{\Ddiag}{D}}{\cocone{\Ddiag}{E}}{c}{\composecocone{t}c}\]
  其定义为：
  \[\composecocone{t}(i,j,h)=(t\circ{}i,t\circ{}j,\mapfunc{t}\circ{}h).\]
而给定 $D$， $E$， $F$，%$:\P$,
函数 $t:D\to{}E$， $u:E\to{}F$ 以及 $c:\cocone{\Ddiag}{D}$，则有
\begin{align}
  \composecocone{\idfunc[D]}c &= c \label{eq:composeconeid}\\
  \composecocone{(u\circ{}t)}c&=\composecocone{u}(\composecocone{t}c). \label{eq:composeconefunc}
\end{align}

\begin{defn}
  给定一个由 $n$-\type{} 构成的伸展 $\Ddiag$，一个 $n$-\type{} $D$，以及一个上锥
  $c:\cocone{\Ddiag}{D}$，如果对每个 $n$-\type{} $E$，映射
  \[\function{(D\to{}E)}{\cocone{\Ddiag}{E}}{t}{\composecocone{t}c}\]
  都是一个等价，则称配对 $(D,c)$ 为 \define{在 $n$-\type{} 中的推出（pushout of $\Ddiag$ in $n$-types）}。
  \indexdef{pushout!in ntypes@in $n$-types}%
\end{defn}

\begin{comment}
We showed in \cref{thm:pushout-ump} that pushouts exist when $\P$ is \type itself, by giving a direct construction in terms of higher
inductive types.
For a general \P, pushouts may or may not exist, but if they do, then they are unique.

\begin{lem}
  If $(D,c)$ and $(D',c')$ are two pushouts of $\Ddiag$ in $\P$, then
  $(D,c)=(D',c')$.
\end{lem}
\begin{proof}
  We first prove that the two types $D$ and $D'$ are equivalent.

  Using the universal property of $D$ with $D'$, we see that the following map is an
  equivalence
  %
  \[
    \function{(D\to{}D')}{\cocone{\Ddiag}{D'}}{t}{\composecocone{t}c}
  \]
  %
  In particular, there is a function $f:D\to{}D'$ satisfying $\composecocone{f}c=c'$. In the
  same way there is a function $g:D'\to{}D$ such that $\composecocone{g}c'=c$.

  In order to prove that $g\circ{}f=\idfunc[D]$ we use the universal property of
  $D$ for $D$, which says that the following map is an equivalence:
  %
  \[
  \function{(D\to{}D)}{\cocone{\Ddiag}{D}}{t}{\composecocone{t}c}
  \]
  %
  Using the functoriality of $t\mapsto{}\composecocone{t}c$ we see that
  \begin{align*}
    \composecocone{(g\circ{}f)}c &= \composecocone{g}(\composecocone{f}c) \\
    &= \composecocone{g}c' \\
    &= c \\
    &= \composecocone{\idfunc[D]}c
  \end{align*}
  hence
  $g\circ{}f=\idfunc[D]$, because equivalences are injective. The same argument
  with $D'$ instead of $D$ shows that $f\circ{}g=\idfunc[D']$.

  Hence $D$ and $D'$ are equal, and the fact that $(D,c)=(D',c')$ follows from
  the fact that the equivalence between $D$ and $D'$ we just defined sends $c$
  to $c'$.
\end{proof}

\begin{cor}
  The type of pushouts of $\Ddiag$ in $\P$ is a mere proposition. In particular if
  pushouts merely exist then they actually exist.
\end{cor}

As in the case of pullbacks, if \P is reflective, then pushouts in \P always exist.
However, unlike the case of pullbacks, pushouts in \P are not the same as the pushouts in \type: they are obtained by applying the
reflector.
\end{comment}

为了构造 $n$-\type{} 的推出，我们需要先说明如何反射伸展与上锥。

\bgroup
\def\reflect(#1){\trunc n{#1}}

\begin{defn}
  设
  \[\Ddiag=\quad\vcenter{\xymatrix{C \ar^g[r] \ar_f[d] & B \\ A & }}\]
  为一个伸展（span）。我们用 $\reflect(\Ddiag)$ 表示如下的 $n$-类型伸展：
  \[\reflect(\Ddiag)\defeq\quad \vcenter{\xymatrix{\reflect(C) \ar^{\reflect(g)}[r]
      \ar_{\reflect(f)}[d] & \reflect(B) \\ \reflect(A) & }}\]
\end{defn}

\begin{defn}
  设 $D:\type$ 且 $c=(i,j,h):\cocone{\Ddiag}{D}$。
  我们定义
  \[\reflect(c)=(\reflect(i),\reflect(j),k):
  \cocone{\reflect(\Ddiag)}{\reflect(D)}\]
  其中 $k$ 为复合同伦
  \[ \reflect(i) \circ \reflect(f) \htpy \reflect(i\circ f) \htpy \reflect(j\circ g) \htpy \reflect(j) \circ \reflect(g) \]
  这里使用了 \cref{thm:trunc-htpy} 以及 $\reflect(\blank)$ 的函子性。
  % \[\reflect(h):\prd{c:\reflect(C)}\reflect(i)(\reflect(f)(c))=\reflect(j)(\reflect(g)(c))\]
  % is defined in the following way:
\end{defn}

\egroup

我们现在观察到，从每个类型到其 $n$-截断的映射，在以下意义下组合成了一个伸展之间的映射。

\begin{defn}
  设
  \[\Ddiag=\quad\vcenter{\xymatrix{C \ar^g[r] \ar_f[d] & B \\ A & }}
  \qquad\text{and}\qquad
  \Ddiag'=\quad\vcenter{\xymatrix{C' \ar^{g'}[r] \ar_{f'}[d] & B' \\ A' & }}
  \]
  为两个伸展。
  一个从 $\Ddiag$ 到 $\Ddiag'$ 的\define{伸展之间的映射}
  \indexdef{map!of spans}%
  由函数 $\alpha:A\to A'$、$\beta:B\to B'$、$\gamma:C\to C'$ 以及同伦 $\phi: \alpha\circ f \htpy f'\circ \gamma$ 和 $\psi:\beta\circ g \htpy g' \circ \gamma$ 构成。
\end{defn}

因此，对任意伸展 $\Ddiag$，我们都有一个由 $\tprojf[A]n$、$\tprojf[B]n$、$\tprojf[C]n$ 以及来自~\eqref{eq:trunc-nat} 的自然性同伦 $\mathsf{nat}^f_n$ 和 $\mathsf{nat}^g_n$ 所构成的伸展之间的映射 $\tprojf[\Ddiag] n : \Ddiag \to \trunc n\Ddiag$。

我们还需要知道伸展之间的映射具有函子性。
即，若 $(\alpha,\beta,\gamma,\phi,\psi):\Ddiag \to \Ddiag'$ 是一个伸展之间的映射且 $D$ 是任意类型，则有
\[ \function{\cocone{\Ddiag'}{D}}{\cocone{\Ddiag}{D}}{(i,j,h)}{(i\circ \alpha,j\circ\beta, k)} \]
其中 $k: \prd{z:C} i(\alpha(f(z))) = j(\beta(g(z)))$ 是复合
\begin{equation}\label{eq:mapofspans-htpy}
\xymatrix{
  i(\alpha(f(z))) \ar@{=}[r]^{\apfunc{i}(\phi)} &
  i(f'(\gamma(z))) \ar@{=}[r]^{h(\gamma(z))} &
  j(g'(\gamma(z))) \ar@{=}[r]^{\apfunc{j}(\psi)} &
  j(\beta(g(z))). }
\end{equation}
我们将这个上锥记为 $(i,j,h) \circ (\alpha,\beta,\gamma,\phi,\psi)$。
此外，该函子作用与上锥的另一函子作用可交换：

\begin{lem}\label{thm:conemap-funct}
  给定 $(\alpha,\beta,\gamma,\phi,\psi):\Ddiag \to \Ddiag'$ 以及 $t:D\to E$，下图交换：
  \begin{equation*}
    \vcenter{\xymatrix{
        \cocone{\Ddiag'}{D}\ar[r]^{t \circ {\blank}}\ar[d] &
        \cocone{\Ddiag'}{E}\ar[d]\\
        \cocone{\Ddiag}{D}\ar[r]_{t \circ {\blank}} &
        \cocone{\Ddiag}{E}
      }}
  \end{equation*}
\end{lem}

\begin{proof}
  给定 $(i,j,h):\cocone{\Ddiag'}{D}$，注意两个复合都给出一个上锥，其前两个分量为 $t\circ i\circ \alpha$ 和 $t\circ j\circ\beta$。
  因此，只需验证二者的同伦是否相同。
  对于右上方的复合，其同伦是将 $(i,j,h)$ 替换为 $(t\circ i, t\circ j, \apfunc{t}\circ h)$ 后的~\eqref{eq:mapofspans-htpy}：
  \begin{equation*}
    \xymatrix@+2.8em{
      {t \, i \, \alpha \, f \, z} \ar@{=}[r]^{\apfunc{t\circ i}(\phi)} &
      {t \, i \, f' \, \gamma \, z} \ar@{=}[r]^{\apfunc{t}(h(\gamma(z)))} &
      {t \, j \, g' \, \gamma \, z} \ar@{=}[r]^{\apfunc{t\circ j}(\psi)} &
      {t \, j \, \beta \, g \, z}
    }
  \end{equation*}
  （为简洁起见，我们省略了函数参数周围的括号。）
  另一方面，对于左下方的复合，其同伦是 $\apfunc{t}$ 作用于~\eqref{eq:mapofspans-htpy} 的结果。
  由于 $\apfunc{}$ 保路径的拼接，这等于
  \begin{equation*}
    \xymatrix@+2.8em{
      {t \, i \, \alpha \, f \, z} \ar@{=}[r]^{\apfunc{t}(\apfunc{i}(\phi))} &
      {t \, i \, f' \, \gamma \, z} \ar@{=}[r]^{\apfunc{t}(h(\gamma(z)))} &
      {t \, j \, g' \, \gamma \, z} \ar@{=}[r]^{\apfunc{t}(\apfunc{j}(\psi))} &
      {t \, j \, \beta \, g \, z}. }
  \end{equation*}
  但是 $\apfunc{t}\circ \apfunc{i} = \apfunc{t\circ i}$，且对 $j$ 类似，因此这两个同伦相等。
\end{proof}

最后注意到，由于我们使用 \cref{thm:trunc-htpy} 定义了 $\trunc nc : \cocone{\trunc n \Ddiag}{\trunc n D}$，附加条件~\eqref{eq:trunc-htpy} 蕴含着
\begin{equation}
  \tprojf[D] n \circ c = \trunc n c \circ \tprojf[\Ddiag]n. \label{eq:conetrunc}
\end{equation}
对任意 $c:\cocone{\Ddiag}{D}$ 成立。
现在我们可以证明期望的定理了。

\begin{thm}
  \label{reflectcommutespushout}
  \index{universal!property!of pushout}%
  设 $\Ddiag$ 为一个伸展，$(D,c)$ 是其推出。
  则 $(\trunc nD,\trunc n c)$ 是 $\trunc n\Ddiag$ 在 $n$-类型中的推出。
\end{thm}

\begin{proof}
  设 $E$ 是一个 $n$-类型，考虑下图：
\bgroup
\def\reflect(#1){\trunc n{#1}}
  \begin{equation*}
  \vcenter{\xymatrix{
      (\trunc nD \to E)\ar[r]^-{\blank\circ \tprojf[D] n}\ar[d]_{\blank\circ \trunc nc} &
      (D\to E)\ar[d]^{\blank\circ c}\\
      \cocone{\trunc n \Ddiag}{E}\ar[r]^-{\blank\circ \tprojf[\Ddiag]n}\ar@{<-}[d]_{\ell_1} &
      \cocone{\Ddiag}{E}\ar@{<-}[d]^{\ell_2}\\
      (\reflect(A)\to{}E)\times_{(\reflect(C)\to{}E)}(\reflect(B)\to{}E)\ar[r] &
      (A\to{}E)\times_{(C\to{}E)}(B\to{}E)
      }}
  \end{equation*}
\egroup
  上方水平箭头是等价，因为 $E$ 是 $n$-类型；而 $\blank\circ c$ 是等价，因为 $c$ 是推出上锥。
  因此，根据 2-out-of-3 性质，要证明 $\blank\circ \trunc nc$ 是等价，只需证明上方方块交换且中间水平箭头是等价。
  为证明上方方块交换，令 $t:\trunc nD \to E$；则有
  \begin{align}
    \big(t \circ \trunc n c\big) \circ \tprojf[\Ddiag] n
    &= t \circ \big(\trunc n c \circ \tprojf[\Ddiag] n\big)
    \tag{by \cref{thm:conemap-funct}}\\
    &= t\circ \big(\tprojf[D]n \circ c\big)
    \tag{by~\eqref{eq:conetrunc}}\\
    &= \big(t\circ \tprojf[D]n\big) \circ c
    \tag{by~\eqref{eq:composeconefunc}}.
  \end{align}
  为证明中间水平箭头是等价，考虑下方方块。
  下方两个垂直箭头就是 $\happly$ 的应用：
  \begin{align*}
    \ell_1(i,j,p) &\defeq (i,j,\happly(p))\\
    \ell_2(i,j,p) &\defeq (i,j,\happly(p))
  \end{align*}
  因此由函数外延性可知它们是等价。
  最下方的水平箭头定义为
  \[ (i,j,p) \mapsto \big( i\circ \tprojf[A]n,\;\; j \circ \tprojf[B] n,\;\; q\big) \]
  其中 $q$ 是复合
  \begin{align}
    i\circ \tprojf[A]n \circ f
    &= i\circ \trunc nf \circ \tprojf[C]n
    \tag{by $\funext(\lam{z} \apfunc{i}(\mathsf{nat}^f_n(z)))$}\\
    &= j\circ \trunc ng \circ \tprojf[C]n
    \tag{by $\apfunc{\blank\circ \tprojf[C] n}(p)$}\\
    &= j\circ \tprojf[B]n \circ g.
    \tag{by $\funext(\lam{z} \apfunc{j}(\mathsf{nat}^g_n(z)))$}
  \end{align}
  这是一个等价，因为它由一个余伸展的等价所诱导。
  因此，再次由 2-out-of-3，只需证明下方方块交换。
  但绕行下方方块的两个复合在前两个分量上依定义相等，故只需证明对左下角的 $(i,j,p)$ 和 $z:C$，路径
  \[ \happly(q,z) : i(\tproj n{f(z)}) = j(\tproj n{g(z)}) \]
  （其中 $q$ 如上所述）
  等于复合
  \begin{align}
    i(\tproj n{f(z)})
    &= i(\trunc nf(\tproj nz))
    \tag{by $\apfunc{i}(\mathsf{nat}^f_n(z))$}\\
    &= j(\trunc ng(\tproj nz))
    \tag{by $\happly(p,\tproj nz)$}\\
    &= j(\tproj n{g(z)}).
    \tag{by $\apfunc{j}(\mathsf{nat}^g_n(z))$}
  \end{align}
  然而，由于 $\happly$ 具有函子性，只需检查三个分量路径相等即可：
  \begin{align*}
    \happly({\funext(\lam{z} \apfunc{i}(\mathsf{nat}^f_n(z)))},z)
    &= {\apfunc{i}(\mathsf{nat}^f_n(z))}\\
    \happly(\apfunc{\blank\circ \tprojf[C] n}(p), z)
    &= {\happly(p,\tproj nz)}\\
    \happly({\funext(\lam{z} \apfunc{j}(\mathsf{nat}^g_n(z)))},z)
    &= {\apfunc{j}(\mathsf{nat}^g_n(z))}.
  \end{align*}
  其中第一和第三条路径正是 $\happly$ 是 $\funext$ 的拟逆这一事实，而第二条路径则是关于 $\happly$ 与预复合的一个简单通用引理。
\end{proof}

\section{连通性}
\label{sec:connectivity}

$n$-类型是指在 $n$ 维以上不包含有趣信息的类型。
与之相对，一个\emph{$n$-连通类型}是指在 $n$ 维\emph{以下}不包含有趣信息的类型。
事实证明，研究函数的更一般概念也是很自然的。

\begin{defn}
一个函数 $f:A\to B$ 被称为是 \define{$n$-连通的（$n$-connected）}
\indexdef{function!n-connected@$n$-connected}%
\indexsee{n-connected@$n$-connected!function}{function, $n$-connected}%
如果对于所有 $b:B$，类型 $\trunc n{\hfiber f b}$ 都是可缩的：
\begin{equation*}
  \mathsf{conn}_n(f)\defeq \prd{b:B}\iscontr(\trunc n{\hfiber{f}b}).
\end{equation*}
一个类型 $A$ 被称为是 \define{$n$-连通的（$n$-connected）}
\indexsee{n-connected@$n$-connected!type}{type, $n$-connected}%
\indexdef{type!n-connected@$n$-connected}%
 如果唯一的函数 $A\to\unit$ 是 $n$-连通的，亦即，如果 $\trunc nA$ 是可缩的。
\end{defn}


\indexsee{connected!function}{function, $n$-connected}

因此，一个函数 $f:A\to B$ 是 $n$-连通的，当且仅当对每个 $b:B$，其纤维 $\hfib{f}b$ 都是 $n$-连通的。
显然，每个函数都是 $(-2)$-连通的。
在下一个层级上，我们有：

\begin{lem}\label{thm:minusoneconn-surjective}
  \index{function!surjective}%
  一个函数 $f$ 是 $(-1)$-连通的，当且仅当它是\cref{sec:mono-surj}意义下的满的（surjective）。
\end{lem}

\begin{proof}
  我们曾定义 $f$ 是满的，若对所有 $b$，类型 $\trunc{-1}{\hfiber f b}$ 都有实例。
  但由于这是一个命题，有实例等价于可缩。
\end{proof}

这样一来，对于 $n\ge 0$，函数的 $n$-连通性可以看作是一种更强的满射性。
从范畴论角度看，$(-1)$-连通性对应于对象上的本质满性，而 $n$-连通性则对应于 $k$-态射（$k\le n+1$）上的本质满性。

\cref{thm:minusoneconn-surjective}还蕴含着，一个类型 $A$ 是 $(-1)$-连通的，当且仅当它仅知有实例。
当一个类型是 $0$-连通的时候，我们可以简称为它是 \define{连通的}，
\indexdef{connected!type}%
\indexdef{type!connected}%
而当它是 $1$-连通的，我们则称之为是 \define{单连通的}。
\indexdef{simply connected type}%
\indexdef{type!simply connected}%

\begin{rmk}\label{rmk:connectedness-indexing}
  尽管我们对于类型的 $n$-连通性定义与同伦论中的标准定义一致，但我们对于\emph{函数}的 $n$-连通性定义，与经典同伦论中一个常见的编号约定相差了 1。
  我们称一个函数 $f$ 是 $n$-连通的，如果它的所有纤维都是 $n$-连通的；而一些经典同伦论学者会称这样的函数为 $(n+1)$-连通的。
  （这是由于历史上更关注\emph{余纤维（cofibers）}而非纤维。）
\end{rmk}

现在我们来观察连通映射的几条闭包性质。
\index{function!n-connected@$n$-connected}

\begin{lem}
\index{retract!of a function}%
设 $g$ 是一个 $n$-连通函数 $f$ 的收缩核。那么 $g$ 也是 $n$-连通的。
\end{lem}

\begin{proof}
这是\cref{lem:func_retract_to_fiber_retract}的直接推论。
\end{proof}

\begin{cor}
如果 $g$ 同伦于一个 $n$-连通函数 $f$，那么 $g$ 也是 $n$-连通的。
\end{cor}

\begin{lem}\label{lem:nconnected_postcomp}
设 $f:A\to B$ 是 $n$-连通的。那么 $g:B\to C$ 是 $n$-连通的，当且仅当复合 $g\circ f$ 是 $n$-连通的。
\end{lem}

\begin{proof}
对任意 $c:C$，我们有
\begin{align*}
  \trunc n{\hfib{g\circ f}c}
  & \eqvsym \Trunc n{ \sm{w:\hfib{g}c}\hfib{f}{\proj1 w}}
  \tag{by \cref{ex:unstable-octahedron}}\\
  & \eqvsym \Trunc n{\sm{w:\hfib{g}c} \trunc n{\hfib{f}{\proj1 w}}}
  \tag{by \cref{thm:trunc-in-truncated-sigma}}\\
  & \eqvsym \trunc n{\hfib{g}c}.
  \tag{since $\trunc n{\hfib{f}{\proj1 w}}$ is contractible}
\end{align*}
由此可见，$\trunc n{\hfib{g}c}$ 是可缩的，当且仅当$\trunc n{\hfib{g\circ f}c}$是可缩的。
\end{proof}

重要的是，$n$-连通函数可以等价地刻画为那些满足关于 $n$-类型的“归纳原理”的函数。\index{induction principle!for connected maps}
这一想法将直接引导我们在\cref{sec:freudenthal}中证明 Freudenthal 纬悬定理。

\begin{lem}\label{prop:nconnected_tested_by_lv_n_dependent types}
对于函数 $f:A\to B$ 和类型族 $P:B\to\type$，考虑以下函数：
\begin{equation*}
\lam{s} s\circ f :\Parens{\prd{b:B} P(b)}\to\Parens{\prd{a:A}P(f(a))}.
\end{equation*}
对于一个固定的 $f$ 和 $n\ge -2$，下列命题等价：
\begin{enumerate}
\item $f$ 是 $n$-连通的。\label{item:conntest1}
\item 对于每个 $P:B\to\ntype{n}$，映射 $\lam{s} s\circ f$ 都是一个等价。\label{item:conntest2}
\item 对于每个 $P:B\to\ntype{n}$，映射 $\lam{s} s\circ f$ 都有一个截面。\label{item:conntest3}
\end{enumerate}
\end{lem}

\begin{proof}
假设 $f$ 是 $n$-连通的，并设 $P:B\to\ntype{n}$。那么我们有如下等价关系：
\begin{align}
  \prd{b:B} P(b) & \eqvsym \prd{b:B} \Parens{\trunc n{\hfib{f}b} \to P(b)}
  \tag{since $\trunc n{\hfib{f}b}$ is contractible}\\
  & \eqvsym \prd{b:B} \Parens{\hfib{f}b\to P(b)}
  \tag{since $P(b)$ is an $n$-type}\\
  & \eqvsym \prd{b:B}{a:A}{p:f(a)= b} P(b)
  \tag{by the left universal property of $\Sigma$-types}\\
  & \eqvsym \prd{a:A} P(f(a)).
  \tag{by the left universal property of path types}
\end{align}
我们略去关于此等价确实由 $\lam{s} s\circ f$ 给出的证明。
因此，有\ref{item:conntest1}$\Rightarrow$\ref{item:conntest2}，且显然有\ref{item:conntest2}$\Rightarrow$\ref{item:conntest3}。
为证\ref{item:conntest3}$\Rightarrow$\ref{item:conntest1}，考虑类型族
\begin{equation*}
P(b)\defeq \trunc n{\hfib{f}b}.
\end{equation*}
那么，由\ref{item:conntest3}可得到一个映射$c:\prd{b:B} \trunc n{\hfib{f}b}$，满足$c(f(a))=\tproj n{\pairr{a,\refl{f(a)}}}$。
为证明每个 $\trunc n{\hfib{f}b}$ 都是可缩的，我们需要构造一个类型如下的函数：
\begin{equation*}
\prd{b:B}{w:\trunc n{\hfib{f}b}} w= c(b).
\end{equation*}
根据\cref{thm:truncn-ind}，为此只需构造一个类型如下的函数：
\begin{equation*}
\prd{b:B}{a:A}{p:f(a)= b} \tproj n{\pairr{a,p}}= c(b).
\end{equation*}
但通过调整变量并使用道路归纳，这等价于如下类型：
\begin{equation*}
\prd{a:A} \tproj n{\pairr{a,\refl{f(a)}}}= c(f(a)).
\end{equation*}
而根据我们对 $c(f(a))$ 的选取，这一性质成立。
\end{proof}

\begin{cor}\label{cor:totrunc-is-connected}
对任意类型 $A$，典范函数 $\tprojf n:A\to\trunc n A$ 是 $n$-连通的。
\end{cor}

\begin{proof}
根据\cref{thm:truncn-ind}以及与之相伴的唯一性原理，\cref{prop:nconnected_tested_by_lv_n_dependent types}的条件成立。
\end{proof}

例如，当 $n=-1$ 时，\cref{cor:totrunc-is-connected}说明从一个类型到其命题截断的映射 $A\to \brck A$ 是满的。

\begin{cor}\label{thm:nconn-to-ntype-const}\label{connectedtotruncated}
一个类型 $A$ 是 $n$-连通的，当且仅当对每个 $n$-类型 $B$，映射
\begin{equation*}
  \lam{b}{a} b: B \to (A\to B)
\end{equation*}
都是等价。
换言之，“从 $A$ 到任何 $n$-类型的映射都是常值映射”。
\end{cor}

\begin{proof}
  将 \cref{prop:nconnected_tested_by_lv_n_dependent types} 应用于一个陪域为 $\unit$ 的函数即得。
\end{proof}

\begin{lem}\label{lem:nconnected_to_leveln_to_equiv}
设 $B$ 是一个 $n$-类型，并设 $f:A\to B$ 是一个函数。那么诱导函数 $g:\trunc n A\to B$ 是等价，当且仅当 $f$ 是 $n$-连通的。
\end{lem}

\begin{proof}
根据 \cref{cor:totrunc-is-connected}，$\tprojf n$ 是 $n$-连通的。
由于 $f = g\circ \tprojf n$，根据
\cref{lem:nconnected_postcomp} 可知，$f$ 是 $n$-连通的当且仅当 $g$ 是 $n$-连通的。
但 $g$ 是两个 $n$-类型之间的函数，所以它的纤维也都是 $n$-类型。
因此，$g$ 是 $n$-连通的当且仅当它是等价。
\end{proof}

我们也可以用基点包含映射的连通性来刻画连通的带点类型。

\begin{lem}\label{thm:connected-pointed}
  \index{basepoint}%
  设 $A$ 是一个类型，$a_0:\unit\to A$ 是一个基点，且 $n\ge -1$。
  那么 $A$ 是 $n$-连通的，当且仅当映射 $a_0$ 是 $(n-1)$-连通的。
\end{lem}

\begin{proof}
  首先假设 $a_0:\unit\to A$ 是 $(n-1)$-连通的，并设 $B$ 是一个 $n$-类型；我们将使用 \cref{thm:nconn-to-ntype-const}。
  映射 $\lam{b}{a} b: B \to (A\to B)$ 有一个收缩，由 $f\mapsto f(a_0)$ 给出，因此只需证明它也有一个截面，即对任意 $f:A\to B$，存在 $b:B$ 使得 $f = \lam{a}b$。
  我们取 $b\defeq f(a_0)$。
  定义 $P:A\to\type$ 为 $P(a) \defeq (f(a)=f(a_0))$。
  那么 $P$ 是一个由 $(n-1)$-类型构成的族，并且我们有 $P(a_0)$；由于 $a_0:\unit\to A$ 是 $(n-1)$-连通的，因此我们有 $\prd{a:A} P(a)$。
  于是，$f = \lam{a} f(a_0)$ 如所需。

  现在假设 $A$ 是 $n$-连通的，设给定 $P:A\to\ntype{(n-1)}$ 和 $u:P(a_0)$。
  根据 \cref{prop:nconnected_tested_by_lv_n_dependent types}，只需构造 $f:\prd{a:A} P(a)$ 使得 $f(a_0)=u$。
  现在 $\ntype{(n-1)}$ 是一个 $n$-类型且 $A$ 是 $n$-连通的，因此由 \cref{thm:nconn-to-ntype-const}，存在一个 $n$-类型 $B$ 使得 $P = \lam{a} B$。
  于是，我们得到一族等价 $g:\prd{a:A} (\eqv{P(a)}{B})$。
  定义 $f(a) \defeq \opp{g_a}(g_{a_0}(u))$；则 $f:\prd{a:A} P(a)$ 且 $f(a_0) = u$ 如所需。
\end{proof}

特别地，一个带点类型 $(A,a_0)$ 是 0-连通的，当且仅当 $a_0:\unit\to A$ 是满的，也就是说 $\prd{x:A} \brck{x=a_0}$。
关于不一定带点的情形的类似结果，参见 \cref{ex:connectivity-inductively}。

\cref{lem:nconnected_postcomp} 的一个有用变体是：

\begin{lem}\label{lem:nconnected_postcomp_variation}
设 $f:A\to B$ 是一个函数，$P:A\to\type$ 和 $Q:B\to\type$ 是类型族。假设 $g:\prd{a:A} P(a)\to Q(f(a))$
是一个逐纤维 $n$-连通%
\index{fiberwise!n-connected family of functions@$n$-connected family of functions}
的函数族，即每个函数 $g_a : P(a) \to Q(f(a))$ 都是 $n$-连通的。
若 $f$ 也是 $n$-连通的，则函数
\begin{align*}
\varphi &:\Parens{\sm{a:A} P(a)}\to\Parens{\sm{b:B} Q(b)}\\
\varphi(a,u) &\defeq \pairr{f(a),g_a(u)}.
\end{align*}
也是 $n$-连通的。
反之，若 $\varphi$ 和每个 $g_a$ 都是 $n$-连通的，且 $Q$ 逐纤维地仅知有实例（即对所有 $b:B$ 有 $\brck{Q(b)}$），则 $f$ 是 $n$-连通的。
\end{lem}

\begin{proof}
对任意 $b:B$ 和 $v:Q(b)$，我们有
{\allowdisplaybreaks
\begin{align*}
\trunc n{\hfib{\varphi}{\pairr{b,v}}} & \eqvsym \Trunc n{\sm{a:A}{u:P(a)}{p:f(a)= b} \trans{p}{g_a(u)}= v}\\
& \eqvsym \Trunc n{\sm{w:\hfib{f}b}{u:P(\proj1(w))} g_{\proj 1 w}(u)= \trans{\opp{\proj2(w)}}{v}}\\
& \eqvsym \Trunc n{\sm{w:\hfib{f}b} \hfib{g(\proj1 w)}{\trans{\opp{\proj 2(w)}}{v}}}\\
& \eqvsym \Trunc n{\sm{w:\hfib{f}b} \trunc n{\hfib{g(\proj1 w)}{\trans{\opp{\proj 2(w)}}{v}}}}\\
& \eqvsym \trunc n{\hfib{f}b}
\end{align*}}%
其中沿 $f(p)$ 和 $f(p)^{-1}$ 的输运是关于 $Q$ 的。
因此，若其中一个是可缩的，则另一个也是可缩的。

特别地，若 $f$ 是 $n$-连通的，则 $\trunc n{\hfib{f}b}$ 对所有 $b:B$ 都是可缩的，从而 $\trunc n{\hfib{\varphi}{\pairr{b,v}}}$ 对所有 $(b,v):\sm{b:B} Q(b)$ 也是可缩的。
另一方面，若 $\varphi$ 是 $n$-连通的，则 $\trunc n{\hfib{\varphi}{\pairr{b,v}}}$ 对所有 $(b,v)$ 都是可缩的，因此对任何存在 $v:Q(b)$ 的 $b:B$，$\trunc n{\hfib{f}b}$ 也是可缩的。
最后，由于可缩性是一个命题，只需仅知存在这样的 $v$ 即可。
\end{proof}

如果 $Q$ 不是逐纤维仅知有实例的，\cref{lem:nconnected_postcomp_variation} 的逆命题可能不成立。
例如，如果 $P$ 和 $Q$ 都恒取 $\emptyt$，那么 $\varphi$ 和每个 $g_a$ 都是等价，但 $f$ 可以是任意的。

在另一个方向上，我们有

\begin{lem}\label{prop:nconn_fiber_to_total}
设 $P,Q:A\to\type$ 是类型族，并考虑一个从 $P$ 到 $Q$ 的逐纤维变换\index{fiberwise!transformation}
\begin{equation*}
f:\prd{a:A} \Parens{P(a)\to Q(a)}
\end{equation*}
。
那么诱导映射 $\total f: \sm{a:A}P(a) \to \sm{a:A} Q(a)$ 是 $n$-连通的，当且仅当每个 $f(a)$ 是 $n$-连通的。
\end{lem}

当然，“仅当”方向也是 \cref{lem:nconnected_postcomp_variation} 的一个特例。

\begin{proof}
根据 \cref{fibwise-fiber-total-fiber-equiv}，对每个 $x:A$ 和 $v:Q(x)$，我们有
$\hfib{\total f}{\pairr{x,v}}\eqvsym\hfib{f(x)}v$
因此 $\trunc n{\hfib{\total f}{\pairr{x,v}}}$ 是可缩的，当且仅当
$\trunc n{\hfib{f(x)}v}$ 是可缩的。
\end{proof}

关于连通映射的另一个有用事实是，它们在 $n$-截断上诱导一个等价：

\begin{lem} \label{lem:connected-map-equiv-truncation}
若 $f : A \to B$ 是 $n$-连通的，则它诱导一个等价
$\eqv{\trunc{n}{A}}{\trunc{n}{B}}$。
\end{lem}

\begin{proof}
令 $c$ 为 $f$ 是 $n$-连通的证明。从左到右，我们使用映射 $\trunc{n}{f} : \trunc{n}{A} \to \trunc{n}{B}$。
为了定义从右到左的映射，根据截断的泛性质，只需给出一个映射 $\mathsf{back} : B \to {\trunc{n}{A}}$。
我们可以如下定义这个映射：
\[
\mathsf{back}(y) \defeq \trunc{n}{\proj{1}}{(\proj{1}{(c(y))})}.
\]
根据定义，$c(y)$ 的类型是 $\iscontr(\trunc n {\hfiber{f}y})$，所以它的第一个分量的类型是 $\trunc n{\hfiber{f}y}$，而我们通过投影可以从中得到一个 $\trunc n A$ 的元素。

接下来，我们证明这两个复合映射都是恒等映射。在两个方向上，由于目标都是在 $n$-截断类型中给出一条道路，只需处理构造子 $\tprojf{n}$ 的情形。

在一个方向上，需证对任意 $x:A$，
\[
\trunc{n}{\proj{1}}{(\proj{1}{(c(f(x)))})} = \tproj{n}{x}.
\]
但 $\tproj{n}{(x, \refl{f(x)})} : \trunc n{\hfiber{f}{f(x)}}$，而 $c(f(x))$ 断言这个类型是可缩的，因此
\[
\proj{1}{(c(f(x)))} = \tproj{n}{(x, \refl{})}.
\]
将 $\trunc{n}{\proj{1}}$ 应用于该等式两边即得结论。

在另一个方向上，需证对任意 $y:B$，
\[
\trunc{n}{f}(\trunc{n}{\proj{1}} (\proj{1}{(c(y))})) = \tproj{n}{y}.
\]
$\proj{1}{(c(y))}$ 的类型是 $\trunc n {\hfiber{f}y}$，我们想要的道路本质上是 $\hfiber{f}y$ 的第二个分量，但要确保截断的部分能对应上。

一般地，假设给定 $p:\trunc{n}{\sm{x:A} B(x)}$ 并希望证明
$P(\trunc{n}{\proj{1}{}}(p))$。根据截断归纳法，只需对所有 $a:A$ 和 $b:B(a)$ 证明 $P(\tproj{n}{a})$。
将这一原理应用于当前情形，只需在给定 $a:A$ 和 $b:f (a) = y$ 时证明
\[
\trunc{n}{f}(\tproj{n}{a}) = \tproj{n}{y}
\]
但左边等于 $\tproj{n}{f (a)}$，所以将 $\tprojf{n}$ 应用于 $b$ 两边即得结论。
\end{proof}

有人可能会猜想这个事实完全刻画了 $n$-连通映射，但实际上 $n$-连通性比这稍强一些。
例如，包含映射 $\bfalse:\unit \to\bool$ 在 $(-1)$-截断上诱导一个等价，但它不是满的（即不是 $(-1)$-连通的）。
在 \cref{sec:long-exact-sequence-homotopy-groups} 中我们将看到，一般的差异在于一种类似的额外满射性。

\section{正交分解}
\label{sec:image-factorization}

\index{unique!factorization system|(}%
\index{orthogonal factorization system|(}%
在集合论中，满射和单射构成一个唯一分解系统：每个函数本质上唯一地分解为一个满射后接一个单射。
我们已经看到，满射可以自然地推广到 $n$-连通映射，因此很自然地会问：这些映射是否也参与一个分解系统？
以下是单射的相应推广。

\begin{defn}
  一个函数 $f:A\to B$ 是\define{$n$-截断的（$n$-truncated）}
  \indexdef{n-truncated@$n$-truncated!function}%
  \indexdef{function!n-truncated@$n$-truncated}%
  如果对任意 $b:B$，纤维 $\hfib f b$ 都是 $n$-类型。
\end{defn}

特别地，$f$ 是 $(-2)$-截断的当且仅当它是等价。
当然，$A$ 是 $n$-类型当且仅当 $A\to\unit$ 是 $n$-截断的。
此外，$n$-截断映射也可以像 $n$-类型那样递归地定义。

\begin{lem}\label{thm:modal-mono}
  对任意 $n\ge -2$，函数 $f:A\to B$ 是 $(n+1)$-截断的当且仅当对任意 $x,y:A$，映射 $\apfunc{f}:(x=y) \to (f(x)=f(y))$ 是 $n$-截断的。
  \index{function!embedding}%
  \index{function!injective}%
  特别地，$f$ 是 $(-1)$-截断的当且仅当它是 \cref{sec:mono-surj} 中意义上的嵌入。
\end{lem}

\begin{proof}
  注意对任意 $(x,p),(y,q):\hfib f b$，有
  \begin{align*}
    \big((x,p) = (y,q)\big)
    &= \sm{r:x=y} (p = \apfunc f(r)\ct q)\\
    &= \sm{r:x=y} (\apfunc f (r) = p\ct \opp q)\\
    &= \hfib{\apfunc{f}}{p\ct \opp q}.
  \end{align*}
  因此，$f$ 的任意纤维中的任意道路空间都是 $\apfunc{f}$ 的纤维。
  另一方面，取 $b\defeq f(y)$ 和 $q\defeq \refl{f(y)}$，我们看到 $\apfunc f$ 的任意纤维都是 $f$ 的纤维中的一个道路空间。
  由于 $f$ 是 $\nplusone$-截断的当且仅当其所有纤维的道路空间都是 $n$-类型，结论得证。
\end{proof}

现在我们可以以一种相当明显的方式构造分解。

\begin{defn}\label{defn:modal-image}
设 $f:A\to B$ 是一个函数。$f$ 的\define{$n$-像（$n$-image）}
\indexdef{image}%
\indexdef{image!n-image@$n$-image}%
\indexdef{n-image@$n$-image}%
\indexdef{function!n-image of@$n$-image of}%
定义为
\begin{equation*}
\im_n(f)\defeq \sm{b:B} \trunc n{\hfib{f}b}.
\end{equation*}
当 $n=-1$ 时，我们简记 $\im(f)$ 并称其为 $f$ 的\define{像（image）}。
\end{defn}

\begin{lem}\label{prop:to_image_is_connected}
对任意函数 $f:A\to B$，典范映射 $\tilde{f}:A\to\im_n(f)$ 是 $n$-连通的。
因此，任意函数可分解为一个 $n$-连通映射后接一个 $n$-截断映射。
\end{lem}

\begin{proof}
注意 $A\eqvsym\sm{b:B}\hfib{f}b$。映射 $\tilde{f}$ 是由典范的逐纤维变换
\begin{equation*}
\prd{b:B} \Parens{\hfib{f}b\to\trunc n{\hfib{f}b}}.
\end{equation*}
诱导出的全空间上的映射。
由于每个映射 $\hfib{f}b\to\trunc n{\hfib{f}b}$ 根据 \cref{cor:totrunc-is-connected} 是 $n$-连通的，故由 \cref{prop:nconn_fiber_to_total} 知 $\tilde{f}$ 是 $n$-连通的。
最后，投影 $\proj1:\im_n(f) \to B$ 是 $n$-截断的，因为它的纤维等价于 $f$ 的纤维的 $n$-截断。
\end{proof}

在接下来的引理中，我们将为证明唯一分解定理准备一些工具。

\begin{lem}\label{prop:factor_equiv_fiber}
假设我们有函数的交换图
\begin{equation*}
  \xymatrix{
    {A} \ar[r]^{g_1} \ar[d]_{g_2} &
    {X_1} \ar[d]^{h_1} &
    \\
    {X_2} \ar[r]_{h_2}
    &
    {B}
  }
\end{equation*}
其中$H:h_1\circ g_1\htpy h_2\circ g_2$，且 $g_1$ 与 $g_2$ 是 $n$-连通的，而 $h_1$ 与 $h_2$ 是 $n$-截断的。
则对任意 $b:B$，存在一个等价
\begin{equation*}
E(H,b):\hfib{h_1}b\eqvsym\hfib{h_2}b
\end{equation*}
使得对任意 $a:A$，有一个相等证明
\[\overline{E}(H,a) :  E(H,h_1(g_1(a)))({g_1(a),\refl{h_1(g_1(a))}}) = \pairr{g_2(a),\opp{H(a)}}.\]
\end{lem}

\begin{proof}
任取 $b:B$。则我们有以下等价：
\begin{align}
\hfib{h_1}b
& \eqvsym \sm{w:\hfib{h_1}b} \trunc n{ \hfib{g_1}{\proj1 w}}
\tag{since $g_1$ is $n$-connected}\\
& \eqvsym \Trunc n{\sm{w:\hfib{h_1}b}\hfib{g_1}{\proj1 w}}
\tag{by \cref{thm:refl-over-ntype-base}, since $h_1$ is $n$-truncated}\\
& \eqvsym \trunc n{\hfib{h_1\circ g_1}b}
\tag{by \cref{ex:unstable-octahedron}}
\end{align}
对 $h_2$ 和 $g_2$ 同理。
此外，由于有同伦$H:h_1\circ g_1\htpy h_2\circ g_2$，存在一个明显的等价$\hfib{h_1\circ g_1}b\eqvsym\hfib{h_2\circ g_2}b$。
因此对任意 $b:B$，我们得到
\begin{equation*}
\hfib{h_1}b\eqvsym\hfib{h_2}b
\end{equation*}
通过分析底层的函数，我们得到元素
$\pairr{g_1(a),\refl{h_1(g_1(a))}}$ 在 $E$ 的各组成等价作用下变化的如下表示。
其中有一些相等证明是定义性的，但另一些（在下方用 $=$ 标出）仅是命题性的；将它们组合起来，便得到 $\overline E(H,a)$。
{\allowdisplaybreaks
\begin{align*}
\pairr{g_1(a),\refl{h_1(g_1(a))}} &
    \overset{=}{\mapsto} \Pairr{\pairr{g_1(a),\refl{h_1(g_1(a))}}, \tproj n{ \pairr{a,\refl{g_1(a)}}}}\\
  & \mapsto \tproj n { \pairr{\pairr{g_1(a),\refl{h_1(g_1(a))}}, \pairr{a,\refl{g_1(a)}} }}\\
  & \mapsto \tproj n { \pairr{a,\refl{h_1(g_1(a))}}}\\
  & \overset{=}{\mapsto} \tproj n { \pairr{a,\opp{H(a)}}}\\
  & \mapsto \tproj n { \pairr{\pairr{g_2(a),\opp{H(a)}},\pairr{a,\refl{g_2(a)}}} }\\
  & \mapsto \Pairr{\pairr{g_2(a),\opp{H(a)}}, \tproj n {\pairr{a,\refl{g_2(a)}}} }\\
  & \mapsto \pairr{g_2(a),\opp{H(a)}}
\end{align*}}
第一个等式成立是因为：对于一般的 $b$，映射
\narrowequation{ \hfib{h_1}b \to \sm{w:\hfib{h_1}b} \trunc n{ \hfib{g_1}{\proj1 w}} }
会插入$\trunc n{ \hfib{g_1}{\proj1 w}}$的收缩中心（由"$g_1$ 是 $n$-截断的"这一假设提供）；而在当前讨论的情形中，该类型有明显的元素$\tproj n{ \pairr{a,\refl{g_1(a)}}}$，由可缩性可知它必须等于那个收缩中心。
第二个命题性等式成立是因为等价$\hfib{h_1\circ g_1}b\eqvsym\hfib{h_2\circ g_2}b$将第二个分量与 $\opp{H(a)}$ 拼接，并且我们有$\opp{H(a)} \ct \refl{} = \opp{H(a)}$。
其余等式均为定义性的，读者可以自行验证（假定\cref{ex:unstable-octahedron}有一个合理的解决方案）。
\end{proof}

% The equivalences $E(H,b)$ are such that $E(H^{-1},b)= E(H,b)^{-1}$.

结合\cref{prop:to_image_is_connected,prop:factor_equiv_fiber}，我们得到如下唯一分解结果：

\begin{thm}\label{thm:orth-fact}
对每个 $f:A\to B$，由下式定义的空间 $\fact_n(f)$
\begin{equation*}
\sm{X:\type}{g:A\to X}{h:X\to B} (h\circ g\htpy f)\times\mathsf{conn}_n(g)\times\mathsf{trunc}_n(h)
\end{equation*}
是可缩的。
其收缩中心是来自\cref{prop:to_image_is_connected}的元素
\begin{equation*}
\pairr{\im_n(f),\tilde{f},\proj1,\theta,\varphi,\psi}:\fact_n(f)
\end{equation*}
其中$\theta:\proj1\circ\tilde{f}\htpy f$是典范同伦，$\varphi$ 是
\cref{prop:to_image_is_connected}的证明，而 $\psi$ 是"$\proj1:\im_n(f)\to B$具有 $n$-截断的纤维"这一事实的显然证明。
\end{thm}

\begin{proof}
根据\cref{prop:to_image_is_connected}我们知道 $\fact_n(f)$ 中存在元素，因此只需证明 $\fact_n(f)$ 是一个命题。
假设 $f$ 有两个 $n$-分解
\begin{equation*}
\pairr{X_1,g_1,h_1,H_1,\varphi_1,\psi_1}\qquad\text{and}\qquad\pairr{X_2,g_2,h_2,H_2,\varphi_2,\psi_2}
\end{equation*}
那么我们有点态拼接而成的同伦
\[ H\defeq (\lam{a} H_1(a) \ct H_2^{-1}(a)) \,:\, (h_1\circ g_1\htpy h_2\circ g_2).\]
根据泛等性以及道路与搬运在 $\Sigma$-类型、函数类型和道路类型中的刻画，只需证明以下四点：
\begin{enumerate}
\item 存在一个等价 $e:X_1\eqvsym X_2$，
\item 存在一个同伦$\zeta:e\circ g_1\htpy g_2$，
% \note{Is it easy enough to see that these elements are the various transports?}
\item 存在一个同伦$\eta:h_2\circ e\htpy h_1$，
\item 对任意 $a:A$ 我们有$\opp{\apfunc{h_2}(\zeta(a))} \ct \eta(g_1(a)) \ct H_1(a) = H_2(a)$。
\end{enumerate}
%where $\underline{e}$ is the function underlying the equivalence.
我们按顺序证明这四个断言。
\begin{enumerate}
\item 由\cref{prop:factor_equiv_fiber}，我们有一个逐纤维等价
% \note{It could be a nice exercise for the book to show
% that if $f_1:A_1\to B$ and $f_2:A_2\to B$ have equivalent fibers, then $A_1\eqvsym A_2$}.
\begin{equation*}
E(H) : \prd{b:B} \eqv{\hfib{h_1}b}{\hfib{h_2}b}.
\end{equation*}
这诱导了全空间的一个等价，亦即我们有
\begin{equation*}
\eqvspaced{\Parens{\sm{b:B} \hfib{h_1}b}}{\Parens{\sm{b:B}\hfib{h_2}b}}.
\end{equation*}
当然，由\cref{thm:total-space-of-the-fibers}我们还有等价$X_1\eqvsym\sm{b:B}\hfib{h_1}b$与$X_2\eqvsym\sm{b:B}
\hfib{h_2}b$。
由此我们得到等价 $e:X_1\eqvsym X_2$；读者可以验证 $e$ 的底层函数由下式给出：
\begin{equation*}
e(x) \jdeq \proj1(E(H,h_1(x))(x,\refl{h_1(x)})).
\end{equation*}
\item 由\cref{prop:factor_equiv_fiber}，我们可以选取
  $\zeta(a) \defeq \apfunc{\proj1}(\overline E(H,a)) : e(g_1(a)) = g_2(a)$。
  \label{item:orth-fact-2}
\item 对每个 $x:X_1$，我们有
\begin{equation*}
\proj2(E(H,h_1(x))({x,\refl{h_1(x)}})) :h_2(e(x))= h_1(x),
\end{equation*}
这给出了同伦$\eta:h_2\circ e\htpy h_1$。
\item 根据纤维中道路的刻画（\cref{lem:hfib}），来自\cref{prop:factor_equiv_fiber}的道路 $\overline E(H,a)$ 给出
  $\eta(g_1(a)) = \apfunc{h_2}(\zeta(a)) \ct \opp{H(a)}$。
  代入 $H$ 的定义并整理道路后，即得所求的等式。\qedhere
\end{enumerate}
\end{proof}

% I can't make sense of this, and it doesn't seem necessary
%
% \begin{cor}
% A function $f:A\to B$ is $n$-connected if and only if
% \begin{equation*}
% \prd C\prd{g:\modalfunc(B\to C)} \iscontr\big(\sm{h:\modalfunc(B\to C)}\underline{h}\circ
% f\htpy\underline{g}\circ f\big).
% \end{equation*}
% \end{cor}

通过标准论证，我们得到如下的正交性原理。

\begin{thm}
  设 $e:A\to B$ 是 $n$-连通的，且 $m:C\to D$ 是 $n$-截断的。
  则映射
  \[ \varphi: (B\to C) \;\to\; \sm{h:A\to C}{k:B\to D} (m\circ h \htpy k \circ e) \]
  是一个等价。
\end{thm}

\begin{proof}[证明概要]
  对于陪域中的任意 $(h,k,H)$，令 $h = h_2 \circ h_1$ 且 $k = k_2 \circ k_1$，其中 $h_1$ 与 $k_1$ 是 $n$-连通的，$h_2$ 与 $k_2$ 是 $n$-截断的。
  那么$f = (m\circ h_2) \circ h_1$与$f = k_2 \circ (k_1\circ e)$都是 $m \circ h = k\circ e$ 的 $n$-分解。
  因此，二者之间存在唯一的等价。
  由此不难得出$\hfib\varphi{(h,k,H)}$是可缩的（尽管略显繁琐）。
\end{proof}

\index{orthogonal factorization system|)}%
\index{unique!factorization system|)}%

最后，我们证明像在拉回下是稳定的。
\index{image!stability under pullback}
\index{factorization!stability under pullback}

\begin{lem}\label{lem:hfiber_wrt_pullback}
假设方块
\begin{equation*}
  \vcenter{\xymatrix{
      A\ar[r]\ar[d]_f &
      C\ar[d]^g\\
      B\ar[r]_-h &
      D
      }}
\end{equation*}
是一个拉回方块，并设 $b:B$。那么$\hfib{f}b\eqvsym\hfib{g}{h(b)}$。
\end{lem}

\begin{proof}
这可由拉回的粘贴性质（\cref{ex:pullback-pasting}）得出，因为下图中的类型 $X$
\begin{equation*}
  \vcenter{\xymatrix{
      X\ar[r]\ar[d] &
      A\ar[r]\ar[d]_f &
      C\ar[d]^g\\
      \unit\ar[r]_b &
      B\ar[r]_h &
      D
      }}
\end{equation*}
是左边方块的拉回当且仅当它是整个外层矩形的拉回，而 $\hfib{f}b$ 是左边方块的拉回，$\hfib{g}{h(b)}$ 是整个外层矩形的拉回。
\end{proof}

\begin{thm}\label{thm:stable-images}
\index{stability!of images under pullback}%
考虑函数 $f:A\to B$, $g:C\to D$ 以及图表
\begin{equation*}
  \vcenter{\xymatrix{
      A\ar[r]\ar[d]_{\tilde{f}_n} &
      C\ar[d]^{\tilde{g}_n}\\
      \im_n(f)\ar[r]\ar[d]_{\proj1} &
      \im_n(g)\ar[d]^{\proj1}\\
      B\ar[r]_h &
      D
      }}
\end{equation*}
如果外层矩形是拉回，那么底部的方块也是拉回（从而由\cref{ex:pullback-pasting}，顶部的方块也是拉回）。因此，像在拉回下是稳定的。
\end{thm}

\begin{proof}
假设外层方块是拉回，我们有等价
\begin{align*}
B\times_D\im_n(g) & \jdeq \sm{b:B}{w:\im_n(g)} h(b)=\proj1 w\\
& \eqvsym \sm{b:B}{d:D}{w:\trunc n{\hfib{g}d}} h(b)= d\\
& \eqvsym \sm{b:B} \trunc n{\hfib{g}{h(b)}}\\
& \eqvsym \sm{b:B} \trunc n{\hfib{f}b} &&
\text{(by \cref{lem:hfiber_wrt_pullback})}\\
& \equiv \im_n(f). && \qedhere
\end{align*}
\end{proof}

\index{n-type@$n$-type|)}%

\section{模态}
\label{sec:modalities}

\index{modality|(}

几乎所有关于 $n$-类型和连通性的理论都可以在更一般的框架下进行。
本节的内容在本书后续部分不会用到。

对于推广 $n$-类型理论，我们最初的想法可能是将\cref{thm:trunc-reflective}作为一个定义。

\begin{defn}\label{defn:reflective-subuniverse}
  一个\define{反射子宇宙（reflective subuniverse）}
  \indexdef{reflective!subuniverse}%
  \indexdef{subuniverse, reflective}%
  是一个谓词 $P:\type\to\prop$，满足对于每个类型 $A:\type$，我们都有一个类型 $\reflect A$ 使得 $P(\reflect A)$ 成立，以及一个映射
  $\project_A:A\to\reflect A$，并具有如下性质：对于每个满足 $P(B)$ 的类型 $B:\type$，下面的映射是一个等价：
  \[\function{(\reflect A\to{}B)}{(A\to{}B)}{f}{f\circ\project_A}.\]
\end{defn}

我们记$\P \defeq \setof{A:\type | P(A)}$，因此 $A:\P$ 表示 $A:\type$ 且 $P(A)$ 成立。
我们也将上述映射的拟逆记为 $\rec{\modal}$。
记号 $\reflect$ 可能看起来有点奇怪，但稍后我们会明白它的意义。

对于任意反射子宇宙，我们可以用通常的方式证明范畴论中关于反射子范畴的所有熟知结论。例如，我们有：

\begin{itemize}
\item 一个类型 $A$ 属于 $\P$ 当且仅当 $\project_A:A\to\reflect A$ 是一个等价。
\item $\P$ 在收缩核下封闭。
  特别地，只要 $\project_A$ 有一个收缩，$A$ 就属于 $\P$。
\item 运算 $\reflect$ 在某种合适的、同伦相干的意义下是一个函子；我们可以在需要的任何高阶层次上给出精确表述。
\item $\P$ 中的类型在所有极限（如积和拉回）下封闭。
  特别地，对于任意 $A:\P$ 和 $x,y:A$，恒等类型 $(x=_A y)$ 也属于 $\P$，因为它是两个函数 $\unit\to A$ 的拉回。
\item $\P$ 中的余极限可以通过将 $\reflect$ 作用于类型的普通余极限来构造。
\end{itemize}

重要的是，对积的封闭性也推广到"无穷积"，即依值函数类型。

\begin{thm}\label{thm:reflsubunv-forall}
  如果 $B:A\to\P$ 是反射子宇宙 \P 中任意一族类型，那么 $\prd{x:A} B(x)$ 也属于 \P。
\end{thm}

\begin{proof}
  对于任意 $x:A$，考虑由$\mathsf{ev}_x(f) \defeq f(x)$定义的函数$\mathsf{ev}_x : (\prd{x:A} B(x)) \to B(x)$。
  因为 $B(x)$ 属于 $P$，这个函数可以延拓为一个函数
  \[ \rec{\modal}(\mathsf{ev}_x) : \reflect\Parens{\prd{x:A} B(x)} \to B(x). \]
  因此我们可以通过$h(z)(x) \defeq \rec{\modal}(\mathsf{ev}_x)(z)$定义$h:\reflect(\prd{x:A} B(x)) \to \prd{x:A} B(x)$。
  那么 $h$ 就是$\project_{\prd{x:A} B(x)}$的一个收缩，从而${\prd{x:A} B(x)}$属于 $\P$。
\end{proof}

特别地，如果 $B:\P$ 且 $A$ 是任意类型，那么 $(A\to B)$ 属于 \P。
用范畴论的语言来说，这意味着任何反射子宇宙都是一个\define{指数理想（exponential ideal）}。
\indexdef{exponential ideal}%
进而，由标准论证可知反射子保持有限积。

\begin{cor}\label{cor:trunc_prod}
  对于任意类型 $A$ 和 $B$ 以及任意反射子宇宙，诱导的映射$\reflect(A\times B) \to \reflect(A) \times \reflect(B)$是一个等价。
\end{cor}

\begin{proof}
  只需证明 $\reflect(A) \times \reflect(B)$ 具有与 $\reflect(A\times B)$ 相同的泛性质。
  根据上述关于 $\P$ 中类型在极限下封闭的说明，$\reflect(A) \times \reflect(B)$ 属于 $\P$。
  现在令 $C:\P$；我们有
  \begin{align*}
    (\reflect(A) \times \reflect(B) \to C)
    &= (\reflect(A) \to (\reflect(B) \to C))\\
    &= (\reflect(A) \to (B \to C))\\
    &= (A \to (B \to C))\\
    &= (A \times B \to C)
  \end{align*}
  这里利用了 $\reflect(B)$ 和 $\reflect(A)$ 的泛性质，以及 $B\to C$ 属于 $\P$ 的事实（因为 $C$ 属于 $\P$）。
  可以直接验证这个等价是通过与 $\mreturn_A \times \mreturn_B$ 复合得到的，符合要求。
\end{proof}

类型的每个反射子范畴都自动成为指数理想，并且具有保持积的反射子，这可能显得奇怪。然而，在经典集合范畴中，出于同样的原因，情况也是如此。只是这一事实通常不被提及，因为经典的集合范畴——与同伦类型范畴相反——并没有许多有趣的反射子范畴。

$n$-类型的两个基本性质不被一般的反射子宇宙所具有：\cref{thm:ntypes-sigma}（在 $\Sigma$-类型下封闭）和\cref{thm:truncn-ind}（截断归纳）。然而，这两个性质的类似命题彼此等价。

\begin{thm}\label{thm:modal-char}
  对于一个反射子宇宙 \P，以下两条逻辑等价。
  \begin{enumerate}
  \item 如果 $A:\P$ 且 $B:A\to \P$，那么 $\sm{x:A} B(x)$ 属于 \P。\label{item:mchr1}
  \item 对于每个 $A:\type$，类型族 $B:\reflect A\to\P$，以及映射$g:\prd{a:A} B(\project(a))$，存在$f:\prd{z:\reflect A} B(z)$使得对所有 $a:A$ 有$f(\project(a)) = g(a)$。\label{item:mchr2}
  \end{enumerate}
\end{thm}

\begin{proof}
  假设\ref{item:mchr1}。
  那么在\ref{item:mchr2}所述的情形中，类型$\sm{z:\reflect A} B(z)$属于 $\P$，并且我们有由$g'(a)\defeq (\project(a),g(a))$定义的$g':A\to \sm{z:\reflect A} B(z)$。
  于是，我们得到$\rec{\modal}(g'):\reflect A \to \sm{z:\reflect A} B(z)$，使得$\rec{\modal}(g')(\project(a)) = (\project(a),g(a))$成立。

  现在考虑函数$\proj1 \circ \rec{\modal}(g') : \reflect A \to \reflect A$ 与 $\idfunc[\reflect A]$。
  根据假设，它们与 $\project$ 预复合后相等。
  因此，由 $\reflect$ 的泛性质，它们已经相等，即对所有 $z$ 我们有$p_z:\proj1(\rec{\modal}(g')(z)) = z$。
  现在我们可以定义
  %
  \narrowequation{f(z) \defeq \trans{p_z}{\proj2(\rec{\modal}(g')(z))},}
  %
  利用定义\ref{defn:reflective-subuniverse}给出的等价的伴随性质，可以证明
  %
  \narrowequation{\rec{\modal}(g')(\project(a)) = (\project(a),g(a))}
  %
  的第一分量等于 $p_{\project(a)}$。因此，它的第二分量给出所需的$f(\project(a)) = g(a)$。

  反过来，假设\ref{item:mchr2}，且 $A:\P$ 以及 $B:A\to\P$。
  令 $h$ 为复合
  \[ \reflect\Parens{\sm{x:A} B(x)} \xrightarrow{\reflect(\proj1)} \reflect A \xrightarrow{\opp{(\project_A)}} A. \]
  那么对于$z:\sm{x:A} B(x)$，我们有
  \begin{align*}
    h(\project(z)) &= \opp\project(\reflect(\proj1)(\project(z)))\\
    &= \opp\project(\project(\proj1(z)))\\
    &= \proj1(z).
  \end{align*}
  记这条道路为 $p_z$。
  现在如果我们将$C:\reflect(\sm{x:A} B(x)) \to \type$定义为$C(w) \defeq B(h(w))$，则有
  \[ g \defeq \lam{z} \trans{p_z}{\proj2(z)} \;:\; \prd{z:\sm{x:A} B(x)} C(\project(z)). \]
  于是，由假设可得
  %
  \narrowequation{f:\prd{w:\reflect(\sm{x:A}B(x))} C(w)}
  %
  使得$f(\project(z)) = g(z)$成立。
  $h$ 和 $f$ 共同给出一个函数
  %
  \narrowequation{k:\reflect(\sm{x:A}B(x)) \to \sm{x:A}B(x)}
  %
  由$k(w) \defeq (h(w),f(w))$定义，而 $p_z$ 与相等$f(\project(z)) = g(z)$表明 $k$ 是$\project_{\sm{x:A}B(x)}$的一个收缩。
  因此，$\sm{x:A}B(x)$ 属于 \P。
\end{proof}

注意这与\cref{sec:htpy-inductive}中的讨论相似。
\index{recursion principle!for a modality}%
\index{induction principle!for a modality}%
\index{uniqueness!principle, propositional!for a modality}%
反射子宇宙的反射子的泛性质类似于带有唯一性性质的递归原理，而\cref{thm:modal-char}\ref{item:mchr2}则类似于相应的归纳原理。
与\cref{sec:htpy-inductive}中不同，由于我们只能消去到属于 $\P$ 的类型中，二者在这里并不等价。
\cref{thm:modal-char} 的条件~\ref{item:mchr1} 正是用来弥合这一脱节的。

当然，毫不意外的是，如果我们有归纳原理，就可以推出递归原理。
只要允许消去到道路类型，我们也能推出其唯一性性质。
这引出了以下定义。
注意，任何反射子宇宙都可以由操作 $\reflect:\type\to\type$ 和函数 $\project_A:A\to \reflect A$ 来刻画，因为$P(A) = \isequiv(\project_A)$。

\begin{defn}\label{defn:modality}
一个 \define{模态（modality）}
\indexdef{modality}
是一个操作 $\modal:\type\to\type$，并附带以下数据：
\begin{enumerate}
\item 对每个类型 $A$，有一个函数 $\mreturn^\modal_A:A\to\modal(A)$。\label{item:modal1}
\item 对每个 $A:\type$ 以及每个类型族 $B:\modal(A)\to\type$，有一个函数\label{item:modal2}
\begin{equation*}
\ind{\modal}:\Parens{\prd{a:A}\modal(B(\mreturn^\modal_A(a)))}\to\prd{z:\modal(A)}\modal(B(z)).
\end{equation*}
\item 对每个 $f:\prd{a:A}\modal(B(\mreturn^\modal_A(a)))$，有一条道路 $\ind\modal(f)(\mreturn^\modal_A(a)) = f(a)$。\label{item:modal3}
\item 对任意 $z,z':\modal(A)$，函数 $\mreturn^\modal_{z=z'} : (z=z') \to \modal(z=z')$ 是一个等价。\label{item:modal4}
\end{enumerate}
如果 $\mreturn^\modal_A:A\to\modal(A)$ 是一个等价，则称 $A$ 对于 $\modal$ 是 \define{模态的}
\indexdef{modal!type}%
\indexdef{type!modal}%
，并记
\begin{equation}
  \modaltype\defeq\setof{X:\type | X \text{ is $\modal$-modal} }\label{eq:modaltype}
\end{equation}
为模态类型的类型。
\end{defn}

条件~\ref{item:modal2} 和~\ref{item:modal3} 与 \cref{thm:modal-char}\ref{item:mchr2} 非常相似，但改用 $\modal B(z)$ 来表述，而非假设 $B$ 取值于 $\P$。
这使得我们可以纯粹用操作 $\modal$ 来陈述条件，而无需预先给定谓词 $P:\type\to\prop$。
（这并非完全令人满意，因为在条款~\ref{item:modal4}中我们仍然不太隐讳地用到了 $P$。
我们不知道~\ref{item:modal4} 是否能从~\ref{item:modal1}--\ref{item:modal3} 推出。）
然而，\cref{thm:modal-char}\ref{item:mchr2} 这个看起来更强的性质可以从 \cref{defn:modality}\ref{item:modal2} 及~\ref{item:modal3} 推出，因为对于任何 $C:\modal A \to \modaltype$ 都有 $C(z) \eqvsym \modal C(z)$，并且我们可以借助这个等价来回转换。

\index{universal!property!of a modality}%
如同其他归纳原理一样，这蕴含了一个泛性质。

\begin{thm}\label{prop:lv_n_deptype_sec_equiv_by_precomp}
设 $A$ 是一个类型，且 $B:\modal(A)\to\modaltype$。则函数
\begin{equation*}
(\blank\circ \mreturn^\modal_A) : \Parens{\prd{z:\modal(A)}B(z)} \to \Parens{\prd{a:A}B(\mreturn^\modal_A(a))}
\end{equation*}
是一个等价。
\end{thm}

\begin{proof}
由定义，操作 $\ind{\modal}$ 是 $(\blank\circ \mreturn^\modal_A)$ 的一个右逆。
因此，我们只需对每个 $s:\prd{z:\modal(A)}B(z)$ 找到一个同伦
\begin{equation*}
\prd{z:\modal(A)}s(z)= \ind{\modal}(s\circ \mreturn^\modal_A)(z)
\end{equation*}
来证明它同时也是一个左逆。
根据假设，每个 $B(z)$ 都是模态的，因而每个类型 $s(z)= R^\modal_X(s\circ \mreturn^\modal_A)(z)$
也是模态的。
于是，只需找到一个类型为
\begin{equation*}
\prd{a:A}s(\mreturn^\modal_A(a))= \ind{\modal}(s\circ \mreturn^\modal_A)(\mreturn^\modal_A(a))
\end{equation*}
的函数即可，而这可由 \cref{defn:modality}\ref{item:modal3} 推出。
\end{proof}

特别地，对每个类型 $A$ 和每个模态类型 $B$，我们都有一个等价 $(\modal A\to B)\eqvsym (A\to B)$。

\begin{cor}
  对于任意模态 $\modal$，$\modal$-模态类型构成一个满足 \cref{thm:modal-char} 的等价条件的反射子宇宙。
\end{cor}

因此，模态可以等同于对 $\Sigma$ 类型封闭的反射子宇宙。
名称 \emph{模态（modality）} 当然来源于 \emph{模态逻辑（modal logic）}\index{modal!logic}，后者研究的是一种逻辑，在其中我们可以形成诸如"可能 $A$"（通常写作 $\diamond A$）或"必然 $A$"（通常写作 $\Box A$）这类陈述。
符号 $\modal$ 通常用于表示任意的模态算子\index{modal!operator}。% (rather than a specific one such as $\diamond$ or $\Box$).
在命题即类型的原则下，模态逻辑意义上的模态对应着对\emph{类型}的一种操作，而 \cref{defn:modality} 似乎是定义此类操作的一种合理候选方案。
（更准确地说，我们或许应该称这些模态为\emph{幂等的、单子的（idempotent, monadic）}模态；见注记。）
\index{idempotent!modality}%
如 \cref{subsec:when-trunc} 所述，我们通常可以使用副词\index{adverb}来非正式地谈论这些模态，例如用"仅仅"\index{merely}表示命题截断，用"纯粹"\index{purely}表示恒等模态
\index{identity!modality}%
\index{modality!identity}%
（即由 $\modal A \defeq A$ 定义的那个模态）。

对于任意模态 $\modal$，我们称映射 $f:A\to B$ 是 \define{$\modal$-连通的（$\modal$-connected）}
\indexdef{function!.circle-connected@$\modal$-connected}%
\indexdef{.circle-connected function@$\modal$-connected function}%
，若对所有 $b:B$，$\modal(\hfib f b)$ 都是可缩的；称其为 \define{$\modal$-截断的（$\modal$-truncated）}
\indexdef{function!.circle-truncated@$\modal$-truncated}%
\indexdef{.circle-truncated function@$\modal$-truncated function}%
，若对所有 $b:B$，$\hfib f b$ 都是模态的。
所有来自 \cref{sec:connectivity,sec:image-factorization} 的理论，只要不涉及关联不同 $n$ 值的 $n$-类型，都可以逐字逐句地适用于这一一般情形。
\index{orthogonal factorization system}%
\index{unique!factorization system}%
特别地，我们有一个正交分解系统。

一类不包含 $n$-截断的重要模态是\emph{左正合（left exact）}模态：即那些其函子 $\modal$ 既保持拉回又保持有限积的模态。
\index{topology!Lawvere-Tierney}%
它们是初等意象\index{topos}论中"Lawvere–Tierney\index{Lawvere}\index{Tierney} 拓扑"的范畴化，并在高阶范畴语义中对应于子 $(\infty,1)$-意象。
\index{.infinity1-topos@$(\infty,1)$-topos}%
然而，这已超出了本书的范围。

除 $n$-截断外，其他模态的一些具体例子可以在习题中找到。

\index{modality|)}

\sectionNotes

同伦 $n$-类型这一概念在经典同伦论中由来已久。
是 Voevodsky 认识到，这个概念可以在同伦类型论中递归地定义，从可缩性出发即可。

\index{axiom!Streicher's Axiom K}%
"Axiom K"这一性质由 Thomas Streicher 如此命名，是因为它是继 J 之后相等类型的一个性质——后者是相等类型消去子的传统名称。
\cref{thm:hedberg} 归功于 Hedberg~\cite{hedberg1998coherence}；\cite{krausgeneralizations} 则包含了更多信息与推广。

$n$-连通空间与函数的概念在同伦论中同样经典，不过如前所述，我们对函数连通性的编号与经典编号相差 1。
由此产生的分解系统的重要性，已被 Rezk、Lurie 及其他学者在高阶意象理论中的近期工作所强调。%
\index{.infinity1-topos@$(\infty,1)$-topos}
特别地，本章的结果应与~\cite[\S6.5.1]{lurie:higher-topoi} 相比较。
在 \cref{sec:freudenthal} 中，$n$-连通映射的理论将对我们证明 Freudenthal 纬悬定理起到关键作用。

对\emph{简单}类型论中的模态算子\index{modal!operator}已有大量研究；例如可见~\cite{modalTT}。而在依值类型论的框架下，\cite{ab:bracket-types} 将命题截断（$(-1)$-截断）这一特例作为模态算子\index{modal!operator}来处理。这里所展示的发展极大地扩展并推广了上述工作，同时也借鉴了意象理论\index{topos}中的思想。

一般来说，模态算子\index{modal!operator}至少可以分为两种：一种是像 $\diamond$（“可能”）那样满足 $A\Rightarrow \diamond A$ 的，另一种是像 $\Box$（“必然”）那样满足 $\Box A \Rightarrow A$ 的。
当它们还是\emph{幂等的}（即 $\diamond A = \diamond{\diamond A}$ 或 $\Box A = \Box{\Box A}$）时，前者可以等同于反射子范畴（或等价地，幂等单子），而后者则等同于余反射子范畴（或幂等余单子）。
\index{monad}
\index{comonad}
然而，在依值类型论中，处理余单子那一类要棘手一些，因为它们在拉回之下更少能保持稳定，从而无法被解释为宇宙 \UU 上的运算。
虽然有时有办法绕过这一点（例如参见\cite{QGFTinCHoTT12}），但为了简单起见，这里我们只讨论单子那一类。

在计算方面，单子（从而模态\index{modality}）被用来为函数式编程中的计算效应建模~\cite{Moggi89}。%
\index{programming}%
\index{computational effect}
如果一个计算的执行不会导致副作用（例如在屏幕上打印消息、播放音乐或通过互联网发送数据），就称该计算是\emph{纯}的。
存在像 Haskell\index{Haskell} 这样的“纯函数式”编程语言，在技术上，其中只能编写纯函数：副作用是通过对输出类型应用“单子”来表示的。
例如，类型为 $\mathsf{Int}\to\mathsf{Int}$ 的函数是纯的，而类型为 $\mathsf{Int}\to \mathsf{IO}(\mathsf{Int})$ 的函数在计算结果的过程中可能会执行输入和输出；运算 $\mathsf{IO}$ 就是一个单子。
\index{purely}%
（这就是我们将副词 “purely” 用于恒等单子的原因，因为后者在计算上对应于没有副作用的纯函数。）
本章考虑的所有模态都是幂等的，而函数式编程中使用的模态很少是幂等的，但两者的思想仍然密切相关。

\sectionExercises

\begin{ex}\label{ex:all-types-sets}\
  \begin{enumerate}
    \item 利用 \cref{thm:h-set-refrel-in-paths-sets} 证明：如果对每个类型 $A$ 都有 $\brck{A}\to A$，那么每个类型都是集合。
    \item 证明：如果每个满射函数都（纯粹地）分裂，即如果对每个 $f:A\to B$ 都有
    %
    \narrowequation{\prd{b:B}\brck{\hfib{f}{b}}\to\prd{b:B}\hfib{f}{b}},
    %
    那么每个类型都是集合。
  \end{enumerate}
\end{ex}

\begin{ex}\label{ex:s2-colim-unit}
  对于本题，我们考虑以下关于余极限的一般概念。
  定义一个\define{图（graph）}\indexdef{graph} $\Gamma$，它由一个类型 $\Gamma_0$ 和一个族 $\Gamma_1 : \Gamma_0 \to \Gamma_0 \to \UU$ 组成。
  在一个图 $\Gamma$ 上的\define{图表（diagram）}\index{diagram}（指类型的图表）包含一个族 $F:\Gamma_0 \to \UU$，以及对每对 $x,y:\Gamma_0$，一个函数 $F_{x,y}:\Gamma_1(x,y) \to F(x) \to F(y)$。
  该图表的\define{余极限（colimit）}\index{colimit!of types} 是由以下构造子生成的高阶归纳类型 $\colim(F)$：
  \begin{itemize}
  \item 对每个 $x:\Gamma_0$，一个函数 $\inc_x:F(x) \to \colim(F)$，以及
  \item 对每个 $x,y:\Gamma_0$、每个 $\gamma:\Gamma_1(x,y)$ 以及每个 $a:F(x)$，一条道路 $\inc_y(F_{x,y}(\gamma,a)) = \inc_x(a)$。
  \end{itemize}
  也存在更一般的余极限（例如参见\cref{ex:s2-colim-unit-2}），但这对于许多目的而言已经足够。
  \begin{enumerate}
  \item 给出一个图 $\Gamma$，使得 $\Gamma$-图表的余极限可以等同于 \cref{sec:colimits} 中定义的推出。
    换言之，每个伸展应该诱导出一个 $\Gamma$-图表，其余极限就是该伸展的推出。
  \item 给出一个图 $\Gamma$ 及其上的图表 $F$，使得对所有 $x$ 有 $F(x)=\unit$，但同时有 $\colim(F)=\Sn^2$。
    注意 $\unit$ 是一个 $(-2)$-类型，而 $\Sn^2$ 预期不是任何有限 $n$ 对应的 $n$-类型。
    另见 \cref{ex:s2-colim-unit-2}。
  \end{enumerate}
\end{ex}

\begin{ex}\label{ex:ntypes-closed-under-wtypes}
  证明：如果 $A$ 是一个 $n$-类型且 $B:A\to \ntype{n}$ 是一个 $n$-类型的族，其中 $n\ge -1$，那么 $W$-类型 $\wtype{a:A} B(a)$（参见 \cref{sec:w-types}）也是一个 $n$-类型。
\end{ex}

\begin{ex}\label{ex:connected-pointed-all-section-retraction}
  利用 \cref{prop:nconn_fiber_to_total} 将 \cref{thm:connected-pointed} 推广到任意的截面-收缩对。
\end{ex}

\begin{ex}\label{ex:ntype-from-nconn-const}
  证明 \cref{thm:nconn-to-ntype-const} 的特征刻画在另一个方向上也成立：$B$ 是 $n$-类型，当且仅当每个从 $n$-连通类型到 $B$ 的映射都是常值的。
  理想情况下，你的证明应该能像在 \cref{sec:modalities} 中那样，适用于任何模态。
\end{ex}

\begin{ex}\label{ex:connectivity-inductively}
  证明：对于 $n\ge -1$，一个类型 $A$ 是 $n$-连通的，当且仅当它仅知有实例，并且对于所有 $a,b:A$，类型 $\id[A]ab$ 是 $(n-1)$-连通的。
  因此，由于每个类型都是 $(-2)$-连通的，类型的 $n$-连通性可以仅使用命题截断来归纳地定义。
  （特别地，$A$ 是 $0$-连通的当且仅当 $\brck{A}$ 且 $\prd{a,b:A} \brck{a=b}$。）
\end{ex}

\begin{ex}\label{ex:lemnm}
  \indexdef{excluded middle!LEMnm@$\LEM{n,m}$}%
  对于 $-1\le n,m \le\infty$，记 $\LEM{n,m}$ 为如下命题
  \[ \prd{A:\ntype{n}} \trunc m{A + \neg A},\]
  这里将 $\ntype{\infty} \defeq \type$，且 $\trunc{\infty}{X}\defeq X$。
  证明：
  \begin{enumerate}
  \item 若 $n=-1$ 或 $m=-1$，则 $\LEM{n,m}$ 等价于 \cref{sec:intuitionism} 中的 $\LEM{}$。
  \item 若 $n\ge 0$ 且 $m\ge 0$，则 $\LEM{n,m}$ 与泛等性矛盾。
  \end{enumerate}
\end{ex}

\begin{ex}\label{ex:acnm}
  \indexdef{axiom!of choice!ACnm@$\choice{n,m}$}%
  对于 $-1\le n,m\le\infty$，记 $\choice{n,m}$ 为如下命题
  \[ \prd{X:\set}{Y:X\to\ntype{n}}
  \Parens{\prd{x:X} \trunc m{Y(x)}}
  \to
  \Trunc m{\prd{x:X} Y(x)},
  \]
  约定同 \cref{ex:lemnm}。
  于是 $\choice{0,-1}$ 就是 \cref{sec:axiom-choice} 中的选择公理，
  而 $\choice{\infty,\infty}$ 就是恒等函数。
  （如果我们按 \eqref{eq:ac} 而非 \eqref{eq:epis-split} 的方式类似地定义 $\choice{n,m}$，
  那么 $\choice{\infty,\infty}$ 就会类似于 \cref{thm:ttac}。）
  已知 $\choice{\infty,-1}$ 与泛等性一致，因为它在 Voevodsky 的单纯集模型中成立。
  \begin{enumerate}
  \item 不使用泛等性，证明对任意 $m$，$\LEM{n,\infty}$ 蕴含 $\choice{n,m}$。
    （另一方面，在 \cref{subsec:emacinsets} 中我们将证明 $\choice{}=\choice{0,-1}$ 蕴含 $\LEM{}=\LEM{-1,-1}$。）
  \item 显然，如果 $k\le n$，则 $\choice{n,m}\Rightarrow \choice{k,m}$。
    那么这些原理 $\choice{n,m}$ 之间还有其他蕴含关系吗？
    对于任意 $m\ge 0$ 和任意 $n$，$\choice{n,m}$ 与泛等性一致吗？
    （这些都是开放问题。）\index{open!problem}
  \end{enumerate}
\end{ex}

\begin{ex}\label{ex:acnm-surjset}
  证明 $\choice{n,-1}$ 蕴含以下结论：对任意 $n$-类型 $A$，都仅存在一个集合 $B$ 和一个满射 $B\to A$。
\end{ex}

\begin{ex}\label{ex:acconn}
  \define{$n$-连通选择公理（$n$-connected axiom of choice）}
  \indexdef{n-connected@$n$-connected!axiom of choice}%
  \indexdef{axiom!of choice!n-connected@$n$-connected}%
  指的是如下命题：
  \begin{quote}
    若 $X$ 是一个集合，且 $Y:X\to \type$ 是一个类型族，使得每个 $Y(x)$ 都是 $n$-连通的，则 $\prd{x:X} Y(x)$ 也是 $n$-连通的。
  \end{quote}
  注意 $(-1)$-连通选择公理就是 \cref{ex:acnm} 中的 $\choice{\infty,-1}$。
  \begin{enumerate}
  \item 证明对所有 $n\ge -1$，$(-1)$-连通选择公理蕴含 $n$-连通选择公理。
  \item $n$-连通选择公理与原理 $\choice{n,m}$ 之间还有其他的蕴含关系吗？
    （这是一个开放问题。）\index{open!problem}
  \end{enumerate}
\end{ex}

\begin{ex}\label{ex:n-truncation-not-left-exact}
  证明对任意 $n\ge -1$，$n$-截断模态都不是左正合的。
  也就是说，请给出一个拉回，使得该模态未能保持它。
\end{ex}

\begin{ex}\label{ex:double-negation-modality}
  证明 $X\mapsto (\neg\neg X)$ 是一个模态。\index{modal!operator}%
\end{ex}

\begin{ex}\label{ex:prop-modalities}
  设 $P$ 是一个命题。
  \begin{enumerate}
  \item 证明 $X\mapsto (P\to X)$ 是一个左正合模态。
    这称为与 $P$ 关联的\define{开模态（open modality）}
    \indexdef{open!modality}%
    \indexdef{modality!open}%
    。
  \item 证明 $X\mapsto P*X$ 是一个左正合模态，其中 $*$ 表示联结（见 \cref{sec:colimits}）。
    这称为与 $P$ 关联的\define{闭模态（closed modality）}
    \indexdef{closed!modality}%
    \indexdef{modality!closed}%
    。
  \end{enumerate}
\end{ex}

\begin{ex}\label{ex:f-local-type}
  设 $f:A\to B$ 是一个映射；如果 $(\blank\circ f):(B\to Z) \to (A\to Z)$ 是一个等价，则称类型 $Z$ 是\define{$f$-局部的（$f$-local）}
  \indexdef{f-local type@$f$-local type}%
  \indexdef{type!f-local@$f$-local}%
  。
  \begin{enumerate}
  \item 证明 $f$-局部类型构成一个反射子宇宙。
    你需要使用一个高阶归纳类型来定义反射子（局部化）。
  \item 证明如果 $B=\unit$，那么这个子宇宙是一个模态。
  \end{enumerate}
\end{ex}

\begin{ex}\label{ex:trunc-spokes-no-hub}
  证明：与 \cref{rmk:spokes-no-hub} 不同，我们可以等价地将 $\trunc nA$ 定义为由一个函数 $\tprojf n : A \to \trunc n A$ 以及对于每个 $r:\Sn^{n+1} \to \trunc n A$ 和每个 $x:\Sn^{n+1}$ 的一条路径 $s_r(x):r(x) = r(\base)$ 所生成。
\end{ex}

\begin{ex}\label{ex:s2-colim-unit-2}
  本练习中，我们考虑一种比 \cref{ex:s2-colim-unit} 中稍复杂一些的余极限概念。
  定义\define{带复合的图（graph with composition）}\indexdef{graph!with composition} $\Gamma$ 为如 \cref{ex:s2-colim-unit} 所述的图，并附带对每个 $x,y,z:\Gamma_0$ 的一个函数 $\Gamma_1(y,z) \to \Gamma_1(x,y) \to \Gamma_1(x,z)$，记作 $\delta\mapsto\gamma \mapsto \delta \circ \gamma$。
  （例如，任何如 \cref{cha:category-theory} 所述的预范畴都是带复合的图。）
  带复合的图 $\Gamma$ 上的\define{图表（diagram）}\index{diagram} $F$ 由底层图上的图表，加上对每个 $x,y,z:\Gamma_0$、$\gamma:\Gamma_1(x,y)$ 和 $\delta:\Gamma_1(y,z)$ 的一个同伦 $\cmp_{x,y,z}(\delta,\gamma) : F_{y,z}(\delta) \circ F_{x,y}(\gamma) \htpy F_{x,z}(\delta\circ\gamma)$ 构成。
  此类图表的\define{余极限（colimit）}\index{colimit!of types} 是由以下构造子生成的高阶归纳类型 $\colim(F)$：
  \begin{itemize}
  \item 对每个 $x:\Gamma_0$，有一个函数 $\inc_x:F(x) \to \colim(F)$，
  \item 对每个 $x,y:\Gamma_0$、$\gamma:\Gamma_1(x,y)$ 和 $a:F(x)$，有一条道路 $\glue_{x,y}(\gamma,a) : \inc_y(F_{x,y}(\gamma,a)) = \inc_x(a)$，以及
  \item 对每个 $x,y,z:\Gamma_0$、$\gamma:\Gamma_1(x,y)$、$\delta:\Gamma_1(y,z)$ 和 $a:F(x)$，有一条道路
    \[ \ap{\inc_z}{\cmp_{x,y,z}(\delta,\gamma,a)} \ct \glue_{x,z}(\delta\circ \gamma,a) = \glue_{y,z}(\delta,F_{x,y}(\gamma,a)) \ct \glue_{x,y}(\gamma,a). \]
  \end{itemize}
  （这是对完全同伦论意义上的图表和余极限概念的一种“二阶近似”；后者在所有更高层级上均应包含此类“相干性路径”。
  在类型论中定义此类概念是一个重要的公开问题。）

  试构造一个带复合的图 $\Gamma$，使得 $\Gamma_0$ 是集合且每个类型 $\Gamma_1(x,y)$ 都是命题，并构造其上的图表 $F$，使得对所有 $x$ 均有 $F(x)=\unit$，且 $\colim(F)=\Sn^2$。
\end{ex}

\begin{ex}\label{ex:fiber-map-not-conn}
  对比 \cref{lem:nconnected_postcomp_variation,prop:nconn_fiber_to_total}，人们可能会猜想：若 $f:A\to B$ 是 $n$-连通的，且 $g:\prd{a:A} P(a) \to Q(f(a))$ 诱导出 $n$-连通映射 $\Parens{\sm{a:A} P(a)} \to \Parens{\sm{b:B} Q(b)}$，则 $g$ 是逐纤维 $n$-连通的。
  请给出一个反例以说明此猜想不成立。
  （事实上，推广到模态时，这一性质刻画了左正合模态；参见 \cref{ex:prop-modalities}。）
\end{ex}

\begin{ex}\label{ex:is-conn-trunc-functor}
  证明：若 $f : A \to B$ 是 $n$-连通的，则 $\trunc kf : \trunc kA \to \trunc kB$ 也是 $n$-连通的。
\end{ex}

\begin{ex}\label{ex:categorical-connectedness}
  我们称一个类型 $A$ 是\define{范畴意义下连通的（categorically connected）}
  \indexdef{connected!categorically}，如果对于任意类型 $B, C$，由
  \begin{align*}
    e_{A,B,C}(\inl(g)) &\defeq \lam{x} \inl(g(x)),\\
    e_{A,B,C}(\inr(g)) &\defeq \lam{x} \inr(g(x))
  \end{align*}
  定义的典范映射 $e_{A, B, C}:((A\to B) + (A\to C)) \to (A \to B + C)$ 是一个等价。
  \begin{enumerate}
  \item 证明任何连通的类型都是范畴意义下连通的。
  \item 证明所有范畴意义下连通的类型都是连通的，当且仅当 $\LEM{}$ 成立。（提示：考虑 $A \defeq \Sigma P$，其中 $\neg \neg P$ 成立。）
  \end{enumerate}
\end{ex}

%%% Local Variables:
%%% mode: latex
%%% TeX-master: "hott-online"
%%% End:

\cleartooddpage[\thispagestyle{empty}] % Needed for correct TOC
\part{数学篇}
\label{part:mathematics}

\chapter{同伦论}
\label{cha:homotopy}

\index{acceptance|(}

在本章中，我们将在类型论内部发展一些同伦论。
我们采用 \cref{cha:basics} 中引入的同伦论的\emph{综合方法}：
空间、点、道路和同伦是基本概念，它们由类型及类型的元素来表示，特别是恒等类型。
道路与同伦的代数结构则由类型上自然的 $\infty$-群胚
\index{.infinity-groupoid@$\infty$-groupoid}%
结构所表示，这一结构由恒等类型的规则生成。
利用 \cref{cha:hits} 中引入的高阶归纳类型，我们可以通过万有性质直接描述空间。

\index{synthetic mathematics}%
这种综合方法有几个有趣的方面。
首先，它结合了具体模型（如拓扑空间\index{topological!space}或单纯集\index{simplicial!sets}）的优势与同伦论抽象范畴框架（如 Quillen 模型范畴\index{Quillen model category}）的优势。
一方面，我们的证明直观上很初等，具体地涉及类型中的点、道路和同伦。
另一方面，我们的方法又并不依赖于这些对象的任何具体呈现方式——例如，道路拼接的结合律是通过道路归纳证明的，而不是通过对映射 $[0,1] \to X$ 的重参数化或通过 horn 填充条件来证明。
类型论可以说是研究 $\infty$-群胚的抽象同伦论的一种非常便利的方式：
通过运用恒等类型的规则，我们可以避开许多 $\infty$-群胚定义中涉及的复杂组合学，并且在任何特定的证明中只需阐明所需的那部分结构。

类型论的抽象性意味着我们的证明能自动适用于多种不同的设定。
具体而言，如前所述，同伦类型论在 Kan\index{Kan complex} 单纯集\index{simplicial!sets}中有一种解释，而 Kan 单纯集是 $\infty$-群胚同伦论的一个模型。
因此，我们的证明适用于这个模型，并且通过沿着从单纯集到拓扑空间的几何实现\index{geometric realization}函子做转移，就能得到经典同伦论中相应定理的证明。
然而，尽管细节方面的工作仍在进行中，我们还可以在多种其他类似于 $\infty$-群胚范畴的范畴中解释类型论，例如 $(\infty,1)$-意象\index{.infinity1-topos@$(\infty,1)$-topos}。
因此，在类型论中证明一个结果，也就表明它在这些设定中同样成立。
这种额外的普适性是普通范畴逻辑的已知特性：泛等基础将它也扩展到了同伦论。

其次，我们的综合方法还启发了一些新的类型论方法与证明。
我们的一些证明是经典证明相当直接的转写。
另一些则更具类型论的特色，主要由 $\infty$-群胚运算的计算构成，其风格与计算机科学家用类型论推理计算机程序的方式非常相似。
使这些新证明成为可能的一个原因似乎是，类型论强调了同伦论中不同于其他方法的侧面：
尽管在 Kan\index{Kan complex} 单纯集这样的设定中，像道路归纳和高阶归纳类型的万有性质这样的工具也是可用的，但类型论提升了它们的重要性，因为它们是可用于处理这些类型的\emph{唯一}原始工具。
专注于这些工具，已经导出了对诸如圆周的万有覆叠和 Hopf 纤维化等熟悉构造的新描述，而这些描述仅仅使用了高阶归纳类型的递归原理。
这些描述非常直接，本章的许多证明都涉及对此类纤维化的计算。
\index{mathematics!constructive}%
我们证明的另一个新特点是它们是构造性的（假设泛等性与高阶归纳类型是构造性的）；
我们将在 \cref{sec:moreresults} 中描述它对于球面同伦群的一个应用。

\index{mathematics!formalized}%
\indexsee{formalization of mathematics}{mathematics, formalized}%
第三，我们的综合方法非常适合利用 \Coq 和 \Agda 等\index{proof!assistant}证明助理进行计算机验证。
本章所描述的证明，几乎全部都已通过计算机验证；其中许多证明最初是在证明助理中给出的，而后才为本书``去形式化''。
计算机验证证明的篇幅和工作量与这里展示的非形式证明不相上下，有时甚至更短、更容易完成。

\mentalpause

在着手展示我们的结果之前，我们为不熟悉同伦论的读者简要复习一些基本概念和定理。
同时，我们也概览一下本章将要证明的结果。

\index{classical!homotopy theory|(}%
同伦论是代数拓扑的一个分支，运用抽象代数的工具（如群论）来研究空间的性质。
同伦论研究者探讨的问题之一是如何判断两个空间是否``相同''，这里的``相同''指的是\emph{同伦等价}
\index{homotopy!equivalence!topological}%
（存在来回的连续映射能在同伦的意义下复合为恒等映射——这为``修正''那些不能严格复合为恒等的映射提供了可能）。
判断两个空间是否相同的常用方法，是计算与空间相关的\emph{代数不变量}，其中包括其\emph{同伦群}、\emph{同调群}与\emph{上同调群}。
\index{homology}%
\index{cohomology}%
等价的空间具有同构的同伦群和（上）同调群，因此，若两个空间的这些群不同，它们便不等价。
于是，这些代数不变量提供了关于空间的整体信息，可以用来区分空间，与连续性等概念所提供的局部信息形成互补。
例如，环面局部上看与 $2$-球面相似，但它有一个整体上的差异，因为它中间有一个洞，而这个差异在这两个空间的同伦群中是可见的。

同伦群最简单的例子是空间的\emph{基本群}
\index{fundamental!group}%
，记为 $\pi_1(X,x_0$)：给定空间 $X$ 及其中的一个点 $x_0$，可以构成一个群，其元素是 $x_0$ 处的环路（即从 $x_0$ 到 $x_0$ 的连续道路），并在同伦的意义下视作相同；群的运算由恒等道路（静止不动）、道路拼接和道路反转给出。
例如，$2$-球面的基本群是平凡的，而环面的基本群则不然，这表明球面和环面不是同伦等价的。
直观上可以这样理解：球面上的每一个环路都同伦于恒等道路，因为它内部可以被``填满''。
相比之下，环面上穿过甜甜圈那个孔洞的环路则无法同伦于恒等道路，因此其基本群中存在非平凡的元素。

\index{homotopy!group}
\emph{高阶同伦群}提供了关于空间的更多信息。
固定空间 $X$ 中的一点 $x_0$，考虑恒等道路 $\refl{x_0}$。
那么，$\refl{x_0}$ 与自身之间的同伦的同伦类构成一个群 $\pi_2(X,x_0)$，它揭示了空间某些二维结构的信息。
进而，$\pi_3(X,x_0$) 是同伦之间的同伦的同伦类构成的群，依此类推。
代数拓扑的基本问题之一就是\emph{计算空间 $X$ 的同伦群}，这意味着要给出 $\pi_k(X,x_0)$ 与某个更直接的群描述（例如通过乘法表或群表示给出的群）之间的群同构。
出人意料的是，这是个非常困难的问题，即使对球面这样简单的空间也是如此。
正如从 \cref{tab:homotopy-groups-of-spheres} 表中可以看到的那样，球面的高阶同伦群中会出现一些模式，但目前尚不存在通用公式，并且许多球面同伦群至今仍是未知的。

\bgroup
% Colors for the table of homotopy groups of spheres
\definecolor{xA}{\OPTcolormodel}{\OPTcolxA}
\definecolor{xB}{\OPTcolormodel}{\OPTcolxB}
\definecolor{xC}{\OPTcolormodel}{\OPTcolxC}
\definecolor{xD}{\OPTcolormodel}{\OPTcolxD}
\definecolor{xE}{\OPTcolormodel}{\OPTcolxE}
\definecolor{xF}{\OPTcolormodel}{\OPTcolxF}
\definecolor{xG}{\OPTcolormodel}{\OPTcolxG}
\definecolor{xH}{\OPTcolormodel}{\OPTcolxH}
\definecolor{xI}{\OPTcolormodel}{\OPTcolxI}
\definecolor{xJ}{\OPTcolormodel}{\OPTcolxJ}
\definecolor{xK}{\OPTcolormodel}{\OPTcolxK}
\definecolor{xL}{\OPTcolormodel}{\OPTcolxL}
\definecolor{xM}{\OPTcolormodel}{\OPTcolxM}

\newcommand{\cA}{\colorbox{xG}{\hbox to 20pt {\hfil$0$\hfil}}}
\newcommand{\cB}{\colorbox{xH}{\hbox to 20pt {\hfil$\Z$\hfil}}}
\newcommand{\cC}{\colorbox{xI}{\hbox to 20pt {\hfil$\Z_{2}$\hfil}}}
\newcommand{\cD}{\colorbox{xJ}{\hbox to 20pt {\hfil$\Z_{2}$\hfil}}}
\newcommand{\cE}{\colorbox{xK}{\hbox to 20pt {\hfil$\Z_{24}$\hfil}}}
\newcommand{\cF}{\colorbox{xL}{\hbox to 20pt {\hfil$0$\hfil}}}
\newcommand{\cG}{\colorbox{xM}{\hbox to 20pt {\hfil$0$\hfil}}}
\newcommand{\cH}{\colorbox{xF}{\hbox to 20pt {\hfil$0$\hfil}}}
\newcommand{\cI}{\colorbox{xE}{\hbox to 20pt {\hfil$0$\hfil}}}
\newcommand{\cJ}{\colorbox{xD}{\hbox to 20pt {\hfil$0$\hfil}}}
\newcommand{\cK}{\colorbox{xC}{\hbox to 20pt {\hfil$0$\hfil}}}
\newcommand{\cL}{\colorbox{xB}{\hbox to 20pt {\hfil$0$\hfil}}}
\newcommand{\cM}{\colorbox{xA}{\hbox to 20pt {\hfil$0$\hfil}}}

\begin{table}[htb]
\centering\small
\begin{tabular}{p{15pt}>{\centering\arraybackslash}p{\OPTspherescolwidth}>{\centering\arraybackslash}p{\OPTspherescolwidth}>{\centering\arraybackslash}p{\OPTspherescolwidth}>{\centering\arraybackslash}p{\OPTspherescolwidth}>{\centering\arraybackslash}p{\OPTspherescolwidth}>{\centering\arraybackslash}p{\OPTspherescolwidth}>{\centering\arraybackslash}p{\OPTspherescolwidth}>{\centering\arraybackslash}p{\OPTspherescolwidth}>{\centering\arraybackslash}p{\OPTspherescolwidth}}
\toprule
           & $\Sn^0$ & $\Sn^1$ & $\Sn^{2}$ & $\Sn^{3}$ & $\Sn^{4}$ & $\Sn^{5}$ & $\Sn^{6}$ & $\Sn^{7}$ & $\Sn^{8}$ \\ \addlinespace[3pt] \midrule
$\pi_{1}$  & $0$     & $\Z$    & \cA       & \cH       & \cI       & \cJ       & \cK       & \cL       & \cM       \\ \addlinespace[3pt]
$\pi_{2}$  & $0$     & $0$     & \cB       & \cA       & \cH       & \cI       & \cJ       & \cK       & \cL       \\ \addlinespace[3pt]
$\pi_{3}$  & $0$     & $0$     & $\Z$      & \cB       & \cA       & \cH       & \cI       & \cJ       & \cK       \\ \addlinespace[3pt]
$\pi_{4}$  & $0$     & $0$     & $\Z_{2}$  & \cC       & \cB       & \cA       & \cH       & \cI       & \cJ       \\ \addlinespace[3pt]
$\pi_{5}$  & $0$     & $0$     & $\Z_{2}$  & $\Z_{2}$  & \cC       & \cB       & \cA       & \cH       & \cI       \\ \addlinespace[3pt]
$\pi_{6}$  & $0$     & $0$     & $\Z_{12}$ & $\Z_{12}$ & \cD       & \cC       & \cB       & \cA       & \cH       \\ \addlinespace[3pt]
$\pi_{7}$  & $0$     & $0$     & $\Z_{2}$  & $\Z_{2}$  & {\footnotesize $\Z {\times} \Z_{12}$} & \cD & \cC & \cB     & \cA    \\ \addlinespace[3pt]
$\pi_{8}$  & $0$     & $0$     & $\Z_{2}$  & $\Z_{2}$  & $\Z_{2}^{2}$ & \cE & \cD & \cC & \cB \\ \addlinespace[3pt]
$\pi_{9}$  & $0$     & $0$     & $\Z_{3}$  & $\Z_{3}$  & $\Z_{2}^{2}$ & $\Z_{2}$ & \cE & \cD & \cC \\ \addlinespace[3pt]
$\pi_{10}$ & $0$     & $0$     & $\Z_{15}$ & $\Z_{15}$ & \footnotesize{$\Z_{24} {\times} \Z_{3}$} & $\Z_{2}$ & \cF & \cE & \cD \\ \addlinespace[3pt]
$\pi_{11}$ & $0$     & $0$     & $\Z_{2}$  & $\Z_{2}$  & $\Z_{15}$ & $\Z_{2}$ & $\Z$ & \cF & \cE \\ \addlinespace[3pt]
$\pi_{12}$ & $0$     & $0$     & $\Z_{2}^{2}$ & $\Z_{2}^{2}$ & $\Z_{2}$ & $\Z_{30}$ & $\Z_{2}$ & \cG & \cF \\ \addlinespace[3pt]
$\pi_{13}$ & $0$     & $0$     & {\footnotesize $\Z_{12} {\times} \Z_{2}$} & {\footnotesize $\Z_{12} {\times} \Z_{2}$} & $\Z_{2}^{3}$ & $\Z_{2}$ & $\Z_{60}$ & $\Z_{2}$ & \cG \\ \addlinespace[3pt]
%$\pi_{14}$ & $0$     & $0$     & $\Z_{84} {\times} \Z_{2}^{2}$ & $\Z_{84} {\times} \Z_{2}^{2}$ & {\footnotesize $\Z_{120} {\times} \Z_{12} {\times} \Z_{2}$} & $\Z_{2}^{3}$ & $\Z_{24} {\times} \Z_{2}$ & $\Z_{120}$ & $\Z_{2}$ \\ \addlinespace[3pt]
%$\pi_{15}$ & $0$     & $0$     & $\Z_{2}^{2}$ & $\Z_{2}^{2}$ & $\Z_{84} {\times} \Z_{2}^{5}$ & $\Z_{72} {\times} \Z_{2}$ & $\Z_{2}^{3}$ & $\Z_{2}^{3}$ & $\Z {\times} \Z_{120}$ \\ \addlinespace[3pt]
\bottomrule
\end{tabular}

\caption{球面的同伦群~\cite{wikipedia-groups}。\index{homotopy!group!of sphere}
表中每一格所列的群即为 $n$ 维球面 $\Sn^n$ 的第 $k$ 个同伦群 $\pi_k$ 所同构的群；其中 $\Z$ 是整数加群，
\index{integers}$\Z_{m}$ 是 $m$ 阶循环群。\index{group!cyclic}\index{cyclic group}
}
\label{tab:homotopy-groups-of-spheres}
\end{table}


\egroup

理解这种复杂性的一种途径，是借助 \cref{cha:basics} 中引入的空间与 $\infty$-群胚（$\infty$-groupoid）之间的对应关系。
\index{.infinity-groupoid@$\infty$-groupoid}%
如 \cref{sec:circle} 所述，2-球面由只包含一个点和一条 2 维环路的高阶归纳类型（higher inductive type）给出。
因此，我们可能会问：既然类型 $\Sn ^2$ 没有任何生成元能产生 3 维单元，为什么 $\pi_3(\Sn ^2)$ 会是 $\Z$ 呢？
原来，$\pi_3(\Sn ^2)$ 的生成元是利用 \cref{thm:EckmannHilton} 的证明中所描述的互换律构造出来的：
$\infty$-群胚的代数结构包含不同层级之间的非平凡交互，而这些交互便产生了更高维同伦群的元素。

\index{classical!homotopy theory|)}%

类型论为研究这一结构提供了自然的框架，因为我们可以方便地定义高阶同伦群。
回忆 \cref{def:loopspace}，对于 $n:\N$，带点类型 $(A,a)$ 的 $n$ 重迭代环路空间\index{loop space!iterated}递归地定义为：
\begin{align*}
  \Omega^0(A,a)&=(A,a)\\
  \Omega^{n+1}(A,a)&=\Omega^n(\Omega(A,a)).
\end{align*}
%
由此我们得到 $n$ 维环路的\emph{空间}（即一个类型），而这个空间本身又含有更高阶的同伦结构。\index{loop!n-@$n$-}
通过截断，我们可以得到 $n$ 维环路的集合（这一构造也在 \cref{sec:free-algebras} 中作为例子定义过）：

\begin{defn}[同伦群]\label{def-of-homotopy-groups}
  给定 $n\ge 1$ 和一个带点类型 $(A,a)$，我们定义 $A$ 在 $a$ 处的 \define{同伦群（homotopy groups）}
  \indexdef{homotopy!group}%
  为：
  \[\pi_n(A,a)\defeq \Trunc0{\Omega^n(A,a)}\]
\end{defn}

\noindent
由于 $n\ge 1$，$\Omega^n(A)$ 上的道路拼接与取逆运算诱导出 $\pi_n(A)$ 上的运算，从而自然地得到群结构。
若 $n\ge 2$，则群 $\pi_n(A)$ 是交换群\index{group!abelian}，这由 Eckmann--Hilton 论证\index{Eckmann--Hilton argument}（\cref{thm:EckmannHilton}）可知。
为了方便，我们也可以记 $\pi_0(A) \defeq \trunc0 A$，
但这个情形有所不同：它不仅不是群，而且在定义时也不依赖于 $A$ 中的任何基点。

\index{presentation!of an infinity-groupoid@of an $\infty$-groupoid}%
这个定义适合研究同伦群，因为类型 $X$ 的（高阶）归纳定义将 $X$ 呈现为一个自由类型，类似于一个自由 $\infty$-群胚，
\index{.infinity-groupoid@$\infty$-groupoid}%
而这种呈现\emph{决定了}但并不\emph{显式描述} $X$ 的高阶相等类型。
相等类型中既包含生成元（例如对于圆周来说是 $\lloop$），也包含对这些生成元施加所有群胚运算（恒等、复合、取逆、结合性、互换、\ldots）所得的结果。
因此，空间的高阶归纳呈现允许我们提出``$X$ 的相等类型实际上是什么？''这样的问题，尽管回答它可能需要相当多的数学论证。
这是类型论中一个熟知事实的高维推广：刻画 $X$ 的相等类型可能需要费一番功夫，即使 $X$ 只是一个普通的归纳类型，比如自然数或布尔值。
例如，$\bfalse$ 与 $\btrue$ 不相等这一定理并不能直接从定义推出；
参见 \cref{sec:compute-coprod}。

\index{univalence axiom}%
泛等公理在同伦群的计算中起着必不可少的作用（若无泛等公理，类型论可以有一种解释，使得所有路径——例如圆周上的环路——都是自反性路径）。这一点在下面计算圆周基本群时就会看到：从 $\Omega(\Sn^1)$ 到 $\Z$ 的映射是这样定义的：先将圆周上的一个环路映为集合 $\Z$ 的一个\index{automorphism!of Z, successor@of $\Z$, successor}自同构，使得例如 $\lloop \ct \opp \lloop$ 被映为 $\mathsf{successor} \ct \mathsf{predecessor}$（这里 $\mathsf{successor}$ 与 $\mathsf{predecessor}$ 是 $\Z$ 的自同构，借助泛等性可被视作宇宙中的路径），然后再将该自同构作用于 $0$。泛等性在宇宙中产生非平凡的道路，我们正是利用这一点从高阶归纳类型中的道路提取信息。

在本章中，我们首先计算球面的一些同伦群，包括 $\pi_k(\Sn ^1)$ (\cref{sec:pi1-s1-intro})、$k<n$ 时的 $\pi_{k}(\Sn ^n)$ (\cref{sec:conn-susp,sec:pik-le-n})、借助 Hopf 纤维化 (\cref{sec:hopf}) 与长正合列论证 (\cref{sec:long-exact-sequence-homotopy-groups}) 计算的 $\pi_2(\Sn ^2)$ 与 $\pi_3(\Sn ^2)$，以及借助 Freudenthal 纬悬定理 (\cref{sec:freudenthal}) 计算的 $\pi_n(\Sn ^n)$。接着，我们讨论 van Kampen 定理 (\cref{sec:van-kampen})——它刻画了推出的基本群——以及 Whitehead（怀特海德）原理的适用情况（即：一个在所有同伦群上诱导出等价的映射，何时本身也是等价）(\cref{sec:whitehead})。最后，我们简要概述一些未收入本书的进一步结果，例如 $n\ge 3$ 时的 $\pi_{n+1}(\Sn ^n)$、Blakers--Massey 定理，以及 Eilenberg（艾伦伯格）--Mac Lane（麦克兰恩）空间的一种构造 (\cref{sec:moreresults})。本章的预备知识包括 \cref{cha:typetheory,cha:basics,cha:hits,cha:hlevels} 以及 \cref{cha:logic} 的部分内容。

\index{fundamental!group!of circle|(}
\section{\texorpdfstring{$\pi_1(S^1)$}{π₁(S¹)}}
\label{sec:pi1-s1-intro}

本节的目标是证明 $\id {\pi_1(\Sn ^1)} {\Z}$。
事实上，我们将证明环路空间\index{loop space} ${\Omega(\Sn ^1)}$ 等价于 $\Z$。
这是一个更强的结论，因为根据定义 $\id{\pi_1(\Sn ^1)} {\trunc 0 {\Omega(\Sn ^1)}}$；因此，若 $\id {\Omega(\Sn ^1)} {\Z}$，则由合同性可得 $\id {\trunc
  0 {\Omega(\Sn ^1)}} {\trunc 0 {\Z}}$，而 $\Z$ 按定义是一个集合（因为它是一个集合商，参见 \cref{defn-Z,Z-quotient-by-canonical-representatives}），于是 $\id {\trunc 0 {\Z}} {\Z}$。
不仅如此，一旦知道 ${\Omega(\Sn ^1)}$ 是一个集合，便可推出当 $n>1$ 时 $\pi_n(\Sn^1)$ 是平凡的，这样一来，我们其实就计算出了 $\Sn^1$ 的\emph{全部}同伦群。

\subsection{开始探索}
\label{sec:pi1s1-initial-thoughts}

在 $\Omega(\Sn^1)$ 与 $\Z$ 之间定义两个方向的函数并不困难。
将 \cref{thm:looptothe} 特化到 $\lloop:\base=\base$ 的情形，便得到一个函数 $\lloop^{\blank} : \Z \rightarrow (\id{\base}{\base})$，其定义（粗略地说）为
\[
  \lloop^n =
  \begin{cases}
    \underbrace{\lloop \ct \lloop \ct \cdots \ct \lloop}_{n}  & \text{if $n > 0$,} \\
    \underbrace{\opp \lloop \ct \opp \lloop \ct \cdots \ct \opp \lloop}_{-n} & \text{if $n < 0$,} \\
    \refl{\base} & \text{if $n = 0$.}
\end{cases}
\]
%
而定义反方向的函数 $g:\Omega(\Sn^1)\to\Z$ 则要棘手一些。
注意到后继函数 $\Zsuc:\Z\to\Z$ 是一个等价，
\index{successor!isomorphism on Z@isomorphism on $\Z$}%
因而它在宇宙 $\type$ 中诱导出一条道路 $\ua(\Zsuc):\Z=\Z$。
于是，$\Sn^1$ 的递归原理诱导出一个映射 $c:\Sn^1\to\type$，由 $c(\base)\defeq \Z$ 与 $\apfunc c (\lloop) \defid \ua(\Zsuc)$ 给出。
这样我们就得到了 $\apfunc{c} : (\base=\base) \to (\Z=\Z)$，进而可以定义 $g(p)\defeq \transfib{X\mapsto X}{\apfunc{c}(p)}{0}$。

利用这些定义，借助 $n$ 的归纳原理 \cref{thm:sign-induction}，我们甚至能对任意 $n:\Z$ 证明 $g(\lloop^n)=n$。
（稍后我们将证明一个更一般的结论。）
然而，另一个等式 $\lloop^{g(p)}=p$ 的证明则要困难得多。
最直接的尝试是道路归纳，但道路归纳并不适用于像 $p:(\base=\base)$ 这样\emph{两个端点都固定}的环路！
这需要一个新的思路，它既可以从经典同伦论的角度，也可以从类型论的角度来阐释。
我们先从前者讲起。

\subsection{经典证明}
\label{sec:pi1s1-classical-proof}

\index{classical!homotopy theory|(}%
在经典同伦论中，证明 $\pi_1(\Sn^1)=\Z$ 有一个标准方法，要用到万有覆叠空间。
我们的证明可以看作此证明的类型论版本，其中覆叠空间体现为纤维是集合的纤维化。
\index{fibration}%
\index{total!space}%
回顾一下，同伦论中空间 $B$ 上的\emph{纤维化}对应于类型论中的类型族 $B\to\type$。
\index{path!fibration}%
\index{fibration!of paths}%
特别地，对于点 $x_0:B$，类型族 $(x\mapsto (x_0=x))$ 对应于\emph{道路纤维化} $P_{x_0} B \to B$：$P_{x_0} B$ 中的点是 $B$ 中始于 $x_0$ 的道路，而到 $B$ 的映射选取的是这条道路的另一个端点。
这个全空间 $P_{x_0} B$ 是可缩的，因为我们可以将任一道路``收缩''回其起始端点 $x_0$ —— 我们之前已经见过这一事实的类型论版本，即 \cref{thm:contr-paths}。
此外，$x_0$ 上的纤维就是环路空间\index{loop space} $\Omega(B,x_0)$ —— 在类型论中，由环路空间的定义即可看出。

\begin{figure}\centering
  \begin{tikzpicture}[xscale=1.4,yscale=.6]
    \node (R) at (2,1) {$\mathbb{R}$};
    \node (S1) at (2,-2) {$S^1$};
    \draw[->] (R) -- node[auto] {$w$} (S1);
    \draw (0,-2) ellipse (1 and .4);
    \draw[dotted] (1,0) arc (0:-30:1 and .8);
    \draw (1,0) arc (0:90:1 and .8) arc (90:270:1 and .3) coordinate (t1);
    \draw[white,line width=4pt] (t1) arc (-90:90:1 and .8);
    \draw (t1) arc (-90:90:1 and .8) arc (90:270:1 and .3) coordinate (t2);
    \draw[white,line width=4pt] (t2) arc (-90:90:1 and .8);
    \draw (t2) arc (-90:90:1 and .8) arc (90:270:1 and .3) coordinate (t3);
    \draw[white,line width=4pt] (t3) arc (-90:90:1 and .8);
    \draw (t3) arc (-90:-30:1 and .8) coordinate (t4);
    \draw[dotted] (t4) arc (-30:0:1 and .8);
    \node[fill,circle,inner sep=1pt,label={below:\scriptsize \base}] at (0,-2.4) {};
    \node[fill,circle,inner sep=1pt,label={above left:\scriptsize 0}] at (0,.2) {};
    \node[fill,circle,inner sep=1pt,label={above left:\scriptsize 1}] at (0,1.2) {};
    \node[fill,circle,inner sep=1pt,label={above left:\scriptsize 2}] at (0,2.2) {};
  \end{tikzpicture}
  \caption{经典拓扑中的环绕映射}\label{fig:winding}
\end{figure}

现在，在经典同伦论里——那里将 $\Sn^1$ 视为一个拓扑空间——我们可以这样进行。
\index{winding!map}%
考虑``环绕''映射 $w:\mathbb{R}\to \Sn ^1$，它看起来像一条投影到圆上的螺旋线（见 \cref{fig:winding}）。
映射 $w$ 将螺旋线\index{helix}上的每个点映为它正下方的圆上的点。
它是一个纤维化，并且每个点上的纤维同构于整数。\index{integers}
如果我们将底部沿逆时针方向绕行环路的那条道路提升，我们在螺旋线上就上升一层，使纤维中的整数加一。
类似地，沿底部环路顺时针绕行，对应于在螺旋线上下降一层，使该数减一。
这个纤维化称为圆的\emph{万有覆叠}。
\index{universal!cover}%
\index{cover!universal}%
\index{covering space!universal}%

经典同伦论中的一个基本事实是：基空间 $B$ 上的两个纤维化 $E_1$ 与 $E_2$ 之间的一个映射，若是 $E_1$ 与 $E_2$ 之间的同伦等价，则它在所有纤维上都诱导出同伦等价。
（我们此前也已经见过这一事实的类型论版本，见 \cref{thm:total-fiber-equiv}。）
由于 $\mathbb{R}$ 和 $P_{\base} S^1$ 都是可缩的拓扑空间，故它们同伦等价；因此，它们在基点处的纤维 $\Z$ 与 $\Omega(\Sn ^1)$ 也同伦等价。

\index{classical!homotopy theory|)}%

\subsection{类型论中的万有覆叠}
\label{sec:pi1s1-universal-cover}

\index{universal!cover|(}%
\index{cover!universal|(}%
\index{covering space!universal|(}%

我们来考虑如何在类型论中表达前述证明。
我们已经提到，$\Sn^1$ 的道路纤维化由类型族 $(x\mapsto (\base=x))$ 表示。
对于 $\Sn^1$ 的万有覆叠，我们也已经有了一个很好的候选：它恰恰就是我们在 \cref{sec:pi1s1-initial-thoughts} 中定义的类型族 $c:\Sn^1\to\type$！
根据定义，这个族在 $\base$ 处的纤维是 $\Z$，而沿着 $\lloop$ 传输的效果是加一 —— 因此它的行为正合我们对 \cref{fig:winding} 的期望。

不过，由于我们尚不知道这个族是否具备万有覆叠应有的性质（例如，其全空间是否单连通），我们暂且给它换一个名字。
为便于查阅，我们在此重复其定义。
\index{integers}

\begin{defn}[$\Sn^1$ 的万有覆叠] \label{S1-universal-cover}
  通过圆周递归定义 $\code : \Sn ^1 \to \type$，满足
  \begin{align*}
    \code(\base) &\defeq \Z \\
    \apfunc{\code}({\lloop}) &\defid \ua(\Zsuc).
  \end{align*}
\end{defn}

我们在此简略强调一下这个类型族的定义，因为它与经典同伦论中通常定义覆叠空间的方式大不相同。
要通过圆周递归定义一个函数，我们需要在值域中找出一个点和一个环路。
在当前情形下，值域是 $\type$，我们选择的点是 $\Z$，这对应着我们的期望：万有覆叠的纤维应该是整数。
我们选择的环路是 $\Z$ 上的后继/前驱
\index{successor!isomorphism on Z@isomorphism on $\Z$}%
\index{predecessor!isomorphism on Z@isomorphism on $\Z$}%
同构，这对应着一个几何事实：在基空间中绕环路一圈，在螺旋线上就上升一层。
这部分证明需要用到泛等性，因为我们需要将 $\Z$ 上的一个\emph{非平凡}等价转化为一个相等。

我们称此纤维化为``码''的纤维化，因为它的元素是组合数据，充当圆周上道路的码：整数 $n$ 编码了绕圆 $n$ 圈的道路。

从这个定义出发，容易计算出：用 $\code$ 传输会将 $\lloop$ 映为后继函数，而将 $\opp{\lloop}$ 映为前驱函数：


\begin{lem} \label{lem:transport-s1-code}
\id{\transfib \code \lloop x} {x + 1} 且
\id{\transfib \code {\opp \lloop} x} {x - 1}。
\end{lem}

\begin{proof}
对于第一个等式，我们计算如下：
\begin{align}
{\transfib \code \lloop x}
&= \transfib {A \mapsto A} {(\ap{\code}{\lloop})} x \tag{by \cref{thm:transport-compose}}\\
&= \transfib {A \mapsto A} {\ua (\Zsuc)} x \tag{by computation for $\rec{\Sn^1}$}\\
&= x + 1 \tag{by computation for \ua}.
\end{align}
第二个等式由第一个推出，因为 $\transfib{B}{p}{\blank}$ 与 $\transfib{B}{\opp p}{\blank}$ 总是互逆的，故 $\transfib\code {\opp \lloop}{\blank}$ 必是 $\Zsuc$ 的逆。
\end{proof}

此刻我们就能看出最初的方法问题何在：我们只定义了纤维 $\Omega(\Sn^1)$ 和 \Z 上的 $f$ 与 $g$，但实际上应该定义整个在 $\Sn^1$ 上的\emph{纤维化的}态射。
在类型论中，这意味着我们本该定义具有类型
\begin{align}
  \prd{x:\Sn^1} &((\base=x) \to \code(x)) \qquad\text{and/or} \label{eq:pi1s1-encode}\\
  \prd{x:\Sn^1} &(\code(x) \to (\base=x))\label{eq:pi1s1-decode}
\end{align}
的函数，而不是只定义这些函数在 $x$ 为 \base 时的特殊情况。
这也印证了类型论中一个常见的观察：当试图证明某个归纳类型的特定实例具有某种性质时，通常更容易先将命题推广到该类型的\emph{所有}实例，然后再用归纳法证明。
从这个角度看，$\Omega(\Sn^1)=\Z$ 的证明与 \cref{sec:compute-coprod,sec:compute-nat} 中刻画余积及自然数的相等类型遵循着相同的模式。

至此，完成证明有两种途径。
我们可以继续模仿经典论证，构造~\eqref{eq:pi1s1-encode} 或~\eqref{eq:pi1s1-decode}（二者选一即可），证明全空间之间的同伦等价会诱导纤维上的等价，然后证明万有覆叠的全空间是可缩的。
第一个关于 $\Omega(\Sn^1)=\Z$ 的类型论证明就遵循了这个模式；我们称之为\emph{同伦论证明}。

然而后来我们发现，还存在另一种证明，它更具有类型论的味道，并且更贴近 \cref{sec:compute-coprod,sec:compute-nat} 中的证明风格。
在这一证明中，我们直接构造出~\eqref{eq:pi1s1-encode} 和~\eqref{eq:pi1s1-decode}，并通过计算证明它们互逆。
我们称此为\emph{编码--解码证明}，因为我们将函数~\eqref{eq:pi1s1-encode} 和~\eqref{eq:pi1s1-decode} 分别称为\emph{编码}与\emph{解码}。
两种证明都使用了前面给出的同一个覆叠构造。
前一种证明从全空间的等价诱导出纤维上的等价，而编码--解码证明则显式地构造出纤维之间的逆映射（即\emph{解码}）。
前一种证明使用了可缩性，而编码--解码证明使用了道路归纳、圆周归纳和整数归纳。
这些工具与证明可缩性所用的工具是一样的——事实上，道路归纳\emph{本质上就是}道路纤维化的可缩性再配上 $\mathsf{transport}$——只是应用方式不同。

既然这是一本关于同伦类型论的书，我们将首先介绍编码--解码证明。
若有同伦论背景的读者感到迷失，我们建议直接跳到同伦论证明（\cref{subsec:pi1s1-homotopy-theory}）。

\index{universal!cover|)}%
\index{cover!universal|)}%
\index{covering space!universal|)}%

\subsection{编码--解码证明}
\label{subsec:pi1s1-encode-decode}

\index{encode-decode method|(}%
\indexsee{encode}{encode-decode method}%
\indexsee{decode}{encode-decode method}%
我们从将道路映射到代码的函数~\eqref{eq:pi1s1-encode} 开始：


\begin{defn}
定义 $\encode : \prd{x : \Sn ^1} (\base=x) \rightarrow  \code(x)$ 为
\[
\encode \: p \defeq \transfib{\code} p 0
\]
（其中参数 $x$ 是隐式的）。
\end{defn}


编码的定义方式是将一条道路提升到万有覆叠，这决定了一个等价，然后将得到的等价应用到 $0$。
这个函数的有趣之处在于，当圆上的环路是用同伦类型论中抽象的群胚框架来表示时，它能从这样一个环路计算出一个具体的数。
为了直观理解它是如何做到的，我们可以根据前面的引理观察到：$\transfib \code \lloop x$ 是后继映射，而 $\transfib \code {\opp
  \lloop} x$ 是前驱映射。
此外，$\mathsf{transport}$ 具有函子性（\cref{cha:basics}），所以 $\transfib{\code} {\lloop \ct \lloop}{\blank}$ 就是
\[(\transfib \code \lloop-) \circ (\transfib \code \lloop-)\]
依此类推。
因此，当 $p$ 是一个像
\[
\lloop \ct \opp \lloop \ct \lloop \ct \cdots
\]
这样的复合时，$\transfib{\code}{p}{\blank}$ 将计算出函数的复合，例如
\[
\Zsuc \circ \Zpred \circ \Zsuc \circ \cdots
\]
将这一组函数的复合应用到 $0$，就会计算出该道路的\index{winding!number}%
\emph{环绕数}——即在消去逆元之后，它绕圆了多少圈，其方向由正负号标记。
由此可见，$\encode$ 的计算行为源于高阶归纳类型与泛等性的归约规则，以及 $\mathsf{transport}$ 对复合与逆的作用。

注意，特例 $\encode' \defeq \encode_{\base}$ 的类型是 $(\id \base \base) \rightarrow \Z$。
这将是我们所需等价的一半；实际上，它正是 \cref{sec:pi1s1-initial-thoughts} 中定义的函数 $g$。

类似地，函数~\eqref{eq:pi1s1-decode} 是 \cref{sec:pi1s1-initial-thoughts} 中函数 $\lloop^{\blank}$ 的推广。

\begin{defn}\label{thm:pi1s1-decode}
通过对 $x$ 作圆周归纳来定义 $\decode : \prd{x : \Sn ^1}\code(x) \rightarrow (\base=x)$。
我们只需给出一个函数
${\code(\base) \rightarrow (\base=\base)}$，为此我们使用 $\lloop^{\blank}$，并证明 $\lloop^{\blank}$ 尊重环路即可。
\end{defn}

\begin{proof}
为了证明 $\lloop^{\blank}$ 尊重环路，只需给出一条从 $\lloop^{\blank}$ 到其自身的道路，且这条道路位于 $\lloop$ 之上。
根据依值道路的定义，这意味着需要一条从
\[\transfib {(x' \mapsto \code(x') \rightarrow (\base=x'))} {\lloop} {\lloop^{\blank}}\]
到 $\lloop^{\blank}$ 的道路。
我们如下构造这样一条道路：
\begin{align*}
  \MoveEqLeft \transfib {(x' \mapsto \code(x') \rightarrow (\base=x'))} \lloop {\lloop^{\blank}}\\
&= \transfibf {x'\mapsto (\base=x')}(\lloop) \circ {\lloop^{\blank}} \circ \transfibf \code ({\opp \lloop}) \\
&= (- \ct \lloop) \circ (\lloop^{\blank}) \circ \transfibf \code ({\opp \lloop}) \\
&= (- \ct \lloop) \circ (\lloop^{\blank}) \circ \Zpred \\
&= (n \mapsto \lloop^{n - 1} \ct \lloop).
\end{align*}
在第一行，我们应用了当纤维化的外层连接词是 $\rightarrow$ 时 $\mathsf{transport}$ 的刻画，它将这个 $\mathsf{transport}$ 化归为在定义域和值域类型上与 $\mathsf{transport}$ 的前复合与后复合。
在第二行，我们应用了当类型族是 $x\mapsto \id{\base}{x}$ 时 $\mathsf{transport}$ 的刻画，即道路的后复合。
在第三行，我们使用了从 \cref{lem:transport-s1-code} 得到的 $\code$ 在 $\opp \lloop$ 上的作用。
在第四行，我们仅仅将函数复合化归了。
因此，只需证明对于所有 $n$，有 $\id{\lloop^{n - 1} \ct \lloop}{\lloop^{n}}$。
这是 \cref{thm:sign-induction} 的一个简单应用，利用了群胚律。
\end{proof}

现在我们可以证明 $\encode$ 和 $\decode$ 互为拟逆。
曾经困难的方向现在变得简单了！

\begin{lem} \label{lem:s1-decode-encode}
对于所有 $x: \Sn ^1$ 和 $p : \id \base x$，有 $\id
{\decode_x({{\encode_x(p)}})} p$。
\end{lem}

\begin{proof}
通过道路归纳，只需证明
\narrowequation{\id {\decode_{\base}({{\encode_{\base}(\refl{\base})}})} {\refl{\base}}.}
但是
\narrowequation{\encode_{\base}(\refl{\base}) \jdeq \transfib{\code}{\refl{\base}} 0 \jdeq 0,}
并且 $\decode_{\base}(0) \jdeq \lloop^ 0 \jdeq \refl{\base}$。
\end{proof}

另一个方向也不难多少。

\begin{lem} \label{lem:s1-encode-decode} 对于所有
$x: \Sn ^1$ 和 $c : \code(x)$，我们有 $\id
{\encode_x({{\decode_x(c)}})} c$。
\end{lem}

\begin{proof}
证明通过圆周归纳进行。只需证明 $\base$ 的情形，因为 $\lloop$ 的情形是 $\Z$ 中两条道路之间的道路，而 $\Z$ 是集合，所以这是直接可得的。

因此，我们需要证明：对于所有 $n : \Z$，有
\[
\id {\encode'(\lloop^n)} {n}.
\]
证明通过对 $n$ 的归纳进行，使用 \cref{thm:sign-induction}。
%
\begin{itemize}

\item 对于 $0$ 的情形，结果按定义成立。

\item 对于 $n+1$ 的情形，
\begin{align}
 {\encode'(\lloop^{n+1})}
&= {\encode'(\lloop^{n} \ct \lloop)} \tag{by definition of $\lloop^{\blank}$} \\
&= \transfib{\code}{(\lloop^{n} \ct \lloop)}{0} \tag{by definition of $\encode$}\\
&= \transfib{\code}{\lloop}{(\transfib{\code}{\lloop^n}{0})} \tag{by functoriality}\\
&= {(\transfib{\code}{\lloop^n}{0})} + 1 \tag{by \cref{lem:transport-s1-code}}\\
&= n + 1. \tag{by the inductive hypothesis}
\end{align}

\item 负数的情形是类似的。\qedhere
\end{itemize}
\end{proof}

最后，我们给出以下定理。

\begin{thm}
存在一族等价 $\prd{x : \Sn ^1} (\eqv {(\base=x)} {\code(x)})$。
\end{thm}

\begin{proof}
映射 $\encode$ 和 $\decode$ 互为拟逆，这由
\cref{lem:s1-decode-encode,lem:s1-encode-decode}
给出。
\end{proof}

在 \base 处实例化，得到


\begin{cor}\label{cor:omega-s1}
$\eqv {\Omega(\Sn^1,\base)} {\Z}$。
\end{cor}

一个简单的归纳表明该等价把加法对应到复合，因此有 $\Omega(\Sn ^1) = \Z$（作为群）。

\begin{cor} \label{cor:pi1s1}
$\id{\pi_1(\Sn ^1)} {\Z}$，而当 $n>1$ 时，$\id{\pi_n(\Sn ^1)}0$。
\end{cor}

\begin{proof}
对于 $n=1$，我们在上面从 \cref{cor:omega-s1} 概述了证明。
对于 $n > 1$，有 $\trunc 0 {\Omega^{n}(\Sn ^1)} = \trunc 0 {\Omega^{n-1}(\Omega{\Sn ^1})} = \trunc 0 {\Omega^{n-1}(\Z)}$。
由于 $\Z$ 是集合，$\Omega^{n-1}(\Z)$ 是可缩的，因此这是平凡的。
\end{proof}

\index{encode-decode method|)}%

\subsection{同伦论的证明}
\label{subsec:pi1s1-homotopy-theory}

在 \cref{sec:pi1s1-universal-cover} 中，我们在类型论里定义了假定的万有覆叠 $\code:\Sn^1\to\type$；在 \cref{subsec:pi1s1-encode-decode} 中，我们定义了一个从道路纤维化到万有覆叠的映射 $\encode : \prd{x:\Sn^1} (\base=x) \to \code(x)$。
对于经典证明，剩下的部分是证明：因为两个全空间都是可缩的，所以这个映射诱导了全空间上的等价，并由此推断它在每个纤维上也必然是等价。

\index{total!space}%
在 \cref{thm:contr-paths} 中我们看到全空间 $\sm{x:\Sn^1} (\base=x)$ 是可缩的。
对于另一个全空间，我们有：

\begin{lem}\label{thm:iscontr-s1cover}
  类型 $\sm{x:\Sn^1}\code(x)$ 是可缩的。
\end{lem}

\begin{proof}
  我们以如下取值应用展平引理（\cref{thm:flattening}）：
  \begin{itemize}
  \item $A\defeq\unit$ 且 $B\defeq\unit$，其中 $f$ 和 $g$ 取为显然的函数。
    因此，展平引理中的基础高阶归纳类型 $W$ 等价于 $\Sn^1$。
  \item $C:A\to\type$ 取常值 \Z。
  \item $D:\prd{b:B} (\eqv{\Z}{\Z})$ 取常值 $\Zsuc$。
  \end{itemize}
  于是，展平引理中定义的类型族 $P:\Sn^1\to\type$ 等价于 $\code:\Sn^1\to\type$。
  因此，展平引理告诉我们，$\sm{x:\Sn^1}\code(x)$ 等价于一个具有以下生成元的高阶归纳类型，我们将其记作 $R$：
  \begin{itemize}
  \item 一个函数 $\mathsf{c}: \Z \to R$。
  \item 对每个 $z:\Z$，一条道路 $\mathsf{p}_z:\mathsf{c}(z) = \mathsf{c}(\Zsuc(z))$。
  \end{itemize}
  我们可以将这个类型称为\define{同伦实数（homotopical reals）}；
  \indexdef{real numbers!homotopical}%
  它在经典证明中扮演的角色与拓扑空间 $\mathbb{R}$ 相同。

  因此，只需再证明 $R$ 是可缩的。
  我们选取 $\mathsf{c}(0)$ 作为收缩中心；现在需要证明对所有 $x:R$ 有 $x=\mathsf{c}(0)$。
  我们对 $R$ 进行归纳来证明。
  首先，当 $x$ 为 $\mathsf{c}(z)$ 时，我们需要给出一条道路 $q_z:\mathsf{c}(0) = \mathsf{c}(z)$，这可以通过对 $z:\Z$ 归纳来完成，使用 \cref{thm:sign-induction}：
  \begin{align*}
    q_0 &\defid \refl{\mathsf{c}(0)}\\
    q_{n+1} &\defid q_n \ct \mathsf{p}_n & &\text{for $n\ge 0$}\\
    q_{n-1} &\defid q_n \ct \opp{\mathsf{p}_{n-1}} & &\text{for $n\le 0$.}
  \end{align*}
  其次，我们需要证明，对任意 $z:\Z$，道路 $q_z$ 沿 $\mathsf{p}_z$ 传输后得到 $q_{z+1}$。
  根据道路传输的性质，这意味着我们需要 $q_z \ct \mathsf{p}_z = q_{z+1}$。
  利用 $q_z$ 的定义，通过对 $z$ 归纳容易完成。
  这就完成了 $R$ 可缩的证明，因而 $\sm{x:\Sn^1}\code(x)$ 也是可缩的。
\end{proof}

\begin{cor}\label{thm:encode-total-equiv}
  由 \encode 诱导的映射：
  \[ \tsm{x:\Sn^1} (\base=x) \to \tsm{x:\Sn^1}\code(x) \]
  是一个等价。
\end{cor}

\begin{proof}
  这两个类型都是可缩的。
\end{proof}

\begin{thm}
  $\eqv {\Omega(\Sn^1,\base)} {\Z}$。
\end{thm}

\begin{proof}
  对 $\encode$ 应用 \cref{thm:total-fiber-equiv}，并使用 \cref{thm:encode-total-equiv}。
\end{proof}

本质上，这两种证明并无太大不同：encode-decode 证明可以看作是同伦论证明的一种``约化''或``展开''。
二者各有优势；这两种观点之间的相互交织，也是本课题趣味的一部分。
\index{fundamental!group!of circle|)}

\subsection{将万有覆叠视作恒等系统}
\label{sec:pi1s1-idsys}

注意，纤维化 $\code:\Sn^1\to\type$ 连同 $0:\code(\base)$ 构成 \cref{defn:identity-systems} 意义下的一个\emph{带点谓词（pointed predicate）}。
从这个角度看，\cref{subsec:pi1s1-encode-decode} 中的编码-解码证明本质上是证明 \code 满足 \cref{thm:identity-systems}\ref{item:identity-systems3}，而 \cref{subsec:pi1s1-homotopy-theory} 中的同伦论证明则是证明它满足 \cref{thm:identity-systems}\ref{item:identity-systems4}。
这提示了第三种途径。

\begin{thm}
  在 \cref{defn:identity-systems} 的意义下，偶对 $(\code,0)$ 是 $\base:\Sn^1$ 处的一个恒等系统。
\end{thm}

\begin{proof}
  给定 $D:\prd{x:\Sn^1} \code(x) \to \type$ 及 $d:D(\base,0)$，我们希望定义函数 $f:\prd{x:\Sn^1}{c:\code(x)} D(x,c)$。
  由圆周归纳，只需指定 $f(\base):\prd{c:\code(\base)} D(\base,c)$ 并验证 $\trans{\lloop}{f(\base)} = f(\base)$。

  显然，$\code(\base)\jdeq \Z$。
  由 \cref{lem:transport-s1-code} 并对 $n$ 归纳，可以对任意整数 $n$ 得到一条道路 $p_n : \transfib{\code}{\lloop^n}{0} = n$。
  因此，由 $\Sigma$-类型中的道路，我们在 $\sm{x:\Sn^1} \code(x)$ 中得到一条道路 $\pairpath(\lloop^n,p_n) : (\base,0) = (\base,n)$。
  将 $d$ 沿此道路在 $D$ 相应的纤维化 $\widehat{D}:(\sm{x:\Sn^1} \code(x)) \to\type$ 中传输，便对任意 $n:\Z$ 得到 $D(\base,n)$ 中的一个元素。
  我们定义此元素为 $f(\base)(n)$：
  \[ f(\base)(n) \defeq \transfib{\widehat{D}}{\pairpath(\lloop^n,p_n)}{d}. \]
  %
  现在我们需要 $\transfib{\lam{x} \prd{c:\code(x)} D(x,c)}{\lloop}{f(\base)} = f(\base)$。
  由 \cref{thm:dpath-forall}，这意味着我们需要证明对任意 $n:\Z$，
  %
  \begin{narrowmultline*}
    \transfib{\widehat D}{\pairpath(\lloop,\refl{\trans\lloop n})}{f(\base)(n)}
    =_{D(\base,\trans\lloop n)} \narrowbreak
    f(\base)(\trans\lloop n).
  \end{narrowmultline*}
  %
  现在我们有一条道路 $q:\trans\lloop n = n+1$，故沿此道路传输后，只需证明
  \begin{multline*}
    \transfib{D(\base)}{q}{\transfib{\widehat D}{\pairpath(\lloop,\refl{\trans\lloop n})}{f(\base)(n)}}\\
    =_{D(\base,n+1)} \transfib{D(\base)}{q}{f(\base)(\trans\lloop n)}.
  \end{multline*}
  由关于 transport 与依值应用的几个引理，上式等价于
  \[ \transfib{\widehat D}{\pairpath(\lloop,q)}{f(\base)(n)} =_{D(\base,n+1)} f(\base)(n+1). \]
  然而，展开 $f(\base)$ 的定义，我们有
  \begin{narrowmultline*}
    \transfib{\widehat D}{\pairpath(\lloop,q)}{f(\base)(n)}
    \narrowbreak
    \begin{aligned}[t]
      &= \transfib{\widehat D}{\pairpath(\lloop,q)}{\transfib{\widehat{D}}{\pairpath(\lloop^n,p_n)}{d}}\\
      &= \transfib{\widehat D}{\pairpath(\lloop^n,p_n) \ct \pairpath(\lloop,q)}{d}\\
      &= \transfib{\widehat D}{\pairpath(\lloop^{n+1},p_{n+1})}{d}\\
      &= f(\base)(n+1).
    \end{aligned}
  \end{narrowmultline*}
  这里我们用到了 transport 的函子性、$\Sigma$-类型中复合的刻画（这是留给读者的练习）、以及一个将 $p_n$、$q$ 与 $p_{n+1}$ 联系起来的引理，该引理的陈述与证明也留给读者。

  这就完成了 $f:\prd{x:\Sn^1}{c:\code(x)} D(x,c)$ 的构造。
  由于
  \[f(\base,0) \jdeq \trans{\pairpath(\lloop^0,p_0)}{d} = \trans{\refl{\base}}{d} = d,\]
  我们已证明 $(\code,0)$ 是一个恒等系统。
\end{proof}

\begin{cor}
  对于任意 $x:\Sn^1$，有 $\eqv{(\base=x)}{\code(x)}$。
\end{cor}

\begin{proof}
  由 \cref{thm:identity-systems}。
\end{proof}

当然，这个证明本质上包含了与前两种证明相同的要素。
粗略地说，它统一了 \cref{thm:pi1s1-decode,lem:s1-encode-decode} 的证明，在一般情形下只需进行一次必要的归纳论证。

\begin{rmk}
  注意，以上所有关于 $\eqv{\pi_1(\Sn^1)}{\Z}$ 的证明都本质地使用了泛等公理。
  这是不可避免的：为了证明 $\eqv{\pi_1(\Sn^1)}{\Z}$，泛等性或其类似物是\emph{必需}的。
  在缺乏泛等性的情况下，可以一致地假设``所有类型都是集合''这一陈述（又称``恒等证明的唯一性''或``K 公理''，如 \cref{sec:hedberg} 中所讨论的），而这一陈述反而蕴含 $\eqv{\pi_1(\Sn^1)}{\unit}$。
  事实上，$\pi_1(\Sn^1)$ 的（非）平凡性恰好可以检测是否所有类型都是集合：\cref{thm:loop-nontrivial} 的证明反过来表明，如果 $\lloop=\refl{\base}$，那么所有类型都是集合。
\end{rmk}

\section{纬悬的连通性}
\label{sec:conn-susp}

回顾 \cref{sec:connectivity}，一个类型 $A$ 称为 \define{$n$-连通的（$n$-connected）}，如果 $\trunc nA$ 是可缩的。
本节的目标是证明 \cref{sec:suspension} 中定义的纬悬运算会提升连通度。

\begin{thm} \label{thm:suspension-increases-connectedness}
  如果 $A$ 是 $n$-连通的，那么 $A$ 的纬悬是 $(n+1)$-连通的。
\end{thm}

\begin{proof}
  我们在 \cref{sec:colimits} 中曾指出，$A$ 的纬悬是推出 $\unit\sqcup^A\unit$，因此需要证明下列类型是可缩的：
  %
  \[\trunc{n+1}{\unit\sqcup^A\unit}.\]
  %
  由 \cref{reflectcommutespushout} 可知
  $\trunc{n+1}{\unit\sqcup^A\unit}$ 是图表
  \[\xymatrix{\trunc{n+1}A \ar[d] \ar[r] & \trunc{n+1}{\unit} \\
    \trunc{n+1}{\unit} & }.\]
  %
  在 $\typelep{n+1}$ 中的推出。
  已知 $\trunc{n+1}{\unit}=\unit$，类型
  $\trunc{n+1}{\unit\sqcup^A\unit}$ 也是下列图表在 $\typelep{n+1}$ 中的推出（因为两个图表相等）：
  %
  \[\Ddiag=\vcenter{\xymatrix{\trunc{n+1}A \ar[d] \ar[r] & \unit \\
    \unit & }}.\]
  %
  我们现在证明 $\unit$ 也是 $\Ddiag$ 在 $\typelep{n+1}$ 中的推出。
  %
  令 $E$ 为一个 $(n+1)$-截断的类型；我们需要证明以下映射是一个等价：
  %
  \[\function{(\unit \to E)}{\cocone{\Ddiag}{E}}{y}
  {(y,y,\lamu{u:{\trunc{n+1}A}} \refl{y(\ttt)})}.\]
  %
  其中回顾 $\cocone{\Ddiag}{E}$ 是类型
  \[\sm{f:\unit \to E}{g:\unit \to E}(\trunc{n+1}A\to
  (f(\ttt)=_E{}g(\ttt))).\]
  %
  映射 $\function{(\unit\to E)}{E}{f}{f(\ttt)}$ 是一个等价，因此我们也有
  \[\cocone{\Ddiag}{E}=\sm{x:E}{y:E}(\trunc{n+1}A\to(x=_Ey)).\]
  %
  现在，$A$ 是 $n$-连通的，故 $\trunc{n+1}A$ 也是 $n$-连通的（因为
  $\trunc n{\trunc{n+1}A}=\trunc nA=\unit$），且 $(x=_Ey)$ 是 $n$-截断的（因为
  $E$ 是 $(n+1)$-截断的）。因此，由 \cref{connectedtotruncated}，以下映射是一个等价：
  %
  \[\function{(x=_Ey)}{(\trunc{n+1}A\to(x=_Ey))}{p}{\lam{z} p}\]
  %
  因此我们有
  %
  \[\cocone{\Ddiag}{E}=\sm{x:E}{y:E}(x=_Ey).\]
  %
  但下列映射是一个等价：
  %
  \[\function{E}{\sm{x:E}{y:E}(x=_Ey)}{x}{(x,x,\refl{x})}.\]
  %
  因此
  %
  \[\cocone{\Ddiag}{E}=E.\]
  %
  最终我们得到一个等价：
  %
  \[(\unit \to E)\eqvsym\cocone{\Ddiag}{E}\]
  %
  现在我们可以展开定义以得到该映射的显式表达式，并且容易看出这正是我们一开始的那个映射。

  因此我们证明了 $\unit$ 是 $\Ddiag$ 在 $\typelep{n+1}$ 中的推出。利用推出的唯一性，得到 $\trunc{n+1}{\unit\sqcup^A\unit}=\unit$，这证明了 $A$ 的纬悬是 $(n+1)$-连通的。
\end{proof}

\begin{cor} \label{cor:sn-connected}
  对所有 $n:\N$，球面 $\Sn^n$ 是 $(n-1)$-连通的。
\end{cor}

\begin{proof}
  我们对 $n$ 归纳证明。
  %
  对于 $n=0$，我们需要证明 $\Sn^0$ 仅知有实例，这是显然的。
  %
  设 $n:\N$ 且 $\Sn^n$ 是 $(n-1)$-连通的。根据定义，$\Sn^{n+1}$ 是 $\Sn^n$ 的纬悬，因此由前述引理可知 $\Sn^{n+1}$ 是 $n$-连通的。
\end{proof}

\section{$n$-连通空间的 \texorpdfstring{$\pi_{k \le n}$}{π\_(k≤n)} 与 \texorpdfstring{$\pi_{k < n}(\Sn ^n)$}{π\_(k<n)(Sⁿ)}}
\label{sec:pik-le-n}

令 $(A,a)$ 为一个带点类型且 $n:\N$。回顾
\cref{thm:homotopy-groups}，若 $n>0$，集合 $\pi_n(A,a)$ 具有群结构；若 $n>1$，该群是交换群\index{group!abelian}。

关于 $n$-截断类型与 $n$-连通类型的同伦群，我们现在可以有所讨论。

\begin{lem}
  如果 $A$ 是 $n$-截断的且 $a:A$，那么对所有 $k>n$ 有 $\pi_k(A,a)=\unit$。
\end{lem}

\begin{proof}
  $n$-类型的环路空间是 $(n-1)$-类型，因此 $\Omega^k(A,a)$ 是 $(n-k)$-类型，且 $(n-k)\le-1$，所以 $\Omega^k(A,a)$ 是一个命题。但 $\Omega^k(A,a)$ 是有实例的，故它实际上是可缩的，并且
  $\pi_k(A,a)=\trunc0{\Omega^k(A,a)}=\trunc0{\unit}=\unit$。
\end{proof}

\begin{lem} \label{lem:pik-nconnected}
  如果 $A$ 是 $n$-连通的且 $a:A$，那么对所有 $k\le{}n$ 有 $\pi_k(A,a)=\unit$。
\end{lem}

\begin{proof}
  我们有下列等式序列：
  %
  \begin{narrowmultline*}
    \pi_k(A,a) = \trunc0{\Omega^k(A,a)}
    = \Omega^k(\trunc k{(A,a)})
    = \Omega^k(\trunc k{\trunc n{(A,a)}})
    = \narrowbreak
      \Omega^k(\trunc k{\unit})
    = \Omega^k(\unit)
    = \unit.
  \end{narrowmultline*}
  %
  第三个等式利用了 $k\le{}n$ 这一事实，从而可以使用 $\truncf k\circ\truncf n=\truncf k$，而第四个等式利用了 $A$ 是 $n$-连通的这一事实。
\end{proof}

\begin{cor}
  当 $k < n$ 时，$\pi_k(\Sn ^n) = \unit$。
\end{cor}

\begin{proof}
  由 \cref{cor:sn-connected}，球面 $\Sn^n$ 是 $(n-1)$-连通的，因此我们可以应用 \cref{lem:pik-nconnected}。
\end{proof}

\section{纤维列与长正合列}
\label{sec:long-exact-sequence-homotopy-groups}

\index{fiber sequence|(}%
\index{sequence!fiber|(}%

如果一个函数 $f:X\to Y$ 的陪域配有一个基点 $y_0:Y$，那么 $f$ 在 $y_0$ 上的纤维 $F\defeq \hfib f {y_0}$ 就称为 \define{$f$ 的纤维 (the fiber of $f$)}\index{fiber}。
（若 $Y$ 是连通的，则 $F$ 仅在相差一个等价的意义下被确定；参见 \cref{ex:unique-fiber}。）
我们现在证明，如果 $X$ 也是带点类型且 $f$ 保持基点，那么 $F$、$X$ 和 $Y$ 的同伦群之间存在一个关系，具体表现为一个\emph{长正合列}。我们通过与这样一个 $f$ 相关联的\emph{纤维列}来推导这一结论。

\begin{defn}\label{def:pointedmap}
  带点类型 $(X,x_0)$ 与 $(Y,y_0)$ 之间的 \define{带点映射 (pointed map)}
  \indexdef{pointed!map}%
  \indexsee{function!pointed}{pointed map}%
  是一个映射 $f:X\to Y$，并附带一条道路 $f_0:f(x_0)=y_0$。
\end{defn}

对于任意带点类型 $(X,x_0)$ 和 $(Y,y_0)$，都存在一个带点映射 $(\lam{x} y_0) : X\to Y$，它是取值为基点的常值映射。
我们称此映射为 \define{零映射 (zero map)}\indexdef{zero!map}\indexdef{function!zero}，有时也将其记作 $0:X\to Y$。

回顾，每个带点类型 $(X,x_0)$ 都有一个环路空间 \index{loop space} $\Omega (X,x_0)$。
现在我们指出，这一运算在带点映射上具有函子性。\index{loop space!functoriality of}\index{functor!loop space}

\begin{defn}\label{def:loopfunctor}
  给定带点类型之间的带点映射 $f:X \to Y$，我们如下定义一个带点映射 $\Omega f:\Omega X \to \Omega Y$：
  \[(\Omega f)(p) \defeq \rev{f_0}\ct\ap{f}{p}\ct f_0.\]
  道路 $(\Omega f)_0 : (\Omega f) (\refl{x_0}) = \refl{y_0}$ 将 $\Omega f$ 展示为带点映射，它是类型如下的显然道路：
  \[\rev{f_0}\ct\ap{f}{\refl{x_0}}\ct f_0=\refl{y_0}.\]
\end{defn}

带点映射上还有另一个函子，它将 $f:X\to Y$ 映到 $\proj1 : \hfib f {y_0} \to X$。
当 $f$ 是带点映射时，我们总是将 $\hfib f {y_0}$ 视为以 $(x_0,f_0)$ 为基点的带点类型；在这种情况下，$\proj1$ 也是一个带点映射，其证据为 $(\proj1)_0 \defeq \refl{x_0}$。
因此，这个操作可以迭代进行。

\begin{defn}
  一个带点映射 $f:X\to Y$ 的 \define{纤维列（fiber sequence）}
  \indexdef{fiber sequence}%
  \indexdef{sequence!fiber}%
  是由带点类型与带点映射组成的无穷序列
  \[\xymatrix{\dots \ar[r]^{f^{(n+1)}} & X^{(n+1)} \ar[r]^{f^{(n)}} & X^{(n)} \ar^-{f^{(n-1)}}[r] & \dots \ar[r] & X^{(2)} \ar^-{f^{(1)}}[r] & X^{(1)} \ar[r]^{f^{(0)}} & X^{(0)}}\]
  递归地由
  \[ X^{(0)} \defeq Y\qquad
    X^{(1)} \defeq X\qquad
    f^{(0)} \defeq f\qquad
    \]
  与
  \begin{alignat*}{2}
    X^{(n+1)} &\defeq \hfib {f^{(n-1)}}{x^{(n-1)}_0}\\
    f^{(n)} &\defeq \proj1 &: X^{(n+1)} \to X^{(n)}.
  \end{alignat*}
  定义；
  其中 $x^{(n)}_0$ 表示 $X^{(n)}$ 的基点，按上述方式递归地选取。
\end{defn}

因此，这个纤维列中任意相邻的一对映射都具有如下形式：
\[ \xymatrix{ X^{(n+1)} \jdeq \hfib{f^{(n-1)}}{x^{(n-1)}_0} \ar[rr]^-{f^{(n)}\jdeq \proj1} && X^{(n)} \ar[r]^{f^{(n-1)}} & X^{(n-1)}. } \]
特别地，我们有 $f^{(n-1)} \circ f^{(n)} = 0$。
现在我们注意到，这个序列中出现的类型分别是底空间 $Y$、全空间\index{total!space} $X$ 以及纤维 $F\defeq \hfib f {y_0}$ 的迭代环路空间\index{loop space!iterated}，映射也类似。

\begin{lem}\label{thm:fiber-of-the-fiber}
  设 $f:X\to Y$ 是带点空间之间的一个带点映射。那么：
  \begin{enumerate}
  \item $f^{(1)}\defeq \proj1 : \hfib f {y_0} \to X$ 的纤维等价于 $\Omega Y$。\label{item:fibseq1}
  \item 类似地，$f^{(2)} : \Omega Y \to \hfib f {y_0}$ 的纤维等价于 $\Omega X$。\label{item:fibseq2}
  \item 在这些等价之下，带点映射 $f^{(3)} : \Omega X\to \Omega Y$ 等同于带点映射 $\Omega f\circ\rev{(\blank)}$。\label{item:fibseq3}
  \end{enumerate}
\end{lem}

\begin{proof}
  对于\ref{item:fibseq1}，我们有
  \begin{align*}
    \hfib{f^{(1)}}{x_0}
    &\defeq \sm{z:\hfib{f}{y_0}} (\proj{1}(z) = x_0)\\
    &\eqvsym \sm{x:X}{p:f(x)=y_0} (x = x_0) &\text{(by \cref{ex:sigma-assoc})}\\
    &\eqvsym (f(x_0) = y_0) &\text{(as $\tsm{x:X} (x=x_0)$ is contractible)}\\
    &\eqvsym (y_0 = y_0) &\text{(by $(f_0 \ct \blank)$)}\\
    &\jdeq \Omega Y.
  \end{align*}。
  追踪计算过程可知，这个等价将 $((x,p),q)$ 映到 $\opp{f_0} \ct \ap{f}{\opp q} \ct p$，而其逆将 $r:y_0=y_0$ 映到 $((x_0, f_0 \ct r), \refl{x_0})$。
  特别地，$\hfib{f^{(1)}}{x_0}$ 的基点 $((x_0,f_0),\refl{x_0})$ 被映到 $\opp{f_0} \ct \ap{f}{\opp {\refl{x_0}}} \ct f_0$，后者等于 $\refl{y_0}$。
  因此这个等价是一个带点映射（参见 \cref{ex:pointed-equivalences}）。
  而且，在这个等价之下，$f^{(2)}$ 被等同于 $\lam{r} (x_0, f_0 \ct r): \Omega Y \to \hfib f {y_0}$。

  \cref{item:fibseq2} 通过将\ref{item:fibseq1} 应用于 $f^{(1)}$（代替 $f$）立即得到。
  % The resulting equivalence sends $(((x,p),p'),q') : \hfib{f^{(2)}}{(x_0,f_0)}$ to $\ap{f}{\opp {q'}} \ct p'$, while its inverse sends $s:x_0=x_0$ to $(((x_0,f_0), s), \refl{(x_0,f_0)})$.
  由于 $(f^{(1)})_0 \defeq \refl{x_0}$，在这个等价下 $f^{(3)}$ 被等同于映射 $\Omega X \to \hfib {f^{(1)}}{x_0}$，后者由 $s \mapsto ((x_0,f_0),s)$ 定义。
  因此，当我们与前述等价 $\hfib {f^{(1)}}{x_0} \eqvsym \Omega Y$ 复合时，可以看到 $s$ 被映到 $\opp{f_0} \ct \ap{f}{\opp s} \ct f_0$，根据定义这就是 $(\Omega f)(\opp s)$。我们略去以下证明：这是带点映射之间的相等，而不仅仅是函数之间的相等。
\end{proof}

于是，$f:X\to Y$ 的纤维列可以图示如下：
\[\xymatrix@C=1.5pc{
  \dots \ar[r] &
  \Omega^2 X \ar[r]^-{\Omega^2 f} &
  \Omega^2 Y \ar[r]^-{-\Omega \partial} &
  \Omega F \ar[r]^-{-\Omega i} &
  \Omega X \ar[r]^-{-\Omega f} &
  \Omega Y \ar[r]^-{\partial} &
  F \ar[r]^-{i} &
  X \ar[r]^{f} & Y.}\]
其中减号表示与道路逆 $\rev{(\blank)}$ 的复合。
注意，由 \cref{ex:ap-path-inversion} 可得
\[ \Omega\left(\Omega f\circ \opp{(\blank)}\right) \circ \opp{(\blank)}
= \Omega^2 f \circ \opp{(\blank)} \circ \opp{(\blank)}
= \Omega^2 f.
\]
因此，当 $k$ 是奇数时，$k$ 重环路映射上就会出现减号。

从这个纤维列我们将推导出一个\emph{带点集的正合列}。
\index{image}%
设 $A$ 和 $B$ 是集合，$f:A\to B$ 是一个函数；回顾 \emph{像} $\im(f)$ 的定义（参见 \cref{defn:modal-image}），它可以视为 $B$ 的一个子集：
\[\im(f) \defeq \setof{b:B | \exis{a:A} f(a)=b}. \]
如果 $A$ 和 $B$ 还分别带有基点 $a_0$ 和 $b_0$，并且 $f$ 是一个带点映射，则定义 $f$ 的\define{核（kernel）}
\indexdef{kernel}%
\indexdef{pointed!map!kernel of}%
为 $A$ 的如下子集：
\symlabel{kernel}
\[\ker(f) \defeq \setof{x:A | f(x) = b_0}. \]
当然，这其实就是 $f$ 在基点 $b_0$ 上的纤维；由于 $B$ 是集合，所以它是 $A$ 的子集。

注意，任何群都是一个带点集，其单位元即为基点；任何群同态都是一个带点映射。
在这种情形下，核与像的定义与通常群论中的概念一致。

\begin{defn}
  一个\define{带点集的正合列（exact sequence of pointed sets）}
  \indexdef{exact sequence}%
  \indexdef{sequence!exact}%
  是一个（可能有界的）带点集与带点映射的序列：
  \[\xymatrix{\dots \ar[r] & A^{(n+1)} \ar[r]^-{f^{(n)}} & A^{(n)} \ar[r]^{f^{(n-1)}} & A^{(n-1)} \ar[r] &
    \dots}\]
  使得对每个 $n$，$f^{(n)}$ 的像作为 $A^{(n)}$ 的子集等于 $f^{(n-1)}$ 的核。
  换言之，对所有 $a:A^{(n)}$，有
  \[ (f^{(n-1)}(a) = a^{(n-1)}_0) \iff \exis{b:A^{(n+1)}} (f^{(n)}(b)=a). \]
  其中 $a^{(n)}_0$ 表示 $A^{(n)}$ 的基点。
\end{defn}

通常，正合列中的大多数或全部带点集都是群，且常为交换群。
当我们提到\define{群的正合列（exact sequence of groups）}时，还假定其中的映射是群同态，而不仅仅是带点映射。

\index{exact sequence}


\begin{thm}\label{thm:les}
  设 $f:X \to Y$ 是带点空间之间的一个带点映射，其纤维为 $F\defeq \hfib f {y_0}$。
  则我们得到如下长正合列，其中除最后三项外都是群；除最后六项外，都是交换群。

  \[
  \xymatrix@R=1.2pc@C=3pc{
    \vdots & \vdots & \vdots \ar[lld] \\
    \pi_k(F) \ar[r] & \pi_k(X) \ar[r] & \pi_k(Y) \ar[lld] \\
    \vdots & \vdots & \vdots \ar[lld] \\
    \pi_2(F) \ar[r] & \pi_2(X) \ar[r] & \pi_2(Y) \ar[lld] \\
    \pi_1(F) \ar[r] & \pi_1(X) \ar[r] & \pi_1(Y) \ar[lld]\\
    \pi_0(F) \ar[r] & \pi_0(X) \ar[r] & \pi_0(Y)}
  \]
\end{thm}

% In US trade we might want a page break here, and extra stretch, otherwise the
% whole page looks ugly.
\vspace*{0pt plus 10ex}
\goodbreak

\begin{proof}
  我们首先证明，一个纤维列的 $0$-截断是一个带点集的正合列。
  因此，对于纤维列中任意相邻的一对映射：
  %
  \[\xymatrix{\hfib{f}{z_0} \ar^-g[r] & W \ar^-f[r] & Z}\]
  %
  其中 $g\defeq \proj1$，我们需要证明列
  %
  \[\xymatrix{\trunc{0}{\hfib{f}{z_0}} \ar^-{\trunc0g}[r] & \trunc0{W}
    \ar^-{\trunc0{f}}[r] & \trunc0{Z}}\]
  %
  是正合的，即证明 $\im(\trunc0g)\subseteq\ker(\trunc0f)$ 与 $\ker(\trunc0f)\subseteq\im(\trunc0g)$。

  第一个包含关系等价于 $\trunc0f\circ\trunc0g=0$，这由 $\truncf0$ 的函子性以及 $f\circ g=0$ 可得。
  对于第二个包含关系，我们假设 $w':\trunc0W$ 且 $p':\trunc0f(w')=\tproj0{z_0}$，并证明仅存在 ${t:\hfib{f}{z_0}}$ 使得 $\tproj0{g(t)}=w'$。
  由于我们的目标是一个命题，可以假定 $w'$ 具有 $\tproj0w$ 的形式，其中 $w:W$。
  现在，根据 \cref{thm:path-truncation}，由 $p' : \tproj0{f(w)}=\tproj0{z_0}$ 可得 $p'': \trunc{-1}{f(w)=z_0}$，因此通过进一步的截断归纳，我们可以假设存在某个 $p:f(w)=z_0$。
  但这样我们就得到了 $\tproj0{(w,p)}:\trunc{0}{\hfib{f}{z_0}}$，其在 $\trunc0g$ 下的像正是 $\tproj0w\jdeq w'$，如所需。

  因此，将 $\truncf0$ 应用于 $f$ 的纤维列，我们得到了一个长正合列，其中按所需顺序包含了带点集 $\pi_k(F)$、$\pi_k(X)$ 和 $\pi_k(Y)$。
  当然，对于 $k\ge1$，$\pi_k$ 是一个群（因为它是环路空间的 $0$-截断），并且对于 $k\ge 2$，由 Eckmann--Hilton 论证可知它是交换群
  \index{Eckmann--Hilton argument} (\cref{thm:EckmannHilton})。
  此外，\cref{thm:fiber-of-the-fiber} 允许我们将这个正合列中的映射 $\pi_k(F) \to \pi_k(X)$ 和 $\pi_k(X) \to \pi_k(Y)$ 分别识别为 $(-1)^k \pi_k(i)$ 和 $(-1)^k \pi_k(f)$。

  更一般地，这个长正合列中除最后三个映射外，每个映射都具有 $\trunc0{\Omega h}$ 或 $\trunc0{-\Omega h}$ 的形式，其中 $h$ 是某个映射。
  在前一种情形下，它是群同态；而在后一种情形下，若所涉及的群是交换群，它也是同态；否则它是一个``反同态''。
  然而，当我们把一个群同态替换为它的负映射时，其核与像保持不变，因此任何涉及该映射的列的正合性也保持不变。
  于是，我们可以修改我们的长正合列，得到一个直接包含 $\pi_k(i)$ 和 $\pi_k(f)$ 的列，并且其中所有映射（除最后三个外）都是群同态。
\end{proof}

交换群的正合列的通常性质\index{group!abelian!exact sequence of}可以按通常方式证明。
特别地，我们有：


\begin{lem}\label{thm:ses}
  给定一个交换群的正合列：
  %
  \[\xymatrix{K \ar[r]& G \ar^f[r] & H \ar[r] & Q.}\]
  %
  \begin{enumerate}
  \item 若 $K=0$，则 $f$ 是单的。\label{item:sesinj}
  \item 若 $Q=0$，则 $f$ 是满的。\label{item:sessurj}
  \item 若 $K=Q=0$，则 $f$ 是同构。\label{item:sesiso}
  \end{enumerate}
\end{lem}

\begin{proof}
  由于 $f$ 的核是 $K\to G$ 的像，若 $K=0$，则 $f$ 的核是 $\{0\}$；
  因此 $f$ 是单的，因为它是群同态。
  类似地，由于 $f$ 的像是 $H\to Q$ 的核，若 $Q=0$，则 $f$ 的像是整个 $H$，所以 $f$ 是满的。
  最后，\ref{item:sesiso} 由 \ref{item:sesinj} 与 \ref{item:sessurj} 根据 \cref{thm:mono-surj-equiv} 得出。
\end{proof}

作为一个直接应用，我们现在可以具体说明，一个映射的 $n$-连通性在何种意义上比它在 $n$-截断上诱导等价更强。

\begin{cor}\label{thm:conn-pik}
  设 $f:A\to B$ 是 $n$-连通的，$a:A$，并定义 $b\defeq f(a)$。那么：
  \begin{enumerate}
  \item 若 $k\le n$，则 $\pi_k(f):\pi_k(A,a) \to \pi_k(B,b)$ 是同构。\label{item:conn-pik1}
  \item 若 $k=n+1$，则 $\pi_k(f):\pi_k(A,a) \to \pi_k(B,b)$ 是满的。\label{item:conn-pik2}
  \end{enumerate}
\end{cor}

\begin{proof}
  当 $k=0$ 时，第~\ref{item:conn-pik1}部分可由~\cref{lem:connected-map-equiv-truncation}得出，只需注意到 $\pi_0(f)\equiv\trunc 0f$。
  当 $k=0$ 时，第~\ref{item:conn-pik2}部分可由~\cref{ex:is-conn-trunc-functor}得出，注意到一个函数是满的当且仅当它是 $(-1)$-连通的，依据~\cref{thm:minusoneconn-surjective}。
  当 $k>0$ 时，作为长正合列的一部分，我们得到一个正合列
  \[\xymatrix{\pi_k(\hfib f b) \ar[r]& \pi_k(A,a) \ar^f[r] & \pi_k(B,b) \ar[r] & \pi_{k-1}(\hfib f b).}\]
  现在，由于 $f$ 是 $n$-连通的，$\trunc n{\hfib f b}$ 是可缩的。
  因此，若 $k\le n$，则~$\pi_k(\hfib f b) = \trunc0{\Omega^k(\hfib f b)} = \Omega^k(\trunc k{\hfib f b})$也是可缩的。
  这样一来，由~\cref{thm:ses}\ref{item:sesiso}可知，当 $k\le n$ 时 $\pi_k(f)$ 是同构；而当 $k=n+1$ 时，由~\cref{thm:ses}\ref{item:sessurj}可知它是满的。
\end{proof}

在~\cref{sec:whitehead}中我们将会看到，\cref{thm:conn-pik}的逆命题也成立。

\index{fiber sequence|)}%
\index{sequence!fiber|)}%

\section{Hopf 纤维化}
\label{sec:hopf}

在本节中，我们将定义\define{Hopf 纤维化（Hopf fibration）}。
\indexdef{Hopf!fibration}%

\begin{thm}[Hopf 纤维化]\label{thm:hopf-fibration}
  存在 $\Sn ^2$ 上的一个纤维化 $H$，它在基点处的纤维是 $\Sn ^1$，全空间是 $\Sn ^3$。
\end{thm}

Hopf 纤维化将帮助我们计算球面的若干同伦群。
事实上，它给出了下面的同伦群长正合列
\index{homotopy!group!of sphere}
\index{sequence!exact}
（参见
\cref{sec:long-exact-sequence-homotopy-groups}）：
%
\[
\xymatrix@R=1.2pc{
  \pi_k(\Sn^1) \ar[r] & \pi_k(\Sn^3) \ar[r] & \pi_k(\Sn^2) \ar[lld] \\
  \vdots & \vdots & \vdots \ar[lld] \\
  \pi_2(\Sn^1) \ar[r] & \pi_2(\Sn^3) \ar[r] & \pi_2(\Sn^2) \ar[lld] \\
  \pi_1(\Sn^1) \ar[r] & \pi_1(\Sn^3) \ar[r] & \pi_1(\Sn^2)}
\]
%
我们已经计算了所有的 $\pi_n(\Sn^1)$，以及当 $k<n$ 时的 $\pi_k(\Sn^n)$，因此上式化为：
%
\[
\xymatrix@R=1.2pc{
  0 \ar[r] & \pi_k(\Sn^3) \ar[r] & \pi_k(\Sn^2) \ar[lld] \\
  \vdots & \vdots & \vdots \ar[lld] \\
  0 \ar[r] & \pi_3(\Sn^3) \ar[r] & \pi_3(\Sn^2) \ar[lld] \\
  0 \ar[r] & 0 \ar[r] & \pi_2(\Sn^2) \ar[lld] \\
  \Z \ar[r] & 0 \ar[r] & 0}
\]
%
特别地，我们得到如下结果：

\begin{cor} \label{cor:pis2-hopf}
  对于每个 $k\ge3$，有 $\eqv{\pi_2(\Sn^2)}{\Z}$ 和 $\eqv{\pi_k(\Sn^3)}{\pi_k(\Sn^2)}$（其中该映射由 Hopf 纤维化诱导，视为从全空间 $\Sn^3$ 到基空间 $\Sn^2$ 的映射）。
\end{cor}

事实上，我们可以更进一步：Hopf 纤维化的纤维列将表明 $\Omega^3(\Sn^3)$ 是从 $\Omega^3(\Sn^2)$ 到 $\Omega^2(\Sn^1)$ 的某个映射的纤维。
由于 $\Omega^2(\Sn^1)$ 是可缩的，我们有 $\eqv{\Omega^3(\Sn^3)}{\Omega^3(\Sn^2)}$。
在经典同伦论中，这个事实会是 \cref{cor:pis2-hopf} 和 Whitehead 定理的推论，但 Whitehead 定理在同伦类型论中不一定成立（参见 \cref{sec:whitehead}）。
不过，我们在这里不会使用更精确的版本。

\subsection{推出上的纤维化}
\label{sec:fib-over-pushout}

我们首先给出一个引理，说明如何在推出上构造纤维化。
\index{pushout}%

\begin{lem}\label{lem:fibration-over-pushout}
  设 $\Ddiag=(Y\xleftarrow{j}X\xrightarrow{k}Z)$ 是一个伸展\index{span}，并假设我们有：
  \begin{itemize}
  \item 两个纤维化 $E_Y:Y\to\type$ 和 $E_Z:Z\to\type$。
  \item $E_Y\circ j:X\to\type$ 与 $E_Z\circ k:X\to\type$ 之间的一个等价 $e_X$，即
    \[e_X:\prd{x:X}\eqv{E_Y(j(x))}{E_Z(k(x))}.\]
  \end{itemize}

  那么我们可以构造一个纤维化 $E:Y\sqcup^XZ\to\type$，使得：
  \begin{itemize}
  \item 对于所有 $y:Y$，$E(\inl(y))\judgeq E_Y(y)$。
  \item 对于所有 $z:Z$，$E(\inr(z))\judgeq E_Z(z)$。
  \item 对于所有 $x:X$，$\map E{\glue(x)}=\ua(e_X(x))$（注意等式两边都是 $\type$ 中从 $E_Y(j(x))$ 到 $E_Z(k(x))$ 的道路）。
  \end{itemize}
  此外，该纤维化的全空间构成如下推出方块：
  \[\xymatrix{ \sm{x:X}E_Y(j(x)) \ar[r]_\sim^{\idfunc\times e_X}
    \ar[d]_{j\times\idfunc} &
    \sm{x:X}E_Z(k(x)) \ar[r]^-{k\times\idfunc}
    & \sm{z:Z}E_Z(z) \ar[d]^{\inr} \\
    \sm{y:Y}E_Y(y) \ar[rr]_{\inl} & & \sm{t:Y\sqcup^XZ}E(t) }\]
\end{lem}

\begin{proof}
  我们利用推出 $Y\sqcup^XZ$ 的递归原理来定义 $E$。
  为此，我们需要指定 $E$ 在形如 $\inl(y)$、$\inr(z)$ 的元素上的值，以及 $E$ 在道路 $\glue(x)$ 上的作用，因此可以直接选取以下值：
  \begin{align*}
    E(\inl(y)) & \defeq E_Y(y),\\
    E(\inr(z)) & \defeq E_Z(z),\\
    \map E{\glue(x)} & \defid \ua(e_X(x)).
  \end{align*}
  %
  为证明该纤维化的全空间是一个推出，我们应用展平引理（\cref{thm:flattening}）并取以下值：\index{flattening lemma}
  \begin{itemize}
  \item $A\defeq Y+Z$，$B\defeq X$，且 $f,g:B\to A$ 定义为
    $f(x)\defeq\inl(j(x))$，$g(x)\defeq\inr(k(x))$，
  \item 类型族 $C:A\to\type$ 定义为
    \begin{equation*}
      C(\inl(y)) \defeq E_Y(y)
      \qquad\text{and}\qquad
      C(\inr(z)) \defeq E_Z(z),
    \end{equation*}
  \item 等价族 $D:\prd{b:B}C(f(b))\eqvsym C(g(b))$ 定义为 $e_X$。
  \end{itemize}
  %
  展平引理中的基础高阶归纳类型 $W$ 等价于推出 $Y\sqcup^XZ$，而类型族 $P:Y\sqcup^XZ\to\type$ 等价于上面定义的 $E$。

  因此展平引理告诉我们 $\sm{t:Y\sqcup^XZ}E(t)$ 等价于具有以下生成元的高阶归纳类型 ${E^{\mathrm{tot}}}'$：
  %
  \begin{itemize}
  \item 一个函数 $\mathsf{z}:\sm{a:Y+Z}C(a)\to {E^{\mathrm{tot}}}'$，
  \item 对每个 $x:X$ 和 $t:E_Y(j(x))$，一条道路
    \narrowequation{\mathsf{z}(\inl(j(x)),t)=\mathsf{z}(\inr(k(x)),e_X(t))。}
  \end{itemize}
  %
  再次使用展平引理或直接计算，容易看出 $\eqv{\sm{a:Y+Z}C(a)}{\sm{y:Y}E_Y(y)+\sm{z:Z}E_Z(z)}$，因此
  ${E^{\mathrm{tot}}}'$ 等价于具有以下生成元的高阶归纳类型 $E^{\mathrm{tot}}$：
  %
  \begin{itemize}
  \item 一个函数 $\inl:\sm{y:Y}E_Y(y)\to E^{\mathrm{tot}}$，
  \item 一个函数 $\inr:\sm{z:Z}E_Z(z)\to E^{\mathrm{tot}}$，
  \item 对每个 $(x,t):\sm{x:X}E_Y(j(x))$，一条道路
    \narrowequation{\glue(x,t):\inl(j(x),t) = \inr(k(x),e_X(t))。}
  \end{itemize}
  %
  因此，$E$ 的全空间是 $E_Y$ 和 $E_Z$ 的全空间的推出，正如所需。
\end{proof}

\subsection{Hopf 构造}

\begin{defn}
  一个 \define{H-空间（H-space）}
  \indexdef{H-space}%
  由以下结构组成：
  \begin{itemize}
  \item 一个类型 $A$，
  \item 一个基点 $e:A$，
  \item 一个二元运算 $\mu:A\times A\to A$，以及
  \item 对于每个 $a:A$，有相等 $\mu(e,a)=a$ 和 $\mu(a,e)=a$。
  \end{itemize}
\end{defn}

\begin{lem}
  设 $A$ 是一个连通的 H-空间。那么对于每个 $a:A$，映射 $\mu(a,\blank):A\to
  A$ 和 $\mu(\blank,a):A\to A$ 都是等价。
\end{lem}

\begin{proof}
  我们来证明对于每个 $a:A$，映射 $\mu(a,\blank)$ 是一个等价。另一个结论同理可证。
  %
  ``$\mu(a,\blank)$ 是等价''这一断言对应于一个类型族
  $P:A\to\prop$，而证明它则对应于寻找该类型族的一个截面。

  类型 $\prop$ 是一个集合（\cref{thm:hleveln-of-hlevelSn}），因此我们可以
  通过 $P'(\tproj0a)\defeq
  P(a)$ 定义一个新的类型族 $P':\trunc0A\to\prop$。但根据假设，$A$ 是连通的，因此 $\trunc0A$
  是可缩的。这蕴含着，要找到 $P'$ 的一个截面，只需在 $P'$ 于 $\tproj0e$ 处的纤维中找到一个点即可。而我们有
  $P'(\tproj0e)=P(e)$，该类型是有实例的，因为根据 H-空间的定义，$\mu(e,\blank)$ 等于恒等映射，因而是一个等价。

  我们已经证明了对于每个 $x:\trunc0A$，命题 $P'(x)$ 为真，
  因此特别地，对于每个 $a:A$，命题 $P(a)$ 也为真，因为
  $P(a)$ 就是 $P'(\tproj0a)$。
\end{proof}

\begin{defn}
  设 $A$ 是一个连通的 H-空间。我们利用
  \cref{lem:fibration-over-pushout} 在 $\susp A$ 上构造一个纤维化。

  鉴于 $\susp A$ 是推出 $\unit\sqcup^A\unit$，我们可以通过指定以下内容来定义 $\susp A$ 上的一个纤维化：
  \begin{itemize}
  \item $\unit$ 上的两个纤维化（即两个类型 $F_1$ 和 $F_2$），以及
  \item 一族 $e:A\to(\eqv{F_1}{F_2})$ $F_1$ 与 $F_2$ 之间的等价，$A$ 中每个元素各对应一个。
  \end{itemize}
  %
  我们取 $F_1$ 和 $F_2$ 都为 $A$，并对 $a:A$ 取等价 $\mu(a,\blank)$ 作为 $e(a)$。
\end{defn}

根据 \cref{lem:fibration-over-pushout}，我们得到如下图表：
%
\[\xymatrix{A \ar@{->>}[d] & A \times A \ar[l]_-{\proj2} \ar@{->>}_{\proj1}[d]
  \ar[r]^-{\mu} & A \ar@{->>}[d] \\
  1 & A \ar[r] \ar[l] & 1}\]
%
我们刚才构造的纤维化是 $\susp A$ 上的一个纤维化，其全空间是上方那一行的推出。

此外，利用 $f(x,y)\defeq(\mu(x,y),y)$ 我们得到如下图表：
%
\[\xymatrix{A \ar_{\idfunc}[d] & A \times A \ar[l]_-{\proj2} \ar^f[d]
  \ar[r]^-{\mu} & A \ar^{\idfunc}[d] \\
  A & A\times A \ar^-{\proj2}[l] \ar_-{\proj1}[r] & A}\]
%
该图表交换，并且三个垂直映射都是等价，其中 $f$ 的逆是函数 $g$，定义为
\[g(u,v)\defeq(\opp{\mu(\blank,v)}(u),v).\]
%
这表明两行是等价的（从而相等的）伸展，因此我们所构造纤维化的全空间等价于下方那一行的推出。
而根据定义，这个推出就是 $A$ 与自身的\emph{联结}（见 \cref{sec:colimits}）。
%
于是我们证明了：

\begin{lem}\label{lem:hopf-construction}
  给定一个连通的 H-空间 $A$，存在一个纤维化，称为
  \define{Hopf 构造（Hopf construction）}，
  \indexdef{Hopf!construction}%
  其底空间为 $\susp A$，纤维为 $A$，全空间为 $A*A$。
\end{lem}

\subsection{Hopf 纤维化}

\index{Hopf fibration|(}
\indexsee{fibration!Hopf}{Hopf fibration}
我们将首先在圆 $\Sn^1$ 上构造一个 H-空间结构，从而由
\cref{lem:hopf-construction} 我们将得到一个底空间为 $\Sn^2$，纤维为
$\Sn^1$ 且全空间为 $\Sn^1*\Sn^1$ 的纤维化。然后我们将证明这个联结等价于 $\Sn^3$。

\begin{lem}\label{lem:hspace-S1}
  在圆 $\Sn^1$ 上存在一个 H-空间结构。
\end{lem}

\begin{proof}
  对于 H-空间结构的基点，我们选择 $\base$。
  %
  现在我们需要定义乘法运算
  $\mu:\Sn^1\times\Sn^1\to\Sn^1$。
  我们将通过对 $\Sn^1$ 作递归来定义 $\mu$ 的柯里化形式 $\widetilde\mu:\Sn^1\to(\Sn^1\to\Sn^1)$：
  \begin{equation*}
    \widetilde\mu(\base) \defeq\idfunc[\Sn^1],
    \qquad\text{and}\qquad
    \ap{\widetilde\mu}{\lloop} \defid\funext(h).
  \end{equation*}
  其中 $h:\prd{x:\Sn^1}(x=x)$ 是 \cref{thm:S1-autohtpy} 中定义的函数，
  它具有性质 $h(\base) \defeq\lloop$。

  现在只需证明对每个 $x:\Sn^1$，有 $\mu(x,\base)=\mu(\base,x)=x$。
  %
  根据定义，对于 $x:\Sn^1$，我们有
  $\mu(\base,x)=\widetilde\mu(\base)(x)=\idfunc[\Sn^1](x)=x$。为了证明等式
  $\mu(x,\base)=x$，我们对 $x:\Sn^1$ 进行归纳：
  \begin{itemize}
  \item 如果 $x$ 是 $\base$，那么根据定义 $\mu(\base,\base)=\base$，因此我们有 $\refl\base:\mu(\base,\base)=\base$。
  \item 当 $x$ 沿 $\lloop$ 变化时，我们需要证明
    \[\refl\base\ct\apfunc{\lam{x}x}({\lloop})=
    \apfunc{\lam{x}\mu(x,\base)}({\lloop})\ct\refl\base.\]
    左边等于 $\lloop$，而对于右边我们有：
    \begin{align*}
      \apfunc{\lam{x}\mu(x,\base)}({\lloop})\ct\refl\base &=
      \apfunc{\lam{x}(\widetilde\mu(x))(\base)}({\lloop})\\
      &=\happly(\apfunc{\lam{x}(\widetilde\mu(x))}({\lloop}),\base)\\
      &=\happly(\funext(h),\base)\\
      &=h(\base)\\
      &=\lloop. \qedhere
    \end{align*}
  \end{itemize}
\end{proof}

现在回顾 \cref{sec:colimits}，类型 $A$ 和 $B$ 的\emph{联结} $A*B$ 是图表
\index{join!of types}%
\[A \xleftarrow{\proj1}A\times B \xrightarrow{\proj2} B. \] 的推出。

\begin{lem}
  \index{associativity!of join}%
  联结运算是结合的：若 $A$、$B$ 和 $C$ 是三个类型，
  则我们有等价 $\eqv{(A*B)*C}{A*(B*C)}$。
\end{lem}

\begin{proof}
  我们通过归纳来定义一个映射 $f:(A*B)*C\to A*(B*C)$。首先需要定义
  $f\circ\inl:A*B\to A*(B*C)$，这也将通过归纳来完成，然后定义
  $f\circ\inr:C\to A*(B*C)$，接着定义 $\apfunc{f}\circ\glue:\prd{t:(A*B)\times
    C}f(\inl(\fst(t)))=f(\inr(\snd(t)))$，后者需要对 $t$ 的第一个分量进行归纳：
  \begin{align*}
    (f\circ\inl)(\inl(a)) &\defeq \inl(a), \\
    (f\circ\inl)(\inr(b)) &\defeq \inr(\inl(b)), \\
    \apfunc{f\circ\inl}(\glue(a,b)) &\defid \glue(a,\inl(b)), \\
    f(\inr(c)) &\defeq \inr(\inr(c)),\\
    \apfunc{f}(\glue(\inl(a),c)) &\defid\glue(a,\inr(c)),\\
    \apfunc{f}(\glue(\inr(b),c)) &\defid\apfunc{\inr}(\glue(b,c)),\\
    \apdfunc{\lam{x}\apfunc{f}(\glue(x,c))}(\glue(a,b)) &\defid
    ``\apdfunc{\lam{x}\glue(a,x)}(\glue(b,c))''.
  \end{align*}
  %
  对于最后一个等式，注意其右侧的类型是
  \[\transfib{\lam{x}\inl(a)=\inr(x)}{\glue(b,c)}{\glue(a,\inl(b))}=
  \glue(a,\inr(c))\]
  而它本应具有类型
  %
  \begin{narrowmultline*}
    \transfib{\lam{x}f(\inl(x))=f(\inr(c))}{\glue(a,b)}{\apfunc{f}(\glue(\inl(a),c))}
    = \narrowbreak
    \apfunc{f}(\glue(\inr(b),c)).
  \end{narrowmultline*}
  %
  但根据定义中前面的几条，这两个类型都等价于以下类型：
  \[\glue(a,\inr(c))=\glue(a,\inl(b))\ct\apfunc\inr(\glue(b,c)),\]
  因此，我们可以通过一个等价进行强制转换，从而得到所需的元素。
  %
  类似地，我们可以定义一个映射 $g:A*(B*C)\to(A*B)*C$，而验证 $f$ 和
  $g$ 互为逆映射，则是一段冗长繁琐但在本质上直截了当的计算。
\end{proof}

一个更具概念性的证明概要如下。

\begin{proof}
  考虑如下图表，其中的映射都是自然的投影：
  \[\xymatrix{
    A & A\times C \ar[l] \ar[r] & A\times C\\
    A\times B \ar[u] \ar[d] & A\times B\times C \ar[l]\ar[u]\ar[r]\ar[d] &
    A\times C \ar[u] \ar[d] \\
    B & B\times C \ar[l] \ar[r] & C}\]
  %
  对各列取余极限，得到如下图表，其余极限是 $(A*B)*C$：
  \[\xymatrix{A*B & (A*B)\times C \ar[l]\ar[r] & C}\]
  %
  另一方面，对各行取余极限，得到一个图表，其余极限是 $A*(B*C)$。

  因此，利用一个关于余极限的类 Fubini 定理（我们尚未证明），便可得到等价 $\eqv{(A*B)*C}{A*(B*C)}$。不过，这个关于余极限的 Fubini 定理
  \index{Fubini theorem for colimits}
  其证明仍需那段冗长繁琐的计算。
\end{proof}

\begin{lem}
  对于任意类型 $A$，存在等价 $\eqv{\susp A}{\bool*A}$。
\end{lem}

\begin{proof}
  定义来回两个方向的映射很容易，证明它们互为逆映射也很直接。详情留给读者作为练习。
\end{proof}

现在我们可以构造 Hopf 纤维化：

\begin{thm}
  存在一个以 $\Sn^2$ 为底、$\Sn^1$ 为纤维、$\Sn^3$ 为全空间的纤维化。
  \index{total!space}%
\end{thm}

\begin{proof}
  我们先前证明了 $\Sn^1$ 具有 H-空间结构（参见 \cref{lem:hspace-S1}），
  因此根据 \cref{lem:hopf-construction}，存在一个以 $\Sn^2$ 为底、$\Sn^1$ 为纤维、$\Sn^1*\Sn^1$ 为全空间的纤维化。
  但由前两个结果以及 \cref{thm:suspbool}，我们有：
  \begin{equation*}
    \Sn^1*\Sn^1 = (\susp\bool)*\Sn^1
    =(\bool*\bool)*\Sn^1
    =\bool*(\bool*\Sn^1)
    =\susp(\susp\Sn^1)
    =\Sn^3. \qedhere
  \end{equation*}
\end{proof}


\index{Hopf fibration|)}

\section{Freudenthal 纬悬定理}
\label{sec:freudenthal}

\index{Freudenthal suspension theorem|(}%
\index{theorem!Freudenthal suspension|(}%

在证明 Freudenthal 纬悬定理之前，我们需要一些关于连通性的辅助引理。
在\cref{cha:hlevels}中，我们针对固定的 $n$ 证明了关于 $n$-连通映射与 $n$-类型的一系列事实；现在我们关心的是当 $n$ 变化时会发生什么。
例如，在\cref{prop:nconnected_tested_by_lv_n_dependent types}中我们证明了 $n$-连通映射可由一个相对于 $n$-类型族的``归纳原理''来刻画。
如果我们想沿着一个 $n$-连通映射，向一个 $k>n$ 的 $k$-类型族作归纳，我们并不能立即知道存在一个满足该归纳原理的函数，但下面的引理表明，至少我们的无知是可以被量化的。

\begin{lem}\label{thm:conn-trunc-variable-ind}
  若 $f:A\to B$ 是 $n$-连通的，并且 $P:B\to \ntype{k}$ 是 $k$-类型族（其中 $k\ge n$），则诱导函数
  \[ (\blank\circ f) : \Parens{\prd{b:B} P(b)} \to \Parens{\prd{a:A} P(f(a)) } \]
  是 $(k-n-2)$-截断的。
\end{lem}

\begin{proof}
  我们对自然数 $k-n$ 进行归纳。
  当 $k=n$ 时，这就是\cref{prop:nconnected_tested_by_lv_n_dependent types}。
  %
  对于归纳步骤，假设 $f$ 是 $n$-连通的且 $P$ 是 $(k+1)$-类型族。
  为证明 $(\blank\circ f)$ 是 $(k-n-1)$-截断的，令$\ell:\prd{a:A} P(f(a))$；则我们有
  \[ \hfib{(\blank\circ f)}{\ell} \eqvsym \sm{g:\prd{b:B} P(b)} \prd{a:A} g(f(a)) = \ell(a).\]
  设 $(g,p)$ 和 $(h,q)$ 属于此类型，即 $p:g\circ f \htpy \ell$ 且 $q:h\circ f \htpy \ell$；则我们还知道
  \[ \big((g,p) = (h,q)\big) \eqvsym
  \Parens{\sm{r:g\htpy h} r\circ f = p \ct \opp{q}}.
  \]
  然而，此处右侧是映射
  \[ (\blank\circ f) : \Parens{\prd{b:B} Q(b)} \to \Parens{\prd{a:A} Q(f(a)) } \]
  的一个纤维，其中$Q(b) \defeq (g(b)=h(b))$。
  由于 $P$ 是 $(k+1)$-类型族，$Q$ 是 $k$-类型族，故归纳假设表明该纤维是 $(k-n-2)$-类型。
  因此，$\hfib{(\blank\circ f)}{\ell}$ 的所有路径空间都是 $(k-n-2)$-类型，从而它是一个 $(k-n-1)$-类型。
\end{proof}

回顾一下，如果 $\pairr{A,a_0}$ 和 $\pairr{B,b_0}$ 是带点类型，那么它们的 \define{楔积}
\index{wedge}%
$A\vee B$ 定义为由 $A\xleftarrow{a_0}
\unit\xrightarrow{b_0} B$ 给出的推出。
存在一个标准映射 $i:A\vee B \to A\times B$，由两个映射 $\lam{a} (a,b_0)$ 和 $\lam{b} (a_0,b)$ 定义；下面的引理本质上说，如果 $A$ 和 $B$ 是高度连通的，那么这个映射也是高度连通的。
不过，如果我们使用 \cref{prop:nconnected_tested_by_lv_n_dependent types} 中连通性的刻画，并将其代入楔积的（推广到类型族的）泛性质，那么无论是证明还是使用起来都更方便一些。

\begin{lem}[楔连通性引理]\label{thm:wedge-connectivity}
  假设 $\pairr{A,a_0}$ 和 $\pairr{B,b_0}$ 分别是 $n$-连通和 $m$-连通的带点类型，其中 $n,m\geq0$，并令
%
\narrowequation{P:A\to B\to \ntype{(n+m)}.}
%
那么对于任意 ${f:\prd{a:A} P(a,b_0)}$ 与 ${g:\prd{b:B} P(a_0,b)}$，其中 $p:f(a_0) = g(b_0)$，存在 $h:\prd{a:A}{b:B} P(a,b)$ 以及同伦
%
\begin{equation*}
  q:\prd{a:A} h(a,b_0)=f(a)
  \qquad\text{and}\qquad
  r:\prd{b:B} h(a_0,b)=g(b)
 \end{equation*}
%
使得 $p = \opp{q(a_0)} \ct r(b_0)$。
\end{lem}

\begin{proof}
  定义 $Q:A\to\type$ 为
  \[ Q(a) \defeq \sm{k:\prd{b:B} P(a,b)} (f(a) = k(b_0)). \]
  则我们有 $(g,p):Q(a_0)$。
  由于 $a_0:\unit\to A$ 是 $(n-1)$-连通的，若 $Q$ 是 $(n-1)$-类型族，则我们将得到 $\ell:\prd{a:A} Q(a)$ 使得 $\ell(a_0) = (g,p)$；此时可以定义 $h(a,b) \defeq \proj1(\ell(a))(b)$。
  然而，对于固定的 $a$，类型 $Q(a)$ 是映射
  \[ \Parens{\prd{b:B} P(a,b) } \to P(a,b_0) \]
  在 $f(a)$ 处的纤维，该映射通过用 $b_0:\unit\to B$ 作前复合得到。
  由于 $b_0:\unit\to B$ 是 $(m-1)$-连通的，根据 \cref{thm:conn-trunc-variable-ind}，为使这个纤维是 $(n-1)$-截断的，只需每个类型 $P(a,b)$ 都是 $(n+m)$-类型即可，而这正是我们的假设。
\end{proof}

令 $(X,x_0)$ 为一个带点类型，并回顾 \cref{sec:suspension} 中纬悬 $\susp X$ 的定义，其构造子为 $\north,\south:\susp X$ 与 $\merid:X \to (\north=\south)$。
我们将 $\susp X$ 视为以 $\north$ 为基点的带点空间，因此有 $\Omega\susp X \defeq (\id[\susp X]\north\north)$。
于是存在一个规范映射
\begin{align*}
  \sigma &: X \to \Omega\susp X\\
  \sigma(x) &\defeq \merid(x) \ct \opp{\merid(x_0)}.
\end{align*}

\begin{rmk}
  在经典代数拓扑中，人们考虑的是\emph{约化纬悬}，其中道路 $\merid(x_0)$ 被坍缩为一个点，从而将 $\north$ 和 $\south$ 等同起来。
  约化与非约化纬悬是同伦等价的，因此在我们纯粹的同伦论视角下，这一区分是不可见的——况且高阶归纳类型本来也只允许我们通过高阶道路来``等同''点，所以在同伦类型论中并无必要考虑约化纬悬。
  然而，我们所用纬悬的``非约化性''正是 $\opp{\merid(x_0)}$ 这一（可能令人意外的）项在 $\sigma$ 的定义中出现的原因。
\end{rmk}

我们现在的目标是证明以下定理。

\begin{thm}[Freudenthal 纬悬定理]\label{thm:freudenthal}
  假设 $X$ 带点且 $n$-连通，其中 $n\geq 0$。
  则映射 $\sigma:X\to \Omega\susp(X)$ 是 $2n$-连通的。
\end{thm}

\index{encode-decode method|(}%

我们将使用编码-解码方法，但以稍有不同的方式应用它。
在之前的大多数情形中，我们通过构造一个满足 $\code(a_0)\defeq B$ 的族 $\code:A\to \type$，以及一族等价 $\decode:\prd{x:A}\code(x) \eqvsym (a_0=x)$，将某个类型的环路空间 $\Omega (A,a_0)$ 刻画为与另一个类型 $B$ 等价。
% We have also generalized it to characterize truncations of loop spaces by way of a family of equivalences $\prd{x:A}\code(x) \eqvsym \trunc n{a_0=x}$.

然而在本例中，我们希望证明 $\sigma:X\to \Omega \susp X$ 是 $2n$-连通的。
我们可以使用先前方法的截断版本（例如我们将在 \cref{sec:van-kampen} 中看到的那样）来证明 $\trunc{2n}X\to \trunc{2n}{\Omega \susp X}$ 是一个等价——但这比映射为 $2n$-连通这一断言稍弱（参见 \cref{thm:conn-pik,thm:pik-conn}）。
不过请注意，在一般情形中，为了证明 $\decode(x)$ 是一个等价，我们也可以等价地证明其纤维是可缩的，并且我们仍然可以对底类型进行归纳。
我们可以将这一点推广，用以证明到环路空间的映射的连通性，即证明其纤维的\emph{截断}是可缩的。
此外，与其分别构造 $\code$ 和 $\decode$，我们可以直接构造一族\emph{纤维之截断的代码（codes for the truncations of the fibers）}。

\begin{defn}\label{thm:freudcode}
  若 $X$ 带点且 $n$-连通（$n\geq 0$），则存在一族
  \begin{equation}
    \code:\prd{y:\susp X} (\north=y) \to \type\label{eq:freudcode}
  \end{equation}
  使得
  \begin{align}
    \code(\north,p) &\defeq \trunc{2n}{\hfib{\sigma}{p}}
    \jdeq \trunc{2n}{\tsm{x:X} (\merid(x) \ct \opp{\merid(x_0)} = p)}\label{eq:freudcodeN}\\
    \code(\south,q) &\defeq \trunc{2n}{\hfib{\merid}{q}}
    \jdeq \trunc{2n}{\tsm{x:X} (\merid(x) = q)}.\label{eq:freudcodeS}
  \end{align}
\end{defn}

我们的最终目标是证明对所有 $y:\susp X$ 和 $p:\north=y$，$\code(y,p)$ 都是可缩的。
取 $y\defeq \north$，这将表明 $\sigma$ 的所有纤维都是 $2n$-连通的，从而 $\sigma$ 是 $2n$-连通的。

\begin{proof}[\cref{thm:freudcode} 的证明]
  我们通过对 $y:\susp X$ 进行归纳来定义 $\code(y,p)$，其中前两种情形是~\eqref{eq:freudcodeN} 与~\eqref{eq:freudcodeS}。
  接下来需要对每个 $x_1:X$ 构造一个依值道路
  \[ \dpath{\lam{y}(\north=y)\to\type}{\merid(x_1)}{\code(\north)}{\code(\south)}. \]
  根据 \cref{thm:dpath-arrow}，这等价于给出一族道路
  \[ \prd{q:\north=\south} \code(\north)(\transfib{\lam{y}(\north=y)}{\opp{\merid(x_1)}}{q}) = \code(\south)(q). \]
  而由泛等性（univalence）及道路类型中的传输，这又等价于给出一族等价
  \[ \prd{q:\north=\south} \code(\north,q \ct \opp{\merid(x_1)}) \eqvsym \code(\south,q). \]
  我们将定义一族映射
  \begin{equation}\label{eq:freudmap}
    \prd{q:\north=\south} \code(\north,q \ct \opp{\merid(x_1)}) \to \code(\south,q).
  \end{equation}
  然后证明它们都是等价的。
  因此，令 $q:\north=\south$；根据截断的泛性质以及 $\code(\north,\blank)$ 和 $\code(\south,\blank)$ 的定义，只需对每个 $x_2:X$ 定义一个映射
  \begin{equation*}
    \big(\merid(x_2)\ct \opp{\merid(x_0)} = q \ct \opp{\merid(x_1)}\big)
    \to \trunc{2n}{\tsm{x:X} (\merid(x) = q)}.
  \end{equation*}
  现在，对于每个 $x_1,x_2:X$，该类型是 $2n$-截断的，而 $X$ 是 $n$-连通的。
  于是，由 \cref{thm:wedge-connectivity}，只需分别在 $x_1$ 为 $x_0$ 和 $x_2$ 为 $x_0$ 的情形下定义此映射，并验证当两者同时为 $x_0$ 时它们一致即可。

  当 $x_1$ 为 $x_0$ 时，假设条件是 $r:\merid(x_2)\ct \opp{\merid(x_0)} = q \ct \opp{\merid(x_0)}$。
  由此，从 $r$ 中消去 $\opp{\merid(x_0)}$ 得到 $r':\merid(x_2)=q$，于是我们可以将其像定义为 $\tproj{2n}{(x_2,r')}$。

  当 $x_2$ 为 $x_0$ 时，假设条件是 $r:\merid(x_0)\ct \opp{\merid(x_0)} = q \ct \opp{\merid(x_1)}$。
  将其重新排列，我们得到 $r'':\merid(x_1)=q$，从而可以将其像定义为 $\tproj{2n}{(x_1,r'')}$。

  最后，当 $x_1$ 和 $x_2$ 均为 $x_0$ 时，只需证明得到的 $r'$ 和 $r''$ 是一致的；这是一个关于道路复合的简单引理。
  至此，我们完成了~\eqref{eq:freudmap} 的定义。
  为了证明它是一族等价，由于作为等价这一性质是一个命题（mere proposition）且 $x_0:\unit\to X$ 是（至少）$(-1)$-连通的，只需假设 $x_1$ 是 $x_0$。
  在这种情况下，检查上述构造可知，它本质上是消去 $\opp{\merid(x_0)}$ 这一函数的 $2n$-截断，而该函数本身是一个等价。
\end{proof}

除了~\eqref{eq:freudcodeN} 与~\eqref{eq:freudcodeS} 之外，我们还需要从 $\code$ 的构造中提取一些关于它如何作用于道路的信息。
为此，我们使用以下引理。

\begin{lem}\label{thm:freudlemma}
  设 $A:\UU$，$B:A\to \UU$，以及 $C:\prd{a:A} B(a)\to\UU$，同时给定 $a_1,a_2:A$ 以及 $m:a_1=a_2$ 和 $b:B(a_2)$。
  那么，函数
  \[\transfib{\widehat{C}}{\pairpath(m,t)}{\blank} : C(a_1,\transfib{B}{\opp m}{b}) \to C(a_2,b),\]
  （其中 $t:\transfib{B}{m}{\transfib{B}{\opp m}{b}} = b$ 是明显的相干道路，而 $\widehat{C}:(\sm{a:A} B(a)) \to\type$ 是 $C$ 的非柯里化形式）等于通过泛等性从以下复合得到的等价：
  \begin{align}
    C(a_1,\transfib{B}{\opp m}{b})
    &= \transfib{\lam{a} B(a)\to \UU}{m}{C(a_1)}(b)
    \tag{by~\eqref{eq:transport-arrow}}\\
    &= C(a_2,b). \tag{by $\happly(\apd{C}{m},b)$}
  \end{align}
\end{lem}

\begin{proof}
  根据道路归纳，我们可以假设 $a_2$ 就是 $a_1$ 且 $m$ 就是 $\refl{a_1}$，此时两个函数都是恒等函数。
\end{proof}

我们将此引理应用于以下情形：令 $A\defeq\susp X$，$B\defeq \lam{y}(\north=y)$，$C\defeq\code$；同时令 $a_1\defeq\north$，$a_2\defeq\south$，并对某个 $x_1:X$ 令 $m\defeq \merid(x_1)$；最后令 $b\defeq q$，即某条道路 $\north=\south$。
对 $\susp X$ 进行归纳的计算规则将 $\apd{C}{m}$ 等同于一条由泛等性和函数外延性以特定方式构造出的道路。
\cref{thm:freudlemma} 中描述的第二个函数，实质上就是撤销这些对泛等性和函数外延性的应用，从而归约到我们利用 \cref{thm:wedge-connectivity} 所定义的特定函数~\eqref{eq:freudmap}。
因此，\cref{thm:freudlemma} 表明，沿着 $\pairpath(q,t)$ 进行传输，实质上就恢复了这些函数。

最后，根据构造，当 $x_1$ 或 $x_2$ 与 $x_0$ 重合，并且输入在 $\tproj{2n}{\blank}$ 的像中时，我们能更明确地知道这些函数是什么。
于是，对任意 $x_2:X$，我们有
\begin{equation}
  \transfib{\hat{\code}}{\pairpath(\merid(x_0),t)}{\tproj{2n}{(x_2,r)}}
  =\tproj{2n}{(x_1,r')}\label{eq:freudcompute1}
\end{equation}
其中 $r:\merid(x_2) \ct \opp{\merid(x_0)} = \transfib{B}{\opp{\merid(x_0)}}{q}$ 如前所述是任意的，而 $r':\merid(x_2)=q$ 是由 $r$ 将其终点与 $q \ct \opp{\merid(x_0)}$ 等同并消去 $\opp{\merid(x_0)}$ 而得到的。
类似地，对任意 $x_1:X$，我们有
\begin{equation}
  \transfib{\hat{\code}}{\pairpath(\merid(x_1),t)}{\tproj{2n}{(x_0,r)}}
  = \tproj{2n}{(x_1,r'')}\label{eq:freudcompute2}
\end{equation}
其中 $r:\merid(x_0) \ct \opp{\merid(x_0)} = \transfib{B}{\opp{\merid(x_1)}}{q}$，而 $r'':\merid(x_1)=q$ 是通过将其终点等同起来并重新排列道路而得到的。

\begin{proof}[\cref{thm:freudenthal} 的证明]
  接下来要证明，对于每个 $y:\susp X$ 和 $p:\north=y$， $\code(y,p)$ 都是可缩的。
  首先，我们必须选定一个收缩中心，记作 $c(y,p):\code(y,p)$。
  这对应于我们先前证明中 $\encode$ 函数的定义，因此我们用传输来定义它。
  注意，在 $y$ 取 $\north$ 且 $p$ 取 $\refl{\north}$ 的特殊情况下，我们有
  \[\code(\north,\refl{\north}) \jdeq \trunc{2n}{\tsm{x:X} (\merid(x) \ct \opp{\merid(x_0)} = \refl{\north})}.\]
  因此，我们可以定义 $c(\north,\refl{\north})\defeq \tproj{2n}{(x_0,\mathsf{rinv}_{\merid(x_0)})}$，其中对于任意道路 $q$，$\mathrm{rinv}_q$ 是明显的道路 $q\ct\opp q = \refl{}$。
  接下来，我们可以通过对 $p$ 做道路归纳来得到 $c:\prd{y:\susp X}{p:\north=y} \code(y,p)$，但下面我们会看到，用传输给出一个具体的定义也很重要：
  \[ c(y,p) \defeq \transfib{\hat{\code}}{\pairpath(p,\mathsf{tid}_p)}{c(\north,\refl{\north})}
  \]
  这里 $\hat{\code}: \big(\sm{y:\susp X} (\north=y)\big) \to \type$ 是 \code 的非柯里化形式，而 $\mathsf{tid}_p:\trans{p}{\refl{}} = p$ 是一个标准引理。

  接着，我们必须证明 $\code(y,p)$ 中的每个元素都等于 $c(y,p)$。
  再次利用道路归纳，我们只需考虑 $y$ 为 $\north$ 且 $p$ 为 $\refl{\north}$ 的情形。
  事实上，我们将在更一般的情形下证明它，即仅假设 $y$ 是 $\north$，而 $p$ 是任意的。
  换言之，我们将证明，对于任意 $p:\north=\north$ 和 $d:\code(\north,p)$，都有 $d = c(\north,p)$。
  由于该等式是一个 $(2n-1)$-类型，我们可以假设 $d$ 具有 $\tproj{2n}{(x_1,r)}$ 的形式，其中 $x_1:X$，且 $r:\merid(x_1) \ct \opp{\merid(x_0)} = p$。

  现在，通过进一步的道路归纳，我们可以假定 $r$ 是自反性，并且 $p$ 是 $\merid(x_1) \ct \opp{\merid(x_0)}$。
  （这就是为什么我们在前面要推广到任意的 $p$。）
  于是，我们需要证明
  \begin{equation}
    \tproj{2n}{(x_1, \refl{\merid(x_1) \ct \opp{\merid(x_0)}})}
    \;=\;
    c\left(\north,\refl{\merid(x_1) \ct \opp{\merid(x_0)}}\right).\label{eq:freudgoal}
  \end{equation}
  根据定义，该等式右边是
  \begin{multline*}
    \Transfib{\hat{\code}}{\pairpath(\merid(x_1) \ct \opp{\merid(x_0)}, \nameless)}{\tproj{2n}{(x_0,\nameless)}} \\
    = \transfibf{\hat{\code}}
    \begin{aligned}[t]
      \Big(
      &{\pairpath(\opp{\merid(x_0)}, \nameless)},\\
      &{\Transfib{\hat{\code}}{\pairpath(\merid(x_1), \nameless)}{\tproj{2n}{(x_0,\nameless)}}}
      \Big)
    \end{aligned}
    \\
    = \Transfib{\hat{\code}}{\pairpath(\opp{\merid(x_0)}, \nameless)}{\tproj{2n}{(x_1,\nameless)}}
    = \tproj{2n}{(x_1,\nameless)}
  \end{multline*}
  其中下划线 $\nameless$ 处应填入合适的相干道路。
  这里，第一步使用了传输的函子性，第二步调用了~\eqref{eq:freudcompute2}，第三步则调用了~\eqref{eq:freudcompute1}（同时将传输移到等式的另一边）。
  这样一来，等式右边的第一分量就与~\eqref{eq:freudgoal} 左边的第一分量相同了。
  我们留给读者验证：所有相干道路都会相消，从而在第二分量中给出自反性。
\end{proof}

% As a corollary, we have the following equivalence.

\begin{cor}[Freudenthal 等价性] \label{cor:freudenthal-equiv}
  设 $X$ 是 $n$-连通且带点的类型，其中 $n\geq 0$。
  那么 $\eqv{\trunc{2n}{X}}{\trunc{2n}{\Omega\susp(X)}}$。
\end{cor}

\begin{proof}
  根据 \cref{thm:freudenthal}，$\sigma$ 是 $2n$-连通的。再由
  \cref{lem:connected-map-equiv-truncation} 可知，它因此在 $2n$-截断上诱导一个等价。
\end{proof}

\index{encode-decode method|)}%

\index{Freudenthal suspension theorem|)}%
\index{theorem!Freudenthal suspension|)}%

\index{homotopy!group!of sphere}%
\index{stability!of homotopy groups of spheres}%
\index{type!n-sphere@$n$-sphere}%
Freudenthal 纬悬定理的一个重要推论是，球面的同伦群在某个范围内是稳定的（这些范围对应于 \cref{tab:homotopy-groups-of-spheres} 中从东北到西南方向的对角线）：

\begin{cor}[球面的稳定性] \label{cor:stability-spheres}
  若 $k \le 2n-2$，则 $\pi_{k+1}(S^{n+1}) = \pi_{k}(S^{n})$。
\end{cor}

\begin{proof}
  假设 $k \le 2n-2$。
  %
  根据 \cref{cor:sn-connected}，$\Sn ^{n}$ 是 $\nminusone$-连通的。因此，
  由 \cref{cor:freudenthal-equiv} 可得，
  \[
\trunc{2(n-1)}{\Omega(\susp(\Sn^{n}))} = \trunc{2(n-1)}{\Sn^{n}}.
\]
  再根据 \cref{lem:truncation-le}，由于 $k \le 2(n-1)$，在等式两边同时应用 $\trunc{k}{\blank}$ 可知此式对 $k$ 也成立：
  \begin{equation}\label{eq:freudenthal-for-spheres}
\trunc{k}{\Omega(\susp(\Sn^{n}))} = \trunc{k}{\Sn^{n}}.
\end{equation}
  %
  于是，证明的主要想法如下；我们略去以下验证：这些等价在相应空间的基点上是否作用恰当，以及当 $k > 0$ 时这些等价是否保持乘法结构：
  %
  \begin{align*}
\pi_{k+1}(\Sn^{n+1}) &\jdeq \trunc{0}{\Omega^{k+1}(\Sn^{n+1})} \\
                     &\jdeq \trunc{0}{\Omega^k(\Omega(\Sn^{n+1}))} \\
                     &\jdeq \trunc{0}{\Omega^k(\Omega(\susp(\Sn^{n})))} \\
                     &= \Omega^k(\trunc{k}{(\Omega(\susp(\Sn^{n})))})
                     \tag{by \cref{thm:path-truncation}}\\
                     &= \Omega^k(\trunc{k}{\Sn^{n}})
                     \tag{by \eqref{eq:freudenthal-for-spheres}}\\
                     &= \trunc{0}{\Omega^k(\Sn^{n})}
                     \tag{by \cref{thm:path-truncation}}\\
                     &\jdeq \pi_k(\Sn^{n}). \qedhere
\end{align*}
  %
\end{proof}

这意味着，一旦我们计算出了某条稳定对角线上的一个条目，我们就知道了该对角线上所有的条目。例如：

\begin{thm}\label{thm:pinsn}
对每个 $n\geq 1$，有 $\pi_n(\Sn^n)=\Z$。
\end{thm}

\begin{proof}
证明对 $n$ 施行归纳。我们已知 $\pi_1(\Sn ^1) = \Z$（\cref{cor:pi1s1}）与 $\pi_2(\Sn ^2) = \Z$（\cref{cor:pis2-hopf}）。
当 $n \ge 2$ 时，$n \le (2n - 2)$。于是，由
\cref{cor:stability-spheres}、$\pi_{n+1}(S^{n+1}) = \pi_{n}(S^{n})$ 以及这一等价，再结合归纳假设，即得结论。
\end{proof}

\begin{cor}
对任意 $n\ge -1$，$\Sn^{n+1}$ 不是 $n$-类型。
\end{cor}

\begin{cor}\label{thm:pi3s2}
$\pi_3(\Sn^2)=\Z$。
\end{cor}

\begin{proof}
由 \cref{cor:pis2-hopf}，$\pi_3(\Sn^2) = \pi_3(\Sn^3)$。
而根据 \cref{thm:pinsn}，有 $\pi_3(\Sn^3) = \Z$。
\end{proof}

\section{van Kampen 定理}
\label{sec:van-kampen}

\index{van Kampen theorem|(}%
\index{theorem!van Kampen|(}%

\index{fundamental!group}%
van Kampen 定理用于计算空间的（同伦）推出的基本群 $\pi_1$。
传统上，这一定理针对一个拓扑空间 $X$，它是两个开子空间 $U$ 和 $V$ 的并集；
但在同伦论的意义下，这只是确保 $X$ 是 $U$ 和 $V$ 沿其交集的推出的一种简便方式。
因此，我们将证明一个适用于任意推出的 van Kampen 定理版本。

在本节中，我们将使用与 \cref{sec:pi1-s1-intro} 中计算 $\pi_1(\Sn^1)$ 时相同的编码-解码方法来证明 van Kampen 定理。
此外，也存在一种更偏向同伦论的处理方法；参见 \cref{ex:rezk-vankampen}。

我们需要编码-解码方法的一个更精细的版本。
在 \cref{sec:pi1-s1-intro}（以及 \cref{sec:compute-coprod,sec:compute-nat}）中，我们用它来刻画一个（高阶）归纳类型 $W$ 的道路空间——从而得到环路空间 $\Omega(W)$ 的刻画，并进而得到其 0-截断 $\pi_1(W)$ 的刻画。
在 van Kampen 定理中，我们的目标只是刻画基本群 $\pi_1(W)$，而我们没有任何环路空间或道路空间的显式描述可供使用。

事实证明，我们可以将同样的技术直接用于道路纤维化的截断版本，从而不仅刻画基本\emph{群} $\pi_1(W)$，还刻画整个基本\emph{群胚}。
\index{fundamental!pregroupoid}%
具体来说，对于一个类型 $X$，记 $\Pi_1 X: X\to X\to \type$ 为其相等类型的 $0$-截断，即$\Pi_1 X(x,y) \defeq \trunc0{x=y}$。
注意我们有诱导的群胚运算
\begin{align*}
  (\blank\ct\blank) &\;:\; \Pi_1X(x,y) \to \Pi_1X(y,z) \to \Pi_1X(x,z)\\
  \opp{(\blank)} &\;:\; \Pi_1X(x,y) \to \Pi_1X(y,x)\\
  \refl{x} &\;:\; \Pi_1X(x,x)\\
  \apfunc{f} &\;:\; \Pi_1X(x,y) \to \Pi_1Y(fx,fy)
\end{align*}
我们对其使用与道路上的相应运算相同的记号。

\subsection{朴素 van Kampen}
\label{sec:naive-vankampen}

我们从 van Kampen 定理的一个``朴素''版本开始，它虽有用，但不及经典版本。
在 \cref{sec:better-vankampen} 中，我们将把它改进为一个更有用的版本。

\index{encode-decode method|(}%

给定类型 $A,B,C$ 和函数 $f:A\to B$ 与 $g:A\to C$，令 $P$ 为它们的推出 $B\sqcup^A C$。
如我们在 \cref{sec:colimits} 中所见，$P$ 是由以下生成元生成的高阶归纳类型：


\begin{itemize}
\item $i:B\to P$，
\item $j:C\to P$，以及
\item 对所有 $x:A$，一条道路 $k x:ifx = jgx$。
\end{itemize}


通过对 $P$ 的双重归纳，定义 $\code:P\to P\to \type$ 如下。


\begin{itemize}
\item $\code(ib,ib')$ 是以下序列类型的一个集合商（见 \cref{sec:set-quotients}）：
  %pairs $(\vec a, \vec a', \vec p, \vec q)$
  \[ (b, p_0, x_1, q_1, y_1, p_1, x_2, q_2, y_2, p_2, \dots, y_n, p_n, b') \]
  其中
  \begin{itemize}
  \item $n:\mathbb{N}$
  \item 对 $0<k \le n$，有 $x_k:A$ 和 $y_k:A$
  \item $p_0:\Pi_1B(b,f x_1)$ 与 $p_n:\Pi_1B(f y_n, b')$（当 $n>0$），以及 $p_0:\Pi_1B(b,b')$（当 $n=0$）
  \item $p_k:\Pi_1B(f y_k, fx_{k+1})$（当 $1\le k < n$）
  \item $q_k:\Pi_1C(gx_k, gy_k)$（当 $1\le k\le n$）
  \end{itemize}
  该商由以下等式生成：
  \begin{align*}
    (\dots, q_k, y_k, \refl{fy_k}, y_k, q_{k+1},\dots)
    &= (\dots, q_k\ct q_{k+1},\dots)\\
    (\dots, p_k, x_k, \refl{gx_k}, x_k, p_{k+1},\dots)
    &= (\dots, p_k\ct p_{k+1},\dots)
  \end{align*}
  （见下文 \cref{rmk:naive}）。
  我们留给读者将这种序列类型精确定义为一个归纳类型。
\item $\code(jc,jc')$ 的定义类似，但交换 $B$ 和 $C$ 的角色。
  我们同样在记号上交换 $x$ 和 $y$ 以及 $p$ 和 $q$ 的角色。
\item $\code(ib,jc)$ 和 $\code(jc,ib)$ 的定义也类似，但改变奇偶性，使其起始于一个类型而终止于另一个类型。
\item 对于 $a:A$ 和 $b:B$，我们需要一个等价
  \begin{equation}
    \code(ib, ifa) \eqvsym \code(ib,jga).\label{eq:bfa-bga}
  \end{equation}
  我们将其定义为由如下定义在序列上的两个函数给出：
  \begin{align*}
    (\dots, y_n, p_n,fa) &\mapsto (\dots,y_n,p_n,a,\refl{ga},ga),\\
    (\dots, x_n, p_n, a, \refl{fa}, fa) &\mapsfrom (\dots, x_n, p_n, ga).
  \end{align*}
  这两个函数都容易看出保持等价关系，从而在代码类型上定义了函数。
  从左到右再回到左的复合为
  \[ (\dots, y_n, p_n,fa) \mapsto
  (\dots,y_n,p_n,a,\refl{ga},a,\refl{fa},fa)
  \]
  根据商的一个生成等式，它等于恒等映射。
  另一个方向的复合类似。
  这样我们就定义了一个等价~\eqref{eq:bfa-bga}。
\item 类似地，我们需要等价
  \begin{align*}
    \code(jc,ifa) &\eqvsym \code(jc,jga)\\
    \code(ifa,ib)&\eqvsym (jga,ib)\\
    \code(ifa,jc)&\eqvsym (jga,jc)
  \end{align*}
  它们都以完全相同的方式定义（后两个是通过在序列开头而非结尾添加自反性项来定义的）。
\item 最后，我们需要知道，对于 $a,a':A$，下图交换：
  \begin{equation}\label{eq:bfa-bga-comm}
  \vcenter{\xymatrix{
      \code(ifa,ifa') \ar[r]\ar[d] &
      \code(ifa,jga')\ar[d]\\
      \code(jga,ifa')\ar[r] &
      \code(jga,jga')
      }}
  \end{equation}
  这相当于说，如果我们先在一个序列的开头添加一些东西，然后在结尾再添加一些东西，那么换一种顺序来做也是一样的。
\end{itemize}

\begin{rmk}\label{rmk:naive}
  人们可能期望在 \code 的定义中看到集合商的一些额外的生成等式，例如
  \begin{align*}
    (\dots, p_{k-1} \ct fw, x_{k}', q_{k}, \dots) &=
    (\dots, p_{k-1}, x_{k}, gw \ct q_{k}, \dots)
    \tag{for $w:\Pi_1A(x_{k},x_{k}')$}\\
    (\dots, q_k \ct gw, y_k', p_k, \dots) &=
    (\dots, q_k, y_k, fw \ct p_k, \dots).
    \tag{for $w:\Pi_1A(y_k, y_k')$}
  \end{align*}
  然而，这些并非必需！
  事实上，它们可以通过对 $w$ 进行道路归纳自动得出。
  这正是``朴素''van Kampen 定理与下一小节将要考虑的更精细版本之间的主要区别。
\end{rmk}

继续下去，我们可以刻画在纤维化 $\code$ 中的 transport 操作：


\begin{itemize}
\item 对于 $p:b=_B b'$ 和 $u:P$，我们有
  \[ \mathsf{transport}^{b\mapsto \code(u,ib)}(p, (\dots, y_n,p_n,b))
  = (\dots,y_n,p_n\ct p,b').
  \]
\item 对于 $q:c=_C c'$ 和 $u:P$，我们有
  \[ \mathsf{transport}^{c\mapsto \code(u,jc)}(q, (\dots, x_n,q_n,c))
  = (\dots,x_n,q_n\ct q,c').
  \]
\end{itemize}


这里我们在记号上有所滥用：对 $X$ 中的一条道路及其在 $\Pi_1X$ 中的像使用了相同的名称。
注意，$\Pi_1X$ 中的 transport 也是通过与一条道路（的像）进行拼接给出的。
由此，我们可以通过对 $u$ 进行归纳来证明上述命题。
我们还有：


\begin{itemize}
\item 对于 $a:A$ 和 $u:P$，
  \[ \mathsf{transport}^{v\mapsto \code(u,v)}(ha, (\dots, y_n,p_n,fa))
  = (\dots,y_n,p_n,a,\refl{ga},ga).
  \]
\end{itemize}


这本质上由 $\code$ 的定义得出。

我们还要通过对 $u$ 进行归纳来构造一个函数
\[ r : \prd{u:P} \code(u,u) \]
具体如下：
\begin{align*}
  rib &\defeq (b,\refl{b},b)\\
  rjc &\defeq (c,\refl{c},c)
\end{align*}
而对于 $rka$，我们取复合相等
\begin{align*}
  (ka,ka)_* (fa,\refl{fa},fa)
  &= (ga,\refl{ga},a,\refl{fa},a,\refl{ga},ga) \\
  &= (ga,\refl{ga},ga)
\end{align*}
其中第一个相等来自上述关于在 $\code$ 中 transport 的观察，第二个相等则是定义 $\code$ 时所使用的集合商关系的一个实例。

接下来我们将证明：


\begin{thm}[朴素 van Kampen 定理]\label{thm:naive-van-kampen}
  对于所有 $u,v:P$，存在一个等价
  \[ \Pi_1P(u,v) \eqvsym \code(u,v). \]
\end{thm}

\begin{proof}

为定义一个函数
\[ \encode : \Pi_1P(u,v) \to \code(u,v) \]
只需定义一个函数 $(u=_P v) \to \code(u,v)$ 即可，因为 $\code(u,v)$ 是一个集合。
我们通过 transport 来定义它：
\[\encode(p) \defeq \mathsf{transport}^{v\mapsto \code(u,v)}(p,r(u)).\]
现在来定义
\[ \decode: \code(u,v) \to \Pi_1P(u,v) \]
我们照常对 $u,v:P$ 进行归纳。
在 $u$ 和 $v$ 的每一种情形中，我们将 $i$ 或 $j$ 适当地作用于所有等式 $p_k$ 和 $q_k$，并在 $P$ 中拼接结果，同时利用 $h$ 来等同端点。
例如，当 $u\jdeq ib$ 且 $v\jdeq ib'$ 时，我们定义
\begin{narrowmultline}\label{eq:decode}
 \decode(b, p_0, x_1, q_1, y_1, p_1, \dots, y_n, p_n, b') \defeq\narrowbreak
 (p_0)\ct h(x_1) \ct j(q_1) \ct \opp{h(y_1)} \ct i(p_1) \ct \cdots \ct \opp{h(y_n)}\ct i(p_n).
\end{narrowmultline}
这个定义保持集合商的等价关系以及诸如 \eqref{eq:bfa-bga} 的等价，因为 $h: fi \htpy gj$ 是自然的，且 $f$ 和 $g$ 具有函子性。

照例，为了证明复合
\[ \Pi_1P(u,v) \xrightarrow{\encode} \code(u,v) \xrightarrow{\decode} \Pi_1P(u,v) \]
是恒等映射，我们首先``剥掉''0-截断（因为值域是集合），然后应用道路归纳。
输入 $\refl{u}$ 映到 $ru$，而 $ru$ 又映回 $\refl u$（通过对 $u$ 的进一步归纳来分解 $\decode(ru)$）。

最后，考虑复合
\[  \code(u,v) \xrightarrow{\decode} \Pi_1P(u,v) \xrightarrow{\encode} \code(u,v). \]
我们对 $u,v:P$ 进行归纳。
当 $u\jdeq ib$ 且 $v\jdeq ib'$ 时，此复合是
%
\begin{narrowmultline*}
(b, p_0, x_1, q_1, y_1, p_1, \dots, y_n, p_n, b')
\narrowbreak
\begin{aligned}[t]
  &\mapsto \Big(ip_0\ct hx_1 \ct jq_1 \ct \opp{hy_1} \ct ip_1 \ct \cdots \ct \opp{hy_n}\ct ip_n\Big)_*(rib)\\
  &= (ip_n)_* \cdots(jq_1)_* (hx_1)_*(ip_0)_*(b,\refl{b},b)\\
  &= (ip_n)_* \cdots(jq_1)_* (hx_1)_*(b,p_0,ifx_1)\\
  &= (ip_n)_* \cdots(jq_1)_* (b,p_0,x_1,\refl{gx_1},jgx_1)\\
  &= (ip_n)_* \cdots (b,p_0,x_1,q_1,jgy_1)\\
  &= \quad\vdots\\
  &= (b, p_0, x_1, q_1, y_1, p_1, \dots, y_n, p_n, b').
\end{aligned}
\end{narrowmultline*}
%
即恒等函数。
（准确地说，这里需要一个隐式的归纳论证。）
另外三种点的情形是类似的，而道路的情形则是平凡的，因为所涉及的类型都是集合。
\end{proof}

\index{encode-decode method|)}%

\cref{thm:naive-van-kampen} 使我们能够计算许多类型的基本群，前提是 $A$ 是一个集合；因为在这种情况下，根据定义，每个 $\code(u,v)$ 都是一个\emph{集合}模掉一个关系所得的集合商。

\begin{eg}\label{eg:circle}
  令 $A\defeq \bool$，$B\defeq\unit$，且 $C\defeq \unit$。
  那么 $P \eqvsym S^1$。
  考察比如 $\code(i(\ttt),i(\ttt))$ 的定义，我们发现所有的道路实际上都是平凡的，因此唯一的信息就是元素的序列 $x_1,y_1,\dots,x_n,y_n: \bool$。
  此外，若对某个 $k$ 我们有 $x_k=y_k$ 或 $y_k=x_{k+1}$，那么集合商关系允许我们同时``切除''这两个元素。
  因此，每个这样的序列都等于一个典范的\emph{约化}序列，其中没有两个相邻元素是相等的。
  显然，这样的约化序列由其长度（一个自然数 $n$）以及（如果 $n>1$）$x_1$ 是 $\bfalse$ 还是 $\btrue$ 这一信息所唯一决定，因为这能唯一确定序列的其余部分。
  当然，这些数据可以等同于一个整数：$n$ 对应绝对值，$x_1$ 编码正负号。
  于是我们重新得到了 $\pi_1(S^1)\cong \Z$。
\end{eg}

由于 \cref{thm:naive-van-kampen} 断言的只是集合族之间的双射，这个同构 $\pi_1(S^1)\cong \Z$ 同样也只是集合间的双射。
然而，我们可以在 $\code$ 上定义一种拼接运算（通过拼接序列），并证明 $\encode$ 和 $\decode$ 构成了一个保持该结构的同构。
（用 \cref{cha:category-theory} 的语言来说，这些将是``预群胚''。）
细节留给读者。

\index{fundamental!group|(}%

\begin{eg}\label{eg:suspension}
  更一般地，令 $B\defeq\unit$ 且 $C\defeq \unit$，但 $A$ 是任意的，于是 $P$ 是 $A$ 的纬悬。
  此时道路 $p_k$ 和 $q_k$ 也是平凡的，因此一个道路编码中唯一的信息是元素的序列 $x_1,y_1,\dots,x_n,y_n: A$。
  前两个生成等式表明相邻的相等元素可以消去，因此将这个序列视为如下形式的群中的字是合理的：
  \[ x_1 y_1^{-1} x_2 y_2^{-1} \cdots x_n y_n^{-1} \]
  实际上，它看起来类似于 $A$（或者说 $\trunc0A$；参见 \cref{thm:freegroup-nonset}）上的自由群，但我们只考虑那些以非逆元素开头、交替出现逆与非逆元素、并以逆元素结尾的字。
  这实际上将生成元集合的规模减少了一个。
  例如，如果 $A$ 有一个点 $a:A$，那么我们可以将 $\pi_1(\susp A)$ 等同于由 $\trunc0A$ 作为生成元且满足关系 $\tproj0a = e$ 所给出的群；详见 \cref{ex:vksusppt,ex:vksuspnopt}。
\end{eg}

\begin{eg}\label{eg:wedge}
  令 $A\defeq\unit$，且 $B$ 和 $C$ 是任意的，于是 $f$ 和 $g$ 只是给 $B$ 和 $C$ 配备了基点，比如 $b$ 和 $c$。
  那么 $P$ 就是 $B$ 和 $C$ 的\emph{楔} $B\vee C$（带基点空间范畴中的余积）。
  在这种情况下，元素 $x_k$ 和 $y_k$ 是平凡的，因此唯一的信息是一系列环路 $(p_0,q_1,p_1,\dots,p_n)$，其中 $p_k:\pi_1(B,b)$ 且 $q_k:\pi_1(C,c)$。
  这样的序列模去我们所施加的等价关系后，很容易与群 $\pi_1(B,b)$ 和 $\pi_1(C,c)$ 的\emph{自由积}的显式描述等同，如 \cref{sec:free-algebras} 中所构造的那样。
  因此，我们有 $\pi_1(B\vee C) \cong \pi_1(B) * \pi_1(C)$。
\end{eg}

\index{fundamental!group|)}%

然而，\cref{thm:naive-van-kampen} 离完整的经典 van Kampen 定理还差一步。经典版本能处理 $A$ 不一定是集合的情形，并且断言 $\pi_1(B\sqcup^A C) \cong \pi_1(B) *_{\pi_1(A)} \pi_1(C)$（其中的基点取自 $A$）。
事实上，\cref{thm:naive-van-kampen} 的结论完全没有涉及 $\pi_1(A)$；$A$ 中的道路以一种类型论的方式被``嵌入''到商化构造当中，这使得提取显式信息变得困难，因为 $\code(u,v)$ 是对一个非集合按某个关系做集合商得到的结果。
因此，在下一小节中，我们将考虑一个更好的 van Kampen 定理版本。

\subsection{带基点集的 van Kampen 定理}
\label{sec:better-vankampen}

\index{basepoint!set of}%
我们将要介绍的 van Kampen 定理的改进，与经典代数拓扑中的类似改进密切相关，即在经典理论中为 $A$ 配备一个\emph{基点集 $S$}。
事实上，我们的证明并不需要假设这个``基点集''是一个\emph{集合}——它完全可以是一个任意类型；假设 $S$ 是集合的用处要到之后应用定理进行计算时才体现出来。
重要的是，$S$ 包含 $A$ 的每个连通分支中至少一个点。
我们在类型论中将其表述为：存在一个类型 $S$ 和一个函数 $k:S \to A$，它是满的，亦即 $(-1)$-连通的。
如果 $S\jdeq A$ 且 $k$ 是恒等函数，那么我们将恢复朴素的 van Kampen 定理。
另一个值得记住的例子是，当 $A$ 是带点且 (0-)连通的，并令 $k:\unit\to A$ 为取基点的映射时：由 \cref{thm:minusoneconn-surjective,thm:connected-pointed}，此映射是满的恰好当 $A$ 是 0-连通的。

令 $A,B,C,f,g,P,i,j,h$ 如前一小节所述。
现在，给定满射 $k:S\to A$，我们定义一个辅助类型，用以改善 $k$ 的连通性。
令 $T$ 为由以下生成元构造的高阶归纳类型：

\begin{itemize}
\item 一个函数 $\ell:S\to T$，以及
\item 对每个 $s,s':S$，一个函数 $m:(\id[A]{ks}{ks'}) \to (\id[T]{\ell s}{\ell s'})$。
\end{itemize}

存在一个明显的诱导函数 $\kbar:T\to A$ 使得 $\kbar \ell = k$，并且任何 $p:ks=ks'$ 都等于复合 $ks = \kbar \ell s \overset{\kbar m p}{=} \kbar \ell s' = k s'$。

\begin{lem}\label{thm:kbar}
  $\kbar$ 是 0-连通的。
\end{lem}

\begin{proof}
  我们需要证明，对任意 $a:A$，类型 $\sm{t:T}(\kbar t = a)$ 的 0-截断是可缩的。
  由于可缩性是一个命题，而 $k$ 是 $(-1)$-连通的，我们不妨设 $a=ks$ 对某个 $s:S$ 成立。
  于是我们可以将收缩中心取为 $\tproj0{(\ell s,q)}$，其中 $q$ 是相等 $\kbar\ell s = k s$。

  剩下还需要证明，对于任意 $\phi:\trunc0{\sm{t:T} (\kbar t = ks)}$，我们有 $\phi = \tproj0{(\ell s,q)}$。
  因为后者是一个命题，特别地也是集合，我们可以进一步假设 $\phi=\tproj0{(t,p)}$ 对某个 $t:T$ 和 $p:\kbar t = ks$ 成立。

  现在我们可以对 $t:T$ 进行归纳。
  如果 $t\jdeq\ell s'$，那么由 $ks' = \kbar \ell s' \overset{p}{=} ks$ 可以通过 $m$ 导出一个相等 $\ell s = \ell s'$。
  于是，根据 $\kbar$ 的定义以及同伦纤维中相等的定义，我们得到相等 $(ks',p) = (ks,q)$，从而得到 $\tproj0{(ks',p)} = \tproj0{(ks,q)}$。
  接下来我们必须证明，当 $t$ 沿着 $m$ 变化时，这些相等彼此相容。
  但这些相等是在一个集合（即 $\trunc0{\sm{t:T} (\kbar t = ks)}$）中的，因此这是自动的。
\end{proof}

\begin{rmk}
  \index{kernel!pair}%
  $T$ 可以被视为 $k$ 的``核偶对''的（同伦）余等化子。
  如果 $S$ 和 $A$ 是集合，那么 $k$ 的 $(-1)$-连通性将蕴含 $A$ 就是这个余等化子的 $0$-截断（参见 \cref{cha:set-math}）。
  对于一般的类型，高阶意象理论则提示，$k$ 的 $(-1)$-连通性将转而蕴含 $A$ 是 $k$ 的``单纯核''的余极限（亦称``几何实现''）。
  \index{.infinity1-topos@$(\infty,1)$-topos}%
  \index{geometric realization}%
  \index{simplicial!kernel}%
  \index{kernel!simplicial}%
  类型 $T$ 是这个单纯核的``1-骨架''的余极限，因此，它能将 $k$ 的连通性提高 $1$ 是合理的。
  更一般地，我们可以期望 $n$-骨架\index{skeleton!of a CW-complex}的余极限将连通性提高 $n$。
\end{rmk}

\index{encode-decode method|(}%

现在，我们通过双重归纳如下定义 $\code:P\to P\to \type$：

\begin{itemize}
\item $\code(ib,ib')$ 现在是如下类型序列的集合商：
  \[ (b, p_0, x_1, q_1, y_1, p_1, x_2, q_2, y_2, p_2, \dots, y_n, p_n, b') \]
  其中
  \begin{itemize}
  \item $n:\mathbb{N}$，
  \item 对 $0<k \le n$，有 $x_k:S$ 和 $y_k:S$，
  \item 当 $n>0$ 时，$p_0:\Pi_1B(b,f k x_1)$ 且 $p_n:\Pi_1B(f k y_n, b')$；当 $n=0$ 时，$p_0:\Pi_1B(b,b')$，
  \item 对 $1\le k < n$，有 $p_k:\Pi_1B(fk y_k, fkx_{k+1})$，
  \item 对 $1\le k\le n$，有 $q_k:\Pi_1C(gkx_k, gky_k)$。
  \end{itemize}
  该商由以下相等生成（参见\cref{rmk:naive}）：
  \begin{align*}
    (\dots, q_k, y_k, \refl{fy_k}, y_k, q_{k+1},\dots)
    &= (\dots, q_k\ct q_{k+1},\dots)\\
    (\dots, p_k, x_k, \refl{gx_k}, x_k, p_{k+1},\dots)
    &= (\dots, p_k\ct p_{k+1},\dots)\\
    (\dots, p_{k-1} \ct fw, x_{k}', q_{k}, \dots) &=
    (\dots, p_{k-1}, x_{k}, gw \ct q_{k}, \dots)
    \tag{for $w:\Pi_1A(kx_{k},kx_{k}')$}\\
    (\dots, q_k \ct gw, y_k', p_k, \dots) &=
    (\dots, q_k, y_k, fw \ct p_k, \dots).
    \tag{for $w:\Pi_1A(ky_k, ky_k')$}
  \end{align*}
  我们后面会用到在此类序列上 $\decode$ 的定义，它和之前一样，将所有路径 $p_k$ 和 $q_k$ 与 $h$ 的实例拼接起来，以得到 $\Pi_1P(ifb,ifb')$ 中的一个元素，参见~\eqref{eq:decode}。
  和之前一样，其他三种点的情况几乎相同。
\item 对于 $a:A$ 和 $b:B$，我们需要一个等价：
  \begin{equation}
    \code(ib, ifa) \eqvsym \code(ib,jga).\label{eq:bfa-bga2}
  \end{equation}
  由于 $\code$ 是集合值的，根据\cref{thm:kbar}，我们可以假设 $a=\kbar t$，其中 $t:T$。
  接下来，我们可以对 $t$ 进行归纳。
  若 $t\jdeq \ell s$，其中 $s:S$，那么我们如\cref{sec:naive-vankampen}那样定义~\eqref{eq:bfa-bga2}：
  \begin{align*}
    (\dots, y_n, p_n,fks) &\mapsto (\dots,y_n,p_n,s,\refl{gks},gks),\\
    (\dots, x_n, p_n, s, \refl{fks}, fks) &\mapsfrom (\dots, x_n, p_n, gks).
  \end{align*}
  这些定义尊重等价关系，并且如前所述定义了拟逆。
  现在假设 $t$ 沿着某个 $w:ks=ks'$ 对应的 $m_{s,s'}(w)$ 变化；我们必须证明~\eqref{eq:bfa-bga2} 尊重沿着 $\kbar mw$ 的传输。
  根据 $\kbar$ 的定义，这本质上可归结为沿 $w$ 本身的传输。
  由道路类型中传输的特征描述，我们需要证明的是：
  \[ w_*(\dots, y_n, p_n,fks) = (\dots,y_n, p_n \ct fw, fks') \]
  被~\eqref{eq:bfa-bga2} 映射到
  \[ w_*(\dots,y_n,p_n,s,\refl{gks},gks) = (\dots, y_n, p_n, s, \refl{gks} \ct gw, gks') \]
  但这直接由定义 \code 的集合商时所施加的新生成元得出。
\item 其他三种必需的等价以类似方式定义。
\item 最后，由于交换性~\eqref{eq:bfa-bga-comm} 是一个命题，根据 $k$ 的 $(-1)$-连通性，我们可以假设 $a=ks$ 且 $a'=ks'$，在这种情况下，其证明与之前完全相同。
\end{itemize}

\begin{thm}[带基点集的 van Kampen 定理]\label{thm:van-Kampen}
  对于所有 $u,v:P$，存在一个等价：
  \[ \Pi_1P(u,v) \eqvsym \code(u,v). \]
  其中 \code 的定义如本节所述。
\end{thm}

\begin{proof}
  基本上和之前一样。
  为了证明 $\decode$ 尊重商关系的新生成元，我们利用了 $h$ 的自然性。
  为了证明 $\decode$ 尊重诸如~\eqref{eq:bfa-bga2}之类的等价，我们需要对 $\kbar$ 和 $T$ 进行归纳，以便将这些等价分解为它们的定义式，但随后这又简单地变为 $f$ 和 $g$ 的函子性。
  其余部分是容易的。
  特别地，对于 $\encode\circ\decode$ 不需要额外的论证，因为目标是在集合中证明一个相等，所以关于 $h$ 的情况是平凡的。
\end{proof}

\index{encode-decode method|)}%

\index{fundamental!group|(}%
\cref{thm:van-Kampen} 允许我们计算空间~$A$ 的基本群，
即使 $A$ 不是集合，只要 $S$ 是集合；因为在此情况下，
根据定义，每个 $\code(u,v)$ 都是对一个\emph{集合}按某个关系做集合商的结果。
在这一点上，它是对
\cref{thm:naive-van-kampen}
的一个改进。

\begin{eg}\label{eg:clvk}
  假设 $S\defeq \unit$，那么 $A$ 有一个基点 $a \defeq k(\ttt)$ 并且是连通的。
  此时，推出的环路的编码可以等同于 $\pi_1(B,f(a))$ 和 $\pi_1(C,g(a))$ 中环路的交错序列，模去一个等价关系；该关系允许我们在它们之间滑动 $\pi_1(A,a)$ 的元素（分别在应用 $f$ 和 $g$ 之后）。
  因此，$\pi_1(P)$ 可以等同于 \emph{融合自由积}
  \index{amalgamated free product}%
  \index{free!product!amalgamated}%
  $\pi_1(B) *_{\pi_1(A)} \pi_1(C)$（群范畴中的推出），如\cref{sec:free-algebras}所构造。
  这（当 $B$ 和 $C$ 是 $P$ 的开子空间且 $A$ 是它们的交集时）大概是 van Kampen 定理最经典的版本。
\end{eg}

\begin{eg}\label{eg:cofiber}
  \index{cofiber}
  作为\cref{eg:clvk}的一个特例，我们再额外假设 $C\defeq\unit$，于是 $P$ 是余纤维 $B/A$。
  那么 $C$ 中的每个环路都等于自反性，所以道路编码上的关系允许我们将所有序列塌缩为 $B$ 中的单个环路。
  附加的关系要求在左边、右边或中间乘以 $\pi_1(A)$ 的像中的一个元素后等于单位元。
  因此，我们可以将 $\pi_1(B/A)$ 等同于群 $\pi_1(B)$ 模去由 $\pi_1(A)$ 的像生成的正规子群所得的商群。
\end{eg}

\begin{eg}\label{eg:torus}
  \index{torus}
  作为\cref{eg:cofiber}的进一步特例，令 $B\defeq S^1 \vee S^1$，令 $A\defeq S^1$，并令 $f:A\to B$ 选取复合环路 $p \ct q \ct \opp p \ct \opp q$，其中 $p$ 和 $q$ 是组成 $B$ 的两个 $S^1$ 拷贝中的生成环路。
  此时 $P$ 就是环面 $T^2$ 的一个表示。
  实际上，不难将 $P$ 与\cref{sec:hubs-spokes}中使用特定环路的锥所描述的 $T^2$ 表示等同起来。
  因此，$\pi_1(T^2)$ 是两个生成元上的自由群\index{generator!of a group}（即 $\pi_1(B)$）模去关系$p \ct q \ct \opp p \ct \opp q = 1$所得的商群。
  这显然给出两个生成元上的自由\emph{交换}群\index{group!abelian}，即 $\Z\times\Z$。
\end{eg}

\begin{eg}
  \index{CW complex}
  \index{hub and spoke}
  更一般地，任何 CW 复形都可以通过反复对球面“取锥”而得到，如\cref{sec:hubs-spokes}所述。
  具体来说，我们从一组点（“0-胞腔”）构成的集合 $X_0$ 出发，它称为 CW 复形的“0-骨架”。
  对某个 1-胞腔集合 $S_1$ 和某个“附着映射”族 $f_1$ 取推出
  \begin{equation*}
    \vcenter{\xymatrix{
        S_1 \times \Sn^0\ar[r]^-{f_1}\ar[d] &
        X_0\ar[d]\\
        \unit \ar[r] &
        X_1
      }}
  \end{equation*}
  得到“1-骨架”\index{skeleton!of a CW-complex} $X_1$。
  \index{attaching map}%
  然后再对某个 2-胞腔集合 $S_2$ 和某个附着映射族 $f_2$ 取推出
  \begin{equation*}
    \vcenter{\xymatrix{
        S_2 \times \Sn^1\ar[r]^{f_2}\ar[d] &
        X_1\ar[d]\\
        \unit \ar[r] &
        X_2
      }}
  \end{equation*}
  得到 2-骨架 $X_2$，依此类推。
  每次取推出所得的基本群都可以用 van Kampen 定理计算：得到一个群表示，其生成元由 1-骨架导出，关系由 $S_2$ 和 $f_2$ 导出。
  此后各步的推出不再改变基本群，因为当 $n>1$ 时 $\pi_1(\Sn^n)$ 是平凡的（参见\cref{sec:pik-le-n}）。
\end{eg}

\begin{eg}\label{eg:kg1}
  特别地，假设给定任意一个（集合意义上的）群 $G = \langle X \mid R \rangle$ 的表示\index{presentation!of a group}，其中 $X$ 是生成元集合，$R$ 是这些生成元\index{generator!of a group}上的字构成的集合。
  令 $B\defeq \bigvee_X S^1$ 且 $A\defeq \bigvee_R S^1$，并将 $f:A\to B$ 定义为将每个 $S^1$ 拷贝映到 $B$ 的生成环路中对应的字。
  由此可得 $\pi_1(P) \cong G$；这样我们就构造了一个基本群为 $G$ 的连通类型。
  由于任何群都有表示，所以任何群都可以成为某个类型的基本群。
  如果我们对这样一个类型取 1-截断，就得到一个类型，其唯一的非平凡同伦群为 $G$；这称为\define{艾伦伯格–麦克兰恩空间} $K(G,1)$。%
  \indexdef{Eilenberg--Mac Lane space}%
\end{eg}

\index{fundamental!group|)}%

\index{van Kampen theorem|)}%
\index{theorem!van Kampen|)}%

\section{Whitehead 定理与 Whitehead 原理}
\label{sec:whitehead}

在经典同伦论中，若映射 $f:A\to B$ 对 $A$ 中的所有点 $a$ 都诱导同构$\pi_n(A,a) \cong \pi_n(B,f(a))$（同时也诱导同构$\pi_0(A)\cong\pi_0(B)$），那么只要空间 $A$ 和 $B$ 性质良好（例如具有 CW 复形的同伦类型），$f$ 就必然是一个同伦等价。
\index{theorem!Whitehead's}%
\index{Whitehead's!theorem}%
这被称为\emph{Whitehead 定理}。
实际上，那些使 Whitehead 定理失效的“性质不良”空间在类型论中是不可见的。
粗略来说，性质良好的拓扑空间已经足以呈现 $\infty$-群胚，%
\index{.infinity-groupoid@$\infty$-groupoid}
而同伦类型论直接处理 $\infty$-群胚，而非具体的拓扑空间。
因此，人们或许会预期 Whitehead 定理在泛等基础中成立。

然而，情况并非如此：Whitehead 定理是不可证明的。
事实上，在已知的某些类型论模型中，该定理不成立，尽管其原因与在性质不良的拓扑空间上失效的原因完全不同。
这些模型是“非超完备 $\infty$-意象”
\index{.infinity1-topos@$(\infty,1)$-topos}%
\index{.infinity1-topos@$(\infty,1)$-topos!non-hypercomplete}%
（参见~\cite{lurie:higher-topoi}）；粗略地说，它们由 $\infty$ 维底空间上的 $\infty$-群胚层构成。

\index{axiom!Whitehead's principle|(}%
\index{Whitehead's!principle|(}%

因此，从基础的观点来看，我们可以将\emph{Whitehead 原理}视为一条“经典性公理”，类似于 \LEM{} 和 \choice{}。
它可以被一致地假设，但并不属于以计算为动机的类型论，也并非在所有自然模型中成立。
但在基于集合论的基础下，这条原理是隐而不显的：在一个 $\infty$-群胚由集合（使用拓扑空间、单纯集合等模型）构建起来的世界里，它不可能不成立。

这听起来或许有些奇怪，但实际上并不意外。
同伦类型论是同伦类型的\emph{抽象}理论，而在集合论框架下，拓扑空间或单纯集合的同伦论只是该理论的一个\emph{具体}模型；
这就像整数是抽象的环论的一个具体模型一样。
可以想见，任何具体模型都会具有一些并非相应抽象理论所固有的特殊性质，但有时我们希望将它们作为额外的公理来假定（例如，整数环是主理想整环，但并非所有环都是如此）。

描述类型论的模型已超出本书范围，因此我们不会解释 Whitehead 原理在某些模型中为何可能不成立。
然而，我们可以证明，只要所涉及的类型对某个有限 $n$ 是 $n$-截断的，该原理便成立，证明方法是对 $n$ 作“向下”归纳。
这一结果本身有其独立的意义（例如，它蕴含了艾伦伯格–麦克兰恩空间的本质唯一性）；
同时，其证明过程也有望提供一些直观解释，帮助我们理解为什么在缺少截断条件的情况下，我们无法指望证明类似的定理。

我们首先给出 \cref{thm:mono-surj-equiv} 的如下变形，它最终将为截断版 Whitehead 原理的归纳步骤提供支持。
它可以看作范畴论中经典结论“全忠实且本质满的函子是范畴等价”在类型论与 $\infty$-群胚语境下的版本。

\begin{thm}\label{thm:whitehead0}
  设 $f:A\to B$ 是一个函数，满足：
  \begin{enumerate}
  \item $\trunc0 f : \trunc0 A \to \trunc0 B$ 是满的，并且\label{item:whitehead01}
  \item 对任意 $x,y:A$，函数 $\apfunc f : (\id[A]xy) \to (\id[B]{f(x)}{f(y)})$ 是一个等价。\label{item:whitehead02}
  \end{enumerate}
  那么 $f$ 是一个等价。
\end{thm}

\begin{proof}
  注意条件~\ref{item:whitehead02} 恰好表达了 $f$ 是嵌入，参见~\cref{sec:mono-surj}。
  因此，由 \cref{thm:mono-surj-equiv}，只需证明 $f$ 是满的，即对任意 $b:B$ 证明 $\trunc{-1}{\hfib f b}$。
  给定 $b$。由于 $\trunc0 f$ 是满的，仅知存在 $a:A$ 使得 $\trunc 0 f(\tproj0a) = \tproj0b$。
  而我们的目标是命题，故可以假设给定了这样一个 $a$。
  于是有 $\tproj0{f(a)} = \trunc 0 f(\tproj0a) =\tproj0b$，从而 $\trunc{-1}{f(a)=b}$。
  同样，因为目标仍是命题，我们可以进一步假设 $f(a)=b$。
  由此 $\hfib f b$ 有实例，因此仅知有实例。
\end{proof}

由于同伦群是环路空间\index{loop space}的截断，而非道路空间的截断，我们需要将上述定理修改为关于环路空间的表述。
回忆从 \cref{def:loopfunctor} 引入的映射 $\Omega f$。

\begin{cor}\label{thm:whitehead1}
  设 $f:A\to B$ 是一个函数，满足：
  \begin{enumerate}
  \item $\trunc0 f : \trunc0 A \to \trunc0 B$ 是一个双射，并且
  \item 对任意 $x:A$，函数 $\Omega f : \Omega(A,x) \to \Omega(B,f(x))$ 是一个等价。
  \end{enumerate}
  那么 $f$ 是一个等价。
\end{cor}

\begin{proof}
  由 \cref{thm:whitehead0}，只需证明对任意 $x,y:A$，$\apfunc f : (\id[A]xy) \to (\id[B]{f(x)}{f(y)})$ 是一个等价。
  而根据 \cref{thm:equiv-inhabcod}，我们可以假设 $\id[B]{f(x)}{f(y)}$。
  特别地，有 $\tproj0{f(x)} = \tproj0{f(y)}$，因此由于 $\trunc0 f$ 是等价，可知 $\tproj0 x = \tproj0y$，故 $\tproj{-1}{x=y}$。
  由于我们所要证明的是命题（即 $\apfunc f$ 为等价），我们可以假设给定 $p:x=y$。
  但此时以下方块在同伦意义下交换：
  \begin{equation*}
  \vcenter{\xymatrix@C=3pc{
      \Omega(A,x)\ar[r]^-{\blank\ct p}\ar[d]_{\Omega f} &
      (\id[A]xy) \ar[d]^{\apfunc f}\\
      \Omega(B,f(x))\ar[r]_-{\blank\ct f(p)} &
      (\id[B]{f(x)}{f(y)}).
      }}
  \end{equation*}
  顶边和底边的映射都是等价，而左边的映射由假设也是等价。
  因此，根据“三中取二”性质，右边的映射也是等价。
\end{proof}

现在我们可以证明截断版 Whitehead 原理。

\begin{thm}\label{thm:whiteheadn}
  设 $A$ 与 $B$ 是 $n$-类型，且 $f:A\to B$ 满足：
  \begin{enumerate}
  \item $\trunc0f:\trunc0A \to \trunc0B$ 是一个双射，并且\label{item:wh0}
  \item $\pi_k(f):\pi_k(A,x) \to \pi_k(B,f(x))$ 对所有 $k\ge 1$ 和所有 $x:A$ 都是双射。\label{item:whk}
  \end{enumerate}
  那么 $f$ 是一个等价。
\end{thm}

\noindent
条件~\ref{item:wh0} 几乎就是~\ref{item:whk} 中 $k=0$ 的情形，只是它没有提及任何基点 $x:A$。

\begin{proof}
  我们通过对 $n$ 的归纳来证明。
  当 $n=-2$ 时，结论是平凡的。
  因此，假设结论对所有 $n$-类型之间的函数都成立，并令 $A$ 和 $B$ 为 $(n+1)$-类型，且 $f:A\to B$ 如上所述。
  \cref{thm:whitehead1} 中的第一个条件由假设成立，因此只需证明对于任意 $x:A$，函数 $\Omega f: \Omega(A,x) \to \Omega(B,f(x))$ 是一个等价。

  由于 $\Omega(A,x)$ 和 $\Omega(B,f(x))$ 是 $n$-类型，我们可以应用归纳假设。
  我们需要验证 $\trunc 0{\Omega f}$ 是双射，并且对于所有 $k\geq1$ 和 $p : x = x$，映射 $\pi_k(\Omega f):\pi_k(x = x,p) \to \pi_k(f(x) = f(x),\Omega f(p))$ 是双射。
  第一个论断由假设成立，因为 $\trunc 0{\Omega f} \jdeq \pi_1(f)$。
  为证明第二个论断，我们首先将其推广：我们证明对于所有 $y : A$ 和 $q : x = y$，有 $\pi_k(\apfunc f):\pi_k(x = y,q) \to \pi_k(f(x) = f(y),\apfunc f(q))$。
  这蕴含了待证结论，因为当 $y\defeq x$ 时，在等同其基点 $\Omega f(p)=\apfunc f(p)$ 的意义下，有 $\pi_k(\Omega f)=\pi_k(\apfunc f)$。
  为证明这一推广，由道路归纳，只需证 $q$ 为 $\refl{a}$ 的情形。
  此时有 $\pi_k(\apfunc f) = \pi_k(\Omega f) = \pi_{k+1}(f)$，而 $\pi_{k+1}(f)$ 由原始假设为双射。
\end{proof}

注意，如果 $A$ 和 $B$ 对任何有限的 $n$ 都不是 $n$-类型，那么归纳法就无从开始。

\begin{cor}\label{thm:whitehead-contr}
  如果 $A$ 是一个 $0$-连通的 $n$-类型，并且对于所有 $k$ 和 $a:A$ 有 $\pi_k(A,a)=0$，那么 $A$ 是可缩的。
\end{cor}

\begin{proof}
  将 \cref{thm:whiteheadn} 应用到映射 $A\to\unit$ 上。
\end{proof}

作为一个应用，我们可以推断出 \cref{thm:conn-pik} 的逆命题。

\begin{cor}\label{thm:pik-conn}
  对于 $n\ge 0$，一个映射 $f:A\to B$ 是 $n$-连通的，当且仅当下述所有条件成立：
  \begin{enumerate}
  \item $\trunc0f:\trunc0A \to \trunc0B$ 是一个同构。
  \item 对于任意 $a:A$ 和 $k\le n$，映射 $\pi_k(f):\pi_k(A,a) \to \pi_k(B,f(a))$ 是一个同构。
  \item 对于任意 $a:A$，映射 $\pi_{n+1}(f):\pi_{n+1}(A,a) \to \pi_{n+1}(B,f(a))$ 是满的。
  \end{enumerate}
\end{cor}

\begin{proof}
  “仅当”方向就是 \cref{thm:conn-pik}。
  反过来，由纤维化的长正合列（\cref{thm:les}），
  \index{sequence!exact}%
  这些前提蕴含了对于所有 $k\le n$ 和 $a:A$ 有 $\pi_k(\hfib f {f(a)})=0$，并且 $\trunc 0{\hfib f{f(a)}}$ 是可缩的。
  由于对于 $k\le n$ 有 $\pi_k(\hfib f {f(a)}) = \pi_k(\trunc n{\hfib f{f(a)}})$，并且 $\trunc n{\hfib f{f(a)}}$ 是 $n$-连通的，根据 \cref{thm:whitehead-contr}，对任意 $a$，它是可缩的。

  余下只需证明对于 $b:B$（不一定具有 $f(a)$ 的形式），$\trunc n{\hfib f{b}}$ 是可缩的。
  然而，由假设，存在 $x:\trunc0A$ 使得 $\tproj 0b = \trunc0f(x)$。
  因为可缩性是命题，我们可以假定 $x$ 具有 $\tproj0a$ 的形式（其中 $a:A$），在这种情况下 $\tproj 0 b = \trunc0f(\tproj0a) = \tproj0{f(a)}$，因此 $\trunc{-1}{b=f(a)}$。
  同样由于可缩性是命题，我们可以假定 $b=f(a)$，结论随之得出。
\end{proof}

一个使得 $\trunc0f$ 是双射且对所有 $k$ 有 $\pi_k(f)$ 是双射的映射 $f$，被称为 \define{$\infty$-连通的（$\infty$-connected）}%
\indexdef{function!.infinity-connected@$\infty$-connected}%
\indexdef{.infinity-connected function@$\infty$-connected function}
或一个\define{弱等价}。%
\indexdef{equivalence!weak}%
\indexdef{weak equivalence!of types}
这等价于要求 $f$ 对所有 $n$ 都是 $n$-连通的。
一个类型 $Z$ 被称为 \define{$\infty$-截断的（$\infty$-truncated）}%
\indexdef{type!.infinity-truncated@$\infty$-truncated}%
\indexdef{.infinity-truncated type@$\infty$-truncated type}
或\define{超完备的}%
\indexdef{type!hypercomplete}%
\indexdef{hypercomplete type}
如果只要 $f$ 是 $\infty$-连通的，诱导映射
\[(\blank\circ f):(B\to Z) \to (A\to Z)\]
就是一个等价——也就是说，如果 $Z$ 认为每个 $\infty$-连通映射都是等价。
\indexdef{axiom!Whitehead's principle}%
那么，如果我们想把 Whitehead 原理当作一条公理来假设，可以使用以下任意一种等价形式。

\begin{itemize}
\item 每个 $\infty$-连通函数都是一个等价。
\item 每个类型都是 $\infty$-截断的。
\end{itemize}

在高阶意象模型中，
\index{.infinity1-topos@$(\infty,1)$-topos}%
$\infty$-截断的类型按照 \cref{sec:modalities} 的意义构成一个反射子宇宙（即一个 $(\infty,1)$-意象的“超完备化”），但我们不知道这一性质在一般情况下是否成立。

\index{axiom!Whitehead's principle|)}%
\index{Whitehead's!principle|)}%

对任何 $n$ 都不是 $n$-类型的类型是否\emph{存在}，这或许并不显然，但事实上确实存在。
在经典同伦论中，对于任何 $n\ge 2$，$\Sn^n$ 都具有这一性质。
我们在同伦类型论中尚未证明这一事实，但存在其他类型，我们可以证明它们具有“无限截断层次”。

\begin{eg}
  假设我们有 $B:\nat\to\type$，使得对每个 $n$，类型 $B(n)$ 都包含一个与 $n$ 重自反性不相等的 $n$-环路\index{loop!n-@$n$-}，例如 $p_n:\Omega^n(B(n),b_n)$ 满足 $p_n \neq \refl{b_n}^n$。
  （举例来说，我们可以定义 $B(n)\defeq \Sn^n$，并令 $p_n$ 为 $1:\Z$ 在同构 $\pi_n(\Sn^n)\cong \Z$ 下的像。）
  考虑 $C\defeq \prd{n:\nat} B(n)$，其中的点 $c:C$ 由 $c(n)\defeq b_n$ 定义。
  因为环路空间与积交换，所以对任意 $m$，我们有
  \[\eqvspaced{\Omega^m (C,c)}{\prd{n:\nat}\Omega^m(B(n),b_n)}.\]
  在此等价下，$\refl{c}^m$ 对应于函数 $(n\mapsto \refl{b_n}^m)$。
  现在在右边的类型中定义 $q_m$ 为
  \[ q_m(n) \defeq
  \begin{cases}
    p_n &\quad m=n\\
    \refl{b_n}^m &\quad m\neq n.
  \end{cases}
  \]
  若 $q_m = (n\mapsto \refl{b_n}^m)$ 成立，则 $p_n = \refl{b_n}^n$ 也成立，但事实并非如此。
  因此，$q_m \neq (n\mapsto \refl{b_n}^m)$，于是在 $\Omega^m(C,c)$ 中存在一个与 $\refl{c}^m$ 不相等的点。
  故 $C$ 对任何 $m:\nat$ 都不是 $m$-类型。
\end{eg}

我们预期，利用宇宙包含所有高阶归纳类型（如对一切 $n$ 包含 $\Sn^n$）这一事实，也应该能够证明一个宇宙 $\UU$ 本身对任何 $n$ 都不是 $n$-类型。
不过，这一点尚未得到证明。

\section{编码-解码方法的一般陈述}
\label{sec:general-encode-decode}

\indexdef{encode-decode method}

我们已经使用编码-解码方法来刻画各种类型的道路空间，包括余积（\cref{thm:path-coprod}）、自然数（\cref{thm:path-nat}）、截断（\cref{thm:path-truncation}）、圆（\cref{{cor:omega-s1}}）、纬悬（\cref{thm:freudenthal}）和推出（\cref{thm:van-Kampen}）。
本章引言中提及的许多其他定理的证明也用到了这一方法的变体，例如直接证明 $\pi_n(\Sn^n)$、Blakers--Massey 定理以及构造艾伦伯格–麦克兰恩空间。
虽然很想将这一方法抽象为一个引理，但这很困难，因为不同的问题需要稍有不同的变体。
例如，同一方法的不同变体可以用来刻画一个环路空间（如 \cref{thm:path-coprod,cor:omega-s1}）或整个道路空间（如 \cref{thm:path-nat}）；可以给出环路空间的完整刻画（例如 $\Omega^1(\Sn ^1)$），也可以仅刻画其某个截断（例如 van Kampen 定理）；既可以用于计算同伦群，也可以用于证明一个映射是 $n$-连通的（例如 Freudenthal 定理和 Blakers--Massey 定理）。

然而，我们可以针对该方法的特定变体来陈述引理。
这些引理的证明几乎是平凡的；主要目的在于通过一般性地陈述来阐明这一方法。
最简单的情况是用编码-解码方法来刻画一个类型的环路空间，如 \cref{thm:path-coprod} 和 \cref{{cor:omega-s1}}。

\begin{lem}[环路空间的编码-解码]\label{lem:encode-decode-loop}
  \index{loop space}%
给定一个带点类型 $(A,a_0)$ 和一个纤维化 $\code : A \to \type$，如果
\begin{enumerate}
\item $c_0 : \code(a_0)$，\label{item:ed1}
\item $\decode : \prd{x:A} \code(x) \to (\id{a_0}{x})$，\label{item:ed2}
\item 对所有 $c : \code(a_0)$ 有 \id{\transfib{\code}{\decode(c)}{c_0}}{c}，并且\label{item:ed3}
\item $\id{\decode(c_0)}{\refl{}}$，\label{item:ed4}
\end{enumerate}
那么 $(\id{a_0}{a_0})$ 等价于 $\code(a_0)$。
\end{lem}

\begin{proof}
定义
$\encode : \prd{x:A} (\id{a_0}{x}) \to \code(x)$ 为
\[
\encode_x(\alpha) = \transfib{\code}{\alpha}{c_0}.
\]
我们证明 $\encode_{a_0}$ 和 $\decode_{a_0}$ 互为拟逆。
复合 $\encode_{a_0} \circ \decode_{a_0}$ 由假设~\ref{item:ed3} 立即可得。
对于另一个复合，我们证明
\[
\prd{x:A}{p : \id{a_0}{x}} \id{\decode_{x} (\encode_{x} p)}{p}.
\]
由道路归纳，只需证明
$\id{{\decode_{{a_0}} (\encode_{{a_o}} \refl{})}}{\refl{}}$。
在归约 $\mathsf{transport}$ 之后，只需证明
$\id{{\decode_{{a_0}} (c_0)}}{\refl{}}$，而这正是假设~\ref{item:ed4}。
\end{proof}

如果想要得到 $(\id{a_0}{\blank})$ 与 $\code$ 之间的逐纤维等价，只需将条件 (iii) 加强为
\[
\prd{x:A}{c : \code(x)} \id{\encode_{x}(\decode_{x}(c))}{c}.
\]
然而，为了计算一个环路空间（例如 $\Omega(\Sn ^1)$），这一更强的假设并非必要。

另一个常见的变体出现在计算同伦群时，它刻画的是环路空间的截断：

\begin{lem}[环路空间截断的编码-解码]
假设有一个带点类型 $(A,a_0)$ 和一个纤维化 $\code : A \to \type$，其中对每个 $x$，$\code(x)$ 都是 $k$-类型。
定义
\[
\encode : \prd{x:A} \trunc{k}{\id{a_0}{x}} \to \code(x)
\]
（利用 $\code(x)$ 是 $k$-类型这一事实）通过截断递归，将 $\alpha : \id{a_0}{x}$ 映到 $\transfib{\code}{\alpha}{c_0}$。假设：
\begin{enumerate}
\item $c_0 : \code(a_0)$，
\item $\decode : \prd{x:A} \code(x) \to \trunc{k}{\id{a_0}{x}}$，
\item \label{item:decode-encode-loop-iii}
  $\id{\encode_{a_0}(\decode_{a_0}(c))}{c}$ 对所有 $c : \code(a_0)$ 成立，并且
\item \label{item:decode-encode-loop-iv}
  $\id{\decode(c_0)}{\tproj{}{\refl{}}}$。
\end{enumerate}
那么 $\trunc{k}{\id{a_0}{a_0}}$ 等价于 $\code(a_0)$。
\end{lem}

\begin{proof}
$\decode \circ \encode$ 为恒等映射，这由 \ref{item:decode-encode-loop-iii} 立即可得。
%
为证明 $\encode \circ \decode$，我们首先进行截断归纳，由此只需证明
\[
\prd{x:A}{p : \id{a_0}{x}} \id{\decode_{x}(\encode_{x}(\tproj{k}{p}))}{\tproj{k}{p}}.
\]
之所以可以进行截断归纳，是因为 $k$-类型中的道路也是 $k$-类型。
为证明此类型，我们进行道路归纳，并在归约 \encode 之后，使用假设~\ref{item:decode-encode-loop-iv}。
\end{proof}

\section{更多结果}
\label{sec:moreresults}

尽管本章没有给出证明，但下列结果也已在同伦类型论中建立。

\begin{thm}
\index{homotopy!group!of sphere}
存在一个 $k$，使得对所有 $n \ge 3$，$\pi_{n+1}(\Sn ^n) =
\Z_k$。
\end{thm}

\begin{proof}[关于证明的注记.]
证明包括对 $\pi_4(\Sn ^3)$ 的计算，以及诉诸稳定性结论（\cref{cor:stability-spheres}）。
在这一结果的经典陈述中，$k$ 等于 $2$。
虽然我们尚未验证 $k$ 确实为 $2$，但我们对 $\pi_4(\Sn ^3)$ 的计算是构造性的\index{mathematics!constructive}，与本章中的所有其他证明一样。
（更准确地说，它没有使用任何额外公理，例如 \LEM{} 或 \choice{}，从而使其构造性与泛等性和高阶归纳类型相当。）
因此，给定同伦类型论的一种计算性解释，我们就能够在计算机上运行该证明，从而验证 $k$ 确实为 $2$。
这个例子颇耐人寻味，因为这是首个无需事先知道答案即可完成的同伦群计算。
\end{proof}

% Recall from \cref{sec:colimits} that $X \sqcup^C Y$ denotes the
% (homotopy) pushout of $X$ and $Y$ along $C$.
\index{pushout}%

\begin{thm}[Blakers--Massey 定理]\label{Blakers-Massey}
  \indexdef{theorem!Blakers--Massey}%
  \indexsee{Blakers--Massey theorem}{theorem, Blakers--Massey}%
  假设给定映射 $f : C  \rightarrow X$ 和 $g : C \rightarrow Y$。
  先取 $f$ 与 $g$ 的推出 $X \sqcup^C Y $，再取其内含映射的拉回 $\inl : X \rightarrow X \sqcup^C Y \leftarrow Y : \inr$，便得到一个诱导映射 $C \to X \times_{(X \sqcup^C Y)} Y$。

  如果 $f$ 是 $i$-连通的而 $g$ 是 $j$-连通的，那么这个诱导映射就是 $(i+j)$-连通的。
  换句话说，对于任意点 $x:X$ 和 $y:Y$，对应的纤维 $C_{x,y}$（即 $(f,g) : C \to X \times Y $ 的纤维）给出了推出中道路空间 $\id[X \sqcup^C Y]{\inl(x)}{\inr(y)}$ 的一个近似。
\end{thm}

需要注意的是，在经典代数拓扑中，Blakers--Massey 定理常常以稍有不同的形式陈述：
映射 $f$ 和 $g$ 被替换为 CW 复形子复形的包含映射，而同伦推出与同伦拉回则分别被替换为并集与交集。
为了在同伦类型论中表达这一定理，我们必须用同伦不变的概念来替换此类概念。
在 van Kampen 定理（\cref{sec:van-kampen}）中我们已经看到了另一个类似的例子：在那里，我们必须用同伦推出来替换开子集的并集。

\begin{thm}[Eilenberg--Mac Lane 空间]\label{Eilenberg-Mac-Lane-Spaces}
\index{Eilenberg--Mac Lane space}
对于任意交换\index{group!abelian}群 $G$ 和正整数 $n$，都存在一个 $n$-类型
$K(G,n)$，使得 $\pi_n(K(G,n)) = G$，并且当 $k\neq n$ 时，$\pi_k(K(G,n)) = 0$。
\end{thm}

\begin{thm}[覆叠空间]\label{thm:covering-spaces}
  \index{covering space}%
  对于连通空间 $A$，$A$ 上的覆叠空间与带有 $\pi_1(A)$ 作用的集合之间存在等价。
\end{thm}

\sectionNotes

%% For the spheres,  two
%% different definitions of the $n$-sphere $\Sn ^n$ have been used: the first as the
%% suspension of $\Sn ^ {n-1}$ (\cref{sec:suspension}), and the second
%% as a higher inductive type with one base point and one loop in
%% $\Omega^n$ (\cref{sec:circle}); we list the status for both
%% definitions.  The equivalence of the two can be deduced from the remarks
%% at the end of \cref{sec:suspension}.

{
\newcommand{\humancheck}{\ding{52}}
\newcommand{\computercheck}{\ding{52}\kern-0.5em\ding{52}}


\begin{table}[htb]
  \centering
\begin{tabular}{lcc}
\toprule
定理         & 状态 \\
\midrule
$\pi_1(\Sn ^1)$                     & \computercheck \\
$\pi_{k<n}(\Sn ^n)$                  & \computercheck \\
%% $\pi_{k<n}(\Sn ^n)$ --- suspension  & \computercheck \\
%% $\pi_2(\Sn ^2)$ --- suspension      & \computercheck \\
同伦群的长正合列 & \computercheck    \\
Hopf 纤维化的全空间是 $\Sn ^3$ & \humancheck    \\
$\pi_2(\Sn ^2)$                     & \computercheck \\
$\pi_3(\Sn ^2)$                     & \humancheck    \\
$\pi_n(\Sn ^n)$                     & \computercheck \\
%% $\pi_n(\Sn ^n)$ --- suspension      & \computercheck \\
$\pi_4(\Sn ^3)$                     & \humancheck    \\
Freudenthal 纬悬定理      & \computercheck \\
Blakers--Massey 定理              & \computercheck \\
Eilenberg--Mac Lane 空间 $K(G,n)$ & \computercheck \\
van Kampen 定理                & \computercheck \\
覆叠空间                     & \computercheck \\
关于 $n$-类型的 Whitehead 原理 & \computercheck \\
\bottomrule
\end{tabular}
\caption{同伦论中定理的手工证明（\humancheck）与计算机证明（\computercheck）情况。}
  \label{tab:theorems}
\end{table}


}

本章所描述的定理都是经典同伦论中的标准结果；其中许多在 \cite{hatcher02topology} 中有所记述。
在这部分附注里，我们回顾这些定理在同伦类型论中新的综合证明是如何发展起来的。
\cref{tab:theorems} 列出了已在同伦类型论中得到证明的同伦论定理，以及它们是否经过了计算机验证。
这些成果几乎全部诞生于 2013 年高等研究院（IAS）的春季学期，是一次重要协作努力的结晶。
许多人为此作出了贡献，例如担任某个证明的主要作者、提出值得研究的问题、参与大量相关问题的讨论与研讨班，或对已有成果提供反馈。
以下人员是上述定理的首批同伦类型论证明的主要作者。
除非另有说明，对于已经过计算机验证的定理，其主要作者也是率先使用计算机证明助理将证明\index{mathematics!formalized}形式化的人。

\begin{itemize}
\item
Shulman 给出了 $\pi_1(\Sn^1)$ 的同伦论计算。Licata 后来发现了编码-解码证明以及编码-解码方法。

\item
Brunerie 计算了 $\pi_{k<n}(\Sn ^ n)$。
Licata 后来给出了一个编码-解码的版本。

\item Voevodsky（沃沃茨基）构造了同伦群的长正合列。

\item Lumsdaine 构造了 Hopf 纤维化。
  Brunerie 证明了其全空间是 $\Sn ^3$，从而算出了 $\pi_2(\Sn ^2)$ 和 $\pi_3(\Sn ^3)$。

\item Licata 和 Brunerie 直接计算了 $\pi_n(\Sn ^n)$。

\item
  Lumsdaine 证明了 Freudenthal 纬悬定理；Licata 和 Lumsdaine 将这一证明形式化。
\item Lumsdaine, Finster 和 Licata 证明了 Blakers--Massey 定理；
  Lumsdaine, Brunerie, Licata 和 Hou 将其形式化。

\item
Licata 使用编码-解码方法计算了 $\pi_2(\Sn ^2)$，并利用 Freudenthal 纬悬定理计算了 $\pi_n(\Sn ^n)$；运用类似技术，他构造了 $K(G,n)$。

\item
Shulman 证明了 van Kampen 定理；Hou 将此证明形式化。

\item
Licata 证明了关于 $n$-类型的 Whitehead 定理。

\item Brunerie 计算了 $\pi_4(\Sn ^3)$。

\item
Hou 建立并形式化了覆叠空间的理论。
\end{itemize}

同伦论与类型论之间的相互交融对于这些成果的发展至关重要。
例如，$\pi_1(\Sn ^1)=\mathbb{Z}$ 的第一个证明是 \cref{subsec:pi1s1-homotopy-theory} 中给出的那个，它沿用了经典同伦论的思路。
对这一证明的类型论分析催生了编码-解码方法。
对 $\pi_2(\Sn ^2)$ 的首次计算也沿用了经典方法，但这很快便引出了该结果的一个编码-解码证明。
编码-解码计算被推广到 $\pi_n(\Sn ^n)$，这进而引出了 Freudenthal 纬悬定理的证明——其方法是将编码-解码论证与关于连通性的经典同伦论推理相结合。
这一证明反过来又催生了 Blakers--Massey 定理与 Eilenberg--Mac Lane 空间。
这一系列成果的快速进展，彰显了我们对这两门学科之间联系的新理解所蕴含的潜力。

\sectionExercises

\begin{ex}\label{ex:homotopy-groups-resp-prod}
  证明同伦群保持积：$\eqv{\pi_n(A\times B)}{\pi_n(A)\times \pi_n(B)}$。
\end{ex}

\begin{ex}\label{ex:decidable-equality-susp}
  \index{decidable!equality}%
  证明：若 $A$ 是一个具备可判定相等性（参见 \cref{defn:decidable-equality}）的集合，则其纬悬 $\susp A$ 是一个 1-类型。
  （这是否可以不借助可判定相等性的假设来证明，仍是一个\index{open!problem}开放问题。）
\end{ex}

\begin{ex}\label{ex:contr-infinity-sphere-colim}
  定义 $\Sn^\infty$ 为序列 $\Sn^0 \to \Sn^1 \to \Sn^2 \to\cdots$ 的余极限。
  证明 $\Sn^\infty$ 是可缩的。
\end{ex}

\begin{ex}\label{ex:contr-infinity-sphere-susp}
  将 $\Sn^\infty$ 定义为由以下生成元生成的高阶归纳类型：
  \begin{itemize}
  \item 两个点 $\north:\Sn^\infty$ 和 $\south:\Sn^\infty$，及
  \item 对每个 $x:\Sn^\infty$，一条道路 $\merid(x):\north=\south$。
  \end{itemize}
  换言之，$\Sn^\infty$ 就是它自身的纬悬。
  证明 $\Sn^\infty$ 是可缩的。
\end{ex}

\begin{ex}\label{ex:unique-fiber}
  设 $f:X\to Y$ 是一个映射且 $Y$ 是连通的。
  证明对任意 $y_1,y_2:Y$，有 $\brck{\eqv{\hfib f {y_1}}{\hfib f {y_2}}}$。
\end{ex}

\begin{ex}\label{ex:ap-path-inversion}
  对任意带点类型 $A$，令 $i_A : \Omega A \to \Omega A$ 表示环路取逆，即 $i_A \defeq \lam{p} \rev{p}$。
  证明 $i_{\Omega A} : \Omega^2 A \to \Omega^2 A$ 等于 $\Omega(i_A)$。
\end{ex}

\begin{ex}\label{ex:pointed-equivalences}
  将\textbf{带点等价}定义为其底层函数为等价的带点映射。
  \begin{enumerate}
  \item 证明带点类型 $(X,x_0)$ 与 $(Y,y_0)$ 之间的带点等价类型等价于 $\id[\pointed\type]{(X,x_0)}{(Y,y_0)}$。
  \item 用带点拟逆和带点同伦，按照 \cref{cha:equivalences} 中的某种相干风格，重新表述带点等价的概念。
  \end{enumerate}
\end{ex}

\begin{ex}\label{ex:HopfJr}
  仿照 \cref{sec:hopf} 中 Hopf 纤维化的例子，将 \define{低阶 Hopf 纤维化}
  \indexdef{Hopf!fibration!junior}%
定义为 $\Sn ^1$ 上的纤维化（即类型族），其在基点处的纤维为 $\Sn ^0$，全空间为 $\Sn ^1$。这也称为圆 $\Sn ^1$ 的``扭曲二重覆叠''。
\end{ex}

\begin{ex}\label{ex:SuperHopf}
再次仿照 \cref{sec:hopf} 中 Hopf 纤维化的例子，定义 $\Sn ^4$ 上的一个类似纤维化，其在基点处的纤维为 $\Sn ^3$，全空间为 $\Sn ^7$。这是同伦类型论中的一个\index{open!problem}未解决问题（在经典同伦论中已知存在这样的纤维化）。
\end{ex}

\begin{ex}\label{ex:vksusppt}
  继续 \cref{eg:suspension} 的讨论，证明若 $A$ 有一个点 $a:A$，则可将 $\pi_1(\susp A)$ 等同于以 $\trunc0A$ 为生成元、以关系 $\tproj0a = e$ 给出的群。
  进而证明，若假设排中律，该群也是 $\trunc0 A \setminus \{\tproj0 a \}$ 上的自由群。
\end{ex}

\begin{ex}\label{ex:vksuspnopt}
  再次继续 \cref{eg:suspension} 的讨论，但这次不假定 $A$ 带点，证明可将 $\pi_1(\susp A)$ 等同于以生成元 $\trunc0A \times \trunc0A$ 和关系
  \begin{equation*}
    (a,b) = \opp{(b,a)},
    \qquad
    (a,c) = (a,b)\cdot (b,c),
    \qquad\text{and}\qquad
    (a,a) = e.
  \end{equation*}
  给出的群。
\end{ex}

\index{acceptance|)}

%%% Local Variables:
%%% mode: latex
%%% TeX-master: "hott-online"
%%% End:

\chapter{范畴论}
\label{cha:category-theory}

数学的各个分支中，范畴论或许是与集合论基础相处得最不融洽的一支。
其原因之一是，范畴论中的大多数内容在一类比"相等"更弱的"相同性"概念（例如范畴中的同构，或者范畴的等价）下保持不变，而集合论却无法把握这一点。
但泛等公理已经为类型解决了同一类问题，做法是将相等与等价等同起来。
因此，在泛等基础下，考虑一种"范畴"的概念是有意义的，其中对象的相等以类似的方式与同构等同起来。

忽略大小问题，在基于集合的数学里，一个范畴由一个对象\emph{集合} $A_0$ 以及对每一对 $x,y\in A_0$ 的一个态射\emph{集合} $\hom_A(x,y)$ 组成。
在泛等基础下，一个"朴素"的范畴定义只需将上述\emph{集合}替换为\emph{类型}即可。
如果我们允许这些类型包含任意的高阶同伦，那么就应当加上更高阶的相容性条件，从而导向某种 $(\infty,1)$-范畴的概念。
\index{.infinity1-category@$(\infty,1)$-category}%
但在现阶段，我们的目标更为保守。
我们只考虑 1-范畴，因此将类型 $\hom_A(x,y)$ 限制为集合，即 0-类型。
如果不加任何进一步的条件，我们就称这个概念为\emph{预范畴（precategory）}。

如果我们添加一个额外条件，要求对象类型 $A_0$ 本身也是一个集合，那么得到的概念其行为将与传统的集合论定义非常相似。
依照 Toby Bartels 的说法，我们称这个概念为\emph{严格范畴（strict category）}。
\index{strict!category}%
或者，我们可以要求某种泛等公理的推广版本，将 $(x=_{A_0} y)$ 与从 $x$ 到 $y$ 的同构类型 $\mathsf{iso}(x,y)$ 等同起来。
鉴于我们通常认为后一个选择才是"正确的"定义，我们将简单地称其为\emph{范畴（category）}。

一个能很好说明这三种范畴概念之间区别的例子是如下陈述："每一个全忠实且本质满的函子都是范畴的等价"。
这一陈述在基于集合的经典范畴论中等价于选择公理。
\index{mathematics!classical}%
\index{axiom!of choice}%
\index{classical!category theory}%


\begin{enumerate}
\item 对于严格范畴而言，该陈述仍然等价于选择公理。
\item 对于预范畴而言，不存在任何一致的选择公理能使其为真。
\item 对于范畴而言，该陈述\emph{无需}任何选择公理即可证明。
\end{enumerate}


我们将在本章证明最后一个论断，以及范畴的其他一些良好性质，例如，等价的范畴（作为范畴类型的元素）是相等的。
我们还将描述一种将预范畴 $A$ "饱和"为范畴 $\widehat A$ 的通用方法，我们称其为 $A$ 的\emph{Rezk 完备化（Rezk completion）}。
\index{completion!Rezk}%
不过，它也可以合理地被称为\emph{层栈完备化（stack completion）}（参见章末注记）。

Rezk 完备化也能进一步阐明范畴的等价这一概念。
例如，函子 $A \to \widehat{A}$ 总是全忠实且本质满的，因而是"弱等价（weak equivalence）"。
由此可以推出，一个预范畴是范畴，当且仅当它能把所有全忠实且本质满的函子都视为等价；
因此，我们的"范畴"概念其实已经蕴含在"全忠实且本质满函子"这一概念之中了。

我们假设读者已具备经典范畴论的基本知识。\index{classical!category theory}
请记住，每当我们书写 \type，它都表示某个类型宇宙，但在不同时候所指的可能不同；
我们所说的一切对于任意一致的宇宙层级选择都成立\index{universe level}。
我们将使用来自 \cref{cha:typetheory,cha:basics} 的同伦类型论基本概念，以及来自 \cref{cha:logic} 的命题截断，除此之外基本不涉及 \cref{part:foundations} 的其他内容——唯一的例外是构造 Rezk 完备化的第二种方法会用到高阶归纳类型。

\section{Categories and precategories}
\label{sec:cats}

在经典数学中，范畴有许多等价定义。
在我们的情形中，由于我们拥有依值类型，最自然的选择是将箭头（态射）定义为一个以对象为指标的类型族。
这与态射类型在范畴论中的使用方式完全吻合：除非我们知道两个箭头的定义域和陪域一致，否则我们根本不会去比较它们。
此外，对于 1-范畴的理论来说，显然所有的态射类型都应该是集合。
这就将我们引向如下的定义。

\begin{defn}\label{ct:precategory}
  一个\define{预范畴}
  \indexdef{precategory}
  $A$ 由以下数据组成。
  \begin{enumerate}
  \item 一个类型 $A_0$，其元素称为\define{对象}。%
    \indexdef{object!in a (pre)category}
    我们用 $a:A$ 表示 $a:A_0$。
  \item 对于每一对 $a,b:A$，一个集合 $\hom_A(a,b)$，其元素称为\define{箭头}或\define{态射}。%
    \indexsee{arrow}{morphism}%
    \indexdef{morphism!in a (pre)category}%
    \indexdef{hom-set}%
  \item 对于每个 $a:A$，有一个态射 $1_a:\hom_A(a,a)$，称为\define{恒同态射}。%
    \indexdef{identity!morphism in a (pre)category}
  \item 对于每组 $a,b,c:A$，有一个函数%
    \indexdef{composition!of morphisms in a (pre)category}
    \[  \hom_A(b,c) \to \hom_A(a,b) \to \hom_A(a,c) \]
    称为\define{复合}，并以中缀形式记作 $g\mapsto f\mapsto g\circ f$，有时也简写为 $gf$。
  \item 对于每一对 $a,b:A$ 以及 $f:\hom_A(a,b)$，有 $\id f {1_b\circ f}$ 和 $\id f {f\circ 1_a}$。
  \item 对于每组 $a,b,c,d:A$ 以及
    \begin{equation*}
      f:\hom_A(a,b), \qquad
      g:\hom_A(b,c), \qquad
      h:\hom_A(c,d),
    \end{equation*}
    有 $\id {h\circ (g\circ f)}{(h\circ g)\circ f}$。
  \end{enumerate}
\end{defn}

预范畴这个概念存在的问题是，对于对象 $a,b:A$，我们有两种可能不同的"相同性"概念。
一方面，我们有类型 $(\id[A_0]{a}{b})$。
但另一方面，还有范畴论中标准的\emph{同构（isomorphism）}概念。

\begin{defn}\label{ct:isomorphism}
  一个态射 $f:\hom_A(a,b)$ 是一个\define{同构}
  \indexdef{isomorphism!in a (pre)category}%
  如果存在一个态射 $g:\hom_A(b,a)$，使得 $\id{g\circ f}{1_a}$ 且 $\id{f\circ g}{1_b}$。
  我们将此类同构的类型记作 $a\cong b$。
\end{defn}

\begin{lem}\label{ct:isoprop}
  对于任意 $f:\hom_A(a,b)$，类型"$f$ 是一个同构"是一个命题。
  因此，对于任意 $a,b:A$，类型 $a\cong b$ 是一个集合。
\end{lem}

\begin{proof}
  假设给定 $g:\hom_A(b,a)$ 以及 $\eta:(\id{1_a}{g\circ f})$ 和 $\epsilon:(\id{f\circ g}{1_b})$，同样地给定 $g'$、$\eta'$ 和 $\epsilon'$。
  我们需要证明 $\id{(g,\eta,\epsilon)}{(g',\eta',\epsilon')}$。
  但由于所有的 hom-集都是集合，它们的相等类型都是命题，因此只需证明 $\id g {g'}$。
  为此，我们有
  \[g' = 1_a\circ g' = (g\circ f)\circ g' = g\circ (f\circ g') = g\circ 1_b = g\]
  这里用到了 $\eta$ 和 $\epsilon'$。
\end{proof}

\symlabel{ct:inv}
\index{inverse!in a (pre)category}%
如果 $f:a\cong b$，我们记 $\inv f$ 为其逆，根据 \cref{ct:isoprop}，这个逆是由 $f$ 唯一决定的。

在预范畴中，关于这两种相同性概念，我们拥有的唯一关系如下。

\begin{lem}[\textsf{idtoiso}]\label{ct:idtoiso}
  如果 $A$ 是一个预范畴且 $a,b:A$，那么
  \[(\id a b)\to (a \cong b).\]
\end{lem}

\begin{proof}
  通过对相等做归纳，我们可以假定 $a$ 和 $b$ 是相同的。
  但此时我们有 $1_a:\hom_A(a,a)$，这显然是一个同构。
\end{proof}

显然，这种情况类似于促使我们引入泛等公理的问题。
事实上，我们有如下例子：

\begin{eg}\label{ct:precatset}
  \index{set}%
  存在一个预范畴 \uset，其对象类型是 \set，并且满足 $\hom_{\uset}(A,B) \defeq (A\to B)$。
  其中的恒同态射是恒等函数，复合是函数复合。
  对于这个预范畴，\cref{ct:idtoiso} 等于（限制在集合上的）来自 \cref{sec:compute-universe} 的映射 $\idtoeqv$。

  当然，更精确地说，我们应该称这个范畴为 $\uset_\UU$，因为它的对象只是相对于某个宇宙 \UU 的\emph{小集合}。
  \index{small!set}%
\end{eg}

因此，很自然地给出如下定义。

\begin{defn}\label{ct:category}
  一个\define{范畴}
  \indexdef{category}
  是指一个预范畴，它满足以下条件：对所有 $a,b:A$，由 \cref{ct:idtoiso} 给出的函数 $\idtoiso_{a,b}$ 都是一个等价。
\end{defn}

特别地，在一个范畴中，如果 $a\cong b$，那么 $a=b$。

\begin{eg}\label{ct:eg:set}
  \index{univalence axiom}%
  泛等公理直接蕴含 \uset 是一个范畴。
  利用泛等性还可以证明，任何由集合层面的结构（例如群、环、拓扑空间等等）构成的预范畴也都是范畴；详见 \cref{sec:sip}。
\end{eg}

我们还注意到下述事实。

\begin{lem}\label{ct:obj-1type}
  在范畴中，对象的类型是一个 1-类型。
\end{lem}

\begin{proof}
  只需证明：对任意 $a,b:A$，类型 $\id a b$ 都是一个集合。
  但 $\id a b$ 等价于 $a \cong b$，而后者是一个集合。
\end{proof}

\symlabel{isotoid}
我们记 $\isotoid$ 为由 \cref{ct:idtoiso} 给出的映射 $\idtoiso$ 的逆 $(a\cong b) \to (\id a b)$。
二者之间存在如下重要关系。

\begin{lem}\label{ct:idtoiso-trans}
  对于 $p:\id a a'$、$q:\id b b'$ 以及 $f:\hom_A(a,b)$，我们有
  \begin{equation}\label{ct:idtoisocompute}
    \id{\trans{(p,q)}{f}}
    {\idtoiso(q)\circ f \circ \inv{\idtoiso(p)}}.
  \end{equation}
\end{lem}

\begin{proof}
  通过归纳，可以假设 $p$ 和 $q$ 分别是 $\refl a$ 和 $\refl b$。
  此时，式~\eqref{ct:idtoisocompute} 的左边就是 $f$。
  而根据定义，$\idtoiso(\refl a)$ 是 $1_a$，且 $\idtoiso(\refl b)$ 是 $1_b$，因此式~\eqref{ct:idtoisocompute} 的右边是 $1_b\circ f\circ 1_a$，它与 $f$ 相等。
\end{proof}

类似地，可以证明
\begin{gather}
  \id{\idtoiso(\rev p)}{\inv {(\idtoiso(p))}}\\
  \id{\idtoiso(p\ct q)}{\idtoiso(q)\circ \idtoiso(p)}\\
  \id{\isotoid(f\circ e)}{\isotoid(e)\ct \isotoid(f)}
\end{gather}
等等。

\begin{eg}\label{ct:orders}
  在一个预范畴中，如果每个集合 $\hom_A(a,b)$ 都是一个命题，则等价地说，它是一个配备了纯粹关系 ``$\le$'' 的类型 $A_0$，该关系是自反的（$a\le a$）且传递的（若 $a\le b$ 且 $b\le c$，则 $a\le c$）。
  我们称其为\define{预序（preorder）}。
  \indexdef{preorder}

  在预序中，一个见证 $f: a\le b$ 是同构当且仅当存在某个见证 $g: b\le a$。
  因此，$a\cong b$ 就是命题 ``$a\le b$ 且 $b\le a$''。
  所以，一个预序 $A$ 是范畴当且仅当：(1) 每个类型 $a=b$ 是一个命题，并且 (2) 对于任意 $a,b:A_0$，存在一个函数 $(a\cong b) \to (a=b)$。
  换言之，$A_0$ 必须是一个集合，并且 $\le$ 必须是反对称的\index{relation!antisymmetric}（若 $a\le b$ 且 $b\le a$，则 $a=b$）。
  我们称其为\define{（偏）序（(partial) order）}或\define{偏序集（poset）}。
  \indexdef{partial order}%
  \indexdef{poset}%
\end{eg}

\begin{eg}\label{ct:gaunt}
  若 $A$ 是范畴，则 $A_0$ 是集合当且仅当对于任意 $a,b:A_0$，类型 $a\cong b$ 是一个命题。
  既然 $A$ 是范畴，这等价于说 $A$ 中的每个自同构都是恒同箭头。另一方面，若 $A$ 是一个 $A_0$ 为集合的预范畴，则 $A$ 恰在它是\emph{骨架的}\index{mathematics!classical}（任意两个同构的对象都相等）\emph{并且}每个自同构都是恒同箭头时，才是一个范畴。
  这类范畴有时被称为 \define{gaunt 范畴（gaunt category）}~\cite{bsp12infncats}。
  \indexdef{category!gaunt}%
  \indexdef{gaunt category}%
  \index{skeletal category}%
  \index{category!skeletal}%
  对我们的范畴而言，并没有真正的 ``骨架性'' 概念，除非把 \cref{ct:category} 本身视作这样一个概念。
\end{eg}

\begin{eg}\label{ct:discrete}
  对于任何 1-类型 $X$，存在一个范畴，以 $X$ 为其对象类型，并以 $\hom(x,y) \defeq (x=y)$ 为态射集。
  如果 $X$ 是集合，我们称之为 $X$ 上的\define{离散（discrete）}
  \indexdef{category!discrete}%
  \indexdef{discrete!category}%
  范畴。
  一般地，我们称之为\define{群胚（groupoid）}
  \indexdef{groupoid}
  （参见 \cref{ct:groupoids}）。
\end{eg}

\begin{eg}\label{ct:fundgpd}
  对于\emph{任何}类型 $X$，存在一个预范畴，以 $X$ 为其对象类型，并以 $\hom(x,y) \defeq \pizero{x=y}$ 为态射集。
  其复合运算
  \[ \pizero{y=z} \to \pizero{x=y} \to \pizero{x=z} \]
  通过基于截断的归纳法，由拼接 $(y=z)\to(x=y)\to(x=z)$ 定义。
  我们称其为 $X$ 的\define{基本预群胚（fundamental pregroupoid）}
  \indexdef{fundamental!pregroupoid}%
  \indexsee{pregroupoid, fundamental}{fundamental pregroupoid}%
  。（事实上，我们已经在 \cref{sec:van-kampen} 中遇到过它；亦可参见 \cref{ex:rezk-vankampen}。）
\end{eg}

\begin{eg}\label{ct:hoprecat}
  存在一个预范畴，其对象类型为 \type，且 $\hom(X,Y) \defeq \pizero{X\to Y}$，复合通过基于截断的归纳法由通常的复合 $(Y\to Z) \to (X\to Y) \to (X\to Z)$ 定义。
  我们称之为\define{类型的同伦预范畴（homotopy precategory of types）}。
  \indexdef{precategory!of types}%
  \index{homotopy!category of types@(pre)category of types}%
\end{eg}

\begin{eg}\label{ct:rel}
  令 \urel 为如下预范畴：
  \begin{itemize}
  \item 其对象是集合。
  \item $\hom_{\urel}(X,Y) = X\to Y\to \prop$。
  \item 对于一个集合 $X$，我们有 $1_X(x,x') \defeq (x=x')$。
  \item 对于 $R:\hom_{\urel}(X,Y)$ 和 $S:\hom_{\urel}(Y,Z)$，它们的复合定义为
    \[ (S\circ R)(x,z) \defeq \Brck{\sm{y:Y} R(x,y) \times S(y,z)}.\]
  \end{itemize}
  假设 $R:\hom_{\urel}(X,Y)$ 是一个同构，其逆为 $S$。
  我们观察到如下事实。
  \begin{enumerate}
  \item 如果 $R(x,y)$ 且 $S(y',x)$，那么 $(R\circ S)(y',y)$，于是 $y'=y$。
    类似地，如果 $R(x,y)$ 且 $S(y,x')$，那么 $x=x'$。\label{item:rel1}
  \item 对于任意 $x$，我们有 $x=x$，因此 $(S\circ R)(x,x)$。
    于是，仅只存在 $y:Y$ 使得 $R(x,y)$ 且 $S(y,x)$。\label{item:rel2}
  \item 假设 $R(x,y)$。
    由~\ref{item:rel2}，仅只存在 $y'$ 使得 $R(x,y')$ 且 $S(y',x)$。
    但由~\ref{item:rel1}，仅有 $y'=y$，而由于 $Y$ 是集合，故 $y'=y$。
    因此，沿此相等传输 $S(y',x)$，我们得到 $S(y,x)$。
    总之，$R(x,y)\to S(y,x)$。
    类似地，$S(y,x) \to R(x,y)$。\label{item:rel3}
  \item 如果 $R(x,y)$ 且 $R(x,y')$，那么由~\ref{item:rel3}，$S(y',x)$，从而由~\ref{item:rel1}，$y=y'$。
    因此，对于任意 $x$，至多存在一个 $y$ 使得 $R(x,y)$。
    又由~\ref{item:rel2}，这样的 $y$ 仅只存在，故而这样的 $y$ 存在。
  \end{enumerate}
  综上所述，如果 $R:\hom_{\urel}(X,Y)$ 是一个同构，那么对每个 $x:X$ 恰好存在一个 $y:Y$ 使得 $R(x,y)$，对偶地亦然。
  因此，存在一个函数 $f:X\to Y$ 将每个 $x$ 映到这个 $y$，并且这是一个等价；于是 $X=Y$。
  再稍加论证，我们得出结论：\urel 是一个范畴。
\end{eg}

现在我们本可以限制自己只考虑范畴而非预范畴。
相反，我们将同时为预范畴和范畴发展许多概念，以强调范畴的性态远优于预范畴，也远优于经典\index{mathematics!classical}数学中的普通范畴。

在\crefrange{sec:strict-categories}{sec:dagger-categories}中我们还会看到，在一些不那么常规的语境下，除了范畴之外，某些特定种类的预范畴也有其用武之地，它们各自以不同的方式``修正''了对象的相等关系。
这正凸显了预范畴的 ``预'' 之含义：它们是定义多种重要范畴结构的基本原料。

\section{函子与变换}
\label{sec:transfors}

接下来的定义相当显然，无需修改。

\begin{defn}\label{ct:functor}
  设 $A$ 和 $B$ 为预范畴。
  一个从 $A$ 到 $B$ 的\define{函子（functor）}
  \indexdef{functor}%
  $F:A\to B$ 包含以下数据：
  \begin{enumerate}
  \item 一个函数 $F_0:A_0\to B_0$，通常也记作 $F$。
  \item 对任意 $a,b:A$，一个函数 $F_{a,b}:\hom_A(a,b) \to \hom_B(Fa,Fb)$，通常也记作 $F$。
  \item 对任意 $a:A$，有 $\id{F(1_a)}{1_{Fa}}$。
  \item 对任意 $a,b,c:A$ 以及 $f:\hom_A(a,b)$ 和 $g:\hom_A(b,c)$，有
    \[\id{F(g\circ f)}{Fg\circ Ff}.\]\label{ct:functor:comp}
  \end{enumerate}
\end{defn}

注意，通过对恒等进行归纳可知，函子也保持 \idtoiso。

\begin{defn}\label{ct:nattrans}
  对于函子 $F,G:A\to B$，一个从 $F$ 到 $G$ 的\define{自然变换（natural transformation）}
  \indexdef{natural!transformation}%
  \indexsee{transformation!natural}{natural transformation}%
  $\gamma:F\to G$ 包含以下数据：
  \begin{enumerate}
  \item 对任意 $a:A$，一个态射 $\gamma_a:\hom_B(Fa,Ga)$（称为 ``分量''）。
  \item 对任意 $a,b:A$ 和 $f:\hom_A(a,b)$，有 $\id{Gf\circ \gamma_a}{\gamma_b\circ Ff}$（称为 ``自然性公理''）。
  \end{enumerate}
\end{defn}

由于每个类型 $\hom_B(Fa,Gb)$ 都是集合，其相等类型是命题。
因此，自然性公理是命题，从而自然变换的相等由其分量的相等决定。
特别地，对任意 $F$ 和 $G$，从 $F$ 到 $G$ 的自然变换类型也是集合。

类似地，函子的相等由其函数 $A_0\to B_0$ 以及（沿此传输的）相应 hom-集上的函数所决定。

\begin{defn}\label{ct:functor-precat}
  \indexdef{precategory!of functors}%
  对于预范畴 $A,B$，存在一个预范畴 $B^A$，称为\define{函子预范畴（functor precategory）}，其定义为：
  \begin{itemize}
  \item $(B^A)_0$ 是从 $A$ 到 $B$ 的函子类型。
  \item $\hom_{B^A}(F,G)$ 是从 $F$ 到 $G$ 的自然变换类型。
  \end{itemize}
\end{defn}

\begin{proof}
  我们定义 $(1_F)_a\defeq 1_{Fa}$。
  其自然性由预范畴的单位公理保证。
  对于 $\gamma:F\to G$ 和 $\delta:G\to H$，我们定义 $(\delta\circ\gamma)_a\defeq \delta_a\circ \gamma_a$。
  其自然性由结合律保证。
  类似地，$B^A$ 的单位律和结合律继承自 $B$ 的相应公理。
\end{proof}

\begin{lem}\label{ct:natiso}
  \index{natural!isomorphism}%
  \index{isomorphism!natural}%
  一个自然变换 $\gamma:F\to G$ 是 $B^A$ 中的同构，当且仅当每个 $\gamma_a$ 都是 $B$ 中的同构。
\end{lem}

\begin{proof}
  若 $\gamma$ 是同构，则存在其逆 $\delta:G\to F$。
  根据 $B^A$ 中复合的定义，有 $(\delta\gamma)_a\jdeq \delta_a\gamma_a$，类似地有 $(\gamma\delta)_a\jdeq \gamma_a\delta_a$。
  因此，由 $\id{\delta\gamma}{1_F}$ 和 $\id{\gamma\delta}{1_G}$ 可推出 $\id{\delta_a\gamma_a}{1_{Fa}}$ 与 $\id{\gamma_a\delta_a}{1_{Ga}}$，所以 $\gamma_a$ 是同构。

  反之，假设每个 $\gamma_a$ 都是同构，其逆记作 $\delta_a$。
  我们定义一个以 $\delta_a$ 为分量的自然变换 $\delta:G\to F$；对其自然性公理，我们有
  \[ Ff\circ \delta_a = \delta_b\circ \gamma_b\circ Ff \circ \delta_a = \delta_b\circ Gf\circ \gamma_a\circ \delta_a = \delta_b\circ Gf. \]
  由于自然变换的复合与恒同由其分量决定，我们有 $\id{\gamma\delta}{1_G}$ 和 $\id{\delta\gamma}{1_F}$。
\end{proof}

下面的结果是基础性的。

\begin{thm}\label{ct:functor-cat}
  \indexdef{category!of functors}%
  \indexdef{functor!category of}%
  若 $A$ 是预范畴而 $B$ 是范畴，则 $B^A$ 也是范畴。
\end{thm}

\begin{proof}
  令 $F,G:A\to B$；我们必须证明 $\idtoiso:(\id{F}{G}) \to (F\cong G)$ 是一个等价。

  为了给出其逆，假设 $\gamma:F\cong G$ 是一个自然同构。
  那么对于任意 $a:A$，我们有一个同构 $\gamma_a:Fa \cong Ga$，从而得到一个相等 $\isotoid(\gamma_a):\id{Fa}{Ga}$。
  根据函数外延性，我们得到一个相等 $\bar{\gamma}:\id[(A_0\to B_0)]{F_0}{G_0}$。

  由于函子的后两条公理仅仅是命题，为证明 $\id{F}{G}$，只需证明对于任意 $a,b:A$，函数
  \begin{align*}
    F_{a,b}&:\hom_A(a,b) \to \hom_B(Fa,Fb)\mathrlap{\qquad\text{and}}\\
    G_{a,b}&:\hom_A(a,b) \to \hom_B(Ga,Gb)
  \end{align*}
  沿 $\bar\gamma$ 传输后变为相等。
  根据函数外延性的计算规则，当作用于 $a$ 时，$\bar\gamma$ 等于 $\isotoid(\gamma_a)$。
  而由 \cref{ct:idtoiso-trans} 可知，将 $Ff:\hom_B(Fa,Fb)$ 沿 $\isotoid(\gamma_a)$ 和 $\isotoid(\gamma_b)$ 传输，等于复合 $\gamma_b\circ Ff\circ \inv{(\gamma_a)}$；再根据 $\gamma$ 的自然性，此复合即等于 $Gf$。

  这就完成了函数 $(F\cong G) \to (\id F G)$ 的定义。
  现在考虑复合
  \[ (\id F G) \to (F\cong G) \to (\id F G). \]
  由于 hom 集都是集合，它们的相等类型是命题，因此要证明两个相等 $p,q:\id F G$ 相等，只需证明 $\id[\id{F_0}{G_0}]{p}{q}$。
  但在 $\bar\gamma$ 的定义中，若 $\gamma$ 具有 $\idtoiso(p)$ 的形式，则 $\gamma_a$ 会等于 $\idtoiso(p_a)$（通过对 $p$ 归纳容易证明）。
  于是，$\isotoid(\gamma_a)$ 将等于 $p_a$，从而根据函数外延性，我们会有 $\id{\bar\gamma}{p}$，这正是所需要的。

  最后，考虑复合
  \[(F\cong G)\to  (\id F G) \to (F\cong G). \]
  由于自然变换的相等性可以逐分量地检验，只需证明对每个 $a$ 有 $\id{\idtoiso(\bar\gamma)_a}{\gamma_a}$。
  但如前所述，我们有 $\id{\idtoiso(\bar\gamma)_a}{\idtoiso((\bar\gamma)_a)}$，同时根据函数外延性的计算规则有 $\id{(\bar\gamma)_a}{\isotoid(\gamma_a)}$。
  因为 $\isotoid$ 和 $\idtoiso$ 互逆，故得 $\id{\idtoiso(\bar\gamma)_a}{\gamma_a}$，如所需。
\end{proof}

特别地，范畴（而非预范畴）之间自然同构的函子是相等的。

\mentalpause

我们现在定义函子与自然变换所有常用的复合方式。

\begin{defn}\label{ct:functor-composition}
  对于函子 $F:A\to B$ 和 $G:B\to C$，它们的复合 $G\circ F:A\to C$ 定义为：
  \begin{itemize}
  \item 对象上的复合 $(G_0\circ F_0) : A_0 \to C_0$。
  \item 对每个 $a,b:A$，态射上的复合
    \[(G_{Fa,Fb}\circ F_{a,b}):\hom_A(a,b) \to \hom_C(GFa,GFb).\]。
  \end{itemize}
  公理容易验证。
\end{defn}

\begin{defn}\label{ct:whisker}
  对于函子 $F:A\to B$ 和 $G,H:B\to C$ 以及自然变换 $\gamma:G\to H$，其复合 $(\gamma F):GF\to HF$ 定义为：
  \begin{itemize}
  \item 对每个 $a:A$，分量 $\gamma_{Fa}$。
  \end{itemize}
  自然性容易验证。
  类似地，对于上述的 $\gamma$ 和函子 $K:C\to D$，其复合 $(K\gamma):KG\to KH$ 定义为：
  \begin{itemize}
  \item 对每个 $b:B$，分量 $K(\gamma_b)$。
  \end{itemize}
\end{defn}

\begin{lem}\label{ct:interchange}
  \index{interchange law}%
  对于函子 $F,G:A\to B$ 和 $H,K:B\to C$ 以及自然变换 $\gamma:F\to G$ 和 $\delta:H\to K$，我们有
  \[\id{(\delta G)(H\gamma)}{(K\gamma)(\delta F)}.\]
\end{lem}

\begin{proof}
  只需逐分量地验证：对 $a:A$，我们有
  \begin{align*}
    ((\delta G)(H\gamma))_a
    &\jdeq (\delta G)_{a}(H\gamma)_a\\
    &\jdeq \delta_{Ga}\circ H(\gamma_a)\\
    &= K(\gamma_a) \circ \delta_{Fa} \tag{by naturality of $\delta$}\\
    &\jdeq (K \gamma)_a\circ (\delta F)_a\\
    &\jdeq ((K \gamma)(\delta F))_a.\qedhere
  \end{align*}
\end{proof}

\index{horizontal composition!of natural transformations}%
\index{classical!category theory}%
在经典范畴论中，人们将 $\gamma:F\to G$ 和 $\delta:H\to K$ 的"水平复合"定义为 ${(\delta G)(H\gamma)}$ 和 ${(K\gamma)(\delta F)}$ 的共同值。
我们将不这样做，因为尽管这两个变换相等，它们并非\emph{定义性}相等。
这也意味着：我们可以无歧义地使用符号 $\circ$（或并置）来表示所有种类的复合——只有一种方式来复合两个自然变换（而非将自然变换与函子在某一侧复合）。

\begin{lem}\label{ct:functor-assoc}
  \index{associativity!of functor composition}
  函子的复合满足结合律：$\id{H(GF)}{(HG)F}$。
\end{lem}

\begin{proof}
  由于函数的复合是结合的，这对于对象和 hom 集上的作用立即成立。
  又因为 hom 集是集合，其余数据自动满足。
\end{proof}

\cref{ct:functor-assoc} 中的等式同样不是定义性相等。
（函数的复合是定义性结合的，但构成函子的公理也必须被复合，这破坏了定义性结合律。）因此，我们还需要知道结合律的\emph{一致性}\index{coherence}。

\begin{lem}\label{ct:pentagon}
  \index{associativity!of functor composition!coherence of}%
  \cref{ct:functor-assoc} 是一致的，即下列五边形\index{pentagon, Mac Lane} 等式交换：
  \[ \xymatrix{ & K(H(GF)) \ar@{=}[dl] \ar@{=}[dr]\\
    (KH)(GF) \ar@{=}[d] && K((HG)F) \ar@{=}[d]\\
    ((KH)G)F && (K(HG))F \ar@{=}[ll].}
  \]
\end{lem}

\begin{proof}
  如 \cref{ct:functor-assoc} 中那样，这对于对象上的作用是显然的，其余部分自动成立。
\end{proof}

此后，我们将滥用记号，用 $H\circ G\circ F$ 或 $HGF$ 兼指 $H(GF)$ 与 $(HG)F$，在需要时沿 \cref{ct:functor-assoc} 做传输。
对于单位，我们也有类似的相干性结果。

\begin{lem}\label{ct:units}
  对于函子 $F:A\to B$，有等式 $\id{(1_B\circ F)}{F}$ 和 $\id{(F\circ 1_A)}{F}$，使得当再给定函子 $G:B\to C$ 时，下面的等式三角形交换：
  \[ \xymatrix{
    G\circ (1_B \circ F) \ar@{=}[rr] \ar@{=}[dr] &&
    (G\circ 1_B)\circ F \ar@{=}[dl] \\
    & G \circ F.}
  \]
\end{lem}

这些思路的进一步发展，可参见 \cref{ct:pre2cat,ct:2cat}。

\section{伴随函子}
\label{sec:adjunctions}

伴随函子的定义是直截了当的；其真正有趣之处来自证明相关性。

\begin{defn}\label{ct:adjoints}
  一个函子 $F:A\to B$ 被称为\define{左伴随（left adjoint）}
  \indexdef{left!adjoint}%
  \indexdef{adjoint!functor}%
  \indexdef{right!adjoint}%
  \indexdef{adjoint!functor}%
  \index{functor!adjoint}%
  如果存在
  \begin{itemize}
  \item 一个函子 $G:B\to A$。
  \item 一个自然变换 $\eta:1_A \to GF$（称为\define{单位（unit）}\indexdef{unit!of an adjunction}）。
  \item 一个自然变换 $\epsilon:FG\to 1_B$（称为\define{余单位（counit）}\indexdef{counit of an adjunction}）。
  \item $\id{(\epsilon F)(F\eta)}{1_F}$。
  \item $\id{(G\epsilon)(\eta G)}{1_G}$。
  \end{itemize}
\end{defn}

最后两个等式被称为\define{三角恒等式（triangle identities）}\indexdef{triangle!identity}或\define{之字形恒等式（zigzag identities）}\indexdef{zigzag identity}。
\indexdef{identity!triangle}\indexdef{identity!zigzag}
我们请读者类似地定义右伴随。

\begin{lem}\label{ct:adjprop}
  若 $A$ 是一个范畴（但 $B$ 可以只是预范畴），则类型"$F$ 是左伴随"是一个命题。
\end{lem}

\begin{proof}
  假设给定满足三角恒等式的 $(G,\eta,\epsilon)$ 以及另一组 $(G',\eta',\epsilon')$。
  定义 $\gamma:G\to G'$ 为 $(G'\epsilon)(\eta' G)$，并定义 $\delta:G'\to G$ 为 $(G\epsilon')(\eta G')$。
  那么，利用 \cref{ct:interchange} 与三角恒等式，有
  \begin{align*}
    \delta\gamma &=
    (G\epsilon')(\eta G')(G'\epsilon)(\eta'G)\\
    &= (G\epsilon')(G F G'\epsilon)(\eta G' F G)(\eta'G)\\
    &= (G\epsilon)(G\epsilon'FG)(G F \eta' G)(\eta G)\\
    &= (G\epsilon)(\eta G)\\
    &= 1_G
  \end{align*}
  类似地，可以证明 $\id{\gamma\delta}{1_{G'}}$，因此 $\gamma$ 是一个自然同构 $G\cong G'$。
  根据 \cref{ct:functor-cat}，我们得到一个相等 $\id G {G'}$。

  现在我们需要知道，当 $\eta$ 和 $\epsilon$ 沿这个相等做传输后，它们分别等于 $\eta'$ 和 $\epsilon'$。
  根据 \cref{ct:idtoiso-trans}，这一传输通过视情况与 $\gamma$ 或 $\delta$ 复合来给出。
  对于 $\eta$，得到
  \begin{equation*}
    (G'\epsilon F)(\eta'GF)\eta
    = (G'\epsilon F)(G'F\eta)\eta'
    = \eta'
  \end{equation*}
  其中用到了 \cref{ct:interchange} 和三角恒等式。
  $\epsilon$ 的情况类似。
  最后，由于 hom-集合都是集合，三角恒等式在传输后自动保持正确。
\end{proof}

在 \cref{sec:yoneda} 中，我们将给出 \cref{ct:adjprop} 的另一个证明。

\section{等价}
\label{sec:equivalences}

在范畴论中，通常将\emph{范畴的等价（equivalence of categories）}定义为一个函子 $F:A\to B$，使得存在一个函子 $G:B\to A$ 以及自然同构 $F G \cong 1_B$ 和 $G F \cong 1_A$。
然而，与"是一个伴随"这一性质不同，若不进行截断，这样定义的"等价"将不是一个命题，其原因与拟逆的类型行为不佳相同（参见 \cref{sec:quasi-inverses}）。
正如 \cref{sec:hae} 中那样，我们可以通过使用\emph{伴随等价（adjoint equivalence）}这一常用概念来避免此问题。
\indexdef{adjoint!equivalence!of (pre)categories}

\begin{defn}\label{ct:equiv}
  一个函子 $F:A\to B$ 是一个\define{(预)范畴的等价（equivalence of (pre)categories）}
  \indexdef{equivalence!of (pre)categories}%
  \indexdef{category!equivalence of}%
  \indexdef{precategory!equivalence of}%
  \index{functor!equivalence}%
  如果它是一个左伴随，并且其单位 $\eta$ 和余单位 $\epsilon$ 都是同构。
  我们记 $\cteqv A B$ 为从 $A$ 到 $B$ 的（预）范畴的等价的类型。
\end{defn}

根据 \cref{ct:adjprop,ct:isoprop}，若 $A$ 是一个范畴，则类型"$F$ 是一个预范畴的等价"是一个命题。

\begin{lem}\label{ct:adjointification}
  如果对于一个函子 $F:A\to B$，存在 $G:B\to A$ 以及同构 $GF\cong 1_A$ 和 $FG\cong 1_B$，那么 $F$ 是一个预范畴的等价。
\end{lem}

\begin{proof}
  其证明类似于 \cref{thm:equiv-iso-adj} 中关于类型的等价的证明。
\end{proof}

\begin{defn}\label{ct:full-faithful}
  我们称一个函子 $F:A\to B$ 是\define{忠实的（faithful）}
  \indexdef{functor!faithful}%
  \index{faithful functor}%a
  如果对于所有 $a,b:A$，函数
  \[F_{a,b}:\hom_A(a,b) \to \hom_B(Fa,Fb)\]
  是单的；称它是\define{全的（full）}
  \indexdef{functor!full}%
  \indexdef{full functor}%
  如果对于所有 $a,b:A$ 这个函数是满的。
  如果它既是忠实的又是全的（从而每个 $F_{a,b}$ 都是一个等价），则称 $F$ 是\define{全忠实的（fully faithful）}。
  \indexdef{functor!fully faithful}%
  \indexdef{fully faithful functor}%
\end{defn}

\begin{defn}\label{ct:split-essentially-surjective}
  我们称一个函子 $F:A\to B$ 是\define{分裂本质满的（split essentially surjective）}
  \indexdef{functor!split essentially surjective}%
  \indexdef{split!essentially surjective functor}%
  如果对于所有 $b:B$，都存在一个 $a:A$ 使得 $Fa\cong b$。
\end{defn}

\begin{lem}\label{ct:ffeso}
  对于任意预范畴 $A$ 和 $B$ 以及函子 $F:A\to B$，下列两个类型等价。
  \begin{enumerate}
  \item $F$ 是预范畴的等价。\label{item:ct:ffeso1}
  \item $F$ 是全忠实的且分裂本质满的。\label{item:ct:ffeso2}
  \end{enumerate}
\end{lem}

\begin{proof}
  假设 $F$ 是预范畴的等价，并指定了 $G,\eta,\epsilon$。
  那么我们有函数
  \begin{align*}
      \hom_B(Fa,Fb) &\to \hom_A(a,b),\\
      g &\mapsto \inv{\eta_b}\circ G(g)\circ \eta_a.
  \end{align*}
  对于 $f:\hom_A(a,b)$，我们有
  \[ \inv{\eta_{b}}\circ G(F(f))\circ \eta_{a}  =
  \inv{\eta_{b}} \circ \eta_{b} \circ f=
  f
  \]
  而对于 $g:\hom_B(Fa,Fb)$，我们有
  \begin{align*}
    F(\inv{\eta_b} \circ G(g)\circ\eta_a)
    &= F(\inv{\eta_b})\circ F(G(g))\circ F(\eta_a)\\
    &= \epsilon_{Fb}\circ F(G(g))\circ F(\eta_a)\\
    &= g\circ\epsilon_{Fa}\circ F(\eta_a)\\
    &= g
  \end{align*}
  其中用到了 $\epsilon$ 的自然性以及两次三角恒等式。
  因此，$F_{a,b}$ 是一个等价，故 $F$ 是全忠实的。
  最后，对于任意 $b:B$，我们有 $Gb:A$ 以及 $\epsilon_b:FGb\cong b$。

  反过来，假设 $F$ 是全忠实且分裂本质满的。
  定义 $G_0:B_0\to A_0$，将 $b:B$ 映至由指定的本质分裂所给出的 $a:A$，并将同样以指定方式给出的同构 $FGb\cong b$ 记为 $\epsilon_b$。

  现在对任意 $g:\hom_B(b,b')$，定义 $G(g):\hom_A(Gb,Gb')$ 为满足 $\id{F(G(g))}{\inv{(\epsilon_{b'})}\circ g \circ \epsilon_b }$ 的唯一态射（该态射存在，因为 $F$ 是全忠实的）。
  最后，对 $a:A$，定义 $\eta_a:\hom_A(a,GFa)$ 为满足 $\id{F\eta_a}{\inv{\epsilon_{Fa}}}$ 的唯一态射。
  容易验证 $G$ 是函子，且 $(G,\eta,\epsilon)$ 将 $F$ 展现为预范畴的等价。

  现考虑复合~\ref{item:ct:ffeso1}$\to$\ref{item:ct:ffeso2}$\to$\ref{item:ct:ffeso1}。
  我们显然恢复了同一个函数 $G_0:B_0 \to A_0$。
  关于 $G$ 在 hom 集上的作用，我们需要证明对 $g:\hom_B(b,b')$，$G(g)$ 是满足 $F(G(g)) = \inv{(\epsilon_{b'})}\circ g \circ \epsilon_b$ 的那个（必然唯一的）态射。
  但此等式由 $\epsilon$ 的自然性（已假设其成立）可得。
  我们也显然恢复了 $\epsilon$，而 $\eta$ 由 $\id{F\eta_a}{\inv{\epsilon_{Fa}}}$ 唯一刻画（这是预范畴的等价之结构中所假设成立的三角恒等式之一）。
  因此，此复合等于恒等映射。

  最后，考虑另一个方向的复合~\ref{item:ct:ffeso2}$\to$\ref{item:ct:ffeso1}$\to$\ref{item:ct:ffeso2}。
  由于"是全忠实的"是一个命题，只需注意到对每个 $b:B$，我们恢复了相同的 $a:A$ 和同构 $Fa \cong b$。
  这是清楚的，因为我们正是用这个函数和同构在~\ref{item:ct:ffeso1} 中定义了 $G_0$ 和 $\epsilon$，而它们又恰好是我们恢复~\ref{item:ct:ffeso2} 时所用的。
  因此，两个方向的复合都等于恒等映射，于是我们得到等价 \eqv{\text{\ref{item:ct:ffeso1}}}{\text{\ref{item:ct:ffeso2}}}。
\end{proof}

然而，如果 $A$ 不是范畴，那么 \cref{ct:ffeso} 中的两个类型都未必是命题。
这提示我们也应考虑以下概念。

\begin{defn}\label{ct:essentially-surjective}
  一个函子 $F:A\to B$ 是\define{本质满的}
  \indexdef{functor!essentially surjective}%
  \indexdef{essentially surjective functor}%
  若对所有 $b:B$，\emph{仅仅}存在 $a:A$ 使得 $Fa\cong b$。
  我们称 $F$ 是\define{弱等价}
  \indexsee{equivalence!of (pre)categories!weak}{weak equivalence}%
  \indexdef{weak equivalence!of precategories}%
  \indexsee{functor!weak equivalence}{weak equivalence}%
  若它是全忠实的且本质满的。
\end{defn}

是弱等价\emph{总是}一个命题。
然而，对于范畴而言，弱等价与等价之间没有区别。

\index{acceptance}


\begin{lem}\label{ct:catweq}
  若 $F:A\to B$ 是全忠实的且 $A$ 是范畴，则对任意 $b:B$，类型 $\sm{a:A} (Fa\cong b)$ 是一个命题。
  因此，范畴之间的函子为等价当且仅当它为弱等价。
\end{lem}

\begin{proof}
  假设在 $\sm{a:A} (Fa\cong b)$ 中给定 $(a,f)$ 和 $(a',f')$。
  那么 $\inv{f'}\circ f$ 是一个同构 $Fa \cong Fa'$。
  由于 $F$ 是全忠实的，我们有 $g:a\cong a'$ 满足 $Fg = \inv{f'}\circ f$。
  又因为 $A$ 是范畴，我们有 $p:a=a'$ 满足 $\idtoiso(p)=g$。
  现在 $Fg = \inv{f'}\circ f$ 意味着 $\trans{(\map{(F_0)}{p})}{f} = f'$，从而（根据依值对类型中相等的刻画）得到 $(a,f)=(a',f')$。

  因此，对于定义域为范畴的全忠实函子，本质满性与分裂本质满性等价，故弱等价性与等价性等价。
\end{proof}

这是我们的范畴论相对于基于集合的方法的一个重要优势。
若采用纯基于集合的范畴定义，陈述"每个全忠实且本质满的函子都是范畴的等价"等价于选择公理 \choice{}。
而在这里，我们可以免费获得它，作为唯一选择原理的范畴论版本（\cref{sec:unique-choice}）。
（事实上，这一性质在预范畴中刻画了范畴；见 \cref{sec:rezk}。）

另一方面，下列关于范畴的等价的刻画或许更为有用。

\begin{defn}\label{ct:isocat}
  一个函子 $F:A\to B$ 是一个\define{（预）范畴的同构（isomorphism of (pre)categories）}
  \indexdef{isomorphism!of (pre)categories}%
  \indexdef{category!isomorphism of}%
  \indexdef{precategory!isomorphism of}%
  ，如果 $F$ 是全忠实的且 $F_0:A_0\to B_0$ 是类型之间的等价。
\end{defn}

这个定义是我们的一般规则（参见 \cref{sec:basics-equivalences}）的一个例外，该规则规定"同构"一词只用于集合和类似集合的对象。
然而，它在这里确实传达了恰当的含义，因为对于一般的预范畴而言，同构强于等价。

注意，作为预范畴的同构总只是一个性质。
记 $A\cong B$ 为从 $A$ 到 $B$ 的（预）范畴同构所构成的类型。

\begin{lem}\label{ct:isoprecat}
  对于预范畴 $A$ 和 $B$ 以及函子 $F:A\to B$，以下条件等价。
  \begin{enumerate}
  \item $F$ 是预范畴的同构。\label{item:ct:ipc1}
  \item 存在 $G:B\to A$, $\eta:1_A = GF$ 以及 $\epsilon:FG=1_B$，使得\label{item:ct:ipc2}
    \begin{equation}
      \apfunc{(\lam{H} F H)}({\eta}) = \apfunc{(\lam{K} K F)}({\opp\epsilon}).\label{eq:ct:isoprecattri}
    \end{equation}
  \item 仅仅存在 $G:B\to A$, $\eta:1_A = GF$ 以及 $\epsilon:FG=1_B$。\label{item:ct:ipc3}
  \end{enumerate}
\end{lem}

注意，如果 $B_0$ 不是 1-类型，那么~\eqref{eq:ct:isoprecattri} 可能不是一个命题。

\begin{proof}
  首先注意，由于同态集是集合，函子等式之间的等式由其对象部分唯一决定。
  因此，由函数外延性，~\eqref{eq:ct:isoprecattri} 等价于
  \begin{equation}
    \map{(F_0)}{\eta_0}_a = \opp{(\epsilon_0)}_{F_0 a}.\label{eq:ct:ipctri}
  \end{equation}
  对所有 $a:A_0$ 成立。注意，这恰好就是 $G_0$、$\eta_0$ 和 $\epsilon_0$ 用以证明 $F_0$ 是类型上的半伴随等价所需的三角恒等式。

  现在假设~\ref{item:ct:ipc1}。
  令 $G_0:B_0 \to A_0$ 为 $F_0$ 的逆，并令 $\eta_0: \idfunc[A_0] = G_0 F_0$ 和 $\epsilon_0:F_0G_0 = \idfunc[B_0]$ 满足三角恒等式，这恰好就是~\eqref{eq:ct:ipctri}。现在如下定义 $G_{b,b'}:\hom_B(b,b') \to \hom_A(G_0b,G_0b')$：
  \[ G_{b,b'}(g) \defeq
  \inv{(F_{G_0b,G_0b'})}\Big(\idtoiso(\opp{(\epsilon_0)}_{b'}) \circ g \circ \idtoiso((\epsilon_0)_b)\Big)
  \]
  （这里用到了 $F$ 全忠实的假设）。由于 \idtoiso 将逆映射为逆，将拼接映射为复合，且 $F$ 是函子，可知 $G$ 也是一个函子。

  根据定义，我们有 $(GF)_0 \jdeq G_0 F_0$，而根据 $\eta_0$，这等于 $\idfunc[A_0]$。
  为了得到 $1_A = GF$，我们需要证明：当沿着 $\eta_0$ 迁移时，$\hom_A(a,a')$ 上的恒等函数变为等于复合 $G_{Fa,Fa'} \circ F_{a,a'}$。
  换句话说，对于任意 $f:\hom_A(a,a')$ 我们必须有
  \begin{multline*}
    \idtoiso((\eta_0)_{a'}) \circ f \circ \idtoiso(\opp{(\eta_0)}_a)\\
    = \inv{(F_{GFa,GFa'})}\Big(\idtoiso(\opp{(\epsilon_0)}_{Fa'})
    \circ F_{a,a'}(f) \circ \idtoiso((\epsilon_0)_{Fa})\Big).
  \end{multline*}
  但这等价于
  \begin{multline*}
    (F_{GFa,GFa'})\Big(\idtoiso((\eta_0)_{a'}) \circ f \circ \idtoiso(\opp{(\eta_0)}_a)\Big)\\
    = \idtoiso(\opp{(\epsilon_0)}_{Fa'})
    \circ F_{a,a'}(f) \circ \idtoiso((\epsilon_0)_{Fa}).
  \end{multline*}
  而这可以从 $F$ 的函子性、$F$ 保持 \idtoiso 的事实以及~\eqref{eq:ct:ipctri} 推出。因此我们得到 $\eta:1_A = GF$。

  在另一边，我们有 $(FG)_0\jdeq F_0 G_0$，根据 $\epsilon_0$，这等于 $\idfunc[B_0]$。
  为了得到 $FG=1_B$，我们需要证明：当沿着 $\epsilon_0$ 迁移时，$\hom_B(b,b')$ 上的恒等函数变为等于复合 $F_{Gb,Gb'} \circ G_{b,b'}$。
  也就是说，对于任意 $g:\hom_B(b,b')$ 我们必须有
  \begin{multline*}
    F_{Gb,Gb'}\Big(\inv{(F_{Gb,Gb'})}\Big(\idtoiso(\opp{(\epsilon_0)}_{b'}) \circ g \circ \idtoiso((\epsilon_0)_b)\Big)\Big)\\
    = \idtoiso((\opp{\epsilon_0})_{b'}) \circ g \circ \idtoiso((\epsilon_0)_b).
  \end{multline*}
  但这正是 $\inv{(F_{Gb,Gb'})}$ 是 $F_{Gb,Gb'}$ 的逆这一事实。并且我们之前已指出~\eqref{eq:ct:isoprecattri} 等价于~\eqref{eq:ct:ipctri}，故而~\ref{item:ct:ipc2} 成立。

  反过来，假设给定~\ref{item:ct:ipc2}；那么 $G$, $\eta$ 和 $\epsilon$ 的对象部分连同~\eqref{eq:ct:ipctri} 表明 $F_0$ 是类型之间的等价。并且对于 $a,a':A_0$，我们通过下式定义 $\overline{G}_{a,a'}: \hom_B(Fa,Fa') \to \hom_A(a,a')$：
  \begin{equation}
    \overline{G}_{a,a'}(g) \defeq \idtoiso(\opp{\eta})_{a'} \circ G(g) \circ \idtoiso(\eta)_a.\label{eq:ct:gbar}
  \end{equation}
  根据 $\idtoiso(\eta)$ 的自然性，对于任意 $f:\hom_A(a,a')$ 我们有
  \begin{align*}
    \overline{G}_{a,a'}(F_{a,a'}(f))
    &= \idtoiso(\opp{\eta})_{a'} \circ G(F(f)) \circ \idtoiso(\eta)_a\\
    &= \idtoiso(\opp{\eta})_{a'} \circ \idtoiso(\eta)_{a'} \circ f \\
    &= f.
  \end{align*}
  另一方面，对于 $g:\hom_B(Fa,Fa')$ 我们有
  \begin{align*}
    F_{a,a'}(\overline{G}_{a,a'}(g))
    &= F(\idtoiso(\opp{\eta})_{a'}) \circ F(G(g)) \circ F(\idtoiso(\eta)_a)\\
    &= \idtoiso(\epsilon)_{Fa'}
    \circ F(G(g))
    \circ \idtoiso(\opp{\epsilon})_{Fa}\\
    &= \idtoiso(\epsilon)_{Fa'}
    \circ \idtoiso(\opp{\epsilon})_{Fa'}
    \circ g\\
    &= g.
  \end{align*}
  （这里需要一些关于 \idtoiso 与须接兼容性的引理，我们留给读者陈述并证明。）
  因此，$F_{a,a'}$ 是一个等价，所以 $F$ 是全忠实的；即~\ref{item:ct:ipc1} 成立。

  现在，复合~\ref{item:ct:ipc1}$\to$\ref{item:ct:ipc2}$\to$\ref{item:ct:ipc1} 等于恒等映射，因为~\ref{item:ct:ipc1} 是一个命题。
  在另一边，追踪上述构造，我们看到复合~\ref{item:ct:ipc2}$\to$\ref{item:ct:ipc1}$\to$\ref{item:ct:ipc2} 本质上保持了对象部分 $G_0$, $\eta_0$, $\epsilon_0$ 以及~\eqref{eq:ct:isoprecattri} 的对象部分。在后三种情形中，对象部分就是全部，因为同态集是集合。

  因此，只需证明我们能恢复 $G$ 在同态集上的作用。
  换句话说，我们必须证明，如果 $g:\hom_B(b,b')$，那么
  \[ G_{b,b'}(g) =
  \overline{G}_{G_0b,G_0b'}\Big(\idtoiso(\opp{(\epsilon_0)}_{b'}) \circ g \circ \idtoiso((\epsilon_0)_b)\Big)
  \]
  其中 $\overline{G}$ 由~\eqref{eq:ct:gbar} 定义。然而，这可以从 $G$ 的函子性以及\emph{另一个}三角恒等式推出，我们在 \cref{cha:equivalences} 中已看到该恒等式等价于~\eqref{eq:ct:ipctri}。

  现在，由于~\ref{item:ct:ipc1} 是一个命题，~\ref{item:ct:ipc2} 亦然，因此只需证明它们与~\ref{item:ct:ipc3} 逻辑等价。
  显然，~\ref{item:ct:ipc2}$\to$\ref{item:ct:ipc3}，故可假设~\ref{item:ct:ipc3}。
  由于~\ref{item:ct:ipc1} 是一个命题，我们可以假设给定了 $G$, $\eta$ 和 $\epsilon$。
  那么 $G_0$ 连同 $\eta$ 和 $\epsilon$ 意味着 $F_0$ 是等价。
  此外，我们还有自然同构 $\idtoiso(\eta):1_A\cong GF$ 和 $\idtoiso(\epsilon):FG\cong 1_B$，因此由 \cref{ct:adjointification}，$F$ 是预范畴的等价，从而是全忠实的。
\end{proof}

从 \cref{ct:isoprecat}\ref{item:ct:ipc2} 以及函子范畴中的 $\idtoiso$，我们立即得出：预范畴的任何同构都是等价。
但对于预范畴而言，反过来不一定成立。

\begin{eg}\label{ct:chaotic}
  设 $X$ 为一个类型，$x_0:X$ 为其中一个元素，并令 $X_{\mathrm{ch}}$ 表示 $X$ 上的\emph{混沌预范畴（chaotic precategory）}\indexdef{chaotic precategory}或\emph{密着预范畴（indiscrete precategory）}\indexdef{indiscrete precategory}。
  根据定义，我们有 $(X_{\mathrm{ch}})_0\defeq X$，并且 $\hom_{X_{\mathrm{ch}}}(x,x') \defeq \unit$ 对所有 $x,x'$ 成立。
  那么唯一的函子 $X_{\mathrm{ch}}\to \unit$ 是预范畴的等价，但除非 $X$ 是可缩的，否则它不是同构。

  这个例子也表明，一个预范畴可以等价于一个范畴，而它本身却不必是范畴。
  当然，如果一个预范畴\emph{同构}于一个范畴，那么它本身必须是一个范畴。
\end{eg}

然而，对于范畴而言，这两个概念是一致的。

\begin{lem}\label{ct:eqv-levelwise}
  对于范畴 $A$ 和 $B$，函子 $F:A\to B$ 是范畴的等价当且仅当它是范畴的同构。
\end{lem}

\begin{proof}
  由于两者都只是命题，只需证明它们逻辑等价即可。
  首先假设 $F$ 是范畴的等价，其中 $(G,\eta,\epsilon)$ 已给定。
  我们已经知道 $F$ 是全忠实的。
  根据 \cref{ct:functor-cat}，自然同构 $\eta$ 和 $\epsilon$ 给出相等 $\id{1_A}{GF}$ 和 $\id{FG}{1_B}$，因而特别地有相等 $\id{\idfunc[A]}{G_0\circ F_0}$ 和 $\id{F_0\circ G_0}{\idfunc[B]}$。
  因此，$F_0$ 是类型的等价。

  反过来，假设 $F$ 是全忠实的，且 $F_0$ 是类型的等价，其逆记为 $G_0$。
  那么对于每个 $b:B$，我们有 $G_0 b:A$ 和相等 $\id{FGb}{b}$，从而得到同构 $FGb\cong b$。
  于是，根据 \cref{ct:ffeso}，$F$ 是范畴的等价。
\end{proof}

当然，（预）范畴之间还有第三种"相同"的概念：相等。
然而，泛等公理意味着它与同构一致。

\begin{lem}\label{ct:cat-eq-iso}
  如果 $A$ 和 $B$ 是预范畴，那么函数
  \[(\id A B) \to (A\cong B)\]
  （由恒同函子归纳定义）是类型的等价。
\end{lem}

\begin{proof}
  与 $\Sigma$ 类型的情形一样，给出 $\id A B$ 的一个元素，等价于给出：
  \begin{itemize}
  \item 一个相等 $P_0:\id{A_0}{B_0}$，
  \item 对每个 $a,b:A_0$，一个相等
    \[P_{a,b}:\id{\hom_A(a,b)}{\hom_B(\trans {P_0} a,\trans {P_0} b)},\]，
  \item 以及相等 $\id{\trans {(P_{a,a})} {1_a}}{1_{\trans {P_0} a}}$ 和
    \narrowequation{\id{\trans {(P_{a,c})} {gf}}{\trans {(P_{b,c})} g \circ \trans {(P_{a,b})} f}.}
  \end{itemize}
  （这里我们再次利用了 Hom 集的相等类型是命题这一事实。）
  然而，根据泛等性，这等价于给出：
  \begin{itemize}
  \item 一个类型的等价 $F_0:\eqv{A_0}{B_0}$，
  \item 对每个 $a,b:A_0$，一个类型的等价
    \[F_{a,b}:\eqv{\hom_A(a,b)}{\hom_B(F_0 (a),F_0 (b))},\]，
  \item 以及相等 $\id{F_{a,a}(1_a)}{1_{F_0 (a)}}$ 和 $\id{F_{a,c}(gf)}{F_{b,c} (g)\circ F_{a,b} (f)}$。
  \end{itemize}
  但这恰好构成一个函子 $F:A\to B$，且它是范畴的同构。
  并且通过对相等的归纳，这一等价 $\eqv{(\id A B)}{(A\cong B)}$ 等于由归纳得到的那个。
\end{proof}

因此，对于范畴，相等也与等价一致。
我们可以这样理解：范畴、函子和自然变换不仅构成一个前2-范畴，还构成一个2-范畴（参见 \cref{ct:pre2cat}）。

\begin{thm}\label{ct:cat-2cat}
  如果 $A$ 和 $B$ 是范畴，那么函数
  \[(\id A B) \to (\cteqv A B)\]
  （由恒同函子归纳定义）是类型的等价。
\end{thm}

\begin{proof}
  根据 \cref{ct:cat-eq-iso,ct:eqv-levelwise}。
\end{proof}

作为推论，范畴的类型是一个 $2$-类型。
因为 $\cteqv A B$ 是从 $A$ 到 $B$ 的函子类型的子类型（这些函子构成一个范畴的对象），所以它是 $1$-类型；因此，相等类型 $\id A B$ 也是 $1$-类型。

\section{米田引理}
\label{sec:yoneda}
\index{Yoneda!lemma|(}

回顾我们有一个范畴 \uset，其对象是集合，态射是函数。
我们现在证明，每个预范畴都有一个取值于 \uset 的 Hom 函子。
首先我们需要定义（预）范畴的反范畴和积。

\begin{defn}\label{ct:opposite-category}
  对于预范畴 $A$，其\define{反范畴（opposite）}
  \indexdef{opposite of a (pre)category}%
  \indexdef{precategory!opposite}%
  \indexdef{category!opposite}%
  $A\op$ 是一个预范畴，其对象类型与 $A$ 相同，其中 $\hom_{A\op}(a,b) \defeq \hom_A(b,a)$，恒同态射和复合运算继承自 $A$。
\end{defn}

\begin{defn}\label{ct:prod-cat}
  对于预范畴 $A$ 和 $B$，其\define{积（product）}
  \index{precategory!product of}%
  \index{category!product of}%
  \index{product!of (pre)categories}%
  $A\times B$ 是一个预范畴，其中 $(A\times B)_0 \defeq A_0 \times B_0$ 且
  \[\hom_{A\times B}((a,b),(a',b')) \defeq \hom_A(a,a') \times \hom_B(b,b').\]
  恒同态射定义为 $1_{(a,b)}\defeq (1_a,1_b)$，复合运算定义为
  \narrowequation{(g,g')(f,f') \defeq ((gf),(g'f')).}
\end{defn}

\begin{lem}\label{ct:functorexpadj}
  对于预范畴 $A,B,C$，以下类型是等价的。
  \begin{enumerate}
  \item 从 $A\times B$ 到 $C$ 的函子。
  \item 从 $A$ 到 $C^B$ 的函子。
  \end{enumerate}
\end{lem}

\begin{proof}
  给定 $F:A\times B\to C$，对任意 $a:A$，我们显然有一个函子 $F_a : B\to C$。
  这就给出了一个函数 $A_0 \to (C^B)_0$。
  其次，对任意 $f:\hom_A(a,a')$，对任意 $b:B$，我们有态射 $F_{(a,b),(a',b)}(f,1_b):F_a(b) \to F_{a'}(b)$。
  这些是自然变换 $F_a \to F_{a'}$ 的分量。
  关于 $a$ 的函子性容易验证，因此我们得到一个函子 $\hat{F}:A\to C^B$。

  反过来，假设给定 $G:A\to C^B$。
  那么对任意 $a:A$ 和 $b:B$，我们有对象 $G(a)(b):C$，从而给出一个函数 $A_0 \times B_0 \to C_0$。
  并且对于 $f:\hom_A(a,a')$ 和 $g:\hom_B(b,b')$，我们有 $\hom_C(G(a)(b), G(a')(b'))$ 中的态射
  \begin{equation*}
     G(a')_{b,b'}(g)\circ G_{a,a'}(f)_b = G_{a,a'}(f)_{b'} \circ  G(a)_{b,b'}(g)
  \end{equation*}。
  其函子性也容易验证，因此我们得到一个函子 $\check{G}:A\times B \to C$。

  最后，显然这两个操作互逆。
\end{proof}

现在，对任意预范畴 $A$，我们有一个 Hom 函子
\indexdef{hom-functor}%
\[\hom_A : A\op \times A \to \uset.\]
它将一对 $(a,b): (A\op)_0 \times A_0 \jdeq A_0 \times A_0$ 送到集合 $\hom_A(a,b)$。
对于态射 $(f,f') : \hom_{A\op\times A}((a,b),(a',b'))$，根据定义我们有 $f:\hom_A(a',a)$ 和 $f':\hom_A(b,b')$，因此可以定义
\begin{align*}
  (\hom_A)_{(a,b),(a',b')}(f,f')
  &\defeq (g \mapsto (f'gf))\\
  &: \hom_A(a,b) \to \hom_A(a',b').
\end{align*}
其函子性容易验证。

于是，由 \cref{ct:functorexpadj}，我们得到一个诱导函子 $\y:A\to \uset^{A\op}$，称之为 \define{米田嵌入（Yoneda embedding）}。
\indexdef{Yoneda!embedding}%
\indexdef{embedding!Yoneda}%

\begin{thm}[米田引理]\label{ct:yoneda}
  \indexdef{Yoneda!lemma}
  对任意预范畴 $A$，任意 $a:A$，以及任意函子 $F:\uset^{A\op}$，我们有一个同构
  \begin{equation}\label{eq:yoneda}
    \hom_{\uset^{A\op}}(\y a, F) \cong Fa.
  \end{equation}
  此外，该同构关于 $a$ 和 $F$ 均是自然的。
\end{thm}

\begin{proof}
  给定一个自然变换 $\alpha:\y a \to F$，我们可以考虑分量 $\alpha_a : \y a(a) \to F a$。
  由于 $\y a(a)\jdeq \hom_A(a,a)$，我们有 $1_a : \y a(a)$，因此 $\alpha_a(1_a) : F a$。
  这就给出了从 \eqref{eq:yoneda} 的左边到右边的函数 $(\alpha \mapsto \alpha_a(1_a))$。

  在另一个方向上，给定 $x:F a$，我们通过下式定义 $\alpha:\y a \to F$：
  \[\alpha_{a'}(f) \defeq F_{a,a'}(f)(x). \]
  其自然性容易验证，因此这给出了从 \eqref{eq:yoneda} 的右边到左边的函数。

  为了证明它们互逆，首先假设给定 $x:F a$。
  那么按上述方式定义 $\alpha$ 后，我们有 $\alpha_a(1_a) = F_{a,a}(1_a)(x) = 1_{F a}(x) = x$。
  另一方面，如果我们假设给定 $\alpha:\y a \to F$ 并按上述方式定义 $x$，那么对任意 $f:\hom_A(a',a)$，我们有
  \begin{align*}
    \alpha_{a'}(f)
    &= \alpha_{a'} (\y a_{a,a'}(f)(1_a))\\
    &= (\alpha_{a'}\circ \y a_{a,a'}(f))(1_a)\\
    &= (F_{a,a'}(f)\circ \alpha_a)(1_a)\\
    &= F_{a,a'}(f)(\alpha_a(1_a))\\
    &= F_{a,a'}(f)(x).
  \end{align*}
  因此，两个复合都等于恒等映射。
  我们将自然性的证明留给读者。
\end{proof}

\begin{cor}\label{ct:yoneda-embedding}
  米田嵌入 $\y :A\to \uset^{A\op}$ 是全忠实的。
\end{cor}

\begin{proof}
  根据 \cref{ct:yoneda}，我们有
  \[ \hom_{\uset^{A\op}}(\y a, \y b) \cong \y b(a) \jdeq \hom_A(a,b). \]
  容易验证，这个同构实际上就是 \y 在态射集上的作用。
\end{proof}

\begin{cor}\label{ct:yoneda-mono}
  如果 $A$ 是一个范畴，那么 $\y_0 : A_0 \to (\uset^{A\op})_0$ 是一个嵌入。
  特别地，如果 $\y a = \y b$，那么 $a=b$。
\end{cor}

\begin{proof}
  根据 \cref{ct:yoneda-embedding}，\y 在同构集上诱导出同构。
  但由于 $A$ 和 $\uset^{A\op}$ 是范畴且 \y 是函子，这等价于在相等类型上诱导出同构，而后者正是嵌入的定义。
\end{proof}

\begin{defn}\label{ct:representable}
  一个函子 $F:\uset^{A\op}$ 被称为 \define{可表的（representable）}
  \indexdef{functor!representable}%
  \indexdef{representable functor}%
  如果存在 $a:A$ 以及一个同构 $\y a \cong F$。
\end{defn}

\begin{thm}\label{ct:representable-prop}
  如果 $A$ 是一个范畴，那么"$F$ 是可表的"这一类型是一个命题。
\end{thm}

\begin{proof}
  根据定义，"$F$ 是可表的"恰好是 $\y_0$ 在 $F$ 上的纤维。
  由于 $\y_0$ 根据 \cref{ct:yoneda-mono} 是一个嵌入，这个纤维是一个命题。
\end{proof}

特别地，在一个范畴中，同一函子的任意两个表示是相等的。
我们可以用这一点来给出 \cref{ct:adjprop} 的另一个证明。
首先，我们用可表性来刻画伴随。

\begin{lem}\label{ct:adj-repr}
  对任意预范畴 $A$ 和 $B$ 以及一个函子 $F:A\to B$，以下类型是等价的。
  \begin{enumerate}
  \item $F$ 是一个左伴随\index{adjoint!functor}。\label{item:ct:ar1}
  \item 对每个 $b:B$，从 $A\op$ 到 \uset 的函子 $(a \mapsto \hom_B(Fa,b))$ 是可表的\index{representable functor}。\label{item:ct:ar2}
  \end{enumerate}
\end{lem}

\begin{proof}
  类型~\ref{item:ct:ar2} 的一个元素由一个函数 $G_0:B_0 \to A_0$ 以及一族同构
  \[ \gamma_{a,b}:\hom_B(Fa,b) \cong \hom_A(a,G_0 b) \]
  组成，使得对于 $f:\hom_{A}(a,a')$ 有 $\gamma_{a,b}(g \circ Ff) = \gamma_{a',b}(g)\circ f$。

  在此基础上，对 $a:A$ 我们定义 $\eta_a \defeq \gamma_{a,Fa}(1_{Fa})$，对 $b:B$ 定义 $\epsilon_b \defeq \inv{(\gamma_{Gb,b})}(1_{Gb})$。
  现在对于 $g:\hom_B(b,b')$，我们定义
  \[ G_{b,b'}(g) \defeq \gamma_{G b, b'}(g \circ \epsilon_b) \]
  验证 $G$ 是函子、$\eta$ 和 $\epsilon$ 是满足三角恒等式的自然变换，与经典情形完全相同；又因为它们都是命题，我们并不关心它们的具体值。
  于是我们得到一个函数~\ref{item:ct:ar2}$\to$\ref{item:ct:ar1}。

  反之，若 $F$ 是一个左伴随，我们当然已经指定了 $G_0$，且可将 $\gamma_{a,b}$ 取为复合
  \[ \hom_B(Fa,b)
  \xrightarrow{G_{Fa,b}} \hom_A(GFa,Gb)
  \xrightarrow{(\blank\circ \eta_a)} \hom_A(a,Gb).
  \]
  由于 $\eta$ 是自然的，这显然也是自然的；并且它有一个逆，由
  \[ \hom_A(a,Gb)
  \xrightarrow{F_{a,Gb}} \hom_B(Fa,FGb)
  \xrightarrow{(\epsilon_b \circ \blank )} \hom_A(Fa,b)
  \]
  给出（利用三角恒等式）。
  因此我们又有~\ref{item:ct:ar1}$\to$~\ref{item:ct:ar2}。

  对于复合~\ref{item:ct:ar2}$\to$\ref{item:ct:ar1}$\to$~\ref{item:ct:ar2}，函数 $G_0$ 显然保持不变，故只需验证 $\gamma$ 得以恢复。
  但新的 $\gamma$ 将 $f:\hom_B(Fa,b)$ 映到
  \begin{align*}
    G(f) \circ \eta_a
    &\jdeq \gamma_{G Fa, b}(f \circ \epsilon_{Fa}) \circ \eta_a\\
    &= \gamma_{G Fa, b}(f \circ \epsilon_{Fa} \circ F\eta_a)\\
    &= \gamma_{G Fa, b}(f)
  \end{align*}
  所以它与旧的 $\gamma$ 一致。

  最后，对于~\ref{item:ct:ar1}$\to$\ref{item:ct:ar2}$\to$~\ref{item:ct:ar1}，在对象层面我们确实恢复了函子 $G$。
  新的 $G_{b,b'}:\hom_B(b,b') \to \hom_A(Gb,Gb')$ 将 $g$ 映到
  \begin{align*}
    \gamma_{G b, b'}(g \circ \epsilon_b)
    &\jdeq G(g \circ \epsilon_b) \circ \eta_{Gb}\\
    &= G(g) \circ G\epsilon_b \circ \eta_{Gb}\\
    &= G(g)
  \end{align*}
  故与旧的一致。
  新的 $\eta_a$ 定义为 $\gamma_{a,Fa}(1_{Fa}) \jdeq G(1_{Fa}) \circ \eta_a$，因此等于旧的 $\eta_a$。
  而新的 $\epsilon_b$ 定义为 $\inv{(\gamma_{Gb,b})}(1_{Gb}) \jdeq \epsilon_b \circ F(1_{Gb})$，同样等于旧的 $\epsilon_b$。
\end{proof}

\begin{cor}\label{ct:adjprop2}[\cref{ct:adjprop}]
  若 $A$ 是一个范畴且 $F:A\to B$，则"$F$ 是左伴随"这一类型是一个命题。
\end{cor}

\begin{proof}
  由 \cref{ct:representable-prop}，若 $A$ 是范畴，则 \cref{ct:adj-repr}\ref{item:ct:ar2} 中的类型是一个命题。
\end{proof}


\index{Yoneda!lemma|)}

\section{严格范畴}
\label{sec:strict-categories}

\index{bargaining|(}%

\begin{defn}\label{ct:strict-category}
  一个\define{严格范畴}
  \indexdef{category!strict}%
  \indexdef{strict!category}%
  是其对象类型为集合的预范畴。
\end{defn}

按照数学红鲱鱼原则\index{red herring principle}，严格范畴未必是范畴。
实际上，一个范畴是严格范畴当且仅当它是 gaunt 范畴（\cref{ct:gaunt}）。
\index{gaunt category}%
\index{category!gaunt}%
大多数时候，范畴论讨论的是范畴而非严格范畴，但有时人们也需要考虑严格范畴。
其主要优势在于，严格范畴具有比等价更严格的"相同"概念，即同构（或由 \cref{ct:cat-eq-iso} 等价地，为相等）。

以下是严格范畴的一个来源。

\begin{eg}\label{ct:mono-cat}
  设 $A$ 是一个预范畴，$x:A$ 是一个对象。
  我们如下定义一个预范畴 $\mathsf{mono}(A,x)$：
  \index{monomorphism}
  \indexsee{mono}{monomorphism}
  \indexsee{monic}{monomorphism}
  \begin{itemize}
  \item 它的对象由 $A$ 中的一个对象 $y$ 以及一个单态射 $m:\hom_A(y,x)$ 组成。
    （照例，若 $(m\circ f = m\circ g) \Rightarrow (f=g)$，则称 $m:\hom_A(y,x)$ 是一个\define{单态射}（或说它是\define{单的}）。）
  \item 从 $(y,m)$ 到 $(z,n)$ 的态射是 $A$ 中从 $y$ 到 $z$ 的任意态射（不必保持 $m$ 和 $n$）。
  \end{itemize}
  $\mathsf{mono}(A,x)$ 中对象的一个相等 $(y,m)=(z,n)$ 由一个等式 $p:y=z$ 和一个等式 $\trans{p}{m}=n$ 组成，而根据 \cref{ct:idtoiso-trans}，这等价于一个等式 $m=n\circ \idtoiso(p)$。
  由于 hom-集是集合，这类等式的类型是一个命题。
  但由于 $m$ 和 $n$ 是单态射，满足 $m = n\circ f$ 的态射 $f$ 的类型也是一个命题。
  因此，如果 $A$ 是一个范畴，那么 $(y,m)=(z,n)$ 是一个命题，从而 $\mathsf{mono}(A,x)$ 是一个严格范畴。
\end{eg}

这个例子可以对偶化，并以各种方式推广。以下是严格范畴的一个有趣应用。

\begin{eg}\label{ct:galois}
  设 $E/F$ 是域的有限Galois 扩张
  \index{Galois!extension}%
  $G$ 是其Galois 群。
  \index{Galois!group}%
  则存在一个严格范畴，其对象是中间域 $F\subseteq K\subseteq E$，态射是逐点固定 $F$ 的域同态\index{homomorphism!field}（但不必与到 $E$ 的包含映射交换）。
  还有一个严格范畴，其对象是子群 $H\subseteq G$，态射是 $G$-集合的态射 $G/H \to G/K$。
  Galois 理论基本定理
  \index{fundamental!theorem of Galois theory}%
  指出这两个预范畴是同构的（而不仅仅是等价）。
\end{eg}

\index{bargaining|)}%

\section{\texorpdfstring{$\dagger$}{†}-范畴}
\label{sec:dagger-categories}

此外，还有一类有用的预范畴值得一提：其对象的类型不是集合，但也不是范畴。

\begin{defn}\label{ct:dagger-precategory}
  一个\define{$\dagger$-预范畴（$\dagger$-precategory）}
  \indexdef{.dagger-precategory@$\dagger$-precategory}%
  \indexdef{precategory!.dagger-@$\dagger$-}%
  是一个预范畴 $A$ 连同以下结构。
  \begin{enumerate}
  \item 对于每对 $x,y:A$，有一个函数 $\dgr{(-)}:\hom_A(x,y) \to \hom_A(y,x)$。
  \item 对所有 $x:A$，有 $\dgr{(1_x)} = 1_x$。
  \item 对所有 $f,g$，有 $\dgr{(g\circ f)} = \dgr f \circ \dgr g$。
  \item 对所有 $f$，有 $\dgr{(\dgr f)} = f$。
  \end{enumerate}
\end{defn}

\begin{defn}\label{ct:unitary}
  $\dagger$-预范畴中的态射 $f:\hom_A(x,y)$ 如果满足 $\dgr f \circ f = 1_x$ 且 $f\circ \dgr f = 1_y$，则称为\define{酉的（unitary）}
  \indexdef{.dagger-precategory@$\dagger$-precategory!unitary morphism in}%
  \indexdef{unitary morphism}%
  \indexdef{morphism!unitary}%
  \indexdef{isomorphism!unitary}%
  。
\end{defn}

显然，每个酉态射都是同构，而且"是酉的"这一性质是一个命题。
因此，对于每对 $x,y:A$，从 $x$ 到 $y$ 的所有酉同构构成一个集合，记为 $(x\unitaryiso y)$。

\begin{lem}\label{ct:idtounitary}
  如果 $p:(x=y)$，那么 $\idtoiso(p)$ 是酉的。
\end{lem}

\begin{proof}
  由归纳法，可假设 $p$ 是 $\refl x$。
  但此时 $\dgr{(1_x)} \circ 1_x = 1_x\circ 1_x = 1_x$ 成立，另一式同理。
\end{proof}

\begin{defn}\label{ct:dagger-category}
  一个\define{$\dagger$-范畴（$\dagger$-category）}
  \indexdef{.dagger-category@$\dagger$-category}%
  是一个 $\dagger$-预范畴，且满足对于所有 $x,y:A$，从 \cref{ct:idtounitary} 出发的函数
  \[ (x=y) \to (x \unitaryiso y) \]
  是一个等价。
\end{defn}

\begin{eg}\label{ct:rel-dagger-cat}
  \cref{ct:rel} 中的范畴 \urel 若定义 $(\dgr R)(y,x) \defeq R(x,y)$，则成为一个 $\dagger$-预范畴。
  证明 \urel 是范畴的过程实际上已表明每个同构都是酉的；因此 \urel 也是 $\dagger$-范畴。
\end{eg}

\begin{eg}\label{ct:groupoid-dagger-cat}
  任何群胚若定义 $\dgr f \defeq \inv{f}$，则成为一个 $\dagger$-范畴。
\end{eg}

\begin{eg}\label{ct:hilb}
  令 \uhilb 为如下预范畴。
  \begin{itemize}
  \item 其对象是带有内积 $\langle \blank,\blank\rangle$ 的\index{finite!-dimensional vector space}有限维\index{vector!space}向量空间。
  \item 其态射是任意的线性映射。
    \index{function!linear}%
    \indexsee{linear map}{function, linear}%
  \end{itemize}
  由标准线性代数，有限维内积空间之间的任意线性映射 $f:V\to W$ 都有唯一确定的\index{adjoint!linear map}伴随 $\dgr f:W\to V$，由条件 $\langle f v,w\rangle = \langle v,\dgr f w\rangle$ 刻画。
  通过这种方式，\uhilb 成为一个 $\dagger$-预范畴。
  此外，一个线性同构是酉的当且仅当它是一个\define{等距映射（isometry）}，
  \indexdef{isometry}%
  即满足 $\langle fv,fw\rangle = \langle v,w\rangle$。
  由此可知，\uhilb 是一个 $\dagger$-范畴，尽管它不是一个范畴（并非每个线性同构都是酉的）。
\end{eg}

在\index{mathematics!classical}\index{classical!category theory}经典数学基础下，$\dagger$-范畴已经发展出了相当多的一般理论。
人们很早就注意到，对于 $\dagger$-范畴的对象而言，正确的"相同性"概念是酉同构而非任意同构；这在范畴学家中曾引起一些困惑。
同伦类型论通过将 $\dagger$-范畴（如同严格范畴一样）视为预范畴的一种不同种类，解决了这个问题。

\section{结构同一性原理}
\label{sec:sip}
 \index{structure!identity principle|(}

\emph{结构同一性原理}是一个非正式的原理，它表达同构的结构是同一的。
我们的目标是证明一个一般的抽象结果，可应用于一大类结构概念；这些结构可以是多类的，甚至是依值类的，可以是无穷元的，甚至是高阶的。

最简单的一类单类结构仅由一个类型构成，没有任何附加结构。
泛等公理以强形式表达了针对这种结构概念的结构同一性原理：对于类型 $A,B$，典范函数 $(A=B)\to (\eqv A B)$ 是一个等价。

我们从一个预范畴 $X$ 出发。
在应用于单类一阶结构时，$X$ 将是 $\bbU$-小\index{set}集合的范畴 %\uset%，其中 $\bbU$ 是一个泛等类型宇宙。

\begin{defn}\label{ct:sig}
  在 $X$ 上的一个\define{结构概念（notion of structure）}
  \indexdef{structure!notion of}%
  $(P,H)$ 由以下数据构成。
  \begin{enumerate}
  \item 一个类型族 $P:X_0 \to \type$。
    对于每个 $x:X_0$，$Px$ 中的元素称为 $x$ 上的\define{$(P,H)$-结构（$(P,H)$-structure）}
    \indexsee{PH-structure@$(P,H)$-structure}{structure}%
    \indexdef{structure!PH@$(P,H)$-}%
    。
  \item 对于 $x,y:X_0$，$f:\hom_X(x,y)$ 以及 $\alpha:Px$ 和 $\beta:Py$，给定一个命题
  \[ H_{\alpha\beta}(f).\]
    若 $H_{\alpha\beta}(f)$ 为真，则称 $f$ 是从 $\alpha$ 到 $\beta$ 的\define{$(P,H)$-同态（$(P,H)$-homomorphism）}
    \indexdef{homomorphism!of structures}%
    \indexdef{structure!homomorphism of}%
    。
  \item 对于 $x:X_0$ 和 $\alpha:Px$，我们有 $H_{\alpha\alpha}(1_x)$。\label{item:sigid}
  \item 对于 $x,y,z:X_0$ 和 $\alpha:Px$、$\beta:Py$、$\gamma:Pz$，
如果 $f:\hom_X(x,y)$ 且 $g:\hom_X(y,z)$，我们有\label{item:sigcmp}
  \[ H_{\alpha\beta}(f)\to H_{\beta\gamma}(g)\to H_{\alpha\gamma}(g\circ   f).\]
   \end{enumerate}
  当 $(P,H)$ 是一个结构概念时，对于 $\alpha,\beta:Px$，我们定义
  \[ (\alpha\leq_x\beta) \defeq H_{\alpha\beta}(1_x).\]
  根据条件~\ref{item:sigid} 和~\ref{item:sigcmp}，这是一个以 $Px$ 为对象类型的预序（\cref{ct:orders}）。
  如果对于所有 $x:X$，这个预序事实上是一个偏序，则称 $(P,H)$ 是一个\define{标准结构概念（standard notion of structure）}
  \indexdef{structure!standard notion of}%
  。
\end{defn}

注意，对于一个标准结构概念，每个类型 $Px$ 实际上必须是一个集合。
现在，对于任意结构概念 $(P,H)$，我们定义\define{$(P,H)$-结构的预范畴（precategory of $(P,H)$-structures）}，
\indexdef{precategory!of PH-structures@of $(P,H)$-structures}%
\indexdef{structure!precategory of PH@precategory of $(P,H)$-}%
$A = \mathsf{Str}_{(P,H)}(X)$
。

\begin{itemize}
\item $A$ 的对象类型是 $A_0 \defeq \sm{x:X_0} Px$。
  若 $a\jdeq (x,\alpha):A_0$，我们可记 $|a| \defeq x$。
\item 对于 $(x,\alpha):A_0$ 和 $(y,\beta):A_0$，我们定义
  \[\hom_A((x,\alpha),(y,\beta)) \defeq \setof{ f:x \to y | H_{\alpha\beta}(f)}.\]
\end{itemize}

其中的复合与恒同态射继承自 $X$；条件~\ref{item:sigid} 与~\ref{item:sigcmp} 保证了这些运算可以提升到 $A$ 中。

\begin{thm}[结构同一性原理]\label{thm:sip}
  \indexdef{structure!identity principle}%
  若 $X$ 是一个范畴，且 $(P,H)$ 是 $X$ 上的一个标准结构概念，则预范畴 $\mathsf{Str}_{(P,H)}(X)$ 也是一个范畴。
\end{thm}

\begin{proof}
  根据依值对类型中相等的定义，给出一个相等 $(x,\alpha)=(y,\beta)$ 相当于给出以下数据：
  \begin{itemize}
  \item 一个相等 $p:x=y$，以及
  \item 一个相等 $\trans{p}{\alpha}=\beta$。
  \end{itemize}
  由于 $P$ 取集合值，后者是一个命题。
另一方面，容易看出，$\mathsf{Str}_{(P,H)}(X)$ 中的一个同构 $(x,\alpha)\cong (y,\beta)$ 由以下数据构成：
  \begin{itemize}
  \item $X$ 中的一个同构 $f:x\cong y$，使得
  \item $H_{\alpha\beta}(f)$ 与 $H_{\beta\alpha}(\inv f)$ 成立。
  \end{itemize}
  当然，后者也是一个命题。
由于 $X$ 是一个范畴，函数 $(x=y) \to (x\cong y)$ 是一个等价。
因此，只需证明对于任意 $p:x=y$ 以及任意 $(\alpha:Px)$ 和 $(\beta:Py)$，$\trans{p}{\alpha}=\beta$ 当且仅当 $H_{\alpha\beta}(\idtoiso (p))$ 与 $H_{\beta\alpha}(\inv{\idtoiso(p)})$ 同时成立。

  “仅当”方向只需用到范畴 $\mathsf{Str}_{(P,H)}(X)$ 的 $\idtoiso$ 函数的存在性。
对于“当”方向，通过对 $p$ 进行归纳，可以假设 $y\jdeq x$ 且 $p\jdeq\refl x$。
不过此时 $\idtoiso (p)\jdeq 1_x$，因此 $\inv{\idtoiso(p)}=1_x$。
于是 $\alpha\leq_x \beta$ 且 $\beta\leq_x \alpha$，又因 $(P,H)$ 是标准结构概念，故 $\alpha=\beta$。
\end{proof}

作为一个例子，这套方法为 \cref{ct:functor-cat} 的证明提供了另一种表述方式。

\begin{eg}\label{ct:sip-functor-cat}
  设 $A$ 为预范畴，$B$ 为范畴。
  存在一个预范畴 $B^{A_0}$，其对象是函数 $A_0 \to B_0$，而从 $F_0:A_0 \to B_0$ 到 $G_0:A_0 \to B_0$ 的态射集是 $\prd{a:A_0} \hom_B(F_0 a, G_0 a)$。
  复合与恒同态射直接继承自 $B$。
  不难证明，$\gamma:\hom_{B^{A_0}}(F_0, G_0)$ 是同构当且仅当每个分量 $\gamma_a$ 都是同构，因此我们有 $\eqv{(F_0 \cong G_0)}{\prd{a:A_0} (F_0 a \cong G_0 a)}$。
  进而，$B^{A_0}$ 中的映射 $\idtoiso : (F_0 = G_0) \to (F_0 \cong G_0)$ 等于复合
  \[ (F_0 = G_0) \longrightarrow \prd{a:A_0} (F_0 a  = G_0 a) \longrightarrow \prd{a:A_0} (F_0 a \cong G_0 a) \longrightarrow (F_0 \cong G_0) \]
  其中第一个映射由函数外延性可知是等价，第二个由于是等价的依值积（因为 $B$ 是范畴）故亦为等价，第三个如前所述。
  因此，$B^{A_0}$ 是一个范畴。

  现在我们在 $B^{A_0}$ 上定义一个结构概念，使得 $P(F_0)$ 是将 $F_0$ 拓展为函子（即保持复合与恒同态射）的运算 $F:\prd{a,a':A_0} \hom_A(a,a') \to \hom_B(F_0 a,F_0 a')$ 的类型。
  由于每个 $\hom_B(\blank,\blank)$ 都是集合，故 $P(F_0)$ 也是一个集合。
  给定这样的 $F$ 和 $G$，我们定义 $\gamma:\hom_{B^{A_0}}(F_0, G_0)$ 为同态，如果它构成一个自然变换。\index{natural!transformation}
  在 \cref{ct:functor-precat} 中，我们本质上已经验证了这是一个结构概念。
  此外，如果 $F$ 和 $F'$ 都是 $F_0$ 上的结构，且恒同态射是从 $F$ 到 $F'$ 的自然变换，那么对于任何 $f:\hom_A(a,a')$，我们有 $F'f = F'f \circ 1_{F_0 a} = 1_{F_0 a}\circ F f = F f$。
  应用函数外延性，可得 $F = F'$。
  因此，我们得到了一个\emph{标准}的结构概念，从而由 \cref{thm:sip} 可知预范畴 $B^A$ 是一个范畴。
\end{eg}

作为另一个例子，我们考虑一阶签名结构的范畴。
我们定义\define{一阶签名（first-order signature）}，
\indexdef{first-order!signature}%
\indexdef{signature!first-order}%
$\Omega$，它由函数符号的集合 $\Omega_0$ 和关系符号的集合 $\Omega_1$ 组成，其中每个符号 $\omega$ 的元\index{arity} $|\omega|$ 是一个集合。
一个\define{$\Omega$-结构（$\Omega$-structure）}
\indexdef{structure!Omega@$\Omega$-}%
\indexsee{omega-structure@$\Omega$-structure}{structure}%
$a$ 由一个集合 $|a|$ 以及如下指派组成：对每个函数符号 $\omega$，指派一个 $|a|$ 上的 $|\omega|$ 元函数 $\omega^a:|a|^{|\omega|}\to |a|$；对每个关系符号 $\omega$，指派一个 $|a|$ 上的 $|\omega|$ 元关系 $\omega^a$，该关系对每个 $x:|a|^{|\omega|}$ 赋予一个命题。
给定 $\Omega$-结构 $a, b$，一个函数 $f: |a| \to |b|$ 称为\define{同态 $a\to b$（homomorphism $a\to b$）}
\indexdef{homomorphism!of Omega-structures@of $\Omega$-structures}%
\indexdef{structure!homomorphism of Omega@homomorphism of $\Omega$-}%
，如果它保持结构；即对于签名中的每个符号 $\omega$ 和每个 $x:|a|^{|\omega|}$，满足：

\begin{enumerate}
\item $f(\omega^ax) = \omega^b(f\circ x)$ 若 $\omega:\Omega_0$，且
\item $\omega^ax\to\omega^b(f\circ x)$ 若 $\omega:\Omega_1$。
\end{enumerate}

注意，每个 $x:|a|^{|\omega|}$ 是一个函数 $x:|\omega|\to |a|$，因此 $f\circ x : b^\omega$。

现在假设给定一个（泛等的）宇宙 $\bbU$ 和一个 $\bbU$-小签名 $\Omega$；即 $|\Omega|$ 是一个 $\bbU$-小集合，且对每个 $\omega:|\Omega|$，集合 $|\omega|$ 是 $\bbU$-小的。
于是我们有 $\bbU$-小集合的范畴 $\uset_\bbU$。我们想要定义 $\uset_\bbU$ 上的 $\bbU$-小 $\Omega$-结构的预范畴，并利用 \cref{thm:sip} 证明它是一个范畴。

我们用一阶签名 $\Omega$ 来给出一个 $\uset_\bbU$ 上的标准结构概念 $(P,H)$。

\begin{defn}\label{defn:fo-notion-of-structure}
\mbox{}
\begin{enumerate}
\item 对每个 $\bbU$-小集合 $x$，定义
  \[ Px \defeq P_0x\times P_1x.\]
  这里
  %
  \begin{align*}
    P_0x &\defeq \prd{\omega:\Omega_0} x^{|\omega|}\to x, \mbox{ and } \\
    P_1x &\defeq \prd{\omega:\Omega_1} x^{|\omega|}\to \propU,
  \end{align*}
\item 对 $\bbU$-小集合 $x,y$ 以及
  $\alpha:P^\omega x,\;\beta:P^\omega y,\; f:x\to y$，定义
  \[ H_{\alpha\beta}(f) \defeq H_{0,\alpha\beta}(f)\wedge H_{1,\alpha\beta}(f).\]
  这里
  \begin{align*}
    H_{0,\alpha\beta}(f) &\defeq
    \fall{\omega:\Omega_0}{u:x^{|\omega|}} f(\alpha u)=\;\beta(f\circ u),
    \mbox{ and }\\
    H_{1,\alpha\beta}(f) &\defeq
    \fall{\omega:\Omega_1}{u:x^{|\omega|}} \alpha u\to\beta(f\circ u).
  \end{align*}
\end{enumerate}
\end{defn}

现在例行验证 $(P,H)$ 是 $\uset_\bbU$ 上的标准结构概念，从而可由 \cref{thm:sip} 得出预范畴 $Str_{(P,H)}(\uset_\bbU)$ 是一个范畴。
最后只需注意，这本质上与 $\uset_\bbU$ 上的 $\bbU$-小 $\Omega$-结构的预范畴相同。
 \index{structure!identity principle|)}

\section{Rezk 完备化}
\label{sec:rezk}

在本节中，我们将给出一种将预范畴替换为范畴的通用方法。
事实上，我们将给出两种。
二者都依赖于以下事实："范畴将弱等价视为等价"。

为了证明这一点，我们先证明几个引理，它们属于完全标准的范畴论内容，但表述时格外谨慎，以确保正确使用了 $\truncf{-1}$ 的消去子。
在经典\index{mathematics!classical}\index{classical!category theory}范畴论中，如果想避免选择公理，也同样需要这样的谨慎：每当我们想定义一个函数时，都需要以某种方式唯一地刻画其值。

\begin{lem}\label{ct:esosurj-postcomp-faithful}
  若 $A,B,C$ 是预范畴，且 $H:A\to B$ 是本质满的函子，则 $(\blank\circ H):C^B \to C^A$ 是忠实的。
\end{lem}

\begin{proof}
  设 $F,G:B\to C$，$\gamma,\delta:F\to G$ 满足 $\gamma H = \delta H$；我们需要证明 $\gamma=\delta$。
  为此，取 $b:B$；我们要证明 $\gamma_b=\delta_b$。
  这是一个命题，而 $H$ 是本质满的，因此我们可以假设给定 $a:A$ 和同构 $f:Ha\cong b$。
  但现在我们有
  \[ \gamma_b = G(f) \circ \gamma_{Ha} \circ F(\inv{f})
  = G(f) \circ \delta_{Ha} \circ F(\inv{f})
  = \delta_b.\qedhere
  \]
\end{proof}

\begin{lem}\label{ct:esofull-precomp-ff}
  若 $A,B,C$ 是预范畴，且 $H:A\to B$ 本质满且全，则 $(\blank\circ H):C^B \to C^A$ 是全忠实的。
\end{lem}

\begin{proof}
  只需再证它是全的。
  为此，取 $F,G:B\to C$ 和 $\gamma:FH \to GH$。
  我们断言，对任意 $b:B$，类型
  \begin{equation}\label{eq:fullprop}
    \sm{g:\hom_C(Fb,Gb)} \prd{a:A}{f:Ha\cong b} (\gamma_a =  \inv{Gf}\circ g\circ Ff)
  \end{equation}
  是可缩的。
  可缩性是命题，而 $H$ 是本质满的，故可以假设给定 $a_0:A$ 和 $h:Ha_0\cong b$。

  现取 $g\defeq Gh \circ \gamma_{a_0} \circ \inv{Fh}$。
  给定任意其他的 $a:A$ 和 $f:Ha\cong b$，我们需要证明 $\gamma_a =  \inv{Gf}\circ g\circ Ff$。
  由于 $H$ 是全的， merely 存在一个态射 $k:\hom_A(a,a_0)$ 使得 $Hk = \inv{h}\circ f$。
  又因我们的目标是命题，可以假设给定这样的 $k$。
  于是有
  \begin{align*}
    \gamma_a &= \inv{GHk}\circ \gamma_{a_0} \circ FHk\\
    &= \inv{Gf} \circ Gh \circ \gamma_{a_0} \circ \inv{Fh} \circ Ff\\
    &= \inv{Gf}\circ g\circ Ff.
  \end{align*}
  因此，~\eqref{eq:fullprop} 是有实例的。
  还需证明它是命题。
  取 $g,g':\hom_C(Fb, Gb)$ 满足：对所有 $a:A$ 和 $f:Ha\cong b$，同时有 $(\gamma_a =  \inv{Gf}\circ g\circ Ff)$ 和 $(\gamma_a =  \inv{Gf}\circ g'\circ Ff)$。
  依值积类型是命题，故只需证明 $g=g'$。
  而这是命题，因此可以假设 $a_0:A$ 和 $h:Ha_0\cong b$，此时有
  \[ g = Gh \circ \gamma_{a_0} \circ \inv{Fh} = g'.\]
  %
  这证明了~\eqref{eq:fullprop} 对所有 $b:B$ 可缩。
  现在定义 $\delta:F\to G$：对每个 $b$，取 $\delta_b$ 为~\eqref{eq:fullprop} 中那个唯一的 $g$。
  为验证其自然性，假设给定 $f:\hom_B(b,b')$；我们需要证明 $Gf \circ \delta_b = \delta_{b'}\circ Ff$。
  如前，可以假设 $a:A$ 和 $h:Ha\cong b$，同样假设 $a':A$ 和 $h':Ha'\cong b'$。
  由于 $H$ 既本质满又全，还可假设 $k:\hom_A(a,a')$ 满足 $Hk = \inv{h'}\circ f\circ h$。

  由 $\gamma$ 的自然性，$GHk\circ \gamma_a = \gamma_{a'} \circ FHk$。
  利用 $\delta$ 的定义，得
  \begin{align*}
    Gf \circ \delta_b
    &= Gf \circ Gh \circ \gamma_a \circ \inv{Fh}\\
    &= Gh' \circ GHk\circ \gamma_a \circ \inv{Fh}\\
    &= Gh' \circ \gamma_{a'} \circ FHk \circ \inv{Fh}\\
    &= Gh' \circ \gamma_{a'} \circ \inv{Fh'} \circ Ff\\
    &= \delta_{b'} \circ Ff.
  \end{align*}
  故 $\delta$ 是自然的。
  最后，对任意 $a:A$，将 $\delta_{Ha}$ 的定义应用于 $a$ 和 $1_a$，即得 $\gamma_a = \delta_{Ha}$。
  因此 $\delta \circ H = \gamma$。
\end{proof}

定理余下部分的证明思路几乎完全相同，只在关键一步插入 $C$ 作为\emph{范畴}的条件，我们在下文用斜体标出以示强调。
正是在这一步，我们试图在不使用选择公理的情况下定义一个以\emph{对象}为值的函数，因此必须谨慎对待"唯一指定"一个对象意味着什么。
在经典\index{mathematics!classical}\index{classical!category theory}范畴论中，我们只能说这个对象在相差一个唯一同构的意义下被指定，但在集合论基础中，这种程度的唯一性不足以在不援引 \choice{} 的情况下给出一个函数。
然而在泛等基础中，如果 $C$ 是范畴，则同构即相等，我们便有了所需的那种唯一性（即处于可缩空间之中）。

\index{weak equivalence!of precategories|(}%

\begin{thm}\label{ct:cat-weq-eq}
  若 $A,B$ 是预范畴，$C$ 是范畴，且 $H:A\to B$ 是弱等价，则 $(\blank\circ H):C^B \to C^A$ 是同构。
\end{thm}

\begin{proof}
  由 \cref{ct:functor-cat}，$C^B$ 和 $C^A$ 都是范畴。
  故由 \cref{ct:eqv-levelwise}，只需证明 $(\blank\circ H)$ 是等价。
  而由前两个引理已知它是全忠实的，故根据 \cref{ct:catweq}，只需证明它是本质满的。
  为此，假设给定 $F:A\to C$；我们要 merely 存在一个 $G:B\to C$ 使得 $GH\cong F$。

  对每个 $b:B$，令 $X_b$ 为由以下数据组成的类型：
  \begin{enumerate}
  \item 一个元素 $c:C$；以及
  \item 对每个 $a:A$ 和 $h:Ha\cong b$，一个同构 $k_{a,h}:Fa\cong c$；使得\label{item:eqvprop2}
  \item 对~\ref{item:eqvprop2} 中那样的每对 $(a,h)$ 和 $(a',h')$，以及每个满足 $h'\circ Hf = h$ 的 $f:\hom_A(a,a')$，有 $k_{a',h'}\circ Ff = k_{a,h}$。\label{item:eqvprop3}
  \end{enumerate}
  我们断言：对任意 $b:B$，类型 $X_b$ 是可缩的。
  由于这是命题，可以假设给定 $a_0:A$ 和 $h_0:Ha_0 \cong b$。
  令 $c^0\defeq Fa_0$。
  接下来，给定 $a:A$ 和 $h:Ha\cong b$，由于 $H$ 是全忠实的，存在唯一的同构 $g_{a,h}:a\to a_0$ 满足 $Hg_{a,h} = \inv{h_0}\circ h$；定义 $k^0_{a,h} \defeq Fg_{a,h}$。
  最后，若 $h'\circ Hf = h$，则 $\inv{h_0}\circ h'\circ Hf = \inv{h_0}\circ h$，从而 $g_{a',h'} \circ f = g_{a,h}$，故 $k^0_{a',h'}\circ Ff = k^0_{a,h}$。
  因此 $X_b$ 是有实例的。

  现假设给定另一个 $(c^1,k^1): X_b$。
  则 $k^1_{a_0,h_0}:c^0 \jdeq Fa_0 \cong c^1$。
  \emph{由于 $C$ 是范畴，我们有 $p:c^0=c^1$ 及 $\idtoiso(p) = k^1_{a_0,h_0}$。}
  而对任意 $a:A$ 和 $h:Ha\cong b$，对 $(c^1,k^1)$ 取 $f\defeq g_{a,h}$ 应用~\ref{item:eqvprop3}，得
  \[k^1_{a,h} = k^1_{a_0,h_0} \circ k^0_{a,h} = \trans{p}{k^0_{a,h}}\]
  这就给出了构造等式 $(c^0,k^0)=(c^1,k^1)$ 所需的数据，从而完成了 $X_b$ 可缩性的证明。

  由于每个 $X_b$ 都可缩，类型 $\prd{b:B} X_b$ 也可缩。
  特别地，它是有实例的，故我们有一个函数，将一个 $c$ 和一个 $k$ 分配给每个 $b:B$。
  定义 $G_0(b)$ 为此 $c$；这便给出函数 $G_0 :B_0 \to C_0$。

  接下来需要定义 $G$ 在态射上的作用。
  对每个 $b,b':B$ 和 $f:\hom_B(b,b')$，令 $Y_f$ 为由以下数据组成的类型：
  \begin{enumerate}[resume]
  \item 一个态射 $g:\hom_C(Gb,Gb')$，使得
  \item 对每个 $a:A$ 和 $h:Ha\cong b$、每个 $a':A$ 和 $h':Ha'\cong b'$，以及任意 $\ell:\hom_A(a,a')$，有\label{item:eqvprop5}
    \[ (h' \circ H\ell = f \circ h)
    \to
    (k_{a',h'} \circ F\ell = g\circ k_{a,h}). \]
  \end{enumerate}
  我们断言：对任意 $b,b'$ 和 $f$，类型 $Y_f$ 是可缩的。
  由于这是命题，可以假设给定 $a_0:A$ 和 $h_0:Ha_0\cong b$，以及 $a'_0:A$ 和 $h'_0:Ha'_0\cong b'$。
  由于 $H$ 是全忠实的，存在唯一的 $\ell_0:\hom_A(a_0,a_0')$ 满足 $h'_0 \circ H\ell_0 = f \circ h_0$。
  定义 $g_0 \defeq k_{a_0',h_0'} \circ F \ell_0 \circ \inv{(k_{a_0,h_0})}$。

  现对任意满足 $(h' \circ H\ell = f \circ h)$ 的 $a,h,a',h',\ell$，有 $\inv{h}\circ h_0:Ha_0\cong Ha$，故存在唯一的 $m:a_0\cong a$ 满足 $Hm = \inv{h}\circ h_0$，从而 $h\circ Hm = h_0$。
  类似地，存在唯一的 $m':a_0'\cong a'$ 满足 $h'\circ Hm' = h_0'$。
  现由~\ref{item:eqvprop3}，得 $k_{a,h}\circ Fm = k_{a_0,h_0}$ 和 $k_{a',h'}\circ Fm' = k_{a_0',h_0'}$。
  又有
  \begin{align*}
    Hm' \circ H\ell_0
    &= \inv{(h')} \circ h_0' \circ H\ell_0\\
    &= \inv{(h')} \circ f \circ h_0\\
    &= \inv{(h')} \circ f \circ h \circ \inv{h} \circ h_0\\
    &= H\ell \circ Hm
  \end{align*}
  由于 $H$ 全忠实，故 $m'\circ \ell_0 = \ell\circ m$。
  最后，计算得
  \begin{align*}
    g_0 \circ k_{a,h}
    &= k_{a_0',h_0'} \circ F \ell_0 \circ \inv{(k_{a_0,h_0})} \circ k_{a,h}\\
    &= k_{a_0',h_0'} \circ F \ell_0 \circ \inv{Fm}\\
    &= k_{a_0',h_0'} \circ \inv{(Fm')} \circ F\ell\\
    &= k_{a',h'}\circ F\ell.
  \end{align*}
  至此证明了 $Y_f$ 是有实例的。
  为证其可缩，由于态射集是集合，只需取另一个满足~\ref{item:eqvprop5} 的 $g_1:\hom_C(Gb,Gb')$ 并证明 $g_0=g_1$。
  但我们手中仍有先前指定的 $a_0,h_0,a_0',h_0',\ell_0$，而~\ref{item:eqvprop5} 蕴含 $g_0$ 和 $g_1$ 都必须等于 $k_{a_0',h_0'} \circ F \ell_0 \circ \inv{(k_{a_0,h_0})}$。

  这就完成了对每个 $b,b':B$ 和 $f:\hom_B(b,b')$ 证明 $Y_f$ 可缩的工作。
  因此，存在一个函数，将唯一的实例分配给每个这样的 $f$；记此函数为 $G_{b,b'}:\hom_B(b,b') \to \hom_C(Gb,Gb')$。
  验证 $G$ 是函子的过程是直接的；在每种情况下，选择 $a,h$ 并应用~\ref{item:eqvprop5} 即可。

  最后，对任意 $a_0:A$，定义 $c\defeq Fa_0$ 和 $k_{a,h}\defeq F g$（其中 $g:\hom_A(a,a_0)$ 是满足 $Hg = h$ 的唯一同构），即给出 $X_{Ha_0}$ 的一个元素。
  因此，它等于先前指定的元素；故 $GHa=Fa$。
  类似地，对 $f:\hom_A(a_0,a_0')$，我们可以沿这些等式做传输来定义 $Y_{Hf}$ 的一个元素，它因此必须等于先前指定的元素。
  故 $GH=F$，进而 $GH\cong F$，如所需。
\end{proof}

\index{universal!property!of Rezk completion}%
因此，若预范畴 $A$ 有一个到范畴的弱等价函子 $A\to \widehat{A}$，则后者就是 $A$ 在范畴中的"反射"：任何从 $A$ 到范畴的函子，都将本质上唯一地通过 $\widehat{A}$ 分解。
现在我们给出两种这样的弱等价的构造。

\indexsee{Rezk 完备化}{completion, Rezk}%
\index{完备化!Rezk|(defstyle}%

\begin{thm}\label{thm:rezk-completion}
  对任何预范畴 $A$，存在一个范畴 $\widehat A$ 和一个弱等价 $A\to\widehat{A}$。
\end{thm}

\begin{proof}[证明一]
  令 $\widehat{A}_0 \defeq \setof{ F:\uset^{A\op} | \exis{a:A} (\y a \cong F)}$，其态射集继承自 $\uset^{A\op}$。
  则包含映射 $\widehat{A} \to \uset^{A\op}$ 是全忠实的，且在对象上是一个嵌入。
  由于 $\uset^{A\op}$ 是一个范畴（根据 \cref{ct:functor-cat}，且由泛等性可知 \uset 也是范畴），$\widehat A$ 也是一个范畴。

  令 $A\to\widehat A$ 为米田嵌入。
  由 \cref{ct:yoneda-embedding} 知它是全忠实的，而由 $\widehat{A}_0$ 的定义知它是本质满的。
  因此它是一个弱等价。
\end{proof}

这个证明十分巧妙，但存在一个缺点：它会提升宇宙层级。
如果 $A$ 是宇宙 \bbU 中的范畴，则在此证明中 \uset 必须至少与 $\uset_\bbU$ 一样大。
于是 $\uset_\bbU$ 和 $(\uset_\bbU)^{A\op}$ 本身不再是 \bbU 中的范畴，而只存在于更高的宇宙中；并且 \emph{先验地}，$\widehat A$ 亦是如此。
可以设想用一条重置公理来处理这个问题，但同样可以用高阶归纳类型给出一个直接构造。

\begin{proof}[第二种证明]
  我们定义一个高阶归纳类型 $\widehat A_0$，其构造子如下：
  \begin{itemize}
  \item 一个函数 $i:A_0 \to \widehat A_0$。
  \item 对每个 $a,b:A$ 和 $e:a\cong b$，一个相等 $je:\id{ia}{ib}$。
  \item 对每个 $a:A$，一个相等 $\id{j(1_a)}{\refl{ia}}$。
  \item 对每个 $(a,b,c:A)$、$(f:a\cong b)$ 和 $(g:b\cong c)$，一个相等 $\id{j(g \circ f)}{j(f)\ct j(g)}$。
  \item 1-截断：对所有 $x,y:\widehat A_0$ 和 $p,q:\id x y$ 以及 $r,s:\id p q$，一个相等 $\id r s$。
  \end{itemize}
  注意，对任意 $a,b:A$ 和 $p:\id a b$，我们有 $\id{j(\idtoiso(p))}{\map i p}$。
  这可以通过对 $p$ 作道路归纳并利用第三个构造子得出。

  类型 $\widehat A_0$ 将作为 $\widehat A$ 的对象类型；现在我们来构建其余的所有结构。
  （下面的证明属于那种能极大受益于计算机证明助手的一类证明：\index{proof!assistant}它广而浅，涉及许多短小的情形需要考虑，而很大一部分工作只是写下需要验证的内容。）

  \mentalpause

  \emph{第一步：}我们通过对 $\widehat A_0$ 的双重归纳来定义一个类型族 $\hom_{\widehat A}:\widehat A_0\to \widehat A_0 \to \set$。
  由于 \set 是一个 1-类型，我们可以忽略 1-截断构造子。
  当 $x$ 和 $y$ 形如 $ia$ 和 $ib$ 时，我们取 $\hom_{\widehat A}(ia,ib) \defeq \hom_A(a,b)$。
  接下来需要考虑所有其他可能的构造子组合。

  首先固定 $x=ia$。
  如果 $y$ 沿相等 $je:\id{ib}{ib'}$ 变化，其中 $e:b\cong b'$，我们需要一个相等 $\id{\hom_A(a,b)}{\hom_A(a,b')}$。
  由泛等性，只需给出一个等价 $\eqv{\hom_A(a,b)}{\hom_A(a,b')}$。
  我们取这个等价为函数 $(e\circ \blank ):\hom_A(a,b)\to \hom_A(a,b')$。
  为说明这是一个等价，我们给出其逆为 $(\inv e\circ \blank )$，而它确为逆的证明来自 $\inv e$ 是 $e$ 在 $A$ 中的逆这一事实。

  当 $y$ 沿相等 $\id{j(1_b)}{\refl{ib}}$ 变化时，我们需要一个相等 $\id{(1_b\circ \blank )}{\refl{\hom_A(a,b)}}$；这由预范畴的单位律 $\id{1_b\circ g}{g}$ 可得。
  类似地，当 $y$ 沿相等 $\id{j(g\circ f)}{j(f)\ct j(g)}$ 变化时，我们需要一个相等 $\id{((g\circ f)\circ \blank )}{(g\circ (f\circ \blank ))}$，这由结合律可得。
  % Finally, as $y$ varies along the 1-truncation constructor, we need only to observe that \set is 1-truncated.

  现在考虑 $x$ 的其他构造子。
  假设 $x$ 沿相等 $j(e):\id{ia}{ia'}$ 变化，其中 $e:a \cong a'$；我们再次需要处理 $y$ 的所有构造子。
  如果 $y$ 是 $ib$，那么我们需要一个相等 $\id{\hom_A(a,b)}{\hom_A(a',b)}$。
  由泛等性，这可由一个等价得出，对此我们可以使用 $(\blank\circ \inv e)$，其逆为 $(\blank\circ e)$。

  仍在 $x$ 沿 $j(e)$ 变化的假设下，现在假设 $y$ 也沿 $j(f)$ 变化，其中 $f:b\cong b'$。
  那么我们需要知道两个拼接的相等
  \begin{gather*}
    \hom_A(a,b) = \hom_A(a',b) = \hom_A(a',b') \mathrlap{\qquad\text{and}}\\
    \hom_A(a,b) = \hom_A(a,b') = \hom_A(a',b')
  \end{gather*}
  是相等的。
  这由结合律可得：$(f\circ \blank)\circ \inv e = f\circ (\blank\circ \inv e)$。
  $y$ 的其他两个构造子是平凡的，因为它们是集合中的二阶相等。

  对于 $x$ 的下两个构造子，除了 $y$ 的第一个构造子外，其余同样是平凡的。
  当 $x$ 沿 $j(1_a)=\refl{ia}$ 变化且 $y$ 是 $ib$ 时，我们再次使用单位律。
  类似地，当 $x$ 沿 $\id{j(g\circ f)}{j(f)\ct j(g)}$ 变化时，我们再次使用结合律。
  这就完成了 $\hom_{\widehat A}:\widehat A_0 \to \widehat A_0 \to \set$ 的构造。

  \mentalpause

  \emph{第二步：}我们通过对 $\widehat A_0$ 的归纳来赋予 $\widehat A$ 预范畴结构。
  % The reader is probably getting bored at this point, so we skip the details.
  我们现在是向集合（即 $\widehat A$ 的 Hom 集）做消去，因此除了前两个构造子外，处理起来都是平凡的。

  对于恒同态射，如果 $x$ 是 $ia$，那么我们有 $\hom_{\widehat A}(x,x) \jdeq \hom_A(a,a)$，我们定义 $1_x \defeq 1_{ia}$。
  如果 $x$ 沿 $je$ 变化，其中 $e:a\cong a'$，我们必须证明 $\transfib{x\mapsto \hom_{\widehat A}(x,x)}{je}{1_{ia}} = 1_{ia'}$。
  但根据 $\hom_{\widehat A}$ 的定义，沿 $je$ 的传输是由与 $e$ 和 $\inv e$ 复合给出的，而我们有 $e\circ 1_{ia} \circ \inv{e} = 1_{ia'}$。

  对于复合，如果 $x,y,z$ 分别是 $ia,ib,ic$，那么 $\hom_{\widehat A}$ 归约为 $\hom_A$，我们可以定义 $\widehat A$ 中的复合为 $A$ 中的复合。
  而当 $x$、$y$ 或 $z$ 沿 $je$ 变化时，我们验证以下等式：
  \begin{align*}
    e \circ (g\circ f) &= (e\circ g) \circ f,\\
    g\circ f &= (g\circ \inv e) \circ (e\circ f),\\
    (g\circ f) \circ \inv e &= g \circ (f\circ \inv e).
  \end{align*}
  最后，结合律和单位律都是命题，因此除了第一个构造子外，所有构造子都是平凡的。
  但在那种情况下，我们有 $A$ 中对应的公理。

  \mentalpause

  \emph{第三步：}我们证明 $\widehat A$ 是一个范畴。
  也就是说，我们必须证明对所有 $x,y:\widehat A$，函数 $\idtoiso:(x=y) \to (x\cong y)$ 是一个等价。
  首先，我们对所有 $x,y:\widehat A$，通过归纳定义一个函数 $k_{x,y}:(x\cong y) \to (x=y)$。
  如前所述，由于我们的目标是一个集合，只需处理前两个构造子即可。

  当 $x$ 和 $y$ 分别是 $ia$ 和 $ib$ 时，我们有 $\hom_{\widehat A}(ia,ib)\jdeq \hom_A(a,b)$，其复合与恒同态射也一并继承，因此 $(ia\cong ib)$ 等价于 $(a\cong b)$。
  但现在我们有了构造子 $j:(a\cong b) \to (ia=ib)$。

  接着，如果 $y$ 沿 $j(e)$ 变化，其中 $e:b\cong b'$，我们必须证明对 $f:a\cong b$ 有 $j(\trans{j(e)}{f}) = j(f) \ct j(e)$。
  但根据 $\hom_{\widehat A}$ 在相等上的定义，沿 $j(e)$ 的传输等价于与 $e$ 后复合，因此这个相等来自 $\widehat A_0$ 的最后一个构造子。
  当 $x$ 沿 $j(e)$ 变化，其中 $e:a\cong a'$ 时，剩余情况类似。
  这就完成了 $k:\prd{x,y:\widehat A_0} (x\cong y) \to (x=y)$ 的定义。

  现在我们需要证明的一点是：如果 $p:x=y$，那么 $k(\idtoiso(p))=p$。
  通过对 $p$ 作归纳，我们可以假定它是 $\refl x$，从而 $\idtoiso(p)\jdeq 1_x$。
  现在我们对 $x:\widehat A_0$ 作归纳，由于我们的目标是一个命题（因为 $\widehat A_0$ 是一个 1-类型），除了第一个构造子外，所有构造子都是平凡的。
  但如果 $x$ 是 $ia$，那么 $k(1_{ia}) \jdeq j(1_a)$，根据 $\widehat A_0$ 的第三个构造子，它等于 $\refl{ia}$。

  为了完成 $\widehat A$ 是范畴的证明，我们必须证明：如果 $f:x\cong y$，那么 $\idtoiso(k(f))=f$。
  通过归纳，我们可以假定 $x$ 和 $y$ 分别是 $ia$ 和 $ib$，在这种情况下 $f$ 必然源自一个同构 $g:a\cong b$，且我们有 $k(f)\jdeq j(g)$。
  然而，对任意 $p$，我们有 $\idtoiso(p) = \trans{p}{1}$，所以特别地 $\idtoiso (j(g)) = \trans{j(g)}{1_{ia}}$。
  而根据 $\hom_{\widehat A}$ 在相等上的定义，这是由 $1_{ia}$ 与等价 $g$ 复合给出的，因此等于 $g$。

  \index{encode-decode method}%
  注意这一步与编码-解码方法\index{encode-decode method}的相似性，后者曾在 \cref{sec:compute-coprod,sec:compute-nat,cha:homotopy} 中使用。
  我们再次通过递归定义一个代码族（这里是 $(x,y)\mapsto (x\cong y)$）并通过对 $\widehat A_0$ 和道路的归纳来定义编码和解码函数，以刻画高阶归纳类型（这里是 $\widehat A_0$）的相等类型。

  \mentalpause

  \emph{第四步：}我们定义一个弱等价 $I:A \to \widehat A$。
  我们取 $I_0 \defeq i : A_0 \to \widehat A_0$，并且根据 $\hom_{\widehat A}$ 的构造，我们有函数 $I_{a,b}:\hom_A(a,b) \to \hom_{\widehat A}(Ia,Ib)$，它们构成了一个函子 $I:A \to \widehat A$。
  由构造可知，这个函子是全忠实的，所以还需证明它是本质满的。
  也就是说，对所有 $x:\widehat A$，我们希望纯存在一个 $a:A$ 使得 $Ia\cong x$。
  一如既往，我们通过对 $x$ 的归纳来论证，并且由于目标是一个命题，除了第一个构造子外，所有构造子都是平凡的。
  但如果 $x$ 是 $ia$，那么我们显然有 $a:A$ 且 $Ia\jdeq ia$，因此 $Ia \cong ia$。
  （注意，如果我们试图证明 $I$ 是\emph{分裂}本质满的，我们会陷入困境，因为我们对 $A_0$ 中的相等毫无所知，从而无法处理任何进一步的构造子。）
\end{proof}

我们将这一构造 $A\mapsto \widehat A$ 称为 \define{Rezk 完备化}，尽管高阶意象语义也提供了将其称为 \define{层栈完备化} 的理由。
\index{.infinity1-topos@$(\infty,1)$-topos}%
\index{stack}%
\index{completion!Rezk|)}%

我们已经看到，实践中出现的大多数预范畴本身就是范畴，因为它们是从 \uset 构造而来的，而根据泛等公理，\uset 本身就是一个范畴。
然而，确实存在少数情况，必须借助 Rezk 完备化才能得到一个范畴。

\begin{eg}\label{ct:rezk-fundgpd-trunc1}
  回忆 \cref{ct:fundgpd}，对任意类型 $X$，都存在一个预群胚，以 $X$ 为对象类型，态射为 $\hom(x,y) \defeq \pizero{x=y}$。
  \indexdef{fundamental!groupoid}%
  \index{fundamental!pregroupoid}%
  \indexsee{groupoid!fundamental}{fundamental group\-oid}%
  它的 Rezk 完备化就是 $X$ 的\emph{基本群胚}。
  回想群胚等价于 1-类型这一事实，不难将此群胚等同于 $\trunc1X$。
\end{eg}

\begin{eg}\label{ct:hocat}
  回忆 \cref{ct:hoprecat}，存在一个预范畴，其对象类型为 \type，态射为 $\hom(X,Y) \defeq \pizero{X\to Y}$。
  它的 Rezk 完备化可以称为\define{类型的同伦范畴}。
  \index{category!of types}%
  \index{homotopy!category of types@(pre)category of types}%
  它的对象类型可以等同于 $\trunc1\type$（参见 \cref{ct:ex:hocat}）。
\end{eg}

Rezk 完备化还使我们能够表明，"范畴"这一概念由"预范畴的弱等价"这一概念所决定。
因此，既然后者（弱等价）是不可避免的，前者（范畴）亦是如此。

\begin{thm}\label{ct:weq-iso-precat-cat}
  一个预范畴 $C$ 是范畴，当且仅当对于每个预范畴的弱等价 $H:A\to B$，诱导出的函子 $(\blank\circ H):C^B \to C^A$ 是预范畴的同构。
\end{thm}

\begin{proof}
  "仅当"方向是 \cref{ct:cat-weq-eq}。
  反之，令 $H$ 为 $I:A\to\widehat A$。
  由于 $(\blank\circ I)_0$ 是等价，存在 $R:\widehat A\to A$ 使得 $RI=1_A$。
  于是 $IRI=I$，但同样由于 $(\blank\circ I)_0$ 是等价，这蕴涵 $IR =1_{\widehat A}$。
  根据 \cref{ct:isoprecat}\ref{item:ct:ipc3}，$I$ 是预范畴的同构。
  而 $\widehat A$ 是范畴，故 $A$ 也是范畴。
\end{proof}

\index{weak equivalence!of precategories|)}%

\newpage

\sectionNotes

范畴的原始定义当然是在集合论基础上给出的，因此一个范畴的对象构成一个集合（或者，对于大范畴而言，构成一个类）。
随着时间的推移，人们逐渐认识到对象的所有"范畴论"性质在同构下都是不变的，并且范畴中对象的相等通常不是一个很有用的概念。
多位作者~\cite{blanc:eqv-log,freyd:invar-eqv,makkai:folds,makkai:comparing}发现，依值类型逻辑使得我们能够在不涉及任何对象相等概念的情况下表述范畴的定义，并且在这种逻辑中可证明的命题恰恰是那些在同构下不变的"范畴论"命题。
\index{evil}%

尽管大部分范畴论内容似乎在对象的同构以及范畴的等价下都是不变的，但也存在一些有趣的例外，这些例外引发了关于何谓"范畴论的"的哲学讨论。
例如，Peter May 在 2010 年 5 月的范畴论邮件列表中提出了 \cref{ct:galois}，作为一个案例，其中两个（按集合论通常方式定义的）范畴是同构而非仅仅等价，这一区别至关重要。
对于倡导范畴论的同构不变版本的人来说，$\dagger$-范畴的情况也多少有些令人困惑，因为 $\dagger$-范畴中对象之间"正确的"相同性概念并非通常的同构，而是\emph{酉的}同构。
\index{isomorphism!invariance under}%

满足 \cref{ct:category} 中"饱和性"或"泛等性"原则的范畴，最初由 Hofmann 和 Streicher~\cite{hs:gpd-typethy} 研究。
几年后，Voevodsky（沃沃茨基）、Shulman 以及可能还有其他研究者大致同时、各自独立地想到了这个条件，并由 Ahrens 和 Kapulkin~\cite{aks:rezk} 加以形式化。
这一框架将上述所有例子统一了起来：有些预范畴是范畴，有些是严格范畴，等等。
一个关于一大类代数结构"同构蕴含相等"（假设泛等公理）的一般性定理由 Coquand（科康）和 Danielsson 证明；而 \cref{sec:sip} 中结构同一性原理的阐述则归功于 Aczel。

撇开关于范畴论的哲学考量不论，Rezk~\cite{rezk01css} 在定义 $(\infty,1)$-范畴的概念时发现，
\index{.infinity1-category@$(\infty,1)$-category}%
有一种做法非常方便：不局限于一个对象\emph{集合}配以对象间的态射空间，而是采用一个对象\emph{空间}，将对象间的所有等价与同伦都囊括其中。
由此得到了一类表现非常良好的 $(\infty,1)$-范畴模型，以特定的单纯空间呈现，Rezk 称之为\emph{完备 Segal 空间（complete Segal spaces）}。
\index{complete!Segal space}%
\index{Segal!space}%
该模型一个特别好的方面在于 \cref{ct:eqv-levelwise} 的类比：完备 Segal 空间之间的映射是等价的，当且仅当它是单纯空间之间的逐层等价。

在 Voevodsky 的泛等基础单纯\index{simplicial!sets}集合模型中加以解释时，我们的预范畴类似于 Rezk 的"Segal 空间"的一个截断类比，而我们的范畴则对应于他的"完备 Segal 空间"。
\index{Segal!category}%
严格范畴则对应于（经过弱化与截断的）所谓的"Segal 范畴"。
已知 Segal 范畴与完备 Segal 空间是 $(\infty,1)$-范畴的等价模型（例如参见~\cite{bergner:infty-one}），因此在单纯集合模型中，范畴与严格范畴所给出的范畴论是"等价"的——尽管如前所述，前者仍然具有许多优点。
然而，在高阶意象\index{.infinity1-topos@$(\infty,1)$-topos}的更一般的范畴语义下，严格范畴对应于相应层 1-意象\index{topos}中的（传统意义上的）内部范畴，而范畴对应于一个\emph{层栈（stack）}。
\index{stack}%
一般来说，相对于一个意象而言，后者是更合适的一类"范畴"。

在 Rezk 的语境中，我们所谓的"Rezk 完备化"对应于完备 Segal 空间的模型范畴中的纤维化替换\index{fibrant replacement}。
由于这是利用超限归纳论证构造的，它与我们的第二种构造——即作为高阶归纳类型给出的那个——最为接近。
然而，在同伦类型论的高阶意象模型中，Rezk 完备化对应于\emph{层栈完备化（stack completion）}，
\index{completion!stack}\index{stack!completion}
后者既可以用超限归纳~\cite{jt:strong-stacks} 来构造，也可以利用米田嵌入 \cite{bunge:stacks-morita-internal} 来构造。

\sectionExercises

\begin{ex}\label{ex:slice-precategory}
  对于一个预范畴 $A$ 和 $a:A$，定义\define{切片预范畴} $A/a$。
  \indexsee{precategory!slice}{category, slice}%
  \indexsee{slice (pre)category}{category, slice}%
  证明若 $A$ 是范畴，则 $A/a$ 也是范畴。
  \indexdef{category!slice}%
\end{ex}

\begin{ex}\label{ex:set-slice-over-equiv-functor-category}
  对任意集合 $X$，证明切片范畴 $\uset/X$ 等价于函子范畴 $\uset^X$，其中在后一种情形中将 $X$ 视为离散范畴。
\end{ex}

\begin{ex}\label{ex:functor-equiv-right-adjoint}
  \index{adjoint!functor}%
  \index{adjoint!equivalence}%
  证明一个函子是范畴的等价，当且仅当它是一个\emph{右}伴随，且其单位与余单位皆为同构。
\end{ex}

\begin{ex}\label{ct:pre2cat}
  定义\define{前2-范畴}的概念。
  \indexdef{pre-2-category}%
  证明在 \cref{sec:transfors} 中定义的预范畴、函子和自然变换构成一个前2-范畴。
  类似地，通过将等式（例如 \cref{ct:functor-assoc,ct:units} 中的那些）替换为满足类似相容条件的自然同构，来定义\define{前双范畴}。
  \indexdef{pre-bicategory}%
  定义一个从前2-范畴到前双范畴的函数，并证明当将其定义域与陪域都限制到那些 Hom 预范畴是范畴的前2-范畴上时，它成为一个等价。
\end{ex}

\begin{ex}\label{ct:2cat}
  定义\define{2-范畴}为满足类似于 \cref{ct:category} 中条件的前2-范畴。
  \indexdef{2-category}%
  验证范畴的前2-范畴 \ucat 是一个2-范畴。本章有多少内容可以在任意2-范畴内部完成？
\end{ex}

\begin{ex}\label{ct:groupoids}
  定义一个2-范畴，其对象是1-类型，态射是函数，2-态射是同伦。
  证明它在适当的意义上等价于 \ucat 中由\emph{群胚}（即其中每个箭头都是同构的范畴）张成的全子2-范畴。
\end{ex}

\begin{ex}\label{ex:2strict-cat}
  \index{strict!category}%
  回顾\emph{严格类别}，它是对象类型为集合的预范畴。
  证明严格范畴的前2-范畴等价于下述前2-范畴。
  \begin{itemize}
  \item 其对象是范畴 $A$，配备一个满射
    % \footnote{Recall that a function $f:X\to Y$ is a \emph{surjection} if for every $y:Y$, there \emph{merely exists} an $x:X$ such that $f(x)=y$.  This is to be distinguished from a \emph{split surjection}, which has the property that for every $y:Y$ there \emph{exists} an $x:X$ such that $f(x)=y$.}
    $p_A:A_0'\to A_0$，其中 $A_0'$ 是一个集合。
  \item 其态射是函子 $F:A\to B$，配备一个函数 $F_0':A_0' \to B_0'$ 使得 $p_B \circ F_0' = F_0 \circ p_A$。
  \item 其 2-态射即为自然变换。
  \end{itemize}
\end{ex}

\begin{ex}\label{ex:pre2dagger-cat}
  定义 $\dagger$-范畴的前2-范畴，其上的态射预范畴具有 $\dagger$-结构。
  证明两个 $\dagger$-范畴相等，当且仅当它们在适当的意义下是“酉等价”的。
\end{ex}

\begin{ex}\label{ct:ex:hocat}
  证明一个函数 $X\to Y$ 是一个等价，当且仅当它在 \cref{ct:hocat} 的同伦范畴中的像是一个同构。
  证明这个范畴的对象类型是 $\trunc1\type$。
\end{ex}

\begin{ex}\label{ex:dagger-rezk}
  构造从 $\dagger$-预范畴到 $\dagger$-范畴的 \define{$\dagger$-Rezk 完备化（$\dagger$-Rezk completion）}，并赋予其适当的泛性质。
\end{ex}

\begin{ex}\label{ex:rezk-vankampen}
  \index{van Kampen theorem}%
  \index{theorem!van Kampen}%
  \index{fundamental!groupoid}%
  \index{fundamental!pregroupoid}%
  利用 \cref{ct:fundgpd,ct:rezk-fundgpd-trunc1} 中的基本（预）群胚和 \cref{sec:rezk} 中的 Rezk 完备化，给出 van Kampen 定理（\cref{sec:van-kampen}）的另一个证明。
\end{ex}

\begin{ex}\label{ex:stack}
  设 $X$ 与 $Y$ 为集合，$p:Y\to X$ 为满射。
  \begin{enumerate}
  \item 对任意预范畴 $A$，定义 $A$ 中相对于 $p$ 的\define{下降数据}范畴 $\mathrm{Desc}(A,p)$。
    \indexdef{descent data}%
  \item 证明任意预范畴 $A$ 都是关于 $p$ 的\define{预层栈}，即典范函子 $A^X \to \mathrm{Desc}(A,p)$ 是全忠实的。
    \indexdef{prestack}%
  \item 证明若 $A$ 是一个范畴，则它是关于 $p$ 的\define{层栈}，即 $A^X \to \mathrm{Desc}(A,p)$ 是一个等价。
    \indexdef{stack}%
  \item 证明命题“每个严格范畴都是关于任意集合满射的层栈”等价于选择公理。
    \index{axiom!of choice}%
    \index{strict!category}%
  \end{enumerate}
\end{ex}

% Local Variables:
% TeX-master: "hott-online"
% End:

\chapter{集合论}
\label{cha:set-math}

\index{set|(}%

我们将集合视为具有特别简单同伦性质的类型，这一观念（参见\cref{sec:basics-sets}）与 Zermelo--Fraenkel\index{set theory!Zermelo--Fraenkel} 集合论中的集合颇为不同，后者构成一个具有复杂嵌套隶属结构的累积层次。
对许多数学目的而言，同伦类型论中的集合与 Zermelo--Fraenkel 集合论中的集合同样适用，但二者之间也存在重要的差异。

本章我们首先在\cref{sec:piw-pretopos}中证明，范畴 $\uset$ 拥有（大多数）集合范畴通常具有的性质。
\index{mathematics!constructive}%
\index{mathematics!predicative}%
在构造性、直谓、泛等的基础中，它是一个“$\Pi\mathsf{W}$-预意象（$\Pi\mathsf{W}$-pretopos）”；
\index{propositional!resizing}%
而如果我们假设命题大小重整（\cref{subsec:prop-subsets}），它就是一个初等意象\index{topos}；
如果我们再假设 \LEM{} 和 \choice{}，那么它便是 Lawvere（劳维尔）的\emph{集合范畴的初等理论}\index{Lawvere} 的一个模型。
\index{Elementary Theory of the Category of Sets}%
这足以保证同伦类型论中的集合，其行为与大多数集合论以外的数学家所使用的集合一致。

在本章余下部分，我们将考察一些传统上归入“集合论”的课题。
在\cref{sec:cardinals,sec:ordinals,sec:wellorderings}中，我们将研究基数与序数。
这些概念在集合论中传统上借助全局隶属关系来定义，但我们将会看到，泛等公理能够提供一种同样便捷、更具“结构性”的方法。

最后，在\cref{sec:cumulative-hierarchy}中，我们将探讨在同伦类型论\emph{内部}构造一个累积层次的可能性，并为其配备一个类似于 Zermelo--Fraenkel 集合论的二元隶属关系。
这将高阶归纳类型与代数集合论的思想结合起来。
\index{algebraic set theory}%
\index{set theory!algebraic}%

在本章中，我们将频繁使用\cref{subsec:prop-trunc}中描述的传统逻辑记号。
除了\cref{cha:basics,cha:logic}的基础理论外，我们还会用到\cref{sec:colimits,sec:set-quotients}中关于余极限和商的高阶归纳类型，以及\cref{cha:hlevels}中的一些截断理论，特别是当 $n=-1$ 时\cref{sec:image-factorization}中的因子分解系统。
在\cref{sec:ordinals}中，我们将用归纳族（\cref{sec:generalizations}）来描述良基性；
而在\cref{sec:cumulative-hierarchy}中，则会使用一个更复杂的高阶归纳类型来呈现累积层次。

%\section{\texorpdfstring{$\set$}{Set} is a \texorpdfstring{$\Pi$}{Π}W-pretopos}
\section{The category of sets}
\label{sec:piw-pretopos}

回顾在\cref{cha:category-theory}中，我们将范畴 \uset 定义为由所有 $0$-类型（位于某个宇宙 \UU 中）及其间的映射构成，并指出它是一个范畴（而不仅仅是预范畴）。
下面我们依次考察 \uset 所具备的各级结构。

\subsection{Limits and colimits}
\label{subsec:limits-sets}

\index{limit!of sets}%
\index{colimit!of sets}%

由于集合对积封闭，\cref{thm:prod-ump}中积的泛性质立即表明 \uset 具有有限积。
事实上，由等价 \[ \Parens{X\to \prd{a:A} B(a)} \eqvsym \Parens{\prd{a:A} (X\to B(a))}.\] 同样容易得到无限积。
并且我们在\cref{ex:pullback}\index{pullback}中看到，$f:A\to C$ 和 $g:B\to C$ 的拉回可以定义为 $\sm{a:A}{b:B} f(a)=g(b)$；当 $A,B,C$ 是集合时，它也是集合，且继承了应有的泛性质。
因此，\uset 在自然的意义上是一个\emph{完备的}范畴。
\index{category!complete}%
\index{complete!category}%

由于集合对 $+$ 封闭且包含 \emptyt，\uset 具有有限余积。
类似地，由于当 $A$ 和每个 $B(a)$ 是集合时，$\sm{a:A}B(a)$ 也是集合，所以它在 \uset 中提供了族 $B$ 的余积。
最后，我们在\cref{sec:pushouts}中证明了 $n$-类型中存在推出，特别地，这包含了 \uset。
因此，\uset 也是\emph{余完备的}。
\index{category!cocomplete}%
\index{cocomplete category}%

\subsection{Images}
\label{sec:image}

%We will show that $\uset$ is a $\Pi$W-pretopos.
接下来，我们证明 $\uset$ 是一个\define{正则范畴（regular category）}，即：
\indexdef{category!regular}%
\indexdef{regular!category}%
%

\begin{enumerate}
\item $\uset$ 是有限完备的。\label{item:reg1}
\item 任何函数 $f : A \to B$ 的核偶对 $\proj1,\proj2: (\sm{x,y:A} f(x)= f(y)) \to A$ 有一个余等子。\label{item:reg2}
  \indexdef{kernel!pair}
\item 正则满态射的拉回仍是正则满态射。\label{item:reg3}
\end{enumerate}

%
回忆一下，一个\define{正则满态射（regular epimorphism）}
\indexdef{epimorphism!regular}%
\indexdef{regular!epimorphism}%
是指可以成为\emph{某个}映射对的余等子的态射。
因此，在~\ref{item:reg3}中，要求余等子的拉回仍然是余等子，但不一定是那个被拉回的映射对的余等子。

\index{set-coequalizer}%
\index{image}
对于函数 $f:A\to B$ 的核偶对的余等子，一个自然的候选是 $f$ 的\emph{像（image）}，正如在 \cref{sec:image-factorization} 中所定义的那样。
回顾，我们曾定义 $\im(f)\defeq \sm{b:B} \brck{\hfib f b}$，配有函数
$\tilde{f}:A\to\im(f)$ 与 $i_f:\im(f)\to B$，由
\begin{align*}
  \tilde{f} & \defeq \lam{a} \Pairr{f(a),\,\bproj{\pairr{a,\refl{f(a)}}}}\\
i_f & \defeq \proj1
\end{align*}
定义，构成如下图表：
\begin{equation*}
  \xymatrix{
    **[l]{\sm{x,y:A} f(x)= f(y)}
    \ar@<0.25em>[r]^{\proj1}
    \ar@<-0.25em>[r]_{\proj2}
    &
    {A}
    \ar[r]^(0.4){\tilde{f}}
    \ar[rd]_{f}
    &
    {\im(f)}
    \ar@{..>}[d]^{i_f}
    \\ & &
    B
  }
\end{equation*}

回忆一下，一个函数 $f:A\to B$ 被称为\emph{满的（surjective）}，如果
\index{function!surjective}%
\narrowequation{\fall{b:B}\brck{\hfib f b},}
或等价地 $\fall{b:B} \exis{a:A} f(a)=b$。
我们也曾提到，集合之间的一个函数 $f:A\to B$ 被称为\emph{单的（injective）}，如果
\index{function!injective}%
$\fall{a,a':A} (f(a) = f(a')) \Rightarrow (a=a')$，或等价地，如果它的每个纤维都是命题。
由于这些分别是 \cref{cha:hlevels} 意义下的 $(-1)$-连通映射与 $(-1)$-截断映射，那里的一般理论表明：上图中的 $\tilde f$ 是满的而 $i_f$ 是单的，并且这个分解在拉回下是稳定的。

现在我们将满射性与单射性同相应的范畴论概念等同起来。
首先我们观察到，范畴论中的单态射与满态射有一个稍强的等价表述。

\begin{lem}\label{thm:mono}
  对于范畴 $A$ 中的一个态射 $f:\hom_A(a,b)$，以下条件等价。
  \begin{enumerate}
  \item $f$ 是一个\define{单态射（monomorphism）}：
    \indexdef{monomorphism}%
    对所有 $x:A$ 与 ${g,h:\hom_A(x,a)}$，若 $f\circ g = f\circ h$ 则 $g=h$。\label{item:mono1}
  \item （若 $A$ 有拉回）对角态射 $a\to a\times_b a$ 是一个同构。\label{item:mono4}
  \item 对所有 $x:A$ 和 $k:\hom_A(x,b)$，类型 $\sm{h:\hom_A(x,a)} (k = f\circ h)$ 是一个命题。\label{item:mono2}
  \item 对所有 $x:A$ 与 ${g:\hom_A(x,a)}$，类型 $\sm{h:\hom_A(x,a)} (f\circ g = f\circ h)$ 是可缩的。\label{item:mono3}
  \end{enumerate}
\end{lem}

\begin{proof}
  条件~\ref{item:mono1} 与条件~\ref{item:mono4} 的等价性是标准的范畴论结果。
  现在考虑集合之间的函数 $(f\circ \blank ):\hom_A(x,a) \to \hom_A(x,b)$。
  条件~\ref{item:mono1} 说它是单的，而条件~\ref{item:mono2} 说它的纤维是命题；因此它们等价。
  条件~\ref{item:mono2} 蕴涵条件~\ref{item:mono3}：取 $k\defeq f\circ g$，并注意一个有实例的命题是可缩的。
  最后，由条件~\ref{item:mono3} 可推出条件~\ref{item:mono1}，因为如果 $p:f\circ g= f\circ h$，那么 $(g,\refl{})$ 和 $(h,p)$ 都是条件~\ref{item:mono3} 中类型的项，因此它们相等，于是 $g=h$。
\end{proof}

\begin{lem}\label{thm:inj-mono}
  集合之间的一个函数 $f:A\to B$ 是单的，当且仅当它是 \uset 中的一个单态射\index{monomorphism}。
\end{lem}

\begin{proof}
  留给读者。
\end{proof}

当然，一个\define{满态射（epimorphism）}
\indexdef{epimorphism}%
\indexsee{epi}{epimorphism}%
就是反范畴中的单态射。
我们现在证明，在 \uset 中，满态射恰好就是满射，并且也恰好是余等子（正则满态射）。

$\uset$ 中一对映射 $f,g:A\to B$ 的余等子被定义为一般（同伦）余等子的 0-截断。
为清晰起见，我们可称之为\define{集合余等子（set-coequalizer）}。
\indexdef{set-coequalizer}%
\indexsee{coequalizer!of sets}{set-coequalizer}%
其泛性质可方便地表述如下。

\begin{lem}
\index{universal!property!of set-coequalizer}%
设 $f,g:A\to B$ 是集合 $A$ 与 $B$ 之间的函数。其集合余等子 $c_{f,g}:B\to Q$ 具有以下性质：对任意集合 $C$ 以及任意满足 $h\circ f = h\circ g$ 的 $h:B\to C$，类型
\begin{equation*}
\sm{k:Q\to C} (k\circ c_{f,g} = h)
\end{equation*}
是可缩的。
\end{lem}

\begin{lem}\label{epis-surj}
对于集合之间的任意函数 $f:A\to B$，以下条件等价：
\begin{enumerate}
\item $f$ 是一个满态射。
\item 考虑 $\uset$ 中定义映射锥\index{cone!of a function}的推出图表
\begin{equation*}
  \xymatrix{
    {A}
    \ar[r]^{f}
    \ar[d]
    &
    {B}
    \ar[d]^{\iota}
    \\
    {\unit}
    \ar[r]_{t}
    &
    {C_f}
  }
\end{equation*}
则类型 $C_f$ 是可缩的。
\item $f$ 是满射。
\end{enumerate}
\end{lem}

\begin{proof}
设 $f:A\to B$ 为集合之间的函数，并假设它是满态射；我们证明 $C_f$ 是可缩的。
$C_f$ 的构造子 $\unit\to C_f$ 给出了一个元素 $t:C_f$。
我们需要证明
\begin{equation*}
\prd{x:C_f} x= t.
\end{equation*}
注意 $x= t$ 是一个命题，因此我们可以对 $C_f$ 使用归纳法。
当然，当 $x$ 就是 $t$ 时，我们有 $\refl{t}:t=t$，所以只需找到
\begin{align*}
I_0 & : \prd{b:B} \iota(b)= t\\
I_1 & : \prd{a:A} \opp{\alpha_1(a)} \ct I_0(f(a))=\refl{t}
\end{align*}
其中 $\iota:B\to C_f$ 和 $\alpha_1:\prd{a:A} \iota(f(a))= t$ 是 $C_f$ 的其他构造子。
注意 $\alpha_1$ 是一个从 $\iota\circ f$ 到 $\mathsf{const}_t\circ f$ 的同伦，因此我们得到元素
\begin{equation*}
\pairr{\iota,\refl{\iota\circ f}},\pairr{\mathsf{const}_t,\alpha_1}:
\sm{h:B\to C_f} \iota\circ f \htpy h\circ f.
\end{equation*}
根据 \cref{thm:mono}\ref{item:mono3} 的对偶（以及函数外延性），存在一条道路
\begin{equation*}
\gamma:\pairr{\iota,\refl{\iota\circ f}}=\pairr{\mathsf{const}_t,\alpha_1}.
\end{equation*}
因此，我们可以定义 $I_0(b)\defeq \happly(\projpath1(\gamma),b):\iota(b)=t$。
我们还有
\[\projpath2(\gamma) : \trans{\projpath1(\gamma)}{\refl{\iota\circ f}} = \alpha_1. \]
这个传送涉及与 $f$ 的预复合，而预复合与 $\happly$ 交换。
因此，由道路类型中的传送，对任意 $a:A$ 我们得到 $I_0(f(a)) = \alpha_1(a)$，这就给出了 $I_1$。

现在假设 $C_f$ 是可缩的；我们证明 $f$ 是满的。
我们首先通过对 $C_f$ 递归定义一个类型族 $P:C_f\to\prop$，这是可行的，因为 \prop 是一个集合。
在点构造子上，我们定义
\begin{align*}
P(t) & \defeq \unit\\
P(\iota(b)) & \defeq \brck{\hfiber{f}b}.
\end{align*}
为了完成 $P$ 的构造，还需要给出一条道路
\narrowequation{\brck{\hfiber{f}{f(a)}} =_\prop \unit}
对所有 $a:A$。
然而，$\brck{\hfiber{f}{f(a)}}$ 有实例 $(f(a),\refl{f(a)})$。
由于它是一个命题，这意味着它是可缩的——从而等价，因此相等，于 \unit。
这就完成了 $P$ 的定义。
现在，既然 $C_f$ 被假设为可缩的，那么对任意 $x:C_f$，$P(x)$ 等价于 $P(t)$。
特别地，对每个 $b:B$，$P(\iota(b))\jdeq \brck{\hfiber{f}b}$ 等价于 $P(t)\jdeq \unit$，因此是可缩的。
于是，$f$ 是满的。

最后，假设 $f:A\to B$ 是满的，并考虑一个集合 $C$ 以及两个函数 $g,h:B\to C$ 满足 $g\circ f = h\circ f$。
由于 $f$ 被假设为满的，对所有 $b:B$，类型 $\brck{\hfib f b}$ 是可缩的。
因此我们有如下等价：
\begin{align*}
\prd{b:B} (g(b)= h(b))
& \eqvsym \prd{b:B} \Parens{\brck{\hfib f b} \to (g(b)= h(b))}\\
& \eqvsym \prd{b:B} \Parens{\hfib f b \to (g(b)= h(b))}\\
& \eqvsym \prd{b:B}{a:A}{p:f(a)= b} g(b)= h(b)\\
& \eqvsym \prd{a:A} g(f(a))= h(f(a))
\end{align*}
在第二行中我们用到了 $g(b)=h(b)$ 是一个命题这一事实，因为 $C$ 是一个集合。
但由假设，后一个类型是有实例的。
\end{proof}

% \begin{rem}
% The above theorem is not true when we replace $\set$ by $\type$
% (replacing it also in the definition of $\mathsf{epi}$ and $\mathsf{epi}'$).
% However, we do
% get the implications $\textit{ii.}\Rightarrow\textit{iii.}\Rightarrow
% \textit{iv.}$
% \end{rem}

\begin{thm}\label{thm:set_regular}\label{lem:images_are_coequalizers}
范畴 $\uset$ 是正则范畴。此外，集合之间的满函数是正则满态射。
\end{thm}

\begin{proof}
范畴论中有一个标准引理：如果一个范畴有有限极限，并且有一个拉回稳定的正交分解系统\index{orthogonal factorization system} $(\mathcal{E},\mathcal{M})$，其中 $\mathcal{M}$ 是单态射，那么该范畴是正则的，此时 $\mathcal{E}$ 自动由正则满态射组成。
（例如参见~\cite[A1.3.4]{elephant}。）
该分解系统的存在性已在 \cref{thm:orth-fact} 中证明。
\end{proof}

\begin{lem}\label{lem:pb_of_coeq_is_coeq}
在 \uset 中，正则满态射的拉回仍是正则满态射。
\end{lem}

\begin{proof}
  我们在 \cref{thm:stable-images} 中证明了 $n$-连通函数的拉回是 $n$-连通的。
  根据 \cref{lem:images_are_coequalizers}，只需取 $n=-1$ 应用此结论即可。
\end{proof}

\indexdef{image!of a subset}
作为 \uset（集合的范畴）是正则范畴的一个推论，我们可以在子集上定义“像”操作。
也就是说，给定 $f:A\to B$，$A$ 的任意子集 $P:\power A$（即一个谓词 $P:A\to \prop$）都有一个\define{像（image）}，它是 $B$ 的一个子集。
这可以直接定义为 $\setof{ y:B | \exis{x:A} f(x)=y \land P(x)}$，或者间接地定义为复合函数
\[ \setof{ x:A | P(x) } \to A \xrightarrow{f} B.\]
\symlabel{subset-image}
的像（按前述意义）。
我们有时也会使用常见的记号 $\setof{f(x) | P(x)}$ 来表示 $P$ 的像。

\subsection{商}\label{subsec:quotients}

\index{set-quotient|(}%
既然我们已经知道 $\uset$ 是正则的，那么为了证明 $\uset$ 是正合的，就需要证明每个等价关系都是有效的。
\index{effective!equivalence relation|(}%
\index{relation!effective equivalence|(}%
换句话说，给定一个等价关系 $R:A\to A\to\prop$，存在一对映射
$\proj1,\proj2:\sm{x,y:A} R(x,y)\to A$
的余等子 $c_R$；并且，$\proj1$ 和 $\proj2$ 构成了 $c_R$ 的核偶\index{kernel!pair}。

我们已经在 \cref{sec:set-quotients} 中见过两种由等价关系 $R:A\to A\to\prop$ 构造集合商的一般方法。
第一种可以描述为两个投影
\[\proj1,\proj2:\Parens{\sm{x,y:A} R(x,y)} \to A.\]
的集合余等子。
这类商的一个重要性质如下所述。

\begin{defn}
  如果下面的方块
\begin{equation*}
  \xymatrix{
    {\sm{x,y:A} R (x,y)}
    \ar[r]^(0.7){\proj1}
    \ar[d]_{\proj2}
    &
    {A}
    \ar[d]^{c_R}
    \\
    {A}
    \ar[r]_{c_R}
    &
    {A/R}
    }
\end{equation*}
是一个拉回，那么关系 $R:A\to A\to\prop$ 就被称为是\define{有效的}。
\indexdef{effective!relation}
\indexdef{effective!equivalence relation}%
\indexdef{relation!effective equivalence}%
\end{defn}

由于 $c_R$ 和自身的标准拉回是 $\sm{x,y:A} (c_R(x)=c_R(y))$，根据 \cref{thm:total-fiber-equiv}，这等价于要求典范变换 $\prd{x,y:A} R(x,y) \to (c_R(x)=c_R(y))$ 是一个逐纤维等价。

\begin{lem}\label{lem:sets_exact}
假设 $\pairr{A,R}$ 是一个等价关系。那么对任意 $x,y:A$，存在一个等价
\begin{equation*}
(c_R(x)= c_R(y))\eqvsym R(x,y)
\end{equation*}
换言之，等价关系是有效的。
\end{lem}

\begin{proof}
我们首先将 $R$ 延拓成一个关系 $\widetilde{R}:A/R\to A/R\to\prop$，随后我们将证明它等价于 $A/R$ 上的相等类型。我们通过对 $A/R$ 进行双重归纳来定义 $\widetilde{R}$（注意，根据命题的泛等公理，$\prop$ 是一个集合）。我们定义 $\widetilde{R}(c_R(x),c_R(y)) \defeq R(x,y)$。对于 $r:R(x,x')$ 和 $s:R(y,y')$，利用 $R$ 的传递性与对称性，我们可以从 $R(x,y)$ 到 $R(x',y')$ 给出一个等价。这就完成了 $\widetilde{R}$ 的定义。

剩下只需证明，对每一对 $w,w':A/R$，都有 $\widetilde{R}(w,w')\eqvsym (w= w')$。
方向 $(w=w')\to \widetilde{R}(w,w')$ 可以通过传送得到，只要我们证明 $\widetilde{R}$ 是自反的，而这可以通过一个简单的归纳完成。
另一个方向 $\widetilde{R}(w,w')\to (w= w')$ 是一个命题，并且由于 $c_R:A\to A/R$ 是满的，我们只需假设 $w$ 和 $w'$ 具有 $c_R(x)$ 和 $c_R(y)$ 的形式。
但在这种情况下，我们已有典范映射 $\widetilde{R}(c_R(x),c_R(y)) \defeq R(x,y) \to (c_R(x)=c_R(y))$。
（再次注意，这里出现了编码-解码方法。\index{encode-decode method}）
\end{proof}

商的第二种构造是作为 $R$ 的等价类组成的集合（即其幂集\index{power set}的一个子集）：
\[ A\sslash R \defeq \setof{ P:A\to\prop | P \text{ is an equivalence class of } R}. \]
为了使该构造与 $A$ 和 $R$ 保持在同一个宇宙中，需要命题大小重整\index{propositional resizing}\index{impredicative!quotient}\index{resizing}。

注意，如果我们把 $R$ 看作从 $A$ 到 $A\to \prop$ 的函数，那么 $A\sslash R$ 就等价于 \cref{sec:image} 中构造的像 $\im(R)$。
现在，在 \cref{lem:images_are_coequalizers} 中我们已经证明了像是余等子。特别地，我们立刻得到余等子图表
\begin{equation*}
  \xymatrix{
    **[l]{\sm{x,y:A} R (x)= R (y)}
    \ar@<0.25em>[r]^{\proj1}
    \ar@<-0.25em>[r]_{\proj2}
    &
    {A}
    \ar[r]
    &
    {A \sslash R.}
  }
\end{equation*}
我们可以利用这个来给出另一个证明：任何等价关系都是有效的，并且两种商的定义是一致的。

\begin{thm}\label{prop:kernels_are_effective}
对于任意两个集合之间的任意函数 $f:A\to B$，由
$\ker(f,x,y)\defeq (f(x)= f(y))$
给出的关系 $\ker(f):A\to A\to\prop$ 是有效的。
\end{thm}

\begin{proof}
我们将使用 $\im(f)$ 是 $\proj1,\proj2:
(\sm{x,y:A} f(x)= f(y))\to A$ 的余等子这一事实。
%we get this equivalence from~\cref{prop:images_are_coequalizers}
注意，函数
\[c_f\defeq\lam{a} \Parens{f(a),\brck{\pairr{a,\refl{f(a)}}}}
: A \to \im(f)
\]
的核偶由两个投影
\begin{equation*}
\proj1,\proj2:\Parens{\sm{x,y:A} c_f(x)= c_f(y)}\to A.
\end{equation*}
组成。
对于任意 $x,y:A$，我们有等价
\begin{align*}
  (c_f(x)= c_f(y))
  & \eqvsym \Parens{\sm{p:f(x)= f(y)} \trans{p}{\brck{\pairr{x,\refl{f(x)}}}} =\brck{\pairr{y,\refl{f(x)}}}}\\
  & \eqvsym (f(x)= f(y)),
\end{align*}
其中最后一个等价成立，是因为对于任意 $b:B$，$\brck{\hfiber{f}b}$ 是一个命题。
因此，我们得到
\begin{equation*}
\Parens{\sm{x,y:A} c_f(x)= c_f(y)}\eqvsym \Parens{\sm{x,y:A} f(x)= f(y)}
\end{equation*}
由此我们可以得出结论：对于任意函数 $f$，其核 $\ker f$ 都是一个有效关系。
\end{proof}

\begin{thm}
等价关系是有效的，并且存在一个等价 $A/R \eqvsym A\sslash  R $。
\end{thm}

\begin{proof}
我们需要分析余等子图表
\begin{equation*}
  \xymatrix{
    **[l]{\sm{x,y:A} R (x)= R (y)}
    \ar@<0.25em>[r]^{\proj1}
    \ar@<-0.25em>[r]_{\proj2}
    &
    {A}
    \ar[r]
    &
    {A \sslash R}
  }
\end{equation*}
根据泛等公理，类型 $R(x) = R(y)$ 等价于从 $R(x)$ 到 $R(y)$ 的同伦的类型，而这个类型又等价于
\narrowequation{\prd{z:A} R (x,z)\eqvsym R (y,z).}
由于 $R$ 是一个等价关系，后一个类型等价于 $R(x,y)$。综上所述，我们得到 $(R(x) = R(y)) \eqvsym R(x,y)$，因此 $R$ 是有效的，因为它等价于一个有效关系。并且，图表
\begin{equation*}
  \xymatrix{
    **[l]{\sm{x,y:A} R(x, y)}
    \ar@<0.25em>[r]^{\proj1}
    \ar@<-0.25em>[r]_{\proj2}
    &
    {A}
    \ar[r]
    &
    {A \sslash R.}
  }
\end{equation*}
是一个余等子图表。由于余等子在等价意义下是唯一的，于是有 $A/R \eqvsym A\sslash  R $。
\end{proof}

我们以商的第三种可能的构造来结束本节。设 $A$ 是一个集合，$R$ 是其上的一个等价关系。考虑以 $A$ 为对象、以 $R$ 为态射集的预范畴；那么该预范畴的 Rezk 完备化
\index{completion!Rezk}%
（参见 \cref{sec:rezk}）的对象类型就是这个商。请读者自行验证其中的细节。

\index{effective!equivalence relation|)}%
\index{relation!effective equivalence|)}%
\index{set-quotient|)}%

\subsection{\texorpdfstring{$\uset$}{Set} 是一个 \texorpdfstring{$\Pi\mathsf{W}$}{ΠW}-预意象}
\label{subsec:piw}

\index{structural!set theory|(}%

所谓 \emph{$\Pi\mathsf{W}$-预意象}
\index{PiW-pretopos@$\Pi\mathsf{W}$-pretopos}%
\indexsee{pretopos}{$\Pi\mathsf{W}$-预意象}
——即一个具有不相交有限余积、有效等价关系以及多项式自函子之初始代数的局部笛卡尔闭范畴
\index{locally cartesian closed category}%
\index{category!locally cartesian closed}%
——旨在作为一种“直谓的”
\index{mathematics!predicative}%
意象概念，亦即一个“直谓集合”的范畴。它可以为构造性数学
\index{mathematics!constructive}%
提供基础，其作用类似于通常的集合范畴之于经典
\index{mathematics!classical}%
数学。

通常，在构造性类型论中，人们需要诉诸于一种外部的 ``setoid'' 构造——即正合完备化——以得到一个具有此类闭包性质的范畴。
\index{setoid}\index{completion!exact}%
  特别是，数学中许多通常涉及（非构造性）幂集的构造都需要性质良好的商。值得注意的是，泛等基础能够\emph{在内部}（通过高阶归纳类型）提供这些构造，而无须借助此类外部构造。这体现了我们方法的一项显著优势，在后续的例子中我们将看到这一点。

\begin{thm}
  \index{PiW-pretopos@$\Pi\mathsf{W}$-预意象}
  范畴 $\uset$ 是一个 $\Pi\mathsf{W}$-预意象（$\Pi\mathsf{W}$-pretopos）。
\end{thm}

\begin{proof}
  我们有一个初始对象
  \index{initial!set}%
  $\emptyt$ 以及有限的、不相交的和 $A+B$。  这些结构在拉回下是稳定的，原因仅在于拉回具有右伴随\index{adjoint!functor}。  事实上，$\uset$ 是局部笛卡尔闭的：因为对于集合间的任意映射 $f:A\to B$，其``纤维化替换''\index{fibrant replacement} $\sm{a:A}f(a)=b$ 在 $B$ 之上等价于 $A$，并且该替换具有依值函数类型。

我们刚刚已经证明了 $\uset$ 是正则的（\cref{thm:set_regular}）并且其商是有效的（\cref{lem:sets_exact}）。因此，我们得到了一个局部笛卡尔闭的预意象。最后，根据\cref{ex:ntypes-closed-under-wtypes}，$n$-类型对 $W$-类型的形成是封闭的，并且根据\cref{thm:w-hinit}，$W$-类型是多项式自函子的初始代数。由此可知，$\uset$ 是一个 $\Pi\mathsf{W}$-预意象。
\end{proof}

\index{topos|(}
人们自然会问，是否有某种因素阻碍 $\uset$ 成为（初等）意象？
除了已经提到的结构外，一个意象还具有一个
\emph{子对象分类子（subobject classifier）}：
\indexdef{subobject classifier}%
\index{classifier!subobject}%
\index{power set}%
即一个带点对象，分类单态射（的等价类）。\index{monomorphism}（事实上，有了子对象分类子后，事情会变得更简单一些：只需笛卡尔闭性就能得到余极限。）

在同伦类型论中，泛等性意味着类型$\prop \defeq \sm{X:\UU}\isprop(X)$确实分类单态射（论证类似于\cref{sec:object-classification}），但一般来说，它和背景宇宙 $\UU$ 一样大。
因此，在"是 $0$-类型"的意义上它是一个"集合"，但在"是 $\UU$ 中对象"的意义上它并非"小的"，故而也不属于范畴 \uset。
然而，如果我们假定一种恰当的命题大小重整形式（参见\cref{subsec:prop-subsets}），那么我们就能找到一个 $\prop$ 的小版本，使得 $\uset$ 成为一个初等意象。

\begin{thm}\label{thm:settopos}
  \index{propositional!resizing}%
  如果存在一个包含所有命题的类型 $\Omega:\UU$，那么范畴 $\uset_\UU$ 是一个初等意象。
\end{thm}


\index{topos|)}

为此的一个充分条件是排中律，即我们称之为 \LEM{} 的"命题"形式；因为此时有 $\prop = \bool$，而 $\bool$ \emph{是}小的，并且它也分类所有命题。
此外，在意象理论中，\LEM{} 的一个众所周知的充分条件是选择公理，而选择公理在经典\index{mathematics!classical}集合论中当然通常被假定为公理。
在下一节中，我们将简要探讨在我们的框架下这些条件之间的关系。

\index{structural!set theory|)}%

%%%%%%%%%%%%%%%%%%%%%%%%%%%%%%%%%%%%%%%%%%%%%%%%
\subsection{选择公理蕴含排中律}
\label{subsec:emacinsets}

% In this section we prove a classic result that the axiom of choice implies excluded
% middle.

我们先给出以下引理。

\begin{lem}\label{prop:trunc_of_prop_is_set}
若 $A$ 是一个命题，则其纬悬 $\susp(A)$ 是一个集合，且 $A$ 等价于$\id[\susp(A)]{\north}{\south}$。
\end{lem}

\begin{proof}
为了证明 $\susp(A)$ 是一个集合，我们定义一个族 $P:\susp(A)\to\susp(A)\to\type$，它具有下述性质：对每个 $x,y:\susp(A)$，$P(x,y)$ 是一个命题，并且它等价于恒等类型 $\idtypevar{\susp(A)}$。
%
我们作如下定义：
\begin{align*}
P(\north,\north) & \defeq \unit &
P(\south,\north) & \defeq A\\
P(\north,\south) & \defeq A &
P(\south,\south) & \defeq \unit.
\end{align*}
我们必须验证该定义保持路径。
给定任意 $a : A$，存在一条经线 $\merid(a) : \north = \south$，因此我们也应该有
%
\begin{equation*}
  P(\north, \north) = P(\north, \south) = P(\south, \north) = P(\south, \south).
\end{equation*}
%
但由于 $A$ 有实例 $a$，它等价于 $\unit$，所以我们有
%
\begin{equation*}
  P(\north, \north) \eqvsym P(\north, \south) \eqvsym P(\south, \north) \eqvsym P(\south, \south).
\end{equation*}
%
泛等公理将这些转化为所需的相等性。同时，$P(x,y)$ 对于所有 $x, y : \susp(A)$ 都是一个命题，这可以通过对 $x$ 和 $y$ 进行归纳，并利用"是命题"本身是一个命题这一事实来证明。

注意 $P$ 是一个自反关系。
因此我们可以应用 \cref{thm:h-set-refrel-in-paths-sets}，于是只需构造 $\tau : \prd{x,y:\susp(A)}P(x,y)\to(x=y)$。我们通过双归纳来完成这一构造。
当 $x$ 是 $\north$ 时，我们通过下式定义 $\tau(\north)$：
%
\begin{equation*}
  \tau(\north,\north,u) \defeq \refl{\north}
  \qquad\text{and}\qquad
  \tau(\north,\south,a) \defeq \merid(a).
\end{equation*}
%
如果 $A$ 有实例 $a$，那么 $\merid(a) : \north = \south$，因此我们还需要
\narrowequation{
  \trans{\merid(a)}{\tau(\north, \north)} = \tau(\north, \south).
}
这可由函数外延性得到，利用如下事实：对于所有 $x : A$，
%
\begin{multline*}
  \trans{\merid(a)}{\tau(\north,\north,x)} =
  \tau(\north,\north,x) \ct \opp{\merid(a)} \jdeq \\
  \refl{\north} \ct \merid(a) =
  \merid(a) =
  \merid(x) \jdeq
  \tau(\north, \south, x).
\end{multline*}
对称地，我们可以通过下式定义 $\tau(\south)$：
%
\begin{equation*}
  \tau(\south,\north, a) \defeq \opp{\merid(a)}
  \qquad\text{and}\qquad
  \tau(\south,\south, u) \defeq \refl{\south}.
\end{equation*}
%
为了完成 $\tau$ 的构造，需要验证 $\trans{\merid(a)}{\tau(\north)} = \tau(\south)$，给定任意 $a : A$。验证过程与前类似，通过对 $\tau$ 的第二个参数进行归纳来完成。

因此，由 \cref{thm:h-set-refrel-in-paths-sets}，我们有 $\susp(A)$ 是一个集合，并且对于所有 $x,y:\susp(A)$，$P(x,y) \eqvsym (\id{x}{y})$。
取 $x\defeq \north$ 和 $y\defeq \south$ 即得到所需的 $\eqv{A}{(\id[\susp(A)]\north\south)}$。
\end{proof}

\begin{thm}[Diaconescu（迪亚科内斯库）]\label{thm:1surj_to_surj_to_pem}
  \index{axiom!of choice}%
  \index{excluded middle}%
  \index{Diaconescu's theorem}\index{theorem!Diaconescu's}%
  选择公理蕴含排中律。
\end{thm}

\begin{proof}
我们使用在 \cref{thm:ac-epis-split} 中给出的选择公理的等价形式。
考虑一个命题 $A$。
由 $f(\bfalse) \defeq \north$ 和 $f(\btrue) \defeq \south$ 定义的函数 $f:\bool\to\susp(A)$ 是满的。
事实上，我们有
$\pairr{\bfalse,\refl{\north}} : \hfiber{f}{\north}$
以及 $\pairr{\btrue,\refl{\south}} :\hfiber{f}{\south}$。
由于 $\bbrck{\hfiber{f}{x}}$ 是一个命题，通过归纳即可得证所述的满射性。

由 \cref{prop:trunc_of_prop_is_set}，纬悬 $\susp(A)$ 是一个集合，因此根据选择公理，仅仅存在 $f$ 的一个截面 $g: \susp(A) \to \bool$。
由于 $\bool$ 上的相等性是可判定的，我们得到
\begin{equation*}
 (g(f(\bfalse))= g(f(\btrue))) +
 \lnot (g(f(\bfalse))= g(f(\btrue))),
\end{equation*}
并且，由于 $g$ 是 $f$ 的截面，从而是单的，
\begin{equation*}
(f(\bfalse) = f(\btrue)) +
\lnot (f(\bfalse) = f(\btrue)).
\end{equation*}
最后，因为由 \cref{prop:trunc_of_prop_is_set} 有 $(f(\bfalse)=f(\btrue)) = (\north=\south) = A$，我们得到 $A+\neg A$。
\end{proof}

% This conclusion needs only \LEM{}, see \cref{ex:lemnm}.

% \begin{cor}\label{cor:ACtoLEM0}
%   If the axiom of choice \choice{} holds then $\brck{A + \neg A}$ for every set $A$.
% \end{cor}

% \begin{proof}
%   There is a surjection
%   \[
%   A + \neg A \epi \brck{A} + \brck{\neg A} \epi
%   \brck{(\brck{A} + \brck{\neg A})} = \brck{A} \vee \brck{\neg A} = \brck{A} \vee \neg \brck{A} = \unit,
%   \]
%   %
%   where in the last step excluded middle is available as a consequence of the axiom of choice.
%   Again by the axiom of choice there merely exists a section of the surjection, but this
%   is none other than an inhabitant of $A + \neg A$. Therefore $\brck{A+\neg A}$.
% \end{proof}

\index{denial}


\begin{thm}\label{thm:ETCS}
  \index{Elementary Theory of the Category of Sets}%
  \index{category!well-pointed}%
  如果选择公理成立，那么范畴 $\uset$ 是一个良点化的、带选择公理的布尔初等意象。
\end{thm}

\begin{proof}
  由于 \choice{} 蕴含 \LEM{}，根据 \cref{thm:settopos} 及其后的注记，我们得到一个带选择公理的布尔初等意象。我们将良点性的证明留给读者作为练习（\cref{ex:well-pointed}）。
\end{proof}

\begin{rmk}
  定理中提到的范畴条件被称为 Lawvere（劳维尔）为\emph{集合范畴的初等理论（Elementary Theory of the Category of Sets）}~\cite{lawvere:etcs-long} 所提出的公理。
\end{rmk}

\section{基数}
\label{sec:cardinals}

\begin{defn}\label{defn:card}
  \define{基数类型（type of cardinal numbers）}
  \indexdef{type!of cardinal numbers}%
  \indexdef{cardinal number}%
  \indexsee{number!cardinal}{cardinal number}%
  是集合类型 \set 的 0-截断：
  \[ \card \defeq \pizero{\set} \]
  因此，一个\define{基数（cardinal number）}，或简称\define{基数（cardinal）}，是 $\card\jdeq \pizero\set$ 的一个实例。当然，从技术上讲，每个宇宙 \type 都关联有一个独立的类型 $\card_\UU$。
\end{defn}

%\begin{rmk}

  % , but with these conventions we can state theorems beginning with ``for all cardinal numbers\dots''\ and give them exactly the same sort of meaning as those beginning ``for all types\dots''.
%\end{rmk}

按照截断的惯例，如果 $A$ 是一个集合，则 $\cd{A}$ 表示其在典范投影 $\set \to \trunc0\set \jdeq \card$ 下的像；我们称 $\cd{A}$ 为 $A$ 的\define{基数（cardinality）}\indexdef{cardinality}。
根据定义，\card 是一个集合。
它还从 \set 继承了半环结构。

\begin{defn}
  \define{基数加法（cardinal addition）}
  \indexdef{addition!of cardinal numbers}%
  \index{cardinal number!addition of}%
  \[ (\blank+\blank) : \card \to \card \to \card \]
  通过截断上的归纳来定义：
  \[ \cd{A} + \cd{B} \defeq \cd{A+B}.\]
\end{defn}

\begin{proof}
  由于 $\card\to\card$ 是一个集合，要为所有 $\alpha:\card$ 定义 $(\alpha+\blank):\card\to\card$，由归纳法，只需假设 $\alpha$ 是 $\cd{A}$，其中 $A:\set$。
  现在我们要定义 $(\cd{A}+\blank) :\card\to\card$，即要为所有 $\beta:\card$ 定义 $\cd{A}+\beta :\card$。
  然而，由于 $\card$ 是一个集合，由归纳法只需假设 $\beta$ 是 $\cd{B}$，其中 $B:\set$。
  此时，我们可以将 $\cd{A}+\cd{B}$ 定义为 $\cd{A+B}$。
\end{proof}

\begin{defn}
  类似地，\define{基数乘法}
  \indexdef{multiplication!of cardinal numbers}%
  \index{cardinal number!multiplication of}%
  \[ (\blank\cdot\blank) : \card \to \card \to \card \]
  运算也由截断归纳定义：
  \[ \cd{A} \cdot \cd{B} \defeq \cd{A\times B} \]
\end{defn}

\begin{lem}\label{card:semiring}
  \card 是一个交换半环\index{semiring}，即对任意 $\alpha,\beta,\gamma:\card$，我们有：
  \begin{align*}
    (\alpha+\beta)+\gamma &= \alpha+(\beta+\gamma)\\
    \alpha+0 &= \alpha\\
    \alpha + \beta &= \beta + \alpha\\
    (\alpha \cdot \beta) \cdot \gamma &= \alpha \cdot (\beta\cdot\gamma)\\
    \alpha \cdot 1 &= \alpha\\
    \alpha\cdot\beta &= \beta\cdot\alpha\\
    \alpha\cdot(\beta+\gamma) &= \alpha\cdot\beta + \alpha\cdot\gamma
  \end{align*}
  其中 $0 \defeq \cd{\emptyt}$，$1\defeq\cd{\unit}$。
\end{lem}

\begin{proof}
  我们证明乘法的交换律 $\alpha\cdot\beta = \beta\cdot\alpha$；其他性质的证明完全类似。
  由于 \card 是一个集合，类型 $\alpha\cdot\beta = \beta\cdot\alpha$ 是一个命题，特别地也是一个集合。
  于是，由归纳法，只需假设存在 $A,B:\set$ 使得 $\alpha$ 和 $\beta$ 分别是 $\cd{A}$ 和 $\cd{B}$。
  此时，$\cd{A}\cdot \cd{B} \jdeq \cd{A\times B}$ 且 $\cd{B}\cdot\cd{A} \jdeq \cd{B\times A}$，因此只需证明 $A\times B = B\times A$。
  最后，由泛等性，只需给出一个等价 $A\times B \eqvsym B\times A$。
  这很简单：取 $(a,b) \mapsto (b,a)$ 及其显然的逆即可。
\end{proof}

\begin{defn}
  \define{基数幂}运算同样由截断归纳定义：
  \indexdef{exponentiation, of cardinal numbers}%
  \index{cardinal number!exponentiation of}%
  \[ \cd{A}^{\cd{B}} \defeq \cd{B\to A}. \]
\end{defn}

\begin{lem}\label{card:exp}
  对任意 $\alpha,\beta,\gamma:\card$，有：
  \begin{align*}
    \alpha^0 &= 1\\
    1^\alpha &= 1\\
    \alpha^1 &= \alpha\\
    \alpha^{\beta+\gamma} &= \alpha^\beta \cdot \alpha^\gamma\\
    \alpha^{\beta\cdot \gamma} &= (\alpha^{\beta})^\gamma\\
    (\alpha\cdot\beta)^\gamma &= \alpha^\gamma \cdot \beta^\gamma
  \end{align*}
\end{lem}

\begin{proof}
  证明方法同 \cref{card:semiring}。
\end{proof}

\begin{defn}
  \define{基数不等式}这一关系
  \index{order!non-strict}%
  \index{cardinal number!inequality of}%
  \[ (\blank\le\blank) : \card\to\card\to\prop \]
  由截断归纳定义：
  \symlabel{inj}
  \[ \cd{A} \le \cd{B} \defeq \brck{\inj(A,B)} \]
  其中 $\inj(A,B)$ 是从 $A$ 到 $B$ 的单射的类型。
  \index{function!injective}%
  换句话说，$\cd{A} \le \cd{B}$ 意味着仅知存在从 $A$ 到 $B$ 的单射。
\end{defn}

\begin{lem}
  基数不等式是一个预序，即对于 $\alpha,\beta:\card$，有
  \index{preorder!of cardinal numbers}%
  \begin{gather*}
    \alpha \le \alpha\\
    (\alpha \le \beta) \to (\beta\le\gamma) \to (\alpha\le\gamma)
  \end{gather*}
\end{lem}

\begin{proof}
  与之前一样，由截断归纳。
  以传递性为例：由于 $(\alpha \le \beta) \to (\beta\le\gamma) \to (\alpha\le\gamma)$ 是一个命题，由 0-截断归纳，可以假设 $\alpha$、$\beta$ 和 $\gamma$ 分别是 $\cd{A}$、$\cd{B}$ 和 $\cd{C}$。
  现在，由于 $\cd{A} \le \cd{C}$ 也是一个命题，由 $(-1)$-截断归纳，可以假设给定了单射 $f:A\to B$ 和 $g:B\to C$。
  那么 $g\circ f$ 就是从 $A$ 到 $C$ 的一个单射，因此 $\cd{A} \le \cd{C}$ 成立。
  自反性的证明则更为简单。
\end{proof}

同样地，可以证明基数不等式与半环运算是相容的。

\begin{lem}\label{thm:injsurj}
  \index{function!injective}%
  \index{function!surjective}%
  考虑以下陈述：
  \begin{enumerate}
  \item 存在一个单射 $A\to B$。\label{item:cle-inj}
  \item 存在一个满射 $B\to A$。\label{item:cle-surj}
  \end{enumerate}
  那么，在假设排中律的前提下：
  \index{excluded middle}%
  \index{axiom!of choice}%
  \begin{itemize}
  \item 给定 $a_0:A$，有 \ref{item:cle-inj} $\to$ \ref{item:cle-surj}。
  \item 因此，如果 $A$ 仅知有实例，有 \ref{item:cle-inj} $\to$ 仅知 \ref{item:cle-surj}。
  \item 假设选择公理，有 \ref{item:cle-surj} $\to$ 仅知 \ref{item:cle-inj}。
  \end{itemize}
\end{lem}

\begin{proof}
  如果 $f:A\to B$ 是一个单射，如下定义 $g:B\to A$。
  对于 $b:B$，由于 $f$ 是单的，$f$ 在 $b$ 处的纤维是一个命题。
  因此，根据排中律，要么存在 $a:A$ 使得 $f(a)=b$，要么不存在。
  在第一种情况下，定义 $g(b)\defeq a$；否则令 $g(b)\defeq a_0$。
  那么对于任意 $a:A$，都有 $a = g(f(a))$，所以 $g$ 是满射。

  第二个陈述由第一个陈述经截断归纳得出。
  对于第三个陈述，如果 $g:B\to A$ 是满射，那么根据选择公理，仅知存在一个函数 $f:A\to B$，使得对所有 $a$ 都有 $g(f(a)) = a$。
  而这个 $f$ 必然是单射。
\end{proof}

\begin{thm}[Schroeder--Bernstein]
  \index{theorem!Schroeder--Bernstein}%
  \index{Schroeder--Bernstein theorem}%
  假设排中律，对于集合 $A$ 和 $B$，有
  \[ \inj(A,B) \to \inj(B,A) \to (A\cong B) \]
\end{thm}

\begin{proof}
  通常的"来回"论证可直接适用，无需显著改动。
  注意，它实际上构造了一个同构 $A\cong B$（这里假设了排中律，以便判定一个给定元素是属于一个循环、一条无限链、一条始于 $A$ 的链，还是一条始于 $B$ 的链）。
\end{proof}

\begin{cor}
  假设排中律，基数不等式是一个偏序，即对于 $\alpha,\beta:\card$，有
  \[ (\alpha\le\beta) \to (\beta\le\alpha) \to (\alpha=\beta). \]
\end{cor}

\begin{proof}
  由于 $\alpha=\beta$ 是一个命题，由截断归纳，可以假设 $\alpha$ 和 $\beta$ 分别是 $\cd{A}$ 和 $\cd{B}$，并且有单射 $f:A\to B$ 和 $g:B\to A$。
  那么由 Schroeder--Bernstein 定理，得到一个同构 $A\cong B$，进而得到相等 $\cd{A}=\cd{B}$。
\end{proof}

最后，我们可以复现 Cantor 定理，它表明对于每个基数，都存在一个更大的基数。

\begin{thm}[Cantor]
  \index{Cantor's theorem}%
  \index{theorem!Cantor's}%
  对于 $A:\set$，不存在满射 $A \to (A\to \bool)$。
\end{thm}

\begin{proof}
  假设 $f:A \to (A\to \bool)$ 是任意函数，定义 $g:A\to \bool$ 为 $g(a) \defeq \neg f(a)(a)$。
  如果 $g = f(a_0)$，则 $g(a_0) = f(a_0)(a_0)$ 但 $g(a_0) = \neg f(a_0)(a_0)$，产生矛盾。
  因此，$f$ 不是满射。
\end{proof}

\begin{cor}
  假设排中律，对于任意 $\alpha:\card$，存在基数 $\beta$ 使得 $\alpha\le\beta$ 且 $\alpha\neq\beta$。
\end{cor}

\begin{proof}
  令 $\beta = 2^\alpha$。
  现在我们要证明的是一个命题，因此由归纳法，可以假设 $\alpha$ 是 $\cd{A}$，于是 $\beta\jdeq \cd{A\to \bool}$。
  利用排中律，定义函数 $f:A\to (A\to \bool)$ 如下：
  \[f(a)(a') \defeq
  \begin{cases}
    \btrue &\quad a=a'\\
    \bfalse &\quad a\neq a'.
  \end{cases}
  \]
  如果 $f(a)=f(a')$，则 $f(a')(a) = f(a)(a) = \btrue$，从而 $a=a'$；因此 $f$ 是单射。
  所以，$\alpha \jdeq \cd{A} \le \cd{A\to \bool} \jdeq 2^\alpha$。

  另一方面，如果 $2^\alpha \le \alpha$，则会存在一个单射 $(A\to\bool)\to A$。
  由 \cref{thm:injsurj}，由于我们有 $(\lam{x} \bfalse):A\to \bool$ 且假设了排中律，则会存在一个满射 $A \to (A\to \bool)$，这与 Cantor 定理矛盾。
\end{proof}

\section{序数}
\label{sec:ordinals}

\index{ordinal|(}%

\begin{defn}\label{defn:accessibility}
  设 $A$ 是一个集合，且
  \[(\blank<\blank):A\to A\to \prop\]
  是 $A$ 上的一个二元关系。
  我们归纳地定义元素 $a:A$ 关于 $<$ 是\define{可达的}的含义：
  \indexdef{accessibility}%
  \indexsee{accessible}{accessibility}%
  \begin{itemize}
  \item 如果每个 $b<a$ 都是可达的，那么 $a$ 就是可达的。
  \end{itemize}
  我们记 $\acc(a)$ 表示 $a$ 是可达的。
\end{defn}

乍看之下，这样的归纳定义似乎永远无法启动，但显然，如果 $a$ 满足\emph{不存在} $b$ 使得 $b<a$ 这一性质，那么 $a$ 就是空虚地可达的。

注意，这是一个类型族的归纳定义，类似于 \cref{sec:generalizations} 中考虑的向量类型。
更准确地说，它有一个构造子，比如说 $\acc_<$，其类型为
\[ \acc_< : \prd{a:A} \Parens{\prd{b:A} (b<a) \to \acc(b)} \to \acc(a). \]
\index{induction principle!for accessibility}%
$\acc$ 的归纳原理是指，对于任意 $P:\prd{a:A} \acc(a) \to \type$，如果我们有
\[f:\prd{a:A}{h:\prd{b:A} (b<a) \to \acc(b)}
\Parens{\prd{b:A}{l:b<a} P(b,h(b,l))} \to
P(a,\acc_<(a,h)),
\]
那么就有一个由归纳定义的 $g:\prd{a:A}{c:\acc(a)} P(a,c)$，满足
\[g(a,\acc_<(a,h)) \jdeq f(a,\,h,\,\lam{b}{l} g(b,h(b,l))).\]
这一表述颇为繁琐，但我们通常只在 $P:A\to\type$ 仅依赖于 $A$ 这一更简单的情况下应用它。
在这种情况下，$f$ 的第二和第三个参数可以合并，因此我们需要证明的是：
\[f:\prd{a:A} \Parens{\prd{b:A} (b<a) \to \acc(b) \times P(b)}
\to P(a).
\]
也就是说，假设每个 $b<a$ 都是可达的且 $g(b):P(b)$ 已定义，由此定义 $g(a):P(a)$。

省略 $P$ 的第二个参数是合理的，这由下面的引理保证，而该引理的证明也是我们使用归纳原理更一般形式的唯一场合。

\begin{lem}
  可达性\index{accessibility}是一个命题。
\end{lem}

\begin{proof}
  我们必须证明，对于任意 $a:A$ 和 $s_1,s_2:\acc(a)$，有 $s_1=s_2$。
  我们通过对 $s_1$ 的归纳来证明这一点，令
  \[P_1(a,s_1) \defeq \prd{s_2:\acc(a)} (s_1=s_2). \]
  因此，我们必须证明，对于任意 $a:A$、${h_1:\prd{b:A} (b<a) \to \acc(b)}$ 以及
  \[ k_1:{\prd{b:A}{l:b<a}{t:\acc(b)} h_1(b,l) = t},\]
  对任意 $s_2:\acc(a)$ 都有 $\acc_<(a,h) = s_2$。
  我们将此陈述视为 $\prd{a:A}{s_2:\acc(a)} P_2(a,s_2)$，其中
  \[P_2(a,s_2) \defeq
  \prd{h_1 : \cdots } %{h_1:\prd{b:A} (b<a) \to \acc(b)}
  {k_1 : \cdots} % \Parens{\prd{b:A}{l:b<a}{t:\acc(b)} h_1(b,l) = t} \to
  (\acc_<(a,h_1) = s_2);
  \]
  因此可以通过对 $s_2$ 的归纳来证明。
  于是，假设 $h_2 : \prd{b:A} (b<a) \to \acc(b)$，以及 $k_2$ 具有一个极其复杂但在此无关紧要的类型，
  % \begin{narrowmultline*}
  %   k_2:\prd{b:A}{l:b<a}
  %   \prd{h_1:\prd{b':A} (b'<b) \to \acc(b')}
  %   \narrowbreak
  %   \Parens{\prd{b':A}{l':b'<b}{t':\acc(b')} h_1(b',l') = t'} \to
  %   (\acc_<(b,h_1) = h_2(b,l)).
  % \end{narrowmultline*}
  并且必须证明，对于任何具有上述类型的 $h_1$ 和 $k_1$，有 $\acc_<(a,h_1) = \acc_<(a,h_2)$。
  根据函数外延性，只需证明对所有 $b:A$ 和 $l:b<a$，有 $h_1(b,l) = h_2(b,l)$。
  而这由 $k_1$ 即得。
\end{proof}

\begin{defn}
  集合 $A$ 上的二元关系 $<$ 称为\define{良基的}
  \indexdef{relation!well-founded}%
  \indexdef{well-founded!relation}%
  如果 $A$ 的每个元素都是可达的。
\end{defn}

良基性的要旨在于：对于 $P:A\to \type$，我们可以利用 $\acc$ 的归纳原理得出 $\prd{a:A} \acc(a) \to P(a)$，然后应用良基性得出 $\prd{a:A} P(a)$。
换言之，如果从 $\fall{b:A} (b<a) \to P(b)$ 可以推出 $P(a)$，那么就有 $\fall{a:A} P(a)$。
这种方法称为\define{良基归纳}\indexdef{well-founded!induction}。

\begin{lem}
  良基性是一个命题。
\end{lem}

\begin{proof}
  $<$ 的良基性是类型 $\prd{a:A} \acc(a)$，由于每个 $\acc(a)$ 都是命题，所以它也是命题。
\end{proof}

\begin{eg}\label{thm:nat-wf}
  或许最常见的良基关系是 $\nat$ 上通常的严格序关系。
  为了证明它是良基的，需要对每个 $n:\nat$ 证明 $n$ 是可达的。
  \index{strong!induction}%
  这正是用 $\nat$ 上的普通归纳法证明"强归纳法"的常见证法。

  具体而言，我们对 $n:\nat$ 作归纳，证明对所有 $k\le n$，$k$ 都是可达的。
  基本情形是 $0$ 可达，这空虚成立，因为没有任何东西严格小于 $0$。
  在归纳步骤中，假设对所有 $k\le n$（即对所有 $k < n+1$），$k$ 都是可达的；于是根据定义，$n+1$ 也是可达的。

  $\nat$ 上的另一种良基关系可以通过如下方式定义：对所有 $n:\nat$，令 $n < \suc(n)$。
  这个关系的良基性几乎就是 $\nat$ 的普通归纳原理。
\end{eg}

\begin{eg}\label{thm:wtype-wf}
  设 $A:\set$ 且 $B : A \to \set$ 为一族集合。
  回顾 \cref{sec:w-types}，$W$ 类型 $\wtype{a:A} B(a)$ 由如下唯一的构造子归纳地生成：
  \begin{itemize}
  \item $\supp : \prd{a:A} (B(a) \to \wtype{x:A} B(x)) \to \wtype{x:A} B(x)$
  \end{itemize}
  我们通过对其第二个参数递归来定义 $\wtype{x:A} B(x)$ 上的关系 $<$：
  \begin{itemize}
  \item 对任意 $a:A$ 和 $f:B(a) \to \wtype{x:A} B(x)$，定义 $w<\supp(a,f)$ 为：仅仅存在 $b:B(a)$ 使得 $w = f(b)$。
  \end{itemize}
  现在我们用 $\wtype{x:A}B(x)$ 的通常归纳原理证明每个 $w:\wtype{x:A} B(x)$ 对此关系都是可达的。
  即假设给定 $a:A$ 和 $f:B(a) \to \wtype{x:A} B(x)$，以及一个提升 $f' : \prd{b:B(a)} \acc(f(b))$。
  于是根据 $<$ 的定义，对所有 $w<\supp(a,f)$ 都有 $\acc(w)$；因此 $\supp(a,f)$ 是可达的。
\end{eg}

良基性允许我们递归地定义函数、通过归纳证明命题，如下面的引理所示。
回顾 \cref{subsec:prop-subsets}，$\power B$ 表示\emph{幂集}\index{power set} $\power B \defeq (B\to\prop)$。

\begin{lem}\label{thm:wfrec}
  设 $B$ 为集合，且我们有一个函数
  \[ g : \power B \to B \]
  那么，如果 $<$ 是 $A$ 上的良基关系，则存在一个函数 $f:A\to B$，使得对所有 $a:A$ 都有
  \begin{equation*}
    f(a) = g\Big(\setof{ f(a') | a'<a }\Big).
  \end{equation*}
\end{lem}


\noindent
（这里使用了 \cref{sec:image} 中子集像的记号。）


\begin{proof}
  我们首先对每个 $a:A$ 和 $s:\acc(a)$ 定义一个元素 $\bar f(a,s):B$。
  由归纳原理，只需假设 $s$ 是一个函数，为每个 $a'<a$ 指派一个见证 $s(a'):\acc(a')$，而且对每个这样的 $a'$ 已有元素 $\bar f(a',s(a')):B$。
  在此情形下，我们定义
  \begin{equation*}
    \bar f(a,s) \defeq g\Big(\setof{ \bar f(a',s(a')) | a'<a }\Big).
  \end{equation*}

  现在，由于 $<$ 是良基的，我们有一个函数 $w:\prd{a:A} \acc(a)$。
  因此，我们可以定义 $f(a)\defeq \bar f (a,w(a))$。
\end{proof}

在经典\index{mathematics!classical}逻辑中，良基性有一种更为人熟知的等价表述。
在下文中，我们称一个子集 $B: \power A$ 是\define{非空的}
\indexdef{nonempty subset}
如果它与空子集 $(\lam{x}\bot) : \power X$ 不相等。
我们留待读者验证：在假设排中律的前提下，这等价于仅被居住，即条件 $\exis{x:A} x\in B$。

\begin{lem}\label{thm:wfmin}
  \index{excluded middle}%
  假设排中律，则 $<$ 是良基的当且仅当每个非空子集 $B: \power A$ 仅仅有一个极小元。
\end{lem}

\begin{proof}
  先假设 $<$ 是良基的，并设 $B\subseteq A$ 是一个没有极小元的子集。
  也就是说，对任意满足 $a\in B$ 的 $a:A$，都仅仅存在某个 $b:A$ 使得 $b<a$ 且 $b\in B$。

  我们断言：对任意 $a:A$ 和 $s:\acc(a)$，都有 $a\notin B$。对此作归纳证明。
  可假设 $s$ 是一个函数，对每个 $a'<a$ 指派一个证明 $s(a'):\acc(a')$，而且对每个这样的 $a'$，已有 $a'\notin B$。
  若 $a\in B$，则由假设，将仅仅存在某个 $b<a$ 使得 $b\in B$，这与归纳假设矛盾。
  因此，$a\notin B$；这就完成了归纳。
  由于 $<$ 是良基的，对所有 $a:A$ 都有 $a\notin B$，即 $B$ 是空集。

  现在反过来假设每个非空子集都仅仅有一个极小元。
  令 $B = \setof{ a:A | \neg \acc(a) }$。
  若 $B$ 非空，则它仅仅有一个极小元。
  于是，仅仅存在某个满足 $a\in B$ 的 $a:A$，使得对所有 $b<a$，都有 $\acc(b)$。
  但根据定义（并通过对截断作归纳），$a$ 仅仅可达，从而可达，这与 $a\in B$ 矛盾。
  因此 $B$ 是空集，故 $<$ 是良基的。
\end{proof}

\begin{defn}
  集合 $A$ 上的良基关系 $<$ 被称为\define{外延的}
  \indexdef{relation!extensional}%
  \indexdef{extensional!relation}%
  若对任意 $a,b:A$，都有
  \[ \Parens{\fall{c:A} (c<a) \Leftrightarrow (c<b)} \to (a=b). \]
\end{defn}

注意，由于 $A$ 是集合，外延性是命题。
这里的"外延性"与函数外延性无关，也与相等类型的外延性无关。
\index{axiom!of extensionality}%
更确切地说，它是经典集合论中外延公理的"局部"对应物。

\begin{thm}
  外延良基关系所构成的类型是一个集合。
\end{thm}

\begin{proof}
  根据泛等公理，只需证明：如果 $(A,<)$ 是外延且良基的，并且 $f:(A,<) \cong (A,<)$，那么 $f=\idfunc[A]$。
  \index{automorphism!of extensional well-founded relations}%
  我们通过对 $<$ 作归纳来证明对所有 $a:A$ 都有 $f(a)=a$。
  归纳假设是：对于所有 $a'<a$，有 $f(a')=a'$。

  现在，由于 $A$ 是外延的，要证 $f(a)=a$，只需证明
  \[\fall{c:A}(c<f(a)) \Leftrightarrow (c<a).\]
  然而，因为 $f$ 是一个自同构，我们有 $(c<a) \Leftrightarrow (f(c)<f(a))$。
  但由归纳假设，$c<a$ 意味着 $f(c)=c$，所以 $(c<a) \to (c<f(a))$。
  另一方面，如果 $c<f(a)$，那么 $f^{-1}(c)<a$，从而再次利用归纳假设得到 $c = f(f^{-1}(c)) = f^{-1}(c)$；于是有 $c<a$。
  因此，对任意 $c:A$，我们有 $(c<a) \Leftrightarrow (c<f(a))$，故 $f(a)=a$。
\end{proof}

\begin{defn}\label{def:simulation}
  如果 $(A,<)$ 和 $(B,<)$ 都是外延且良基的，则称满足以下条件的函数 $f:A\to B$ 为\define{模拟}
  \indexdef{simulation}%
  \indexsee{function!simulation}{simulation}%
  ：
  \begin{enumerate}
  \item 若 $a<a'$，则 $f(a)<f(a')$；并且\label{item:sim1}
  \item 对所有 $a:A$ 和 $b:B$，若 $b<f(a)$，则仅仅存在 $a'<a$ 满足 $f(a')=b$。\label{item:sim2}
  \end{enumerate}
\end{defn}

\begin{lem}
  任意模拟都是单的。
\end{lem}

\begin{proof}
  我们通过双重良基归纳法证明：对任意 $a,b:A$，若 $f(a)=f(b)$，则 $a=b$。
  关于 $a:A$ 的归纳假设是：对任意 $a'<a$ 和任意 $b:B$，若 $f(a')=f(b)$，则 $a=b$。
  关于 $b:A$ 的内层归纳假设是：对任意 $b'<b$，若 $f(a)=f(b')$，则 $a=b'$。

  假设 $f(a)=f(b)$；我们需要证明 $a=b$。
  由外延性，只需证明对任意 $c:A$ 我们有 $(c<a)\Leftrightarrow (c<b)$。
  如果 $c<a$，那么由 \cref{def:simulation}\ref{item:sim1} 可得 $f(c)<f(a)$。
  故 $f(c)<f(b)$，于是由 \cref{def:simulation}\ref{item:sim2} 可知，仅仅存在 $c':A$ 使得 $c'<b$ 且 $f(c)=f(c')$。
  由关于 $a$ 的归纳假设，我们有 $c=c'$，因此 $c<b$。
  对偶的论证同理可证。
\end{proof}

特别地，这表明在 \cref{def:simulation}\ref{item:sim2} 中，"仅仅"一词可以省略而不改变意思。

\begin{cor}
  如果 $f:A\to B$ 是模拟，那么对所有 $a:A$ 和 $b:B$，若 $b<f(a)$，则\emph{纯粹存在} $a'<a$ 使得 $f(a')=b$。
\end{cor}

\begin{proof}
  由于 $f$ 是单的，$\sm{a:A} (f(a)=b)$ 是命题。
\end{proof}

我们称子集 $C :\power B$ 为\define{初始段}
\indexdef{initial!segment}%
\indexsee{segment, initial}{initial segment}%
，如果 $c\in C$ 且 $b<c$ 蕴含 $b\in C$。
模拟的像必为初始段，而任意初始段的包含映射都是模拟。
因此，由泛等性，每个模拟 $A\to B$ 都\emph{相等于} $B$ 的某个初始段的包含映射。

\begin{thm}
  对一个集合 $A$，令 $P(A)$ 为 $A$ 上外延良基关系的类型。
  若 $\mathord{<_A} : P(A)$ 与 $\mathord{<_B} : P(B)$ 且 $f:A\to B$，令 $H_{\mathord{<_A}\mathord{<_B}}(f)$ 为命题"$f$ 是模拟"。
  则 $(P,H)$ 是 \cref{sec:sip} 意义下 \uset 上的标准结构概念。
\end{thm}

\begin{proof}
  恒等映射是模拟、模拟的复合仍是模拟，这些留给读者验证。
  因此，模拟构成了一个结构概念。
  为验证标准性，我们须证：若 $<$ 和 $\prec$ 是 $A$ 上的两个外延良基关系，且 $\idfunc[A]$ 在两个方向都是模拟，则 $<$ 与 $\prec$ 相等。
  由于外延性和良基性都是命题，为此只需 $\fall{a,b:A} (a<b) \Leftrightarrow (a\prec b)$ 成立。
  而这由 $\idfunc[A]$ 的 \cref{def:simulation}\ref{item:sim1} 即可推出。
\end{proof}

\begin{cor}\label{thm:wfcat}
  存在一个范畴，其对象为配备外延良基关系的集合，态射为模拟。
\end{cor}

实际上，这个范畴是一个偏序集。

\begin{lem}
  对于外延良基的 $(A,<)$ 和 $(B,<)$，至多存在一个模拟 $f:A\to B$。
\end{lem}

\begin{proof}
  假设 $f,g:A\to B$ 都是模拟。
  由于"是模拟"是一个命题，只需证明 $\fall{a:A}(f(a)=g(a))$。
  对 $<$ 作归纳，可假设对所有 $a'<a$ 均有 $f(a')=g(a')$。
  而由 $B$ 的外延性，要证 $f(a)=g(a)$，只需 $\fall{b:B}(b<f(a)) \Leftrightarrow (b<g(a))$ 成立。

  由于 $f$ 是模拟，若 $b<f(a)$，则存在 $a'<a$ 使得 $f(a')=b$。
  由归纳假设亦有 $g(a')=b$，从而 $b<g(a)$。
  对偶论证同理可证。
\end{proof}

因此，若 $A$ 和 $B$ 配备了外延良基关系，可记 $A\le B$ 表示存在模拟 $f:A\to B$。
\cref{thm:wfcat} 表明：若 $A\le B$ 且 $B\le A$，则 $A=B$。

\begin{defn}
  一个\define{序数}
  \indexdef{ordinal}%
  \indexsee{number!ordinal}{ordinal}%
  是一个集合 $A$ 连同其上的一个外延良基关系，且该关系是\emph{传递的}，即满足 $\fall{a,b,c:A}(a<b)\to (b<c) \to (a<c)$。
\end{defn}

\begin{eg}
  当然，$\nat$ 上通常的严格序是传递的。
  不难看出它也是外延的；因此它是一个序数。
  依照惯例，我们记这个序数为 $\omega$。
\end{eg}

\symlabel{ord}
令 $\ord$ 表示序数的类型。
根据之前的结果，$\ord$ 是一个集合并具有自然的偏序。
我们现在证明 $\ord$ 也容许一个良基关系。

\symlabel{initial-segment}
如果 $A$ 是一个序数且 $a:A$，令 $\ordsl A a \defeq \setof{ b:A | b<a}$ 表示\emph{初始段}。
\index{initial!segment}%
注意，若 $\ordsl A a = \ordsl A b$ 作为序数相等，则该同构必须保持它们到 $A$ 的包含映射（因为模拟构成偏序集），从而它们作为 $A$ 的子集也是相等的。
因此，由 $A$ 的外延性知 $a=b$。
于是，函数 $a\mapsto \ordsl A a$ 是一个单射 $A\to \ord$。

\begin{defn}
  对于序数 $A$ 和 $B$，一个模拟 $f:A\to B$ 被称为\define{有界的}
  \indexdef{simulation!bounded}%
  \indexdef{bounded!simulation}%
  如果存在 $b:B$ 使得 $A = \ordsl B b$。
\end{defn}

上述讨论表明，这样的 $b$ 若存在则唯一，因此有界性是一个命题。

我们记 $A<B$ 表示存在一个从 $A$ 到 $B$ 的有界模拟。
由于模拟是唯一的，$A<B$ 也是一个命题。

\begin{thm}\label{thm:ordord}
  $(\ord,<)$ 是一个序数。
\end{thm}

\noindent
更准确地说，该定理断言：某个宇宙\index{universe level}中的序数类型 $\ord_{\UU_i}$ 本身是更高一层宇宙中的序数，即 $(\ord_{\UU_i},<):\ord_{\UU_{i+1}}$。

\begin{proof}
  设 $A$ 是一个序数；首先证明对所有 $a:A$，$\ordsl A a$ 在 $\ord$ 中是可达的。
  对 $A$ 作良基归纳，假设对所有 $b<a$，$\ordsl A b$ 是可达的。
  根据可达性的定义，须证对所有 $B<\ordsl A a$，$B$ 在 $\ord$ 中是可达的。
  然而，若 $B<\ordsl A a$，则存在某个 $b<a$ 使得 $B = \ordsl{(\ordsl A a)}{b} = \ordsl A b$，而由归纳假设知后者是可达的。
  因此，对所有 $a:A$，$\ordsl A a$ 是可达的。

  现在证明 $A$ 在 $\ord$ 中是可达的。根据定义，须证对所有 $B<A$，$B$ 是可达的。
  但如前所述，$B<A$ 意味着存在某个 $a:A$ 使得 $B=\ordsl A a$，而由刚才的证明知这是可达的。
  因此，$\ord$ 是良基的。

  对于外延性，假设 $A$ 和 $B$ 是满足下式的序数：
  \narrowequation{\prd{C:\ord} (C<A) \Leftrightarrow (C<B).}
  那么对每个 $a:A$，由于 $\ordsl A a<A$，我们有 $\ordsl A a<B$，从而存在 $b:B$ 使得 $\ordsl A a = \ordsl B b$。
  定义 $f:A\to B$ 将每个 $a$ 映到相应的 $b$；可以直接验证 $f$ 是一个同构。
  于是 $A\cong B$，由泛等性知 $A=B$。

  最后，容易看出 $<$ 是传递的。
\end{proof}

将 $\ord$ 视为序数通常非常方便，但也有其隐患。
例如，考虑以下引理，我们需留意其中宇宙的使用方式。

\begin{lem}\label{thm:ordsucc}
  设 \bbU 为一个宇宙。
  对于任意 $A:\ord_\bbU$，存在 $B:\ord_\bbU$ 使得 $A<B$。
\end{lem}

\begin{proof}
  令 $B=A+\unit$，并规定 $\unit$ 中的元素 $\ttt$ 大于 $A$ 的所有元素。
  则 $B$ 是一个序数，且容易看出 $A\cong \ordsl B \ttt$。
\end{proof}

在 \cref{thm:ordsucc} 证明中所构造的序数 $B$ 称为 $A$ 的\define{后继}\indexdef{successor!of an ordinal}。

这条引理揭示了使用"类型歧义"\index{typical ambiguity}风格——即用 \UU 表示一个任意且未指定的宇宙——可能带来的隐患。
考虑如下另一种证明。

\begin{proof}[\cref{thm:ordsucc} 的另一种拟议证明]
  注意到 $C<A$ 当且仅当存在某个 $a:A$ 使得 $C=\ordsl A a$。
  这给出了一个同构 $A \cong \ordsl \ord A$，从而 $A<\ord$。
  于是我们可以取 $B\defeq\ord$。
\end{proof}

如果我们用类型歧义的风格陈述 \cref{thm:ordsucc}，则第二种证明是有效的。
但这样得到的引理用处较小，因为第二种证明会强制引理陈述中第二个"\ord"所指的宇宙层级高于第一个"\ord"。
而第一种证明则允许两个宇宙是同一个。

类似的讨论也适用于下一条引理：该引理也可通过观察"对任意 $A:\ord$ 皆有 $A\le \ord$"来证明，但那样得到的结论用处较小。

\begin{lem}\label{thm:ordunion}
  设 \bbU 为一个宇宙。
  对于任意 $X:\type$ 与 $F:X\to \ord_\bbU$，存在 $B:\ord_\bbU$ 使得对所有 $x:X$ 均有 $Fx\le B$。
\end{lem}

\begin{proof}
  在 $\sm{x:X} Fx$ 上定义等价关系 $\eqr$ 如下，并令 $B$ 为其集合商（参见 \cref{rmk:quotient-of-non-set}）：
  \[ (x,y) \eqr (x',y')
  \;\defeq\;
  \Big(\ordsl{(Fx)}{y} \cong \ordsl{(Fx')}{y'}\Big).
  \]
  定义 $(x,y)<(x',y')$ 当且仅当 $\ordsl{(Fx)}{y} < \ordsl{(Fx')}{y'}$。
  该关系显然下降到商上，并可验证使 $B$ 成为一个序数。
  此外，对每个 $x:X$，诱导的映射 $Fx\to B$ 是一个模拟。
\end{proof}

\section{经典良序}
\label{sec:wellorderings}

\index{denial|(}%
现在我们证明，我们的序数与更为人熟知的经典\index{mathematics!classical}良序等价。

\begin{lem}
  \index{excluded middle}%
  假设排中律，每个序数都是三分的：
  \index{trichotomy of ordinals}%
  \index{ordinal!trichotomy of}%
  \[ \fall{a,b:A} (a<b) \vee (a=b) \vee (b<a). \]
\end{lem}

\begin{proof}
  对 $a$ 归纳，可假设对任意 $a'<a$ 及任意 $b':A$，有 $(a'<b') \vee (a'=b') \vee (b'<a')$。
  现在对 $b$ 归纳，可假设对任意 $b'<b$，有 $(a<b') \vee (a=b') \vee (b'<a)$。

  根据排中律，要么仅存在 $b'<b$ 使得 $a<b'$，要么仅存在 $b'<b$ 使得 $a=b'$，要么对任意 $b'<b$ 皆有 $b'<a$。
  在第一种情形，由传递性仅知 $a<b$，而 $a<b$ 是命题，故确实有 $a<b$。
  类似地，在第二种情形，由迁移可得 $a<b$。
  因此，我们假设 $\fall{b':A}(b'<b)\to (b'<a)$。

  现在类似地，要么仅存在 $a'<a$ 使得 $b<a'$，要么仅存在 $a'<a$ 使得 $a'=b$，要么对任意 $a'<a$ 皆有 $a'<b$。
  在第一种和第二种情形，$b<a$，因此我们可假设 $\fall{a':A}(a'<a)\to (a'<b)$。
  然而，由外延性，这两个假设现在蕴涵 $a=b$。
\end{proof}

\begin{lem}
  良基关系不包含循环，即
  \[ \fall{n:\mathbb{N}}{a:\mathbb{N}_n\to A} \neg\Big((a_0<a_1) \wedge \dots \wedge (a_{n-1}<a_n)\wedge (a_n<a_0)\Big). \]
\end{lem}

\begin{proof}
  我们对 $a:A$ 归纳证明：不存在包含 $a$ 的循环。
  于是，由归纳可假设对所有 $a'<a$，不存在包含 $a'$ 的循环。
  但在任何包含 $a$ 的循环中，都存在某个小于 $a$ 且属于该循环的元素。
\end{proof}

\indexdef{relation!irreflexive}%
\index{irreflexivity!of well-founded relation}%
特别地，良基关系必须是\define{不自反的}，即对任意 $a$ 有 $\neg(a<a)$。

\begin{thm}\label{thm:wellorder}
  假设排中律，则 $(A,<)$ 是一个序数当且仅当 $A$ 的每个非空子集 $B\subseteq A$ 都有最小元。
\end{thm}

\begin{proof}
  若 $A$ 是序数，则由 \cref{thm:wfmin}，其每个非空子集都仅有一个极小元。
  但三分性保证了任何极小元都是最小元。
  此外，最小元若存在则唯一，因此仅存在一个最小元与实际拥有一个效果相同。

  反之，若每个非空子集都有最小元，则由 \cref{thm:wfmin} 知 $A$ 是良基的。
  我们也有三分性，因为对任意 $a,b$，子集
  $ \setof{a,b} \defeq \setof{x:A | x=a \lor x=b} $
  仅有一个最小元，而该最小元必定是 $a$ 或 $b$。
  这蕴含传递性：若 $a<b$ 且 $b<c$，则 $a=c$ 或 $c<a$ 都会产生循环。
  类似地，它也蕴含外延性：若 $\fall{c:A}(c<a)\Leftrightarrow (c<b)$，则 $a<b$ 将导致（取 $c$ 为 $a$ 时）$a<a$，这构成循环；$b<a$ 的情形同理；故必有 $a=b$。
\end{proof}

在经典\index{mathematics!classical}数学中，\cref{thm:wellorder} 这一刻画被直接当作\emph{良序}的定义，而\emph{序数}则构成良序同构类的一组标准代表元。
在我们的语境中，结构同一性原理意味着无需寻找这样的代表元：任何一个良序都不逊于任何其他良序。

现在我们转而考虑选择公理的推论。
对任意集合 $X$，令 $\powerp X$ 表示由 $X$ 的\emph{仅知有实例的子集}构成的类型：
\symlabel{inhabited-powerset}
\[ \powerp X \defeq \setof{ Y : \power X | \exis{x:X} x\in Y}. \]
若假设排中律，这便等价于由 $X$ 的\emph{非空}\index{nonempty subset}子集构成的类型，且有 $\power X \eqvsym (\powerp X) + \unit$。

\begin{thm}\label{thm:wop}
  \index{axiom!of choice}%
  \index{excluded middle}%
  在假设排中律的前提下，以下两条陈述等价：
  \begin{enumerate}
  \item 对每个集合 $X$，仅存在一个函数 $f: \powerp X \to X$，使得对所有 $Y:\powerp X$ 都有 $f(Y)\in Y$。\label{item:wop1}
  \item 每个集合仅具有一个序数结构。\label{item:wop2}
  \end{enumerate}
\end{thm}

\noindent
当然，~\ref{item:wop1} 是选择公理的经典\index{mathematics!classical}标准形式；参见 \cref{ex:choice-function}。

\begin{proof}
  一个方向的证明很简单：假设~\ref{item:wop2}。
  由于我们要证明的是命题~\ref{item:wop1}，可以假定 $A$ 是一个序数。
  于是就可以把 $f(B)$ 定义为 $B$ 中的最小元。

  现在假设~\ref{item:wop1}。
  同前，由于~\ref{item:wop2} 是一个命题，可以假定已经给定了这样的 $f$。
  我们以显然的方式将 $f$ 延拓为函数
  \[ \bar f:\power X \eqvsym (\powerp X) + \unit \longrightarrow X+\unit
  \]

  现在，对任意序数 $A$，可通过良基递归定义 $g_A : A \to X + \unit$：
  \[ g_A(a) \defeq
    \bar f\Big(X \setminus \setof{ g_A(b) | \strut (b<a) \wedge (g_A(b) \in X) }\Big)
  \]
  （显然地将 $X$ 视作 $X + \unit$ 的子集）。

  令 $A'\defeq \setof{a:A | g_A(a) \in X}$ 为 $X \subseteq X + \unit$ 的原像；则我们断言限制映射 $g_A' : A' \to X$ 是单射。
  因为若 $a, a' : A$ 且 $a \neq a'$，则由三分性且不失一般性，可假设 $a' < a$。
  于是 $g_A(a') \in \setof{ g_A(b) | b<a }$，再由对所有 $Y$ 都有 $f(Y) \in Y$ 便知 $g_A(a) \neq g_A(a')$。

  此外，$A'$ 是 $A$ 的一个初始段。
  因为 $g_A(a)$ 属于 \unit 当且仅当 $\setof{g_A(b)|b<a} = X$，且若此条件成立，则它对任意 $a' > a$ 亦成立。
  因此，$A'$ 本身也是一个序数。

  最后，由于 \ord 是一个序数，我们可以取 $A \defeq \ord$。
  令 $X'$ 为 $g_\ord' : \ord' \to X$ 的像；则 $g_\ord'$ 的逆给出一个单射 $H : X' \to \ord$。
  由 \cref{thm:ordunion}，存在一个序数 $C$，使得对所有 $x : X'$ 都有 $Hx \le C$。
  再由 \cref{thm:ordsucc}，存在另一个序数 $D$ 使得 $C < D$，从而对所有 $x : X'$ 都有 $Hx < D$。
  现在我们有
  \begin{align*}
    g_{\ord}(D) &= \bar f\Big( X \setminus \setof{ g_\ord(B) | \rule{0pt}{1em} B<D \wedge (g_\ord(B) \in X)} \Big)\\
    &=\bar f\Big( X \setminus \setof{ g_\ord(B) | \rule{0pt}{1em} g_\ord(B) \in X} \Big)
  \end{align*}
  因为若 $B : \ord$ 且 $g_\ord(B) \in X$，则存在某个 $x : X'$ 使得 $B = Hx$，从而 $B < D$。
  现在，如果
  \[\setof{ g_\ord(B) | \rule{0pt}{1em} g_\ord(B) \in X}\]
  不等于整个 $X$，则 $g_\ord(D)$ 将属于 $X$ 却不在此子集中，而由于 $D$ 本身也是 $B$ 的一个可能取值，这将导致矛盾。
  因此这个子集必须是整个 $X$，从而 $g_\ord'$ 不仅是单射，还是满射。
  于是，我们就可以将 $\ord'$ 上的序数结构迁移到 $X$ 上。
\end{proof}

\begin{rmk}
  如果我们之前对 \cref{thm:ordsucc} 或 \cref{thm:ordunion} 给出了错误的证明，那么由此得到的 \cref{thm:wop} 的证明将是无效的：因为那样将无法一致地指定宇宙层级\index{universe level}。
  事实上，我们需要命题大小重整（它可由 \LEM{} 推出）来确保 $X'$ 在等价意义下与 $X$ 处于同一宇宙。
\end{rmk}

\begin{cor}
  假设选择公理，函数 $\ord\to\set$（它忘记序结构）是一个满射。
\end{cor}

注意，\ord 是一个集合，而 \set 是一个 $1$-类型。
一般而言，没有理由指望 $1$-类型能接受来自某个集合的满射函数。
即使选择公理似乎也不蕴含\emph{每个} $1$-类型都能如此（尽管参见 \cref{ex:acnm-surjset}），但它可以直接推出这一点对于由 \set 构造出来的 $1$-类型是成立的，例如 \cref{sec:sip} 中结构范畴的对象类型。
下面的推论也适用于这类范畴。

\begin{cor}
  \index{weak equivalence!of precategories}%
  假设 \choice{}，\uset 接受一个来自严格范畴的弱等价函子。
\end{cor}

\begin{proof}
  令 $X_0\defeq \ord$，并且对于 $A,B:X_0$，令 $\hom_X(A,B) \defeq (A\to B)$。
  由于 \ord 是一个集合，$X$ 是一个严格范畴，且上述满射 $X_0 \to \set$ 可扩张为弱等价函子 $X\to \uset$。
\end{proof}

现在回顾 \cref{sec:cardinals}，我们还有一个进一步的满射 $\cd{\blank}:\set\to\card$，从而得到一个复合的满射 $\ord\to\card$，它把每个序数映到它的基数。

\begin{thm}
  假设 \choice{}，满射 $\ord\to\card$ 具有一个截面。
\end{thm}

\begin{proof}
  对此有一个简单但错误的证明：因为 \ord 和 \card 都是集合，\choice{} 蕴含它们之间的任何满射\emph{仅}存在一个截面。
  然而，我们实际上拥有一个\emph{具体的}典范截面：因为 \ord 是一个序数，其每个非空子集都有一个唯一确定的最小元。
  因此，我们可以将每个基数映射到其对应纤维中的那个最小元。
\end{proof}

集合论中的传统做法是将基数等同于它们在 \ord 中的像，即具有该基数的那个最小序数。

由此可知，\card 也典范地具有序数结构：事实上，是一个与 \ord 同构的结构。
具体来说，我们通过良基递归定义一个函数 $\aleph:\ord\to\ord$，使得 $\aleph(A)$ 是满足如下条件的最小序数：其基数大于所有 $a:A$ 对应的 $\aleph({\ordsl A a})$。
那么（在假设 \choice{} 下）$\aleph$ 的像恰好就是 \card 的像。

\index{denial|)}%

\index{ordinal|)}%

\section{累积层次}
\label{sec:cumulative-hierarchy}

\index{bargaining|(}%
对于给定的宇宙 $\UU$，我们可以将其中所有集合的累积层次 $V$ 定义为一个高阶归纳类型。这种定义方式使得 $V$ 本身也是一个集合（位于更大的宇宙 $\UU'$ 中），并配备一个满足集合论常见定律的二元“隶属”关系 $x\in y$。

\begin{defn}\label{defn:V}
  相对于类型宇宙 $\UU$ 的\define{累积层次}
  \indexdef{cumulative!hierarchy, set-theoretic}%
  \indexdef{hierarchy!cumulative, set-theoretic}%
  $V$ 是由以下构造子生成的高阶归纳类型。
  %
  \begin{enumerate}
  \item 对于每一个 $A : \UU$ 和 $f : A \to V$，有一个元素 $\vset(A, f) : V$。
  \item 对所有 $A, B : \UU$，以及 $f : A \to V$ 和 $g : B \to V$，若满足
    %
    \begin{narrowmultline} \label{eq:V-path}
      \big(\fall{a:A} \exis{b:B} \id[V]{f(a)}{g(b)}\big) \land \narrowbreak
      \big(\fall{b:B} \exis{a:A} \id[V]{f(a)}{g(b)}\big)
    \end{narrowmultline}
    %
    则存在一条路径 $\id[V]{\vset(A,f)}{\vset(B,g)}$。
  \item 0-截断构造子：对所有 $x,y:V$ 和 $p,q:x=y$，有 $p=q$。
  \end{enumerate}
\end{defn}

用集合论的语言来说，$\vset(A,f)$ 可以理解为经典集合论意义下 $A$ 在 $f$ 下的像所构成的集合，即 $\setof{ f(a) | a \in A }$。
但是，我们将避免使用这种记法，因为它会与我们用于子类型的记法冲突（不过可以参见下面的 \eqref{eq:class-notation} 和 \cref{def:TypeOfElements}）。

层次 $V$ 是从空映射 $\rec\emptyt(V) : \emptyt \to V$（由空类型的递归原理给出）起步的，它给出的空集是 $\emptyset = \vset(\emptyt,\rec\emptyt(V))$。
接着，单点集 $\{\emptyset\}$ 通过映射 $\unit \to V$（定义为 $\ttt \mapsto \emptyset$）进入 $V$，以此类推。
（通过添加额外的点构造子，该定义也可调整以包含任意由“原子”或“本元”构成的集合。）
类型 $V$ 与基础宇宙 $\UU$ 处于同一宇宙层级。

$V$ 的第二个构造子的形式与我们之前见过的任何构造子都不同：它不仅仅涉及 $V$ 中的路径（我们在 \cref{sec:hittruncations} 中曾说过这有点奇怪），还涉及对其求和的截断。
它当然不符合 \cref{sec:naturality} 中描述的一般模式，因此它的归纳原理应该是什么可能并不显然。
幸运的是，就像我们在 \cref{sec:hittruncations} 中对 0-截断的第一个定义那样，它可以利用辅助的高阶归纳类型重新表达。
我们将其细节留给读者探索（参见 \cref{ex:cumhierhit}）。

\index{induction principle!for cumulative hierarchy}%
最终，$V$ 的归纳原理（用模式匹配的语言写出来）是说：给定 $P:V\to \set$，为了构造 $h:\prd{x:V} P(x)$，只需给出以下内容。

\begin{enumerate}
\item 对于任意 $f:A\to V$，在假设已对所有 $a:A$ 给出 $h(f(a))$ 的前提下，构造 $h(\vset(A,f))$。
\item 验证若 $f : A \to V$ 和 $g : B \to V$ 满足~\eqref{eq:V-path}，则 $\dpath{P}{q}{h(\vset(A,f))}{h(\vset(B,g))}$，其中 $q$ 是由 $V$ 的第二个构造子及~\eqref{eq:V-path} 产生的路径；这里需归纳地假设 $h(f(a))$ 和 $h(g(b))$ 对所有 $a:A$ 和 $b:B$ 均有定义，并且满足以下条件：
\begin{eqnarray*}
    &       & \big(\fall{a:A} \exis{b:B} \exis{p:f(a)=g(b)} \dpath{P}{p}{h(f(a))}{h(g(b))}\big) \\
    & \land & \big(\fall{b:B} \exis{a:A} \exis{p:f(a)=g(b)} \dpath{P}{p}{h(f(a))}{h(g(b))}\big)
\end{eqnarray*}
\end{enumerate}

第二个条款检查正在定义的映射必须尊重在 \eqref{eq:V-path} 中引入的路径。
正如我们通常在使用模式匹配陈述高阶归纳原理时那样，这看起来像是同义反复，其实不然。
关键在于，“$h(f(a))$”本质上是一个我们无法窥探其内部的形式符号，而 $h(\vset(A,f))$ 必须依据它来定义。
因此，在第二个条款中，我们在适当的时候假设这些形式符号相等，然后验证由第一个条款的构造所产生的元素也相等。
当然，如果 $P$ 是一个命题族，那么第二个条款会自动满足。

注意到，通过归纳可知，对于每个 $v:V$，仅有 $A:\UU$ 和 $f:A\to V$ 使得 $v=\vset(A,f)$。
因此，尝试如下定义 $V$ 上的 \define{隶属关系}
\indexdef{membership, for cumulative hierarchy}%
$x\in v$ 是合理的：
%
% Note: "membership" rather than "elementhood", because "element" is taken.
%
\symlabel{V-membership}
\begin{equation*}
  (x \in \vset(A,f)) \defeq (\exis{a : A} x = f(a)).
\end{equation*}
%
要验证该定义是有效的，我们必须使用 $V$ 的递归原理。假设我们有一条通过~\eqref{eq:V-path} 构造的道路 $\vset(A, f) = \vset(B, g)$。若 $x \in \vset(A,f)$，则仅有 $a : A$ 使得 $x = f(a)$；但由~\eqref{eq:V-path}，仅有 $b : B$ 使得 $f(a) = g(b)$，因此 $x = g(b)$ 且 $x \in \vset(B,g)$。反之亦然。

$V$ 上的 \define{子集关系}
\indexdef{subset!relation on the cumulative hierarchy}%
$x\subseteq y$ 按通常方式定义为：
%
\begin{equation*}
  (x \subseteq y) \defeq \fall{z : V} z \in x \Rightarrow z \in y.
\end{equation*}

一个 \define{类}
\indexdef{class}%
可以视作 $V$ 上的一个纯命题谓词。我们说一个类 $C : V \to \prop$ 是一个 \define{$V$-集（$V$-set）}
\indexdef{set!in the cumulative hierarchy}%
，如果仅存在 $v:V$ 使得
%
\begin{equation*}
  \fall{x : V} C(x) \Leftrightarrow x \in v.
\end{equation*}
我们也可以使用类的常规记法，这与我们对于子类型的标准记法一致：
\begin{equation}
  \setof{ x | C(x) } \defeq \lam{x}C(x).\label{eq:class-notation}
\end{equation}
%
一个类 $C: V\to \prop$ 被称为 \define{$\UU$-小的（$\UU$-small）}
\indexdef{class!small}%
\indexdef{small!class}%
，如果它的所有值 $C(x)$ 都位于 $\UU$ 中，具体而言即 $C: V\to \prop_{\UU}$。
由于 $V$ 与构建它所基于的基础宇宙 $\UU$ 位于同一宇宙 $\UU'$ 中，这对于任意 $v,w:V$ 的相等类型 $v=_V w$ 同样成立。因此，为了在缺乏命题大小重整的情况下获得一个行为良好的理论，
\index{propositional!resizing}%
\index{resizing}%
拥有一个 $\UU$-小的相等关系的“重整”会很方便，我们可以如下通过归纳来定义它。

\begin{defn}\label{def:bisimulation}
  通过对 $V$ 进行双重归纳来定义 \define{双模拟}
  \indexdef{bisimulation}%
  关系
  %
  \begin{equation*}
    \mathord\bisim : V \times V \longrightarrow \prop_{\UU}
  \end{equation*}
  %
  ，其中对于 $\vset(A,f)$ 和 $\vset(B,g)$，我们令：
  \begin{narrowmultline*}
    \vset(A,f)  \bisim \vset(B,g) \defeq \narrowbreak
    \big(\fall{a:A}\exis{b:B} f(a)  \bisim g(b)\big) \land
    \big(\fall{b:B}\exis{a:A} f(a) \bisim g(b)\big).
  \end{narrowmultline*}
\end{defn}


%
为验证该定义是正确的，我们只需检查它与通过~\eqref{eq:V-path} 构造的道路 $\vset(A, f) = \vset(B, g)$ 相容，而这是显然的；此外还需验证 $\prop_{\UU}$ 是一个集合，而它确实是。注意，由构造可知 $u \bisim v$ 位于 $\propU$ 中。

\begin{lem}\label{lem:BisimEqualsId}
对于任意 $u,v:V$，我们有 $(u=_V v) = (u \bisim v)$。
\end{lem}

\begin{proof}
容易通过归纳证明 $\bisim$ 是自反的，因此通过传输我们有 $(u=_V v)\to (u \bisim v)$。
因此，只需再证 $(u \bisim v)\to (u=_V v)$。
通过对 $u$ 和 $v$ 进行归纳，我们可以假设它们分别是 $\vset(A,f)$ 和 $\vset(B,g)$。
（我们可以忽略 $V$ 的道路构造子，因为 $(u \bisim v)\to (u=_V v)$ 是一个命题。）
那么根据定义，$\vset(A,f)\bisim\vset(B,g)$ 蕴含 $(\fall{a:A}\exis{b:B}f(a)  \bisim g(b))$，反之亦然。
而归纳假设告诉我们 $(\fall{a:A}\exis{b:B}f(a) = g(b))$，反之亦然。
因此，由 $V$ 的道路构造子，我们有 $\vset(A,f) =_V \vset(B,g)$。
\end{proof}

人们也许认为可以省略 $V$ 的 0-截断构造子，并通过将 \cref{thm:h-set-refrel-in-paths-sets} 应用于双模拟来\emph{证明} $V$ 是 0-截断的。然而，在 \cref{lem:BisimEqualsId} 的证明中，我们使用了 $V$ 是 0-截断的这一事实，从而得出 $(u \bisim v)\to (u=_V v)$ 是一个命题，使得在归纳中只需假设 $u$ 和 $v$ 是 $\vset(A,f)$ 和 $\vset(B,g)$ 即可。

现在，我们可以利用重整后的相等关系来得到以下有用的原理。

\begin{lem}\label{lem:MonicSetPresent}
对于每个 $u:V$，都存在给定的 $A_u:\UU$ 和单射 $m_u: A_u \mono V$，使得 $u = \vset(A_u, m_u)$。
\end{lem}

\begin{proof}
  任取一个表示 $u = \vset(A,f)$，并将 $f:A\to V$ 分解为一个满射后接一个单射：
  %
  \begin{equation*}
    f = m_u\circ e_u : A \epi A_u \mono V.
  \end{equation*}
  %
  显然 $u = \vset(A_u, m_u)$，只要 $A_u$ 仍在 $\UU$ 中即可；而若 $e_u : A \epi A_u$ 的核在 $\UU$ 中，则此条件成立。但 $e_u : A \epi A_u$ 的核是 $V$ 上恒等映射沿 $f : A\to V$ 的拉回，而由我们刚刚所证，在等价意义下它是 $\UU$-小的。现在，从 $u:V$ 构造满足 $m_u :A_u \mono V$ 和 $u = \vset(A_u, m_u)$ 的对 $(A_u, m_u)$ 在 $V$ 之上等价意义下唯一，从而由泛等性知在相等意义下唯一。于是，由唯一选择原理~\eqref{cor:UC}，存在映射 $c : V\to\sm{A:\UU}(A\to V)$ 使得 $c(u) = (A_u, m_u)$，其中 $m_u :A_u \mono V$ 且 $u = \vset(c(u))$，如所断言。
\end{proof}

\begin{defn}\label{def:TypeOfElements}
  对于 $u:V$，刚才构造的满足 $u = \vset(A_u, m_u)$ 的单的表示 $m_u: A_u \mono V$ 可以称为 $u$ 的 \define{成员类型}
  \indexdef{type!of members}%
  ，并记作 $m_u : [u] \mono V$，甚至简记为 $[u] \mono V$。我们可以将 $[u]$ 视作“由 $u$ 的成员组成的 $V$ 的子类”。
\end{defn}

\begin{thm}\label{thm:VisCST}
  \index{axiom!of set theory, for the cumulative hierarchy}%
  对于 $(V, {\in})$，以下性质成立：
  %
  \begin{enumerate}
  \item \emph{外延性（extensionality）:}
    %
    \begin{equation*}
      \fall{x, y : V} x \subseteq y \land y \subseteq x \Leftrightarrow x = y.
    \end{equation*}
    %
     \item \emph{空集（empty set）:} 对所有 $x:V$，均有 $\neg (x\in \emptyset)$。
    %
    \item \emph{无序对（pairing）:} 对所有 $u, v:V$，类 $\{u, v\} \defeq \setof{ x | x = u \vee x = v}$ 是一个 $V$-集。
      %
    \item \emph{无穷（infinity）:}\index{axiom!of infinity} 存在一个 $v:V$，满足 $\emptyset\in v$，且 $x\in v$ 蕴含 $x\cup \{x\}\in v$。
    %
  \item \emph{并集（union）:} 对所有 $v:V$，类 $\cup v\defeq \setof{ x | \exis{u:V} x \in u \in v}$ 是一个 $V$-集。
    %
    \item \emph{函数集（function set）:} 对所有 $u, v:V$，类 $v^u \defeq \setof{ x | x : u\to v}$ 是一个 $V$-集。%
      \footnote{此处 $x:u\to v$ 意指 $x$ 是一个适当的、由有序对所组成的集合，这是集合论中编码函数的惯常方式。}
    %
   \item \emph{$\in$-归纳（$\in$-induction）:} 如果 $C : V \to \prop$ 是一个类，使得只要对所有 $x\in a$ 都有 $C(x)$ 成立，则 $C(a)$ 便成立，那么对所有 $v:V$，均有 $C(v)$。
   %
     \item \emph{替换（replacement）:}\index{axiom!of replacement} 给定任意 $r : V \to V$ 和 $x : V$，类
       %
       \begin{equation*}
         \setof{ y | \exis{z : V} z \in x \land y = r(z)}
       \end{equation*}
       %
       是一个 $V$-集。
  %
   \item \emph{分离（separation）:}\index{axiom!of separation} 给定任意 $a : V$ 和 $\UU$-小的 $C : V \to \propU$，类
     %
     \begin{equation*}
       \setof{ x | x \in a \land C(x)}
     \end{equation*}
     %
     是一个 $V$-集。
  \end{enumerate}
\end{thm}

\begin{proof}[证明概要]
  \mbox{}
  %
  \begin{enumerate}
  \item 外延性（extensionality）：如果 $\vset(A,f) \subseteq \vset(B, g)$，那么对每个 $a : A$ 都有 $f(a) \in \vset(B, g)$，因此对每个 $a : A$ 仅有 $b:B$ 使得
    $f(a) = g(b)$。由假设 $\vset(B, g) \subseteq \vset(A, f)$ 可得~\eqref{eq:V-path} 的另一半，因此 $\vset(A,f) = \vset(B,g)$。

  \item 空集（empty set）：假设 $x\in \emptyset = \vset(\emptyt,\rec\emptyt(V))$。那么 $\exis{a:\emptyt}x=\, \rec\emptyt(V,a)$，这是荒谬的。

  \item 无序对（pairing）：给定 $u$ 和 $v$，令 $w=\vset(\bool,\rec\bool(V,u,v))$。
    \index{pair!unordered}

  \item 无穷（infinity）：取 $w = \vset(\nat,I)$，其中 $I: \nat \to V$ 由递归式 $I(0) \defeq \emptyset$ 及 $I(n+1) \defeq I(n)\cup \{I(n)\}$ 给出。

  \item 并集（union）：任取 $v:V$ 及其任意一个表示 $f :A\to V$ 满足 $v=\vset(A,f)$。然后令 $\tilde{A} \defeq \sm{a:A}[fa]$，其中 $m_{fa} : [fa] \mono V$ 是来自 \cref{def:TypeOfElements} 的成员类型。显然 $\tilde{A}$ 是 $\UU$-小的，并且我们有 $\cup v \defeq \vset(\tilde{A}, \lam{x} m_{f(\proj1(x))}(\proj2(x)))$。

  \item 函数集（function set）：给定 $u, v:V$，取成员类型 $[u] \mono V$ 和 $[v] \mono V$，以及函数类型 $[u]\to [v]$。我们想定义一个映射
  \[
 r: ([u]\to [v])\ \longrightarrow\ V
  \]
   使其满足“$r(f) = \setof{ \pairr{x, f(x)} | x : [u] }$”；但为使该式有意义，我们必须先定义有序对 $\pairr{x, y}$，然后取映射 $r': x \mapsto \pairr{x, f(x)}$，接着便可令 $r(f)\defeq \vset([u], r')$。我们可按惯常的方式，用无序对来定义有序对。

  \item $\in$-归纳（$\in$-induction）：令 $C : V \to \prop$ 为一个类，使得只要对所有 $x\in a$ 有 $C(x)$ 成立，则 $C(a)$ 成立；并任取 $v=\vset(B,g)$。为了通过归纳证明 $C(v)$，假设对所有 $b:B$ 有 $C(g(b))$。对任意 $x\in v$，仅有某个 $b:B$ 使得 $x = g(b)$，因此 $C(x)$ 成立。从而 $C(v)$ 得证。

  \item 替换（replacement）：用 $C$ 表示所讨论的类。
    断言“$C$ 是一个 $V$-集”是一个命题，因此我们可以如下进行归纳证明。
    假设 $x$ 是 $\vset(A, f)$，我们断言 $w \defeq \vset(A, r \circ f)$ 即为我们所寻找的集合。
    若 $C(y)$，则仅有 $z : V$ 和 $a : A$ 使得 $z = f(a)$ 且 $y = r(z)$，因此 $y \in w$。
    反之，若 $y \in w$，则仅有 $a : A$ 使得 $y = r(f(a))$，所以若取 $z \defeq f(a)$，便知 $C(y)$ 成立。

  \item 我们称一个类 $C: V\to\prop$ 是 \define{可分的}
    \indexdef{class!separable}%
    \indexdef{separable class}%
    如果对任意 $a:V$，类
  %
  \symlabel{class-intersection}
  \begin{equation*}
    a \cap C \defeq\setof{x | x\in a \wedge C(x)}
  \end{equation*}
  %
  是一个 $V$-集。
我们需要证明任何 $\UU$-小的 $C: V \to \propU$ 都是可分的。实际上，给定 $a=\vset(A,f)$，令 $A' = \sm{x:A}C(fx)$，并取 $f' = f\circ i$，其中 $i : A' \to A$ 是显然的包含映射。然后我们可以取 $a' = \vset(A',f')$，并且 $x\in a\wedge C(x) \Leftrightarrow x\in a'$ 如所断言。我们需要 $C$ 取值于 $\UU$ 的假设，以便 $A' = \sm{x:A}C(fx)$ 能落在 $\UU$ 中。\qedhere
\end{enumerate}
\end{proof}

另外，拥有一个关于可分性的严格句法判据也很方便，这样从类的表达式就能直接看出它产生的是一个 $V$-集。一个熟悉的此类条件就是“$\Delta_0$”，它意味着表达式仅由相等关系 $x=_V y$ 和隶属关系 $x\in y$ 经由纯命题联结词 $\neg$、$\land$、$\lor$、$\Rightarrow$ 以及关于特定集合的量词（即形如 $\exists(x\in a)$ 和 $\forall(y\in b)$ 的量词，这些称为\define{有界的（bounded）}量词\index{bounded!quantifier}\index{quantifier!bounded}）构造而成。\indexdef{separation!.Delta0@$\Delta_0$}%

\begin{cor}\label{cor:Delta0sep}
如果类 $C: V \to \prop$ 是上述意义上的 $\Delta_0$ 类，那么它是可分的。
\end{cor}


\index{axiom!of $\Delta_0$-separation}%

\begin{proof}
回顾一下，$x = y$ 有一个 $\UU$-小的调整 $x \bisim y$。由于 $x\in y$ 是通过 $x=y$ 定义的，我们同样有隶属关系的一个 $\UU$-小的调整
%
\symlabel{resized-membership}
\begin{equation*}
  x\bin\vset(A,f) \defeq \exis{a:A} x \bisim f(a).
\end{equation*}
%
现在，令 $\Phi$ 为 $C$ 的一个 $\Delta_0$ 表达式，使得作为类有 $\Phi = C$（严格来说，我们应当区分表达式与其含义，但这里我们模糊这一区别）。令 $\widetilde{\Phi}$ 为将其中所有 $=$ 和 $\in$ 的出现替换为调整后的对应物 $\bisim$ 和 $\bin$ 所得的结果。显然 $\widetilde{\Phi}$ 也表示 $C$，即对所有 $x:V$ 有 $\widetilde{\Phi}(x) \Leftrightarrow C(x)$，因而由泛等性得 $\widetilde{\Phi}=C$。现在只需证明 $\widetilde{\Phi}$ 是 $\UU$-小的，那么根据定理，它便是可分的。

我们通过对表达式的构造进行归纳来证明 $\widetilde{\Phi}$ 是 $\UU$-小的。基础情形是 $x \bisim y$ 和 $x\bin y$，它们已被调整到 $\UU$ 中。$\UU$ 显然在单纯命题运算（以及 $(-1)$-截断）下封闭，于是只剩下有界量词 $\exists(x\in a)$ 和 $\forall(y\in b)$ 需要检查。根据定义，
\begin{align*}
\exists(x\in a) P(x) &\defeq \Brck {\sm{x:V}(x\bin a \land P(x))},\\
\forall(y\in b) P(x) &\defeq  \prd{x:V}(x\bin a \to P(x)).
\end{align*}
考虑 $\brck {\sm{x:V}(x\bin a \land P(x))}$。根据归纳假设 $P(x)$ 是 $\UU$-小的，从而主体 $(x\bin a \land P(x))$ 是 $\UU$-小的，但对 $V$ 的量化不一定保持在 $\UU$ 内。然而，在当前情形下，我们可以将其替换为对 $a$ 的成员类型 $[a]\mono V$ 的量化，且不难证明
\begin{equation*}
  \sm{x:V}(x\bin a \land P(x)) = \sm{x:[a]} P(x).
\end{equation*}
由于 $[a]$ 和 $P(x)$ 都在 $\UU$ 中，右端确实仍在 $\UU$ 中。$\prd{x:V}(x\bin a \to P(x))$ 的情形类似，利用 $\prd{x:V}(x\bin a \to P(x)) = \prd{x:[a]}P(x)$ 即可。
\end{proof}

我们已经证明，在带有一个宇宙 $\UU$ 的类型论中，累积层次 $V$ 是一个包含许多标准公理的"构造性集合论"
\index{constructive!set theory}%
的模型。
然而，据我们所知，它缺少 \emph{强收集}
\index{axiom!strong collection}%
\index{collection!strong}%
\index{strong!collection}%
与 \emph{子集收集}
\index{axiom!subset collection}%
\index{collection!subset}%
\index{subset!collection}%
公理，而这两条公理包含在 \CZF{}~\cite{AczelCZF} 中。
在这种集合论到类型论的通常解释中，这两条公理是类似集囿的相等定义的推论；而在其他构造出的集合论模型中，强收集可能因其他原因而成立。
我们不知道这两条公理是否在我们的模型 $(V,\in)$ 中成立，但似乎不太可能。
由于 $V$ 是系统\emph{内}的高阶归纳类型，而非\emph{外部}构造，它在某些方面与以往的解释不同并不令人意外。

最后，考虑将集合的选择公理以 $\choice{}$ 的形式（来自上文的 \cref{subsec:emacinsets}）添加到我们的类型论中。根据 \cref{thm:1surj_to_surj_to_pem}，可得 $\LEM{}$ 也成立，进而由 \cref{thm:settopos} 知 $\set$ 是一个带有子对象分类子 $\bool$ 的意象\index{topos}。此时，有 $\prop = \bool:\UU$，因而\emph{所有类都是可分的}。
于是我们证明了：

\begin{lem}\label{lem:fullsep}
  在带有 $\choice{}$ 的类型论中，$V$ 满足 \define{（完全）分离}
  \indexdef{separation!full}%
  律：给定\emph{任意}类 $C : V \to \prop$ 和 $a : V$，类 $a \cap C$ 是一个 $V$-集。
\end{lem}

\begin{thm}\label{thm:zfc}
在带选择公理 $\choice{}$ 和宇宙 $\UU$ 的类型论中，累积层次 $V$ 是带选择公理的 Zermelo--Fraenkel\index{set theory!Zermelo--Fraenkel} 集合论（ZFC）的一个模型。
\end{thm}

\begin{proof}
我们已经具备了 \cref{thm:VisCST} 所列的所有公理，以及完全分离律，因此只需证明对每个 $a:V$ 都存在幂集\index{power set} $\power a:V$。
但由于我们有排中律 $\LEM{}$，这些幂集即为函数类型 $\power a = (a\to\bool)$。
于是 $V$ 便成为 Zermelo--Fraenkel 集合论（ZF）的一个模型。
从 $\choice{}$ 验证集合论的选择公理，我们留作一个简单的练习。
\end{proof}

\index{bargaining|)}%

%%%%%%%%%%%%%%%%%%%%%%%%%%%%%%%%%%%%%%%%%%%%%%%%%%%%%%%%%%%%%%%%%%%%%%
\sectionNotes

人们对集合范畴所期望的基本性质，可以追溯到初等意象理论的早期。
在 \cref{subsec:emacinsets} 中提及的\emph{集合范畴的初等理论（Elementary theory of the category of sets）}是由 Lawvere\index{Lawvere} 在 \cite{lawvere:etcs-long} 中引入的，作为对集合论的一种范畴论公理化。
\index{Elementary Theory of the Category of Sets}%
$\Pi W$-预意象（$\Pi W$-pretopos）的概念，被视为初等意象的一个直谓版本，是在~\cite{MoerdijkPalmgren2002} 中提出的；亦可参阅~\cite{palmgren:cetcs}。

\cref{sec:piw-pretopos} 中对集合范畴的处理大体遵循了~\cite{RijkeSpitters}。
满态射是满射这一事实（\cref{epis-surj}）在经典数学中广为人知，但要\emph{直谓地}证明它却并不像看起来那么平凡。
\index{数学!直谓的}%
\cite{Mines/R/R:1988} 中的证明用到了幂集运算（这一运算本身是非直谓的），不过，该证明也可以看作如下较弱命题的一个直谓证明：若一个映射在宇宙 $\UU_i$ 中是满态射，则它在下一个宇宙 $\UU_{i+1}$ 中是满射。
Wilander~\cite{Wilander2010} 给出了关于集囿（setoids）的一个直谓证明。
我们的证明与 Wilander 的类似，但通过使用推出和泛等性避开了集囿。

\cref{thm:1surj_to_surj_to_pem} 中从 $\choice{}$ 推出 $\LEM{}$ 的蕴含关系，是 Diaconescu~\cite{Diaconescu} 在意象理论中证明的一个定理在同伦类型论中的改编；Bishop~\cite[Problem~2]{Bishop1967} 早已将其作为一个问题提出。

关于序数的直觉主义理论，参见~\cite{taylor:ordinals,Taylor99} 以及 \cite{JoyalMoerdijk1995}。
在类型论中通过归纳原理（包括可达性\index{accessibility}这一归纳谓词）定义良基性的研究，见于~\cite{Huet80,Paulson86,Nordstrom88}，尽管这一思想可以追溯到 Gentzen 对算术一致性\index{一致性!算术的}的证明~\cite{Gentzen36}。

代数集合论（algebraic set theory）的思想启发了我们在 \cref{sec:cumulative-hierarchy} 中对累积层次的发展，这一思想归功于~\cite{JoyalMoerdijk1995}，但它源自更早的工作~\cite{AczelCZF}。
\index{algebraic set theory}%
\index{集合论!代数的}%

%%%%%%%%%%%%%%%%%%%%%%%%%%%%%%%%%%%%%%%%%%%%%%%%%%%%%%%%%%%%%%%%%%%%%%
\sectionExercises

\begin{ex}\label{ex:utype-ct}
  仿照 $\uset$ 的模式，我们希望构造一个由所有类型（在给定宇宙 $\UU$ 中）以及它们之间的映射组成的范畴 $\utype$。
  然而，为了使它成为 \cref{sec:cats} 意义上的范畴，我们必须先声明 $\hom(X,Y) \defeq \pizero{X\to Y}$，其上的复合由截断归纳从普通复合 $(Y\to Z) \to (X\to Y) \to (X\to Z)$ 来定义。
  这在 \cref{ct:hoprecat} 中被定义为\emph{类型的同伦预范畴（homotopy precategory of types）}。
  但它仍然不是一个范畴，而只是一个预范畴（其对象类型 $\UU$ 甚至连 $0$-类型都不是）。
  经过 Rezk 完备化
  \index{completion!Rezk}%
  （参见 \cref{ct:hocat}），它成为一个范畴，并且根据 \cref{ct:ex:hocat}，其对象类型可以等同于 $\trunc1\type$。
  证明所得到的范畴 $\utype$ 与 $\uset$ 不同，它不是一个预意象。
\end{ex}

\begin{ex}\label{ex:surjections-have-sections-impl-ac}
  证明：如果在范畴 $\uset$ 中每个满射都有一个截面，那么选择公理成立。
\end{ex}

\begin{ex}\label{ex:well-pointed}
  证明：在 $\LEM{}$ 下，范畴 $\uset$ 是良点的，
  \indexdef{category!well-pointed}%
  即以下陈述成立：对于任意 $f, g : A\to B$，如果 $f \neq g$，则存在函数 $a : 1\to A$ 使得 $f(a) \neq g(a)$。
  证明切片范畴
  \index{category!slice}%
  $\uset/\bool$（由函数 $A\to \bool$ 和交换三角形构成）不具备此性质。
  （提示：$\uset/\bool$ 中的终对象是恒等函数 $\bool \to \bool$，因此在这个范畴中，存在没有元素 $1\to X$ 的对象 $X$。）
\end{ex}

\begin{ex}\label{ex:add-ordinals}
  \index{addition!of ordinal numbers}%
  证明：如果 $(A,<_A)$ 和 $(B,<_B)$ 是良基的、外延的，或者是序数，那么 $A+B$ 同样如此，其中 $<$ 定义为
  \begin{align*}
    (a<a') &\defeq (a<_A a') & \text{for }& a,a':A\\
    (b<b') &\defeq (b<_B b') & \text{for }& b,b':B\\
    (a<b) &\defeq \unit      & \text{for }& (a:A),(b:B)\\
    (b<a) &\defeq \emptyt    & \text{for }& (a:A),(b:B).
  \end{align*}
\end{ex}


% \begin{proof}
%   We first prove by induction on $<_A$ that every element of $A$ is accessible in $A+B$.
%   This is easy since the only elements less than $a:A$ in $A+B$ are also in $A$.
%   We then prove by induction on $<_B$ that every element of $B$ is accessible in $A+B$.
%   This is easy since we have already proven that every element of $A$ is accessible.
% \end{proof}

\begin{ex}\label{ex:multiply-ordinals}
  \index{multiplication!of ordinal numbers}%
  证明：如果 $(A,<_A)$ 和 $(B,<_B)$ 是良基的、外延的，或者是序数，那么 $A\times B$ 同样如此，其中 $<$ 定义为
  \[ ((a,b) <(a',b')) \defeq (a<_A a') \vee ((a=a') \wedge (b<_B b')). \]
\end{ex}


% \begin{proof}
%   We prove by induction on $<_A$ that for every $a:A$, every element of the form $(a,b)$ is accessible in $A\times B$.
%   The inductive hypothesis is that for all $a'<_A a$, every pair $(a',b)$ is accessible.
%   Inside this induction, we prove by induction on $<_B$ that for every $b:B$, the element $(a,b)$ is accessible.
%   The nested inductive hypothesis is that for every $b'<_B b$, the element $(a,b')$ is accessible.
%   But now, if $(a',b')< (a,b)$, then either $a<_A a'$ in which case $(a',b')$ is accessible by the first inductive hypothesis, or $a=a'$ and $b'<_B b$, in which case $(a,b')$ is accessible by the second inductive hypothesis.
%   Thus, by definition of accessibility, $(a,b)$ is accessible.
%   This completes both inductions.
% \end{proof}

\begin{ex}\label{ex:algebraic-ordinals}
  在序数上定义通常的代数运算，并证明它们满足通常的性质。
\end{ex}

\begin{ex}\label{ex:prop-ord}
  注意 $\bool$ 是一个序数，它带有一个显然的 $<$ 关系，其中仅 $\bfalse<\btrue$ 成立。
  \begin{enumerate}
  \item 在 $\prop$ 上定义一个 $<$ 关系，使其成为一个序数。
  \item 证明 $\id[\ord]\bool\prop$ 当且仅当 \LEM{} 成立。
  \end{enumerate}
\end{ex}

\begin{ex}\label{ex:ninf-ord}
  回顾，当把 $\nat$ 视为序数时，我们称其为 $\omega$；因此我们也有序数 $\omega+1$。
  另一方面，我们定义
  \[ \nat_\infty \defeq \setof{a:\nat\to\bool | \fall{n:\nat} (a_n \le a_{\suc(n)}) } \]
  其中 $\le$ 表示 $\bool$ 上显然的偏序关系，满足 $\bfalse\le\btrue$。
  \begin{enumerate}
  \item 在 $\nat_\infty$ 上定义一个 $<$ 关系，使其成为一个序数。
  \item 证明 $\id[\ord]{\omega+1}{\nat_\infty}$ 当且仅当有限全知原理~\eqref{eq:lpo} 成立。%
    \index{limited principle of omniscience}%
  \end{enumerate}
\end{ex}

\begin{ex}\label{ex:well-founded-extensional-simulation}
  证明：若 $(A,<)$ 是良基且外延的，并且 $A:\UU$，则存在一个模拟 $A\to V$，其中 $(V,\in)$ 是从宇宙~\UU 构建的累积层次（见 \cref{sec:cumulative-hierarchy}）。
\end{ex}

\begin{ex}\label{ex:choice-function}
  证明 \cref{thm:wop}\ref{item:wop1} 等价于选择公理~\eqref{eq:ac}。
\end{ex}

\begin{ex}\label{ex:cumhierhit}
  给定类型 $A$ 和 $B$，定义一个\define{双完全关系}
  \indexsee{bitotal relation}{relation, bitotal}%
  \indexdef{relation!bitotal}%
  为 $R:A\to B\to \prop$，使得
  \[ \Big(\fall{a:A}\exis{b:B} R(a,b) \Big) \land \Big(\fall{b:B}\exis{a:A} R(a,b) \Big). \]
  对于这样的 $A$、$B$、$R$，令 $A\sqcup^R B$ 为由以下构造子生成的高阶归纳类型：
  \begin{itemize}
  \item $i:A\to A\sqcup^R B$，
  \item $j:B\to A\sqcup^R B$，
  \item 对每一对满足 $R(a,b)$ 的 $a:A$ 和 $b:B$，有一条道路 $i(a)=j(b)$。
  \end{itemize}
  证明累积层次 $V$ 可以用以下更直接的构造子列表来定义，并且所得的归纳原理就是 \cref{sec:cumulative-hierarchy} 中给出的那个：
  \begin{itemize}
  \item 对每个 $A : \UU$ 和 $f : A \to V$，有一个元素 $\vset(A, f) : V$。
  \item 对任意 $A,B:\UU$ 和双完全关系
    \index{relation!bitotal}%
    $R:A\to B\to \prop$，以及任意映射 $h:A\sqcup^R B \to V$，存在一条道路 $\id{\vset(A,h\circ i)}{\vset(B,h\circ j)}$。
  \item 0-截断构造子。
  \end{itemize}
\end{ex}

\begin{ex}\label{ex:strong-collection}
  在 \CZF 中，\define{强收集公理}
  \indexdef{axiom!strong collection}%
  \indexdef{collection!strong}%
  \indexdef{strong!collection}%
  的形式为：
   \begin{multline*}
   \Parens{\fall{x\in v}\exis{y} R(x,y)} \Rightarrow \\
   \exis{w}\big[\big(\fall{x\in v}\exis{y\in w}R(x,y)\big)\land \big(\fall{y\in w}\exis{x\in v}R(x,y) \big)\big]
   \end{multline*}
   它在累积层次 $V$ 中成立吗？（我们不知道这个问题的答案。）
\end{ex}

\begin{ex}\label{ex:choice-cumulative-hierarchy-choice}
验证：若假定 $\choice{}$，则累积层次 $V$ 满足通常集合论意义上的选择公理，该公理可以表述为：
  \[
   \fall{x:V} \Parens{(\fall{y\in x}\exis{z:V} z\in y) \Rightarrow  \exis{c\in(\cup x)^x}\fall{y\in x} c(y)\in y}
   \]
\end{ex}

\begin{ex}\label{ex:plump-ordinals}
  假定命题大小重整，证明存在一个单纯的谓词 $\mathsf{isPlump}:\ord\to\prop$，使得对于任意 $A:\ord$ 有
  \begin{multline*}\label{eq:plump}
    \mathsf{isPlump}(A) = \Parens{\fall{B<A} \mathsf{isPlump}(B)} \wedge\narrowbreak
    \Parens{\fall{C,B:\ord} C\le B < A \wedge \mathsf{isPlump}(C) \Rightarrow C < A}.
  \end{multline*}
  注意，$\mathsf{isPlump}$ 不能简单地通过关于 \ord 的良基归纳来定义；你必须使用一个不同的良基关系。
  我们说一个序数 $A$ 是\define{丰满的}~\cite{taylor:ordinals,Taylor99}，如果 $\mathsf{isPlump}(A)$。
  \index{plump!ordinal}\index{ordinal!plump}%
\end{ex}

\begin{ex}\label{ex:not-plump}
  证明 \LEM{} 等价于"所有序数都是丰满的"这一陈述。
\end{ex}

\begin{ex}\label{ex:plump-successor}
  定义序数 $A$ 的\define{丰满后继}\index{plump!successor}\indexdef{successor!plump}为
  \[ t(A) \defeq \setof{ B:\ord | (B\le A) \wedge \mathsf{isPlump}(B) } \]
  \begin{enumerate}
  \item 根据定义，$t(A)$ 属于下一个更高的宇宙。证明：在假定命题大小重整的条件下，它等于一个与 $A$ 在同一宇宙中的序数。
  \item 再次假定命题大小重整，证明如果 $A$ 是丰满的（\cref{ex:plump-ordinals}），那么 $t(A)$ 也是丰满的。
  \end{enumerate}
\end{ex}

\begin{ex}\label{ex:ZF-algebras}
  相对于宇宙 $\UU_i$ 的一个\define{ZF-代数}~\cite{JoyalMoerdijk1995}
  \index{ZF-algebra}%
  是一个偏序集（见 \cref{ct:orders}）$V:\UU_{i+1}$，它具有所有以 $\UU_i$ 中类型为指标的上确界，并且配备了一个"后继"函数 $s:V\to V$（不一定在任何意义上保持 $\le$）。
  \begin{enumerate}
  \item 证明累积层次 $(V_{\UU_i},\subseteq,s)$ 是初始 ZF-代数，其中 $s(x)$ 是单点集 $\setof{x}$。
  \item 证明 $(\ord_{\UU_i},\le,s)$ 是满足"对所有 $x$ 有 $x\le s(x)$"这一性质的初始 ZF-代数，其中 $s(A)=A+\unit$ 是来自 \cref{thm:ordsucc} 的后继\index{successor!of an ordinal}。
  \item 假定命题大小重整，证明 $\Parens{\setof{A:\ord_{\UU_i} | \mathsf{isPlump}(A) },\le,t}$ 是满足"对所有 $x,y$ 有 $(x\le y) \Rightarrow (t(x)\le t(y))$"这一性质的初始 ZF-代数，其中 $t$ 是来自 \cref{ex:plump-successor} 的丰满后继。
  \end{enumerate}
\end{ex}

\begin{ex}\label{ex:monos-are-split-monos-iff-LEM-holds}
  对于范畴 $A$，态射 $f: \hom_A(a,b)$ 被称为是一个\define{可分裂单态射}，如果存在一个态射 $g: \hom_A(b,a)$ 使得 $\id{g \circ f}{1_a}$。（这样的 $g$ 称为 $f$ 的一个\define{收缩}。）
  证明以下陈述逻辑等价：
  \begin{enumerate}
  \item $\LEM{}$。
  \item 对所有集合 $A$ 和 $B$，如果 $A$ 是有实例的，那么对于 $\uset$ 中的每一个单态射 $f: A \to B$，$f$ 在 $\uset$ 中也是可分裂单态射。
  \end{enumerate}
\end{ex}

\index{set|)}%

% Local Variables:
% TeX-master: "hott-online"
% End:

\chapter{实数}
\label{cha:real-numbers}

\index{real numbers|(}%
任何一个名副其实的数学基础，最终都必须能够构造出数学分析中所理解的那种实数，也就是作为一个完备的 Archimedean 有序域。
\index{ordered field}%
完备性有两种概念。Cauchy（柯西）的版本要求实数对柯西序列\index{Cauchy!sequence}的极限封闭，而更强的 Dedekind版本则要求对Dedekind 分割\index{cut!Dedekind}封闭。
这引出了两种构造实数的方法，我们将分别在\cref{sec:dedekind-reals}和\cref{sec:cauchy-reals}中研究。在\cref{RD-final-field,RC-initial-Cauchy-complete}中，我们用泛性质来刻画这两种构造：Dedekind 实数是终 Archimedean 有序域，而 Cauchy 实数是始 Cauchy 完备 Archimedean 有序域。

在传统的构造性数学\index{mathematics!constructive}中，实数似乎总是需要某些妥协。
例如，Dedekind 实数与幂集或某些其他形式的非直谓性配合得更好，而 Cauchy 实数在可数选择公理\index{axiom!of choice!countable}成立时表现良好。
然而，我们给出了 Cauchy 实数的一种新构造，将它定义为一个高阶归纳-归纳类型；这似乎提供了第三种可能，它既不需要幂集，也不需要可数选择公理。

在~\cref{sec:comp-cauchy-dedek}中，我们比较这两种实数构造。
Cauchy 实数包含在Dedekind 实数中。
若排中律或可数选择公理成立，则二者一致；但在一般情况下，这一包含可能是真包含。

在~\cref{sec:compactness-interval}中，我们考察闭区间~$[0,1]$ 的三种紧性概念。
我们首先证明 $[0,1]$ 是\emph{度量紧}\indexdef{metricly compact}\indexdef{compactness!metric}的，意思是它既完备又完全有界，并且度量紧空间上的一致连续映射\index{uniformly continuous}的行为符合预期。
与之形成对比的是，Bolzano（波尔查诺）--Weierstra\ss{}（魏尔斯特拉斯）性质——即每个序列都有收敛子列——蕴涵有限全知原理，而后者是排中律的一个特例。
最后，我们讨论 Heine（海涅）--Borel（博雷尔）紧性。
有限子覆盖性质的朴素表述行不通，但一种证明相关的归纳覆盖\index{inductive cover}概念则可行。
本节内容基本属于标准的构造性分析。

在同伦类型论中发展起来的实数与分析理论，可以轻松地与经典数学兼容。
只要假设排中律与选择公理，我们便得到标准的经典分析\index{mathematics!classical}\index{classical!analysis}：此时Dedekind 实数与 Cauchy 实数一致，关于Dedekind 实数非直谓本质的基础问题随之消失，区间也尽可能地紧了。

在 \cref{sec:surreals} 中，我们通过将 Conway（康威）的超现实数构造为高阶归纳-归纳类型来结束本章；
这一构造在泛等类型论中比在经典集合论中更为自然。

除了 \cref{cha:basics,cha:logic} 的基本理论之外，如上所述，我们对 Cauchy 实数和超现实数使用了``高阶归纳-归纳类型''：它们结合了 \cref{cha:hits} 的思想与 \cref{sec:generalizations} 中提到的归纳-归纳类型概念。
我们也将频繁使用 \cref{subsec:prop-trunc} 中描述的传统逻辑记号，以及（在 \cref{sec:piw-pretopos} 中证明的）事实：我们的``集合''的行为方式符合预期。

注意，圆周的万有覆叠的全空间（它在 \cref{subsec:pi1s1-homotopy-theory} 中扮演了类似于经典代数拓扑中``实数''的角色）\emph{不是}我们要找的实数类型。
该类型是可缩的，因此等价于单点类型，所以它不能配备非平凡的代数结构。

\section{有理数域}
\label{sec:field-rati-numb}

\indexdef{rational numbers}%
\indexsee{number!rational}{rational numbers}%
我们首先构造有理数 \Q，因为实数可以看作 \Q 的完备化。
专家会指出，\Q 可以被任何近似域替代，
\indexdef{field!approximate}%
即 \Q 的一个子环，其中存在任意精确的近似逆
\index{inverse!approximate}%
。
一个例子是2进有理数环，
\index{rational numbers!dyadic}%
即形如 $n/2^k$ 的有理数。
如果我们在计算机上实现构造性数学，近似域会更合适，但我们将这种精细之处留给那些在意~$\pi$ 的数字的人。

我们在 \cref{sec:set-quotients} 中将整数 \Z 构造为 $\N\times \N$ 的一个商集，并观察到这个商集由一个幂等元生成。
在 \cref{sec:free-algebras} 中我们看到 \Z 是 \unit 上的自由群；类似地，可以证明它也是 \emptyt 上的自由交换环\index{ring}。
有理数域 \Q 的构造也遵循同样的思路，即作为商集
%
\[ \Q \defeq (\Z \times \N)/{\approx} \]
%
其中
\[ (u,a) \approx (v,b) \defeq (u (b + 1) = v (a + 1)). \]
%
换言之，对子 $(u, a)$ 代表有理数 $u / (1 + a)$。这里不可能出现除以零的情况，因为我们巧妙地给分母~$a$ 加了 1。
这里同样有典范的代表元选择方式，即最简分数。因此，我们可以应用 \cref{lem:quotient-when-canonical-representatives} 来得到一个集合 \Q，它同样具有可判定相等性。
\index{decidable!equality}%

我们不专门写出 \Q 上的算术运算，相信读者即使面对分母加了 1 的情形，也知道如何计算分数。
我们仅记录如下结论：有理数 \Q 的有序域构造是完全不成问题的，它具有可判定的相等性与可判定的序关系。
它也可刻画为初始有序域。
\index{initial!ordered field}%

\symlabel{positive-rationals}
\indexdef{rational numbers!positive}%
\indexdef{positive!rational numbers}%
最后，我们将用 $\Qp \defeq \setof{ q : \Q | q > 0 }$ 表示正有理数的类型。

\section{Dedekind 实数}
\label{sec:dedekind-reals}

\index{real numbers!Dedekind|(}%
我们首先回顾 Dedekind构造的基本思想。
我们采用双侧Dedekind 分割，而非常用的单侧版本，因为对称性使构造更为优雅，且在构造性数学与经典数学中均适用。
\index{mathematics!constructive}%
一个\emph{Dedekind 分割}\index{cut!Dedekind} 由 \Q 的一对子集 $(L, U)$ 组成，分别称为\emph{下}分割与\emph{上}分割，它们满足：
%

\begin{enumerate}
\item \emph{有实例的（inhabited）：} 存在 $q \in L$ 和 $r \in U$，
\item \emph{舍入的（rounded）：} $q \in L \Leftrightarrow \exis {r \in \Q} q < r \land r \in L$
  且 $r \in U \Leftrightarrow \exis {q \in \Q} q \in U \land q < r$，
  \index{rounded!Dedekind cut}
\item \emph{不相交的（disjoint）：} $\lnot (q \in L \land q \in U)$，以及
\item \emph{定位的（located）：} $q < r \Rightarrow q \in L \lor r \in U$。
  \index{locatedness}%
\end{enumerate}

%
从两个方向理解``舍入性''条件：自左向右读，它告诉我们分割是\emph{开的（open）}，
\index{open!cut}%
而从右向左读，它说明 $L$ 和 $U$ 分别是\emph{下集（lower set）}与\emph{上集（upper set）}。
定位性条件指出，$L$ 与 $U$ 之间不存在大的间隙。
因为分割总是开的，它们从不包含``中间的那个点''——即使该点为有理数时亦然。
一个典型的Dedekind 分割如下图所示：
%

\begin{center}
  \begin{tikzpicture}[x=\textwidth]
    \draw[<-),line width=0.75pt] (0,0) -- (0.297,0) node[anchor=south east]{$L\ $};
    \draw[(->,line width=0.75pt] (0.300, 0) node[anchor=south west]{$\ U$} -- (0.9, 0) ;
  \end{tikzpicture}
\end{center}

%
我们或许会朴素地将非形式化的定义翻译到类型论中，即认为一个分割就是一对映射 $L, U : \Q \to \prop$。但在 \cref{subsec:prop-subsets} 中我们已经看到，$\prop$ 是 $\prop_{\UU_i}$ 的一种带有\index{typical ambiguity}类型歧义的记法，其中 $\UU_i$ 是某个宇宙。
一旦我们使用某个特定的 $\UU_i$ 来定义分割，实数类型便位于下一个宇宙 $\UU_{i+1}$ 中，实数的一个性质位于更高两层的宇宙 $\UU_{i+2}$ 中，实数子集的一个性质位于 $\UU_{i+3}$ 中，依此类推。
原则上，我们完全可以跟踪这些\index{universe level}宇宙层级，尤其是在证明助手的帮助下；但在此处这样做只会带来我们宁可避免的繁琐事务。
因此，我们将做出一个简化假设：就我们的所有目的而言，一个单一的命题类型 $\Omega$ 就已足够。

事实上，Dedekind 实数的构造对逻辑层面的处理相当稳健。我们可以通过几种方式来理解这里使用单一类型 $\Omega$ 的做法：
%

\begin{enumerate}

\item 我们可以将 $\Omega$ 等同于那个有歧义的 $\prop$，并追踪定义和构造中出现的所有宇宙。

\item 我们可以假设命题大小重整公理，
  \index{propositional!resizing}%
  如 \cref{subsec:prop-subsets} 所示，
  它本质上将各层的 $\prop_{\UU_i}$ “坍缩”到最低的
  宇宙层级\index{universe level}，我们称其为 $\Omega$。

\item 一位对类型论宇宙或计算的复杂细节不感兴趣的经典数学家，可以直接假设命题的排中律~\eqref{eq:lem}，从而令 $\Omega \jdeq \bool$。
  \index{excluded middle}
  这不仅根除了关于 $\prop$ 的层级\index{universe level}问题，还将我们所做的一切转化为标准的经典\index{mathematics!classical}实数构造。

\item 在谱系的另一端，我们或许会问：让这些构造得以成立的最低要求是什么。一个谓词成为Dedekind 分割的条件，可以只用合取、析取以及对可数集~\Q 的存在量词\index{quantifier!existential}来表达。因此，我们可以取 $\Omega$ 为初始 \emph{$\sigma$-框架}，
  \index{initial!sigma-frame@$\sigma$-frame}%
  \index{sigma-frame@$\sigma$-frame!initial|defstyle}%
  即一个带有可数并\index{join!in a lattice}的格\index{lattice}，其中二元交运算对可数并满足分配律。（初始 $\sigma$-框架不可能是双点格 $\bool$，因为除非假设排中律，否则 $\bool$ 对可数并不封闭。）这将导致把~$\Omega$ 构造为一个高阶归纳-归纳类型，但像 \cref{sec:cauchy-reals} 那样的一个实验已经足够了。
\end{enumerate}

在上述所有情况下，$\Omega$ 都是一个集合。
%
事不宜迟，我们现在就将这个非形式的定义翻译为类型论的形式。在本章中，我们使用的逻辑记号遵循 \cref{defn:logical-notation}。

\begin{defn} \label{defn:dedekind-reals}
  一个 \define{Dedekind 分割（Dedekind cut）}
  \indexsee{Dedekind!cut}{cut, Dedekind}%
  \indexdef{cut!Dedekind}%
  是一对谓词 $L : \Q \to \Omega$ 与 $U : \Q \to \Omega$ 构成的 $(L, U)$，满足以下条件：
  %
  \begin{enumerate}
  \item \label{defn:dedekind-reals-inhabited}
    \emph{有实例的（即有界的）：} $\exis{q : \Q} L(q)$ 且 $\exis{r : \Q} U(r)$，
  \item \emph{舍入的（rounded）：} 对所有 $q, r : \Q$，
    \index{rounded!Dedekind cut}
    %
    \begin{align*}
      L(q) &\Leftrightarrow \exis{r : \Q} (q < r) \land L(r)
      \qquad\text{and}\\
      U(r) &\Leftrightarrow \exis{q : \Q} (q < r) \land U(q),
    \end{align*}
  \item \emph{不相交的（disjoint）：} $\lnot (L(q) \land U(q))$ 对所有 $q : \Q$，
  \item \emph{定位的（located）：} $(q < r) \Rightarrow L(q) \lor U(r)$ 对所有 $q, r : \Q$。
  \index{locatedness}%
  \end{enumerate}
  %
  我们用 $\dcut(L, U)$ 表示这些条件的合取。\define{Dedekind 实数（Dedekind reals）}的类型定义为
  \indexsee{Dedekind!real numbers}{real numbers, De\-de\-kind}%
  \indexdef{real numbers!Dedekind}%
  %
  \begin{equation*}
    \RD \defeq \setof{ (L, U) : (\Q \to \Omega) \times (\Q \to \Omega) | \dcut(L,U)}.
  \end{equation*}
\end{defn}

显然 $\dcut(L, U)$ 是一个命题，并且由于 $\Q \to \Omega$ 是一个集合，Dedekind 实数也构成一个集合。关于Dedekind 分割的变体（这些变体会导出扩充实数、下实数、上实数以及区间域），请参见
\cref{ex:RD-extended-reals,ex:RD-lower-cuts,ex:RD-interval-arithmetic}。

存在一个嵌入 $\Q \to \RD$，它将每个有理数 $q : \Q$ 与分割 $(L_q, U_q)$ 关联起来，其中
%
\begin{equation*}
  L_q(r) \defeq (r < q)
  \qquad\text{and}\qquad
  U_q(r) \defeq (q < r).
\end{equation*}
%
对于与有理数关联的分割 $(L_q, U_q)$，我们简记为 $q$。

\subsection{Dedekind 实数的代数结构}
\label{sec:algebr-struct-dedek}

Dedekind 实数的代数结构及序结构的构造，在直觉主义逻辑中照常进行。我们不再赘述细节，而是指出经典\index{mathematics!classical}构造与直觉主义构造之间的差异。用 $L_x$ 和 $U_x$ 表示实数 $x : \RD$ 的下分割与上分割，我们将加法定义为%
%
\indexdef{addition!of Dedekind reals}%
\begin{align*}
  L_{x + y}(q) &\defeq \exis{r, s : \Q} L_x(r) \land L_y(s) \land q = r + s, \\
  U_{x + y}(q) &\defeq \exis{r, s : \Q} U_x(r) \land U_y(s) \land q = r + s,
\end{align*}
%
加法逆元则定义为
%
\begin{align*}
  L_{-x}(q) &\defeq \exis{r : \Q} U_x(r) \land q = - r, \\
  U_{-x}(q) &\defeq \exis{r : \Q} L_x(r) \land q = - r.
\end{align*}
%
在这些运算下，$(\RD, 0, {+}, {-})$ 构成一个交换\index{group!abelian}群。乘法的定义则略显繁琐：
%
\indexdef{multiplication!of Dedekind reals}%
\begin{align*}
  L_{x \cdot y}(q) &\defeq
  \begin{aligned}[t]
    \exis{a, b, c, d : \Q} & L_x(a) \land U_x(b) \land L_y(c) \land U_y(d) \land {}\\
                           & \qquad q < \min (a \cdot c, a \cdot d, b \cdot c, b \cdot d),
  \end{aligned} \\
  U_{x \cdot y}(q) &\defeq
  \begin{aligned}[t]
    \exis{a, b, c, d : \Q} & L_x(a) \land U_x(b) \land L_y(c) \land U_y(d) \land {}\\
                           & \qquad \max (a \cdot c, a \cdot d, b \cdot c, b \cdot d) < q.
  \end{aligned}
\end{align*}
%
\index{interval!arithmetic}%
这些公式与区间算术中的区间乘法有关。在区间算术中，端点均为有理数的区间 $[a,b]$ 和 $[c,d]$ 相乘得到区间
%
\begin{equation*}
  [a,b] \cdot [c,d] =
  [\min(a c, a d, b c, b d), \max(a c, a d, b c, b d)].
\end{equation*}
%
例如，下分割的公式可以理解为：当存在分别包含 $x$ 和 $y$ 的区间 $[a,b]$ 和 $[c,d]$，使得 $q$ 位于 $[a,b] \cdot [c,d]$ 的左侧时，就有 $q < x \cdot y$。一般而言，将一个满足 $L_x(a)$ 和 $U_x(b)$ 的区间 $[a,b]$ 视为~$x$ 的一种近似是很有用的，参见
\cref{ex:RD-interval-arithmetic}。

我们现在已经得到了一个含幺交换环\index{ring}
\index{unit!of a ring}%
$(\RD, 0, 1, {+}, {-}, {\cdot})$。为了处理乘法逆元，我们必须先引入序。定义 $\leq$ 和 $<$ 为
%
\begin{align*}
  (x \leq y) &\ \defeq \ \fall{q : \Q} L_x(q) \Rightarrow L_y(q), \\
  (x < y)    &\ \defeq \ \exis{q : \Q} U_x(q) \land L_y(q).
\end{align*}

\begin{lem} \label{dedekind-in-cut-as-le}
  对于所有 $x : \RD$ 和 $q : \Q$，$L_x(q) \Leftrightarrow (q < x)$ 且 $U_x(q)
  \Leftrightarrow (x < q)$。
\end{lem}

\begin{proof}
  如果 $L_x(q)$，则由舍入性，仅存在 $r > q$ 使得 $L_x(r)$ 成立；而 $U_q(r)$ 成立，由此可得 $q < x$。
  反之，如果 $q < x$，则存在 $r : \Q$ 使得 $U_q(r)$ 和 $L_x(r)$ 同时成立；又因为 $L_x$ 是下集，所以 $L_x(q)$ 成立。
  另一半证明是对称的。
\end{proof}

\index{partial order}%
\index{transitivity!of . for reals@of $<$ for reals}
\index{transitivity!of . for reals@of $\leq$ for reals}
\index{relation!irreflexive}
\index{irreflexivity!of . for reals@of $<$ for reals}
关系 $\leq$ 是偏序，而 $<$ 传递且不自反。线性序
\index{order!linear}%
\index{linear order}%
%
\begin{equation*}
  (x < y) \lor (y \leq x)
\end{equation*}
%
在假设排中律时成立，但若不假设排中律，则得到弱线性序
%
\index{order!weakly linear}
\index{weakly linear order}
\begin{equation} \label{eq:RD-linear-order}
  (x < y) \Rightarrow (x < z) \lor (z < y).
\end{equation}
%
乍一看，或许不清楚~\eqref{eq:RD-linear-order} 与线性序有什么关系。但如果我们取 $x \jdeq u - \epsilon$ 和 $y \jdeq u + \epsilon$，其中 $\epsilon > 0$，则得到
%
\begin{equation*}
  (u - \epsilon < z) \lor (z < u + \epsilon).
\end{equation*}
%
这就是``在微小数值误差意义下''的线性序。换句话说，既然我们实际上无法期望以无限精度进行计算，那么 $<$ 只能在我们已计算出的有限精度内判定，也就不足为奇了。

为了证明~\eqref{eq:RD-linear-order} 成立，假设 $x < y$。则仅存在 $q : \Q$ 使得 $U_x(q)$ 且 $L_y(q)$。
由舍入性，仅存在 $r, s : \Q$ 使得 $r < q < s$，$U_x(r)$ 且 $L_y(s)$。
于是，由定位性可得 $L_z(r)$ 或 $U_z(s)$。
在第一种情形得到 $x < z$，在第二种情形得到 $z < y$。

在经典数学中，所有非零的数都存在乘法逆元。然而，在没有排中律的情况下，需要一个更强的条件。称两个实数 $x, y : \RD$ 是\define{分离的（apart）}
\indexdef{apartness}%
记作 $x \apart y$，当$(x < y) \lor (y < x)$：
%
\symlabel{apart}
\begin{equation*}
  (x \apart y) \defeq (x < y) \lor (y < x).
\end{equation*}
%
如果 $x \apart y$，则 $\lnot (x = y)$。
其逆命题在假设排中律时成立，但在构造性数学中不可证。
\index{mathematics!constructive}%
事实上，如果 $\lnot (x = y)$ 蕴含 $x \apart y$，则可推出一小部分排中律；参见\cref{ex:reals-apart-neq-MP}。

\begin{thm} \label{RD-inverse-apart-0}
  一个实数可逆，当且仅当，它与 $0$ 分离。
\end{thm}

\begin{rmk}
  我们注意到，一个实数可逆当且仅当它仅可逆。实际上，这在任何环\index{ring}中都成立，因为环是一个集合，而乘法逆元若存在则唯一。相关讨论参见\cref{cor:UC}之后的内容。
\end{rmk}

\begin{proof}
  假设 $x \cdot y = 1$。则仅存在 $a, b, c, d : \Q$ 使得 $a < x < b$，$c < y < d$ 且$0 < \min (a c, a d, b c, b d)$。由 $0 < a c$ 和 $0 < b c$ 可知，$a$、$b$ 和 $c$ 要么全为正，要么全为负。
  因此，要么 $0 < a < x$，要么 $x < b < 0$，于是 $x \apart 0$。

  反过来，如果 $x \apart 0$，则要么 $x > 0$，要么 $x < 0$。
  若 $x > 0$，我们如下定义 $x^{-1}$：
  %
  \begin{align*}
    L_{x^{-1}}(q) &\defeq
    (q > 0) \Rightarrow \exis{r : \Q} U_x(r) \land (q r < 1),
    \\
    U_{x^{-1}}(q) &\defeq
    (q > 0) \land \exis{r : \Q} L_x(r) \land (q r > 1).
  \end{align*}
  %
  若 $x < 0$，则如下定义：
  %
  \begin{align*}
    L_{x^{-1}}(q) &\defeq
    (q < 0) \land \exis{r : \Q} U_x(r) \land (q r > 1),
    \\
    U_{x^{-1}}(q) &\defeq
    (q < 0) \Rightarrow \exis{r : \Q} L_x(r) \land (q r < 1). \qedhere
  \end{align*}
\end{proof}

\index{ordered field!archimedean}%
\index{dense}%
\indexsee{order-dense}{dense}%
Archimedean 原理可以有多种陈述方式。我们发现，以``$\Q$ 在 $\RD$ 中稠密''这种形式来陈述，最具启发性。

\begin{thm}[$\RD$ 的 Archimedean 原理] \label{RD-archimedean}
  %
  对于所有 $x, y : \RD$，如果 $x < y$，则仅存在 $q : \Q$ 使得 $x < q < y$。
\end{thm}

\begin{proof}
  由 $<$ 的定义直接得出。
\end{proof}

在讨论Dedekind 实数的完备性之前，让我们先精确地说明它们所具有的代数结构。在下面的定义中，我们的目标不是追求公理化的最小化，而是希望获得实用的结构与性质。

\begin{defn} \label{ordered-field} 一个\define{有序域}
  \indexdef{ordered field}%
  \indexsee{field!ordered}{ordered field}%
  是一个集合 $F$，连同常量 $0$、$1$，运算 $+$、$-$、$\cdot$、$\min$、$\max$，以及纯粹关系 $\leq$、$<$、$\apart$，使得下列条件成立：
  %
  \begin{enumerate}
  \item $(F, 0, 1, {+}, {-}, {\cdot})$ 构成一个带单位的交换环；
    \index{unit!of a ring}%
    \index{ring}%
  \item $x : F$ 可逆当且仅当 $x \apart 0$；
  \item $(F, {\leq}, {\min}, {\max})$ 构成一个格；
  \item 严格序 $<$ 是传递的、不自反的，
    \index{relation!irreflexive}
    \index{irreflexivity!of . in a field@of $<$ in a field}%
    且弱线性的（$x < y \Rightarrow x < z \lor z < y$）；\index{transitivity!of . in a field@of $<$ in a field}
    \index{order!weakly linear}
    \index{weakly linear order}
    \index{strict!order}%
    \index{order!strict}%
  \item 分离关系 $\apart$ 是不自反的、对称的且协传递的（$x \apart y \Rightarrow x \apart z \lor y \apart z$）；
    \index{relation!irreflexive}
    \index{irreflexivity!of apartness}%
    \indexdef{relation!cotransitive}%
    \index{cotransitivity of apartness}%
  \item 对所有 $x, y, z : F$：
    %
    \begin{align*}
      x \leq y &\Leftrightarrow \lnot (y < x), &
      x < y \leq z &\Rightarrow x < z, \\
      x \apart y &\Leftrightarrow (x < y) \lor (y < x), &
      x \leq y < z &\Rightarrow x < z, \\
      x \leq y &\Leftrightarrow x + z \leq y + z, &
      x \leq y \land 0 \leq z &\Rightarrow x z \leq y z, \\
      x < y &\Leftrightarrow x + z < y + z, &
      0 < z \Rightarrow (x < y &\Leftrightarrow x z < y z), \\
      0 < x + y &\Rightarrow 0 < x \lor 0 < y, &
      0 &< 1.
    \end{align*}
  \end{enumerate}
  %
  每个这样的域都有一个标准嵌入 $\Q \to F$。一个有序域是\define{Archimedean 的}
  \indexdef{ordered field!archimedean}%
  \indexsee{archimedean property}{ordered field, archi\-mede\-an}%
  ，如果对所有 $x, y : F$，当 $x < y$ 时纯粹存在 $q :
  \Q$ 使得 $x < q < y$。
\end{defn}

\begin{thm} \label{RD-archimedean-ordered-field}
  Dedekind 实数构成一个有序的 Archimedean 域。
\end{thm}

\begin{proof}
  我们省略证明，希望迄今为止所展示的内容足以使这一定理显得可信。
\end{proof}

\subsection{Dedekind 实数是 Cauchy 完备的}
\label{sec:RD-cauchy-complete}

回忆一下，$x : \N \to \Q$ 称为一个\emph{柯西序列}\indexdef{Cauchy!sequence}，如果它满足
%
\begin{equation} \label{eq:cauchy-sequence}
  \prd{\epsilon : \Qp} \sm{n : \N} \prd{m, k \geq n} |x_m - x_k| < \epsilon.
\end{equation}
%
请注意，我们\emph{没有}截断内部的存在量词，因为我们实际上想要计算收敛速率——一个没有误差估计的逼近几乎不提供任何有用信息。由 \cref{thm:ttac}，\eqref{eq:cauchy-sequence} 给出一个函数 $M
: \Qp \to \N$，称为\emph{收敛模（modulus of convergence）}\indexdef{modulus!of convergence}，使得 $m, k \geq M(\epsilon)$
蕴含 $|x_m - x_k| < \epsilon$。由此可得，对所有 $\delta, \epsilon : \Qp$ 均有 $|x_{M(\delta/2)} - x_{M(\epsilon/2)}|<
\delta + \epsilon$。事实上，映射 $(\epsilon \mapsto
x_{M(\epsilon/2)}) : \Qp \to \Q$ 承载了与原始 Cauchy 条件~\eqref{eq:cauchy-sequence} 相同的关于极限的信息。后续我们将使用这些逼近函数而非柯西序列。

\begin{defn} \label{defn:cauchy-approximation}
  一个\define{Cauchy 逼近（Cauchy approximation）}
  \indexdef{Cauchy!approximation}%
  是一个映射 $x : \Qp \to \RD$，它满足
  %
  \begin{equation}
    \label{eq:cauchy-approx}
    \fall{\delta, \epsilon :\Qp} |x_\delta - x_\epsilon| < \delta + \epsilon.
  \end{equation}
  %
  一个 Cauchy 逼近 $x : \Qp \to \RD$ 的\define{极限（limit）}
  \index{limit!of a Cauchy approximation}%
  是一个数 $\ell : \RD$，使得
  %
  \begin{equation*}
    \fall{\epsilon, \theta : \Qp} |x_\epsilon - \ell| < \epsilon + \theta.
  \end{equation*}
\end{defn}

\begin{thm} \label{RD-cauchy-complete}
  $\RD$ 中的每个 Cauchy 逼近都有一个极限。
\end{thm}

\begin{proof}
  注意，我们这里证明的是极限的存在性，而不仅仅是纯粹存在性。
  给定一个 Cauchy 逼近 $x : \Qp \to \RD$，定义
  %
  \begin{align*}
    L_y(q) &\defeq \exis{\epsilon, \theta : \Qp} L_{x_\epsilon}(q + \epsilon + \theta),\\
    U_y(q) &\defeq \exis{\epsilon, \theta : \Qp} U_{x_\epsilon}(q - \epsilon - \theta).
  \end{align*}
  %
  显然，$L_y$ 和 $U_y$ 都是有实例且舍入的。互不相交性可由 Cauchy 逼近条件得出。为证明定位性，考虑任意
  满足 $q < r$ 的 $q, r : \Q$。存在 $\epsilon : \Qp$ 使得
  $4 \epsilon < r - q$。由于 $q + 2 \epsilon < r - 2 \epsilon$，纯粹地
  $L_{x_\epsilon}(q + 2 \epsilon)$ 或者 $U_{x_\epsilon}(r - 2 \epsilon)$。在第一种情况下，
  我们有 $L_y(q)$，而在第二种情况下有 $U_y(r)$。

  为了证明 $y$ 是 $x$ 的极限，考虑任意 $\epsilon, \theta : \Qp$。因为
  $\Q$ 在 $\RD$ 中稠密，所以纯粹存在 $q, r : \Q$ 使得
  %
  \begin{narrowmultline*}
    x_\epsilon - \epsilon - \theta < q < x_\epsilon - \epsilon - \theta/2
    < x_\epsilon < \\
    x_\epsilon + \epsilon + \theta/2 < r < x_\epsilon + \epsilon + \theta,
  \end{narrowmultline*}
  %
  因此 $q < y < r$。由此可得 $|y - x_\epsilon| < \epsilon + \theta$。
\end{proof}

为了完备性，我们也记录下经典的表述。

\begin{cor}
  假设 $x : \N \to \RD$ 满足 Cauchy 条件~\eqref{eq:cauchy-sequence}。则
  存在 $y : \RD$ 使得
  %
  \begin{equation*}
    \prd{\epsilon : \Qp} \sm{n : \N} \prd{m \geq n} |x_m - y| < \epsilon.
  \end{equation*}
\end{cor}

\begin{proof}
  由 \cref{thm:ttac}，存在 $M : \Qp \to \N$ 使得 $\bar{x}(\epsilon) \defeq
  x_{M(\epsilon/2)}$ 是一个 Cauchy 逼近。令 $y$ 为其极限（该极限的存在性由
  \cref{RD-cauchy-complete} 保证）。给定任意的 $\epsilon : \Qp$，令 $n \defeq M(\epsilon/4)$
  并观察到，对于任意 $m \geq n$，
  %
  \begin{narrowmultline*}
    |x_m - y| \leq |x_m - x_n| + |x_n - y| =
    |x_m - x_n| + |\bar{x}(\epsilon/2) - y| < \narrowbreak
    \epsilon/4 + \epsilon/2 + \epsilon/4 = \epsilon.\qedhere
  \end{narrowmultline*}
\end{proof}

\subsection{Dedekind 实数是Dedekind 完备的}
\label{sec:RD-dedekind-complete}

我们通过取 $\Q$ 上Dedekind 分割的类型得到了 $\RD$。但我们本也可以从任一 Archimedean 有序域 $F$ 出发，构造 $F$ 上的\index{cut!Dedekind}Dedekind 分割。这些分割同样构成一个 Archimedean 有序域 $\bar{F}$，即\define{$F$ 的Dedekind 完备化}，%
\index{completion!Dedekind}%
\indexsee{Dedekind!completion}{completion, Dedekind}
其中 $F$ 作为子域被包含。如果将这一构造应用于 $\RD$ 本身，我们会得到更多的实数吗？答案是否定的。事实上，我们将证明一个更强的结论：$\RD$ 是终极的。

一个有序域~$F$ 称为\define{对 $\Omega$ 可容许的}，
\indexsee{admissible!ordered field}{ordered field, admissible}%
\indexdef{ordered field!admissible}%
如果 $F$ 上的严格序 $<$ 是一个映射 ${<} : F \to F \to \Omega$。

\begin{thm} \label{RD-final-field}
  每个对 $\Omega$ 可容许的 Archimedean 有序域都是 $\RD$ 的子域。
\end{thm}

\begin{proof}
  设 $F$ 为一个 Archimedean 有序域。对每个 $x : F$，由下式定义 $L_x, U_x : \Q \to
  \Omega$：
  %
  \begin{equation*}
    L_x(q) \defeq (q < x)
    \qquad\text{and}\qquad
    U_x(q) \defeq (x < q).
  \end{equation*}
  %
  （此处用到了 $F$ 对 $\Omega$ 可容许的假设。）
  则 $(L_x, U_x)$ 是一个\index{cut!Dedekind}Dedekind 分割。事实上，由于 $F$ 是 Archimedean 的且 $<$ 具有传递性，该分割是有实例的且为舍入的；由于 $<$ 不自反，它是不交的；由于 $<$ 是弱线性序，它是定位的。令 $e : F \to \RD$ 为映射 $e(x) \defeq (L_x,
  U_x)$。

  我们断言 $e$ 是一个保持且反映序的域嵌入。首先注意到，对于有理数 $q$，有 $e(q) = q$。其次，对所有 $x, y : F$，有如下等价：
  %
  \begin{narrowmultline*}
    x < y \Leftrightarrow
    (\exis{q : \Q} x < q < y) \Leftrightarrow \narrowbreak
    (\exis{q : \Q} U_x(q) \land L_y(q)) \Leftrightarrow
    e(x) < e(y),
  \end{narrowmultline*}
  %
  因此 $e$ 确实保持并反映序关系。$e(x + y) = e(x) + e(y)$ 成立是因为，对所有 $q : \Q$：
  %
  \begin{equation*}
    q < x + y \Leftrightarrow
    \exis{r, s : \Q} r < x \land s < y \land q = r + s.
  \end{equation*}
  %
  从右到左的蕴含是显然的。至于另一个方向，若 $q < x + y$，则仅仅存在 $r : \Q$ 使得 $q - y < r < x$；取 $s \defeq q - r$，即得所需的 $r$ 和 $s$。我们将 $e$ 保持乘法的验证留作练习。
\end{proof}

为了证明 $\RD$ 上的Dedekind 分割不会给出任何新东西，我们只需再有一个引理。

\begin{lem} \label{lem:cuts-preserve-admissibility}
  如果 $F$ 对 $\Omega$ 可容许，那么它的Dedekind 完备化也对 $\Omega$ 可容许。
  \index{completion!Dedekind}%
\end{lem}

\begin{proof}
  设 $\bar{F}$ 为 $F$ 的Dedekind 完备化。$\bar{F}$ 上的严格序定义为
  %
  \begin{equation*}
    ((L,U) < (L',U')) \defeq \exis{q : \Q} U(q) \land L'(q).
  \end{equation*}
  %
  由于 $U(q)$ 和 $L'(q)$ 都是 $\Omega$ 的元素，只要 $\Omega$ 对合取与可数存在量化封闭——而我们从一开始就做了这一假设——引理即成立。
\end{proof}

\begin{cor} \label{RD-dedekind-complete}
  %
  \indexdef{complete!ordered field, Dedekind}%
  \indexdef{Dedekind!completeness}%
  Dedekind 实数是Dedekind 完备的：对每个取实数值的Dedekind 分割 $(L, U)$，都存在唯一的 $x : \RD$，使得对所有 $y : \RD$，有 $L(y) = (y < x)$ 且 $U(y) = (x < y)$。
\end{cor}

\begin{proof}
  由 \cref{lem:cuts-preserve-admissibility}，$\RD$ 的Dedekind 完备化 $\barRD$ 对 $\Omega$ 可容许，故由 \cref{RD-final-field} 可得嵌入 $\barRD \to \RD$，以及嵌入 $\RD \to \barRD$。但这两个嵌入必为同构，因为它们的复合是固定稠密子域~$\Q$ 的保序域\index{homomorphism!field}同态，从而就是恒等映射。由 $\barRD \to \RD$ 为同构，推论立即得证。
\end{proof}

\index{real numbers!Dedekind|)}%

\section{Cauchy 实数}
\label{sec:cauchy-reals}

\index{real numbers!Cauchy|(}%
\index{completion!Cauchy|(}%
\indexsee{Cauchy!completion}{completion, Cauchy}%
Cauchy 实数按其意图是 \Q 在柯西序列极限下的完备化。\index{Cauchy!sequence}
在 Cauchy 实数的经典构造中，我们考虑 \Q 中所有柯西序列的集合 $\mathcal{C}$，然后形成适当的商 $\mathcal{C}/{\approx}$。
然后，为证明 $\mathcal{C}/{\approx}$ 是 Cauchy 完备的，我们考虑一个柯西序列 $x : \N \to \mathcal{C}/{\approx}$，将其提升为一个序列的序列 $\bar{x} : \N \to \mathcal{C}$，并用 $\bar{x}$ 构造 $x$ 的极限。然而，将~$x$ 提升为 $\bar{x}$ 需要用到可数选择公理（即~\eqref{eq:ac} 在 $X=\N$ 时的特例）或排中律，而这是我们可能希望避免的。
\indexdef{axiom!of choice!countable}%
每个以取商为最后一步的实数构造都有此缺陷。
在构造性数学中，摆脱这一困境有三种常见途径：
\index{mathematics!constructive}%
%
\index{bargaining}%

\begin{enumerate}
\item 假装实数是一个 setoid $(\mathcal{C}, {\approx})$，即柯西序列的类型 $\mathcal{C}$ 附带一个通过行政命令加上去的重合\index{coincidence, of Cauchy approximations}关系。如此一来，一个实数序列就\emph{只不过}是一个代表这些实数的柯西序列的序列。
\item 向诱惑低头，接受可数选择公理。毕竟，该公理在大多数基于计算观点的构造性数学模型（如可实现性模型）中都是成立的。
\item 宣布 Cauchy 实数不足取，转而构造Dedekind 实数。在某些语境下，例如在构造性数学的层论模型中，这样的判决完全成立。然而，正如我们在 \cref{sec:dedekind-reals} 所见，构造性的Dedekind 实数也有其自身的问题。
\end{enumerate}

然而，利用高阶归纳类型，我们还有第四种解决方案。我们相信它比以上任何一种都更可取，甚至对经典数学家来说也很有趣。
这个想法是：Cauchy 实数应当是由~\Q 生成的\emph{自由完备度量空间}\index{free!complete metric space}。
一般来说，构造任何种类的自由结构，都需要将该结构的操作反复应用于生成元。
例如，集合 $X$ 上自由群的元素，并不仅仅是 $X$ 中元素的二元乘积和逆，还包括通过迭代乘积和求逆构造而得到的``词''。
因此，我们自然可以预期 Cauchy 完备化也是如此，这里相关的\emph{操作}就是\emph{取一个柯西序列的极限}。
（在这种情况下，迭代必须超限地进行，因为即使在无穷多步之后，也仍然会有新的柯西序列需要取极限。）

前文提到的论证表明，如果排中律或可数选择公理成立，那么 Cauchy 完备化会变得非常特殊：在构造一个空间的完备化时，只需在应用该操作\emph{一步}之后即可停止。
这可以类比于自由幺半群和自由群能够用（约化）词来显式描述的事实。
然而，我们在 \cref{sec:free-algebras} 看到，高阶归纳类型允许我们\emph{直接}构造自由结构，无论是否还有可用的显式描述。
在本节中我们将说明，这对于 Cauchy 实数同样成立（类似的技术可用于构造任何度量空间的 Cauchy 完备化；参见 \cref{ex:metric-completion}）。
具体来说，高阶归纳类型允许我们\emph{同时}添加柯西序列的极限并对重合关系作商，从而避免了将实数序列提升为代表元序列的问题。
\index{completion!Cauchy|)}%

\subsection{Cauchy 实数的构造}
\label{sec:constr-cauchy-reals}

将 Cauchy 实数 $\RC$ 构造为高阶归纳类型，比 \cref{sec:free-algebras} 中考虑的自由代数结构的构造更为微妙。
我们打算引入一个``取极限''构造子，它以实数的柯西序列作为输入；但``实数的柯西序列''这一概念本身，又依赖于某种度量实数间``距离''的方法。
当然，一般来说，两个实数之间的距离本身也是一个实数，这可能带来有问题的循环。

然而，对于``实数的柯西序列''这一概念，我们实际需要的并非一般的``距离''概念，而是一种能够表达``两个实数之间的距离\index{distance}小于 $\epsilon$''的方法，对任意 $\epsilon:\Qp$ 成立。
这可以用一个二元关系族来表示，记作 $\mathord{\close\epsilon} : \RC\to\RC\to \prop$。
关系 $x \close\epsilon y$ 的预期含义是 $|x - y| < \epsilon$；但由于我们尚未定义减法、绝对值或不等式等概念（毕竟我们才刚刚在定义 $\RC$），我们就不得不在定义 $\RC$ 本身的同时，一并定义这些关系 $\close\epsilon$。
又因为 $\close\epsilon$ 是一个以两份 $\RC$ 为指标的类型族，所以我们无法用普通的相互（高阶）归纳定义来完成；取而代之，我们必须使用一种\emph{高阶归纳-归纳定义}。
\index{inductive-inductive type!higher}

回顾 \cref{sec:generalizations} 可知，普通意义上的归纳-归纳定义允许我们通过同时归纳的方式，定义一个类型以及一个以该类型为指标的类型族。
当然，它的``高阶''版本则允许该类型和该类型族都既包含点构造子，也包含道路构造子。
本书不会尝试建立高阶归纳-归纳定义的一般性理论，但我们希望接下来对 $\RC$ 和 $\close\epsilon$ 的描述，能够使这个想法清晰明了。

\begin{rmk}
  我们也可以考虑\emph{高阶归纳-递归定义（higher inductive-recursive definition）}，其中 $\close\epsilon$ 是利用 $\RC$ 的\emph{递归}原理来定义的，而这一定义又与 $\RC$ 本身的\emph{归纳}定义同时进行。
  但我们选择了归纳-归纳的路线，原因有二。
  首先，高阶归纳-递归定义似乎更难在同伦语义中给出证成。
  其次，也是更重要的一点，归纳-归纳的定义能给出更强的归纳原理，而发展 Cauchy 实数的基本理论就已经需要用到它。
\end{rmk}

最后，如同我们在讨论 Dedekind 实数的 Cauchy 完备性时（见 \cref{sec:RD-cauchy-complete}）所做的那样，我们将使用\emph{Cauchy 逼近（Cauchy approximations）}（\cref{defn:cauchy-approximation}）而不是 Cauchy 序列。
当然，现在的 Cauchy 逼近将是由 Cauchy 实数构成的，而不是 Dedekind 实数或有理数。

\begin{defn}\label{defn:cauchy-reals}
  令 $\RC$ 以及关系 $\closesym:\RC \times \RC \times \Qp \to \type$ 为以下的高阶归纳-归纳类型族。
  \define{Cauchy 实数（Cauchy reals）}\indexdef{real numbers!Cauchy}%\indexsee{Cauchy!real numbers}{real numbers, Cau\-chy}的类型 $\RC$ 由以下构造子生成：
  \begin{itemize}
  \item \emph{有理点（rational points）}：
    对任意 $q : \Q$，存在一个实数 $\rcrat(q)$。
    \index{rational numbers!as Cauchy real numbers}%
  \item \emph{极限点（limit points）}：
    对任意 $x : \Qp \to \RC$，若满足
    %
    \begin{equation}
      \label{eq:RC-cauchy}
      \fall{\delta, \epsilon : \Qp} x_\delta \close{\delta + \epsilon} x_\epsilon
    \end{equation}
    %
    则存在一点 $\rclim(x) : \RC$。我们称 $x$ 为一个\define{Cauchy 逼近（Cauchy approximation）}\indexdef{Cauchy!approximation}%
\index{limit!of a Cauchy approximation}。%
  \item \emph{道路（paths）}：
    对 $u, v : \RC$，若满足
    %
    \begin{equation}
      \label{eq:RC-path}
      \fall{\epsilon : \Qp} u \close\epsilon v
    \end{equation}
    %
    则存在一条道路 $\rceq(u, v) : \id[\RC]{u}{v}$。
  \end{itemize}
  同时，类型族 $\closesym:\RC\to\RC\to\Qp \to\type$ 由以下构造子生成。
  此处 $q$ 与 $r$ 表示有理数；$\delta$、$\epsilon$ 与 $\eta$ 表示正有理数；$u$ 与 $v$ 表示 Cauchy 实数；$x$ 与 $y$ 表示 Cauchy 逼近：
  \begin{itemize}
  \item 对任意 $q,r,\epsilon$，若 $-\epsilon < q - r < \epsilon$，则 $\rcrat(q) \close\epsilon \rcrat(r)$，
  \item 对任意 $q,y,\epsilon,\delta$，若 $\rcrat(q) \close{\epsilon - \delta} y_\delta$，则 $\rcrat(q) \close{\epsilon} \rclim(y)$，
  \item 对任意 $x,r,\epsilon,\delta$，若 $x_\delta \close{\epsilon - \delta} \rcrat(r)$，则 $\rclim(x) \close\epsilon \rcrat(r)$，
  \item 对任意 $x,y,\epsilon,\delta,\eta$，若 $x_\delta \close{\epsilon - \delta - \eta} y_\eta$，则 $\rclim(x) \close\epsilon \rclim(y)$，
  \item 对任意 $u,v,\epsilon$，若 $\xi,\zeta : u \close{\epsilon} v$，则 $\xi=\zeta$（命题截断）。
  \end{itemize}
\end{defn}

\mentalpause

$\RC$ 的第一个构造子说明任何有理数都可以看作一个实数。
第二个构造子说明，从任意 Cauchy 逼近出发，我们可以得到一个新的实数，称为它的“极限”。
第三个构造子则表达了这样一个想法：如果两个 Cauchy 逼近相一致，那么它们的极限相等。

$\closesym$ 的前四个构造子分别规定：两个有理数何时接近、一个有理数与一个极限何时接近，以及两个极限何时接近。
对于两个有理数的情况，这正好是通常意义下有理数的 $\epsilon$-接近；而其他情况则可以通过以下观察导出：每个逼近项 $x_\delta$ 都应当在其极限 $\rclim(x)$ 的 $\delta$ 范围之内。

我们提醒自己注意证明相关性：通过 $\rclim$ 得到的实数，其表示不仅包含一个 Cauchy 逼近 $x$，还包含 \eqref{eq:RC-cauchy} 的一个证明 $p$，所以严格来说我们应该写成 $\rclim(x,p)$ 而不是 $\rclim(x)$。
类似的观察也适用于 $\rceq$ 和 \eqref{eq:RC-path}，但我们将仅写作 $\rceq : u = v$ 而非 $\rceq(u, v, p) : u = v$。
这些记法滥用之所以可以接受，是因为我们省略了那些命题以及可以轻易推断出的信息。
同样地，$\mathord{\close\epsilon}$ 的最后一个构造子也使得我们可以不给其他四个构造子命名。

我们立刻就能用许多实数来填充 $\RC$。假设 $x : \N \to
\Q$ 是一个传统的有理数 Cauchy 序列\index{Cauchy!sequence}，并令 $M : \Qp \to \N$ 为其收敛模。
那么 $\rcrat \circ x \circ M : \Qp \to \RC$ 就是一个 Cauchy 逼近（利用 $\closesym$ 的第一个构造子来生成必要的见证）。
因此，$\rclim(\rcrat \circ x \circ M)$ 就是一个实数。
像 $\sqrt{2}$、$\pi$、$e$…… 这样的著名实数，都是这类有理数 Cauchy 序列的极限。

\subsection{Cauchy 实数上的归纳与递归}
\label{sec:induct-recurs-cauchy}

当然，要想对 $\RC$ 做些有用的事情，我们需要给出它的归纳原理。
正如同时归纳定义两个或更多对象时的情况一样，$\RC$ 和 $\closesym$ 的基本归纳原理要求对两者同时进行归纳。
因此，我们预期该原理是说：假设有两个类型族分别居于 $\RC$ 和 $\closesym$ 之上，并给定对应于每个构造子的数据，那么这两个类型族就都存在截面。
然而，由于 $\closesym$ 以 $\RC$ 的两个拷贝为指标，这些类型族的精确依赖关系就有些微妙了。
归纳原理将适用于任意一对类型族：
\begin{align*}
A&:\RC\to\type\\
B&:\prd{x,y:\RC} A(x) \to A(y) \to \prd{\epsilon:\Qp} (x\close\epsilon y) \to \type.
\end{align*}
$A$ 的类型是明显的，但 $B$ 的类型就需要稍加思索了。
既然 $B$ 必须依赖于 $\closesym$，而 $\closesym$ 本身又依赖于两个 $\RC$ 的拷贝和一个 $\Qp$ 的拷贝，那么很明显，$B$ 也必须依赖于变量 $x,y:\RC$、$\epsilon:\Qp$ 以及 $(x\close\epsilon y)$ 的一个元素。
稍不明显的是，$B$ 还必须依赖于 $A(x)$ 和 $A(y)$。

如果我们考虑非依值情形（即递归原理），这一点可能会更清楚；在那种情形下，$A$ 是一个简单的类型（而不是类型族）。
此时，我们会希望 $B$ 不依赖于 $x,y:\RC$ 或 $x\close\epsilon y$。
但递归原理（连同其相关的唯一性原理）意在表明，带有关联关系 $\close\epsilon$ 的 $\RC$ 是某个范畴中的``始对象''。
因此，在这种情况下，$A$ 和 $B$ 的依赖结构应当反映出 $\RC$ 和 $\close\epsilon$ 的依赖结构：也就是说，我们应有 $B:A\to A\to \Qp \to \type$。
将这一观察与依值情形下 $B$ 还必须依赖于 $x,y:\RC$ 和 $x\close\epsilon y$ 的事实结合起来，便不可避免地导出了上面给出的 $B$ 的类型。

\symlabel{RC-recursion}
将 $B$ 看作一个以 $\epsilon$ 为索引的、介于类型 $A(x)$ 与 $A(y)$ 之间的关系族，是很有帮助的。
据此，我们可以将 $B(x,y,a,b,\epsilon,\xi)$ 写作 $(x,a) \bsim_\epsilon^\xi (y,b)$。
由于 $\xi:x\close\epsilon y$ 在存在时是唯一的，我们通常将其从记号中省略，而写作 $(x,a) \bsim_\epsilon (y,b)$；只要我们记住这个关系仅在 $x\close\epsilon y$ 成立时有定义，这样做就不会出问题。
我们有时还会进一步简化，写作 $a\bsim_\epsilon b$，此时 $x$ 和 $y$ 可以从 $a$ 和 $b$ 的类型推断出来；但为了清晰起见，有时仍有必要将它们包含在记号中。

\index{induction principle!for Cauchy reals}%
现在，给定类型族 $A:\RC\to\type$ 以及如上所述的关系族 $\bsim$，归纳原理的前提由以下数据组成，每一条对应 $\RC$ 或 $\closesym$ 的一个构造子：

\begin{itemize}
\item 对于任意 $q : \Q$，有一个元素 $f_q:A(\rcrat(q))$。
\item 对于任意 Cauchy 逼近 $x$，以及任意满足 $a:\prd{\epsilon:\Qp} A(x_\epsilon)$ 的 \begin{equation}
    \fall{\delta, \epsilon : \Qp}
    (x_\delta,a_\delta) \bsim_{\delta+\epsilon} (x_\epsilon,a_\epsilon),
    \label{eq:depCauchyappx}
  \end{equation}，有一个元素 $f_{x,a}:A(\rclim(x))$。
  我们将这样的 $a$ 称为一个\define{依值 Cauchy 逼近（dependent Cauchy approximation）}
  \indexdef{Cauchy!approximation!dependent}%
  \indexsee{approximation, Cauchy}{Cauchy approximation}%
  \indexdef{dependent!Cauchy approximation}%
  在 $x$ 上。
\item 对于满足 $h:\fall{\epsilon : \Qp} u \close\epsilon v$ 的 $u, v : \RC$，以及所有满足 $\fall{\epsilon:\Qp} (u,a) \bsim_\epsilon (v,b)$ 的 $a:A(u)$ 与 $b:A(v)$，有一条依值道路 $\dpath{A}{\rceq(u,v)}{a}{b}$。
\item 对于 $q,r:\Q$ 和 $\epsilon:\Qp$，如果 $-\epsilon < q - r < \epsilon$，则有
  \narrowequation{(\rcrat(q),f_q) \bsim_\epsilon (\rcrat(r),f_r).}
\item 对于 $q:\Q$、$\delta,\epsilon:\Qp$、Cauchy 逼近 $y$，以及 $y$ 上的依值 Cauchy 逼近 $b$，如果 $\rcrat(q) \close{\epsilon - \delta} y_\delta$，那么
  \[(\rcrat(q),f_q) \bsim_{\epsilon-\delta} (y_\delta,b_\delta)
  \;\Rightarrow\;
  (\rcrat(q),f_q) \bsim_\epsilon (\rclim(y),f_{y,b}).\]
\item 类似地，对于 $r:\Q$、$\delta,\epsilon:\Qp$、Cauchy 逼近 $x$，以及 $x$ 上的依值 Cauchy 逼近 $a$，如果 $x_\delta \close{\epsilon - \delta} \rcrat(r)$，那么
  \[(x_\delta,a_\delta) \bsim_{\epsilon-\delta} (\rcrat(r),f_r)
  \;\Rightarrow\;
  (\rclim(x),f_{x,a}) \bsim_\epsilon (\rcrat(q),f_r).
  \]
\item 对于 $\epsilon,\delta,\eta:\Qp$、Cauchy 逼近 $x,y$，以及分别在 $x$ 和 $y$ 上的依值 Cauchy 逼近 $a$ 与 $b$，如果我们有 $x_\delta \close{\epsilon - \delta - \eta} y_\eta$，那么
  \[ (x_\delta,a_\delta) \bsim_{\epsilon - \delta - \eta} (y_\eta,b_\eta)
  \;\Rightarrow\;
  (\rclim(x),f_{x,a}) \bsim_\epsilon (\rclim(y),f_{y,b}).\]
\item 对于 $\epsilon:\Qp$、$x,y:\RC$、$\xi,\zeta:x\close{\epsilon} y$，以及 $a:A(x)$ 和 $b:A(y)$，$(x,a) \bsim_\epsilon^\xi (y,b)$ 与 $(x,a) \bsim_\epsilon^\zeta (y,b)$ 中的任意两个元素沿 $\xi=\zeta$ 依值相等。
  注意，按照惯例，这等价于要求 $\bsim$ 取值于命题。
\end{itemize}

在这些前提下，我们推导出函数
\begin{align*}
  f&:\prd{x:\RC} A(x)\\
  g&:\prd{x,y:\RC}{\epsilon:\Qp}{\xi:x\close{\epsilon} y}
  (x,f(x)) \bsim_\epsilon^\xi (y,f(y))
\end{align*}
它们按照预期计算：
\begin{align}
  f(\rcrat(q)) &\defeq f_q, \label{eq:rcsimind1}\\
  f(\rclim(x)) &\defeq f_{x,(f,g)[x]}. \label{eq:rcsimind2}
\end{align}
这里 $(f,g)[x]$ 表示将 $f$ 和 $g$ 应用于 Cauchy 逼近 $x$ 以得到 $x$ 上的依值 Cauchy 逼近。
也就是说，我们定义 $(f,g)[x]_\epsilon \defeq f(x_\epsilon) : A(x_\epsilon)$，然后对于任意 $\epsilon,\delta:\Qp$，我们有 $g(x_\epsilon,x_\delta,\epsilon+\delta,\xi)$ 来见证 $(f,g)[x]$ 是一个依值 Cauchy 逼近，其中 $\xi: x_\epsilon \close{\epsilon+\delta} x_\delta$ 来自 $x$ 是 Cauchy 逼近的假设。

我们之后再也不会用到这个记号了，所以不必费心去记它。
通常我们采用模式匹配的约定，其中 $f$ 通过形如~\eqref{eq:rcsimind1} 与~\eqref{eq:rcsimind2} 的等式来定义，而~\eqref{eq:rcsimind2} 的右侧可能出现符号 $f(x_\epsilon)$，以及它们构成依值 Cauchy 逼近这一假设。

不过，这条归纳原理的表述确实还是相当繁琐。
为了便于理解，我们注意到它作为特例包含了两个独立的归纳原理，分别关于~$\RC$ 和~$\closesym$。
首先，假设只给定一个类型族 $A:\RC\to\type$，并将 $\bsim$ 定义为取常值 \unit。
这样一来，许多所需的数据就变得平凡，余下仅需提供：

\begin{itemize}
\item 对任意 $q : \Q$，有一个元素 $f_q:A(\rcrat(q))$，
\item 对任意 Cauchy 逼近 $x$，以及任意 $a:\prd{\epsilon:\Qp} A(x_\epsilon)$，有一个元素 $f_{x,a}:A(\rclim(x))$，
\item 对 $u, v : \RC$ 和 $h:\fall{\epsilon : \Qp} u \close\epsilon v$，以及 $a:A(u)$ 与 $b:A(v)$，有 $\dpath{A}{\rceq(u,v)}{a}{b}$。
\end{itemize}

给定这些数据，归纳原理就给出一个函数 $f:\prd{x:\RC} A(x)$，使得
\begin{align*}
  f(\rcrat(q)) &\defeq f_q,\\
  f(\rclim(x)) &\defeq f_{x,f(x)}.
\end{align*}
我们称此原理为 \define{$\RC$-归纳}；它本质上说的是，若将 $\close\epsilon$ 视为给定的，则 $\RC$ 就是由其构造子归纳生成的。

请注意，若 $A$ 是命题，则 $\RC$-归纳中的第三个假设自动成立（稍后我们会看到，这二者实际上是等价的）。
因此，要证明关于实数的命题，只需对有理数以及 Cauchy 逼近的极限加以证明即可。
下面是一个例子。

\begin{lem} \label{lem:close-reflexive}
  对任意 $u:\RC$ 与 $\epsilon:\Qp$，有 $u\close\epsilon u$。
\end{lem}

\begin{proof}
  定义 $A(u) \defeq \fall{\epsilon:\Qp} (u\close\epsilon u)$。
  由于这是一个命题（根据 $\closesym$ 的最后一个构造子），由 $\RC$-归纳，只需对 $u$ 为 $\rcrat(q)$ 和 $u$ 为 $\rclim(x)$ 两种情形加以证明。
  在第一种情形，显然对任意 $\epsilon$ 都有 $|q-q|<\epsilon$，于是根据 $\closesym$ 的第一个构造子，有 $\rcrat(q) \close\epsilon \rcrat(q)$。
  %
  在第二种情形，我们可以归纳地假设对所有 $\delta,\epsilon:\Qp$ 都有 $x_\delta \close\epsilon x_\delta$。
  那么特别地，我们有 $x_{\epsilon/3} \close{\epsilon/3} x_{\epsilon/3}$，从而根据 $\closesym$ 的第四个构造子，得到 $\rclim(x) \close{\epsilon} \rclim(x)$。
\end{proof}

由 \cref{lem:close-reflexive} 可知，只有当族 $A:\RC\to\type$ 是命题时，直接应用 $\RC$-归纳才有可能成功。
为了看清这一点，固定 $u:\RC$。
取 $v$ 为 $u$，则 $\RC$-归纳的第三个假设告诉我们，对任意 $a : A(u)$，有 $\dpath{A}{\rceq(u,u)}{a}{a}$。
此外，若再给定一个点 $b : A(u)$，我们还可得到 $\dpath{A}{\rceq(u,u)}{a}{b}$。
根据依值道路类型的定义，我们得出结论：这两条道路合起来意味着 $a = b$，即 $A(u)$ 中所有的点都相等。

\begin{thm}\label{thm:Cauchy-reals-are-a-set}
  $\RC$ 是一个集合。
\end{thm}

\begin{proof}
  我们刚刚证明了，纯粹关系
  \narrowequation{P(u,v) \defeq \fall{\epsilon:\Qp} (u\close\epsilon v)}
  是自反的。
  由于它蕴含相等关系（根据 $\RC$ 的道路构造子），结论由 \cref{thm:h-set-refrel-in-paths-sets} 可得。
\end{proof}

我们还可以证明，尽管 $\RC$ 可能不是\emph{有理数} Cauchy 序列集合的商集，但它仍然是\emph{实数} Cauchy 序列集合的商集。
（当然，这不能作为 $\RC$ 的有效\emph{定义}，但它是一个有用的性质。）
我们将 Cauchy 逼近的类型定义为
%
\symlabel{cauchy-approximations}%
\index{Cauchy!approximation!type of}%
\begin{equation*}
  \CAP \defeq
  \setof{ x : \Qp \to \RC |
    \fall{\epsilon, \delta : \Qp} x_\delta \close{\delta + \epsilon} x_\epsilon
  }.
\end{equation*}
$\RC$ 的第二个构造子给出了一个函数 $\rclim:\CAP\to\RC$。

\begin{lem} \label{RC-lim-onto}
  每个实数都仅仅是极限点：$\fall{u : \RC} \exis{x : \CAP} u = \rclim(x)$。
  换言之，$\rclim:\CAP\to\RC$ 是满的。
\end{lem}

\begin{proof}
  根据 $\RC$-归纳，我们可以对 $u$ 分情况讨论。
  显然，如果 $u$ 是一个极限 $\rclim(x)$，结论是平凡的。
  因此假设 $u$ 是一个有理点 $\rcrat(q)$；我们断言 $u$ 等于 $\rclim(\lam{\epsilon} \rcrat(q))$。
  根据 $\RC$ 的道路构造子，只需证明对所有 $\epsilon:\Qp$ 有 $\rcrat(q) \close\epsilon \rclim(\lam{\epsilon} \rcrat(q))$。
  而根据 $\closesym$ 的第二个构造子，为此只需找到 $\delta:\Qp$ 使得 $\rcrat(q)\close{\epsilon-\delta} \rcrat(q)$。
  但根据 $\closesym$ 的第一个构造子，我们可以取任意满足 $\delta<\epsilon$ 的 $\delta:\Qp$。
\end{proof}

%

\begin{lem} \label{RC-lim-factor}
  如果 $A$ 是一个集合，且函数 $f : \CAP \to A$ 保持 Cauchy 逼近的\index{coincidence!of Cauchy approximations}重合，意指
  %
  \begin{equation*}
    \fall{x, y : \CAP} \rclim(x) = \rclim(y) \Rightarrow f(x) = f(y),
  \end{equation*}
  %
  则 $f$ 沿 $\rclim : \CAP \to \RC$ 唯一分解。
\end{lem}

\begin{proof}
  由于 $\rclim$ 是满的，根据 \cref{lem:images_are_coequalizers}，$\RC$ 是 $\CAP$ 关于 $\rclim$ 的\index{kernel!pair}核对的商集。
  而这正是本引理的断言。
\end{proof}

现在来看归纳原理的第二个特殊情况。假设我们将 $A$ 取为取值恒为 $\unit$ 的常值类型族。
在这种情况下，$\bsim$ 就只是一个以 $\epsilon$ 为指标、定义在 $\epsilon$-接近的实数对上的关系族，因此我们可以记 $u\bsim_\epsilon v$ 来代替 $(u,\ttt)\bsim_\epsilon (v,\ttt)$。
此时所需的数据简化为如下几项，其中 $q, r$ 表示有理数，$\epsilon, \delta, \eta$ 表示正有理数，$x, y$ 表示 Cauchy 逼近：

\begin{itemize}
\item 若 $-\epsilon < q - r < \epsilon$，则
  $\rcrat(q) \bsim_\epsilon \rcrat(r)$，
\item 若 $\rcrat(q) \close{\epsilon - \delta} y_\delta$ 且
  $\rcrat(q)\bsim_{\epsilon-\delta} y_\delta$，
  则 $\rcrat(q) \bsim_\epsilon \rclim(y)$，
\item 若 $x_\delta \close{\epsilon - \delta} \rcrat(r)$ 且
  $x_\delta \bsim_{\epsilon-\delta} \rcrat(r)$，
  则 $\rclim(y) \bsim_\epsilon \rcrat(q)$，
\item 若 $x_\delta \close{\epsilon - \delta - \eta} y_\eta$ 且
  $x_\delta\bsim_{\epsilon - \delta - \eta} y_\eta$，
  则 $\rclim(x) \bsim_\epsilon \rclim(y)$。
\end{itemize}

最终得到的结论是 $\fall{u,v:\RC}{\epsilon:\Qp} (u\close\epsilon v) \to (u \bsim_\epsilon v)$。
我们称这一原理为 \define{$\closesym$-归纳（$\closesym$-induction）}；它本质上说的是，如果我们把 $\RC$ 视为给定的，那么 $\close\epsilon$ 作为类型族，是由\emph{它自己的}构造子归纳生成的。
例如，我们可以用它来证明 $\closesym$ 是对称的。

\begin{lem}\label{thm:RCsim-symmetric}
  对于任意 $u,v:\RC$ 和 $\epsilon:\Qp$，有 $(u\close\epsilon v) = (v\close\epsilon u)$。
\end{lem}

\begin{proof}
  由于两边都是命题，利用对称性只需证明其中一个蕴含关系即可。
  因此，设 $(u\bsim_\epsilon v) \defeq (v\close\epsilon u)$。
  根据 $\closesym$-归纳，我们可以将问题归约到下述情形：$u\close\epsilon v$ 是由 $\closesym$ 的四个有趣的构造子之一推导出来的。
  在第一种情形，即 $u$ 和 $v$ 均为有理数时，结果是显然的（我们可以再次应用第一个构造子）。
  在另外三种情形中，归纳假设（结合 $\Q$ 中加法的交换性）恰好给出了 $\closesym$ 的另一个构造子所需的输入（第二和第三个构造子互换角色，而第四个则保持不变）。
\end{proof}

因此，一般的归纳原理——我们可称之为 \define{$(\RC,\closesym)$-归纳（$(\RC,\closesym)$-induction）}——是 $\RC$-归纳与 $\closesym$-归纳的一种联合形式。
例如，考虑它的非依值版本，我们称之为 \define{$(\RC,\closesym)$-递归（$(\RC,\closesym)$-recursion）}，这也是我们最常用到的形式。
\index{recursion principle!for Cauchy reals}%
普通的 $\RC$-递归告诉我们，要定义一个函数 $f : \RC \to A$，只需满足：

\begin{enumerate}
\item 对每个 $q : \Q$，构造 $f(\rcrat(q)) : A$，
\item 对每个 Cauchy 逼近 $x : \Qp \to \RC$，构造 $f(x) : A$，
  这里假定对所有的 $\epsilon : \Qp$，$f(x_\epsilon)$ 已被定义，
\item 对所有满足 $\fall{\epsilon:\Qp} u\close\epsilon v$ 的 $u, v : \RC$，证明 $f(u) = f(v)$。\label{item:rcrec3}
\end{enumerate}

然而，在不知道 $f$ 如何作用在 $\epsilon$-接近的 Cauchy 实数上的情况下，通常很难证明~\ref{item:rcrec3}。
增强的 $(\RC,\closesym)$-递归原理弥补了这一缺陷，它允许我们指定一个\emph{任意的} ``$f$ 在 $\epsilon$-接近的 Cauchy 实数上的作用方式''，然后我们可以通过与 $f$ 的定义同步进行的归纳来证明这一点确实成立。
这个``作用方式''就是关系族 $\bsim$。
由于 $A$ 独立于 $\RC$，为了简化，我们可以假设 $\bsim$ 只依赖于 $A$ 和 $\Qp$，因此用 $a\bsim_\epsilon b$ 代替 $(u,a) \bsim_\epsilon (v,b)$ 不会产生歧义。
在这种情况下，使用 $(\RC,\closesym)$-递归定义一个函数 $f:\RC\to A$ 需要处理以下几种情形（现在我们使用模式匹配约定的写法来书写）。

\begin{itemize}
\item 对每个 $q : \Q$，构造 $f(\rcrat(q)) : A$。
\item 对每个 Cauchy 逼近 $x : \Qp \to \RC$，构造 $f(\rclim(x)) : A$，这里归纳地假设对所有 $\epsilon : \Qp$，$f(x_\epsilon)$ 已经定义，并且构成一个``关于 $\bsim$ 的 Cauchy 逼近''，即 $\fall{\epsilon,\delta:\Qp} (f(x_\epsilon) \bsim_{\epsilon+\delta} f(x_\delta))$。
\item 证明关系 $\bsim$ 是\emph{分离的}，即对任意 $a,b:A$，
  \indexdef{relation!separated family of}%
  \indexdef{separated family of relations}%
\narrowequation{(\fall{\epsilon:\Qp} a\bsim_\epsilon b) \Rightarrow (a=b).}
\item 证明对 $q,r:\Q$，若 $-\epsilon< q-r <\epsilon$，则 $f(\rcrat(q))\bsim_\epsilon f(\rcrat(r))$。
\item 对任意 $q:\Q$ 与任意 Cauchy 逼近 $y$，证明
\narrowequation{f(\rcrat(q)) \bsim_\epsilon f(\rclim(y)),} 这里归纳地假设 $\rcrat(q)\close{\epsilon-\delta} y_\delta$ 和 $f(\rcrat(q)) \bsim_{\epsilon-\delta} f(y_\delta)$ 对某个 $\delta:\Qp$ 成立，并且 $\eta \mapsto f(x_\eta)$ 是关于 $\bsim$ 的 Cauchy 逼近。
\item 对任意 Cauchy 逼近 $x$ 与任意 $r:\Q$，证明
\narrowequation{f(\rclim(x)) \bsim_\epsilon f(\rcrat(r)),}
这里归纳地假设 $x_\delta \close{\epsilon-\delta} \rcrat(r)$ 和 $f(x_\delta) \bsim_{\epsilon-\delta} f(\rcrat(r))$ 对某个 $\delta:\Qp$ 成立，并且 $\eta\mapsto f(x_\eta)$ 是关于 $\bsim$ 的 Cauchy 逼近。
\item 对任意 Cauchy 逼近 $x,y$，证明
\narrowequation{f(\rclim(x)) \bsim_\epsilon f(\rclim(y)),}
这里归纳地假设 $x_\delta \close{\epsilon-\delta-\eta} y_\eta$ 和 $f(x_\delta) \bsim_{\epsilon-\delta-\eta} f(y_\eta)$ 对某个 $\delta,\eta:\Qp$ 成立，并且 $\theta\mapsto f(x_\theta)$ 与 $\theta\mapsto f(y_\theta)$ 都是关于 $\bsim$ 的 Cauchy 逼近。
\end{itemize}


注意，在后四个证明中，我们可以自由使用前两条数据中给出的 $f(\rcrat(q))$ 和 $f(\rclim(x))$ 的具体定义。
然而，分离性的证明必须适用于 $A$ 的\emph{任意}两个元素，而与 $f$ 无关：它是关系族 $\bsim$ 的一种``可接受性''条件。
因此，我们通常在定义 $\bsim$ 之后立即验证它，然后再去定义 $f(\rcrat(q))$ 和 $f(\rclim(x))$。

在上述假设下，$(\RC,\closesym)$-递归给出一个函数 $f:\RC\to A$，使得 $f(\rcrat(q))$ 和 $f(\rclim(x))$ 与前两条中给出的定义判断相等。
此外，我们还可以得出
\begin{equation}
  \fall{u,v:\RC}{\epsilon:\Qp} (u\close\epsilon v) \to (f(u) \bsim_\epsilon f(v)).\label{eq:RC-sim-recursion-extra}
\end{equation}

作为一个典型例子，$(\RC,\closesym)$-递归允许我们将定义在 $\Q$ 上的函数延拓到整个 $\RC$，只要这些函数是充分连续的。
\index{function!continuous}%

\begin{defn}\label{defn:lipschitz}
  函数 $f:\Q\to\RC$ 称为 \define{Lipschitz 的}
  \indexdef{function!Lipschitz}%
  \indexdef{Lipschitz!function}%
  \indexdef{Lipschitz!constant}%
  \indexdef{constant!Lipschitz}%
  ，如果存在 $L:\Qp$（称为 \define{Lipschitz 常数}）使得
  \[ |q - r|<\epsilon \Rightarrow (f(q) \close{L\epsilon} f(r)) \]
  对所有 $\epsilon:\Qp$ 和 $q,r:\Q$ 成立。
  %
  类似地，$g:\RC\to\RC$ 是 \define{Lipschitz 的}，如果存在 $L:\Qp$ 使得
  \[ (u\close\epsilon v) \Rightarrow (g(u) \close{L\epsilon} g(v)) \]
  对所有 $\epsilon:\Qp$ 和 $u,v:\RC$ 成立。
\end{defn}

特别地，注意根据 $\closesym$ 的第一个构造子，如果 $f:\Q\to\Q$ 在显然的意义上是 Lipschitz 的，那么复合 $\Q\xrightarrow{f} \Q \to \RC$ 也是 Lipschitz 的。

\begin{lem}\label{RC-extend-Q-Lipschitz}
  设 $f : \Q \to \RC$ 是 Lipschitz 的，常数为 $L : \Qp$。
  则存在一个 Lipschitz 映射 $\bar{f} : \RC \to \RC$，同样以 $L$ 为常数，使得 $\bar{f}(\rcrat(q)) \jdeq f(q)$ 对所有 $q:\Q$ 成立。
\end{lem}

\begin{proof}
  % Uniqueness follows directly from \cref{RC-continuous-eq}.
  我们通过 $(\RC,\closesym)$-递归来定义 $\bar{f}$，陪域为 $A\defeq \RC$。
  我们将关系 $\mathord{\bsim}: \RC \to \RC \to \Qp \to \prop$ 定义为
  \begin{align*}
    (u \bsim_\epsilon v) &\defeq (u \close{L\epsilon} v).
  \end{align*}
  对于 $q : \Q$，我们定义
  %
  \begin{equation*}
    \bar{f}(\rcrat(q)) \defeq \rcrat(f(q)).
  \end{equation*}
  %
  对于一个 Cauchy 逼近 $x : \Qp \to \RC$，我们定义
  %
  \begin{equation*}
    \bar{f}(\rclim(x)) \defeq \rclim(\lamu{\epsilon : \Qp} \bar{f}(x_{\epsilon/L})).
  \end{equation*}
  %
  为使此定义有意义，我们必须验证 $y \defeq \lamu{\epsilon : \Qp} \bar{f}(x_{\epsilon/L})$ 确实是一个 Cauchy 逼近。
  而这步的归纳假设是：对任意 $\delta,\epsilon:\Qp$ 都有 $\bar{f}(x_\delta) \bsim_{\delta+\epsilon} \bar{f}(x_\epsilon)$，即 $\bar{f}(x_\delta) \close{L\delta+L\epsilon} \bar{f}(x_\epsilon)$。
  由此我们得到
  \[y_\delta \jdeq f(x_{\delta/L}) \close{\delta + \epsilon} f(x_{\epsilon/L})   \jdeq y_\epsilon. \]

  为证明分离性，我们只须观察到：$\fall{\epsilon:\Qp} a\bsim_\epsilon b$ 意味着 $\fall{\epsilon:\Qp} a\close{L\epsilon} b$，这蕴含着 $\fall{\epsilon:\Qp}a\close\epsilon b$，从而 $a=b$。

  为了完成 $(\RC,\closesym)$-递归，剩下还需验证 $\bsim$ 所满足的四个条件。
  这基本相当于证明 $\bar f$ 对于 $\closesym$ 的所有四个构造子都是 Lipschitz 的。
  \begin{enumerate}
  \item 当 $u$ 是 $\rcrat(q)$ 而 $v$ 是 $\rcrat(r)$，且满足 $-\epsilon < |q-r| <\epsilon$ 时，由 $f$ 是 Lipschitz 的假设可得 $f(q) \close{L\epsilon} f(r)$，从而根据定义有 $\bar{f}(\rcrat(q)) \bsim_\epsilon \bar{f}(\rcrat(r))$。
  \item 当 $u$ 是 $\rclim(x)$ 而 $v$ 是 $\rcrat(q)$，且满足 $x_\eta \close{\epsilon - \eta} \rcrat(q)$ 时，归纳假设是 $\bar{f}(x_\eta) \close{L \epsilon - L \eta} \rcrat(f(q))$；这结合 $\closesym$ 的第三个构造子，即证明了
      \narrowequation{\bar{f}(\rclim(x)) \close{L \epsilon} \bar{f}(\rcrat(q)).}
  \item 当 $u$ 为有理数而 $v$ 为极限的对称情形，证明本质上相同。
  \item 当 $u$ 是 $\rclim(x)$ 而 $v$ 是 $\rclim(y)$，并且存在 $\delta, \eta : \Qp$ 使得 $x_\delta \close{\epsilon - \delta - \eta} y_\eta$ 时，归纳假设是 $\bar{f}(x_\delta) \close{L \epsilon - L \delta - L \eta} \bar{f}(y_\eta)$；这结合 $\closesym$ 的第四个构造子，即证明了 $\bar{f}(\rclim(x)) \close{L
        \epsilon} \bar{f}(\rclim(y))$。
  \end{enumerate}
  至此完成了 $(\RC,\closesym)$-递归，从而完成了 $\bar f$ 的构造。
  所期望的等式 $\bar f(\rcrat(q))\jdeq f(q)$ 正是 $(\RC,\closesym)$-递归的第一个计算规则，而附加条件~\eqref{eq:RC-sim-recursion-extra} 则正好说明 $\bar f$ 是常量为 $L$ 的 Lipschitz 函数。
\end{proof}

在对 $\closesym$ 有更好的刻画之前，我们目前能做的也就只有这些了。
我们在 $\closesym$ 的构造子中规定了两种不同形式的 Cauchy 实数在何种条件下 $\epsilon$-接近。
但是，我们如何知道，在最终得到的归纳-归纳类型族中，这些就是该事实的\emph{唯一}见证呢？
我们已经看到，归纳类型族（例如相等类型，参见 \cref{sec:identity-systems}）和高阶归纳类型往往包含“比最初放入的还要多”的东西，因此这并非一个无关紧要的问题。

为了更精确地刻画 $\closesym$，我们将在 $\RC$ 上\emph{递归地}定义一个关系族 $\approx_\epsilon$，使其能在构造子上进行计算，然后证明这个关系族等价于 $\close\epsilon$。

\begin{thm}\label{defn:RC-approx}
  存在一族纯粹关系 $\mathord\approx:\RC\to\RC\to\Qp\to\prop$，使得
  \begin{align}
    (\rcrat(q) \approx_\epsilon \rcrat(r))  &\defeq
    (-\epsilon < q - r < \epsilon)\label{eq:RCappx1}\\
    (\rcrat(q) \approx_\epsilon \rclim(y)) &\defeq
    \exis{\delta : \Qp} \rcrat(q) \approx_{\epsilon - \delta} y_\delta\label{eq:RCappx2}\\
    (\rclim(x) \approx_\epsilon \rcrat(r)) &\defeq
    \exis{\delta : \Qp} x_\delta \approx_{\epsilon - \delta} \rcrat(r)\label{eq:RCappx3}\\
    (\rclim(x) \approx_\epsilon \rclim(y)) &\defeq
    \exis{\delta, \eta : \Qp} x_\delta \approx_{\epsilon - \delta - \eta} y_\eta.\label{eq:RCappx4}
  \end{align}
  此外，我们还有
  \begin{gather}
    (u \approx_\epsilon v) \Leftrightarrow \exis{\theta:\Qp} (u \approx_{\epsilon-\theta} v) \label{RC-sim-rounded}\\
    (u \approx_\epsilon v) \to (v\close\delta w) \to (u\approx_{\epsilon+\delta} w)\label{eq:RC-sim-rtri}\\
    (u \close\epsilon v) \to (v\approx_\delta w) \to (u\approx_{\epsilon+\delta} w)\label{eq:RC-sim-ltri}.
  \end{gather}
\end{thm}

附加条件~\eqref{RC-sim-rounded}--\eqref{eq:RC-sim-ltri} 是使归纳定义得以成立所必需的。
条件~\eqref{RC-sim-rounded} 称为\define{舍入的（rounded）}。
\indexsee{relation!rounded}{rounded relation}%
\indexdef{rounded!relation}%
从右向左读，便得出 $\approx$ 的\define{单调性（monotonicity）}，
\index{monotonicity}%
\index{relation!monotonic}%
%
\begin{equation*}
  (\delta < \epsilon) \land (u \approx_\delta v) \Rightarrow (u \approx_\epsilon v)
\end{equation*}
%
而从左向右读，则给出 $\approx$ 的\define{开性（openness）}，
\index{open!relation}%
\index{relation!open}%
%
\begin{equation*}
  (u \approx_\epsilon v) \Rightarrow \exis{\delta : \Qp} (\delta < \epsilon) \land (u \approx_\delta v).
\end{equation*}
%
条件~\eqref{eq:RC-sim-rtri} 与~\eqref{eq:RC-sim-ltri} 是三角不等式的两种表现形式，它们表明 $\approx$ 是 $\closesym$ 两侧的``模''。

\begin{proof}
  我们将通过双重 $(\RC,\closesym)$-递归来定义 $\mathord\approx:\RC\to\RC\to\Qp\to\prop$。
  首先，我们将目标域为 $\RC\to\Qp\to\prop$ 的某个子集进行 $(\RC,\closesym)$-递归，该子集由满足舍入性及某一种适当形式的三角形不等式的谓词族构成。
  若将这些谓词视作某个二元关系的一半，我们便将其记作 $(u,\epsilon) \mapsto (\hapx_\epsilon u)$，并用符号 $\hapname$ 指代整个关系。
  现在我们可以将 $A$ 精确地写为
  \begin{multline*}
    A \defeq\; \Bigg\{ \hapname :\RC\to\Qp\to\prop \;\bigg|\; \\
    \Big(\fall{u:\RC}{\epsilon:\Qp}
    \big((\hapx_\epsilon u) \Leftrightarrow \exis{\theta:\Qp} (\hapx_{\epsilon-\theta} u)\big)\Big)  \\
    \land \Big(\fall{u,v:\RC}{\eta,\epsilon:\Qp} (u\close\epsilon v) \to\\
    \big((\hapx_\eta u) \to (\hapx_{\eta+\epsilon} v) \big) \land \big((\hapx_\eta v) \to (\hapx_{\eta+\epsilon} u) \big)\Big)\Bigg\}
  \end{multline*}
  与子集的惯常记法一致，我们对 $A$ 的一个居民与其第一个分量 $\hapname$ 使用相同的记号。
  作为 $(\RC,\closesym)$-递归所需的关系族，我们考虑以下关系，它将确保另一种形式的三角形不等式：
  \begin{narrowmultline*}
    (\hapname \bsim_\epsilon \hapbname ) \defeq \narrowbreak
    \fall{u:\RC}{\eta:\Qp} ((\hapx_\eta u) \to (\hapxb_{\epsilon+\eta} u))
    \land \narrowbreak
    ((\hapxb_\eta u) \to (\hapx_{\epsilon+\eta} u)).
  \end{narrowmultline*}
  我们观察到这些关系是分离的。
  因为假设
  \narrowequation{\fall{\epsilon:\Qp} (\hapname \bsim_\epsilon \hapbname),}
  要证明 $\hapname = \hapbname$，只需对所有 $u:\RC$ 证明 $(\hapx_\epsilon u) \Leftrightarrow (\hapxb_\epsilon u)$。
  但由舍入性，$\hapx_\epsilon u$ 蕴含存在某个 $\theta$ 使得 $\hapx_{\epsilon-\theta} u$，这结合 $\hapname \bsim_\epsilon \hapbname$ 便蕴含 $\hapxb_\epsilon u$；反之亦然。

  现在递归原理所需的前两个数据如下。
  \begin{itemize}
  \item 对于任意 $q:\Q$，我们必须给出 $A$ 的一个元素，记作 $(\rcrat(q)\approx_{(\blank)} \blank)$。
  \item 对于任意 Cauchy 逼近 $x$，若假设已定义了一个函数 $\Qp \to A$（记作 $\epsilon \mapsto (x_\epsilon \approx_{(\blank)} \blank)$）满足性质
    % \[ \fall{u,v:\RC}{\delta,\epsilon,\eta:\Qp} (x_\delta \approx_\eta u) \to (u\close{\delta+\epsilon} v) \to (x_\epsilon \approx_{\eta+\delta+\epsilon} v) \]
    \begin{equation}
      \fall{u:\RC}{\delta,\epsilon,\eta:\Qp} (x_\delta \approx_\eta u) \to (x_\epsilon \approx_{\eta+\delta+\epsilon} u),\label{eq:appxrec2}
    \end{equation}
    我们就必须给出 $A$ 的一个元素，记作 $(\rclim(x)\approx_{(\blank)} \blank)$。
  \end{itemize}
  在这两种情形下，我们都通过嵌套的 $(\RC,\closesym)$-递归来给出所需的定义，此递归的目标域是 $\Qp\to\prop$ 中由舍入的命题族构成的子集。
  若将这些命题视作某个二元关系的“零半边”，我们便将其记作 $\epsilon \mapsto (\tap{\epsilon})$，并用符号 $\tapname$ 指代整个族。
  现在我们可以将这些内层递归的目标域精确地写为
  \begin{narrowmultline*}
    C \defeq
    \bigg\{ \tapname :\Qp\to\prop \;\;\Big|\;\; \narrowbreak
    \fall{\epsilon:\Qp} \Big((\tap\epsilon) \Leftrightarrow \exis{\theta:\Qp} (\tap{\epsilon-\theta})\Big)\bigg\}
  \end{narrowmultline*}
  我们取所需的关系族为三角形不等式的剩余部分：
  \begin{narrowmultline*}
    (\tapname \bbsim_\epsilon \tapbname) \defeq
    \fall{\eta:\Qp} ((\tap\eta) \to (\tapb{\epsilon+\eta})) \land
    \narrowbreak
    ((\tapb\eta) \to (\tap{\epsilon+\eta})).
  \end{narrowmultline*}
  利用 $C$ 中所有元素的舍入性，通过与 $\bsim$ 相同的论证可知这些关系是分离的。

  注意，若这样一个内层递归成功，它将产生一个谓词族 $\hapname : \RC\to\Qp\to \prop$，这些谓词是舍入的（因为它们在 $\Qp\to\prop$ 中的像落在 $C$ 中）并且满足
  \[ \fall{u,v:\RC}{\epsilon:\Qp} (u\close\epsilon v) \to \big((\hapx_{(\blank)} u) \bbsim_\epsilon (\hapx_{(\blank)} u)\big). \]
  展开 $\bbsim$ 的定义，这恰是 $\hapname$ 属于 $A$ 所需的第三个条件；因此这正是我们需要的。

  此时，我们可以给出定义~\eqref{eq:RCappx1}--\eqref{eq:RCappx4}，作为两个内层递归各自对应于有理点和极限的前两个子句。
  在每种情形下，我们都必须验证该关系是舍入的，从而位于 $C$ 中。
  在有理数-有理数情形~\eqref{eq:RCappx1} 中，这是显然的；而在其他情形中，则由归纳假设得出。（在~\eqref{eq:RCappx2} 中，相关的归纳假设是 $(\rcrat(q) \approx_{(\blank)} y_\delta) : C$，而在~\eqref{eq:RCappx3} 和~\eqref{eq:RCappx4} 中则是 $(x_\delta \approx_{(\blank)} \blank) : A$。）

  子递归剩余的数据在于证明 \eqref{eq:RCappx1}--\eqref{eq:RCappx4} 对 $\closesym$ 的构造子满足右侧的三角形不等式。
  共有八种情形——每个子递归四种——对应于~\eqref{eq:RC-sim-rtri} 中 $u$、$v$ 和 $w$ 选为有理点或极限的八种可能组合。
  首先我们考虑 $u$ 是 $\rcrat(q)$ 的情形。
  \begin{enumerate}
  \item 假设 $\rcrat(q)\approx_\phi \rcrat(r)$ 且 $-\epsilon<|r-s|<\epsilon$，我们必须证明 $\rcrat(q)\approx_{\phi+\epsilon} \rcrat(s)$。
    但由 $\approx$ 的定义，这归约为有理数的三角形不等式。
  \item 我们假设 $\phi,\epsilon,\delta:\Qp$ 满足 $\rcrat(q)\approx_\phi \rcrat(r)$ 与 $\rcrat(r) \close{\epsilon-\delta} y_\delta$，并归纳地假设
    \begin{equation}
      \fall{\psi:\Qp}(\rcrat(q) \approx_{\psi} \rcrat(r)) \to (\rcrat(q) \approx_{\psi+\epsilon-\delta} y_\delta).\label{eq:RCappx-rtri-rrl1}
    \end{equation}
    我们还假设 $\psi,\delta\mapsto (\rcrat(q) \approx_{\psi} y_\delta)$ 是一个关于 $\bbsim$ 的 Cauchy 逼近，即
    \begin{equation}
      \fall{\psi,\xi,\zeta:\Qp} (\rcrat(q) \approx_{\psi} y_\xi) \to (\rcrat(q) \approx_{\psi+\xi+\zeta} y_\zeta),\label{eq:RCappx-rtri-rrl2}
    \end{equation}
    不过在此情形中我们并不需要这个假设。
    事实上，令 $\psi\defeq \phi$，\eqref{eq:RCappx-rtri-rrl1} 立即给出 $\rcrat(q) \approx_{\phi+\epsilon-\delta} y_\delta$，从而由 $\approx$ 的定义即得 $\rcrat(q) \approx_{\phi+\epsilon} \rclim(y)$。
  \item 我们假设 $\phi,\epsilon,\delta:\Qp$ 满足 $\rcrat(q)\approx_\phi \rclim(y)$ 与 $y_\delta \close{\epsilon-\delta} \rcrat(r)$，并归纳地假设
    \begin{gather}
      \fall{\psi:\Qp}(\rcrat(q) \approx_{\psi} y_\delta) \to (\rcrat(q) \approx_{\psi+\epsilon-\delta} \rcrat(r)).\label{eq:RCappx-rtri-rlr1}\\
      \fall{\psi,\xi,\zeta:\Qp} (\rcrat(q) \approx_{\psi} y_\xi) \to (\rcrat(q) \approx_{\psi+\xi+\zeta} y_\zeta).\label{eq:RCappx-rtri-rlr2}
    \end{gather}
    根据定义，$\rcrat(q)\approx_\phi \rclim(y)$ 意味着存在 $\theta:\Qp$ 使得 $\rcrat(q) \approx_{\phi-\theta} y_\theta$。
    因此，由假设~\eqref{eq:RCappx-rtri-rlr2}，我们也有 $\rcrat(q) \approx_{\phi+\delta} y_\delta$，然后由~\eqref{eq:RCappx-rtri-rlr1} 可知 $\rcrat(q) \approx_{\phi+\epsilon} \rcrat(r)$，正如所求。
  \item 我们假设 $\phi,\epsilon,\delta,\eta:\Qp$ 满足 $\rcrat(q)\approx_\phi \rclim(y)$ 与 $y_\delta \close{\epsilon-\delta-\eta} z_\eta$，并归纳地假设
    \begin{gather}
      \fall{\psi:\Qp}(\rcrat(q) \approx_{\psi} y_\delta) \to (\rcrat(q) \approx_{\psi+\epsilon-\delta-\eta} z_\eta), \label{eq:RCappx-rtri-rll1}\\
      \fall{\psi,\xi,\zeta:\Qp} (\rcrat(q) \approx_{\psi} y_\xi) \to (\rcrat(q) \approx_{\psi+\xi+\zeta} y_\zeta), \label{eq:RCappx-rtri-rll2}\\
      \fall{\psi,\xi,\zeta:\Qp} (\rcrat(q) \approx_{\psi} z_\xi) \to (\rcrat(q) \approx_{\psi+\xi+\zeta} z_\zeta). \label{eq:RCappx-rtri-rll3}
    \end{gather}
    同样，$\rcrat(q)\approx_\phi \rclim(y)$ 意味着存在 $\xi:\Qp$ 使得 $\rcrat(q) \approx_{\phi-\xi} y_\xi$，而~\eqref{eq:RCappx-rtri-rll2} 蕴含 $\rcrat(q) \approx_{\phi+\delta} y_\delta$，~\eqref{eq:RCappx-rtri-rll1} 蕴含 $\rcrat(q) \approx_{\phi+\epsilon-\eta} z_\eta$。
    但由 $\approx$ 的定义，这蕴含了即得所求的 $\rcrat(q) \approx_{\phi+\epsilon} \rclim(z)$。
  \end{enumerate}
  现在我们考虑 $u$ 是 $\rclim(x)$ 的情形，其中 $x$ 是一个 Cauchy 逼近。
  在此情形下，定义 $(\rclim(x) \approx_{(\blank)} {\blank}) : A$ 时的外围归纳假设为 ${(x_\delta \approx_{(\blank)} {\blank})}: A$，因此这些关系除了舍入性之外还满足右侧的三角形不等式。
  \begin{enumerate}\setcounter{enumi}{4}
  \item 假设 $\rclim(x)\approx_\phi \rcrat(r)$ 且 $-\epsilon<|r-s|<\epsilon$，我们必须证明 $\rclim(x)\approx_{\phi+\epsilon} \rcrat(s)$。
    根据 $\approx$ 的定义，前者意味着 $x_\delta \approx_{\phi-\delta} \rcrat(r)$，因此上述三角形不等式蕴含 $x_\delta \approx_{\epsilon+\phi-\delta} \rcrat(s)$，从而即得所求的 $\rclim(x)\approx_{\phi+\epsilon} \rcrat(s)$。
  \item 我们假设 $\phi,\epsilon,\delta:\Qp$ 满足 $\rclim(x)\approx_\phi \rcrat(r)$ 与 $\rcrat(r) \close{\epsilon-\delta} y_\delta$，以及两个多余的归纳假设。
    %
    根据定义，存在 $\eta:\Qp$ 使得 $x_\eta \approx_{\phi-\eta} \rcrat(r)$，因此归纳的三角形不等式给出 $x_\eta \approx_{\phi+\epsilon-\eta-\delta} y_\delta$。
    由 $\approx$ 的定义立得 $\rclim(x) \approx_{\phi+\epsilon} \rclim(y)$。
  \item 我们假设 $\phi,\epsilon,\delta:\Qp$ 满足 $\rclim(x)\approx_\phi \rclim(y)$ 与 $y_\delta \close{\epsilon-\delta} \rcrat(r)$，以及两个多余的归纳假设。
    根据定义，存在 $\xi,\theta:\Qp$ 使得 $x_\xi \approx_{\phi-\xi-\theta} y_\theta$。
    由于 $y$ 是一个 Cauchy 逼近，我们有 $y_\theta \close{\theta+\delta} y_\delta$，因此由归纳的三角形不等式可得 $x_\xi \approx_{\phi+\delta-\xi} y_\delta$，进而 $x_\xi \close{\phi+\epsilon-\xi} \rcrat(r)$。
    再由 $\approx$ 的定义即得 $\rclim(x) \approx_{\phi+\epsilon}\rcrat(r)$，正如所求。
  \item 最后，我们假设 $\phi,\epsilon,\delta,\eta:\Qp$ 满足 $\rclim(x)\approx_\phi \rclim(y)$ 与 $y_\delta \close{\epsilon-\delta-\eta} z_\eta$。
    那么如前所述，存在 $\xi,\theta:\Qp$ 使得 $x_\xi \approx_{\phi-\xi-\theta} y_\theta$，如前一样两次应用三角形不等式足矣。
  \end{enumerate}

  这就完成了两个内层递归，从而定义了关系族 $(\rcrat(q)\approx_{(\blank)}\blank)$ 和 $(\rclim(x)\approx_{(\blank)}\blank)$。
  因为它们都是 $A$ 的元素，所以是舍入的，并且对 $\closesym$ 满足右侧的三角形不等式。
% , and satisfy~\eqref{eq:appxrec2}.
  余下的工作是验证与 $\bsim$ 相关的条件，即这些关系对 $\closesym$ 的构造子满足\emph{左侧}的三角形不等式。
  四种情形对应于~\eqref{eq:RC-sim-ltri} 中 $u$ 和 $v$ 取有理点或极限的四种选择，且由于它们均为命题，我们可以应用 $\RC$-归纳法并假设 $w$ 也是有理点或极限。
  这又产生八种情形，其证明与刚给出的那些本质相同；因此我们就不再拿它们来折磨读者了。
\end{proof}

现在我们可以证明：

\begin{thm}\label{thm:RC-sim-characterization}
  对任意 $u,v:\RC$ 和 $\epsilon:\Qp$，有 $(u\close\epsilon v) = (u\approx_\epsilon v)$。
\end{thm}

\begin{proof}
  由于两边都是命题，只需证明双向蕴含。
  对于从左到右的方向，我们对 $C(u,v,\epsilon)\defeq (u\approx_\epsilon v)$ 使用 $\closesym$-归纳。
  因此，只需考虑 $\closesym$ 的四个构造子。
  在每种情形中，$u$ 和 $v$ 都被特化为有理点或极限，从而 $\approx$ 的定义可直接计算，归纳假设也总是适用。

  对于从右到左的方向，我们使用 $\RC$-归纳，假设 $u$ 和 $v$ 是有理点或极限，从而 $\approx$ 可求值。
  此时 $\approx$ 的定义以及归纳假设恰好提供了 $\closesym$ 相应构造子所需的数据。
\end{proof}

\index{encode-decode method}%
略加引申，我们可以将 $\approx$ 称为 $\closesym$ 的``代码''纤维化，而上述证明的两个方向分别对应 \encode 和 \decode。
根据 $\approx$ 的定义，由 \cref{thm:RC-sim-characterization} 我们得到等价
\begin{align*}
  (\rcrat(q) \close\epsilon \rcrat(r))  &=
  (-\epsilon < q - r < \epsilon)\\
  (\rcrat(q) \close\epsilon \rclim(y)) &=
  \exis{\delta : \Qp} \rcrat(q) \close{\epsilon - \delta} y_\delta\\
  (\rclim(x) \close\epsilon \rcrat(r)) &=
  \exis{\delta : \Qp} x_\delta \close{\epsilon - \delta} \rcrat(r)\\
  (\rclim(x) \close\epsilon \rclim(y)) &=
  \exis{\delta, \eta : \Qp} x_\delta \close{\epsilon - \delta - \eta} y_\eta.
\end{align*}
我们的证明还给出以下附加信息。

\begin{cor}
  \index{triangle!inequality for R@inequality for $\RC$}%
  \indexsee{inequality!triangle}{triangle inequality}%
  $\closesym$ 是\index{rounded!relation}舍入的，并且满足三角不等式：
    \begin{gather}
      \eqvspaced{
        (u \close\epsilon v)
      }{
        \exis{\theta : \Qp} u \close{\epsilon - \theta} v
      }\\
      (u\close\epsilon v) \to (v\close\delta w) \to (u\close{\epsilon+\delta} w). \label{item:RC-sim-triangle}
    \end{gather}
\end{cor}


% \begin{proof}
%   The construction of $\approx$ showed simultaneously that it is rounded, and satisfies ``triangle inequalities'' such as
%   \[ (u\approx_\epsilon v) \to (v\close\delta w) \to (u\approx_{\epsilon+\delta} w). \]
%   Thus, both properties follow from \cref{thm:RC-sim-characterization}.
% \end{proof}

有了三角不等式，我们就可以证明 Cauchy 逼近的``极限''确实表现得像极限。

\begin{lem}\label{thm:RC-sim-lim}
  对任意 $u:\RC$，任意 Cauchy 逼近 $y$，以及任意 $\epsilon,\delta:\Qp$，若 $u\close\epsilon y_\delta$，则 $u\close{\epsilon+\delta} \rclim(y)$。
\end{lem}

\begin{proof}
  我们对 $u$ 使用 $\RC$-归纳。
  若 $u$ 是 $\rcrat(q)$，这正是 $\closesym$ 的第二个构造子。
  现在假设 $u$ 是 $\rclim(x)$，并且每个 $x_\eta$ 都满足如下性质：对任意 $y,\epsilon,\delta$，若 $x_\eta\close\epsilon y_\delta$，则 $x_\eta \close{\epsilon+\delta} \rclim(y)$。
  特别地，在此假设中取 $y\defeq x$ 且 $\delta\defeq\eta$，可得：对任意 $\eta,\theta:\Qp$，有 $x_\eta \close{\eta+\theta} \rclim(x)$。

  现在任取 $y,\epsilon,\delta$ 并假设 $\rclim(x) \close\epsilon y_\delta$。
  由舍入性，存在 $\theta$ 使得 $\rclim(x) \close{\epsilon-\theta} y_\delta$。
  那么根据上面的观察，对任意 $\eta$ 有 $x_\eta \close{\eta+\theta/2} \rclim(x)$，进而由三角不等式得到 $x_\eta \close{\epsilon+\eta-\theta/2} y_\delta$。
  因此，$\closesym$ 的第四个构造子给出 $\rclim(x) \close{\epsilon+2\eta+\delta-\theta/2} \rclim(y)$。
  于是，取 $\eta \defeq \theta/4$，即得所求。
\end{proof}

\begin{lem}\label{thm:RC-sim-lim-term}
  对任意 Cauchy 逼近 $y$ 及任意 $\delta,\eta:\Qp$，有 $y_\delta \close{\delta+\eta} \rclim(y)$。
\end{lem}

\begin{proof}
  在前一个引理中取 $u\defeq y_\delta$ 和 $\epsilon\defeq \eta$ 即可。
\end{proof}

\begin{rmk}
  我们可能会期望 $y_\delta \close{\delta} \rclim(y)$ 成立，但这一性质在某些例子中并不成立。
  例如，考虑由 $x_\epsilon \defeq \epsilon$ 定义的 $x$。
  其极限显然是 $0$，但 $|\epsilon - 0 |<\epsilon$ 并不成立，只有 $\le$ 成立。
\end{rmk}

作为应用，\cref{thm:RC-sim-lim-term} 使我们能够证明 Lipschitz 函数从 \cref{RC-extend-Q-Lipschitz} 出发的扩展是唯一的。

\begin{lem}\label{RC-continuous-eq}
  \index{function!continuous}%
  设 $f,g:\RC\to\RC$ 连续，即满足
  \[ \fall{u:\RC}{\epsilon:\Qp}\exis{\delta:\Qp}\fall{v:\RC} (u\close\delta v) \to (f(u) \close\epsilon f(v)) \]
  对 $g$ 亦同。
  若对所有 $q:\Q$ 有 $f(\rcrat(q))=g(\rcrat(q))$，则 $f=g$。
\end{lem}

\begin{proof}
  我们通过 $\RC$-归纳证明对所有 $u$ 有 $f(u)=g(u)$。
  有理数的情形就是假设本身。
  于是，假设对所有 $\delta$ 有 $f(x_\delta)=g(x_\delta)$。
  我们将证明对所有 $\epsilon$ 有 $f(\rclim(x))\close\epsilon g(\rclim(x))$，从而可应用 $\RC$ 的道路构造子。

  由于 $f$ 和 $g$ 连续，故存在 $\theta,\eta$ 使得对所有 $v$ 有
  \begin{align*}
    (\rclim(x)\close\theta v) &\to (f(\rclim(x)) \close{\epsilon/2} f(v))\\
    (\rclim(x)\close\eta v) &\to (g(\rclim(x)) \close{\epsilon/2} g(v)).
  \end{align*}
  取 $\delta < \min(\theta,\eta)$，由 \cref{thm:RC-sim-lim-term} 我们既有 $\rclim(x)\close\theta y_\delta$ 也有 $\rclim(x)\close\eta y_\delta$。
  因此
  \[ f(\rclim(x)) \close{\epsilon/2} f(y_\delta) = g(y_\delta) \close{\epsilon/2} g(\rclim(x))\]
  于是由三角不等式可得 $f(\rclim(x))\close\epsilon g(\rclim(x))$。
\end{proof}

\subsection{Cauchy 实数的代数结构}
\label{sec:algebr-struct-cauchy}

我们首先定义加法结构 $(\RC, 0, {+}, {-})$。显然，加法单位元 $0$ 就是 $\rcrat(0)$，而加法逆元 ${-} : \RC \to \RC$ 则是利用 \cref{RC-extend-Q-Lipschitz} 将 $\Q$ 上的加法逆元 ${-} : \Q \to \Q$ 扩展而得，Lipschitz 常数为 $1$。加法本身的定义则需要多花一些功夫。

\begin{lem} \label{RC-binary-nonexpanding-extension}
  假设 $f : \Q \times \Q \to \Q$ 满足，对所有 $q, r, s : \Q$，
  %
  \begin{equation*}
    |f(q, s) - f(r, s)| \leq |q - r|
    \qquad\text{and}\qquad
    |f(q, r) - f(q, s)| \leq |r - s|.
  \end{equation*}
  %
  那么存在函数 $\bar{f} : \RC \times \RC \to \RC$，使得对所有 $q, r : \Q$ 有 $\bar{f}(\rcrat(q), \rcrat(r)) = f(q,r)$。此外，对所有 $u, v, w : \RC$ 和 $q : \Qp$，
  %
  \begin{equation*}
    u \close\epsilon v \Rightarrow \bar{f}(u,w) \close\epsilon \bar{f}(v,w)
    \quad\text{and}\quad
    v \close\epsilon w \Rightarrow \bar{f}(u,v) \close\epsilon \bar{f}(u,w).
  \end{equation*}
\end{lem}

\begin{proof}
  我们使用 $(\RC, {\closesym})$-递归来构造 $\bar{f}$ 的柯里化形式，即一个映射
  $\RC \to A$，其中 $A$ 是取值于实数的非扩张\index{function!non-expanding}\index{non-expanding function}函数构成的空间：
  %
  \begin{equation*}
    A \defeq
    \setof{ h : \RC \to \RC |
      \fall{\epsilon : \Qp} \fall{u, v : \RC}
      u \close\epsilon v \Rightarrow h(u) \close\epsilon h(v)
    }.
  \end{equation*}
  %
  我们还需要 $A$ 上的一个合适的 $\bsim_\epsilon$，其定义为
  %
  \begin{equation*}
    (h \bsim_\epsilon k) \defeq \fall{u : \RC} h(u) \close\epsilon k(u).
  \end{equation*}
  %
  显然，若 $\fall{\epsilon : \Qp} h \bsim_\epsilon k$ 则对所有 $u : \RC$ 都有 $h(u) = k(u)$，因此 $\bsim$ 是分离的。

  对于基础情况，我们将 $\bar{f}(\rcrat(q)) : A$（其中 $q : \Q$）定义为 Lipschitz 映射 $\lam{r} f(q,r)$ 从 $\Q \to \Q$ 到 $\RC \to \RC$ 的延拓，这由 \cref{RC-extend-Q-Lipschitz} 以 Lipschitz 常数~$1$ 构造而成。
  接着，对于一个 Cauchy 逼近 $x$，我们将 $\bar{f}(\rclim(x)) : \RC \to \RC$ 定义为
  %
  \begin{equation*}
    \bar{f}(\rclim(x))(v) \defeq \rclim (\lam{\epsilon} \bar{f}(x_\epsilon)(v)).
  \end{equation*}
  %
  为使此定义有效，$\lam{\epsilon} \bar{f}(x_\epsilon)(v)$ 应当是一个 Cauchy 逼近。为此，考虑任意 $\delta, \epsilon : \Q$。
  则由假设有 $\bar{f}(x_\delta) \bsim_{\delta + \epsilon} \bar{f}(x_\epsilon)$，于是
  $\bar{f}(x_\delta)(v) \close{\delta + \epsilon} \bar{f}(x_\epsilon)(v)$。
  此外，$\bar{f}(\rclim(x))$ 是非扩张的，因为由归纳假设，$\bar{f}(x_\epsilon)$ 是非扩张的。
  事实上，若 $u \close\epsilon v$，则对所有 $\epsilon : \Q$，
  %
  \begin{equation*}
    \bar{f}(x_{\epsilon/3})(u) \close{\epsilon/3} \bar{f}(x_{\epsilon/3})(v),
  \end{equation*}
  %
  因此，由 $\closesym$ 的第四个构造子可得 $\bar{f}(\rclim(x))(u) \close\epsilon \bar{f}(\rclim(x))(v)$。

  我们还需要验证另外四个条件，这里仅举例说明其中一个。
  假设 $\epsilon : \Qp$，且对某个 $\delta : \Qp$ 有 $\rcrat(q) \close{\epsilon - \delta}
  y_\delta$ 与 $\bar{f}(\rcrat(q)) \bsim_{\epsilon - \delta} \bar{f}(y_\delta)$。
  为证明 $\bar{f}(\rcrat(q)) \bsim_\epsilon \bar{f}(\rclim(y))$，考虑任意 $v : \RC$ 并注意到
  %
  \begin{equation*}
    \bar{f}(\rcrat(q))(v) \close{\epsilon - \delta} \bar{f}(y_\delta)(v).
  \end{equation*}
  %
  因此，由 $\closesym$ 的第二个构造子，我们得到
  \narrowequation{\bar{f}(\rcrat(q))(v) \close\epsilon \bar{f}(\rclim(y))(v)}
  如所欲证。
\end{proof}

我们可以将 \cref{RC-binary-nonexpanding-extension} 应用于任何在每个变量上分别非扩张的二元有理函数。
加法就是这样一个函数，因此我们得到了 ${+} : \RC \times \RC \to \RC$。
\indexdef{addition!of Cauchy reals}%
此外，只要要求延拓在每个变量上非扩张，延拓就是唯一的；并且正如单变量情形一样，有理数上的等式可延拓为实数上的等式。
由于非扩张映射的复合仍是非扩张的，我们可以断定加法满足通常的性质，如交换律与结合律。
\index{associativity!of addition!of Cauchy reals}%
因此，$(\RC, 0, {+}, {-})$ 是一个交换群。

我们也可以将 \cref{RC-binary-nonexpanding-extension} 应用于函数 $\min : \Q \times
\Q \to \Q$ 和 $\max : \Q \times \Q \to \Q$，从而使 $\RC$ 成为一个格。
$\RC$ 上的偏序 $\leq$ 通过 $\max$ 定义为
%
\symlabel{leq-RC}
\index{order!non-strict}%
\index{non-strict order}%
\begin{equation*}
  (u \leq v) \defeq (\max(u, v) = v).
\end{equation*}
%
关系 $\leq$ 是偏序，因为它在 $\Q$ 上就是偏序，而偏序公理可以用关于 $\min$ 和 $\max$ 的等式来表达，所以这些性质传递到了 $\RC$。

\index{absolute value}%
另一个通过同样方法可延拓到 $\RC$ 的函数是绝对值 $|{\blank}|$。
同样，它具有预期的性质，因为这些性质从 $\Q$ 传递到了 $\RC$。

\symlabel{lt-RC}
由 $\leq$ 定义严格序 $<$ 如下：
\index{strict!order}%
\index{order!strict}%
%
\begin{equation*}
  (u < v) \defeq \exis{q, r : \Q} (u \leq \rcrat(q)) \land (q < r) \land (\rcrat(r) \leq v).
\end{equation*}
%
也就是说，当 merely 存在一对有理数 $q < r$ 使得 $x \leq
\rcrat(q)$ 且 $\rcrat(r) \leq v$ 时，$u < v$ 成立。
不难验证 $<$ 不自反且传递，并满足有序域的其他预期性质。
Archimedean 原理直接由 $<$ 的定义得出。

\index{ordered field!archimedean}%


\begin{thm}[$\RC$ 的Archimedean 原理] \label{RC-archimedean}
  %
  对于所有满足 $u < v$ 的 $u, v : \RC$，mere 存在 $q : \Q$ 使得 $u < q < v$。
\end{thm}

\begin{proof}
  由 $u < v$，我们 merely 得到 $r, s : \Q$ 使得 $u \leq r < s \leq v$，此时可取 $q
  \defeq (r + s) / 2$。
\end{proof}

现在我们在 $\RC$ 上已具备足够的结构，可以用标准概念来表达 $u \close\epsilon v$。

\begin{lem}\label{thm:RC-le-grow}
  如果 $q:\Q$ 和 $u:\RC$ 满足 $u\le \rcrat(q)$，那么对于任意 $v:\RC$ 和 $\epsilon:\Qp$，若 $u\close\epsilon v$ 则 $v\le \rcrat(q+\epsilon)$。
\end{lem}

\begin{proof}
  注意到函数 $\max(\rcrat(q),\blank):\RC\to\RC$ 是 Lipschitz 的，其常数为 $1$。
  先考虑 $u=\rcrat(r)$ 为有理数的情形。
  为此，我们对 $v$ 施行归纳。
  如果 $v$ 是有理数，那么结论是显然的。
  如果 $v$ 是 $\rclim(y)$，我们归纳地假设：对任意 $\epsilon,\delta$，如果 $\rcrat(r)\close\epsilon y_\delta$，则 $y_\delta \le \rcrat(q+\epsilon)$，即 $\max(\rcrat(q+\epsilon),y_\delta)=\rcrat(q+\epsilon)$。

  现在假设给定 $\epsilon$ 且 $\rcrat(r)\close\epsilon \rclim(y)$，那么存在 $\theta$ 使得 $\rcrat(r)\close{\epsilon-\theta} \rclim(y)$ 成立，从而当 $\delta<\theta$ 时就有 $\rcrat(r)\close\epsilon y_\delta$。
  于是由归纳假设，对这样的 $\delta$ 有 $\max(\rcrat(q+\epsilon),y_\delta)=\rcrat(q+\epsilon)$。
  但根据定义，
  \[\max(\rcrat(q+\epsilon),\rclim(y)) \jdeq \rclim(\lam{\delta} \max(\rcrat(q+\epsilon),y_\delta)).\]
  由于一个最终常值的 Cauchy 逼近的极限就是该常值，我们有
  \[\max(\rcrat(q+\epsilon),\rclim(y)) = \rcrat(q+\epsilon),\]从而得到 $\rclim(y)\le \rcrat(q+\epsilon)$。

  现在考虑一般的 $u:\RC$。
  因为 $u\le \rcrat(q)$ 意味着 $\max(\rcrat(q),u)=\rcrat(q)$，结合假设 $u\close\epsilon v$ 以及 $\max(\rcrat(q),-)$ 的 Lipschitz 性质，可推出 $\max(\rcrat(q),v) \close\epsilon \rcrat(q)$。
  这样，既然 $\rcrat(q)\le \rcrat(q)$，由第一种情形得 $\max(\rcrat(q),v) \le \rcrat(q+\epsilon)$，再由 $\le$ 的传递性即得 $v\le \rcrat(q+\epsilon)$。
\end{proof}

\begin{lem}\label{thm:RC-lt-open}
  设 $q:\Q$ 和 $u:\RC$ 满足 $u<\rcrat(q)$。则：
  \begin{enumerate}
  \item 对任意 $v:\RC$ 和 $\epsilon:\Qp$，若 $u\close\epsilon v$，则 $v< \rcrat(q+\epsilon)$。\label{item:RCltopen1}
  \item 存在 $\epsilon:\Qp$，使得对任意 $v:\RC$，只要 $u\close\epsilon v$，就有 $v<\rcrat(q)$。\label{item:RCltopen2}
  \end{enumerate}
\end{lem}

\begin{proof}
  根据定义，$u<\rcrat(q)$ 意味着存在 $r:\Q$，满足 $r<q$ 且 $u\le \rcrat(r)$。
  于是由 \cref{thm:RC-le-grow}，对任意 $\epsilon$，如果 $u\close\epsilon v$，则 $v\le \rcrat(r+\epsilon)$。
  结论~\ref{item:RCltopen1} 直接可得，因为 $r+\epsilon<q+\epsilon$；至于~\ref{item:RCltopen2}，只需取任意 $\epsilon <q-r$ 即可。
\end{proof}

现在我们可以证明，辅助关系 $\closesym$ 正是我们所预期的。

\begin{thm} \label{RC-sim-eqv-le}
  \index{distance}%
  $\eqv{(u \close\epsilon v)}{(|u - v| < \rcrat(\epsilon))}$
  对所有 $u, v : \RC$ 和 $\epsilon : \Qp$ 成立。
\end{thm}

\begin{proof}
  减法与绝对值的 Lipschitz 性质表明：如果 $u\close\epsilon v$，那么 $|u-v| \close\epsilon |u-u| = 0$。
  因此，对于从左到右的方向，只需证明：如果 $u\close\epsilon 0$，那么 $|u|<\rcrat(\epsilon)$。
  我们对 $u$ 进行 $\RC$-归纳。

  如果 $u$ 是有理数，那么结论直接成立，因为绝对值与序延拓了 $\Qp$ 上的标准定义。
  如果 $u$ 是 $\rclim(x)$，那么由舍入性，存在 $\theta:\Qp$ 使得 $\rclim(x)\close{\epsilon-\theta} 0$。
  于是，由三角不等式，我们有 $x_{\theta/3} \close{\epsilon-2\theta/3} 0$，从而归纳假设给出 $|x_{\theta/3}|<\rcrat(\epsilon-2\theta/3)$。
  但是 $x_{\theta/3} \close{2\theta/3} \rclim(x)$，因此由 Lipschitz 性质得到 $|x_{\theta/3}| \close{2\theta/3} |\rclim(x)|$，这样一来，\cref{thm:RC-lt-open}\ref{item:RCltopen1} 蕴涵 $|\rclim(x)|<\rcrat(\epsilon)$。

  对于另一个方向，我们对 $u$ 和 $v$ 进行 $\RC$-归纳。
  如果两者都是有理数，那么这就是 $\closesym$ 的第一个构造子。

  如果 $u$ 是 $\rcrat(q)$ 而 $v$ 是 $\rclim(y)$，我们归纳地假设：对于任意 $\epsilon,\delta$，如果 $|\rcrat(q)-y_\delta|<\rcrat(\epsilon)$ 那么 $\rcrat(q) \close{\epsilon} y_\delta$。
  取定一个满足 $|\rcrat(q) - \rclim(y)|<\rcrat(\epsilon)$ 的 $\epsilon$。
  由于 $\Q$ 在 $\RC$ 中序稠密，存在 $\theta<\epsilon$ 使得 $|\rcrat(q) - \rclim(y)|<\rcrat(\theta)$。
  现在，对任意 $\delta,\eta$ 我们有 $\rclim(y)\close{2\delta} y_\delta$，因此由 Lipschitz 性质，
  \[ |\rcrat(q) - \rclim(y)| \close{\delta+\eta} |\rcrat(q) - y_\delta|. \]
  于是，由 \cref{thm:RC-lt-open}\ref{item:RCltopen1}，我们得到 $|\rcrat(q) - y_\delta| < \rcrat(\theta+2\delta)$。
  那么由归纳假设，$\rcrat(q) \close{\theta+2\delta} y_\delta$，从而由三角不等式得到 $\rcrat(q)\close{\theta+4\delta} \rclim(y)$。
  因此，只需取 $\delta \defeq (\epsilon-\theta)/4$ 即可。

  剩下的两种情况完全类似。
\end{proof}

\indexdef{multiplication!of Cauchy reals}%
接下来，我们希望为 $\RC$ 配备乘法结构。
对每个 $q : \Q$，映射 $r \mapsto q \cdot r$ 是 Lipschitz 的，其常数为\footnote{我们定义 Lipschitz 常数为\emph{正的}有理数。} $|q| + 1$，因此我们可以将其延拓为实数上乘以 $q$ 的乘法。
所以 $\RC$ 是 $\Q$ 上的向量空间\index{vector!space}。
一般来说，我们可以将实数的乘法定义为
%
\begin{equation}
  u \cdot v \defeq
  {\textstyle \frac{1}{2}} \cdot ((u + v)^2 - u^2 - v^2),\label{mult-from-square}
\end{equation}
%
因此我们只需要平方函数\index{squaring function} $u \mapsto u^2$ 作为映射 $\RC \to \RC$。
平方函数不是 Lipschitz 映射，但在每个有界域上它是 Lipschitz 的，据此可以将其拼接起来。
定义开区间和闭区间为
%
\indexdef{interval!open and closed}%
\indexdef{open!interval}%
\indexdef{closed!interval}%
\begin{equation*}
  [u,v] \defeq \setof{ x : \RC | u \leq x \leq v }
  \qquad\text{and}\qquad
  (u,v) \defeq \setof{ x : \RC | u < x < v }.
\end{equation*}
%
虽然严格来说，$[u,v]$ 或 $(u,v)$ 的元素是一个 Cauchy 实数连同其一个证明，但由于后者仅居于一个命题，故无关紧要。
因此，正如处理子集类型时的常见做法，每当 $x:\RC$ 满足 $u\leq x \leq v$ 时，我们通常就简记为 $x:[u,v]$，其余类似。

\begin{thm} \label{RC-squaring}
  %
  存在唯一的函数 ${(\blank)}^2 : \RC \to \RC$，它延拓有理数的平方映射 $q \mapsto
  q^2$，并且满足
  %
  \begin{equation*}
    \fall{n : \N}
    \fall{u, v : [-n, n]}
    |u^2 - v^2| \leq 2 \cdot n \cdot |u - v|.
  \end{equation*}
\end{thm}

\begin{proof}
  我们首先观察到，对于每个 $u : \RC$，总存在 $n : \N$ 使得 $-n \leq u \leq n$，参见 \cref{ex:traditional-archimedean}，因此映射
  %
  \begin{equation*}
    e : \Parens{\sm{n : \N} [-n, n]} \to \RC
    \qquad\text{defined by}\qquad
    e(n, x) \defeq x
  \end{equation*}
  %
  是满射。接着，对于每个 $n : \N$，平方映射
  %
  \begin{equation*}
    s_n : \setof{ q : \Q | -n \leq q \leq n } \to \Q
    \qquad\text{defined by}\qquad
    s_n(q) \defeq q^2
  \end{equation*}
  %
  具有 Lipschitz 常数 $2 n$，于是我们可以利用 \cref{RC-extend-Q-Lipschitz} 将其延拓为一个具有 Lipschitz 常数 $2 n$ 的映射 $\bar{s}_n : [-n, n] \to \RC$，具体细节参见 \cref{RC-Lipschitz-on-interval}。这些映射 $\bar{s}_n$ 是相容的：如果对某个 $m, n : \N$ 有 $m < n$，那么 $s_n$ 在 $[-m, m]$ 上的限制必须与 $s_m$ 一致，因为二者都是 Lipschitz 的，因此在~\cref{RC-continuous-eq} 的意义下是连续的。因此，根据 \cref{lem:images_are_coequalizers}，映射
  %
  \begin{equation*}
    \Parens{\sm{n : \N} [-n, n]} \to \RC,
    \qquad\text{given by}\qquad
    (n, x) \mapsto s_n(x)
  \end{equation*}
  %
  可以唯一地通过 $\RC$ 分解，从而给出我们想要的函数。
\end{proof}

至此，我们已经有了实数的环结构和 Archimedean 序。为了确立 $\RC$ 是一个 Archimedean 有序域，还需要证明逆元的存在性。

\begin{thm}
  \index{apartness}%
  一个 Cauchy 实数可逆当且仅当它与零分离。
\end{thm}

\begin{proof}
  首先，假设 $u : \RC$ 有一个逆元 $v : \RC$。根据 Archimedean 原理，存在 $q : \Q$ 使得 $|v| < q$。那么 $1 = |u v| < |u| \cdot v < |u| \cdot q$，因此 $|u| > 1/q$，也就是说 $u \apart 0$。

  反之，与 \cref{RC-squaring} 中平方函数的构造类似，我们通过粘合函数来构造逆映射
  %
  \begin{equation*}
    ({\blank})^{-1} : \setof{ u : \RC | u \apart 0 } \to \RC
  \end{equation*}
  %
  。我们仅概述主要步骤。对于每个 $q : \Q$，令
  %
  \begin{equation*}
    [q, \infty) \defeq \setof{u : \RC | q \leq u}
    \qquad\text{and}\qquad
    (-\infty, q] \defeq \setof{u : \RC | u \leq -q}.
  \end{equation*}
  %
  那么，当 $q$ 取遍 $\Qp$ 时，类型 $(-\infty, q]$ 和 $[q, \infty)$ 共同覆盖了 $\setof{u : \RC | u \apart 0}$。在每个这样的 $[q, \infty)$ 和 $(-\infty, q]$ 上，应用 \cref{RC-extend-Q-Lipschitz} 即可得到具有 Lipschitz 常数 $1/q^2$ 的逆函数。最后，\cref{lem:images_are_coequalizers} 保证逆函数可唯一地经由 $\setof{ u : \RC | u
    \apart 0 }$ 分解。
\end{proof}

我们用一个定理来总结 $\RC$ 的代数结构。

\begin{thm} \label{RC-archimedean-ordered-field}
  Cauchy 实数构成一个 Archimedean 有序域。
\end{thm}

\subsection{Cauchy 实数是 Cauchy 完备的}
\label{sec:cauchy-reals-cauchy-complete}

我们通过在有理数 $\Q$ 上对 Cauchy 逼近取极限来构造 $\RC$，因此 $\RC$ 理应是 Cauchy 完备的。得益于 \cref{RC-sim-eqv-le}，在 $\RC$ 构造中定义的 Cauchy 逼近 $x : \Qp \to \RC$，与 \cref{defn:cauchy-approximation} 中定义的（适应于 $\RC$ 的）Cauchy 逼近之间没有区别。

这样一来，给定一个 Cauchy 逼近 $x : \Qp \to \RC$，自然期望 $\rclim(x)$ 就是它的极限，这里极限的概念是按 \cref{defn:cauchy-approximation} 定义的。而根据 \cref{RC-sim-eqv-le} 和 \cref{thm:RC-sim-lim-term}，情况正是如此。我们已经证明了：

\begin{thm}
  $\RC$ 中的每个 Cauchy 逼近都有一个极限。
\end{thm}

若某 Archimedean 有序域中每个 Cauchy 逼近都有极限，则称之为
\define{Cauchy 完备的（Cauchy complete）}。
\indexdef{Cauchy!completeness}%
\indexdef{complete!ordered field, Cauchy}%
\index{ordered field}%
Cauchy 实数是这样的最小域。

\begin{thm} \label{RC-initial-Cauchy-complete}
  Cauchy 实数可以嵌入到每个 Cauchy 完备的 Archimedean 有序域中。
\end{thm}

\begin{proof}
  \index{limit!of a Cauchy approximation}%
  设 $F$ 是一个 Cauchy 完备的 Archimedean 有序域。由于极限是唯一的，
  存在一个算子 $\lim$，将 $F$ 中的 Cauchy 逼近映为其极限。我们
  通过 $(\RC, {\closesym})$-递归定义嵌入 $e : \RC \to F$ 如下
  %
  \begin{equation*}
    e(\rcrat(q)) \defeq q
    \qquad\text{and}\qquad
    e(\rclim(x)) \defeq \lim (e \circ x).
  \end{equation*}
  %
  $F$ 上的一个合适的 $\bsim$ 为
  %
  \begin{equation*}
    (a \bsim_\epsilon b) \defeq |a - b| < \epsilon.
  \end{equation*}
  %
  因为 $F$ 是 Archimedean 的，所以这是一个分离关系。$(\RC, {\closesym})$-递归的其余条款
  很容易验证。此外还需要验证 $e$ 是一个固定有理数的有序域嵌入。
\end{proof}

\index{real numbers!Cauchy|)}%

\section{Cauchy 实数与 Dedekind 实数的比较}
\label{sec:comp-cauchy-dedek}

\index{real numbers!Dedekind|(}%
\index{real numbers!Cauchy|(}%
\index{depression|(}

下面我们也讨论一下 Cauchy 实数与 Dedekind 实数之间的关系。由
\cref{RC-archimedean-ordered-field}，$\RC$ 是一个 Archimedean 有序域。它对于 $\Omega$ 也是
\define{可容许的}\index{ordered field!admissible}，这很容易验证。（若 $\Omega$ 为初始
$\sigma$-帧
\index{initial!sigma-frame@$\sigma$-frame}%
\index{sigma-frame@$\sigma$-frame!initial}%
则只需一个简单的归纳，而在其他情况下则是显然的。）
因此，由 \cref{RD-final-field}，存在一个有序域嵌入
%
\begin{equation*}
  \RC \to \RD
\end{equation*}
%
该嵌入固定有理数。
（我们也可以从 \cref{RC-initial-Cauchy-complete,RD-cauchy-complete} 得到这个嵌入。）
一般来说，如果没有进一步的假设，我们不期望 $\RC$ 与 $\RD$ 重合。

\begin{lem} \label{lem:untruncated-linearity-reals-coincide}
  %
  若对每个 $x : \RD$，都仅仅存在
  %
  \begin{equation}
    \label{eq:untruncated-linearity}
    c : \prd{q, r : \Q} (q < r) \to (q < x) + (x < r)
  \end{equation}
  %
  则 Cauchy 实数与Dedekind 实数重合。
\end{lem}

\begin{proof}
  注意，\eqref{eq:untruncated-linearity} 中的类型是 \eqref{eq:RD-linear-order} 的一个未截断变体，后者断言 $<$ 是弱线性序。
  我们已知 $\RC$ 可嵌入 $\RD$，因此只需证明每个Dedekind 实数仅仅是有理数柯西序列\index{Cauchy!sequence}的极限。

  考虑任意的 $x : \RD$。由假设，仅仅存在满足引理陈述中条件的 $c$；又由分割\index{cut!Dedekind}的非空性，仅仅存在 $a, b : \Q$ 使得 $a < x < b$。
  我们用递归构造序列\index{sequence} $f : \N \to \setof{ \pairr{q, r} \in \Q \times \Q | q < r }$：
  %
  \begin{enumerate}
  \item 令 $f(0) \defeq \pairr{a, b}$。
  \item 假设 $f(n)$ 已定义为 $\pairr{q_n, r_n}$ 且 $q_n < r_n$。
    定义 $s \defeq (2 q_n + r_n)/3$ 与 $t \defeq (q_n + 2 r_n)/3$。于是 $c(s,t)$ 可判定 $s < x$ 与 $x < t$ 二者孰是。
    若判定为 $s < x$，则令 $f(n+1) \defeq
    \pairr{s, r_n}$；否则令 $f(n+1) \defeq \pairr{q_n, t}$。
  \end{enumerate}
  %
  将序列~$f$ 的第 $n$ 项记作 $\pairr{q_n, r_n}$。不难看出，对所有 $n : \N$，有 $q_n < x < r_n$ 及 $|q_n - r_n| \leq (2/3)^n \cdot |q_0 - r_0|$。
  因此 $q_0, q_1, \ldots$ 和 $r_0, r_1, \ldots$ 都是收敛到Dedekind 分割~$x$ 的柯西序列。
  我们已证明，对每个 $x : \RD$，都仅仅存在收敛到 $x$ 的柯西序列。
\end{proof}

该引理表明，可数选择公理或排中律均足以保证 $\RC$ 与 $\RD$ 重合。

\begin{cor} \label{when-reals-coincide}
  \index{axiom!of choice!countable}%
  \index{excluded middle}%
  若排中律或可数选择公理成立，则 $\RC$ 与 $\RD$ 等价。
\end{cor}

\begin{proof}
  若排中律成立，则 $(x < y) \to (x < z) + (z < y)$ 可证：要么 $x < z$，要么 $\lnot (x < z)$。
  前一种情况即已完成，而后一种情况下，由 $z \leq x < y$ 得 $z < y$。
  因此，我们得到~\eqref{eq:untruncated-linearity}，从而可以应用 \cref{lem:untruncated-linearity-reals-coincide}。

  假设可数选择公理成立。集合 $S = \setof{ \pairr{q, r} \in \Q \times \Q | q <
    r }$ 与 $\N$ 等价，因此可以对``$x$ 可定位''这一陈述
  %
  \begin{equation*}
    \fall{\pairr{q, r} : S} (q < x) \lor (x < r).
  \end{equation*}
  %
  应用可数选择公理。
  注意，$(q < x) \lor (x < r)$ 可表达为一个存在性陈述 $\exis{b :
    \bool} (b = \bfalse \to q < x) \land (b = \btrue \to x < r)$。该选择函数的柯里化形式
  恰好就是~\eqref{eq:untruncated-linearity}，从而
  \cref{lem:untruncated-linearity-reals-coincide} 再次适用。
\end{proof}

\index{real numbers!Dedekind|)}%
\index{real numbers!Cauchy|)}%
\index{real numbers!agree}%

\index{depression|)}

\section{区间的紧性}
\label{sec:compactness-interval}

\index{mathematics!classical|(}%
\index{mathematics!constructive|(}%

我们此前已指出，对实数的构造与经典逻辑完全兼容。
因此，通过假设排中律~\eqref{eq:lem} 和选择公理~\eqref{eq:ac}，我们可以发展经典分析\index{classical!analysis}\index{analysis!classical}，其结果实质上等同于照搬任何一本标准分析教科书。

\index{analysis!constructive}%
\index{constructive!analysis}%
尽管如此，任何关注计算的人——例如数值分析师——都应当对在计算上有意义的环境中发展分析学感到好奇。
在构造性框架下进行分析甚至是可能的，这一点已由~\cite{Bishop1967} 所证明。
作为经典分析与构造性分析之异同的一个示例，我们将仅简要讨论一个主题——闭区间 $[0,1]$ 的紧性以及围绕这一概念的若干定理。

在构造性数学中，经典等价的概念往往会产生分叉，这是一个常见现象，而紧性也不例外。
最常用的三种紧性概念是：
%
\indexdef{compactness}%


\begin{enumerate}
\item \define{度量紧的（metrically compact）}：“Cauchy 完备且完全有界”，
  \indexdef{metrically compact}%
  \indexdef{compactness!metric}%
\item \define{Bolzano--Weierstra\ss{} 紧的（Bolzano--Weierstra\ss{} compact）}：“每个序列都有收敛子列”，
  \index{compactness!Bolzano--Weierstrass@Bolzano--Weierstra\ss{}}%
  \indexsee{Bolzano--Weierstrass@Bolzano--Weierstra\ss{}}{compactness}%
  \index{sequence}%
\item \define{Heine--Borel 紧的（Heine--Borel compact）}：“每个开覆盖都有有限子覆盖”。
  \index{compactness!Heine--Borel}%
  \indexsee{Heine--Borel}{compactness}%
\end{enumerate}


%
它们在经典数学中都是等价的。
让我们看看它们在同伦类型论中表现如何。我们可以使用Dedekind 实数或 Cauchy 实数，因此直接将实数记作~$\R$。我们先回顾几个基本定义。

\indexsee{space!metric}{metric space}
\index{metric space|(}%

\begin{defn} \label{defn:metric-space}
  一个\define{度量空间（metric space）}
  \indexdef{metric space}%
  $(M, d)$ 由一个集合 $M$ 连同一个映射 $d : M \times M \to \R$ 组成，满足对所有的 $x, y, z : M$，
  %
  \begin{align*}
    d(x,y) &\geq 0, &
    d(x,y) &= d(y,x), \\
    d(x,y) &= 0 \Leftrightarrow x = y, &
    d(x,z) &\leq d(x,y) + d(y,z).
  \end{align*}
  %
\end{defn}

\begin{defn} \label{defn:complete-metric-space}
  $M$ 中的一个\define{Cauchy 逼近（Cauchy approximation）}
  \index{Cauchy!approximation}%
  是一个序列 $x : \Qp \to M$，满足
  %
  \begin{equation*}
    \fall{\delta, \epsilon} d(x_\delta, x_\epsilon) < \delta + \epsilon.
  \end{equation*}
  %
  \index{limit!of a Cauchy approximation}%
  一个 Cauchy 逼近 $x : \Qp \to M$ 的\define{极限（limit）}是一个点 $\ell : M$，满足
  %
  \begin{equation*}
    \fall{\epsilon, \theta : \Qp} d(x_\epsilon, \ell) < \epsilon + \theta.
  \end{equation*}
  %
  \indexdef{metric space!complete}%
  \indexdef{complete!metric space}%
  若一个度量空间中每个 Cauchy 逼近都有极限，则称其为\define{完备度量空间（complete metric space）}。
\end{defn}

\begin{defn} \label{defn:total-bounded-metric-space}
  对于一个正有理数 $\epsilon$，度量空间 $(M, d)$ 中的一个\define{$\epsilon$-网（$\epsilon$-net）}
  \indexdef{epsilon-net@$\epsilon$-net}%
  是
  %
  \begin{equation*}
    \sm{n : \N}{x_1, \ldots, x_n : M}
    \fall{y : M} \exis{k \leq n} d(x_k, y) < \epsilon.
  \end{equation*}
  %
  的一个元素。换言之，这是一组有限点列 $x_1, \ldots, x_n$，使得 $M$ 中的每一点仅位于某个 $x_k$ 的 $\epsilon$ 邻域内。

  一个度量空间 $(M, d)$ 是\define{完全有界的（totally bounded）}
  \indexdef{totally bounded metric space}%
  \indexdef{metric space!totally bounded}%
  ，若它对每个 $\epsilon$ 都有 $\epsilon$-网：
  %
  \begin{equation*}
    \prd{\epsilon : \Qp}
    \sm{n : \N}{x_1, \ldots, x_n : M}
    \fall{y : M} \exis{k \leq n} d(x_k, y) < \epsilon.
  \end{equation*}
\end{defn}

\begin{rmk}
  在完全有界性的定义中，我们使用了一种不够严格的记法 $\sm{n : \N}{x_1, \ldots, x_n : M}$。严格来说，我们应该写成 $\sm{x : \lst{M}}$，其中 $\lst{M}$ 是从 \cref{sec:bool-nat} 引入的有限列表\index{type!of lists}的归纳类型。不过，那样会使其余部分的表述稍显繁琐。
\end{rmk}

请注意，在完全有界性的定义中，我们要求的是 $\epsilon$-网的纯粹存在性，而非仅存在性。这样我们就能得到一个函数，为每个 $\epsilon : \Qp$ 指派一个具体的 $\epsilon$-网。这样的函数或许可以称为``完全有界模''。一般而言，在将经典的度量概念移植到同伦类型论时，我们应当节制地使用命题截断，通常是为了避免要求一个从 $\R$ 到 $\Q$ 或 $\N$ 的非常值映射。例如，以下是``一致连续''的``正确''定义。

\begin{defn} \label{defn:uniformly-continuous}
  度量空间上的映射 $f : M \to \R$ 称为\define{一致连续的（uniformly continuous）}
  \indexdef{function!uniformly continuous}%
  \indexdef{uniformly continuous function}%
  ，若满足
  %
  \begin{equation*}
    \prd{\epsilon : \Qp}
    \sm{\delta : \Qp}
    \fall{x, y : M}
    d(x,y) < \delta \Rightarrow |f(x) - f(y)| < \epsilon.
  \end{equation*}
  %
  特别地，一致连续映射有一个\define{一致连续模（modulus of uniform continuity）}\indexdef{modulus!of uniform continuity}，这是一个为每个 $\epsilon$ 指派对应 $\delta$ 的函数。
\end{defn}

现在我们来证明 $[0,1]$ 在第一种意义下是紧的。

\begin{thm} \label{analysis-interval-ctb}
  \index{compactness!metric}%
  \index{interval!open and closed}%
  闭区间 $[0,1]$ 是完备且完全有界的。
\end{thm}

\begin{proof}
  给定 $\epsilon : \Qp$，存在 $k : \N$ 使得 $2/k < \epsilon$，因此我们可以取 $\epsilon$-网 $x_i = i/k$，其中 $i = 0, \ldots, k$。这是一个 $\epsilon$-网，因为对于每个 $y : [0,1]$，仅存在 $i$ 使得 $0 \leq i \leq k$ 且 $(i -
  1)/k < y < (i+1)/k$，从而 $|y - x_i| < 2/k < \epsilon$。

  为证 $[0,1]$ 的完备性，考虑一个 Cauchy 逼近 $x : \Qp \to [0,1]$，并令 $\ell$ 为其在 $\R$ 中的极限。由于 $\max$ 和 $\min$ 是 Lipschitz 映射，由 $r(x) \defeq \max(0, \min(1, x))$ 定义的收缩 $r : \R \to [0,1]$ 与 Cauchy 逼近的极限交换，因此有
  %
  \begin{equation*}
    r(\ell) =
    r (\lim x) =
    \lim (r \circ x) =
    \lim x =
    \ell,
  \end{equation*}
  %
  这意味着 $0 \leq \ell \leq 1$，如所需。
\end{proof}

于是，我们至少在同伦类型论中获得了一个好的紧性概念。遗憾的是，它局限于度量空间，因为完全有界是一个度量概念。我们稍后会考虑另外两个概念，但首先我们证明，在完全有界空间上的一致连续映射具有\define{上确界（supremum）}，
\indexsee{least upper bound}{supremum}%
即一个小于或等于所有其他上界的上界。

\begin{thm} \label{ctb-uniformly-continuous-sup}
  %
  \indexdef{supremum!of uniformly continuous function}%
  在完全有界度量空间 $(M, d)$ 上的一致连续映射 $f : M \to \R$ 有一个上确界 $m : \R$。对于每个 $\epsilon : \Qp$，都存在 $u : M$ 使得 $|m - f(u)| < \epsilon$。
\end{thm}

\begin{proof}
  令 $h : \Qp \to \Qp$ 为 $f$ 的一致连续模。
  我们如下定义一个逼近 $x : \Qp \to \R$：对于任意的 $\epsilon : \Q$，由 $M$ 的完全有界性可以得到一个 $h(\epsilon)$-网 $y_0, \ldots, y_n$。定义
  %
  \begin{equation*}
    x_\epsilon \defeq \max (f(y_0), \ldots, f(y_n)).
  \end{equation*}
  %
  我们断言 $x$ 是一个 Cauchy 逼近。考虑任意的 $\epsilon, \eta : \Q$，此时
  %
  \begin{equation*}
    x_\epsilon \jdeq \max (f(y_0), \ldots, f(y_n))
    \quad\text{and}\quad
    x_\eta \jdeq \max (f(z_0), \ldots, f(z_m))
  \end{equation*}
  %
  其中的 $y_0, \ldots, y_n$ 是某个 $h(\epsilon)$-网，$z_0, \ldots, z_m$ 是某个 $h(\eta)$-网。
  每一个 $z_i$ 仅仅与某个 $y_j$ 的距离在 $h(\epsilon)$ 以内，因此 $|f(z_i) - f(y_j)| <
  \epsilon$，由此可得
  %
  \begin{equation*}
    f(z_i) < \epsilon + f(y_j) \leq \epsilon + x_\epsilon,
  \end{equation*}
  %
  故 $x_\eta < \epsilon + x_\epsilon$。对称地，我们得到 $x_\eta < \eta +
  x_\eta$，因此 $|x_\eta - x_\epsilon| < \eta + \epsilon$。

  我们断言 $m \defeq \lim x$ 是 $f$ 的上确界。要证明对所有 $x : M$ 有 $f(x) \leq m$，只需证明 $\lnot (m < f(x))$。于是，反设 $m <
  f(x)$。存在 $\epsilon : \Qp$ 使得 $m + \epsilon < f(x)$。但此时，仅由参与定义 $x_\epsilon$ 的某个 $y_i$ 便得 $|f(x) - f(y_i)| <
  \epsilon$，因此 $m < f(x) - \epsilon < f(y_i) \leq m$，矛盾。

  最后，我们通过证明 $m$ 满足定理的第二部分来完成证明，因为
  这样 $m$ 就自动成为最小上界。给定任意的 $\epsilon : \Qp$，一方面
  $|m - f(x_{\epsilon/2})| < 3 \epsilon/4$，另一方面，仅由参与定义 $x_{\epsilon/2}$ 的某个 $y_i$ 便得 $|f(x_{\epsilon/2}) - f(y_i)| <
  \epsilon/4$，于是取 $u \defeq y_i$，由三角不等式即得 $|m - f(u)| < \epsilon$。
\end{proof}

现在，如果在 \cref{ctb-uniformly-continuous-sup} 中我们还知道 $M$ 是完备的，我们或许可以寄望于将一致连续的假设弱化为连续，并将结论加强为存在一个使上确界得以达到的点。这些改进的通常证明依赖于如下事实：在一个完备且完全有界的空间中
%


\begin{enumerate}
\item 连续蕴涵一致连续，并且
\item 每个序列都有收敛的子序列。
\end{enumerate}


%
第一条陈述可由 Heine--Borel 紧性容易地推出，第二条则正是 Bolzano--Weierstra\ss{} 紧性。
\index{compactness!Bolzano--Weierstrass@Bolzano--Weierstra\ss{}}%
遗憾的是，这两条都或多或少存在问题。我们先证明 Bolzano--Weierstra\ss{} 紧性蕴含一种被称为\define{有限全知原理}的排中律特例：
\indexsee{axiom!limited principle of omniscience}{limited principle of omniscience}%
\indexdef{limited principle of omniscience}%
即对于每个 $\alpha : \N \to \bool$，
%
\begin{equation} \label{eq:lpo}
  \Parens{\sm{n : \N} \alpha(n) = \btrue} +
  \Parens{\prd{n : \N} \alpha(n) = \bfalse}.
\end{equation}
%
从计算的角度来说，我们并不期望这条原理成立，因为它要求我们判断一个函数是否有无穷多个取值为 $\bfalse$ 的值。

\begin{thm} \label{analysis-bw-lpo}
  %
  区间 $[0,1]$ 的 Bolzano--Weierstra\ss{} 紧性蕴涵有限全知原理。
  \index{compactness!Bolzano--Weierstrass@Bolzano--Weierstra\ss{}}%
\end{thm}

\begin{proof}
  给定任意 $\alpha : \N \to \bool$，定义序列\index{sequence} $x : \N \to [0,1]$ 如下：
  %
  \begin{equation*}
    x_n \defeq
    \begin{cases}
      0 & \text{if $\alpha(k) = \bfalse$ for all $k < n$,}\\
      1 & \text{if $\alpha(k) = \btrue$ for some $k < n$}.
    \end{cases}
  \end{equation*}
  %
  如果 Bolzano--Weierstra\ss{} 性质成立，则存在严格递增的 $f : \N \to
  \N$ 使得 $x \circ f$ 是柯西序列\index{Cauchy!sequence}。对于充分大的 $n:
  \N$，第 $n$ 项 $x_{f(n)}$ 与其极限的距离在 $1/6$ 以内。要么 $x_{f(n)} < 2/3$，
  要么 $x_{f(n)} > 1/3$。若 $x_{f(n)} < 2/3$，则 $x_n$ 收敛于 $0$，于是 $\prd{n : \N}
  \alpha(n) = \bfalse$。若 $x_{f(n)} > 1/3$，则 $x_{f(n)} = 1$，因此 $\sm{n : \N}
  \alpha(n) = \btrue$。
\end{proof}

也许 Bolzano--Weierstra\ss{} 紧性的丧失并不会让我们太过惋惜，但要舍弃 Heine--Borel 紧性则似乎要困难得多——经典数学和 Brouwer（布劳威尔）的直觉主义都接受了它，这一事实便足以印证。由于我们不想过多涉入一般拓扑，我们将采用基本开集来展开讨论。对于 $\R$ 而言，这些就是以有理数为端点的开区间。一个由某类型 $I$ 索引的此类区间族，就是一个映射
%
\begin{equation*}
  \mathcal{F} : I \to \setof{(q, r) : \Q \times \Q | q < r},
\end{equation*}
%
其想法是：满足 $q < r$ 的有理数对 $(q, r)$ 确定了类型 $\setof{ x : \R | q < x < r}$。为稍作方便起见，我们也允许退化区间，因此我们将\define{基本区间族}
\indexdef{family!of basic intervals}%
\indexdef{interval!family of basic}%
取为一个映射
%
\begin{equation*}
  \mathcal{F} : I \to \Q \times \Q.
\end{equation*}
%
更确切地说，一个族是一个依值对 $(I, \mathcal{F})$，而不只是
$\mathcal{F}$。一个\define{有限基本区间族}是由 $\setof{ m :
  \N | m < n}$（对于某个 $n : \N$）索引的族。我们通常用有限列表 $[(q_0, r_0), \ldots,
(q_{n-1}, r_{n-1})]$ 来表示它。最后，$(I, \mathcal{F})$ 的一个\define{有限子族}\indexdef{subfamily, finite, of intervals} 由一个指标列表 $[i_1, \ldots, i_n]$ 给出，该列表随后确定了有限族
$[\mathcal{F}(i_1), \ldots, \mathcal{F}(i_n)]$。

只要我们清楚有序对 $(q, r)$ 与它所对应的区间 $\setof{ x : \R | q < x < r}$ 之间的区别，便可以放心地用同一个记号 $(q, r)$ 来表示二者。
区间之间的\define{交}与\define{包含}关系可以用它们的端点来表示：
%
\symlabel{interval-intersection}
\symlabel{interval-subset}
\begin{align*}
  (q, r) \cap (s, t) &\ \defeq\  (\max(q, s), \min(r, t)),\\
  (q, r) \subseteq (s, t) &\ \defeq\ (q < r \Rightarrow s \leq q < r \leq t).
\end{align*}
%
我们说 $\intfam{i}{I}{(q_i, r_i)}$ \define{（逐点）覆盖} $[a,b]$，若
\indexdef{interval!pointwise cover}%
\indexdef{cover!pointwise}%
\indexdef{pointwise!cover}%
%
\begin{equation} \label{eq:cover-pointwise-truncated}
  \fall{x : [a,b]} \exis{i : I} q_i < x < r_i.
\end{equation}
%
$[0,1]$ 的 \define{Heine–Borel 紧性}
\indexdef{compactness!Heine--Borel}%
是指：对于 $[0,1]$ 的每一个覆盖族，都仅仅存在一个仍能覆盖 $[0,1]$ 的有限子族。

\index{depression}

\begin{thm} \label{classical-Heine-Borel}
  \index{excluded middle}%
  若排中律成立，则 $[0,1]$ 是 Heine–Borel 紧的。
\end{thm}

\begin{proof}
  为推出矛盾，假定有一个族 $\intfam{i}{I}{(a_i,
    b_i)}$ 覆盖了 $[0,1]$，但其任意有限子族都不能覆盖 $[0,1]$。
  我们构造一系列闭区间 $[q_n, r_n]$，它们是嵌套的，长度趋于 $0$，并且其中任何一个都不能被 $\intfam{i}{I}{(a_i, b_i)}$ 的有限子族覆盖。

  我们设 $[q_0, r_0] \defeq [0,1]$。
  假定已经构造好 $[q_n, r_n]$，令 $s
  \defeq (2 q_n + r_n)/3$ 与 $t \defeq (q_n + 2 r_n)/3$。
  区间 $[q_n, t]$ 和 $[s, r_n]$ 都被 $\intfam{i}{I}{(a_i, b_i)}$ 覆盖，但它们不可能同时拥有有限子覆盖，否则 $[q_n, r_n]$ 也会有有限子覆盖。
  要么 $[q_n, t]$ 有有限子覆盖，要么没有。
  若有，则设 $[q_{n+1}, r_{n+1}] \defeq [s, r_n]$；否则设 $[q_{n+1},
  r_{n+1}] \defeq [q_n, t]$。

  序列 $q_0, q_1, \ldots$ 和 $r_0, r_1, \ldots$ 都是柯西序列，它们收敛到同一个点 $x : [0,1]$，该点包含于每一个 $[q_n, r_n]$ 中。
  仅仅存在 $i : I$ 使得 $a_i < x < b_i$。
  由于区间 $[q_n, r_n]$ 的长度趋于零，存在 $n : \N$ 使得 $a_i < q_n \leq x
  \leq r_n < b_i$，但这就意味着 $[q_n, r_n]$ 被单个区间 $(a_i, b_i)$ 所覆盖，同时却又没有有限子覆盖。矛盾。
\end{proof}

没有了排中律，或是缺少一点布劳威尔直觉主义，我们似乎就束手无策了。
然而，$[0,1]$ 的 Heine–Borel 紧性\emph{是能够}在构造性框架下恢复的，而且其方式仍然与经典数学兼容！
为此，我们需要重新审视``覆盖''这一概念。
\eqref{eq:cover-pointwise-truncated} 的问题在于，截断的存在量词允许空间以任意杂乱的方式被覆盖，因此从计算的角度来说，我们便毫无机会``仅仅''提取出一个有限的子覆盖。
去掉截断，我们得到
%
\begin{equation} \label{eq:cover-pointwise}
  \prd{x : [0,1]} \sm{i : I} q_i < x < r_i,
\end{equation}
%
这或许能有所帮助，但它对覆盖的要求实在太高了。
按照这个定义，我们甚至无法证明 $(0,3)$ 和 $(2,5)$ 覆盖了 $[1,4]$，因为那相当于给出一个非常值的映射 $[1,4] \to \bool$，参见
\cref{ex:reals-non-constant-into-Z}。
这里，我们可以从``无点拓扑''
\index{pointfree topology}%
\index{topology!pointfree}%
（即 locale 理论）
\index{locale}%
中汲取经验：覆盖的概念应当基于开集来表达，而不涉及点。
这种对空间的``整体性''观点将使我们能够分析覆盖的概念，进而恢复 Heine–Borel 紧性。
locale 理论使用幂集，
\index{power set}%
我们可以通过假设命题大小重整来获得它；
\index{propositional!resizing}%
不过，我们可以转而借鉴 locale 理论的直谓近亲的思想，
\index{mathematics!predicative}%
它被称为``形式拓扑''。
\index{formal!topology}%

\index{acceptance|(}

假设我们有一个族 $\pairr{I, \mathcal{F}}$ 和一个区间 $(a, b)$。
我们如何在不涉及点的情况下表达``$(a,b)$ 被该族覆盖''这一事实呢？
这里有一个：如果 $(a, b)$ 等于某个 $\mathcal{F}(i)$，那么它被该族覆盖。
还有一个：如果 $(a,b)$ 被另一个族 $(J, \mathcal{G})$ 覆盖，并且每个 $\mathcal{G}(j)$ 又都被 $\pairr{I, \mathcal{F}}$ 覆盖，那么 $(a,b)$ 就被 $\pairr{I, \mathcal{F}}$ 覆盖。
请注意，我们正在列举可以用来\emph{推导}出 $\pairr{I, \mathcal{F}}$ 覆盖 $(a,b)$ 的\emph{规则}。
我们需要找到一组足够好的规则，并将它们转化为一个归纳定义。

\begin{defn} \label{defn:inductive-cover}
  %
  所谓\define{归纳覆盖 $\cover$}
  \indexdef{inductive!cover}%
  \indexdef{cover!inductive}%
  是一个纯粹关系
  %
  \begin{equation*}
    {\cover} : (\Q \times \Q) \to \Parens{\sm{I : \type} (I \to \Q \times \Q)} \to \prop
  \end{equation*}
  %
  它由以下规则归纳地定义，其中 $q, r, s, t$ 是有理数，而
  $\pairr{I, \mathcal{F}}$ 与 $\pairr{J, \mathcal{G}}$ 是基本区间族：
  %
  \begin{enumerate}

  \item \emph{自反性：}
    \index{reflexivity!of inductive cover}%
    $\mathcal{F}(i) \cover \pairr{I, \mathcal{F}}$ 对所有 $i : I$ 成立，

  \item \emph{传递性：}
    \index{transitivity!of inductive cover}%
    若 $(q, r) \cover \pairr{J, \mathcal{G}}$ 且 $\fall{j : J} \mathcal{G}(j) \cover \pairr{I,\mathcal{F}}$，
    则 $(q, r) \cover \pairr{I, \mathcal{F}}$，

  \item \emph{单调性：}
    \index{monotonicity!of inductive cover}%
    若 $(q, r) \subseteq (s, t)$ 且 $(s,t) \cover \pairr{I, \mathcal{F}}$ 则 $(q, r) \cover
    \pairr{I, \mathcal{F}}$，

  \item \emph{局部化：}
    \index{localization of inductive cover}%
    若 $(q, r) \cover (I, \mathcal{F})$ 则 $(q, r) \cap (s, t) \cover
    \intfam{i}{I}{(\mathcal{F}(i) \cap (s, t))}$。

  \item \label{defn:inductive-cover-interval-1}
    若 $q < s < t < r$，则 $(q, r) \cover [(q, t), (s, r)]$，

  \item \label{defn:inductive-cover-interval-2}
    $(q, r) \cover \intfam{u}{\setof{ (s,t) : \Q \times \Q | q < s < t < r}}{u}$。
  \end{enumerate}
\end{defn}

此定义应理解为一个高阶归纳类型，其中所列规则为点构造子，且该类型是 $(-1)$-截断的。
前四条规则具有一般性，直观上应不难理解。
后两条规则则是针对实直线特有的：一条是说，若两个区间有重叠，则它们可以覆盖一个区间；另一条是说，一个区间可以从内部被覆盖。
顺便指出，若 $r \leq q$，则根据最后一条规则，区间 $(q, r)$ 被空族覆盖。

归纳覆盖具有 Heine--Borel 性质，其证明需要下面一个引理。

\begin{lem} \label{reals-formal-topology-locally-compact}
  假设 $q < s < t < r$ 且 $(q, r) \cover \pairr{I, \mathcal{F}}$。那么纯粹存在 $\pairr{I, \mathcal{F}}$ 的一个有限子族，它归纳地覆盖 $(s, t)$。
\end{lem}

\begin{proof}
  我们对 $(q, r) \cover \pairr{I, \mathcal{F}}$ 进行归纳来证明此命题。共有六种情况：
  %
  \begin{enumerate}

  \item 自反性：若 $(q, r) = \mathcal{F}(i)$，则由单调性知，$(s, t)$ 被有限子族 $[\mathcal{F}(i)]$ 覆盖。

  \item 传递性：
    假设 $(q, r) \cover \pairr{J, \mathcal{G}}$ 且 $\fall{j : J} \mathcal{G}(j) \cover
    \pairr{I, \mathcal{F}}$。根据归纳假设，纯粹存在 $[\mathcal{G}(j_1), \ldots, \mathcal{G}(j_n)]$ 覆盖 $(s, t)$。
    再次应用归纳假设，每一个 $\mathcal{G}(j_k)$ 都被 $\pairr{I, \mathcal{F}}$ 的一个有限子族覆盖，我们可以将这些有限子族收集起来，构成一个覆盖 $(s, t)$ 的有限子族。

  \item 单调性：
    若 $(q, r) \subseteq (u, v)$ 且 $(u, v) \cover \pairr{I, \mathcal{F}}$，则可以对 $(u, v) \cover \pairr{I, \mathcal{F}}$ 应用归纳假设，因为此时有 $u <
    s < t < v$。

  \item 局部化：
    假设 $(q', r') \cover \pairr{I, \mathcal{F}}$ 且 $(q, r) = (q', r') \cap (a, b)$。
    由于 $q' < s < t < r'$，根据归纳假设，存在 $(s, t)$ 的一个有限子覆盖 $[\mathcal{F}(i_1), \ldots, \mathcal{F}(i_n)]$。同时我们也知道 $a < s <
    t < b$，因此 $(s, t) = (s, t) \cap (a, b)$ 被
    $[\mathcal{F}(i_1) \cap (a,b), \ldots, \mathcal{F}(i_n) \cap (a,b)]$ 覆盖，而后者正是
    $\intfam{i}{I}{(\mathcal{F}(i) \cap (a, b))}$ 的一个有限子族。

  \item 若对某个 $q < u < v < r$ 有 $(q, r) \cover [(q, v), (u, r)]$，则由单调性可得
    $(s, t) \cover [(q, v), (u, r)]$。

  \item 最后，由自反性即得 $(s, t) \cover \intfam{z}{\setof{ (u,v):\Q \times \Q | q < u < v < r}}{z}$。\qedhere
  \end{enumerate}
\end{proof}

若纯粹存在 $\epsilon : \Qp$ 使得 $(a - \epsilon, b +
\epsilon) \cover \pairr{I, \mathcal{F}}$，则称 \define{$\pairr{I, \mathcal{F}}$ 归纳地覆盖 $[a, b]$}。

\begin{cor} \label{interval-Heine-Borel}
  \index{compactness!Heine-Borel}%
  \index{interval!open and closed}%
  对于归纳覆盖而言，闭区间是 Heine--Borel 紧的。
\end{cor}

\begin{proof}
  假设 $[a, b]$ 被 $\pairr{I, \mathcal{F}}$ 归纳覆盖，即纯粹存在 $\epsilon : \Qp$ 使得 $(a - \epsilon, b + \epsilon) \cover \pairr{I, \mathcal{F}}$。
  由 \cref{reals-formal-topology-locally-compact}，存在
  $(a - \epsilon/2, b + \epsilon/2)$ 的一个有限子覆盖，因而它也是 $[a, b]$ 的有限子覆盖。
\end{proof}

从形式拓扑学\index{topology!formal}的经验可知，归纳覆盖的这些规则足以开展构造性的无点拓扑学。
不过，我们也能自行提供证据表明，它们确实是一个合理的概念。

\begin{thm} \label{inductive-cover-classical}
  \mbox{}
  %
  \begin{enumerate}
  \item 归纳覆盖同时也是逐点覆盖。
  \item 若假设排中律，则逐点覆盖同时也是归纳覆盖。
  \end{enumerate}
\end{thm}

\begin{proof}
  \mbox{}
  %
  \begin{enumerate}

  \item
    考虑一个基本区间族 $\pairr{I, \mathcal{F}}$，其中我们记 $(q_i,
    r_i) \defeq \mathcal{F}(i)$，$(a,b)$ 是一个由 $\pairr{I,
      \mathcal{F}}$ 归纳覆盖的区间，$x$ 满足 $a < x < b$。
    %
    我们通过对 $(a,b) \cover \pairr{I, \mathcal{F}}$ 作归纳来证明，仅仅存在指标 $i : I$ 使得 $q_i < x < r_i$。大部分情形都相当显然，因此我们只展示其中两种。
    若 $(a,b) \cover \pairr{I, \mathcal{F}}$ 由自反性得到，则仅仅存在某个 $i : I$ 使得 $(a,b) = (q_i, r_i)$，从而有 $q_i < x < r_i$。
    若 $(a,b)
    \cover \pairr{I, \mathcal{F}}$ 经 $\intfam{j}{J}{(s_j, t_j)}$ 由传递性得到，则由归纳假设，仅仅存在 $j : J$ 使得 $s_j < x < t_j$；
    继而由于 $(s_j, t_j) \cover \pairr{I, \mathcal{F}}$，再次利用归纳假设可知，仅仅存在某个 $i : I$ 使得 $q_i < x < r_i$。其他情形也同样直截了当。

  \item 假设 $\intfam{i}{I}{(q_i, r_i)}$ 点态覆盖 $(a, b)$。根据
    \cref{defn:inductive-cover} 的 \cref{defn:inductive-cover-interval-2}，只需证明每当
    $a < c < d < b$ 时，$\intfam{i}{I}{(q_i, r_i)}$ 都归纳覆盖 $(c, d)$，因此考虑这样的 $c$ 和 $d$。
    由 \cref{classical-Heine-Borel}
    可知，存在一个有限子族 $[i_1, \ldots, i_n]$ 已经点态覆盖了 $[c,
    d]$，从而也点态覆盖 $(c,d)$。设 $\epsilon : \Qp$ 为关于 $(q_{i_1}, r_{i_1}), \ldots, (q_{i_n}, r_{i_n})$ 的一个Lebesgue 数
    \index{Lebesgue number}
    ，如 \cref{ex:finite-cover-lebesgue-number} 所述。取正数 $k : \N$ 使得 $2 (d - c)/k
    < \min(1, \epsilon)$。对于 $0 \leq i \leq k$，令
    %
    \begin{equation*}
      c_k \defeq ((k - i) c + i d) / k.
    \end{equation*}
    %
    区间 $(c_0, c_2)$, $(c_1, c_3)$, \dots, $(c_{k-2}, c_k)$ 通过反复运用 \cref{defn:inductive-cover} 中的传递性及~\cref{defn:inductive-cover-interval-1}
    归纳覆盖
    $(c,d)$。由于它们的宽度都小于 $\epsilon$，每一个这样的区间都被包含在某个 $(q_i, r_i)$ 之中；
    于是利用传递性和单调性即可断定，$\intfam{i}{I}{(q_i, r_i)}$ 归纳覆盖了 $(c, d)$。 \qedhere
  \end{enumerate}
\end{proof}

前一定理的要旨在于，就经典数学而言，点态覆盖与归纳覆盖之间并无区别。
特别地，由于在同伦类型论中假定排中律是一致的，我们不可能构造出一个不是点态覆盖的归纳覆盖。
或者换一种说法，点态覆盖与归纳覆盖之间的区别，不在于它们覆盖了什么，而在于覆盖的\emph{证明}。

我们大可以像这样继续写下去，再写出一本书来，不过还是就此止步吧；希望以上讨论已为分析学可在同伦类型论中发展这一论断提供了充分的依据。
好奇的读者可以查阅 \cref{ex:mean-value-theorem} 以了解构造性版本的中间值定理。

\index{acceptance|)}

\index{mathematics!classical|)}%
\index{mathematics!constructive|)}%

\section{超现实数}
\label{sec:surreals}

\index{surreal numbers|(}%

本节我们将考察另一个高阶归纳-归纳类型的例子，它汇集了我们之前的许多线索：Conway 的域 $\NO$，即\emph{超现实数（surreal numbers）}~\cite{conway:onag}。
超现实数是（Dedekind ）实数（\cref{sec:dedekind-reals}）与序数（\cref{sec:ordinals}）的自然共同推广。
Conway 在运用排中律和选择公理的经典\index{mathematics!classical}数学框架下，将超现实数定义为一对超现实数的\emph{集合}，记为 $\surr L R$，并满足 $L$ 中的每个元素都严格小于 $R$ 中的每个元素。
这显然像是一个归纳定义，但若视其为归纳定义，会面临三个问题。

首先，该定义需要用到超现实数之间的（严格）序关系，因此这个关系必须与超现实数的类型 $\NO$ 同时定义。
（Conway 通过先定义\emph{游戏（game）}\index{game!Conway}避开了这个问题，游戏类似于超现实数，但省略了 $L$ 与 $R$ 之间的相容性条件。）
正如 Cauchy 实数中的关系 $\closesym$ 一样，这个同时定义\emph{先验地}既可以是归纳-归纳的，也可以是归纳-递归的。出于与为 $\closesym$ 做出该选择时相同的理由，我们在此也选择归纳-归纳定义。

此外，我们将分别（并相互归纳地）定义超现实数的严格序 $<$ 与非严格序 $\le$。
Conway 用 $\le$ 来定义 $<$，这种方式在经典意义下是合理的，但在构造性意义下则不然。
\index{mathematics!constructive}%
再者，一个否定形式的 $<$ 定义，将使其不适合作为高阶归纳类型构造子的假设（参见 \cref{sec:strictly-positive}）。

其次，Conway 说 $\surr L R$ 中的 $L$ 和 $R$ 应该是“超现实数的集合”，但若将其朴素地理解为一个谓词 $\NO\to\prop$，则它不是正的，因此不能用作归纳构造子的输入。
不过，这本身也不是对 Conway 原意的良好类型论翻译，因为在集合论中，超现实数构成一个真类，而集合 $L$ 和 $R$ 是真正的（小）集合，并非 $\NO$ 的任意子类。
在类型论中，这意味着 $\NO$ 将相对于一个宇宙 $\UU$ 来定义，但其本身属于下一个更高的宇宙 $\UU'$，就像序数的集合 $\ord$、基数的集合 $\card$、累积层次 $V$，甚至（在没有命题大小重整时的）Dedekind 实数一样。
\index{propositional!resizing}%
这样一来，我们就要求超现实数的“集合” $L$ 和 $R$ 是 $\UU$-小的，因此用某个以 $\UU$-小类型为指标的\emph{族}来表示它们是很自然的。
（这与我们在 \cref{sec:cumulative-hierarchy} 中对累积层次所做的完全一样。）
也就是说，超现实数的构造子将具有类型
\[ \prd{\LL,\RR:\UU} (\LL\to\NO) \to (\RR\to \NO) \to (\text{some condition}) \to \NO \]
这确实是严格正的。\index{strict!positivity}

最后，在给出 $\NO$ 及其序关系的相互定义后，Conway 规定，若 $x\le y$ 且 $y\le x$，则两个超现实数 $x$ 和 $y$ 是\emph{相等}的。
这自然可以理解为对一个“预超现实数”的集合按等价关系取商。
%(In set-theoretic foundations, one has to us an additional trick to deal with large equivalence classes.)
然而，在没有选择公理的情况下，这样的商构造会遇到与通常构造 Cauchy 实数时的商相同的问题：一对\emph{超现实数}的族将不再能给出一个新的超现实数 $\surr L R$，因为我们不一定能将 $L$ 和 $R$ “提升”为预超现实数的族。
当然，我们可以用与处理 Cauchy 实数相同的方式来解决这个问题，即使用一个\emph{高阶}归纳-归纳定义。

\begin{defn}\label{defn:surreals}
  \define{超现实数}的类型 $\NO$，
  \indexdef{surreal numbers}%
  \indexsee{number!surreal}{surreal numbers}%
  连同关系 $\mathord<:\NO\to\NO\to\type$ 和 $\mathord\le:\NO\to\NO\to\type$，由以下的高阶归纳-归纳方式定义。
  类型 $\NO$ 具有以下构造子。
  \begin{itemize}
  \item 对于任何 $\LL,\RR:\UU$ 以及函数 $\LL\to \NO$ 和 $\RR\to \NO$，若满足条件 $\fall{L:\LL}{R:\RR} x^L<x^R$，则存在一个超现实数 $x$。对于 $L:\LL$ 和 $R:\RR$，我们将这两个函数的值分别记作 $x^L$ 和 $x^R$。
  \item 对于任何满足 $x\le y$ 且 $y\le x$ 的 $x,y:\NO$，我们有 $\noeq(x,y):x=y$。
  \end{itemize}
  我们将第一个构造子的输入称为一个\define{分割}。
  \indexdef{cut!of surreal numbers}%
  如果 $x$ 是由一个分割构造出的超现实数，那么记号 $x^L$ 将隐含 $L:\LL$，类似地 $x^R$ 隐含 $R:\RR$。
  这样我们通常可以避免为索引类型 $\LL$ 和 $\RR$ 显式命名，当讨论涉及多个不同的分割时，这会很方便。
  遵循 Conway 的术语，我们称 $x^L$ 为 $x$ 的一个\emph{左选项（left option）}\indexdef{option of a surreal number}，称 $x^R$ 为一个\emph{右选项（right option）}。

  路径构造子意味着不同的分割可以定义出同一个超现实数。
  因此，谈论任意一个超现实数 $x$ 的左选项或右选项是无意义的，除非我们还知道 $x$ 是由某个特定分割定义的。
  因此，在下文中，我们会比如用“给定一个定义超现实数 $x$ 的分割”这样的说法，以区别于“给定一个超现实数 $x$”。

  关系 $\le$ 具有以下构造子。
  \index{non-strict order}%
  \index{order!non-strict}%
  \begin{itemize}
  \item 给定分别定义两个超现实数 $x$ 和 $y$ 的分割，如果对于所有 $L$ 有 $x^L<y$，并且对于所有 $R$ 有 $x<y^R$，那么 $x\le y$。
  \item 命题截断：
    对于任何 $x,y:\NO$，如果 $p,q:x\le y$，那么 $p=q$。
  \end{itemize}
  关系 $<$ 具有以下构造子。
  \index{strict!order}%
  \index{order!strict}%
  \begin{itemize}
    % Don't technically need x in the first one and y in the second one to be defined by cuts?
  \item 给定分别定义两个超现实数 $x$ 和 $y$ 的分割，如果存在某个 $L$ 使得 $x\le y^L$，那么 $x<y$。
  \item 给定分别定义两个超现实数 $x$ 和 $y$ 的分割，如果存在某个 $R$ 使得 $x^R\le y$，那么 $x<y$。
  \item 命题截断：对于任何 $x,y:\NO$，如果 $p,q:x<y$，那么 $p=q$。
  \end{itemize}
\end{defn}

\noindent
我们将其与 Conway 的定义进行比较：

\begin{itemize}\footnotesize
\item[-] 如果 $L,R$ 是任意两个数集，并且 $L$ 中没有任何成员 $\ge$ $R$ 中的任何成员，那么就存在一个数 $\surr L R$。
  所有的数都是以这种方式构造出来的。
\item[-] $x\ge y$ 当且仅当（没有 $x^R\le y$，且 $x\le$ 没有 $y^L$）。
\item[-] $x=y$ 当且仅当（$x \ge y$ 且 $y\ge x$）。
\item[-] $x>y$ 当且仅当（$x\ge y$ 且 $y\not\ge x$）。
\end{itemize}


若所有对象都是 [超现实] 数而非``游戏''\index{game!Conway}，则在 $x>y$ 的定义中纳入 $x\ge y$ 便非必要。
因此，Conway 的 $<$ 不过就是他的 $\ge$ 的否定，从而他关于 $\surr L R$ 成为超现实数的条件与我们的相同。
对 Conway 的 $\le$ 取否定，再消去双重否定，便得到了我们对 $<$ 的定义；之后我们便可以用 $<$ 来重新表述他的 $\le$，而无需使用否定。

我们立刻就能用许多超现实数来填充 $\NO$。
如同 Conway 那样，我们记
\symlabel{surreal-cut}
\[\surr{x,y,z,\dots}{u,v,w,\dots}\]
表示由分割所定义的超现实数，其中 $\LL\to\NO$ 和 $\RR\to\NO$ 是由 $x,y,z,\dots$ 和 $u,v,w,\dots$ 所描述的类型族。
当然，如果 $\LL$ 或 $\RR$ 是 $\emptyt$，我们就在记号中留空对应的部分。
这里与表示子集的标准记号 $\setof{x:A | P(x)}$ 发生了令人遗憾的冲突，但本节中我们不会使用后者。


\begin{itemize}
\item 我们递归地定义 $\iota_{\nat}:\nat\to\NO$ 如下：
  \begin{align*}
    \iota_{\nat}(0) &\defeq \surr{}{},\\
    \iota_{\nat}(\suc(n)) &\defeq \surr{\iota_{\nat}(n)}{}.
  \end{align*}
  也就是说，$\iota_{\nat}(0)$ 由分割 $\emptyt\to\NO$ 和 $\emptyt\to\NO$ 定义。
  类似地，$\iota_{\nat}(\suc(n))$ 由 $\unit\to\NO$（选出 $\iota_{\nat}(n)$）和 $\emptyt\to\NO$ 定义。
\item 类似地，我们利用符号情形递归原理（\cref{thm:sign-induction}）定义 $\iota_{\Z}:\Z\to\NO$：
  \begin{align*}
    \iota_{\Z}(0) &\defeq \surr{}{},\\
    \iota_{\Z}(n+1) &\defeq \surr{\iota_{\Z}(n)}{} & &\text{$n\ge 0$,}\\
    \iota_{\Z}(n-1) &\defeq \surr{}{\iota_{\Z}(n)} & &\text{$n\le 0$.}
  \end{align*}
\item 所谓\define{2进有理数（dyadic rational）}
  \indexdef{rational numbers!dyadic}%
  \indexsee{dyadic rational}{rational numbers, dyadic}%
  是指一对 $(a,n)$，其中 $a:\Z$ 且 $n:\nat$，并且满足若 $n>0$ 则 $a$ 是奇数。
  我们将其记作 $a/2^n$，并视其为对应的有理数。
  若用 $\Q_D$ 表示所有2进有理数的集合，我们通过对 $n$ 归纳来定义 $\iota_{\Q_D}:\Q_D\to\NO$：
  \begin{align*}
    \iota_{\Q_D}(a/2^0) &\defeq \iota_{\Z}(a),\\
    \iota_{\Q_D}(a/2^n) &\defeq \surr{\iota_{\Q_D}(a/2^n - 1/2^n)}{\iota_{\Q_D}(a/2^n + 1/2^n)},
    \quad \text{for $n>0$.}
  \end{align*}
  这里我们用到了如下事实：如果 $n>0$ 且 $a$ 是奇数，那么 $a/2^n \pm 1/2^n$ 是一个分母小于 $a/2^n$ 的2进有理数。
\item 我们定义 $\iota_{\RD}:\RD\to\NO$,\label{reals-into-surreals} 其中 $\RD$ 是 \cref{sec:dedekind-reals} 中（任意版本的）Dedekind 实数，如下：
  \begin{align*}
    \iota_{\RD}(x) &\defeq
    \surr{q\in\Q_D \text{ such that } q<x}{q\in\Q_D \text{ such that } x<q}.
  \end{align*}
  与前几种情形不同，当我们将2进有理数视为Dedekind 实数时，上式延拓 $\iota_{\Q_D}$ 这一点并非显然。
  这可由简单性定理（\cref{thm:NO-simplicity}）推出。
\item 回顾 \cref{sec:ordinals} 中由关系 $<$ 良序的\emph{序数（ordinal）}\index{ordinal} 类型 $\ord$，其中 $A<B$ 意味着对某个 $b:B$ 有 $A = \ordsl B b$。
  我们通过关于 $\ord$ 的良基递归（\cref{thm:wfrec}）定义 $\iota_{\ord}:\ord\to\NO$\label{ord-into-surreals}：
  \begin{equation*}
    \iota_{\ord}(A) \defeq
    \surr{\iota_{\ord}(\ordsl A a) \text{ for all } a:A}{}.
  \end{equation*}
  同样由简单性定理可知，将 $\iota_{\ord}$ 限制在有限序数上与 $\iota_{\nat}$ 一致。
  （但我们提醒读者，与上面的例子不同，除非将其限制到更小的一类序数上，否则 $\iota_{\ord}$ 并非构造性地单的；参见 \cref{ex:ord-into-surreals,ex:hiit-plump}。）
\item 从 Conway 那里再取几个有趣的例子：
  \begin{align*}
    \omega &\defeq \surr{0,1,2,3,\dots}{} \qquad\text{(also an ordinal)}\\
    -\omega &\defeq \surr{}{\dots,-3,-2,-1,0}\\
    1/\omega &\defeq \textstyle\surr{0}{1,\frac12,\frac14,\frac18,\dots}\\
    \omega-1 &\defeq \surr{0,1,2,3,\dots}{\omega}\\
    \omega/2 &\defeq \surr{0,1,2,3,\dots}{\dots,\omega-2,\omega-1,\omega}.
  \end{align*}
\end{itemize}

在识别由不同分割所给出的超现实数时，下面这个简单的观察很有用。

\begin{thm}[Conway 的简单性定理]\label{thm:NO-simplicity}
  \index{simplicity theorem}%
  \index{theorem!Conway's simplicity}%
  假设 $x$ 和 $z$ 是由分割定义的超现实数，并且满足以下条件。
  \begin{itemize}
  \item 对于所有的 $L$ 和 $R$，有 $x^L < z < x^R$。
  \item 对于 $z$ 的每一个左选项 $z^L$，存在一个左选项 $x^{L'}$ 使得 $z^L\le x^{L'}$。
  \item 对于 $z$ 的每一个右选项 $z^R$，存在一个右选项 $x^{R'}$ 使得 $x^{R'}\le z^R$。
  \end{itemize}
  那么 $x=z$。
\end{thm}

\begin{proof}
  应用 $\NO$ 的路径构造子，我们必须证明 $x\le z$ 且 $z\le x$。
  前者需要证明对所有 $L$ 有 $x^L<z$（这正是我们的假设），并且对所有 $R$ 有 $x<z^R$。
  但根据假设，对于任意 $z^R$，存在一个 $x^{R'}$ 满足 $x^{R'}\le z^R$，从而 $x<z^R$，如所需。
  因此 $x\le z$；$z\le x$ 的证明是对称的。
\end{proof}

\index{induction principle!for surreal numbers}
然而，要想对超现实数进行更深入的讨论，我们需要它们的归纳原理。
适用于 $(\NO,\le,<)$ 的相互归纳原理涉及三个类型族：
\begin{align*}
  A &: \NO\to\type\\
  B &: \prd{x,y:\NO}{a:A(x)}{b:A(y)} (x\le y) \to \type\\
  C &: \prd{x,y:\NO}{a:A(x)}{b:A(y)} (x<y) \to \type.
\end{align*}
与 Cauchy 实数的归纳原理类似，将 $B$ 和 $C$ 视为类型 $A(x)$ 与 $A(y)$ 之间的关系族，会更容易理解。
\symlabel{NO-recursion}
因此，我们将 $B(x,y,a,b,\xi)$ 写作 $(x,a) \ble^\xi (y,b)$，将 $C(x,y,a,b,\xi)$ 写作 $(x,a) \blt^\xi (y,b)$。
类似地，我们通常省略 $\xi$，因为它居于一个命题，因而无关紧要；我们也常常省略 $x$ 和 $y$，直接简写为 $a\ble b$ 或 $a\blt b$。
采用这些记号后，归纳原理的前提如下。

\begin{itemize}
\item 对于定义超现实数 $x$ 的任何分割，以及
  \begin{enumerate}
  \item 对每个 $L$，给定元素 $a^L:A(x^L)$，和
  \item 对每个 $R$，给定元素 $a^R:A(x^R)$，使得
  \item 对所有 $L$ 和 $R$ 有 $(x^L,a^L) \blt (x^R,a^R)$
  \end{enumerate}
  存在一个指定的元素 $f_a:A(x)$。
  我们称这样的数据为定义 $x$ 的分割上的一个\define{依值分割（dependent cut）}
  \indexdef{cut!of surreal numbers!dependent}%
  \indexdef{dependent!cut}。%
\item 对于任意 $x,y:\NO$ 及 $a:A(x)$ 和 $b:A(y)$，若 $x\le y$ 且 $y\le x$，并且还有 $(x,a) \ble (y,b)$
  与 $(y,b) \ble (x,a)$，
  则有 $\dpath{A}{\noeq}{a}{b}$。
\item 给定定义两个超现实数 $x$ 和 $y$ 的分割，以及 $x$ 上的依值分割 $a$ 和 $y$ 上的依值分割 $b$，使得对所有 $L$ 有 $x^L<y$ 且 $(x^L,a^L)\blt (y,f_b)$，
  并且对所有 $R$ 有 $x<y^R$ 且 $(x,f_a) \blt (y^R,b^R)$，
  则有 $(x,f_a) \ble (y,f_b)$。
\item $\ble$ 取值于命题。
\item 给定定义两个超现实数 $x$ 和 $y$ 的分割、$x$ 上的依值分割 $a$ 和 $y$ 上的依值分割 $b$，以及一个满足 $x\le y^{L_0}$ 和 $(x,f_a) \ble (y^{L_0},b^{L_0})$ 的 $L_0$，
  则有 $(x,f_a) \blt (y,f_b)$。
\item 给定定义两个超现实数 $x$ 和 $y$ 的分割、$x$ 上的依值分割 $a$ 和 $y$ 上的依值分割 $b$，以及一个满足 $x^{R_0}\le y$ 和 $(x^{R_0},a^{R_0}),\ble (y,f_b)$ 的 ${R_0}$，
  则有 $(x,f_a) \blt (y,f_b)$。
\item $\blt$ 取值于命题。
\end{itemize}

在这些前提下，我们推导出一个函数 $f:\prd{x:\NO} A(x)$，使得
\begin{align}
  f(x) &\;\jdeq\; f_{f[x]} \label{eq:noind1}\\
  (x\le y) &\;\Rightarrow\; (x,f(x)) \ble (y,f(y)) \notag\\
  (x< y) &\;\Rightarrow\; (x,f(x)) \blt (y,f(y)). \notag
\end{align}
在点构造子的计算规则~\eqref{eq:noind1} 中，$x$ 是由一个分割定义的超现实数，而 $f[x]$ 表示通过应用 $f$ 得到的 $x$ 上的依值分割（并利用 $f$ 将 $<$ 映到 $\blt$ 这一事实）。
通常，我们会使用模式匹配记法，其中 $f$ 在分割 $\surr{x^L}{x^R}$ 上的定义可能会用到符号 $f(x^L)$ 和 $f(x^R)$，以及它们构成一个依值分割的假设。

与 Cauchy 实数的情况类似，通过平凡化 $A$、$\ble$ 和 $\blt$ 中的某些部分，我们可以得到一些特例。
将 $\ble$ 和 $\blt$ 取为常值 \unit，我们就得到了 \define{$\NO$-归纳法（$\NO$-induction）}，为简洁起见，我们只针对命题来陈述它：

\begin{itemize}
\item 给定 $P:\NO\to\prop$，若对每个由分割定义的超现实数 $x$，只要对所有 $L$ 和 $R$ 有 $P(x^L)$ 和 $P(x^R)$ 成立，就有 $P(x)$ 成立，则 $P(x)$ 对所有 $x:\NO$ 成立。
\end{itemize}

这可以与 Conway 的下述评论相对照：

\begin{quote}\footnotesize
  一般来说，当我们希望对所有数 $x$ 确立一个命题 $P(x)$ 时，我们将用归纳法来证明它：由所有 $P(x^L)$ 和 $P(x^R)$ 成立来推导 $P(x)$。
  我们认为``所有数都是以这种方式构造的''这一说法为这一做法的合法性提供了依据。
\end{quote}

利用 $\NO$-归纳法，我们可以证明

\begin{thm}[Conway 定理 0]\label{thm:NO-refl-opt}\
  \index{theorem!Conway's 0}%
  \begin{enumerate}
  \item 对任意 $x:\NO$，有 $x\le x$。\label{item:NO-le-refl}
  \item 对任意由分割定义的 $x:\NO$，对所有 $L$ 和 $R$，有 $x^L <x$ 和 $x<x^R$。\label{item:NO-lt-opt}
  \end{enumerate}
\end{thm}

\begin{proof}
  首先注意，若 $x\le x$，则每当 $x$ 作为某个分割 $y$ 的左选项出现时，由 $<$ 的第一个构造子可得 $x<y$；类似地，每当 $x$ 作为某个分割 $y$ 的右选项出现时，由 $<$ 的第二个构造子可得 $y<x$。
  特别地，有 \ref{item:NO-le-refl}$\Rightarrow$\ref{item:NO-lt-opt}。

  我们通过对 $x$ 作 $\NO$-归纳法来证明 \ref{item:NO-le-refl}。
  因此，假设 $x$ 是由一个分割定义的，且对所有 $L$ 和 $R$ 有 $x^L\le x^L$ 和 $x^R \le x^R$。
  但根据我们上面的观察，这些假设意味着对所有 $L$ 和 $R$ 有 $x^L<x$ 和 $x<x^R$，从而由 $\le$ 的构造子可得 $x\le x$。
\end{proof}

\begin{cor}\label{thm:NO-set}
  $\NO$ 是一个 $0$-类型。
%  (As with $V$, it might be confusing to say that it is a ``set''.)
\end{cor}

\begin{proof}
  纯粹关系 $R(x,y)\defeq (x\le y) \land (y\le x)$ 通过 $\NO$ 的道路构造子蕴含相等性，并且由 \cref{thm:NO-refl-opt}\ref{item:NO-le-refl} 可知它包含对角线。
  因此，\cref{thm:h-set-refrel-in-paths-sets} 适用。
\end{proof}

相比之下，Conway 的定理 1（$\le$ 的传递性）用我们的定义来证明则要困难一些；参见 \cref{thm:NO-unstrict-transitive}。

% Of course, we also have:

% \begin{lem}
%   Every surreal number is merely defined by a cut.
% \end{lem}
% \begin{proof}
%   Obvious by $\NO$-induction.
% \end{proof}

我们还需要联合递归原理，即 \define{$(\NO,\le,<)$-递归（$({\NO,\le,<})$-recursion）}。
为方便起见，我们如下陈述。
假设 $A$ 是一个配备了关系 $\mathord\ble:A\to A\to\prop$ 和 $\mathord\blt:A\to A\to\prop$ 的类型。
那么我们可以通过以下步骤来定义 $f:\NO\to A$。

\begin{enumerate}
\item 对于由分割定义的任意 $x$，假设 $f(x^L)$ 和 $f(x^R)$ 已定义且对所有 $L$ 和 $R$ 满足 $f(x^L)\blt f(x^R)$，我们必须定义 $f(x)$。（我们称此为递归的\emph{主条款}。）\label{item:NO-rec-primary}
\item 证明 $\ble$ 是\emph{反对称的}\index{relation!antisymmetric}：若 $a\ble b$ 且 $b\ble a$，则 $a=b$。
\item 对于由分割定义的 $x,y$，若对所有 $L$ 有 $x^L<y$ 且对所有 $R$ 有 $x<y^R$，并且归纳地假设对所有 $L$ 有 $f(x^L)\blt f(y)$，对所有 $R$ 有 $f(x)\blt f(y^R)$，以及对所有 $L$ 和 $R$ 还有 $f(x^L)\blt f(x^R)$ 和 $f(y^L)\blt f(y^R)$，则我们必须证明 $f(x)\ble f(y)$。
\item 对于由分割定义的 $x,y$ 以及某个满足 $x\le y^{L_0}$ 的 $L_0$，并且归纳地假设 $f(x)\ble f(y^{L_0})$，以及对所有 $L$ 和 $R$ 还有 $f(x^L)\blt f(x^R)$ 和 $f(y^L)\blt f(y^R)$，则我们必须证明 $f(x)\blt f(y)$。
\item 对于由分割定义的 $x,y$ 以及某个满足 $x^{R_0}\le y$ 的 $R_0$，并且归纳地假设 $f(x^{R_0})\ble f(y)$，以及对所有 $L$ 和 $R$ 还有 $f(x^L)\blt f(x^R)$ 和 $f(y^L)\blt f(y^R)$，则我们必须证明 $f(x)\blt f(y)$。\label{item:NO-rec-last}
\end{enumerate}

后三个条款可以更简洁地表述为：我们必须证明 $f$（如第一个条款所定义）将 $\le$ 映为 $\ble$，将 $<$ 映为 $\blt$。
我们将这些性质称为\emph{$f$ 保持不等式}。
此外，在证明 $f$ 保持不等式时，我们可以假设所考虑的 $\le$ 或 $<$ 的具体实例是由其某个构造子生成的，且可使用归纳假设：$f$ 保持了出现在该构造子输入中的所有不等式。

如果我们完成了上述~\ref{item:NO-rec-primary}--\ref{item:NO-rec-last}，就得到 $f:\NO\to A$，它在分割上的计算如~\ref{item:NO-rec-primary} 所指定，且保持所有不等式：
%
\begin{narrowmultline*}
  \fall{x,y:\NO}\Big((x\le y) \to (f(x)\ble f(y))\Big) \land
  \narrowbreak
  \Big((x< y) \to (f(x)\blt f(y))\Big).
\end{narrowmultline*}
%
与 Cauchy 实数的 $(\RC,\closesym)$-递归类似，这个递归原理对于定义 $\NO$ 上的函数至关重要，因为我们无法先在``预超现实数''上定义函数，事后再证明它尊重相等的概念。

\begin{eg}
  我们来定义\emph{取负}函数 $\NO\to\NO$。
  我们应用联合递归原理，取 $A\defeq\NO$，$(x\ble y)\defeq (y\le x)$，以及 $(x\blt y)\defeq (y< x)$。
  显然此 $\ble$ 是反对称的。

  对于定义中的主条款，假设 $x$ 由分割定义，且 $-x^L$ 与 $-x^R$ 已定义并满足对所有 $L$ 和 $R$ 有 $-x^L \blt -x^R$。
  按定义，这意味着对所有 $L$ 和 $R$ 有 $-x^R < -x^L$，因此我们可以通过分割 $\surr{-x^R}{-x^L}$ 来定义 $-x$。
  此记号遵循 Conway 的约定，指的是这样一个分割：其左选项以类型 $\RR$ 为指标——$\RR$ 正是 $x$ 的右选项的指标类型；其右选项以类型 $\LL$ 为指标——$\LL$ 正是 $x$ 的左选项的指标类型；相应的族 $\RR\to\NO$ 和 $\LL\to\NO$ 由 $x$ 的对应族与取负复合而成。

  现在我们必须验证 $f$ 保持不等式。
  \begin{itemize}
  \item 对于 $x\le y$，我们可以假设对所有 $L$ 有 $x^L<y$ 且对所有 $R$ 有 $x < y^R$，并证明 $-y\le -x$。
    但由归纳假设可得 $-y < -x^L$ 且 $-y^R < -x$，再由 $-y$、$-x$ 的定义及 $\le$ 的构造子，即得所需结果。
  \item 对于 $x<y$，在第一种情形中，即它由某个 $x\le y^{L_0}$ 生成时，我们可以归纳地假设 $-y^{L_0} \le -x$，此时由 $<$ 的构造子可得 $-y<-x$。
  \item 类似地，若 $x<y$ 由 $x^{R_0}\le y$ 生成，归纳假设为 $-y \le -x^R$，同样可得 $-y<-x$。
  \end{itemize}
\end{eg}

然而，要在此基础上做进一步的讨论，我们需要像在 \cref{thm:RC-sim-characterization} 中对 Cauchy 实数所做的那样，更明确地刻画关系 $\le$ 和 $<$。
同样地，为了使归纳得以进行，我们必须同时证明这些关系的若干本质性质。

\begin{thm}\label{defn:No-codes}
  存在关系 $\mathord\preceq:\NO\to\NO\to\prop$ 和 $\mathord\prec:\NO\to\NO\to\prop$，使得如果 $x$ 和 $y$ 是由分割定义的超现实数，则
  \begin{align*}
    (x\preceq y) &\defeq
    \big(\fall{L} x^L\prec y\big) \land \big(\fall{R} x\prec y^R\big)\\
    (x\prec y) &\defeq
    \big(\exis{L} x\preceq y^L\big) \lor \big(\exis{R} x^R \preceq y\big).
  \end{align*}
  此外，我们有
  \begin{equation}\label{eq:NO-codes-unstrict}
    (x\prec y) \to (x\preceq y)
  \end{equation}
  以及所有合理的传递性，使 $\prec$ 和 $\preceq$ 成为 $\le$ 和 $<$ 上的``双模''\index{bimodule}：
  \begin{equation}\label{eq:NO-codes-transitivity}
    \begin{array}{c@{\hspace{1cm}}c}
      (x \le y) \to (y\preceq z) \to (x\preceq z) &
      (x \preceq y) \to (y\le z) \to (x\preceq z) \\
      (x \le y) \to (y\prec z) \to (x\prec z) &
      (x \preceq y) \to (y< z) \to (x\prec z) \\
      (x < y) \to (y\preceq z) \to (x\prec z) &
      (x \prec y) \to (y\le z) \to (x\prec z).
  \end{array}
  \end{equation}
\end{thm}

\begin{proof}
  我们通过对 $x, y$ 的双重 $(\NO,\le,<)$-归纳来定义 $\preceq$ 和 $\prec$。
  第一次归纳是一个简单的递归，其值域是 $(\NO\to\prop)\times (\NO\to\prop)$ 的子集 $A$，该子集由一对谓词组成，其中一个蕴含另一个，并且满足“右侧传递性”，即 \eqref{eq:NO-codes-unstrict} 和 \eqref{eq:NO-codes-transitivity} 的右列，其中 $(x\preceq \blank)$ 和 $(x\prec \blank)$ 被替换为两个给定的谓词。
  如同在 \cref{defn:RC-approx} 的证明中，我们将这些谓词视为二元关系的一半来书写，记为 $y\mapsto (\hle y)$ 和 $y\mapsto (\hlt y)$，并用 $\hlname$ 表示这一对关系。
  % The precise definition of $A$ is
  % \begin{align*}
  %   A\defeq \bigg\{ \hlname : (\NO\to\prop)\times (\NO\to\prop) \;\bigg|\;\\
  %   \begin{split}
  %     \fall{y,z:\NO}
  %     &\Big( (\hle y) \to (y\le z) \to (\hle z) \Big)\\
  %     \land\; &\Big( (\hle y) \to (y< z) \to (\hlt z) \Big)\\
  %     \land\; &\Big( (\hlt y) \to (y\le z) \to (\hlt z) \Big)\\
  %     \land\; &\Big( (\hlt y) \to (y< z) \to (\hlt z) \Big) \bigg\}
  %   \end{split}
  % \end{align*}
  我们在 $A$ 上装备以下两个关系：
  \begin{align*}
    (\hlname \ble \hlbname) &\defeq
    \fall{y:\NO} \Big( (\hleb y) \to (\hle y) \Big) \land
    \Big( (\hltb y) \to (\hlt y) \Big),\\
    (\hlname \blt \hlbname) &\defeq
    \fall{y:\NO} \Big( (\hleb y) \to (\hlt y) \Big).
    %\land \Big( (\hltb y) \to (\hlt y) \Big)
  \end{align*}
  注意 $\ble$ 是反对称的，因为如果 $\hlname \ble \hlbname$ 且 $\hlbname \ble \hlname$，那么对于所有 $y$ 有 $(\hleb y) \Leftrightarrow (\hle y)$ 且 $(\hltb y) \Leftrightarrow (\hlt y)$，因此由命题的泛等性和函数外延性可得 $\hlname=\hlbname$。
  此外，说一个函数 $\NO\to A$ 保持不等式，精确地说就是：当将其视为 $\NO$ 上的一对二元关系时，它满足“左侧传递性”（\eqref{eq:NO-codes-transitivity} 的左列）。

  现在处理递归的主情况，我们假设给定一个由分割定义的 $x$，以及对于所有 $L$ 和 $R$ 的关系 $(x^L \prec \blank)$, $(x^R \prec \blank)$, $(x^L \preceq \blank)$, 和 $(x^R \preceq \blank)$。其中严格的关系蕴含非严格的关系，满足右侧传递性，并且使得
  \begin{equation}\label{eq:NO-prec-outer-IH}
    \fall{L,R}{y:\NO}\Big( (x^R\preceq y) \to (x^L \prec y) \Big).
    % \land\Big( (x^R \prec y) \to (x^L \prec y) \Big)
  \end{equation}
  现在我们需要对所有 $y$ 定义 $(x\prec y)$ 和 $(x\preceq y)$。
  这里与 \cref{defn:RC-approx} 相反，我们使用一个嵌套的归纳而非嵌套的递归，以便能够归纳地使用关于不等式 $x^L<x$ 和 $x<x^R$ 的左侧传递性。
  定义 $A':\NO\to\type$，其中 $A'(y)$ 是 $\prop\times\prop$ 的子集，由两个命题组成，记作 $\tle y$ 和 $\tlt y$（对应的元素记为 $\tlname:A'(y)$），使得
  \begin{gather}
    (\tlt y) \to (\tle y)\\
    \fall{L} (\tle y)\to (x^L\prec y) \label{eq:NO-prec-IHL}\\
    \fall{R} (x^R \preceq y) \to (\tlt y) \label{eq:NO-prec-IHR}.
  \end{gather}
  使用类似于 $\ble$ 和 $\blt$ 的记号，我们在 $A'$ 上装备两个关系，对于 $\tlname:A'(y)$ 和 $\tlbname:A'(z)$ 定义如下：
  \begin{align*}
    (\tlname \bble \tlbname) &\defeq
    \Big((\tle y) \to (\tleb z)\Big) \land \Big((\tlt y) \to (\tltb z)\Big)\\
    (\tlname \bblt \tlbname) &\defeq
    \Big((\tle y) \to (\tltb z)\Big). % \land \Big(\tlt \to \tltb\Big).
  \end{align*}
  % (These are the type families $B$ and $C$ in the general induction principle.)
  同样，$\bble$ 在适当的意义下显然是反对称的。
  此外，一个能保持不等式的函数 $\prd{y:\NO} A'(y)$ 恰好是这样一对谓词：其中一个蕴含另一个，满足右侧传递性，并且关于不等式 $x^L<x$ 和 $x<x^R$ 满足左侧传递性。
  因此，这个内层归纳将提供我们完成外层递归主情况所需的内容。

  对于内层归纳的主情况，我们还假设给定一个由分割定义的 $y$，以及对于所有 $L$ 和 $R$ 的性质 $(x\prec y^L)$, $(x\prec y^R)$, $(x\preceq y^L)$, 和 $(x\preceq y^R)$。其中严格的关系蕴含非严格的关系，关于 $x^L<x$ 和 $x<x^R$ 满足左侧传递性，关于 $y^L<y^R$ 满足右侧传递性。
  % \begin{equation}
  %   \fall{L,R}\Big((x \preceq y^L) \to (x \prec y^R)\Big) % \land \Big((x \prec y^L) \to (x\prec y^R)\Big).
  %   \label{eq:NO-prec-inner-IH}
  % \end{equation}
  我们现在可以给出定理陈述中指定的定义：
  \begin{align}
    (x\preceq y) &\defeq
    (\fall{L} x^L\prec y) \land (\fall{R} x\prec y^R), \label{eq:NO-preceq-def}\\
    (x\prec y) &\defeq
    (\exis{L} x\preceq y^L) \lor (\exis{R} x^R \preceq y).\label{eq:NO-prec-def}
  \end{align}
  为了让这一定义出 $A'(y)$ 的一个元素，我们必须首先证明 $(x\prec y) \to (x\preceq y)$。
  假设 $x\prec y$ 有两种情况。
  一方面，如果存在 $L_0$ 使得 $x\preceq y^{L_0}$，那么通过关于 $y^{L_0}<y^R$ 的右侧传递性，我们得到对于所有 $R$ 有 $x\prec y^R$。
  此外，通过关于 $x^L<x$ 的左侧传递性，对于任何 $L$ 我们有 $x^L \prec y^{L_0}$，从而由右侧传递性得到 $x^L\prec y$。
  因此，$x\preceq y$。

  另一方面，如果存在 $R_0$ 使得 $x^{R_0}\preceq y$，那么由 \eqref{eq:NO-prec-outer-IH}，我们得到对于所有 $L$ 有 $x^L \prec y$。
  并且通过关于 $x<x^{R_0}$ 和 $y<y^R$ 的左侧与右侧传递性，对于任何 $R$ 我们有 $x\prec y^R$。
  因此，$x\preceq y$。

  我们还需要证明这些定义关于 $x^L<x$ 和 $x<x^R$ 满足左侧传递性。
  但如果 $x\preceq y$，根据定义对于所有 $L$ 有 $x^L\prec y$；而如果 $x^R\preceq y$，同样根据定义有 $x\prec y$。

  因此，\eqref{eq:NO-preceq-def} 和 \eqref{eq:NO-prec-def} 确实定义了 $A'(y)$ 的一个元素。
  我们现在需要验证，这个定义作为进入 $A'$ 的依值函数是保持不等式的，即这些关系满足右侧传递性。
  记住在每种情况下，我们可以归纳地假设它们对于不等式构造器中产生的所有不等式都满足右侧传递性。
  \begin{itemize}
  \item 假设 $x\preceq y$ 且 $y\le z$，后者源于对于所有 $L$ 和 $R$ 有 $y^L<z$ 和 $y<z^R$。
    那么将内层递归的归纳假设应用于 $y<z^R$，得到对于任何 $R$ 有 $x\prec z^R$。
    此外，根据定义 $x\preceq y$ 意味着对于任何 $L$ 有 $x^L \prec y$，因此通过外层递归的归纳假设我们得到 $x^L \prec z$。
    于是，$x\preceq z$。
  \item 假设 $x\preceq y$ 且 $y<z$。
    首先，假设 $y<z$ 源于 $y\le z^{L_0}$。
    那么将内层归纳假设应用于 $y\le z^{L_0}$，得到 $x \preceq z^{L_0}$，因此 $x\prec z$。

    其次，假设 $y<z$ 源于 $y^{R_0}\le z$。
    那么根据定义，$x\preceq y$ 意味着 $x\prec y^{R_0}$，然后将内层归纳假设应用于 $y^{R_0}\le z$，得到 $x\prec z$。
  \item 假设 $x\prec y$ 且 $y\le z$，后者源于对于所有 $L$ 和 $R$ 有 $y^L<z$ 和 $y<z^R$。
    根据定义，$x\prec y$ 意味着仅仅存在 $R_0$ 使得 $x^{R_0}\preceq y$ 或存在 $L_0$ 使得 $x\preceq y^{L_0}$。
    如果 $x^{R_0}\preceq y$，那么外层归纳假设给出 $x^{R_0}\preceq z$，因此 $x\prec z$。
    如果 $x\preceq y^{L_0}$，那么将内层归纳假设应用于 $y^{L_0}<z$（这由 $y\le z$ 的构造子保证）得到 $x\prec z$。
  % \item Suppose $x\prec y$ and $y<z$.
  %   First, suppose $y<z$ arises from $y\le z^{L_0}$.
  %   Then the inner inductive hypothesis for $y\le z^{L_0}$ yields $x\prec z^{L_0}$, hence $x\preceq z^{L_0}$; thus $x\prec z$.

  %   Second, suppose $y<z$ arises from $y^{R_0}\le z$.
  %   Then by definition, $x\prec y$ implies there merely exists $R_1$ with $x^{R_1}\preceq y$ or $L_1$ with $x\preceq y^{L_1}$.
  %   If $x^{R_1}\preceq y$, then the outer inductive hypothesis implies $x^{R_1}\prec z$, hence $x^{R_1}\preceq z$, and thus $x\prec z$.
  %   And if $x\preceq y^{L_1}$, then the inner inductive hypothesis applied to $y^{L_1}<y^{R_0}$ (which comes from $y$ being defined as a cut) and $y^{R_0}\le z$ yields $x\prec z$.
  \end{itemize}
  这就完成了内层归纳。
  因此，对于任何由分割定义的 $x$，我们由 \eqref{eq:NO-preceq-def} 和 \eqref{eq:NO-prec-def} 定义了 $(x\prec \blank)$ 和 $(x\preceq \blank)$，并且它们满足右侧传递性。

  为了完成外层递归，我们需要验证这些定义满足左侧传递性。
  在对 z 进行 $\NO$-归纳之后，我们最终会得到与上面为右侧传递性描述的三种情况本质上相同的三种情况。
  因此，我们省略它们。
\end{proof}

\begin{thm}\label{thm:NO-encode-decode}
  对于任意 $x,y:\NO$，我们有 $(x<y)=(x\prec y)$ 与 $(x\le y)=(x\preceq y)$。
\end{thm}

\begin{proof}
  从左到右，我们使用 $(\NO,\le,<)$-归纳，其中 $A(x)\defeq\unit$，并用 $\preceq$ 和 $\prec$ 来提供关系 $\ble$ 和 $\blt$。
  在所有构造子情形中，$x$ 和 $y$ 都由分割定义，因此 $\preceq$ 和 $\prec$ 的定义可以直接计算，归纳前提适用。

  从右到左，我们使用 $\NO$-归纳，假设 $x$ 和 $y$ 由分割定义。
  但此时 $\preceq$ 和 $\prec$ 的定义以及归纳前提，恰好提供了 $\le$ 和 $<$ 的相关构造子所需的数据。
  % From right to left, we first prove by $\NO$-induction on $x$ that for any $y:\NO$ we have $(x\prec y) \to (x<y)$ and $(x\preceq y) \to (x\le y)$.
  % Thus, we assume this to be true for all $x^L$ and $x^R$ in a cut, and show it for the resulting $x:\NO$.
  % Next, we prove by $\NO$-induction on $y$ that $(x\prec y) \to (x<y)$ and $(x\preceq y) \to (x\le y)$, hence we assume it to be true for all $y^L$ and $y^R$ in a cut, and show it for the resulting $y:\NO$.
  % Now since $x$ and $y$ are both defined by cuts, $x\preceq y$ means that $x^L\prec y$ and $x\prec y^R$ for all $L$ and $R$.
  % By the inductive hypotheses, this gives $x^L<y$ and $x<y^R$, hence $x\le y$ by the constructor of $\le$.
  % Similarly, $x\prec y$ yields merely an $R_0$ with $x^{R_0}\preceq y$ or an $L_0$ with $x\preceq y^{L_0}$.
  % Hence merely $x^{R_0}\le y$ or $x\le y^{L_0}$ by the inductive hypothesis, so $x<y$ by a constructor.
\end{proof}

\begin{cor}\label{thm:NO-unstrict-transitive}
  $\NO$ 上的关系 $\le$ 和 $<$ 满足
  \[ \fall{x,y:\NO} (x<y) \to (x\le y) \]
  且是传递的：
  \index{transitivity!of . for surreals@of $<$ for surreals}
  \index{transitivity!of . for surreals@of $\leq$ for surreals}
  \begin{gather*}
    (x\le y) \to (y\le z) \to (x\le z)\\
    (x\le y) \to (y< z) \to (x< z)\\
    (x< y) \to (y\le z) \to (x< z).
  \end{gather*}
\end{cor}

与 Cauchy 实数类似，在定义 $\NO$ 上所有运算时，联合 $(\NO,\le,<)$-递归原理仍然是必不可少的。

\begin{eg}\label{eg:surreal-addition}
\index{addition!of surreal numbers}%
我们通过先对第一个变元递归、再对第二个变元归纳，来定义 $\mathord+:\NO\to\NO\to\NO$。
对于外层递归，我们取余域为 $\NO\to\NO$ 的一个子集，该子集由那些对所有 $x,y$ 满足 $(x<y) \to (g(x)<g(y))$ 和 $(x\le y) \to (g(x)\le g(y))$ 的函数 $g$ 组成。
对于这样的 $g,h$，我们定义 $(g\ble h)\defeq \fall{x:\NO} g(x)\le h(x)$ 与 $(g\blt h)\defeq \fall{x:\NO} g(x)< h(x)$。
显然 $\ble$ 是反对称的。

对于递归的主条款，我们假设 $x$ 由分割定义，函数 $(x^L+\blank)$ 和 $(x^R+\blank)$ 已定义、保序，且满足 $x^L+y<x^R+y$，然后定义 $(x+\blank)$。
如同在 \cref{defn:No-codes} 中那样，我们不进行内层递归，而是向类型族 $A:\NO\to\type$ 进行内层归纳，其中 $A(y)$ 是 $\NO$ 中满足以下条件的 $z$ 构成的子集：对所有 $x^L$ 有 $x^L + y < z$，且对所有 $x^R$ 有 $x^R + y > z$。
我们为 $A$ 赋予从 $\NO$ 诱导的关系 $\le$ 和 $<$，因此反对称性是显然的。
对于内层归纳的主条款，我们进一步假设 $y$ 也由分割定义，且每个 $x+y^L$ 和 $x+y^R$ 已定义并满足 $x^L+y^L < x+y^L$、$x^L+y^R < x+y^R$、$x+y^L < x^R + y^L$ 以及 $x+y^R < x^R+y^R$（这些来自对 $A(y)$ 元素施加的额外条件），同时也满足 $x+y^L < x+y^R$（因为 $A(y)$ 中的元素 $x+y^L$ 和 $x+y^R$ 构成了一个依值分割）。
现在我们给出 Conway 的定义：
\[ x+y \defeq \surr{x^L+y, x+y^L}{x^R+y,x+y^R}. \]
换句话说，$x+y$ 的左选项是所有形如 $x^L+y$（对某个左选项 $x^L$）或 $x+y^L$（对某个左选项 $y^L$）的数。
我们必须证明这些左选项中的每一个都小于每一个右选项：
\begin{itemize}
\item $x^L+y < x^R+y$，由外层归纳前提。
\item $x^L+y < x^L + y^R < x + y^R$，第一个因为 $(x^L+\blank)$ 保序，第二个因为 $x+y^R : A(y^R)$。
\item $x+y^L < x^R+ y^L < x^R + y$，第一个因为 $x+y^L : A(y^L)$，第二个因为 $(x^R+\blank)$ 保序。
\item $x+y^L < x+y^R$，由内层归纳前提（具体来说，由我们有一个依值分割这一事实）。
\end{itemize}
我们还必须证明如此定义的 $x+y$ 属于 $A(y)$，即 $x^L + y < x+y$ 且 $x+y < x^R + y$；但这由 \cref{thm:NO-refl-opt}\ref{item:NO-lt-opt} 即得。

接下来我们必须验证 $(x+\blank)$ 的定义保序：
\begin{itemize}
\item 若 $y\le z$ 是由已知对所有 $L$ 和 $R$ 有 $y^L<z$ 且 $y<z^R$ 而得到的，则内层归纳前提给出 $x+y^L<x+z$ 和 $x+y < x+z^R$，而外层归纳前提给出 $x^L+y \le x^L+z$ 和 $x^R+ y \le x^R+z$。
  此外，由于 $x^R+y$ 依定义是 $x+y$ 的一个右选项，我们有 $x+y < x^R+y$。
  类似地，$x^L+z$ 是 $x+z$ 的一个左选项，因此 $x^L+z < x+z$。
  于是，利用传递性，我们得到 $x^L+y < x+z$ 和 $x+y < x^R+z$；因此由 $\le$ 的构造子可得 $x+y \le x+z$。
\item 若 $y<z$ 来自某个 $L_0$ 使得 $y\le z^{L_0}$，则由归纳有 $x+y \le x+z^{L_0}$，因此 $x+y<x+z$，因为 $x+z^{L_0}$ 是 $x+z$ 的一个右选项。
\item 类似地，若 $y<z$ 来自 $y^{R_0}\le z$，则 $x+y<x+z$，因为 $x+y^{R_0}\le x+z$。
\end{itemize}
这就完成了内层归纳。
对于外层递归，我们必须验证 $+$ 在左侧也保序。
进行一次 $\NO$-归纳后，论证过程完全相同。
\end{eg}

\index{acceptance|(}%
\index{mathematics!formalized}%
在~\cite{conway:onag} 第零部分的附录中，Conway 讨论了如何在 ZFC 集合论中形式化超现实数：沿序数迭代，并对每个等价类取秩最低的代表元所构成的集合；或者用``符号展开''来表示数。
他随后指出

\begin{quote}\footnotesize
  这些构造出奇复杂的性质，更多地揭示了 ZF 框架内形式化的本质，而非我们的数系本身……
\end{quote}

接着，他主张建立一种关于``容许的构造种类''的一般理论，其中应当包括：

\begin{enumerate}\footnotesize
\item 对象可以以任何具有合理构造性的方式从先前的对象中创建。\label{item:conway1}
\item 所创建对象之间的相等关系可以是任意的等价关系。\label{item:conway2}
\end{enumerate}

\noindent
条件~\ref{item:conway1} 可以很自然地理解为对\emph{归纳定义（inductive definition）}一般原则的证成，例如在 \cref{sec:strictly-positive,sec:generalizations} 中阐述的那些。
特别地，构造子的严格正性（strict positivity）条件，可以看作是``具有合理构造性''这一要求的形式化表述。
条件~\ref{item:conway2} 则提示我们，应将此扩展到各种\emph{高阶}归纳定义（higher inductive definitions），在其中我们可以引入道路构造子，以任意合理的方式规定对象之间的相等。
例如，Conway 在下一段中写道：

\begin{quote}\footnotesize
  ……我们也可以，比方说，自由地创造一个新对象 $(x,y)$，并将其称为 $x$ 和 $y$ 的有序对。
  我们还可以创造一个与 $(x,y)$ 不同但与之共存的有序对 $[x,y]$……
  而如果我们想把 $(x,y)$ 做成无序对，我们可以通过等价关系来定义相等，即：当且仅当 $x=z,y=t$ \emph{或} $x=t,y=z$ 时，有 $(x,y)=(z,t)$。
\end{quote}

通过特定形式的构造子来引入具有新名称的新对象——这种自由，恰恰是我们在归纳定义理论中所拥有的。
正如我们在 \cref{sec:appetizer-univalence} 中处理自然数的两个副本 $\nat$ 和 $\nat'$ 时一样，如果我们写下与 Cartesian 积类型 $A\times B$ 完全相同的定义，我们会得到一个不同的积类型 $A\times' B$，其典范元素可自由地写作 $[x,y]$。
并且，我们可以通过添加适当的道路构造子，使其成为无序对的类型。% (and perhaps 0-truncating).

诚然，Conway 的重点并非专门抱怨 ZF，而是一并反对所有的基础理论：

\begin{quote}\footnotesize
  ……本提议并非要用某个特定理论来替代 ZF……
  相反，提议的是：我们赋予自己自由，去创造任意此类数学理论，但要证明一个元定理，它一劳永逸地确保任何此类理论都可以用任何标准的基础理论来形式化。
\end{quote}

或许有人会回应说，事实上，泛等基础并非 Conway 心目中的那些``标准基础理论''之一，而是这样一种\emph{元理论}：在其中我们可以表达创造新理论的能力，也可以据此证明 Conway 的元定理。
例如，超现实数就是 Conway 心目中的``数学理论''之一，而我们已看到，它们可以在泛等基础内部构造出来并证明其合理性。
类似地，Conway 早前曾指出：

\begin{quote}\footnotesize
  ……集合论就是这样一种理论，集合通过对应于通常公理的构造过程从更早的集合构造出来，而相等关系即具有相同成员。
\end{quote}

这一描述与 \cref{sec:cumulative-hierarchy} 中对集合论累积层次的高阶归纳构造紧密契合。
Conway 的元定理便对应于我们多次提及的那个事实：我们可以在 ZFC 内部构造一个泛等基础的模型（这超出了本书的范围）。

然而，泛等基础本身是如此丰富和强大，若仅仅将它降格为一种用于构造类集合理论的元理论，那将是愚蠢的。
我们看到，即使在集合（0-类型）的层面，泛等基础中的高阶归纳类型也能通过万有性质直接产生对象的构造（\cref{sec:free-algebras}），例如 Cauchy 完备化的构造性理论（\cref{sec:cauchy-reals}）。
但最重要的是，在基础系统中直接为同伦论和范畴论建模的潜力（\cref{cha:homotopy,cha:category-theory}），赋予了泛等基础任何集合论基础都无法比拟的优势。
\index{acceptance|)}%

\index{surreal numbers|)}%

\sectionNotes

在Dedekind 实数上定义代数运算，尤其是乘法，既有些棘手又相当繁琐。
有好几种方法可以建立算术运算：每种都有其优势，但它们似乎都需要一些技术性工作。
例如，Richman~\cite{Richman:reals} 先在正分割上定义乘法，然后再代数地延拓到所有Dedekind 分割；
而 Conway~\cite{conway:onag} 则观察到，超现实数的乘法定义同样可以很好地适用于Dedekind 实数。

我们对Dedekind 实数的处理借鉴了~\cite{BauerTaylor09} 中的许多想法，该文献在 Abstract Stone Duality 的背景下构造了Dedekind 实数。
\index{Abstract Stone Duality}%
这是一种（受限的）简单类型 $\lambda$-演算，带有一个特定的对象 $\Sigma$，该对象分类开集，并由对偶性也分类闭集。
在~\cite{BauerTaylor09} 中也可以找到算术运算基本性质的详细证明。

$\RC$ 是最小的 Cauchy 完备的 Archimedean 的有序域，这一事实（已在 \cref{RC-initial-Cauchy-complete} 中证明）表明，我们的 Cauchy 实数很可能与 Escard{\'o}--Simpson 实数~\cite{EscardoSimpson:01} 相一致。
\index{real numbers!Escardo--Simpson@Escard\'o--Simpson}%
验证这是否确实如此\index{open!problem}，将是一个有趣的问题。
Escard{\'o}--Simpson 实数的概念，更准确地说，是相应的闭区间概念，之所以有趣，是因为它可以在任何具有有限积的范畴中陈述。

在增加了``正则扩张公理''的构造性集合论中，人们也可以尝试通过超限迭代、对 Cauchy 序列的极限取闭包来定义 Cauchy 完备化。
验证这种构造是否与我们的构造一致，同样是一个有趣的问题。

构造性数学中一个广为流传的结论是：Cauchy 实数与Dedekind 实数的一致需要依赖选择公理。
但较少人知道的是，可数选择公理就足够了。
回想一下，\define{依赖选择公理}
\indexdef{axiom!of choice!dependent}%
\index{axiom!of choice!countable}%
\index{total!relation}%
说的是：对于 $A$ 上的一个全关系 $R$（这里全关系是指 $\fall{x : A} \exis{y : A} R(x,y)$），以及任意 $a : A$，都仅存在 $f : \N \to A$ 使得 $f(0) = a$，并且对于所有 $n : \N$ 都有 $R(f(n), f(n+1))$。
我们的 \cref{when-reals-coincide} 使用了典型的技巧，将一个依赖选择公理的应用转化为使用可数选择公理。
具体来说，我们先用一次可数选择公理，提前做好所有可能出现的那些选择，然后利用这个选择函数来避免做出依赖性的选择。

各种紧性概念在构造性背景下的复杂关系在 \cite{bridges2002compactness} 中有所讨论。Palmgren~\cite{Palmgren:FT} 则对逐点分析和无点拓扑进行了很好的比较。

超现实数由~\cite{conway:onag} 定义，使用了一种归纳定义，但并未在任何基础系统下明确证明其合理性。
正因如此，一些后来的作者倾向于使用符号展开式或其他更明确的表示法，以便更明显地编码到集合论中。
在类型论中表示它们的想法最初由 Hancock 考虑，而 Setzer 和 Forsberg~\cite{forsbergfinite} 指出，超现实数及其不等式关系 $<$ 和 $\le$ 自然地构成了一个归纳-归纳定义。
这里介绍的\emph{高阶}归纳-归纳版本，内置了超现实数的正确相等概念，是新的。

\sectionExercises

\begin{ex}\label{ex:alt-dedekind-reals}
  通过先定义平方再利用 \cref{mult-from-square}，给出Dedekind 实数的另一种定义。验证由此得到的是一个交换环。
\end{ex}

\begin{ex} \label{ex:RD-extended-reals}
  %
  假设我们去掉 \cref{defn:dedekind-reals} 中的有界性条件 \ref{defn:dedekind-reals-inhabited}。
  那么我们就得到了 \define{扩充实数}
  \indexdef{real numbers!extended}%
  \indexdef{extended real numbers}%
  它们包含 $-\infty \defeq (\emptyt, \Q)$ 和 $\infty \defeq (\Q, \emptyt)$。哪些分割上的算术运算定义对扩充实数仍有意义？我们会得到什么样的代数结构？
\end{ex}

\begin{ex} \label{ex:RD-lower-cuts}
  %
  通过考虑单侧分割，我们分别得到\define{下Dedekind 实数（lower Dedekind reals）}和\define{上Dedekind 实数（upper Dedekind reals）}。
  \indexdef{real numbers!Dedekind!upper}%
  \indexdef{real numbers!Dedekind!lower}%
  \indexdef{lower Dedekind reals}%
  \indexdef{upper Dedekind reals}%
  \index{cut!Dedekind}%
  例如，一个下实数可由一个谓词 $L : \Q \to \Omega$ 给出，它满足以下条件：
  %
  \begin{enumerate}
  \item \emph{有实例的：} $\exis{q : \Q} L(q)$；
  \item \emph{舍入的：} $L(q) = \exis{r : \Q} q < r \land L(r)$。
    \index{rounded!Dedekind cut}
  \end{enumerate}
  %
  （我们也可以要求 $\exis{r : \Q} \lnot L(r)$，以排除分割 $\infty \defeq
  \Q$。）在下实数上可以定义哪些算术运算？特别是，加法逆元的情况如何？
\end{ex}

\begin{ex} \label{ex:RD-interval-arithmetic}
  %
  \index{interval!arithmetic}%
  假设我们去掉 \cref{defn:dedekind-reals} 中的位于性条件。
  那么我们就得到了\define{区间域（interval domain）}
  \indexdef{interval!domain}%
  $\mathbb{I}$，因为此时分割允许存在``间隙''，而这些间隙就是区间。通过
  %
  \begin{narrowmultline*}
    ((L, U) \sqsubseteq (L', U'))
    \defeq \narrowbreak
    (\fall{q : \Q} L(q) \Rightarrow L'(q)) \land
    (\fall{q : \Q} U(q) \Rightarrow U'(q)).
  \end{narrowmultline*}
  %
  定义 $\mathbb{I}$ 上的偏序 $\sqsubseteq$。
  $\mathbb{I}$ 中的极大元有哪些？定义``端点''运算，它将区间域中的一个元素映为其下端点与上端点。这些端点是实数、下实数还是上实数（参见
  \cref{ex:RD-lower-cuts}）？分割上的哪些算术运算定义对区间域仍然有意义？
\end{ex}

\begin{ex} \label{ex:RD-lt-vs-le}
  证明对于所有 $x, y : \RD$，
  %
  \begin{equation*}
    \lnot (x < y) \Rightarrow y \leq x
  \end{equation*}
  %
  以及
  %
  \begin{equation*}
    \eqv{(x \leq y)}{\Parens{\prd{\epsilon : \Qp} x < y + \epsilon}}.
  \end{equation*}
  %
  成立。
  $\lnot (x \leq y)$ 是否蕴含 $y < x$？
\end{ex}

\begin{ex} \label{ex:reals-non-constant-into-Z}
  \mbox{}
  %
  \begin{enumerate}
  \item
    假设排中律，构造一个从 $\RD$ 到 $\Z$ 的非常值映射。
  \item
    假设 $f : \RD \to \Z$ 是一个映射，满足 $f(0) = 0$ 且对所有 $x >
    0$ 有 $f(x) \neq 0$。试由此推导出有限全知原理~\eqref{eq:lpo}。
\index{limited principle of omniscience}%
  \end{enumerate}
\end{ex}

\begin{ex} \label{ex:traditional-archimedean}
  \index{ordered field!archimedean}%
  证明在有序域 $F$ 中，$\Q$ 的稠密性与传统的 Archimedean 公理是等价的：
  %
  \begin{equation*}
    (\fall{x, y : F} x < y \Rightarrow \exis{q : \Q} x < q < y)
    \Leftrightarrow
    (\fall{x : F} \exis{k : \Z} x < k).
  \end{equation*}
\end{ex}

\begin{ex} \label{RC-Lipschitz-on-interval} 假设 $a, b : \Q$ 且 $f : \setof{ q : \Q |
    a \leq q \leq b } \to \RC$ 是 Lipschitz 的，常数为~$L$。证明存在 $f$ 的唯一延拓 $\bar{f} : [a,b] \to \RC$，且该延拓是 Lipschitz 的，常数为~$L$。提示：不必为闭区间重做 \cref{RC-extend-Q-Lipschitz}，可以注意到存在一个收缩 $r : \RC \to [-n,n]$，然后将 \cref{RC-extend-Q-Lipschitz} 应用于 $f \circ r$ 即可。
\end{ex}

\begin{ex} \label{ex:metric-completion}
  \index{completion!of a metric space}%
  将 $\RC$ 的构造推广，以构造任意度量空间的 Cauchy 完备化。首先，思考哪一种实数概念作为度量空间距离函数\index{distance}的陪域最为自然。这有影响吗？接下来，详细完成两种构造：
  %
  \begin{enumerate}
  \item 遵循 Cauchy 实数的构造，将度量空间的完备化定义为一个在柯西序列极限下封闭的归纳-归纳类型。\index{Cauchy!sequence}
  \item 使用由 Lawvere（劳维尔）~\cite{lawvere:metric-spaces}\index{Lawvere} 和 Richman（里奇曼）~\cite{Richman00thefundamental} 提出的构造，其中度量空间 $(M, d)$ 的完备化由\define{位置（locations）}的类型给出。
    \indexdef{location}%
    一个位置是一个函数 $f : M \to \R$，使得
    %
    \begin{enumerate}
    \item $f(x) \geq |f(y) - d(x,y)|$ 对所有 $x, y : M$ 成立；
    \item $\inf_{x \in M} f(x) = 0$，即 $\fall{\epsilon : \Qp} \exis{x : M} |f(x)| < \epsilon$ 与 $\fall{x : M} f(x) \geq 0$。
    \end{enumerate}
    %
    这里的想法是，$f$ 看起来像是在测量到某个点的距离。
  \end{enumerate}
  %
  \index{universal!property!of metric completion}%
  最后，证明度量空间完备化的如下泛性质：从一个度量空间到 Cauchy 完备度量空间的局部一致连续映射，可以唯一地延拓到该完备化上。（我们称一个映射是\define{局部一致连续的（locally uniformly continuous）}
  \indexdef{function!locally uniformly continuous}%
  \indexdef{locally uniformly continuous map}%
  如果它在每个开球上都是一致连续的。）
\end{ex}

\index{metric space|)}%

\begin{ex} \label{ex:reals-apart-neq-MP}
  \define{Markov 原理（Markov's principle）}
  \indexdef{axiom!Markov's principle}%
  \indexdef{Markov's principle}%
  说的是：对于所有 $f : \nat \to \bool$，
  %
  \begin{equation*}
    (\lnot \lnot \exis{n : \nat} f(n) = \btrue)
    \Rightarrow
    \exis{n : \nat} f(n) = \btrue.
  \end{equation*}
  %
  这是双重否定律~\eqref{eq:ldn} 的一个特例。证明
  $\fall{x, y: \RD} x \neq y \Rightarrow x \apart y$ 蕴含 Markov 原理。反之是否也成立？
\end{ex}

\begin{ex} \label{ex:reals-apart-zero-divisors}
  \index{apartness}%
  验证实数满足下述“无零因子”性质：
  $x y \apart 0 \Leftrightarrow x \apart 0 \land y \apart 0$。
\end{ex}

\begin{ex} \label{ex:finite-cover-lebesgue-number}
  %
  设 $(q_1, r_1), \ldots, (q_n, r_n)$ 逐点覆盖 $(a, b)$。则存在 $\epsilon : \Qp$，使得只要 $a < x < y < b$ 且 $|x - y| < \epsilon$，就仅仅存在 $i$ 使得 $q_i < x < r_i$ 和 $q_i < y < r_i$。这样的 $\epsilon$ 称为该覆盖的一个 \define{Lebesgue 数}
  \indexdef{Lebesgue number}。%
\end{ex}

\begin{ex} \label{ex:mean-value-theorem}
  %
  证明下述中值定理的近似版本：
  %
  \begin{quote}
    \emph{
      若 $f : [0,1] \to \R$ 是一致连续的，且 $f(0) < 0 < f(1)$，则对每个 $\epsilon : \Qp$，都仅仅存在 $x : [0,1]$ 使得 $|f(x)| < \epsilon$。
    }
  \end{quote}
  %
  提示：不要尝试使用二分法，因为它会引入选择公理。而是用一个分段线性映射来逼近 $f$。如何构造一个分段线性映射？
\end{ex}

\begin{ex}\label{ex:knuth-surreal-check}
  检查~\cite{knuth74:_surreal_number} 中的所有内容是否都能使用 \cref{sec:surreals} 中定义的高阶归纳-归纳超现实数来完成。
\end{ex}

\begin{ex}\label{ex:reals-into-surreals}
  回忆在第~\pageref{reals-into-surreals} 页定义的函数 $\iota_{\RD}:\RD\to\NO$。
  \begin{enumerate}
  \item 证明 $\iota_{\RD}$ 是单的。
  \item $\iota_{\RD}$ 可以显然地扩展到扩充实数（\cref{ex:RD-extended-reals}）和区间域（\cref{ex:RD-interval-arithmetic}）。它们是单的吗？
  \end{enumerate}
\end{ex}

\begin{ex}\label{ex:ord-into-surreals}
  证明在第~\pageref{ord-into-surreals} 页定义的函数 $\iota_{\ord}:\ord\to\NO$ 是单的当且仅当 \LEM{} 成立。
\end{ex}

\begin{ex}\label{ex:hiit-plump}
  模仿 $\NO$ 的定义（但仅使用左选项），定义一个配备二元关系 $\le$ 和 $<$ 的类型 $\mathsf{POrd}$。
  \begin{enumerate}
  \item 构造一个映射 $j:\mathsf{POrd} \to \NO$，并证明它是一个嵌入。
  \item 证明在关系 $<$ 下，$\mathsf{POrd}$ 是一个序数（在更高一层的宇宙中，如 $\ord$）。
  \item 假设命题大小重整，证明 $\mathsf{POrd}$ 等价于 \cref{ex:plump-ordinals} 中 \ord 的子集 \[\setof{A:\ord | \mathsf{isPlump}(A)}\]由此推断，$\iota_{\ord}:\ord\to\NO$ 限制在丰满序数上时是单的。
  \end{enumerate}
  在没有命题大小重整的情况下，我们依然可以称 $\mathsf{POrd}$ 的元素（或它们在 $\NO$ 中的像）为 \define{丰满序数}。\index{ordinal!plump}\index{plump!ordinal}
\end{ex}

\begin{ex}\label{ex:pseudo-ordinals}
  如果一个超现实数等于一个没有右选项的分割 $\surr{x^L}{}$（但其左选项本身可以包含右选项），则称之为 \define{伪序数}\index{pseudo-ordinal}\index{ordinal!pseudo-}。证明“每个伪序数都是丰满序数”这一陈述等价于 \LEM{}。
\end{ex}

\begin{ex}\label{ex:double-No-recursion}
  注意到 \cref{defn:No-codes} 和 \cref{eg:surreal-addition} 都使用了相似的模式来定义函数 $\NO \to \NO \to B$：一个外层的 $\NO$-递归，其陪域是保序函数 $\NO\to B$ 的集合；随后是一个内层的 $\NO$-归纳，进入一个类型族 $A:\NO\to\type$，其中 $A(y)$ 是 $B$ 的子集，用以确保不等式 $x^L<x$ 和 $x<x^R$ 也得到保持。请表述并证明一个“双重 $\NO$-递归”的一般原理，用以概括这些证明。
\end{ex}

\index{real numbers|)}%

%%% Local Variables:
%%% mode: latex
%%% TeX-master: "hott-online"
%%% End:

%%%% Appendix

\cleartooddpage[\thispagestyle{empty}] % Needed for correct TOC
\phantomsection % Needed for correct TOC also
\part*{附录}

% We use magic to get the appendix look like Bibliography and Index

\appendix

\renewcommand{\chaptermark}[1]{\markboth{\textsc{Appendix. \thechapter. #1}}{}}
\renewcommand{\sectionmark}[1]{\markright{\textsc{\thesection\ #1}}}

% !TeX root = hott-online.tex

\titleformat{\chapter}[display]{\fontsize{23}{25}\fontseries{m}\fontshape{it}\selectfont}{\chaptertitlename}{20pt}{\fontsize{35}{35}\fontseries{b}\fontshape{n}\selectfont}
\chapter{类型论的形式化表述}
\label{cha:rules}

\index{formal!type theory|(}%
\index{type theory!formal|(}%
\index{rules of type theory|(}%

正如我们可以在集合论中发展数学而无需显式用到 Zermelo--Fraenkel 集合论的公理一样，
在本书中，我们也是在泛等基础中发展数学，而并未显式地参照同伦类型论的形式系统。
尽管如此，将同伦类型论作为一个形式系统加以精确描述仍然是重要的，例如，这使我们能够：
%

\begin{itemize}
\item 陈述并证明其元理论性质，包括逻辑一致性；
\item 构造其模型，例如在单纯集合、模型范畴、高阶意象等模型中；以及
\item 在诸如 \Coq 或 \Agda 之类的证明助理中将其实现。
  \index{proof!assistant}
\end{itemize}

%
甚至同伦类型论的逻辑一致性\index{consistency}——即在空上下文中不存在项 $a:\emptyt$——也并非显而易见：如果我们错误地选择了一个使得 $\eqv{\emptyt}{\unit}$ 的等价定义，
那么泛等公理将推出 $\emptyt$ 也有元素，因为 $\unit$ 确实有元素。
例如，我们作为高阶归纳类型定义的 $\Sn^1$ 产生了一个行为如同通常圆周的类型，这一点同样并非显而易见。

在讨论这些问题之前，我们必须先明确类型论的两个方面。
回顾引言，类型论包含一组规则，用以规定判断 $a:A$ 与 $a\jdeq a':A$ 何时成立——例如，积类型由这样一条规则刻画：只要 $a:A$ 且 $b:B$，就有 $(a,b):A\times B$。
为使上述说法精确化，我们必须首先精确定义\emph{项的语法}——即这些判断所关联的对象 $a,a',A,\dots$；然后，精确定义\emph{判断及其推导规则}——即如何从已有判断推导出新判断。

在本附录中，我们将给出 Martin-L\"{o}f 类型论及构成同伦类型论的各项扩展的两种表述。
第一种表述（\cref{sec:syntax-informally}）将项的语法和判断形式描述为无类型 $\lambda$-演算的一个扩展，而推导规则则仅作非形式化的描述。
第二种表述（\cref{sec:syntax-more-formally}）则以自然演绎的风格，按类型论文献中的通行做法，归纳地定义项、判断以及推导规则。

\section*{预备知识}
\label{sec:formal-prelim}

在 \cref{cha:typetheory} 中，我们给出了类型论的两种基本\define{判断}。
\index{judgment}
第一种，$a:A$，断言项 $a$ 具有类型 $A$。
第二种，$a\jdeq b:A$，则陈述两个项 $a$ 与 $b$ 在类型 $A$ 上是\define{判断相等}的。%
\index{equality!judgmental}
\index{judgmental equality}
这些判断由一组在 \cref{sec:syntax-more-formally} 中描述的推导规则归纳地定义。

构造类型 $A$ 的一个元素 $a$，就是推导出判断 $a:A$；在本书中，我们以非形式化的论证来描述 $a$ 的构造方式，
但在形式层面，必须给定一个精确的项 $a$ 以及 $a:A$ 的一个完整推导。

不过，本书正文与本附录在类型论表述上的主要区别在于：此处的判断明确地在一个环境\define{上下文}中表述，
\index{context}
该上下文即形如
\[
  x_1:A_1, x_2:A_2,\dots,x_n:A_n.
\]
的假设列表。
上下文中的一个元素 $x_i : A_i$ 表达了一个假设：\index{variable}变量 $x_i$ 具有类型 $A_i$。
出现在上下文中的变量 $x_1, \ldots, x_n$ 必须互不相同。
我们用字母 $\Gamma$ 和 $\Delta$ 来缩写上下文。

在上下文 $\Gamma$ 下的判断 $a:A$ 写作
\[ \oftp\Gamma aA \]
这意味着在 $\Gamma$ 中所列假设下，有 $a:A$。
当假设列表为空时，我们简写为
\[ \oftp{}aA \]
或
\[ \oftp\emptyctx aA \]
其中 $\emptyctx$ 表示空上下文。对相等判断亦然：
\[
  \jdeqtp\Gamma{a}{b}{A}
\]

然而，这类判断仅在上下文为\define{良构}时才有意义，
\index{context!well-formed}%
这一概念由我们的第三个、也是最后一个判断\[
  \wfctx{(x_1:A_1, x_2:A_2,\dots,x_n:A_n)}
\]所刻画，该判断表达的是每个 $A_i$ 在上下文 $x_1:A_1,
x_2:A_2,\dots,x_{i-1}:A_{i-1}$ 中都是一个类型。
因此特别地，如果 $\oftp\Gamma aA$ 且 $\wfctx\Gamma$，则我们可知每个 $A_i$ 仅包含变量
$x_1,\dots,x_{i-1}$，且 $a$ 和 $A$ 仅包含变量
$x_1,\dots,x_n$。
\index{variable!in context}

在非正式的数学表述中，上下文是隐含的。
在证明的每一步，数学家都知道有哪些变量可用以及它们各自具有什么类型；这要么源于历史惯例
（例如，$n$ 通常表示数字，$f$ 表示函数等等），要么是因为变量通过诸如``设 $x$ 为一个实数''之类的语句被显式地引入。
我们将在 \cref{sec:more-formal-pi,sec:more-formal-sigma} 中讨论使用显式上下文的一些好处。

我们用 $B[a/x]$ 表示\define{替换}操作：以项 $a$ 代入项 $B$ 中变量 $x$ 的所有自由出现处，
\index{substitution}%
其间可能对约束变量进行避免捕获的重命名，
\index{variable!and substitution}%
详见 \cref{sec:function-types}。
替换的一般形式
%
\[
   B[a_1,\dots,a_n/x_1,\dots,x_n]
\]
%
是将表达式 $a_1,\dots,a_n$ 同时代入变量 $x_1,\dots,x_n$ 的位置。

在表达式 $B$ 中\define{绑定变量} $x$，
\indexdef{variable!bound}%
意味着将二者一并纳入一个更大的表达式，称为\define{抽象}，
\indexdef{abstraction}%
其目的在于表达 $x$ 相对于 $B$ 是``局部的''这一事实，即它不应与出现在别处的其他 $x$ 相混淆。
约束变量对程序员而言十分熟悉，对数学家来说则不那么熟悉。
根据具体情况，绑定变量的记法有多种，例如 $x \mapsto B$、$\lam x B$ 和 $x \,.\, B$。
对于用项 $a$ 替换抽象表达式中变量这一操作，我们可以记作 $C[a]$；即可以定义 $(x.B)[a]$ 为 $B[a/x]$。
如 \cref{sec:function-types} 所述，在表达式内部统一更改约束变量的名称（即``$\alpha$-变换''）
\index{alpha-conversion@$\alpha $-conversion}%
不会改变该表达式。
因此，严格地说，一个表达式是若干仅在约束变量名称上有所不同的句法形式所构成的等价类。

我们也可以将判断
\[
  x_1:A_1, x_2:A_2,\dots,x_n:A_n \vdash a : A
\]
中的每个变量 $x_i$ 视为在其\define{作用域}内被绑定，
\indexdef{variable!scope of}%
\index{scope}%
该作用域由表达式 $A_{i+1},
\ldots, A_n$、$a$ 和 $A$ 组成。

\section{第一种表述}
\label{sec:syntax-informally}

我们类型论中的对象与类型都可以写成项，所用语法如下，它是 $\lambda$-演算的扩展，引入了\emph{变量} $x, x',\dots$、
\index{variable}%
\emph{原始常量}
\index{primitive!constant}%
\index{constant!primitive}%
$c,c',\dots$、\emph{定义常量}\index{constant!defined} $f,f',\dots$，以及项的构造操作：
%
\[
  t \production x \mid \lam{x} t \mid t(t') \mid c \mid f
\]
%
这里所使用的记号意味着，一个项 $t$ 要么是一个变量 $x$，要么形如 $\lam{x} t$（其中 $x$ 是变量，$t$ 是项），
要么形如 $t(t')$（其中 $t$ 和 $t'$ 是项），要么是一个原始常量 $c$，要么是一个定义常量 $f$。
语法标记``$\lambda$''、``(''、``)''与``.''都是用于引导人眼阅读的标点符号。

我们把 $t(t_1,\dots,t_n)$ 用作重复应用 $t(t_1)(t_2)\dots (t_n)$ 的缩写。
\index{infix notation}%
当 $\star$ 是一个原始常量或定义常量时，我们也可以使用\emph{中缀}记号，将 $\star(t_1,t_2)$ 写作 $t_1\;
\star\; t_2$。

每个定义常量具有零条、一条或多条\define{定义方程}。
\index{equation, defining}%
\index{defining equation}%
定义常量分为两种。
一种\emph{显式定义常量}
\index{constant!explicit}
$f$ 具有唯一一条定义方程
  \[ f(x_1,\dots,x_n)\defeq t,\]
其中 $t$ 不涉及 $f$。
%
例如，我们可以引入显式定义常量 $\circ$，其定义方程为
  \[ \circ (x,y)(z) \defeq x(y(z)),\]
并采用中缀记号 $x\circ y$ 来表示 $\circ(x,y)$。这当然就是函数的复合。

第二种定义常量用于指定一个（带参数的）映射
$f(x_1,\dots,x_n,x)$，其中 $x$ 在一个由零个或多个原始常量生成的类型上变化。
对于每一个这样的原始常量 $c$，都有一条如下形式的定义方程
\[
  f(x_1,\dots,x_n,c(y_1,\dots,y_m)) \defeq t,
\]
其中 $f$ 可以出现在 $t$ 中，但其出现方式须使得这些方程显然能够确定一个全函数。
这类已定义函数的典型例子是自然数上由原始递归定义的函数。我们可以将这种函数定义方式称为\emph{全递归定义}。
\index{total!recursive definition}%
在计算机科学与逻辑学中，函数在递归数据类型上的这种定义方式被称为\define{结构递归定义}。
\index{definition!by structural recursion}%
\index{structural!recursion}%
\index{recursion!structural}%

\define{可转换性}
\index{convertibility of terms}%
\index{term!convertibility of}%
项 $t$ 与 $t'$ 之间的可转换性 $t \conv t'$，是由常量的定义方程、
函数类型的计算规则\index{computation rule!for function types}
%
\[
  (\lam{x} t)(u) \defeq t[u/x],
\]
%
以及使它关于应用和 $\lambda$-抽象\index{lambda abstraction@$\lambda$-abstraction}成为\emph{同余关系}的规则所生成的等价关系：
%


\begin{itemize}
\item 如果 $t \conv t'$ 且 $s \conv s'$，那么 $t(s) \conv t'(s')$，并且
\item 如果 $t \conv t'$，那么 $(\lam{x} t) \conv (\lam{x} t')$。
\end{itemize}


\noindent
进而，判断相等性 $t \jdeq u : A$ 可由如下单条规则导出：
%


\begin{itemize}
\item 如果 $t:A$，$u:A$，且 $t \conv u$，那么 $t \jdeq u : A$。
\end{itemize}


%
判断相等性是一个等价关系。

请注意，此处的类型论表述与正文所使用的有所不同：它不包含函数的判断唯一性原理 $f \jdeq (\lam{x} f(x))$。
这样的相等性要求判断相等性对所涉项的类型敏感（因为仅当已知 $f$ 是函数时，此相等关系才有意义），而在此处的表述中，可转换性关系与类型无关。
\cref{sec:syntax-more-formally} 中的第二种表述则包含了这条唯一性原理。

\subsection{类型宇宙}

我们假设存在一个由原始常量表示的\define{宇宙}层次：
\index{type!universe}
%
\begin{equation*}
  \UU_0, \quad \UU_1, \quad  \UU_2, \quad \ldots
\end{equation*}
%
宇宙的前两条规则表明它们构成类型的累积层次：
%


\begin{itemize}
\item 对 $m < n$，有 $\UU_m : \UU_n$，
\item 如果 $A:\UU_m$ 且 $m \le n$，那么 $A:\UU_n$，
\end{itemize}


%
第三条规则表达了这样一个意思：宇宙中的对象可以作为类型，出现在判断中冒号的右侧：
%


\begin{itemize}
\item 如果 $\Gamma \vdash A : \UU_n$，并且 $x$ 是一个新变量，%
\footnote{所谓“新”是指它不在 $\Gamma$ 或 $A$ 中出现。}
那么 $\vdash (\Gamma, x:A)\; \ctx$。
\end{itemize}


%
在本书正文中，类型 $A$ 与 $B$ 之间的判断相等性 $A \jdeq B : \UU_n$ 通常简写为 $A \jdeq B$。
这是类型歧义\index{typical ambiguity}的一个实例：我们总可以切换到更大的宇宙，但这并不影响判断的有效性。

下面的转换规则允许我们在类型判断中用一个相等的类型替换原类型：
%


\begin{itemize}
\item 如果 $a:A$ 且 $A \jdeq B$，那么 $a:B$。
\end{itemize}

\subsection{依值函数类型（\texorpdfstring{$\Pi$}{Π}-类型）}

我们引入原始常量 $c_\Pi$，但将
$c_\Pi(A,\lam{x} B)$ 写作 $\tprd{x:A}B$。
关于此类表达式及形如 $\lam{x} b$ 的表达式的判断，由以下规则引入：
%


\begin{itemize}
\item 若 $\Gamma \vdash A:\UU_n$ 且 $\Gamma,x:A \vdash B:\UU_n$，则 $\Gamma \vdash \tprd{x:A}B : \UU_n$
\item 若 $\Gamma, x:A \vdash b:B$ 则 $\Gamma \vdash (\lam{x} b) : (\tprd{x:A} B)$
\item 若 $\Gamma\vdash g:\tprd{x:A} B$ 且 $\Gamma\vdash t:A$ 则 $\Gamma\vdash g(t):B[t/x]$
\end{itemize}


%
如果 $x$ 在 $B$ 中不自由出现，我们将 $\tprd{x:A} B$ 缩写为非依值函数类型
$A\rightarrow B$，并导出以下规则：
%


\begin{itemize}
\item 若 $\Gamma\vdash g:A \rightarrow B$ 且 $\Gamma\vdash t:A$ 则 $\Gamma\vdash g(t):B$
\end{itemize}


%
使用非依值函数类型并隐去上下文 $\Gamma$，上述规则可以写成如下另一种形式，我们在本附录后续部分将采用这种形式：
%


\begin{itemize}
\item 若 $A:\UU_n$ 且 $B:A\to\UU_n$，则 $\tprd{x:A}B(x) : \UU_n$
\item 若 $x:A \vdash b:B(x)$ 则 $ \lam{x} b : \tprd{x:A} B(x)$
\item 若 $g:\tprd{x:A} B(x)$ 且 $t:A$ 则 $g(t):B(t)$
\end{itemize}


%

\subsection{依值对类型（\texorpdfstring{$\Sigma$}{Σ}-类型）}

我们引入原始常量 $c_\Sigma$ 和 $c_{\mathsf{pair}}$。An
expression of the form $c_\Sigma(A,\lam{a} B)$ is written as $\sm{a:A}B$,
and an expression of the form $c_{\mathsf{pair}}(a,b)$ is written as $\tup
a b$. We write $A\times B$ instead of $\sm{x:A} B$ if $x$ is not free in $B$.

Judgments concerning such expressions are introduced by the following
rules:
%


\begin{itemize}
\item if $A:\UU_n$ and $B: A \rightarrow \UU_n$, then $\sm{x:A}B(x) : \UU_n$
\item if, in addition, $a:A$ and $b:B(a)$, then $\tup a b:\sm{x:A}B(x)$
\end{itemize}


%
If we have $A$ and $B$ as above, $C : (\sm{x:A}B(x)) \rightarrow \UU_m$, and
\[
  d:\tprd{x:A}{y:B(x)} C(\tup x y)
\]
we can introduce a defined constant
\[
  f:\tprd{p:\sm{x:A}B(x)} C(p)
\]
with the defining equation
\[
  f(\tup x y)\defeq d(x,y).
\]
%
Note that $C$, $d$, $x$, and $y$ may contain extra implicit parameters $x_1,\ldots,x_n$ if they were obtained in some non-empty context; therefore, the fully explicit recursion schema is
%
\begin{narrowmultline*}
 f(x_1,\dots,x_n,\tup{x(x_1,\dots,x_n)}{y(x_1,\dots,x_n)}) \defeq
 \narrowbreak
 d(x_1,\dots,x_n,\tup{x(x_1,\dots,x_n)}{y(x_1,\dots,x_n)}).
\end{narrowmultline*}

\subsection{余积类型}

我们引入原始常量 $c_+$、$c_\inlsym$ 与 $c_\inrsym$。
我们将 $c_+(A,B)$ 写作 $A+B$，将 $c_\inlsym(a)$ 写作 $\inl(a)$，并将 $c_\inrsym(a)$ 写作 $\inr(a)$：
%


\begin{itemize}
\item 若 $A,B : \UU_n$，则 $A + B : \UU_n$
\item 此外，$\inl: A \rightarrow A+B$ 且 $\inr: B \rightarrow A+B$
\end{itemize}


%
对于如上所述的 $A$ 与 $B$，如果有 $C : A+B \rightarrow \UU_m$、$d:\tprd{x:A} C(\inl(x))$ 以及 $e:\tprd{y:B} C(\inr(y))$，
那么可以引入一个定义常量 $f:\tprd{z:A+B}C(z)$，其定义方程为
%
\begin{equation*}
  f(\inl(x)) \defeq d(x)
  \qquad\text{and}\qquad
  f(\inr(y)) \defeq e(y).
\end{equation*}

\subsection{有限类型}

我们引入满足以下规则的原始常量 $\ttt$、$\emptyt$ 与 $\unit$：
%


\begin{itemize}
\item $\emptyt : \UU_0$，$\unit : \UU_0$
\item $\ttt:\unit$
\end{itemize}

给定 $C : \emptyt \rightarrow \UU_n$，我们可以引入一个定义常量 $f:\tprd{x:\emptyt} C(x)$，不带定义方程。

给定 $C : \unit \rightarrow \UU_n$ 与 $d : C(\ttt)$，我们可以引入一个定义常量 $f:\tprd{x:\unit} C(x)$，其定义方程为 $f(\ttt) \defeq d$。

\subsection{自然数}

自然数类型是通过引入满足以下规则的原始常量 $\N$、$0$ 和 $\suc$ 得到的：
%


\begin{itemize}
  \item $\N : \UU_0$，
  \item $0:\N$，
  \item $\suc:\N\rightarrow \N$。
\end{itemize}


%
此外，我们可以用原始递归来定义函数。
如果有 $C : \N \rightarrow \UU_k $，那么只要具备
%
\begin{align*}
  d & : C(0) \\
  e & : \tprd{x:\N}(C(x)\rightarrow C(\suc (x)))
\end{align*}
%
就可以引入一个定义常量 $f:\tprd{x:\N}C(x)$，其定义方程为
%
\begin{equation*}
  f(0) \defeq d
  \qquad\text{and}\qquad
  f(\suc (x)) \defeq e(x,f(x)).
\end{equation*}

\subsection{\texorpdfstring{$W$}{W}-类型}

对于 $W$-类型，我们引入原始常量 $c_\wtypesym$ 与 $c_\suppsym$。
我们将形如 $c_\wtypesym(A,\lam{x} B)$ 的表达式写作 $\wtype{x:A}B$，并将形如 $c_\suppsym(x,u)$ 的表达式写作 $\supp(x,u)$：
%


\begin{itemize}
\item 若 $A:\UU_n$ 且 $B: A \rightarrow \UU_n$，则 $\wtype{x:A}B(x) : \UU_n$
\item 此外，若 $a:A$ 且 $u:B(a)\rightarrow \wtype{x:A}B(x)$，则 $\supp(a,u):\wtype{x:A}B(x)$。
\end{itemize}


% 
这里同样可以通过全递归来定义函数。如果有如上所述的 $A$ 与 $B$ 以及 $C : (\wtype{x:A}B(x)) \rightarrow \UU_m$，那么只要具备
\[
  d:\tprd{a:A}{u:B(a) \rightarrow \wtype{x:A}B(x)}((\tprd{y:B(a)}C(u(y))) \rightarrow C(\supp(a,u)))
\]
就可以引入一个定义常量
$f:\tprd{z:\wtype{x:A}B(x)} C(z)$，
其定义方程为
\[
  f(\supp(a,u)) \defeq d(a,u,f\circ u).
\]

\subsection{相等类型}

我们引入原始常量 $c_\idsym$ 和 $c_{\refl{}}$。我们将 $c_\idsym(A,a,b)$ 记作 $\id[A] a b$，并将 $c_{\refl{}}(A,a)$ 记作 $\refl a$，此时默认 $a:A$：
%


\begin{itemize}
\item 若 $A : \UU_n$，$a:A$，且 $b:A$，则 $\id[A] a b : \UU_n$。
\item 若 $a:A$，则 $\refl a :\id[A] a a$。
\end{itemize}


%
给定 $a:A$，如果 $y:A, z:\id[A] a y \vdash C : \UU_m$ 且 $\vdash d:C[a,\refl{a}/y,z]$，那么可以引入一个定义常量
\[
  f:\tprd{y:A}{z:\id[A] a y} C
\]
其定义方程为
\[
  f(a,\refl{a})\defeq d.
\]

\section{第二套表示法}
\label{sec:syntax-more-formally}

本节中有三种判断


\begin{mathpar}
\wfctx\Gamma
\and
\oftp\Gamma{a}{A}
\and
\jdeqtp\Gamma{a}{a'}{A}
\end{mathpar}


我们通过给出推导它们的规则来规定这些判断。一条典型的 \define{推导规则（inference rule）}
\indexsee{inference rule}{rule}%
\indexdef{rule}%
具有如下形式
%
\begin{equation*}
  \inferrule*[right=\textsc{Name}]
  {\mathcal{J}_1 \\ \cdots \\ \mathcal{J}_k}
  {\mathcal{J}}
\end{equation*}
%
它表示：只要我们已经推导出 \define{前提} $\mathcal{J}_1, \ldots, \mathcal{J}_k$，就可以推导出 \define{结论} $\mathcal{J}$。
（注意，由于这些是判断而非类型，因此它们并非 \cref{sec:types-vs-sets} 意义上类型论\emph{内部}的假设；相反，它们是演绎系统中的假设，即元理论中的假设。）
在规则右侧，我们写出规则的 \textsc{Name}，且可能附有在应用规则前需要检查的附加条件。

一个判断的 \define{推导}
\index{derivation}%
是一棵由这些推导规则构造而成的树，该判断位于树的根部。
例如，利用下面给出的规则，以下是 $\oftp{\emptyctx}{\lamu{x:\unit} x}{\unit\to\unit}$ 的一个推导。
%


\begin{mathpar}
\inferrule*[right=$\Pi$-\rintro]
  {\inferrule*[right=$\Vble$]
    {\inferrule*[right=\ctx-\textsc{ext}]
      {\inferrule*[right=$\unit$-\rform]
        {\inferrule*[right=\ctx-\textsc{emp}]
          {\ }
          {\wfctx {\emptyctx}}}
        {\oftp{}{\unit}{\UU_0}}}
      {\wfctx {\tmtp x\unit}}}
   {\oftp{\tmtp x\unit}{x}{\unit}}}
 {\oftp{\emptyctx}{\lamu{x:\unit} x}{\unit\to\unit}}
\end{mathpar}

\subsection{上下文}
\label{subsec:contexts}

\index{context}%
上下文是一个列表
%
\begin{equation*}
  \tmtp{x_1}{A_1}, \tmtp{x_2}{A_2}, \ldots, \tmtp{x_n}{A_n}
\end{equation*}
%
它表示互不相同的变量 \indexdef{variable}% $x_1, \ldots, x_n$ 被假定分别具有类型 $A_1, \ldots, A_n$。该列表可以为空。我们用字母 $\Gamma$ 和 $\Delta$ 简写上下文，并可将它们并置以构成更大的上下文。

判断 $\wfctx{\Gamma}$ 形式上表达了 $\Gamma$ 是良构上下文这一事实，并由推导规则
%
所规定。

\begin{mathpar}
  \inferrule*[right=\ctx-\textsc{emp}]
  {\ }
  {\wfctx\emptyctx}
\and
  \inferrule*[right=\ctx-\textsc{ext}]
  {\oftp{\tmtp{x_1}{A_1}, \ldots, \tmtp{x_{n-1}}{A_{n-1}}}{A_n}{\UU_i}}
  {\wfctx{(\tmtp{x_1}{A_1}, \ldots, \tmtp{x_n}{A_n})}}
\end{mathpar}

%
第二条规则有一个附带条件：变量 $x_n$ 必须与变量 $x_1, \ldots, x_{n-1}$ 不同。
注意，规则 $\ctx$-\textsc{ext} 的前提和结论是不同形式的判断：前提说的是在变量 $x_1, \ldots, x_{n-1}$ 的上下文中，表达式 $A_n$ 具有类型 $\UU_i$；而结论说的是扩展后的上下文 $(\tmtp{x_1}{A_1}, \ldots, \tmtp{x_n}{A_n})$ 是良构的。

该系统具有以下元理论性质：如果任何形如 $\oftp{\Gamma}{a}{A}$ 或 $\jdeqtp\Gamma{a}{a'}{A}$ 的判断是可推导的，那么判断 $\wfctx\Gamma$（即上下文 $\Gamma$ 是良构的）也是可推导的。
所有规则的前提的选取，只包含恰好足以使该性质可证的良构性假设，但不再多。
例如，$\ctx$-\textsc{ext} 没有必要假设 $(\tmtp{x_1}{A_1}, \ldots, \tmtp{x_{n-1}}{A_{n-1}})$ 的良构性，因为这可由其前提的可推导性推出；但对于下一节中的 $\Vble$ 规则，假设其上下文的良构性则是必要的。
这样的选择只是精确表述类型理论的众多可能方式之一，但对此类问题的详细探讨已超出本附录的范围。

\subsection{结构性规则}

\index{structural!rules|(}%
\index{rule!structural|(}%

上下文包含假设这一事实由下面这条规则来表达：我们可以推导出那些在上下文中列出的类型判断：
%

\begin{mathpar}
  \inferrule*[right=$\Vble$]
  {\wfctx {(\tmtp{x_1}{A_1}, \ldots, \tmtp{x_n}{A_n})} }
  {\oftp{\tmtp{x_1}{A_1}, \ldots, \tmtp{x_n}{A_n}}{x_i}{A_i}}
\end{mathpar}

%
与 $\ctx$-\textsc{ext} 类似，规则 $\Vble$ 的前提和结论也是不同形式的判断，只不过现在方向反了过来：我们从一个良构的上下文出发，推导出一个类型判断。

下面这些重要的原则，称为 \define{替换（substitution）}
\indexdef{rule!of substitution}%
和
\define{弱化（weakening）}，
\indexdef{rule!of weakening}%
并不需要显式地假设。
相反，通过对所有可能推导的结构进行归纳，可以证明：只要这些规则的前提是可推导的，它们的结论也是可推导的。\footnote{这样的规则称为 \define{可容许的（admissible）}\indexdef{rule!admissible}\indexsee{admissible!rule}{rule, admissible}。}
对于类型判断，这些原则表现为：
%

\begin{mathpar}
  \inferrule*[right=$\Subst_1$]
  {\oftp\Gamma{a}{A} \\ \oftp{\Gamma,\tmtp xA,\Delta}{b}{B}}
  {\oftp{\Gamma,\Delta[a/x]}{b[a/x]}{B[a/x]}}
\and
  \inferrule*[right=$\Weak_1$]
  {\oftp\Gamma{A}{\UU_i} \\ \oftp{\Gamma,\Delta}{b}{B}}
  {\oftp{\Gamma,\tmtp xA,\Delta}{b}{B}}
\end{mathpar}

而对于判断相等性，它们则变为：

\begin{mathpar}
  \inferrule*[right=$\Subst_2$]
  {\oftp\Gamma{a}{A} \\ \jdeqtp{\Gamma,\tmtp xA,\Delta}{b}{c}{B}}
  {\jdeqtp{\Gamma,\Delta[a/x]}{b[a/x]}{c[a/x]}{B[a/x]}}
\and
  \inferrule*[right=$\Subst_3$]
  {\jdeqtp\Gamma{a}{b}{A} \\ \oftp{\Gamma,\tmtp xA,\Delta}{c}{C}}
  {\jdeqtp{\Gamma,\Delta[a/x]}{c[a/x]}{c[b/x]}{C[a/x]}}
\and
  \inferrule*[right=$\Weak_2$]
  {\oftp\Gamma{A}{\UU_i} \\ \jdeqtp{\Gamma,\Delta}{b}{c}{B}}
  {\jdeqtp{\Gamma,\tmtp xA,\Delta}{b}{c}{B}}
\end{mathpar}

%
除了为每个类型形成子给出的判断相等性规则之外，我们还假设判断相等性是一个等价关系，且为类型判断所保持。

\begin{mathparpagebreakable}
  \inferrule*{\oftp\Gamma{a}{A}}{\jdeqtp\Gamma{a}{a}{A}}
\and
  \inferrule*{\jdeqtp\Gamma{a}{b}{A}}{\jdeqtp\Gamma{b}{a}{A}}
\and
  \inferrule*{\jdeqtp\Gamma{a}{b}{A} \\ \jdeqtp\Gamma{b}{c}{A}}{\jdeqtp\Gamma{a}{c}{A}}
\and
  \inferrule*{\oftp\Gamma{a}{A} \\ \jdeqtp\Gamma{A}{B}{\UU_i}}{\oftp\Gamma{a}{B}}
\and
  \inferrule*{\jdeqtp\Gamma{a}{b}{A} \\ \jdeqtp\Gamma{A}{B}{\UU_i}}{\jdeqtp\Gamma{a}{b}{B}}
\end{mathparpagebreakable}

%
最后，我们假设判断相等性是一个同余关系且为类型判断所保持，即每个类型形成子和项形成子都在各自的每个参数中保持判断相等性。
例如，除了 $\Pi$-\rintro\ 规则之外，我们还假设规则
\[
  \inferrule*[right=$\Pi$-\rintro-eq]
  {\oftp\Gamma{A}{\UU_i} \\
   \oftp{\Gamma,\tmtp xA}{B}{\UU_i} \\
   \jdeqtp{\Gamma,\tmtp xA}{b}{b'}{B}}
  {\jdeqtp\Gamma{\lamu{x:A} b}{\lamu{x:A'} b'}{\tprd{x:A} B}}
\]
为补全依值函数类型的情形，我们还假设了两个类似的规则：$\Pi$-\textsc{form-eq}\ 和 $\Pi$-\textsc{elim-eq}。
这些局部原则（在每个类型上）合在一起，就蕴含了上面给出的全局同余原则 $\Subst_2$ 和 $\Subst_3$。
为简洁起见，我们将省略这些局部规则。

\index{rule!structural|)}%
\index{structural!rules|)}%

\subsection{类型宇宙}

\index{type!universe}%

我们假设存在一个无限的类型宇宙层级
%
\begin{equation*}
  \UU_0, \quad \UU_1, \quad  \UU_2, \quad \ldots
\end{equation*}
%
每个宇宙都包含于下一个宇宙之中，并且任何在 $\UU_i$ 中的类型也在 $\UU_{i+1}$ 中：
%


\begin{mathpar}
\inferrule*[right=\UU-\textsc{intro}]
  {\wfctx \Gamma }
  {\oftp\Gamma{\UU_i}{\UU_{i+1}}}
\and
\inferrule*[right=\UU-\textsc{cumul}]
  {\oftp\Gamma{A}{\UU_i}}
  {\oftp\Gamma{A}{\UU_{i+1}}}
\end{mathpar}


%
我们将如此设定类型论的规则，使得 $\oftp\Gamma{a}{A}$ 蕴含对于某个 $i$ 有 $\oftp\Gamma{A}{\UU_i}$。
换言之，如果 $A$ 扮演类型的角色，那么它就在某个宇宙中。
我们类型系统的另一个性质是，$\jdeqtp\Gamma{a}{b}{A}$
蕴含 $\oftp\Gamma{a}{A}$ 和 $\oftp\Gamma{b}{A}$。

\subsection{依值函数类型 (\texorpdfstring{$\Pi$}{Π}-类型)}
\label{sec:more-formal-pi}

\index{type!dependent function}%
\index{type!function}%

在\cref{sec:function-types}中，我们引入了非依值函数 $A\to B$，以便将类型族定义为函数 $\lam{x:A} B:A\to\UU_i$，
进而引出依值函数类型 $\tprd{x:A} B$。
但有了显式上下文，我们可以将 $\lam{x:A} B:A\to\UU_i$ 替换为判断
%
\begin{equation*}
  \oftp{\tmtp xA}{B}{\UU_i}.
\end{equation*}
%
这样一来，我们就可以直接定义依值函数，而不再需要借助非依值函数。
这遵循了一条一般原则：每个类型形成子及其相关的常量和规则，都应当独立于所有其他类型形成子而引入。
%
事实上，从今往后，每一个类型形成子都将通过以下方式被系统地引入：

\begin{itemize}
\item 一条\define{形成规则（formation rule）}，说明该类型形成子在何时可以应用；\index{formation rule}\index{rule!formation}
\item 若干条\define{引入规则（introduction rules）}，说明如何构造该类型的元素；\index{introduction rule}\index{rule!introduction}
\item \define{消去规则（elimination rules）}，或称归纳原理，说明如何使用该类型的元素；
  \index{induction principle}\index{eliminator}
\item \define{计算规则（computation rules）}，即解释当消去规则作用于引入规则的结果时会发生什么的判断相等性；
  \index{computation rule}
  \indexsee{rule!computation}{computation rule}
\item 可选的\define{唯一性原理（uniqueness principles）}，即解释该类型的每个元素如何由作用于其上的消去规则的结果唯一确定的判断相等性。
  \index{uniqueness!principle}
  \indexsee{principle!uniqueness}{uniqueness principle}
\end{itemize}

（另见\cref{rmk:introducing-new-concepts}。）

对于依值函数类型，这些规则是：
%

\begin{mathparpagebreakable}
  \def\premise{\oftp{\Gamma}{A}{\UU_i} \and \oftp{\Gamma,\tmtp xA}{B}{\UU_i}}
  \inferrule*[right=$\Pi$-\rform]
    \premise
    {\oftp\Gamma{\tprd{x:A}B}{\UU_i}}
\and
  \inferrule*[right=$\Pi$-\rintro]
  {\oftp{\Gamma,\tmtp xA}{b}{B}}
  {\oftp\Gamma{\lam{x:A} b}{\tprd{x:A} B}}
\and
  \inferrule*[right=$\Pi$-\relim]
  {\oftp\Gamma{f}{\tprd{x:A} B} \\ \oftp\Gamma{a}{A}}
  {\oftp\Gamma{f(a)}{B[a/x]}}
\and
  \inferrule*[right=$\Pi$-\rcomp]
  {\oftp{\Gamma,\tmtp xA}{b}{B} \\ \oftp\Gamma{a}{A}}
  {\jdeqtp\Gamma{(\lam{x:A} b)(a)}{b[a/x]}{B[a/x]}}
\and
  \inferrule*[right=$\Pi$-\runiq]
  {\oftp\Gamma{f}{\tprd{x:A} B}}
  {\jdeqtp\Gamma{f}{(\lamu{x:A}f(x))}{\tprd{x:A} B}}
\end{mathparpagebreakable}

表达式 $\lam{x:A} b$ 绑定了 $b$ 中 $x$ 的自由出现，$\tprd{x:A} B$ 对 $B$ 中的 $x$ 也是如此。

当 $x$ 不在 $B$ 中自由出现，因而 $B$ 不依赖于 $A$ 时，我们便得到普通函数类型 $A\to B \defeq \tprd{x:A} B$ 作为特例。
我们将其作为 $\to$ 的\emph{定义}。

我们可以将表达式 $\lam{x:A} b$ 缩写为 $\lamu{x:A} b$，在此约定：被省略的类型 $A$ 应在类型检查之前被适当地填入。

\subsection{依值对类型 (\texorpdfstring{$\Sigma$}{Σ}-类型)}
\label{sec:more-formal-sigma}

\index{type!dependent pair}%
\index{type!product}%

在\cref{sec:sigma-types}中，我们需要 $\to$ 和 $\prdsym$ 类型来定义 $\smsym$ 的引入与消去规则；而正如 $\prdsym$ 的情况一样，上下文允许我们独立地陈述 $\smsym$ 的规则。
回顾一下，像 $\Sigma$ 这样的正类型（positive type）的消去规则被称为\emph{归纳（induction）}，记作 $\ind{}$。
%

\begin{mathparpagebreakable}
  \def\premise{\oftp{\Gamma}{A}{\UU_i} \and \oftp{\Gamma,\tmtp xA}{B}{\UU_i}}
  \inferrule*[right=$\Sigma$-\rform]
    \premise
    {\oftp\Gamma{\tsm{x:A} B}{\UU_i}}
  \and
  \inferrule*[right=$\Sigma$-\rintro]
    {\oftp{\Gamma, \tmtp x A}{B}{\UU_i} \\
     \oftp\Gamma{a}{A} \\ \oftp\Gamma{b}{B[a/x]}}
    {\oftp\Gamma{\tup ab}{\tsm{x:A} B}}
  \and
  \inferrule*[right=$\Sigma$-\relim]
    {\oftp{\Gamma, \tmtp z {\tsm{x:A} B}}{C}{\UU_i} \\
     \oftp{\Gamma,\tmtp x A,\tmtp y B}{g}{C[\tup x y/z]} \\
     \oftp\Gamma{p}{\tsm{x:A} B}}
    {\oftp\Gamma{\ind{\tsm{x:A} B}(z.C,x.y.g,p)}{C[p/z]}}
  \and
  \inferrule*[right=$\Sigma$-\rcomp]
    {\oftp{\Gamma, \tmtp z {\tsm{x:A} B}}{C}{\UU_i} \\
     \oftp{\Gamma, \tmtp x A, \tmtp y B}{g}{C[\tup x y/z]} \\\\
     \oftp\Gamma{a}{A} \\ \oftp\Gamma{b}{B[a/x]}}
    {\jdeqtp\Gamma{\ind{\tsm{x:A} B}(z.C,x.y.g,\tup{a}{b})}{g[a,b/x,y]}{C[\tup {a} {b}/z]}}
\end{mathparpagebreakable}

%
表达式 $\tsm{x:A} B$ 绑定 $x$ 在 $B$ 中的自由出现。
此外，由于 $\ind{\tsm{x:A} B}$ 的某些参数含有上下文 $\Gamma$ 之外的自由变量，我们（沿用上文的变量名）在 $C$ 中绑定 $z$，在 $g$ 中绑定 $x$ 与 $y$。
这些绑定记作 $z.C$ 与 $x.y.g$，用以指明被绑定的变量名。
\index{variable!bound}%
特别地，我们将 $\ind{\tsm{x:A} B}$ 视为一个原语，其中两个参数含有绑定器；这表面上类似于把 $\ind{\tsm{x:A} B}$ 看作一个以函数为参数的函数，但实则不同。

当 $B$ 不含 $x$ 的自由出现时，我们便得到 Cartesian 积这一特例 $A \times B \defeq \tsm{x:A} B$。
我们就以此作为 Cartesian 积的\emph{定义}。

请注意，我们并没有为 $\Sigma$ 类型预设判断唯一性原理（尽管我们本可以这么做）；相应的命题唯一性原理的证明参见 \cref{thm:eta-sigma}。

\subsection{余积类型}

\index{type!coproduct}%

\begin{mathparpagebreakable}
  \inferrule*[right=$+$-\rform]
  {\oftp\Gamma{A}{\UU_i} \\ \oftp\Gamma{B}{\UU_i}}
  {\oftp\Gamma{A+B}{\UU_i}}
\\
  \inferrule*[right=$+$-\rintro${}_1$]
  {\oftp\Gamma{A}{\UU_i} \\ \oftp\Gamma{B}{\UU_i} \\\\ \oftp\Gamma{a}{A}}
  {\oftp\Gamma{\inl(a)}{A+B}}
\and
  \inferrule*[right=$+$-\rintro${}_2$]
  {\oftp\Gamma{A}{\UU_i} \\ \oftp\Gamma{B}{\UU_i} \\\\ \oftp\Gamma{b}{B}}
  {\oftp\Gamma{\inr(b)}{A+B}}
\\
  \inferrule*[right=$+$-\relim]
  {\oftp{\Gamma,\tmtp z{(A+B)}}{C}{\UU_i} \\\\
   \oftp{\Gamma,\tmtp xA}{c}{C[\inl(x)/z]} \\
   \oftp{\Gamma,\tmtp yB}{d}{C[\inr(y)/z]} \\\\
   \oftp\Gamma{e}{A+B}}
  {\oftp\Gamma{\ind{A+B}(z.C,x.c,y.d,e)}{C[e/z]}}
\and
  \inferrule*[right=$+$-\rcomp${}_1$]
  {\oftp{\Gamma,\tmtp z{(A+B)}}{C}{\UU_i} \\
   \oftp{\Gamma,\tmtp xA}{c}{C[\inl(x)/z]} \\
   \oftp{\Gamma,\tmtp yB}{d}{C[\inr(y)/z]} \\\\
   \oftp\Gamma{a}{A}}
  {\jdeqtp\Gamma{\ind{A+B}(z.C,x.c,y.d,\inl(a))}{c[a/x]}{C[\inl(a)/z]}}
\and
  \inferrule*[right=$+$-\rcomp${}_2$]
  {\oftp{\Gamma,\tmtp z{(A+B)}}{C}{\UU_i} \\
   \oftp{\Gamma,\tmtp xA}{c}{C[\inl(x)/z]} \\
   \oftp{\Gamma,\tmtp yB}{d}{C[\inr(y)/z]} \\\\
   \oftp\Gamma{b}{B}}
  {\jdeqtp\Gamma{\ind{A+B}(z.C,x.c,y.d,\inr(b))}{d[b/y]}{C[\inr(b)/z]}}
\end{mathparpagebreakable}

%
在 $\ind{A+B}$ 中，$z$ 在 $C$ 中被绑定，$x$ 在 $c$ 中被绑定，$y$ 在 $d$ 中被绑定。

\subsection{空类型 \texorpdfstring{$\emptyt$}{0}}

\index{type!empty|(}%

\begin{mathparpagebreakable}
  \inferrule*[right=$\emptyt$-\rform]
  {\wfctx\Gamma}
  {\oftp\Gamma\emptyt{\UU_i}}
\and
  \inferrule*[right=$\emptyt$-\relim]
  {\oftp{\Gamma,\tmtp x\emptyt}{C}{\UU_i} \\ \oftp\Gamma{a}{\emptyt}}
  {\oftp\Gamma{\ind{\emptyt}(x.C,a)}{C[a/x]}}
\end{mathparpagebreakable}

%
在 $\ind{\emptyt}$ 中，$x$ 在 $C$ 中被绑定。空类型没有引入规则，也没有计算规则。

\index{type!empty|)}%

\subsection{单位类型 \texorpdfstring{$\unit$}{1}}
\label{sec:more-formal-unit}

\index{type!unit|(}%

\begin{mathparpagebreakable}
  \inferrule*[right=$\unit$-\rform]
  {\wfctx\Gamma}
  {\oftp\Gamma\unit{\UU_i}}
\and
  \inferrule*[right=$\unit$-\rintro]
  {\wfctx\Gamma}
  {\oftp\Gamma{\ttt}{\unit}}
\and
  \inferrule*[right=$\unit$-\relim]
  {\oftp{\Gamma,\tmtp x\unit}{C}{\UU_i} \\
   \oftp{\Gamma}{c}{C[\ttt/x]} \\
   \oftp\Gamma{a}{\unit}}
  {\oftp\Gamma{\ind{\unit}(x.C,c,a)}{C[a/x]}}
\and
  \inferrule*[right=$\unit$-\rcomp]
  {\oftp{\Gamma,\tmtp x\unit}{C}{\UU_i} \\
   \oftp{\Gamma}{c}{C[\ttt/x]}}
  {\jdeqtp\Gamma{\ind{\unit}(x.C,c,\ttt)}{c}{C[\ttt/x]}}
\end{mathparpagebreakable}

%
在 $\ind{\unit}$ 中，变量 $x$ 在 $C$ 中被绑定。

请注意，我们也没有为单位类型预设判断唯一性原理；相应的命题唯一性陈述的证明参见 \cref{sec:finite-product-types}。

\index{type!unit|)}%

\subsection{自然数类型}

\index{natural numbers|(}%

我们遵循 \cref{sec:inductive-types}，给出自然数的规则。

\begin{mathparpagebreakable}
  \def\premise{
     \oftp{\Gamma,\tmtp x{\N}}{C}{\UU_i} \\
     \oftp\Gamma{c_0}{C[0/x]} \\
     \oftp{\Gamma,\tmtp{x}\N,\tmtp y C}{c_s}{C[\suc(x)/x]}}
  %
  \inferrule*[right=$\N$-\rform]
  {\wfctx\Gamma}
  {\oftp\Gamma{\N}{\UU_i}}
\and
  \inferrule*[right=$\N$-\rintro${}_1$]
  {\wfctx\Gamma}
  {\oftp\Gamma{0}{\N}}
\and
  \inferrule*[right=$\N$-\rintro${}_2$]
  {\oftp\Gamma{n}{\N}}
  {\oftp\Gamma{\suc(n)}{\N}}
\and
  \inferrule*[right=$\N$-\relim]
  {\premise \\ \oftp\Gamma{n}{\N}}
  {\oftp\Gamma{\ind{\N}(x.C,c_0,x.y.c_s,n)}{C[n/x]}}
\and
  \inferrule*[right=$\N$-\rcomp${}_1$]
  {\premise}
  {\jdeqtp\Gamma{\ind{\N}(x.C,c_0,x.y.c_s,0)}{c_0}{C[0/x]}}
\and
  \inferrule*[right=$\N$-\rcomp${}_2$]
  {\premise \\ \oftp\Gamma{n}{\N}}
  {\Gamma\vdash
    {\begin{aligned}[t]
      &\ind{\N}(x.C,c_0,x.y.c_s,\suc(n)) \\
      &\quad \jdeq c_s[n,\ind{\N}(x.C,c_0,x.y.c_s,n)/x,y] : C[\suc(n)/x]
    \end{aligned}}}
\end{mathparpagebreakable}


%
在 $\ind{\N}$ 中，$x$ 在 $C$ 中被绑定，$x$ 和 $y$ 在 $c_s$ 中被绑定。

其他归纳定义的类型遵循相同的一般模式。

\index{natural numbers|)}%

\subsection{相等类型}

\label{sec:more-formal-identity}

\index{type!identity|(}%

此处的表述对应于 \cref{sec:identity-types} 中相等类型的（非基化的）道路归纳原理。

\begin{mathparpagebreakable}
  \inferrule*[right=$\idsym$-\rform]
  {\oftp\Gamma{A}{\UU_i} \\ \oftp\Gamma{a}{A} \\ \oftp\Gamma{b}{A}}
  {\oftp\Gamma{\id[A]{a}{b}}{\UU_i}}
\and
  \inferrule*[right=$\idsym$-\rintro]
  {\oftp\Gamma{A}{\UU_i} \\ \oftp\Gamma{a}{A}}
  {\oftp\Gamma{\refl a}{\id[A]aa}}
\and
  \inferrule*[right=$\idsym$-\relim]
  {\oftp{\Gamma,\tmtp xA,\tmtp yA,\tmtp p{\id[A]xy}}{C}{\UU_i} \\
   \oftp{\Gamma,\tmtp zA}{c}{C[z,z,\refl z/x,y,p]} \\
   \oftp\Gamma{a}{A} \\ \oftp\Gamma{b}{A} \\ \oftp\Gamma{p'}{\id[A]ab}}
  {\oftp\Gamma{\indid{A}(x.y.p.C,z.c,a,b,p')}{C[a,b,p'/x,y,p]}}
\and
  \inferrule*[right=$\idsym$-\rcomp]
  {\oftp{\Gamma,\tmtp xA,\tmtp yA,\tmtp p{\id[A]xy}}{C}{\UU_i} \\
   \oftp{\Gamma,\tmtp zA}{c}{C[z,z,\refl z/x,y,p]} \\
   \oftp\Gamma{a}{A}}
  {\jdeqtp\Gamma{\indid{A}(x.y.p.C,z.c,a,a,\refl a)}{c[a/z]}{C[a,a,\refl a/x,y,p]}}
\end{mathparpagebreakable}


%
在 $\indid{A}$ 中，$x$、$y$ 和 $p$ 在 $C$ 中被绑定，而 $z$ 在 $c$ 中被绑定。

\index{type!identity|)}%

\subsection{定义}

\index{definition}%

尽管目前列出的规则已允许我们直接构造所需的一切，但为了方便起见，我们仍然希望能够使用命名的常量，例如 $\isequiv$。非形式地说，我们可以把这些常量简单地看作缩写，但在形式化中，情况要稍微微妙一些。

举例来说，考虑函数复合 $\circ$，它将 $f:A\to B$ 和 $g:B\to C$ 映为 $g\circ f:A\to C$。有些出人意料的是，为了在形式化中实现这一点，$\circ$ 必须接受的参数不仅是 $f$ 和 $g$，还包括它们的类型 $A$、$B$、$C$：
%
\begin{narrowmultline*}
  {\circ} \defeq \lam{A:\UU_i}{B:\UU_i}{C:\UU_i}
  \narrowbreak
  \lam{g:B\to C}{f:A\to B}{x:A} g(f(x)).
\end{narrowmultline*}
%
从实用角度出发，我们并不希望每次应用 $\circ$ 时都标注上 $A$、$B$ 和 $C$，因为它们通常很容易从上下文中推断出来。我们希望简单地写作 $g\circ f$。那么，严格来说，$g \circ f$ 并不是 $\lam{x : A} g(f(x))$ 的缩写，因为它涉及了我们想要隐藏的额外 \define{隐式参数（implicit arguments）}。
\index{implicit argument}

推断隐式参数、处理类型歧义（typical ambiguity）\index{typical ambiguity}（\cref{sec:universes}）、确保符号只被定义一次等工作，统称为 \define{精细化（elaboration）}。\index{elaboration, in type theory}
精细化必须在检查推导之前完成，因此通常不作为核心类型论的一部分来呈现。然而，任何不进行精细化的类型论实现几乎都无法使用；更多讨论参见 \cite{Coq,norell2007towards}。

\section{同伦类型论}
\label{sec:hott-features}

在本节中，我们将阐述使同伦类型论区别于标准 Martin-L\"{o}f 类型论的额外公理：函数外延性、泛等公理，以及高阶归纳类型。我们以第二种表述（\cref{sec:syntax-more-formally}）的风格来叙述它们，尽管第一种表述（\cref{sec:syntax-informally}）也同样适用。

\subsection{函数外延性与泛等性}

有两种基本方式可以引入那些不增添新语法或判断相等性的公理（函数外延性与泛等性即属此类）：
要么添加一个原始常量来居留该公理，要么通过假设一个居留该公理的变量来证明所有依赖于该公理的定理（参见~\cref{sec:axioms}）。
虽然这两种方式本质上是等价的，但我们选择了前者，因为我们认为同伦类型论的公理是其核心理论不可或缺的一部分。

\index{function extensionality}%
\cref{axiom:funext} 通过引入一个常量 $\funext$ 来形式化，该常量断言 $\happly$ 是一个等价：
%


\begin{mathparpagebreakable}
  \inferrule*[right=$\Pi$-外延]
  {\oftp\Gamma{f}{\tprd{x:A} B} \\
   \oftp\Gamma{g}{\tprd{x:A} B}}
  {\oftp\Gamma{\funext(f,g)}{\isequiv(\happly_{f,g})}}
\end{mathparpagebreakable}


%
$\happly$ 与 $\isequiv$ 的定义可分别在~\eqref{eq:happly} 与~\cref{sec:concluding-remarks} 中找到。

\index{univalence axiom}%
\cref{axiom:univalence} 也以类似的方式形式化：
%


\begin{mathparpagebreakable}
  \inferrule*[right=$\UU_i$-泛等]
  {\oftp\Gamma{A}{\UU_i} \\
   \oftp\Gamma{B}{\UU_i}}
  {\oftp\Gamma{\univalence(A,B)}{\isequiv(\idtoeqv_{A,B})}}
\end{mathparpagebreakable}


%
$\idtoeqv$ 的定义可在~\eqref{eq:uidtoeqv} 中找到。

\subsection{圆类型}

\index{type!circle}%

这里我们给出一个基本的高阶归纳类型的例子；其他类型也遵循相同的一般模式，只是会更为复杂精细。

需要注意，下面的规则并不严格遵循\cref{sec:syntax-more-formally}中普通归纳类型的模式：这些规则涉及传输和映射的函子性（\cref{sec:functors}）的概念，并且第二条计算规则是命题相等，而非判断相等。这些差异在\cref{sec:dependent-paths}中讨论。

\begin{mathparpagebreakable}
  \inferrule*[right=$\Sn^1$-\rform]
  {\wfctx\Gamma}
  {\oftp\Gamma{\Sn^1}{\UU_i}}
\and
  \inferrule*[right=$\Sn^1$-\rintro${}_1$]
  {\wfctx\Gamma}
  {\oftp\Gamma{\base}{\Sn^1}}
\and
  \inferrule*[right=$\Sn^1$-\rintro${}_2$]
  {\wfctx\Gamma}
  {\oftp\Gamma{\lloop}{\id[\Sn^1]{\base}{\base}}}
\and
  \inferrule*[right=$\Sn^1$-\relim]
  {\oftp{\Gamma,\tmtp x{\Sn^1}}{C}{\UU_i} \\
   \oftp{\Gamma}{b}{C[\base/x]} \\
   \oftp{\Gamma}{\ell}{\dpath C \lloop b b} \\
   \oftp\Gamma{p}{\Sn^1}}
  {\oftp\Gamma{\ind{\Sn^1}(x.C,b,\ell,p)}{C[p/x]}}
\and
  \inferrule*[right=$\Sn^1$-\rcomp${}_1$]
  {\oftp{\Gamma,\tmtp x{\Sn^1}}{C}{\UU_i} \\
   \oftp{\Gamma}{b}{C[\base/x]} \\
   \oftp{\Gamma}{\ell}{\dpath C \lloop b b}}
  {\jdeqtp\Gamma{\ind{\Sn^1}(x.C,b,\ell,\base)}{b}{C[\base/x]}}
\and
  \inferrule*[right=$\Sn^1$-\rcomp${}_2$]
  {\oftp{\Gamma,\tmtp x{\Sn^1}}{C}{\UU_i} \\
   \oftp{\Gamma}{b}{C[\base/x]} \\
   \oftp{\Gamma}{\ell}{\dpath C \lloop b b}}
  {\oftp\Gamma{\Sn^1\text{-}\mathsf{loopcomp}}
    {\id {\apd{(\lamu{y:\Sn^1} \ind{\Sn^1}(x.C,b,\ell,y))}{\lloop}} {\ell}}}
\end{mathparpagebreakable}


%
在 $\ind{\Sn^1}$ 中，$x$ 在 $C$ 中被绑定。依值道路的记号 ${\dpath C \lloop b b}$ 是在 \cref{sec:dependent-paths} 中引入的。
\index{rules of type theory|)}%

\section{基本元理论}
\index{metatheory|(}%

本节讨论\cref{sec:syntax-informally}中给出的类型论的元理论性质，类似的结果对\cref{sec:syntax-more-formally}也成立。当我们加入\cref{sec:hott-features}的特性时，要弄清这些性质中哪些仍然成立，很快就会引向一些开放问题\index{open!problem}，正如本节末尾所讨论的那样。

回想一下，\cref{sec:syntax-informally}将类型论的项定义为无类型 $\lambda$-演算的扩展。$\lambda$-演算有其自身的计算概念，即函数类型的计算规则\index{computation rule!for function types}：
\[
  (\lam{x} t)(u) \defeq t[u/x].
\]
这条规则，连同已定义常量的定义方程，构成了\emph{重写规则}\index{rewriting rule}\index{rule!rewriting}，它们决定了一个重写系统的约化步骤。这些步骤之所以能产生一种计算概念，是因为每条规则都有一个天然的方向：我们通过将函数应用于其参数来求值，从而简化 $(\lam{x} t)(u)$。

此外，这个系统是\emph{合流的}\index{confluence}，也就是说，如果 $a$ 经过若干步骤简化到 $a'$ 和 $a''$，那么存在某个 $b$，使得 $a'$ 和 $a''$ 最终都能简化到 $b$。因此，我们可以定义 $t\conv u$ 来表示 $t$ 和 $u$ 能简化到同一个项。

（在\cref{sec:syntax-more-formally}中情况类似：尽管我们在那里将计算规则呈现为无向的等式 $\jdeq$，但我们可以赋予它一种操作语义，即，将消去子应用于引入形式会简化为与之相等的项，而不是反过来。）

运用类型论中的标准技术，可以证明 \cref{sec:syntax-informally} 中的系统具有如下性质：

\begin{thm}\label{thm:conversion-preserves-typing}
如果 $A : \UU$ 且 $A \conv A'$，那么 $A' : \UU$。
如果 $t:A$ 且 $t \conv t'$，那么 $t':A$。
\end{thm}

我们称一个项是\define{可停机的（normalizable）}
\indexdef{term!normalizable}%
\index{normalization}%
\indexdef{normalizable term}%
（相应地，\define{强可停机的（strongly
normalizable）}）
\indexdef{term!strongly normalizable}%
\index{normalization!strong}%
\index{strong!normalization}%
，如果从它出发的某个（相应地，每一个）重写序列会终止。

\begin{thm}\label{thm:strong-normalization}
如果 $A : \UU$，那么 $A$ 是强可停机的。
如果 $t:A$，那么 $A$ 和 $t$ 都是强可停机的。
\end{thm}

我们称一个项处于\define{既约形式（normal form）}
\index{normal form}%
\index{term!normal form of}%
，如果它不能再被简化；称一个项是\define{封闭的（closed）}
\index{closed!term}%
\index{term!closed}%
，如果没有自由变量在其中出现。
一个封闭的既约类型必定是一个原始类型，即形如 $c(\vec{v})$，其中 $c$ 是某个原始常量（封闭既约项的列表 $\vec{v}$ 若为空则可省略，例如 $\N$）。
事实上，我们可以显式地描述所有既约形式：

\begin{lem}\label{lem:normal-forms}
  既约项可以通过以下语法描述：
  %
  \begin{align*}
    v & \production  k \mid \lam{x} v \mid c(\vec{v}) \mid f(\vec{v}), \\
    k &\production x \mid k(v) \mid f(\vec{v})(k),
  \end{align*}
  %
  其中 $f(\vec{v})$ 代表已定义函数 $f$ 的部分应用。
  特别地，一个既约类型具有形式 $k$ 或 $c(\vec{v})$。
\end{lem}

\begin{thm}
  如果 $A$ 是既约的，那么判断 $A : \UU$ 是可判定的。如果 $A : \UU$ 且 $t$ 是既约的，那么判断 $t:A$ 是可判定的。
\end{thm}

\cref{sec:syntax-informally} 中系统的逻辑一致性\index{consistency}可立即得出：如果我们在空上下文中有一个 $a:\emptyt$，那么根据 \cref{thm:conversion-preserves-typing,thm:strong-normalization}，$a$ 可简化为一个既约项 $a':\emptyt$。但根据 \cref{lem:normal-forms}，这样的项并不存在。

\begin{cor}
 \cref{sec:syntax-informally} 中的系统是逻辑一致的。
\end{cor}

类似地，我们有\emph{典范性（canonicity）}\indexdef{canonicity}性质：如果空上下文中有一个 $a:\N$，那么 $a$ 可简化为一个既约项 $\suc^k(0)$，其中 $k$ 是某个数字。

\begin{cor}
 \cref{sec:syntax-informally} 中的系统具有典范性。
\end{cor}

最后，如果 $a,A$ 都是既约的，那么 $a:A$ 是否成立是\emph{可判定的}；换句话说，因为类型检查相当于验证一个证明的正确性，这意味着``看到一个正确的证明时，我们总能将其识别出来''。

\begin{cor}
 在 \cref{sec:syntax-informally} 的系统中，一个项是否构成证明是可判定的。
\end{cor}

\mentalpause

上述结果并不适用于同伦类型论的扩展系统（即由 \cref{sec:hott-features} 扩展后的上述系统），因为泛等公理和高阶归纳类型构造子的出现永远不会被简化，从而破坏了 \cref{lem:normal-forms}。是否存在对这些常量的应用进行简化的方法以恢复典范性，这是一个开放问题\index{open!problem}。我们也没有一个能描述所有合法高阶归纳类型的模式，且不确定如何正确地形式化它们的规则（例如，高阶构造子上的计算规则是否应该是判断相等）。

通过设计适当的归约过程，或许能证明添加了泛等性和高阶归纳类型的 Martin-L\"{o}f 类型论的一致性，但目前证明这些系统一致性的唯一途径是通过语义模型——对泛等性，有 Voevodsky \cite{klv:ssetmodel} 在 Kan\index{Kan complex} 复形中给出的模型；对高阶归纳类型，有 Lumsdaine 和 Shulman \cite{ls:hits} 给出的模型。

其他的元理论问题，以及我们当前成果的总结，在本书导言的``构造性''与``待解问题''两节中有更详细的讨论。

\index{metatheory|)}%

\sectionNotes\label{subsec:general-remarks}

% This presentation is strongly inspired by two  Martin-L\"of 1972 and 1973.

这套包含引入规则（原始常量）、消去规则与计算规则（已定义常量）的规则系统，其灵感来源于 Gentzen 的自然演绎。强化存在量词消去规则的可能性在 \cite{howard:pat} 中有所提示。析取的公理的强化出现在 \cite{Martin-Lof-1972} 中，荒谬消去和相等类型的强化则在 \cite{Martin-Lof-1973} 中。$W$-类型在 \cite{Martin-Lof-1979} 中被引入。它们推广了由 \cite{Tait-1968} 引入的一种树的概念。
\index{Martin-L\"of}%

%inspired from unpublished work of Spector.

自然数和序数的原始递归（primitive recursion）的推广形式见 \cite{Hilbert-1925}。这促成了 Gödel 的系统 $T$（\cite{Goedel-T-1958}）。\cite{Tait-1966} 分析了该系统，并遵循 \cite{Goedel-T-1958} 的做法，用``定义相等（definitional equality）''这一术语来指称转换（conversion）：如果两个项能通过反复应用归约规则（reduction rules）而归约到同一项，则称它们是\emph{判断相等（judgmentally equal）}的。de Bruijn \cite{deBruijn-1973} 在 \emph{AUTOMATH} 的表述中也使用了这一术语。\index{AUTOMATH}

本书的第二种表述包含了一种相当标准的内涵（intensional）Martin-L\"{o}f 类型论，并附带同伦类型论所需的一些额外特性。与 \cite{hofmann:syntax-and-semantics} 的参考表述相比，本书的类型论有几点非关键差异：
%

\begin{itemize}
\item \cite{martin-lof:bibliopolis} 意义下的 \`{a} la Russell（罗素）风格宇宙；以及
\item $\Pi$ 类型上的判断性 $\eta$ 规则与函数外延性（function extensionality）；
\end{itemize}


以及几点对同伦类型论至关重要的特性：


\begin{itemize}
\item 泛等公理（univalence axiom）；以及
\item 高阶归纳类型（higher inductive types）。
\end{itemize}


%
为方便起见，本书定义函数时主要采用基于\emph{模式匹配（pattern matching）}的归纳定义。
\index{pattern matching}%
\index{definition!by pattern matching}%
模式匹配这一概念是可以形式化的，\cref{sec:syntax-informally} 即是一例。然而，在 \cref{sec:syntax-more-formally} 采用的标准类型论表述中，为每个类型构造子（type former）仅引入一个\emph{依值消去子（dependent eliminator）}，所有从该类型出发的函数都必须通过这个消去子来定义。这种方法无论在句法上还是语义上都更容易形式化，因为它归结为类型构造子的泛性质（universal property）。这两种方法是等价的；更详细的讨论见 \cref{sec:pattern-matching}。

\index{type theory!formal|)}%
\index{formal!type theory|)}%

%%% Local Variables: 
%%% mode: latex
%%% TeX-master: "hott-online"
%%% End: 

% Joke
\nocite{Angiuli13}
\nocite{BauerAcceptanceVideo}

%%%% Bibliography
\bibliographystyle{halpha}
\phantomsection % black magic to get TOC to point to correct page
\addcontentsline{toc}{part}{\bibname}
\markboth{}{\textsc{Bibliography}}
{\renewcommand{\markboth}[2]{} % Prevent bibliography from resetting the header to something silly
  \OPTbibliographyfont
\bibliography{references}}

\cleartooddpage[\thispagestyle{empty}]

%%%% Index of symbols

\phantomsection % black magic to get TOC to point to correct page
\markboth{}{\textsc{符号表}}
\addcontentsline{toc}{part}{符号表}
\chapter*{符号表}

% Shorthand for \pageref, we have lots of these.
\newcommand{\pg}[1]{p.~\pageref{#1}}
% In the next macro whitespace matters, so be careful
\newcommand{\symbolindex}[2]{\hbox{\makebox[0.2\textwidth][s]{#1\hfill}\hspace*{2.5em}\parbox[t]{0.65\textwidth}{#2\hfill}}\\[1.5pt]}

% The entries in this table are sorted "alphabetically" whatever that means.

{\OPTindexfont % Set same font size as for the index

% MAKE SURE THERE ARE NO EMPTY LINES, OR ELSE WE GET A PARAGRAPH BREAK.
% ALSO, DO NOT INSERT ANY WHITESPACE AT THE BEGINNING OF ARGUMENTS OF \symbolindex
\noindent
%%%%%%%%%% Equalities %%%%%%%%%%
\symbolindex{$x \defeq a$	}{定义，\pg{defn:defeq}}
\symbolindex{$a \jdeq b$	}{判断相等，\pg{defn:judgmental-equality}}
\symbolindex{$a =_A b$	}{相等类型，\pg{sec:identity-types}}
\symbolindex{$a = b$	}{相等类型，\pg{sec:identity-types}}
\symbolindex{$x \defid b$	}{按定义成立的命题相等，\pg{rmk:defid}}
\symbolindex{$\idtypevar{A}(a,b)$	}{相等类型，\pg{sec:identity-types}}
\symbolindex{$\dpath{P}{p}{a}{b}$	}{依值道路类型，\pg{eq:dpath}}
\symbolindex{$a \neq b$	}{不等，\pg{sec:disequality}}
%%%%%%%%%% Stuff that's hard to alphabetize %%%%%%%%%%
\symbolindex{$\refl{x}$	}{$x$ 处的自反道路，\pg{sec:identity-types}}
\symbolindex{$\opp{p}$	}{道路的逆，\pg{lem:opp}}
\symbolindex{$p\ct q$	}{道路拼接，\pg{lem:concat}}
\symbolindex{$p\leftwhisker r$	}{左须接，\pg{thm:EckmannHilton}}
\symbolindex{$r\rightwhisker q$	}{右须接，\pg{thm:EckmannHilton}}
\symbolindex{$r\hct s$	}{2-道路的水平复合，\pg{thm:EckmannHilton}}
\symbolindex{$g\circ f$	}{函数的复合，\pg{ex:composition}}
\symbolindex{$g\circ f$	}{预范畴中态射的复合，\pg{ct:precategory}}
\symbolindex{$\inv{f}$	}{等价的拟逆，\pg{thm:equiv-eqrel}}
\symbolindex{$\inv{f}$	}{预范畴中同构的逆，\pg{ct:inv}}
%%%%%%%%%% Numbers etc. %%%%%%%%%%
\symbolindex{$\emptyt$	}{空类型，\pg{sec:coproduct-types}}
\symbolindex{$\unit$	}{单位类型，\pg{sec:finite-product-types}}
\symbolindex{$\ttt$	}{$\unit$ 的典范元素，\pg{sec:finite-product-types}}
\symbolindex{$\bool$	}{布尔类型，\pg{sec:type-booleans}}
\symbolindex{$\btrue$, $\bfalse$	}{$\bool$ 的构造子，\pg{sec:type-booleans}}
\symbolindex{$\izero$, $\ione$	}{区间 $\interval$ 的点构造子，\pg{sec:interval}}
%%%%%%%%%% A %%%%%%%%%%
\symbolindex{$\choice{}$	}{选择公理，\pg{eq:ac}}
\symbolindex{$\choice{\infty}$	}{``类型论选择公理''，\pg{thm:ttac}}
\symbolindex{$\acc(a)$	}{可达性谓词，\pg{defn:accessibility}}
\symbolindex{$P \land Q$	}{逻辑合取（``且''），\pg{defn:logical-notation}}
\symbolindex{$\apfunc{f}(p)$ or $\ap{f}{p}$	}{将 $f:A\to B$ 作用于 $p:\id[A]xy$，\pg{lem:map}}
\symbolindex{$\apd{f}{p}$	}{将 $f:\prd{a:A} B(a)$ 作用于 $p:\id[A]xy$，\pg{lem:mapdep}}
% This is not currently a distinctive notation that someone would look up
% $\aptwo{f}{p}$ & two-dimensional $\apfunc{}$, \pg{thm:ap2}
% \\
\symbolindex{$\apdtwo{f}{p}$	}{二维依值 $\apfunc{}$，\pg{thm:apd2}}
\symbolindex{$x\apart y$	}{实数的分离关系，\pg{apart}}
%%%%%%%%%% B %%%%%%%%%%
\symbolindex{$\base$	}{$\Sn^1$ 的基点，\pg{sec:intro-hits}}
\symbolindex{$\base$            }{$\Sn^2$ 的基点，\pg{s2a} 与 \pg{s2b}}
\symbolindex{$\biinv(f)$        }{``$f$ 是双可逆的''这一命题，\pg{defn:biinv}}
\symbolindex{$x \bisim y$       }{双模拟，\pg{def:bisimulation}}
\symbolindex{$\blank$           }{用于隐式 $\lambda$-抽象的空白记号，\pg{blank}}
%%%%%%%%%% C %%%%%%%%%%
\symbolindex{$\CAP$             }{Cauchy 逼近的类型，\pg{cauchy-approximations}}
\symbolindex{$\card$            }{基数的类型，\pg{defn:card}}
\symbolindex{$\modal A$         }{$A$ 上的反射子或模态，\pg{defn:reflective-subuniverse} 与 \pg{defn:modality}}
\symbolindex{$\cocone{X}{Y}$    }{上锥的类型，\pg{defn:cocone}}
% Not used
%\symbolindex{
%  \cone{X}{Y}$
%}{
%  type of cones
%}
\symbolindex{$\code$            }{道路的编码族，\pg{sec:compute-coprod}、\pg{S1-universal-cover}、\pg{sec:general-encode-decode}}
\symbolindex{$A \setminus B$    }{子集的补，\pg{complement}}
\symbolindex{$\cons(x,\ell)$    }{列表的拼接构造子，\pg{lst} 与 \pg{lst-freemonoid}}
\symbolindex{$\contr_x$         }{到达收缩中心的道路，\pg{defn:contractible}}
\symbolindex{$\mathcal{F}\cover\pairr{J, \mathcal{G}}$          }{归纳覆盖，\pg{defn:inductive-cover}}
\symbolindex{$\dcut(L,U)$       }{``是 Dedekind 分割''这一性质，\pg{defn:dedekind-reals}}
\symbolindex{$\surr{L}{R}$      }{定义超现实数的分割，\pg{surreal-cut}}
%%%%%%%%%% D %%%%%%%%%%
\symbolindex{$\dgr{X}$          }{$\dagger$-范畴中的态射反转，\pg{sec:dagger-categories}}
\symbolindex{$\decode$          }{道路的译码函数，\pg{sec:compute-coprod}、\pg{S1-universal-cover}、\pg{sec:general-encode-decode}}
%%%%%%%%%% E %%%%%%%%%%
\symbolindex{$\encode$	}{道路的编码函数，\pg{sec:compute-coprod}、\pg{S1-universal-cover}、\pg{sec:general-encode-decode}}
\symbolindex{$\mreturn^\modal_A$ or $\mreturn_A$	}{函数 $A\to\modal A$，\pg{defn:reflective-subuniverse} 与 \pg{defn:modality}}
\symbolindex{$A \epi B$	}{满态射或满射}
\symbolindex{$\noeq(x,y)$	}{超现实数的道路构造子，\pg{defn:surreals}}
\symbolindex{$\rceq(u,v)$	}{Cauchy 实数的道路构造子，\pg{defn:cauchy-reals}}
\symbolindex{$a \eqr b$	}{等价关系，\pg{equivalencerelation}}
\symbolindex{$\eqv{X}{Y}$	}{等价的类型，\pg{eq:eqv}}
\symbolindex{$\texteqv{X}{Y}$	}{等价的类型（与 $\eqv{X}{Y}$ 相同）}
\symbolindex{$\cteqv{A}{B}$	}{范畴之间等价的类型，\pg{ct:equiv}}
\symbolindex{$P ⇔ Q$	}{逻辑等价，\pg{defn:logical-notation}}
\symbolindex{$\exis{x:A} B(x)$	}{存在命题的逻辑记号，\pg{defn:logical-notation}}
\symbolindex{$\extend f$	}{沿 $\eta_A$ 对 $f:A\to B$ 的延拓，\pg{extend}}
%%%%%%%%%% F %%%%%%%%%%
\symbolindex{$\bot$	}{逻辑假，\pg{defn:logical-notation}}
\symbolindex{$\hfib{f}{b}$	}{$f:A\to B$ 在 $b:B$ 处的纤维，\pg{defn:homotopy-fiber}}
\symbolindex{$\Fin(n)$  }{标准有限类型，\pg{fin}}
\symbolindex{$\fall{x:A} B(x)$	}{依值函数类型的逻辑记号，\pg{defn:logical-notation}}
\symbolindex{$\funext$	}{函数外延性，\pg{axiom:funext}}
\symbolindex{$A\to B$	}{函数类型，\pg{sec:function-types}}
\symbolindex{$B^A$      }{函子预范畴，\pg{ct:functor-precat}}
%%%%%%%%%% G %%%%%%%%%%
\symbolindex{$\glue$	}{$A \sqcup^C B$ 的道路构造子，\pg{sec:colimits}}
%%%%%%%%%% H %%%%%%%%%%
\symbolindex{$\happly$	}{把函数间的道路变为同伦的函数，\pg{eq:happly}}
\symbolindex{$\hom_A(a,b)$	}{预范畴中的态射集，\pg{ct:precategory}}
\symbolindex{$f \htpy g$	}{函数间的同伦，\pg{defn:homotopy}}
%%%%%%%%%% I %%%%%%%%%%
\symbolindex{$\interval$	}{区间类型，\pg{sec:interval}}
\symbolindex{$\idfunc[A]$	}{$A$ 的恒等函数，\pg{idfunc}}
\symbolindex{$1_a$	}{预范畴中的恒等态射，\pg{ct:precategory}}
\symbolindex{$\idtoeqv$	}{泛等性所要取逆的函数 $(A=B)\to(\eqv A B)$，\pg{eq:uidtoeqv}}
\symbolindex{$\idtoiso$	}{预范畴中的函数 $(a=b) \to (a\cong b)$，\pg{ct:idtoiso}}
\symbolindex{$\im(f)$	}{映射 $f$ 的像，\pg{defn:modal-image}}
\symbolindex{$\im_n(f)$	}{映射 $f$ 的 $n$-像，\pg{defn:modal-image}}
\symbolindex{$P \Rightarrow Q$	}{逻辑蕴含（``蕴含''），\pg{defn:logical-notation}}
\symbolindex{$a \in P$	}{属于某子集或子类型，\pg{membership}}
\symbolindex{$x\in v$	}{属于累积层次，\pg{V-membership}}
\symbolindex{$x\bin v$	}{缩小规模后的属于关系，\pg{resized-membership}}
\symbolindex{$\ind{\emptyt}$	}{$\emptyt$ 的归纳，\pg{defn:induction-emptyt},}
\symbolindex{$\ind{\unit}$	}{$\unit$ 的归纳，\pg{defn:induction-unit},}
\symbolindex{$\ind{\bool}$	}{$\bool$ 的归纳，\pg{defn:induction-bool},}
\symbolindex{$\ind{\nat}$	}{$\nat$ 的归纳，\pg{defn:induction-nat}，及}
\symbolindex{$\indid{A}$	}{$=_A$ 的道路归纳，\pg{defn:induction-ML-id},}
\symbolindex{$\indidb{A}$	}{$=_A$ 的基于基点的道路归纳，\pg{defn:induction-PM-id},}
\symbolindex{$\ind{A \times B}$	}{$A \times B$ 的归纳，\pg{defn:induction-times},}
\symbolindex{$\ind{\sm{x:A} B(x)}$	}{$\sm{x:A} B$ 的归纳，\pg{defn:induction-sm},}
\symbolindex{$\ind{A + B}$	}{$A + B$ 的归纳，\pg{defn:induction-plus},}
\symbolindex{$\ind{\wtype{x:A} B(x)}$	}{$\wtype{x:A} B$ 的归纳，\pg{defn:induction-wtype}}
\symbolindex{$\ordsl{A}{a}$	}{序数的初始段，\pg{initial-segment}}
\symbolindex{$\inj(A,B)$	}{单射的类型，\pg{inj}}
\symbolindex{$\inl$	}{余积的第一个内射，\pg{sec:coproduct-types}}
\symbolindex{$\inr$	}{余积的第二个内射，\pg{sec:coproduct-types}}
\symbolindex{$A \cap B$	}{子集之交，\pg{intersection}，类之交，\pg{class-intersection}，或区间之交，\pg{interval-intersection}}
\symbolindex{$\iscontr(A)$	}{``$A$ 是可缩的''这一命题，\pg{defn:contractible}}
\symbolindex{$\isequiv(f)$	}{``$f$ 是一个等价''这一命题，\pg{basics-isequiv}、\pg{cha:equivalences} 与 \pg{sec:concluding-remarks}}
\symbolindex{$\ishae(f)$	}{``$f$ 是半伴随等价''这一命题，\pg{defn:ishae}}
\symbolindex{$a\cong b$	}{（预）范畴中同构的类型，\pg{ct:isomorphism}}
\symbolindex{$A\cong B$	}{预范畴之间同构的类型，\pg{ct:isocat}}
\symbolindex{$A\cong B$	}{集合之间同构的类型，\pg{basics:iso}}
\symbolindex{$a\unitaryiso b$	}{酉同构的类型，\pg{ct:unitary}}
\symbolindex{$\isotoid$	}{范畴中 $\idtoiso$ 的逆，\pg{isotoid}}
\symbolindex{$\istype{n}(X)$	}{``$X$ 是一个 $n$-类型''这一命题，\pg{def:hlevel}}
\symbolindex{$\isprop(A)$	}{``$A$ 是一个命题''这一命题，\pg{defn:isprop}}
\symbolindex{$\isset(A)$	}{``$A$ 是一个集合''这一命题，\pg{defn:set}}
%%%%%%%%%% J %%%%%%%%%%
\symbolindex{$A*B$	}{$A$ 与 $B$ 的联结，\pg{join}}
%%%%%%%%%% K %%%%%%%%%%
\symbolindex{$\ker(f)$	}{带点集合间映射的核，\pg{kernel}}
%%%%%%%%%% L %%%%%%%%%%
\symbolindex{$\lam{x} b(x)$	}{$\lambda$-抽象，\pg{eq:lambda-abstraction}}
\symbolindex{$\lcoh{f}{g}{\eta}$	}{左伴随相干数据的类型，\pg{defn:lcoh-rcoh}}
\symbolindex{$\LEM{}$	}{排中律，\pg{eq:lem}}
\symbolindex{$\LEM{\infty}$	}{不相容的命题即类型式排中律，\pg{thm:not-lem} 与 \pg{lem-infty}}
\symbolindex{$x < y$	}{自然数、\pg{leq-nat}，序数、\pg{sec:ordinals}，Cauchy 实数、\pg{lt-RC}，超现实数、\pg{defn:surreals} 等上的严格不等式}
\symbolindex{$x ≤ y$	}{自然数、\pg{leq-nat}，Cauchy 实数、\pg{leq-RC}，超现实数、\pg{defn:surreals} 等上的非严格不等式}
\symbolindex{$⪯$, $≺$	}{超现实数上 $\le$ 与 $<$ 的递归版本，\pg{defn:No-codes}}
\symbolindex{$\ble$, $\blt$, $⊑$, $\bblt$	}{$\NO$-递归陪域上的序，\pg{NO-recursion}}
\symbolindex{$\rclim(x)$	}{Cauchy 逼近的极限，\pg{defn:cauchy-reals}}
\symbolindex{$\linv(f)$	}{$f$ 的左逆的类型，\pg{defn:linv-rinv}}
\symbolindex{$\lst{X}$	}{$X$ 的元素构成的列表的类型，\pg{lst} 与 \pg{lst-freemonoid}}
\symbolindex{$\lloop$	}{$\Sn^1$ 的道路构造子，\pg{sec:intro-hits}}
%%%%%%%%%% M %%%%%%%%%%
\symbolindex{$\Map_*(A,B)$	}{带点映射的类型，\pg{based-maps}}
\symbolindex{$x ↦ b$	}{$\lambda$-抽象的另一种记法，\pg{mapsto}}
\symbolindex{$\max(x,y)$	}{某序下的最大值，例如 \pg{ordered-field} 与 \pg{leq-RC}}
\symbolindex{$\merid(a)$	}{$\susp A$ 在 $a:A$ 处的经线，\pg{sec:suspension}}
\symbolindex{$\min(x,y)$	}{某序下的最小值，例如 \pg{ordered-field} 与 \pg{leq-RC}}
\symbolindex{$A \mono B$	}{单态射或嵌入}
%%%%%%%%%% N %%%%%%%%%%
\symbolindex{$\N$	}{自然数类型，\pg{sec:inductive-types}}
\symbolindex{$\north$	}{$\susp A$ 的北极，\pg{sec:suspension}}
\symbolindex{$\natw$, $\zerow$, $\sucw$	}{编码为 $W$-类型的自然数，\pg{natw}}
\symbolindex{$\nalg$	}{$\nat$-代数的类型，\pg{defn:nalg}}
\symbolindex{$\nhom(C,D)$	}{$\nat$-同态的类型，\pg{defn:nhom}}
\symbolindex{$\nil$	}{空列表，\pg{lst} 与 \pg{lst-freemonoid}}
\symbolindex{$\NO$	}{超现实数的类型，\pg{defn:surreals}}
\symbolindex{$\neg P$	}{逻辑否定（``非''），\pg{defn:logical-notation}}
\symbolindex{$\typele{n}$, $\typeleU{n}$	}{$n$-类型的宇宙，\pg{universe-of-ntypes}}
%%%%%%%%%% O %%%%%%%%%%
\symbolindex{$\Omega(A,a)$, $\Omega A$		}{带点类型的环路空间，\pg{def:loopspace}}
\symbolindex{$\Omega^k(A,a)$, $\Omega^k A$		}{迭代环路空间，\pg{def:loopspace}}
\symbolindex{$A\op$	}{反范畴，\pg{ct:opposite-category}}
\symbolindex{$P \lor Q$	}{逻辑析取（``或''），\pg{defn:logical-notation}}
\symbolindex{$\ord$	}{序数的类型，\pg{ord}}
%%%%%%%%%% P %%%%%%%%%%
\symbolindex{$\pairr{a,b}$	}{（依值）对，\pg{sec:finite-product-types} 与 \pg{defn:dependent-pair}}
\symbolindex{$\pairpath$	}{$=_{A \times B}$ 的构造子，\pg{defn:pairpath}}
\symbolindex{$\pi_n(A)$	}{$A$ 的第 $n$ 个同伦群，\pg{thm:homotopy-groups} 与 \pg{def-of-homotopy-groups}}
\symbolindex{$\power A$	}{幂集，\pg{powerset}}
\symbolindex{$\powerp A$	}{有实例的幂集，\pg{inhabited-powerset}}
\symbolindex{$\Zpred$	}{前驱函数 $\Z\to\Z$，\pg{subsec:pi1s1-encode-decode}}
\symbolindex{$A\times B$	}{Cartesian 积类型，\pg{sec:finite-product-types}}
\symbolindex{$\prd{x:A} B(x)$	}{依值函数类型，\pg{sec:pi-types}}
\symbolindex{$\proj1(t)$	}{对的第一个投影，\pg{defn:proj} 与 \pg{defn:dependent-proj1}}
\symbolindex{$\proj2(t)$	}{对的第二个投影，\pg{defn:proj} 与 \pg{defn:dependent-proj1}}
\symbolindex{$\prop$, $\propU$	}{命题的宇宙，\pg{propU}}
\symbolindex{$A \times_C B$	}{$A$ 与 $B$ 在 $C$ 上的拉回，\pg{eq:defn-pullback}}
\symbolindex{$A \sqcup^C B$	}{$A$ 与 $B$ 在 $C$ 下的推出，\pg{sec:colimits}}
%%%%%%%%%% Q %%%%%%%%%%
\symbolindex{$\Q$	}{有理数的类型，\pg{sec:field-rati-numb}}
\symbolindex{$\Qp$	}{正有理数的类型，\pg{positive-rationals}}
\symbolindex{$\qinv(f)$	}{$f$ 的拟逆的类型，\pg{qinv}}
\symbolindex{$A/R$	}{集合按等价关系取的商，\pg{sec:set-quotients}}
\symbolindex{$A\sslash R$	}{商的另一种定义，\pg{def:VVquotient}}
%%%%%%%%%% R %%%%%%%%%%
\symbolindex{$\R$	}{实数的类型（两种之一），\pg{sec:compactness-interval}}
\symbolindex{$\RC$	}{Cauchy 实数的类型，\pg{defn:cauchy-reals}}
\symbolindex{$\RD$	}{Dedekind 实数的类型，\pg{defn:dedekind-reals}}
\symbolindex{$\rcrat(q)$	}{视为 Cauchy 实数的有理数，\pg{defn:cauchy-reals}}
\symbolindex{$\rcoh{f}{g}{\epsilon}$	}{右伴随相干数据的类型，\pg{defn:lcoh-rcoh}}
\symbolindex{$\rec{\emptyt}$	}{$\emptyt$ 的递归子，\pg{defn:recursor-emptyt}}
\symbolindex{$\rec{\unit}$	}{$\unit$ 的递归子，\pg{defn:recursor-unit}}
\symbolindex{$\rec{\bool}$	}{$\bool$ 的递归子，\pg{defn:recursor-bool}}
\symbolindex{$\rec{\nat}$	}{$\nat$ 的递归子，\pg{defn:recursor-nat}}
\symbolindex{$\rec{A \times B}$	}{$A \times B$ 的递归子，\pg{defn:recursor-times}}
\symbolindex{$\rec{\sm{x:A} B(x)}$	}{$\sm{x:A} B$ 的递归子，\pg{defn:recursor-sm}}
\symbolindex{$\rec{A + B}$	}{$A + B$ 的递归子，\pg{defn:recursor-plus}}
\symbolindex{$\rec{\wtype{x:A} B(x)}$	}{$\wtype{x:A} B$ 的递归子，\pg{defn:recursor-wtype}}
\symbolindex{$\rinv$	}{$f$ 的右逆的类型，\pg{defn:linv-rinv}}
%%%%%%%%%% S %%%%%%%%%%
\symbolindex{$\south$	}{$\susp A$ 的南极，\pg{sec:suspension}}
\symbolindex{$\Sn^n$	}{$n$-维球面，\pg{sec:circle}}
\symbolindex{$\seg$	}{区间 $\interval$ 的道路构造子，\pg{sec:interval}}
\symbolindex{$\set$, $\setU$	}{集合的宇宙，\pg{setU}}
\symbolindex{$\uset$	}{集合的范畴，\pg{ct:precatset}}
\symbolindex{$\vset(A,f)$	}{累积层次的构造子，\pg{defn:V}}
\symbolindex{$x\sim_\epsilon y$	}{$\RC$ 上的 $\epsilon$-接近关系，\pg{defn:cauchy-reals}}
\symbolindex{$x\approx_\epsilon y$	}{$\sim_\epsilon$ 的递归版本，\pg{defn:RC-approx}}
\symbolindex{$\bsim_\epsilon$ or $\bbsim_\epsilon$	}{$\RC$-递归陪域上的接近关系，\pg{RC-recursion}}
\symbolindex{$A\wedge B$	}{$A$ 与 $B$ 的缩积，\pg{smash}}
\symbolindex{$\setof{x : A | P(x)}$	}{子集类型，\pg{defn:setof}}
\symbolindex{$\setof{ f(x) | P(x)}$	}{子集的像，\pg{subset-image}}
\symbolindex{$B \subseteq C$	}{子集类型间的包含，\pg{subset}}
\symbolindex{$(q,r)\subseteq (s,t)$	}{区间的包含，\pg{interval-subset}}
\symbolindex{$\suc$	}{后继函数 $\N\to\N$，\pg{sec:inductive-types}}
\symbolindex{$\Zsuc$	}{后继函数 $\Z\to\Z$，\pg{sec:pi1s1-initial-thoughts}}
\symbolindex{$A+B$	}{余积类型，\pg{sec:coproduct-types}}
\symbolindex{$\sm{x:A} B(x)$	}{依值对类型，\pg{sec:sigma-types}}
\symbolindex{$\supp(a, f)$	}{$W$-类型的构造子，\pg{defn:supp}}
\symbolindex{$\surf$	}{$\Sn^2$ 的 2-道路构造子，\pg{s2a} 与 \pg{s2b}}
\symbolindex{$\susp A$	}{$A$ 的纬悬，\pg{sec:suspension}}
%%%%%%%%%% T %%%%%%%%%%
\symbolindex{$\total{f}$	}{全空间上的诱导映射，\pg{defn:total-map}}
\symbolindex{$\trans{p}{u}$	}{$u:P(x)$ 沿 $p:x=y$ 的迁移，\pg{lem:transport}}
\symbolindex{$\transfib{P}{p}{u}$	}{$u:P(x)$ 沿 $p:x=y$ 的迁移，\pg{lem:transport}}
\symbolindex{$\transtwo{X}{Y}$	}{二维迁移，\pg{thm:transport2}}
\symbolindex{$\transconst{X}{Y}{Z}$	}{常值族中的迁移，\pg{thm:trans-trivial}}
\symbolindex{$\trunc{n}{A}$	}{$A$ 的 $n$-截断，\pg{sec:truncations}}
\symbolindex{$\tproj[A]{n}{a}$, $\tproj{n}{a}$	}{$a:A$ 在 $\trunc n A$ 中的像，\pg{sec:truncations}}
\symbolindex{$\brck{A}$	}{$A$ 的命题截断，\pg{subsec:prop-trunc} 与 \pg{sec:hittruncations}}
\symbolindex{$\bproj{a}$	}{$a:A$ 在 $\brck A$ 中的像，\pg{subsec:prop-trunc} 与 \pg{sec:hittruncations}}
\symbolindex{$\top$	}{逻辑真，\pg{defn:logical-notation}}
%%%%%%%%%% U %%%%%%%%%%
\symbolindex{$\nameless$	}{未命名的对象或变量}
\symbolindex{$A \cup B$	}{子集之并，\pg{union}}
\symbolindex{$\uniq{A\times B}$ }{积 $A\times B$ 的唯一性原理，\pg{uniquenessproduct}}
\symbolindex{$\uniq{\unit}$     }{$\unit$ 的唯一性原理，\pg{uniquenessunit}}
\symbolindex{$\UU$	}{宇宙类型，\pg{sec:universes}}
\symbolindex{$\modaltype$	}{模态类型的宇宙，\pg{eq:modaltype}}
\symbolindex{$\pointed\type$	}{带点类型的宇宙，\pg{def:pointedtype}}
\symbolindex{$\ua$	}{泛等性给出的 $\idtoeqv$ 的逆，\pg{ua}}
%%%%%%%%%% V %%%%%%%%%%
\symbolindex{$V$	}{累积层次，\pg{defn:V}}
%%%%%%%%%% W %%%%%%%%%%
\symbolindex{$\walg(A,B)$	}{$w$-代数的类型，\pg{walg}}
\symbolindex{$\whom_{A,B}(C,D)$	}{$\w$-同态的类型，\pg{whom}}
\symbolindex{$\wtype{x:A} B(x)$	}{$W$-类型（归纳类型），\pg{sec:w-types}}
\symbolindex{$A\vee B$	}{$A$ 与 $B$ 的楔，\pg{wedge}}
%%%%%%%%%% X %%%%%%%%%%
%%%%%%%%%% Y %%%%%%%%%%
\symbolindex{$\y$	}{米田嵌入，\pg{ct:yoneda}}
%%%%%%%%%% Z %%%%%%%%%%
\symbolindex{$\Z$	}{整数的类型，\pg{defn-Z}}}

%%% Local Variables:
%%% mode: latex
%%% TeX-master: "hott-online"
%%% End:


\cleartooddpage[\thispagestyle{empty}]

%%%%% Index of terms

%% Global cross-references for index
\indexsee{principle}{axiom}
\indexsee{number!real}{real numbers}
\indexsee{abelian group}{group, abelian}
\indexsee{sequence!Cauchy}{Cauchy sequence}
\indexsee{adjunction}{adjoint functor}
\indexsee{higher topos}{$(\infty,1)$-topos}
\indexsee{topos!higher}{$(\infty,1)$-topos}
\indexsee{source!of a function}{domain}
\indexsee{target!of a function}{codomain}
\indexsee{type!truncation of}{truncation}
\indexsee{propositional!truncation}{truncation}
\indexsee{proof-relevant mathematics}{mathematics, proof-relevant}
\indexsee{classical!mathematics}{mathematics, classical}
\indexsee{classical!logic}{logic}
\indexsee{constructive!mathematics}{mathematics, constructive}
\indexsee{constructive!logic}{logic}
\indexsee{intuitionistic logic}{logic}
\indexsee{definition!inductive}{type, inductive}
\indexsee{inductive!definition}{type, inductive}
\indexsee{bounded!totally}{totally bounded}
\indexsee{sum!of numbers}{addition}
\indexsee{continuous map}{function, continuous}
\indexsee{function!continuity of@``continuity'' of}{``continuity''}
\indexsee{function!functoriality of@``functoriality'' of}{``functoriality''}
\indexsee{codes}{encode-decode method}
\indexsee{inequality}{order}
\indexsee{Coq@\Coq}{proof assistant}
\indexsee{Agda@\Agda}{proof assistant}
\indexsee{NuPRL@\NuPRL}{proof assistant}
\indexsee{generator!of an inductive type}{constructor}
\indexsee{groupoid!.infinity-@$\infty$-}{$\infty$-groupoid}
\indexsee{higher groupoid}{$\infty$-groupoid}
\indexsee{hierarchy!of n-types@of $n$-types}{$n$-type}
\indexsee{homotopy!n-type@$n$-type}{$n$-type}
\indexsee{homotopy!theory, classical}{classical homotopy theory}
\indexsee{homotopy!fiber}{fiber}
\indexsee{homotopy!limit}{limit of types}
\indexsee{homotopy!colimit}{colimit of types}
\indexsee{implementation}{proof assistant}
\indexsee{notation, abuse of}{abuse of notation}
\indexsee{language, abuse of}{abuse of language}
\indexsee{operator!induction}{induction principle}
\indexsee{operator!modal}{modality}
\indexsee{commutative!group}{group, abelian}
\indexsee{countable axiom of choice}{axiom of choice, countable}

% tell the index to get itself into the table of contents
\phantomsection % black magic to get TOC to point to correct page
\addcontentsline{toc}{part}{索引}
\markboth{}{\textsc{索引}}
\renewcommand{\markboth}[2]{}
{\OPTindexfont
  \setlength{\columnsep}{\OPTindexcolumnsep}
\printindex}

% The back cover
\include{back}

\end{document}

%%% Local Variables:
%%% mode: latex
%%% TeX-master: "hott-online"
%%% End:
